\documentclass[11pt,oneside]{book}
\usepackage[T1]{fontenc}
\usepackage[utf8]{inputenc}
\usepackage{amsmath,amssymb,amsthm}
\usepackage[margin=1in]{geometry}
\usepackage[hidelinks]{hyperref}
\usepackage{needspace}

\newbox\exerciseintrobox
\newdimen\exerciseintroheight
\newenvironment{exerciseintro}
  {\par\setbox\exerciseintrobox=\vbox\bgroup\hsize=\linewidth}
  {\par\egroup
   \exerciseintroheight=\ht\exerciseintrobox
   \advance\exerciseintroheight by \dp\exerciseintrobox
   \ifdim\exerciseintroheight<\textheight
     \Needspace{\exerciseintroheight}%
   \else
     \Needspace{16\baselineskip}%
   \fi
   \unvbox\exerciseintrobox}

\newenvironment{context}
  {\begin{quote}\small\noindent\textbf{Permitted context.}\par\smallskip}
  {\end{quote}}
\newenvironment{problem}
  {\par\medskip\noindent\textbf{Problem.}\ }
  {\par\medskip}
\newenvironment{solution}
  {\begin{proof}[Solution]}
  {\end{proof}}
\newcommand{\corpuschapter}[1]{%
  \clearpage\phantomsection
  \markboth{#1}{#1}%
  {\Huge\bfseries #1\par}\vspace{1.5em}%
}
\newcommand{\corpussection}[1]{%
  \par\bigskip\phantomsection
  \markright{#1}%
  {\Large\bfseries #1\par}\medskip
}
\newcommand{\exerciseheading}[1]{%
  \par\bigskip\phantomsection
  {\large\bfseries #1\par}\smallskip
}

\setcounter{tocdepth}{2}
\setlength{\emergencystretch}{2em}
\title{Hartshorne Exercise Solutions\\[0.5em]\large Finalized draft}
\author{Unofficial scholarly corpus}
\date{20 September 2026}

\begin{document}
\hypersetup{pageanchor=false}
\frontmatter
\maketitle
\hypersetup{pageanchor=true}
\tableofcontents
\mainmatter

\corpuschapter{Appendix A}
\addcontentsline{toc}{chapter}{Appendix A}

\begin{exerciseintro}
\corpussection{A.6 Appendix A Exercises}
\addcontentsline{toc}{section}{A.6 Appendix A Exercises}

% record: H-A-06-001
\exerciseheading{Exercise A.6.1}
\addcontentsline{toc}{subsection}{Exercise A.6.1}

\begin{context}
Work over an algebraically closed field \(k\).
Chow groups are cycles modulo the subgroup generated by pushforwards
of principal divisors on normalizations of integral subvarieties.
Proper pushforward, pullback between smooth varieties, the projection
formula, localization, and intersections with Cartier divisors are
available. A proper generically finite map pushes a principal divisor
to the divisor of the field norm; if its image has smaller dimension,
the corresponding pushforward is zero, using degree zero of a principal
divisor on the generic proper curve in the dimension-one fiber case.
\end{context}

\begin{problem}
Show that rational equivalence is generated by relations between the
fibers at \(0,1\) of cycles \(W\) on \(X\times\mathbb A^1\)
meeting those fibers properly.
\end{problem}
\end{exerciseintro}

\begin{solution}
First take an integral component of \(W\) dominating \(\mathbb A^1\).
Close it in \(X\times\mathbb P^1\) and normalize.
The rational function \(t/(t-1)\) has divisor the difference of its
fibers at \(0,1\), with their Cartier intersection multiplicities.
It has neither zero nor pole over infinity.
Projection to \(X\) is proper, so pushing this principal divisor
shows that the two fiber cycles are rationally equivalent.
Components over other fixed parameters meet neither specified
fiber; components inside those fibers are excluded by proper
intersection. Additivity proves this direction for arbitrary cycles.

Conversely a generating relation is \(\nu_*\operatorname{div}(f)\),
where \(\nu:\widetilde V\to V\subseteq X\) is normalization and
\(f\in k(V)^\times\).
Constant \(f\) gives the zero relation.
Otherwise take the closure of the graph of
\(t=f/(1+f)\) in \(\widetilde V\times\mathbb P^1\), and push
its cycle to \(X\times\mathbb P^1\).
The graph is birational to \(\widetilde V\).
At codimension-one points, normality gives a discrete valuation
ring, and the rational map to \(\mathbb P^1\) extends there.
The fiber at \(0\) is the zero divisor of \(f\), and the fiber
at \(1\) is its pole divisor, with their orders.
Any additional divisorial components over indeterminacy map
into codimension at least two in \(V\), so have zero pushforward
in the required dimension.
Restrict the resulting cycle to \(X\times\mathbb A^1\);
its fibers meet properly and their difference is exactly
\(\nu_*\operatorname{div}(f)\).
Adding constant families and taking sums generates all the original
rational-equivalence relations.
\end{solution}

\begin{exerciseintro}
% record: H-A-06-002
\exerciseheading{Exercise A.6.2}
\addcontentsline{toc}{subsection}{Exercise A.6.2}

\begin{context}
Let \(f:X\to X'\) be a proper dominant generically finite map of
normal varieties over \(k\).
Normalization in a finite function-field extension is finite.
A proper birational map to a normal variety is an isomorphism over
codimension-one points. A finite torsion-free module over a discrete
valuation ring is free. The norm is the determinant of multiplication
on the finite-dimensional extension field.
\end{context}

\begin{problem}
Prove that \(f_*\) carries linearly equivalent Weil divisors to
linearly equivalent Weil divisors.
\end{problem}
\end{exerciseintro}

\begin{solution}
It suffices to prove
\[
f_*\operatorname{div}_X(g)
=\operatorname{div}_{X'}N_{k(X)/k(X')}(g)
\quad(g\in k(X)^\times).
\]
Let \(Z\) be the normalization of \(X'\) in \(k(X)\), and factor
\[
X\xrightarrow{h}Z\xrightarrow{g_0}X'.
\]
Here \(h\) is proper birational and \(g_0\) is finite.
Over a neighborhood of every codimension-one point of \(X'\),
\(h\) is an isomorphism at the relevant codimension-one
points; exceptional divisors with smaller-dimensional image
contribute zero. Removing closed subsets of codimension at least
two therefore reduces the calculation to a finite map, flat at
these points by torsion-freeness over a discrete valuation ring.

Fix a prime divisor \(D\) of \(X'\).
Over its local discrete valuation ring \(R\), the finite normalization
algebra is a free \(R\)-module. For an integral nonzero \(g\),
the order of its determinant is the length of the multiplication
cokernel. Decomposing that finite-length module at the primes \(E\)
above \(D\) gives
\[
v_D(N(g))=\sum_{E\mapsto D}[k(E):k(D)]\,v_E(g).
\]
Write a general \(g\) as a quotient to obtain the same equality.
The right side is exactly the coefficient in divisor pushforward.
This proof uses no separability assumption.
Equality at every prime divisor proves the formula, and hence the claim.
\end{solution}

\begin{exerciseintro}
% record: H-A-06-003
\exerciseheading{Exercise A.6.3}
\addcontentsline{toc}{subsection}{Exercise A.6.3}

\begin{context}
Work over an algebraically closed field \(k\).
Chow groups are cycles modulo the subgroup generated by pushforwards
of principal divisors on normalizations of integral subvarieties.
Proper pushforward, pullback between smooth varieties, the projection
formula, localization, and intersections with Cartier divisors are
available. A proper generically finite map pushes a principal divisor
to the divisor of the field norm; if its image has smaller dimension,
the corresponding pushforward is zero, using degree zero of a principal
divisor on the generic proper curve in the dimension-one fiber case.

A projective \(r\)-fold admits a finite linear projection to
\(\mathbb P^r\) with center disjoint from it.
Over the principal ideal domain \(k[t]\), torsion-free modules are flat.
The Hilbert polynomial is constant in a projective flat family.
Fibers of a one-parameter cycle are rationally equivalent.
\end{context}

\begin{problem}
Use projection to prove directly that a degree-\(d\) integral
subvariety \(Y\subseteq\mathbb P^n\) is rationally equivalent
to \(d\) times a linear subspace of its dimension.
\end{problem}
\end{exerciseintro}

\begin{solution}
Choose coordinates \([u_0:\cdots:u_r:v_1:\cdots:v_{n-r}]\)
so that the \(u_i\) have no common zero on \(Y\) and define
a finite surjection \(Y\to L=\mathbb P^r\).
The map
\[
Y\times\mathbb A^1\longrightarrow\mathbb P^n\times\mathbb A^1,
\qquad ([u:v],t)\longmapsto([u:tv],t)
\]
is everywhere defined and proper. Give its image \(W\) the reduced
scheme structure. It is integral and dominates the parameter line,
so its structure sheaf is torsion-free, hence flat over \(k[t]\).
For \(t\ne0\), its fiber is an automorphic image of \(Y\).
For \(t=0\), its support is precisely \(L\), by finiteness and
surjectivity of the projection.
Its associated fiber cycle is consequently \(m[L]\) for a positive
integer \(m\). Constancy of the Hilbert polynomial makes its degree
equal to \(d\); as \(\deg L=1\), \(m=d\).
The relation between the fibers at \(1\) and \(0\) is
\([Y]\sim d[L]\), as required.
\end{solution}

\begin{exerciseintro}
% record: H-A-06-004
\exerciseheading{Exercise A.6.4}
\addcontentsline{toc}{subsection}{Exercise A.6.4}

\begin{context}
Work over an algebraically closed field \(k\).
Chow groups are cycles modulo the subgroup generated by pushforwards
of principal divisors on normalizations of integral subvarieties.
Proper pushforward, pullback between smooth varieties, the projection
formula, localization, and intersections with Cartier divisors are
available. A proper generically finite map pushes a principal divisor
to the divisor of the field norm; if its image has smaller dimension,
the corresponding pushforward is zero, using degree zero of a principal
divisor on the generic proper curve in the dimension-one fiber case.

Let \(\pi:X=\mathbb P(\mathcal E)\to C\) be a ruled surface over
a smooth projective curve. It has a section \(s\): a rank-two
bundle on a smooth curve has a line quotient after saturating
a rational quotient. Every fiber is \(\mathbb P^1\), on which
any two points are rationally equivalent. On a smooth curve,
zero-cycles modulo rational equivalence form its Picard group.
\end{context}

\begin{problem}
Prove \(A^2(X)\cong\operatorname{Pic}C\).
\end{problem}
\end{exerciseintro}

\begin{solution}
Proper maps give \(\pi_*:A^2(X)\to A^1(C)\) and
\(s_*:A^1(C)\to A^2(X)\), with \(\pi_*s_*=\mathrm{id}\).
For each closed point \(x\in X\), the fiber over \(\pi(x)\)
contains both \(x\) and \(s(\pi(x))\); a rational function on
that projective line has divisor their difference.
Thus \([x]=[s(\pi(x))]\) in \(A^2(X)\).
Point cycles generate, so \(s_*\pi_*=\mathrm{id}\) as well.
The resulting inverse isomorphisms identify this group with
\(A^1(C)=\operatorname{Pic}C\).
\end{solution}

\begin{exerciseintro}
% record: H-A-06-005
\exerciseheading{Exercise A.6.5}
\addcontentsline{toc}{subsection}{Exercise A.6.5}

\begin{context}
Work over an algebraically closed field \(k\).
Chow groups are cycles modulo the subgroup generated by pushforwards
of principal divisors on normalizations of integral subvarieties.
Proper pushforward, pullback between smooth varieties, the projection
formula, localization, and intersections with Cartier divisors are
available. A proper generically finite map pushes a principal divisor
to the divisor of the field norm; if its image has smaller dimension,
the corresponding pushforward is zero, using degree zero of a principal
divisor on the generic proper curve in the dimension-one fiber case.

Let \(\pi:\widetilde X\to X\) blow up a point \(P\) of a smooth
projective surface. Its exceptional curve \(E\cong\mathbb P^1\)
has normal bundle \(\mathcal O_E(-1)\).
The map is an isomorphism off \(E\).
Divisor classes satisfy
\(\operatorname{Pic}\widetilde X=\pi^*\operatorname{Pic}X\oplus\mathbb Z[E]\).
One may choose two very ample divisors meeting transversely at \(P\)
and at their other intersection points.
\end{context}

\begin{problem}
Prove that \(A(\widetilde X)=\pi^*A(X)\oplus\mathbb Z[E]\)
as graded groups, and describe multiplication, including
\(E^2=-\pi^*[P]\).
\end{problem}
\end{exerciseintro}

\begin{solution}
Let \(q\) be the class of a point of \(E\).
The normal-bundle self-intersection calculation gives \(E^2=-q\).
Also \(E\cdot\pi^*D=0\) for every divisor \(D\) on \(X\),
because its pulled-back line bundle restricts trivially to \(E\).

To identify \(\pi^*[P]\), choose the two transverse divisors
in the context, with distinct tangent directions at \(P\).
Their total transforms are their strict transforms plus \(E\).
Over \(P\) the strict transforms are disjoint, each meeting
\(E\) transversely in a point. Expanding their intersection
therefore gives \(q+q+E^2=q\) above \(P\).
Away from \(P\) the map is an isomorphism, so subtracting the
other intersection points proves \(\pi^*[P]=q\).
Hence \(E^2=-\pi^*[P]\).

The projection formula gives \(\pi_*\pi^*=\mathrm{id}\).
For any zero-cycle class \(\alpha\) on \(\widetilde X\), the class
\(\alpha-\pi^*\pi_*\alpha\) restricts to zero off \(E\).
By localization it is a multiple \(nq\).
Pushing down gives \(n[P]=0\), and pulling back then gives
\(nq=0\); thus \(\pi^*\) is an isomorphism on zero-cycles.
It is an isomorphism in degree zero as well.
The Picard splitting handles degree one, proving the graded
group decomposition. The relations just computed and the
fact that \(\pi^*\) is a ring map determine multiplication;
all products of total codimension greater than two vanish.
\end{solution}

\begin{exerciseintro}
% record: H-A-06-006
\exerciseheading{Exercise A.6.6}
\addcontentsline{toc}{subsection}{Exercise A.6.6}

\begin{context}
For a smooth projective variety, Chern classes multiply in an exact
sequence of vector bundles. The self-intersection formula for a
regular immersion \(i:Y\to X\) of codimension \(r\) is
\(i^*i_*1=c_r(N_{Y/X})\).
Hirzebruch--Riemann--Roch says
\(\chi(\mathcal F)=\deg(\operatorname{ch}(\mathcal F)\operatorname{td}(T_X))_{\dim X}\).
The degree-at-most-three Todd expansion is
\[
\operatorname{td}(T_X)=1+\frac{c_1}{2}
+\frac{c_1^2+c_2}{12}+\frac{c_1c_2}{24},
\]
and \(\operatorname{ch}(\mathcal O(D))=\exp(D)\).
\end{context}

\begin{problem}
Identify the self-intersection of the diagonal of a smooth
projective \(n\)-fold with \(c_n(T_X)\).
\end{problem}
\end{exerciseintro}

\begin{solution}
For \(i:X\cong\Delta\hookrightarrow X\times X\), restriction
of the ambient tangent bundle is \(T_X\oplus T_X\).
The tangent map is \(v\mapsto(v,v)\), whose quotient is
canonically \(T_X\) by \((v,w)\mapsto w-v\).
Thus \(N_{\Delta/(X\times X)}\cong T_X\).
Self-intersection gives \(i^*i_*1=c_n(T_X)\) in \(A^n(X)\).
Equivalently in the ambient Chow ring,
\([\Delta]^2=i_*c_n(T_X)\).
The latter formula makes explicit which group contains the
ordinary product of the two diagonal classes.
\end{solution}

\begin{exerciseintro}
% record: H-A-06-007
\exerciseheading{Exercise A.6.7}
\addcontentsline{toc}{subsection}{Exercise A.6.7}

\begin{context}
For a smooth projective variety, Chern classes multiply in an exact
sequence of vector bundles. The self-intersection formula for a
regular immersion \(i:Y\to X\) of codimension \(r\) is
\(i^*i_*1=c_r(N_{Y/X})\).
Hirzebruch--Riemann--Roch says
\(\chi(\mathcal F)=\deg(\operatorname{ch}(\mathcal F)\operatorname{td}(T_X))_{\dim X}\).
The degree-at-most-three Todd expansion is
\[
\operatorname{td}(T_X)=1+\frac{c_1}{2}
+\frac{c_1^2+c_2}{12}+\frac{c_1c_2}{24},
\]
and \(\operatorname{ch}(\mathcal O(D))=\exp(D)\).

Write \(K\) for the canonical divisor, so \(c_1(T_X)=-K\).
For a threefold, \(\chi(\mathcal O_X)=1-p_a(X)\).
Products of classes of total degree three below denote their degrees.
\end{context}

\begin{problem}
For a smooth projective threefold, prove
\[
1-p_a=\frac{c_1c_2}{24},\qquad
\chi(\mathcal O(D))=
\frac{D(D-K)(2D-K)}{12}+\frac{Dc_2}{12}+1-p_a.
\]
\end{problem}
\end{exerciseintro}

\begin{solution}
Apply Riemann--Roch to the structure sheaf: its character is one,
so only the degree-three Todd term contributes, giving
\(\chi(\mathcal O_X)=c_1c_2/24=1-p_a\).
For \(\mathcal O(D)\), multiply the Todd expansion by
\(1+D+D^2/2+D^3/6\). The top-degree term is
\[
\frac{D^3}{6}+\frac{D^2c_1}{4}
+\frac{D(c_1^2+c_2)}{12}+\frac{c_1c_2}{24}.
\]
Substituting \(c_1=-K\) and the first identity gives the formula,
since \(2D^3-3D^2K+DK^2=D(D-K)(2D-K)\).
\end{solution}

\begin{exerciseintro}
% record: H-A-06-008
\exerciseheading{Exercise A.6.8}
\addcontentsline{toc}{subsection}{Exercise A.6.8}

\begin{context}
For a smooth projective variety, Chern classes multiply in an exact
sequence of vector bundles. The self-intersection formula for a
regular immersion \(i:Y\to X\) of codimension \(r\) is
\(i^*i_*1=c_r(N_{Y/X})\).
Hirzebruch--Riemann--Roch says
\(\chi(\mathcal F)=\deg(\operatorname{ch}(\mathcal F)\operatorname{td}(T_X))_{\dim X}\).
The degree-at-most-three Todd expansion is
\[
\operatorname{td}(T_X)=1+\frac{c_1}{2}
+\frac{c_1^2+c_2}{12}+\frac{c_1c_2}{24},
\]
and \(\operatorname{ch}(\mathcal O(D))=\exp(D)\).

On \(\mathbb P^3\), \(A(\mathbb P^3)=\mathbb Z[h]/(h^4)\),
\(\deg h^3=1\), and the Euler sequence gives
\(\operatorname{td}(T_{\mathbb P^3})=1+2h+11h^2/6+h^3\).
For a rank-two bundle with \(c_1=ah,c_2=bh^2\),
\[
\operatorname{ch}(\mathcal E)
=2+ah+(a^2-2b)h^2/2+(a^3-3ab)h^3/6.
\]
\end{context}

\begin{problem}
Prove that the two integer Chern numbers of a rank-two bundle
on \(\mathbb P^3\) satisfy \(ab\equiv0\pmod2\).
\end{problem}
\end{exerciseintro}

\begin{solution}
The coefficient of \(h^3\) in the product gives
\[
\chi(\mathcal E)=2+\frac{11a}{6}+a^2-2b+\frac{a^3-3ab}{6}
=1+\binom{a+3}{3}-2b-\frac{ab}{2}.
\]
The Euler characteristic and the binomial expression are integers
for every integer \(a\), including negative ones.
Therefore \(ab/2\) is an integer, proving the parity assertion.
\end{solution}

\begin{exerciseintro}
% record: H-A-06-009
\exerciseheading{Exercise A.6.9}
\addcontentsline{toc}{subsection}{Exercise A.6.9}

\begin{context}
For a smooth projective variety, Chern classes multiply in an exact
sequence of vector bundles. The self-intersection formula for a
regular immersion \(i:Y\to X\) of codimension \(r\) is
\(i^*i_*1=c_r(N_{Y/X})\).
Hirzebruch--Riemann--Roch says
\(\chi(\mathcal F)=\deg(\operatorname{ch}(\mathcal F)\operatorname{td}(T_X))_{\dim X}\).
The degree-at-most-three Todd expansion is
\[
\operatorname{td}(T_X)=1+\frac{c_1}{2}
+\frac{c_1^2+c_2}{12}+\frac{c_1c_2}{24},
\]
and \(\operatorname{ch}(\mathcal O(D))=\exp(D)\).

For a smooth surface in \(\mathbb P^4\), with \(d=H^2\), the normal
sequence and self-intersection give
\[
d^2-10d-5HK-2K^2+12\chi(\mathcal O_X)=0.
\]
On the rational ruled surface \(\mathbb F_e\),
\(C_0^2=-e,\ C_0F=1,\ F^2=0,\ K=-2C_0-(e+2)F\),
\(K^2=8,\ \chi=1\).
A divisor \(aC_0+bF\) is very ample exactly when
\(a\geq1,\ b\geq ae+1\).
The quadric realizes \(\mathbb F_0\) in \(\mathbb P^3\);
\(|C_0+2F|\) realizes \(\mathbb F_1\) as the cubic scroll in \(\mathbb P^4\).
A K3 surface has \(K=0,\chi=2\); an abelian surface has \(K=0,\chi=0\).
\end{context}

\begin{problem}
(a) Check the surface formula on the rational cubic scroll.

(b) Determine the possible degrees of a K3 surface embedded in
\(\mathbb P^4\).

(c) Determine the degree of an abelian surface embedded in \(\mathbb P^4\).

(d) Determine which \(\mathbb F_e\), \(e\geq0\), embed in \(\mathbb P^4\).
\end{problem}
\end{exerciseintro}

\begin{solution}
(a) On \(\mathbb F_1\), take \(H=C_0+2F\).
Then \(d=3,HK=-5,K^2=8,\chi=1\), so the left side is
\(9-30+25-16+12=0\).

(b) The formula becomes \(d^2-10d+24=0\), hence \(d=4\) or \(6\).
These possibilities are realized by a smooth quartic in a linear
\(\mathbb P^3\), and a smooth \((2,3)\) complete intersection in
\(\mathbb P^4\), respectively.

(c) It becomes \(d^2-10d=0\). Since \(d>0\), it gives \(d=10\).

(d) Let the hyperplane divisor of an embedding be \(H=aC_0+bF\).
Put \(t=2b-ae\). Very ampleness gives \(a\geq1,t\geq ae+2\geq2\).
Then \(d=at\), \(HK=-t-2a\), and the formula is
\[
a^2t^2-10at+5t+10a-4=0.
\]
For \(a\geq2\), this polynomial increases for \(t\geq2\), since
its derivative there is at least \(4a^2-10a+5>0\).
At \(t=2\) its value is \(2(2a-3)(a-1)>0\), a contradiction.
Thus \(a=1\), and the equation is \((t-2)(t-3)=0\).
If \(t=2\), \(2b-e=2\) and \(b\geq e+1\) imply \(e=0,b=1\).
If \(t=3\), the same inequalities and parity imply \(e=1,b=2\).
Both occur by the stated embeddings, with the quadric included in
a hyperplane of \(\mathbb P^4\). Thus exactly \(\mathbb F_0,\mathbb F_1\)
admit such embeddings.
\end{solution}

\begin{exerciseintro}
% record: H-A-06-010
\exerciseheading{Exercise A.6.10}
\addcontentsline{toc}{subsection}{Exercise A.6.10}

\begin{context}
For a smooth projective variety, Chern classes multiply in an exact
sequence of vector bundles. The self-intersection formula for a
regular immersion \(i:Y\to X\) of codimension \(r\) is
\(i^*i_*1=c_r(N_{Y/X})\).
Hirzebruch--Riemann--Roch says
\(\chi(\mathcal F)=\deg(\operatorname{ch}(\mathcal F)\operatorname{td}(T_X))_{\dim X}\).
The degree-at-most-three Todd expansion is
\[
\operatorname{td}(T_X)=1+\frac{c_1}{2}
+\frac{c_1^2+c_2}{12}+\frac{c_1c_2}{24},
\]
and \(\operatorname{ch}(\mathcal O(D))=\exp(D)\).

The tangent bundle of an abelian variety is trivial.
For a smooth \(r\)-fold in \(\mathbb P^n\), its normal bundle
has rank \(n-r\). Chern classes above the rank vanish.
The Euler sequence gives \(c(T_{\mathbb P^n})=(1+h)^{n+1}\).
\end{context}

\begin{problem}
Prove that an abelian threefold cannot be embedded in \(\mathbb P^5\).
\end{problem}
\end{exerciseintro}

\begin{solution}
If it were embedded, its normal bundle \(N\) would have rank two.
The tangent-normal exact sequence and triviality of the tangent
bundle of \(X\) would give
\[
c(N)=c(T_{\mathbb P^5}|_X)=(1+H)^6.
\]
The degree-three term would be \(c_3(N)=20H^3\).
But rank two forces \(c_3(N)=0\). Taking degrees yields
\(20\deg X=0\) in \(\mathbb Z\), impossible since an embedded
projective threefold has positive degree.
This is an integer Chow-group argument, so it does not disappear
in characteristics two or five.
\end{solution}

\corpuschapter{Appendix B}
\addcontentsline{toc}{chapter}{Appendix B}

\begin{exerciseintro}
\corpussection{B.6 Appendix B Exercises}
\addcontentsline{toc}{section}{B.6 Appendix B Exercises}

% record: H-B-06-001
\exerciseheading{Exercise B.6.1}
\addcontentsline{toc}{subsection}{Exercise B.6.1}

\begin{context}
All algebraic schemes are of finite type over \(\mathbb C\).
Analytification preserves and detects reducedness, smoothness and
connectedness; algebraic separatedness is equivalent to analytic
Hausdorffness, and algebraic properness to analytic properness.
The analytic local rings are faithfully flat over the corresponding
algebraic local rings; analytification is faithful and exact on
coherent sheaves. A smooth separated algebraic curve is its smooth
projective completion minus finitely many points.

A bounded holomorphic function on a punctured disc extends across
the puncture. Holomorphic functions on a compact connected Riemann
surface are constant.
\end{context}

\begin{problem}
Show that the open unit disc is not the analytification of an
algebraic scheme.
\end{problem}
\end{exerciseintro}

\begin{solution}
If \(X_h\) were the unit disc, the comparison facts would make
\(X\) a smooth connected separated curve.
Its smooth projective completion \(\overline X\) adds finitely
many points. The coordinate of the disc pulls back to a bounded
nonconstant holomorphic function on \(X_h\).
Near every added point, it is a bounded function on a punctured
coordinate disc, so extends holomorphically.
It would then be a nonconstant holomorphic function on the compact
connected Riemann surface \(\overline X_h\), a contradiction.
The same contradiction applies if there were no added points.
\end{solution}

\begin{exerciseintro}
% record: H-B-06-002
\exerciseheading{Exercise B.6.2}
\addcontentsline{toc}{subsection}{Exercise B.6.2}

\begin{context}
All algebraic schemes are of finite type over \(\mathbb C\).
Analytification preserves and detects reducedness, smoothness and
connectedness; algebraic separatedness is equivalent to analytic
Hausdorffness, and algebraic properness to analytic properness.
The analytic local rings are faithfully flat over the corresponding
algebraic local rings; analytification is faithful and exact on
coherent sheaves. A smooth separated algebraic curve is its smooth
projective completion minus finitely many points.

Weierstrass's product theorem supplies an entire function with any
prescribed discrete zero set and positive integral multiplicities.
Every ideal in \(\mathbb C[t]\) is principal.
\end{context}

\begin{problem}
Let \(|z_n|\to\infty\), and let \(\mathcal J\subseteq\mathcal O_{\mathbb C}\)
be the sheaf of holomorphic functions vanishing at these points.
Show that it is not the analytification of an algebraic ideal in
\(\mathcal O_{\mathbb A^1}\), but is abstractly isomorphic to the
analytification of a coherent algebraic sheaf.
\end{problem}
\end{exerciseintro}

\begin{solution}
Use the set of distinct \(z_n\); it is infinite and discrete,
regardless of any repetitions in the sequence.
A nonzero algebraic ideal is generated by a nonzero polynomial,
whose zero set is finite, so its analytic ideal cannot equal
\(\mathcal J\). The zero ideal cannot work either, since
\(\mathcal J\) is nonzero.
Choose an entire function \(g\) having simple zeros exactly at
the specified set. Locally the vanishing ideal is \((g)\):
at a zero both are generated by a local parameter; elsewhere
\(g\) is a unit.
Multiplication by \(g\) therefore gives an isomorphism
\(\mathcal O_{\mathbb C}\cong\mathcal J\).
Taking the algebraic sheaf \(\mathcal F=\mathcal O_{\mathbb A^1}\)
gives \(\mathcal F_h\cong\mathcal J\) as sheaves.
This isomorphism does not identify their specified inclusions
as algebraic ideals.
\end{solution}

\begin{exerciseintro}
% record: H-B-06-003
\exerciseheading{Exercise B.6.3}
\addcontentsline{toc}{subsection}{Exercise B.6.3}

\begin{context}
All algebraic schemes are of finite type over \(\mathbb C\).
Analytification preserves and detects reducedness, smoothness and
connectedness; algebraic separatedness is equivalent to analytic
Hausdorffness, and algebraic properness to analytic properness.
The analytic local rings are faithfully flat over the corresponding
algebraic local rings; analytification is faithful and exact on
coherent sheaves. A smooth separated algebraic curve is its smooth
projective completion minus finitely many points.

The punctured algebraic plane \(U=\mathbb A^2\setminus\{0\}\)
has \(\operatorname{Pic}U=0\): removing codimension two preserves
the divisor class group of the normal factorial affine plane,
and on a smooth variety Picard and divisor class groups agree.
A nowhere-zero holomorphic function on \(\mathbb C^*\times\mathbb C\)
has a holomorphic logarithm exactly when its winding number around
the \(\mathbb C^*\) factor is zero. Cauchy's integral theorem applies
to each variable separately.
\end{context}

\begin{problem}
On \(\mathbb C^2\setminus\{0\}\), glue trivial analytic line bundles
on \(U_1=\{z\ne0\}\), \(U_2=\{w\ne0\}\) by
\(\exp(-1/(zw))\) on the overlap.
Prove that the resulting bundle is not algebraizable.
\end{problem}
\end{exerciseintro}

\begin{solution}
Any coherent algebraization of the analytic line bundle would itself
be invertible: faithful flatness of analytic local rings descends
local freeness of rank one. Such an algebraic line bundle would be
trivial by \(\operatorname{Pic}U=0\).
Suppose the analytic bundle were trivial. Its transition would be
a quotient \(a/b\), with \(a\) a holomorphic unit on \(U_1\)
and \(b\) a holomorphic unit on \(U_2\).
Its winding along each coordinate circle is zero, since it is
the exponential of the single-valued function \(-1/(zw)\).
Along a \(z\)-circle, \(b\) has winding zero because it is a unit
on a full \(z\)-disc; hence \(a\) has winding zero too.
Similarly \(b\) has zero winding in its \(w\)-direction.
The logarithm criterion gives \(a=e^A,b=e^B\) with \(A,B\)
holomorphic on their respective opens.
Connectedness of the overlap then gives
\[
-\frac1{zw}=A-B+2\pi i n
\]
for a fixed integer \(n\).
Integrate against \(dz\,dw\) over two coordinate circles.
The left side has integral \(-(2\pi i)^2\).
The \(A\)-term integrates to zero first in \(w\), since \(A\)
is holomorphic across \(w=0\); the \(B\)-term integrates to zero
first in \(z\); the constant term also gives zero.
This contradiction proves nontriviality and hence nonalgebraizability.
\end{solution}

\begin{exerciseintro}
% record: H-B-06-004
\exerciseheading{Exercise B.6.4}
\addcontentsline{toc}{subsection}{Exercise B.6.4}

\begin{context}
All algebraic schemes are of finite type over \(\mathbb C\).
Analytification preserves and detects reducedness, smoothness and
connectedness; algebraic separatedness is equivalent to analytic
Hausdorffness, and algebraic properness to analytic properness.
The analytic local rings are faithfully flat over the corresponding
algebraic local rings; analytification is faithful and exact on
coherent sheaves. A smooth separated algebraic curve is its smooth
projective completion minus finitely many points.

Holomorphic functions on a reduced compact connected complex
analytic space are constant (apply the maximum principle on
normalizations of its irreducible components).
A separated finite-type scheme has a proper compactification.
In an integral variety, an affine neighborhood of a boundary point
can be cut by hypersurfaces to obtain a curve through that point
not contained in a prescribed proper closed subset.
A nonproper separated integral algebraic curve is affine.
Projective GAGA identifies global sections of coherent sheaves.
\end{context}

\begin{problem}
Prove that for reduced proper \(X\), algebraic and analytic global
regular functions agree. Prove the converse detection statement
for separated \(X\): if \(X\) is not proper, some coherent
\(\mathcal F\) has a nonsurjective comparison map on global sections.
Explain why separatedness is necessary in the converse.
\end{problem}
\end{exerciseintro}

\begin{solution}
If \(X\) is reduced and proper, both algebraic and analytic functions
are constant on each connected component. Algebraically, properness
forces the image of a morphism to the affine line to be finite;
reducedness identifies the ring with these componentwise constants.
Analytically use compactness and the maximum principle.
The connected components correspond under analytification and are
finite in number. Both rings are thus the same finite product of
copies of \(\mathbb C\), with the comparison map its identity.

Now let \(X\) be separated but not proper. Choose a proper
compactification \(j:X\hookrightarrow\overline X\), and let
\(D\) be a reduced irreducible component of the closure of
\(X_{\mathrm{red}}\) meeting \(\overline X\setminus X\).
Select a boundary point of \(D\). Curve selection inside \(D\),
followed by closure in \(\overline X\), gives a proper integral
curve \(\overline C\) through that point meeting \(X\) but
not contained in it. Its intersection \(C=\overline C\cap X\)
is closed in \(X\), nonempty, and a proper open of \(\overline C\);
it is a nonproper affine integral curve.
Choose a nonconstant \(h\in\Gamma(C,\mathcal O_C)\).
The analytic function \(e^h\) is not algebraic regular:
on the smooth projective model of \(C\), the rational function
\(h\) has a pole somewhere outside \(C\), and \(e^h\) has an
essential singularity there. An algebraic function would be meromorphic.
For the closed inclusion \(i:C\hookrightarrow X\), take
\(\mathcal F=i_*\mathcal O_C\), a coherent sheaf.
Its analytic global sections contain \(e^h\), absent from the
image of its algebraic global sections, proving the claim.

Without separatedness the converse as stated is false.
Glue two copies of \(\mathbb P^1\) along their identical complements
of one point, leaving that point doubled. The resulting finite-type
scheme is not separated, hence not proper.
For any coherent sheaf, sections on each of the two projective opens
algebraize by GAGA. An analytic global section gives two such
algebraic sections agreeing analytically on the overlap; faithfulness
makes them agree algebraically, so they glue.
Thus every coherent global-section comparison is an isomorphism
on this nonproper scheme.
\end{solution}

\begin{exerciseintro}
% record: H-B-06-005
\exerciseheading{Exercise B.6.5}
\addcontentsline{toc}{subsection}{Exercise B.6.5}

\begin{context}
All algebraic schemes are of finite type over \(\mathbb C\).
Analytification preserves and detects reducedness, smoothness and
connectedness; algebraic separatedness is equivalent to analytic
Hausdorffness, and algebraic properness to analytic properness.
The analytic local rings are faithfully flat over the corresponding
algebraic local rings; analytification is faithful and exact on
coherent sheaves. A smooth separated algebraic curve is its smooth
projective completion minus finitely many points.

Chow's theorem makes a closed analytic subspace of complex projective
space algebraic. A finite birational map onto a normal curve is an
isomorphism. Proper maps of algebraic curves with finite fibers are
finite. A holomorphic map to a coordinate disc extends across a
puncture if its coordinate is bounded.
\end{context}

\begin{problem}
If two nonsingular affine algebraic curves have isomorphic
analytic spaces, prove that the algebraic curves are isomorphic.
\end{problem}
\end{exerciseintro}

\begin{solution}
Complete the curves to smooth projective curves
\(\overline X,\overline X'\), adding finite nonempty sets \(S,S'\).
Let \(f:X_h\to X'_h\) be the biholomorphism.
At \(p\in S\), its cluster set in compact \(\overline X'_h\)
is nonempty and connected: it is the intersection of the nested
compact connected closures of images of punctured small discs.
It contains no point of \(X'_h\), since convergence to such a point
and continuity of \(f^{-1}\) would force convergence inside \(X_h\).
It is therefore a connected subset of finite \(S'\), hence a single
point \(p'\). Thus \(f\) extends continuously at \(p\).
In a coordinate disc at \(p'\), the extended map is holomorphic
by the removable-singularity theorem. Apply the same argument to
\(f^{-1}\). The extensions are inverse biholomorphisms of the
compactifications and identify their missing sets.

The graph of this biholomorphism is a closed analytic subspace
of the projective product. By Chow it is algebraic.
Each projection of its reduced algebraic graph is a proper
bijective map of smooth curves, generically of degree one in
characteristic zero; thus it is finite birational, hence an
isomorphism by normality. The resulting algebraic isomorphism
restricts to \(X\cong X'\), since it matches \(S,S'\).
\end{solution}

\begin{exerciseintro}
% record: H-B-06-006
\exerciseheading{Exercise B.6.6}
\addcontentsline{toc}{subsection}{Exercise B.6.6}

\begin{context}
All algebraic schemes are of finite type over \(\mathbb C\).
Analytification preserves and detects reducedness, smoothness and
connectedness; algebraic separatedness is equivalent to analytic
Hausdorffness, and algebraic properness to analytic properness.
The analytic local rings are faithfully flat over the corresponding
algebraic local rings; analytification is faithful and exact on
coherent sheaves. A smooth separated algebraic curve is its smooth
projective completion minus finitely many points.

Projective GAGA gives an equivalence of coherent algebraic and analytic
sheaves, compatible with global sections and morphisms.
Faithful flatness detects invertibility and surjectivity of coherent
maps. A line bundle with \(n+1\) generating sections defines a morphism
to \(\mathbb P^n\), uniquely up to isomorphism of that data.
\end{context}

\begin{problem}
Prove that an analytic morphism \(f:X_h\to Y_h\) between
analytifications of projective schemes has a unique algebraic
origin, allowing nonreduced schemes.
\end{problem}
\end{exerciseintro}

\begin{solution}
Embed \(Y\) in \(\mathbb P^n\). The composite analytic map
pulls back \(\mathcal O(1)\) with its \(n+1\) standard sections.
GAGA algebraizes this coherent sheaf as \(\mathcal L\) on \(X\)
and the sections as \(s_0,\ldots,s_n\).
Faithful flatness at closed points makes \(\mathcal L\) invertible
and the evaluation map \(\mathcal O_X^{n+1}\to\mathcal L\) surjective:
its analytic cokernel is zero, so the algebraic cokernel is zero.
The resulting morphism \(g:X\to\mathbb P^n\) has analytification
equal to the composite with \(f\).

Let \(\mathcal I_Y\) be the homogeneous ideal sheaf of \(Y\).
The map \(g^*\mathcal I_Y\to\mathcal O_X\) analytifies to zero
because the analytic morphism factors through \(Y_h\).
Faithfulness forces the original map to be zero as well.
Thus \(g\) factors through the closed subscheme \(Y\), giving
the desired algebraic morphism, including its nilpotent structure.
For uniqueness, an analytic identification of the pulled-back
line bundles and generating sections algebraizes uniquely by
full faithfulness of GAGA. The associated maps to \(\mathbb P^n\)
coincide; their factorizations through the fixed closed immersion
of \(Y\) coincide as well.
\end{solution}

\corpuschapter{Appendix C}
\addcontentsline{toc}{chapter}{Appendix C}

\begin{exerciseintro}
\corpussection{C.5 Appendix C Exercises}
\addcontentsline{toc}{section}{C.5 Appendix C Exercises}

% record: H-C-05-001
\exerciseheading{Exercise C.5.1}
\addcontentsline{toc}{subsection}{Exercise C.5.1}

\begin{context}
For a scheme of finite type over \(\mathbb F_q\), put
\(N_r(X)=\#X(\mathbb F_{q^r})\) and
\[
Z(X,t)=\exp\left(\sum_{r\geq1}N_r(X)t^r/r\right)\in1+t\mathbb Q[[t]].
\]
Formal logarithm and exponential are inverse over \(\mathbb Q\), and
\(-\log(1-at)=\sum_{r\geq1}a^r t^r/r\).
\end{context}

\begin{problem}
For a disjoint decomposition into locally closed
\(\mathbb F_q\)-subschemes \(X_i\), prove
\(Z(X,t)=\prod_i Z(X_i,t)\).
\end{problem}
\end{exerciseintro}

\begin{solution}
Every \(\mathbb F_{q^r}\)-point lies in exactly one stratum,
so \(N_r(X)=\sum_i N_r(X_i)\).
For a finite decomposition, substitute this in the defining logarithm
and exponentiate to obtain the product.
The same proof allows an infinite decomposition: for each bounded
range of \(r\), only finitely many strata can contribute points,
since each \(X(\mathbb F_{q^r})\) is finite.
Thus the sums and product are well defined coefficient by coefficient
in the \(t\)-adic topology.
\end{solution}

\begin{exerciseintro}
% record: H-C-05-002
\exerciseheading{Exercise C.5.2}
\addcontentsline{toc}{subsection}{Exercise C.5.2}

\begin{context}
For a scheme of finite type over \(\mathbb F_q\), put
\(N_r(X)=\#X(\mathbb F_{q^r})\) and
\[
Z(X,t)=\exp\left(\sum_{r\geq1}N_r(X)t^r/r\right)\in1+t\mathbb Q[[t]].
\]
Formal logarithm and exponential are inverse over \(\mathbb Q\), and
\(-\log(1-at)=\sum_{r\geq1}a^r t^r/r\).

For smooth projective dimension \(n\), the Weil properties ask for
rationality, a functional equation
\(Z(1/(q^nt))=\pm q^{nE/2}t^E Z(t)\), and factors \(P_i\)
with reciprocal roots of absolute value \(q^{i/2}\).
Their degrees should equal the complex Betti numbers for a smooth
lift, and \(E\) is the top tangent Chern number.
The Euler sequence gives \(c(T_{\mathbb P^n})=(1+h)^{n+1}\).
Complex projective \(n\)-space has one cell in each even real
dimension \(0,2,\ldots,2n\).
\end{context}

\begin{problem}
Compute \(Z(\mathbb P^n,t)\) and verify all the Weil properties
for projective space.
\end{problem}
\end{exerciseintro}

\begin{solution}
Counting nonzero vectors modulo nonzero scalar gives
\[
N_r=1+q^r+\cdots+q^{nr}.
\]
The defining logarithm therefore yields
\[
Z(\mathbb P^n,t)=\prod_{j=0}^n(1-q^jt)^{-1},
\]
a rational function with integer factors.
The top tangent Chern number is \(E=n+1\).
Direct substitution gives
\[
Z(\mathbb P^n,1/(q^nt))
=(-1)^{n+1}q^{n(n+1)/2}t^{n+1}Z(\mathbb P^n,t),
\]
the required equation.
Take \(P_{2j}=1-q^jt\) and \(P_{2j+1}=1\).
Each reciprocal root has absolute value \(q^j=q^{(2j)/2}\);
there are no odd-degree roots. The factor degrees give one
Betti number in each even degree and zero in odd degrees,
matching the cell decomposition of \(\mathbb P^n(\mathbb C)\).
Their alternating sum is \(n+1=E\).
\end{solution}

\begin{exerciseintro}
% record: H-C-05-003
\exerciseheading{Exercise C.5.3}
\addcontentsline{toc}{subsection}{Exercise C.5.3}

\begin{context}
For a scheme of finite type over \(\mathbb F_q\), put
\(N_r(X)=\#X(\mathbb F_{q^r})\) and
\[
Z(X,t)=\exp\left(\sum_{r\geq1}N_r(X)t^r/r\right)\in1+t\mathbb Q[[t]].
\]
Formal logarithm and exponential are inverse over \(\mathbb Q\), and
\(-\log(1-at)=\sum_{r\geq1}a^r t^r/r\).
\end{context}

\begin{problem}
Prove \(Z(X\times\mathbb A^1,t)=Z(X,qt)\).
\end{problem}
\end{exerciseintro}

\begin{solution}
For every \(r\), the affine line has \(q^r\) rational points,
so \(N_r(X\times\mathbb A^1)=q^rN_r(X)\).
The logarithm of the left side is
\(\sum N_r(X)(qt)^r/r\), exactly that of the right side.
Exponentiating proves the identity.
\end{solution}

\begin{exerciseintro}
% record: H-C-05-004
\exerciseheading{Exercise C.5.4}
\addcontentsline{toc}{subsection}{Exercise C.5.4}

\begin{context}
For a scheme of finite type over \(\mathbb F_q\), put
\(N_r(X)=\#X(\mathbb F_{q^r})\) and
\[
Z(X,t)=\exp\left(\sum_{r\geq1}N_r(X)t^r/r\right)\in1+t\mathbb Q[[t]].
\]
Formal logarithm and exponential are inverse over \(\mathbb Q\), and
\(-\log(1-at)=\sum_{r\geq1}a^r t^r/r\).

A closed point of degree \(d\) over \(\mathbb F_q\) has residue
field \(\mathbb F_{q^d}\) and contributes \(d\) points over
\(\mathbb F_{q^r}\) if \(d\mid r\), and none otherwise.
The arithmetic Euler product is
\(\zeta_X(s)=\prod_{x\ \text{closed}}(1-\#k(x)^{-s})^{-1}\).
\end{context}

\begin{problem}
For \(X\) of finite type over \(\mathbb F_q\), prove
\(\zeta_X(s)=Z(X,q^{-s})\), interpreting the equality first
as a formal Euler product and then where it converges.
\end{problem}
\end{exerciseintro}

\begin{solution}
Let \(b_d\) be the number of closed points of degree \(d\).
It is finite, and \(N_r=\sum_{d\mid r}d b_d\).
Consequently
\[
\log\prod_{d\geq1}(1-t^d)^{-b_d}
=\sum_{d,m\geq1}b_d t^{dm}/m
=\sum_{r\geq1}N_r t^r/r.
\]
Each coefficient uses finitely many divisors, so the formal
rearrangement is valid. Exponentiating gives the closed-point
Euler product for \(Z(X,t)\).
Substitution \(t=q^{-s}\) gives \(\#k(x)^{-s}=t^{\deg x}\)
at each factor, hence the desired identity.
It is also an analytic identity for sufficiently large
\(\operatorname{Re}s\): finite affine covers and embeddings
bound \(N_r\) by \(Cq^{Mr}\), ensuring absolute convergence
when \(|t|<q^{-M}\).
\end{solution}

\begin{exerciseintro}
% record: H-C-05-005
\exerciseheading{Exercise C.5.5}
\addcontentsline{toc}{subsection}{Exercise C.5.5}

\begin{context}
For a scheme of finite type over \(\mathbb F_q\), put
\(N_r(X)=\#X(\mathbb F_{q^r})\) and
\[
Z(X,t)=\exp\left(\sum_{r\geq1}N_r(X)t^r/r\right)\in1+t\mathbb Q[[t]].
\]
Formal logarithm and exponential are inverse over \(\mathbb Q\), and
\(-\log(1-at)=\sum_{r\geq1}a^r t^r/r\).

Let \(C\) be smooth, projective, geometrically connected of genus \(g\).
Assume its Weil properties:
\[
Z(C,t)=\frac{P(t)}{(1-t)(1-qt)},\quad
P(t)=\prod_{i=1}^{2g}(1-\alpha_i t),\quad
P(t)=q^g t^{2g}P(1/(qt)).
\]
In particular \(P(0)=1\).
\end{context}

\begin{problem}
Prove that \(N_1,\ldots,N_g\), together with \(q,g\), determine
every \(N_r\).
\end{problem}
\end{exerciseintro}

\begin{solution}
Taking logarithms gives
\(s_r:=\sum_i\alpha_i^r=1+q^r-N_r\).
Write \(P(t)=\sum_{j=0}^{2g}c_jt^j\), \(c_0=1\).
Logarithmic differentiation, or Newton's identities, gives
\[
jc_j=-\sum_{i=1}^j s_i c_{j-i}.
\]
Thus \(s_1,\ldots,s_g\) determine \(c_1,\ldots,c_g\) recursively
over \(\mathbb Q\). Reciprocity gives
\(c_{2g-j}=q^{g-j}c_j\) for \(0\leq j\leq g\), determining
every remaining coefficient.
The rational function \(Z(C,t)\) is now determined, and its
formal logarithm recovers every \(N_r\).
If \(g=0\), no initial counts are needed: \(P=1\), so
\(N_r=1+q^r\).
\end{solution}

\begin{exerciseintro}
% record: H-C-05-006
\exerciseheading{Exercise C.5.6}
\addcontentsline{toc}{subsection}{Exercise C.5.6}

\begin{context}
For a scheme of finite type over \(\mathbb F_q\), put
\(N_r(X)=\#X(\mathbb F_{q^r})\) and
\[
Z(X,t)=\exp\left(\sum_{r\geq1}N_r(X)t^r/r\right)\in1+t\mathbb Q[[t]].
\]
Formal logarithm and exponential are inverse over \(\mathbb Q\), and
\(-\log(1-at)=\sum_{r\geq1}a^r t^r/r\).

For an elliptic curve \(E/\mathbb F_q\), Frobenius \(F\) and
its dual satisfy \(F\overline F=[q]\) and
\(F+\overline F=[a]\) for an integer \(a\), with
\(|a|\leq2\sqrt q\).
The degree formula is
\(\deg(1-F^r)=1+q^r-(F^r+\overline F^{\,r})\), where the last
sum is identified with its integer scalar.
The isogeny \(1-F^r\) is separable because its tangent map is the
identity. Its degree equals its number of geometric kernel points.
\end{context}

\begin{problem}
Use these elliptic-curve results to compute the zeta function
and verify the Weil properties for \(E\). Show explicitly that
the Hasse bound is equivalent to the required absolute values
of the two reciprocal roots.
\end{problem}
\end{exerciseintro}

\begin{solution}
The kernel of \(1-F^r\) is \(E(\mathbb F_{q^r})\), so its degree
is \(N_r\). Let \(\alpha,\beta\) be the complex roots of
\(T^2-aT+q\). The identities for Frobenius imply the same
recurrence for \(F^r+\overline F^{\,r}\) as for
\(\alpha^r+\beta^r\), with initial values \(2,a\).
Therefore
\[
N_r=1+q^r-\alpha^r-\beta^r,\qquad
Z(E,t)=\frac{(1-\alpha t)(1-\beta t)}{(1-t)(1-qt)}
=\frac{1-at+qt^2}{(1-t)(1-qt)}.
\]
This proves rationality and integrality.
Substitution verifies \(Z(E,1/(qt))=Z(E,t)\), the functional
equation with Euler number zero.
The inequality \(|a|\leq2\sqrt q\) says the real-coefficient
quadratic has either conjugate nonreal roots or a repeated real root.
Since their product is \(q\), both have modulus \(\sqrt q\).
Conversely these two moduli imply
\(|a|=|\alpha+\beta|\leq2\sqrt q\).
Thus the Hasse bound gives exactly the Riemann-hypothesis condition.
Finally the degrees of the factors are \(1,2,1\), the Betti
numbers of a complex genus-one curve; their alternating sum is zero.
\end{solution}

\begin{exerciseintro}
% record: H-C-05-007
\exerciseheading{Exercise C.5.7}
\addcontentsline{toc}{subsection}{Exercise C.5.7}

\begin{context}
For a scheme of finite type over \(\mathbb F_q\), put
\(N_r(X)=\#X(\mathbb F_{q^r})\) and
\[
Z(X,t)=\exp\left(\sum_{r\geq1}N_r(X)t^r/r\right)\in1+t\mathbb Q[[t]].
\]
Formal logarithm and exponential are inverse over \(\mathbb Q\), and
\(-\log(1-at)=\sum_{r\geq1}a^r t^r/r\).

For a smooth projective geometrically connected genus-\(g\) curve,
assume
\[
Z(C,t)=\frac{\prod_{i=1}^{2g}(1-\alpha_i t)}{(1-t)(1-qt)},
\quad N_r=1-a_r+q^r,\quad |a_r|\leq2g q^{r/2}.
\]
The functional equation makes the multiset of nonzero roots
\(\{\alpha_i\}\) invariant under \(\alpha\mapsto q/\alpha\).
\end{context}

\begin{problem}
(a) Show \(a_r=\sum_i\alpha_i^r\).

(b) Prove that the bounds \(|a_r|\leq2gq^{r/2}\) for all \(r\)
are equivalent to \(|\alpha_i|\leq\sqrt q\) for all \(i\).

(c) Use the functional equation to obtain equality of every modulus.
\end{problem}
\end{exerciseintro}

\begin{solution}
(a) Expand the formal logarithms of the numerator and denominator
and compare the coefficient of \(t^r/r\). This gives
\(N_r=1+q^r-\sum_i\alpha_i^r\), as required.

(b) If all root moduli are at most \(\sqrt q\), the triangle
inequality gives the bound immediately.
Conversely the assumed bounds make
\[
A(t)=\sum_{r\geq1}a_r t^{r-1}
\]
holomorphic for \(|t|<q^{-1/2}\).
Near zero, part (a) identifies it with the rational function
\[
\sum_{i=1}^{2g}\frac{\alpha_i}{1-\alpha_i t}.
\]
If some \(|\alpha_i|>\sqrt q\), this function has a pole inside
that disk. Terms with the same \(\alpha_i\) add a positive integer
multiple of the same nonzero pole; different values have different
poles, so cancellation is impossible.
This contradicts the holomorphic power series, by uniqueness of
analytic continuation. All moduli are therefore at most \(\sqrt q\).

(c) The paired root \(q/\alpha_i\) also satisfies that upper bound.
It gives \(|\alpha_i|\geq\sqrt q\), hence equality.
For genus zero the multiset is empty and the conclusion is vacuous.
\end{solution}

\corpuschapter{Chapter I}
\addcontentsline{toc}{chapter}{Chapter I}

\begin{exerciseintro}
\corpussection{I.1 Affine Varieties}
\addcontentsline{toc}{section}{I.1 Affine Varieties}

% record: H-I-01-001
\exerciseheading{Exercise I.1.1}
\addcontentsline{toc}{subsection}{Exercise I.1.1}

\begin{context}
Let \(k\) be algebraically closed, of any characteristic.
For an affine algebraic set \(W\), its coordinate ring is
\(A(W)=k[x,y]/I(W)\). The Nullstellensatz identifies \(I(V(f))\)
with \(\sqrt{(f)}\); an irreducible polynomial in \(k[x,y]\)
generates a prime ideal. Changes of affine coordinates induce
isomorphisms of coordinate rings. A homogeneous binary quadratic
over \(k\) is a product of two linear forms, possibly proportional.
\end{context}

\begin{problem}
(a) For \(Y=V(y-x^2)\), show that \(A(Y)\cong k[t]\).

(b) For \(Z=V(xy-1)\), prove that \(A(Z)\) is not isomorphic as
a \(k\)-algebra to \(k[t]\).

(c) Classify \(A(V(f))\) for every irreducible polynomial \(f\)
of total degree two in \(k[x,y]\): show it is one of the two
rings in (a),(b), and give a criterion deciding which.
\end{problem}
\end{exerciseintro}

\begin{solution}
(a) Substitution \(x\mapsto t,\ y\mapsto t^2\) is a surjective
homomorphism with kernel \((y-x^2)\), as division by the monic
linear polynomial in \(y\) shows. Its quotient is \(k[t]\), a
domain, so the ideal is prime and is the vanishing ideal of \(Y\).
Hence \(A(Y)\cong k[t]\).

(b) There are inverse homomorphisms
\[
k[x,y]/(xy-1)\longleftrightarrow k[t,t^{-1}],
\qquad x\longmapsto t,\quad y\longmapsto t^{-1}.
\]
Thus this quotient is a domain and equals \(A(Z)\).
It has the nonconstant unit \(x\), with inverse \(y\).
In \(k[t]\), degrees add on nonzero products, so every unit is a
nonzero constant. A \(k\)-algebra isomorphism fixes \(k\) and
preserves units, giving the contradiction.

(c) Write \(f=f_2+f_1+f_0\) by total degree.
If the two factors of \(f_2\) are linearly independent, choose
linear coordinates and scale \(f\) to obtain
\[
f=xy+ax+by+c=(x+b)(y+a)+(c-ab).
\]
Irreducibility forces \(c-ab\ne0\). Translating and then scaling
one coordinate gives the equation \(uv=1\), so \(A(V(f))\cong A(Z)\).

If those factors are proportional, linear coordinates and a scalar
put \(f\) in the form \(x^2+ax+by+c\).
Here \(b\ne0\): otherwise this quadratic in \(x\) splits over \(k\)
and is reducible. Solving for \(y\) identifies the quotient with
\(k[x]\), so \(A(V(f))\cong A(Y)\).
This criterion uses only the factorization of the leading
homogeneous part: a repeated factor gives \(A(Y)\), and two
distinct factors give \(A(Z)\). The argument uses no division
by two, and applies also in characteristic two.
\end{solution}

\begin{exerciseintro}
% record: H-I-01-002
\exerciseheading{Exercise I.1.2}
\addcontentsline{toc}{subsection}{Exercise I.1.2}

\begin{context}
Let \(k\) be an algebraically closed field. For an algebraic set \(X\subseteq\mathbb A_k^n\), write \(I(X)\) for its vanishing ideal and \(A(X)=k[x_1,\ldots,x_n]/I(X)\) for its affine coordinate ring. One may use that an affine algebraic set is irreducible exactly when its coordinate ring is a domain, and that \(\dim k[t]=1\).
\end{context}

\begin{problem}
Let
\[
Y=\{(t,t^2,t^3):t\in k\}\subseteq\mathbb A_k^3.
\]
Prove that \(Y\) is a one-dimensional affine variety, determine generators for \(I(Y)\), and identify \(A(Y)\) with a polynomial ring in one variable.
\end{problem}
\end{exerciseintro}

\begin{solution}
Consider the homomorphism
\[
\phi:k[x,y,z]\longrightarrow k[t],\qquad
x\longmapsto t,\quad y\longmapsto t^2,\quad z\longmapsto t^3.
\]
It is surjective because its image contains \(t\). The polynomials \(y-x^2\) and \(z-x^3\) lie in its kernel. Conversely,
\[
k[x,y,z]/(y-x^2,z-x^3)\cong k[x],
\]
and under this isomorphism the map induced by \(\phi\) is \(k[x]\to k[t]\), \(x\mapsto t\), which is an isomorphism. Hence
\[
\ker\phi=(y-x^2,z-x^3).
\]

Since \(k\) is infinite, a polynomial \(f\in k[x,y,z]\) vanishes at every
point of \(Y\) if and only if the one-variable polynomial
\(\phi(f)=f(t,t^2,t^3)\) vanishes at every \(t\in k\), which is equivalent to
\(\phi(f)=0\). Hence \(I(Y)=\ker\phi\).

The common zero set of these two generators consists precisely of the points \((t,t^2,t^3)\), so \(Y=V(y-x^2,z-x^3)\) is closed. Its coordinate ring is
\[
A(Y)\cong k[x,y,z]/(y-x^2,z-x^3)\cong k[t].
\]
This ring is a domain, so \(Y\) is irreducible, and its Krull dimension is \(1\). Thus \(Y\) is a one-dimensional affine variety and the displayed two polynomials generate \(I(Y)\).
\end{solution}

\begin{exerciseintro}
% record: H-I-01-003
\exerciseheading{Exercise I.1.3}
\addcontentsline{toc}{subsection}{Exercise I.1.3}

\begin{context}
Let \(k\) be algebraically closed. Closed irreducible subsets of affine space correspond to prime vanishing ideals. A finite union of irreducible closed subsets, none contained in another, is the irreducible-component decomposition.
\end{context}

\begin{problem}
In \(\mathbb A_k^3\) with coordinates \(x,y,z\), let
\[
Y=V(x^2-yz,xz-x).
\]
Express \(Y\) as the union of its irreducible components and give the prime ideal of each component.
\end{problem}
\end{exerciseintro}

\begin{solution}
At a point of \(Y\), the second equation is \(x(z-1)=0\). If \(x=0\), the first equation becomes \(yz=0\). This gives the two lines
\[
L_y=V(x,z),\qquad L_z=V(x,y).
\]
If \(z=1\), the first equation gives \(y=x^2\), producing the curve
\[
C=V(z-1,y-x^2).
\]
These cases exhaust all points, so
\[
Y=L_y\cup L_z\cup C.
\]

The corresponding coordinate rings are
\[
k[x,y,z]/(x,z)\cong k[y],\qquad
k[x,y,z]/(x,y)\cong k[z],
\]
and
\[
k[x,y,z]/(z-1,y-x^2)\cong k[x].
\]
They are domains, so the three defining ideals are prime and the three closed sets are irreducible. The lines are distinct, and neither line is contained in \(C\); likewise \(C\) is not contained in either line. Therefore these are exactly the irreducible components of \(Y\), with prime ideals
\[
I(L_y)=(x,z),\qquad I(L_z)=(x,y),\qquad I(C)=(z-1,y-x^2).
\]
\end{solution}

\begin{exerciseintro}
% record: H-I-01-004
\exerciseheading{Exercise I.1.4}
\addcontentsline{toc}{subsection}{Exercise I.1.4}

\begin{context}
Let \(k\) be algebraically closed. The affine line \(\mathbb A_k^1\) is irreducible, so any two nonempty Zariski-open subsets of it intersect. The product topology on \(\mathbb A_k^1\times\mathbb A_k^1\) has a basis consisting of products \(U\times V\) of open subsets.
\end{context}

\begin{problem}
Identify the set \(\mathbb A_k^2\) with \(\mathbb A_k^1\times\mathbb A_k^1\). Prove that the Zariski topology on \(\mathbb A_k^2\) is not the product of the Zariski topologies on the two affine-line factors.
\end{problem}
\end{exerciseintro}

\begin{solution}
The diagonal
\[
\Delta=\{(a,a):a\in k\}=V(x-y)
\]
is closed in the Zariski topology of \(\mathbb A_k^2\).

It is not closed in the product topology. Indeed, choose \(a,b\in k\) with \(a\ne b\). If the complement of \(\Delta\) were product-open, there would be nonempty Zariski-open neighborhoods \(U\) of \(a\) and \(V\) of \(b\) such that
\[
(U\times V)\cap\Delta=\varnothing.
\]
This condition is equivalent to \(U\cap V=\varnothing\), contradicting irreducibility of \(\mathbb A_k^1\). Thus \(\Delta\) is Zariski-closed in \(\mathbb A_k^2\) but not closed in the product topology, so the two topologies differ.
\end{solution}

\begin{exerciseintro}
% record: H-I-01-005
\exerciseheading{Exercise I.1.5}
\addcontentsline{toc}{subsection}{Exercise I.1.5}

\begin{context}
Let \(k\) be algebraically closed. The affine coordinate ring of
\(X\subseteq\mathbb A_k^n\) is
\[
A(X)=k[x_1,\ldots,x_n]/I(X).
\]
The Hilbert Nullstellensatz gives \(I(V(I))=\sqrt I\) for every ideal
\(I\subseteq k[x_1,\ldots,x_n]\).

The empty algebraic set and the zero ring are allowed in this correspondence.
\end{context}

\begin{problem}
Characterize the \(k\)-algebras that occur as affine coordinate rings of algebraic subsets of some affine space. Prove that these are exactly the finitely generated reduced \(k\)-algebras.
\end{problem}
\end{exerciseintro}

\begin{solution}
Suppose first that \(B=A(X)\) for an algebraic set \(X\subseteq\mathbb A_k^n\). Then \(B\) is generated over \(k\) by the residue classes of the coordinate functions. Moreover, \(I(X)\) is radical: if \(f^m\) vanishes on \(X\), then \(f\) vanishes on \(X\). Therefore the quotient \(k[x_1,\ldots,x_n]/I(X)\) has no nonzero nilpotents, so \(B\) is reduced.

Conversely, let \(B\) be a finitely generated reduced \(k\)-algebra. Choose generators and a surjection
\[
k[x_1,\ldots,x_n]\longrightarrow B
\]
with kernel \(I\). Since \(B\) is reduced, \(I\) is radical. Set \(X=V(I)\subseteq\mathbb A_k^n\). The Nullstellensatz yields
\[
I(X)=I(V(I))=\sqrt I=I.
\]
Consequently
\[
A(X)=k[x_1,\ldots,x_n]/I(X)\cong k[x_1,\ldots,x_n]/I\cong B.
\]
Thus finite generation and reducedness are both necessary and sufficient.
\end{solution}

\begin{exerciseintro}
% record: H-I-01-006
\exerciseheading{Exercise I.1.6}
\addcontentsline{toc}{subsection}{Exercise I.1.6}

\begin{context}
A nonempty topological space is irreducible if it is not the union of two proper closed subsets. Equivalently, every two nonempty open subsets intersect.
\end{context}

\begin{problem}
Prove both statements below.

\begin{enumerate}
\item Every nonempty open subset of an irreducible space is dense and is irreducible in its subspace topology.
\item If a subset is irreducible in its subspace topology, then its closure in the ambient space is irreducible.
\end{enumerate}
\end{problem}
\end{exerciseintro}

\begin{solution}
Let \(X\) be irreducible and let \(U\subseteq X\) be nonempty and open. Every nonempty open subset \(V\subseteq X\) meets \(U\), which is precisely the assertion that \(\overline U=X\).

To prove that \(U\) is irreducible, let \(U_1,U_2\) be nonempty open subsets of \(U\). Write \(U_i=U\cap V_i\) with \(V_i\) open in \(X\). Each \(U\cap V_i\) is a nonempty open subset of \(X\). Irreducibility of \(X\) therefore gives
\[
U_1\cap U_2=(U\cap V_1)\cap(U\cap V_2)\ne\varnothing.
\]
Hence \(U\) is irreducible.

Now let \(Y\subseteq X\) be irreducible and suppose
\[
\overline Y=F_1\cup F_2
\]
with \(F_1,F_2\) closed in \(\overline Y\). Intersecting with \(Y\) gives
\[
Y=(Y\cap F_1)\cup(Y\cap F_2).
\]
Since \(Y\) is irreducible, it is contained in one of the \(F_i\), say \(F_1\). The set \(F_1\) is closed in \(\overline Y\), so it contains the closure of \(Y\) inside \(\overline Y\), which is all of \(\overline Y\). Thus \(F_1=\overline Y\), proving that \(\overline Y\) is irreducible.
\end{solution}

\begin{exerciseintro}
% record: H-I-01-007
\exerciseheading{Exercise I.1.7}
\addcontentsline{toc}{subsection}{Exercise I.1.7}

\begin{context}
A topological space is noetherian if every descending chain of closed subsets stabilizes. A space is quasi-compact if every open cover has a finite subcover. One may use that a quasi-compact subset of a Hausdorff space is closed.

A finite union may have no members; that union is the empty set.
\end{context}

\begin{problem}
(a) Prove that the following conditions on a topological space \(X\) are equivalent: descending chains of closed subsets stabilize; every nonempty family of closed subsets has a minimal member; ascending chains of open subsets stabilize; every nonempty family of open subsets has a maximal member.

(b) Prove that every noetherian space is quasi-compact.

(c) Prove that every subspace of a noetherian space is noetherian.

(d) Prove that a noetherian Hausdorff space is a finite discrete space.
\end{problem}
\end{exerciseintro}

\begin{solution}
(a) Assume descending chains of closed sets stabilize. If a nonempty family of closed sets had no minimal member, one could begin with one member and repeatedly choose a strictly smaller member, producing a nonstabilizing descending chain. Thus every such family has a minimal member. Conversely, a minimal member of the family formed by the terms of a descending chain forces the chain to stabilize.

Taking complements reverses inclusion and exchanges closed sets with open sets. It therefore carries the first condition to the ascending-chain condition for open sets and the minimal-member condition for closed sets to the maximal-member condition for open sets. All four conditions are equivalent.

(b) Let \(\{U_i\}_{i\in I}\) cover \(X\), and consider the family of all finite unions of the \(U_i\). By part (a), this family has a maximal member \(V\). If \(V\ne X\), choose \(x\in X\setminus V\) and then \(U_j\) containing \(x\). The finite union \(V\cup U_j\) strictly contains \(V\), a contradiction. Hence \(V=X\), giving a finite subcover.

(c) Let \(Y\subseteq X\), and let
\[
F_1\supseteq F_2\supseteq\cdots
\]
be closed subsets of \(Y\). For each \(i\), because \(F_i\) is closed in \(Y\),
\[
F_i=Y\cap\overline{F_i}^{X}.
\]
The closures in \(X\) form a descending chain and hence stabilize. Intersecting with \(Y\) shows that the original chain stabilizes. Thus \(Y\) is noetherian.

(d) By parts (b) and (c), every subset of \(X\) is quasi-compact. Because \(X\) is Hausdorff, every quasi-compact subset is closed. Hence every subset is closed, so \(X\) is discrete. Part (b) applied to the cover of \(X\) by singleton subsets now shows that \(X\) is finite.
\end{solution}

\begin{exerciseintro}
% record: H-I-01-008
\exerciseheading{Exercise I.1.8}
\addcontentsline{toc}{subsection}{Exercise I.1.8}

\begin{context}
Let \(k\) be algebraically closed and let \(Y\) be an affine variety with coordinate ring \(A=A(Y)\). One may use that irreducible closed subsets \(Z\subseteq Y\) correspond to prime ideals \(I(Z)/I(Y)\) of \(A(Y)\), that primes minimal over a nonzero principal ideal in a finitely generated integral \(k\)-algebra have height one, and that
\[
\dim(A/\mathfrak p)+\operatorname{ht}(\mathfrak p)=\dim A.
\]
\end{context}

\begin{problem}
Let \(Y\subseteq\mathbb A_k^n\) be an affine variety of dimension \(r\), and let \(H\) be a hypersurface that does not contain \(Y\). Prove that every irreducible component of \(Y\cap H\) has dimension \(r-1\).
\end{problem}
\end{exerciseintro}

\begin{solution}
Write \(A=A(Y)\), and let \(f\in A\) be the restriction to \(Y\) of an equation defining \(H\). Since \(Y\not\subseteq H\), the element \(f\) is nonzero. If \(Y\cap H\) is empty there is nothing to prove, so assume it is nonempty. Then \(f\) is not a unit.

An irreducible component \(Z\) of \(Y\cap H=V_Y(f)\) corresponds, through the affine variety--coordinate ring correspondence, to a prime ideal \(\mathfrak p\subseteq A\) minimal over \((f)\). The ring \(A\) is a domain, and the principal ideal theorem therefore gives
\[
\operatorname{ht}(\mathfrak p)=1.
\]
The dimension formula now yields
\[
\dim Z=\dim(A/\mathfrak p)=\dim A-1=r-1.
\]
This applies to every irreducible component.
\end{solution}

\begin{exerciseintro}
% record: H-I-01-009
\exerciseheading{Exercise I.1.9}
\addcontentsline{toc}{subsection}{Exercise I.1.9}

\begin{context}
Let \(k\) be algebraically closed. One may use the hypersurface-section theorem: if an affine variety \(W\) of dimension \(d\) is not contained in the zero set \(H=V(f)\) of a nonzero polynomial, every irreducible component of \(W\cap H\) has dimension \(d-1\). The irreducible components of a finite union are the maximal members among the irreducible components of the sets in that union.
\end{context}

\begin{problem}
Suppose an ideal \(I\subseteq k[x_1,\ldots,x_n]\) can be generated by \(r\) elements. Prove that every irreducible component of \(V(I)\) has dimension at least \(n-r\).
\end{problem}
\end{exerciseintro}

\begin{solution}
Write \(I=(f_1,\ldots,f_r)\), and argue by induction on \(r\). For \(r=0\), the zero set is \(\mathbb A_k^n\), of dimension \(n\).

Let
\[
X=V(f_1,\ldots,f_{r-1})=W_1\cup\cdots\cup W_m
\]
be its irreducible-component decomposition. By induction, \(\dim W_i\ge n-r+1\) for every \(i\). Since
\[
V(I)=X\cap V(f_r)=\bigcup_{i=1}^m\bigl(W_i\cap V(f_r)\bigr),
\]
every irreducible component \(Z\) of \(V(I)\) is an irreducible component of some \(W_i\cap V(f_r)\).

If \(W_i\subseteq V(f_r)\), then \(W_i\cap V(f_r)=W_i\), so
\[
\dim Z=\dim W_i\ge n-r+1.
\]
Otherwise the hypersurface-section theorem gives
\[
\dim Z=\dim W_i-1\ge n-r.
\]
In either case \(\dim Z\ge n-r\), which proves the assertion.
\end{solution}

\begin{exerciseintro}
% record: H-I-01-010
\exerciseheading{Exercise I.1.10}
\addcontentsline{toc}{subsection}{Exercise I.1.10}

\begin{context}
The dimension of a nonempty topological space is the supremum of
lengths \(n\) of chains \(Z_0\subsetneq\cdots\subsetneq Z_n\) of
nonempty irreducible closed subsets. Put \(\dim\varnothing=-1\).
A space is noetherian if every descending sequence of closed
subsets stabilizes.
A subset of an irreducible space that is nonempty and open is
irreducible and dense; the closure of an irreducible subset is
irreducible.
\end{context}

\begin{problem}
(a) Prove \(\dim Y\leq\dim X\) for every subspace \(Y\subseteq X\).

(b) For an open cover \(X=\bigcup_i U_i\), prove
\(\dim X=\sup_i\dim U_i\).

(c) Give a space with a dense open subset of strictly smaller dimension.

(d) If \(X\) is irreducible of finite dimension and \(Y\) is closed
with \(\dim Y=\dim X\), prove \(Y=X\).

(e) Construct a noetherian space of infinite dimension.
\end{problem}
\end{exerciseintro}

\begin{solution}
(a) Given a chain of irreducible closed subsets \(Z_j\) of \(Y\),
take their closures in \(X\).
They are irreducible closed subsets of \(X\).
Since \(\overline{Z_j}^{\,X}\cap Y=Z_j\), all inclusions remain
strict. Thus every chain length in \(Y\) occurs in \(X\).

(b) Part (a) gives \(\sup_i\dim U_i\leq\dim X\).
Conversely, take a chain \(Z_0\subsetneq\cdots\subsetneq Z_n\)
in \(X\), and choose \(U_i\) meeting \(Z_0\).
Each \(Z_j\cap U_i\) is nonempty, closed in \(U_i\), and irreducible.
It is dense in \(Z_j\), so equality of any two such intersections
would imply equality of their closures \(Z_j\).
They therefore give a strict chain of length \(n\) in \(U_i\).
Taking the supremum over all chains proves the other inequality.

(c) Let \(X=\{\eta,s\}\), with closed subsets
\(\varnothing,\{s\},X\). Both \(\{s\}\) and \(X\) are irreducible,
so \(\dim X=1\). The singleton \(U=\{\eta\}\) is open and dense,
but has dimension zero.

(d) Put \(d=\dim X<\infty\).
If \(\dim Y=d\), a chain of length \(d\) exists in \(Y\), since
the supremum is over integers. Because \(Y\) is closed, its chain
members are closed in \(X\). If \(Y\ne X\), adjoining the irreducible
closed set \(X\) at the top gives a chain of length \(d+1\),
a contradiction.

(e) Take a point \(\eta\) and disjoint finite sets
\[
C_n=\{p_{n,1},\ldots,p_{n,n}\}\quad(n\geq1).
\]
Let \(X=\{\eta\}\cup\bigcup_{n\geq1}C_n\).
Declare the proper closed subsets to be finite unions of initial
segments
\(\{p_{n,1},\ldots,p_{n,j}\}\), with \(0\leq j\leq n\);
also declare \(X\) closed.
These sets are closed under finite unions and arbitrary intersections.
For the latter, if some intersected set is proper, the intersection
is contained in its finite support, and within each chain remains
an initial segment. This defines a topology.

Every proper closed set is finite. A descending chain therefore
stabilizes as soon as it leaves \(X\), by decreasing cardinalities.
Thus \(X\) is noetherian.
The closure of \(p_{n,j}\) is its initial segment, which is irreducible:
any closed subset containing that point contains the whole segment.
The nested initial segments of \(C_n\) give a chain of length
\(n-1\), for every \(n\). Consequently \(\dim X=\infty\).
\end{solution}

\begin{exerciseintro}
% record: H-I-01-011
\exerciseheading{Exercise I.1.11}
\addcontentsline{toc}{subsection}{Exercise I.1.11}

\begin{context}
Let \(k\) be algebraically closed. The vanishing ideal of a
parametrized subset of affine space is the kernel of substitution
into \(k[t]\), because \(k\) is infinite.
For a finitely generated domain over \(k\), dimension equals the
transcendence degree of its fraction field; in \(k[x,y,z]\),
\(\operatorname{ht}\mathfrak p+\dim k[x,y,z]/\mathfrak p=3\).
If an ideal \(I\) is generated by \(r\) elements, then
\(\dim_k I/\mathfrak m I\leq r\) for a maximal ideal with residue
field \(k\). The same statement holds after localization.
\end{context}

\begin{problem}
Let \(Y=\{(t^3,t^4,t^5):t\in k\}\subseteq\mathbb A^3\).
Show that \(Y\) is a closed irreducible curve and that \(I(Y)\)
is a prime ideal of height two requiring at least three generators,
even locally at the origin.
\end{problem}
\end{exerciseintro}

\begin{solution}
Let \(\phi:k[x,y,z]\to k[t]\) send \(x,y,z\) to \(t^3,t^4,t^5\).
The three polynomials
\[
f_1=y^2-xz,\qquad f_2=yz-x^3,\qquad f_3=z^2-x^2y
\]
lie in its kernel. Set \(J=(f_1,f_2,f_3)\).
Using these relations, every monomial reduces modulo \(J\) to
a \(k[x]\)-linear combination of \(1,y,z\).
Indeed each replacement of \(y^2,yz,z^2\) strictly decreases
the total exponent of \(y,z\), so the reductions terminate.
The images of these normal forms are combinations of
\[
t^{3a},\qquad t^{3a+4},\qquad t^{3a+5}\quad(a\geq0),
\]
whose exponents are pairwise distinct, as seen modulo three.
They are linearly independent in \(k[t]\). Thus the induced map
on the quotient is injective, proving \(\ker\phi=J\).
It follows that \(J\) is prime.

The common zero set of \(J\) equals the parametrized set.
If \(x=0\), its equations force \(y=z=0\).
If \(x\ne0\), put \(t=y/x\). Then
\[
t^2=z/x,\qquad t^3=yz/x^2=x,\qquad t^4=y,\qquad t^5=z.
\]
Hence every common zero has the required form.
So \(Y=V(J)\), \(I(Y)=J\), and
\(A(Y)=k[t^3,t^4,t^5]\).
Its fraction field is \(k(t)\), because \(t=t^4/t^3\), so \(Y\)
has dimension one and \(J\) has height two.

Let \(\mathfrak m=(x,y,z)\).
Since \(J\subseteq\mathfrak m^2\), we have
\(\mathfrak mJ\subseteq\mathfrak m^3\).
The quadratic initial forms of \(f_1,f_2,f_3\) are
\[
y^2-xz,\qquad yz,\qquad z^2,
\]
which are linearly independent.
Thus no nonzero constant linear combination of the \(f_i\)
belongs to \(\mathfrak mJ\).
Their classes form a basis of \(J/\mathfrak mJ\), of dimension
three. The same quotient computes
\(J_{\mathfrak m}/\mathfrak m_{\mathfrak m}J_{\mathfrak m}\),
so neither \(J\) nor its localization can be generated by two elements.
This gives the claimed failure of a local complete intersection
at the origin.
\end{solution}

\begin{exerciseintro}
% record: H-I-01-012
\exerciseheading{Exercise I.1.12}
\addcontentsline{toc}{subsection}{Exercise I.1.12}

\begin{context}
Work with real points and the real Zariski topology.
A nonempty space is irreducible if it is not the union of two
proper closed subsets. A monic polynomial in \(x\) over
\(\mathbb R[y]\) that factors in \(\mathbb R(y)[x]\) has its monic
factors in \(\mathbb R[y][x]\), by Gauss's lemma.
\end{context}

\begin{problem}
Find an irreducible polynomial in \(\mathbb R[x,y]\) whose real
zero set is nonempty and reducible.
\end{problem}
\end{exerciseintro}

\begin{solution}
Take
\[
f=x^2+(y^2-1)^2.
\]
If it were reducible, its degree two and monicity in \(x\),
together with Gauss's lemma, would give a polynomial root
\(a(y)\in\mathbb R[y]\). That root would satisfy
\[
a(y)^2=-(y^2-1)^2.
\]
At \(y=0\) this says \(a(0)^2=-1\), which is impossible over
\(\mathbb R\). A factor independent of \(x\) is also impossible
because \(f\) is monic in \(x\). Thus \(f\) is irreducible.

For real \(x,y\), both summands of \(f\) are nonnegative.
Their sum vanishes precisely when \(x=0\) and \(y^2=1\).
Therefore
\[
V_{\mathbb R}(f)=\{(0,1),(0,-1)\}.
\]
Each singleton is algebraic and closed, so this set is the
union of two proper nonempty closed subsets and is reducible.
\end{solution}

\begin{exerciseintro}
\corpussection{I.2 Projective Varieties}
\addcontentsline{toc}{section}{I.2 Projective Varieties}

% record: H-I-02-001
\exerciseheading{Exercise I.2.1}
\addcontentsline{toc}{subsection}{Exercise I.2.1}

\begin{context}
Let \(k\) be algebraically closed and \(S=k[x_0,\ldots,x_n]\) with its standard grading. For a homogeneous ideal \(I\), its affine zero set is a cone, and every nonzero point of that cone determines a point of the projective zero set \(Z_+(I)\subseteq\mathbb P_k^n\). One may use the affine Hilbert Nullstellensatz.
\end{context}

\begin{problem}
Let \(I\subseteq S\) be homogeneous, and let \(f\in S\) be homogeneous of positive degree. Assume \(f\) vanishes at every point of \(Z_+(I)\). Prove that \(f^q\in I\) for some integer \(q>0\).
\end{problem}
\end{exerciseintro}

\begin{solution}
Let \(C=V_{\mathbb A^{n+1}}(I)\) be the affine cone defined by \(I\). If \(a\in C\) is nonzero, its projective class \([a]\) belongs to \(Z_+(I)\), so the hypothesis gives \(f(a)=0\). At the origin, \(f(0)=0\) because \(f\) has positive degree. Hence \(f\) vanishes on all of \(C\).

The affine Nullstellensatz gives
\[
f\in I(C)=\sqrt I.
\]
By the definition of the radical, there is an integer \(q>0\) such that \(f^q\in I\). This is the homogeneous Nullstellensatz.
\end{solution}

\begin{exerciseintro}
% record: H-I-02-002
\exerciseheading{Exercise I.2.2}
\addcontentsline{toc}{subsection}{Exercise I.2.2}

\begin{context}
Let \(k\) be algebraically closed, let \(S=k[x_0,\ldots,x_n]\), and write \(S_+=(x_0,\ldots,x_n)\) for the irrelevant ideal. Let \(S_d\) denote the vector space of homogeneous forms of degree \(d\). One may use the homogeneous Nullstellensatz.
\end{context}

\begin{problem}
For a homogeneous ideal \(I\subseteq S\), prove that the following conditions are equivalent:

\begin{enumerate}
\item its projective zero set is empty;
\item \(\sqrt I\) is either \(S\) or \(S_+\);
\item \(I\) contains \(S_d\) for some \(d>0\).
\end{enumerate}
\end{problem}
\end{exerciseintro}

\begin{solution}
Assume first that the projective zero set is empty. If \(I=S\), then \(\sqrt I=S\). Otherwise \(I\) is proper, and homogeneity gives \(I\subseteq S_+\), hence \(\sqrt I\subseteq S_+\). By the homogeneous Nullstellensatz, each homogeneous polynomial of positive degree that vanishes on the empty projective zero set belongs to \(\sqrt I\). In particular, every \(x_i\) belongs to \(\sqrt I\), so
\[
S_+\subseteq\sqrt I\subseteq S_+.
\]
Thus \(\sqrt I=S_+\), proving the second condition.

Suppose next that \(\sqrt I=S\) or \(S_+\). In the first case \(I=S\), so it contains \(S_d\) for every \(d\). In the second case, for each \(i\) there is an integer \(N_i>0\) with \(x_i^{N_i}\in I\). Choose
\[
d>\sum_{i=0}^n(N_i-1).
\]
Every monomial of degree \(d\) is divisible by some \(x_i^{N_i}\), and hence belongs to \(I\). Since these monomials span \(S_d\), we have \(S_d\subseteq I\).

Finally, suppose \(S_d\subseteq I\) for some \(d>0\). At any projective point at least one coordinate, say \(x_i\), is nonzero. But \(x_i^d\in S_d\subseteq I\), so that point cannot be a common zero of \(I\). Therefore the projective zero set is empty, completing the cycle of implications.
\end{solution}

\begin{exerciseintro}
% record: H-I-02-003
\exerciseheading{Exercise I.2.3}
\addcontentsline{toc}{subsection}{Exercise I.2.3}

\begin{context}
Let \(k\) be algebraically closed and \(S=k[x_0,\ldots,x_n]\),
graded by total degree. For homogeneous polynomials \(T\), let
\(V_+(T)\) be their common projective zero set. For a subset
\(Y\subseteq\mathbb P^n\), let \(I(Y)\) be the homogeneous ideal
generated by homogeneous polynomials vanishing on \(Y\).
The irrelevant ideal is \(S_+=(x_0,\ldots,x_n)\).
We may use the affine Nullstellensatz and noetherianity of \(S\).
An ideal generated by homogeneous elements contains every
homogeneous component of each of its elements; radicals and
intersections of homogeneous ideals are homogeneous.
\end{context}

\begin{problem}
(a) Show \(T_1\subseteq T_2\) implies \(V_+(T_1)\supseteq V_+(T_2)\).

(b) Show \(Y_1\subseteq Y_2\) implies \(I(Y_1)\supseteq I(Y_2)\).

(c) Prove \(I(Y_1\cup Y_2)=I(Y_1)\cap I(Y_2)\).

(d) If \(\mathfrak a\) is homogeneous and \(V_+(\mathfrak a)\ne
\varnothing\), prove \(I(V_+(\mathfrak a))=\sqrt{\mathfrak a}\).

(e) Prove \(V_+(I(Y))=\overline Y\) for every projective subset \(Y\).
\end{problem}
\end{exerciseintro}

\begin{solution}
(a) Every point annihilating every element of \(T_2\) annihilates
every element of its subset \(T_1\).

(b) A homogeneous polynomial vanishing on \(Y_2\) vanishes on \(Y_1\);
passing to the ideals they generate gives the asserted inclusion.

(c) In each degree, vanishing on the union is equivalent to
vanishing on both subsets. The two homogeneous ideals therefore
have identical components in every degree and are equal.
This also covers the empty subset, whose ideal is \(S\).

(d) Put \(Y=V_+(\mathfrak a)\).
Its nonemptiness implies that \(\mathfrak a\) contains no nonzero
constant. Its affine zero set in \(\mathbb A^{n+1}\) is exactly
the origin together with all nonzero representatives of points of \(Y\).
Every positive-degree homogeneous polynomial vanishing on \(Y\)
therefore vanishes on this affine zero set; the affine Nullstellensatz
puts it in \(\sqrt{\mathfrak a}\).
The only degree-zero polynomial vanishing on nonempty \(Y\) is zero.
Thus \(I(Y)\subseteq\sqrt{\mathfrak a}\).
Conversely every homogeneous element whose power lies in
\(\mathfrak a\) vanishes on \(Y\). Since the radical is homogeneous,
its homogeneous components generate it, giving
\(\sqrt{\mathfrak a}\subseteq I(Y)\).

(e) The set \(V_+(I(Y))\) is closed and contains \(Y\), so it contains
\(\overline Y\).
If \(W=V_+(\mathfrak b)\) is any closed set containing \(Y\),
each homogeneous element of \(\mathfrak b\) vanishes on \(Y\).
Thus \(\mathfrak b\subseteq I(Y)\), and part (a) gives
\(V_+(I(Y))\subseteq W\).
Intersecting all such \(W\) proves the reverse inclusion.
\end{solution}

\begin{exerciseintro}
% record: H-I-02-004
\exerciseheading{Exercise I.2.4}
\addcontentsline{toc}{subsection}{Exercise I.2.4}

\begin{context}
Let \(k\) be algebraically closed and \(S=k[x_0,\ldots,x_n]\),
graded by total degree. For homogeneous polynomials \(T\), let
\(V_+(T)\) be their common projective zero set. For a subset
\(Y\subseteq\mathbb P^n\), let \(I(Y)\) be the homogeneous ideal
generated by homogeneous polynomials vanishing on \(Y\).
The irrelevant ideal is \(S_+=(x_0,\ldots,x_n)\).
We may use the affine Nullstellensatz and noetherianity of \(S\).
An ideal generated by homogeneous elements contains every
homogeneous component of each of its elements; radicals and
intersections of homogeneous ideals are homogeneous.

For homogeneous \(\mathfrak a\) with nonempty projective zero set,
\(I(V_+(\mathfrak a))=\sqrt{\mathfrak a}\).
The projective zero set is empty precisely when
\(\sqrt{\mathfrak a}=S\) or \(S_+\).
For every subset \(Y\), \(V_+(I(Y))=\overline Y\).
Primality of a homogeneous ideal can be checked on products of
homogeneous elements.
\end{context}

\begin{problem}
(a) Establish the inclusion-reversing correspondence between
projective algebraic sets and homogeneous radical ideals other
than \(S_+\), including the empty set.

(b) Prove that an algebraic set \(Y\) is irreducible exactly when
\(I(Y)\) is prime.

(c) Prove that \(\mathbb P^n\) is irreducible.
\end{problem}
\end{exerciseintro}

\begin{solution}
(a) For nonempty algebraic \(Y\), its vanishing ideal is homogeneous
and radical: if a power vanishes, the polynomial itself vanishes,
and this can be applied to homogeneous components.
It is not \(S_+\), since no projective point annihilates all
coordinates. The two identities in the context show that \(I\)
and \(V_+\) are inverse on nonempty algebraic sets and homogeneous
radical ideals other than \(S,S_+\).
The empty set has ideal \(S\), and \(V_+(S)=\varnothing\).
This extends the bijection as stated. Both operations reverse
inclusion directly from their definitions.

(b) Suppose first that \(Y\) is irreducible.
Its ideal is proper because \(Y\ne\varnothing\).
If homogeneous \(f,g\) satisfy \(fg\in I(Y)\), then
\[
Y=(Y\cap V_+(f))\cup(Y\cap V_+(g)).
\]
Irreducibility forces one term to be all of \(Y\); hence
\(f\) or \(g\) lies in \(I(Y)\). The homogeneous primality test
therefore makes \(I(Y)\) prime.

Conversely suppose \(I(Y)\) is prime. In particular \(Y\ne\varnothing\).
If \(Y=Y_1\cup Y_2\) with both \(Y_i\) proper and closed in \(Y\),
then they are algebraic, since \(Y\) is closed.
The correspondence in (a) gives
\(I(Y)\subsetneq I(Y_i)\).
Choose homogeneous \(f_i\in I(Y_i)\setminus I(Y)\).
Their product vanishes on \(Y_1\cup Y_2\) but neither factor
belongs to \(I(Y)\), contradicting primality.

(c) A homogeneous polynomial vanishing on all projective points
also vanishes on every nonzero vector; if it has positive degree
it vanishes at zero as well. A polynomial vanishing on all
of \(k^{n+1}\) is zero, by induction in the variables over the
infinite field \(k\).
Nonzero constants do not vanish on projective space.
Thus \(I(\mathbb P^n)=(0)\), which is prime, and (b) applies.
\end{solution}

\begin{exerciseintro}
% record: H-I-02-005
\exerciseheading{Exercise I.2.5}
\addcontentsline{toc}{subsection}{Exercise I.2.5}

\begin{context}
Let \(k\) be algebraically closed and \(S=k[x_0,\ldots,x_n]\),
graded by total degree. For homogeneous polynomials \(T\), let
\(V_+(T)\) be their common projective zero set. For a subset
\(Y\subseteq\mathbb P^n\), let \(I(Y)\) be the homogeneous ideal
generated by homogeneous polynomials vanishing on \(Y\).
The irrelevant ideal is \(S_+=(x_0,\ldots,x_n)\).
We may use the affine Nullstellensatz and noetherianity of \(S\).
An ideal generated by homogeneous elements contains every
homogeneous component of each of its elements; radicals and
intersections of homogeneous ideals are homogeneous.

Every projective algebraic set \(Y\) satisfies \(V_+(I(Y))=Y\).
A topological space is noetherian when descending chains of
closed subsets stabilize.
\end{context}

\begin{problem}
(a) Prove that projective space is noetherian.

(b) Prove that every projective algebraic set is uniquely a
finite union of irreducible closed subsets, none contained in
another. Include the empty set.
\end{problem}
\end{exerciseintro}

\begin{solution}
(a) A descending chain of closed sets \(Y_1\supseteq Y_2\supseteq
\cdots\) gives an ascending chain \(I(Y_1)\subseteq I(Y_2)\subseteq
\cdots\). It stabilizes because \(S\) is noetherian.
Applying \(V_+\) shows that the original chain stabilizes.

(b) Suppose some nonempty closed set has no finite irreducible
decomposition. By the descending chain condition, choose a
minimal such set \(Y\). It cannot be irreducible, so it is the
union of two proper closed subsets. By minimality each of those
is a finite union of irreducible closed sets, giving such a
decomposition of \(Y\), a contradiction.
In a finite decomposition, remove any member contained in another.

For uniqueness, let \(Y=\bigcup_{i=1}^r A_i=\bigcup_{j=1}^s B_j\)
be two decompositions without containments.
The irreducible set \(A_i\) is covered by the finitely many
closed subsets \(A_i\cap B_j\); it must lie in some \(B_j\).
In turn that \(B_j\) lies in some \(A_\ell\).
Then \(A_i\subseteq A_\ell\), so the no-containment condition
forces \(\ell=i\), and \(A_i=B_j\).
Interchanging the two decompositions proves equality of the
collections. For \(Y=\varnothing\), the unique collection is
the empty one.
\end{solution}

\begin{exerciseintro}
% record: H-I-02-006
\exerciseheading{Exercise I.2.6}
\addcontentsline{toc}{subsection}{Exercise I.2.6}

\begin{context}
Let \(k\) be algebraically closed and \(S=k[x_0,\ldots,x_n]\),
graded by total degree. For homogeneous polynomials \(T\), let
\(V_+(T)\) be their common projective zero set. For a subset
\(Y\subseteq\mathbb P^n\), let \(I(Y)\) be the homogeneous ideal
generated by homogeneous polynomials vanishing on \(Y\).
The irrelevant ideal is \(S_+=(x_0,\ldots,x_n)\).
We may use the affine Nullstellensatz and noetherianity of \(S\).
An ideal generated by homogeneous elements contains every
homogeneous component of each of its elements; radicals and
intersections of homogeneous ideals are homogeneous.

The standard chart \(U_i=\{x_i\ne0\}\) is affine space with
coordinates \(x_j/x_i\), \(j\ne i\).
For a finitely generated \(k\)-domain, dimension equals the
transcendence degree of its fraction field.
The dimension of a topological space covered by open subsets is
the supremum of their dimensions. A nonempty open subset of an
irreducible space is dense.
\end{context}

\begin{problem}
Let \(Y\subseteq\mathbb P^n\) be a projective variety and
\(R=S/I(Y)\). Prove \(\dim R=\dim Y+1\).
For every nonempty affine chart \(Y_i=Y\cap U_i\), identify
its coordinate ring with \((R_{x_i})_0\), and prove
\(\dim Y_i=\dim Y\).
\end{problem}
\end{exerciseintro}

\begin{solution}
Choose \(i\) such that \(Y_i\ne\varnothing\); then \(x_i\ne0\)
in the graded domain \(R\). The degree-zero subring
\[
B=(R_{x_i})_0=k[x_0/x_i,\ldots,x_n/x_i]
\]
is generated by the indicated ratios, omitting \(x_i/x_i=1\).
Its relations are exactly the polynomials vanishing on the
affine algebraic set \(Y_i\).
Indeed dehomogenization of an element of \(I(Y)\) vanishes there.
Conversely, homogenize a polynomial relation that vanishes on
\(Y_i\). The resulting homogeneous polynomial vanishes on the
dense subset \(Y_i\) of \(Y\), hence on \(Y\), so belongs to
\(I(Y)\). This proves \(B=A(Y_i)\).

Every homogeneous element of \(R_{x_i}\) of degree \(e\in\mathbb Z\)
is \(x_i^e\) times an element of \(B\).
The grading is a direct sum, so no nontrivial Laurent polynomial
in \(x_i\) with coefficients in \(B\) can vanish.
Consequently
\[
R_{x_i}\cong B[x_i,x_i^{-1}],\qquad
\operatorname{Frac}(R)=\operatorname{Frac}(B)(x_i),
\]
where \(x_i\) is transcendental over \(\operatorname{Frac}(B)\).
The equality of fraction fields also uses that localization of
a domain does not change its fraction field.
Taking transcendence degrees gives
\(\dim R=\dim B+1=\dim Y_i+1\).
Thus every nonempty chart has the same dimension \(\dim R-1\).
Since those charts form an open cover of \(Y\), their supremum
equals \(\dim Y\), proving both claims.
\end{solution}

\begin{exerciseintro}
% record: H-I-02-007
\exerciseheading{Exercise I.2.7}
\addcontentsline{toc}{subsection}{Exercise I.2.7}

\begin{context}
Let \(k\) be algebraically closed and \(S=k[x_0,\ldots,x_n]\),
graded by total degree. For homogeneous polynomials \(T\), let
\(V_+(T)\) be their common projective zero set. For a subset
\(Y\subseteq\mathbb P^n\), let \(I(Y)\) be the homogeneous ideal
generated by homogeneous polynomials vanishing on \(Y\).
The irrelevant ideal is \(S_+=(x_0,\ldots,x_n)\).
We may use the affine Nullstellensatz and noetherianity of \(S\).
An ideal generated by homogeneous elements contains every
homogeneous component of each of its elements; radicals and
intersections of homogeneous ideals are homogeneous.

For a projective variety \(Z\), every nonempty standard affine
chart has dimension \(\dim S(Z)-1=\dim Z\).
Every nonempty open subset of an affine variety has the same
dimension as that variety.
The closure of an irreducible set is irreducible, and subspaces
cannot have dimension greater than the ambient space.
\end{context}

\begin{problem}
(a) Prove \(\dim\mathbb P^n=n\).

(b) If \(Y\subseteq\mathbb P^n\) is a quasi-projective variety,
prove \(\dim Y=\dim\overline Y\).
\end{problem}
\end{exerciseintro}

\begin{solution}
(a) The homogeneous coordinate ring of projective space is
\(k[x_0,\ldots,x_n]\), of dimension \(n+1\).
The chart dimension formula therefore gives \(\dim\mathbb P^n=n\).

(b) Put \(Z=\overline Y\). It is a projective variety, and \(Y\)
is a nonempty open subset of \(Z\).
Choose a standard chart meeting \(Y\). Then \(Y\cap U_i\) is
a nonempty open subset of the affine variety \(Z\cap U_i\).
The two dimension statements in the context give
\[
\dim(Y\cap U_i)=\dim(Z\cap U_i)=\dim Z.
\]
Subspace monotonicity now gives
\(\dim Z\leq\dim Y\leq\dim Z\), proving equality.
\end{solution}

\begin{exerciseintro}
% record: H-I-02-008
\exerciseheading{Exercise I.2.8}
\addcontentsline{toc}{subsection}{Exercise I.2.8}

\begin{context}
Let \(k\) be algebraically closed and \(S=k[x_0,\ldots,x_n]\),
graded by total degree. For homogeneous polynomials \(T\), let
\(V_+(T)\) be their common projective zero set. For a subset
\(Y\subseteq\mathbb P^n\), let \(I(Y)\) be the homogeneous ideal
generated by homogeneous polynomials vanishing on \(Y\).
The irrelevant ideal is \(S_+=(x_0,\ldots,x_n)\).
We may use the affine Nullstellensatz and noetherianity of \(S\).
An ideal generated by homogeneous elements contains every
homogeneous component of each of its elements; radicals and
intersections of homogeneous ideals are homogeneous.

For a projective variety \(Y\), \(\dim S/I(Y)=\dim Y+1\).
For a prime \(\mathfrak p\subset S\),
\(\operatorname{ht}\mathfrak p+\dim S/\mathfrak p=n+1\).
The ring \(S\) is a unique factorization domain, in which each
height-one prime is generated by an irreducible element.
The homogeneous radical ideals with empty projective zero set
are \(S,S_+\).
\end{context}

\begin{problem}
For \(n\geq1\), prove that a projective variety
\(Y\subseteq\mathbb P^n\) has dimension \(n-1\) exactly when
it is the zero set of one irreducible homogeneous polynomial
of positive degree.
\end{problem}
\end{exerciseintro}

\begin{solution}
Suppose \(\dim Y=n-1\).
Then \(\dim S/I(Y)=n\), so the prime \(I(Y)\) has height one.
It is therefore \((f)\) for an irreducible polynomial \(f\).
Because this ideal is homogeneous, the top homogeneous component
\(f_d\) belongs to \((f)\).
It has the same total degree as \(f\), so divisibility forces
\(f_d=cf\) for some nonzero scalar \(c\).
Thus \(f\) is homogeneous. It is nonconstant because \(I(Y)\)
is proper, and \(Y=V_+(f)\).

Conversely, let \(f\) be irreducible, homogeneous, and of positive
degree. Its ideal is a height-one prime, with quotient dimension \(n\).
It is neither \(S\) nor \(S_+\), the latter having quotient
dimension zero. Thus \(V_+(f)\) is nonempty.
The projective Nullstellensatz gives \(I(V_+(f))=(f)\);
its primality makes this zero set irreducible, and the dimension
formula gives dimension \(n-1\).
\end{solution}

\begin{exerciseintro}
% record: H-I-02-009
\exerciseheading{Exercise I.2.9}
\addcontentsline{toc}{subsection}{Exercise I.2.9}

\begin{context}
Work over an algebraically closed field \(k\).
For a projective algebraic set \(Y\), \(I(Y)\) denotes its
homogeneous vanishing ideal. For homogeneous ideals with nonempty
zero set, \(I(V_+(\mathfrak a))=\sqrt{\mathfrak a}\).
An algebraic set is irreducible exactly when its vanishing ideal
is prime. For a projective variety,
\(\dim Y=\dim k[x_0,\ldots,x_n]/I(Y)-1\).
Standard charts are affine spaces, and their topology is obtained
by dehomogenizing polynomials.

For \(f\in k[x_1,\ldots,x_n]\) of degree \(d\), write
\(f^h=x_0^d f(x_1/x_0,\ldots,x_n/x_0)\).
The chart \(x_0\ne0\) identifies affine space with an open
subset of projective space.
\end{context}

\begin{problem}
(a) For an affine variety \(Y\subseteq\mathbb A^n\), prove
that the homogeneous ideal of its projective closure is generated
by \(f^h\) for all \(f\in I(Y)\).

(b) Determine the ideals of the affine twisted cubic
\(\{(t,t^2,t^3):t\in k\}\) and its projective closure.
Show why homogenizing just a chosen generating set of the
affine ideal need not generate the closure's ideal.
\end{problem}
\end{exerciseintro}

\begin{solution}
(a) Every \(f^h\) vanishes on \(Y\) and hence on its closure,
since its projective zero set is closed.
Conversely, let \(F\) be homogeneous of degree \(e\) and vanish
on the closure. Then \(f=F(1,x_1,\ldots,x_n)\) belongs to \(I(Y)\).
If \(F\ne0\), its dehomogenization is nonzero, of some degree \(d\leq e\),
and
\[
F=x_0^{e-d}f^h.
\]
Thus every homogeneous element of the closure's ideal lies in
the ideal generated by all these homogenizations, proving equality.

(b) Write the projective coordinates as \(w,x,y,z\), with \(w=1\)
the affine chart. Substitution gives the affine ideal
\[
I_{\mathrm{aff}}=(y-x^2,z-x^3);
\]
division by these two relations identifies the quotient with \(k[x]\).
The projective ideal is
\[
J=(wy-x^2,\ wz-xy,\ xz-y^2).
\tag{1}
\]
To prove this assertion directly, label \(w,x,y,z\) by indices
\(0,1,2,3\), and map them to
\(s^3,s^2t,st^2,t^3\), respectively.
The relations in (1) replace any product of two indices whose
difference is at least two by the product of the two closer
indices, preserving their sum.
Each replacement strictly decreases the sum of squares of all
indices in the monomial, so the process terminates.
The resulting monomial uses at most two adjacent indices and is
uniquely determined by its degree \(d\) and index sum \(j\).
Its image is \(s^{3d-j}t^j\).
These images are distinct monomials, hence linearly independent.
Consequently \(J\) is exactly the substitution kernel and is prime.

For homogeneous \(F\) of degree \(d\), vanishing at every
\((1,t,t^2,t^3)\) is equivalent to the zero polynomial
\(F(s^3,s^2t,st^2,t^3)\): the latter is homogeneous, and its
specialization \(s=1\) vanishes for every \(t\) in an infinite field.
Thus \(J\) is the homogeneous ideal of the affine set and its closure.

Homogenizing the two displayed affine generators yields only
\[
wy-x^2,\qquad w^2z-x^3.
\]
Both vanish along the whole line \(w=x=0\). But the point
\([0:0:1:0]\) on this line does not satisfy \(xz-y^2=0\).
These two homogenizations therefore do not generate \(J\), and
even their projective zero set is too large.
\end{solution}

\begin{exerciseintro}
% record: H-I-02-010
\exerciseheading{Exercise I.2.10}
\addcontentsline{toc}{subsection}{Exercise I.2.10}

\begin{context}
Work over an algebraically closed field \(k\).
For a projective algebraic set \(Y\), \(I(Y)\) denotes its
homogeneous vanishing ideal. For homogeneous ideals with nonempty
zero set, \(I(V_+(\mathfrak a))=\sqrt{\mathfrak a}\).
An algebraic set is irreducible exactly when its vanishing ideal
is prime. For a projective variety,
\(\dim Y=\dim k[x_0,\ldots,x_n]/I(Y)-1\).
Standard charts are affine spaces, and their topology is obtained
by dehomogenizing polynomials.

An affine algebraic set is irreducible exactly when its vanishing
ideal is prime, and its dimension is that of its coordinate ring.
In a finite union of closed subsets, dimension is the maximum of
their dimensions. A polynomial vanishing at every scalar
multiple of a vector over an infinite field has each homogeneous
component vanishing at that vector.
\end{context}

\begin{problem}
For nonempty projective algebraic \(Y\subseteq\mathbb P^n\),
let \(C(Y)\subseteq\mathbb A^{n+1}\) consist of zero and all
nonzero vectors representing points of \(Y\).

(a) Show that it is algebraic and has affine ideal \(I(Y)\).

(b) Show that \(C(Y)\) is irreducible exactly when \(Y\) is.

(c) Prove \(\dim C(Y)=\dim Y+1\), allowing reducible \(Y\).
\end{problem}
\end{exerciseintro}

\begin{solution}
(a) The affine zero set of the homogeneous ideal \(I(Y)\)
is exactly \(C(Y)\): away from zero its equations express membership
in \(Y\), and at zero every positive-degree equation vanishes.
There is no nonzero constant in \(I(Y)\) because \(Y\ne\varnothing\).
Thus the cone is closed.
If \(f=\sum f_d\) vanishes on \(C(Y)\), fix a representative
vector \(a\) of a point of \(Y\) and evaluate
\[
0=f(\lambda a)=\sum_d\lambda^d f_d(a)\quad(\lambda\in k).
\]
Since \(k\) is infinite, all coefficients vanish.
Therefore every \(f_d\) belongs to \(I(Y)\), so \(f\in I(Y)\).
The reverse inclusion is immediate, proving equality of ideals.

(b) The affine and projective irreducibility criteria now both
amount to primality of the same ideal \(I(Y)\).

(c) Write \(Y=Y_1\cup\cdots\cup Y_r\) as its finite nonempty
irreducible components. Their cones are irreducible by (b);
no one contains another, because such a containment would imply
the corresponding containment after projectivizing nonzero vectors.
Hence they are exactly the components of \(C(Y)\).
For each \(i\), the coordinate-ring dimension formula gives
\[
\dim C(Y_i)=\dim S/I(Y_i)=\dim Y_i+1.
\]
Taking the maximum over \(i\) proves the assertion.
\end{solution}

\begin{exerciseintro}
% record: H-I-02-011
\exerciseheading{Exercise I.2.11}
\addcontentsline{toc}{subsection}{Exercise I.2.11}

\begin{context}
Work over an algebraically closed field \(k\).
For a projective algebraic set \(Y\), \(I(Y)\) denotes its
homogeneous vanishing ideal. For homogeneous ideals with nonempty
zero set, \(I(V_+(\mathfrak a))=\sqrt{\mathfrak a}\).
An algebraic set is irreducible exactly when its vanishing ideal
is prime. For a projective variety,
\(\dim Y=\dim k[x_0,\ldots,x_n]/I(Y)-1\).
Standard charts are affine spaces, and their topology is obtained
by dehomogenizing polynomials.
\end{context}

\begin{problem}
(a) Prove that a projective variety is an intersection of hyperplanes
if and only if its homogeneous ideal is generated by linear forms.

(b) If such a variety has dimension \(r\) in \(\mathbb P^n\),
show that its ideal is minimally generated by \(n-r\) linear forms.

(c) If linear varieties \(Y,Z\subseteq\mathbb P^n\) have dimensions
\(r,s\), show that \(r+s\geq n\) implies \(Y\cap Z\ne\varnothing\).
If the intersection is nonempty, show it is linear and has
dimension at least \(r+s-n\).
\end{problem}
\end{exerciseintro}

\begin{solution}
(a) If \(I(Y)\) is generated by linear forms, its zero set is the
intersection of their hyperplanes.
Conversely, let \(Y\) be such an intersection, and choose a basis
\(\ell_1,\ldots,\ell_c\) of the span of its defining linear forms.
An invertible linear coordinate change makes them \(x_0,\ldots,x_{c-1}\).
Their ideal is prime, since the quotient is the polynomial ring
in the remaining \(n+1-c\) variables.
Because \(Y\) is nonempty, the projective Nullstellensatz gives
\(I(Y)=(\ell_1,\ldots,\ell_c)\).

(b) The same quotient shows
\(\dim Y=(n+1-c)-1=n-c\), so \(c=n-r\).
These \(c\) independent linear forms generate the ideal minimally:
any homogeneous generating collection must span its degree-one
part, of dimension \(c\).
Even allowing nonhomogeneous generators cannot lower the number,
because reduction modulo \(S_+I(Y)\) leaves precisely that
\(c\)-dimensional vector space.

(c) Let \(V,W\subseteq k^{n+1}\) be the vector subspaces defining
\(Y,Z\). Their dimensions are \(r+1,s+1\), and linear algebra gives
\[
\dim(V\cap W)\geq(r+1)+(s+1)-(n+1)=r+s-n+1.
\]
If \(r+s\geq n\), this is positive, so its projectivization is
nonempty. In every nonempty case
\[
Y\cap Z=\mathbb P(V\cap W),\qquad
\dim(Y\cap Z)=\dim(V\cap W)-1\geq r+s-n.
\]
The intersection is therefore a linear variety with the claimed bound.
\end{solution}

\begin{exerciseintro}
% record: H-I-02-012
\exerciseheading{Exercise I.2.12}
\addcontentsline{toc}{subsection}{Exercise I.2.12}

\begin{context}
Work over an algebraically closed field \(k\).
For a projective algebraic set \(Y\), \(I(Y)\) denotes its
homogeneous vanishing ideal. For homogeneous ideals with nonempty
zero set, \(I(V_+(\mathfrak a))=\sqrt{\mathfrak a}\).
An algebraic set is irreducible exactly when its vanishing ideal
is prime. For a projective variety,
\(\dim Y=\dim k[x_0,\ldots,x_n]/I(Y)-1\).
Standard charts are affine spaces, and their topology is obtained
by dehomogenizing polynomials.
\end{context}

\begin{problem}
Let \(n,d>0\), let \(\alpha\) range over exponent vectors of
total degree \(d\) in \(n+1\) variables, and put
\(N=\binom{n+d}{d}-1\). Define
\[
v_d:\mathbb P^n\to\mathbb P^N,\qquad [a]\longmapsto[a^\alpha]_\alpha,
\]
and \(\theta:k[y_\alpha]\to k[x_0,\ldots,x_n]\) by
\(\theta(y_\alpha)=x^\alpha\).

(a) Show that \(\mathfrak a=\ker\theta\) is homogeneous and prime,
with nonempty irreducible projective zero set.

(b) Prove \(\operatorname{im}v_d=V_+(\mathfrak a)\).

(c) Prove that \(v_d\) is a homeomorphism onto its image.

(d) Identify the projective twisted cubic with \(v_3(\mathbb P^1)\).
\end{problem}
\end{exerciseintro}

\begin{solution}
(a) The map multiplies degrees by \(d\), so homogeneous components
of different degrees cannot cancel in its image.
Its kernel is therefore homogeneous. The image is a subring of
a domain, so the kernel is prime.
Substitution at any nonzero vector gives a projective common zero;
at least one pure power \(a_i^d\) is nonzero.
Thus its projective zero set is nonempty and irreducible.

(b) Substitution immediately gives \(\operatorname{im}v_d
\subseteq V_+(\mathfrak a)\).
For the reverse inclusion, write \(Y_i=y_{de_i}\) for a pure-power
coordinate, where \(e_i\) is the \(i\)-th unit vector.
The kernel contains the relations
\[
y_\alpha^d=\prod_iY_i^{\alpha_i}.
\tag{1}
\]
Thus at a projective common zero not all \(Y_i\) can vanish.
Choose \(i\) with \(Y_i\ne0\), and normalize all coordinates by
dividing by \(Y_i\), so that \(Y_i=1\).
Put \(a_i=1\) and \(a_j=y_{(d-1)e_i+e_j}\) for \(j\ne i\).
There are further kernel relations
\[
y_\alpha Y_i^{d-1}
=\prod_j y_{(d-1)e_i+e_j}^{\alpha_j},
\tag{2}
\]
because both monomials map to \(x_i^{d(d-1)}x^\alpha\).
After normalization (2) gives \(y_\alpha=\prod_j a_j^{\alpha_j}\).
The given point is therefore \(v_d([a_0:\cdots:a_n])\).
No extraction of \(d\)-th roots was used, so this also works
when the characteristic divides \(d\).

(c) Pulling back a homogeneous equation in the \(y_\alpha\)
gives a homogeneous polynomial in the \(x_i\). Hence \(v_d\)
is continuous.
On the open set \(Y_i\ne0\), the inverse is given in the chart
\(x_i\ne0\) by
\[
\frac{x_j}{x_i}
=\frac{y_{(d-1)e_i+e_j}}{Y_i}\quad(j\ne i).
\]
These formulas prove injectivity, and are continuous in affine
chart coordinates: inverse images of polynomial zero sets are
zero sets obtained by clearing a nonvanishing denominator.
They agree on overlaps, so they define a continuous inverse
on the whole image. Thus \(v_d\) is a homeomorphism.

(d) With the coordinates ordered as \(s^3,s^2t,st^2,t^3\),
the chart \(s\ne0\) of \(v_3(\mathbb P^1)\) is
\([1:u:u^2:u^3]\). It is the affine twisted cubic.
It is dense in the image, since this chart is dense in
\(\mathbb P^1\) and (c) is a homeomorphism.
The image is closed by (b), so it is precisely the projective closure.
\end{solution}

\begin{exerciseintro}
% record: H-I-02-013
\exerciseheading{Exercise I.2.13}
\addcontentsline{toc}{subsection}{Exercise I.2.13}

\begin{context}
Work over an algebraically closed field \(k\).
For a projective algebraic set \(Y\), \(I(Y)\) denotes its
homogeneous vanishing ideal. For homogeneous ideals with nonempty
zero set, \(I(V_+(\mathfrak a))=\sqrt{\mathfrak a}\).
An algebraic set is irreducible exactly when its vanishing ideal
is prime. For a projective variety,
\(\dim Y=\dim k[x_0,\ldots,x_n]/I(Y)-1\).
Standard charts are affine spaces, and their topology is obtained
by dehomogenizing polynomials.

The degree-two Veronese map \(v_2:\mathbb P^2\to Y\subseteq
\mathbb P^5\) is a homeomorphism onto a projective variety.
Every irreducible plane curve is \(V_+(f)\) for an irreducible
homogeneous polynomial of positive degree.
The homomorphism
\(\theta:k[y_0,\ldots,y_5]\to k[x_0,x_1,x_2]\) sends the six
coordinates to the six quadratic monomials.
Every even-degree monomial is a product of quadratic monomials.
\end{context}

\begin{problem}
For every irreducible closed curve \(Z\) on the Veronese surface
\(Y\subseteq\mathbb P^5\), construct a hypersurface \(V\) of
\(\mathbb P^5\) such that \(V\cap Y=Z\) as sets.
\end{problem}
\end{exerciseintro}

\begin{solution}
The inverse image \(W=v_2^{-1}(Z)\) is an irreducible closed
curve in \(\mathbb P^2\), because homeomorphisms preserve
closedness, irreducibility, and dimension.
Write \(W=V_+(f)\), where \(f\) is homogeneous and irreducible
of degree \(m>0\).
Every monomial of \(f^2\), of degree \(2m\), can be grouped into
\(m\) quadratic monomials. Hence there is a homogeneous
\(G\in k[y_0,\ldots,y_5]\) of degree \(m\) with
\(\theta(G)=f^2\).
Its zero set meets \(Y\) exactly in \(Z\), since
\(G(v_2(a))=f(a)^2\).

If a hypersurface is required to be irreducible, factor
\(G=c\prod_j G_j\) into homogeneous irreducible factors.
None has zero pullback, because the product pulls back to the
nonzero polynomial \(f^2\).
Unique factorization in \(k[x_0,x_1,x_2]\) gives
\(\theta(G_j)=c_j f^{a_j}\) with \(a_j\geq1\):
the exponent is positive because its degree is \(2\deg G_j>0\).
Thus \(V=V_+(G_j)\), for any factor, is an irreducible
hypersurface with \(V\cap Y=Z\).
The construction also applies in characteristic two.
\end{solution}

\begin{exerciseintro}
% record: H-I-02-014
\exerciseheading{Exercise I.2.14}
\addcontentsline{toc}{subsection}{Exercise I.2.14}

\begin{context}
Work over an algebraically closed field \(k\).
For a projective algebraic set \(Y\), \(I(Y)\) denotes its
homogeneous vanishing ideal. For homogeneous ideals with nonempty
zero set, \(I(V_+(\mathfrak a))=\sqrt{\mathfrak a}\).
An algebraic set is irreducible exactly when its vanishing ideal
is prime. For a projective variety,
\(\dim Y=\dim k[x_0,\ldots,x_n]/I(Y)-1\).
Standard charts are affine spaces, and their topology is obtained
by dehomogenizing polynomials.
\end{context}

\begin{problem}
For nonnegative integers \(r,s\), put \(N=(r+1)(s+1)-1\).
Define the Segre map of sets
\[
\sigma:\mathbb P^r\times\mathbb P^s\to\mathbb P^N,\qquad
([a_i],[b_j])\longmapsto[a_i b_j]_{i,j}.
\]
Prove that its image is a projective variety.
\end{problem}
\end{exerciseintro}

\begin{solution}
Rescaling either input rescales all output coordinates by a
common nonzero scalar, so the map is well-defined.
Let
\[
\theta:k[z_{ij}]\to k[x_0,\ldots,x_r,y_0,\ldots,y_s],
\qquad z_{ij}\longmapsto x_i y_j.
\]
Its kernel \(\mathfrak a\) is prime because its image is a domain,
and homogeneous because a degree-\(d\) polynomial maps to one
of total degree \(2d\).
Substitution gives \(\operatorname{im}\sigma\subseteq V_+(\mathfrak a)\).

The kernel contains all two-by-two minors
\[
z_{ij}z_{i'j'}-z_{ij'}z_{i'j}.
\]
For a projective common zero, choose a nonzero entry \(z_{i_0j_0}\).
Set \(a_i=z_{ij_0}\) and \(b_j=z_{i_0j}/z_{i_0j_0}\).
The minor relations imply \(z_{ij}=a_i b_j\) for every pair \(i,j\).
Both vectors are nonzero since their indicated entries are nonzero.
This point is therefore in the Segre image.
Thus \(\operatorname{im}\sigma=V_+(\mathfrak a)\).
It is nonempty, closed, and irreducible by primality, as required.

The same calculation also proves injectivity: the projective
column and row factors of a nonzero rank-one matrix are uniquely
determined.
\end{solution}

\begin{exerciseintro}
% record: H-I-02-015
\exerciseheading{Exercise I.2.15}
\addcontentsline{toc}{subsection}{Exercise I.2.15}

\begin{context}
Work over an algebraically closed field \(k\).
For a projective algebraic set \(Y\), \(I(Y)\) denotes its
homogeneous vanishing ideal. For homogeneous ideals with nonempty
zero set, \(I(V_+(\mathfrak a))=\sqrt{\mathfrak a}\).
An algebraic set is irreducible exactly when its vanishing ideal
is prime. For a projective variety,
\(\dim Y=\dim k[x_0,\ldots,x_n]/I(Y)-1\).
Standard charts are affine spaces, and their topology is obtained
by dehomogenizing polynomials.

Every proper closed subset of \(\mathbb P^1\) is finite; in
particular any two nonempty open subsets intersect.
The Segre image consists of projective rank-one matrices.
Its inverse on the chart with one nonzero entry is obtained by
taking row and column ratios, which are continuous in affine charts.
\end{context}

\begin{problem}
Let \(Q=V_+(xy-zw)\subseteq\mathbb P^3\).

(a) Identify \(Q\) with the Segre image of \(\mathbb P^1\times\mathbb P^1\).

(b) Describe two families of lines parametrized by \(\mathbb P^1\),
with different lines in the same family disjoint and every line
in one family meeting each line of the other in one point.

(c) Exhibit a different irreducible curve on \(Q\), and show that
the Segre bijection is not a homeomorphism from the product of
the two Zariski topologies.
\end{problem}
\end{exerciseintro}

\begin{solution}
(a) Use the coordinates
\[
([a_0:a_1],[b_0:b_1])
\longmapsto[x:y:z:w]
=[a_0b_0:a_1b_1:a_0b_1:a_1b_0].
\]
The equation \(xy-zw=0\) is the determinant of
\(\left(\begin{smallmatrix}x&z\\ w&y\end{smallmatrix}\right)\).
Its nonzero matrices of determinant zero are exactly the
rank-one matrices, hence exactly this Segre image.
The equation is irreducible: as a polynomial in \(y\), it is
primitive and linear over \(k[x,z,w]\), with coprime coefficients
\(x,-zw\), so Gauss's lemma applies.
Thus \(Q\) is an irreducible surface.

(b) Fixing \([a]\) and varying \([b]\) gives a line \(L_a\),
since its coordinates span a two-dimensional vector subspace.
Fixing \([b]\) gives the line \(M_b\).
Uniqueness of the two projective factors of a rank-one matrix
shows that distinct \(L_a\)'s are disjoint, as are distinct
\(M_b\)'s, and that
\[
L_a\cap M_b=\{\sigma([a],[b])\}.
\]

(c) The diagonal pair \(([s:t],[s:t])\) has image
\[
[s^2:t^2:st:st].
\]
Its image is \(Q\cap V_+(z-w)\): a rank-one matrix is symmetric
exactly when its row and column factors are proportional.
In the plane \(z=w\), this is the irreducible conic \(xy=z^2\).
For example it is the degree-two Veronese image of \(\mathbb P^1\),
so is an irreducible closed curve, and neither factor is constant.
It is therefore none of the ruling lines in (b).

The diagonal \(\Delta\subseteq\mathbb P^1\times\mathbb P^1\)
is dense for the product topology: every nonempty basic open
\(U\times V\) meets it, since \(U\cap V\ne\varnothing\).
It is a proper subset, hence not closed in that topology.
Its image on \(Q\) is the closed conic just described.
Therefore the Segre bijection from the product topological space
cannot be a homeomorphism.
\end{solution}

\begin{exerciseintro}
% record: H-I-02-016
\exerciseheading{Exercise I.2.16}
\addcontentsline{toc}{subsection}{Exercise I.2.16}

\begin{context}
Work over an algebraically closed field \(k\).
For a projective algebraic set \(Y\), \(I(Y)\) denotes its
homogeneous vanishing ideal. For homogeneous ideals with nonempty
zero set, \(I(V_+(\mathfrak a))=\sqrt{\mathfrak a}\).
An algebraic set is irreducible exactly when its vanishing ideal
is prime. For a projective variety,
\(\dim Y=\dim k[x_0,\ldots,x_n]/I(Y)-1\).
Standard charts are affine spaces, and their topology is obtained
by dehomogenizing polynomials.

The projective twisted cubic in coordinates \(w,x,y,z\) is the
closure of \([1:t:t^2:t^3]\).
A prime ideal defines an irreducible algebraic set, and
nonzero linear forms define hyperplanes.
\end{context}

\begin{problem}
(a) Show that the intersection in \(\mathbb P^3\) of the quadrics
\[
Q_1=V_+(x^2-yw),\qquad Q_2=V_+(xy-zw)
\]
is the union of a twisted cubic and a line.

(b) For \(C=V_+(x^2-yz)\) and \(L=V_+(y)\) in \(\mathbb P^2\),
show that their intersection is a single point \(P\), but
\(I(C)+I(L)\ne I(P)\).
\end{problem}
\end{exerciseintro}

\begin{solution}
(a) On \(w\ne0\), normalize \(w=1\). The two equations then
give \(y=x^2,\ z=x^3\), exactly the affine twisted cubic.
On \(w=0\), the first equation forces \(x=0\), and the second
imposes no further restriction. These points form the line
\(M=V_+(w,x)\).
The projective twisted cubic \(T\) lies on both quadrics, by
substituting \([w:x:y:z]=[s^3:s^2t:st^2:t^3]\).
We have therefore proved
\[
Q_1\cap Q_2=T\cup M.
\]
Neither curve contains the other: \(T\) has points with \(w\ne0\),
whereas \([0:0:1:0]\in M\) is not on \(T\).
They meet at \([0:0:0:1]\).
Both quadrics are irreducible by Gauss's lemma, applied to the
primitive linear polynomials in \(w\) with coefficients
\(-y,x^2\) and \(-z,xy\), respectively.
Thus this is an intersection of varieties that is reducible.

(b) Setting \(y=0\) in the conic equation gives \(x^2=0\).
The only projective common zero is \(P=[0:0:1]\), with ideal
\(I(P)=(x,y)\).
The conic equation is irreducible by the same primitive-linear
argument, so \(I(C)=(x^2-yz)\). Consequently
\[
I(C)+I(L)=(x^2-yz,y)=(x^2,y)\subsetneq(x,y).
\]
The containment is strict because the image of \(x\) in
\(k[x,y,z]/(x^2,y)\cong k[x,z]/(x^2)\) is nonzero.
Its square vanishes, explaining why the sum ideal has the right
zero set but is not the vanishing ideal of that point.
\end{solution}

\begin{exerciseintro}
% record: H-I-02-017
\exerciseheading{Exercise I.2.17}
\addcontentsline{toc}{subsection}{Exercise I.2.17}

\begin{context}
Work over an algebraically closed field \(k\).
For a projective algebraic set \(Y\), \(I(Y)\) denotes its
homogeneous vanishing ideal. For homogeneous ideals with nonempty
zero set, \(I(V_+(\mathfrak a))=\sqrt{\mathfrak a}\).
An algebraic set is irreducible exactly when its vanishing ideal
is prime. For a projective variety,
\(\dim Y=\dim k[x_0,\ldots,x_n]/I(Y)-1\).
Standard charts are affine spaces, and their topology is obtained
by dehomogenizing polynomials.

If an ideal of a polynomial ring is generated by \(q\) elements,
every prime minimal over it has height at most \(q\).
For a projective variety \(Y\subseteq\mathbb P^n\),
\(\operatorname{ht}I(Y)=n-\dim Y\).
A strict complete intersection has its homogeneous ideal
generated by \(n-\dim Y\) elements; a set-theoretic complete
intersection is the common zero set of that many hypersurfaces.

In coordinates \(w,x,y,z\), the projective twisted cubic \(T\)
has ideal
\[
I(T)=(wy-x^2,\ wz-xy,\ xz-y^2).
\]
These generators follow from the parametrization
\([s^3:s^2t:st^2:t^3]\): straightening pairs of indices of
difference at least two gives one monomial for each degree
and index sum, whose images are linearly independent.
\end{context}

\begin{problem}
(a) If a projective variety \(Y\subseteq\mathbb P^n\) is the
zero set of a homogeneous ideal generated by \(q\) elements,
prove \(\dim Y\geq n-q\).

(b) Prove that every strict complete intersection is a
set-theoretic complete intersection.

(c) Show that the twisted cubic is not a strict complete
intersection, but is the set-theoretic intersection of a
quadric and a cubic surface.

(d) Explain the distinction illustrated by (c) in relation
to the source's research question whether all irreducible
projective space curves are set-theoretic intersections of
two surfaces.
\end{problem}
\end{exerciseintro}

\begin{solution}
(a) Write \(Y=V_+(\mathfrak a)\).
Since \(Y\) is nonempty, the projective Nullstellensatz gives
\(\sqrt{\mathfrak a}=I(Y)\), a prime ideal.
It is the unique minimal prime over \(\mathfrak a\).
The height theorem gives \(n-\dim Y=\operatorname{ht}I(Y)\leq q\),
as desired.

(b) If \(I(Y)\) is generated by \(n-r\) homogeneous polynomials,
where \(r=\dim Y\), then their common zero set is \(Y\).
If each hypersurface is required to be irreducible, choose from
each generator an irreducible homogeneous factor vanishing on \(Y\).
Such a factor exists because \(I(Y)\) is prime.
The zero set of the chosen factor is contained in that of the
original generator and contains \(Y\). Their intersection is
therefore still exactly \(Y\).
For \(Y=\mathbb P^n\), use the empty intersection.

(c) The three displayed quadrics are linearly independent,
and \(I(T)\) has no degree-one elements, since the four
parametrizing monomials are linearly independent.
Put \(\mathfrak m=(w,x,y,z)\).
Then \(\mathfrak m I(T)\) has no nonzero degree-two element,
so the three quadrics remain independent in
\(I(T)/\mathfrak m I(T)\).
This quotient cannot be spanned by two generators.
Hence \(I(T)\) cannot be generated by two polynomials.

On the other hand, take
\[
Q=wy-x^2,\qquad F=wz^2-2xyz+y^3.
\]
Both vanish on \(T\).
At a common zero with \(w\ne0\), the equation \(Q=0\) gives
\(y=x^2/w\), and substitution in \(F\) gives
\[
F=w\left(z-\frac{x^3}{w^2}\right)^2.
\]
Thus \(z=x^3/w^2\), so this point is on \(T\).
At \(w=0\), the equations force \(x=0\) and then \(y=0\),
leaving the single point \([0:0:0:1]\), also on \(T\).
Therefore \(V_+(Q)\cap V_+(F)=T\) set-theoretically, in every
characteristic, including two.
The quadric is irreducible. The cubic is irreducible as well:
view it as a polynomial in \(w\) over \(k[x,y,z]\). Its
coefficients \(z^2\) and \(y^3-2xyz\) are coprime, so Gauss's
lemma reduces irreducibility to that of a linear polynomial
over the fraction field.

(d) Part (c) concerns two different requirements.
Two equations can have radical ideal \(I(T)\) even though
\(I(T)\) itself requires three generators.
Thus failure of a strict complete intersection supplies no
counterexample to the set-theoretic question.
The general question asks for equations for every irreducible
projective space curve; the explicit equations above establish
only the twisted-cubic case.
\end{solution}

\begin{exerciseintro}
\corpussection{I.3 Morphisms}
\addcontentsline{toc}{section}{I.3 Morphisms}

% record: H-I-03-001
\exerciseheading{Exercise I.3.1}
\addcontentsline{toc}{subsection}{Exercise I.3.1}

\begin{context}
Work over an algebraically closed field \(k\); varieties are irreducible
and carry the Zariski topology. A regular function is locally a quotient
of polynomials with nonvanishing denominator. A morphism is a continuous
map pulling back regular functions to regular functions.
Morphisms to an affine variety correspond contravariantly to
\(k\)-algebra homomorphisms of its coordinate ring into the ring of
global regular functions. Morphisms and regular functions can be
checked and glued on open covers. The local ring at a point consists
of germs of regular functions, with maximal ideal the germs vanishing
there. On an affine variety with coordinate ring \(A\), it is
\(A_{\mathfrak m}\). The function field is \(\operatorname{Frac}A\).

A projective variety has only constant global regular functions.
A plane curve is the zero set of one irreducible polynomial, homogeneous
in the projective case. A quotient of \(k[x,y,z]\) by a proper ideal
generated by two elements has dimension at least one. The homogeneous Nullstellensatz says that
an empty projective zero set has affine cone consisting only of the origin.
Conics mean irreducible degree-two curves. Affine linear changes put an
affine conic into either a graph of a quadratic polynomial or \(xy=1\):
the leading quadratic splits into either two different linear factors
or a repeated factor. Polynomial rings are unique factorization domains.
\end{context}

\begin{problem}
(a) Classify affine plane conics up to isomorphism as
\(\mathbb A^1\) or \(\mathbb A^1\setminus\{0\}\).

(b) Show that no proper open subset of \(\mathbb A^1\) is
isomorphic to \(\mathbb A^1\).

(c) Show that every projective plane conic is isomorphic to
\(\mathbb P^1\).

(d) Distinguish \(\mathbb A^2\) and \(\mathbb P^2\) topologically.

(e) Show that a variety which is both affine and projective up
to isomorphism is a point.
\end{problem}
\end{exerciseintro}

\begin{solution}
(a) Projection onto the parameter coordinate identifies a quadratic
graph with \(\mathbb A^1\), with polynomial inverse.
The maps \(t\mapsto(t,t^{-1})\) and \((x,y)\mapsto x\)
identify \(xy=1\) with \(\mathbb A^1\setminus\{0\}\).
These are morphisms in both directions.

(b) A nonempty proper open subset is
\(U=\mathbb A^1\setminus\{a_1,\ldots,a_s\}\), where \(s\geq1\).
The regular function \(t-a_1\) is a nonconstant unit on \(U\).
An isomorphism would carry it to a nonconstant unit in \(k[t]\),
whose units are precisely \(k^\times\), a contradiction.
The empty open subset cannot be isomorphic to a nonempty variety.

(c) Choose a point on the conic and make it \([0:0:1]\).
The equation is
\[
z(ax+by)+q(x,y)=0.
\]
The linear form \(ax+by\) is nonzero; otherwise the binary quadratic
\(q\) splits over \(k\), contradicting irreducibility.
Change \(x,y\) so that this form is \(x\).
The coefficient of \(y^2\) in \(q\) is nonzero, since otherwise
the whole equation is divisible by \(x\).
Replacing \(z\) by \(z+\alpha x+\beta y\), and rescaling, gives
\(xz=y^2\). No division by \(2\) was used.
The morphism
\[
[s:t]\longmapsto[s^2:st:t^2]
\]
has image this conic. On \(x\ne0\) its inverse is \([x:y]\);
on \(z\ne0\) its inverse is \([y:z]\).
These charts cover the conic and the inverses agree on their overlap.
Thus it is an isomorphism in every characteristic.

(d) Any two projective plane curves meet. Otherwise the cone of their
intersection would be just the origin, so the quotient by their two
homogeneous equations would have dimension zero, contrary to the
dimension bound in the context.
But the two closed affine lines \(x=0\) and \(x=1\) are disjoint.
A homeomorphism would send these to disjoint closed irreducible
one-dimensional subsets of \(\mathbb P^2\), hence to projective
plane curves, giving a contradiction.

(e) The affine coordinate ring of such a variety is its ring of global
regular functions. Projectivity makes that ring \(k\), whose affine
variety is one point.
\end{solution}

\begin{exerciseintro}
% record: H-I-03-002
\exerciseheading{Exercise I.3.2}
\addcontentsline{toc}{subsection}{Exercise I.3.2}

\begin{context}
Work over an algebraically closed field \(k\); varieties are irreducible
and carry the Zariski topology. A regular function is locally a quotient
of polynomials with nonvanishing denominator. A morphism is a continuous
map pulling back regular functions to regular functions.
Morphisms to an affine variety correspond contravariantly to
\(k\)-algebra homomorphisms of its coordinate ring into the ring of
global regular functions. Morphisms and regular functions can be
checked and glued on open covers. The local ring at a point consists
of germs of regular functions, with maximal ideal the germs vanishing
there. On an affine variety with coordinate ring \(A\), it is
\(A_{\mathfrak m}\). The function field is \(\operatorname{Frac}A\).

A proper closed subset of an irreducible curve is finite.
Thus its Zariski topology on points is the cofinite topology.
\end{context}

\begin{problem}
(a) Prove that \(t\mapsto(t^2,t^3)\) is a homeomorphic
morphism from \(\mathbb A^1\) onto \(V(y^2-x^3)\), but is not
an isomorphism.

(b) In characteristic \(p>0\), prove the same assertions
for the morphism \(\mathbb A^1\to\mathbb A^1\), \(t\mapsto t^p\).
\end{problem}
\end{exerciseintro}

\begin{solution}
(a) Polynomial coordinates define a morphism. Away from the origin,
a point of the cusp has \(x\ne0\), and \(t=y/x\) satisfies
\(t^2=x,t^3=y\). The origin has the unique preimage \(0\).
Hence the map is bijective, and it is a homeomorphism because both
curves have the cofinite topology.
Its pullback has image \(k[t^2,t^3]\subset k[t]\).
This image does not contain \(t\), since its monomials have exponents
zero or at least two. The pullback is not an isomorphism, so neither
is the morphism.

(b) Algebraic closure gives a root of \(T^p-a\) for each \(a\).
The identity \(T^p-a=(T-b)^p\) for any such root makes it unique.
Thus the polynomial morphism is bijective and again a homeomorphism.
Its pullback \(k[u]\to k[t]\), \(u\mapsto t^p\), has image \(k[t^p]\),
which does not contain \(t\).
Here this is the \(k\)-morphism defined by that polynomial; it fixes
constant functions under pullback.
\end{solution}

\begin{exerciseintro}
% record: H-I-03-003
\exerciseheading{Exercise I.3.3}
\addcontentsline{toc}{subsection}{Exercise I.3.3}

\begin{context}
Work over an algebraically closed field \(k\); varieties are irreducible
and carry the Zariski topology. A regular function is locally a quotient
of polynomials with nonvanishing denominator. A morphism is a continuous
map pulling back regular functions to regular functions.
Morphisms to an affine variety correspond contravariantly to
\(k\)-algebra homomorphisms of its coordinate ring into the ring of
global regular functions. Morphisms and regular functions can be
checked and glued on open covers. The local ring at a point consists
of germs of regular functions, with maximal ideal the germs vanishing
there. On an affine variety with coordinate ring \(A\), it is
\(A_{\mathfrak m}\). The function field is \(\operatorname{Frac}A\).
\end{context}

\begin{problem}
Let \(\varphi:X\to Y\) be a morphism.

(a) Construct the induced local homomorphism
\(\varphi_P^*:\mathcal O_{\varphi(P),Y}\to\mathcal O_{P,X}\).

(b) Prove that \(\varphi\) is an isomorphism exactly when it is
a homeomorphism and all these local homomorphisms are isomorphisms.

(c) If \(\varphi(X)\) is dense in \(Y\), prove that every
\(\varphi_P^*\) is injective.
\end{problem}
\end{exerciseintro}

\begin{solution}
(a) A germ represented by \(f\) on a neighborhood \(V\) of \(\varphi(P)\)
is sent to the germ of \(f\circ\varphi\) on \(\varphi^{-1}(V)\).
Agreement on a smaller neighborhood is preserved by inverse image,
so this is well defined. It respects addition, multiplication, and \(k\).
The pulled-back value at \(P\) is \(f(\varphi(P))\).
Consequently the inverse image of the maximal ideal is exactly the
maximal ideal at \(\varphi(P)\), making it a local homomorphism.

(b) An inverse morphism supplies inverse maps on germs and on points.
Conversely suppose the stated conditions hold, and write
\(\psi=\varphi^{-1}\) for the continuous inverse.
For a regular function \(g\) on an open \(U\subseteq X\) and \(P\in U\),
surjectivity on local rings gives a regular function \(h\) near
\(\varphi(P)\) whose pullback agrees with \(g\) near \(P\).
Shrink the latter neighborhood to \(W\subseteq U\) inside the
inverse image of the domain of \(h\).
Because \(\varphi\) is a homeomorphism, \(\varphi(W)\) is open.
On it, \(g\circ\psi=h\). Thus \(g\circ\psi\) is locally regular
everywhere on \(\varphi(U)\), hence regular.
This proves that \(\psi\) is a morphism.

(c) Let a germ represented by \(f\) on \(V\) pull back to zero at \(P\).
Then \(f\circ\varphi\) vanishes on a nonempty open subset of
\(\varphi^{-1}(V)\). The latter is irreducible, so a regular function
zero on that dense open subset is zero everywhere there.
The subset \(\varphi(X)\cap V\) is dense in \(V\).
Thus \(f\) vanishes on a dense subset of \(V\), and hence identically.
Its germ is zero, as required.
\end{solution}

\begin{exerciseintro}
% record: H-I-03-004
\exerciseheading{Exercise I.3.4}
\addcontentsline{toc}{subsection}{Exercise I.3.4}

\begin{context}
Work over an algebraically closed field \(k\); varieties are irreducible
and carry the Zariski topology. A regular function is locally a quotient
of polynomials with nonvanishing denominator. A morphism is a continuous
map pulling back regular functions to regular functions.
Morphisms to an affine variety correspond contravariantly to
\(k\)-algebra homomorphisms of its coordinate ring into the ring of
global regular functions. Morphisms and regular functions can be
checked and glued on open covers. The local ring at a point consists
of germs of regular functions, with maximal ideal the germs vanishing
there. On an affine variety with coordinate ring \(A\), it is
\(A_{\mathfrak m}\). The function field is \(\operatorname{Frac}A\).

For \(d\geq1\), the degree-\(d\) Veronese map has coordinates
\(y_\alpha=x^\alpha\), one for each multi-index \(\alpha\) of
total degree \(d\). It is a bijection onto a closed projective
subvariety \(V\). Write \(e_i\) for the \(i\)-th coordinate vector.
\end{context}

\begin{problem}
Prove that the degree-\(d\) Veronese map
\(\nu_d:\mathbb P^n\to V\) is an isomorphism.
\end{problem}
\end{exerciseintro}

\begin{solution}
On \(x_i\ne0\), the ratios \(y_\alpha/y_{de_i}\) are monomials
in \(x_j/x_i\), so the map is a morphism on these covering charts.
On \(V\cap\{y_{de_i}\ne0\}\) recover the affine coordinates by
\[
\frac{x_j}{x_i}=
\frac{y_{(d-1)e_i+e_j}}{y_{de_i}}.
\]
These are regular functions. The charts cover \(V\) because every
image point has some \(x_i^d\ne0\). The displayed formula gives the
set-theoretic inverse on each chart, so these inverse morphisms agree
on overlaps and glue. This proves the assertion, including \(d=1\),
without a separability or characteristic restriction.
\end{solution}

\begin{exerciseintro}
% record: H-I-03-005
\exerciseheading{Exercise I.3.5}
\addcontentsline{toc}{subsection}{Exercise I.3.5}

\begin{context}
Work over an algebraically closed field \(k\); varieties are irreducible
and carry the Zariski topology. A regular function is locally a quotient
of polynomials with nonvanishing denominator. A morphism is a continuous
map pulling back regular functions to regular functions.
Morphisms to an affine variety correspond contravariantly to
\(k\)-algebra homomorphisms of its coordinate ring into the ring of
global regular functions. Morphisms and regular functions can be
checked and glued on open covers. The local ring at a point consists
of germs of regular functions, with maximal ideal the germs vanishing
there. On an affine variety with coordinate ring \(A\), it is
\(A_{\mathfrak m}\). The function field is \(\operatorname{Frac}A\).

The degree-\(d\) Veronese map is an isomorphism onto a closed
subvariety of projective space. The complement of a hyperplane
in \(\mathbb P^N\) is isomorphic to \(\mathbb A^N\).
Closed subvarieties of affine space are affine.
\end{context}

\begin{problem}
If \(H\subseteq\mathbb P^n\) is a hypersurface, prove that
\(\mathbb P^n\setminus H\) is affine.
\end{problem}
\end{exerciseintro}

\begin{solution}
Write \(H=V_+(F)\), where \(F\) is nonzero homogeneous of positive
degree \(d\). In the degree-\(d\) Veronese coordinates, \(F\) becomes
a nonzero linear form \(L\). If \(V\) is the Veronese image, then
\[
\nu_d(\mathbb P^n\setminus H)=V\cap(\mathbb P^N\setminus V_+(L)).
\]
The right side is closed in the affine hyperplane complement,
because \(V\) is closed in \(\mathbb P^N\). It is therefore affine.
Restricting the Veronese isomorphism proves the claim.
\end{solution}

\begin{exerciseintro}
% record: H-I-03-006
\exerciseheading{Exercise I.3.6}
\addcontentsline{toc}{subsection}{Exercise I.3.6}

\begin{context}
Work over an algebraically closed field \(k\); varieties are irreducible
and carry the Zariski topology. A regular function is locally a quotient
of polynomials with nonvanishing denominator. A morphism is a continuous
map pulling back regular functions to regular functions.
Morphisms to an affine variety correspond contravariantly to
\(k\)-algebra homomorphisms of its coordinate ring into the ring of
global regular functions. Morphisms and regular functions can be
checked and glued on open covers. The local ring at a point consists
of germs of regular functions, with maximal ideal the germs vanishing
there. On an affine variety with coordinate ring \(A\), it is
\(A_{\mathfrak m}\). The function field is \(\operatorname{Frac}A\).

For \(A=k[x,y]\), the principal open \(D(g)\) has global
regular-function ring \(A_g\). The ring \(A\) is a unique factorization
domain. Isomorphisms of affine coordinate rings correspond to
isomorphisms of affine varieties.
\end{context}

\begin{problem}
Prove that \(X=\mathbb A^2\setminus\{(0,0)\}\) is not affine.
\end{problem}
\end{exerciseintro}

\begin{solution}
The opens \(D(x)\) and \(D(y)\) cover \(X\), so the regular functions
on \(X\), viewed in \(k(x,y)\), form
\[
\mathcal O(X)=A_x\cap A_y=A.
\]
For the last equality, a reduced fraction in \(A_x\) has denominator
dividing a power of \(x\); membership in \(A_y\) forces it also to
divide a power of \(y\). Unique factorization makes that denominator
a unit.

The inclusion \(j:X\to\mathbb A^2\) induces the identity isomorphism
\(A\to\mathcal O(X)\). If \(X\) were affine, the affine
coordinate-ring equivalence would make \(j\) an isomorphism.
But \(j\) misses the origin. This contradiction proves nonaffineness.
\end{solution}

\begin{exerciseintro}
% record: H-I-03-007
\exerciseheading{Exercise I.3.7}
\addcontentsline{toc}{subsection}{Exercise I.3.7}

\begin{context}
Work over an algebraically closed field \(k\); varieties are irreducible
and carry the Zariski topology. A regular function is locally a quotient
of polynomials with nonvanishing denominator. A morphism is a continuous
map pulling back regular functions to regular functions.
Morphisms to an affine variety correspond contravariantly to
\(k\)-algebra homomorphisms of its coordinate ring into the ring of
global regular functions. Morphisms and regular functions can be
checked and glued on open covers. The local ring at a point consists
of germs of regular functions, with maximal ideal the germs vanishing
there. On an affine variety with coordinate ring \(A\), it is
\(A_{\mathfrak m}\). The function field is \(\operatorname{Frac}A\).

The complement of a projective hypersurface is affine.
A projective variety has only constant global regular functions.
An affine variety whose coordinate ring is \(k\) is a point.
Every projective plane curve is a hypersurface.
\end{context}

\begin{problem}
(a) Show that two curves in \(\mathbb P^2\) always meet.

(b) Show that a positive-dimensional projective variety
\(Y\subseteq\mathbb P^n\) meets every hypersurface \(H\).
\end{problem}
\end{exerciseintro}

\begin{solution}
(b) If \(Y\cap H\) were empty, \(Y\) would be a closed subvariety
of the affine variety \(\mathbb P^n\setminus H\); hence \(Y\)
would be affine. Its global regular functions are \(k\), since
it is projective. Thus it would be a point, contradicting
\(\dim Y\geq1\).

(a) Apply (b) with one curve as \(Y\) and the other as \(H\).
If the word curve allows a reducible algebraic set, choose a
one-dimensional irreducible component from each and apply this result.
\end{solution}

\begin{exerciseintro}
% record: H-I-03-008
\exerciseheading{Exercise I.3.8}
\addcontentsline{toc}{subsection}{Exercise I.3.8}

\begin{context}
Work over an algebraically closed field \(k\); varieties are irreducible
and carry the Zariski topology. A regular function is locally a quotient
of polynomials with nonvanishing denominator. A morphism is a continuous
map pulling back regular functions to regular functions.
Morphisms to an affine variety correspond contravariantly to
\(k\)-algebra homomorphisms of its coordinate ring into the ring of
global regular functions. Morphisms and regular functions can be
checked and glued on open covers. The local ring at a point consists
of germs of regular functions, with maximal ideal the germs vanishing
there. On an affine variety with coordinate ring \(A\), it is
\(A_{\mathfrak m}\). The function field is \(\operatorname{Frac}A\).
\end{context}

\begin{problem}
For different coordinate hyperplanes \(H_i,H_j\subseteq\mathbb P^n\),
prove that every regular function on
\(\mathbb P^n\setminus(H_i\cap H_j)\) is constant.
\end{problem}
\end{exerciseintro}

\begin{solution}
Relabel so that \(i=0,j=1\), and put \(U_\ell=\{x_\ell\ne0\}\).
The indicated open set is \(U_0\cup U_1\).
On \(U_0\), write \(t_\ell=x_\ell/x_0\). A regular function is
a polynomial \(P(t_1,\ldots,t_n)\).
On \(U_1\), use \(s_0=x_0/x_1\) and \(s_\ell=x_\ell/x_1\)
for \(\ell\geq2\). On the overlap,
\[
t_1=s_0^{-1},\qquad t_\ell=s_\ell s_0^{-1}\quad(\ell\geq2).
\]
A monomial \(t_1^{a_1}\cdots t_n^{a_n}\) therefore becomes
\[
s_0^{-(a_1+\cdots+a_n)}s_2^{a_2}\cdots s_n^{a_n}.
\]
Different exponent vectors give different Laurent monomials,
so no cancellation can remove their negative powers of \(s_0\).
For \(P\) to be polynomial in the \(U_1\) coordinates, all
nonzero-degree monomials must have coefficient zero.
Thus \(P\) and the original function are constant.
For \(n=1\) this is the same argument with no \(s_\ell\), \(\ell\geq2\).
\end{solution}

\begin{exerciseintro}
% record: H-I-03-009
\exerciseheading{Exercise I.3.9}
\addcontentsline{toc}{subsection}{Exercise I.3.9}

\begin{context}
Work over an algebraically closed field \(k\); varieties are irreducible
and carry the Zariski topology. A regular function is locally a quotient
of polynomials with nonvanishing denominator. A morphism is a continuous
map pulling back regular functions to regular functions.
Morphisms to an affine variety correspond contravariantly to
\(k\)-algebra homomorphisms of its coordinate ring into the ring of
global regular functions. Morphisms and regular functions can be
checked and glued on open covers. The local ring at a point consists
of germs of regular functions, with maximal ideal the germs vanishing
there. On an affine variety with coordinate ring \(A\), it is
\(A_{\mathfrak m}\). The function field is \(\operatorname{Frac}A\).

For a \(k\)-algebra \(R\) and a maximal ideal \(\mathfrak m\) with
residue field \(k\), the vector space \(\mathfrak m/\mathfrak m^2\)
and its dimension are invariants of that local \(k\)-algebra.
Every maximal ideal of \(k[s,t]\) has this dimension two.
\end{context}

\begin{problem}
Let \(X=\mathbb P^1\) and let \(Y\subseteq\mathbb P^2\) be its
degree-two Veronese image. Although \(X\cong Y\), show that their
homogeneous coordinate rings are not isomorphic, even as
ungraded \(k\)-algebras.
\end{problem}
\end{exerciseintro}

\begin{solution}
The rings are
\[
S(X)=k[s,t],\qquad S(Y)=k[a,b,c]/(ac-b^2).
\]
The second equality follows from the parametrization
\([s:t]\mapsto[s^2:st:t^2]\): its image is exactly the irreducible
conic \(ac=b^2\).
At the maximal ideal \(\mathfrak m=(a,b,c)\) of the second ring,
the relation has degree two, so it imposes no linear relation
modulo \(\mathfrak m^2\).
Thus \(\dim_k\mathfrak m/\mathfrak m^2=3\).
An isomorphism with \(k[s,t]\) would give a maximal ideal in
that polynomial ring with the same dimension, which is impossible.
The Veronese map nevertheless gives \(X\cong Y\); hence the
homogeneous coordinate ring depends on the projective embedding.
\end{solution}

\begin{exerciseintro}
% record: H-I-03-010
\exerciseheading{Exercise I.3.10}
\addcontentsline{toc}{subsection}{Exercise I.3.10}

\begin{context}
Work over an algebraically closed field \(k\); varieties are irreducible
and carry the Zariski topology. A regular function is locally a quotient
of polynomials with nonvanishing denominator. A morphism is a continuous
map pulling back regular functions to regular functions.
Morphisms to an affine variety correspond contravariantly to
\(k\)-algebra homomorphisms of its coordinate ring into the ring of
global regular functions. Morphisms and regular functions can be
checked and glued on open covers. The local ring at a point consists
of germs of regular functions, with maximal ideal the germs vanishing
there. On an affine variety with coordinate ring \(A\), it is
\(A_{\mathfrak m}\). The function field is \(\operatorname{Frac}A\).

A subvariety is an irreducible locally closed subset with its induced
variety structure. A regular function on a subvariety is locally
the restriction of a quotient of ambient affine-chart polynomials
with denominator nonzero at the point in question.
\end{context}

\begin{problem}
Let \(\varphi:X\to Y\) be a morphism and let
\(X'\subseteq X\), \(Y'\subseteq Y\) be subvarieties with
\(\varphi(X')\subseteq Y'\). Prove that
\(\varphi|_{X'}:X'\to Y'\) is a morphism.
\end{problem}
\end{exerciseintro}

\begin{solution}
The restricted map is continuous for the subspace topologies.
Let \(g\) be regular on an open subset \(V'\subseteq Y'\).
Near any point \(Q\in V'\), choose an ambient affine chart and
write \(g=a/b\) on a neighborhood in \(Y'\), with \(b(Q)\ne0\).
After shrinking an open neighborhood \(V\) of \(Q\) in \(Y\),
the same quotient defines a regular function \(h\) on \(V\)
whose restriction agrees with \(g\) on \(Y'\cap V\subseteq V'\).
The pullback \(h\circ\varphi\) is regular on \(\varphi^{-1}(V)\).
Restricting it to \(X'\) remains regular, and there it equals
\(g\circ\varphi|_{X'}\).
This works near every point of the inverse image of \(V'\).
The pullback of \(g\) is therefore regular, proving the assertion.
\end{solution}

\begin{exerciseintro}
% record: H-I-03-011
\exerciseheading{Exercise I.3.11}
\addcontentsline{toc}{subsection}{Exercise I.3.11}

\begin{context}
Work over an algebraically closed field \(k\); varieties are irreducible
and carry the Zariski topology. A regular function is locally a quotient
of polynomials with nonvanishing denominator. A morphism is a continuous
map pulling back regular functions to regular functions.
Morphisms to an affine variety correspond contravariantly to
\(k\)-algebra homomorphisms of its coordinate ring into the ring of
global regular functions. Morphisms and regular functions can be
checked and glued on open covers. The local ring at a point consists
of germs of regular functions, with maximal ideal the germs vanishing
there. On an affine variety with coordinate ring \(A\), it is
\(A_{\mathfrak m}\). The function field is \(\operatorname{Frac}A\).

Prime ideals in an affine coordinate ring correspond to irreducible
closed subvarieties. The primes in \(A_{\mathfrak m}\) correspond,
by extension and contraction, to the primes of \(A\) contained in
\(\mathfrak m\). A nonempty open subset of an irreducible space is dense.
\end{context}

\begin{problem}
For \(P\in X\), give a bijection between the prime ideals of
\(\mathcal O_{P,X}\) and the closed subvarieties of \(X\) containing \(P\).
\end{problem}
\end{exerciseintro}

\begin{solution}
Choose an affine open neighborhood \(U\) of \(P\), with ring \(A\)
and maximal ideal \(\mathfrak m_P\).
A prime of \(\mathcal O_{P,X}=A_{\mathfrak m_P}\) contracts to
a prime \(\mathfrak p\subseteq\mathfrak m_P\).
The subvariety \(V_U(\mathfrak p)\) contains \(P\); take its
closure in \(X\).

Conversely, if \(Z\subseteq X\) is closed irreducible and contains \(P\),
then \(Z\cap U\) is nonempty, closed in \(U\), and irreducible.
Its prime ideal is contained in \(\mathfrak m_P\); extend this
ideal to \(A_{\mathfrak m_P}\).

The two constructions are inverse: taking the closure in \(X\)
and intersecting with \(U\) recovers a closed subset of \(U\),
and \(\overline{Z\cap U}=Z\) by density.
Extension and contraction of primes are inverse as well.
Intrinsically, the prime associated to \(Z\) consists of germs
whose restriction vanishes on a neighborhood of \(P\) in \(Z\).
This description also proves independence of the affine neighborhood.
\end{solution}

\begin{exerciseintro}
% record: H-I-03-012
\exerciseheading{Exercise I.3.12}
\addcontentsline{toc}{subsection}{Exercise I.3.12}

\begin{context}
Work over an algebraically closed field \(k\); varieties are irreducible
and carry the Zariski topology. A regular function is locally a quotient
of polynomials with nonvanishing denominator. A morphism is a continuous
map pulling back regular functions to regular functions.
Morphisms to an affine variety correspond contravariantly to
\(k\)-algebra homomorphisms of its coordinate ring into the ring of
global regular functions. Morphisms and regular functions can be
checked and glued on open covers. The local ring at a point consists
of germs of regular functions, with maximal ideal the germs vanishing
there. On an affine variety with coordinate ring \(A\), it is
\(A_{\mathfrak m}\). The function field is \(\operatorname{Frac}A\).

For a finitely generated integral \(k\)-algebra \(A\), dimension
equals the transcendence degree of its fraction field, and
\(\operatorname{ht}\mathfrak p+\dim A/\mathfrak p=\dim A\).
Maximal ideals have residue field \(k\). Localization identifies
the primes of \(A_{\mathfrak p}\) with the primes contained in
\(\mathfrak p\), so \(\dim A_{\mathfrak p}=\operatorname{ht}\mathfrak p\).
Every nonempty open subset of a variety has its dimension and
function field.
\end{context}

\begin{problem}
Prove that \(\dim\mathcal O_{P,X}=\dim X\) for every point \(P\)
of a variety \(X\).
\end{problem}
\end{exerciseintro}

\begin{solution}
Choose an affine open neighborhood \(U\) of \(P\), with coordinate
ring \(A\) and maximal ideal \(\mathfrak m\).
Then \(\mathcal O_{P,X}=A_{\mathfrak m}\).
Since \(A/\mathfrak m=k\) has dimension zero, the affine dimension
formula gives
\[
\dim\mathcal O_{P,X}
=\operatorname{ht}\mathfrak m
=\dim A
=\dim U
=\dim X.
\]
Here points are the closed points of the classical variety.
The assertion is not a statement that every scheme-theoretic point
has local-ring dimension \(\dim X\).
\end{solution}

\begin{exerciseintro}
% record: H-I-03-013
\exerciseheading{Exercise I.3.13}
\addcontentsline{toc}{subsection}{Exercise I.3.13}

\begin{context}
Work over an algebraically closed field \(k\); varieties are irreducible
and carry the Zariski topology. A regular function is locally a quotient
of polynomials with nonvanishing denominator. A morphism is a continuous
map pulling back regular functions to regular functions.
Morphisms to an affine variety correspond contravariantly to
\(k\)-algebra homomorphisms of its coordinate ring into the ring of
global regular functions. Morphisms and regular functions can be
checked and glued on open covers. The local ring at a point consists
of germs of regular functions, with maximal ideal the germs vanishing
there. On an affine variety with coordinate ring \(A\), it is
\(A_{\mathfrak m}\). The function field is \(\operatorname{Frac}A\).

For a finitely generated integral \(k\)-algebra \(A\), dimension
equals the transcendence degree of its fraction field, and
\(\operatorname{ht}\mathfrak p+\dim A/\mathfrak p=\dim A\).
Maximal ideals have residue field \(k\). Localization identifies
the primes of \(A_{\mathfrak p}\) with the primes contained in
\(\mathfrak p\), so \(\dim A_{\mathfrak p}=\operatorname{ht}\mathfrak p\).
Every nonempty open subset of a variety has its dimension and
function field.

For an irreducible locally closed subset \(Y\), its closure has
the same function field and dimension. The localization
\(A_{\mathfrak p}\) is local with residue field
\(\operatorname{Frac}(A/\mathfrak p)\).
Eisenstein's criterion applies to polynomials over
\(k[t_1,\ldots,t_d]\).
\end{context}

\begin{problem}
For a subvariety \(Y\subseteq X\), consider pairs \((U,f)\),
where \(U\subseteq X\) is open, \(U\cap Y\ne\varnothing\), and \(f\)
is regular on \(U\); identify pairs whose functions agree on their
overlap. Prove that their ring \(\mathcal O_{Y,X}\) is local,
has residue field \(K(Y)\), and has dimension
\(\dim X-\dim Y\). Identify the cases \(Y=\{P\}\) and \(Y=X\),
and explain why \(K(Y)\) is not algebraically closed when
\(\dim Y>0\).
\end{problem}
\end{exerciseintro}

\begin{solution}
Each representative defines a rational function on the irreducible
variety \(X\). Two representatives are equivalent exactly when
they give the same rational function. Intersections of their domains
still meet \(Y\), since their intersections with irreducible \(Y\)
are nonempty opens. Thus sum and product of representatives define
a subring of \(K(X)\).

Choose an affine open \(W\) meeting \(Y\), put \(A=A(W)\), and
let \(\mathfrak p\) be the prime ideal of
\(\overline{Y\cap W}\) in \(W\).
We claim that this subring is \(A_{\mathfrak p}\).
For \(a/b\in A_{\mathfrak p}\), \(b\) does not vanish identically
on \(Y\cap W\). Thus \(D_W(b)\) meets \(Y\), and \(a/b\) is
an allowed representative.
Conversely take \((U,f)\) and choose a point in \(U\cap W\cap Y\).
Near it \(f=a/b\), with \(a,b\in A\) and \(b\) nonzero at that
point. Then \(b\notin\mathfrak p\), and the rational function of
this representative lies in \(A_{\mathfrak p}\).

This identifies the maximal ideal with
\(\mathfrak pA_{\mathfrak p}\), and the residue field with
\[
\operatorname{Frac}(A/\mathfrak p)=K(Y).
\]
The dimension formula gives
\[
\dim\mathcal O_{Y,X}
=\operatorname{ht}\mathfrak p
=\dim X-\dim Y.
\]
For \(Y=\{P\}\) this is the ordinary ring of germs at \(P\).
For \(Y=X\), \(\mathfrak p=0\), and it is \(K(X)\).

Finally let \(d=\dim Y>0\) and take a transcendence basis
\(t_1,\ldots,t_d\) of \(K(Y)/k\).
The extension of \(F=k(t_1,\ldots,t_d)\) to \(K(Y)\) is finite,
say of degree \(e\).
For every integer \(N>e\), the polynomial \(T^N-t_1\) is
irreducible over \(F\) by Eisenstein at \(t_1\).
If \(K(Y)\) were algebraically closed, it would contain a root,
and hence an \(F\)-subextension of degree \(N>e\), impossible.
\end{solution}

\begin{exerciseintro}
% record: H-I-03-014
\exerciseheading{Exercise I.3.14}
\addcontentsline{toc}{subsection}{Exercise I.3.14}

\begin{context}
Work over an algebraically closed field \(k\); varieties are irreducible
and carry the Zariski topology. A regular function is locally a quotient
of polynomials with nonvanishing denominator. A morphism is a continuous
map pulling back regular functions to regular functions.
Morphisms to an affine variety correspond contravariantly to
\(k\)-algebra homomorphisms of its coordinate ring into the ring of
global regular functions. Morphisms and regular functions can be
checked and glued on open covers. The local ring at a point consists
of germs of regular functions, with maximal ideal the germs vanishing
there. On an affine variety with coordinate ring \(A\), it is
\(A_{\mathfrak m}\). The function field is \(\operatorname{Frac}A\).
\end{context}

\begin{problem}
Let \(H\cong\mathbb P^n\) be a hyperplane in \(\mathbb P^{n+1}\)
and \(P\notin H\). Project a point \(Q\ne P\) onto the intersection
of the line \(PQ\) with \(H\).

(a) Prove that this projection is a morphism.

(b) For the twisted cubic
\([t:u]\mapsto[t^3:t^2u:tu^2:u^3]\), project from
\([0:0:1:0]\) onto \(z=0\), and find its plane image.
\end{problem}
\end{exerciseintro}

\begin{solution}
(a) Choose coordinates with \(H=\{x_{n+1}=0\}\) and
\(P=[0:\cdots:0:1]\).
The map is
\[
[x_0:\cdots:x_{n+1}]\longmapsto[x_0:\cdots:x_n].
\]
Away from \(P\) these coordinates are not all zero.
On every chart \(x_i\ne0\), \(i\leq n\), the target affine
coordinates are the regular ratios \(x_j/x_i\).
These charts cover the domain, proving that the map is a morphism.

(b) The specified point is not on the twisted cubic: \(x=w=0\)
would force \(t=u=0\).
After dropping \(z\), the parametrization becomes
\[
[t:u]\longmapsto[x:y:w]=[t^3:t^2u:u^3].
\]
It satisfies \(y^3=x^2w\).
Conversely on \(x\ne0\), normalize \(x=1\); every point of this
equation is \([1:a:a^3]\), obtained from \([t:u]=[1:a]\).
On \(x=0\) the equation forces \(y=0\), leaving \([0:0:1]\),
obtained from \([0:1]\). Thus the image is exactly
\(V_+(y^3-x^2w)\).
At \([0:0:1]\), the chart \(w=1\) has equation \(x^2=y^3\),
the cuspidal cubic, with normalization parameter
\(x=v^3,y=v^2\). This description holds in characteristics
two and three as well.
\end{solution}

\begin{exerciseintro}
% record: H-I-03-015
\exerciseheading{Exercise I.3.15}
\addcontentsline{toc}{subsection}{Exercise I.3.15}

\begin{context}
Work over an algebraically closed field \(k\); varieties are irreducible
and carry the Zariski topology. A regular function is locally a quotient
of polynomials with nonvanishing denominator. A morphism is a continuous
map pulling back regular functions to regular functions.
Morphisms to an affine variety correspond contravariantly to
\(k\)-algebra homomorphisms of its coordinate ring into the ring of
global regular functions. Morphisms and regular functions can be
checked and glued on open covers. The local ring at a point consists
of germs of regular functions, with maximal ideal the germs vanishing
there. On an affine variety with coordinate ring \(A\), it is
\(A_{\mathfrak m}\). The function field is \(\operatorname{Frac}A\).

The polynomial functions on an affine variety are its coordinate
ring. Tensoring vector spaces over a field preserves injections.
Noether normalization gives a polynomial subalgebra of dimension
\(\dim A\) over which a finitely generated integral \(k\)-algebra
\(A\) is a finite module. Integral finite extensions preserve dimension.
\end{context}

\begin{problem}
For affine varieties \(X\subseteq\mathbb A^n\) and
\(Y\subseteq\mathbb A^m\), equip \(X\times Y\subseteq\mathbb A^{n+m}\)
with its induced algebraic-set structure.

(a) Prove that it is irreducible.

(b) Identify its coordinate ring with \(A(X)\otimes_k A(Y)\).

(c) Verify the categorical product property among varieties.

(d) Prove \(\dim(X\times Y)=\dim X+\dim Y\).
\end{problem}
\end{exerciseintro}

\begin{solution}
(a) The equations of \(X\) in the first coordinates and \(Y\)
in the second show that \(X\times Y\) is closed.
Suppose \(X\times Y=Z_1\cup Z_2\) with \(Z_i\) closed, and put
\[
X_i=\{x\in X:\{x\}\times Y\subseteq Z_i\}.
\]
For fixed \(x\), the fiber is isomorphic to irreducible \(Y\),
so it lies in one \(Z_i\). Hence \(X=X_1\cup X_2\).
To see \(X_i\) is closed, take each polynomial defining \(Z_i\)
and express its restriction as
\(\sum_{\ell=1}^r a_\ell(x)b_\ell(y)\), choosing the finitely
many \(b_\ell\) linearly independent in \(A(Y)\).
The polynomial vanishes identically on \(\{x\}\times Y\)
exactly when all \(a_\ell(x)\) vanish.
Thus \(X_i\) is cut out by polynomial equations in \(X\).
Irreducibility of \(X\) gives \(X=X_i\) for one \(i\), and then
\(X\times Y=Z_i\).

(b) The homomorphism
\[
A(X)\otimes_k A(Y)\longrightarrow A(X\times Y),\qquad
a\otimes b\longmapsto((x,y)\mapsto a(x)b(y))
\]
is surjective because the coordinate functions generate the target.
If an element in its kernel is written
\(\sum a_\ell\otimes b_\ell\) with the \(b_\ell\) independent,
evaluation at every \(x\) shows \(a_\ell(x)=0\) for every \(\ell\).
Thus all \(a_\ell=0\); the homomorphism is injective.

(c) The projections are polynomial maps. Given morphisms
\(f:Z\to X\), \(g:Z\to Y\), the only possible map is
\(h(z)=(f(z),g(z))\).
Its affine coordinates are the regular coordinate functions
of \(f\) and \(g\), so it is a morphism to \(\mathbb A^{n+m}\)
with image in the closed subvariety \(X\times Y\).
It is therefore the required unique morphism.

(d) Choose normalization subalgebras
\(k[u_1,\ldots,u_r]\subseteq A(X)\) and
\(k[v_1,\ldots,v_s]\subseteq A(Y)\), with \(r=\dim X,s=\dim Y\).
Tensoring embeds
\(k[u_1,\ldots,u_r,v_1,\ldots,v_s]\) in the ring of (b).
Tensor products of finite module generating sets show that the
latter is finite over this polynomial subalgebra.
Therefore its dimension is \(r+s\), proving the formula.
\end{solution}

\begin{exerciseintro}
% record: H-I-03-016
\exerciseheading{Exercise I.3.16}
\addcontentsline{toc}{subsection}{Exercise I.3.16}

\begin{context}
Work over an algebraically closed field \(k\); varieties are irreducible
and carry the Zariski topology. A regular function is locally a quotient
of polynomials with nonvanishing denominator. A morphism is a continuous
map pulling back regular functions to regular functions.
Morphisms to an affine variety correspond contravariantly to
\(k\)-algebra homomorphisms of its coordinate ring into the ring of
global regular functions. Morphisms and regular functions can be
checked and glued on open covers. The local ring at a point consists
of germs of regular functions, with maximal ideal the germs vanishing
there. On an affine variety with coordinate ring \(A\), it is
\(A_{\mathfrak m}\). The function field is \(\operatorname{Frac}A\).

The Segre embedding identifies \(\mathbb P^n\times\mathbb P^m\)
with the closed variety of rank-one matrices in
\(\mathbb P^{(n+1)(m+1)-1}\). Its affine chart where the
\((i,j)\)-entry is nonzero is the product of the standard
affine charts \(x_i\ne0,y_j\ne0\), with coordinates obtained by
row and column ratios.
Affine products are irreducible and satisfy the categorical
product property. A quasi-projective variety is a locally
closed irreducible subset of projective space.
\end{context}

\begin{problem}
For quasi-projective \(X\subseteq\mathbb P^n\) and
\(Y\subseteq\mathbb P^m\), use the Segre embedding to give
\(X\times Y\) its induced structure.

(a) Show that \(X\times Y\) is quasi-projective.

(b) Show that it is projective if both factors are projective.

(c) Prove the categorical product property.
\end{problem}
\end{exerciseintro}

\begin{solution}
The Segre charts show that both projections are morphisms:
their coordinates on each chart are the row and column ratios.

(a) Write \(X=C\cap U\) and \(Y=D\cap V\), with \(C,D\) closed
and \(U,V\) open in their projective spaces.
Inside the Segre variety, \(X\times Y\) is the intersection of
the closed set \(p_1^{-1}(C)\cap p_2^{-1}(D)\) with the open set
\(p_1^{-1}(U)\cap p_2^{-1}(V)\), hence is locally closed.

Cover \(X\) and \(Y\) by nonempty affine opens \(X_a,Y_b\).
Their products are open in \(X\times Y\) by the projections.
The Segre chart structure on each agrees with the affine
product structure: the charts on the ambient product have
the ordinary coordinate ratios, and restriction gives those
on the subvarieties. Thus each \(X_a\times Y_b\) is irreducible.
Any two such opens intersect, since any two nonempty opens
of either factor intersect.
An open cover by irreducible sets with pairwise nonempty
intersections has irreducible union: two nonempty open subsets
of the union meet by passing through a common covering set.
Consequently \(X\times Y\) is irreducible and quasi-projective.

(b) If \(X,Y\) are closed, their inverse images under the
projections have closed intersection. Thus the product is
closed in the projective Segre variety and is projective.

(c) The projections restrict to morphisms on \(X\times Y\).
Given \(f:Z\to X\) and \(g:Z\to Y\), define
\(h(z)=(f(z),g(z))\).
The opens \(f^{-1}(X_a)\cap g^{-1}(Y_b)\) cover \(Z\).
On each, the affine product property says that \(h\) is a
morphism to \(X_a\times Y_b\).
These morphisms agree as maps on overlaps and glue to a
morphism \(Z\to X\times Y\). Its projections are \(f,g\),
and uniqueness follows pointwise.
\end{solution}

\begin{exerciseintro}
% record: H-I-03-017
\exerciseheading{Exercise I.3.17}
\addcontentsline{toc}{subsection}{Exercise I.3.17}

\begin{context}
Work over an algebraically closed field \(k\); varieties are irreducible
and carry the Zariski topology. A regular function is locally a quotient
of polynomials with nonvanishing denominator. A morphism is a continuous
map pulling back regular functions to regular functions.
Morphisms to an affine variety correspond contravariantly to
\(k\)-algebra homomorphisms of its coordinate ring into the ring of
global regular functions. Morphisms and regular functions can be
checked and glued on open covers. The local ring at a point consists
of germs of regular functions, with maximal ideal the germs vanishing
there. On an affine variety with coordinate ring \(A\), it is
\(A_{\mathfrak m}\). The function field is \(\operatorname{Frac}A\).

A domain is normal if it is integrally closed in its fraction field.
Polynomial rings over a field are unique factorization domains
and therefore normal. Localizations of normal domains are normal:
if \(z\) is integral over \(S^{-1}A\), clearing denominators in
a monic equation makes \(sz\) integral over \(A\) for some \(s\in S\).
For any domain \(A\),
\(A=\bigcap_{\mathfrak m\ \text{maximal}}A_{\mathfrak m}\)
inside its fraction field. Indeed, if \(z\notin A\), the proper
ideal \(\{a\in A:az\in A\}\) is contained in a maximal ideal
at which \(z\) is still absent.
A variety is normal at a point when its local ring is normal.

Every projective conic is isomorphic to \(\mathbb P^1\).
The integral closure of a finitely generated integral
\(k\)-algebra in its fraction field is finite as a module
over that algebra, and hence finitely generated over \(k\).
\end{context}

\begin{problem}
(a) Prove that projective plane conics are normal.

(b) Prove that the projective quadrics \(xy=zw\) and \(xy=z^2\)
in \(\mathbb P^3\) are normal.

(c) Show that the affine cusp \(y^2=x^3\) is not normal.

(d) For affine \(Y\), prove that \(Y\) is normal exactly when
\(A(Y)\) is normal.

(e) Construct the normalization \(\pi:\widetilde Y\to Y\)
of an affine variety and show that every dominant morphism
\(f:Z\to Y\) from a normal variety has a unique morphism
\(\widetilde f:Z\to\widetilde Y\) satisfying
\(\pi\widetilde f=f\).
\end{problem}
\end{exerciseintro}

\begin{solution}
(d) Put \(A=A(Y)\). If \(A\) is normal, so is each localization
\(A_{\mathfrak m}\), proving normality at each point.
Conversely, an element of \(\operatorname{Frac}A\) integral over
\(A\) is integral over every \(A_{\mathfrak m}\).
If all these local rings are normal, it belongs to their
intersection \(A\). This proves the equivalence.

(a) The two standard affine charts of \(\mathbb P^1\) have
coordinate ring \(k[t]\), a normal domain.
Normality is a condition on local rings, so \(\mathbb P^1\)
and every variety isomorphic to it are normal.

(b) For \(xy=zw\), each nonzero coordinate chart eliminates
one other coordinate from the equation and is \(\mathbb A^2\).
The charts cover, and all their local rings are normal.

For \(xy=z^2\), the charts \(x\ne0\) and \(y\ne0\) are likewise
\(\mathbb A^2\).
The chart \(w\ne0\) has ring
\[
R=k[x,y,z]/(xy-z^2)\cong k[s^2,t^2,st]\subseteq k[s,t]=B.
\]
To verify the isomorphism, reduce every polynomial modulo
\(z^2=xy\) to \(a(x,y)+z b(x,y)\). The images of its monomials
have respectively even-even and odd-odd exponent pairs in
\(s,t\), and are linearly independent. Thus the displayed
surjection to the subring has zero kernel.
The subring \(R\) consists of the polynomials all of whose
monomials have even total degree.
It satisfies
\[
B\cap\operatorname{Frac}R=R.
\]
Indeed, if \(b=p/q\in B\) with \(p,q\in R\), \(q\ne0\),
split \(b=b_0+b_1\) into its even- and odd-degree parts.
In \(qb=p\), the odd-degree part is \(qb_1=0\), so \(b_1=0\).
If an element of \(\operatorname{Frac}R\) is integral over \(R\),
it is integral over normal \(B\), hence lies in \(B\), and
the intersection identity puts it in \(R\).
Thus \(R\) is normal in every characteristic.
The three charts cover the quadric, since \(x=y=w=0\) also
forces \(z=0\). Part (d) proves normality everywhere.

(c) The cusp has coordinate ring \(A=k[t^2,t^3]\).
The element \(t=t^3/t^2\in\operatorname{Frac}A\) is integral,
since it satisfies \(T^2-t^2=0\), but it is not in \(A\).
Thus (d) proves that the cusp is not normal.
More precisely, it is not normal at its origin: if
\(t=a/b\in A_{(t^2,t^3)}\), the denominator \(b\) has a
nonzero constant term. Then \(tb\) has a nonzero coefficient
of \(t\), whereas \(a\in A\) has none, a contradiction.

(e) Let \(A=A(Y)\), \(K=\operatorname{Frac}A\), and let
\(\widetilde A\) be its integral closure in \(K\).
The finiteness theorem makes \(\widetilde A\) the coordinate
ring of an affine variety \(\widetilde Y\).
It is normal, so (d) makes \(\widetilde Y\) normal.
The inclusion \(A\subseteq\widetilde A\) induces \(\pi\).

Dominance of \(f\) gives an inclusion \(K\hookrightarrow K(Z)\).
For \(b\in\widetilde A\), its image is integral over
\(\mathcal O_{z,Z}\) for every \(z\in Z\), since the coefficients
of a monic equation for \(b\) come from \(A\) and pull back
to regular functions. Normality of each local ring puts this
rational function in every \(\mathcal O_{z,Z}\).
It is therefore regular on \(Z\).
We obtain a homomorphism
\(\widetilde A\to\mathcal O(Z)\), and hence a morphism
\(\widetilde f:Z\to\widetilde Y\) with \(\pi\widetilde f=f\).

For uniqueness, write any \(b\in\widetilde A\subseteq K\)
as \(c/a\) with \(a,c\in A\), \(a\ne0\).
Every lift must pull \(b\) back to \(f^*(c)/f^*(a)\) on the
nonempty open set where \(f^*(a)\ne0\); this set is nonempty
by dominance. Regular functions equal on a dense open set
are equal everywhere. Thus all coordinate pullbacks, and
hence the lifts, agree.
\end{solution}

\begin{exerciseintro}
% record: H-I-03-018
\exerciseheading{Exercise I.3.18}
\addcontentsline{toc}{subsection}{Exercise I.3.18}

\begin{context}
Work over an algebraically closed field \(k\); varieties are irreducible
and carry the Zariski topology. A regular function is locally a quotient
of polynomials with nonvanishing denominator. A morphism is a continuous
map pulling back regular functions to regular functions.
Morphisms to an affine variety correspond contravariantly to
\(k\)-algebra homomorphisms of its coordinate ring into the ring of
global regular functions. Morphisms and regular functions can be
checked and glued on open covers. The local ring at a point consists
of germs of regular functions, with maximal ideal the germs vanishing
there. On an affine variety with coordinate ring \(A\), it is
\(A_{\mathfrak m}\). The function field is \(\operatorname{Frac}A\).

A domain is normal if it is integrally closed in its fraction field.
Polynomial rings over a field are unique factorization domains
and therefore normal. Localizations of normal domains are normal:
if \(z\) is integral over \(S^{-1}A\), clearing denominators in
a monic equation makes \(sz\) integral over \(A\) for some \(s\in S\).
For any domain \(A\),
\(A=\bigcap_{\mathfrak m\ \text{maximal}}A_{\mathfrak m}\)
inside its fraction field. Indeed, if \(z\notin A\), the proper
ideal \(\{a\in A:az\in A\}\) is contained in a maximal ideal
at which \(z\) is still absent.
A variety is normal at a point when its local ring is normal.

If \(S\) is a homogeneous coordinate ring, the affine chart
\(x_i\ne0\) has ring \((S_{x_i})_0\), the degree-zero subring
of the graded localization. An affine variety is normal exactly
when its coordinate ring is normal.
Projective normality means integral closedness of \(S\).
\end{context}

\begin{problem}
(a) Prove that a projectively normal variety is normal.

(b) For the twisted quartic
\[
Y=\{[t^4:t^3u:tu^3:u^4]:[t:u]\in\mathbb P^1\}\subseteq\mathbb P^3,
\]
prove that \(Y\) is normal but not projectively normal.

(c) Show that \(Y\cong\mathbb P^1\), and conclude that
projective normality depends on the embedding.
\end{problem}
\end{exerciseintro}

\begin{solution}
(a) Suppose \(S\) is normal and put \(B=S_{x_i}\), \(A=B_0\)
for a nonempty chart.
The localization \(B\) is normal. If
\(f\in\operatorname{Frac}A\) is integral over \(A\), it is
integral over \(B\), so \(f\in B\).
Write \(f=p/q\) with \(p,q\in A\), \(q\ne0\), and decompose
\(f\) into homogeneous components in \(B\).
Since \(q\) has degree zero, \(qf=p\) forces all nonzero-degree
components of \(f\) to vanish. Hence \(f\in B_0=A\).
Every chart ring is normal, so the variety is normal.

(b) Denote the target coordinates by \(x,y,z,w\).
The image above is closed: it is exactly the common zero set of
\[
y^4=x^3w,\qquad z^4=xw^3,\qquad x^2z=y^3,\qquad w^2y=z^3.
\]
These relations hold on the parametrization.
They force at least one of \(x,w\) to be nonzero.
On \(x\ne0\), putting \(a=y/x\) gives the point
\([1:a:a^3:a^4]\); on \(w\ne0\), putting \(b=z/w\)
gives \([b^4:b^3:b:1]\).
These are precisely the two parametrizing charts.
Their inverse maps to \(\mathbb P^1\) are respectively
\([t:u]=[x:y]\) and \([t:u]=[z:w]\).
They are regular and agree on the overlap, proving
\(Y\cong\mathbb P^1\).
The latter has normal affine charts \(k[a]\) and \(k[b]\),
so \(Y\) is normal.

Its homogeneous coordinate ring is the image of the homogeneous
substitution homomorphism, namely
\[
R=k[t^4,t^3u,tu^3,u^4].
\]
Indeed a homogeneous polynomial vanishes on all parametrized
points exactly when its substituted polynomial vanishes on
all of \(k^2\), equivalently when that polynomial is zero,
since \(k\) is infinite.
The element \(v=t^2u^2\) belongs to \(\operatorname{Frac}R\),
because \(v=(t^3u)^2/t^4\).
It is integral over \(R\), since \(v^2=t^4u^4\in R\).
But \(v\notin R\): in total degree four, \(R\) is spanned by
the four displayed generators, none having exponent pair \((2,2)\).
Thus this embedding is not projectively normal.

(c) The inverse chart maps in (b) give \(Y\cong\mathbb P^1\).
In contrast, the standard embedding of \(\mathbb P^1\) has
normal homogeneous coordinate ring \(k[t,u]\).
This proves the embedding dependence.
\end{solution}

\begin{exerciseintro}
% record: H-I-03-019
\exerciseheading{Exercise I.3.19}
\addcontentsline{toc}{subsection}{Exercise I.3.19}

\begin{context}
Work over an algebraically closed field \(k\); varieties are irreducible
and carry the Zariski topology. A regular function is locally a quotient
of polynomials with nonvanishing denominator. A morphism is a continuous
map pulling back regular functions to regular functions.
Morphisms to an affine variety correspond contravariantly to
\(k\)-algebra homomorphisms of its coordinate ring into the ring of
global regular functions. Morphisms and regular functions can be
checked and glued on open covers. The local ring at a point consists
of germs of regular functions, with maximal ideal the germs vanishing
there. On an affine variety with coordinate ring \(A\), it is
\(A_{\mathfrak m}\). The function field is \(\operatorname{Frac}A\).

The Jacobian matrix of a polynomial map is the matrix of its
formal partial derivatives, and satisfies the chain rule over
any field. The units in \(k[x_1,\ldots,x_n]\) are the nonzero
constants.
\end{context}

\begin{problem}
Let \(\varphi:\mathbb A^n\to\mathbb A^n\) have polynomial
coordinates \(f_1,\ldots,f_n\), and write
\(J_\varphi=\det(\partial f_i/\partial x_j)\).

(a) If \(\varphi\) is an automorphism, prove that
\(J_\varphi\in k^\times\).

(b) Discuss the converse research question posed in the source,
including its necessary characteristic restriction.
\end{problem}
\end{exerciseintro}

\begin{solution}
(a) An inverse morphism \(\psi\) also has polynomial coordinates.
The chain rule applied to \(\psi\circ\varphi\) yields
\[
D\psi(\varphi(x))D\varphi(x)=I_n.
\]
Taking determinants gives
\[
J_\psi(\varphi(x))J_\varphi(x)=1
\]
in the polynomial ring. Thus \(J_\varphi\) is a unit,
and hence a nonzero constant.

(b) In characteristic \(p>0\), the converse is false.
For example
\[
(x_1,\ldots,x_n)\longmapsto(x_1-x_1^p,x_2,\ldots,x_n)
\]
has Jacobian determinant \(1\), but is not injective:
replacing \(x_1\) by \(x_1+a\), \(a\in\mathbb F_p\),
leaves its image unchanged.
In characteristic zero and dimension one, a polynomial with
nonzero constant derivative is linear, so the converse does hold.

The source's research question, called the Jacobian conjecture,
asks whether a polynomial map in characteristic zero with
nonzero constant Jacobian determinant must be an automorphism,
in dimensions at least two.
Part (a) establishes the necessary condition. Sufficiency in
dimension one and its failure in positive characteristic do
not settle the characteristic-zero question in higher dimension.
\end{solution}

\begin{exerciseintro}
% record: H-I-03-020
\exerciseheading{Exercise I.3.20}
\addcontentsline{toc}{subsection}{Exercise I.3.20}

\begin{context}
Work over an algebraically closed field \(k\); varieties are irreducible
and carry the Zariski topology. A regular function is locally a quotient
of polynomials with nonvanishing denominator. A morphism is a continuous
map pulling back regular functions to regular functions.
Morphisms to an affine variety correspond contravariantly to
\(k\)-algebra homomorphisms of its coordinate ring into the ring of
global regular functions. Morphisms and regular functions can be
checked and glued on open covers. The local ring at a point consists
of germs of regular functions, with maximal ideal the germs vanishing
there. On an affine variety with coordinate ring \(A\), it is
\(A_{\mathfrak m}\). The function field is \(\operatorname{Frac}A\).

A noetherian normal domain \(R\) satisfies the height-one intersection
theorem
\[
R=\bigcap_{\operatorname{ht}\mathfrak q=1}R_{\mathfrak q}
\quad\text{inside }\operatorname{Frac}R.
\]
For an affine variety with ring \(A\),
\(\operatorname{ht}\mathfrak p+\dim A/\mathfrak p=\dim A\);
localization preserves the chains of primes lying below a given
prime. In particular all closed-point local rings have dimension
equal to that of the variety. A positive-dimensional closed
subvariety cannot consist of just one point.
\end{context}

\begin{problem}
Let \(\dim Y\geq2\), and suppose \(Y\) is normal at \(P\).

(a) Show that every regular function on \(Y\setminus\{P\}\)
extends to a regular function on \(Y\).

(b) Give a counterexample in dimension one.
\end{problem}
\end{exerciseintro}

\begin{solution}
(a) Let \(f\) be the given function, viewed as an element of \(K(Y)\).
Choose an affine open \(U\) containing \(P\), put \(A=A(U)\),
and let \(\mathfrak m\) correspond to \(P\).
The local ring \(R=A_{\mathfrak m}\) is a noetherian normal
domain of dimension \(\dim Y\geq2\).

Consider a height-one prime of \(R\), corresponding to
\(\mathfrak p\subseteq\mathfrak m\) in \(A\).
It has \(\operatorname{ht}\mathfrak p=1\).
Consequently \(V_U(\mathfrak p)\) has dimension
\(\dim Y-1\geq1\), and contains a point \(Q\ne P\).
The regularity of \(f\) at \(Q\) supplies a representation
\(f=a/b\) with \(a,b\in A\) and \(b(Q)\ne0\).
Thus \(b\notin\mathfrak p\), and \(f\in A_{\mathfrak p}\).
This is exactly the localization of \(R\) at the selected
height-one prime.
The intersection theorem therefore implies \(f\in R\).

Represent this germ by a regular function \(g\) on a neighborhood
of \(P\). As rational functions, \(g=f\), so they agree on the
punctured part of that neighborhood. Gluing \(g\) to the given
function extends it to all of \(Y\).
The extension is unique because it agrees with \(f\) on the
dense open \(Y\setminus\{P\}\).

(b) On the normal curve \(\mathbb A^1\), the function \(1/t\)
is regular away from \(0\), but is not in the local ring
\(k[t]_{(t)}\). Indeed \(1/t=a/b\) with \(b(0)\ne0\)
would imply \(b=ta\), contradicting \(b(0)\ne0\).
It has no regular extension across \(0\).
\end{solution}

\begin{exerciseintro}
% record: H-I-03-021
\exerciseheading{Exercise I.3.21}
\addcontentsline{toc}{subsection}{Exercise I.3.21}

\begin{context}
Work over an algebraically closed field \(k\); varieties are irreducible
and carry the Zariski topology. A regular function is locally a quotient
of polynomials with nonvanishing denominator. A morphism is a continuous
map pulling back regular functions to regular functions.
Morphisms to an affine variety correspond contravariantly to
\(k\)-algebra homomorphisms of its coordinate ring into the ring of
global regular functions. Morphisms and regular functions can be
checked and glued on open covers. The local ring at a point consists
of germs of regular functions, with maximal ideal the germs vanishing
there. On an affine variety with coordinate ring \(A\), it is
\(A_{\mathfrak m}\). The function field is \(\operatorname{Frac}A\).

Products of varieties have their categorical product structure:
a pair of morphisms to \(X,Y\) gives a morphism to \(X\times Y\).
A group variety has multiplication and inversion morphisms
whose values give a group on points.
The variety \(\mathbb G_m=\mathbb A^1\setminus\{0\}\) has ring
\(k[t,t^{-1}]\); its products have the analogous Laurent coordinate
rings.
\end{context}

\begin{problem}
(a) Verify that \(\mathbb G_a=\mathbb A^1\), with addition,
is a group variety.

(b) Verify that \(\mathbb G_m=\mathbb A^1\setminus\{0\}\),
with multiplication, is a group variety.

(c) Give \(\operatorname{Hom}(X,G)\) a natural group structure
when \(G\) is a group variety.

(d) Identify \(\operatorname{Hom}(X,\mathbb G_a)\) with
the additive group of \(\mathcal O(X)\).

(e) Identify \(\operatorname{Hom}(X,\mathbb G_m)\) with
the group \(\mathcal O(X)^\times\).
\end{problem}
\end{exerciseintro}

\begin{solution}
(a) Addition \((a,b)\mapsto a+b\) and inversion \(a\mapsto-a\)
are polynomial morphisms. The field addition axioms give the
group laws, with identity \(0\).

(b) Multiplication \((a,b)\mapsto ab\) is a morphism on
\(\mathbb G_m\times\mathbb G_m\). The inverse \(a\mapsto a^{-1}\)
is regular because \(t^{-1}\in k[t,t^{-1}]\).
The field multiplication axioms on \(k^\times\) give the group
laws, with identity \(1\).

(c) For morphisms \(f,g:X\to G\), define
\[
(fg)(x)=f(x)g(x).
\]
This is a morphism: compose the paired map \((f,g)\) with
the multiplication morphism of \(G\).
The identity is the constant morphism with value \(e\), and
the inverse of \(f\) is the composition of \(f\) with inversion.
Associativity and the other laws follow at each point from
the laws in \(G\). Precomposition with a morphism \(X'\to X\)
is a group homomorphism, expressing the naturality.

(d) A morphism \(X\to\mathbb A^1\) is uniquely its coordinate
regular function. Pointwise addition of morphisms gives the
sum of these functions, identifying the group with
\((\mathcal O(X),+)\).

(e) A morphism to \(\mathbb G_m\) gives a regular function \(f\)
with regular reciprocal, so \(f\in\mathcal O(X)^\times\).
Conversely a unit \(f\) has no zeros and defines a morphism
to \(\mathbb A^1\) with image in \(\mathbb G_m\); locally its
reciprocal is regular, giving the morphism to this open subvariety.
Equivalently, it is the homomorphism
\(k[t,t^{-1}]\to\mathcal O(X)\), \(t\mapsto f\).
The two constructions are inverse and carry pointwise
multiplication to multiplication of units.
\end{solution}

\begin{exerciseintro}
\corpussection{I.4 Rational Maps}
\addcontentsline{toc}{section}{I.4 Rational Maps}

% record: H-I-04-001
\exerciseheading{Exercise I.4.1}
\addcontentsline{toc}{subsection}{Exercise I.4.1}

\begin{context}
Work over an algebraically closed field \(k\). Varieties are irreducible.
Regularity and the morphism property are local on the source.
Two morphisms to a variety agreeing on a dense open subset agree
everywhere on their common domain. A rational map is an equivalence
class of morphisms on nonempty open subsets, with this agreement relation.
Dominant rational maps correspond contravariantly to embeddings of
function fields over \(k\); a function-field isomorphism is equivalent
to birationality and to isomorphism of suitable nonempty open subsets.
Every point has an affine open neighborhood.
\end{context}

\begin{problem}
Glue regular functions on two open subsets when they agree on the
overlap. Deduce that a rational function has a largest open domain
on which it is regular.
\end{problem}
\end{exerciseintro}

\begin{solution}
For \(f\in\mathcal O(U)\), \(g\in\mathcal O(V)\) agreeing on
\(U\cap V\), define \(h\) by these functions on \(U\cup V\).
Agreement makes this well defined. At each point it agrees on a
neighborhood with either \(f\) or \(g\), so is regular.
The same reasoning glues any compatible family.
Representatives of one rational function agree on their overlaps:
their difference vanishes on a dense open subset, hence everywhere.
They glue on the open union \(D\) of their domains. This union contains
every other such domain and is the unique largest domain.
\end{solution}

\begin{exerciseintro}
% record: H-I-04-002
\exerciseheading{Exercise I.4.2}
\addcontentsline{toc}{subsection}{Exercise I.4.2}

\begin{context}
Work over an algebraically closed field \(k\). Varieties are irreducible.
Regularity and the morphism property are local on the source.
Two morphisms to a variety agreeing on a dense open subset agree
everywhere on their common domain. A rational map is an equivalence
class of morphisms on nonempty open subsets, with this agreement relation.
Dominant rational maps correspond contravariantly to embeddings of
function fields over \(k\); a function-field isomorphism is equivalent
to birationality and to isomorphism of suitable nonempty open subsets.
Every point has an affine open neighborhood.
\end{context}

\begin{problem}
Prove that a rational map \(X\dashrightarrow Y\) has a unique
largest open subset on which it is represented by a morphism.
\end{problem}
\end{exerciseintro}

\begin{solution}
Let \(D\) be the union of all representative domains. Representatives
agree on overlaps by the agreement theorem for morphisms.
They therefore glue to a map \(D\to Y\). Its inverse images of open
subsets are open on each member of the open cover, hence open.
Its pullbacks of regular functions are regular locally, hence regular.
Thus the glued map is a morphism. It represents the rational map
and contains every other domain, proving the assertion.
\end{solution}

\begin{exerciseintro}
% record: H-I-04-003
\exerciseheading{Exercise I.4.3}
\addcontentsline{toc}{subsection}{Exercise I.4.3}

\begin{context}
Work over an algebraically closed field \(k\). Varieties are irreducible.
Regularity and the morphism property are local on the source.
Two morphisms to a variety agreeing on a dense open subset agree
everywhere on their common domain. A rational map is an equivalence
class of morphisms on nonempty open subsets, with this agreement relation.
Dominant rational maps correspond contravariantly to embeddings of
function fields over \(k\); a function-field isomorphism is equivalent
to birationality and to isomorphism of suitable nonempty open subsets.
Every point has an affine open neighborhood.

The local rings of the affine plane are localizations of \(k[u,v]\).
In \(k[u,v]_{(u,v)}\), neither \(v/u\) nor \(u/v\) is present:
clearing a denominator nonzero at the origin contradicts unique
factorization.
\end{context}

\begin{problem}
(a) Determine the regular domain of \(x_1/x_0\) on \(\mathbb P^2\).

(b) Determine its domain as a rational map to \(\mathbb P^1\),
using the usual inclusion \(\mathbb A^1\subseteq\mathbb P^1\).
\end{problem}
\end{exerciseintro}

\begin{solution}
(a) On \(x_0\ne0\), this is the affine coordinate \(x_1/x_0\).
At \(x_0=0,x_1\ne0\), use \(x_1=1\); the function is \(1/u\)
with \(u=x_0/x_1\), absent from the local ring at \(u=0\).
At \([0:0:1]\), use \(u=x_0/x_2,v=x_1/x_2\); the function
is \(v/u\), excluded by the divisibility test.
Its full domain is therefore \(D_+(x_0)\).

(b) The map is \([x_0:x_1:x_2]\mapsto[x_0:x_1]\), a morphism
away from \(P=[0:0:1]\). It maps \(x_0=0,x_1\ne0\) to infinity.
If it extended to \(P\), one of the two target affine charts would
contain the value. Pulling back its coordinate would make either
\(v/u\) or \(u/v\) regular near \(P\), impossible.
Thus its largest domain is \(\mathbb P^2\setminus\{P\}\).
\end{solution}

\begin{exerciseintro}
% record: H-I-04-004
\exerciseheading{Exercise I.4.4}
\addcontentsline{toc}{subsection}{Exercise I.4.4}

\begin{context}
Work over an algebraically closed field \(k\). Varieties are irreducible.
Regularity and the morphism property are local on the source.
Two morphisms to a variety agreeing on a dense open subset agree
everywhere on their common domain. A rational map is an equivalence
class of morphisms on nonempty open subsets, with this agreement relation.
Dominant rational maps correspond contravariantly to embeddings of
function fields over \(k\); a function-field isomorphism is equivalent
to birationality and to isomorphism of suitable nonempty open subsets.
Every point has an affine open neighborhood.

Every irreducible projective conic is isomorphic to \(\mathbb P^1\).
A curve is rational when its function field is \(k(t)\).
Gauss's lemma tests polynomial irreducibility over a fraction field.
\end{context}

\begin{problem}
(a) Show that projective conics are rational.

(b) Show that \(y^2=x^3\) is rational.

(c) For \(Y=V_+(y^2z-x^2(x+z))\), show that projection from
\([0:0:1]\) onto \(z=0\) is birational to \(\mathbb P^1\).
\end{problem}
\end{exerciseintro}

\begin{solution}
(a) Its isomorphism to \(\mathbb P^1\) is birational.

(b) Put \(t=y/x\). The equation gives \(t^2=x,t^3=y\), so
the function field is \(k(t)\).

(c) On \(z=1,x\ne0\), put \(t=y/x\). The equation gives
\[
x=t^2-1,\qquad y=t(t^2-1).
\]
Projection is \([x:y]=[1:t]\), with inverse on a nonempty open
given by \(t\mapsto[t^2-1:t(t^2-1):1]\).
The cubic is irreducible: a factor of
\(y^2-x^2(x+1)\) over \(k(x)\) would make \(x+1\) a square
there, contrary to its odd order at \(x=-1\).
Gauss's lemma and homogenization apply.
Thus this is a birational map of varieties in every characteristic.
In characteristic two the tangent cone is a repeated line;
the designation \emph{nodal} requires characteristic different from two.
\end{solution}

\begin{exerciseintro}
% record: H-I-04-005
\exerciseheading{Exercise I.4.5}
\addcontentsline{toc}{subsection}{Exercise I.4.5}

\begin{context}
Work over an algebraically closed field \(k\). Varieties are irreducible.
Regularity and the morphism property are local on the source.
Two morphisms to a variety agreeing on a dense open subset agree
everywhere on their common domain. A rational map is an equivalence
class of morphisms on nonempty open subsets, with this agreement relation.
Dominant rational maps correspond contravariantly to embeddings of
function fields over \(k\); a function-field isomorphism is equivalent
to birationality and to isomorphism of suitable nonempty open subsets.
Every point has an affine open neighborhood.

The quadric \(xy=zw\) has rulings by lines, with distinct lines
of one ruling disjoint. Any two projective plane curves intersect.
\end{context}

\begin{problem}
Prove that \(Q=V_+(xy-zw)\subseteq\mathbb P^3\) is birational
to \(\mathbb P^2\), but not isomorphic to it.
\end{problem}
\end{exerciseintro}

\begin{solution}
On \(x=1\), its equation is \(y=zw\), giving an open
\(\mathbb A^2\), also an open of \(\mathbb P^2\).
Hence the varieties are birational.
An isomorphism would take two disjoint ruling lines to disjoint
closed irreducible curves in \(\mathbb P^2\), contradicting their
intersection property.
\end{solution}

\begin{exerciseintro}
% record: H-I-04-006
\exerciseheading{Exercise I.4.6}
\addcontentsline{toc}{subsection}{Exercise I.4.6}

\begin{context}
Work over an algebraically closed field \(k\). Varieties are irreducible.
Regularity and the morphism property are local on the source.
Two morphisms to a variety agreeing on a dense open subset agree
everywhere on their common domain. A rational map is an equivalence
class of morphisms on nonempty open subsets, with this agreement relation.
Dominant rational maps correspond contravariantly to embeddings of
function fields over \(k\); a function-field isomorphism is equivalent
to birationality and to isomorphism of suitable nonempty open subsets.
Every point has an affine open neighborhood.
\end{context}

\begin{problem}
For \(c:[x:y:z]\mapsto[yz:xz:xy]\):

(a) Prove that \(c\) is birational and its own inverse.

(b) Give open subsets on which it is an isomorphism.

(c) Find the full domains of \(c,c^{-1}\) and describe their morphisms.
\end{problem}
\end{exerciseintro}

\begin{solution}
(a) On \(xyz\ne0\), the second iterate is
\([x^2yz:xy^2z:xyz^2]=[x:y:z]\). Thus \(c\) is its own
rational inverse.

(b) Take \(U=V=\{xyz\ne0\}\). The formula and its identical
inverse are morphisms on this torus.

(c) The coordinates vanish simultaneously at just the three
coordinate vertices. Elsewhere they define a morphism, contracting
each punctured coordinate line to the opposite vertex.
At the vertex \(x=1,y=z=0\), the expression is \([yz:z:y]\).
An extension with first target coordinate nonzero would make
\(1/y,1/z\) regular at the origin. With second coordinate nonzero
it would make \(y/z\) regular; with third coordinate nonzero it
would make \(z/y\) regular. Unique factorization in
\(k[y,z]_{(y,z)}\) excludes all three alternatives.
Permuting coordinates handles the other vertices.
Thus both full domains are the plane minus its three coordinate
vertices. The contractions explain why these full-domain morphisms
are not isomorphisms.
\end{solution}

\begin{exerciseintro}
% record: H-I-04-007
\exerciseheading{Exercise I.4.7}
\addcontentsline{toc}{subsection}{Exercise I.4.7}

\begin{context}
Work over an algebraically closed field \(k\). Varieties are irreducible.
Regularity and the morphism property are local on the source.
Two morphisms to a variety agreeing on a dense open subset agree
everywhere on their common domain. A rational map is an equivalence
class of morphisms on nonempty open subsets, with this agreement relation.
Dominant rational maps correspond contravariantly to embeddings of
function fields over \(k\); a function-field isomorphism is equivalent
to birationality and to isomorphism of suitable nonempty open subsets.
Every point has an affine open neighborhood.
\end{context}

\begin{problem}
If \(\mathcal O_{Q,Y}\cong\mathcal O_{P,X}\) as \(k\)-algebras,
prove that some open neighborhoods of \(P,Q\) are isomorphic,
with \(P\) corresponding to \(Q\).
\end{problem}
\end{exerciseintro}

\begin{solution}
Choose affine neighborhoods with rings \(A,B\) and write the
isomorphism \(\theta:B_{\mathfrak n}\to A_{\mathfrak m}\).
It extends to a function-field isomorphism.
Images of finitely many algebra generators of \(B\) are regular
on a common neighborhood \(U_0\) of \(P\), defining a morphism
\(f:U_0\to Y\). The isomorphism carries maximal ideal to maximal
ideal and fixes \(k\), so evaluation gives \(f(P)=Q\).
The inverse construction gives \(g:V_0\to X\) with \(g(Q)=P\).
Shrink within the original affine neighborhoods.
Set \(U_1=U_0\cap f^{-1}(V_0)\),
\(V_1=V_0\cap g^{-1}(U_0)\).
The maps \(gf\) on \(U_1\) and \(fg\) on \(V_1\) induce identity
function-field maps and therefore equal the identity on dense opens,
hence everywhere. For \(x\in U_1\), the equality \(g(f(x))=x\in U_0\)
implies \(f(x)\in V_1\); the analogous inclusion holds for \(g\).
Thus these restrictions are inverse neighborhood isomorphisms.
\end{solution}

\begin{exerciseintro}
% record: H-I-04-008
\exerciseheading{Exercise I.4.8}
\addcontentsline{toc}{subsection}{Exercise I.4.8}

\begin{context}
Work over an algebraically closed field \(k\). Varieties are irreducible.
Regularity and the morphism property are local on the source.
Two morphisms to a variety agreeing on a dense open subset agree
everywhere on their common domain. A rational map is an equivalence
class of morphisms on nonempty open subsets, with this agreement relation.
Dominant rational maps correspond contravariantly to embeddings of
function fields over \(k\); a function-field isomorphism is equivalent
to birationality and to isomorphism of suitable nonempty open subsets.
Every point has an affine open neighborhood.

Noether normalization makes an affine \(d\)-dimensional coordinate
ring finite integral over a polynomial ring in \(d\) variables.
Lying over then gives a surjection on closed points over algebraically
closed \(k\). Proper closed subsets of an irreducible curve are finite.
An infinite cardinal equals each of its positive finite powers.
\end{context}

\begin{problem}
(a) Prove that every positive-dimensional variety over \(k\)
has cardinality \(|k|\).

(b) Prove that any two curves over \(k\) are homeomorphic.
\end{problem}
\end{exerciseintro}

\begin{solution}
(a) The field \(k\) is infinite. Projective space has finitely many
affine charts, each of cardinality \(|k|\), so an embedding in
projective space gives \(|X|\leq|k|\).
Take an affine open \(U\) of dimension \(d\geq1\).
Normalization and lying over give a surjection \(U\to\mathbb A^d(k)\):
the residue field above a maximal ideal is finite over \(k\), hence \(k\).
Therefore \(|X|\geq|U|\geq|k^d|=|k|\).

(b) The topology on a curve is cofinite. By (a) two curves have
a bijection; it preserves finite sets in both directions and is
therefore a homeomorphism.
\end{solution}

\begin{exerciseintro}
% record: H-I-04-009
\exerciseheading{Exercise I.4.9}
\addcontentsline{toc}{subsection}{Exercise I.4.9}

\begin{context}
Work over an algebraically closed field \(k\). Varieties are irreducible.
Regularity and the morphism property are local on the source.
Two morphisms to a variety agreeing on a dense open subset agree
everywhere on their common domain. A rational map is an equivalence
class of morphisms on nonempty open subsets, with this agreement relation.
Dominant rational maps correspond contravariantly to embeddings of
function fields over \(k\); a function-field isomorphism is equivalent
to birationality and to isomorphism of suitable nonempty open subsets.
Every point has an affine open neighborhood.

Over a perfect field, finitely generated function fields are separably
generated, and any generating set contains a separating transcendence
basis. Finite separable extensions have distinct embeddings into an
algebraic closure and satisfy the primitive element theorem.
Finite integral extensions have closed surjective maps on prime spectra.
A proper closed subset of a projective line is finite.
Height-one primes in a polynomial ring are principal.
\end{context}

\begin{problem}
For projective \(X\subseteq\mathbb P^n\) of dimension \(r\), with
\(n\geq r+2\), find a point outside \(X\) such that projection
to a hyperplane is birational onto its image.
Deduce that successive projections give a hypersurface in
\(\mathbb P^{r+1}\) birational to \(X\).
\end{problem}
\end{exerciseintro}

\begin{solution}
Choose coordinates so every coordinate vertex lies outside \(X\)
and \(X\) is not contained in \(x_0=0\).
These conditions can be imposed together: they are nonempty open
conditions on invertible matrices, an irreducible variety.
The vertex conditions are nonempty because \(X\) is proper, and
the hyperplane condition is nonempty because \(X\) is nonempty.
On \(x_0\ne0\), the ratios \(t_i=x_i/x_0\) generate \(K=K(X)\).
After permutation, \(t_1,\ldots,t_r\) are a separating basis.
Choose two remaining generators \(\beta_1,\beta_2\), and let \(L\)
be generated over \(k\) by all the other ratios.
Then \(K=L(\beta_1,\beta_2)\) is finite separable over \(L\).

For all but finitely many \(c\in k\),
\(\alpha=\beta_1+c\beta_2\) generates \(K/L\).
Indeed distinct \(L\)-embeddings agree on this element for at most
one \(c\), since agreement on both \(\beta_i\) makes them equal.
There are only finitely many embeddings.
Projection retaining \(x_0\), the coordinates defining \(L\),
and the linear combination defining \(\alpha\), has point center
\(P_c\) on the line joining the two remaining coordinate vertices.
This line is not contained in \(X\), so only finitely many of the
distinct \(P_c\) belong to \(X\).
Since \(k\) is infinite, choose \(c\) avoiding both exceptional sets.
Projection is then a morphism on \(X\); its image field is
\(L(\alpha)=K\), so it is birational to its image closure.

Its image is closed as well. Put the chosen center at the last
coordinate vertex. An equation of \(X\) nonvanishing there has
a nonzero \(x_n^d\) coefficient, and after scaling is monic in \(x_n\).
On each target chart \(x_i\ne0\), this makes \(x_n/x_i\) integral
over the algebra of retained coordinate ratios.
Thus the affine image on this chart is closed by lying over.
These charts cover the target, so the projective image is closed.
Iterate until the ambient dimension is \(r+1\); the image still
has dimension \(r\), and its height-one homogeneous ideal is
principal. It is the required hypersurface.
\end{solution}

\begin{exerciseintro}
% record: H-I-04-010
\exerciseheading{Exercise I.4.10}
\addcontentsline{toc}{subsection}{Exercise I.4.10}

\begin{context}
Work over an algebraically closed field \(k\). Varieties are irreducible.
Regularity and the morphism property are local on the source.
Two morphisms to a variety agreeing on a dense open subset agree
everywhere on their common domain. A rational map is an equivalence
class of morphisms on nonempty open subsets, with this agreement relation.
Dominant rational maps correspond contravariantly to embeddings of
function fields over \(k\); a function-field isomorphism is equivalent
to birationality and to isomorphism of suitable nonempty open subsets.
Every point has an affine open neighborhood.

The blow-up of \(\mathbb A^2\) at the origin is given by
\(xb=ya\) in \(\mathbb A^2\times\mathbb P^1\).
The strict transform is the closure of the inverse image away
from the origin. Irreducible curves have cofinite point topology.
\end{context}

\begin{problem}
Blow up \(Y=V(y^2-x^3)\) at the origin. Prove that its strict
transform is \(\mathbb A^1\), meets the exceptional curve in one
point, and maps homeomorphically but not isomorphically to \(Y\).
\end{problem}
\end{exerciseintro}

\begin{solution}
On \(a=1\), put \(v=b/a\); then \(y=xv\), and the total equation
is \(x^2(v^2-x)=0\). The strict transform is
\(x=v^2,y=v^3\), with coordinate ring \(k[v]\).
Its intersection with \(x=0\) is the single point \(v=0\).
Scheme-theoretically this intersection is defined by \(v^2=0\)
and has length two; the intersection is not transverse.
On \(b=1\), put \(u=a/b\); then \(x=yu\), and the strict transform
is \(1-yu^3=0\). It has \(u\ne0\), so lies entirely in the first
chart, with \(v=u^{-1}\). Thus the first chart covers the transform.
The projection is \(v\mapsto(v^2,v^3)\), bijective with inverse
\(v=y/x\) off the origin and one point over the origin.
Its bijection is a homeomorphism for the cofinite topologies.
The pullback ring has image \(k[v^2,v^3]\), not containing \(v\);
hence it is not an isomorphism.
\end{solution}

\begin{exerciseintro}
\corpussection{I.5 Nonsingular Varieties}
\addcontentsline{toc}{section}{I.5 Nonsingular Varieties}

% record: H-I-05-001
\exerciseheading{Exercise I.5.1}
\addcontentsline{toc}{subsection}{Exercise I.5.1}

\begin{context}
Work over an algebraically closed field \(k\).
For a reduced affine hypersurface \(f=0\), singular points are the
common zeros of \(f\) and all its first partial derivatives.
The multiplicity of a plane curve at the origin is the lowest
nonzero total degree of its equation; the corresponding homogeneous
form is its tangent cone. A regular local ring is characterized
by equality of its dimension and the dimension of its cotangent space.

Assume \(\operatorname{char}k\ne2\). A node has two distinct tangent
directions, a cusp has local model \(y^2=x^3\), and a tacnode has
two smooth branches with the same tangent and contact order two.
An ordinary triple point has three distinct tangent directions.
Real sketches refer to these equations over \(\mathbb R\).
In characteristic different from two every formal power-series unit
has a square root after choosing a square root of its constant term;
its remaining coefficients are determined recursively.
\end{context}

\begin{problem}
Find all singular points and sketch the real curves:

(a) \(x^2=x^4+y^4\).

(b) \(xy=x^6+y^6\).

(c) \(x^3=y^2+x^4+y^4\).

(d) \(x^2y+xy^2=x^4+y^4\).
Identify the four types of singularity at the origin.
\end{problem}
\end{exerciseintro}

\begin{solution}
(a) For \(f=x^2-x^4-y^4\), the derivatives are
\(2x(1-2x^2),-4y^3\).
Thus a singular point has \(y=0\); its equation and first derivative
then force \(x=0\). The leading term is \(x^2\), and
\(x^2(1-x^2)=y^4\) gives the two formal smooth branches
\(x=\pm y^2+\text{higher terms}\). The origin is a tacnode,
with common tangent \(x=0\).

(b) The derivatives are \(y-6x^5,x-6y^5\).
In characteristic three these immediately force the origin.
Otherwise a nonzero critical point has
\(xy=6x^6=6y^6\); the equation then gives
\(xy=xy/3\), impossible since \(2,3\ne0\).
If one coordinate is zero, the other is zero as well.
The leading term \(xy\) makes the origin a node with the coordinate
axes as tangents.

(c) The derivatives are \(x^2(3-4x),-2y(1+2y^2)\).
The origin is singular, with leading term \(-y^2\), and
\[
y^2(1+y^2)=x^3(1-x).
\]
Taking square roots of units in the completed local ring converts
this to a cusp \(Y^2=x^3\); its tangent is \(y=0\).
The remaining critical candidates have \(x=3/4\) and
\(y^2=-1/2\). The defining equation there equals \(91/256\).
Consequently there are exactly two additional singular points
in characteristics \(7\) or \(13\), and none otherwise.
At either extra point, the quadratic Taylor form is
\(-9u^2/8+2v^2\). Its two factors are distinct in these
characteristics, so both extra points are nodes.
Candidates with only one coordinate nonzero do not lie on the curve;
in characteristic three \(3/4=0\) adds no further candidate.

(d) Write \(f=g_3-g_4\), where
\(g_3=xy(x+y)\) and \(g_4=x^4+y^4\).
At a singular point Euler's identity gives \(3g_3-4g_4=0\),
while \(f=0\) gives \(g_3=g_4\); hence both vanish.
If \(xy(x+y)=0\), either a coordinate is zero or \(y=-x\).
The equation \(x^4+y^4=0\) then forces the origin, using
\(\operatorname{char}k\ne2\).
Its three distinct tangents are \(x=0,y=0,x+y=0\), so it is
an ordinary triple point.

\noindent\begin{minipage}{\linewidth}
Here are qualitative sketches over \(\mathbb R\); a dot marks the
origin in each panel. They illustrate the real branches, separately
from the singularity assertions over general \(k\):
\begin{center}
\begin{picture}(360,95)
\put(40,45){\makebox(0,0){\(\bullet\)}}
\put(130,45){\makebox(0,0){\(\bullet\)}}
\put(190,45){\makebox(0,0){\(\bullet\)}}
\put(310,40){\makebox(0,0){\(\bullet\)}}
\put(5,0){(a) tacnode}
\qbezier(40,45)(40,80)(15,75)
\qbezier(15,75)(-5,45)(15,15)
\qbezier(15,15)(40,10)(40,45)
\qbezier(40,45)(40,80)(65,75)
\qbezier(65,75)(85,45)(65,15)
\qbezier(65,15)(40,10)(40,45)
\put(105,0){(b) node}
\qbezier(130,45)(130,80)(155,80)
\qbezier(155,80)(180,45)(130,45)
\qbezier(130,45)(95,45)(95,20)
\qbezier(95,20)(130,-5)(130,45)
\put(193,0){(c) cusp}
\qbezier(190,45)(225,45)(235,70)
\qbezier(235,70)(260,90)(265,45)
\qbezier(265,45)(260,0)(235,20)
\qbezier(235,20)(225,45)(190,45)
\put(285,0){(d) triple}
\qbezier(310,40)(310,85)(335,85)
\qbezier(335,85)(370,40)(310,40)
\qbezier(310,40)(290,60)(285,50)
\qbezier(285,50)(280,40)(310,40)
\qbezier(310,40)(330,20)(320,15)
\qbezier(320,15)(310,10)(310,40)
\end{picture}
\end{center}
\end{minipage}
The first has two lobes on opposite sides of the vertical tangent;
the second has lobes in the first and third quadrants; the cusp has
one lobe on the positive \(x\)-side. The triple point has one larger
lobe in the first quadrant and two smaller lobes in the second and
fourth. These real sketches do not represent the extra points that
occur only in characteristics \(7,13\).
\end{solution}

\begin{exerciseintro}
% record: H-I-05-002
\exerciseheading{Exercise I.5.2}
\addcontentsline{toc}{subsection}{Exercise I.5.2}

\begin{context}
Work over an algebraically closed field \(k\).
For a reduced affine hypersurface \(f=0\), singular points are the
common zeros of \(f\) and all its first partial derivatives.
The multiplicity of a plane curve at the origin is the lowest
nonzero total degree of its equation; the corresponding homogeneous
form is its tangent cone. A regular local ring is characterized
by equality of its dimension and the dimension of its cotangent space.

Assume \(\operatorname{char}k\ne2\).
A nondegenerate quadratic cone gives a conical double point.
A transverse product of two distinct linear factors describes
an ordinary double curve locally.
\end{context}

\begin{problem}
Find and describe the singular loci in \(\mathbb A^3\):

(a) \(xy^2=z^2\).

(b) \(x^2+y^2=z^2\).

(c) \(xy+x^3+y^3=0\).
\end{problem}
\end{exerciseintro}

\begin{solution}
(a) The gradient is \((y^2,2xy,-2z)\), so the singular locus is
the line \(y=z=0\). At \((a,0,0)\), \(a\ne0\), the transverse
quadratic form is \(a y^2-z^2\), with distinct factors.
The parametrization
\((u,v)\mapsto(u^2,v,uv)\) has two points above a general point
of this singular line and only one above the origin.
At the origin the branches coalesce in the displayed equation:
this is the pinch point. Over \(\mathbb R\), for \(x>0\) the two
sheets are \(z=\pm\sqrt{x}\,y\); at \(x=0\) they meet, while
for \(x<0\) only the singular axis remains.

(b) The gradient \((2x,2y,-2z)\) vanishes only at the origin.
The equation is itself a nondegenerate quadratic form in three
variables, giving a conical double point (the real circular cone).

(c) The derivatives are \(y+3x^2,x+3y^2,0\).
In characteristic three they force \(x=y=0\).
Otherwise at a nonzero critical pair,
\(xy=-3x^3=-3y^3\), and the equation becomes \(xy/3=0\),
a contradiction. Thus the singular locus is the whole \(z\)-axis.
The transverse lowest form is \(xy\), two distinct lines.
Since the equation is independent of \(z\), the surface is the
product of a nodal plane curve with an affine line, giving the
double-line picture.
\end{solution}

\begin{exerciseintro}
% record: H-I-05-003
\exerciseheading{Exercise I.5.3}
\addcontentsline{toc}{subsection}{Exercise I.5.3}

\begin{context}
Work over an algebraically closed field \(k\).
For a reduced affine hypersurface \(f=0\), singular points are the
common zeros of \(f\) and all its first partial derivatives.
The multiplicity of a plane curve at the origin is the lowest
nonzero total degree of its equation; the corresponding homogeneous
form is its tangent cone. A regular local ring is characterized
by equality of its dimension and the dimension of its cotangent space.
\end{context}

\begin{problem}
(a) For a plane curve, prove that multiplicity one at a point
is equivalent to nonsingularity there.

(b) Compute the multiplicities at every singular point of the
four curves \(x^2-x^4-y^4=0\), \(xy-x^6-y^6=0\),
\(x^3-y^2-x^4-y^4=0\), and \(xy(x+y)-x^4-y^4=0\),
assuming characteristic different from two.
\end{problem}
\end{exerciseintro}

\begin{solution}
(a) Translate the point to the origin and expand
\(f=f_1+f_2+\cdots\), with \(f_i\) homogeneous of degree \(i\).
The coefficient vector of \(f_1\) is \((f_x(0),f_y(0))\).
Thus \(f_1\ne0\) exactly when the Jacobian criterion says the
point is nonsingular, and exactly when its multiplicity is one.

(b) At the origin the respective lowest forms are
\(x^2,xy,-y^2,xy(x+y)\).
The multiplicities are therefore \(2,2,2,3\).
The derivative equations give no other singularities except on
the third curve in characteristics \(7,13\), where the two extra
points have \(x=3/4,y^2=-1/2\).
After translation their quadratic forms are
\(-9u^2/8+2v^2\), nonzero, so their multiplicities are two.
\end{solution}

\begin{exerciseintro}
% record: H-I-05-004
\exerciseheading{Exercise I.5.4}
\addcontentsline{toc}{subsection}{Exercise I.5.4}

\begin{context}
Work over an algebraically closed field \(k\).
For a reduced affine hypersurface \(f=0\), singular points are the
common zeros of \(f\) and all its first partial derivatives.
The multiplicity of a plane curve at the origin is the lowest
nonzero total degree of its equation; the corresponding homogeneous
form is its tangent cone. A regular local ring is characterized
by equality of its dimension and the dimension of its cotangent space.

The ambient local ring at a plane point is
\(R=k[x,y]_{(x,y)}\) after translation.
Length of a finite-length \(R\)-module equals its dimension over \(k\).
The associated graded ring of \(R\) is \(k[x,y]\).
For a principal ideal \((f)\), its initial ideal is generated by
the lowest homogeneous form of \(f\), because orders add under
multiplication. Distinct irreducible plane curves have no common
one-dimensional component.
\end{context}

\begin{problem}
For distinct irreducible plane curves \(Y,Z\) through \(P\), with local
equations \(f,g\), define \((Y\cdot Z)_P=\ell(R/(f,g))\).

(a) Prove finiteness and the inequality
\((Y\cdot Z)_P\geq\mu_P(Y)\mu_P(Z)\).

(b) Show that all but finitely many lines through \(P\) meet
\(Y\) with local multiplicity \(\mu_P(Y)\).

(c) If a projective plane curve has degree \(d\), prove that
its total intersection multiplicity with a line not a component
is \(d\).
\end{problem}
\end{exerciseintro}

\begin{solution}
(a) The absence of a common component makes \((f,g)\) primary
to the maximal ideal after localization at \(P\).
Its quotient is zero-dimensional noetherian, hence Artinian
and of finite length.
Put \(m=\operatorname{ord}f,n=\operatorname{ord}g\) and \(A=R/(f)\),
with maximal ideal \(\mathfrak a\).
Since \(\operatorname{gr}A=k[x,y]/(f_m)\), direct counting gives
\[
\ell(A/\mathfrak a^N)=mN-\frac{m(m-1)}2\qquad(N\geq m).
\]
Indeed the degree-\(j\) piece has dimension \(j+1\) for \(j<m\)
and \(m\) for \(j\geq m\).
Multiplication by \(g\) induces
\[
A/\mathfrak a^{N-n}\longrightarrow A/\mathfrak a^N
\longrightarrow A/(g,\mathfrak a^N)\longrightarrow0.
\]
For large \(N\), the last term is \(A/(g)\).
Its length is at least the difference of the first two lengths,
which is \(mn\), proving the bound without assuming the first map
injective.

(b) A line of direction \((a,b)\) is parametrized by \((at,bt)\).
If \(f_m(a,b)\ne0\), the restricted series has order \(m\), and
the quotient of \(k[t]_{(t)}\) has length \(m\).
Only the finitely many zeros of the nonzero binary form \(f_m\)
in \(\mathbb P^1\) are excluded.

(c) Restrict the degree-\(d\) homogeneous equation to the line
\(\mathbb P^1\). It is a nonzero binary form of degree \(d\),
whose linear factors over \(k\), counted with multiplicity,
number \(d\). In each affine chart, the order of a root is
exactly the local quotient length defining the intersection.
Summing gives \(d\).
\end{solution}

\begin{exerciseintro}
% record: H-I-05-005
\exerciseheading{Exercise I.5.5}
\addcontentsline{toc}{subsection}{Exercise I.5.5}

\begin{context}
Work over an algebraically closed field \(k\).
For a reduced affine hypersurface \(f=0\), singular points are the
common zeros of \(f\) and all its first partial derivatives.
The multiplicity of a plane curve at the origin is the lowest
nonzero total degree of its equation; the corresponding homogeneous
form is its tangent cone. A regular local ring is characterized
by equality of its dimension and the dimension of its cotangent space.

For a projective plane hypersurface, singularity is detected by
vanishing of its homogeneous gradient at a point on the hypersurface.
Two positive-degree projective plane curves meet; consequently a
reducible or nonreduced plane equation has a singular point.
\end{context}

\begin{problem}
For each integer \(d>0\) and each characteristic, give a
nonsingular projective plane curve of degree \(d\).
\end{problem}
\end{exerciseintro}

\begin{solution}
If the characteristic does not divide \(d\) (including characteristic
zero), take
\[
F=x^d+y^d+z^d.
\]
For \(d>1\), its gradient can vanish only when all coordinates vanish;
for \(d=1\), the gradient is a nonzero constant vector.

If the characteristic divides \(d\), then \(d\geq2\). Take
\[
F=x^d+y^{d-1}z+xz^{d-1}.
\]
Its derivative \(F_x=z^{d-1}\) forces \(z=0\) at a possible
singular point. For \(d>2\), \(F_z=y^{d-1}\) there forces \(y=0\),
and \(F=0\) then forces \(x=0\).
For \(d=2\) in characteristic two, \(F_z=y+x\);
together with \(z=0,F=0\) it again forces \(x=y=0\).
No projective singular point exists.
The intersection observation in the context proves irreducibility
of each displayed equation, so these are nonsingular curves.
\end{solution}

\begin{exerciseintro}
% record: H-I-05-006
\exerciseheading{Exercise I.5.6}
\addcontentsline{toc}{subsection}{Exercise I.5.6}

\begin{context}
Work over an algebraically closed field \(k\).
For a reduced affine hypersurface \(f=0\), singular points are the
common zeros of \(f\) and all its first partial derivatives.
The multiplicity of a plane curve at the origin is the lowest
nonzero total degree of its equation; the corresponding homogeneous
form is its tangent cone. A regular local ring is characterized
by equality of its dimension and the dimension of its cotangent space.

The blow-up of the plane at the origin has charts \(y=xt\) and
\(x=ys\). The strict-transform equation is obtained by removing
the largest common exceptional power from the substituted equation.
It agrees with the original curve away from the origin.
Assume characteristic different from two for the named node and tacnode.
Here resolution means making the strict transform nonsingular.
It does not require the total transform to have simple normal crossings;
that embedded-resolution condition can require further blow-ups.
\end{context}

\begin{problem}
(a) Blow up the node \(xy=x^6+y^6\) and cusp
\(x^3=y^2+x^4+y^4\) at the origin and show the points above the
origin are nonsingular. Determine when their whole strict transforms
are nonsingular.

(b) Prove that blowing up an ordinary double point gives two
distinct nonsingular points above it.

(c) Show that blowing up the tacnode \(x^2=x^4+y^4\) gives a node,
which one more blow-up resolves.

(d) Resolve the triple cusp \(y^3=x^5\) by two blow-ups.
\end{problem}
\end{exerciseintro}

\begin{solution}
(a) For the node, the \(x\)-chart has equation
\(t-x^4(1+t^6)=0\). On the exceptional line \(x=0\), \(t=0\)
and the \(t\)-derivative is one.
The \(y\)-chart gives a second smooth point in the other tangent
direction. The original curve has no singularity away from the
origin, so this strict transform is nonsingular.
For the cusp, the \(x\)-chart gives
\[
x-t^2-x^2-x^2t^4=0.
\]
There is one point above the origin, \(x=t=0\), with \(x\)-derivative
one. The other chart has no additional exceptional point.
The whole strict transform is nonsingular except in characteristics
\(7,13\): the two other original nodes
\(x=3/4,y^2=-1/2\) are unaffected by this blow-up.
This qualifies the global nonsingularity assertion for those fields.

(b) Choose coordinates so neither tangent is the vertical direction.
If \(f=f_2+f_3+\cdots\), the \(x\)-chart equation is
\[
f_2(1,t)+x f_3(1,t)+\cdots=0.
\]
The two exceptional points are the distinct roots of \(f_2(1,t)\).
At each the \(t\)-derivative is nonzero, proving nonsingularity.
The other chart gives no further point, by the choice of coordinates.

(c) In the \(y\)-chart, put \(x=yu\). The equation becomes
\[
u^2-y^2(u^4+1)=0.
\]
Its unique exceptional point is \(u=y=0\), with leading form
\(u^2-y^2\), two distinct lines. This is a node, resolved by (b).
The \(x\)-chart adds no exceptional point.

(d) The original lowest degree is three.
Putting \(y=xt\) gives \(t^3=x^2\), a double cusp at the sole
exceptional point. Blow that point up in the chart \(x=ts\);
the strict transform is \(t=s^2\), a smooth curve.
In each other chart there is no additional exceptional point.
These calculations work in every characteristic; the final equation
has a partial derivative equal to one.
\end{solution}

\begin{exerciseintro}
% record: H-I-05-007
\exerciseheading{Exercise I.5.7}
\addcontentsline{toc}{subsection}{Exercise I.5.7}

\begin{context}
Work over an algebraically closed field \(k\).
For a reduced affine hypersurface \(f=0\), singular points are the
common zeros of \(f\) and all its first partial derivatives.
The multiplicity of a plane curve at the origin is the lowest
nonzero total degree of its equation; the corresponding homogeneous
form is its tangent cone. A regular local ring is characterized
by equality of its dimension and the dimension of its cotangent space.

Let \(F(x,y,z)\) be homogeneous of degree \(d>1\), defining a
nonsingular projective plane curve. The blow-up of affine space
at its origin records the projective direction \([a:b:c]\);
on \(a\ne0\) it has coordinates \((x,v,w)\) with \(y=xv,z=xw\).
\end{context}

\begin{problem}
Let \(X=V(F)\subseteq\mathbb A^3\) and blow up its vertex \(O\).

(a) Show that \(O\) is its only singular point.

(b) Show that the strict transform \(\widetilde X\) is nonsingular.

(c) Identify the exceptional fiber with \(Y=V_+(F)\).
\end{problem}
\end{exerciseintro}

\begin{solution}
(a) All first derivatives vanish at the origin since \(d>1\).
At a nonzero point of the cone, projective nonsingularity says at
least one partial derivative is nonzero. Thus just the vertex is singular.

(b) On the direction chart \(a\ne0\), substitution gives
\(x^dF(1,v,w)\); the strict transform is \(F(1,v,w)=0\).
It is \(Y\cap\{a\ne0\}\) times \(\mathbb A^1_x\), nonsingular
by the affine Jacobian criterion. The charts \(b\ne0,c\ne0\)
give the same description and cover the transform.

(c) On the first chart the exceptional fiber has \(x=0\) and
\(F(1,v,w)=0\). Its direction coordinates are \([1:v:w]\).
On overlaps these are exactly the usual transitions of \(Y\),
so the chart identifications glue to an isomorphism of the
exceptional fiber with \(Y\).
\end{solution}

\begin{exerciseintro}
% record: H-I-05-008
\exerciseheading{Exercise I.5.8}
\addcontentsline{toc}{subsection}{Exercise I.5.8}

\begin{context}
Work over an algebraically closed field \(k\).
For a reduced affine hypersurface \(f=0\), singular points are the
common zeros of \(f\) and all its first partial derivatives.
The multiplicity of a plane curve at the origin is the lowest
nonzero total degree of its equation; the corresponding homogeneous
form is its tangent cone. A regular local ring is characterized
by equality of its dimension and the dimension of its cotangent space.

An affine \(r\)-dimensional variety in \(\mathbb A^n\) is nonsingular
where the Jacobian of generators of its ideal has rank \(n-r\).
For homogeneous \(f\) of degree \(d\), Euler's identity is
\(\sum x_i f_{x_i}=df\), in every characteristic.
\end{context}

\begin{problem}
Let homogeneous \(f_1,\ldots,f_s\) generate the ideal of an
\(r\)-dimensional \(Y\subseteq\mathbb P^n\). Prove that \(P=[a]\)
is nonsingular exactly when \((f_{i,x_j}(a))\) has rank \(n-r\),
and show that its rank is independent of the representative \(a\).
\end{problem}
\end{exerciseintro}

\begin{solution}
Replacing \(a\) by \(\lambda a\) multiplies row \(i\) by
\(\lambda^{\deg f_i-1}\), a nonzero scalar, preserving rank.
Choose \(a_0\ne0\) and normalize it to one.
On this affine chart the equations are \(g_i(t)=f_i(1,t)\);
they generate the affine ideal by localization and dehomogenization.
Their Jacobian is the matrix formed by columns \(1,\ldots,n\)
of the homogeneous Jacobian. Euler's identity on \(Y\) gives
\[
f_{i,x_0}(a)=-\sum_{j=1}^n a_j f_{i,x_j}(a).
\]
Thus column zero lies in the span of the other columns, and the
two Jacobians have equal rank. The affine criterion proves the claim.
No division by \(\deg f_i\) occurs, so the proof includes degrees
divisible by the characteristic.
\end{solution}

\begin{exerciseintro}
% record: H-I-05-009
\exerciseheading{Exercise I.5.9}
\addcontentsline{toc}{subsection}{Exercise I.5.9}

\begin{context}
Work over an algebraically closed field \(k\).
For a reduced affine hypersurface \(f=0\), singular points are the
common zeros of \(f\) and all its first partial derivatives.
The multiplicity of a plane curve at the origin is the lowest
nonzero total degree of its equation; the corresponding homogeneous
form is its tangent cone. A regular local ring is characterized
by equality of its dimension and the dimension of its cotangent space.

Two projective plane curves of positive degree intersect.
The homogeneous gradient criterion detects nonsingularity on
a projective plane hypersurface.
\end{context}

\begin{problem}
Let \(f\) be a nonzero homogeneous polynomial of positive degree
in three variables. If its gradient is nonzero at every point
of \(V_+(f)\), prove that \(f\) is irreducible.
\end{problem}
\end{exerciseintro}

\begin{solution}
If \(f=gh\) with both factors of positive degree, choose a
projective common zero of \(g,h\). The product rule
\(f_{x_i}=g h_{x_i}+h g_{x_i}\) makes every partial vanish there.
This contradicts the hypothesis. The same argument includes repeated
factors, since their common zero set is nonempty.
Thus \(f\) is irreducible, and its projective zero set is nonsingular
by the stated gradient criterion.
\end{solution}

\begin{exerciseintro}
% record: H-I-05-010
\exerciseheading{Exercise I.5.10}
\addcontentsline{toc}{subsection}{Exercise I.5.10}

\begin{context}
Work over an algebraically closed field \(k\).
For a reduced affine hypersurface \(f=0\), singular points are the
common zeros of \(f\) and all its first partial derivatives.
The multiplicity of a plane curve at the origin is the lowest
nonzero total degree of its equation; the corresponding homogeneous
form is its tangent cone. A regular local ring is characterized
by equality of its dimension and the dimension of its cotangent space.

At a closed point of a variety, the local-ring dimension equals the
variety dimension. Nakayama's lemma identifies the minimum number of
generators of the maximal ideal \(\mathfrak m\) with
\(\dim_k\mathfrak m/\mathfrak m^2\), which is at least the ring
dimension by the height theorem. Nonsingularity is equivalent to
regularity of the local ring.
\end{context}

\begin{problem}
Define \(T_PX=(\mathfrak m_P/\mathfrak m_P^2)^\vee\).

(a) Prove \(\dim T_PX\geq\dim X\), with equality exactly at
nonsingular points.

(b) Construct the natural linear map \(T_P\varphi:T_PX\to
T_{\varphi(P)}Y\) for a morphism \(\varphi:X\to Y\).

(c) Show that the projection of \(x=y^2\) onto the \(x\)-axis
has zero tangent map at the origin.
\end{problem}
\end{exerciseintro}

\begin{solution}
(a) Duality preserves finite vector-space dimension.
The inequality follows from the generator bound in the context.
Equality is exactly the defining condition for the local ring to
be regular, hence for the point to be nonsingular.

(b) The local homomorphism
\(\mathcal O_{\varphi(P),Y}\to\mathcal O_{P,X}\) sends maximal
ideal to maximal ideal and its square into the square.
It induces a \(k\)-linear map on their quotients; dualizing gives
\(T_P\varphi\). Functoriality follows from composition of these
local homomorphisms.

(c) The coordinate ring of the parabola is \(k[y]\), with \(x=y^2\).
Pullback sends the cotangent generator \(x\bmod x^2\) of the axis
to \(y^2\bmod(y)^2=0\). The cotangent map and its dual are zero,
including in characteristic two.
\end{solution}

\begin{exerciseintro}
% record: H-I-05-011
\exerciseheading{Exercise I.5.11}
\addcontentsline{toc}{subsection}{Exercise I.5.11}

\begin{context}
Work over an algebraically closed field \(k\).
For a reduced affine hypersurface \(f=0\), singular points are the
common zeros of \(f\) and all its first partial derivatives.
The multiplicity of a plane curve at the origin is the lowest
nonzero total degree of its equation; the corresponding homogeneous
form is its tangent cone. A regular local ring is characterized
by equality of its dimension and the dimension of its cotangent space.

Assume \(\operatorname{char}k\ne2\).
For two equations in affine three-space, independent gradients give
a nonsingular one-dimensional zero set locally. A plane curve
with nonvanishing homogeneous gradient is irreducible.
\end{context}

\begin{problem}
Let \(Y\subseteq\mathbb P^3\) be cut out by
\(x^2-xz-yw=0,\ yz-xw-zw=0\), and put \(P=[0:0:0:1]\).
Show that projection from \(P\) identifies \(Y\setminus\{P\}\)
with \(C\setminus\{S\}\), where
\[
C=V_+(y^2z-x^3+xz^2),\qquad S=[1:0:-1].
\]
Prove that \(Y\) is an irreducible nonsingular curve.
Explain why characteristic two must be excluded.
\end{problem}
\end{exerciseintro}

\begin{solution}
Eliminating \(w\) gives the displayed cubic equation.
For a point \([x:y:z]\) of \(C\), the desired value of \(w\)
satisfies
\[
yw=x(x-z),\qquad (x+z)w=yz.
\]
These are compatible precisely because it lies on \(C\).
On \(y\ne0\) use \(w=x(x-z)/y\), and on \(x+z\ne0\)
use \(w=yz/(x+z)\), interpreted in homogeneous coordinates.
They agree on the overlap and are regular inverse maps.
The two opens cover \(C\setminus\{S\}\).
Over \(S\) the first equation would say \(0=2\), so no point
of \(Y\setminus\{P\}\) maps there.
This proves the claimed isomorphism.

The cubic \(C\) is nonsingular. On \(z=1\), its equation is
\(y^2=x^3-x\). A critical point has \(y=0,3x^2=1\);
the equation requires \(x=0,1,-1\), none satisfying both
conditions in characteristic different from two.
At infinity its only point is \([0:1:0]\), where the \(z\)-partial
is nonzero. Thus \(C\) is irreducible and nonsingular.
Consequently \(Y\setminus\{P\}\) is an irreducible nonsingular curve.

At \(P\), set \(w=1\). The two gradients in \(x,y,z\) are
\((0,-1,0)\) and \((-1,0,-1)\), linearly independent.
Hence \(Y\) is a nonsingular curve near \(P\); in particular
\(P\) is not an isolated component. It belongs to the closure of
\(Y\setminus\{P\}\), proving irreducibility of the whole \(Y\)
and nonsingularity everywhere.

In characteristic two, the line \(x=z,y=0\) lies in \(Y\),
and its points other than \(P\) project to \(S=[1:0:1]\).
The proposed isomorphism fails. Thus the characteristic
qualification is essential.
\end{solution}

\begin{exerciseintro}
% record: H-I-05-012
\exerciseheading{Exercise I.5.12}
\addcontentsline{toc}{subsection}{Exercise I.5.12}

\begin{context}
Work over an algebraically closed field \(k\).
For a reduced affine hypersurface \(f=0\), singular points are the
common zeros of \(f\) and all its first partial derivatives.
The multiplicity of a plane curve at the origin is the lowest
nonzero total degree of its equation; the corresponding homogeneous
form is its tangent cone. A regular local ring is characterized
by equality of its dimension and the dimension of its cotangent space.

Assume \(\operatorname{char}k\ne2\) and \(k\) algebraically closed.
A nonzero quadratic form corresponds to a symmetric bilinear form;
successive orthogonal decompositions diagonalize it.
Nonzero scalars in \(k\) have square roots.
\end{context}

\begin{problem}
For a nonzero homogeneous quadratic \(f\) in \(n+1\) variables:

(a) Put it in the form \(x_0^2+\cdots+x_r^2\) by linear coordinates.

(b) Determine when it is irreducible.

(c) In the irreducible case determine its singular locus.

(d) Retain the irreducible case and assume \(2\leq r<n\).
Describe the quadric as a cone with linear axis
over a nonsingular quadric in \(\mathbb P^r\).
\end{problem}
\end{exerciseintro}

\begin{solution}
(a) A nonzero symmetric bilinear form in characteristic different
from two has a vector \(v\) with nonzero self-pairing: otherwise
polarization would make the whole form zero.
Split off its orthogonal one-dimensional summand and repeat on
the complement. Rescale each nonzero diagonal coefficient by a
square root. The number of nonzero coefficients is \(r+1\).

(b) With one coefficient the form is a square; with two it factors
over \(k\). Conversely a product of two linear forms has bilinear
matrix of rank at most two, since its image lies in the span of
their coefficient vectors. Thus rank at least three, equivalently
\(r\geq2\), is exactly irreducibility.

(c) Its gradient is \((2x_0,\ldots,2x_r,0,\ldots,0)\).
The singular locus is
\(Z=V_+(x_0,\ldots,x_r)\), of dimension \(n-r-1\).
For \(r=n\) it is empty (dimension convention \(-1\)); otherwise
it is the stated linear space.

(d) Since \(r\geq2\), the base quadric below is nonempty and
irreducible. In the complementary subspace
\(\mathbb P^r=\{x_{r+1}=\cdots=x_n=0\}\), the same equation
defines a nonsingular quadric \(Q'\).
Every point of the full quadric outside \(Z\) has nonzero
first coordinate block representing a point of \(Q'\), and its
remaining block either is zero or represents a point of \(Z\).
It lies on their joining line. Conversely every such line
satisfies the equation, since it involves only the first block.
Including its endpoints gives exactly the cone with axis \(Z\).
\end{solution}

\begin{exerciseintro}
% record: H-I-05-013
\exerciseheading{Exercise I.5.13}
\addcontentsline{toc}{subsection}{Exercise I.5.13}

\begin{context}
Work over an algebraically closed field \(k\).
For a reduced affine hypersurface \(f=0\), singular points are the
common zeros of \(f\) and all its first partial derivatives.
The multiplicity of a plane curve at the origin is the lowest
nonzero total degree of its equation; the corresponding homogeneous
form is its tangent cone. A regular local ring is characterized
by equality of its dimension and the dimension of its cotangent space.

The normalization \(B\) of a finitely generated integral \(k\)-algebra
\(A\) in its fraction field is a finite \(A\)-module.
Integral closure commutes with localization (clear denominators
in a monic integral equation).
The support of a finite module \(M\) is \(V(\operatorname{Ann}M)\).
An affine variety is normal at its maximal ideal precisely when
that localization is integrally closed.
\end{context}

\begin{problem}
Without using existence of nonsingular points, prove that the
nonnormal points of a variety form a proper closed subset.
\end{problem}
\end{exerciseintro}

\begin{solution}
On an affine open with ring \(A\), let \(B\) be its normalization.
By localization, \(A_{\mathfrak m}\) is normal exactly when
\(B_{\mathfrak m}=A_{\mathfrak m}\), equivalently
\((B/A)_{\mathfrak m}=0\).
The module \(B/A\) is finite, so the nonnormal set on this open
is the closed set \(V(\operatorname{Ann}_A(B/A))\).
Choose finitely many module generators of \(B\) and clear their
denominators in \(\operatorname{Frac}A\).
Their product yields nonzero \(a\in A\) with \(aB\subseteq A\).
Thus the annihilator is nonzero, and its zero set is proper
(or empty).
These descriptions on affine opens agree because normality is
intrinsic to the local ring. A subset closed on an open cover is
closed globally. The nonempty normal open \(D(a)\) in any one
affine chart also proves that this global closed subset is proper.
\end{solution}

\begin{exerciseintro}
% record: H-I-05-014
\exerciseheading{Exercise I.5.14}
\addcontentsline{toc}{subsection}{Exercise I.5.14}

\begin{context}
Work over an algebraically closed field \(k\).
For a reduced affine hypersurface \(f=0\), singular points are the
common zeros of \(f\) and all its first partial derivatives.
The multiplicity of a plane curve at the origin is the lowest
nonzero total degree of its equation; the corresponding homogeneous
form is its tangent cone. A regular local ring is characterized
by equality of its dimension and the dimension of its cotangent space.

Analytic isomorphism here means isomorphism of completed local
\(k\)-algebras. For a plane equation of order \(m\), its associated
graded local ring is \(k[x,y]/(f_m)\), with Hilbert function eventually
equal to \(m\). A formal substitution with invertible linear part is
an automorphism. Weierstrass preparation replaces a series regular
of order two in \(y\), up to a unit, by \(y^2+a(x)y+b(x)\).
Completions of reduced algebraic plane curve local rings are reduced.
The number of branches and
\(\delta=\dim_k(\widetilde R/R)\), with \(\widetilde R\) the
normalization of a reduced complete curve ring, are ring invariants.
\end{context}

\begin{problem}
(a) Prove that analytically isomorphic plane singularities have
the same multiplicity.

(b) If the leading form of \(f\in k[[x,y]]\) is \(g_sh_t\),
where the homogeneous factors are coprime, lift the factorization
to \(f=gh\) with those leading forms.

(c) Show that all ordinary double points are analytically
isomorphic, as are all ordinary triple points; explain the
one-parameter variation of ordinary four-fold points.

(d) In characteristic different from two, classify a reduced
plane double point as \(y^2=x^r\), \(r\geq2\), and prove uniqueness
of \(r\).
\end{problem}
\end{exerciseintro}

\begin{solution}
(a) An isomorphism of local rings preserves the maximal ideal,
its powers, and their successive quotient dimensions.
The eventual value of this Hilbert function is the order of
the plane equation, so the multiplicities coincide.

(b) Put \(m=s+t\). Fix the leading pieces \(g_s,h_t\), whose
product already equals the degree-\(m\) term of \(f\).
Construct higher homogeneous pieces successively, starting at
\(j=1\). At degree \(m+j\), the known lower pieces leave
an error \(E_{m+j}\), to be written
\[
E_{m+j}=g_s h_{t+j}+h_t g_{s+j}.
\]
The homogeneous map from forms of degrees \(t+j,s+j\) onto
forms of degree \(s+t+j\) is surjective for \(j\geq0\).
Indeed its kernel consists exactly of
\((h_t q,-g_s q)\), with \(q\) of degree \(j\), by coprimality.
Its image dimension is
\((t+j+1)+(s+j+1)-(j+1)=s+t+j+1\), the target dimension.
Thus the correction exists at every step; the series converge
in the formal topology and their product is \(f\).

(c) Applying (b) repeatedly factors an ordinary multiple point
into smooth branches with the prescribed distinct tangent lines.
For two branches, their equations form formal coordinates, giving
the ring \(k[[x,y]]/(xy)\).
For three, straighten the first two to \(x=0,y=0\).
The third has equation \(xA+yB\), where \(A,B\) are units because
its tangent is distinct from both axes.
The change \(x'=xA,y'=yB\) gives the form
\(x'y'(x'+y')\), up to a unit.

For four, already the cones
\[
xy(x-y)(x-\lambda y)=0,\qquad \lambda\notin\{0,1\},
\]
give a one-parameter collection. An analytic isomorphism induces
a linear isomorphism on cotangent spaces and therefore a projective
transformation of their four unordered tangent directions.
Cross-ratios imply that \(\lambda,\mu\) can be equivalent only if
\[
\mu\in\left\{\lambda,1-\lambda,\lambda^{-1},
(1-\lambda)^{-1},\frac{\lambda}{\lambda-1},
\frac{\lambda-1}{\lambda}\right\}.
\]
Conversely these changes arise by permuting four directions.
Thus parameters modulo these finite identifications form a
one-dimensional family of distinct analytic types; choosing one
representative per orbit gives mutually nonisomorphic examples.
This does not assert that every distinct raw value of \(\lambda\)
gives a distinct unordered configuration, or classify the effect of
all higher-order terms of an arbitrary ordinary four-fold point.

(d) Choose a \(y\)-direction where the quadratic leading form
is nonzero. Preparation gives \(y^2+a(x)y+b(x)\), with
\(\operatorname{ord}a\geq1,\operatorname{ord}b\geq2\).
Completing the square gives \(y'^2-c(x)\).
Reducedness implies \(c\ne0\); otherwise the completed equation
would be a square. Write \(c=x^r u(x)\), with \(r\geq2\)
and \(u\) a unit. It has a square root in \(k[[x]]\), obtained
coefficient by coefficient since two is invertible.
Replacing \(y'\) by \(y'/\sqrt{u}\) yields \(y^2=x^r\).

If \(r=2m+1\), the normalization is \(k[[t]]\) with
\(x=t^2,y=t^{2m+1}\). Its quotient by the original ring has
basis \(t,t^3,\ldots,t^{2m-1}\), so \(\delta=m\) and there is
one branch. If \(r=2m\), the two branches are \(y=\pm x^m\).
The ring embeds in \(k[[x]]\oplus k[[x]]\) as pairs congruent
modulo \(x^m\); hence \(\delta=m\) and there are two branches.
In both cases \(r=2\delta-\text{(number of branches)}+2\).
These invariants prove uniqueness. The values \(r=2,3,4\)
give respectively node, cusp, and tacnode.
\end{solution}

\begin{exerciseintro}
% record: H-I-05-015
\exerciseheading{Exercise I.5.15}
\addcontentsline{toc}{subsection}{Exercise I.5.15}

\begin{context}
Work over an algebraically closed field \(k\).
For a reduced affine hypersurface \(f=0\), singular points are the
common zeros of \(f\) and all its first partial derivatives.
The multiplicity of a plane curve at the origin is the lowest
nonzero total degree of its equation; the corresponding homogeneous
form is its tangent cone. A regular local ring is characterized
by equality of its dimension and the dimension of its cotangent space.

Nonzero degree-\(d\) ternary forms up to scalar form
\(\mathbb P^N\), \(N=\binom{d+2}{2}-1\).
Projection of a closed subset of \(\mathbb P^N\times\mathbb P^2\)
to \(\mathbb P^N\) is closed, by projective elimination.
The gradient criterion is valid in every characteristic.
For each \(d>0\) there is a smooth degree-\(d\) plane curve:
use \(x^d+y^d+z^d\) if the characteristic does not divide \(d\),
and \(x^d+y^{d-1}z+xz^{d-1}\) otherwise.
\end{context}

\begin{problem}
(a) Describe the correspondence from \(\mathbb P^N\) to plane
zero sets defined by degree-\(d\) equations, including its failures
of injectivity.

(b) Show that nonsingular irreducible degree-\(d\) curves correspond
bijectively to a nonempty Zariski-open subset of \(\mathbb P^N\).
\end{problem}
\end{exerciseintro}

\begin{solution}
(a) A nonzero coefficient vector defines a homogeneous form and
its projective zero set; scalar multiples define the same set.
Conversely every set admitting a degree-\(d\) equation is obtained.
If \(f=c\prod h_i^{a_i}\) is its factorization into distinct
irreducible homogeneous forms, its zero set remembers exactly
the factors \(h_i\), not their positive multiplicities \(a_i\).
Thus fibers consist precisely of exponent choices satisfying
\(\sum a_i\deg h_i=d\), modulo a common scalar.
For a squarefree \(f\), all these exponents are forced to be one,
so its point is uniquely determined.

(b) In parameter space times \(\mathbb P^2\), impose the four equations
\[
F=F_x=F_y=F_z=0,
\]
where \(F\) is the universal degree-\(d\) form.
Their incidence set is closed and its projection is closed by
elimination. The complement consists exactly of forms whose zero
sets satisfy the nonvanishing-gradient condition.
Such a form is irreducible: a nontrivial factorization would yield
a common projective zero of two factors where all derivatives vanish.
Consequently the complement parametrizes precisely the stated curves,
and (a) gives bijectivity.
The explicit smooth examples in the context prove nonemptiness.
The equation \(F=0\) must be retained in the incidence condition:
when the characteristic divides \(d\), Euler's identity does not
recover it from the partial derivatives.
\end{solution}

\begin{exerciseintro}
\corpussection{I.6 Nonsingular Curves}
\addcontentsline{toc}{section}{I.6 Nonsingular Curves}

% record: H-I-06-001
\exerciseheading{Exercise I.6.1}
\addcontentsline{toc}{subsection}{Exercise I.6.1}

\begin{context}
Work over an algebraically closed field \(k\). A nonsingular curve is
an open subvariety of its unique nonsingular projective model.
Rational maps from a nonsingular curve to a projective variety extend
uniquely to morphisms. Function-field isomorphisms give birational
maps and inverse isomorphisms of nonsingular projective models.
Nonsingular curves are normal. Principal affine opens have localized
coordinate rings, and localizations of unique factorization domains
are unique factorization domains.
\end{context}

\begin{problem}
Let \(Y\) be a nonsingular rational curve not isomorphic to
\(\mathbb P^1\).

(a) Show that \(Y\) is isomorphic to an open subset of \(\mathbb A^1\).

(b) Show that \(Y\) is affine.

(c) Show that \(A(Y)\) is a unique factorization domain.
\end{problem}
\end{exerciseintro}

\begin{solution}
(a) Its nonsingular projective model is \(\mathbb P^1\), since
its function field is \(k(t)\). The open embedding in that model
has nonempty complement. Send one missing point to infinity by a
projective linear change of coordinates. Then \(Y\) lies openly in
\(\mathbb A^1\).

(b) A proper closed subset of \(\mathbb A^1\) is finite, so the
image is \(D(h)\), where \(h=\prod(t-a_i)\) ranges over the finite
missing affine points; take \(h=1\) if none are missing.
The graph \(hu=1\) realizes it as an affine variety.

(c) Its coordinate ring is \(k[t,h^{-1}]\), a localization of
the principal ideal domain \(k[t]\), hence a unique factorization domain.
\end{solution}

\begin{exerciseintro}
% record: H-I-06-002
\exerciseheading{Exercise I.6.2}
\addcontentsline{toc}{subsection}{Exercise I.6.2}

\begin{context}
Work over an algebraically closed field \(k\). A nonsingular curve is
an open subvariety of its unique nonsingular projective model.
Rational maps from a nonsingular curve to a projective variety extend
uniquely to morphisms. Function-field isomorphisms give birational
maps and inverse isomorphisms of nonsingular projective models.
Nonsingular curves are normal. Principal affine opens have localized
coordinate rings, and localizations of unique factorization domains
are unique factorization domains.

Assume \(\operatorname{char}k\ne2\).
Gauss's lemma tests polynomial irreducibility over a fraction field.
An affine variety is normal exactly when its coordinate ring is
integrally closed. An affine nonsingular rational curve has a
unique factorization coordinate ring, since it is a principal
open subset of \(\mathbb A^1\).
\end{context}

\begin{problem}
For \(Y=V(y^2-x^3+x)\), put \(A=k[x,y]/(y^2-x^3+x)\), \(K=K(Y)\).

(a) Prove that \(Y\) is nonsingular and \(A\) is a normal domain.

(b) Show that \(k[x]\) is polynomial and that \(A\) is its integral
closure in \(K\).

(c) Define \(\sigma(x)=x,\sigma(y)=-y\); prove that
\(N(a)=a\sigma(a)\) lies in \(k[x]\), with \(N(1)=1\) and
\(N(ab)=N(a)N(b)\).

(d) Determine the units, prove \(x,y\) irreducible, and show \(A\)
is not a unique factorization domain.

(e) Deduce that \(Y\) is not rational.
\end{problem}
\end{exerciseintro}

\begin{solution}
(a) The polynomial \(x^3-x\) is not a square in \(k(x)\), because
its order at \(x=0\) is one. Thus the quadratic in \(y\) is irreducible.
A singular point would have \(y=0,3x^2=1\), whereas the equation
then requires \(x=0,\pm1\); these cannot occur together when
two is nonzero. This includes characteristic three.
The curve is nonsingular, hence normal, so \(A\) is integrally closed.

(b) Division by the monic relation gives the free \(k[x]\)-basis
\(1,y\). In particular \(k[x]\) injects and has no polynomial relation.
The element \(y\) is integral over it, so \(A\) is integral over \(k[x]\).
An element of \(K\) integral over \(k[x]\) is also integral over \(A\);
normality puts it in \(A\). This proves the integral-closure statement.

(c) The substitution preserves the equation and squares to the identity.
For the unique expression \(a=f(x)+yg(x)\),
\[
N(a)=f(x)^2-(x^3-x)g(x)^2.
\]
It lies in \(k[x]\), and the two norm identities follow from
commutativity and the automorphism property.

(d) If both summands are nonzero, their degrees are respectively
even and odd, so their leading terms cannot cancel. Thus for nonzero \(a\),
\[
\deg N(a)=\max\{2\deg f,\ 3+2\deg g\},
\]
omitting zero summands. A unit has norm a unit in \(k[x]\), hence
degree zero, so is a nonzero constant. Conversely those are units.
Every nonunit has norm degree at least two.
Since \(N(x)=x^2\) and \(N(y)=-(x^3-x)\) have degrees two and three,
neither can be a product of norms of two nonunits.
Thus \(x,y\) are irreducible.
But \(x\mid y^2=x(x-1)(x+1)\) and \(x\nmid y\): the basis
\(1,y\) would require a polynomial coefficient \(1/x\) if \(y=xa\).
Thus the irreducible \(x\) is not prime, excluding unique factorization.

(e) The affine curve is not \(\mathbb P^1\), since it has
nonconstant global regular functions. If it were rational,
its nonsingularity and the open-subset result would make \(A\)
a unique factorization domain, contradicting (d).
\end{solution}

\begin{exerciseintro}
% record: H-I-06-003
\exerciseheading{Exercise I.6.3}
\addcontentsline{toc}{subsection}{Exercise I.6.3}

\begin{context}
Work over an algebraically closed field \(k\). A nonsingular curve is
an open subvariety of its unique nonsingular projective model.
Rational maps from a nonsingular curve to a projective variety extend
uniquely to morphisms. Function-field isomorphisms give birational
maps and inverse isomorphisms of nonsingular projective models.
Nonsingular curves are normal. Principal affine opens have localized
coordinate rings, and localizations of unique factorization domains
are unique factorization domains.
\end{context}

\begin{problem}
Give counterexamples to extension of rational maps if
(a) the nonsingular source has dimension at least two, or
(b) the target need not be projective.
\end{problem}
\end{exerciseintro}

\begin{solution}
(a) The rational map \(\mathbb A^2\dashrightarrow\mathbb P^1\),
\((x,y)\mapsto[x:y]\), has no extension at the origin.
Either affine target chart at a proposed value would make \(y/x\)
or \(x/y\) regular in \(k[x,y]_{(x,y)}\), impossible by divisibility.
The source is nonsingular of dimension two.

(b) The coordinate \(t\) gives a rational map
\(\mathbb P^1\dashrightarrow\mathbb A^1\).
At infinity it is \(1/u\) in the parameter \(u=1/t\), absent
from \(k[u]_{(u)}\). Thus it cannot extend there, although the
source is a nonsingular projective curve.
\end{solution}

\begin{exerciseintro}
% record: H-I-06-004
\exerciseheading{Exercise I.6.4}
\addcontentsline{toc}{subsection}{Exercise I.6.4}

\begin{context}
Work over an algebraically closed field \(k\). A nonsingular curve is
an open subvariety of its unique nonsingular projective model.
Rational maps from a nonsingular curve to a projective variety extend
uniquely to morphisms. Function-field isomorphisms give birational
maps and inverse isomorphisms of nonsingular projective models.
Nonsingular curves are normal. Principal affine opens have localized
coordinate rings, and localizations of unique factorization domains
are unique factorization domains.

A morphism from a projective variety has closed image
(projective elimination applied to its closed graph).
Proper closed subsets of an irreducible curve are finite.
\end{context}

\begin{problem}
Show that a nonconstant rational function on a nonsingular
projective curve defines a surjective morphism to \(\mathbb P^1\)
with finite fibers.
\end{problem}
\end{exerciseintro}

\begin{solution}
Regard the function as a rational map to \(\mathbb P^1\).
The extension theorem supplies its morphism on the entire curve.
Its image is irreducible and closed. A proper such subset of
\(\mathbb P^1\) is a single point, which would make the function
constant. Thus the image is all of \(\mathbb P^1\).
Each fiber is closed and proper because the morphism is nonconstant,
so is finite. Surjectivity also says it is nonempty.
\end{solution}

\begin{exerciseintro}
% record: H-I-06-005
\exerciseheading{Exercise I.6.5}
\addcontentsline{toc}{subsection}{Exercise I.6.5}

\begin{context}
Work over an algebraically closed field \(k\). A nonsingular curve is
an open subvariety of its unique nonsingular projective model.
Rational maps from a nonsingular curve to a projective variety extend
uniquely to morphisms. Function-field isomorphisms give birational
maps and inverse isomorphisms of nonsingular projective models.
Nonsingular curves are normal. Principal affine opens have localized
coordinate rings, and localizations of unique factorization domains
are unique factorization domains.

Varieties have closed diagonals. Projection
\(\mathbb P^n\times T\to T\) is closed for a variety \(T\),
by projective elimination on affine charts of \(T\).
\end{context}

\begin{problem}
If a nonsingular projective curve \(X\) is a locally closed
subvariety of a variety \(Y\), prove that \(X\) is closed in \(Y\).
\end{problem}
\end{exerciseintro}

\begin{solution}
The inclusion \(i:X\to Y\) is a morphism, so its graph is closed
in \(X\times Y\): it is the inverse image of the diagonal of \(Y\)
under \((x,y)\mapsto(i(x),y)\).
Embed the projective \(X\) as a closed subset of \(\mathbb P^n\).
The graph is then closed in \(\mathbb P^n\times Y\).
Its projection to \(Y\) is closed by elimination, and is exactly
the given subset \(X\). This proof in fact uses only projectivity,
not nonsingularity or the dimension-one hypothesis.
\end{solution}

\begin{exerciseintro}
% record: H-I-06-006
\exerciseheading{Exercise I.6.6}
\addcontentsline{toc}{subsection}{Exercise I.6.6}

\begin{context}
Work over an algebraically closed field \(k\). A nonsingular curve is
an open subvariety of its unique nonsingular projective model.
Rational maps from a nonsingular curve to a projective variety extend
uniquely to morphisms. Function-field isomorphisms give birational
maps and inverse isomorphisms of nonsingular projective models.
Nonsingular curves are normal. Principal affine opens have localized
coordinate rings, and localizations of unique factorization domains
are unique factorization domains.

An invertible two-by-two matrix acts on homogeneous coordinates
of \(\mathbb P^1\); scalar matrices act trivially.
For relatively prime nonproportional \(P,Q\in k[T]\),
\(P(T)-uQ(T)\) is irreducible in \(k(u)[T]\):
apply Gauss's lemma to this primitive polynomial linear in \(u\).
\end{context}

\begin{problem}
(a) Show that \(t\mapsto(at+b)/(ct+d)\), \(ad-bc\ne0\), gives
an automorphism of \(\mathbb P^1\).

(b) Identify its automorphism group with the group of
\(k\)-automorphisms of \(k(t)\).

(c) Show that all these automorphisms are fractional linear,
so the projective linear group is the full automorphism group.
\end{problem}
\end{exerciseintro}

\begin{solution}
(a) Homogenizing gives
\([T:U]\mapsto[aT+bU:cT+dU]\).
The determinant condition prevents simultaneous zero coordinates.
The inverse matrix supplies an inverse morphism.

(b) Pullback gives the contravariant correspondence. Every field
automorphism and its inverse extend to inverse morphisms of the
nonsingular projective model \(\mathbb P^1\).
To state a group isomorphism, send a geometric automorphism
\(\varphi\) to \((\varphi^{-1})^*\); ordinary pullback alone
reverses composition.

(c) A field automorphism sends \(t\) to a nonconstant \(P(t)/Q(t)\)
with coprime polynomials \(P,Q\).
Put \(d=\max(\deg P,\deg Q)\).
The polynomial \(P(T)-uQ(T)\) has degree \(d\) in \(T\) and is
irreducible over \(k(u)\). Substituting the transcendental
\(u=P(t)/Q(t)\) proves
\([k(t):k(P(t)/Q(t))]=d\).
For an automorphism this degree is one.
Thus \(P,Q\) are linear or constant and give
\((at+b)/(ct+d)\), with nonzero determinant because the ratio is
nonconstant. Two matrices give the same function exactly when
they differ by scalar, proving the asserted projective linear description.
\end{solution}

\begin{exerciseintro}
% record: H-I-06-007
\exerciseheading{Exercise I.6.7}
\addcontentsline{toc}{subsection}{Exercise I.6.7}

\begin{context}
Work over an algebraically closed field \(k\). A nonsingular curve is
an open subvariety of its unique nonsingular projective model.
Rational maps from a nonsingular curve to a projective variety extend
uniquely to morphisms. Function-field isomorphisms give birational
maps and inverse isomorphisms of nonsingular projective models.
Nonsingular curves are normal. Principal affine opens have localized
coordinate rings, and localizations of unique factorization domains
are unique factorization domains.

Every automorphism of \(\mathbb P^1\) is fractional linear.
Four unordered distinct points have cross-ratio defined up to
the six values \(\lambda,1-\lambda,\lambda^{-1},(1-\lambda)^{-1},
\lambda/(\lambda-1),(\lambda-1)/\lambda\).
Fractional linear maps act transitively on ordered triples of distinct points.
\end{context}

\begin{problem}
If \(\mathbb A^1\setminus\{P_1,\ldots,P_r\}\) and
\(\mathbb A^1\setminus\{Q_1,\ldots,Q_s\}\) are isomorphic,
prove \(r=s\). Is equality sufficient?
\end{problem}
\end{exerciseintro}

\begin{solution}
An isomorphism induces a function-field automorphism and hence
extends uniquely to an automorphism of their projective completions
\(\mathbb P^1\). Its inverse extends too, so it bijects the missing
sets, including the point at infinity on each side.
Therefore \(r+1=s+1\), giving \(r=s\).

Equality is not sufficient in general. For \(r=s=3\), take missing
affine sets \(\{0,1,\lambda\}\), \(\{0,1,\mu\}\).
Choose \(\lambda\notin\{0,1\}\) and choose \(\mu\) outside
\(\{0,1\}\) and its finite cross-ratio orbit, possible since \(k\)
is infinite. The four missing projective points cannot be carried
to one another by an automorphism, so the punctured curves are
not isomorphic. For \(r=s\leq2\), equality does suffice, by the
transitivity on at most three missing projective points.
\end{solution}

\begin{exerciseintro}
\corpussection{I.7 Intersections in Projective Space}
\addcontentsline{toc}{section}{I.7 Intersections in Projective Space}

% record: H-I-07-001
\exerciseheading{Exercise I.7.1}
\addcontentsline{toc}{subsection}{Exercise I.7.1}

\begin{context}
Work over an algebraically closed field \(k\). The Hilbert polynomial
of an \(r\)-dimensional projective variety has leading term
\((\deg Y)m^r/r!\).
Hilbert polynomials are additive in exact sequences of graded modules.
For a reduced pure-dimensional algebraic set, degree is the sum of
the degrees of its irreducible components.
A proper hyperplane section has degree equal to that of the variety
when component multiplicities are included; discarding multiplicities
can only decrease degree.
\end{context}

\begin{problem}
(a) Compute the degree of the degree-\(d\) Veronese image of
\(\mathbb P^n\).

(b) Compute the degree of the Segre image of
\(\mathbb P^r\times\mathbb P^s\).
\end{problem}
\end{exerciseintro}

\begin{solution}
(a) Degree \(m\) in the Veronese coordinate ring consists of
all forms of degree \(dm\) in \(n+1\) variables, because every
such monomial factors into \(m\) degree-\(d\) monomials.
Thus its Hilbert function is \(\binom{dm+n}{n}\), whose
leading coefficient is \(d^n/n!\). The degree is \(d^n\).

(b) Degree \(m\) in the Segre ring consists of all bihomogeneous
forms of bidegree \((m,m)\), since each monomial pairs its
\(m\) first-factor variables with its \(m\) second-factor variables.
Its Hilbert function is
\(\binom{m+r}{r}\binom{m+s}{s}\).
The leading coefficient is \(1/(r!s!)\) and the dimension is \(r+s\),
so the degree is \((r+s)!/(r!s!)=\binom{r+s}{r}\).
\end{solution}

\begin{exerciseintro}
% record: H-I-07-002
\exerciseheading{Exercise I.7.2}
\addcontentsline{toc}{subsection}{Exercise I.7.2}

\begin{context}
Work over an algebraically closed field \(k\). The Hilbert polynomial
of an \(r\)-dimensional projective variety has leading term
\((\deg Y)m^r/r!\).
Hilbert polynomials are additive in exact sequences of graded modules.
For a reduced pure-dimensional algebraic set, degree is the sum of
the degrees of its irreducible components.
A proper hyperplane section has degree equal to that of the variety
when component multiplicities are included; discarding multiplicities
can only decrease degree.

Define \(p_a(Y)=(-1)^{\dim Y}(P_Y(0)-1)\).
The polynomial \(\binom{m+n}{n}\) is the Hilbert polynomial of
\(\mathbb P^n\); binomial expressions below are evaluated as
polynomials even at negative integers.
Two homogeneous equations generating the ideal of a codimension-two
complete intersection form a regular sequence.
\end{context}

\begin{problem}
Compute arithmetic genera:

(a) \(\mathbb P^n\).

(b) A plane curve of degree \(d\).

(c) A degree-\(d\) hypersurface in \(\mathbb P^n\).

(d) A strict complete-intersection curve of surfaces of degrees
\(a,b\) in \(\mathbb P^3\).

(e) The Segre product of projective varieties \(Y,Z\) of dimensions
\(r,s\), in terms of their arithmetic genera.
\end{problem}
\end{exerciseintro}

\begin{solution}
(a) Since \(\binom{n}{n}=1\), the answer is zero.

(b) Multiplication by the nonzero defining equation gives
\(P_Y(m)=\binom{m+2}{2}-\binom{m+2-d}{2}\).
Thus \(p_a(Y)=1-P_Y(0)=(d-1)(d-2)/2\).

(c) The same exact sequence gives
\[
P_H(m)=\binom{m+n}{n}-\binom{m+n-d}{n}.
\]
Consequently
\(p_a(H)=(-1)^n\binom{n-d}{n}=\binom{d-1}{n}\).

(d) The two regular-sequence exact sequences, equivalently their
Koszul resolution, give
\[
P_Y(m)=\binom{m+3}{3}-\binom{m+3-a}{3}
-\binom{m+3-b}{3}+\binom{m+3-a-b}{3}.
\]
Expanding at zero and using \(p_a=1-P_Y(0)\) yields
\[
p_a(Y)=1+\frac{ab(a+b-4)}2.
\]
Ideal generation, not merely set-theoretic intersection, is the
hypothesis needed for this Hilbert polynomial.

(e) The degree-\(m\) component of the Segre coordinate ring is
the tensor product of the two degree-\(m\) components.
Therefore \(P_{Y\times Z}(m)=P_Y(m)P_Z(m)\).
Substituting \(P_Y(0)=1+(-1)^r p_a(Y)\) and the analogous
formula for \(Z\) gives
\[
p_a(Y\times Z)=p_a(Y)p_a(Z)+(-1)^s p_a(Y)+(-1)^r p_a(Z).
\]
\end{solution}

\begin{exerciseintro}
% record: H-I-07-003
\exerciseheading{Exercise I.7.3}
\addcontentsline{toc}{subsection}{Exercise I.7.3}

\begin{context}
Work over an algebraically closed field \(k\). The Hilbert polynomial
of an \(r\)-dimensional projective variety has leading term
\((\deg Y)m^r/r!\).
Hilbert polynomials are additive in exact sequences of graded modules.
For a reduced pure-dimensional algebraic set, degree is the sum of
the degrees of its irreducible components.
A proper hyperplane section has degree equal to that of the variety
when component multiplicities are included; discarding multiplicities
can only decrease degree.

If a plane curve is smooth at the origin, its local equation starts
with a nonzero linear form. Intersection multiplicity with a line
not a component is the order of its restricted equation.
Homogeneous coordinates of a morphism may be given by forms of
equal degree having no common zero on its domain.
\end{context}

\begin{problem}
For a plane curve \(Y\), show that each nonsingular point has a unique
tangent line with intersection multiplicity greater than one.
Show that sending the point to this line is a morphism to the dual
projective plane. Address the case when \(Y\) itself is a line.
\end{problem}
\end{exerciseintro}

\begin{solution}
Translate a smooth point to the affine origin and write \(f=\ell+h\),
where \(\ell\ne0\) is linear and \(h\) has order at least two.
On a line with direction \((a,b)\), the restricted equation has
linear coefficient \(\ell(a,b)\). Its intersection multiplicity
is one unless \(\ell(a,b)=0\), which selects exactly one line;
on that line the order is greater than one.
For a curve of degree greater than one the restricted equation
cannot vanish identically, because that would make the line a component.
If \(Y\) is a line, its tangent is itself: its intersection quotient
has infinite length, while every other line through the point has
multiplicity one.

For homogeneous equation \(F\), the tangent line at a smooth
point \(P\) has coefficients
\[
[F_{x_0}(P):F_{x_1}(P):F_{x_2}(P)].
\]
The coefficients have equal degree and are not all zero on
\(\operatorname{Reg}Y\); Euler's identity shows the line passes
through \(P\), and dehomogenization identifies it with the line
just found. Ratios on the three target charts are regular,
so this is a morphism. Its image closure is the dual locus;
for a line it is a point rather than a one-dimensional curve.
\end{solution}

\begin{exerciseintro}
% record: H-I-07-004
\exerciseheading{Exercise I.7.4}
\addcontentsline{toc}{subsection}{Exercise I.7.4}

\begin{context}
Work over an algebraically closed field \(k\). The Hilbert polynomial
of an \(r\)-dimensional projective variety has leading term
\((\deg Y)m^r/r!\).
Hilbert polynomials are additive in exact sequences of graded modules.
For a reduced pure-dimensional algebraic set, degree is the sum of
the degrees of its irreducible components.
A proper hyperplane section has degree equal to that of the variety
when component multiplicities are included; discarding multiplicities
can only decrease degree.

The singular set of an irreducible plane curve is finite.
The closure of an image of a curve under a morphism has dimension
at most one. The tangent assignment on the smooth locus is a morphism.
A line not equal to a plane curve of degree \(d\) has total
intersection multiplicity \(d\); transverse smooth intersections
have multiplicity one.
\end{context}

\begin{problem}
Show that a nonempty open set of lines in the dual plane meets
a degree-\(d\) plane curve in exactly \(d\) distinct points.
\end{problem}
\end{exerciseintro}

\begin{solution}
For each singular point, lines through it form a projective line
in the dual plane. The closure of all smooth-point tangent lines
has dimension at most one.
Their union is a finite union of proper closed subsets of the
irreducible dual plane, hence is proper closed.
Let \(U\) be its nonempty open complement.
A line in \(U\) passes through no singular point and is nowhere
tangent to the curve. If the curve itself is a line, that one
dual point was also excluded. All intersections therefore have
multiplicity one. Their total multiplicity is \(d\), so there
are exactly \(d\) distinct points.
\end{solution}

\begin{exerciseintro}
% record: H-I-07-005
\exerciseheading{Exercise I.7.5}
\addcontentsline{toc}{subsection}{Exercise I.7.5}

\begin{context}
Work over an algebraically closed field \(k\). The Hilbert polynomial
of an \(r\)-dimensional projective variety has leading term
\((\deg Y)m^r/r!\).
Hilbert polynomials are additive in exact sequences of graded modules.
For a reduced pure-dimensional algebraic set, degree is the sum of
the degrees of its irreducible components.
A proper hyperplane section has degree equal to that of the variety
when component multiplicities are included; discarding multiplicities
can only decrease degree.

The multiplicity at \([0:0:1]\) is the lowest degree of a nonzero
term of \(F(x,y,1)\). A binary homogeneous form factors into
linear factors over \(k\). Birational inverse maps on nonempty
opens prove rationality.
\end{context}

\begin{problem}
(a) Prove that an irreducible plane curve of degree \(d>1\)
has no point of multiplicity at least \(d\).

(b) Prove that such a curve with a point of multiplicity \(d-1\)
is rational.
\end{problem}
\end{exerciseintro}

\begin{solution}
(a) Put the point at \([0:0:1]\). Multiplicity at least \(d\)
forces the degree-\(d\) equation to be \(F=f_d(x,y)\).
This binary form splits into \(d\) linear factors, contrary to
irreducibility and \(d>1\).

(b) The equation now has the form
\[
F=z f_{d-1}(x,y)+f_d(x,y),\qquad f_{d-1}\ne0.
\]
Projection from the point sends \([x:y:z]\) to \([x:y]\).
On the nonempty open \(f_{d-1}(u,v)\ne0\), its inverse is
\[
[u:v]\longmapsto
[u f_{d-1}(u,v):v f_{d-1}(u,v):-f_d(u,v)].
\]
These degree-\(d\) coordinates satisfy \(F=0\), and direct
substitution shows both compositions are the identity where
defined. Thus the curve is birational to \(\mathbb P^1\).
\end{solution}

\begin{exerciseintro}
% record: H-I-07-006
\exerciseheading{Exercise I.7.6}
\addcontentsline{toc}{subsection}{Exercise I.7.6}

\begin{context}
Work over an algebraically closed field \(k\). The Hilbert polynomial
of an \(r\)-dimensional projective variety has leading term
\((\deg Y)m^r/r!\).
Hilbert polynomials are additive in exact sequences of graded modules.
For a reduced pure-dimensional algebraic set, degree is the sum of
the degrees of its irreducible components.
A proper hyperplane section has degree equal to that of the variety
when component multiplicities are included; discarding multiplicities
can only decrease degree.

For a standard graded domain \(S\) over an infinite field,
graded Noether normalization provides algebraically independent
linear forms \(l_0,\ldots,l_r\) such that \(S\) is finite over
\(B=k[l_0,\ldots,l_r]\).
A finite graded module of rank \(e\) over this polynomial ring
has Hilbert polynomial with leading coefficient \(e/r!\):
choose homogeneous generically independent generators; the
remaining torsion modules have smaller dimension.
Polynomial rings are integrally closed.
\end{context}

\begin{problem}
Prove that a reduced projective algebraic set of pure dimension
\(r\) has degree one exactly when it is a linear variety.
\end{problem}
\end{exerciseintro}

\begin{solution}
For a linear \(\mathbb P^r\), the coordinate ring is a polynomial
ring in \(r+1\) linear variables, of degree one.
Conversely, additivity of degree and positivity of each component
degree show that a pure-dimensional degree-one set has exactly
one irreducible component. Its homogeneous coordinate ring \(S\)
is a domain. Choose the normalization algebra \(B\) in the context.
The Hilbert leading coefficient shows that
\([\operatorname{Frac}S:\operatorname{Frac}B]=\deg Y=1\).
Thus \(B\subseteq S\subseteq\operatorname{Frac}B\) and \(S\)
is integral over \(B\). Integral closedness gives \(S=B\).
In particular all ambient degree-one coordinates restrict to linear
combinations of \(l_0,\ldots,l_r\). Their linear relations define
a linear \(\mathbb P^r\) containing \(Y\); the coordinate-ring
equality says that \(Y\) is exactly that linear space.
\end{solution}

\begin{exerciseintro}
% record: H-I-07-007
\exerciseheading{Exercise I.7.7}
\addcontentsline{toc}{subsection}{Exercise I.7.7}

\begin{context}
Work over an algebraically closed field \(k\). The Hilbert polynomial
of an \(r\)-dimensional projective variety has leading term
\((\deg Y)m^r/r!\).
Hilbert polynomials are additive in exact sequences of graded modules.
For a reduced pure-dimensional algebraic set, degree is the sum of
the degrees of its irreducible components.
A proper hyperplane section has degree equal to that of the variety
when component multiplicities are included; discarding multiplicities
can only decrease degree.

A general codimension-\(r\) linear section of an \(r\)-dimensional
projective variety consists of \(\deg Y\) distinct smooth points.
More generally, a finite intersection with a codimension-\(r\)
linear subspace has at most \(\deg Y\) distinct points: cut
successively with a flag of hyperplanes, preserving proper
intersection at each step, and use the degree bound for each cut.
An irreducible cone with point vertex over \(Z\) has dimension
\(\dim Z+1\) and degree \(\deg Z\), since its coordinate ring is
the polynomial extension of the coordinate ring of \(Z\).
The projective tangent space at a smooth point of an \(r\)-fold
is a linear \(\mathbb P^r\). Degree-one varieties are linear.
Blowing up a smooth point has exceptional fiber \(\mathbb P^{r-1}\)
and resolves projection from that point. A projective dominant morphism
of equal-dimensional varieties has finite fibers over a nonempty open
subset of the target, by the dimension-of-fibers theorem.
\end{context}

\begin{problem}
Let \(Y\subseteq\mathbb P^n\) have dimension \(r\) and degree \(d>1\),
and let \(P\) be a smooth point. Let \(X\) be the closure of the
union of lines \(PQ\), \(Q\in Y\setminus\{P\}\).

(a) Prove that \(X\) is irreducible of dimension \(r+1\).

(b) Prove that \(\deg X<d\).
\end{problem}
\end{exerciseintro}

\begin{solution}
(a) Let \(Z\) be the closure of the projection of \(Y\setminus\{P\}\)
from \(P\) to a complementary hyperplane.
It is irreducible, and \(X\) is exactly its cone with vertex \(P\).
The union-of-lines parametrization bounds \(\dim X\leq r+1\),
while \(Y\subseteq X\) gives \(\dim X\geq r\).
If \(\dim X=r\), irreducibility makes \(X=Y\).
Then every line \(PQ\) is contained in \(Y\), so its direction
at \(P\) lies in \(T_PY\). Consequently \(Y\) lies in the
projective tangent \(\mathbb P^r\), and having the same dimension
equals it. This gives degree one, a contradiction.
Hence \(\dim X=r+1\), and \(\dim Z=r\).

(b) The projection has zero-dimensional generic fiber, since
\(\dim Y=\dim Z\).
More explicitly, the closed incidence graph in the blow-up at \(P\)
is projective over \(Z\); over a nonempty open of \(Z\) its fibers
are finite. Its exceptional image lies among the tangent directions
at \(P\), of dimension at most \(r-1\).
Choose a general codimension-\(r\) linear subspace \(\Lambda\)
of the target, so that \(Z\cap\Lambda\) consists of \(\deg Z\)
distinct points, all in this finite-fiber open and outside the
exceptional image. Their inverse images in \(Y\setminus\{P\}\)
are finite and nonempty.
The span \(L=\langle P,\Lambda\rangle\) has codimension \(r\)
in \(\mathbb P^n\), and \(L\cap Y\) is these finite fibers
together with \(P\). Therefore it has at least \(1+\deg Z\)
distinct points. The linear-section bound gives
\[
1+\deg Z\leq d.
\]
Since \(\deg X=\deg Z\) for a cone, the required strict inequality follows.
\end{solution}

\begin{exerciseintro}
% record: H-I-07-008
\exerciseheading{Exercise I.7.8}
\addcontentsline{toc}{subsection}{Exercise I.7.8}

\begin{context}
Work over an algebraically closed field \(k\). The Hilbert polynomial
of an \(r\)-dimensional projective variety has leading term
\((\deg Y)m^r/r!\).
Hilbert polynomials are additive in exact sequences of graded modules.
For a reduced pure-dimensional algebraic set, degree is the sum of
the degrees of its irreducible components.
A proper hyperplane section has degree equal to that of the variety
when component multiplicities are included; discarding multiplicities
can only decrease degree.

A variety has a nonempty smooth locus.
For a projective \(r\)-fold of degree greater than one, the cone
swept out by lines through a smooth point has dimension \(r+1\)
and strictly smaller degree. A degree-one projective variety
is linear. A codimension-one projective variety in a linear
space has a principal homogeneous ideal, whose generator's
degree equals the variety's degree.
\end{context}

\begin{problem}
Prove that a degree-two \(r\)-dimensional projective variety
lies in a linear \(\mathbb P^{r+1}\), where it is a quadric hypersurface.
\end{problem}
\end{exerciseintro}

\begin{solution}
Choose a smooth point and form the cone of joining lines.
It contains \(Y\), has dimension \(r+1\), and has positive
degree strictly less than two. Its degree is therefore one,
so it is a linear space \(L\cong\mathbb P^{r+1}\).
The variety \(Y\) is a hypersurface in \(L\).
The generator of its homogeneous prime ideal has degree
\(\deg Y=2\), giving the asserted quadric equation.
There is no missing zero-dimensional case: over algebraically
closed \(k\) an irreducible zero-dimensional variety is one
point, of degree one.
\end{solution}

\corpuschapter{Chapter II}
\addcontentsline{toc}{chapter}{Chapter II}

\begin{exerciseintro}
\corpussection{II.1 Sheaves}
\addcontentsline{toc}{section}{II.1 Sheaves}

% record: H-II-01-001
\exerciseheading{Exercise II.1.1}
\addcontentsline{toc}{subsection}{Exercise II.1.1}

\begin{context}
Let \(X\) be a topological space and \(A\) an abelian group. The constant presheaf assigns \(A\) to every nonempty open set, assigns \(0\) to the empty set, and uses the identity maps between nonempty opens and the unique restriction maps to the empty set. The constant sheaf \(\mathcal A\) is the sheaf of locally constant \(A\)-valued functions. One may use the universal property of sheafification, the fact that sheafification preserves stalks, and the criterion that a morphism of sheaves is an isomorphism exactly when it is an isomorphism on every stalk.
\end{context}

\begin{problem}
Prove that the sheaf associated to the constant presheaf with value \(A\) is the constant sheaf \(\mathcal A\).
\end{problem}
\end{exerciseintro}

\begin{solution}
There is a natural morphism from the constant presheaf to \(\mathcal A\): an element \(a\in A\) is sent on each nonempty open set to the constant function with value \(a\). At a point \(x\in X\), every germ of a locally constant function is represented on some neighborhood of \(x\) by a constant function. Consequently the induced map on the stalk at \(x\) is an isomorphism
\[
A\longrightarrow \mathcal A_x.
\]

Let \(A^+\) be the sheafification of the constant presheaf. By the universal property, the natural morphism to \(\mathcal A\) factors uniquely through a morphism
\[
A^+\longrightarrow\mathcal A.
\]
Sheafification does not change stalks, so this morphism is an isomorphism on every stalk. The stalk criterion now shows that \(A^+\to\mathcal A\) is an isomorphism of sheaves.
\end{solution}

\begin{exerciseintro}
% record: H-II-01-002
\exerciseheading{Exercise II.1.2}
\addcontentsline{toc}{subsection}{Exercise II.1.2}

\begin{context}
Sheaves are sheaves of abelian groups on a fixed topological space \(X\). The kernel of a sheaf morphism is computed sectionwise. Its image is the sheafification of the sectionwise image presheaf. Stalks are filtered colimits over neighborhoods, and filtered colimits of abelian groups preserve kernels, images, and exact sequences.
\end{context}

\begin{problem}
Let \(\phi:\mathcal F\to\mathcal G\) be a morphism of sheaves.

(a) Prove that \((\ker\phi)_x=\ker(\phi_x)\) and \((\operatorname{im}\phi)_x=\operatorname{im}(\phi_x)\) for every \(x\in X\).

(b) Prove that \(\phi\) is injective, respectively surjective, exactly when every map on stalks \(\phi_x\) has that property.

(c) Prove that a sequence of sheaves is exact exactly when its sequence of stalks is exact at every point.
\end{problem}
\end{exerciseintro}

\begin{solution}
(a) Taking the stalk of the sectionwise kernel gives
\[
(\ker\phi)_x
=\varinjlim_{x\in U}\ker\bigl(\mathcal F(U)\to\mathcal G(U)\bigr)
\cong\ker\bigl(\mathcal F_x\to\mathcal G_x\bigr),
\]
because filtered colimits are exact. For the image, first take the sectionwise image presheaf and then sheafify. Sheafification preserves stalks, and filtered colimits commute with images, so
\[
(\operatorname{im}\phi)_x\cong\operatorname{im}(\phi_x).
\]

(b) The morphism \(\phi\) is injective exactly when \(\ker\phi=0\). A sheaf is zero exactly when all of its stalks are zero, and part (a) identifies these stalks with \(\ker(\phi_x)\). This proves the injective assertion. Similarly, \(\phi\) is surjective exactly when its cokernel is zero. Applying the same argument to the cokernel, or equivalently using the image calculation in part (a), proves the surjective assertion.

(c) For consecutive maps \(\mathcal F^{i-1}\to\mathcal F^i\to\mathcal F^{i+1}\), exactness at \(\mathcal F^i\) means
\[
\operatorname{im}(\mathcal F^{i-1}\to\mathcal F^i)
=\ker(\mathcal F^i\to\mathcal F^{i+1}).
\]
Two subsheaves are equal exactly when their stalks agree. Part (a) identifies the two stalks in this equality with the image and kernel of the corresponding stalk maps. Thus exactness can be checked point by point.
\end{solution}

\begin{exerciseintro}
% record: H-II-01-003
\exerciseheading{Exercise II.1.3}
\addcontentsline{toc}{subsection}{Exercise II.1.3}

\begin{context}
For sheaves of abelian groups, a morphism is surjective if and only if it is surjective on every stalk. Equality of two germs means that their representatives agree after restriction to a sufficiently small neighborhood.

A nowhere-zero holomorphic function admits local logarithms, and a holomorphic differential has local primitives on small disks.
\end{context}

\begin{problem}
(a) Prove that a morphism \(\phi:\mathcal F\to\mathcal G\) is surjective if and only if every section \(s\in\mathcal G(U)\) can be lifted locally: there is an open cover \(U=\bigcup_i U_i\) and sections \(t_i\in\mathcal F(U_i)\) satisfying \(\phi(t_i)=s|_{U_i}\).

(b) Give a surjective morphism of sheaves for which the map on sections over some open set is not surjective.
\end{problem}
\end{exerciseintro}

\begin{solution}
(a) Suppose \(\phi\) is surjective. For each \(x\in U\), surjectivity of \(\phi_x\) gives a germ \(\tau_x\in\mathcal F_x\) mapping to \(s_x\). Choose an open neighborhood \(V_x\subseteq U\) of \(x\) and a section \(t_x\in\mathcal F(V_x)\) representing \(\tau_x\). The germs of \(\phi(t_x)\) and \(s|_{V_x}\) agree at \(x\), so after shrinking to a neighborhood \(U_x\subseteq V_x\), the two sections agree on \(U_x\). The sets \(U_x\) cover \(U\), and the restrictions \(t_x|_{U_x}\) are the required local lifts.

Conversely, apply the local-lifting property to a representative of an arbitrary germ \(s_x\in\mathcal G_x\). Some member of the resulting cover contains \(x\), and the lift on that member produces a preimage of \(s_x\) in \(\mathcal F_x\). Thus every \(\phi_x\) is surjective, hence so is \(\phi\).

(b) Let \(X=\mathbb C^{\times}\) with its usual complex topology. On \(X\), consider the exponential morphism from the additive sheaf of holomorphic functions to the multiplicative sheaf of nowhere-vanishing holomorphic functions,
\[
\exp:\mathcal O_X\longrightarrow\mathcal O_X^{\times}.
\]
Every nowhere-vanishing holomorphic function has a holomorphic logarithm locally, so this is a surjective morphism of sheaves. The global section \(z\mapsto z\) has no global holomorphic logarithm on \(\mathbb C^{\times}\): integrating its logarithmic derivative around the unit circle gives
\[
\frac{1}{2\pi i}\int_{|z|=1}\frac{dz}{z}=1,
\]
whereas the logarithmic derivative of an exponential \(e^h\) is \(dh\) and has integral zero around every closed curve. Hence the map on global sections is not surjective.
\end{solution}

\begin{exerciseintro}
% record: H-II-01-004
\exerciseheading{Exercise II.1.4}
\addcontentsline{toc}{subsection}{Exercise II.1.4}

\begin{context}
For a presheaf \(\mathcal P\), write \(\mathcal P^+\) for its associated sheaf. Sheafification preserves stalks. A morphism of sheaves is injective exactly when it is injective on every stalk.

The presheaves here take values in abelian groups.
\end{context}

\begin{problem}
(a) Let \(\phi:\mathcal F\to\mathcal G\) be a morphism of presheaves such that \(\mathcal F(U)\to\mathcal G(U)\) is injective for every open \(U\). Prove that \(\phi^+:\mathcal F^+\to\mathcal G^+\) is injective.

(b) Deduce that the image sheaf of a morphism of sheaves \(\mathcal F\to\mathcal G\) is naturally a subsheaf of \(\mathcal G\).
\end{problem}
\end{exerciseintro}

\begin{solution}
(a) The map on stalks \(\phi_x:\mathcal F_x\to\mathcal G_x\) is injective. Indeed, if a germ represented by \(s\in\mathcal F(U)\) maps to zero, then after restricting to some neighborhood \(V\subseteq U\) of \(x\), the section \(\phi(s|_V)\) is zero. Injectivity of \(\mathcal F(V)\to\mathcal G(V)\) implies \(s|_V=0\), so the original germ is zero. Because sheafification preserves stalks, \(\phi^+\) is injective on every stalk and therefore injective as a sheaf morphism.

(b) Let \(\mathcal I\) be the presheaf defined by
\[
\mathcal I(U)=\operatorname{im}\bigl(\mathcal F(U)\to\mathcal G(U)\bigr).
\]
The evident morphism \(\mathcal I\to\mathcal G\) is injective on every open set. Part (a) therefore gives an injection \(\mathcal I^+\to\mathcal G^+\). Since \(\mathcal G\) is already a sheaf, \(\mathcal G^+\cong\mathcal G\), while \(\mathcal I^+\) is by definition the image sheaf. Hence the image sheaf is naturally a subsheaf of \(\mathcal G\).
\end{solution}

\begin{exerciseintro}
% record: H-II-01-005
\exerciseheading{Exercise II.1.5}
\addcontentsline{toc}{subsection}{Exercise II.1.5}

\begin{context}
A morphism of sheaves is injective or surjective exactly when every induced map on stalks is injective or surjective, respectively. A morphism of sheaves is an isomorphism exactly when it is an isomorphism on every stalk.
\end{context}

\begin{problem}
Prove that a morphism of sheaves is an isomorphism if and only if it is both injective and surjective.
\end{problem}
\end{exerciseintro}

\begin{solution}
Every isomorphism is injective and surjective. Conversely, let \(\phi:\mathcal F\to\mathcal G\) be both injective and surjective. For every \(x\in X\), the map
\[
\phi_x:\mathcal F_x\longrightarrow\mathcal G_x
\]
is both injective and surjective, hence is an isomorphism of abelian groups. The stalk criterion for isomorphisms then implies that \(\phi\) is an isomorphism of sheaves.

Alternatively, local lifts give a direct construction of the inverse. For an open set \(U\) and a section \(s\in\mathcal G(U)\), surjectivity provides an open cover \(U=\bigcup_i U_i\) and sections \(t_i\in\mathcal F(U_i)\) with \(\phi(t_i)=s|_{U_i}\). On an overlap \(U_i\cap U_j\), injectivity gives \(t_i|_{U_i\cap U_j}=t_j|_{U_i\cap U_j}\). The sheaf axiom therefore glues the \(t_i\) to a unique section \(t\in\mathcal F(U)\). This construction commutes with restriction and defines a morphism \(\mathcal G\to\mathcal F\) inverse to \(\phi\).
\end{solution}

\begin{exerciseintro}
% record: H-II-01-006
\exerciseheading{Exercise II.1.6}
\addcontentsline{toc}{subsection}{Exercise II.1.6}

\begin{context}
If \(\mathcal F'\subseteq\mathcal F\) is a subsheaf, the quotient sheaf \(\mathcal F/\mathcal F'\) is the sheafification of \(U\mapsto\mathcal F(U)/\mathcal F'(U)\), and
\[
(\mathcal F/\mathcal F')_x\cong\mathcal F_x/\mathcal F'_x.
\]
Exactness of sheaves can be checked on stalks.
\end{context}

\begin{problem}
(a) Prove that the canonical map \(\mathcal F\to\mathcal F/\mathcal F'\) is surjective with kernel \(\mathcal F'\).

(b) Conversely, given an exact sequence
\[
0\longrightarrow\mathcal F'\longrightarrow\mathcal F\longrightarrow\mathcal F''\longrightarrow0,
\]
prove that \(\mathcal F'\) is identified with a subsheaf of \(\mathcal F\) and \(\mathcal F''\cong\mathcal F/\mathcal F'\).
\end{problem}
\end{exerciseintro}

\begin{solution}
(a) At every point \(x\), the canonical sequence is
\[
0\longrightarrow\mathcal F'_x\longrightarrow\mathcal F_x
\longrightarrow\mathcal F_x/\mathcal F'_x\longrightarrow0.
\]
It is exact as a sequence of abelian groups. The displayed identification of the quotient stalk and the stalkwise exactness criterion therefore give an exact sequence of sheaves
\[
0\longrightarrow\mathcal F'\longrightarrow\mathcal F
\longrightarrow\mathcal F/\mathcal F'\longrightarrow0.
\]

(b) Exactness at \(\mathcal F'\) makes \(\mathcal F'\to\mathcal F\) injective, and exactness at \(\mathcal F\) identifies its image with \(\ker(\mathcal F\to\mathcal F'')\). We may therefore identify \(\mathcal F'\) with this subsheaf of \(\mathcal F\). The map \(\mathcal F\to\mathcal F''\) annihilates \(\mathcal F'\), so the universal property of the quotient produces a morphism
\[
\mathcal F/\mathcal F'\longrightarrow\mathcal F''.
\]
On each stalk this is the canonical isomorphism
\[
\mathcal F_x/\mathcal F'_x\longrightarrow\mathcal F''_x
\]
coming from exactness. Hence it is an isomorphism of sheaves.
\end{solution}

\begin{exerciseintro}
% record: H-II-01-007
\exerciseheading{Exercise II.1.7}
\addcontentsline{toc}{subsection}{Exercise II.1.7}

\begin{context}
For a morphism of sheaves \(\phi:\mathcal F\to\mathcal G\), kernels, images, and cokernels have stalks equal to the corresponding kernels, images, and cokernels of \(\phi_x\). Quotient sheaves satisfy \((\mathcal F/\mathcal K)_x=\mathcal F_x/\mathcal K_x\).
\end{context}

\begin{problem}
For a morphism \(\phi:\mathcal F\to\mathcal G\), prove the natural isomorphisms

(a) \(\operatorname{im}\phi\cong\mathcal F/\ker\phi\);

(b) \(\operatorname{coker}\phi\cong\mathcal G/\operatorname{im}\phi\).
\end{problem}
\end{exerciseintro}

\begin{solution}
(a) The morphism \(\mathcal F\to\operatorname{im}\phi\) is surjective and has kernel \(\ker\phi\), so it induces a natural map
\[
\mathcal F/\ker\phi\longrightarrow\operatorname{im}\phi.
\]
At each point \(x\), this is the usual first-isomorphism map
\[
\mathcal F_x/\ker(\phi_x)\longrightarrow\operatorname{im}(\phi_x),
\]
which is an isomorphism. The induced sheaf morphism is therefore an isomorphism.

(b) By definition, the cokernel sheaf is the sheafification of the presheaf cokernel. The composite \(\mathcal F\to\mathcal G\to\mathcal G/\operatorname{im}\phi\) is zero, so it induces a natural morphism
\[
\operatorname{coker}\phi\longrightarrow\mathcal G/\operatorname{im}\phi.
\]
Its stalk at \(x\) is the canonical isomorphism
\[
\operatorname{coker}(\phi_x)
=\mathcal G_x/\operatorname{im}(\phi_x).
\]
Hence the morphism is an isomorphism of sheaves.
\end{solution}

\begin{exerciseintro}
% record: H-II-01-008
\exerciseheading{Exercise II.1.8}
\addcontentsline{toc}{subsection}{Exercise II.1.8}

\begin{context}
Sheaves are sheaves of abelian groups on \(X\). For a morphism \(\mathcal F\to\mathcal G\), the kernel sheaf is computed sectionwise:
\[
(\ker(\mathcal F\to\mathcal G))(U)
=\ker(\mathcal F(U)\to\mathcal G(U)).
\]
\end{context}

\begin{problem}
For an open subset \(U\subseteq X\), prove that the section functor \(\Gamma(U,-)\) is left exact. Thus, from an exact sequence
\[
0\longrightarrow\mathcal F'\longrightarrow\mathcal F\longrightarrow\mathcal F'',
\]
deduce an exact sequence
\[
0\longrightarrow\Gamma(U,\mathcal F')\longrightarrow\Gamma(U,\mathcal F)
\longrightarrow\Gamma(U,\mathcal F'').
\]
\end{problem}
\end{exerciseintro}

\begin{solution}
Exactness of the sheaf sequence says that \(\mathcal F'\) is the kernel of the morphism \(\mathcal F\to\mathcal F''\). Since kernels of sheaves are computed sectionwise,
\[
\Gamma(U,\mathcal F')
=\ker\bigl(\Gamma(U,\mathcal F)\to\Gamma(U,\mathcal F'')\bigr).
\]
The map \(\Gamma(U,\mathcal F')\to\Gamma(U,\mathcal F)\) is therefore injective, and its image is exactly the kernel of the following map. This is precisely the claimed exactness. No surjectivity onto \(\Gamma(U,\mathcal F'')\) is asserted; local lifts of a section need not glue globally.
\end{solution}

\begin{exerciseintro}
% record: H-II-01-009
\exerciseheading{Exercise II.1.9}
\addcontentsline{toc}{subsection}{Exercise II.1.9}

\begin{context}
Let \(\mathcal F\) and \(\mathcal G\) be sheaves of abelian groups on \(X\). In an additive category, a biproduct is an object that satisfies the universal properties of both product and coproduct.
\end{context}

\begin{problem}
Prove that the presheaf
\[
U\longmapsto\mathcal F(U)\oplus\mathcal G(U)
\]
is a sheaf and that it is simultaneously the direct product and direct sum of \(\mathcal F\) and \(\mathcal G\) in the category of sheaves of abelian groups.
\end{problem}
\end{exerciseintro}

\begin{solution}
Restrictions are defined componentwise. If pairs of sections \((s_i,t_i)\) agree on all overlaps of an open cover, then the \(s_i\) glue uniquely to a section of \(\mathcal F\), and the \(t_i\) glue uniquely to a section of \(\mathcal G\). Their pair is the unique glued section. Thus the displayed presheaf is a sheaf; denote it by \(\mathcal F\oplus\mathcal G\).

The component projections define morphisms
\[
\mathcal F\oplus\mathcal G\longrightarrow\mathcal F,
\qquad
\mathcal F\oplus\mathcal G\longrightarrow\mathcal G.
\]
Given morphisms \(\mathcal H\to\mathcal F\) and \(\mathcal H\to\mathcal G\), pairing them sectionwise produces a unique morphism \(\mathcal H\to\mathcal F\oplus\mathcal G\). Hence this sheaf is the product.

Similarly, the component inclusions define morphisms from \(\mathcal F\) and \(\mathcal G\) into \(\mathcal F\oplus\mathcal G\). Given morphisms \(\mathcal F\to\mathcal H\) and \(\mathcal G\to\mathcal H\), their sectionwise sum produces the unique morphism \(\mathcal F\oplus\mathcal G\to\mathcal H\) extending both. Hence it is also the coproduct, and therefore the direct sum.
\end{solution}

\begin{exerciseintro}
% record: H-II-01-010
\exerciseheading{Exercise II.1.10}
\addcontentsline{toc}{subsection}{Exercise II.1.10}

\begin{context}
Let \(\{\mathcal F_i\}\) be a directed system of sheaves on \(X\). Colimits of presheaves are computed sectionwise. If \(\mathcal P\) is a presheaf and \(\mathcal G\) is a sheaf, every presheaf morphism \(\mathcal P\to\mathcal G\) factors uniquely through the sheafification \(\mathcal P^+\).

All sheaves in this exercise are sheaves of abelian groups.
\end{context}

\begin{problem}
Define \(\varinjlim_i\mathcal F_i\) to be the sheaf associated to the presheaf
\[
U\longmapsto\varinjlim_i\mathcal F_i(U).
\]
Prove that this sheaf satisfies the universal property of the direct limit in the category of sheaves.
\end{problem}
\end{exerciseintro}

\begin{solution}
Let \(\mathcal P\) be the sectionwise direct-limit presheaf. The canonical maps \(\mathcal F_i(U)\to\mathcal P(U)\) commute with restriction and with the transition maps, so after sheafification they give a compatible family
\[
\mathcal F_i\longrightarrow\mathcal P^+.
\]

Now let \(\mathcal G\) be a sheaf equipped with compatible morphisms \(\psi_i:\mathcal F_i\to\mathcal G\). For every open \(U\), the universal property of the direct limit of abelian groups gives a unique homomorphism
\[
\mathcal P(U)=\varinjlim_i\mathcal F_i(U)\longrightarrow\mathcal G(U)
\]
compatible with the \(\psi_i(U)\). These homomorphisms commute with restriction, so they form a morphism of presheaves \(\mathcal P\to\mathcal G\). Because \(\mathcal G\) is a sheaf, this morphism factors uniquely through a sheaf morphism
\[
\mathcal P^+\longrightarrow\mathcal G.
\]
The factorization is compatible with every \(\psi_i\), and uniqueness follows first from the presheaf colimit and then from sheafification. Thus \(\mathcal P^+\) is the direct limit in the category of sheaves.
\end{solution}

\begin{exerciseintro}
% record: H-II-01-011
\exerciseheading{Exercise II.1.11}
\addcontentsline{toc}{subsection}{Exercise II.1.11}

\begin{context}
A topological space is noetherian if every ascending chain of open subsets
stabilizes; equivalently, every open subset is quasi-compact.  A filtered
direct limit of abelian groups commutes with finite products and finite
equalizers.  For a finite open cover \(U=\bigcup_{a=1}^n U_a\), the sheaf
axiom for \(\mathcal F\) says that
\[
0\longrightarrow \mathcal F(U)\longrightarrow
\prod_a\mathcal F(U_a)
\rightrightarrows
\prod_{a,b}\mathcal F(U_a\cap U_b)
\]
is an equalizer diagram.
\end{context}

\begin{problem}
Let \(\{\mathcal F_i\}_{i\in I}\) be a filtered direct system of sheaves of
abelian groups on a noetherian topological space \(X\).  Prove that the
sectionwise direct-limit presheaf
\[
\mathcal P(U)=\varinjlim_i\mathcal F_i(U)
\]
already satisfies the sheaf axioms.  Deduce that
\[
\Gamma\!\left(X,\varinjlim_i\mathcal F_i\right)
\cong \varinjlim_i\Gamma(X,\mathcal F_i).
\]
\end{problem}
\end{exerciseintro}

\begin{solution}
First consider a finite open cover \(U=U_1\cup\cdots\cup U_n\).  For every
\(i\), the sheaf axiom gives an equalizer diagram
\[
0\longrightarrow \mathcal F_i(U)\longrightarrow
\prod_a\mathcal F_i(U_a)
\rightrightarrows
\prod_{a,b}\mathcal F_i(U_a\cap U_b).
\]
Taking the filtered direct limit preserves this diagram because all products
displayed here are finite.  Hence \(\mathcal P\) has unique gluing for every
finite open cover.

Now let \(U=\bigcup_{\alpha\in A}U_\alpha\) be an arbitrary cover and let
\(s_\alpha\in\mathcal P(U_\alpha)\) be compatible sections.  Since \(U\) is
quasi-compact, choose a finite subcover
\(U_{\alpha_1},\ldots,U_{\alpha_n}\).  The finite-cover result produces a
unique \(s\in\mathcal P(U)\) whose restriction to each selected member is
\(s_{\alpha_a}\).

For any \(\alpha\), the finite family
\[
U_\alpha\cap U_{\alpha_1},\ldots,
U_\alpha\cap U_{\alpha_n}
\]
covers \(U_\alpha\).  On each member of this cover, compatibility gives
\(s|_{U_\alpha\cap U_{\alpha_a}}
=s_\alpha|_{U_\alpha\cap U_{\alpha_a}}\).  Uniqueness for finite covers
therefore implies \(s|_{U_\alpha}=s_\alpha\).  The same argument, applied to
the difference of two putative gluings, proves uniqueness.  Thus
\(\mathcal P\) is a sheaf.

Consequently no sheafification changes the sectionwise direct limit.  Taking
\(U=X\) in its definition gives
\[
\Gamma\!\left(X,\varinjlim_i\mathcal F_i\right)
=\varinjlim_i\mathcal F_i(X)
=\varinjlim_i\Gamma(X,\mathcal F_i).
\]
\end{solution}

\begin{exerciseintro}
% record: H-II-01-012
\exerciseheading{Exercise II.1.12}
\addcontentsline{toc}{subsection}{Exercise II.1.12}

\begin{context}
Limits of abelian groups are computed as compatible families inside products.
A presheaf of abelian groups is a sheaf precisely when every compatible
family of local sections has a unique gluing.
\end{context}

\begin{problem}
For an inverse system \(\{\mathcal F_i\}\) of sheaves on a topological space
\(X\), prove that
\[
U\longmapsto\varprojlim_i\mathcal F_i(U)
\]
is a sheaf.  Show that the resulting sheaf, denoted
\(\varprojlim_i\mathcal F_i\), has the universal property of the inverse
limit in the category of sheaves.
\end{problem}
\end{exerciseintro}

\begin{solution}
Write \(\mathcal L(U)=\varprojlim_i\mathcal F_i(U)\), with restrictions
defined componentwise.  Suppose \(U=\bigcup_\alpha U_\alpha\) and that
\(\ell_\alpha\in\mathcal L(U_\alpha)\) agree on overlaps.  For every \(i\),
let \(\ell_{\alpha,i}\in\mathcal F_i(U_\alpha)\) be the \(i\)-th component.
The sheaf axiom for \(\mathcal F_i\) gives a unique
\(\ell_i\in\mathcal F_i(U)\) restricting to all the
\(\ell_{\alpha,i}\).

If \(i\to j\) is a transition arrow in the inverse system, the image of
\(\ell_i\) and the section \(\ell_j\) have the same restriction to every
\(U_\alpha\).  They are therefore equal by uniqueness in
\(\mathcal F_j\).  Thus the family \((\ell_i)_i\) is compatible and defines
\(\ell\in\mathcal L(U)\).  Componentwise uniqueness proves uniqueness of
\(\ell\), so \(\mathcal L\) is a sheaf.

The projections \(\mathcal L\to\mathcal F_i\) form a compatible cone.  If
\(\mathcal G\) is a sheaf with a compatible cone
\(\psi_i:\mathcal G\to\mathcal F_i\), then for every open \(U\) the maps
\(\psi_i(U)\) induce a unique homomorphism
\[
\mathcal G(U)\longrightarrow\varprojlim_i\mathcal F_i(U)=\mathcal L(U).
\]
These homomorphisms commute with restriction because all the \(\psi_i\) do,
and hence define a sheaf morphism \(\mathcal G\to\mathcal L\).  Its
components are the prescribed maps, and componentwise uniqueness makes it
the only such morphism.  This is exactly the inverse-limit universal
property.
\end{solution}

\begin{exerciseintro}
% record: H-II-01-013
\exerciseheading{Exercise II.1.13}
\addcontentsline{toc}{subsection}{Exercise II.1.13}

\begin{context}
For a presheaf \(\mathcal F\), its associated sheaf may be realized as the
sheaf whose sections over \(U\) are functions \(t\) assigning to each
\(x\in U\) a germ \(t(x)\in\mathcal F_x\), subject to the condition that
every point has a neighborhood on which these germs are represented by one
section of \(\mathcal F\).  The topology on a set receiving maps
\(\rho_a:T_a\to E\) that is strongest among those making every \(\rho_a\)
continuous is characterized by: \(W\subseteq E\) is open if and only if
\(\rho_a^{-1}(W)\) is open for every \(a\).
\end{context}

\begin{problem}
Let \(\mathcal F\) be a presheaf on \(X\).  Form
\[
E(\mathcal F)=\coprod_{x\in X}\mathcal F_x
\]
and let \(\pi:E(\mathcal F)\to X\) send a germ in \(\mathcal F_x\) to \(x\).
For \(s\in\mathcal F(U)\), define
\(\bar s:U\to E(\mathcal F)\) by \(\bar s(x)=s_x\).  Give
\(E(\mathcal F)\) the strongest topology for which every such map \(\bar s\)
is continuous.  Prove that the associated sheaf \(\mathcal F^+\) is the
sheaf of continuous sections of \(\pi\).  In particular, characterize when
\(\mathcal F\) itself is a sheaf in these terms.
\end{problem}
\end{exerciseintro}

\begin{solution}
For \(s\in\mathcal F(U)\), put
\[
B(s,U)=\{s_x:x\in U\}=\bar s(U).
\]
For the specified topology, if an open \(W\) contains \(s_x\), then
\(\bar s^{-1}(W)\) is an open neighborhood of \(x\).
The image of the restriction of \(s\) to that neighborhood is
a set \(B(s|_{\bar s^{-1}(W)},\bar s^{-1}(W))\) contained in \(W\).
Thus the basic sets just defined refine every open of that topology.

We first show that these sets form a basis for the topology on
\(E(\mathcal F)\).  Given \(t\in\mathcal F(V)\), the inverse image of
\(B(s,U)\) under \(\bar t\) is
\[
\{x\in U\cap V:t_x=s_x\}.
\]
Equality of two germs holds on a neighborhood of the point where it holds,
so this set is open.  Hence every \(B(s,U)\) is open by the definition of
the topology.  These sets cover \(E(\mathcal F)\), since every germ has a
representative.  If a germ \(e\) belongs to both \(B(s,U)\) and \(B(t,V)\),
then \(s\) and \(t\) have equal germs at \(\pi(e)\), and after shrinking to
an open neighborhood \(W\subseteq U\cap V\) their restrictions agree.
Thus \(e\in B(s|_W,W)\subseteq B(s,U)\cap B(t,V)\), proving the basis claim.

The projection \(\pi\) is continuous: for open \(W\subseteq X\), the
preimage \(\pi^{-1}(W)\) has open inverse image under every map \(\bar s\).
Moreover, \(\pi:B(s,U)\to U\) is a homeomorphism with inverse \(\bar s\).

Let \(t:U\to E(\mathcal F)\) satisfy \(\pi\circ t\) equal to the identity.
If \(t\) is continuous and \(x\in U\), choose a representative
\(s\in\mathcal F(V)\) of the germ \(t(x)\), with \(x\in V\).  The basic
open set \(B(s,V)\) contains \(t(x)\), so
\[
W=t^{-1}(B(s,V))
\]
is a neighborhood of \(x\) in \(U\).  For \(y\in W\), both \(t(y)\) and
\(s_y\) lie over \(y\) in \(B(s,V)\); hence \(t(y)=s_y\).  Thus \(t\) is
locally represented by sections of \(\mathcal F\).

Conversely, a section of \(\pi\) that is locally of the form \(\bar s\) is
continuous, because continuity is local on its domain and every \(\bar s\)
is continuous by construction.  Therefore continuous sections of \(\pi\)
are exactly the locally representable germ functions, which are the
sections of \(\mathcal F^+\) by the stated realization of sheafification.
The identifications commute with restriction, so they give an isomorphism
of sheaves.

The canonical map \(\mathcal F(U)\to\mathcal F^+(U)\) sends \(s\) to
\(\bar s\).  Hence \(\mathcal F\) is a sheaf if and only if, for every open
\(U\), this map identifies \(\mathcal F(U)\) with all continuous sections
of \(\pi\) over \(U\).
\end{solution}

\begin{exerciseintro}
% record: H-II-01-014
\exerciseheading{Exercise II.1.14}
\addcontentsline{toc}{subsection}{Exercise II.1.14}

\begin{context}
For a section \(s\in\mathcal F(U)\), the equality \(s_x=0\) in the stalk
\(\mathcal F_x\) means that \(s\) restricts to zero on some neighborhood of
\(x\) contained in \(U\).  Stalks commute with direct sums of sheaves.  On
a \(T_1\)-space, the skyscraper sheaf with value \(A\ne0\) at a point \(p\)
has nonzero stalk precisely at \(p\).
\end{context}

\begin{problem}
For a sheaf \(\mathcal F\) on \(X\) and a section
\(s\in\mathcal F(U)\), define
\[
\operatorname{Supp}(s)=\{x\in U:s_x\ne0\}.
\]
Prove that \(\operatorname{Supp}(s)\) is closed in \(U\).  The support of
the sheaf is
\[
\operatorname{Supp}(\mathcal F)=\{x\in X:\mathcal F_x\ne0\}.
\]
Give an example showing that the latter set need not be closed.
\end{problem}
\end{exerciseintro}

\begin{solution}
If \(x\notin\operatorname{Supp}(s)\), then \(s_x=0\).  By the definition of
a stalk, there is an open neighborhood \(V\subseteq U\) of \(x\) on which
\(s|_V=0\).  Every germ \(s_y\) with \(y\in V\) is therefore zero, so
\[
V\subseteq U\setminus\operatorname{Supp}(s).
\]
The complement of the support is open in \(U\), and the support is closed.

For the second assertion, take \(X=\mathbb R\) with its usual topology and,
for every \(q\in\mathbb Q\), let \(\mathcal S_q\) be the skyscraper sheaf
with value \(\mathbb Z\) at \(q\).  Set
\[
\mathcal F=\bigoplus_{q\in\mathbb Q}\mathcal S_q.
\]
Because stalks commute with direct sums,
\[
\mathcal F_x\cong
\bigoplus_{q\in\mathbb Q}(\mathcal S_q)_x.
\]
This group is \(\mathbb Z\) when \(x\in\mathbb Q\) and is zero when
\(x\notin\mathbb Q\).  Thus
\(\operatorname{Supp}(\mathcal F)=\mathbb Q\), which is not closed in
\(\mathbb R\).
\end{solution}

\begin{exerciseintro}
% record: H-II-01-015
\exerciseheading{Exercise II.1.15}
\addcontentsline{toc}{subsection}{Exercise II.1.15}

\begin{context}
Morphisms of sheaves of abelian groups may be added pointwise on sections.
Sections of a sheaf that agree on an open cover are equal, and compatible
local sections glue uniquely.
\end{context}

\begin{problem}
Let \(\mathcal F\) and \(\mathcal G\) be sheaves of abelian groups on \(X\).
Show that
\[
U\longmapsto
\operatorname{Hom}(\mathcal F|_U,\mathcal G|_U)
\]
is naturally a sheaf of abelian groups, with restriction given by restricting
a sheaf morphism.  This is the sheaf of local morphisms from
\(\mathcal F\) to \(\mathcal G\).
\end{problem}
\end{exerciseintro}

\begin{solution}
For a fixed open \(U\), define the sum of two morphisms by
\[
(\varphi+\psi)_V(s)=\varphi_V(s)+\psi_V(s)
\qquad
(V\subseteq U,\ s\in\mathcal F(V)).
\]
The zero and additive inverses are defined in the same way.  Naturality with
respect to restrictions follows from that of \(\varphi\) and \(\psi\), so
this makes
\(\operatorname{Hom}(\mathcal F|_U,\mathcal G|_U)\) an abelian group.
Restricting morphisms is a group homomorphism and gives a presheaf.

Let \(U=\bigcup_\alpha U_\alpha\), and suppose morphisms
\[
\varphi_\alpha:\mathcal F|_{U_\alpha}\longrightarrow
\mathcal G|_{U_\alpha}
\]
agree after restriction to every \(U_\alpha\cap U_\beta\).  For an open
\(V\subseteq U\) and \(s\in\mathcal F(V)\), the sections
\[
\varphi_{\alpha,V\cap U_\alpha}
  (s|_{V\cap U_\alpha})\in\mathcal G(V\cap U_\alpha)
\]
agree on pairwise overlaps.  They therefore glue uniquely to a section
\(\varphi_V(s)\in\mathcal G(V)\).

The uniqueness of these gluings shows that \(\varphi_V\) is a homomorphism,
that the maps \(\varphi_V\) commute with restriction as \(V\) varies, and
that the resulting sheaf morphism \(\varphi:\mathcal F|_U\to\mathcal G|_U\)
restricts to every \(\varphi_\alpha\).  If two global morphisms have those
restrictions, their values on every section agree locally and hence agree
globally in \(\mathcal G\).  Thus the gluing is unique, proving the sheaf
axiom.
\end{solution}

\begin{exerciseintro}
% record: H-II-01-016
\exerciseheading{Exercise II.1.16}
\addcontentsline{toc}{subsection}{Exercise II.1.16}

\begin{context}
A sheaf \(\mathcal F\) of abelian groups is flasque if every restriction
\(\mathcal F(U)\to\mathcal F(V)\), for open \(V\subseteq U\), is
surjective.  A short exact sequence of sheaves is exact on stalks.  In
particular, an epimorphism of sheaves has local lifts of sections, while the
kernel sheaf is computed sectionwise.  Every set can be well ordered.
\end{context}

\begin{problem}
(a) Prove that a constant sheaf on an irreducible topological space is
flasque.

(b) Let
\[
0\longrightarrow\mathcal F'\longrightarrow\mathcal F
\longrightarrow\mathcal F''\longrightarrow0
\]
be exact, with \(\mathcal F'\) flasque.  Prove that for every open \(U\)
the induced sequence
\[
0\longrightarrow\mathcal F'(U)\longrightarrow\mathcal F(U)
\longrightarrow\mathcal F''(U)\longrightarrow0
\]
is exact.

(c) In the sequence of part (b), assume both \(\mathcal F'\) and
\(\mathcal F\) are flasque.  Prove that \(\mathcal F''\) is flasque.

(d) If \(f:X\to Y\) is continuous and \(\mathcal F\) is flasque on \(X\),
prove that \(f_*\mathcal F\) is flasque on \(Y\).

(e) For any sheaf \(\mathcal F\), define
\[
\mathcal G(U)=\prod_{x\in U}\mathcal F_x,
\]
with restriction given by forgetting components.  Thus a section of
\(\mathcal G\) is an arbitrary, not necessarily continuous, choice of a germ
over every point.  Prove that \(\mathcal G\) is a flasque sheaf and that the
germ maps give a natural injection \(\mathcal F\to\mathcal G\).
\end{problem}
\end{exerciseintro}

\begin{solution}
(a) Let \(\mathcal A\) be the constant sheaf associated to an abelian group
\(A\).  Every nonempty open subset of an irreducible space is irreducible,
hence connected.  Therefore a locally constant map from a nonempty open set
to \(A\) is constant, so
\[
\mathcal A(U)=
\begin{cases}
A,&U\ne\varnothing,\\
0,&U=\varnothing.
\end{cases}
\]
Every restriction is either the identity of \(A\) or the unique map to the
zero group, and is consequently surjective.

(b) Left exactness on sections gives exactness through
\(\mathcal F(U)\); only surjectivity onto \(\mathcal F''(U)\) remains.
Take \(s\in\mathcal F''(U)\).  Stalkwise surjectivity gives an open cover
\(U=\bigcup_{\alpha<\kappa}U_\alpha\), indexed by an ordinal, and lifts
\(u_\alpha\in\mathcal F(U_\alpha)\) of \(s|_{U_\alpha}\).

We construct compatible lifts on the increasing open sets
\[
V_\alpha=\bigcup_{\beta<\alpha}U_\beta.
\]
Start with \(t_0=0\) on \(V_0=\varnothing\).
At a limit ordinal, the previously constructed compatible sections glue
uniquely on \(V_\alpha\).  At a successor stage, suppose
\(t_\alpha\in\mathcal F(V_\alpha)\) has already been constructed and put
\(W=V_\alpha\cap U_\alpha\).  The difference
\[
t_\alpha|_W-u_\alpha|_W
\]
maps to zero in \(\mathcal F''(W)\).  Exactness at \(\mathcal F\) identifies
it with a section \(a\in\mathcal F'(W)\).  Since \(\mathcal F'\) is
flasque, extend \(a\) to
\(\widetilde a\in\mathcal F'(U_\alpha)\).  The sections
\(t_\alpha\) and \(u_\alpha+\widetilde a\) now agree on \(W\), so they glue
to a lift on \(V_{\alpha+1}\).  Transfinite induction yields a lift
\(t_\kappa\in\mathcal F(U)\) of \(s\), proving surjectivity.

(c) Let \(V\subseteq U\) be open and \(s\in\mathcal F''(V)\).  Applying
part (b) on the open subspace \(V\) lifts \(s\) to
\(t\in\mathcal F(V)\).  Flasqueness of \(\mathcal F\) extends \(t\) to a
section \(\widetilde t\in\mathcal F(U)\).  Its image in
\(\mathcal F''(U)\) restricts to \(s\).  Thus every restriction map of
\(\mathcal F''\) is surjective.

(d) For open \(V\subseteq U\subseteq Y\), the restriction in the direct
image is
\[
(f_*\mathcal F)(U)=\mathcal F(f^{-1}U)
\longrightarrow
\mathcal F(f^{-1}V)=(f_*\mathcal F)(V).
\]
It is surjective because \(f^{-1}V\subseteq f^{-1}U\) are open in \(X\).

(e) An arbitrary family of functions into the pointwise groups glues
uniquely, point by point, so \(U\mapsto\prod_{x\in U}\mathcal F_x\) is a
sheaf.  If \(V\subseteq U\), a member of \(\mathcal G(V)\) extends to
\(\mathcal G(U)\) by assigning zero to every point of \(U\setminus V\).
Thus \(\mathcal G\) is flasque.

For each open \(U\), define
\[
\iota_U:\mathcal F(U)\longrightarrow\mathcal G(U),
\qquad
s\longmapsto(s_x)_{x\in U}.
\]
These homomorphisms commute with restriction.  If \(\iota_U(s)=0\), every
point \(x\in U\) has a neighborhood on which \(s\) is zero.  Those
neighborhoods cover \(U\), and uniqueness in the sheaf axiom gives \(s=0\).
Hence \(\iota\) is injective.
\end{solution}

\begin{exerciseintro}
% record: H-II-01-017
\exerciseheading{Exercise II.1.17}
\addcontentsline{toc}{subsection}{Exercise II.1.17}

\begin{context}
For \(P\in X\), a point \(Q\) lies in \(\overline{\{P\}}\) precisely when
every open neighborhood of \(Q\) contains \(P\).  The closure of a single
point is irreducible.  The constant sheaf associated to an abelian group
\(A\) has section group \(A\) on each nonempty open subset of an
irreducible space and zero on the empty set.
\end{context}

\begin{problem}
Let \(X\) be a topological space, \(P\in X\), and \(A\) an abelian group.
Define
\[
i_P(A)(U)=
\begin{cases}
A,&P\in U,\\
0,&P\notin U,
\end{cases}
\]
with the evident restriction maps.  Verify that this is a sheaf and compute
its stalk at every \(Q\in X\).  Then identify it with \(i_*\mathcal A\),
where \(i:\overline{\{P\}}\to X\) is the inclusion and \(\mathcal A\) is the
constant sheaf associated to \(A\) on \(\overline{\{P\}}\).
\end{problem}
\end{exerciseintro}

\begin{solution}
Let \(U=\bigcup_\alpha U_\alpha\).  If \(P\notin U\), every group in the
sheaf axiom is zero.  If \(P\in U\), at least one \(U_\alpha\) contains
\(P\).  A compatible family is then uniquely determined by its value in
\(A\) on any member containing \(P\), and that value gives its unique
gluing in \(i_P(A)(U)=A\).  Thus \(i_P(A)\) is a sheaf.

Its stalk at \(Q\) is
\[
i_P(A)_Q=\varinjlim_{Q\in V}i_P(A)(V).
\]
If \(Q\notin\overline{\{P\}}\), some neighborhood \(V\) of \(Q\) omits
\(P\); such neighborhoods are cofinal after intersecting with \(V\), so
the stalk is zero.  If \(Q\in\overline{\{P\}}\), every neighborhood of
\(Q\) contains \(P\).  The directed system is then constantly \(A\) with
identity transition maps, and the stalk is \(A\).  Hence
\[
i_P(A)_Q\cong
\begin{cases}
A,&Q\in\overline{\{P\}},\\
0,&Q\notin\overline{\{P\}}.
\end{cases}
\]

Put \(Z=\overline{\{P\}}\).  Since \(Z\) is irreducible,
\(\mathcal A(W)=A\) for nonempty open \(W\subseteq Z\), and
\(\mathcal A(\varnothing)=0\).  For open \(U\subseteq X\),
\[
(i_*\mathcal A)(U)=\mathcal A(U\cap Z).
\]
An open set \(U\) meets \(Z\) if and only if it contains \(P\): if
\(Q\in U\cap Z\), the neighborhood \(U\) of \(Q\) must contain \(P\).
Consequently \((i_*\mathcal A)(U)\) is \(A\) exactly when \(P\in U\), and
is zero otherwise.  These identifications respect restrictions, giving
\(i_P(A)\cong i_*\mathcal A\).
\end{solution}

\begin{exerciseintro}
% record: H-II-01-018
\exerciseheading{Exercise II.1.18}
\addcontentsline{toc}{subsection}{Exercise II.1.18}

\begin{context}
For a continuous map \(f:X\to Y\), the direct image is
\((f_*\mathcal F)(V)=\mathcal F(f^{-1}V)\).  For a sheaf \(\mathcal G\) on
\(Y\), the inverse image \(f^{-1}\mathcal G\) is the sheafification of the
presheaf
\[
U\longmapsto
\varinjlim_{V\supseteq f(U)}\mathcal G(V).
\]
A morphism from a presheaf to a sheaf factors uniquely through
sheafification.
\end{context}

\begin{problem}
For a continuous map \(f:X\to Y\), construct natural morphisms
\[
\varepsilon_{\mathcal F}:f^{-1}f_*\mathcal F\longrightarrow\mathcal F
\quad\text{and}\quad
\eta_{\mathcal G}:\mathcal G\longrightarrow f_*f^{-1}\mathcal G.
\]
Use them to prove a natural bijection
\[
\operatorname{Hom}_X(f^{-1}\mathcal G,\mathcal F)
\cong
\operatorname{Hom}_Y(\mathcal G,f_*\mathcal F)
\]
for sheaves \(\mathcal F\) on \(X\) and \(\mathcal G\) on \(Y\).
\end{problem}
\end{exerciseintro}

\begin{solution}
For an open \(U\subseteq X\), every open \(V\subseteq Y\) containing
\(f(U)\) gives a restriction map
\[
(f_*\mathcal F)(V)=\mathcal F(f^{-1}V)\longrightarrow\mathcal F(U).
\]
These maps are compatible as \(V\) shrinks, so they induce a morphism from
the inverse-image presheaf of \(f_*\mathcal F\) to \(\mathcal F\).
Because \(\mathcal F\) is a sheaf, it factors uniquely through
sheafification and gives
\[
\varepsilon_{\mathcal F}:f^{-1}f_*\mathcal F\longrightarrow\mathcal F.
\]
The construction uses only restriction maps, so it is natural in
\(\mathcal F\).

For open \(V\subseteq Y\), a section \(s\in\mathcal G(V)\) determines, by
the canonical map into the inverse-image presheaf and then sheafification,
a section of \(f^{-1}\mathcal G\) over \(f^{-1}V\).  This defines
\[
\eta_{\mathcal G,V}:\mathcal G(V)\longrightarrow
(f_*f^{-1}\mathcal G)(V)
\]
compatibly with restrictions, and hence a natural morphism
\(\eta_{\mathcal G}:\mathcal G\to f_*f^{-1}\mathcal G\).

These maps satisfy the two triangle identities
\[
\varepsilon_{f^{-1}\mathcal G}\circ
f^{-1}(\eta_{\mathcal G})
=\operatorname{id}_{f^{-1}\mathcal G},
\qquad
f_*(\varepsilon_{\mathcal F})\circ
\eta_{f_*\mathcal F}
=\operatorname{id}_{f_*\mathcal F}.
\]
Indeed, before sheafification a section representing a germ is first sent
to the same section over its inverse-image open set and then restricted
back; the result is its original germ.  The first identity follows after
sheafification.  The second is checked directly on
\((f_*\mathcal F)(V)=\mathcal F(f^{-1}V)\), where the composite is the
identity restriction map.

Now define
\[
\Phi(\alpha)=f_*(\alpha)\circ\eta_{\mathcal G}
\]
for \(\alpha:f^{-1}\mathcal G\to\mathcal F\), and
\[
\Psi(\beta)=\varepsilon_{\mathcal F}\circ f^{-1}(\beta)
\]
for \(\beta:\mathcal G\to f_*\mathcal F\).  Naturality of \(\eta\) and
\(\varepsilon\), followed by the two triangle identities, gives
\[
\Psi(\Phi(\alpha))=\alpha,
\qquad
\Phi(\Psi(\beta))=\beta.
\]
Thus \(\Phi\) and \(\Psi\) are inverse bijections.  Every construction is
natural in both sheaves, so \(f^{-1}\) is left adjoint to \(f_*\).
\end{solution}

\begin{exerciseintro}
% record: H-II-01-019
\exerciseheading{Exercise II.1.19}
\addcontentsline{toc}{subsection}{Exercise II.1.19}

\begin{context}
Let \(Z\) be closed in \(X\), let \(U=X\setminus Z\), and write
\(i:Z\to X\) and \(j:U\to X\) for the inclusions.  Direct image is defined
by \((i_*\mathcal F)(V)=\mathcal F(V\cap Z)\).  Sheafification preserves
stalks.  A morphism of sheaves of abelian groups is an isomorphism, or a
sequence of such sheaves is exact, if and only if this holds at every stalk.
\end{context}

\begin{problem}
(a) For a sheaf \(\mathcal F\) on \(Z\), show that
\[
(i_*\mathcal F)_x\cong
\begin{cases}
\mathcal F_x,&x\in Z,\\
0,&x\notin Z.
\end{cases}
\]

(b) For a sheaf \(\mathcal F\) on \(U\), let \(j_!\mathcal F\) be the
sheafification of the presheaf
\[
\mathcal P(V)=
\begin{cases}
\mathcal F(V),&V\subseteq U,\\
0,&V\not\subseteq U.
\end{cases}
\]
Prove that its restriction to \(U\) is \(\mathcal F\), that its stalk is
\(\mathcal F_x\) for \(x\in U\) and zero for \(x\notin U\), and that these
properties characterize it up to a unique isomorphism compatible with the
given identification over \(U\).

(c) For a sheaf \(\mathcal F\) on \(X\), prove that there is a natural exact
sequence
\[
0\longrightarrow j_!(\mathcal F|_U)
\longrightarrow\mathcal F
\longrightarrow i_*(\mathcal F|_Z)
\longrightarrow0.
\]
\end{problem}
\end{exerciseintro}

\begin{solution}
(a) By definition,
\[
(i_*\mathcal F)_x
=\varinjlim_{x\in V}\mathcal F(V\cap Z).
\]
If \(x\in Z\), the sets \(V\cap Z\), as \(V\) ranges through neighborhoods
of \(x\) in \(X\), are exactly a cofinal system of neighborhoods of \(x\)
in \(Z\); the limit is \(\mathcal F_x\).  If \(x\notin Z\), the open
neighborhood \(X\setminus Z\) of \(x\) has empty intersection with \(Z\).
After intersecting all neighborhoods with it, the system is constantly
\(\mathcal F(\varnothing)=0\), so the stalk is zero.

(b) If \(x\in U\), neighborhoods of \(x\) contained in \(U\) are cofinal
among all neighborhoods of \(x\).  On these neighborhoods \(\mathcal P\)
equals \(\mathcal F\), and hence \(\mathcal P_x=\mathcal F_x\).  If
\(x\notin U\), no neighborhood of \(x\) is contained in \(U\), so every
term defining \(\mathcal P_x\) is zero.  Sheafification preserves these
stalks.  Also, \(\mathcal P|_U=\mathcal F\) is already a sheaf, so
\((j_!\mathcal F)|_U\cong\mathcal F\).

For uniqueness, let \(\mathcal G\) be a sheaf on \(X\), supplied with an
isomorphism \(\theta:\mathcal F\to\mathcal G|_U\), and suppose
\(\mathcal G_x=0\) outside \(U\).  Define a presheaf morphism
\(\mathcal P\to\mathcal G\): on \(V\subseteq U\) use
\(\theta_V\), and on every other \(V\) use the zero map.  These maps commute
with restriction, including a restriction from an open not contained in
\(U\) to one contained in \(U\), because the source section is then zero.
The map factors uniquely through a sheaf morphism
\[
j_!\mathcal F\longrightarrow\mathcal G.
\]
It induces \(\theta\) on stalks in \(U\), and is the unique map between zero
stalks outside \(U\); it is therefore an isomorphism on every stalk and
hence an isomorphism of sheaves.  The factorization property also gives
uniqueness among isomorphisms compatible with \(\theta\).

(c) If \(V\subseteq U\), restriction gives
\(\mathcal F(V)\to\mathcal F(V)\); together with the zero maps on opens not
contained in \(U\), these maps induce
\[
j_!(\mathcal F|_U)\longrightarrow\mathcal F.
\]
Restriction to \(Z\), equivalently the unit for inverse and direct image,
gives
\[
\mathcal F\longrightarrow i_*(\mathcal F|_Z).
\]
At \(x\in U\), the resulting stalk sequence is
\[
0\longrightarrow\mathcal F_x
\longrightarrow\mathcal F_x\longrightarrow0\longrightarrow0,
\]
where the middle map is the identity.  At \(x\in Z\), it is
\[
0\longrightarrow0\longrightarrow\mathcal F_x
\longrightarrow\mathcal F_x\longrightarrow0,
\]
again with the nonzero map equal to the identity under the natural stalk
identification.  The sequence is exact at every stalk, hence exact as a
sequence of sheaves.
\end{solution}

\begin{exerciseintro}
% record: H-II-01-020
\exerciseheading{Exercise II.1.20}
\addcontentsline{toc}{subsection}{Exercise II.1.20}

\begin{context}
For \(s\in\mathcal F(V)\), its support is
\(\{x\in V:s_x\ne0\}\).  If \(j:U\to X\) is the inclusion of an open
subset, then
\[
(j_*(\mathcal F|_U))(V)=\mathcal F(V\cap U).
\]
Kernels of morphisms of sheaves of abelian groups are computed sectionwise.
\end{context}

\begin{problem}
Let \(Z\) be a closed subset of \(X\) and let \(\mathcal F\) be a sheaf on
\(X\).

(a) Prove that
\[
V\longmapsto
\{s\in\mathcal F(V):\operatorname{Supp}(s)\subseteq Z\cap V\}
\]
is a sheaf.  Denote it by \(\mathcal H_Z^0(\mathcal F)\).

(b) Put \(U=X\setminus Z\) and let \(j:U\to X\) be the inclusion.  Prove
that
\[
0\longrightarrow\mathcal H_Z^0(\mathcal F)
\longrightarrow\mathcal F
\longrightarrow j_*(\mathcal F|_U)
\]
is exact.  If \(\mathcal F\) is flasque, prove in addition that the last
map is surjective.
\end{problem}
\end{exerciseintro}

\begin{solution}
(a) Restriction preserves the support condition: if \(W\subseteq V\), then
\[
\operatorname{Supp}(s|_W)
\subseteq\operatorname{Supp}(s)\cap W
\subseteq Z\cap W.
\]
Thus the displayed assignment is a presheaf.  Suppose
\(V=\bigcup_\alpha V_\alpha\) and compatible sections
\(s_\alpha\in\mathcal F(V_\alpha)\) all have support in \(Z\cap V_\alpha\).
They have a unique gluing \(s\in\mathcal F(V)\).  If \(x\in V\setminus Z\),
choose \(\alpha\) with \(x\in V_\alpha\).  Then
\((s_\alpha)_x=0\), so \(s_x=0\).  Hence
\(\operatorname{Supp}(s)\subseteq Z\cap V\), and the gluing belongs to the
presheaf.  Uniqueness is inherited from \(\mathcal F\), proving the sheaf
axiom.

(b) The natural morphism
\(\mathcal F\to j_*(\mathcal F|_U)\) is, on an open \(V\), the restriction
\[
\mathcal F(V)\longrightarrow\mathcal F(V\cap U).
\]
Its kernel consists of sections whose restriction to \(V\cap U\) is zero.
This is equivalent to the vanishing of every germ at every point of
\(V\setminus Z\), and hence to
\(\operatorname{Supp}(s)\subseteq Z\cap V\).  The sectionwise kernel is
therefore \(\mathcal H_Z^0(\mathcal F)\), which proves exactness.

If \(\mathcal F\) is flasque, each displayed restriction map is surjective
because \(V\cap U\subseteq V\) are open.  Thus the sheaf morphism
\(\mathcal F\to j_*(\mathcal F|_U)\) is surjective even on sections over
every open set, and in particular is an epimorphism of sheaves.
\end{solution}

\begin{exerciseintro}
% record: H-II-01-021
\exerciseheading{Exercise II.1.21}
\addcontentsline{toc}{subsection}{Exercise II.1.21}

\begin{context}
Let \(X\) be a variety over an algebraically closed field \(k\), equipped
with its sheaf \(\mathcal O_X\) of regular functions.  Affine open subsets
form a basis, and if \(Y\) is closed in an affine open \(W\), restriction
identifies \(k[Y]\) with \(k[W]/I(Y)\).  A sequence of sheaves of abelian
groups is exact if and only if it is exact on stalks.  On
\(\mathbb P_k^1\), every global regular function is constant, every rational
function has only finitely many poles, and
\(\Gamma(\mathbb P_k^1,\mathcal O)=k\).
Sections on a noetherian space commute with filtered colimits (II.1.11).
\end{context}

\begin{problem}
Let \(X\) be a variety over \(k\).

(a) For a closed subset \(Y\subseteq X\), define
\[
\mathcal I_Y(U)=
\{f\in\mathcal O_X(U):f(y)=0\text{ for every }y\in Y\cap U\}.
\]
Show that \(\mathcal I_Y\) is a sheaf of ideals in \(\mathcal O_X\).

(b) If \(Y\) is a closed reduced algebraic subset and \(i:Y\to X\) is the inclusion,
prove
\[
\mathcal O_X/\mathcal I_Y\cong i_*\mathcal O_Y.
\]

(c) Take \(X=\mathbb P_k^1\) and \(Y=\{P,Q\}\) for distinct points.  With
\(\mathcal S=i_{P*}\mathcal O_P\oplus i_{Q*}\mathcal O_Q\), show that
\[
0\longrightarrow\mathcal I_Y\longrightarrow\mathcal O_X
\longrightarrow\mathcal S\longrightarrow0
\]
is exact, but that the induced map
\(\Gamma(X,\mathcal O_X)\to\Gamma(X,\mathcal S)\) is not surjective.

(d) Still on \(X=\mathbb P_k^1\), let \(K=k(X)\) and let \(\mathcal K\) be
the constant sheaf associated to the additive group of \(K\).  Construct a
natural injection \(\mathcal O_X\to\mathcal K\).  For
\(I_P=K/\mathcal O_{X,P}\), prove
\[
\mathcal K/\mathcal O_X
\cong\bigoplus_{P\in X} i_P(I_P).
\]

(e) Prove that in part (d) the sequence on global sections
\[
0\longrightarrow\Gamma(X,\mathcal O_X)
\longrightarrow\Gamma(X,\mathcal K)
\longrightarrow\Gamma(X,\mathcal K/\mathcal O_X)
\longrightarrow0
\]
is exact.
\end{problem}
\end{exerciseintro}

\begin{solution}
(a) If \(V\subseteq U\) and \(f\) vanishes on \(Y\cap U\), then \(f|_V\)
vanishes on \(Y\cap V\), so the assignment is a presheaf of ideals.  Now
let \(U=\bigcup_\alpha U_\alpha\) and let compatible
\(f_\alpha\in\mathcal I_Y(U_\alpha)\) be given.  They glue uniquely in
\(\mathcal O_X\) to \(f\in\mathcal O_X(U)\).  If \(y\in Y\cap U\), choose
\(\alpha\) with \(y\in U_\alpha\); then
\(f(y)=f_\alpha(y)=0\).  Hence \(f\in\mathcal I_Y(U)\).  This proves the
sheaf axiom, and the ideal property is immediate sectionwise.

(b) Restriction of regular functions defines
\[
\rho:\mathcal O_X\longrightarrow i_*\mathcal O_Y.
\]
Its sectionwise kernel is \(\mathcal I_Y\).  To check local surjectivity,
take an affine open neighborhood \(W\subseteq X\).  Then \(Y\cap W\) is
closed in \(W\), and the coordinate-ring description gives a surjection
\[
\mathcal O_X(W)=k[W]\longrightarrow
k[Y\cap W]=\mathcal O_Y(Y\cap W)
\]
with kernel \(I(Y\cap W)=\mathcal I_Y(W)\).  Since affine opens form a
basis, \(\rho\) is surjective on every stalk.  Thus
\[
0\longrightarrow\mathcal I_Y\longrightarrow\mathcal O_X
\longrightarrow i_*\mathcal O_Y\longrightarrow0
\]
is exact, and taking the quotient gives the claimed isomorphism.

(c) The reduced algebraic set \(Y\) is the disjoint union of its two reduced points, so
\(i_*\mathcal O_Y\cong
i_{P*}\mathcal O_P\oplus i_{Q*}\mathcal O_Q=\mathcal S\).
Part (b) gives the asserted exact sequence.  On global sections,
\[
\Gamma(X,\mathcal O_X)=k,\qquad
\Gamma(X,\mathcal S)=k\oplus k,
\]
and the map is \(c\mapsto(c,c)\).  For example, \((0,1)\) is not in its
image, so the map is not surjective.  This exhibits the failure of the
global-section functor to preserve epimorphisms.

(d) Every regular function on a nonempty open subset of the irreducible
curve \(X\) determines a rational function.  These inclusions commute with
restriction and are injective, giving
\(\mathcal O_X\hookrightarrow\mathcal K\).

For a nonempty open \(V\) and \(f\in\mathcal K(V)=K\), send \(f\) to its
family of principal parts
\[
\bigl([f]_P\bigr)_{P\in V},
\qquad [f]_P\in I_P=K/\mathcal O_{X,P}.
\]
Only the finitely many poles of \(f\) give nonzero entries, so this defines
a section of \(\bigoplus_P i_P(I_P)\).  The construction commutes with
restriction and gives a sheaf morphism
\[
\delta:\mathcal K\longrightarrow\bigoplus_{P\in X}i_P(I_P).
\]
At a point \(P\), its map on stalks is the quotient map
\[
K\longrightarrow K/\mathcal O_{X,P}=I_P:
\]
all other skyscraper summands have zero stalk there.  This map is
surjective with kernel \(\mathcal O_{X,P}\).  The resulting sequence is
therefore exact on every stalk:
\[
0\longrightarrow\mathcal O_X\longrightarrow\mathcal K
\xrightarrow{\delta}\bigoplus_{P\in X}i_P(I_P)
\longrightarrow0.
\]
Taking its cokernel identification proves the assertion.

(e) Sections on a noetherian space commute with filtered colimits
by Exercise II.1.11.  Applying this to the finite subsums gives
\[
\Gamma(X,\mathcal K)=K,\qquad
\Gamma\!\left(X,\bigoplus_P i_P(I_P)\right)=\bigoplus_P I_P.
\]
It remains only to prove that \(K\to\bigoplus_P I_P\) is surjective.  First
fix \(P=\infty\), choose the affine coordinate \(t\) on
\(X\setminus\{\infty\}\), and take a class in \(K/\mathcal O_{X,\infty}\)
represented by \(a(t)/b(t)\).  Euclidean division gives
\[
\frac{a(t)}{b(t)}=q(t)+\frac{r(t)}{b(t)},
\qquad \deg r<\deg b.
\]
The second term is regular at \(\infty\), so the prescribed class is
represented by the polynomial \(q(t)\).  This polynomial is regular at
every finite point and therefore has no principal part except possibly at
\(\infty\).

For an arbitrary \(P\), apply the same argument after a projective linear
change of coordinate carrying \(P\) to \(\infty\).  Thus every element of
the single summand \(I_P\) is the principal-part family of a rational
function with no poles away from \(P\).  An element of
\(\bigoplus_P I_P\) has only finitely many nonzero components; summing the
corresponding rational functions realizes that element.  Hence the final
map is surjective.  Its kernel is
\(\Gamma(X,\mathcal O_X)=k\) by the exact sequence of part (d), completing
the proof.
\end{solution}

\begin{exerciseintro}
% record: H-II-01-022
\exerciseheading{Exercise II.1.22}
\addcontentsline{toc}{subsection}{Exercise II.1.22}

\begin{context}
Compatible local sections of a sheaf glue uniquely.  An isomorphism of
sheaves may be checked locally on an open cover.  Uniqueness of an object
constructed from gluing data means uniqueness up to a unique isomorphism
that respects the specified local identifications.
\end{context}

\begin{problem}
Let \(X=\bigcup_iU_i\) be an open cover.  Suppose a sheaf
\(\mathcal F_i\) is given on each \(U_i\), together with isomorphisms
\[
\varphi_{ij}:\mathcal F_i|_{U_i\cap U_j}
\longrightarrow\mathcal F_j|_{U_i\cap U_j}
\]
such that \(\varphi_{ii}\) is the identity and
\[
\varphi_{ik}=\varphi_{jk}\circ\varphi_{ij}
\]
on every triple overlap.  Prove that the \(\mathcal F_i\) glue to a sheaf
\(\mathcal F\) on \(X\): there are isomorphisms
\(\psi_i:\mathcal F|_{U_i}\to\mathcal F_i\) satisfying
\(\psi_j=\varphi_{ij}\circ\psi_i\) on overlaps.  Prove the corresponding
uniqueness statement.
\end{problem}
\end{exerciseintro}

\begin{solution}
For an open set \(V\subseteq X\), define \(\mathcal F(V)\) to be the group
of families
\[
(s_i)_i,\qquad s_i\in\mathcal F_i(V\cap U_i),
\]
that satisfy
\[
\varphi_{ij}\bigl(s_i|_{V\cap U_i\cap U_j}\bigr)
=s_j|_{V\cap U_i\cap U_j}
\]
for every \(i,j\).  Restriction is componentwise.

This is a sheaf.  Indeed, for a cover \(V=\bigcup_\alpha V_\alpha\), a
compatible family of elements of \(\mathcal F(V_\alpha)\) can be glued in
each component \(\mathcal F_i\), producing
\(s_i\in\mathcal F_i(V\cap U_i)\).  The transition equalities hold locally
on the cover and hence globally by uniqueness in \(\mathcal F_j\).
Componentwise uniqueness proves uniqueness of the resulting family.

For \(V\subseteq U_i\), projection to the \(i\)-th component defines
\[
\psi_{i,V}:\mathcal F(V)\longrightarrow\mathcal F_i(V).
\]
To construct its inverse, start with \(s\in\mathcal F_i(V)\) and, for each
\(j\), set
\[
s_j=\varphi_{ij}
\bigl(s|_{V\cap U_j}\bigr)
\in\mathcal F_j(V\cap U_j).
\]
The cocycle identity implies that the family \((s_j)_j\) satisfies all
compatibility relations.  Since \(\varphi_{ii}\) is the identity, this
construction is inverse to projection.  It also commutes with restriction,
so \(\psi_i:\mathcal F|_{U_i}\to\mathcal F_i\) is an isomorphism.  The
defining compatibility relation gives
\(\psi_j=\varphi_{ij}\circ\psi_i\) on \(U_i\cap U_j\).

Finally, let \(\mathcal G\) be another sheaf with compatible
isomorphisms \(\theta_i:\mathcal G|_{U_i}\to\mathcal F_i\).  For
\(s\in\mathcal G(V)\), the family
\[
\bigl(\theta_i(s|_{V\cap U_i})\bigr)_i
\]
is compatible with the \(\varphi_{ij}\), and hence defines an element of
\(\mathcal F(V)\).  These maps form a sheaf morphism
\(\Theta:\mathcal G\to\mathcal F\).  On each \(U_i\) it is
\(\psi_i^{-1}\circ\theta_i\), so it is an isomorphism.  Any morphism
respecting the local identifications has these same restrictions on the
cover \(\{U_i\}\), and is therefore equal to \(\Theta\).  This proves
uniqueness up to a unique compatible isomorphism.
\end{solution}

\begin{exerciseintro}
\corpussection{II.2 Schemes}
\addcontentsline{toc}{section}{II.2 Schemes}

% record: H-II-02-001
\exerciseheading{Exercise II.2.1}
\addcontentsline{toc}{subsection}{Exercise II.2.1}

\begin{context}
For a ring \(A\) and \(f\in A\), write
\[
D(f)=\{\mathfrak p\in\operatorname{Spec}A:f\notin\mathfrak p\}.
\]
Prime ideals of \(A_f\) correspond to prime ideals of \(A\) disjoint from
\(\{1,f,f^2,\ldots\}\).  Principal open subsets form a basis of
\(\operatorname{Spec}A\), and
\(\mathcal O_{\operatorname{Spec}A}(D(g))=A_g\).
\end{context}

\begin{problem}
Let \(X=\operatorname{Spec}A\).  Prove that the locally ringed space
\[
\bigl(D(f),\mathcal O_X|_{D(f)}\bigr)
\]
is naturally isomorphic to \(\operatorname{Spec}A_f\).
\end{problem}
\end{exerciseintro}

\begin{solution}
The localization map \(A\to A_f\) induces a continuous map
\[
\theta:\operatorname{Spec}A_f\longrightarrow\operatorname{Spec}A,
\qquad
\mathfrak q\longmapsto\{a\in A:a/1\in\mathfrak q\}.
\]
Extension and contraction give inverse bijections between the prime ideals
of \(A_f\) and the prime ideals of \(A\) not containing \(f\), so the image
of \(\theta\) is \(D(f)\).

For \(g\in A\), the preimage of \(D(g)\cap D(f)\) is
\(D(g/1)\subseteq\operatorname{Spec}A_f\).  Conversely these opens form a
basis on \(\operatorname{Spec}A_f\).  Thus \(\theta\) is a homeomorphism
onto \(D(f)\).

It remains to compare the structure sheaves.  On the corresponding basic
opens there are natural isomorphisms
\[
\mathcal O_X(D(f)\cap D(g))
=A_{fg}
\cong (A_f)_{g/1}
=\mathcal O_{\operatorname{Spec}A_f}(D(g/1)).
\]
They are the canonical localization isomorphisms, so they commute with all
restriction maps between basic opens.  Hence they glue to an isomorphism
between \(\mathcal O_X|_{D(f)}\) and the structure sheaf of
\(\operatorname{Spec}A_f\).  On a prime \(\mathfrak p\in D(f)\), the
induced stalk map is the natural isomorphism
\[
A_{\mathfrak p}\cong
(A_f)_{\mathfrak p A_f},
\]
which is local.  Therefore the homeomorphism with this sheaf isomorphism is
an isomorphism of locally ringed spaces.
\end{solution}

\begin{exerciseintro}
% record: H-II-02-002
\exerciseheading{Exercise II.2.2}
\addcontentsline{toc}{subsection}{Exercise II.2.2}

\begin{context}
A locally ringed space is a scheme if every point has an open neighborhood
isomorphic to an affine scheme.  Principal opens \(D(f)\) form a basis of
an affine scheme \(\operatorname{Spec}A\), and with the restricted structure
sheaf each \(D(f)\) is isomorphic to \(\operatorname{Spec}A_f\).
\end{context}

\begin{problem}
Let \((X,\mathcal O_X)\) be a scheme and let \(U\subseteq X\) be open.
Prove that
\[
\bigl(U,\mathcal O_X|_U\bigr)
\]
is a scheme.
\end{problem}
\end{exerciseintro}

\begin{solution}
Fix \(x\in U\).  Since \(X\) is a scheme, there is an affine open
neighborhood \(V\) of \(x\), say \(V\cong\operatorname{Spec}A\).  The set
\(U\cap V\) is open in \(V\) and contains \(x\).  Because principal opens
form a basis, there is an \(f\in A\) such that
\[
x\in D(f)\subseteq U\cap V.
\]
With the structure sheaf restricted from \(X\), this neighborhood is the
same locally ringed space as \(D(f)\) inside
\(\operatorname{Spec}A\), and hence is isomorphic to
\(\operatorname{Spec}A_f\).

Thus every point of \(U\) has an affine open neighborhood for the restricted
sheaf.  This proves that \((U,\mathcal O_X|_U)\) is a scheme, called the
open subscheme induced by \(X\).
\end{solution}

\begin{exerciseintro}
% record: H-II-02-003
\exerciseheading{Exercise II.2.3}
\addcontentsline{toc}{subsection}{Exercise II.2.3}

\begin{context}
For a ring \(A\), write
\(A_{\mathrm{red}}=A/\sqrt{(0)}\).  The nilradical is contained in every
prime ideal, localization commutes with passage to the reduced quotient,
and \(\operatorname{Spec}A_{\mathrm{red}}\to\operatorname{Spec}A\) is a
homeomorphism.  Sheafification preserves stalks, and a section of a sheaf
is zero if all of its germs are zero.
\end{context}

\begin{problem}
A scheme \(X\) is called reduced if \(\mathcal O_X(U)\) is a reduced ring
for every open \(U\subseteq X\).

(a) Prove that \(X\) is reduced if and only if every local ring
\(\mathcal O_{X,x}\) is reduced.

(b) Let \((\mathcal O_X)_{\mathrm{red}}\) be the sheafification of
\[
U\longmapsto\mathcal O_X(U)_{\mathrm{red}}.
\]
Prove that \(\bigl(X,(\mathcal O_X)_{\mathrm{red}}\bigr)\) is a scheme,
denoted \(X_{\mathrm{red}}\), and that the natural morphism
\(X_{\mathrm{red}}\to X\) is a homeomorphism on underlying spaces.

(c) If \(X\) is reduced, prove that every morphism \(f:X\to Y\) factors
uniquely through \(Y_{\mathrm{red}}\to Y\).
\end{problem}
\end{exerciseintro}

\begin{solution}
(a) Suppose first that \(X\) is reduced.  Given \(x\in X\), choose an affine
neighborhood \(\operatorname{Spec}A\).  The ring
\(A=\mathcal O_X(\operatorname{Spec}A)\) is reduced, and every localization
\(A_{\mathfrak p}\) is reduced: if
\((a/s)^n=0\), then some \(t\notin\mathfrak p\) satisfies
\(ta^n=0\), so \((ta)^n=0\), whence \(ta=0\) and \(a/s=0\).
Thus \(\mathcal O_{X,x}\cong A_{\mathfrak p}\) is reduced.

Conversely, suppose every local ring is reduced.  If
\(s\in\mathcal O_X(U)\) satisfies \(s^n=0\), then
\((s_x)^n=0\) in \(\mathcal O_{X,x}\) for every \(x\in U\).  Hence
\(s_x=0\) for every \(x\), and the sheaf axiom gives \(s=0\).  Therefore
every \(\mathcal O_X(U)\) is reduced.

(b) The assertion is local on \(X\).  Let
\(V=\operatorname{Spec}A\) be an affine open.  On a principal open
\(D(f)\subseteq V\),
\[
\mathcal O_X(D(f))_{\mathrm{red}}
=(A_f)_{\mathrm{red}}
\cong(A_{\mathrm{red}})_f.
\]
These isomorphisms commute with restriction.  Since the principal opens
form a basis, after sheafification they identify
\[
\bigl(V,(\mathcal O_X)_{\mathrm{red}}|_V\bigr)
\cong\operatorname{Spec}A_{\mathrm{red}}.
\]
Thus the indicated locally ringed space is covered by affine schemes and is
a scheme.

The quotient homomorphisms \(A\to A_{\mathrm{red}}\) induce, on these
charts, morphisms
\(\operatorname{Spec}A_{\mathrm{red}}\to\operatorname{Spec}A\).
They are compatible on overlaps and therefore glue to
\(X_{\mathrm{red}}\to X\).  Every prime of \(A\) contains its nilradical,
so contraction gives a homeomorphism
\(\operatorname{Spec}A_{\mathrm{red}}\cong\operatorname{Spec}A\).
These local homeomorphisms are the identity on the common underlying point
set, and hence the global morphism is a homeomorphism on underlying spaces.

(c) Use the homeomorphism \(Y_{\mathrm{red}}\to Y\) to take the underlying
continuous map of the desired factorization to be the map underlying \(f\).
For every open \(V\subseteq Y\), the homomorphism
\[
f^\#_V:\mathcal O_Y(V)\longrightarrow
\mathcal O_X(f^{-1}V)
\]
kills every nilpotent because its target is reduced by part (a) and the
definition of a reduced scheme.  It therefore factors uniquely through
\(\mathcal O_Y(V)_{\mathrm{red}}\).  These factorizations commute with
restriction, so by the universal property of sheafification they give a
unique sheaf morphism
\[
(\mathcal O_Y)_{\mathrm{red}}\longrightarrow f_*\mathcal O_X.
\]

At \(x\in X\), with \(y=f(x)\), the resulting local-ring homomorphism is
\[
\mathcal O_{Y,y}/\sqrt{(0)}
\longrightarrow\mathcal O_{X,x}.
\]
It is local: the inverse image of the maximal ideal of
\(\mathcal O_{X,x}\) is the maximal ideal of
\(\mathcal O_{Y,y}\) modulo its nilradical, because the original composite
from \(\mathcal O_{Y,y}\) was local.  Hence the factorization is a morphism
of schemes.  Uniqueness follows from the uniqueness of every quotient
factorization and of sheafification.
\end{solution}

\begin{exerciseintro}
% record: H-II-02-004
\exerciseheading{Exercise II.2.4}
\addcontentsline{toc}{subsection}{Exercise II.2.4}

\begin{context}
A ring homomorphism \(A\to B\) determines a unique morphism
\(\operatorname{Spec}B\to\operatorname{Spec}A\), and this gives a
contravariant bijection for affine schemes.  Morphisms of schemes whose
restrictions agree on an open cover of the source glue uniquely.
\end{context}

\begin{problem}
Let \(A\) be a ring and \(X\) a scheme.  Sending a morphism
\(f:X\to\operatorname{Spec}A\) to its map on global sections gives
\[
\alpha:
\operatorname{Hom}_{\mathrm{Sch}}(X,\operatorname{Spec}A)
\longrightarrow
\operatorname{Hom}_{\mathrm{Ring}}
\bigl(A,\Gamma(X,\mathcal O_X)\bigr).
\]
Prove that \(\alpha\) is a bijection.
\end{problem}
\end{exerciseintro}

\begin{solution}
Let
\(\varphi:A\to\Gamma(X,\mathcal O_X)\) be a ring homomorphism.  Choose an
affine open cover \(X=\bigcup_iU_i\), with
\(U_i=\operatorname{Spec}B_i\).  Restriction of global sections gives
\[
A\xrightarrow{\varphi}\Gamma(X,\mathcal O_X)
\longrightarrow\Gamma(U_i,\mathcal O_X)=B_i,
\]
and hence a morphism
\(f_i:U_i\to\operatorname{Spec}A\).

The morphisms \(f_i\) agree on overlaps.  Indeed, cover
\(U_i\cap U_j\) by affine opens \(W=\operatorname{Spec}C\).  Both
restrictions to \(W\) correspond to the same composite
\[
A\xrightarrow{\varphi}\Gamma(X,\mathcal O_X)
\longrightarrow\Gamma(W,\mathcal O_X)=C.
\]
The affine correspondence makes the two restrictions equal on every such
\(W\), hence on \(U_i\cap U_j\).  The \(f_i\) therefore glue uniquely to a
morphism \(f_\varphi:X\to\operatorname{Spec}A\).

The map on global sections induced by \(f_\varphi\) is \(\varphi\): after
restriction to every \(U_i\), both maps are the displayed composite, and
sections agreeing on an open cover are equal.  Thus \(\alpha\) is
surjective.

If \(f,g:X\to\operatorname{Spec}A\) induce the same homomorphism on global
sections, then on each affine \(U_i=\operatorname{Spec}B_i\) their
restrictions correspond to the same ring map
\[
A\longrightarrow\Gamma(X,\mathcal O_X)\longrightarrow B_i.
\]
Hence \(f|_{U_i}=g|_{U_i}\) for every \(i\), and gluing uniqueness gives
\(f=g\).  Therefore \(\alpha\) is injective as well.
\end{solution}

\begin{exerciseintro}
% record: H-II-02-005
\exerciseheading{Exercise II.2.5}
\addcontentsline{toc}{subsection}{Exercise II.2.5}

\begin{context}
The prime ideals of \(\mathbb Z\) are \((0)\) and \((p)\) for positive
prime numbers \(p\).  A morphism from a scheme \(X\) to
\(\operatorname{Spec}A\) is naturally equivalent to a unital ring
homomorphism \(A\to\Gamma(X,\mathcal O_X)\).
\end{context}

\begin{problem}
Describe \(\operatorname{Spec}\mathbb Z\) as a topological space, and prove
that it is a final object in the category of schemes.
\end{problem}
\end{exerciseintro}

\begin{solution}
The points are the generic point \((0)\) and one closed point \((p)\) for
each positive prime number \(p\).  If \(n\ne0\), then
\[
V((n))=\{(p):p\text{ divides }n\},
\]
which is finite; \(V((0))\) is the entire space.  Conversely, every finite
set of closed points is \(V((p_1\cdots p_r))\).  Thus the nonempty open sets
are exactly the subsets containing \((0)\) whose complements consist of
finitely many closed points.

For any scheme \(X\), the affine-target correspondence gives
\[
\operatorname{Hom}_{\mathrm{Sch}}
(X,\operatorname{Spec}\mathbb Z)
\cong
\operatorname{Hom}_{\mathrm{Ring}}
\bigl(\mathbb Z,\Gamma(X,\mathcal O_X)\bigr).
\]
There is exactly one unital homomorphism from \(\mathbb Z\) to any unital
ring, namely \(n\mapsto n\cdot1\).  Hence there is exactly one morphism
\(X\to\operatorname{Spec}\mathbb Z\), proving that
\(\operatorname{Spec}\mathbb Z\) is final.
\end{solution}

\begin{exerciseintro}
% record: H-II-02-006
\exerciseheading{Exercise II.2.6}
\addcontentsline{toc}{subsection}{Exercise II.2.6}

\begin{context}
Prime ideals are proper ideals.  Ring homomorphisms preserve the
multiplicative identity; the zero ring has \(0=1\).  There is one continuous
map from the empty space to any topological space, and local conditions on
stalks are vacuous when the source has no points.
\end{context}

\begin{problem}
Describe the spectrum of the zero ring and prove that it is an initial
object in the category of schemes.
\end{problem}
\end{exerciseintro}

\begin{solution}
The zero ring has no proper ideals, hence no prime ideals.  Therefore
\[
\operatorname{Spec}(0)=\varnothing.
\]
Its structure sheaf is the unique sheaf of rings on the empty space, and the
resulting locally ringed space is a scheme because the condition of having
affine neighborhoods is vacuous.

Let \(X\) be any scheme.  There is exactly one continuous map
\(f:\varnothing\to X\).  For every open \(V\subseteq X\),
\[
(f_*\mathcal O_{\varnothing})(V)
=\mathcal O_{\varnothing}(\varnothing)=0,
\]
and there is exactly one unital homomorphism
\(\mathcal O_X(V)\to0\).  These maps form the unique possible morphism of
sheaves
\(\mathcal O_X\to f_*\mathcal O_{\varnothing}\).  There are no source
stalks on which to check locality, so this is a morphism of locally ringed
spaces.  It is unique, and consequently
\(\operatorname{Spec}(0)\) is initial in the category of schemes.
\end{solution}

\begin{exerciseintro}
% record: H-II-02-007
\exerciseheading{Exercise II.2.7}
\addcontentsline{toc}{subsection}{Exercise II.2.7}

\begin{context}
For \(x\in X\), write
\[
k(x)=\mathcal O_{X,x}/\mathfrak m_x
\]
for its residue field.  The spectrum of a field \(K\) has one point and
local ring \(K\).  A homomorphism of local rings is local precisely when
the inverse image of the maximal ideal of the target is the maximal ideal
of the source.
\end{context}

\begin{problem}
Let \(X\) be a scheme and \(K\) a field.  Prove that giving a morphism
\(\operatorname{Spec}K\to X\) is equivalent to giving a point \(x\in X\)
together with an injective field homomorphism \(k(x)\to K\).
\end{problem}
\end{exerciseintro}

\begin{solution}
Let \(g:\operatorname{Spec}K\to X\) be a morphism and let \(x\) be the image
of the unique point.  The morphism on stalks is a local homomorphism
\[
g_x^\#:\mathcal O_{X,x}\longrightarrow K.
\]
The maximal ideal of \(K\) is zero, so locality says
\[
\ker(g_x^\#)=\mathfrak m_x.
\]
Consequently \(g_x^\#\) factors through an injective homomorphism
\[
k(x)=\mathcal O_{X,x}/\mathfrak m_x\longrightarrow K.
\]

Conversely, suppose \(x\in X\) and an injection
\(\iota:k(x)\to K\) are given.  The composite
\[
\mathcal O_{X,x}\longrightarrow k(x)\xrightarrow{\iota}K
\]
is a local homomorphism.  Define the underlying continuous map from the
one-point space \(\operatorname{Spec}K\) to send its point to \(x\).
For an open \(U\subseteq X\), map a section \(s\in\mathcal O_X(U)\) to its
germ \(s_x\) and then through the displayed homomorphism when \(x\in U\);
when \(x\notin U\), the target section ring over the empty inverse image is
the zero ring and the map is unique.  These maps define a morphism of
sheaves whose only map on source stalks is the specified local
homomorphism.  Hence they define a scheme morphism.

The two constructions recover both the image point and the stalk map, and
are inverse to one another.
\end{solution}

\begin{exerciseintro}
% record: H-II-02-008
\exerciseheading{Exercise II.2.8}
\addcontentsline{toc}{subsection}{Exercise II.2.8}

\begin{context}
For \(x\in X\), the Zariski tangent space is
\[
T_xX=\operatorname{Hom}_{k(x)}
\bigl(\mathfrak m_x/\mathfrak m_x^2,k(x)\bigr).
\]
The ring of dual numbers \(D=k[\epsilon]/(\epsilon^2)\) is local with
maximal ideal \((\epsilon)\), residue field \(k\), and one prime ideal.
A \(k\)-morphism induces homomorphisms of local \(k\)-algebras.
\end{context}

\begin{problem}
Let \(X\) be a scheme over a field \(k\).  Prove that \(k\)-morphisms
\[
\operatorname{Spec}k[\epsilon]/(\epsilon^2)\longrightarrow X
\]
are naturally equivalent to pairs consisting of a \(k\)-rational point
\(x\in X\) and a tangent vector in \(T_xX\).
\end{problem}
\end{exerciseintro}

\begin{solution}
Let \(h:\operatorname{Spec}D\to X\) be a \(k\)-morphism, and let \(x\) be
the image of its unique point.  Composing \(h\) with the closed point
\(\operatorname{Spec}k\to\operatorname{Spec}D\) gives a \(k\)-morphism
\(\operatorname{Spec}k\to X\) with image \(x\).  Equivalently, there is a
\(k\)-embedding \(k(x)\to k\).  Since the structural map embeds \(k\) into
\(k(x)\) and the composite is the identity on \(k\), this forces
\(k(x)=k\); thus \(x\) is \(k\)-rational.

The stalk map of \(h\) is a local \(k\)-algebra homomorphism
\[
\varphi:\mathcal O_{X,x}\longrightarrow D.
\]
Its reduction modulo \(\epsilon\) is the residue map
\(\mathcal O_{X,x}\to k\).  Therefore
\(\varphi(\mathfrak m_x)\subseteq(\epsilon)\).  For \(u\in\mathfrak m_x\)
there is a unique scalar \(\lambda(u)\in k\) with
\(\varphi(u)=\lambda(u)\epsilon\).  The map \(\lambda\) is \(k\)-linear,
and
\[
\varphi(uv)=\lambda(u)\lambda(v)\epsilon^2=0
\qquad(u,v\in\mathfrak m_x),
\]
so it vanishes on \(\mathfrak m_x^2\).  It consequently defines an element
\[
\overline\lambda\in
\operatorname{Hom}_k(\mathfrak m_x/\mathfrak m_x^2,k)=T_xX.
\]

Conversely, let \(x\) be \(k\)-rational and let
\(\overline\lambda\in T_xX\) be a tangent vector, viewed as a functional
\(\mathfrak m_x/\mathfrak m_x^2\to k\).  The structural copy of \(k\) in
\(\mathcal O_{X,x}\) splits the residue map, so every element is written
uniquely as \(c+u\), with \(c\in k\) and \(u\in\mathfrak m_x\).  Define
\[
\varphi(c+u)=c+\overline\lambda(u)\epsilon.
\]
Because \(\overline\lambda\) vanishes on \(\mathfrak m_x^2\), this formula
is multiplicative; it is plainly additive, preserves \(k\), and is local.
The correspondence between a point plus a local stalk homomorphism and a
morphism from the one-point scheme \(\operatorname{Spec}D\) now produces
the desired \(k\)-morphism.  The two constructions are inverse and natural.
\end{solution}

\begin{exerciseintro}
% record: H-II-02-009
\exerciseheading{Exercise II.2.9}
\addcontentsline{toc}{subsection}{Exercise II.2.9}

\begin{context}
A generic point of a closed subset \(Z\) is a point \(\zeta\) for which
\(\overline{\{\zeta\}}=Z\).  In \(\operatorname{Spec}A\), the closure of
the point \(\mathfrak p\) is \(V(\mathfrak p)\), and \(V(I)\) is
irreducible precisely when \(\sqrt I\) is prime.  Every nonempty open subset
of an irreducible space is dense.
\end{context}

\begin{problem}
Prove that every nonempty irreducible closed subset of a scheme has a
unique generic point.
\end{problem}
\end{exerciseintro}

\begin{solution}
Let \(Z\subseteq X\) be nonempty, irreducible, and closed.  Choose an affine
open \(U=\operatorname{Spec}A\) meeting \(Z\).  The subset \(Z\cap U\) is a
nonempty irreducible closed subset of \(U\), so
\[
Z\cap U=V(\mathfrak p)
\]
for a unique prime ideal \(\mathfrak p\subseteq A\).  Let \(\zeta\) be the
corresponding point.  Its closure inside \(U\) is \(Z\cap U\).

Because \(Z\cap U\) is a nonempty open subset of the irreducible space
\(Z\), it is dense in \(Z\).  The closure of \(\zeta\) in \(X\) is
contained in \(Z\), since \(Z\) is closed and contains \(\zeta\), but it
contains the closure of \(Z\cap U\), because it contains
\(\overline{\{\zeta\}}^{\,U}=Z\cap U\).  Therefore
\[
\overline{\{\zeta\}}^{\,X}=Z,
\]
so \(\zeta\) is a generic point.

For uniqueness, suppose \(\zeta'\) is another generic point of \(Z\).
Since \(\zeta\in\overline{\{\zeta'\}}\), every open neighborhood of
\(\zeta\), in particular \(U\), contains \(\zeta'\).  Inside
\(\operatorname{Spec}A\), the two points have the same closure
\(Z\cap U\).  Their corresponding prime ideals are therefore equal, since
the closure of a prime \(\mathfrak q\) is \(V(\mathfrak q)\).  Thus
\(\zeta'=\zeta\).
\end{solution}

\begin{exerciseintro}
% record: H-II-02-010
\exerciseheading{Exercise II.2.10}
\addcontentsline{toc}{subsection}{Exercise II.2.10}

\begin{context}
The ring \(\mathbb R[x]\) is a principal ideal domain.  Its nonconstant
monic irreducible polynomials are the linear polynomials \(x-a\), with
\(a\in\mathbb R\), and the quadratic polynomials
\((x-z)(x-\overline z)\), with \(z\in\mathbb C\setminus\mathbb R\), where
\(z\) and \(\overline z\) determine the same polynomial.
\end{context}

\begin{problem}
Describe \(\operatorname{Spec}\mathbb R[x]\), including its topology, and
compare its points and topology with \(\mathbb R\) and with \(\mathbb C\).
\end{problem}
\end{exerciseintro}

\begin{solution}
There is one nonclosed point, the generic point \((0)\).  Every other prime
is maximal and is generated by a monic irreducible polynomial.  Thus the
closed points are
\[
(x-a)\quad(a\in\mathbb R)
\]
and
\[
\bigl((x-z)(x-\overline z)\bigr)
\quad
(z\in\mathbb C\setminus\mathbb R),
\]
with the same point obtained from \(z\) and \(\overline z\).  The first
kind has residue field \(\mathbb R\); the second has residue field
\(\mathbb C\).

For a nonzero polynomial \(f\), the closed set \(V(f)\) consists of the
finitely many closed points corresponding to the distinct irreducible
factors of \(f\).  The only closed set containing the generic point is the
whole spectrum.  Hence every nonempty open set contains \((0)\) and all but
finitely many closed points, and every such cofinite subset containing
\((0)\) is open.

The usual real line accounts only for the closed points of the form
\((x-a)\); on this subset the induced Zariski topology is cofinite, not the
usual Euclidean topology.  All closed points together are naturally
identified with the quotient set
\(\mathbb C/(z\sim\overline z)\), while the spectrum has in addition the
generic point \((0)\).  Thus it is neither the usual real line nor the
usual complex line as a topological space.
\end{solution}

\begin{exerciseintro}
% record: H-II-02-011
\exerciseheading{Exercise II.2.11}
\addcontentsline{toc}{subsection}{Exercise II.2.11}

\begin{context}
For \(k=\mathbb F_p\), the ring \(k[x]\) is a principal ideal domain.  For
each \(n\ge1\), the polynomial \(x^{p^n}-x\) is the product, without
repetition, of all monic irreducible polynomials over \(k\) whose degrees
divide \(n\).  Every finite field of order \(p^d\) is isomorphic to
\(\mathbb F_{p^d}\).
We use the elementary Mobius inversion formula for sums over divisors of an integer.
\end{context}

\begin{problem}
Describe \(\operatorname{Spec}\mathbb F_p[x]\).  Determine the residue
field at each point and the number of points having each possible residue
field.
\end{problem}
\end{exerciseintro}

\begin{solution}
The spectrum has the generic point \((0)\), whose residue field is the
rational function field \(\mathbb F_p(x)\).  Its closed points are the
ideals \((f)\), where \(f\in\mathbb F_p[x]\) is monic and irreducible.  If
\(\deg f=d\), then
\[
k((f))=\mathbb F_p[x]/(f)\cong\mathbb F_{p^d}.
\]
As for every polynomial ring in one variable over a field, the nonempty
open sets contain the generic point and omit only finitely many closed
points.

Let \(N_p(d)\) denote the number of monic irreducible polynomials of degree
\(d\).  Factoring \(x^{p^n}-x\) as in the context and comparing degrees
gives
\[
p^n=\sum_{d\mid n}d\,N_p(d).
\]
Möbius inversion therefore yields
\[
N_p(d)=\frac1d\sum_{e\mid d}\mu(e)\,p^{d/e}.
\]
Consequently there are exactly \(N_p(d)\) closed points whose residue field
is isomorphic to \(\mathbb F_{p^d}\), and there is one generic point with
residue field \(\mathbb F_p(x)\).
\end{solution}

\begin{exerciseintro}
% record: H-II-02-012
\exerciseheading{Exercise II.2.12}
\addcontentsline{toc}{subsection}{Exercise II.2.12}

\begin{context}
Sheaves on an open cover glue along sheaf isomorphisms satisfying the
cocycle condition.  A locally ringed space covered by open subspaces that
are schemes is a scheme.  An isomorphism of schemes identifies both the
underlying open spaces and their structure sheaves by local
homomorphisms.
\end{context}

\begin{problem}
Let \(\{X_i\}\) be a family of schemes.  For \(i\ne j\), let
\(U_{ij}\subseteq X_i\) be open and let
\(\varphi_{ij}:U_{ij}\to U_{ji}\) be an isomorphism.  Assume
\(\varphi_{ji}=\varphi_{ij}^{-1}\), and on triple overlaps assume
\[
\varphi_{ij}(U_{ij}\cap U_{ik})=U_{ji}\cap U_{jk},
\qquad
\varphi_{ik}=\varphi_{jk}\circ\varphi_{ij}.
\]
Construct a scheme \(X\) and morphisms \(\psi_i:X_i\to X\) such that:
each \(\psi_i\) is an isomorphism onto an open subscheme; their images
cover \(X\); the intersection of the \(i\)-th and \(j\)-th images is the
image of \(U_{ij}\); and
\(\psi_i=\psi_j\circ\varphi_{ij}\) on \(U_{ij}\).  Explain the case in
which all the \(U_{ij}\) are empty.
\end{problem}
\end{exerciseintro}

\begin{solution}
Set \(U_{ii}=X_i\) and \(\varphi_{ii}\) equal to the identity.  On the
disjoint union of the underlying sets, declare
\[
x\in X_i\sim y\in X_j
\quad\Longleftrightarrow\quad
x\in U_{ij}\ \text{and}\ y=\varphi_{ij}(x).
\]
The inverse condition gives symmetry, and the triple-overlap condition
gives transitivity, so this is an equivalence relation.  It never identifies
two distinct points of the same \(X_i\).

Let \(X\) be the quotient set and let \(\psi_i:X_i\to X\) be the canonical
injections.  Give \(X\) the topology in which \(W\subseteq X\) is open if
and only if every \(\psi_i^{-1}(W)\) is open in \(X_i\).  Each
\(\psi_i(X_i)\) is open: its inverse image in \(X_j\) is \(U_{ji}\).
More generally, if \(V\subseteq X_i\) is open, then
\[
\psi_j^{-1}(\psi_i(V))
=\varphi_{ji}^{-1}(V\cap U_{ij}),
\]
which is open.  Hence \(\psi_i\) is an open continuous bijection from
\(X_i\) to \(\psi_i(X_i)\), and therefore a homeomorphism.  The equivalence
relation also gives
\[
\psi_i(X_i)\cap\psi_j(X_j)=\psi_i(U_{ij})
=\psi_j(U_{ji})
\]
and \(\psi_i=\psi_j\circ\varphi_{ij}\) on \(U_{ij}\).

Transport \(\mathcal O_{X_i}\) by \(\psi_i\) to a sheaf on the open set
\(\psi_i(X_i)\).  On the intersection of the \(i\)-th and \(j\)-th open
sets, the scheme isomorphism \(\varphi_{ij}\) gives an isomorphism of these
transported sheaves.  The assumed cocycle identity for the
\(\varphi_{ij}\) gives the cocycle identity for the sheaf isomorphisms.
The sheaf-gluing theorem therefore produces a sheaf \(\mathcal O_X\) whose
restriction to \(\psi_i(X_i)\) is identified with \(\mathcal O_{X_i}\).
These identifications make every \(\psi_i\) an isomorphism of locally
ringed spaces onto an open subspace.

The images \(\psi_i(X_i)\) cover \(X\), and each is a scheme.  Thus
\((X,\mathcal O_X)\) is a scheme and has all the required properties.  If
every \(U_{ij}\) is empty for \(i\ne j\), no distinct summands are
identified; the construction is the disjoint union
\(\coprod_iX_i\).
\end{solution}

\begin{exerciseintro}
% record: H-II-02-013
\exerciseheading{Exercise II.2.13}
\addcontentsline{toc}{subsection}{Exercise II.2.13}

\begin{context}
A topological space is noetherian if its open subsets satisfy the ascending
chain condition.  A space is quasi-compact if every open cover has a finite
subcover.  Principal opens form a basis of \(\operatorname{Spec}A\), and
\[
\bigcup_\alpha D(f_\alpha)=\operatorname{Spec}A
\quad\Longleftrightarrow\quad
(f_\alpha)_\alpha=A.
\]
Closed subsets of an affine spectrum are the sets \(V(I)\).
\end{context}

\begin{problem}
(a) Prove that a topological space is noetherian if and only if every open
subset is quasi-compact.

(b) Prove that every affine scheme is quasi-compact, but that its underlying
space need not be noetherian.

(c) Prove that \(\operatorname{Spec}A\) is noetherian when \(A\) is a
noetherian ring.

(d) Give an example of a nonnoetherian ring whose spectrum is nevertheless
a noetherian topological space.
\end{problem}
\end{exerciseintro}

\begin{solution}
(a) Suppose \(X\) is noetherian and let \(U=\bigcup_{\alpha}U_\alpha\) be
an open cover of an open subset \(U\).  The finite unions of the
\(U_\alpha\) form a family of open sets.  If no one of them equals \(U\),
successively adjoining a member that adds a new point produces a strictly
ascending chain of open subsets, contrary to noetherianity.  Hence finitely
many \(U_\alpha\) cover \(U\).

Conversely, suppose every open subset is quasi-compact and let
\[
U_1\subseteq U_2\subseteq\cdots
\]
be an ascending chain.  The open set \(U=\bigcup_nU_n\) is covered by the
\(U_n\), so finitely many suffice.  Since the cover is nested,
\(U=U_N\) for the largest index occurring, and the chain stabilizes.

(b) Let \(\operatorname{Spec}A=\bigcup_\alpha W_\alpha\) be an open cover.
Refining by principal opens gives a cover by \(D(f_\lambda)\).  Its
complement is
\[
V((f_\lambda)_\lambda)=\varnothing,
\]
so the ideal generated by all the \(f_\lambda\) is the unit ideal.  An
expression \(1=\sum_{\lambda\in F}a_\lambda f_\lambda\) uses only finitely
many of them, and the corresponding \(D(f_\lambda)\), hence finitely many
of the original \(W_\alpha\), cover the spectrum.

Affineness does not imply noetherianity of the underlying space.  For
\[
A=k[x_1,x_2,x_3,\ldots],
\]
the closed subsets
\[
V(x_1)\supsetneq V(x_1,x_2)\supsetneq
V(x_1,x_2,x_3)\supsetneq\cdots
\]
form a strictly descending chain.  Strictness follows because the prime
ideal \((x_1,\ldots,x_n)\) lies in the \(n\)-th set but not the next.

(c) A descending chain \(V(I_1)\supseteq V(I_2)\supseteq\cdots\) may be
written using radical ideals, in which case it corresponds to an ascending
chain
\[
\sqrt{I_1}\subseteq\sqrt{I_2}\subseteq\cdots.
\]
If \(A\) is noetherian, this chain stabilizes, so the chain of closed sets
stabilizes.  Equivalently, the open subsets satisfy the ascending chain
condition.

(d) Let
\[
A=k[x_1,x_2,\ldots]\big/
(x_ix_j:i,j\ge1),
\]
and let \(\mathfrak m=(x_1,x_2,\ldots)\) denote the image of the indicated
ideal.  Then \(\mathfrak m^2=0\), so every prime ideal contains
\(\mathfrak m\).  Since \(A/\mathfrak m\cong k\), the ideal
\(\mathfrak m\) itself is prime and is the only point of
\(\operatorname{Spec}A\).  The spectrum is therefore noetherian.  But
\[
(x_1)\subsetneq(x_1,x_2)\subsetneq\cdots
\]
is a strictly ascending chain of ideals, so \(A\) is not noetherian.
\end{solution}

\begin{exerciseintro}
% record: H-II-02-014
\exerciseheading{Exercise II.2.14}
\addcontentsline{toc}{subsection}{Exercise II.2.14}

\begin{context}
For a nonnegatively graded ring \(S\), let
\(S_+=\bigoplus_{d>0}S_d\).  The standard opens
\[
D_+(f)=\{\mathfrak p\in\operatorname{Proj}S:f\notin\mathfrak p\},
\qquad f\in S_+\text{ homogeneous},
\]
form a basis, and
\(D_+(f)\cong\operatorname{Spec}(S_f)_0\).  A ring has empty spectrum if
and only if it is the zero ring.  Let \(t(V)\) denote the scheme associated
to a variety \(V\).
\end{context}

\begin{problem}
(a) Prove that \(\operatorname{Proj}S=\varnothing\) if and only if every
element of \(S_+\) is nilpotent.

(b) Let \(\varphi:S\to T\) be a degree-preserving homomorphism and set
\[
U=\{\mathfrak q\in\operatorname{Proj}T:
\mathfrak q\not\supseteq\varphi(S_+)\}.
\]
Show that \(U\) is open and that \(\varphi\) induces a natural morphism
\(U\to\operatorname{Proj}S\).

(c) Suppose there is an integer \(d_0\) such that
\(\varphi_d:S_d\to T_d\) is an isomorphism for every \(d\ge d_0\).  Prove
that \(U=\operatorname{Proj}T\) and that the morphism in part (b) is an
isomorphism.

(d) If \(V\) is a projective variety with homogeneous coordinate ring
\(S\), prove that \(t(V)\cong\operatorname{Proj}S\).
\end{problem}
\end{exerciseintro}

\begin{solution}
(a) If every element of \(S_+\) is nilpotent, every prime ideal contains
\(S_+\); hence no homogeneous prime qualifies as a point of
\(\operatorname{Proj}S\).

Conversely, suppose \(\operatorname{Proj}S\) is empty.  For every homogeneous
\(f\in S_+\), the standard open
\[
D_+(f)\cong\operatorname{Spec}(S_f)_0
\]
is empty, so \((S_f)_0\) is the zero ring.  In particular \(1=0\) in
\(S_f\), which says that some power of \(f\) is zero.  Every homogeneous
element of positive degree is therefore nilpotent.  The nilpotent elements
of a commutative ring form an ideal, so every finite sum of such homogeneous
elements, and hence every element of \(S_+\), is nilpotent.

(b) One has
\[
U=\bigcup_{\substack{f\in S_+\\ f\ \mathrm{homogeneous}}}
D_+(\varphi(f)),
\]
so \(U\) is open.  On each member of this cover, localization of
\(\varphi\) gives a ring homomorphism
\[
(S_f)_0\longrightarrow(T_{\varphi(f)})_0
\]
and thus a scheme morphism
\[
D_+(\varphi(f))\longrightarrow D_+(f).
\]
These morphisms agree on pairwise intersections because further
localization is functorial, so they glue to \(U\to\operatorname{Proj}S\).
On points it is
\(\mathfrak q\mapsto\varphi^{-1}(\mathfrak q)\); membership in \(U\)
ensures that this homogeneous prime does not contain \(S_+\).

(c) If \(\mathfrak q\in\operatorname{Proj}T\) were outside \(U\), it would
contain \(\varphi(S_+)\), and therefore \(T_d\) for every \(d\ge d_0\).
For any homogeneous \(t\in T_+\), a sufficiently high power \(t^n\) has
degree at least \(d_0\), so \(t^n\in\mathfrak q\), and primality gives
\(t\in\mathfrak q\).  This would imply
\(\mathfrak q\supseteq T_+\), a contradiction.  Thus
\(U=\operatorname{Proj}T\).

The opens \(D_+(f)\) with \(f\) homogeneous of positive degree and
\(\deg f\ge d_0\) cover \(\operatorname{Proj}S\): replace any homogeneous
element avoided by a given prime by a sufficiently high power.  For such
an \(f\), of degree \(d\), there are direct-limit descriptions
\[
(S_f)_0=\varinjlim
\bigl(S_0\xrightarrow{\,f\,}S_d\xrightarrow{\,f\,}S_{2d}
\xrightarrow{\,f\,}\cdots\bigr)
\]
and the analogous expression for \(T\).  The map between these systems is
an isomorphism from a sufficiently late term onward, so it induces
\[
(S_f)_0\cong(T_{\varphi(f)})_0.
\]
The morphism is consequently an isomorphism on an open cover, and hence is
an isomorphism globally.

(d) Embed \(V\) in \(\mathbb P_k^n\) and write
\[
S=k[x_0,\ldots,x_n]/I(V).
\]
On the standard affine chart \(x_i\ne0\), the coordinate ring of
\(V\cap\mathbb A_i^n\) is obtained by dehomogenizing at \(x_i\), namely
\[
k[V\cap\mathbb A_i^n]\cong(S_{x_i})_0.
\]
Therefore
\[
t(V\cap\mathbb A_i^n)
\cong\operatorname{Spec}(S_{x_i})_0
=D_+(x_i)\subseteq\operatorname{Proj}S.
\]
On overlaps these identifications are the same localization
identifications used in the constructions of both \(t(V)\) and
\(\operatorname{Proj}S\).  They consequently glue, and the charts cover
both schemes.  This yields \(t(V)\cong\operatorname{Proj}S\).
\end{solution}

\begin{exerciseintro}
% record: H-II-02-015
\exerciseheading{Exercise II.2.15}
\addcontentsline{toc}{subsection}{Exercise II.2.15}

\begin{context}
Let \(k\) be algebraically closed and let \(t(V)\) be the scheme associated
to a \(k\)-variety \(V\).  Affine charts of \(t(V)\) are spectra of
finitely generated reduced \(k\)-algebras.  By the weak Nullstellensatz,
the residue field at every maximal ideal of such an algebra is \(k\).
Every nonempty closed subset of a finite-type \(k\)-scheme has a closed
point.
\end{context}

\begin{problem}
(a) Prove that a point \(P\in t(V)\) is closed if and only if
\(k(P)=k\).

(b) If \(f:X\to Y\) is a morphism of schemes over \(k\) and
\(k(P)=k\), prove that \(k(f(P))=k\).

(c) For \(k\)-varieties \(V,W\), prove that the natural map
\[
\operatorname{Hom}_{\mathrm{Var}}(V,W)
\longrightarrow
\operatorname{Hom}_{\mathrm{Sch}/k}(t(V),t(W))
\]
is a bijection.
\end{problem}
\end{exerciseintro}

\begin{solution}
(a) If \(P\) is closed, choose an affine neighborhood
\(\operatorname{Spec}A\) of \(P\).  Its corresponding prime
\(\mathfrak p\) is maximal, and the weak Nullstellensatz gives
\[
k(P)=A/\mathfrak p=k.
\]

Conversely, suppose \(k(P)=k\), and let \(Q\) lie in the closure of
\(\{P\}\).  Any affine open neighborhood
\(\operatorname{Spec}B\) of \(Q\) also contains \(P\).  If
\(\mathfrak p\subseteq B\) represents \(P\), then
\[
k(P)=\operatorname{Frac}(B/\mathfrak p)=k.
\]
Since \(B/\mathfrak p\) is a \(k\)-subalgebra of its fraction field \(k\)
and contains \(k\), it equals \(k\).  Thus \(\mathfrak p\) is maximal.
The specialization \(Q\) corresponds to a prime containing
\(\mathfrak p\), so it equals \(\mathfrak p\), and \(Q=P\).  The closure of
\(\{P\}\) is therefore \(\{P\}\).

(b) A morphism induces an injective homomorphism of residue fields
\[
k(f(P))\longrightarrow k(P)=k.
\]
Because \(f\) is a \(k\)-morphism, this is a homomorphism of field
extensions of \(k\), and its image contains the structural copy of \(k\).
It follows that \(k(f(P))=k\).

(c) The construction \(t\) plainly preserves composition.  A morphism of
varieties is determined by its maps on affine coordinate rings, so if two
such morphisms induce the same scheme morphism, they are equal.  Thus the
displayed map is injective.

For surjectivity, let \(F:t(V)\to t(W)\) be a \(k\)-morphism.  By parts (a)
and (b), \(F\) sends the closed points corresponding to \(V\) into those
corresponding to \(W\), and hence determines a set map \(g:V\to W\).
We verify that it is regular locally.  Around any closed point of \(t(V)\),
choose affine opens
\[
\operatorname{Spec}A\subseteq t(V),\qquad
\operatorname{Spec}B\subseteq t(W)
\]
such that \(F(\operatorname{Spec}A)\subseteq\operatorname{Spec}B\).
The restriction of \(F\) is induced by a \(k\)-algebra homomorphism
\(B\to A\).  On closed points this is exactly the regular map between the
corresponding affine varieties.  These local maps agree because they all
come from \(F\), so they glue to a morphism of varieties \(g:V\to W\).

The affine neighborhoods used above cover every closed point of \(t(V)\).
They in fact cover all of \(t(V)\): otherwise their open union would have
a nonempty closed complement, which would contain a closed point.  On each
of these affine opens, \(t(g)\) and \(F\) arise from the same ring
homomorphism \(B\to A\).  Hence they agree locally and therefore globally.
Thus every \(F\) comes from a unique morphism \(g\), proving bijectivity.
\end{solution}

\begin{exerciseintro}
% record: H-II-02-016
\exerciseheading{Exercise II.2.16}
\addcontentsline{toc}{subsection}{Exercise II.2.16}

\begin{context}
For \(f\in\Gamma(X,\mathcal O_X)\), define
\[
X_f=\{x\in X:f_x\notin\mathfrak m_x\}.
\]
On an affine scheme \(\operatorname{Spec}B\), a section \(g\in B\) is a
unit in \(B_{\mathfrak p}\) exactly when \(g\notin\mathfrak p\), and
\[
\Gamma(D(g),\mathcal O)=B_g.
\]
A section whose germs vanish everywhere is zero.
\end{context}

\begin{problem}
(a) If \(U=\operatorname{Spec}B\) is an affine open of \(X\) and
\(\bar f=f|_U\in B\), prove
\[
U\cap X_f=D(\bar f).
\]
Deduce that \(X_f\) is open.

(b) Assume \(X\) is quasi-compact.  Put
\(A=\Gamma(X,\mathcal O_X)\).  If \(a\in A\) restricts to zero on \(X_f\),
prove that \(f^na=0\) for some \(n>0\).

(c) Suppose \(X\) has a finite affine open cover \(\{U_i\}\) such that
every \(U_i\cap U_j\) is quasi-compact.  If
\(b\in\Gamma(X_f,\mathcal O_{X_f})\), prove that \(f^nb\) extends to a
section in \(A\) for some \(n>0\).

(d) Under the hypothesis of part (c), prove
\[
\Gamma(X_f,\mathcal O_{X_f})\cong A_f.
\]
\end{problem}
\end{exerciseintro}

\begin{solution}
(a) A point \(x\in U\), represented by
\(\mathfrak p\in\operatorname{Spec}B\), belongs to \(X_f\) precisely when
\(\bar f/1\) is a unit in \(B_{\mathfrak p}\).  This holds precisely when
\(\bar f\notin\mathfrak p\).  Hence
\[
U\cap X_f=D(\bar f).
\]
Affine opens cover \(X\), and the right side is open in \(U\), so \(X_f\)
is open.

(b) Choose a finite affine cover
\(X=\bigcup_{i=1}^rU_i\), possible because \(X\) is quasi-compact.  Write
\[
U_i=\operatorname{Spec}B_i,\qquad
f_i=f|_{U_i},\qquad a_i=a|_{U_i}.
\]
By part (a), \(U_i\cap X_f=D(f_i)\).  The hypothesis says that \(a_i\)
maps to zero in
\((B_i)_{f_i}\), so some \(n_i>0\) satisfies
\(f_i^{n_i}a_i=0\) in \(B_i\).  If \(n=\max_i n_i\), then
\((f^na)|_{U_i}=0\) for every \(i\).  The sheaf axiom gives \(f^na=0\)
on \(X\).

(c) Write \(U_i=\operatorname{Spec}B_i\).  On
\(U_i\cap X_f=D(f_i)\), the section \(b\) is represented by an element of
\((B_i)_{f_i}\).  Since there are finitely many \(i\), after increasing one
integer \(n\) we may write
\[
b|_{U_i\cap X_f}=\frac{a_i}{f_i^n}
\qquad\text{with }a_i\in B_i
\]
for every \(i\).

On \(W_{ij}=U_i\cap U_j\), the section
\[
d_{ij}=a_i|_{W_{ij}}-a_j|_{W_{ij}}
\]
restricts to zero on \((W_{ij})_f=W_{ij}\cap X_f\), because both fractions
represent \(b\).  The scheme \(W_{ij}\) is quasi-compact by hypothesis.
Applying part (b) to \(W_{ij}\), with \(f\) and \(d_{ij}\) restricted
there, gives an exponent \(m_{ij}\) for which
\[
f^{m_{ij}}d_{ij}=0\quad\text{on }W_{ij}.
\]
There are only finitely many pairs.  Taking \(m\) at least every
\(m_{ij}\), the sections \(f^ma_i\in\mathcal O_X(U_i)\) agree on all
overlaps.  They glue to \(a\in A\), and on \(X_f\) one has
\[
a|_{X_f}=f^{n+m}b.
\]
Thus a power of \(f\) times \(b\) extends globally.

(d) The section \(f|_{X_f}\) is invertible: it is invertible in every
stalk, so its local inverses glue.  Restriction \(A\to
\Gamma(X_f,\mathcal O_{X_f})\) therefore factors through a homomorphism
\[
\rho:A_f\longrightarrow\Gamma(X_f,\mathcal O_{X_f}).
\]
If \(\rho(a/f^r)=0\), then \(a|_{X_f}=0\).  Part (b), applicable because
the finite affine cover makes \(X\) quasi-compact, gives \(f^ma=0\) for
some \(m\); hence \(a/f^r=0\) in \(A_f\).  Thus \(\rho\) is injective.
For any section \(b\) of the target, part (c) gives
\(f^nb=a|_{X_f}\) for some \(a\in A\), so \(b=\rho(a/f^n)\).
Therefore \(\rho\) is an isomorphism.
\end{solution}

\begin{exerciseintro}
% record: H-II-02-017
\exerciseheading{Exercise II.2.17}
\addcontentsline{toc}{subsection}{Exercise II.2.17}

\begin{context}
Morphisms of schemes may be glued on an open cover of the source.  For
\(A=\Gamma(X,\mathcal O_X)\), the identity \(A\to A\) determines a
canonical morphism \(X\to\operatorname{Spec}A\).  If \(X\) has a finite
affine cover with quasi-compact pairwise intersections, then
\[
\Gamma(X_f,\mathcal O_X|_{X_f})\cong A_f.
\]
\end{context}

\begin{problem}
(a) Let \(h:X\to Y\) be a morphism.  Suppose \(Y\) has an open cover
\(\{U_i\}\) for which every induced morphism
\[
h^{-1}(U_i)\longrightarrow U_i
\]
is an isomorphism.  Prove that \(h\) is an isomorphism.

(b) Prove that a scheme \(X\) is affine if and only if there are finitely
many \(f_1,\ldots,f_r\in A=\Gamma(X,\mathcal O_X)\) such that every
\(X_{f_i}\) is affine and \((f_1,\ldots,f_r)=A\).
\end{problem}
\end{exerciseintro}

\begin{solution}
(a) Let
\(g_i:U_i\to h^{-1}(U_i)\subseteq X\) be the inverse of the displayed
isomorphism, followed by the open immersion into \(X\).  On
\(U_i\cap U_j\), both \(g_i\) and \(g_j\) are inverses to the restriction
of \(h\); uniqueness of an inverse makes them equal.  They glue to a
morphism \(g:Y\to X\).  The identities
\[
h\circ g=\operatorname{id}_Y,\qquad
g\circ h=\operatorname{id}_X
\]
hold on the respective open covers and hence globally.  Thus \(h\) is an
isomorphism.

(b) If \(X\) is affine, take \(r=1\) and \(f_1=1\).  Then
\(X_{f_1}=X\) and \(f_1\) generates the unit ideal.

Conversely, suppose such \(f_i\) are given.  Because they generate \(A\),
at every \(x\in X\) at least one germ \((f_i)_x\) is a unit; otherwise all
would belong to the maximal ideal of \(\mathcal O_{X,x}\), as would their
linear combination \(1\).  Hence the affine opens \(X_{f_i}\) cover \(X\).
Moreover,
\[
X_{f_i}\cap X_{f_j}=X_{f_if_j}
\]
is a principal open in the affine scheme \(X_{f_i}\), and is therefore
affine and quasi-compact.  The global-section localization result from the
context applies to this finite cover and gives
\[
\Gamma(X_{f_i},\mathcal O_X)=A_{f_i}.
\]

Let \(h:X\to\operatorname{Spec}A\) be the canonical morphism.  The
preimage of \(D(f_i)\) is \(X_{f_i}\).  On this preimage, \(h\) is the
affine morphism induced by the natural map
\[
A_{f_i}\longrightarrow
\Gamma(X_{f_i},\mathcal O_X)=A_{f_i},
\]
which is the identity under the displayed identification.  Therefore
\[
X_{f_i}\longrightarrow D(f_i)
\]
is an isomorphism.  Since the \(f_i\) generate the unit ideal, the
\(D(f_i)\) cover \(\operatorname{Spec}A\).  Part (a) now shows that
\(h\) is an isomorphism, so \(X\) is affine.
\end{solution}

\begin{exerciseintro}
% record: H-II-02-018
\exerciseheading{Exercise II.2.18}
\addcontentsline{toc}{subsection}{Exercise II.2.18}

\begin{context}
Let \(\varphi:A\to B\) induce
\(f:\operatorname{Spec}B\to\operatorname{Spec}A\).  Then
\[
f^{-1}(D(a))=D(\varphi(a)),
\]
and the map of structure sheaves on this basic open is
\[
A_a\longrightarrow B_{\varphi(a)}.
\]
A morphism of sheaves is injective or surjective if and only if its stalk
maps have that property.  The nilradical of a ring is the intersection of
all its prime ideals.
\end{context}

\begin{problem}
(a) For \(a\in A\), prove that \(a\) is nilpotent if and only if
\(D(a)=\varnothing\).

(b) Prove that \(\varphi\) is injective if and only if
\[
f^\#:\mathcal O_{\operatorname{Spec}A}
\longrightarrow f_*\mathcal O_{\operatorname{Spec}B}
\]
is injective.  If these conditions hold, prove that \(f\) is dominant.

(c) If \(\varphi\) is surjective, prove that \(f\) is a homeomorphism onto
a closed subset and that \(f^\#\) is surjective.

(d) Conversely, suppose \(f\) is a homeomorphism onto a closed subset and
\(f^\#\) is surjective.  Prove that \(\varphi\) is surjective.
\end{problem}
\end{exerciseintro}

\begin{solution}
(a) The set \(D(a)\) is empty precisely when every prime ideal contains
\(a\).  Since the intersection of all prime ideals is the nilradical, this
is equivalent to \(a\) being nilpotent.

(b) Suppose \(\varphi\) is injective.  For every \(a\in A\), the
localization map
\[
A_a\longrightarrow B_{\varphi(a)}
\]
is injective.  Indeed, if the image of \(c/a^n\) is zero, then
\(\varphi(a)^m\varphi(c)=0\) for some \(m\); injectivity gives
\(a^mc=0\), so \(c/a^n=0\).  Since the \(D(a)\) form a basis, \(f^\#\) is
injective.  Conversely, if \(f^\#\) is injective, its map on sections over
the whole affine scheme is
\[
A\longrightarrow B,
\]
namely \(\varphi\), so \(\varphi\) is injective.

To prove dominance under these conditions, let \(c\in A\) vanish at every
point of \(f(\operatorname{Spec}B)\).  Then \(\varphi(c)\) belongs to every
prime of \(B\), so it is nilpotent.  Injectivity implies that \(c\) is
nilpotent in \(A\), and hence it vanishes at every point of
\(\operatorname{Spec}A\).  Thus no proper closed subset contains the image,
and \(f\) is dominant.

(c) Write \(B=A/I\).  Extension and contraction identify
\(\operatorname{Spec}B\) homeomorphically with the closed subset
\[
V(I)\subseteq\operatorname{Spec}A.
\]
At a point \(\mathfrak p\in V(I)\), the relevant stalk map is the
surjection
\[
A_{\mathfrak p}\longrightarrow
(A/I)_{\mathfrak p/I}.
\]
At a point outside \(V(I)\), the stalk of \(f_*\mathcal O_{\operatorname{Spec}B}\)
is zero, because that point has a neighborhood disjoint from the image.
Every stalk map of \(f^\#\) is therefore surjective, so \(f^\#\) is
surjective.

(d) Put \(I=\ker\varphi\) and factor
\[
A\twoheadrightarrow A'=A/I\hookrightarrow B.
\]
Correspondingly, factor \(f\) as
\[
\operatorname{Spec}B\xrightarrow{g}\operatorname{Spec}A'
\xrightarrow{i}\operatorname{Spec}A.
\]
By part (b), \(g\) is dominant.  The closure of the image of \(f\) is
\(V(I)\); since the image is assumed closed, it equals \(V(I)\).  As
\(i\) is a homeomorphism from \(\operatorname{Spec}A'\) to \(V(I)\) and
\(f\) is a homeomorphism onto that subset, \(g\) is a homeomorphism.

The sheaf map for \(f\) factors as
\[
\mathcal O_{\operatorname{Spec}A}
\longrightarrow i_*\mathcal O_{\operatorname{Spec}A'}
\longrightarrow i_*g_*\mathcal O_{\operatorname{Spec}B}.
\]
The first arrow is surjective by part (c), and the composite is surjective
by hypothesis.  It follows stalkwise on \(V(I)\) that the second arrow is
surjective.  The injection \(A'\hookrightarrow B\), together with part
(b), says that this second arrow is also injective after identifying
\(\operatorname{Spec}A'\) with its image.  Thus
\[
\mathcal O_{\operatorname{Spec}A'}\cong
g_*\mathcal O_{\operatorname{Spec}B}.
\]
Taking global sections yields \(A'\cong B\), and the isomorphism is the
injection induced by \(\varphi\).  Hence \(A\to B\) is surjective.
\end{solution}

\begin{exerciseintro}
% record: H-II-02-019
\exerciseheading{Exercise II.2.19}
\addcontentsline{toc}{subsection}{Exercise II.2.19}

\begin{context}
For \(X=\operatorname{Spec}A\), one has
\(\Gamma(X,\mathcal O_X)=A\).  Sections on disjoint open subsets glue
uniquely.  Every prime ideal of a product \(A_1\times A_2\) is of the form
\(\mathfrak p_1\times A_2\) or \(A_1\times\mathfrak p_2\).
\end{context}

\begin{problem}
For a ring \(A\), prove that the following conditions are equivalent:

\(\mathrm{(i)}\) \(\operatorname{Spec}A\) is disconnected.

\(\mathrm{(ii)}\) There are nonzero \(e_1,e_2\in A\) such that
\[
e_1^2=e_1,\qquad e_2^2=e_2,\qquad
e_1e_2=0,\qquad e_1+e_2=1.
\]

\(\mathrm{(iii)}\) There are nonzero rings \(A_1,A_2\) with
\(A\cong A_1\times A_2\).
\end{problem}
\end{exerciseintro}

\begin{solution}
Assume \(\mathrm{(i)}\).  Write
\(X=\operatorname{Spec}A=U_1\coprod U_2\), where both \(U_i\) are nonempty
and open.  On \(U_1\) take the section \(1\) and on \(U_2\) take the
section \(0\).  They glue to
\(e_1\in\Gamma(X,\mathcal O_X)=A\).  Interchanging \(0\) and \(1\) gives
\(e_2\).  All four identities in \(\mathrm{(ii)}\) hold after restriction
to both members of the cover and hence hold globally.  Neither element is
zero because it restricts to \(1\) on a nonempty open set.  Thus
\(\mathrm{(i)}\) implies \(\mathrm{(ii)}\).

Assume \(\mathrm{(ii)}\).  The ideals \(A_1=e_1A\) and \(A_2=e_2A\) are
rings with identities \(e_1\) and \(e_2\), respectively, and both are
nonzero.  Define
\[
\Phi:A\longrightarrow A_1\times A_2,\qquad
a\longmapsto(e_1a,e_2a).
\]
Orthogonality makes \(\Phi\) a unital ring homomorphism.  Its inverse is
\[
(u,v)\longmapsto u+v:
\]
the identities \(e_1+e_2=1\) and \(e_1e_2=0\) verify both composites.
Hence \(\mathrm{(ii)}\) implies \(\mathrm{(iii)}\).

Finally, if \(A\cong A_1\times A_2\) with both factors nonzero, then
\[
\operatorname{Spec}A
\cong\operatorname{Spec}A_1\coprod\operatorname{Spec}A_2.
\]
Both summands are nonempty and are open and closed, so the spectrum is
disconnected.  Thus \(\mathrm{(iii)}\) implies \(\mathrm{(i)}\).
\end{solution}

\begin{exerciseintro}
\corpussection{II.3 First Properties of Schemes}
\addcontentsline{toc}{section}{II.3 First Properties of Schemes}

% record: H-II-03-001
\exerciseheading{Exercise II.3.1}
\addcontentsline{toc}{subsection}{Exercise II.3.1}

\begin{context}
A morphism \(f:X\to Y\) is locally of finite type if \(Y\) has an affine
open cover \(V_i=\operatorname{Spec}B_i\) such that every
\(f^{-1}(V_i)\) has an affine open cover
\(\operatorname{Spec}A_{ij}\) with each \(A_{ij}\) a finitely generated
\(B_i\)-algebra.  Localizing a finitely generated algebra at finitely many elements
preserves finite generation.  Base change to any localization of
the base also preserves finite generation over that localized base.  Principal
opens form a basis on an affine scheme.
\end{context}

\begin{problem}
Prove that \(f:X\to Y\) is locally of finite type if and only if, for every
affine open \(V=\operatorname{Spec}B\subseteq Y\), the open set
\(f^{-1}(V)\) can be covered by affine opens
\(U_j=\operatorname{Spec}A_j\) for which every \(A_j\) is a finitely
generated \(B\)-algebra.
\end{problem}
\end{exerciseintro}

\begin{solution}
The stated condition immediately implies the definition by choosing any
affine open cover of \(Y\).  Conversely, suppose the condition in the
definition holds for a fixed affine cover \(\{V_i\}\), and let
\(V=\operatorname{Spec}B\) be any affine open.

Fix \(x\in f^{-1}(V)\) and choose \(i\) with \(f(x)\in V_i\), together
with one of the prescribed affine neighborhoods
\[
U_0=\operatorname{Spec}A_0\subseteq f^{-1}(V_i)
\]
of \(x\), where \(A_0\) is of finite type over
\(B_i=\Gamma(V_i,\mathcal O_Y)\).  Inside \(V\cap V_i\), choose a
principal neighborhood
\[
W=D_V(g)
\]
of \(f(x)\).  It is affine, with coordinate ring \(B_g\), and its
inclusion in the affine scheme \(V_i\) corresponds to a homomorphism
\(B_i\to B_g\).

The open \(U_0\cap f^{-1}(W)\) is the fiber product of the affine schemes
\(U_0\to V_i\) and \(W\to V_i\).  It is therefore affine with coordinate
ring
\[
A_0\otimes_{B_i}B_g.
\]
This ring is finitely generated over \(B_g\), hence over \(B\).
It contains \(x\) and lies in \(f^{-1}(V)\).  As \(x\) was arbitrary, these
affine neighborhoods cover \(f^{-1}(V)\), proving the desired condition.
\end{solution}

\begin{exerciseintro}
% record: H-II-03-002
\exerciseheading{Exercise II.3.2}
\addcontentsline{toc}{subsection}{Exercise II.3.2}

\begin{context}
A morphism \(f:X\to Y\) is quasi-compact if \(Y\) has an affine open cover
\(\{V_i\}\) such that every \(f^{-1}(V_i)\) is quasi-compact.  Every
quasi-compact scheme has a finite affine open cover.  Principal opens form
a basis in an affine scheme, and a finite union of quasi-compact subspaces
is quasi-compact.
\end{context}

\begin{problem}
Prove that \(f:X\to Y\) is quasi-compact if and only if
\(f^{-1}(V)\) is quasi-compact for every affine open \(V\subseteq Y\).
\end{problem}
\end{exerciseintro}

\begin{solution}
Only the forward implication needs proof.  Choose an affine cover
\(V_i=\operatorname{Spec}B_i\) as in the definition.  Since
\(f^{-1}(V_i)\) is quasi-compact, choose a finite affine cover
\[
f^{-1}(V_i)=\bigcup_{j=1}^{r_i}U_{ij}.
\]
For \(b\in B_i\), the inverse image of \(D(b)\) meets each \(U_{ij}\) in a
principal open: it is defined there by the restriction of \(b\).  Hence
\[
f^{-1}(D(b))
=\bigcup_{j=1}^{r_i}\bigl(U_{ij}\cap f^{-1}(D(b))\bigr)
\]
is a finite union of affine, and therefore quasi-compact, opens.

Now let \(V\subseteq Y\) be affine.  The principal opens
\[
D(b)\subseteq V_i\cap V
\]
as \(i\) and \(b\in B_i\) vary form an open cover of \(V\).  Since \(V\)
is quasi-compact, finitely many, say \(D(b_1),\ldots,D(b_n)\), cover it.
Their inverse images are quasi-compact by the first paragraph, and
\[
f^{-1}(V)=\bigcup_{\ell=1}^n f^{-1}(D(b_\ell))
\]
is consequently quasi-compact.  The converse is immediate by taking an
affine cover of \(Y\).
\end{solution}

\begin{exerciseintro}
% record: H-II-03-003
\exerciseheading{Exercise II.3.3}
\addcontentsline{toc}{subsection}{Exercise II.3.3}

\begin{context}
A morphism is of finite type if the affine covers in the definition of
locally finite type can be chosen finite over each member of an affine
cover of the target.  A morphism is locally of finite type, respectively
quasi-compact, exactly when the corresponding condition holds over every
affine open of the target.

We may use the following algebraic patching lemma.  If
\(D(g_1),\ldots,D(g_r)\) cover \(\operatorname{Spec}A\) and every
\(A_{g_i}\) is a finitely generated \(B\)-algebra, then \(A\) is a
finitely generated \(B\)-algebra.
\end{context}

\begin{problem}
(a) Prove that a morphism is of finite type if and only if it is locally of
finite type and quasi-compact.

(b) Deduce that \(f:X\to Y\) is of finite type if and only if, for every
affine open \(V=\operatorname{Spec}B\subseteq Y\), the inverse image
\(f^{-1}(V)\) has a finite affine cover
\(\operatorname{Spec}A_j\) with all \(A_j\) finitely generated over \(B\).

(c) If \(f\) is of finite type, \(V=\operatorname{Spec}B\subseteq Y\) is
affine, and \(U=\operatorname{Spec}A\subseteq f^{-1}(V)\) is any affine
open, prove that \(A\) is a finitely generated \(B\)-algebra.
\end{problem}
\end{exerciseintro}

\begin{solution}
First note why the patching lemma in the context holds.  Enlarge a
finitely generated \(B\)-subalgebra \(C\subseteq A\) so that it contains
the \(g_i\), coefficients of a relation
\(1=\sum_i c_i g_i\), and numerators of finite generating sets for all
\(A_{g_i}\).  Then \(C_{g_i}=A_{g_i}\) for every \(i\).  For any
\(a\in A\), suitable powers \(g_i^N a\) all lie in \(C\).  The elements
\(g_i^N\) generate the unit ideal in \(C\), so a \(C\)-linear combination
of the \(g_i^Na\) equals \(a\).  Thus \(a\in C\), and \(A=C\).

(a) A finite-type morphism is locally of finite type by definition, and
the finite affine covers of inverse images show that it is quasi-compact.
Conversely, suppose \(f\) is locally of finite type and quasi-compact.
Choose an affine target cover on which the first condition gives affine
source covers by finitely generated algebras.  Each inverse image is
quasi-compact, so finitely many members of its affine cover suffice.  These
are exactly the covers required for \(f\) to be of finite type.

(b) By the affine-independent criteria for local finite type and
quasi-compactness, a finite-type morphism has both properties over every
affine \(V\).  The locally finite-type affine cover of \(f^{-1}(V)\) has a
finite subcover because that inverse image is quasi-compact.  This proves
the forward direction; the reverse direction is the definition applied to
any affine cover of \(Y\).

(c) By part (b), write
\[
f^{-1}(V)=\bigcup_{j=1}^r U_j,
\qquad U_j=\operatorname{Spec}A_j,
\]
with every \(A_j\) finitely generated over \(B\).  The intersections
\(U\cap U_j\) cover the affine scheme \(U\).  Refine this cover by
principal opens of \(U\), and use quasi-compactness to retain finitely many:
\[
U=D(g_1)\cup\cdots\cup D(g_m),
\qquad D(g_\ell)\subseteq U_{j(\ell)}.
\]
We use a small consequence of the patching lemma.  If an affine open
\(W=\operatorname{Spec}C\) lies in
\(\operatorname{Spec}A_{j(\ell)}\), then quasi-compactness of \(W\) gives
a finite cover by principal opens
\(D_{A_{j(\ell)}}(h_q)\) contained in \(W\).  These same sets are the
principal opens \(D_C(h_q)\), and their rings are
\((A_{j(\ell)})_{h_q}\), which are finitely generated over \(B\).
Patching shows that \(C\) is finitely generated over \(B\).  Applying this
to \(W=D(g_\ell)\) proves that every \(A_{g_\ell}\) is a finitely
generated \(B\)-algebra.  The patching lemma now shows that \(A\) itself
is finitely generated over \(B\).
\end{solution}

\begin{exerciseintro}
% record: H-II-03-004
\exerciseheading{Exercise II.3.4}
\addcontentsline{toc}{subsection}{Exercise II.3.4}

\begin{context}
A morphism \(f:X\to Y\) is finite if \(Y\) has an affine open cover
\(\operatorname{Spec}B_i\) whose inverse images are affine
\(\operatorname{Spec}A_i\), with every \(A_i\) finite as a
\(B_i\)-module.  Finiteness of a module is local on the base: if
\(D(b_1),\ldots,D(b_r)\) cover \(\operatorname{Spec}B\) and every
\(M_{b_i}\) is a finite \(B_{b_i}\)-module, then \(M\) is a finite
\(B\)-module.  We may use the affine criterion based on a finite cover by
distinguished affine opens.
\end{context}

\begin{problem}
Prove that \(f:X\to Y\) is finite if and only if, for every affine open
\(V=\operatorname{Spec}B\subseteq Y\), its inverse image is affine,
\[
f^{-1}(V)=\operatorname{Spec}A,
\]
with \(A\) finite as a \(B\)-module.
\end{problem}
\end{exerciseintro}

\begin{solution}
The stated condition plainly implies the definition.  Conversely, assume
finiteness is known over an affine cover of \(Y\), and fix an arbitrary
affine open \(V=\operatorname{Spec}B\).

As in the target-locality argument, choose finitely many principal opens
\(D(b_i)\subseteq V\), each lying in a member of the original cover, such
that the \(D(b_i)\) cover \(V\).  Put
\[
W=f^{-1}(V),\qquad W_i=f^{-1}(D(b_i)).
\]
If \(D_V(b_i)\subseteq V_j=\operatorname{Spec}B_j\) is in the
original cover, with \(f^{-1}(V_j)=\operatorname{Spec}A_j\),
affine base change gives
\(W_i=\operatorname{Spec}(A_j\otimes_{B_j}B_{b_i})\).
The tensor product is finite over \(B_{b_i}\).  Thus
\[
W_i=\operatorname{Spec}A_i
\]
with \(A_i\) finite over \(B_{b_i}\).  The pullbacks of the \(b_i\) are
global functions on \(W\), generate the unit ideal, and have
\(W_i=W_{b_i}\).  Their pairwise intersections are distinguished opens in
the affine schemes \(W_i\), hence are affine and quasi-compact.  The
affineness criterion therefore gives
\[
W=\operatorname{Spec}A,\qquad
A=\Gamma(W,\mathcal O_W).
\]
The localization theorem for global sections identifies
\[
A_{b_i}\cong\Gamma(W_i,\mathcal O_W)=A_i.
\]
Each \(A_{b_i}\) is a finite \(B_{b_i}\)-module.  Since the \(D(b_i)\)
cover \(\operatorname{Spec}B\), finiteness of modules is local on the base,
and \(A\) is a finite \(B\)-module.  This proves the assertion for every
affine \(V\).
\end{solution}

\begin{exerciseintro}
% record: H-II-03-005
\exerciseheading{Exercise II.3.5}
\addcontentsline{toc}{subsection}{Exercise II.3.5}

\begin{context}
A morphism is quasi-finite here if all of its set-theoretic fibers are
finite.  Finite morphisms are stable under base change and composition.
A finite-dimensional algebra over a field is Artinian and has only
finitely many prime ideals.  For an integral ring extension
\(A\subseteq B\), the lying-over theorem makes
\(\operatorname{Spec}B\to\operatorname{Spec}A\) surjective.
The book uses quasi-finite here to mean finite fibers; the common modern definition also requires local finite type.
\end{context}

\begin{problem}
(a) Prove that every finite morphism is quasi-finite.

(b) Prove that every finite morphism is a closed map.

(c) Give a surjective morphism which is of finite type and quasi-finite but
not finite.
\end{problem}
\end{exerciseintro}

\begin{solution}
(a) Let \(f:X\to Y\) be finite and \(y\in Y\).  After restricting to an
affine neighborhood of \(y\), write
\[
Y=\operatorname{Spec}A,\qquad X=\operatorname{Spec}B
\]
with \(B\) a finite \(A\)-module.  The fiber is
\[
X_y=\operatorname{Spec}\bigl(B\otimes_A k(y)\bigr).
\]
The tensor product is finite-dimensional over \(k(y)\), hence Artinian,
and its spectrum is finite.  Its underlying set is the set-theoretic fiber
of \(f\), so \(f\) is quasi-finite.

(b) It is enough to work over an affine open of \(Y\).  A closed subset
\(Z\subseteq X\), with its reduced closed subscheme structure, maps to
\(Y\) by a finite morphism: the closed immersion \(Z\to X\) and \(f\) are
both finite.  Thus it suffices to show that the image of a finite affine
morphism is closed.

Let \(A\to B\) be finite and put \(I=\ker(A\to B)\).  The induced
injection \(A/I\hookrightarrow B\) is integral.  Lying over shows that
\(\operatorname{Spec}B\to\operatorname{Spec}(A/I)\) is surjective, while
\(\operatorname{Spec}(A/I)\to\operatorname{Spec}A\) has image \(V(I)\).
Hence the image is exactly the closed set \(V(I)\).  Applying this to
\(Z\to Y\) proves that \(f(Z)\) is closed.

(c) Let
\[
X=\mathbb A_k^1\coprod
\bigl(\mathbb A_k^1\setminus\{0\}\bigr)
\]
and map both components to \(Y=\mathbb A_k^1\) by their natural maps.  The
first component makes the morphism surjective.  It is of finite type:
one component is an isomorphism and the other is the quasi-compact open
immersion \(D(t)\to\operatorname{Spec}k[t]\).  The fiber over \(0\) has one
point and every other fiber has two, so the map is quasi-finite.

The second component is closed in the disjoint union \(X\), but its image
\(\mathbb A_k^1\setminus\{0\}\) is not closed in \(Y\).  By part (b), the
morphism cannot be finite.
\end{solution}

\begin{exerciseintro}
% record: H-II-03-006
\exerciseheading{Exercise II.3.6}
\addcontentsline{toc}{subsection}{Exercise II.3.6}

\begin{context}
An integral scheme is nonempty, irreducible, and reduced.  Every nonempty
open subset of an irreducible space contains its generic point.  If
\(U=\operatorname{Spec}A\) is an affine open in an integral scheme, then
\(A\) is a domain, and the generic point corresponds to the prime
\((0)\).
\end{context}

\begin{problem}
Let \(X\) be an integral scheme with generic point \(\xi\).  Prove that
\(\mathcal O_{X,\xi}\) is a field.  If
\(U=\operatorname{Spec}A\) is any nonempty affine open of \(X\), identify
this field with the fraction field of \(A\).
\end{problem}
\end{exerciseintro}

\begin{solution}
Every nonempty affine open \(U=\operatorname{Spec}A\) contains \(\xi\).
The ring \(A\) is a domain, and under the affine identification \(\xi\)
corresponds to its zero prime.  Therefore
\[
\mathcal O_{X,\xi}
\cong A_{(0)}
=\operatorname{Frac}(A).
\]
This is a field.  Since the left side is intrinsic to \(X\), the fraction
fields obtained from all nonempty affine opens are canonically identified.
It is called the function field \(K(X)\).
\end{solution}

\begin{exerciseintro}
% record: H-II-03-007
\exerciseheading{Exercise II.3.7}
\addcontentsline{toc}{subsection}{Exercise II.3.7}

\begin{context}
A dominant morphism between integral schemes sends the generic point to
the generic point and induces an inclusion of function fields.  A
zero-dimensional domain of finite type over a field is a finite field
extension.  If a finite set of elements is algebraic over
\(\operatorname{Frac}(A)\), then after inverting one nonzero element of
\(A\) they are integral over the resulting localization.  Integral
finite-type algebras are finite modules.
A finite-type algebra over a field is a Jacobson ring.
\end{context}

\begin{problem}
Let \(f:X\to Y\) be a dominant, generically finite morphism of finite type
between integral schemes.  Prove that there is a dense open
\(V\subseteq Y\) for which
\[
f^{-1}(V)\longrightarrow V
\]
is finite.
\end{problem}
\end{exerciseintro}

\begin{solution}
Let \(\xi\) and \(\eta\) be the generic points of \(X\) and \(Y\).
Dominance gives \(f(\xi)=\eta\).  The generic fiber is a finite-type scheme
over \(K(Y)\) with finitely many points.  It has only the point \(\xi\).
Indeed, an affine neighborhood of any point of that fiber also contains
\(\xi\), because \(X\) is irreducible.  Its intersection with the generic
fiber is the spectrum of a finite-type domain over \(K(Y)\) with finitely
many primes.  A finite-type algebra over a field is Jacobson.  In a domain
with finitely many primes the intersection of its maximal ideals
is zero; if every maximal ideal were nonzero, their product would
be nonzero in that intersection.  Thus the domain is a field, so its
only prime is the one corresponding to \(\xi\).

Choose an affine neighborhood \(U=\operatorname{Spec}B\) of \(\xi\),
lying over an affine \(\operatorname{Spec}A\subseteq Y\).  The coordinate
ring of its generic fiber is
\[
(A\setminus\{0\})^{-1}B.
\]
It is a zero-dimensional domain of finite type over \(K(Y)\), hence a
finite field extension.  It is \(K(X)\), so
\[
[K(X):K(Y)]<\infty.
\]

We next shrink \(Y\) so that its whole inverse image lies in \(U\).
After shrinking the affine target around \(\eta\), quasi-compactness of
\(f\) gives a finite affine cover
\[
f^{-1}(\operatorname{Spec}A)=W_1\cup\cdots\cup W_r
\]
in which \(U\) may be included as one member.  Each nonempty
\(W_i\cap U\) contains \(\xi\).  Choose a nonzero
\(b_i\in\Gamma(W_i,\mathcal O_X)\) such that
\[
\xi\in D_{W_i}(b_i)\subseteq W_i\cap U.
\]
The element \(b_i\in K(X)^\times\) is algebraic over \(K(Y)\).  If
\[
b_i^n+c_{n-1}b_i^{n-1}+\cdots+c_0=0
\qquad(c_j\in K(Y),\ c_0\ne0),
\]
clearing denominators gives a nonzero \(a_i\in A\) with
\(a_i\in b_i\Gamma(W_i,\mathcal O_X)\).  Consequently
\[
W_i\cap f^{-1}(D(a_i))\subseteq D_{W_i}(b_i)\subseteq U.
\]
For \(a=\prod_i a_i\ne0\), it follows from the cover by the \(W_i\) that
\[
f^{-1}(D(a))\subseteq U.
\]
Replace \(A\) by \(A_a\) and \(B\) by \(B_a\).  We have now reduced to
\[
Y=\operatorname{Spec}A,\qquad X=\operatorname{Spec}B,
\]
where \(A\subseteq B\) are domains, \(B\) is a finitely generated
\(A\)-algebra, and \(\operatorname{Frac}B\) is finite over
\(\operatorname{Frac}A\).

Choose \(A\)-algebra generators \(b_1,\ldots,b_r\) of \(B\).  Each \(b_i\)
is algebraic over \(\operatorname{Frac}A\).  Clearing the denominators in
monic equations for all the \(b_i\), there is a nonzero \(a\in A\) such
that every \(b_i\) is integral over \(A_a\).  Hence \(B_a\) is integral
and of finite type over \(A_a\), so it is a finite \(A_a\)-module.  With
\(V=D(a)\subseteq\operatorname{Spec}A\), one has
\[
f^{-1}(V)=\operatorname{Spec}B_a\longrightarrow
\operatorname{Spec}A_a=V,
\]
which is finite.  Since \(a\ne0\) in the domain \(A\), \(V\) is dense in
\(Y\).
\end{solution}

\begin{exerciseintro}
% record: H-II-03-008
\exerciseheading{Exercise II.3.8}
\addcontentsline{toc}{subsection}{Exercise II.3.8}

\begin{context}
Localization commutes with integral closure: if \(A\) is a domain with
fraction field \(K\), \(\widetilde A\) is its integral closure in \(K\),
and \(S\subseteq A\) is multiplicative, then
\[
S^{-1}\widetilde A
\]
is the integral closure of \(S^{-1}A\) in \(K\).  Localizations of an
integrally closed domain are integrally closed.  We may also use the
finiteness theorem: the integral closure of a finitely generated domain
over a field in a finite extension of its fraction field is finite as a
module over that domain.
\end{context}

\begin{problem}
Let \(X\) be an integral scheme.  For every nonempty affine open
\(U=\operatorname{Spec}A\), let \(\widetilde A\) be the integral closure
of \(A\) in \(K(X)\).  Prove that the schemes
\(\widetilde U=\operatorname{Spec}\widetilde A\) glue to a normal integral
scheme \(\widetilde X\), with a natural morphism
\(\nu:\widetilde X\to X\).  Prove the universal property: every dominant
morphism from a normal integral scheme \(Z\) to \(X\) factors uniquely
through \(\nu\).  Finally, prove that \(\nu\) is finite when \(X\) is of
finite type over a field.
\end{problem}
\end{exerciseintro}

\begin{solution}
If \(D(f)\subseteq U=\operatorname{Spec}A\), the localization fact in the
context gives
\[
\widetilde{A_f}=(\widetilde A)_f
\]
inside \(K(X)\).  Hence the normalizations constructed from two affine
opens agree over every principal affine contained in their intersection.
These principal opens cover the intersection, and the resulting
identifications satisfy the cocycle condition because all are identities
inside the same field \(K(X)\).  The schemes
\(\operatorname{Spec}\widetilde A\) therefore glue to a scheme
\(\widetilde X\).

The overlaps are covered by common distinguished affine
neighborhoods, so the gluing identifies the fraction-field
generic points.  Each \(\widetilde A\) is a domain, all charts
share that generic point, and their localizations are integrally
closed.  Thus \(\widetilde X\) is integral and all its local rings are
integrally closed domains; it is normal.  The inclusions
\(A\hookrightarrow\widetilde A\) induce morphisms
\[
\operatorname{Spec}\widetilde A\longrightarrow\operatorname{Spec}A.
\]
They agree on overlaps by the same localization identity, so they glue to
\(\nu:\widetilde X\to X\).

Let \(g:Z\to X\) be dominant, with \(Z\) normal and integral.  Dominance
induces an inclusion \(K(X)\hookrightarrow K(Z)\).  Fix
\(U=\operatorname{Spec}A\subseteq X\) and put \(W=g^{-1}(U)\).  If
\(\alpha\in\widetilde A\), its image in \(K(Z)\) is integral over the image
of \(A\).  For every \(z\in W\), it is therefore integral over the local
ring \(\mathcal O_{Z,z}\).  Normality makes that local ring integrally
closed in \(K(Z)\), so the image of \(\alpha\) lies in every
\(\mathcal O_{Z,z}\).  A rational function on an integral scheme that lies
in all local rings over an open set is a regular section there.  We obtain
a homomorphism
\[
\widetilde A\longrightarrow\Gamma(W,\mathcal O_Z),
\]
and hence a morphism \(W\to\operatorname{Spec}\widetilde A\).  These
morphisms agree on overlaps and glue to
\(\widetilde g:Z\to\widetilde X\), with
\(\nu\circ\widetilde g=g\).  The map on function fields is forced by
\(g\), and regular maps into each affine chart are then forced on the
dense generic point; thus the factorization is unique.

Finally suppose \(X\) is of finite type over a field.  On
\(U=\operatorname{Spec}A\), the ring \(A\) is a finitely generated domain.
Here the relevant field extension in the finiteness theorem is
\(K(X)/\operatorname{Frac}A\), which is the identity extension, so
\(\widetilde A\) is a finite \(A\)-module.  Therefore
\[
\nu^{-1}(U)=\operatorname{Spec}\widetilde A\longrightarrow U
\]
is finite for every affine chart \(U\), and \(\nu\) is finite.
\end{solution}

\begin{exerciseintro}
% record: H-II-03-009
\exerciseheading{Exercise II.3.9}
\addcontentsline{toc}{subsection}{Exercise II.3.9}

\begin{context}
Fiber products of affine schemes satisfy
\[
\operatorname{Spec}B\times_{\operatorname{Spec}A}\operatorname{Spec}C
\cong\operatorname{Spec}(B\otimes_A C).
\]
Points of a localization correspond to primes disjoint from the
multiplicative set.
\end{context}

\begin{problem}
(a) Show that
\[
\mathbb A_k^1\times_{\operatorname{Spec}k}\mathbb A_k^1
\cong\mathbb A_k^2,
\]
but that the underlying point set is not the Cartesian product of the
underlying point sets of the two affine lines, even when \(k\) is
algebraically closed.

(b) Let \(s,t\) be indeterminates over \(k\).  Describe
\[
\operatorname{Spec}k(s)\times_{\operatorname{Spec}k}
\operatorname{Spec}k(t),
\]
although each of the three displayed factor and base spaces has one point.
\end{problem}
\end{exerciseintro}

\begin{solution}
(a) The affine fiber-product formula gives
\[
\operatorname{Spec}k[x]\times_{\operatorname{Spec}k}
\operatorname{Spec}k[y]
\cong\operatorname{Spec}(k[x]\otimes_k k[y])
\cong\operatorname{Spec}k[x,y]=\mathbb A_k^2.
\]
On points, the two projections send a prime
\(\mathfrak q\subseteq k[x,y]\) to its contractions to \(k[x]\) and
\(k[y]\).  The distinct primes
\[
(x-y)\qquad\text{and}\qquad(x-y^2)
\]
both contract to \((0)\) in each polynomial subring.  Thus two different
points of the scheme product map to the same pair of generic points.  The
underlying set cannot be the Cartesian product, regardless of whether
\(k\) is algebraically closed.

(b) The product is
\[
\operatorname{Spec}R,\qquad
R=k(s)\otimes_k k(t).
\]
Equivalently,
\[
R=S^{-1}k[s,t],
\quad
S=\{g(s)h(t):0\ne g\in k[s],\ 0\ne h\in k[t]\}.
\]
Hence its points correspond exactly to primes
\(\mathfrak p\subseteq k[s,t]\) with
\[
\mathfrak p\cap k[s]=(0),\qquad
\mathfrak p\cap k[t]=(0).
\]
They include the generic point \((0)\) and the height-one primes generated
by irreducible polynomials involving both variables; there are no
height-two primes with these contraction properties.
Indeed, the residue field of a height-two prime is finite algebraic
over \(k\) by Zariski's lemma, making both contractions nonzero.
The distinct primes \((s-t^j)\), \(j\geq1\), all survive.
Thus the product has infinitely many points even though both
factors are one-point schemes.
Its topology and structure sheaf are those of the displayed localization.
\end{solution}

\begin{exerciseintro}
% record: H-II-03-010
\exerciseheading{Exercise II.3.10}
\addcontentsline{toc}{subsection}{Exercise II.3.10}

\begin{context}
For \(f:X\to Y\) and \(y\in Y\), the scheme-theoretic fiber is
\[
X_y=X\times_Y\operatorname{Spec}k(y).
\]
If \(Y=\operatorname{Spec}A\), \(X=\operatorname{Spec}B\), and
\(y\) corresponds to \(\mathfrak p\subseteq A\), then
\[
X_y=\operatorname{Spec}\bigl(B\otimes_A k(\mathfrak p)\bigr)
=\operatorname{Spec}\bigl((A\setminus\mathfrak p)^{-1}
(B/\mathfrak pB)\bigr).
\]
\end{context}

\begin{problem}
(a) Prove that the underlying space of \(X_y\) is homeomorphic to the
set-theoretic fiber \(f^{-1}(y)\) with its subspace topology.

(b) Let \(k\) be algebraically closed, let
\[
X=\operatorname{Spec}k[s,t]/(s-t^2),\qquad
Y=\operatorname{Spec}k[s],
\]
and let \(f:X\to Y\) be induced by the inclusion \(k[s]\to
k[s,t]/(s-t^2)\).  Assuming \(\operatorname{char}k\ne2\), describe the
fibers over a nonzero \(a\in k\), over \(0\), and over the generic point
of \(Y\).  Also state what changes in characteristic \(2\).
\end{problem}
\end{exerciseintro}

\begin{solution}
(a) Work first with affine neighborhoods
\(Y=\operatorname{Spec}A\) and \(X=\operatorname{Spec}B\).  The primes of
\[
(A\setminus\mathfrak p)^{-1}(B/\mathfrak pB)
\]
correspond to the primes \(\mathfrak q\subseteq B\) which contain
\(\mathfrak pB\) and avoid the image of \(A\setminus\mathfrak p\).
Equivalently, they are exactly those satisfying
\(\mathfrak q\cap A=\mathfrak p\), namely the points of \(f^{-1}(y)\).
The quotient-spectrum correspondence gives the subspace topology on
\(V(\mathfrak pB)\), and localization then gives the subspace consisting
of primes disjoint from \(A\setminus\mathfrak p\).  Thus the bijection is
a homeomorphism.  These affine homeomorphisms agree on overlaps and glue
for general \(X\) and \(Y\).

(b) For the closed point \(s=a\), the fiber is
\[
X_a\cong\operatorname{Spec}k[t]/(t^2-a).
\]
If \(a\ne0\) and \(\operatorname{char}k\ne2\), choose
\(\alpha\in k\) with \(\alpha^2=a\).  The roots
\(\alpha\) and \(-\alpha\) are distinct, so the Chinese remainder theorem
gives
\[
k[t]/(t^2-a)\cong k\times k.
\]
The fiber is two reduced points, both with residue field \(k\).  At
\(a=0\), the ring is \(k[t]/(t^2)\), a nonreduced local ring with one
point.

At the generic point \(\eta\), whose residue field is \(k(s)\),
\[
X_\eta\cong
\operatorname{Spec}\bigl(k(s)[t]/(t^2-s)\bigr).
\]
The element \(s\) is not a square in \(k(s)\), so \(t^2-s\) is irreducible.
The fiber is one point with residue field a degree-two extension of
\(k(s)\).

In characteristic \(2\), the generic description remains a purely
inseparable degree-two field extension, but every closed fiber is
\(k[t]/((t-\alpha)^2)\).  Thus even for \(a\ne0\) it is a nonreduced
one-point scheme; the asserted two-point conclusion requires the stated
characteristic hypothesis.
\end{solution}

\begin{exerciseintro}
% record: H-II-03-011
\exerciseheading{Exercise II.3.11}
\addcontentsline{toc}{subsection}{Exercise II.3.11}

\begin{context}
An affine closed immersion has the form
\(\operatorname{Spec}(A/I)\to\operatorname{Spec}A\).  Closed subschemes
of a scheme correspond to quasi-coherent ideal sheaves, with a larger
ideal giving a smaller closed subscheme.  Arbitrary sums of
quasi-coherent ideal sheaves are quasi-coherent.  A closed immersion may
be checked on an affine open cover of the target.
\end{context}

\begin{problem}
(a) Prove that a base change of a closed immersion is a closed immersion.

(b) If \(Y\to X=\operatorname{Spec}A\) is a closed immersion, prove that
\(Y\) is affine and is isomorphic over \(X\) to
\(\operatorname{Spec}(A/\mathfrak a)\) for an ideal
\(\mathfrak a\subseteq A\).

(c) Give a closed subset \(C\subseteq X\) its reduced induced subscheme
structure.  If \(Y'\subseteq X\) is any closed subscheme with underlying
space \(C\), prove that \(C_{\mathrm{red}}\to X\) factors through
\(Y'\to X\).

(d) Prove that every morphism \(f:Z\to X\) has a unique smallest closed
subscheme of \(X\) through which it factors.  This is the
scheme-theoretic image.  If \(Z\) is reduced, identify it with the reduced
induced structure on \(\overline{f(Z)}\).
\end{problem}
\end{exerciseintro}

\begin{solution}
(a) The assertion is local on the new base.  In affine coordinates, base
change of
\[
\operatorname{Spec}(A/I)\longrightarrow\operatorname{Spec}A
\]
along \(\operatorname{Spec}B\to\operatorname{Spec}A\) is
\[
\operatorname{Spec}\bigl((A/I)\otimes_A B\bigr)
\cong\operatorname{Spec}(B/IB)
\longrightarrow\operatorname{Spec}B,
\]
which is a closed immersion.

(b) For every \(y\in Y\), choose an affine open neighborhood in \(Y\).
Because the topology on \(Y\) is the subspace topology, there is a
principal open \(D(f)\subseteq X\) containing \(y\) such that
\(Y\cap D(f)\) lies in that affine neighborhood.  It is a distinguished
open there, and hence is affine.  Since \(Y\) is closed in the
quasi-compact space \(X\), finitely many such \(Y\cap D(f_i)\) cover
\(Y\).

The union of the corresponding \(D(f_i)\) in \(X\) is an open neighborhood of \(Y\) in \(X\).  Its closed
complement is quasi-compact and lies in \(X\setminus Y\); add finitely many
principal opens contained in \(X\setminus Y\) to obtain a finite principal
cover of all of \(X\).  The restrictions of the corresponding \(f_i\) to
\(Y\) generate the unit ideal, and each distinguished open \(Y_{f_i}\) is
affine, including the empty ones.  The affine criterion shows that
\(Y=\operatorname{Spec}B\).

The closed immersion is now induced by a ring homomorphism \(A\to B\).
Its underlying map is a homeomorphism onto a closed subset and its sheaf
map is surjective, so the affine closed-immersion criterion makes
\(A\to B\) surjective.  With
\(\mathfrak a=\ker(A\to B)\), this gives
\(Y\cong\operatorname{Spec}(A/\mathfrak a)\) over \(X\).

(c) Work on an affine open \(\operatorname{Spec}A\subseteq X\).  By part
(b), the two closed subschemes have the forms
\[
\operatorname{Spec}(A/I),\qquad
\operatorname{Spec}(A/I').
\]
For the reduced induced structure, \(I\) is the radical ideal of functions
vanishing on \(C\).  Equality of the underlying closed sets gives
\(\sqrt{I'}=I\), and in particular \(I'\subseteq I\).  The quotient map
\[
A/I'\longrightarrow A/I
\]
defines the required factorization locally.  These maps are canonical and
therefore agree on overlaps, giving the global factorization.

(d) For every closed subscheme \(Y_\alpha\subseteq X\) through which \(f\)
factors, let \(\mathcal I_\alpha\subseteq\mathcal O_X\) be its ideal
sheaf.  Set
\[
\mathcal I=\sum_\alpha\mathcal I_\alpha.
\]
This is a quasi-coherent ideal sheaf and defines a closed subscheme
\(Y\subseteq X\).  Each \(\mathcal I_\alpha\) maps to zero under
\(\mathcal O_X\to f_*\mathcal O_Z\), so their sum does as well; hence
\(f\) factors through \(Y\).  Since
\(\mathcal I_\alpha\subseteq\mathcal I\), the inclusion \(Y\to X\)
factors through every \(Y_\alpha\).  This proves existence and the
universal property, which also gives uniqueness.

Now suppose \(Z\) is reduced and put \(C=\overline{f(Z)}\).  Locally, a
function vanishing on \(C\) pulls back to a function lying in every prime
of every affine open of \(Z\).  It is therefore nilpotent, and hence zero
because \(Z\) is reduced.  Thus \(f\) factors through
\(C_{\mathrm{red}}\).  Every closed subscheme through which \(f\) factors
has underlying space containing \(C\); part (c) shows that
\(C_{\mathrm{red}}\) factors through it.  Hence
\(C_{\mathrm{red}}\) has the defining universal property of the
scheme-theoretic image.
\end{solution}

\begin{exerciseintro}
% record: H-II-03-012
\exerciseheading{Exercise II.3.12}
\addcontentsline{toc}{subsection}{Exercise II.3.12}

\begin{context}
A graded homomorphism \(S\to T\) induces a morphism from
\[
U=\{\mathfrak q\in\operatorname{Proj}T:
\mathfrak q\not\supseteq\varphi(S_+)\}
\]
to \(\operatorname{Proj}S\).  On a standard open it is induced by
\[
(S_f)_0\longrightarrow(T_{\varphi(f)})_0.
\]
A morphism covered by affine closed immersions is a closed immersion.
\end{context}

\begin{problem}
(a) Let \(\varphi:S\to T\) be a surjective degree-preserving homomorphism
of graded rings.  Prove that \(U=\operatorname{Proj}T\) and that the
induced morphism
\[
\operatorname{Proj}T\longrightarrow\operatorname{Proj}S
\]
is a closed immersion.

(b) Let \(I\subseteq S\) be homogeneous.  Show that different homogeneous
ideals can define the same closed subscheme of \(\operatorname{Proj}S\).
Specifically, for an integer \(d_0\), put
\[
I'=\bigoplus_{d\ge d_0}I_d.
\]
Prove that \(I\) and \(I'\) define the same closed subscheme.
\end{problem}
\end{exerciseintro}

\begin{solution}
(a) Surjectivity in each degree gives
\(\varphi(S_+)=T_+\).  No point of \(\operatorname{Proj}T\) contains
\(T_+\), so \(U=\operatorname{Proj}T\).  For every homogeneous
\(f\in S_+\), localization and then passage to degree zero give a
surjection
\[
(S_f)_0\twoheadrightarrow(T_{\varphi(f)})_0.
\]
Thus the morphism over \(D_+(f)\) is an affine closed immersion.  The
standard opens cover \(\operatorname{Proj}S\), so the global morphism is a
closed immersion.

(b) The quotients \(S/I\) and \(S/I'\) have identical graded pieces in all
degrees \(d\ge d_0\).  More explicitly, on \(D_+(f)\) an element of the
ideal defining \(\operatorname{Proj}(S/I)\) has the form
\[
\frac{a}{f^n},\qquad a\in I_{n\deg f}.
\]
Multiplying numerator and denominator by a sufficiently high power of
\(f\) makes the numerator have degree at least \(d_0\), without changing
the fraction.  It then belongs to \(I'\).  Hence \(I\) and \(I'\) induce
the same ideal in every ring \((S_f)_0\).  Their affine closed subschemes
agree on the standard cover and therefore agree globally.
\end{solution}

\begin{exerciseintro}
% record: H-II-03-013
\exerciseheading{Exercise II.3.13}
\addcontentsline{toc}{subsection}{Exercise II.3.13}

\begin{context}
A morphism is of finite type exactly when it is locally of finite type and
quasi-compact.  These properties may be checked over arbitrary affine
opens of the target.  A finitely generated algebra over a noetherian ring
is noetherian, and fiber products of affine schemes correspond to tensor
products.
\end{context}

\begin{problem}
Prove the following properties.

(a) Every closed immersion is of finite type.

(b) Every quasi-compact open immersion is of finite type.

(c) A composite of finite-type morphisms is of finite type.

(d) Finite-type morphisms are stable under base change.

(e) If \(X\) and \(Y\) are of finite type over \(S\), then
\(X\times_SY\) is of finite type over \(S\).

(f) If \(X\xrightarrow{f}Y\xrightarrow{g}Z\), \(f\) is quasi-compact, and
\(g\circ f\) is of finite type, then \(f\) is of finite type.

(g) If \(f:X\to Y\) is of finite type and \(Y\) is noetherian, then \(X\)
is noetherian.
\end{problem}
\end{exerciseintro}

\begin{solution}
(a) Over \(\operatorname{Spec}A\), a closed immersion has inverse image
\(\operatorname{Spec}(A/I)\).  The quotient is a finitely generated
\(A\)-algebra, so the morphism is of finite type.

(b) Let \(j:U\to Y\) be a quasi-compact open immersion and
\(V=\operatorname{Spec}A\subseteq Y\).  The open \(U\cap V\) is
quasi-compact, so finitely many principal opens \(D(f_i)\subseteq V\)
cover it.  Each has coordinate ring
\[
A_{f_i}=A[t]/(f_it-1),
\]
a finitely generated \(A\)-algebra.  The affine criterion proves that
\(j\) is of finite type.

(c) Work over an affine \(\operatorname{Spec}A\) in the final target.
The first inverse image has a finite affine cover
\(\operatorname{Spec}B_i\), with each \(B_i\) finitely generated over
\(A\).  The inverse image of each of these has a finite affine cover
\(\operatorname{Spec}C_{ij}\), with \(C_{ij}\) finitely generated over
\(B_i\).  Each \(C_{ij}\) is finitely generated over \(A\), and the
combined cover is finite.

(d) Affine locally, base change replaces a finite-type map \(A\to B\) by
\[
C\longrightarrow B\otimes_A C.
\]
If \(b_1,\ldots,b_n\) generate \(B\) over \(A\), their tensors generate
the latter algebra over \(C\).  The affine descriptions glue, proving
stability under arbitrary base change.

(e) The projection \(X\times_SY\to Y\) is the base change of
\(X\to S\), hence is of finite type by part (d).  Composing with the
finite-type morphism \(Y\to S\) and using part (c) proves the claim.

(f) It remains to show that \(f\) is locally of finite type.  Work with
affine neighborhoods
\[
W=\operatorname{Spec}A\subseteq Z,\qquad
V=\operatorname{Spec}B\subseteq g^{-1}(W),\qquad
U=\operatorname{Spec}C\subseteq f^{-1}(V).
\]
By the arbitrary-affine criterion for the finite-type morphism
\(g\circ f\), \(C\) is finitely generated over \(A\).  Thus the relevant
ring maps have the form
\[
A\longrightarrow B\longrightarrow C.
\]
If \(C\) is generated over \(A\) by \(c_1,\ldots,c_n\), the same elements
generate it over \(B\), since the image of \(A\) is contained in the image
of \(B\).  Thus \(f\) is locally of finite type.  It is quasi-compact by
hypothesis, so it is of finite type.

(g) A noetherian scheme has a finite affine cover
\(\operatorname{Spec}A_i\) with all \(A_i\) noetherian.  Finite type and
quasi-compactness give a finite affine cover of every
\(f^{-1}(\operatorname{Spec}A_i)\) by spectra of finitely generated
\(A_i\)-algebras.  Those algebras are noetherian by the Hilbert basis
theorem.  Hence \(X\) has a finite affine cover by noetherian rings and is
a noetherian scheme.
\end{solution}

\begin{exerciseintro}
% record: H-II-03-014
\exerciseheading{Exercise II.3.14}
\addcontentsline{toc}{subsection}{Exercise II.3.14}

\begin{context}
If \(A\) is a finitely generated algebra over a field \(k\), the residue
field at every maximal ideal is a finite extension of \(k\).  Conversely,
for a prime \(\mathfrak p\) of a finitely generated \(k\)-algebra, if
\(\operatorname{Frac}(A/\mathfrak p)\) is algebraic over \(k\), then
\(A/\mathfrak p\) is a field.  Every nonzero ring has a maximal ideal.
\end{context}

\begin{problem}
Prove that the closed points of a scheme of finite type over a field are
dense.  Give an example showing that this need not hold for an arbitrary
scheme.
\end{problem}
\end{exerciseintro}

\begin{solution}
Let \(X\) be of finite type over \(k\), and let \(U\subseteq X\) be a
nonempty open set.  Choose a nonempty affine open
\(V=\operatorname{Spec}A\subseteq U\).  A maximal ideal of \(A\) gives a
point \(x\in V\) whose residue field \(k(x)\) is finite over \(k\).

We check that \(x\) is closed in all of \(X\), not merely in \(V\).  Let
\(z\in\overline{\{x\}}\) and choose an affine neighborhood
\(\operatorname{Spec}B\) of \(z\).  It also contains \(x\), since every
neighborhood of a specialization contains its generizations.  If
\(\mathfrak p\subseteq B\) represents \(x\), then
\[
\operatorname{Frac}(B/\mathfrak p)=k(x)
\]
is algebraic over \(k\).  The fact in the context makes \(B/\mathfrak p\)
a field, so \(\mathfrak p\) is maximal.  The point \(z\), represented by a
prime containing \(\mathfrak p\), must equal \(x\).  Thus \(x\) is closed
in \(X\).  Every nonempty open contains such a point, proving density.

For a counterexample, take the spectrum of the discrete valuation ring
\[
R=k[t]_{(t)}.
\]
It has a generic point \((0)\) and one closed point \((t)\).  The nonempty
open \(D(t)=\{(0)\}\) contains no closed point, so the closed points are
not dense.
\end{solution}

\begin{exerciseintro}
% record: H-II-03-015
\exerciseheading{Exercise II.3.15}
\addcontentsline{toc}{subsection}{Exercise II.3.15}

\begin{context}
Let \(X\) be of finite type over a field \(k\).  We may use these standard
base-change facts.

First, a field extension is faithfully flat, so reducedness and
irreducibility descend from \(X_L\) to \(X_K\) for extensions
\(L/K/k\).  Second, if the base field is algebraically closed, every
irreducible finite-type scheme remains irreducible after arbitrary field
extension.  Third, a reduced finite-type scheme over a perfect field is
geometrically reduced.  Finally, base change by an algebraic purely
inseparable field extension is a universal homeomorphism.
\end{context}

\begin{problem}
Let \(\overline k\), \(k_s\), and \(k_p\) denote an algebraic closure, a
separable closure, and a perfect closure of \(k\), respectively.

(a) Prove the equivalence of:
\[
X_{\overline k}\text{ is irreducible};\qquad
X_{k_s}\text{ is irreducible};\qquad
X_K\text{ is irreducible for every field extension }K/k.
\]

(b) Prove the equivalence of:
\[
X_{\overline k}\text{ is reduced};\qquad
X_{k_p}\text{ is reduced};\qquad
X_K\text{ is reduced for every field extension }K/k.
\]

(c) Give an integral finite-type \(k\)-scheme which is neither
geometrically irreducible nor geometrically reduced.
\end{problem}
\end{exerciseintro}

\begin{solution}
(a) The third condition implies the first, and the first implies the second
by faithfully flat descent along
\(\overline k/k_s\).  Conversely, \(\overline k/k_s\) is purely
inseparable.  The projection
\[
X_{\overline k}\longrightarrow X_{k_s}
\]
is therefore a universal homeomorphism, so the second condition implies
that \(X_{\overline k}\) is irreducible.

It remains to pass from the first condition to an arbitrary \(K/k\).
Choose an algebraically closed field \(\Omega\) containing both \(K\) and
an embedded algebraic closure of \(k\).  Since an irreducible finite-type
scheme over the algebraically closed field \(\overline k\) stays
irreducible after extending scalars, \(X_\Omega\) is irreducible.  The
projection \(X_\Omega\to X_K\) is surjective, and a continuous image of an
irreducible space is irreducible.  Hence \(X_K\) is irreducible.

(b) Again, the third condition implies the first, and the first descends
to the second.  Conversely, the field \(k_p\) is perfect.  If \(X_{k_p}\)
is reduced, the perfect-field result in the context says it remains
reduced over \(\overline k\), proving the first condition.

Assume now that \(X_{\overline k}\) is reduced and let \(K/k\) be any
extension.  Embed \(K\) and \(\overline k\) in a common algebraically
closed extension \(\Omega\).  A reduced finite-type scheme over the
perfect field \(\overline k\) is geometrically reduced, so \(X_\Omega\) is
reduced.  Reducedness descends through the faithfully flat extension
\(\Omega/K\), giving that \(X_K\) is reduced.  This proves the third
condition.

(c) Take
\[
k=\mathbb F_3(t),\qquad
L=\mathbb F_9(t^{1/3}),\qquad X=\operatorname{Spec}L.
\]
The scheme \(X\) is integral and finite over \(k\).  Over an algebraic
closure \(\overline k\), the separable constant-field extension
\(\mathbb F_9/\mathbb F_3\) gives two factors, while the purely inseparable
equation \(u^3-t\) becomes \((u-t^{1/3})^3\).  Consequently
\[
L\otimes_k\overline k
\cong
\overline k[\epsilon]/(\epsilon^3)
\times
\overline k[\epsilon]/(\epsilon^3).
\]
Its spectrum is reducible and its ring is nonreduced.  Thus \(X\) has
neither geometric property.
\end{solution}

\begin{exerciseintro}
% record: H-II-03-016
\exerciseheading{Exercise II.3.16}
\addcontentsline{toc}{subsection}{Exercise II.3.16}

\begin{context}
In a noetherian topological space, descending chains of closed subsets
stabilize.  Equivalently, every nonempty collection of closed subsets has
a minimal member.
\end{context}

\begin{problem}
Let \(X\) be noetherian and let \(\mathcal P\) be a property of closed
subsets.  Suppose that whenever \(\mathcal P\) holds for every proper
closed subset of a closed set \(Y\), it also holds for \(Y\).  Prove that
\(\mathcal P\) holds for every closed subset of \(X\), in particular for
\(X\).
\end{problem}
\end{exerciseintro}

\begin{solution}
If the conclusion were false, the collection of closed subsets for which
\(\mathcal P\) fails would be nonempty.  Choose a minimal member \(Y\).
Every proper closed subset of \(Y\) then satisfies \(\mathcal P\), so the
hypothesis implies that \(Y\) satisfies \(\mathcal P\), a contradiction.
Thus there is no counterexample.  This argument is called noetherian
induction.
\end{solution}

\begin{exerciseintro}
% record: H-II-03-017
\exerciseheading{Exercise II.3.17}
\addcontentsline{toc}{subsection}{Exercise II.3.17}

\begin{context}
A Zariski space is a noetherian topological space in which every nonempty
irreducible closed subset has a unique generic point.  In a noetherian
space, every closed subset is a finite union of irreducible components.
For a noetherian space \(X\), let \(t(X)\) be the set of its nonempty
irreducible closed subsets, with closed subsets
\[
t(C)=\{Z\in t(X):Z\subseteq C\}
\]
for closed \(C\subseteq X\).  Define
\(\alpha:X\to t(X)\) by \(\alpha(x)=\overline{\{x\}}\).
\end{context}

\begin{problem}
(a) Prove that the underlying space of a noetherian scheme is a Zariski
space.

(b) Prove that a minimal nonempty closed subset of a Zariski space is a
single point.

(c) Prove that every Zariski space is \(T_0\).

(d) Prove that the generic point of an irreducible Zariski space lies in
every nonempty open subset.

(e) Order points by declaring \(x_1\geq x_0\) when \(x_1\) specializes to
\(x_0\), meaning \(x_0\in\overline{\{x_1\}}\).  Identify the minimal and
maximal points.  Prove that closed sets are stable under specialization
and open sets under generization.

(f) Prove that \(t(X)\) is a Zariski space for every noetherian \(X\), and
that \(X\) is a Zariski space if and only if
\(\alpha:X\to t(X)\) is a homeomorphism.
\end{problem}
\end{exerciseintro}

\begin{solution}
(a) The underlying space of a noetherian scheme is noetherian.  Every
nonempty irreducible closed subset of a scheme has a unique generic point,
so it is a Zariski space.

(b) Let \(C\) be minimal among nonempty closed subsets.  It is irreducible:
a decomposition into two proper closed subsets would contradict
minimality.  Let \(\eta\) be its generic point.  For any \(x\in C\), the
nonempty closed set \(\overline{\{x\}}\) is contained in \(C\), hence
equals \(C\).  Both \(x\) and \(\eta\) are then generic points of \(C\),
so uniqueness gives \(x=\eta\).  Thus \(C=\{\eta\}\).

(c) If distinct points \(x,y\) had the same closure, they would be two
generic points of the same irreducible closed subset.  Therefore their
closures differ.  If, for example,
\(y\notin\overline{\{x\}}\), then
\(X\setminus\overline{\{x\}}\) contains \(y\) but not \(x\).  This is the
\(T_0\) axiom.

(d) Let \(\eta\) be the generic point.  If a nonempty open \(U\) omitted
\(\eta\), its closed complement would contain \(\eta\), hence also
\(\overline{\{\eta\}}=X\).  This would make \(U\) empty, a contradiction.

(e) The \(T_0\) property makes specialization a partial order.  A point is
minimal exactly when it has no proper specialization, equivalently when
its closure is a singleton; these are the closed points from part (b).
If \(x\) is maximal, choose an irreducible component \(C\)
containing \(x\).  Its generic point \(\eta_C\) generalizes \(x\),
so maximality forces \(x=\eta_C\).
Conversely, a proper generization of \(\eta_C\) would have
irreducible closure properly containing \(C\), contradicting
maximality of an irreducible component.
Thus the maximal points are precisely the generic points of
irreducible components.

If a closed set contains \(x\), it contains
\(\overline{\{x\}}\), and hence every specialization of \(x\).  If an
open set contains \(x\) and \(y\) is a generization of \(x\), every open
neighborhood of \(x\) contains \(y\); hence the open set contains \(y\).

(f) A descending chain
\[
t(C_1)\supseteq t(C_2)\supseteq\cdots
\]
gives a descending chain \(C_1\supseteq C_2\supseteq\cdots\), because a
closed set in a noetherian space is the union of its finitely many
irreducible components.  The latter chain stabilizes, so \(t(X)\) is
noetherian.  Moreover, \(t(C)\) is irreducible exactly when \(C\) is
irreducible.  In that case the point \(C\in t(X)\) has closure \(t(C)\),
and it is the unique generic point.  Hence \(t(X)\) is a Zariski space.

For every closed \(C\subseteq X\),
\[
\alpha^{-1}(t(C))=C,
\]
so \(\alpha\) is continuous.  It is injective exactly when distinct points
have distinct closures, and it is surjective exactly when every
irreducible closed subset has a generic point.  If \(X\) is a Zariski
space, both hold; furthermore
\(\alpha(C)=t(C)\), so \(\alpha\) is closed and is a homeomorphism.
Conversely, if \(\alpha\) is a homeomorphism, \(X\) is homeomorphic to the
Zariski space \(t(X)\), and is therefore itself a Zariski space.
\end{solution}

\begin{exerciseintro}
% record: H-II-03-018
\exerciseheading{Exercise II.3.18}
\addcontentsline{toc}{subsection}{Exercise II.3.18}

\begin{context}
A constructible subset belongs to the Boolean algebra generated by the open
subsets.  A locally closed subset is an intersection of an open and a
closed subset.  In a Zariski space, closed subsets have finitely many
irreducible components, and the generic point of an irreducible space lies
in every nonempty open subset.
The specialization criterion in (c) retains the Zariski-space hypothesis.
\end{context}

\begin{problem}
(a) Prove that a subset is constructible if and only if it is a finite
disjoint union of locally closed subsets.

(b) In an irreducible Zariski space, prove that a constructible subset is
dense if and only if it contains the generic point.  Prove that a dense
constructible subset contains a nonempty open subset.

(c) Prove that a subset is closed if and only if it is constructible and
stable under specialization.  State and prove the corresponding criterion
for openness.

(d) If \(f:X\to Y\) is continuous between Zariski spaces, prove that the
inverse image of a constructible subset is constructible.
\end{problem}
\end{exerciseintro}

\begin{solution}
(a) Any particular constructible set is obtained by finitely many Boolean
operations from finitely many open sets \(U_1,\ldots,U_n\).  The disjoint
atoms
\[
\left(\bigcap_{i\in I}U_i\right)
\cap
\left(\bigcap_{j\notin I}(X\setminus U_j)\right),
\qquad I\subseteq\{1,\ldots,n\},
\]
partition \(X\).  Each atom is the intersection of an open set with a
closed set and is therefore locally closed.  The constructible set is the
union of some of these finitely many atoms.  Conversely, locally closed
sets are constructible, and constructible sets are closed under finite
unions.

(b) Let \(\eta\) be the generic point.  Any subset containing \(\eta\) is
dense because its closure contains
\(\overline{\{\eta\}}=X\).

Conversely, write a dense constructible subset as a finite union
\[
C=\bigcup_{i=1}^r(U_i\cap F_i),
\]
with \(U_i\) open and \(F_i\) closed.  Refine the \(F_i\) into their
irreducible components and discard empty intersections.  The closure of
each remaining piece is contained in its irreducible closed component.
Since the finite union of these closures is \(X\), irreducibility of \(X\)
forces one component to be all of \(X\).  The corresponding piece is then
a nonempty open subset of \(X\).  It contains \(\eta\), so \(C\) contains
both the generic point and a nonempty open subset.

(c) Closed subsets are constructible and stable under specialization.
Conversely, suppose \(C\) has these two properties.  Decompose \(C\) into
finitely many locally closed pieces, and refine their closed parts into
irreducible components.  A nonempty piece open in an irreducible closed
set \(F\) contains the generic point of \(F\).  Stability under
specialization then makes \(C\) contain all of \(F\).  Consequently \(C\)
is the finite union of these closed sets \(F\), and is closed.

A subset is open if and only if it is constructible and stable under
generization.  This follows by applying the closed criterion to its
complement: complements preserve constructibility and interchange the two
stability conditions.

(d) The inverse image of a locally closed set \(U\cap F\) is
\[
f^{-1}(U)\cap f^{-1}(F),
\]
again locally closed.  Apply part (a) to a finite locally closed
decomposition.
\end{solution}

\begin{exerciseintro}
% record: H-II-03-019
\exerciseheading{Exercise II.3.19}
\addcontentsline{toc}{subsection}{Exercise II.3.19}

\begin{context}
We may use the following commutative-algebra lemma.  Let
\(A\subseteq B\) be noetherian domains with \(B\) a finitely generated
\(A\)-algebra.  For every \(0\ne b\in B\), there is \(0\ne a\in A\) such
that any homomorphism \(A\to K\) to an algebraically closed field, with
\(a\) mapping nonzero, extends to \(B\to K\) with \(b\) mapping nonzero.

Open subschemes of noetherian schemes are quasi-compact, and
quasi-compact open immersions and closed immersions are of finite type.
\end{context}

\begin{problem}
Let \(f:X\to Y\) be a finite-type morphism of noetherian schemes.  Prove
Chevalley's theorem: the image of every constructible subset of \(X\) is
constructible in \(Y\), using these steps.

(a) Reduce to the case in which \(X\) and \(Y\) are affine integral
noetherian schemes and \(f\) is dominant.

(b) In that case, use the lemma in the context to show that \(f(X)\)
contains a nonempty open subset of \(Y\).

(c) Complete the proof by noetherian induction on the target.

(d) Give a morphism of varieties whose image is neither open nor closed.
\end{problem}
\end{exerciseintro}

\begin{solution}
(a) By the finite locally closed decomposition, it is enough to treat one
locally closed subset \(C\subseteq X\).  Give its closure the reduced
closed structure and \(C\) its open subscheme structure.  Since \(X\) is
noetherian, \(C\) is quasi-compact, and the composite
\[
C\longrightarrow X\longrightarrow Y
\]
is of finite type.  Thus it suffices to prove that the image of an entire
finite-type source is constructible.

Cover \(Y\) by finitely many affine opens and each inverse image by finitely
many affine opens.  Images and constructible sets may be united finitely,
so reduce to \(X\) and \(Y\) affine.  They have finitely many irreducible
components.  Give each component its reduced structure and treat its image
separately, reducing to \(X\) integral.  The closure of its image is an
irreducible closed subset of \(Y\); replace \(Y\) by that closure with its
reduced structure.  A closed subscheme of an affine scheme is affine, and
the resulting map is dominant.  We have reached the stated case.

(b) Write \(X=\operatorname{Spec}B\) and
\(Y=\operatorname{Spec}A\).  Dominance and integrality make
\(A\to B\) injective.  Apply the algebra lemma with \(b=1\), obtaining
\(0\ne a\in A\).  We claim
\[
D(a)\subseteq f(X).
\]
For \(\mathfrak p\in D(a)\), embed \(k(\mathfrak p)\) in an algebraic
closure \(K\).  The map \(A\to k(\mathfrak p)\to K\) does not kill \(a\),
so it extends to a homomorphism \(B\to K\).  Its kernel is a prime
\(\mathfrak q\) whose contraction to \(A\) is \(\mathfrak p\).
Therefore \(f(\mathfrak q)=\mathfrak p\).  Since \(A\) is a domain and
\(a\ne0\), \(D(a)\) is a nonempty dense open subset of \(Y\).

(c) Use noetherian induction on closed subsets of \(Y\).  The reductions
in part (a), applied within any closed subset, show that an irreducible
closed target either contains no dense image component or has a dense open
subset \(U\) contained in the image. In the first case the closures
of the images of the finitely many source components are proper closed
subsets. Apply the induction hypothesis to the maps into these closures
and take the finite union of their constructible images.
In the latter case,
\[
f(X)=U\cup
f\bigl(X\setminus f^{-1}(U)\bigr).
\]
The second source is closed in \(X\), its image lies in the proper closed
set \(Y\setminus U\), and its morphism to that closed set is of finite
type.  The induction hypothesis makes the second image constructible.
Thus \(f(X)\) is constructible.  Reducible targets are finite unions of
their irreducible components, so the same conclusion holds there.  Part
(a) then gives the theorem for every constructible subset.

(d) Over an algebraically closed field, consider
\[
\mathbb A^2\longrightarrow\mathbb A^2,\qquad
(x,y)\longmapsto(x,xy).
\]
Its image on variety points is
\[
D(x)\cup\{(0,0)\}.
\]
It is not closed because it contains the dense open \(D(x)\) but is not
all of \(\mathbb A^2\).  It is not open because every open neighborhood of
the origin meets the punctured line
\(\{(0,z):z\ne0\}\), which the image omits.
\end{solution}

\begin{exerciseintro}
% record: H-II-03-020
\exerciseheading{Exercise II.3.20}
\addcontentsline{toc}{subsection}{Exercise II.3.20}

\begin{context}
Let \(A\) be a finitely generated domain over a field \(k\).  We may use
the dimension results from affine algebraic geometry:
\[
\dim A=\operatorname{trdeg}_k\operatorname{Frac}A,
\qquad
\dim A_{\mathfrak p}=\operatorname{ht}\mathfrak p,
\qquad
\operatorname{ht}\mathfrak p+\dim A/\mathfrak p=\dim A.
\]
Every maximal ideal has height \(\dim A\).  Noether normalization gives a
finite injection \(k[x_1,\ldots,x_d]\hookrightarrow A\), with
\(d=\dim A\).  The dimension of a scheme is the supremum of the dimensions
of the affine opens in any affine open cover.
\end{context}

\begin{problem}
Let \(X\) be an integral scheme of finite type over \(k\).

(a) For every closed point \(P\), prove
\(\dim X=\dim\mathcal O_{X,P}\).

(b) Prove
\[
\dim X=\operatorname{trdeg}_kK(X).
\]

(c) For a nonempty closed subset \(Y\subseteq X\), prove
\[
\operatorname{codim}(Y,X)
=\inf_{P\in Y}\dim\mathcal O_{X,P}.
\]

(d) For the nonempty closed subset in (c), prove
\[
\dim Y+\operatorname{codim}(Y,X)=\dim X.
\]

(e) If \(U\subseteq X\) is nonempty and open, prove
\(\dim U=\dim X\).

(f) For a field extension \(k'/k\), prove that every irreducible component
of \(X_{k'}=X\times_k k'\) has dimension \(\dim X\).
\end{problem}
\end{exerciseintro}

\begin{solution}
(a) Choose an affine neighborhood
\(\operatorname{Spec}A\) of \(P\), represented by a maximal ideal
\(\mathfrak m\).  Every nonempty affine open of \(X\) has fraction field
\(K(X)\), so the dimension formula in the context gives the same dimension
\(\operatorname{trdeg}_kK(X)\) on every affine chart.  A finite affine
cover therefore shows
\[
\dim X=\dim A.
\]
Since maximal ideals of the finitely generated domain \(A\) have full
height,
\[
\dim\mathcal O_{X,P}
=\dim A_{\mathfrak m}
=\operatorname{ht}\mathfrak m
=\dim A=\dim X.
\]

(b) On any nonempty affine open \(\operatorname{Spec}A\),
\[
K(X)=\operatorname{Frac}A.
\]
The first equality in the context and part (a)'s affine-cover argument give
\[
\dim X=\dim A
=\operatorname{trdeg}_k\operatorname{Frac}A
=\operatorname{trdeg}_kK(X).
\]

(c) Let \(Y_1,\ldots,Y_r\) be the irreducible components of \(Y\), with
generic points \(\eta_i\).  On an affine neighborhood, \(\eta_i\)
corresponds to a prime \(\mathfrak p_i\), and
\[
\operatorname{codim}(Y_i,X)
=\operatorname{ht}\mathfrak p_i
=\dim\mathcal O_{X,\eta_i}.
\]
If \(P\in Y_i\), its prime contains \(\mathfrak p_i\), so
\(\dim\mathcal O_{X,P}\ge
\dim\mathcal O_{X,\eta_i}\).  Hence the infimum over \(Y\) is attained at
the generic point of a component and equals
\[
\min_i\operatorname{codim}(Y_i,X)
=\operatorname{codim}(Y,X).
\]

(d) For each irreducible component \(Y_i\), the height formula gives
\[
\dim Y_i+\operatorname{codim}(Y_i,X)=\dim X.
\]
Now \(\dim Y\) is the maximum of the \(\dim Y_i\), while
\(\operatorname{codim}(Y,X)\) is the minimum of their codimensions.
The displayed equality shows that these extrema occur for the same
components, and proves the formula for \(Y\).

(e) The open \(U\) is noetherian and therefore quasi-compact, so it remains
of finite type over \(k\).  It is integral and has the same generic point
and function field as \(X\).  Part (b) gives
\[
\dim U=\operatorname{trdeg}_kK(U)
=\operatorname{trdeg}_kK(X)=\dim X.
\]

(f) It is enough to work on an affine open
\(X=\operatorname{Spec}A\), because every irreducible component meets the
base change of an affine open and dimension is unchanged on nonempty open
subsets.  Put \(d=\dim A\) and choose a finite injection
\[
R=k[x_1,\ldots,x_d]\hookrightarrow A.
\]
After scalar extension,
\[
R'=k'[x_1,\ldots,x_d]\longrightarrow
B=A\otimes_k k'
\]
is finite.

We claim that every minimal prime \(\mathfrak q\) of \(B\) contracts to
\((0)\) in \(R'\).  The finite \(R\)-module \(A\) is torsion-free, and it
embeds in a finite free \(R\)-module after clearing denominators in
\(A\otimes_R\operatorname{Frac}R\).  Flat base change to \(R'\) embeds
\(B\) in a finite free \(R'\)-module, so every nonzero element of \(R'\)
is a non-zero-divisor on \(B\).  In a noetherian ring every element of a
minimal prime is a zero-divisor.  Thus
\(\mathfrak q\cap R'=(0)\).

It follows that \(B/\mathfrak q\) is an integral finite extension of
\(R'\).  Integral extensions preserve dimension, so
\[
\dim B/\mathfrak q=\dim R'=d.
\]
These quotients are the irreducible components of \(X_{k'}\), and every
one has dimension \(\dim X\).
\end{solution}

\begin{exerciseintro}
% record: H-II-03-021
\exerciseheading{Exercise II.3.21}
\addcontentsline{toc}{subsection}{Exercise II.3.21}

\begin{context}
Let \(R\) be a discrete valuation ring with uniformizer \(\pi\), fraction
field \(K\), and a coefficient field isomorphic to its residue field
\(k\).  Then \(\dim R[t]=2\), while
\[
R[t]_\pi=K[t]
\]
has dimension \(1\).  A principal prime of a domain has height at most
one, and a nonzero principal prime here has height exactly one.
\end{context}

\begin{problem}
For \(X=\operatorname{Spec}R[t]\), show that conclusions (a), (d), and (e)
of Exercise II.3.20 fail when the finite-type-over-a-field hypothesis is
removed.
\end{problem}
\end{exerciseintro}

\begin{solution}
Consider the principal ideal
\[
\mathfrak p=(\pi t-1)\subseteq R[t].
\]
There is an isomorphism
\[
R[t]/(\pi t-1)\cong R[1/\pi]=K,
\]
so \(\mathfrak p\) is maximal and defines a closed point \(P\in X\).
It is a nonzero principal prime and hence has height one.  Therefore
\[
\dim\mathcal O_{X,P}=1
\quad\text{but}\quad
\dim X=2,
\]
contradicting conclusion (a).

Let \(Y=\{P\}\).  Then \(\dim Y=0\) and
\(\operatorname{codim}(Y,X)=1\), so
\[
\dim Y+\operatorname{codim}(Y,X)=1\ne2=\dim X.
\]
This contradicts (d).

Finally, the nonempty open subset
\[
U=D(\pi)=\operatorname{Spec}K[t]
\]
has dimension \(1\), whereas \(X\) has dimension \(2\).  Thus (e) also
fails.
\end{solution}

\begin{exerciseintro}
% record: H-II-03-022
\exerciseheading{Exercise II.3.22}
\addcontentsline{toc}{subsection}{Exercise II.3.22}

\begin{context}
Let \(f:X\to Y\) be a dominant morphism of integral finite-type
\(k\)-schemes, and put \(e=\dim X-\dim Y\).  We may use:

\[
\operatorname{trdeg}_k k(x)
=\dim X-\dim\mathcal O_{X,x}
\]
for \(x\in X\), and the analogous formula on \(Y\).  In a noetherian
local ring of dimension \(d\), there are \(d\) elements whose generated
ideal has maximal-ideal radical.  Krull's height theorem bounds the height
of a minimal prime over an ideal generated by \(d\) elements by \(d\).
Finite morphisms and open immersions are stable under base change, and
finite-type images of constructible sets are constructible.
\end{context}

\begin{problem}
(a) Let \(Y'\subseteq Y\) be irreducible closed with generic point
\(\eta'\in f(X)\).  If \(Z\) is an irreducible component of
\(f^{-1}(Y')\) whose image contains \(\eta'\), prove
\[
\operatorname{codim}(Z,X)\le
\operatorname{codim}(Y',Y).
\]

(b) For \(y\in f(X)\), prove that every irreducible component of \(X_y\)
has dimension at least \(e\).

(c) Prove that there is a dense open \(U\subseteq X\) such that
\(\dim U_y=e\) for every \(y\in f(U)\).

(d) For an integer \(h\), let \(E_h\) consist of points \(x\in X\) lying
on an irreducible component of \(X_{f(x)}\) of dimension at least \(h\).
Prove that \(E_e=X\); that \(E_h\) is not dense for \(h>e\); and that every
\(E_h\) is closed.

(e) Let
\[
C_h=\{y\in Y:\dim X_y=h\}.
\]
Prove that every \(C_h\) is constructible and that \(C_e\) contains a
dense open subset of \(Y\).
\end{problem}
\end{exerciseintro}

\begin{solution}
(a) Let \(\zeta\) be the generic point of \(Z\).  The hypothesis that
\(\eta'\in f(Z)\) forces \(f(\zeta)=\eta'\).  Work in affine
neighborhoods and let \(\mathfrak p\) and \(\mathfrak q\) be the primes
corresponding to \(\eta'\) and \(\zeta\).  Since \(Z\) is an irreducible
component of the inverse image, \(\mathfrak q\) is minimal over
\(\mathfrak p\mathcal O_X\).

Put \(d=\dim\mathcal O_{Y,\eta'}\).  Choose
\(a_1,\ldots,a_d\) in that local ring whose ideal has radical its maximal
ideal.  In \(\mathcal O_{X,\zeta}\), the radical of the extended ideal is
the maximal ideal: a smaller prime containing it would give a prime
strictly below \(\mathfrak q\) over \(\mathfrak p\), contradicting the
minimality of \(\mathfrak q\).  Krull's height theorem therefore gives
\[
\dim\mathcal O_{X,\zeta}\le d.
\]
Using the codimension formula for generic points,
\[
\operatorname{codim}(Z,X)
=\dim\mathcal O_{X,\zeta}
\le\dim\mathcal O_{Y,\eta'}
=\operatorname{codim}(Y',Y).
\]

(b) Let \(T\) be an irreducible component of \(X_y\), with generic point
\(x\).  The closure \(\overline{\{x\}}\) is an irreducible component of
\(f^{-1}(\overline{\{y\}})\) that dominates
\(\overline{\{y\}}\).  Part (a) gives
\[
\dim\mathcal O_{X,x}\le\dim\mathcal O_{Y,y}.
\]
Since \(T\) is an integral finite-type scheme over \(k(y)\),
\[
\begin{aligned}
\dim T
&=\operatorname{trdeg}_{k(y)}k(x)\\
&=\operatorname{trdeg}_k k(x)-\operatorname{trdeg}_k k(y)\\
&=(\dim X-\dim\mathcal O_{X,x})
  -(\dim Y-\dim\mathcal O_{Y,y})\\
&\ge\dim X-\dim Y=e.
\end{aligned}
\]

(c) Shrink to dense affine opens in \(X\) and \(Y\), without changing
their dimensions.  Choose
\(t_1,\ldots,t_e\in K(X)\) forming a transcendence basis over \(K(Y)\);
after shrinking once more, take them to be regular on \(X\).  They define
a dominant, generically finite morphism
\[
g:X\longrightarrow\mathbb A_Y^e.
\]
By the generically finite result, there is a dense open
\(W\subseteq\mathbb A_Y^e\) such that
\[
g^{-1}(W)\longrightarrow W
\]
is finite.  Dominance and closedness of a finite morphism allow \(W\) to
be chosen so that this map is surjective.  Put \(U=g^{-1}(W)\).

For \(y\in f(U)\), base change gives a finite surjection
\[
U_y\longrightarrow W_y.
\]
The latter is a nonempty open subset of
\(\mathbb A_{k(y)}^e\), so it has dimension \(e\).  A finite morphism
makes every component of \(U_y\) have dimension at most \(e\), while part
(b), or its application to the open \(U\), gives the reverse inequality.
Thus \(\dim U_y=e\).

(d) Part (b) immediately gives \(E_e=X\), and also \(E_h=X\) for
\(h\le e\).  For \(h>e\), the open \(U\) from part (c) is disjoint from
\(E_h\): a component of \(X_y\) meeting \(U_y\) has a nonempty open subset
inside a component of \(U_y\), and therefore has dimension \(e\).  Hence
\(E_h\) is not dense.

We prove closedness by noetherian induction on \(X\).  For \(h>e\), put
\(F=X\setminus U\).  A fiber component of dimension at least \(h\) cannot
meet \(U\), so
\[
E_h(X\to Y)=E_h(F\to Y).
\]
Write the reduced closed scheme \(F\) as the finite union of its
irreducible components \(F_i\), and let \(Y_i\) be the reduced closure of
the image of \(F_i\).  Then \(F_i\to Y_i\) is dominant and
\(\dim F_i<\dim X\).  By induction, every
\(E_h(F_i\to Y_i)\) is closed in \(F_i\), and
\[
E_h(X\to Y)=\bigcup_iE_h(F_i\to Y_i).
\]
This finite union is closed in \(X\).  Together with the case \(h\le e\),
this proves closedness for all \(h\).

(e) Put
\[
F_h=\{y\in Y:\dim X_y\ge h\}.
\]
By definition and part (d),
\[
F_h=f(E_h).
\]
Chevalley's theorem makes \(F_h\) constructible, and hence
\[
C_h=F_h\setminus F_{h+1}
\]
is constructible.  The generic fiber has function field \(K(X)\) over
\(K(Y)\), so
\[
\dim X_\eta
=\operatorname{trdeg}_{K(Y)}K(X)
=\dim X-\dim Y=e.
\]
Thus the generic point \(\eta\) of \(Y\) belongs to \(C_e\).  A
constructible subset of an irreducible Zariski space that contains the
generic point contains a nonempty, hence dense, open subset.  Therefore
\(C_e\) contains a dense open subset of \(Y\).
\end{solution}

\begin{exerciseintro}
% record: H-II-03-023
\exerciseheading{Exercise II.3.23}
\addcontentsline{toc}{subsection}{Exercise II.3.23}

\begin{context}
For affine \(k\)-varieties \(V\) and \(W\),
\[
A(V\times_k W)\cong A(V)\otimes_k A(W).
\]
The functor \(t\) sends an affine variety with coordinate ring \(A\) to
\(\operatorname{Spec}A\).  Products and the functor \(t\) are obtained
globally by gluing their affine-chart constructions.
\end{context}

\begin{problem}
Let \(V,W\) be varieties over an algebraically closed field \(k\).  Prove
that
\[
t(V\times_k W)
\cong t(V)\times_{\operatorname{Spec}k}t(W).
\]
\end{problem}
\end{exerciseintro}

\begin{solution}
First suppose \(V\) and \(W\) are affine varieties, with coordinate
rings \(A=A(V)\) and \(B=A(W)\). Then
\[
\begin{aligned}
t(V\times_k W)
&\cong\operatorname{Spec}(A\otimes_k B)\\
&\cong\operatorname{Spec}A
\times_{\operatorname{Spec}k}\operatorname{Spec}B\\
&=t(V)\times_{\operatorname{Spec}k}t(W).
\end{aligned}
\]
The isomorphism is natural in restrictions to affine open subsets.

For general \(V\) and \(W\), the products of affine charts
\(V_i\times_k W_j\) cover \(V\times_k W\), while
\(t(V_i)\times_{\operatorname{Spec}k}t(W_j)\) cover the scheme fiber
product.  The affine isomorphisms above agree on all overlaps by their
naturality, so they glue to the asserted global isomorphism.
\end{solution}

\begin{exerciseintro}
\corpussection{II.4 Separated and Proper Morphisms}
\addcontentsline{toc}{section}{II.4 Separated and Proper Morphisms}

% record: H-II-04-001
\exerciseheading{Exercise II.4.1}
\addcontentsline{toc}{subsection}{Exercise II.4.1}

\begin{context}
A morphism is proper if it is separated, of finite type, and universally
closed.  Finite morphisms are affine, are stable under arbitrary base
change, and correspond locally to ring maps for which the target is a
finite module.  An integral extension satisfies lying over.
\end{context}

\begin{problem}
Prove that every finite morphism of schemes is proper.
\end{problem}
\end{exerciseintro}

\begin{solution}
Let \(f:X\to Y\) be finite.  It is of finite type.  It is also separated:
over an affine open \(V=\operatorname{Spec}A\subseteq Y\), write
\(f^{-1}(V)=\operatorname{Spec}B\).  The restriction of the diagonal is
the morphism associated to the surjection
\[
B\otimes_A B\longrightarrow B,\qquad b_1\otimes b_2\longmapsto b_1b_2,
\]
so it is a closed immersion.  Since this can be checked over an open
cover of \(Y\), the diagonal of \(f\) is a closed immersion.

It remains to prove universal closedness.  Because every base change of
\(f\) is finite, it is enough to show that a finite morphism is closed.
Again work with a finite ring map \(A\to B\).  If \(J\subseteq B\), set
\[
I=\ker(A\longrightarrow B/J).
\]
The induced injection \(A/I\hookrightarrow B/J\) is integral, so lying
over gives a surjection
\[
\operatorname{Spec}(B/J)\longrightarrow\operatorname{Spec}(A/I).
\]
Consequently the image of the closed subset \(V(J)\subseteq\operatorname{Spec}B\)
is exactly \(V(I)\), hence is closed.  Thus \(f\), and every base change
of \(f\), is closed.  Therefore \(f\) is proper.
\end{solution}

\begin{exerciseintro}
% record: H-II-04-002
\exerciseheading{Exercise II.4.2}
\addcontentsline{toc}{subsection}{Exercise II.4.2}

\begin{context}
For two \(S\)-morphisms \(f,g:X\to Y\), their equalizer is the inverse
image of the diagonal under
\[
(f,g):X\longrightarrow Y\times_S Y.
\]
If \(Y\to S\) is separated, this equalizer is a closed subscheme of
\(X\).  A closed subscheme whose underlying set is all of a reduced
scheme is the whole scheme.
\end{context}

\begin{problem}
Let \(X\) be a reduced \(S\)-scheme and \(Y\) a separated \(S\)-scheme.
Suppose two \(S\)-morphisms \(f,g:X\to Y\) agree on an open dense subset
of \(X\).  Prove that \(f=g\).  Give examples showing that the conclusion
can fail if either \(X\) is nonreduced or \(Y\) is nonseparated.
\end{problem}
\end{exerciseintro}

\begin{solution}
Let
\[
E=X\times_{(f,g),\,Y\times_S Y,\,\Delta_Y}Y.
\]
Because \(Y\to S\) is separated, \(\Delta_Y\) is a closed immersion, so
\(E\hookrightarrow X\) is a closed immersion.  The dense open set on
which \(f\) and \(g\) agree is contained in \(\lvert E\rvert\).
Therefore \(\lvert E\rvert=X\).  On an affine open
\(\operatorname{Spec}A\subseteq X\), the ideal defining \(E\) has
vanishing set all of \(\operatorname{Spec}A\), and hence is contained in
the nilradical of \(A\).  Since \(X\) is reduced, that ideal is zero.
Thus \(E=X\), which says precisely that \(f=g\).

For failure without reducedness, take
\[
X=\operatorname{Spec}k[x,\varepsilon]/(\varepsilon^2,x\varepsilon),
\qquad Y=\mathbb A_k^1,
\]
and let \(f,g:X\to Y\) correspond respectively to
\[
t\longmapsto x,\qquad t\longmapsto x+\varepsilon.
\]
They are distinct, but on the dense open \(D(x)\) one has
\(\varepsilon=0\), so their restrictions agree.

For failure without separatedness, let \(Y\) be the affine line with
doubled origin, obtained by gluing two copies \(Y_1,Y_2\) of
\(\mathbb A_k^1\) along their punctured lines.  The two open immersions
\[
\mathbb A_k^1=Y_i\longrightarrow Y
\]
agree on \(\mathbb A_k^1\setminus\{0\}\) but send the origin to the two
different origins.  Here \(X=\mathbb A_k^1\) is reduced, while \(Y\) is
not separated.
\end{solution}

\begin{exerciseintro}
% record: H-II-04-003
\exerciseheading{Exercise II.4.3}
\addcontentsline{toc}{subsection}{Exercise II.4.3}

\begin{context}
The diagonal of a separated morphism is a closed immersion.  Closed
subschemes of affine schemes are affine, and a fiber product of affine
schemes over an affine scheme is affine.
\end{context}

\begin{problem}
Let \(X\) be separated over an affine scheme \(S\), and let \(U,V\) be
affine open subsets of \(X\).  Prove that \(U\cap V\) is affine.  Give an
example showing that separatedness is necessary.
\end{problem}
\end{exerciseintro}

\begin{solution}
The inverse image of \(U\times_S V\) under the diagonal
\[
\Delta_{X/S}:X\longrightarrow X\times_S X
\]
is \(U\cap V\): a morphism to \(X\) factors through this inverse image
exactly when it factors through both \(U\) and \(V\).  Thus there is a
cartesian square
\[
\begin{array}{ccc}
U\cap V&\longrightarrow&U\times_S V\\
\big\downarrow&&\big\downarrow\\
X&\xrightarrow{\Delta_{X/S}}&X\times_S X.
\end{array}
\]
The lower map is a closed immersion.  Hence the upper map is a closed
immersion.  Since \(S,U,V\) are affine, \(U\times_S V\) is affine, and
therefore so is \(U\cap V\).

For a counterexample, glue two copies \(U,V\) of \(\mathbb A_k^2\) by
the identity on \(\mathbb A_k^2\setminus\{(0,0)\}\).  The resulting
scheme \(X\) is not separated and
\[
U\cap V\cong\mathbb A_k^2\setminus\{(0,0)\}.
\]
This intersection is not affine.  Indeed, with coordinates \(x,y\),
\[
\Gamma(\mathbb A_k^2\setminus\{0\},\mathcal O)
=k[x,y]_x\cap k[x,y]_y=k[x,y]
\]
inside \(k(x,y)\).  If the punctured plane were affine, its inclusion
in \(\mathbb A_k^2=\operatorname{Spec}k[x,y]\) would, after identifying
both rings of global functions, be an isomorphism, which is impossible
because the origin is missing.
\end{solution}

\begin{exerciseintro}
% record: H-II-04-004
\exerciseheading{Exercise II.4.4}
\addcontentsline{toc}{subsection}{Exercise II.4.4}

\begin{context}
The graph of a morphism to a separated \(S\)-scheme is a closed
immersion.  Proper morphisms are stable under base change and
composition.  For a quasi-compact morphism into a noetherian scheme, the
underlying set of its scheme-theoretic image is the closure of its
set-theoretic image.
\end{context}

\begin{problem}
Let \(f:X\to Y\) be an \(S\)-morphism of separated schemes of finite type over
a noetherian scheme \(S\).  Let \(Z\subseteq X\) be a closed subscheme
proper over \(S\).  Prove that \(f(Z)\) is closed in \(Y\), and that its
scheme-theoretic image, with the image subscheme structure, is proper
over \(S\).
\end{problem}
\end{exerciseintro}

\begin{solution}
Write \(q:Z\to Y\) for the restriction of \(f\).  Its graph gives a
factorization
\[
Z\xrightarrow{\Gamma_q}Z\times_S Y
\xrightarrow{\operatorname{pr}_2}Y.
\]
Because \(Y\to S\) is separated, \(\Gamma_q\) is a closed immersion.
The second map is the base change of the proper morphism \(Z\to S\).
Thus \(q\) is proper.  In particular it is closed, so \(q(Z)=f(Z)\) is
closed in \(Y\).

Let \(W\hookrightarrow Y\) be the scheme-theoretic image of \(q\).
Since \(q(Z)\) is already closed, the underlying space of \(W\) is
exactly \(q(Z)\).  To see this locally, for
\(V=\operatorname{Spec}A\subseteq Y\), the subscheme \(W\cap V\) is
defined by
\[
\ker\bigl(A\longrightarrow
\Gamma(q^{-1}(V),\mathcal O_Z)\bigr),
\]
and \(q(q^{-1}(V))\) is dense in that affine closed subscheme.  Its
global image is closed, so equality follows.  Hence the induced map
\(p:Z\to W\) is surjective.

The morphism \(W\to S\) is separated and of finite type because
\(W\hookrightarrow Y\) is a closed immersion.  It remains to prove
universal closedness.  Let \(T\to S\) be arbitrary and let
\(C\subseteq W_T\) be closed.  Surjectivity is preserved by base
change, so \(p_T:Z_T\to W_T\) is surjective.  Therefore the image of
\(C\) in \(T\) equals the image of the closed subset
\(p_T^{-1}(C)\subseteq Z_T\).  Since \(Z_T\to T\) is closed, this image
is closed.  Thus every base change of \(W\to S\) is closed, and
\(W\to S\) is proper.
\end{solution}

\begin{exerciseintro}
% record: H-II-04-005
\exerciseheading{Exercise II.4.5}
\addcontentsline{toc}{subsection}{Exercise II.4.5}

\begin{context}
Let \(X\) be an integral scheme of finite type over a field \(k\), with
function field \(K\).  A valuation ring \(R\) of \(K/k\) has center
\(x\in X\) precisely when
\[
\mathcal O_{X,x}\subseteq R,\qquad
\mathfrak m_R\cap\mathcal O_{X,x}=\mathfrak m_x.
\]
Equivalently, the generic-point morphism
\(\operatorname{Spec}K\to X\) extends to
\(\operatorname{Spec}R\to X\) with closed point mapping to \(x\).

We may use the valuative criteria and the following form of Chow's
lemma: a separated integral finite-type \(k\)-scheme \(X\) admits a
projective, surjective, birational morphism \(X'\to X\), with \(X'\)
quasi-projective over \(k\).
At a specialization one may choose a valuation ring of the relevant function field dominating the local domain there.
\end{context}

\begin{problem}
(a) If \(X\) is separated over \(k\), prove that a valuation of \(K/k\)
has at most one center on \(X\).

(b) If \(X\) is proper over \(k\), prove that every valuation of \(K/k\)
has exactly one center on \(X\).

(c) Prove the converses of (a) and (b).

(d) If \(X\) is proper over an algebraically closed field \(k\), prove
that \(\Gamma(X,\mathcal O_X)=k\).
\end{problem}
\end{exerciseintro}

\begin{solution}
(a) A center is the closed-point image of an extension in the diagram
\[
\begin{array}{ccc}
\operatorname{Spec}K&\longrightarrow&X\\
\big\downarrow&&\big\downarrow\\
\operatorname{Spec}R&\longrightarrow&\operatorname{Spec}k.
\end{array}
\]
The uniqueness part of the valuative criterion for separatedness gives
at most one such extension, hence at most one center.

(b) Properness gives both existence and uniqueness in the valuative
criterion.  The displayed diagram therefore has exactly one extension,
so \(R\) has exactly one center on \(X\).

(c) First suppose every valuation ring of \(K/k\) has at most one center.
We show that the image of
\(\Delta:X\to X\times_kX\) is closed.  It is irreducible with generic
point \(\Delta(\eta)\), where \(\eta\) is the generic point of \(X\).
Let \(z\) be a specialization of \(\Delta(\eta)\).  The specialization
lemma supplies a valuation ring \(R\subseteq K\) and a morphism
\[
\operatorname{Spec}R\longrightarrow X\times_kX
\]
whose generic point maps to \(\Delta(\eta)\) and whose closed point maps
to \(z\).  The two projections give two extensions of the same generic
map \(\operatorname{Spec}K\to X\).  Their closed-point images are centers
of \(R\), so they are equal by hypothesis.  Once the center is fixed,
the extension is unique: on any affine neighborhood its ring map is
the inclusion into \(R\subseteq K\).  Thus the two projections coincide,
the morphism factors through \(\Delta\), and \(z\in\Delta(X)\).
Therefore \(\Delta(X)\) is closed.  Since a diagonal is an immersion, it
is a closed immersion, and \(X\) is separated.

Now assume every valuation ring of \(K/k\) has exactly one center.  We
already know that \(X\) is separated.  Apply Chow's lemma to obtain a
projective, surjective, birational morphism
\[
p:X'\longrightarrow X
\]
with \(X'\) integral and quasi-projective.  Embed \(X'\) as an open
subscheme of an integral projective scheme \(\overline X'\).
We claim that \(X'=\overline X'\).  Otherwise choose
\(z\in\overline X'\setminus X'\) and a valuation ring \(R\) of \(K/k\)
dominating \(\mathcal O_{\overline X',z}\).  By hypothesis \(R\) has a
center on \(X\).  Properness of \(p\) then lifts this center uniquely to
a center on \(X'\).  But \(\overline X'\) is separated, so \(R\) cannot
have both that center and the distinct center \(z\) on \(\overline X'\).
This contradiction proves the claim.  Hence \(X'\) is projective, and
in particular proper over \(k\).

It remains to descend universal closedness along the surjection \(p\).
For any \(k\)-scheme \(T\) and any closed \(C\subseteq X_T\), the inverse
image \(p_T^{-1}(C)\) is closed in the proper \(T\)-scheme \(X'_T\).
Its image in \(T\) is closed, and, because \(p_T\) is surjective, this
image equals the image of \(C\).  Thus \(X_T\to T\) is closed for every
\(T\).  Since \(X\) is separated and of finite type, it is proper.

(d) A section \(a\in\Gamma(X,\mathcal O_X)\) defines a morphism
\[
u:X\longrightarrow\mathbb A_k^1\hookrightarrow\mathbb P_k^1.
\]
Because \(X\) is proper and \(\mathbb P_k^1\) is separated over \(k\),
the graph factorization shows that \(u:X\to\mathbb P_k^1\) is proper.
Its image is therefore closed.  It is irreducible because \(X\) is
integral, and it cannot be all of \(\mathbb P_k^1\), since it misses
infinity.  Hence its image is one closed point.  As \(k\) is
algebraically closed, that point is \(k\)-rational, so \(a\) is a
constant in \(k\).  The reverse inclusion is automatic.
\end{solution}

\begin{exerciseintro}
% record: H-II-04-006
\exerciseheading{Exercise II.4.6}
\addcontentsline{toc}{subsection}{Exercise II.4.6}

\begin{context}
A finitely generated algebra that is integral over its base ring is
finite as a module.  Surjective morphisms of schemes remain surjective
after base change.
\end{context}

\begin{problem}
Let \(f:X\to Y\) be a proper morphism of affine varieties over a field
\(k\).  Prove that \(f\) is finite.
\end{problem}
\end{exerciseintro}

\begin{solution}
Write \(Y=\operatorname{Spec}A\), \(X=\operatorname{Spec}B\).  The
\(A\)-algebra \(B\) is finitely generated because \(f\) is of finite
type.  We prove that each \(b\in B\) is integral over \(A\).

Base-changing \(f\) along
\(\operatorname{Spec}A[T]\to\operatorname{Spec}A\) gives a closed map
\[
\operatorname{Spec}B[T]\longrightarrow\operatorname{Spec}A[T].
\]
Inside its source consider
\[
Z=V(bT-1)=\operatorname{Spec}\bigl(B[T]/(bT-1)\bigr).
\]
Its image is closed.  Put
\[
J=\ker\bigl(A[T]\longrightarrow B[T]/(bT-1)\bigr).
\]
The image of \(Z\) is dense in \(V(J)\) by the definition of \(J\), and
it is closed, hence it is all of \(V(J)\).  Thus
\(Z\to V(J)\) is surjective.

After base change by \(V(T)\hookrightarrow\operatorname{Spec}A[T]\),
the source is empty because \(bT-1\) and \(T\) generate the unit ideal.
Surjectivity after base change therefore shows
\[
\operatorname{Spec}\bigl(A[T]/(J,T)\bigr)=\varnothing,
\qquad\text{so}\qquad J+(T)=A[T].
\]
Consequently \(J\) contains a polynomial
\[
F(T)=1+a_1T+\cdots+a_nT^n.
\]
In \(B_b\), substituting \(T=b^{-1}\) gives
\[
b^n+a_1b^{n-1}+\cdots+a_n=0.
\]
The equality in \(B_b\) means that, for some \(N\geq0\),
\[
b^{N+n}+a_1b^{N+n-1}+\cdots+a_nb^N=0
\]
in \(B\).  This is a monic equation for \(b\), so \(b\) is integral over
\(A\).  Therefore \(B\) is both integral and finitely generated as an
\(A\)-algebra, hence finite as an \(A\)-module.  Thus \(f\) is finite.
\end{solution}

\begin{exerciseintro}
% record: H-II-04-007
\exerciseheading{Exercise II.4.7}
\addcontentsline{toc}{subsection}{Exercise II.4.7}

\begin{context}
For a complex scheme obtained by base change from \(\mathbb R\), complex
conjugation induces a semilinear involution.  We may use elementary
faithfully flat descent for the finite Galois extension
\(\mathbb C/\mathbb R\): affine schemes, morphisms, open immersions,
finite type, and separatedness descend.  If a complex algebra \(A\)
has a conjugate-linear involution \(\sigma\), then
\[
A^\sigma\otimes_{\mathbb R}\mathbb C\xrightarrow{\ \sim\ }A,
\]
because every \(a\in A\) has the unique decomposition
\[
a=\frac{a+\sigma(a)}2
+i\,\frac{a-\sigma(a)}{2i}.
\]
\end{context}

\begin{problem}
Let \(X\) be a separated finite-type scheme over \(\mathbb C\), equipped
with a semilinear involution \(\sigma\).  Assume that any two points of
\(X\) lie in a common affine open.

(a) Prove that \(X\) has a unique separated finite-type real form
\(X_0\), compatible with \(\sigma\).

For such real forms prove:

(b) \(X_0\) is affine if and only if \(X\) is affine.

(c) Morphisms \(X_0\to Y_0\) correspond bijectively to complex
morphisms \(X\to Y\) commuting with the involutions.

(d) If \(X\cong\mathbb A_{\mathbb C}^1\), then
\(X_0\cong\mathbb A_{\mathbb R}^1\).

(e) If \(X\cong\mathbb P_{\mathbb C}^1\), then \(X_0\) is isomorphic
either to \(\mathbb P_{\mathbb R}^1\) or to the conic
\[
x_0^2+x_1^2+x_2^2=0
\]
in \(\mathbb P_{\mathbb R}^2\).
\end{problem}
\end{exerciseintro}

\begin{solution}
(a) For each \(x\in X\), choose an affine open \(W\) containing both
\(x\) and \(\sigma(x)\).  The intersection
\[
U=W\cap\sigma(W)
\]
is affine because \(X\) is separated, and it is invariant under
\(\sigma\).  Thus invariant affine opens form a cover of \(X\).

Write such an open as \(U=\operatorname{Spec}A\), and put
\(A_0=A^\sigma\).  The isomorphism in the context identifies
\[
U\cong\operatorname{Spec}A_0\times_{\mathbb R}\mathbb C
\]
and identifies \(\sigma\) with conjugation on the second factor.
Intersections of two invariant affine opens are again invariant and
affine.  The corresponding real affine schemes glue: the overlap maps
are obtained by faithfully flat descent from the complex open
immersions.  Denote the resulting real scheme by \(X_0\).  Its base
change to \(\mathbb C\) is \(X\), with the prescribed involution.
Finite type and separatedness descend along the faithfully flat map
\(\operatorname{Spec}\mathbb C\to\operatorname{Spec}\mathbb R\).

The same affine calculation shows uniqueness.  On every invariant
affine, a compatible real form must be
\(\operatorname{Spec}(A^\sigma)\), and these identifications agree on
overlaps.  Hence any two forms are uniquely isomorphic in a way inducing
the given identification after complex base change.

(b) If \(X_0\) is affine, its base change \(X\) is affine.  Conversely,
if \(X=\operatorname{Spec}A\), the whole scheme is an invariant affine,
so part (a) constructs the form
\[
\operatorname{Spec}(A^\sigma).
\]
Uniqueness identifies this form with \(X_0\), proving that \(X_0\) is
affine.

(c) Base change sends a real morphism to an equivariant complex
morphism.  Conversely, on invariant affine opens an equivariant map is
given by a semilinear-compatible ring homomorphism, which restricts to
a homomorphism between invariant real subrings.  These restrictions
glue and give a real morphism.  Faithful flatness shows that the two
constructions are inverse, so the correspondence is bijective.

(d) Choose a coordinate \(t\) on \(X\).  A semilinear automorphism of
\(\mathbb C[t]\) has the form
\[
\sigma(t)=at+b,\qquad a\in\mathbb C^\times,\quad b\in\mathbb C.
\]
The condition \(\sigma^2=1\) says
\[
\overline a\,a=1,\qquad \overline a\,b+\overline b=0.
\]
Choose \(c\in\mathbb C^\times\) with \(c/\overline c=a\).  Then
\(\overline c\,b\) is purely imaginary, so choose \(d\in\mathbb C\)
with \(d-\overline d=\overline c\,b\).  For \(u=ct+d\) one obtains
\(\sigma(u)=u\).  Hence
\[
\mathbb C[t]^\sigma=\mathbb R[u],
\]
and part (b) gives \(X_0\cong\mathbb A_{\mathbb R}^1\).

(e) First suppose \(\sigma\) fixes a closed complex point of
\(\mathbb P_{\mathbb C}^1\).  Move it to infinity by a projective change
of coordinates.  Its affine complement is invariant.  The coordinate
change constructed in (d) is affine-linear, so it extends across
infinity to a projective coordinate \(u\) satisfying
\(\sigma(u)=u\).  In this coordinate the involution is ordinary
conjugation on all of \(\mathbb P_{\mathbb C}^1\), and hence
\(X_0\cong\mathbb P_{\mathbb R}^1\).

Now suppose there is no fixed closed point.  Choose a closed complex point \(p\), and choose a
coordinate taking \(p\) to \(0\) and \(\sigma(p)\) to infinity.  In this
coordinate the action on complex points is
\[
z\longmapsto\frac{c}{\overline z}
\]
for some \(c\in\mathbb R^\times\).  It has a fixed complex point exactly when
\(\lvert z\rvert^2=c\), so the fixed-point-free case has \(c<0\).
Rescaling the coordinate reduces the action to
\[
z\longmapsto-\frac1{\overline z}.
\]

The morphism
\[
\begin{array}{rcl}
\mathbb P_{\mathbb C}^1&\longrightarrow&
V(x_0^2+x_1^2+x_2^2)\subseteq\mathbb P_{\mathbb C}^2,\\
{[s:t]}&\longmapsto&
[s^2-t^2:i(s^2+t^2):2st]
\end{array}
\]
is an isomorphism.  Complex conjugation on the conic corresponds under
it to
\[
[s:t]\longmapsto[-\overline t:\overline s],
\]
which is \(z\mapsto-1/\overline z\).  Therefore its real descent is
exactly the conic
\(x_0^2+x_1^2+x_2^2=0\).  These two cases exhaust all involutions and
give the asserted classification.
\end{solution}

\begin{exerciseintro}
% record: H-II-04-008
\exerciseheading{Exercise II.4.8}
\addcontentsline{toc}{subsection}{Exercise II.4.8}

\begin{context}
For \(X\to Y\to Z\), the graph \(X\to X\times_ZY\) is a base
change of \(\Delta_{Y/Z}\), and is a closed immersion if \(Y\to Z\)
is separated.
\end{context}

\begin{problem}
Let \(\mathcal P\) be a property such that:

(a) closed immersions have \(\mathcal P\);

(b) composites of morphisms having \(\mathcal P\) have \(\mathcal P\);

(c) base changes of morphisms having \(\mathcal P\) have
\(\mathcal P\).

Under these assumptions, prove:

(d) a product of morphisms having \(\mathcal P\) has \(\mathcal P\);

(e) if \(f:X\to Y\), \(g:Y\to Z\), \(g\circ f\) has \(\mathcal P\),
and \(g\) is separated, then \(f\) has \(\mathcal P\);

(f) if \(f:X\to Y\) has \(\mathcal P\), then the induced morphism
\(f_{\mathrm{red}}:X_{\mathrm{red}}\to Y_{\mathrm{red}}\) has
\(\mathcal P\).
\end{problem}
\end{exerciseintro}

\begin{solution}
(d) Let \(f:X\to Y\) and \(f':X'\to Y'\) be \(S\)-morphisms having
\(\mathcal P\).  Factor their product as
\[
X\times_SX'\longrightarrow Y\times_SX'
\longrightarrow Y\times_SY'.
\]
The first arrow is the base change of \(f\), and the second is the base
change of \(f'\).  They have \(\mathcal P\) by (c), so their composite
has \(\mathcal P\) by (b).

(e) Factor \(f\) as
\[
X\xrightarrow{\gamma}X\times_ZY
\xrightarrow{\operatorname{pr}_2}Y,\qquad
\gamma(x)=(x,f(x)).
\]
The second map is the base change along \(g\) of
\(g\circ f:X\to Z\), and hence has \(\mathcal P\).  The first map is the
base change of the diagonal
\(\Delta_g:Y\to Y\times_ZY\); since \(g\) is separated, it is a closed
immersion and has \(\mathcal P\).  Their composite \(f\) therefore has
\(\mathcal P\).

(f) The composite
\[
X_{\mathrm{red}}\hookrightarrow X\xrightarrow{f}Y
\]
has \(\mathcal P\) by (a) and (b).  Because \(X_{\mathrm{red}}\) is
reduced, this composite factors uniquely as
\[
X_{\mathrm{red}}\xrightarrow{f_{\mathrm{red}}}
Y_{\mathrm{red}}\hookrightarrow Y.
\]
The second arrow is a closed immersion, hence separated.  Apply (e) to
this factorization: the composite has \(\mathcal P\), so
\(f_{\mathrm{red}}\) has \(\mathcal P\).
\end{solution}

\begin{exerciseintro}
% record: H-II-04-009
\exerciseheading{Exercise II.4.9}
\addcontentsline{toc}{subsection}{Exercise II.4.9}

\begin{context}
A morphism \(X\to Y\) is projective if it factors as a closed immersion
into \(\mathbb P_Y^n\) followed by the projection.  Projective morphisms
are stable under base change.  The Segre morphism is a closed immersion
\[
\mathbb P_Z^m\times_Z\mathbb P_Z^n
\hookrightarrow\mathbb P_Z^{(m+1)(n+1)-1}.
\]
\end{context}

\begin{problem}
Prove that a composite of projective morphisms is projective.  Deduce
that projective morphisms satisfy properties (a)--(f) of Exercise
II.4.8.
\end{problem}
\end{exerciseintro}

\begin{solution}
Let \(f:X\to Y\) and \(g:Y\to Z\) be projective.  Choose closed
immersions
\[
i:X\hookrightarrow\mathbb P_Y^m,\qquad
j:Y\hookrightarrow\mathbb P_Z^n.
\]
Base-changing \(j\) along \(\mathbb P_Z^m\to Z\) gives a closed
immersion
\[
\mathbb P_Y^m
\hookrightarrow
\mathbb P_Z^m\times_Z\mathbb P_Z^n.
\]
Composing this with \(i\), and then with the Segre embedding, gives a
closed immersion
\[
X\hookrightarrow
\mathbb P_Z^m\times_Z\mathbb P_Z^n
\hookrightarrow
\mathbb P_Z^{(m+1)(n+1)-1}.
\]
Its projection to \(Z\) is \(g\circ f\), so \(g\circ f\) is projective.

Now take \(\mathcal P\) to mean projective.  A closed immersion is
projective by using \(\mathbb P_Y^0=Y\), so property (a) holds.  The
result just proved is property (b), and projectivity is stable under
base change, giving (c).  Exercise II.4.8 then supplies properties
(d), (e), and (f).
\end{solution}

\begin{exerciseintro}
% record: H-II-04-010
\exerciseheading{Exercise II.4.10}
\addcontentsline{toc}{subsection}{Exercise II.4.10}

\begin{context}
A finite-type scheme over a noetherian scheme is noetherian.  A locally
closed subscheme of a projective \(S\)-scheme is quasi-projective over
\(S\).  A quasi-compact morphism has a scheme-theoretic image that can
be computed locally on the target.  Finite products of projective
\(S\)-schemes are projective.
\end{context}

\begin{problem}
Let \(X\) be proper over a noetherian scheme \(S\).  Prove Chow's lemma:
there are an \(S\)-projective scheme \(X'\), a morphism \(g:X'\to X\),
and an open dense subset \(U\subseteq X\) such that
\[
g^{-1}(U)\longrightarrow U
\]
is an isomorphism.  Proceed by:

(a) reducing to irreducible \(X\);

(b) finding a finite cover \(X=\bigcup U_i\) by \(S\)-quasi-projective
opens, with open immersions \(U_i\hookrightarrow P_i\) into projective
\(S\)-schemes;

(c) for \(U=\bigcap U_i\), taking the scheme-theoretic closure \(X'\) of
the map
\[
U\longrightarrow X\times_SP_1\times_S\cdots\times_SP_n
\]
and proving that the projection \(h:X'\to\prod_iP_i\) is a closed
immersion;

(d) proving \(g^{-1}(U)\to U\) is an isomorphism.
\end{problem}
\end{exerciseintro}

\begin{solution}
(a) Let \(Z_1,\ldots,Z_r\) be the irreducible components of \(X\).
Give
\[
V_i=X\setminus\bigcup_{j\ne i}Z_j
\]
its original open subscheme structure, and let \(X_i\) be its
scheme-theoretic closure in \(X\).  The underlying space of \(X_i\)
is \(Z_i\), and \(X_i|_{V_i}=V_i\), including its nilpotents.
Each \(X_i\) is irreducible and proper over \(S\).
Apply the irreducible case to obtain projective \(X_i'\) mapping
isomorphically over a dense open \(W_i\subseteq X_i\).
Replacing \(W_i\) by \(W_i\cap V_i\) makes it open in \(X\).
The disjoint union of these opens is dense in \(X\), and
\(X'=\coprod_iX_i'\), projective over \(S\), maps isomorphically over
their union.  Thus the reduction preserves the required open
subscheme even when \(X\) is generically nonreduced.

(b) The schemes \(S\) and \(X\) are noetherian and hence
quasi-compact.  Cover \(S\) by finitely many affine opens and refine a
finite affine cover of \(X\) so that each member
\(U_i=\operatorname{Spec}B_i\) maps into some
\(\operatorname{Spec}A_i\subseteq S\).  Since \(B_i\) is a finitely
generated \(A_i\)-algebra, there is a closed immersion followed by an open immersion
\[
U_i\hookrightarrow\mathbb A_{A_i}^{n_i}
\hookrightarrow\mathbb P_{A_i}^{n_i}.
\]
The last scheme is an open subscheme of \(\mathbb P_S^{n_i}\).  Hence
\(U_i\) is locally closed in \(\mathbb P_S^{n_i}\), and thus is open in
its scheme-theoretic closure \(P_i\).  The scheme \(P_i\) is closed in
\(\mathbb P_S^{n_i}\), so it is projective over \(S\).  This gives the
required finite quasi-projective cover.

(c) Assume \(X\) irreducible.  Every nonempty \(U_i\) is dense, so
\[
U=\bigcap_{i=1}^nU_i
\]
is open dense.  Put \(P=\prod_{i=1}^nP_i\), let
\[
j:U\longrightarrow X\times_SP
\]
be the product of the inclusion \(U\hookrightarrow X\) and the maps
\(U\hookrightarrow U_i\hookrightarrow P_i\), and let
\(X'\hookrightarrow X\times_SP\) be its scheme-theoretic image.  Write
\(g:X'\to X\) and \(h:X'\to P\) for the projections.

For each \(i\), set
\[
V_i=U_i\times_S\prod_{j\ne i}P_j\subseteq P.
\]
We claim
\[
h^{-1}(V_i)=g^{-1}(U_i)
\]
and that \(h\) restricts there to a closed immersion into \(V_i\).
Indeed, over \(U_i\times_SP\), the image of \(U\) lies in the graph of
\(U_i\hookrightarrow P_i\), which is closed because \(P_i\to S\) is
separated.  Hence the scheme-theoretic closure is contained in that
graph and projects as a closed subscheme of \(V_i\).  Conversely, over
\(X\times_SV_i\), the image lies in the graph of the open immersion
\(U_i\hookrightarrow X\).  This graph is closed because \(X\to S\) is
separated.  Thus a point of \(X'\) whose \(P_i\)-coordinate lies in
\(U_i\) has its \(X\)-coordinate equal to it, proving the claim.

Since the \(U_i\) cover \(X\), the image \(h(X')\) is contained in
\(V=\bigcup_iV_i\), and \(h:X'\to V\) is a closed immersion, as this is
local on the target.  On the other hand,
\[
X'\hookrightarrow X\times_SP\longrightarrow P
\]
is proper: the first map is closed and the second is the base change of
the proper morphism \(X\to S\).  Hence \(h(X')\) is closed in all of
\(P\).  The two opens \(V\) and \(P\setminus h(X')\) cover \(P\); on
the first, \(h\) is a closed immersion, and on the second its source is
empty.  Therefore \(h:X'\to P\) is a closed immersion.  Since \(P\) is
projective over \(S\), \(X'\) is projective over \(S\).

(d) Restrict the construction to \(U\times_SP\).  There the original
map \(j\) is the graph of the morphism \(U\to P\).  Since \(P\to S\) is
separated, this graph is already closed.  Its scheme-theoretic closure
is therefore itself, and consequently
\[
g^{-1}(U)\xrightarrow{\ \sim\ }U.
\]
Together with (a)--(c), this proves Chow's lemma.
\end{solution}

\begin{exerciseintro}
% record: H-II-04-011
\exerciseheading{Exercise II.4.11}
\addcontentsline{toc}{subsection}{Exercise II.4.11}

\begin{context}
We use the Krull--Akizuki theorem: the integral closure of a
one-dimensional noetherian domain in a finite extension of its fraction
field is noetherian of dimension one.  An integrally closed
one-dimensional noetherian local domain is a discrete valuation ring.

As customary in the valuative criterion, a field may be treated as the
degenerate rank-zero valuation ring.  If the term “discrete valuation
ring” is required to have a nonzero maximal ideal, part (a) needs the
minor qualification that the original local ring is not a field, except
when the extension has positive transcendence degree.
We also use the finite-type test for universal closedness: a
finite-type morphism \(T\to Y\) is universally closed if all
\(\mathbb A_Y^n\times_Y T\to\mathbb A_Y^n\) are closed.
This test incorporates arbitrary base changes; it follows by
finite-presentation approximation (Stacks Project, Lemma 32.14.2).
\end{context}

\begin{problem}
(a) Let \((\mathcal O,\mathfrak m)\) be a noetherian local domain with
fraction field \(K\), and let \(L/K\) be a finitely generated field
extension.  Prove that there is a discrete valuation ring \(R\) of \(L\)
dominating \(\mathcal O\).  Reduce first to \(L/K\) finite; then use a
suitable generator \(x_1\) of \(\mathfrak m\), the ring
\[
\mathcal O'=\mathcal O
\left[\frac{x_2}{x_1},\ldots,\frac{x_n}{x_1}\right],
\]
a prime minimal over \((x_1)\), and Krull--Akizuki.

(b) Let \(f:X\to Y\) be a finite-type morphism of noetherian schemes.
Prove that \(f\) is separated, respectively proper, if and only if the
corresponding valuative criterion holds for all discrete valuation
rings.
\end{problem}
\end{exerciseintro}

\begin{solution}
(a) Choose a transcendence basis \(t_1,\ldots,t_r\) for \(L/K\).  Then
\[
\mathcal O_1=
\mathcal O[t_1,\ldots,t_r]_
{(\mathfrak m,t_1,\ldots,t_r)}
\]
is a noetherian local domain with fraction field
\(K(t_1,\ldots,t_r)\), and it dominates \(\mathcal O\).  Since \(L\) is
finite over this fraction field, replacing \(\mathcal O\) by
\(\mathcal O_1\) reduces the problem to the case where \(L/K\) is
finite.

Assume \(\mathfrak m\ne0\), and choose generators
\(\mathfrak m=(x_1,\ldots,x_n)\).  After renumbering, we may arrange
that
\[
x_1^{N-1}\notin\mathfrak m^N\qquad\text{for every }N\geq1.
\]
Indeed, otherwise each \(x_i\) would satisfy the reverse inclusion for
all sufficiently large \(N\).  Choosing one common \(N\), a
pigeonhole argument on monomials would give
\(\mathfrak m^d=\mathfrak m^{d+1}\) for some \(d\); Nakayama would imply
\(\mathfrak m^d=0\), impossible in a domain with
\(\mathfrak m\ne0\).

Form
\[
\mathcal O'=\mathcal O
\left[\frac{x_2}{x_1},\ldots,\frac{x_n}{x_1}\right]\subseteq K.
\]
The ideal \(x_1\mathcal O'\) is proper.  Otherwise an expression for
\(x_1^{-1}\) as a polynomial in the \(x_i/x_1\), after clearing
denominators, would put \(x_1^{N-1}\) in \(\mathfrak m^N\) for some
\(N\).  Moreover,
\[
\mathfrak m\subseteq x_1\mathcal O'\cap\mathcal O.
\]
The contraction on the right is proper and \(\mathfrak m\) is maximal,
so equality holds.

Let \(\mathfrak p\) be a prime minimal over
\(x_1\mathcal O'\).  The noetherian ring \(\mathcal O'\) is a domain,
and the principal ideal theorem gives
\(\operatorname{ht}\mathfrak p=1\).  Also
\(\mathfrak p\cap\mathcal O=\mathfrak m\).  Thus
\[
A=\mathcal O'_{\mathfrak p}
\]
is a one-dimensional noetherian local domain dominating
\(\mathcal O\).

Let \(\widetilde A\) be the integral closure of \(A\) in \(L\).
Krull--Akizuki says that \(\widetilde A\) is a noetherian
one-dimensional domain.  Choose a maximal ideal
\(\mathfrak n\subseteq\widetilde A\).  By integrality its contraction is
the maximal ideal of \(A\).  Therefore
\[
R=\widetilde A_{\mathfrak n}
\]
dominates \(A\), hence \(\mathcal O\), and is an integrally closed
one-dimensional noetherian local domain with fraction field \(L\).
It is therefore a discrete valuation ring.

If \(\mathfrak m=0\), the stated rank-zero convention handles the
finite algebraic case.  When \(L/K\) has positive transcendence degree,
one may instead begin with
\(K[t_1,\ldots,t_r]_{(t_1)}\) and apply the finite-extension argument, obtaining a
nontrivial discrete valuation ring.

(b) Necessity follows immediately from the usual valuative criteria.
For sufficiency, recall the specialization step in their proofs.  If
\(q\) lies in the closure of the image of a finite-type morphism
\(T\to Y\), with \(Y\) noetherian, choose a generic point \(\xi\) of an
irreducible component of \(T\) whose image specializes to \(q\).  The
image of
\[
\mathcal O_{Y,q}\longrightarrow k(\xi)
\]
is a noetherian local domain, and \(k(\xi)\) is finitely generated over
its fraction field.  Part (a) therefore supplies a discrete valuation
ring in \(k(\xi)\) dominating that local domain.  Its generic point maps
to \(\xi\) and its closed point maps to \(q\).

For separatedness, apply this argument to the diagonal into
\(X\times_YX\).  A point in its closure is reached by a DVR whose
generic point lies in the diagonal.  The two induced maps into \(X\)
agree generically, so the uniqueness hypothesis makes them equal.
Thus the diagonal has closed image, and a diagonal with closed image
is a closed immersion.

For properness, fix \(n\geq0\) and put
\(T=\mathbb A_Y^n\times_YX\).  For every closed subscheme \(Z\subseteq T\),
the preceding specialization argument applies to \(Z\to\mathbb A_Y^n\),
whose target is noetherian.  The DVR lifting hypothesis for \(f\),
together with the given map to \(\mathbb A_Y^n\), gives a map to \(T\).
Its generic point lies in \(Z\), so its whole image lies in \(Z\):
the defining ideal pulls back to zero in the domain \(R\).
Consequently the image of \(Z\) is closed.  All these affine-space
base changes are therefore closed.  The universal-closedness test
in the context now handles arbitrary, possibly nonnoetherian base
changes.  Combined with finite type and separatedness, it proves
properness.
\end{solution}

\begin{exerciseintro}
% record: H-II-04-012
\exerciseheading{Exercise II.4.12}
\addcontentsline{toc}{subsection}{Exercise II.4.12}

\begin{context}
Let \(k\) be algebraically closed.  A nontrivial valuation ring of a
one-variable function field has a unique center on its nonsingular
projective curve.  On a nonsingular surface, the local ring at the
generic point of a curve is a discrete valuation ring.

In part (b)(2), the local ring at the generic point of \(Y'\) is a
discrete valuation ring provided \(X'\) is normal there, as is the case
when \(X'\) is nonsingular or is obtained by blowing up nonsingular
points.  Without this normality condition, one must pass to the
normalization; the unqualified statement is otherwise false.
\end{context}

\begin{problem}
Let \(k\) be algebraically closed.

(a) If \(K/k\) is finitely generated of transcendence degree one, prove that every valuation
ring of \(K/k\), other than \(K\), is discrete and that these rings are
the local rings of the abstract nonsingular projective curve \(C_K\).

(b) Let \(K/k\) have transcendence degree two and let \(X\) be a complete
nonsingular surface with function field \(K\).

(1) For an irreducible curve \(Y\subseteq X\) with generic point \(x_1\),
show that \(\mathcal O_{X,x_1}\) is a discrete valuation ring centered
at \(x_1\).

(2) If \(f:X'\to X\) is birational and an irreducible curve
\(Y'\subseteq X'\) is contracted to a closed point \(x_0\in X\), show
that the local ring at the generic point of \(Y'\) gives a discrete
valuation centered at \(x_0\), under the normality qualification in the
context.

(3) Starting at a closed point \(x_0\), repeatedly blow up \(x_i\) and
choose a closed point \(x_{i+1}\) on the new exceptional curve.  Put
\[
R_0=\bigcup_{i\geq0}\mathcal O_{X_i,x_i}.
\]
Show that \(R_0\) is local, that a valuation ring \(R\) of \(K\)
dominating \(R_0\) is a valuation ring of \(K/k\) centered at \(x_0\),
and explain when such an \(R\) is discrete.
\end{problem}
\end{exerciseintro}

\begin{solution}
(a) Let \(C_K\) be the nonsingular projective curve with function field
\(K\), and let \(R\ne K\) be a valuation ring of \(K/k\).  Properness
gives a unique center \(x\in C_K\).  It cannot be the generic point,
because then \(R\) would contain \(K\) and equal \(K\).  Thus \(x\) is
closed.  Since \(C_K\) is nonsingular,
\(\mathcal O_{C_K,x}\) is a discrete valuation ring, and \(R\) dominates
it.  An overring of a valuation ring inside its fraction field is a
localization at a prime.  A discrete valuation ring has only the primes
\((0)\) and its maximal ideal; domination rules out localization at
\((0)\).  Hence
\[
R=\mathcal O_{C_K,x}.
\]
Conversely each such local ring is a valuation ring of \(K/k\), so the
nontrivial valuation rings are precisely the closed points of \(C_K\).

(b) We consider the three kinds of surface valuations in turn.

(1) The point \(x_1\) has codimension one on the regular surface
\(X\).  Hence \(\mathcal O_{X,x_1}\) is a one-dimensional regular local
domain and therefore a discrete valuation ring with fraction field
\(K\).  Its center on \(X\) is \(x_1\) by construction.

(2) Let \(\eta'\) be the generic point of \(Y'\).  Birationality
identifies the function fields of \(X'\) and \(X\) with \(K\).  Under
the stated normality hypothesis,
\(\mathcal O_{X',\eta'}\) is an integrally closed one-dimensional
noetherian local domain, hence a discrete valuation ring of \(K/k\).
The local homomorphism
\[
\mathcal O_{X,x_0}\longrightarrow\mathcal O_{X',\eta'}
\]
sends \(\mathfrak m_{x_0}\) into the maximal ideal, because the residue
map factors through the constant map \(Y'\to x_0\); its contraction is
exactly \(\mathfrak m_{x_0}\).  Thus the valuation ring dominates
\(\mathcal O_{X,x_0}\), and its center on \(X\) is \(x_0\).

(3) The blowup map induces an inclusion
\[
\mathcal O_{X_i,x_i}\subseteq\mathcal O_{X_{i+1},x_{i+1}},
\]
and the latter local ring dominates the former.  Consequently
\[
\mathfrak m_0=\bigcup_{i\geq0}\mathfrak m_{X_i,x_i}
\]
is an ideal of \(R_0\).  Every element outside \(\mathfrak m_0\) is a
unit in some, hence every later, local ring.  Thus \(R_0\) is local with
maximal ideal \(\mathfrak m_0\), and its fraction field is \(K\).

If a valuation ring \(R\subseteq K\) dominates \(R_0\), it dominates
each \(\mathcal O_{X_i,x_i}\).  It contains \(k\), and no nonzero
element of \(k\) lies in its maximal ideal, so its valuation is trivial
on \(k\).  It is therefore a valuation ring of \(K/k\), and already its
domination of \(\mathcal O_{X,x_0}\) says that its center on \(X\) is
\(x_0\).

Intrinsically, \(R\) is discrete exactly when its value group has a
least positive element and is generated by that element; equivalently,
\(R\) is noetherian.  A principal maximal ideal alone does not
suffice: the lexicographic valuation of \(k(x,y)\) with
\(v(x)=(1,0)\), \(v(y)=(0,1)\) has maximal ideal \((y)\) but value
group \(\mathbb Z^2\).  It dominates all the blowup local rings
with parameters \(y,x/y^i\), so it occurs in this construction.  There is also
a useful geometric formulation for this construction.  A discrete
\(R\), after normalizing its value, has completion
\[
\widehat R\cong\kappa(R)[[t]].
\]
The inclusions of the local rings then give a dominant formal arc
\[
\operatorname{Spec}\kappa(R)[[t]]\longrightarrow X
\]
whose successive closed-point images lie over the \(x_i\).  Conversely,
any such dominant formal arc, meaning that it induces an injection
\(K\hookrightarrow\kappa(R)((t))\), defines a discrete valuation ring
\[
\{a\in K:\operatorname{ord}_t(a)\geq0\}.
\]
It dominates \(R_0\) exactly when its successive centers are the chosen
points.  In the important residue-field-\(k\) case this is simply a
dominant arc \(\operatorname{Spec}k[[t]]\to X\).

This condition is genuinely restrictive.  For example, on an affine
chart with parameters \(x,y\), a choice
\[
x\longmapsto t,\qquad y\longmapsto\phi(t),
\]
where \(\phi(t)\in tk[[t]]\) is transcendental over \(k(t)\), gives an
injective map \(k(x,y)\hookrightarrow k((t))\) and hence a discrete
example.  In contrast, a monomial valuation with
\[
v(x)=1,\qquad v(y)=\alpha
\]
for irrational \(\alpha>0\) has value group
\(\mathbb Z+\alpha\mathbb Z\), which has no least positive element.
Its successive centers give an infinite sequence of closed infinitely
near points, but its valuation ring is not discrete.
\end{solution}

\begin{exerciseintro}
\corpussection{II.5 Sheaves of Modules}
\addcontentsline{toc}{section}{II.5 Sheaves of Modules}

% record: H-II-05-001
\exerciseheading{Exercise II.5.1}
\addcontentsline{toc}{subsection}{Exercise II.5.1}

\begin{context}
Let \((X,\mathcal O_X)\) be a ringed space.  For an
\(\mathcal O_X\)-module \(\mathcal E\), write
\[
\mathcal E^\vee=\mathcal{H}om_{\mathcal O_X}(\mathcal E,\mathcal O_X).
\]
A morphism of sheaves is an isomorphism if it is so on an open cover.
On an open set where a finite locally free sheaf has basis
\(e_1,\ldots,e_r\), its dual has the dual basis
\(e_1^\vee,\ldots,e_r^\vee\).
\end{context}

\begin{problem}
Let \(\mathcal E\) be finite locally free.

(a) Prove that the evaluation map
\(\mathcal E\to(\mathcal E^\vee)^\vee\) is an isomorphism.

(b) For every \(\mathcal O_X\)-module \(\mathcal F\), prove
\[
\mathcal{H}om(\mathcal E,\mathcal F)
\cong\mathcal E^\vee\otimes\mathcal F.
\]

(c) For all \(\mathcal F,\mathcal G\), prove the natural adjunction
\[
\operatorname{Hom}(\mathcal E\otimes\mathcal F,\mathcal G)
\cong
\operatorname{Hom}\bigl(\mathcal F,\mathcal{H}om(\mathcal E,\mathcal G)\bigr).
\]

(d) If \(f:X\to Y\) is a morphism of ringed spaces,
\(\mathcal F\) is an \(\mathcal O_X\)-module, and \(\mathcal E\) is
finite locally free on \(Y\), prove the projection formula
\[
f_*(\mathcal F\otimes f^*\mathcal E)
\cong f_*\mathcal F\otimes\mathcal E.
\]
\end{problem}
\end{exerciseintro}

\begin{solution}
(a) Define the evaluation morphism by
\[
e\longmapsto\bigl(\lambda\mapsto\lambda(e)\bigr).
\]
On an open set where \(\mathcal E\cong\mathcal O_X^r\), it sends the
basis \(e_i\) to the basis dual to \(e_i^\vee\), and is therefore an
isomorphism.  Such opens cover \(X\), so it is an isomorphism globally.

(b) There is a natural morphism
\[
\mathcal E^\vee\otimes\mathcal F
\longrightarrow\mathcal{H}om(\mathcal E,\mathcal F),\qquad
\lambda\otimes s\longmapsto(e\mapsto\lambda(e)s).
\]
Where \(\mathcal E\cong\mathcal O_X^r\), both sides identify with
\(\mathcal F^r\), and this morphism is the identity under those
identifications.  It is consequently a global isomorphism.

(c) Given \(u:\mathcal E\otimes\mathcal F\to\mathcal G\), send a local
section \(s\) of \(\mathcal F\) to the homomorphism
\[
e\longmapsto u(e\otimes s).
\]
This gives a morphism
\(\mathcal F\to\mathcal{H}om(\mathcal E,\mathcal G)\).
Conversely, evaluate such a morphism on \(\mathcal E\) to obtain
\(\mathcal E\otimes\mathcal F\to\mathcal G\).  The two constructions
are inverse on local sections and are compatible with restriction;
they therefore give the stated natural bijection.

(d) The adjunction map \(f^{-1}\mathcal E\to f^{-1}f_*\!f^*\mathcal E\)
and multiplication define a natural morphism
\[
f_*\mathcal F\otimes_{\mathcal O_Y}\mathcal E
\longrightarrow
f_*(\mathcal F\otimes_{\mathcal O_X}f^*\mathcal E).
\]
Restrict to an open \(V\subseteq Y\) on which
\(\mathcal E|_V\cong\mathcal O_V^r\).  Then
\[
f^*\mathcal E|_{f^{-1}(V)}\cong\mathcal O_{f^{-1}(V)}^r,
\]
and both sides of the displayed morphism restrict to
\((f_*\mathcal F|_V)^r\).  The map is the evident identity there.
These \(V\) cover \(Y\), so the projection morphism is an isomorphism.
\end{solution}

\begin{exerciseintro}
% record: H-II-05-002
\exerciseheading{Exercise II.5.2}
\addcontentsline{toc}{subsection}{Exercise II.5.2}

\begin{context}
If \(R\) is a discrete valuation ring with fraction field \(K\), then
\(\operatorname{Spec}R=\{\eta,s\}\) has precisely the open subsets
\[
\varnothing,\qquad\{\eta\},\qquad X.
\]
Moreover,
\(\mathcal O_X(X)=R\) and
\(\mathcal O_X(\{\eta\})=K\).
\end{context}

\begin{problem}
Let \(X=\operatorname{Spec}R\).

(a) Prove that an \(\mathcal O_X\)-module is equivalent to the data of
an \(R\)-module \(M\), a \(K\)-vector space \(L\), and a \(K\)-linear
map
\[
\rho:M\otimes_RK\longrightarrow L.
\]

(b) Prove that the corresponding sheaf is quasi-coherent if and only
if \(\rho\) is an isomorphism.
\end{problem}
\end{exerciseintro}

\begin{solution}
(a) For an \(\mathcal O_X\)-module \(\mathcal F\), put
\[
M=\mathcal F(X),\qquad L=\mathcal F(\{\eta\}).
\]
Then \(M\) is an \(R\)-module, \(L\) is a \(K\)-vector space, and the
restriction \(M\to L\) is \(R\)-linear.  It therefore extends uniquely
to a \(K\)-linear map
\[
\rho:M\otimes_RK\to L.
\]
There are no other nontrivial restriction maps and no nontrivial open
covers to check.  Conversely, the given data define the values and the
single restriction map of an \(\mathcal O_X\)-module.  These
constructions are mutually inverse, also on morphisms.

(b) The quasi-coherent sheaf associated to \(M\) has
\[
\widetilde M(X)=M,\qquad
\widetilde M(\{\eta\})=M\otimes_RK.
\]
There is a natural morphism \(\widetilde M\to\mathcal F\), equal to the
identity on \(X\) and to \(\rho\) on \(\{\eta\}\).  Hence it is an
isomorphism exactly when \(\rho\) is.  Since every quasi-coherent sheaf
on the affine scheme \(X\) is associated to its module of global
sections, this proves both directions.
\end{solution}

\begin{exerciseintro}
% record: H-II-05-003
\exerciseheading{Exercise II.5.3}
\addcontentsline{toc}{subsection}{Exercise II.5.3}

\begin{context}
For \(X=\operatorname{Spec}A\) and an \(A\)-module \(M\),
\[
\widetilde M(D(f))=M_f.
\]
The distinguished opens form a basis, and
\(\mathcal F(D(f))\) is naturally an \(A_f\)-module for every
\(\mathcal O_X\)-module \(\mathcal F\).
\end{context}

\begin{problem}
Prove that for every \(A\)-module \(M\) and every
\(\mathcal O_X\)-module \(\mathcal F\) there is a natural isomorphism
\[
\operatorname{Hom}_A\bigl(M,\Gamma(X,\mathcal F)\bigr)
\cong
\operatorname{Hom}_{\mathcal O_X}(\widetilde M,\mathcal F).
\]
\end{problem}
\end{exerciseintro}

\begin{solution}
Taking global sections sends a sheaf morphism
\(\widetilde M\to\mathcal F\) to an \(A\)-linear map
\[
M=\Gamma(X,\widetilde M)\longrightarrow\Gamma(X,\mathcal F).
\]

Conversely, let \(u:M\to\Gamma(X,\mathcal F)\) be \(A\)-linear.  For
each \(f\in A\), define
\[
u_f:M_f\longrightarrow\mathcal F(D(f)),\qquad
\frac{m}{f^n}\longmapsto
\frac{u(m)|_{D(f)}}{f^n}.
\]
This is well-defined because the target is an \(A_f\)-module.  The maps
\(u_f\) commute with restriction to smaller distinguished opens.
They therefore define a unique morphism
\(\widetilde M\to\mathcal F\), since the \(D(f)\) form a basis.

The resulting sheaf morphism has global map \(u\).  Conversely, an
\(\mathcal O_X\)-linear sheaf morphism is forced on every \(D(f)\) to be
the localization of its map on \(X\), so the first construction followed
by the second returns it.  The two assignments are inverse and are
natural in both \(M\) and \(\mathcal F\).
\end{solution}

\begin{exerciseintro}
% record: H-II-05-004
\exerciseheading{Exercise II.5.4}
\addcontentsline{toc}{subsection}{Exercise II.5.4}

\begin{context}
On \(U=\operatorname{Spec}A\), the functor
\(M\mapsto\widetilde M\) is exact and commutes with arbitrary direct
sums.  Every module has a presentation by free modules.  If \(A\) is
noetherian, every finite \(A\)-module is finitely presented.
\end{context}

\begin{problem}
Prove that an \(\mathcal O_X\)-module \(\mathcal F\) is quasi-coherent
if and only if locally it is the cokernel of a morphism of free sheaves.
If \(X\) is noetherian, prove that \(\mathcal F\) is coherent if and
only if locally it is the cokernel of a morphism between free sheaves
of finite rank.
\end{problem}
\end{exerciseintro}

\begin{solution}
Suppose first that \(\mathcal F\) is quasi-coherent.  Near any point
choose an affine open \(U=\operatorname{Spec}A\) with
\(\mathcal F|_U\cong\widetilde M\).  Choose a module presentation
\[
A^{(I)}\longrightarrow A^{(J)}\longrightarrow M\longrightarrow0.
\]
Exactness and compatibility with direct sums give
\[
\mathcal O_U^{(I)}\longrightarrow\mathcal O_U^{(J)}
\longrightarrow\mathcal F|_U\longrightarrow0.
\]
Thus \(\mathcal F\) has the required local presentation.

Conversely, free sheaves are quasi-coherent, and on an affine open the
cokernel of
\[
\mathcal O_U^{(I)}\longrightarrow\mathcal O_U^{(J)}
\]
is the sheaf associated to the cokernel of
\(A^{(I)}\to A^{(J)}\).  It is therefore quasi-coherent.  Since
quasi-coherence is local, \(\mathcal F\) is quasi-coherent.

Now assume \(X\) noetherian.  If \(\mathcal F\) is coherent, then on an
affine open it is associated to a finite module \(M\).  The noetherian
property gives a finite presentation
\[
A^m\longrightarrow A^n\longrightarrow M\longrightarrow0,
\]
and hence a presentation by finite-rank free sheaves.  Conversely, the
cokernel of a morphism between finite-rank free sheaves is locally
associated to a finitely presented, hence finite, module.  It is
coherent.
\end{solution}

\begin{exerciseintro}
% record: H-II-05-005
\exerciseheading{Exercise II.5.5}
\addcontentsline{toc}{subsection}{Exercise II.5.5}

\begin{context}
For an affine morphism
\(f:\operatorname{Spec}B\to\operatorname{Spec}A\) and a \(B\)-module
\(M\), the direct image of \(\widetilde M\) is the quasi-coherent
\(\mathcal O_{\operatorname{Spec}A}\)-module associated to \(M\),
viewed as an \(A\)-module.  A module finite over a finite
\(A\)-algebra is finite over \(A\).
\end{context}

\begin{problem}
Let \(f:X\to Y\) be a morphism of schemes.

(a) Give an example in which \(\mathcal F\) is coherent on \(X\) but
\(f_*\mathcal F\) is not coherent on \(Y\), even though \(X,Y\) are
varieties over a field.

(b) Prove that every closed immersion is finite.

(c) If \(f\) is finite between noetherian schemes and \(\mathcal F\) is
coherent on \(X\), prove that \(f_*\mathcal F\) is coherent on \(Y\).
\end{problem}
\end{exerciseintro}

\begin{solution}
(a) Take the structure morphism
\[
f:\mathbb A_k^1\longrightarrow\operatorname{Spec}k
\]
and \(\mathcal F=\mathcal O_{\mathbb A_k^1}\).  This sheaf is coherent,
but
\[
f_*\mathcal F=\Gamma(\mathbb A_k^1,\mathcal O)=k[t]
\]
as a sheaf on the one-point scheme.  It is not a finite \(k\)-module,
so it is not coherent.

(b) Let \(i:Z\hookrightarrow Y\) be a closed immersion.  For every
affine \(V=\operatorname{Spec}A\subseteq Y\), its inverse image is
\[
i^{-1}(V)=\operatorname{Spec}(A/I)
\]
for some ideal \(I\).  The \(A\)-module \(A/I\) is generated by \(1\),
so the restriction of \(i\) over \(V\) is finite.  Hence \(i\) is
finite.

(c) Work over \(V=\operatorname{Spec}A\subseteq Y\).  Because \(f\) is
finite,
\[
f^{-1}(V)=\operatorname{Spec}B
\]
with \(B\) a finite \(A\)-algebra.  The restriction of \(\mathcal F\)
is \(\widetilde M\) for a finite \(B\)-module \(M\).  For every
\(a\in A\),
\[
(f_*\mathcal F)(D(a))
=\mathcal F(\operatorname{Spec}B_a)=M_a,
\]
so \(f_*\mathcal F|_V\cong\widetilde M\), with \(M\) regarded as an
\(A\)-module.  It is finite over \(A\), because it is finite over the
finite \(A\)-module \(B\).  Thus the direct image is coherent on each
such \(V\), and hence on \(Y\).
\end{solution}

\begin{exerciseintro}
% record: H-II-05-006
\exerciseheading{Exercise II.5.6}
\addcontentsline{toc}{subsection}{Exercise II.5.6}

\begin{context}
For \(X=\operatorname{Spec}A\), \(\mathcal F=\widetilde M\), and
\(\mathfrak p\in X\), the stalk is \(M_{\mathfrak p}\).  If
\(\mathfrak a\subseteq A\), write
\[
\Gamma_{\mathfrak a}(M)
=\{m\in M:\mathfrak a^n m=0\text{ for some }n>0\}.
\]
For a closed set \(Z\), the sheaf \(\mathcal H_Z^0(\mathcal F)\)
consists of local sections whose support is contained in \(Z\).
\end{context}

\begin{problem}
(a) For \(m\in M\), prove
\[
\operatorname{supp}(m)=V(\operatorname{Ann}(m)).
\]

(b) If \(A\) is noetherian and \(M\) is finite, prove
\[
\operatorname{Supp}(\widetilde M)=V(\operatorname{Ann}(M)).
\]

(c) Deduce that the support of a coherent sheaf on a noetherian scheme
is closed.

(d) If \(A\) is noetherian, \(Z=V(\mathfrak a)\), and
\(\mathcal F=\widetilde M\), prove
\[
\widetilde{\Gamma_{\mathfrak a}(M)}
\cong\mathcal H_Z^0(\mathcal F).
\]

(e) On a noetherian scheme, prove that
\(\mathcal H_Z^0(\mathcal F)\) is quasi-coherent when \(\mathcal F\) is
quasi-coherent, and coherent when \(\mathcal F\) is coherent.
\end{problem}
\end{exerciseintro}

\begin{solution}
(a) The germ \(m_{\mathfrak p}\) is zero exactly when some
\(s\notin\mathfrak p\) annihilates \(m\).  This is equivalent to
\(\operatorname{Ann}(m)\not\subseteq\mathfrak p\).  Therefore the
primes at which the germ is nonzero are precisely the primes containing
\(\operatorname{Ann}(m)\), which proves the formula.

(b) Let \(m_1,\ldots,m_r\) generate \(M\).  Then
\[
M_{\mathfrak p}=0
\iff (m_i)_{\mathfrak p}=0\text{ for every }i
\iff\mathfrak p\notin\bigcup_iV(\operatorname{Ann}(m_i)).
\]
Since
\[
\operatorname{Ann}(M)=\bigcap_i\operatorname{Ann}(m_i),
\]
the union on the right is
\(V(\operatorname{Ann}(M))\).  Taking complements yields the claim.

(c) On every affine open \(U=\operatorname{Spec}A\), a coherent sheaf
is \(\widetilde M\) for a finite module \(M\), so part (b) shows that
its support meets \(U\) in a closed subset of \(U\).  Closedness is
local on the ambient space, hence the global support is closed.

(d) A section \(m\in M\) has support contained in \(V(\mathfrak a)\)
if and only if
\[
V(\operatorname{Ann}(m))\subseteq V(\mathfrak a),
\]
or equivalently
\(\mathfrak a\subseteq\sqrt{\operatorname{Ann}(m)}\).  Since
\(\mathfrak a\) is finitely generated, this holds exactly when
\(\mathfrak a^n m=0\) for some \(n\).  Thus the two sides have the same
global sections.

The identification is compatible with localization.  More precisely,
for every \(f\in A\),
\[
\Gamma_{\mathfrak a}(M)_f
=\Gamma_{\mathfrak a_f}(M_f).
\]
One inclusion is immediate.  For the other, use finite generation of a
power of \(\mathfrak a\) to clear a common power of \(f\) from the
annihilation relations.  Applying the global calculation over
\(\operatorname{Spec}A_f\) now identifies the two sheaves on every
distinguished open, proving the isomorphism.

(e) Locally write \(Z=V(\mathfrak a)\) and
\(\mathcal F=\widetilde M\).  Part (d) expresses the subsheaf with
support as the sheaf associated to the module
\(\Gamma_{\mathfrak a}(M)\), so it is quasi-coherent.  If
\(\mathcal F\) is coherent, then \(M\) is finite.  Because \(A\) is
noetherian, its submodule \(\Gamma_{\mathfrak a}(M)\) is finite, and
the resulting sheaf is coherent.
\end{solution}

\begin{exerciseintro}
% record: H-II-05-007
\exerciseheading{Exercise II.5.7}
\addcontentsline{toc}{subsection}{Exercise II.5.7}

\begin{context}
For a finite module over a noetherian ring, kernels and cokernels of
maps between finite modules are finite, and their supports are closed.
Nakayama's lemma says that lifts of a basis of
\(M/\mathfrak mM\) generate a finite module \(M\) over a local ring.
\end{context}

\begin{problem}
Let \(X\) be noetherian and \(\mathcal F\) coherent.

(a) If \(\mathcal F_x\) is a free \(\mathcal O_{X,x}\)-module, prove
that \(\mathcal F\) is free on a neighborhood of \(x\).

(b) Prove that \(\mathcal F\) is locally free if and only if all its
stalks are free.

(c) Prove that \(\mathcal F\) is invertible if and only if there is a
coherent sheaf \(\mathcal G\) such that
\[
\mathcal F\otimes\mathcal G\cong\mathcal O_X.
\]
\end{problem}
\end{exerciseintro}

\begin{solution}
(a) Work on \(U=\operatorname{Spec}A\), write
\(\mathcal F|_U=\widetilde M\), and let \(x\) correspond to
\(\mathfrak p\).  Choose elements \(m_1,\ldots,m_r\in M\) whose images
form a basis of \(M_{\mathfrak p}\), after multiplying basis elements
by denominators that are units in \(A_{\mathfrak p}\).  The map
\[
u:A^r\longrightarrow M
\]
is an isomorphism after localization at \(\mathfrak p\).  Its kernel
and cokernel are finite and have supports not containing
\(\mathfrak p\).  Some distinguished neighborhood \(D(f)\) of
\(\mathfrak p\) avoids both supports.  Hence \(u_f\) is an isomorphism,
so \(\mathcal F|_{D(f)}\) is free.

(b) One direction follows by taking stalks.  Conversely, if every
stalk is free, part (a) supplies a free neighborhood of every point.
Thus \(\mathcal F\) is locally free.

(c) If \(\mathcal F\) is invertible, its dual is coherent and evaluation
gives
\[
\mathcal F\otimes\mathcal F^\vee\cong\mathcal O_X.
\]

Conversely, suppose
\(\mathcal F\otimes\mathcal G\cong\mathcal O_X\).
At \(x\), put
\[
B=\mathcal O_{X,x},\qquad
P=\mathcal F_x,\qquad Q=\mathcal G_x.
\]
These are finite modules over the noetherian local ring
\((B,\mathfrak m)\), and \(P\otimes_BQ\cong B\).  Reducing modulo
\(\mathfrak m\) gives
\[
(P/\mathfrak mP)\otimes_{B/\mathfrak m}
(Q/\mathfrak mQ)\cong B/\mathfrak m.
\]
Both vector spaces therefore have dimension one.  Nakayama shows that
\(P\) and \(Q\) are cyclic, say \(P\cong B/I\) and
\(Q\cong B/J\).  Then
\[
B\cong P\otimes_BQ\cong B/(I+J).
\]
The annihilators of the two ends are equal, so \(I+J=0\), and hence
\(I=J=0\).  Thus \(P\cong B\) for every \(x\).  Part (b) now shows
that \(\mathcal F\) is locally free of rank one.
\end{solution}

\begin{exerciseintro}
% record: H-II-05-008
\exerciseheading{Exercise II.5.8}
\addcontentsline{toc}{subsection}{Exercise II.5.8}

\begin{context}
Let \(X\) be noetherian and \(\mathcal F\) coherent.  Define
\[
\varphi(x)=
\dim_{k(x)}\bigl(\mathcal F_x\otimes_{\mathcal O_{X,x}}k(x)\bigr).
\]
For a finite module over a local ring, this dimension is the minimal
number of generators.  A function is upper semicontinuous when each
set \(\{x:\varphi(x)\geq n\}\) is closed.
\end{context}

\begin{problem}
(a) Prove that \(\varphi\) is upper semicontinuous.

(b) If \(\mathcal F\) is locally free and \(X\) is connected, prove
that \(\varphi\) is constant.

(c) Conversely, if \(X\) is reduced and \(\varphi\) is constant, prove
that \(\mathcal F\) is locally free.
\end{problem}
\end{exerciseintro}

\begin{solution}
(a) Let \(\varphi(x)=r\).  On an affine neighborhood
\(\operatorname{Spec}A\), write \(\mathcal F=\widetilde M\) and let
\(\mathfrak p\) correspond to \(x\).  Lift a basis of
\(M_{\mathfrak p}/\mathfrak pM_{\mathfrak p}\).  Nakayama and clearing
denominators give, after shrinking to a distinguished neighborhood of
\(\mathfrak p\), a surjection
\[
A^r\longrightarrow M.
\]
Every fiber on this neighborhood is generated by at most \(r\)
elements, so \(\varphi\leq r\) there.  Thus the set
\(\{x:\varphi(x)<n\}\) is open for every \(n\), and its complement
\(\{x:\varphi(x)\geq n\}\) is closed.

(b) On an open set where
\(\mathcal F\cong\mathcal O_X^r\), the function \(\varphi\) equals
\(r\).  Hence it is locally constant.  A locally constant
integer-valued function on a connected space is constant.

(c) Suppose \(\varphi\equiv r\).  On a reduced affine open
\(\operatorname{Spec}A\), write \(\mathcal F=\widetilde M\).  As in
(a), near any chosen point there is a surjection
\[
A^r\longrightarrow M.
\]
After shrinking, its kernel \(K\) is finite, so choose a finite
presentation
\[
A^m\xrightarrow{v}A^r\longrightarrow M\longrightarrow0
\]
whose image is \(K\).  For every prime \(\mathfrak q\) in this open,
the induced surjection
\[
k(\mathfrak q)^r\longrightarrow M\otimes_Ak(\mathfrak q)
\]
is an isomorphism, since both sides have dimension \(r\).  Therefore
the matrix \(v\otimes k(\mathfrak q)\) is zero.  Every entry of the
matrix \(v\) lies in every prime ideal of \(A\), hence in the
nilradical.  Because \(A\) is reduced, all entries vanish.  Thus
\(K=0\), \(M\cong A^r\), and \(\mathcal F\) is free near the chosen
point.  It is therefore locally free.
\end{solution}

\begin{exerciseintro}
% record: H-II-05-009
\exerciseheading{Exercise II.5.9}
\addcontentsline{toc}{subsection}{Exercise II.5.9}

\begin{context}
Let \(S\) be generated by \(S_1\) over \(S_0\), let \(M\) be graded,
and put \(X=\operatorname{Proj}S\).  Write
\[
\Gamma_*(\mathcal F)=
\bigoplus_{d\in\mathbb Z}\Gamma(X,\mathcal F(d)).
\]
Under the finite-generation hypotheses on \(S\) in part (b),
if a finite graded \(S\)-module \(N\) has
\(\widetilde N=0\), then \(N_d=0\) for all sufficiently large \(d\):
for finitely many degree-one generators \(f_i\) of \(S_+\), powers of
the \(f_i\) kill a finite set of homogeneous generators of \(N\).
Under those same hypotheses, every coherent sheaf on \(X\) is a
quotient of a finite sum of sufficiently negative twists.
\end{context}

\begin{problem}
(a) Construct a natural graded homomorphism
\[
\alpha:M\longrightarrow\Gamma_*(\widetilde M).
\]

(b) Suppose \(S_0=A\) is a finite-type \(k\)-algebra, \(S_1\) is finite
over \(A\), and \(M\) is finite over \(S\).  Prove that
\[
\alpha_d:M_d\longrightarrow
\Gamma(X,\widetilde M(d))
\]
is an isomorphism for every sufficiently large \(d\).

(c) Retain the hypotheses on \(S\) in (b).  Call two graded modules equivalent if their sufficiently high
truncations are isomorphic, and call a module quasi-finitely generated
if it is equivalent to a finite module.  Prove that sheafification and
\(\Gamma_*\) induce inverse equivalences between quasi-finitely
generated graded modules modulo this equivalence and coherent sheaves
on \(X\).
\end{problem}
\end{exerciseintro}

\begin{solution}
(a) If \(m\in M_d\), then on \(D_+(f)\) its image in
\[
\bigl(M(d)_f\bigr)_0
\]
defines a section of \(\widetilde M(d)\).  These local sections agree
on overlaps, so they glue to \(\alpha_d(m)\).  Multiplication of these
fractions shows that the maps \(\alpha_d\) form a natural graded
\(S\)-homomorphism.

(b) Choose degree-one generators \(f_1,\ldots,f_r\) of \(S\) over
\(A\).  For a graded module \(N\), set
\[
E(N)_d=
\left\{(n_i)\in\prod_iN_{d+1}:
f_jn_i=f_in_j\text{ for all }i,j\right\}.
\]
There is a graded map
\[
\beta:N\longrightarrow E(N),\qquad
n\longmapsto(f_1n,\ldots,f_rn).
\]
When \(N=M\), both \(M\) and \(E(M)\) are finite \(S\)-modules because
\(S\) is noetherian.  On \(D_+(f_i)\), the inverse to the sheafified
map sends a compatible tuple to \(n_i/f_i\).  Hence
\(\widetilde\beta\) is an isomorphism.  Its finite kernel and cokernel
have zero associated sheaf, so the observation in the context gives an
integer \(d_0\) such that
\[
\beta_d:M_d\xrightarrow{\ \sim\ }E(M)_d
\qquad(d\geq d_0).
\]

First \(\alpha_d\) is injective for \(d\geq d_0\).
If a nonzero homogeneous \(m\) of such a degree has zero associated
section, some power of \(S_+\) kills \(m\).  Choose a nonzero
homogeneous element in the last nonzero \(S_+^q m\).
It is killed by every \(f_i\), contradicting injectivity of
\(\beta\) in its degree.  Identify \(M_{\geq d_0}\) with its image.

Let \(s\in H^0(X,\widetilde M(d))\), with \(d\geq d_0\).
Clearing denominators on the finite standard cover gives an \(N\)
such that \(S_+^N s\) lies in the image of \(M\).  More explicitly,
for each \(i\) a sufficiently large \(f_i^{q_i}s\) is represented
by an element of \(M\); any degree-\(N\) monomial is divisible by
one of these powers once \(N>\sum_i(q_i-1)\).

Descend on \(N\).  For any degree-\((N-1)\) monomial \(h\), put
\(t=hs\), and lift \(f_it\) uniquely to
\(m_i\in M_{d+N}\).  The relations \(f_jm_i=f_im_j\) hold in \(M\)
by the injectivity already proved.  Surjectivity of
\(\beta_{d+N-1}\) gives an \(m\) with \(f_im=m_i\) for all \(i\).
Then \(\alpha(m)=t\), since their difference is zero on each
\(D_+(f_i)\).  Thus \(S_+^{N-1}s\) is also in the image.
Continuing to \(N=0\) proves surjectivity of \(\alpha_d\).

(c) Define morphisms in the category of tails by
\[
\operatorname{Hom}_{\mathrm{tail}}(M,N)
=\varinjlim_q\operatorname{Hom}_{\mathrm{gr},S}
(M_{\geq q},N_{\geq q}),
\]
with transition maps given by restriction and composition on a
common truncation.  This specifies the category, rather than only
an equivalence relation on objects.

Truncating a graded module does not change its associated sheaf.
Part (b) says that for every finite \(M\), the unit
\[
M\longrightarrow\Gamma_*(\widetilde M)
\]
is an isomorphism in all sufficiently high degrees.  The same holds
for a quasi-finitely generated module after replacing it by an
equivalent finite one.

Conversely, for every quasi-coherent \(\mathcal F\) there is a natural
isomorphism
\[
\widetilde{\Gamma_*(\mathcal F)}\cong\mathcal F.
\]
If \(\mathcal F\) is coherent, choose finitely many global sections of
one twist \(\mathcal F(n)\) that generate it.  They generate a finite
graded submodule
\[
N\subseteq\Gamma_*(\mathcal F)
\]
whose associated sheaf is \(\mathcal F\).  Applying (b) to \(N\) shows
that \(N\to\Gamma_*(\mathcal F)\) is an isomorphism in high degrees.
Thus \(\Gamma_*(\mathcal F)\) is quasi-finitely generated.  The unit
and counit just described are isomorphisms in the quotient category,
so the two functors are inverse equivalences.
\end{solution}

\begin{exerciseintro}
% record: H-II-05-010
\exerciseheading{Exercise II.5.10}
\addcontentsline{toc}{subsection}{Exercise II.5.10}

\begin{context}
Let \(S=A[x_0,\ldots,x_r]\) and \(X=\operatorname{Proj}S\).  For a
homogeneous ideal \(I\), define
\[
\overline I=
\{s\in S:\text{ for every }i\text{ some }n_i\text{ satisfies }
x_i^{n_i}s\in I\}.
\]
The closed subscheme defined by \(I\) has ideal sheaf
\(\widetilde I\), and the opens \(D_+(x_i)\) cover \(X\).
Here
\(\Gamma_{\geq0}(\mathcal F)=\bigoplus_{d\geq0}H^0(X,\mathcal F(d))\).
For all \(r\geq0\), \(\Gamma_{\geq0}(\mathcal O_X)=S\), and
sheafifying \(\Gamma_{\geq0}(\mathcal F)\) recovers \(\mathcal F\).
Restricting to nonnegative degrees is essential when \(r=0\).
\end{context}

\begin{problem}
(a) Prove that \(\overline I\) is a homogeneous ideal.

(b) Prove that two homogeneous ideals define the same closed subscheme
of \(X\) if and only if their saturations agree.

(c) If \(Y\subseteq X\) is closed, prove that
\(\Gamma_{\geq0}(\mathcal I_Y)\) is saturated and is the largest homogeneous
ideal defining \(Y\).

(d) Deduce a bijection between saturated homogeneous ideals of \(S\)
and closed subschemes of \(X\).
\end{problem}
\end{exerciseintro}

\begin{solution}
(a) If \(s,t\in\overline I\), then for each \(i\) one common sufficiently
large power of \(x_i\) sends both \(s\) and \(t\) into \(I\); it
therefore sends \(s+t\) into \(I\).  If \(a\in S\), any power that sends
\(s\) into \(I\) also sends \(as\) into \(I\).  Hence
\(\overline I\) is an ideal.

Write \(s=\sum_ds_d\) with \(s_d\) homogeneous.  Since \(I\) is
homogeneous, \(x_i^{n_i}s\in I\) implies
\(x_i^{n_i}s_d\in I\) for every \(d\).  Thus every homogeneous
component \(s_d\) lies in \(\overline I\), so the ideal is homogeneous.

(b) For homogeneous \(s\in S\),
\[
s\in\overline I
\iff s/1\in I_{x_i}\text{ for every }i.
\]
Consequently \(\overline I\) is the intersection, inside \(S\), of the
contractions of the homogeneous localizations \(I_{x_i}\).

Two ideals define the same closed subscheme precisely when
\[
(I_1)_{(x_i)}=(I_2)_{(x_i)}
\]
for every \(i\).  Since \(x_i\) is invertible in the full graded
localization, its degree-zero part determines every homogeneous degree.
Thus these equalities are equivalent to
\((I_1)_{x_i}=(I_2)_{x_i}\) for every \(i\), which by the preceding
intersection formula is equivalent to
\(\overline I_1=\overline I_2\).

(c) Put \(J=\Gamma_{\geq0}(\mathcal I_Y)\), viewed inside
\(\Gamma_{\geq0}(\mathcal O_X)\).  Suppose a homogeneous section \(s\) has
the property that, for each \(i\), some \(x_i^{n_i}s\) belongs to \(J\).
On \(D_+(x_i)\), the section \(x_i\) is invertible, so
\[
s|_{D_+(x_i)}\in\mathcal I_Y(\deg s).
\]
These opens cover \(X\), and membership in a subsheaf is local.
Therefore \(s\in J\), proving that \(J\) is saturated.

If \(I\) defines \(Y\), the natural map
\(\widetilde I\to\mathcal I_Y\) sends each homogeneous element of \(I\)
to a global twisted section of \(\mathcal I_Y\).  Hence \(I\subseteq J\).
Sheafification gives \(\widetilde J=\mathcal I_Y\), so \(J\)
does define \(Y\).  Thus it is the largest ideal defining \(Y\).

(d) Send a saturated ideal \(I\) to the subscheme defined by
\(\widetilde I\), and send \(Y\) to \(\Gamma_{\geq0}(\mathcal I_Y)\).
Part (c) says the second output is saturated.  Part (b), together with
the maximality in (c), says that starting from a saturated \(I\)
returns \(I\), while the ideal sheaf of the subscheme obtained from
\(Y\) is again \(\mathcal I_Y\).  The two constructions are inverse.
\end{solution}

\begin{exerciseintro}
% record: H-II-05-011
\exerciseheading{Exercise II.5.11}
\addcontentsline{toc}{subsection}{Exercise II.5.11}

\begin{context}
For graded \(A\)-algebras \(S,T\) generated in degree one, their Segre
product is
\[
S\mathbin{\#_A}T=\bigoplus_{n\geq0}S_n\otimes_AT_n.
\]
Degree-one standard affine charts cover all three Proj schemes.
\end{context}

\begin{problem}
For \(X=\operatorname{Proj}S\) and \(Y=\operatorname{Proj}T\), prove
\[
\operatorname{Proj}(S\mathbin{\#_A}T)
\cong X\times_{\operatorname{Spec}A}Y.
\]
Under this isomorphism prove
\[
\mathcal O(1)\cong
p_1^*\mathcal O_X(1)\otimes p_2^*\mathcal O_Y(1).
\]
When degree-one generators \(x_i\) and \(y_j\) give projective
embeddings of \(X\) and \(Y\), show that the generators
\(x_i\otimes y_j\) give their Segre embedding.
\end{problem}
\end{exerciseintro}

\begin{solution}
Let \(R=S\mathbin{\#_A}T\).  The products
\[
D_+(f)\times_AD_+(g)
=\operatorname{Spec}\bigl(S_{(f)}\otimes_AT_{(g)}\bigr)
\]
for \(f\in S_1\), \(g\in T_1\), form an affine cover of
\(X\times_AY\).  There is a natural ring
isomorphism
\[
R_{(f\otimes g)}
\xrightarrow{\ \sim\ }
S_{(f)}\otimes_AT_{(g)}.
\]
Indeed, a degree-zero fraction on the left gives degree-zero fractions
in the two factors.  Conversely, for a tensor of fractions on the
right, multiply numerator and denominator by sufficiently high powers
of \(f\) and \(g\) so that both denominators have the same exponent; it
then comes from a power of \(f\otimes g\).

These isomorphisms commute with further homogeneous localization.
They therefore glue the affine charts of \(\operatorname{Proj}R\) to
the displayed product, proving the first assertion.

On the same chart, the standard generator of
\(\mathcal O_{\operatorname{Proj}R}(1)\) corresponds to the tensor
product of the standard generators of
\(p_1^*\mathcal O_X(1)\) and \(p_2^*\mathcal O_Y(1)\).  The transition
functions are likewise the products of the two transition functions.
Hence the local identifications glue to
\[
\mathcal O_{\operatorname{Proj}R}(1)
\cong p_1^*\mathcal O_X(1)\otimes p_2^*\mathcal O_Y(1).
\]

If \(S_1\) and \(T_1\) are generated by \(x_0,\ldots,x_r\) and
\(y_0,\ldots,y_s\), then \(R_1\) is generated by
\(x_i\otimes y_j\).  The resulting homogeneous coordinates on the
product are exactly all pairwise products of the coordinates of the
two factors.  Thus the induced map to
\(\mathbb P_A^{(r+1)(s+1)-1}\) is the Segre embedding.
\end{solution}

\begin{exerciseintro}
% record: H-II-05-012
\exerciseheading{Exercise II.5.12}
\addcontentsline{toc}{subsection}{Exercise II.5.12}

\begin{context}
An invertible sheaf \(\mathcal L\) on \(X\) is very ample relative to
\(Y\) if there is an immersion
\[
i:X\hookrightarrow\mathbb P_Y^n
\]
such that \(\mathcal L\cong i^*\mathcal O(1)\).  Under the Segre
embedding,
\[
\mathcal O_{\mathbb P_Y^n\times_Y\mathbb P_Y^m}(1,1)
\cong
p_1^*\mathcal O(1)\otimes p_2^*\mathcal O(1)
\]
is the pullback of \(\mathcal O(1)\).
\end{context}

\begin{problem}
(a) If \(\mathcal L\) and \(\mathcal M\) are very ample on \(X\)
relative to \(Y\), prove that \(\mathcal L\otimes\mathcal M\) is very
ample relative to \(Y\).

(b) Let \(f:X\to Y\) and \(g:Y\to Z\).  If \(\mathcal L\) is very ample
for \(f\) and \(\mathcal M\) is very ample for \(g\), prove that
\[
\mathcal L\otimes f^*\mathcal M
\]
is very ample for \(g\circ f\).
\end{problem}
\end{exerciseintro}

\begin{solution}
(a) Choose immersions
\[
i:X\hookrightarrow\mathbb P_Y^n,\qquad
j:X\hookrightarrow\mathbb P_Y^m
\]
realizing \(\mathcal L\) and \(\mathcal M\).  The product map
\[
(i,j):X\longrightarrow
\mathbb P_Y^n\times_Y\mathbb P_Y^m
\]
is an immersion: factor through the closed graph of \(j\) in
\(X\times_Y\mathbb P_Y^m\), followed by the base change of \(i\).  Compose
with the Segre closed immersion into \(\mathbb P_Y^N\).  The pullback
of \(\mathcal O(1)\) is
\[
i^*\mathcal O(1)\otimes j^*\mathcal O(1)
\cong\mathcal L\otimes\mathcal M.
\]
This realizes the tensor product as very ample.

(b) Choose immersions
\[
X\hookrightarrow\mathbb P_Y^n,\qquad
Y\hookrightarrow\mathbb P_Z^m
\]
realizing \(\mathcal L\) and \(\mathcal M\).  Base change of the second
one gives an immersion
\[
\mathbb P_Y^n\hookrightarrow
\mathbb P_Z^n\times_Z\mathbb P_Z^m.
\]
After composing with the first immersion and the Segre embedding, we
obtain an immersion \(X\hookrightarrow\mathbb P_Z^N\).
The pulled-back hyperplane sheaf is
\[
\mathcal L\otimes f^*\mathcal M.
\]
Hence this sheaf is very ample relative to \(Z\).
\end{solution}

\begin{exerciseintro}
% record: H-II-05-013
\exerciseheading{Exercise II.5.13}
\addcontentsline{toc}{subsection}{Exercise II.5.13}

\begin{context}
Let \(S\) be generated by \(S_1\) over \(S_0\).  For \(d>0\), its
\(d\)-th Veronese ring is
\[
S^{(d)}=\bigoplus_{n\geq0}S_{nd},
\]
where \(S_{nd}\) is assigned new degree \(n\).  The degree-\(d\)
monomials in degree-one generators generate \(S^{(d)}\) in new degree
one.
\end{context}

\begin{problem}
For \(X=\operatorname{Proj}S\), prove
\[
\operatorname{Proj}S^{(d)}\cong X
\]
and prove that \(\mathcal O(1)\) on the left corresponds to
\(\mathcal O_X(d)\).  Show that the projective embedding defined by
degree-\(d\) monomials is the \(d\)-uple embedding.
\end{problem}
\end{exerciseintro}

\begin{solution}
Let \(f\in S_d=S^{(d)}_1\).  The same element defines standard opens in
both Proj schemes, and
\[
\bigl(S^{(d)}_f\bigr)_0=(S_f)_0.
\]
Indeed, a degree-zero fraction for the regraded ring has numerator in
\(S_{nd}\) and denominator \(f^n\), which is exactly a degree-zero
fraction in the original grading.  Degree-\(d\) monomials cover
\(\operatorname{Proj}S\), because the degree-one generators generate
the irrelevant ideal.  The displayed affine isomorphisms agree after
further localization and hence glue to
\[
\operatorname{Proj}S^{(d)}\xrightarrow{\ \sim\ }\operatorname{Proj}S.
\]

On \(D_+(f)\), the standard trivialization of
\(\mathcal O_{\operatorname{Proj}S^{(d)}}(1)\) uses the new
degree-one element \(f\).  In the original grading that element has
degree \(d\), so the same transition functions are those of
\(\mathcal O_X(d)\).  This proves the assertion about twisting sheaves.

If \(x_0,\ldots,x_r\) generate \(S_1\), the new degree-one generators
are the monomials \(x_0^{a_0}\cdots x_r^{a_r}\) with
\(\sum a_i=d\).  The associated homogeneous-coordinate map lists all
degree-\(d\) monomials, which is precisely the \(d\)-uple embedding.
\end{solution}

\begin{exerciseintro}
% record: H-II-05-014
\exerciseheading{Exercise II.5.14}
\addcontentsline{toc}{subsection}{Exercise II.5.14}

\begin{context}
Let \(k\) be algebraically closed and let
\(X\subseteq\mathbb P_k^r\) be connected and normal.  Put
\[
S=S(X),\qquad
S'=\bigoplus_{n\geq0}\Gamma(X,\mathcal O_X(n)).
\]
A connected normal noetherian scheme is integral.  We use two standard
graded facts: the integral closure of a graded domain is graded, and a
Veronese subring of a normal graded domain is normal.

For the final criterion, normality must be accompanied by integrality
or connectedness, since projective normality requires the homogeneous
coordinate ring to be a domain.  Over \(k\), surjectivity in degree zero
forces this connectedness.
\end{context}

\begin{problem}
(a) Prove that \(S\) is a domain and that \(S'\) is the integral closure
of \(S\) in its fraction field.

(b) Prove that \(S_n=S'_n\) for every sufficiently large \(n\).

(c) Prove that \(S^{(d)}\) is integrally closed for all sufficiently
large \(d\), and deduce that a sufficiently high \(d\)-uple embedding
of \(X\) is projectively normal.

(d) In the same integral projective setting, prove that an embedding is
projectively normal if and only if \(X\) is normal and every restriction
map
\[
\Gamma(\mathbb P_k^r,\mathcal O(n))
\longrightarrow\Gamma(X,\mathcal O_X(n))
\]
is surjective for \(n\geq0\).
\end{problem}
\end{exerciseintro}

\begin{solution}
(b) We prove this first, since it supplies one ingredient for (a).
The natural graded map
\[
S\longrightarrow\Gamma_*(\widetilde S)
=\bigoplus_{n\in\mathbb Z}\Gamma(X,\mathcal O_X(n))
\]
is the map of Exercise II.5.9.  Its nonnegative target is \(S'\), and
that exercise shows that it is an isomorphism in all sufficiently large
degrees.  Thus
\[
S_n=S'_n\qquad(n\gg0),
\]
which proves (b).

(a) Since \(X\) is connected and normal, it is integral.  Restriction
of homogeneous polynomials gives an injection \(S\hookrightarrow S'\):
its kernel is exactly the saturated homogeneous ideal defining \(X\).
The ring \(S'\) is a domain, because a nonzero section of a line bundle
is nonzero at the generic point, and the product of two such germs is
nonzero.  Hence \(S\) is a domain.

The eventual equality from (b) implies that \(S'\) lies in
\(\operatorname{Frac}S\).  Indeed, for homogeneous \(u\in S'_m\), choose
a nonzero homogeneous \(a\in S\) of positive degree and a power \(a^q\) such that
\(\deg(a^qu)\) is in the equality range; then
\(a^qu\in S\).  It also implies that \(S'\) is integral over \(S\).
For \(m>0\), a sufficiently high power of \(u\) belongs to
\(S_{qm}=S'_{qm}\), so \(u\) satisfies
\[
T^q-u^q=0.
\]
In degree zero, properness and connectedness give
\(S'_0=\Gamma(X,\mathcal O_X)=k=S_0\).  Thus every homogeneous element,
and hence every element, of \(S'\) is integral over \(S\).

It remains to show that \(S'\) is integrally closed.  The integral
closure of a graded domain is graded, so it is enough to consider a
homogeneous element \(z\in\operatorname{Frac}S'\) integral over \(S'\).
For every nonempty standard affine
\[
U_i=X\cap D_+(x_i),
\]
localization and the trivialization by \(x_i\) identify
\[
S'_{x_i}\cong\Gamma(U_i,\mathcal O_X)[x_i,x_i^{-1}].
\]
The affine ring \(\Gamma(U_i,\mathcal O_X)\) is normal, and so is its
Laurent polynomial ring.  Hence \(z\in S'_{x_i}\) for every \(i\).
If \(z\) has degree \(m\), these local elements are compatible sections
of \(\mathcal O_X(m)\), and therefore glue to a global section.  An
integral homogeneous element cannot have negative degree in a
nonnegatively graded domain, so \(m\geq0\), and \(z\in S'_m\).
Therefore \(S'\) is normal and is exactly the integral closure of \(S\).

(c) Choose \(N\) such that \(S_n=S'_n\) for \(n\geq N\), and take
\(d\geq N\).  Then
\[
S^{(d)}=S'^{(d)},
\]
including degree zero.  A Veronese subring of the normal graded domain
\(S'\) is normal: a homogeneous element of its fraction field integral
over the Veronese is integral over \(S'\), hence lies in \(S'\), and
its degree is divisible by \(d\).  Thus \(S^{(d)}\) is integrally
closed.  Exercise II.5.13 identifies it with the homogeneous coordinate
ring of the \(d\)-uple embedding, which is therefore projectively
normal.

(d) Suppose first that \(X\) is normal and all restriction maps are
surjective.  The degree-zero map shows that
\(\Gamma(X,\mathcal O_X)=k\), so \(X\) is connected; normality then
makes it integral.  Surjectivity in every degree says precisely that
\[
S=S'.
\]
Part (a) now shows that this ring is an integrally closed domain, so the
embedding is projectively normal.

Conversely, suppose \(S\) is an integrally closed domain.  The standard
affine charts of \(X=\operatorname{Proj}S\) have rings
\((S_{x_i})_0\).  Each is an integrally closed domain: an element of
its fraction field integral over the degree-zero ring is integral over
the normal localization \(S_{x_i}\), and homogeneity puts it back in
degree zero.  Hence \(X\) is integral and normal.  Applying (a), the
ring \(S'\) is the integral closure of \(S\) in their common fraction
field.  Since \(S\) is already normal, \(S=S'\), which says exactly
that every displayed restriction map is surjective.
\end{solution}

\begin{exerciseintro}
% record: H-II-05-015
\exerciseheading{Exercise II.5.15}
\addcontentsline{toc}{subsection}{Exercise II.5.15}

\begin{context}
On a noetherian scheme, every open subset is quasi-compact and coherent
sheaves form an abelian category.  If
\(i:U\hookrightarrow X\) is a quasi-compact open immersion and
\(\mathcal F\) is quasi-coherent on \(U\), then \(i_*\mathcal F\) is
quasi-coherent.  A quasi-coherent sheaf is the union of subsheaves
\(\mathcal F_\lambda\) when every local section belongs locally, and
hence on each quasi-compact open globally, to one member of the directed
family.
\end{context}

\begin{problem}
Let \(X\) be noetherian.

(a) If \(X\) is affine, prove that every quasi-coherent sheaf is the
union of its coherent subsheaves.

(b) If \(X\) is affine, \(U\subseteq X\) is open, and \(\mathcal F\) is
coherent on \(U\), prove that \(\mathcal F\) extends to a coherent sheaf
on \(X\).

(c) In (b), suppose
\(\mathcal F\subseteq\mathcal G|_U\) for a quasi-coherent
\(\mathcal G\) on \(X\).  Prove that the extension can be chosen as a
coherent subsheaf of \(\mathcal G\).

(d) Prove (c) for an arbitrary noetherian \(X\).  Deduce the extension
theorem by taking \(\mathcal G=i_*\mathcal F\).

(e) Prove that on any noetherian scheme every quasi-coherent sheaf is
the union of its coherent subsheaves.
\end{problem}
\end{exerciseintro}

\begin{solution}
(a) Write \(X=\operatorname{Spec}A\) and
\(\mathcal F=\widetilde M\).  The module \(M\) is the directed union of
its finite submodules \(N\).  Since \(A\) is noetherian,
\(\widetilde N\) is coherent.  On a distinguished open,
\[
\widetilde M(D(f))=M_f=\bigcup_NN_f.
\]
Every open of the noetherian space \(X\) is quasi-compact and hence has
a finite distinguished cover.  Filtered unions commute with the finite
equalizers expressing the sheaf gluing condition on such a cover.
Therefore
\[
\mathcal F(V)=\bigcup_N\widetilde N(V)
\]
for every open \(V\), as required.

(b) Let \(i:U\hookrightarrow X\).  The sheaf \(i_*\mathcal F\) is
quasi-coherent.  By (a) it is the directed union of coherent subsheaves.
Choose a finite affine cover of the quasi-compact open \(U\) on which
\(\mathcal F\) has finitely many generators.  All those finitely many
generators lie in one coherent subsheaf \(\mathcal F'\) of
\(i_*\mathcal F\), after enlarging within the directed family.  Since
\(\mathcal F'|_U\subseteq\mathcal F\) and contains local generating
sets, equality holds:
\[
\mathcal F'|_U=\mathcal F.
\]

(c) Form the fiber product of quasi-coherent sheaves
\[
\mathcal H=
\mathcal G\times_{\,i_*i^*\mathcal G}i_*\mathcal F.
\]
It is a quasi-coherent subsheaf of \(\mathcal G\), and restriction to
\(U\) identifies \(\mathcal H|_U\) with \(\mathcal F\).  Apply the
argument of (b), now to the union of coherent subsheaves of
\(\mathcal H\).  A coherent member containing finite local generators
of \(\mathcal F\) restricts exactly to \(\mathcal F\), and its inclusion
in \(\mathcal H\subseteq\mathcal G\) gives the desired extension.

(d) Choose a finite affine cover
\[
X=X_1\cup\cdots\cup X_m.
\]
Starting with \(\mathcal F\subseteq\mathcal G|_U\), proceed
inductively.  Suppose a coherent subsheaf has been constructed on
\[
W_i=U\cup X_1\cup\cdots\cup X_i.
\]
On the open \(W_i\cap X_{i+1}\) it is a coherent subsheaf of
\(\mathcal G|_{X_{i+1}}\).  Part (c), applied to the affine scheme
\(X_{i+1}\), extends it to a coherent subsheaf there.  The two
subsheaves agree on the overlap as subsheaves of \(\mathcal G\), so
they glue to a coherent subsheaf on \(W_{i+1}\).  After finitely many
steps this gives \(\mathcal F'\subseteq\mathcal G\) on \(X\), with
\(\mathcal F'|_U=\mathcal F\).

For the announced extension theorem, \(i_*\mathcal F\) is
quasi-coherent and restricts to \(\mathcal F\); take
\(\mathcal G=i_*\mathcal F\).

(e) Let \(\mathcal G\) be quasi-coherent, let \(V\subseteq X\) be open,
and let \(s\in\mathcal G(V)\).  The image of
\(\mathcal O_V\to\mathcal G|_V\), \(1\mapsto s\), is coherent because
\(V\) is noetherian.  By (d) it extends to a coherent subsheaf of
\(\mathcal G\) containing \(s\).  The sums of finitely many such
coherent subsheaves are coherent, so they form a directed family.
Every local section lies in one of them, proving that their union is
\(\mathcal G\).
\end{solution}

\begin{exerciseintro}
% record: H-II-05-016
\exerciseheading{Exercise II.5.16}
\addcontentsline{toc}{subsection}{Exercise II.5.16}

\begin{context}
Tensor, symmetric, and exterior powers of sheaves are obtained by
sheafifying the corresponding module constructions.  Exact sequences
of locally free sheaves split locally, because a local basis of the
quotient can be lifted.  Pullback is a tensor functor and a left
adjoint, so it commutes with tensor products and cokernels.
\end{context}

\begin{problem}
Let \(\mathcal F\) be an \(\mathcal O_X\)-module.

(a) If \(\mathcal F\) is locally free of rank \(n\), compute the ranks
of \(T^r\mathcal F\), \(S^r\mathcal F\), and
\(\bigwedge^r\mathcal F\).

(b) If \(\mathcal F\) is locally free of rank \(n\) and \(0\leq r\leq n\), prove that wedge multiplication
is a perfect pairing
\[
\bigwedge^r\mathcal F\otimes
\bigwedge^{n-r}\mathcal F\longrightarrow
\bigwedge^n\mathcal F.
\]

(c) For an exact sequence of locally free sheaves
\[
0\to\mathcal F'\to\mathcal F\to\mathcal F''\to0,
\]
construct a filtration of \(S^r\mathcal F\) with quotients
\[
S^p\mathcal F'\otimes S^{r-p}\mathcal F''.
\]

(d) Prove the analogous statement for exterior powers and deduce
\[
\bigwedge^n\mathcal F\cong
\bigwedge^{n'}\mathcal F'\otimes
\bigwedge^{n''}\mathcal F'',
\qquad n=n'+n''.
\]

(e) Prove that pullback commutes with tensor, symmetric, and exterior
powers.
\end{problem}
\end{exerciseintro}

\begin{solution}
(a) The assertion is local, so take
\(\mathcal F\cong\mathcal O_X^n\) with basis
\(e_1,\ldots,e_n\).  Pure tensors give a basis of \(T^r\mathcal F\);
nondecreasing \(r\)-tuples give a basis of \(S^r\mathcal F\); and
strictly increasing \(r\)-tuples give a basis of
\(\bigwedge^r\mathcal F\).  Their ranks are respectively
\[
n^r,\qquad
\binom{n+r-1}{r},\qquad
\binom nr.
\]

(b) Again work in a local basis.  For every \(r\)-element subset
\(I\subseteq\{1,\ldots,n\}\), the basis vector
\(e_I\) wedges to \(\pm e_1\wedge\cdots\wedge e_n\) with
\(e_{I^c}\), and wedges to zero with every other complementary-degree
basis vector.  The matrix of the pairing is therefore a signed
permutation matrix.  Hence it induces an isomorphism
\[
\bigwedge^r\mathcal F
\cong
\left(\bigwedge^{n-r}\mathcal F\right)^\vee
\otimes\bigwedge^n\mathcal F.
\]
For \(n=2,r=1\), this is
\(\mathcal F\cong\mathcal F^\vee\otimes\bigwedge^2\mathcal F\).

(c) Define \(F^p\subseteq S^r\mathcal F\) to be the image of
\[
S^p\mathcal F'\otimes S^{r-p}\mathcal F
\longrightarrow S^r\mathcal F.
\]
Then
\[
S^r\mathcal F=F^0\supseteq F^1\supseteq\cdots
\supseteq F^r\supseteq F^{r+1}=0.
\]
Locally split \(\mathcal F\cong\mathcal F'\oplus\mathcal F''\).  The
standard decomposition
\[
S^r(\mathcal F'\oplus\mathcal F'')
\cong
\bigoplus_{i=0}^r
S^i\mathcal F'\otimes S^{r-i}\mathcal F''
\]
identifies \(F^p\) with the sum of the terms for \(i\geq p\).
Consequently
\[
F^p/F^{p+1}
\cong S^p\mathcal F'\otimes S^{r-p}\mathcal F''.
\]
This description is independent of the local splitting, so the local
isomorphisms glue.

(d) Define the filtration in the same way using wedge multiplication.
Under a local splitting,
\[
\bigwedge^r(\mathcal F'\oplus\mathcal F'')
\cong
\bigoplus_{i=0}^r
\bigwedge^i\mathcal F'\otimes
\bigwedge^{r-i}\mathcal F'',
\]
and the same argument gives the required quotients.  For
\(r=n'+n''\), all the associated graded pieces vanish except the one
with \(i=n'\).  The filtration therefore gives the asserted canonical
determinant isomorphism.

(e) There are natural maps
\[
f^*(\mathcal F^{\otimes r})\longrightarrow
(f^*\mathcal F)^{\otimes r},
\]
which are isomorphisms by associativity of tensor product.  Symmetric
and exterior powers are cokernels of the permutation and alternating
relations in tensor powers.  Since \(f^*\) commutes with cokernels, the
tensor isomorphism descends to natural isomorphisms
\[
f^*(S^r\mathcal F)\cong S^r(f^*\mathcal F),\qquad
f^*(\bigwedge^r\mathcal F)\cong\bigwedge^r(f^*\mathcal F).
\]
\end{solution}

\begin{exerciseintro}
% record: H-II-05-017
\exerciseheading{Exercise II.5.17}
\addcontentsline{toc}{subsection}{Exercise II.5.17}

\begin{context}
A morphism \(f:X\to Y\) is affine when \(Y\) has an affine open cover
whose inverse images are affine.  We may use the affine communication
lemma: if a scheme over \(\operatorname{Spec}A\) is covered by finitely
many affine principal opens defined by elements generating the unit
ideal of \(A\), then the scheme is affine.
An affine structural morphism need not have an affine total space unless its base is affine.
\end{context}

\begin{problem}
(a) Prove that \(f\) is affine if and only if \(f^{-1}(V)\) is affine
for every affine open \(V\subseteq Y\).

(b) Prove that affine morphisms are quasi-compact and separated, and
that finite morphisms are affine.

(c) For a quasi-coherent \(\mathcal O_Y\)-algebra \(\mathcal A\),
construct a unique \(Y\)-scheme
\[
\mathbf{Spec}_Y\mathcal A
\]
whose inverse image over every affine \(V\subseteq Y\) is
\(\operatorname{Spec}\mathcal A(V)\), compatibly with restriction.

(d) Prove that the structural morphism of this relative spectrum is affine and
\[
f_*\mathcal O_{\mathbf{Spec}_Y\mathcal A}\cong\mathcal A.
\]
Conversely, prove that for affine \(f:X\to Y\),
\(\mathcal A=f_*\mathcal O_X\) is quasi-coherent and
\(X\cong\mathbf{Spec}_Y\mathcal A\).

(e) For affine \(f\) and \(\mathcal A=f_*\mathcal O_X\), prove that
\(f_*\) is an equivalence from quasi-coherent \(\mathcal O_X\)-modules
to quasi-coherent \(\mathcal A\)-modules.
\end{problem}
\end{exerciseintro}

\begin{solution}
(a) Only one direction needs proof.  Assume an affine cover
\(\{V_i\}\) of \(Y\) has affine inverse images, and let
\(V=\operatorname{Spec}A\subseteq Y\) be arbitrary.  Refine the cover
\(\{V\cap V_i\}\) by distinguished opens of \(V\).  Quasi-compactness
of \(V\) gives finitely many
\[
D(a_1),\ldots,D(a_m)
\]
covering \(V\), each contained in some \(V_i\).  Their inverse images
are affine: each is the fiber product of the two affine schemes
\(f^{-1}(V_i)\) and \(D(a_j)\) over the affine scheme \(V_i\).
Inside \(f^{-1}(V)\) they are the principal opens defined by the
pullbacks of the \(a_j\).  The \(a_j\) generate the unit ideal of
\(A\), so the affine
communication lemma applied to \(f^{-1}(V)\) shows that it is affine.

(b) The inverse image of every affine open is affine and therefore
quasi-compact; hence \(f\) is quasi-compact.  Separatedness is local on
the base.  Over \(V=\operatorname{Spec}A\), write
\[
f^{-1}(V)=\operatorname{Spec}B.
\]
The diagonal is defined by the surjection
\[
B\otimes_AB\longrightarrow B,
\]
so it is a closed immersion.  Thus \(f\) is separated.  A finite
morphism is affine by its local definition.

(c) Cover \(Y\) by affine opens \(V_i=\operatorname{Spec}A_i\), and put
\[
X_i=\operatorname{Spec}\mathcal A(V_i).
\]
If \(D(a)\subseteq V_i\), quasi-coherence gives
\[
\mathcal A(D(a))\cong\mathcal A(V_i)_a,
\]
so \(\operatorname{Spec}\mathcal A(D(a))\) is the corresponding
principal open of \(X_i\).  Covering every overlap \(V_i\cap V_j\) by
such distinguished opens produces canonical identifications between
open pieces of \(X_i\) and \(X_j\).  They satisfy the cocycle condition
because all maps are restriction maps of \(\mathcal A\).  Gluing yields
a scheme \(X\) and a morphism
\[
f:X=\mathbf{Spec}_Y\mathcal A\longrightarrow Y.
\]
The construction gives the required description on every affine open
by (a).  That description also forces all the gluing maps, proving
uniqueness up to unique \(Y\)-isomorphism.

(d) The construction and (a) show that \(f\) is affine.  On an affine
\(V\subseteq Y\),
\[
(f_*\mathcal O_X)(V)
=\Gamma(\operatorname{Spec}\mathcal A(V),\mathcal O)
=\mathcal A(V),
\]
compatibly with localization.  These identifications give
\(f_*\mathcal O_X\cong\mathcal A\).

Conversely, for affine \(f\), write over \(V=\operatorname{Spec}A\)
\[
f^{-1}(V)=\operatorname{Spec}B.
\]
Then \(f_*\mathcal O_X|_V\cong\widetilde B\), so
\(\mathcal A=f_*\mathcal O_X\) is a quasi-coherent algebra.  The affine
charts of \(\mathbf{Spec}_Y\mathcal A\) are \(\operatorname{Spec}B\),
with the same localization maps as those of \(X\); they glue to a
canonical isomorphism \(X\cong\mathbf{Spec}_Y\mathcal A\).

(e) On \(V=\operatorname{Spec}A\) as above, a quasi-coherent sheaf on
\(f^{-1}(V)=\operatorname{Spec}B\) is the same thing as a \(B\)-module.
Its direct image is that module viewed as a module over
\(\mathcal A|_V=\widetilde B\).  Conversely, a quasi-coherent
\(\mathcal A\)-module restricts on \(V\) to a \(B\)-module and hence
defines a quasi-coherent sheaf on \(\operatorname{Spec}B\).
Localization of modules shows that these local constructions agree on
overlaps.  They glue to mutually inverse functors, and the local affine
equivalences show that the unit and counit are isomorphisms.
\end{solution}

\begin{exerciseintro}
% record: H-II-05-018
\exerciseheading{Exercise II.5.18}
\addcontentsline{toc}{subsection}{Exercise II.5.18}

\begin{context}
A geometric vector bundle of rank \(n\) over \(Y\) is locally
\(\mathbb A_Y^n\), with transition automorphisms linear in the fiber
coordinates.  For a locally free sheaf \(\mathcal E\), put
\[
\mathbf V(\mathcal E)
=\mathbf{Spec}_Y\bigl(\operatorname{Sym}\mathcal E\bigr).
\]
The symmetric algebra has the universal property that
\(\mathcal O_Y\)-algebra maps
\(\operatorname{Sym}\mathcal E\to\mathcal B\) correspond to
\(\mathcal O_Y\)-linear maps \(\mathcal E\to\mathcal B\).
\end{context}

\begin{problem}
(a) Prove that \(\mathbf V(\mathcal E)\) is a geometric vector bundle
of rank \(n\), independent up to isomorphism of the local bases chosen.

(b) If \(f:X\to Y\) is a geometric vector bundle, prove that its sheaf
of sections \(\mathcal S(X/Y)\) is naturally a locally free
\(\mathcal O_Y\)-module of rank \(n\).

(c) For \(X=\mathbf V(\mathcal E)\), prove naturally that
\[
\mathcal S(X/Y)\cong\mathcal E^\vee.
\]

(d) Deduce a bijection between isomorphism classes of rank-\(n\)
locally free sheaves and rank-\(n\) geometric vector bundles over
\(Y\).
\end{problem}
\end{exerciseintro}

\begin{solution}
(a) On an open \(U\) where
\(\mathcal E|_U\cong\mathcal O_U^n\),
\[
\operatorname{Sym}(\mathcal E|_U)
\cong\mathcal O_U[x_1,\ldots,x_n],
\]
so
\[
\mathbf V(\mathcal E)|_U\cong\mathbb A_U^n.
\]
On an overlap, changing a basis of \(\mathcal E\) by an invertible
matrix changes the generators \(x_i\) by the corresponding linear
substitution.  Thus the transition functions are linear.  A different
choice of bases merely changes the local trivializations by such
linear maps and hence gives an isomorphic vector bundle.

(b) Over a trivializing open \(U\), a section of
\(\mathbb A_U^n\to U\) is an \(n\)-tuple of sections of
\(\mathcal O_U\).  Therefore
\[
\mathcal S(X/Y)|_U\cong\mathcal O_U^n.
\]
Addition and scalar multiplication are defined coordinatewise.  The
transition matrices of the vector bundle are linear, so they preserve
these operations; consequently the local module structures glue to a
natural \(\mathcal O_Y\)-module structure.  The displayed local
identifications prove that it is locally free of rank \(n\).

(c) For every open \(V\subseteq Y\), a section of
\(\mathbf V(\mathcal E)\) over \(V\) is a \(V\)-morphism
\[
V\longrightarrow
\mathbf{Spec}_V\operatorname{Sym}(\mathcal E|_V).
\]
By relative affine duality and the universal property of the symmetric
algebra, these are naturally equivalent to
\[
\operatorname{Hom}_{\mathcal O_V\text{-alg}}
\bigl(\operatorname{Sym}(\mathcal E|_V),\mathcal O_V\bigr)
\cong
\operatorname{Hom}_{\mathcal O_V}(\mathcal E|_V,\mathcal O_V).
\]
The latter is \(\Gamma(V,\mathcal E^\vee)\).  Compatibility with
restriction and with the module operations gives the claimed
isomorphism of sheaves.

(d) Starting with \(\mathcal E\), parts (a)--(c) produce the bundle
\(\mathbf V(\mathcal E)\), whose sheaf of sections is
\(\mathcal E^\vee\).  Dualizing again recovers \(\mathcal E\), because
finite locally free sheaves are canonically isomorphic to their double
duals.

Conversely, start with a geometric vector bundle \(X\to Y\), put
\(\mathcal S=\mathcal S(X/Y)\), and form
\(\mathbf V(\mathcal S^\vee)\).  On every bundle chart both schemes
are \(\mathbb A_U^n\), and the transition matrices agree by the
construction of \(\mathcal S\).  The local identifications therefore
glue to a vector-bundle isomorphism
\[
X\cong\mathbf V(\mathcal S^\vee).
\]
The two constructions are inverse on isomorphism classes.
\end{solution}

\begin{exerciseintro}
\corpussection{II.6 Divisors}
\addcontentsline{toc}{section}{II.6 Divisors}

% record: H-II-06-001
\exerciseheading{Exercise II.6.1}
\addcontentsline{toc}{subsection}{Exercise II.6.1}

\begin{context}
Let \(n\geq1\).  Condition \((*)\) means that a scheme is noetherian, integral,
separated, and regular in codimension one.  For such a scheme the Weil
divisor class group is defined.  We use
\(K=K(X)\) and write
\[
p:X\times\mathbb P^n\longrightarrow X
\]
for the projection.  The standard affine-bundle calculation gives
\(\operatorname{Cl}(X\times\mathbb A^n)\cong\operatorname{Cl}(X)\).
The product with projective space means \(\mathbb P_X^n\), with the indicated structural base.
\end{context}

\begin{problem}
Prove that, if \(X\) satisfies \((*)\), then so does
\(X\times\mathbb P^n\), and prove a natural decomposition
\[
\operatorname{Cl}(X\times\mathbb P^n)
\cong \operatorname{Cl}(X)\oplus\mathbb Z.
\]
Identify a generator of the second summand.
\end{problem}
\end{exerciseintro}

\begin{solution}
Put \(Y=X\times\mathbb P^n\).  It is noetherian and separated because
\(p\) is projective and \(X\) has those properties.  If
\(W=\operatorname{Spec}A\subseteq X\) is nonempty affine, the inverse
image of \(W\) is covered by spectra of polynomial rings over the
domain \(A\).  These charts have the same generic point and glue along
nonempty opens, so \(Y\) is integral.

We check regularity in codimension one.  Let \(y\in Y\) have
codimension one and let \(x=p(y)\).  Since \(p\) is flat, the dimension
formula leaves two possibilities.  If \(x\) is the generic point of
\(X\), then \(y\) is a codimension-one point of
\(Y_\eta=\mathbb P_K^n\), whose local ring is a discrete valuation
ring.  Otherwise \(x\) has codimension one in \(X\) and \(y\) is the
generic point of the fiber over \(\overline{\{x\}}\).  On a standard
affine chart its local ring is
\[
\mathcal O_{X,x}[t_1,\ldots,t_n]_{(\pi)},
\]
where \(\mathcal O_{X,x}\) is a discrete valuation ring with
uniformizer \(\pi\).  This localization is again a discrete valuation
ring.  Hence \(Y\) satisfies \((*)\).

Let \(H=X\times\mathbb P^{n-1}\) be a relative hyperplane.  We now
describe the class group in a way that also proves there are no hidden
relations.  Restrict a divisor to the generic fiber and take its class:
\[
r:\operatorname{Cl}(Y)\longrightarrow
\operatorname{Cl}(\mathbb P_K^n)\cong\mathbb Z.
\]
This is well defined, and \(r([H])=1\).

Pullback of prime divisors gives a homomorphism
\[
p^*:\operatorname{Cl}(X)\longrightarrow\operatorname{Cl}(Y).
\]
We claim that
\[
0\longrightarrow\operatorname{Cl}(X)
\xrightarrow{p^*}\operatorname{Cl}(Y)
\xrightarrow{r}\mathbb Z\longrightarrow0
\tag{1}
\]
is exact.  Surjectivity on the right follows from the relative
hyperplane.  Suppose \(r([D])=0\).  On \(\mathbb P_K^n\), the divisor
\(D|_{Y_\eta}\) is therefore the divisor of some
\(f\in K(\mathbb P_K^n)^*\).  Subtracting \(\operatorname{div}_Y(f)\)
from \(D\) removes every horizontal component.  Every remaining prime
divisor is vertical; flatness and the dimension formula show that it
is \(p^{-1}(Z)\) for a unique prime divisor \(Z\subseteq X\).  Thus
\(D\) differs from \(p^*E\) by a principal divisor.

It remains to see that \(p^*\) is injective.  If
\(p^*E=\operatorname{div}_Y(g)\), then the restriction of \(g\) to
\(\mathbb P_K^n\) has zero divisor.  A rational function on projective
space with neither zeros nor poles is an element of \(K^*\).  Hence
\(g\in K(X)^*\), and comparison at the vertical discrete valuations
gives \(E=\operatorname{div}_X(g)\).  This proves (1).

Finally, \(1\mapsto[H]\) splits \(r\).  Consequently
\[
\operatorname{Cl}(X\times\mathbb P^n)
\cong p^*\operatorname{Cl}(X)\oplus\mathbb Z[H],
\]
as required.
\end{solution}

\begin{exerciseintro}
% record: H-II-06-002
\exerciseheading{Exercise II.6.2}
\addcontentsline{toc}{subsection}{Exercise II.6.2}

\begin{context}
Let \(k\) be algebraically closed and let
\(X\subseteq\mathbb P_k^n\) be an integral closed subvariety that is
regular in codimension one and has positive dimension.  If
\(D=\sum m_iY_i\) is a Weil divisor on \(X\), set
\[
\deg D=\sum_i m_i\deg Y_i.
\]
At the generic point of each prime divisor \(Y_i\), the local ring of
\(X\) is a discrete valuation ring; its valuation is denoted
\(v_{Y_i}\).
\end{context}

\begin{problem}
(a) If a hypersurface \(V\subseteq\mathbb P^n\) does not contain
\(X\), define an intersection divisor \(V\mathbin{.}X\) by restricting
a local equation of \(V\) to \(X\) and taking its order along every
component of \(V\cap X\).  Prove that this is well defined and extends
linearly to ambient divisors having no component that contains \(X\).

(b) Prove that intersection with \(X\) carries principal ambient
divisors to principal divisors on \(X\), and hence gives a homomorphism
\(\operatorname{Cl}(\mathbb P^n)\to\operatorname{Cl}(X)\).

(c) Identify the coefficient of a component \(Y\) in
\(V\mathbin{.}X\) with the intersection multiplicity
\(i(X,V;Y)\).  Deduce
\[
\deg(D\mathbin{.}X)=\deg(D)\deg(X).
\]

(d) Show that every principal divisor on \(X\) is obtained by
intersecting \(X\) with a principal divisor on \(\mathbb P^n\).
Deduce that degree descends to
\(\operatorname{Cl}(X)\to\mathbb Z\), and prove that the diagram
\[
\begin{array}{ccc}
\operatorname{Cl}(\mathbb P^n)&\longrightarrow&\operatorname{Cl}(X)\\
\deg\big\downarrow&&\big\downarrow\deg\\
\mathbb Z&\xrightarrow{\,m\mapsto m\deg X\,}&\mathbb Z
\end{array}
\]
commutes.  In particular, prove that the top map is injective.
\end{problem}
\end{exerciseintro}

\begin{solution}
(a) Let \(Y\) be an irreducible component of \(V\cap X\), with
generic point \(\eta_Y\).  The projective dimension theorem makes
\(Y\) a prime divisor on \(X\).  Choose an open neighborhood \(U\) of
\(\eta_Y\) in \(\mathbb P^n\) on which \(V\) has one equation \(f\).
Define
\[
n_Y=v_Y(f|_X).
\]
If \(f'\) is another local equation, then near \(\eta_Y\) one has
\(f'=uf\), where \(u\) is a unit.  Its restriction to \(X\) is still a
unit, so \(v_Y(f'|_X)=v_Y(f|_X)\).  Thus
\[
V\mathbin{.}X=\sum_{Y\subseteq V\cap X}n_YY
\]
is independent of all choices.  Valuations are additive under
multiplication of equations, which gives the asserted linear
extension.

(b) Let \(D=\operatorname{div}_{\mathbb P^n}(g)\), with no component
containing \(X\).  At the generic point \(\eta_X\), the function \(g\)
is a unit in the ambient local ring, so it has a nonzero image
\(\bar g\in K(X)^*\).  Near the generic point of a prime divisor
\(Y\subseteq X\), write the ambient divisor of \(g\) as a sum of local
equations of its prime components.  Restricting this factorization to
\(X\) and applying \(v_Y\) gives
\[
\operatorname{coeff}_Y(D\mathbin{.}X)=v_Y(\bar g).
\]
Therefore
\[
D\mathbin{.}X=\operatorname{div}_X(\bar g).
\tag{1}
\]
Every class on \(\mathbb P^n\) has a representative supported on
hyperplanes that do not contain \(X\).  Equation (1) shows that changing
such a representative by a principal divisor does not change the
class on \(X\), so the map on class groups is well defined.

(c) Work in the local ring \(R=\mathcal O_{X,\eta_Y}\).  This is a
discrete valuation ring, and the restriction \(\bar f\) of a local
equation of \(V\) is nonzero.  The local definition of intersection
multiplicity gives
\[
i(X,V;Y)=\operatorname{length}_R(R/(\bar f)).
\]
If \(\bar f=u\pi^a\), with \(u\) a unit and \(\pi\) a uniformizer,
then this length is \(a=v_Y(\bar f)=n_Y\).  Hence the divisor
coefficients are precisely the intersection multiplicities.

The generalized Bezout theorem now yields, for an irreducible
hypersurface \(V\),
\[
\deg(V\mathbin{.}X)
=\sum_Y i(X,V;Y)\deg Y
=\deg V\deg X.
\]
Linearity gives the displayed formula for every allowed divisor
\(D\).

(d) Take \(h\in K(X)^*\).  The residue map
\[
\mathcal O_{\mathbb P^n,\eta_X}\longrightarrow
\mathcal O_{X,\eta_X}=K(X)
\]
is surjective, so choose a lift
\(g\in\mathcal O_{\mathbb P^n,\eta_X}\) of \(h\).  Because its residue
is nonzero, \(g\) is a unit at \(\eta_X\); thus no component of
\(\operatorname{div}_{\mathbb P^n}(g)\) contains \(X\).  The same
valuation calculation as in (1) gives
\[
\operatorname{div}_X(h)
=\operatorname{div}_{\mathbb P^n}(g)\mathbin{.}X.
\]
Part (c) therefore implies that every principal divisor on \(X\) has
degree zero.  Degree consequently defines a homomorphism on
\(\operatorname{Cl}(X)\), and part (c) is exactly the commutativity of
the stated diagram.  Since
\(\operatorname{Cl}(\mathbb P^n)\cong\mathbb Z\) by degree and
\(\deg X>0\), multiplication by \(\deg X\) is injective.  The top
horizontal map must therefore be injective as well.
\end{solution}

\begin{exerciseintro}
% record: H-II-06-003
\exerciseheading{Exercise II.6.3}
\addcontentsline{toc}{subsection}{Exercise II.6.3}

\begin{context}
Let \(V\subseteq\mathbb P_k^n\) be an integral projective variety of
positive dimension, regular in codimension one.  Let
\(X=C(V)\subseteq\mathbb A_k^{n+1}\) be its affine cone,
\(\overline X\subseteq\mathbb P_k^{n+1}\) its projective cone, and
\(P\) the vertex.  Write \(H\) for a hyperplane of \(\mathbb P^n\)
not containing \(V\).  All class groups below are Weil divisor class
groups.
In (d), remove Laurent monomial units so that \(g,h\) are ordinary polynomials in the generic fiber coordinate, normalized to constant coefficient one.
\end{context}

\begin{problem}
(a) Show that projection away from \(P\) makes
\(\overline X\setminus\{P\}\to V\) a Zariski-locally trivial
\(\mathbb A^1\)-bundle, and prove
\[
\operatorname{Cl}(V)\cong
\operatorname{Cl}(\overline X\setminus\{P\})
\cong\operatorname{Cl}(\overline X).
\]

(b) Under this identification, prove that the divisor at infinity
\(V\subseteq\overline X\) corresponds to the hyperplane-section class
\([V\mathbin{.}H]\).  Deduce an exact sequence
\[
0\longrightarrow\mathbb Z
\xrightarrow{\,1\mapsto[V\mathbin{.}H]\,}
\operatorname{Cl}(V)
\longrightarrow\operatorname{Cl}(X)\longrightarrow0.
\]

(c) Prove that the homogeneous coordinate ring \(S(V)\) is a unique
factorization domain exactly when \(V\) is projectively normal and
\(\operatorname{Cl}(V)\cong\mathbb Z\), generated by the hyperplane
class.

(d) Prove that localization at the vertex induces an isomorphism
\[
\operatorname{Cl}(X)\xrightarrow{\sim}
\operatorname{Cl}(\operatorname{Spec}\mathcal O_{X,P}).
\]
\end{problem}
\end{exerciseintro}

\begin{solution}
(a) Use homogeneous coordinates \([x_0:\cdots:x_n:z]\), with
\(P=[0:\cdots:0:1]\) and with \(V\) in \(z=0\).  On
\(U_i=V\cap D_+(x_i)\), every point of the projective cone away from
\(P\) has coordinates
\[
([x_0:\cdots:x_n],z/x_i),
\]
which gives an isomorphism
\[
\pi^{-1}(U_i)\cong U_i\times\mathbb A^1.
\]
On overlaps the fiber coordinate changes by multiplication by a unit,
so these charts form the total space of a line bundle.

The affine-bundle class-group argument applies chartwise and at the
generic fiber: pullback gives
\[
\pi^*: \operatorname{Cl}(V)\xrightarrow{\sim}
\operatorname{Cl}(\overline X\setminus\{P\}).
\tag{1}
\]
Indeed, a divisor can first be made vertical by subtracting a
principal divisor on the generic \(\mathbb A^1\), after which it is the
pullback of a divisor on \(V\); the same generic-fiber argument proves
injectivity.  Since \(\dim V\geq1\), the point \(P\) has codimension at
least two in \(\overline X\).  Removing it changes neither prime
divisors nor principal-divisor relations.  Thus the second map in (1),
from \(\operatorname{Cl}(\overline X)\), is also an isomorphism.

(b) Let \(L\subseteq\mathbb P^{n+1}\) be the hyperplane obtained by
joining \(H\) to \(P\), and let \(L_\infty=(z=0)\).  The two ambient
hyperplanes are linearly equivalent.  Their intersections with
\(\overline X\) are respectively
\(\pi^*(V\mathbin{.}H)\) and the divisor at infinity \(V\).  Hence
\[
[V]=\pi^*[V\mathbin{.}H]
\quad\text{in }\operatorname{Cl}(\overline X).
\tag{2}
\]

The complement of the divisor at infinity is the affine cone \(X\).
The localization sequence for class groups and (2) give
\[
\mathbb Z[V\mathbin{.}H]\longrightarrow
\operatorname{Cl}(V)\longrightarrow\operatorname{Cl}(X)
\longrightarrow0.
\]
The first map is injective: the degree homomorphism of Exercise
II.6.2 sends \([V\mathbin{.}H]\) to the positive integer \(\deg V\).
This proves the asserted short exact sequence.

(c) The cone is \(X=\operatorname{Spec}S(V)\).  A noetherian domain is
a unique factorization domain if and only if it is normal and its
divisor class group is zero.  Normality of \(S(V)\) is precisely
projective normality of the given embedding.  By (b),
\[
\operatorname{Cl}(X)=0
\quad\Longleftrightarrow\quad
\operatorname{Cl}(V)=\mathbb Z[V\mathbin{.}H],
\]
and the injection in (b) says that this cyclic group is infinite.
Combining the two criteria proves (c).

(d) For a normal noetherian domain \(A\), localization induces a
surjection
\[
\operatorname{Cl}(A)\longrightarrow\operatorname{Cl}(A_{\mathfrak m})
\]
whose kernel is generated by height-one primes not contained in
\(\mathfrak m\).  Here \(A=S(V)\) and
\(\mathfrak m=A_+\) is the vertex.  The same statement remains valid
under the present regular-in-codimension-one hypothesis because the
argument uses only the codimension-one valuation rings.  It is
therefore enough to show that every prime divisor of \(X\) avoiding
\(P\) is principal.

We record the graded argument.  Let \(D\) be such a prime divisor.
On the punctured cone
\[
X^\circ=X\setminus\{P\}\longrightarrow V
\]
it must dominate \(V\): every divisor lying over a proper divisor of
\(V\) is a cone and contains \(P\).  Over the generic point of \(V\),
the punctured cone is
\(\operatorname{Spec}K[t,t^{-1}]\), after choosing a basis of the
one-dimensional generic fiber.  Thus \(D\) is defined there by an
irreducible Laurent polynomial \(g\), which we normalize to have
constant coefficient \(1\).

Because \(D\) avoids the vertex, its prime ideal contains an element
with nonzero constant term.  After scaling, write it as
\[
f=1+f_1+\cdots+f_a,
\qquad f_i\in A_i.
\]
In \(K[t,t^{-1}]\), the polynomial \(g\) divides \(f\); write
\(f=gh\), normalizing \(h\) also to have constant coefficient \(1\).
For a prime divisor \(Z\subseteq V\), the valuation along its cone is
the Gauss extension of \(v_Z\): after trivializing
\(\mathcal O_V(1)\) at \(Z\), it sends a Laurent polynomial to the
minimum valuation of its coefficients.  All coefficients \(f_i\) are
regular at \(Z\), while the constant coefficient is \(1\), so
\[
v_{C(Z)}(f)=0.
\]
Both \(g\) and \(h\) have constant coefficient \(1\), hence their
Gauss valuations are at most zero.  Since their sum is zero, both are
zero.  Thus \(g\) has no vertical zero or pole.  On the generic fiber,
irreducibility gives valuation one along \(D\) and zero along every
other horizontal prime.  Consequently
\[
\operatorname{div}_X(g)=D.
\]
Every generator of the localization kernel is therefore zero, so the
surjection on class groups is injective and hence an isomorphism.
\end{solution}

\begin{exerciseintro}
% record: H-II-06-004
\exerciseheading{Exercise II.6.4}
\addcontentsline{toc}{subsection}{Exercise II.6.4}

\begin{context}
Let \(k\) be a field of characteristic different from \(2\), put
\[
R=k[x_1,\ldots,x_n],
\]
and let \(f\in R\) be nonconstant and square-free.  Thus every
irreducible factor of \(f\) occurs with exponent one.
\end{context}

\begin{problem}
Prove that
\[
A=R[z]/(z^2-f)
\]
is integrally closed in its fraction field.
\end{problem}
\end{exerciseintro}

\begin{solution}
Because \(f\) is square-free and nonconstant, it is not a square in
\(F=\operatorname{Frac}R\).  Hence \(T^2-f\) is irreducible over
\(F\), and Gauss's lemma shows that \(A\) is a domain.  Its fraction
field is
\[
L=F(z),\qquad z^2=f,
\]
and every element of \(L\) has a unique expression \(a+bz\) with
\(a,b\in F\).

Let \(\alpha=a+bz\in L\) be integral over \(A\).  Since \(A\) is finite
and integral over \(R\), transitivity makes \(\alpha\) integral over
\(R\).  The nontrivial \(F\)-automorphism of \(L\) sends \(z\) to
\(-z\) and preserves \(A\), so \(\bar\alpha=a-bz\) is integral over
\(R\) as well.  Therefore
\[
\operatorname{Tr}_{L/F}(\alpha)=2a,\qquad
\operatorname{N}_{L/F}(\alpha)=a^2-b^2f
\]
are integral over \(R\).  They belong to \(F\), and \(R\) is a unique
factorization domain, hence integrally closed.  Thus
\[
2a\in R,\qquad a^2-b^2f\in R.
\]
The characteristic assumption gives \(a\in R\), and consequently
\(b^2f\in R\).

It remains to show \(b\in R\).  For every irreducible
\(p\in R\), let \(v_p\) be its exponent valuation.  From
\(b^2f\in R\) we have
\[
2v_p(b)+v_p(f)\geq0.
\]
Square-freeness says \(v_p(f)\) is either \(0\) or \(1\).  If the
integer \(v_p(b)\) were negative, the left side would be at most
\(-1\), a contradiction.  Hence \(v_p(b)\geq0\) for every irreducible
\(p\), so \(b\in R\).  We have proved
\(\alpha=a+bz\in A\); therefore \(A\) is integrally closed.
\end{solution}

\begin{exerciseintro}
% record: H-II-06-005
\exerciseheading{Exercise II.6.5}
\addcontentsline{toc}{subsection}{Exercise II.6.5}

\begin{context}
Let \(k\) be algebraically closed of characteristic different from
\(2\), and let \(2\leq r\leq n\).  Consider the affine and projective
quadrics
\[
X=\operatorname{Spec}
k[x_0,\ldots,x_n]/(x_0^2+\cdots+x_r^2),
\qquad
Q=V_+(x_0^2+\cdots+x_r^2)\subseteq\mathbb P_k^n.
\]
The unused variables \(x_{r+1},\ldots,x_n\) make these quadrics cones
when \(r<n\).
\end{context}

\begin{problem}
(a) Prove that \(X\) is normal.

(b) After a linear change of coordinates, write the equation as
\[
x_0x_1=x_2^2+\cdots+x_r^2.
\]
Compute \(\operatorname{Cl}(X)\), obtaining respectively
\(\mathbb Z/2\mathbb Z\), \(\mathbb Z\), and \(0\) when
\(r=2\), \(r=3\), and \(r\geq4\).

(c) Compute \(\operatorname{Cl}(Q)\).  Show that it is
\(\mathbb Z\), with a hyperplane equal to twice a generator, for
\(r=2\); that it is \(\mathbb Z\oplus\mathbb Z\) for \(r=3\); and
that it is \(\mathbb Z\), generated by a hyperplane section, for
\(r\geq4\).

(d) For \(r\geq4\), prove Klein's theorem: every prime divisor
\(Y\subseteq Q\) is the scheme-theoretic intersection of \(Q\) with
an irreducible hypersurface of \(\mathbb P^n\), with multiplicity one.
\end{problem}
\end{exerciseintro}

\begin{solution}
(a) Regard the equation as
\[
x_0^2=-\bigl(x_1^2+\cdots+x_r^2\bigr).
\]
The polynomial on the right is square-free when \(r\geq2\): in the
case \(r=2\) it is a product of two distinct linear forms, and for
\(r\geq3\) it has no repeated irreducible factor.  Exercise II.6.4
therefore shows that the coordinate ring of \(X\) is integrally
closed.

(b) Choose \(i\in k\) with \(i^2=-1\).  Replacing \(x_0,x_1\) by
\(x_0+ix_1,x_0-ix_1\), and rescaling, gives the stated equation.  Put
\[
A=k[x_0,\ldots,x_n]/(x_0x_1-q),\qquad
q=x_2^2+\cdots+x_r^2.
\]
After inverting \(x_0\),
\[
A_{x_0}\cong
k[x_0,x_0^{-1},x_2,\ldots,x_n],
\]
a unique factorization domain whose units are \(k^*x_0^{\mathbb Z}\).
The divisor localization sequence says that \(\operatorname{Cl}(A)\)
is generated by the prime components of \(V(x_0)\), with precisely the
relations supplied by these units.

If \(r=2\), the scheme \(V(x_0)\) has the single prime component
\(D=V(x_0,x_2)\), occurring twice, so
\(\operatorname{div}(x_0)=2D\).  Hence
\[
\operatorname{Cl}(X)\cong\mathbb Z/2\mathbb Z.
\]
If \(r=3\), write
\(q=(x_2+ix_3)(x_2-ix_3)\).  There are two components \(D_+\) and
\(D_-\), and
\(\operatorname{div}(x_0)=D_++D_-\).  Thus
\[
\operatorname{Cl}(X)
\cong\mathbb Z^2/\mathbb Z(1,1)\cong\mathbb Z.
\]
For \(r\geq4\), the quadratic form \(q\) in at least three variables
is irreducible.  Consequently \(V(x_0)\) is one reduced prime divisor
\(D\), and \(\operatorname{div}(x_0)=D\); its class generates and is
zero.  Therefore \(\operatorname{Cl}(X)=0\).

(c) The affine quadric \(X\) is the affine cone over \(Q\).
Exercise II.6.3 gives
\[
0\longrightarrow\mathbb Z[H]\longrightarrow
\operatorname{Cl}(Q)\longrightarrow\operatorname{Cl}(X)
\longrightarrow0,
\tag{1}
\]
where \(H\) is a hyperplane section.

For \(r=2\), first take \(n=2\).  The smooth conic is
\(\mathbb P^1\), so its class group is \(\mathbb Z\), generated by a
point, while \(H\) has degree \(2\).  Increasing \(n\) forms successive
projective cones, and Exercise II.6.3(a) leaves the class group
unchanged.  Thus the same conclusion holds for every \(n\geq2\).

For \(r=3\), the nondegenerate quadric surface in \(\mathbb P^3\) is
\(\mathbb P^1\times\mathbb P^1\) under the Segre embedding.  Its class
group is \(\mathbb Z\oplus\mathbb Z\), generated by the two rulings,
and \(H\) has class \((1,1)\).  Successive projective cones again do
not change the class group.  Finally, if \(r\geq4\), part (b) and (1)
show directly that \(\operatorname{Cl}(Q)=\mathbb Z[H]\).

(d) Let
\[
B=k[x_0,\ldots,x_n]/(x_0^2+\cdots+x_r^2)
\]
be the homogeneous coordinate ring of \(Q\).  By (a) it is normal,
and by (b) its class group is zero for \(r\geq4\).  Hence \(B\) is a
unique factorization domain.

The cone over a prime divisor \(Y\subseteq Q\) corresponds to a
homogeneous height-one prime \(\mathfrak p\subseteq B\).  Factoriality
gives \(\mathfrak p=(g)\).  Since the ideal is homogeneous, its
generator may be chosen homogeneous; primality makes \(g\)
irreducible.  Lift \(g\) to a homogeneous polynomial
\(G\in k[x_0,\ldots,x_n]\).  The lift is irreducible: a nontrivial
homogeneous factorization of \(G\) would give a nontrivial
factorization of \(g\) in \(B\), because no positive-degree
homogeneous factor can become a unit there.  Thus \(V(G)\) is an
irreducible hypersurface, and
\[
Q\cap V(G)=\operatorname{Proj}(B/(g))=Y
\]
scheme-theoretically.  Since \((g)\) is prime, the multiplicity is
one.
\end{solution}

\begin{exerciseintro}
% record: H-II-06-006
\exerciseheading{Exercise II.6.6}
\addcontentsline{toc}{subsection}{Exercise II.6.6}

\begin{context}
Work over an algebraically closed field \(k\) of characteristic
different from two.  Let
\[
E:\ y^2z=x^3-xz^2
\]
be the smooth plane cubic with origin \(O=[0:1:0]\).  Its group law is
the one obtained from
\[
E\longrightarrow\operatorname{Cl}^0(E),\qquad
P\longmapsto[P-O].
\]
Intersections with a line are always counted with multiplicity.
When we say that a point is annihilated by \(m\), the identity is
included; every nonidentity point so obtained has exact order \(m\)
in the cases below.
\end{context}

\begin{problem}
(a) Prove that a line cuts out the divisor \(P+Q+R\) exactly when
\(P+Q+R=O\).

(b) Prove that \(2P=O\) exactly when the tangent line at \(P\) passes
through \(O\).

(c) Prove that \(3P=O\) exactly when \(P\) is an inflection point.

(d) Over \(\mathbb C\), prove that the points whose homogeneous
coordinates can be chosen in \(\mathbb Q\) form a subgroup, and
determine that subgroup explicitly.
\end{problem}
\end{exerciseintro}

\begin{solution}
(a) The line at infinity \(z=0\) cuts out the divisor \(3O\) on
\(E\).  If a line \(L\) cuts out \(P+Q+R\), the quotient of equations
of \(L\) and \(z=0\) is a rational function on \(E\), so
\[
P+Q+R\sim3O.
\]
This is equivalent to
\([P-O]+[Q-O]+[R-O]=0\), hence to \(P+Q+R=O\).

Conversely, suppose \(P+Q+R=O\).  Then the effective divisor
\(P+Q+R\) is linearly equivalent to \(3O\).  Hence there is a rational
function \(h\) with \(\operatorname{div}(h)+3O=P+Q+R\).
In particular \(h\in L(3O)\).  On \(E\setminus\{O\}\), write the affine
coordinates as \(x,y\).  Its coordinate ring is
\[
k[x,y]/(y^2-x^3+x),
\]
so every regular function has a unique form \(p(x)+yq(x)\).  Since
\(\operatorname{ord}_O(x)=-2\) and
\(\operatorname{ord}_O(y)=-3\), the pole bound forces
\(p\) to have degree at most one and \(q\) to be constant.  The
space \(L(3O)\) is therefore spanned by \(1,x,y\).
The homogenization of \(h=a+bx+cy\) is a line section whose divisor
is \(\operatorname{div}(h)+3O=P+Q+R\), proving the converse.

(b) Let the tangent at \(P\) meet \(E\) a third time at \(R\).  Part
(a), with multiplicity, gives
\[
2P+R=O.
\]
The tangent passes through \(O\) exactly when \(R=O\), and this is
equivalent to \(2P=O\).

(c) With the same notation, \(P\) is an inflection point exactly when
\(R=P\).  By \(2P+R=O\), this is equivalent to \(3P=O\).

(d) The inverse map and the chord-tangent construction are defined
over \(\mathbb Q\).  More explicitly, a rational line through two
rational points has rational coefficients, and after two rational
roots are removed from its intersection polynomial with \(E\), the
third root is rational.  Tangents are handled by the same calculation.
Thus \(E(\mathbb Q)\) is a subgroup.

We show that it has no points beyond
\[
O,\qquad (0,0),\qquad (1,0),\qquad (-1,0)
\tag{1}
\]
on the affine chart \(z=1\).  Suppose
\(P=(x,y)\in E(\mathbb Q)\) has \(y\neq0\).  Then
\[
a=\frac{x^2-1}{y},\qquad
b=\frac{2x}{y},\qquad
c=\frac{x^2+1}{y}
\]
satisfy
\[
a^2+b^2=c^2,\qquad \frac{ab}{2}=1,
\]
up to changing signs of the sides.  They would therefore give a
rational right triangle of area \(1\).

For completeness, Fermat's descent rules this out.  If a rational
right triangle has square area, clear denominators and remove the
common divisor to obtain a primitive integral right triangle of square
area, chosen with minimal hypotenuse.  Write its sides as
\[
2mn,\qquad m^2-n^2,\qquad m^2+n^2,
\]
where \(m>n\), \(\operatorname{gcd}(m,n)=1\), and \(m,n\) have opposite
parity.  Its area is
\[
mn(m+n)(m-n).
\]
The four displayed factors are pairwise coprime, so their product can
be a square only if
\[
m=e^2,\qquad n=f^2,\qquad m+n=g^2,\qquad m-n=h^2.
\]
The even parameter must be \(n\).  Indeed, if \(m\) were even, then
\(n=f^2\) would be odd, whereas
\(g^2-h^2=2n\) is impossible: the difference of two odd squares is
divisible by \(8\), while \(2n\) is not.  Thus \(f\) is even.  Set
\[
u=\frac{g+h}{2},\qquad v=\frac{g-h}{2}.
\]
Now
\[
u^2+v^2=e^2,\qquad
\frac{uv}{2}=\left(\frac f2\right)^2.
\]
This is a smaller integral right triangle with square area and
hypotenuse \(e<m^2+n^2\), contradicting minimality.  Hence no rational
right triangle has area \(1\), and a rational point on \(E\) must have
\(y=0\).

The equation then gives \(x(x-1)(x+1)=0\), producing exactly the three
affine points in (1), together with \(O\).  Each nonzero affine point
is annihilated by \(2\), so
\[
E(\mathbb Q)\cong
\mathbb Z/2\mathbb Z\oplus\mathbb Z/2\mathbb Z.
\]
\end{solution}

\begin{exerciseintro}
% record: H-II-06-007
\exerciseheading{Exercise II.6.7}
\addcontentsline{toc}{subsection}{Exercise II.6.7}

\begin{context}
Let \(k\) be algebraically closed of characteristic different from
\(2\), and let
\[
X:\ y^2z=x^3+x^2z
\]
be the nodal plane cubic.  Its node is \(P=[0:0:1]\).  We write
\(\operatorname{CaCl}^0(X)\) for Cartier divisors of degree zero modulo
principal Cartier divisors, and identify \(\mathbb G_m(k)\) with
\(k^*\).
For a required branch-value ratio \(1\), use the constant function \(1\); fractional linear interpolation is needed only for unequal values.
\end{context}

\begin{problem}
Compute the degree-zero Cartier divisor class group of \(X\), and
prove naturally that
\[
\operatorname{CaCl}^0(X)\cong\mathbb G_m(k).
\]
\end{problem}
\end{exerciseintro}

\begin{solution}
On the chart \(z=1\), put \(t=y/x\) in the function field.  The
equation gives
\[
x=t^2-1,\qquad y=t(t^2-1),
\]
so \(K(X)=k(t)\).  The normalization is \(\mathbb P^1\), and its map to
\(X\) identifies the two points \(t=1\) and \(t=-1\) to form the node.
More precisely,
\[
\mathcal O_{X,P}
=k[t^2-1,t(t^2-1)]_{(t^2-1,\,t(t^2-1))}
\subseteq k(t).
\]
Inside the semilocal normalization ring at \(t=1,-1\), this is exactly
the set of rational functions that are regular at both points and have
the same value there.  Indeed, the inclusion from left to right is
immediate.  Conversely, if \(a(t)/b(t)\) has equal values at
\(1\) and \(-1\), multiply numerator and denominator by \(b(-t)\);
both resulting polynomials have equal values at the two points, hence
belong to \(k+(t^2-1)k[t]\).

Every Cartier divisor is linearly equivalent to one supported away
from \(P\): subtract a rational function that is a local equation for
the divisor at \(P\).  Two degree-zero divisors supported in
\(X\setminus\{P\}\) represent the same Cartier class precisely when
their difference is the divisor of an element of
\(\mathcal O_{X,P}^*\).  Under normalization we therefore obtain
\[
\operatorname{CaCl}^0(X)
\cong
\frac{\operatorname{Div}^0(\mathbb P^1\setminus\{1,-1\})}
{\{\operatorname{div}(u):u(1)=u(-1)\neq0\}}.
\tag{1}
\]

For a degree-zero divisor \(D\) in the numerator, choose
\(r\in k(t)^*\) with \(D=\operatorname{div}_{\mathbb P^1}(r)\).
Because the support avoids \(1\) and \(-1\), the two values
\(r(1),r(-1)\) are defined and nonzero.  The function \(r\) is unique
up to a nonzero scalar, so
\[
\Phi(D)=\frac{r(1)}{r(-1)}\in k^*
\]
is well defined and multiplicative.  Its kernel consists exactly of
the divisors in the denominator of (1).  It is surjective: for any
\(\lambda\in k^*\), a fractional linear function can be chosen with
neither a zero nor a pole at \(\pm1\) and with ratio of values
\(\lambda\).  Thus \(\Phi\) induces the required natural isomorphism
\[
\operatorname{CaCl}^0(X)\xrightarrow{\sim}k^*
=\mathbb G_m(k).
\]
\end{solution}

\begin{exerciseintro}
% record: H-II-06-008
\exerciseheading{Exercise II.6.8}
\addcontentsline{toc}{subsection}{Exercise II.6.8}

\begin{context}
The group law on \(\operatorname{Pic}(Y)\) is tensor product.  For a
morphism \(f:X\to Y\), pullback of an \(\mathcal O_Y\)-module is
\[
f^*\mathcal F=f^{-1}\mathcal F
\otimes_{f^{-1}\mathcal O_Y}\mathcal O_X.
\]
On a nonsingular integral scheme, divisor classes and invertible
sheaves correspond by \(D\mapsto\mathcal O(D)\).
\end{context}

\begin{problem}
(a) Prove that pullback induces a homomorphism
\[
f^*:\operatorname{Pic}(Y)\longrightarrow\operatorname{Pic}(X).
\]

(b) If \(f\) is a finite morphism of nonsingular curves, prove that
this homomorphism agrees, under the divisor--line-bundle
correspondence, with pullback of divisors including ramification
multiplicities.

(c) Let \(X\subseteq\mathbb P_k^n\) be an integral locally factorial
closed subscheme.  Prove that restriction of invertible sheaves along
the inclusion agrees with the intersection homomorphism
\(\operatorname{Cl}(\mathbb P^n)\to\operatorname{Cl}(X)\) of Exercise
II.6.2.
\end{problem}
\end{exerciseintro}

\begin{solution}
(a) An invertible sheaf is locally isomorphic to the structure sheaf,
so its pullback is again invertible.  There are natural isomorphisms
\[
f^*(\mathcal L\otimes\mathcal M)
\cong f^*\mathcal L\otimes f^*\mathcal M,
\qquad
f^*(\mathcal L^{-1})\cong(f^*\mathcal L)^{-1}.
\]
They show that the isomorphism class of \(f^*\mathcal L\) depends only
on that of \(\mathcal L\), preserves the tensor-product law, and
preserves inverses.  Hence it defines the claimed group homomorphism.

(b) It is enough to treat a prime divisor \(P\in Y\).  Let \(t\) be a
uniformizer at \(P\).  For every \(Q\in X\) above \(P\), the local rings
are discrete valuation rings and
\[
v_Q(f^\#t)=e(Q/P),
\]
the ramification index.  Thus the pulled-back Cartier divisor has
associated Weil divisor
\[
f^*P=\sum_{Q\mapsto P}e(Q/P)Q.
\]
Locally \(\mathcal O_Y(P)\) is the fractional ideal generated by
\(t^{-1}\).  Tensoring with \(\mathcal O_X\) makes it the fractional
ideal generated by \((f^\#t)^{-1}\), whose divisor is the displayed
pullback.  Therefore, for every divisor \(D\),
\[
f^*\mathcal O_Y(D)\cong\mathcal O_X(f^*D).
\]
This proves compatibility on class groups.

(c) Local factoriality identifies Weil and Cartier divisors on \(X\).
Choose a hyperplane \(H\) not containing \(X\).  Its defining linear
form trivializes \(\mathcal O_{\mathbb P^n}(H)\) away from \(H\);
after restriction to \(X\), the same local equation defines the
Cartier divisor \(H\mathbin{.}X\).  Consequently
\[
\mathcal O_{\mathbb P^n}(H)|_X
\cong\mathcal O_X(H\mathbin{.}X).
\]
Both homomorphisms in question are additive, and
\(\operatorname{Cl}(\mathbb P^n)=\mathbb Z[H]\).  Agreement on \(H\)
therefore proves agreement for every class.
\end{solution}

\begin{exerciseintro}
% record: H-II-06-009
\exerciseheading{Exercise II.6.9}
\addcontentsline{toc}{subsection}{Exercise II.6.9}

\begin{context}
Let \(X\) be an integral projective curve over an algebraically closed
field \(k\), let
\(\pi:\widetilde X\to X\) be its normalization, and for \(P\in X\)
write
\[
\widetilde{\mathcal O}_P
\]
for the integral closure of \(\mathcal O_{X,P}\) in \(K(X)\).  This is
a semilocal ring when \(P\) has more than one branch.  Only finitely
many of the local quotients below are nonzero.
An invertible module over a semilocal ring is free.  Thus for a
finite normalization \(\pi\), a line bundle becomes trivial above
a sufficiently small neighborhood of any point, and
\(R^1\pi_*\mathcal O_{\widetilde X}^*=0\).
The Leray sequence then identifies the two first cohomology groups
used below.
\end{context}

\begin{problem}
(a) Prove the exact sequence
\[
0\longrightarrow
\bigoplus_{P\in X}
\widetilde{\mathcal O}_P^*/\mathcal O_{X,P}^*
\longrightarrow\operatorname{Pic}(X)
\xrightarrow{\pi^*}\operatorname{Pic}(\widetilde X)
\longrightarrow0.
\]

(b) Apply it to prove
\[
0\longrightarrow\mathbb G_a(k)\longrightarrow
\operatorname{Pic}(X)\longrightarrow\mathbb Z\longrightarrow0
\]
for a cuspidal plane cubic, and
\[
0\longrightarrow\mathbb G_m(k)\longrightarrow
\operatorname{Pic}(X)\longrightarrow\mathbb Z\longrightarrow0
\]
for a nodal plane cubic.
\end{problem}
\end{exerciseintro}

\begin{solution}
(a) Inside the constant sheaf \(K(X)^*\), there is an exact sequence
of sheaves of abelian groups
\[
1\longrightarrow\mathcal O_X^*
\longrightarrow\pi_*\mathcal O_{\widetilde X}^*
\longrightarrow\mathcal Q\longrightarrow1.
\tag{1}
\]
Localization commutes with integral closure for this finite
normalization, so
\[
\mathcal Q_P
\cong\widetilde{\mathcal O}_P^*/\mathcal O_{X,P}^*.
\]
The normalization is an isomorphism over the regular locus, whose
complement is a finite set.  Hence \(\mathcal Q\) is the finite direct
sum of the corresponding skyscraper sheaves.  In particular,
\[
H^0(X,\mathcal Q)
=\bigoplus_{P\in X}
\widetilde{\mathcal O}_P^*/\mathcal O_{X,P}^*,
\qquad H^1(X,\mathcal Q)=0.
\tag{2}
\]

Both \(X\) and \(\widetilde X\) are integral and projective over
\(k\), so their global units are \(k^*\).  The map between these two
unit groups is an isomorphism.  Taking cohomology in (1), using (2)
and
\[
H^1(X,\pi_*\mathcal O_{\widetilde X}^*)
=H^1(\widetilde X,\mathcal O_{\widetilde X}^*),
\]
therefore gives exactly
\[
0\longrightarrow H^0(X,\mathcal Q)
\longrightarrow\operatorname{Pic}(X)
\xrightarrow{\pi^*}\operatorname{Pic}(\widetilde X)
\longrightarrow0.
\]

(b) The normalization of either irreducible cubic is
\(\mathbb P^1\), whose Picard group is \(\mathbb Z\).

At the cusp one may choose a parameter \(t\) for which
\[
\mathcal O_{X,P}=k[t^2,t^3]_{(t^2,t^3)},\qquad
\widetilde{\mathcal O}_P=k[t]_{(t)}.
\]
The conductor contains \(t^2k[t]_{(t)}\).  Consequently the quotient
of unit groups is measured by the coefficient of \(t\):
\[
\widetilde{\mathcal O}_P^*/\mathcal O_{X,P}^*
\xrightarrow{\sim}(k,+),\qquad
u\longmapsto\frac{u'(0)}{u(0)}.
\]
The map is a homomorphism by the product rule, is onto using
\(1+ct\), and has kernel precisely the units with no linear term,
which are the units of \(\mathcal O_{X,P}\).  Part (a) gives the
cuspidal sequence.

At the node the normalization has two points above \(P\).  With a
parameter \(t\) taking the distinct values \(0\) and \(1\) there, the local
ring consists of functions regular at both points and having equal
values.  Thus
\[
\widetilde{\mathcal O}_P^*/\mathcal O_{X,P}^*
\xrightarrow{\sim}k^*,\qquad
u\longmapsto\frac{u(0)}{u(1)}.
\]
Its kernel is exactly the local unit group downstairs, and
surjectivity follows by using the polynomial
\(u(t)=\lambda+(1-\lambda)t\) for the ratio \(\lambda\in k^*\).
It is a unit at both branch points; for \(\lambda=1\) it is constant.  Applying (a)
again gives the nodal sequence.
\end{solution}

\begin{exerciseintro}
% record: H-II-06-010
\exerciseheading{Exercise II.6.10}
\addcontentsline{toc}{subsection}{Exercise II.6.10}

\begin{context}
For a noetherian scheme \(X\), let \(K(X)\) be the abelian group
generated by symbols \([\mathcal F]\) for coherent
\(\mathcal O_X\)-modules, subject to
\[
[\mathcal F]=[\mathcal F']+[\mathcal F'']
\]
for every exact sequence
\(0\to\mathcal F'\to\mathcal F\to\mathcal F''\to0\).
This is the Grothendieck group of coherent sheaves, often denoted
\(G_0(X)\).
The category of coherent sheaves on an open \(U\subseteq X\)
is the Serre quotient of \(\operatorname{Coh}(X)\) by coherent
sheaves supported in \(X\setminus U\).  Morphisms in this quotient
allow kernels and cokernels supported in that complement.
\end{context}

\begin{problem}
(a) Compute \(K(\mathbb A_k^1)\).

(b) If \(X\) is integral, prove that generic rank induces a surjective
homomorphism
\[
\operatorname{rank}:K(X)\longrightarrow\mathbb Z.
\]

(c) For a closed subscheme \(i:Y\hookrightarrow X\), with open
complement \(j:U\hookrightarrow X\), prove exactness of
\[
K(Y)\xrightarrow{i_*}K(X)\xrightarrow{j^*}K(U)\longrightarrow0.
\]
\end{problem}
\end{exerciseintro}

\begin{solution}
(a) Write \(\mathbb A_k^1=\operatorname{Spec}k[t]\).  Every finitely
generated \(k[t]\)-module \(M\) admits an exact sequence
\[
0\longrightarrow k[t]^s\longrightarrow k[t]^r
\longrightarrow M\longrightarrow0,
\]
because the kernel of a map between finite free modules over the
principal ideal domain \(k[t]\) is finite free.  Therefore
\[
[M]=(r-s)[k[t]],
\]
so \(K(\mathbb A_k^1)\) is cyclic.  Generic rank sends
\([k[t]]\) to \(1\), and hence proves
\[
K(\mathbb A_k^1)\cong\mathbb Z.
\]

(b) Let \(\eta\) be the generic point and put
\(K_\eta=\mathcal O_{X,\eta}\), the function field.  Taking the
stalk at \(\eta\) is exact, so any exact sequence of coherent sheaves
gives an exact sequence of finite-dimensional \(K_\eta\)-vector spaces.
Dimension is additive, and hence
\[
[\mathcal F]\longmapsto
\dim_{K_\eta}\mathcal F_\eta
\]
respects all defining relations of \(K(X)\).  It is onto because
\(\mathcal O_X\) has rank one.

(c) Both maps are well defined: extension by zero along a closed
immersion and restriction to an open subscheme are exact and preserve
coherence.  Their composite is zero.  Moreover, every coherent sheaf
on \(U\) extends to a coherent sheaf on the noetherian scheme \(X\),
so \(j^*\) is surjective on Grothendieck groups.

Let \(\operatorname{Coh}_Y(X)\) be the full subcategory of coherent
sheaves whose set-theoretic support lies in \(Y\).  It is a Serre
subcategory: it is closed under subobjects, quotients, and
extensions.  Restriction realizes
\[
\operatorname{Coh}(U)
\cong\operatorname{Coh}(X)/\operatorname{Coh}_Y(X).
\tag{1}
\]
Essential surjectivity is the extension theorem just used.  A
morphism on \(U\) extends after replacing its source and target by
subsheaves with quotient supported on \(Y\); taking the coherent
closure of its graph gives the morphism statement in (1).

For any abelian category \(\mathcal A\) and Serre subcategory
\(\mathcal B\), the generators-and-relations definition gives an
exact sequence
\[
K_0(\mathcal B)\longrightarrow K_0(\mathcal A)
\longrightarrow K_0(\mathcal A/\mathcal B)\longrightarrow0.
\tag{2}
\]
Indeed, objects of the quotient have the same generators, while a
sequence that becomes exact in the quotient can be corrected by
kernels and cokernels belonging to \(\mathcal B\).  Applying (2) to
(1), the kernel of \(K(X)\to K(U)\) is generated by classes of
coherent sheaves supported on \(Y\).

It remains to see that all these classes come from \(K(Y)\).  Let
\(\mathcal I\) define \(Y\), and let \(\mathcal F\) be coherent with
support in \(Y\).  Noetherianity and the support condition imply
\(\mathcal I^N\mathcal F=0\) for some \(N\).  The finite filtration
\[
\mathcal F\supseteq\mathcal I\mathcal F\supseteq\cdots
\supseteq\mathcal I^N\mathcal F=0
\]
has quotients annihilated by \(\mathcal I\), hence coherent
\(\mathcal O_Y\)-modules.  Additivity expresses \([\mathcal F]\) as a
sum of their extensions by zero.  Thus the kernel is exactly the
image of \(K(Y)\), proving the sequence.
\end{solution}

\begin{exerciseintro}
% record: H-II-06-011
\exerciseheading{Exercise II.6.11}
\addcontentsline{toc}{subsection}{Exercise II.6.11}

\begin{context}
Let \(X\) be a nonsingular integral curve over an algebraically closed
field.  Write \(K(X)\) for the Grothendieck group of coherent sheaves,
and \([\mathcal F]\) for the class of \(\mathcal F\).  On this regular
curve, \(\operatorname{Cl}(X)=\operatorname{Pic}(X)\), and every local
ring at a closed point is a discrete valuation ring.
\end{context}

\begin{problem}
(a) For a divisor \(D=\sum n_PP\), define
\[
\psi(D)=\sum_P n_P[k(P)]\in K(X).
\]
If \(D\) is effective, prove
\(\psi(D)=[\mathcal O_D]\).  Show that \(\psi\) depends only on the
linear-equivalence class and hence defines
\(\psi:\operatorname{Pic}(X)\to K(X)\).

(b) Prove that every coherent sheaf has a locally free resolution
\[
0\longrightarrow\mathcal E_1\longrightarrow\mathcal E_0
\longrightarrow\mathcal F\longrightarrow0.
\]
Define
\[
\det\mathcal F
=\bigwedge^{r_0}\mathcal E_0
\otimes
\left(\bigwedge^{r_1}\mathcal E_1\right)^{-1}.
\]
Prove independence of the resolution, additivity in short exact
sequences, and
\[
\det(\psi(D))\cong\mathcal O_X(D).
\]

(c) If \(\mathcal F\) has rank \(r\), construct a divisor \(D\), a
torsion sheaf \(\mathcal T\), and an exact sequence
\[
0\longrightarrow\mathcal O_X(D)^{\oplus r}
\longrightarrow\mathcal F\longrightarrow\mathcal T
\longrightarrow0.
\]
Deduce that
\([\mathcal F]-r[\mathcal O_X]\) lies in the image of \(\psi\).

(d) Use rank, determinant, and \(\psi\) to prove
\[
K(X)\cong\operatorname{Pic}(X)\oplus\mathbb Z.
\]
\end{problem}
\end{exerciseintro}

\begin{solution}
(a) At a point \(P\), the local ring is a discrete valuation ring.
If \(D=\sum n_PP\) is effective, the stalk of \(\mathcal O_D\) at
\(P\) has length \(n_P\) and a filtration whose successive quotients
are \(k(P)\).  Additivity in \(K(X)\) gives
\[
[\mathcal O_D]=\sum_Pn_P[k(P)]=\psi(D).
\tag{1}
\]
The ideal sequence of an effective divisor is
\[
0\longrightarrow\mathcal O_X(-D)\longrightarrow\mathcal O_X
\longrightarrow\mathcal O_D\longrightarrow0,
\]
so
\[
\psi(D)=[\mathcal O_X]-[\mathcal O_X(-D)].
\tag{2}
\]
If \(D\sim D'\), choose an effective divisor \(E\) such that both
\(D+E\) and \(D'+E\) are effective.  The corresponding invertible
sheaves \(\mathcal O_X(-D-E)\) and
\(\mathcal O_X(-D'-E)\) are isomorphic, so (2) gives
\(\psi(D+E)=\psi(D'+E)\), and cancellation of \(\psi(E)\) gives
\(\psi(D)=\psi(D')\).  Thus \(\psi\) descends to the Picard group.

(b) A nonsingular curve is regular and quasi-projective.  By clearing
poles of finitely many generators on a finite affine cover, one gets
a surjection from a finite direct sum of invertible sheaves
\(\mathcal E_0\twoheadrightarrow\mathcal F\).  Let
\(\mathcal E_1\) be its kernel.  At each closed point,
\((\mathcal E_1)_P\) is a submodule of a finite free module over a
discrete valuation ring, hence is free.  Coherence then makes
\(\mathcal E_1\) locally free.  This gives the required resolution.

Suppose two such resolutions use
\((\mathcal E_1,\mathcal E_0)\) and
\((\mathcal E'_1,\mathcal E'_0)\).  Their fiber product over
\(\mathcal F\) is a locally free sheaf \(\mathcal E\), fitting into
\[
0\to\mathcal E_1\to\mathcal E\to\mathcal E'_0\to0,
\qquad
0\to\mathcal E'_1\to\mathcal E\to\mathcal E_0\to0.
\]
Both sequences split locally because their quotients are locally
free.  Multiplicativity of top exterior powers gives
\[
\det\mathcal E_1\otimes\det\mathcal E'_0
\cong\det\mathcal E
\cong\det\mathcal E'_1\otimes\det\mathcal E_0,
\]
which proves independence of the resolution.

Now let
\(0\to\mathcal F'\to\mathcal F\to\mathcal F''\to0\) be exact.
Choose a locally free surjection
\(\mathcal E\twoheadrightarrow\mathcal F\), and put
\[
\mathcal K=\ker(\mathcal E\to\mathcal F),\qquad
\mathcal K''=\ker(\mathcal E\to\mathcal F'').
\]
Both kernels are locally free, and there are resolutions
\[
0\to\mathcal K\to\mathcal K''\to\mathcal F'\to0,\quad
0\to\mathcal K\to\mathcal E\to\mathcal F\to0,\quad
0\to\mathcal K''\to\mathcal E\to\mathcal F''\to0.
\]
Substitution in the definition yields
\[
\det\mathcal F
\cong\det\mathcal F'\otimes\det\mathcal F''.
\]
Thus determinant defines a homomorphism
\(\det:K(X)\to\operatorname{Pic}(X)\).

For effective \(D\), the ideal sequence above is a locally free
resolution of \(\mathcal O_D\), so
\[
\det(\psi(D))=\det\mathcal O_D
\cong\mathcal O_X(D).
\]
Additivity extends this identity to every divisor.

(c) Choose a basis \(s_1,\ldots,s_r\) of the generic fiber
\(\mathcal F_\eta\).  These are rational sections of \(\mathcal F\).
Their poles occur at finitely many closed points, so a sufficiently
large effective divisor \(E\) clears all of them and produces a map
\[
\mathcal O_X(-E)^{\oplus r}\longrightarrow\mathcal F
\]
that is an isomorphism at the generic point.  Its kernel has generic
rank zero but is a subsheaf of a locally free sheaf, which is
torsion-free; hence the kernel is zero.  Its cokernel
\(\mathcal T\) has generic stalk zero and is therefore a torsion
sheaf.  Taking \(D=-E\) proves the asserted sequence.

For the divisor \(D=-E\) just constructed, equation (2) for the effective \(E\) gives
\[
[\mathcal O_X(D)]-[\mathcal O_X]=\psi(D).
\tag{3}
\]
Moreover, a torsion sheaf has finite length at finitely many points,
so a composition series gives
\[
[\mathcal T]
=\psi\left(\sum_P\operatorname{length}_{\mathcal O_{X,P}}
(\mathcal T_P)P\right).
\]
The sequence in (c) and (3) now show that
\[
[\mathcal F]-r[\mathcal O_X]
=r\psi(D)+[\mathcal T]\in\operatorname{im}\psi.
\]

(d) Define
\[
\Theta:K(X)\longrightarrow\operatorname{Pic}(X)\oplus\mathbb Z,
\qquad
\alpha\longmapsto(\det\alpha,\operatorname{rank}\alpha),
\]
and
\[
\Psi:\operatorname{Pic}(X)\oplus\mathbb Z\longrightarrow K(X),
\qquad
(\mathcal O_X(D),n)\longmapsto\psi(D)+n[\mathcal O_X].
\]
The map \(\Psi\) is well defined by (a).  Part (b) and the fact that
\(\psi(D)\) has rank zero show
\(\Theta\Psi=\operatorname{id}\).  Part (c) says every class in
\(K(X)\) is \(n[\mathcal O_X]+\psi(D)\); applying determinant identifies
\(\mathcal O_X(D)\) with its first component.  Hence
\(\Psi\Theta=\operatorname{id}\) as well, proving the claimed
isomorphism.
\end{solution}

\begin{exerciseintro}
% record: H-II-06-012
\exerciseheading{Exercise II.6.12}
\addcontentsline{toc}{subsection}{Exercise II.6.12}

\begin{context}
Let \(X\) be a complete nonsingular integral curve over an
algebraically closed field.  Exercise II.6.11 gives
\[
K(X)\cong\operatorname{Pic}(X)\oplus\mathbb Z
\]
through determinant and rank.
\end{context}

\begin{problem}
Prove that there is a unique integer-valued degree for coherent
sheaves on \(X\) satisfying all three properties:

(1) \(\deg\mathcal O_X(D)=\deg D\) for every divisor \(D\);

(2) if \(\mathcal T\) is torsion, then
\[
\deg\mathcal T
=\sum_{P\in X}\operatorname{length}_{\mathcal O_{X,P}}(\mathcal T_P);
\]

(3) degree is additive in every short exact sequence of coherent
sheaves.
\end{problem}
\end{exerciseintro}

\begin{solution}
Additivity means precisely that the desired function must factor
through a homomorphism
\[
\deg:K(X)\longrightarrow\mathbb Z.
\]
Under the decomposition of Exercise II.6.11, define
\[
\deg(\alpha)=\deg(\det\alpha).
\tag{1}
\]
Here the degree on the right is the usual degree homomorphism
\(\operatorname{Pic}(X)\to\mathbb Z\).  Since determinant is
multiplicative in short exact sequences, (1) is additive.

For a line bundle \(\mathcal O_X(D)\), its determinant is itself, so
(1) gives property (1).  If \(\mathcal T\) is torsion, a composition
series over the discrete valuation rings gives
\[
[\mathcal T]
=\sum_P
\operatorname{length}_{\mathcal O_{X,P}}(\mathcal T_P)[k(P)].
\]
Exercise II.6.11 identifies
\(\det[k(P)]\) with \(\mathcal O_X(P)\), of degree one.  Formula (1)
therefore gives property (2).

Finally, uniqueness follows from the same decomposition.  Every class
has the form
\[
\psi(D)+r[\mathcal O_X].
\]
Properties (1) and (2) force its degree to be \(\deg D\), while
property (1) with \(D=0\) forces
\(\deg\mathcal O_X=0\).  Additivity then determines the value on every
class, so no other degree function can satisfy the three conditions.
\end{solution}

\begin{exerciseintro}
\corpussection{II.7 Projective Morphisms}
\addcontentsline{toc}{section}{II.7 Projective Morphisms}

% record: H-II-07-001
\exerciseheading{Exercise II.7.1}
\addcontentsline{toc}{subsection}{Exercise II.7.1}

\begin{context}
Let \((X,\mathcal O_X)\) be a locally ringed space.  An invertible
sheaf has stalk isomorphic to \(\mathcal O_{X,x}\) at every point
\(x\).
\end{context}

\begin{problem}
Prove that every surjective morphism
\[
\mathcal L\longrightarrow\mathcal M
\]
between invertible \(\mathcal O_X\)-modules is an isomorphism.
\end{problem}
\end{exerciseintro}

\begin{solution}
It is enough to check the morphism on stalks.  Fix \(x\in X\), choose
bases of the free rank-one modules \(\mathcal L_x\) and
\(\mathcal M_x\), and identify the stalk map with multiplication by
some \(a\in\mathcal O_{X,x}\).  Surjectivity gives
\(b\in\mathcal O_{X,x}\) with \(ab=1\).  Thus \(a\) is a unit and
multiplication by \(a\) is an isomorphism.  Every stalk map is
therefore an isomorphism, so the original sheaf morphism is an
isomorphism.
\end{solution}

\begin{exerciseintro}
% record: H-II-07-002
\exerciseheading{Exercise II.7.2}
\addcontentsline{toc}{subsection}{Exercise II.7.2}

\begin{context}
Let \(X\) be a \(k\)-scheme and \(\mathcal L\) an invertible sheaf.
Suppose
\[
s_0,\ldots,s_n,\qquad t_0,\ldots,t_m
\]
are two base-point-free lists of global sections, both spanning the
same vector subspace \(V\subseteq\Gamma(X,\mathcal L)\), with
\(n\leq m\).  They define morphisms
\(\phi:X\to\mathbb P_k^n\) and
\(\psi:X\to\mathbb P_k^m\).
When \(m=n\), the projection center of dimension \(-1\) means the empty set.
\end{context}

\begin{problem}
Show that, after an automorphism of \(\mathbb P^n\), the morphism
\(\phi\) is the composite of \(\psi\) with a linear projection
\[
\mathbb P^m\setminus L\longrightarrow\mathbb P^n,
\]
where \(L\) is a linear subspace of dimension \(m-n-1\), and prove
that \(\psi(X)\cap L=\varnothing\).
\end{problem}
\end{exerciseintro}

\begin{solution}
Let
\[
q_s:k^{n+1}\twoheadrightarrow V,\qquad
q_t:k^{m+1}\twoheadrightarrow V
\]
send the standard bases to the two lists of sections.  Choose
complements to their kernels.  Since
\[
\dim\ker q_t=m+1-\dim V\geq n+1-\dim V=\dim\ker q_s,
\]
we can lift \(q_s\) through \(q_t\) by an injective linear map
\[
B:k^{n+1}\hookrightarrow k^{m+1},
\qquad q_tB=q_s.
\tag{1}
\]
Equivalently, the \(s_i\) can be written as linear combinations
\[
s_i=\sum_{j=0}^m b_{ji}t_j
\]
whose \(n+1\) coefficient vectors are linearly independent.  A change
of basis in \(k^{n+1}\) accounts for the allowed automorphism of the
target \(\mathbb P^n\).

Let \(\ell_0,\ldots,\ell_n\) be the corresponding independent linear
forms on the homogeneous coordinates of \(\mathbb P^m\).  Their common
zero locus
\[
L=V_+(\ell_0,\ldots,\ell_n)
\]
has dimension \(m-n-1\), and the tuple
\([\ell_0:\cdots:\ell_n]\) defines the linear projection away from
\(L\).  Pulling these forms back along \(\psi\) gives exactly the
sections \(s_0,\ldots,s_n\).  They generate \(\mathcal L\) at every
point, so they never vanish simultaneously; hence
\(\psi(X)\cap L=\varnothing\).  The projected morphism is therefore
defined on all of \(\psi(X)\), and its homogeneous coordinates are
those defining \(\phi\).  This proves the assertion.
\end{solution}

\begin{exerciseintro}
% record: H-II-07-003
\exerciseheading{Exercise II.7.3}
\addcontentsline{toc}{subsection}{Exercise II.7.3}

\begin{context}
Let \(n\geq1\) and let
\(\phi:\mathbb P_k^n\to\mathbb P_k^m\) be a morphism.  If the image
lies in a proper linear subspace of \(\mathbb P^m\), a factorization
by a linear projection naturally ends with the inclusion of that
linear span.  Thus the usual direct projection statement is read
after replacing \(\mathbb P^m\) by the linear span of the image.
\end{context}

\begin{problem}
(a) Prove that either \(\phi\) is constant, or
\[
m\geq n,\qquad \dim\phi(\mathbb P^n)=n.
\]

(b) In the nonconstant case, prove that there is a unique integer
\(d\geq1\) such that \(\phi\) is induced by base-point-free degree
\(d\) forms.  Factor it through the \(d\)-uple Veronese embedding,
a linear projection, and a change of projective coordinates, with
the linear-span qualification stated above.  Prove also that
\(\phi\) is finite and hence has finite fibers.
\end{problem}
\end{exerciseintro}

\begin{solution}
(a) The morphism is represented by an invertible sheaf
\[
\phi^*\mathcal O_{\mathbb P^m}(1)\cong\mathcal O_{\mathbb P^n}(d)
\]
and \(m+1\) global sections \(p_0,\ldots,p_m\) having no common zero.
There are no nonzero sections when \(d<0\).  If \(d=0\), all sections
are constants and the morphism is constant.  We henceforth assume
\(d>0\).

Put \(S=k[y_0,\ldots,y_n]\).  Since the homogeneous ideal
\(I=(p_0,\ldots,p_m)\) has empty projective zero locus, its radical
contains \(S_+=(y_0,\ldots,y_n)\).  Consequently
\[
S_q\subseteq I\qquad(q\geq N)
\tag{1}
\]
for some \(N\).  On the target chart \(D_+(T_i)\), the inverse image is
\[
D_+(p_i)=\operatorname{Spec}(S_{p_i})_0.
\]
Let
\[
B_i=k[p_0/p_i,\ldots,p_m/p_i]\subseteq(S_{p_i})_0.
\]
If \(h\in S_{qd}\) and \(qd\geq N\), (1) gives homogeneous
\(g_j\in S_{(q-1)d}\) such that \(h=\sum_jp_jg_j\).  Hence
\[
\frac{h}{p_i^q}
=\sum_j\frac{p_j}{p_i}\frac{g_j}{p_i^{q-1}}.
\tag{2}
\]
Repeated use of (2) shows that \((S_{p_i})_0\) is generated as a
\(B_i\)-module by the finitely many fractions
\[
\frac{y^\alpha}{p_i^q},\qquad |\alpha|=qd<N.
\]
Thus every affine restriction of \(\phi\) is finite, and \(\phi\) is
a finite morphism.

Its image is therefore closed, and the induced finite dominant map
from \(\mathbb P^n\) to that image preserves dimension.  Hence
\(\dim\phi(\mathbb P^n)=n\), which forces \(m\geq n\).  This also
proves the finite-fiber assertion needed in (b).

(b) For the factorization, let
\[
N_d=\binom{n+d}{d}-1.
\]
All degree-\(d\) monomials define the Veronese closed immersion
\[
v_d:\mathbb P^n\hookrightarrow\mathbb P^{N_d}.
\]
The forms \(p_i\) are linear combinations of those monomials.  If
\(W\subseteq H^0(\mathbb P^n,\mathcal O(d))\) is their span and
\(\dim W=r\), these combinations define a linear projection
\[
\mathbb P^{N_d}\setminus L\longrightarrow\mathbb P^{r-1}
\]
whose center misses the Veronese image, because the \(p_i\) have no
common zero.  After choosing a basis of \(W\), the original map is
this projection followed by the linear inclusion
\(\mathbb P^{r-1}\hookrightarrow\mathbb P^m\) and an automorphism of
\(\mathbb P^m\).  When the image spans \(\mathbb P^m\), one has
\(r=m+1\), so this is the stated direct linear projection to
\(\mathbb P^m\).

Finally, \(d\) is unique because it is determined by the isomorphism
class of \(\phi^*\mathcal O_{\mathbb P^m}(1)\) in
\(\operatorname{Pic}(\mathbb P^n)\cong\mathbb Z\).
\end{solution}

\begin{exerciseintro}
% record: H-II-07-004
\exerciseheading{Exercise II.7.4}
\addcontentsline{toc}{subsection}{Exercise II.7.4}

\begin{context}
Let \(A\) be noetherian.  We use the criterion that an invertible
sheaf on a finite-type \(A\)-scheme is ample exactly when some positive
power gives an immersion into a projective space over \(A\).

The affine line with doubled origin is obtained by gluing two copies
\(U,V\cong\mathbb A_k^1\) by the identity on
\(U\cap V\cong\mathbb G_m\).
\end{context}

\begin{problem}
(a) Prove that a finite-type \(A\)-scheme possessing an ample
invertible sheaf is separated over \(A\).

(b) For the affine line with doubled origin, compute the Picard group
and determine exactly which invertible sheaves are generated by global
sections.  Use this computation directly to prove that the scheme has
no ample invertible sheaf.
\end{problem}
\end{exerciseintro}

\begin{solution}
(a) If \(\mathcal L\) is ample, the cited criterion supplies an
immersion
\[
X\longrightarrow\mathbb P_A^N.
\]
An immersion is separated, projective space is separated over \(A\),
and a composite of separated morphisms is separated.  Hence
\(X\to\operatorname{Spec}A\) is separated.

(b) Every invertible sheaf on either \(U\) or \(V\) is trivial.
After choosing trivializations, gluing is multiplication on
\(\mathbb G_m\) by a unit of
\[
k[x,x^{-1}]^*=k^*x^{\mathbb Z}.
\]
Changing either trivialization multiplies this unit by an element of
\(k^*\).  Therefore
\[
\operatorname{Pic}(X)
\cong k[x,x^{-1}]^*/(k[x]^*k[x]^*)
\cong\mathbb Z.
\]
Let \(\mathcal L_r\) denote the class represented by \(x^r\).

Choose local bases so that \(e_U=x^re_V\) on the overlap.  A global
section is a pair \(f e_U,g e_V\), with \(f,g\in k[x]\), satisfying
\[
g=x^rf
\quad\text{in }k[x,x^{-1}].
\]
If \(r>0\), every such \(g\) vanishes at the origin in \(V\), so global
sections do not generate the fiber at that origin.  If \(r<0\), every
\(f\) vanishes at the origin in \(U\).  For \(r=0\), the sheaf is
trivial and is generated by \(1\).  Thus
\[
\mathcal L_r\text{ is globally generated}
\quad\Longleftrightarrow\quad r=0.
\tag{1}
\]

Suppose \(\mathcal L_r\) were ample.  Apply the global-generation
definition of ampleness to \(\mathcal L_1\).  It would require
\[
\mathcal L_1\otimes\mathcal L_r^{\otimes n}
=\mathcal L_{1+nr}
\]
to be globally generated for every sufficiently large \(n\).
By (1), this would force \(1+nr=0\) for every sufficiently large
\(n\), which is impossible.  Hence no invertible sheaf on the doubled
line is ample.
\end{solution}

\begin{exerciseintro}
% record: H-II-07-005
\exerciseheading{Exercise II.7.5}
\addcontentsline{toc}{subsection}{Exercise II.7.5}

\begin{context}
Let \(X\) be noetherian and let \(\mathcal L,\mathcal M\) be invertible
sheaves.  In parts (d) and (e), assume that \(X\) is of finite type
over a noetherian ring \(A\), and interpret very ample relative to
\(\operatorname{Spec}A\), using an immersion into a single
\(\mathbb P_A^N\) to define very ampleness.

We repeatedly use the elementary fact that the tensor product of two
globally generated sheaves is globally generated: tensor products of
generating global sections generate every stalk.
\end{context}

\begin{problem}
Prove:

(a) if \(\mathcal L\) is ample and \(\mathcal M\) is globally
generated, then \(\mathcal L\otimes\mathcal M\) is ample;

(b) if \(\mathcal L\) is ample and \(\mathcal M\) is arbitrary, then
\(\mathcal M\otimes\mathcal L^n\) is ample for all sufficiently large
\(n\);

(c) the tensor product of two ample invertible sheaves is ample;

(d) if \(\mathcal L\) is very ample and \(\mathcal M\) is globally
generated, then \(\mathcal L\otimes\mathcal M\) is very ample;

(e) if \(\mathcal L\) is ample, then \(\mathcal L^n\) is very ample
for every sufficiently large \(n\).
\end{problem}
\end{exerciseintro}

\begin{solution}
(a) Let \(\mathcal F\) be coherent.  For \(q\gg0\),
\(\mathcal F\otimes\mathcal L^q\) is globally generated.  So is
\(\mathcal M^q\), because \(\mathcal M\) is globally generated.
Therefore
\[
\mathcal F\otimes(\mathcal L\otimes\mathcal M)^q
\cong
(\mathcal F\otimes\mathcal L^q)\otimes\mathcal M^q
\]
is globally generated for \(q\gg0\).  This is the
global-generation criterion for ampleness.

(b) Ampleness of \(\mathcal L\), applied to \(\mathcal M\), gives an
\(N\) such that
\(\mathcal M\otimes\mathcal L^q\) is globally generated for every
\(q\geq N\).  For \(n\geq N+1\), write
\[
\mathcal M\otimes\mathcal L^n
=\mathcal L\otimes
(\mathcal M\otimes\mathcal L^{n-1}).
\]
Part (a) makes this sheaf ample.

(c) For a coherent \(\mathcal F\), the sheaves
\(\mathcal F\otimes\mathcal L^q\) and \(\mathcal M^q\) are both
globally generated once \(q\) is large.  Their tensor product is
\(\mathcal F\otimes(\mathcal L\otimes\mathcal M)^q\), so the same
criterion proves that \(\mathcal L\otimes\mathcal M\) is ample.

(d) Choose sections \(s_0,\ldots,s_r\) realizing a very ample
immersion for \(\mathcal L\), and finite generating sections
\(t_0,\ldots,t_s\) for \(\mathcal M\).  They define
an immersion and a morphism
\[
i:X\longrightarrow\mathbb P_A^r,\qquad
j:X\longrightarrow\mathbb P_A^s.
\]
The products \(s_at_b\) define the composite
\[
X\xrightarrow{(i,j)}
\mathbb P_A^r\times_A\mathbb P_A^s
\xrightarrow{\mathrm{Segre}}
\mathbb P_A^{(r+1)(s+1)-1}.
\tag{1}
\]
The map \((i,j)\) is an immersion: it factors through the closed graph
of \(j\) in \(X\times_A\mathbb P_A^s\), followed by the base change of
the immersion \(i\).  The Segre map is a closed immersion, so (1) is
an immersion.  Its pulled-back \(\mathcal O(1)\) is
\(\mathcal L\otimes\mathcal M\), proving very ampleness.

(e) Some power \(\mathcal L^a\) is very ample by the ample-sheaf
criterion.  Also, there is \(b\) such that \(\mathcal L^q\) is
globally generated for every \(q\geq b\).  If \(n\geq a+b\), then
\[
\mathcal L^n=\mathcal L^a\otimes\mathcal L^{n-a},
\]
and part (d) makes it very ample.
\end{solution}

\begin{exerciseintro}
% record: H-II-07-006
\exerciseheading{Exercise II.7.6}
\addcontentsline{toc}{subsection}{Exercise II.7.6}

\begin{context}
Let \(X\) be a nonsingular projective integral variety over an
algebraically closed field, let \(D\) be a divisor, and put
\(\mathcal L=\mathcal O_X(D)\).  We use the convention
\[
\dim|qD|=h^0(X,\mathcal L^q)-1,
\]
so an empty linear system has dimension \(-1\).
\end{context}

\begin{problem}
(a) If \(D\) is very ample and gives an embedding
\(X\hookrightarrow\mathbb P^N\), prove that
\[
\dim|qD|=P_X(q)-1
\]
for all sufficiently large \(q\), where \(P_X\) is the Hilbert
polynomial of this embedded variety.

(b) If the class of \(D\) has exact order \(r\) in
\(\operatorname{Pic}(X)\), prove
\[
\dim|qD|=
\begin{cases}
0,&r\text{ divides }q,\\
-1,&r\text{ does not divide }q.
\end{cases}
\]
Thus the function is periodic with period \(r\).
\end{problem}
\end{exerciseintro}

\begin{solution}
(a) Under the given embedding,
\(\mathcal L\cong\mathcal O_X(1)\), and hence
\(\mathcal L^q\cong\mathcal O_X(q)\).  If \(S\) is the homogeneous
coordinate ring of \(X\), the natural map
\[
S_q\longrightarrow\Gamma(X,\mathcal O_X(q))
\]
is an isomorphism for \(q\gg0\).  In the same range the Hilbert
function \(\dim_kS_q\) equals \(P_X(q)\).  Therefore
\[
\dim|qD|
=h^0(X,\mathcal O_X(q))-1
=P_X(q)-1.
\]

(b) If \(r\) divides \(q\), then
\(\mathcal L^q\cong\mathcal O_X\).  Since \(X\) is projective and
integral,
\(\Gamma(X,\mathcal O_X)=k\), so the linear system has dimension zero.

Conversely, suppose \(s\neq0\) is a global section of
\(\mathcal L^q\).  The line bundle \(\mathcal L^q\) is torsion.  If
its order is \(a\), then, after choosing a trivialization
\((\mathcal L^q)^a\cong\mathcal O_X\), the tensor power \(s^a\) is a
nonzero global regular function.  It is therefore a nonzero constant.
Thus \(s\) vanishes nowhere and trivializes \(\mathcal L^q\).  Exact
order \(r\) now implies \(r\mid q\).  Hence there are no nonzero
sections when \(r\not\mid q\), giving dimension \(-1\).
\end{solution}

\begin{exerciseintro}
% record: H-II-07-007
\exerciseheading{Exercise II.7.7}
\addcontentsline{toc}{subsection}{Exercise II.7.7}

\begin{context}
Let \(k\) be algebraically closed.  In part (b) assume
\(\operatorname{char}k\neq2\), which is needed for the stated
five-dimensional subsystem to be an immersion.  On
\(\mathbb P_k^2\), with coordinates \(x,y,z\), the complete system of
conics is \(H^0(\mathcal O_{\mathbb P^2}(2))\).
\end{context}

\begin{problem}
(a) Show that the six quadratic monomials give the Veronese closed
immersion \(\mathbb P^2\hookrightarrow\mathbb P^5\).

(b) Prove that the five sections
\[
x^2,\quad y^2,\quad z^2,\quad y(x-z),\quad (x-y)z
\]
give a closed immersion \(\mathbb P^2\hookrightarrow\mathbb P^4\).

(c) Fix \(P\in\mathbb P^2(k)\).  Show that the conics through \(P\)
embed \(\mathbb P^2\setminus\{P\}\) in \(\mathbb P^4\), extend to a
closed immersion of
\(\widetilde X=\operatorname{Bl}_P\mathbb P^2\), and make
\(\widetilde X\) a degree-three surface.  Prove that strict transforms
of lines through \(P\) are mutually disjoint projective lines covering
\(\widetilde X\).
\end{problem}
\end{exerciseintro}

\begin{solution}
(a) The complete linear system
\(\lvert\mathcal O_{\mathbb P^2}(2)\rvert\) gives the second Veronese
map
\[
[x:y:z]\longmapsto
[x^2:y^2:z^2:xy:xz:yz].
\]
The general \(d\)-uple embedding theorem, with \(d=2\), says that this
is a closed immersion.

(b) View the map as projection of the Veronese surface.  In coordinates
\[
[X_0:X_1:X_2:X_3:X_4:X_5]
=[x^2:y^2:z^2:xy:xz:yz],
\]
the five sections are
\[
X_0,\ X_1,\ X_2,\ X_3-X_5,\ X_4-X_5.
\]
Thus the center of projection is
\[
c=[0:0:0:1:1:1]\in\mathbb P^5.
\]
Under the symmetric-matrix description of the quadratic Veronese
embedding, its secant variety consists of symmetric matrices of rank
at most \(2\).  The point \(c\) corresponds, up to harmless nonzero
scalars in the off-diagonal entries, to
\[
\begin{pmatrix}
0&1&1\\
1&0&1\\
1&1&0
\end{pmatrix},
\]
whose determinant is \(2\).  Hence \(c\) is outside the secant variety
when \(\operatorname{char}k\neq2\).

Projection from a point outside the secant variety restricts to a
closed immersion: it separates every length-two subscheme, since
failure to separate two points or a tangent vector would put the
center on the corresponding secant or tangent line.  Therefore the
five displayed sections give a closed immersion into \(\mathbb P^4\).

(c) Move \(P\) to \([0:0:1]\).  The conics through \(P\) are spanned by
\[
x^2,\quad xy,\quad y^2,\quad xz,\quad yz.
\tag{1}
\]
The blowup has the incidence presentation
\[
\widetilde X
=\{([x:y:z],[s:t])\in\mathbb P^2\times\mathbb P^1:
xt=ys\}.
\tag{2}
\]
Away from \(P\), the second coordinate is the direction
\([s:t]=[x:y]\); over \(P\), it records the exceptional direction.
The Segre embedding of the product in (2) has coordinates
\[
xs,\ xt,\ ys,\ yt,\ zs,\ zt.
\]
On \(\widetilde X\), the relation \(xt=ys\) places the image in a
linear \(\mathbb P^4\), and substituting \([s:t]=[x:y]\) recovers,
up to order, exactly the five sections in (1).  Since (2) and the
Segre map are closed immersions, the induced map
\(\widetilde X\hookrightarrow\mathbb P^4\) is a closed immersion.  On
the complement of the exceptional curve it is the map defined by the
conics through \(P\).

Let \(H\) be the pullback of a line and \(E\) the exceptional curve.
The hyperplane class of this embedding is \(2H-E\).  Using
\[
H^2=1,\qquad H\mathbin{.}E=0,\qquad E^2=-1,
\]
its degree is
\[
(2H-E)^2=4-1=3.
\]
A line through \(P\) has strict-transform class \(H-E\), and
\[
(2H-E)\mathbin{.}(H-E)=1.
\]
It therefore maps isomorphically to a line in \(\mathbb P^4\).
Distinct such transforms are disjoint because their only intersection
downstairs was \(P\), and they meet \(E\) at distinct direction
points.  Every point of (2) has a direction \([s:t]\), so these lines
cover \(\widetilde X\).  Thus the embedded cubic surface is ruled.
\end{solution}

\begin{exerciseintro}
% record: H-II-07-008
\exerciseheading{Exercise II.7.8}
\addcontentsline{toc}{subsection}{Exercise II.7.8}

\begin{context}
Let \(X\) be noetherian, let \(\mathcal E\) be a finite locally free
\(\mathcal O_X\)-module, and use the quotient convention
\[
\mathbb P(\mathcal E)
=\mathbf{Proj}_X(\operatorname{Sym}\mathcal E).
\]
If \(\pi:\mathbb P(\mathcal E)\to X\), there is a tautological
surjection
\[
\pi^*\mathcal E\twoheadrightarrow\mathcal O_{\mathbb P(\mathcal E)}(1).
\]
\end{context}

\begin{problem}
Prove a natural bijection between sections
\(\sigma:X\to\mathbb P(\mathcal E)\) of \(\pi\) and isomorphism
classes of surjections
\[
\mathcal E\twoheadrightarrow\mathcal L
\]
with \(\mathcal L\) invertible.
\end{problem}
\end{exerciseintro}

\begin{solution}
Given a section \(\sigma\), pull back the tautological quotient:
\[
\mathcal E
\cong\sigma^*\pi^*\mathcal E
\twoheadrightarrow
\sigma^*\mathcal O_{\mathbb P(\mathcal E)}(1).
\]
The target
\(\mathcal L=\sigma^*\mathcal O_{\mathbb P(\mathcal E)}(1)\) is
invertible, so this produces a quotient line bundle on \(X\).

Conversely, a surjection
\(\mathcal E\twoheadrightarrow\mathcal L\) induces a surjective map of
graded algebras
\[
\operatorname{Sym}\mathcal E
\longrightarrow\operatorname{Sym}\mathcal L.
\]
The relative-Proj universal property gives an \(X\)-morphism
\[
\sigma:X\cong\mathbb P(\mathcal L)
\longrightarrow\mathbb P(\mathcal E).
\]
Because it is an \(X\)-morphism, \(\pi\sigma=\operatorname{id}_X\).

These constructions commute with restriction to open sets.  Locally,
where \(\mathcal E\cong\mathcal O_X^{r+1}\), they are the familiar
bijection between a point of \(\mathbb P^r\) and a rank-one quotient
of the free module.  The local constructions are inverse and hence so
are the global ones.  Replacing a quotient by an isomorphic quotient
does not change the section, giving the claimed natural bijection.
\end{solution}

\begin{exerciseintro}
% record: H-II-07-009
\exerciseheading{Exercise II.7.9}
\addcontentsline{toc}{subsection}{Exercise II.7.9}

\begin{context}
Let \(X\) be a connected regular noetherian scheme and let
\(\mathcal E\) be finite locally free of rank \(r\geq2\).  Write
\(\pi:\mathbb P(\mathcal E)\to X\), using the quotient convention.
Connectedness is necessary for a single \(\mathbb Z\)-summand; for a
disconnected regular noetherian scheme, the result applies on each
open-and-closed connected component and the integer is a locally
constant function on \(X\).
\end{context}

\begin{problem}
(a) Prove
\[
\operatorname{Pic}(\mathbb P(\mathcal E))
\cong\operatorname{Pic}(X)\oplus\mathbb Z,
\qquad
(\mathcal M,d)\longmapsto
\pi^*\mathcal M\otimes\mathcal O_{\mathbb P(\mathcal E)}(d).
\]

(b) For another finite locally free sheaf \(\mathcal E'\), prove that
\[
\mathbb P(\mathcal E)\cong_X\mathbb P(\mathcal E')
\]
if and only if
\(\mathcal E'\cong\mathcal E\otimes\mathcal M\) for some invertible
sheaf \(\mathcal M\) on \(X\).
\end{problem}
\end{exerciseintro}

\begin{solution}
(a) First let \(U=\operatorname{Spec}A\) be connected, regular, and
affine, and suppose \(\mathcal E|_U\) is free.  Then
\(\mathbb P(\mathcal E)|_U\cong\mathbb P_A^{r-1}\).  Both \(U\) and
\(\mathbb P_A^{r-1}\) are regular and hence locally factorial.  The
class-group computation for projective space over an integral regular
base, together with \(\operatorname{Pic}=\operatorname{Cl}\), gives
\[
\operatorname{Pic}(\mathbb P_A^{r-1})
\cong\operatorname{Pic}(A)\oplus\mathbb Z,
\tag{1}
\]
where the second generator is \(\mathcal O(1)\).

Cover \(X\) by connected affine opens \(U_i\) trivializing
\(\mathcal E\).  For
\(\mathcal N\in\operatorname{Pic}(\mathbb P(\mathcal E))\), equation
(1) gives
\[
\mathcal N|_{\pi^{-1}(U_i)}
\cong\pi^*\mathcal M_i\otimes\mathcal O(k_i).
\]
On a fiber over a point of \(U_i\cap U_j\), both expressions have
degree \(k_i\) and \(k_j\); hence these integers agree whenever the
overlap is nonempty.  Connectedness makes them a single integer \(k\).

Put \(\mathcal N_0=\mathcal N\otimes\mathcal O(-k)\).  Locally it is a
pullback from \(U_i\).  Therefore
\[
\mathcal M=\pi_*\mathcal N_0
\]
is invertible, and the adjunction map
\(\pi^*\mathcal M\to\mathcal N_0\) is an isomorphism; both assertions
can be checked on the trivializing opens, using
\(\pi_*\mathcal O_{\mathbb P^{r-1}_{U_i}}=\mathcal O_{U_i}\).
This proves surjectivity of the displayed homomorphism.

For uniqueness, if
\(\pi^*\mathcal M\otimes\mathcal O(k)\cong\mathcal O\), restriction to
any fiber gives \(k=0\).  Pushing forward then gives
\(\mathcal M\cong\mathcal O_X\).  Thus the homomorphism is an
isomorphism.

(b) Let
\(f:\mathbb P(\mathcal E)\xrightarrow{\sim}\mathbb P(\mathcal E')\)
be an \(X\)-isomorphism.  Part (a) gives
\[
f^*\mathcal O_{\mathbb P(\mathcal E')}(1)
\cong\pi^*\mathcal M\otimes\mathcal O_{\mathbb P(\mathcal E)}(k).
\]
On every geometric fiber, \(f\) is an automorphism of projective
space and hence pulls \(\mathcal O(1)\) back to \(\mathcal O(1)\).
Thus \(k=1\).  Since
\[
\pi_*\mathcal O_{\mathbb P(\mathcal E)}(1)\cong\mathcal E,
\qquad
\pi'_*\mathcal O_{\mathbb P(\mathcal E')}(1)\cong\mathcal E',
\]
pushforward and the projection formula yield
\[
\mathcal E'
\cong\pi_*f^*\mathcal O_{\mathbb P(\mathcal E')}(1)
\cong\mathcal M\otimes\mathcal E.
\]

Conversely, if \(\mathcal E'\cong\mathcal E\otimes\mathcal M\), then
tensoring the degree-\(d\) part of
\(\operatorname{Sym}\mathcal E\) by \(\mathcal M^d\) does not change
relative Proj.  Hence
\[
\mathbb P(\mathcal E\otimes\mathcal M)
\cong_X\mathbb P(\mathcal E),
\]
which proves the converse.
\end{solution}

\begin{exerciseintro}
% record: H-II-07-010
\exerciseheading{Exercise II.7.10}
\addcontentsline{toc}{subsection}{Exercise II.7.10}

\begin{context}
Let \(X\) be noetherian.  A transition automorphism of
\(\mathbb P^n_U\) is called linear if, locally on \(U\), it is induced
by an invertible linear change of the degree-one homogeneous
coordinates.  Scalar matrices induce the same projective
automorphism.
\end{context}

\begin{problem}
(a) Define a projective \(n\)-space bundle over \(X\) by Zariski-local
trivializations with linear transition automorphisms.

(b) Prove that \(\mathbb P(\mathcal E)\) is such a bundle whenever
\(\mathcal E\) is locally free of rank \(n+1\).

(c) If \(X\) is regular, prove that every projective \(n\)-space bundle
\(P\to X\) is isomorphic to \(\mathbb P(\mathcal E)\) for some locally
free \(\mathcal E\).  Give a weaker sufficient hypothesis.

(d) For regular \(X\), deduce a bijection between isomorphism classes
of projective \(n\)-space bundles and rank-\(n+1\) locally free
sheaves modulo
\[
\mathcal E\sim\mathcal E\otimes\mathcal M,
\qquad \mathcal M\in\operatorname{Pic}(X).
\]
\end{problem}
\end{exerciseintro}

\begin{solution}
If \(n=0\), a \(\mathbb P^0\)-bundle is the base \(X\) itself.
Every rank-one locally free sheaf has this projectivization,
and any two differ by tensoring with an invertible sheaf.
This proves the assertion in that case.  Assume henceforth \(n\geq1\).

(a) A projective \(n\)-space bundle is a morphism
\(\pi:P\to X\) with an open cover \(X=\bigcup U_i\) and
\(U_i\)-isomorphisms
\[
\pi^{-1}(U_i)\cong U_i\times\mathbb P^n
\]
such that every transition map over \(U_i\cap U_j\) is locally induced
by an element of
\(\operatorname{GL}_{n+1}(\mathcal O_X)\).  Composition and inversion
of coordinate changes preserve this condition, so it is independent
of refinements and defines a gluing category.

(b) Choose frames of \(\mathcal E\) on an open cover.  They identify
\(\mathbb P(\mathcal E)|_{U_i}\) with \(\mathbb P^n_{U_i}\).
Changes of frame are invertible \((n+1)\)-by-\((n+1)\) matrices, and
their projectivizations are exactly the required linear transition
maps.

(c) Work on one connected component; a connected regular scheme is
integral.  Choose a dense open \(U\subseteq X\) over which \(P\) is
trivial, and let \(\mathcal O(1)\) be the usual invertible sheaf on
\(P_U=\pi^{-1}(U)\).  The total space \(P\) is regular, because it is
locally a projective space over \(X\), and hence it is locally
factorial.  Represent \(\mathcal O(1)\) by a Cartier divisor on
\(P_U\), take the closures in \(P\) of its prime components, and use
local factoriality to obtain a Cartier divisor on \(P\).  Its
invertible sheaf is an extension \(\mathcal L\) of
\(\mathcal O_{P_U}(1)\).

Over a trivializing open \(V\subseteq X\), the projective-space Picard
calculation gives
\[
\mathcal L|_{P_V}
\cong\pi^*\mathcal M_V\otimes\mathcal O_{\mathbb P^n_V}(a_V).
\]
On the dense overlap with \(U\), its fiber degree is \(1\), so
\(a_V=1\) on every component meeting \(U\).  Thus
\[
\mathcal E=\pi_*\mathcal L
\]
is locally free of rank \(n+1\).  The adjunction evaluation map
\[
\pi^*\mathcal E\twoheadrightarrow\mathcal L
\]
is surjective: on \(P_V\) it is, after the harmless factor
\(\mathcal M_V\), the standard evaluation
\(\mathcal O_V^{n+1}\twoheadrightarrow\mathcal O_{\mathbb P^n_V}(1)\).
The quotient universal property produces
\[
g:P\longrightarrow\mathbb P(\mathcal E).
\]
On every \(V\), this is the standard isomorphism
\(\mathbb P^n_V\to\mathbb P(\mathcal O_V^{n+1}\otimes\mathcal M_V)\).
Hence \(g\) is globally an isomorphism.

The proof used local factoriality, rather than the full force of
regularity: it extended the divisor and used the projective-space
Picard calculation.  Thus a noetherian locally factorial base is a
weaker sufficient hypothesis.

(d) Part (b) constructs a bundle from every locally free sheaf, and
part (c) proves surjectivity on isomorphism classes.  Exercise II.7.9
shows that
\[
\mathbb P(\mathcal E)\cong_X\mathbb P(\mathcal E')
\quad\Longleftrightarrow\quad
\mathcal E'\cong\mathcal E\otimes\mathcal M
\]
for an invertible \(\mathcal M\).  This is precisely the stated
bijection.
\end{solution}

\begin{exerciseintro}
% record: H-II-07-011
\exerciseheading{Exercise II.7.11}
\addcontentsline{toc}{subsection}{Exercise II.7.11}

\begin{context}
For a coherent ideal sheaf \(\mathcal I\) on a noetherian scheme
\(X\), write
\[
\operatorname{Bl}_{\mathcal I}X
=\mathbf{Proj}_X\left(\bigoplus_{n\geq0}\mathcal I^n\right).
\]
In part (c), start with a morphism \(f:Z\to X\) represented as such a
blowup by the theorem being strengthened.
\end{context}

\begin{problem}
(a) Prove that blowing up \(\mathcal I^d\), for any \(d\geq1\), gives
the same \(X\)-scheme as blowing up \(\mathcal I\).

(b) If \(\mathcal J\) is an invertible ideal sheaf, prove
\[
\operatorname{Bl}_{\mathcal I\mathcal J}X
\cong_X\operatorname{Bl}_{\mathcal I}X.
\]

(c) Suppose \(X\) is regular, and let \(U\subseteq X\) be the largest
open set over which \(f\) is an isomorphism.  Show that the blowup
ideal may be chosen so that its zero scheme has support exactly
\(X\setminus U\).
\end{problem}
\end{exerciseintro}

\begin{solution}
(a) The Rees algebra of \(\mathcal I^d\) is
\[
\bigoplus_{n\geq0}\mathcal I^{dn},
\]
the \(d\)-th Veronese subalgebra of
\(\bigoplus_{n\geq0}\mathcal I^n\).  Relative Proj is unchanged on
passing to a positive Veronese subalgebra.  Hence the two blowups are
canonically isomorphic over \(X\).

(b) Since \(\mathcal J\) is invertible,
\[
(\mathcal I\mathcal J)^n
\cong\mathcal I^n\otimes\mathcal J^n.
\]
Degreewise tensoring a graded algebra by the powers of an invertible
sheaf does not change relative Proj.  Equivalently, on an open where
\(\mathcal J\) is generated by one nonzerodivisor \(a\), multiplication
by \(a^n\) identifies the degree-\(n\) pieces, and these local Proj
isomorphisms agree on overlaps.  This proves (b).

(c) Begin with \(Z\cong\operatorname{Bl}_{\mathcal I}X\).  The
tautological pullback of \(\mathcal I\) to the blowup is invertible.
Over \(U\), where the blowup is the identity, this says that
\(\mathcal I|_U\) is an invertible ideal.  Thus
\[
\mathcal I|_U=\mathcal O_U(-D)
\]
for an effective Cartier divisor \(D\) on \(U\).

Work componentwise, so \(X\) is integral.  Take the closures of the
prime components of \(D\) in \(X\).  Regularity makes the resulting
effective Weil divisor \(\overline D\) Cartier.  We claim
\[
\mathcal I\subseteq\mathcal O_X(-\overline D).
\tag{1}
\]
This is local.  In a regular local ring, which is a unique
factorization domain, each element of \(\mathcal I\) has along every
height-one component of \(\overline D\) at least the order prescribed
by \(D\), because equality holds at that component's generic point in
\(U\).  Unique factorization then gives divisibility by a local
equation of \(\overline D\), proving (1).

Set
\[
\mathcal K=\mathcal I\otimes\mathcal O_X(\overline D).
\]
By (1), \(\mathcal K\) is a coherent ideal sheaf, and
\(\mathcal K|_U=\mathcal O_U\).  Moreover
\[
\mathcal I=\mathcal K\mathcal O_X(-\overline D),
\]
where \(\mathcal O_X(-\overline D)\) is an invertible ideal.  Part (b)
therefore gives
\[
\operatorname{Bl}_{\mathcal K}X
\cong_X\operatorname{Bl}_{\mathcal I}X\cong_X Z.
\]
The support of \(\mathcal O_X/\mathcal K\) lies in \(X\setminus U\).
If it omitted a point \(x\in X\setminus U\), then
\(\mathcal K=\mathcal O_X\) on a neighborhood of \(x\), and its
blowup would be the identity there.  This would enlarge \(U\),
contrary to maximality.  Hence the support is exactly \(X\setminus U\).
\end{solution}

\begin{exerciseintro}
% record: H-II-07-012
\exerciseheading{Exercise II.7.12}
\addcontentsline{toc}{subsection}{Exercise II.7.12}

\begin{context}
Let \(X\) be noetherian, and let \(Y,Z\subseteq X\) be closed
subschemes with ideal sheaves \(\mathcal I\) and \(\mathcal J\).
Assume neither contains the other.  Their scheme-theoretic
intersection is defined by
\(\mathcal K=\mathcal I+\mathcal J\).
\end{context}

\begin{problem}
Let
\[
\widetilde X=\operatorname{Bl}_{Y\cap Z}X.
\]
Prove that the strict transforms \(\widetilde Y\) and
\(\widetilde Z\) are disjoint.
\end{problem}
\end{exerciseintro}

\begin{solution}
The question is local on \(X\).  On an affine open write the three
ideals as \(I,J,K=I+J\), and let
\[
R=\bigoplus_{d\geq0}K^d
\]
be the Rees algebra.  The strict transform of \(Y\) is cut out in
\(\operatorname{Proj}R\) by the homogeneous ideal whose degree-\(d\)
piece is
\[
K^d\cap I.
\tag{1}
\]
Indeed, its homogeneous coordinate algebra is the image of \(R\) in
the Rees algebra of the induced ideal on \(A/I\):
\[
\bigoplus_{d\geq0}
\frac{K^d+I}{I}
\cong
\bigoplus_{d\geq0}\frac{K^d}{K^d\cap I}.
\]
The analogous ideal for \(\widetilde Z\) has degree-\(d\) piece
\(K^d\cap J\).

For every \(d\geq1\),
\[
K^d=(I+J)^d
\subseteq (K^d\cap I)+(K^d\cap J).
\tag{2}
\]
To see this, expand a product of \(d\) factors from \(I+J\): a term
containing an \(I\)-factor belongs to \(I\cap K^d\), while the remaining
all-\(J\) term belongs to \(J\cap K^d\).  The reverse inclusion in
(2) is automatic, so equality holds.

Thus the sum of the two homogeneous ideals defining the strict
transforms contains every positive-degree piece of \(R\), namely the
irrelevant ideal.  No homogeneous prime representing a point of
\(\operatorname{Proj}R\) can contain that irrelevant ideal.
Consequently the scheme-theoretic intersection
\(\widetilde Y\cap\widetilde Z\) is empty.
\end{solution}

\begin{exerciseintro}
% record: H-II-07-013
\exerciseheading{Exercise II.7.13}
\addcontentsline{toc}{subsection}{Exercise II.7.13}

\begin{context}
Let \(k\) be algebraically closed with
\(\operatorname{char}k\neq2\), and let
\[
C:\ y^2z=x^3+x^2z
\]
be the nodal cubic.  Its smooth locus is \(\mathbb G_m\).  Multiplication
by \(a\in k^*\) on that group extends to an automorphism
\(\varphi_a:C\to C\).

Glue
\[
X_0=C\times(\mathbb P^1\setminus\{\infty\}),\qquad
X_\infty=C\times(\mathbb P^1\setminus\{0\})
\]
over \(C\times\mathbb G_m\) by
\[
(P,u)\longmapsto(\varphi_u(P),u).
\]
The second projections glue to
\(\pi:X\to\mathbb P^1\).
For the finite normalization, local inverse images are spectra
of semilocal rings with trivial Picard groups.  Hence
\(R^1\nu_*\mathcal O_{\widetilde Y}^*=0\), and the Leray sequence
justifies the first-cohomology comparison for the unit sheaves.
\end{context}

\begin{problem}
(a) Prove that \(\pi\) is proper and that \(X\) is a complete variety.

(b) Compute
\[
\operatorname{Pic}(C\times\mathbb A^1)\cong k^*\oplus\mathbb Z,
\qquad
\operatorname{Pic}(C\times\mathbb G_m)
\cong k^*\oplus\mathbb Z\oplus\mathbb Z.
\]

(c) Choose these coordinates so that restriction is
\[
(t,n)\longmapsto(t,0,n)
\]
and pullback by
\(\varphi(P,u)=(\varphi_u(P),u)\) is
\[
(t,d,n)\longmapsto(t,d+n,n).
\]

(d) Prove that every invertible sheaf on \(X\) restricts to degree
zero on \(C\times\{0\}\).  Deduce that \(X\) is proper but not
projective and that \(\pi\) is not projective.
\end{problem}
\end{exerciseintro}

\begin{solution}
(a) Properness is local on the base in the Zariski topology.  Over
each of the two standard affine opens of \(\mathbb P^1\), the morphism
\(\pi\) is the projection
\[
C\times\mathbb A^1\longrightarrow\mathbb A^1,
\]
which is proper because it is the base change of the proper curve
\(C\to\operatorname{Spec}k\).  Hence \(\pi\) is proper.  The composite
\[
X\xrightarrow{\pi}\mathbb P^1\longrightarrow\operatorname{Spec}k
\]
is proper, so \(X\) is complete.  The pieces are integral schemes of
finite type and are glued along a dense open, so \(X\) is an integral
variety.

(b) Choose a coordinate on the normalization
\(\nu:\mathbb P^1\to C\) for which \(0\) and \(\infty\) are the two
points above the node.  For
\(T=\mathbb A^1\) or \(\mathbb G_m\), put
\[
Y=C\times T,\qquad \widetilde Y=\mathbb P^1\times T.
\]
The normalization map gives an exact sequence of unit sheaves on
\(Y\):
\[
1\longrightarrow\mathcal O_Y^*
\longrightarrow\nu_*\mathcal O_{\widetilde Y}^*
\longrightarrow i_*\mathcal O_T^*
\longrightarrow1,
\tag{1}
\]
where \(i:T\to Y\) is the node section.  The last arrow takes the ratio
of the values on the two branches.  Its kernel consists exactly of
units descending to the node, and it is locally surjective.

The global units on \(Y\) and \(\widetilde Y\) are both the units on
\(T\); their branch ratio is \(1\).  Moreover
\[
H^1(T,\mathcal O_T^*)=\operatorname{Pic}(T)=0.
\]
The cohomology sequence of (1) therefore becomes
\[
0\longrightarrow\Gamma(T,\mathcal O_T^*)
\longrightarrow\operatorname{Pic}(C\times T)
\longrightarrow\operatorname{Pic}(\mathbb P^1\times T)
\longrightarrow0.
\tag{2}
\]
The last group is \(\mathbb Z\), measured by degree on the
\(\mathbb P^1\)-fiber.  A divisor \(q\times T\), for a smooth point
\(q\in C\), splits (2).  Since
\[
\Gamma(\mathbb A^1,\mathcal O^*)=k^*,\qquad
\Gamma(\mathbb G_m,\mathcal O^*)=k^*u^{\mathbb Z},
\]
the two asserted Picard-group formulas follow.

(c) Sequence (2) also gives concrete coordinates.  A line bundle on
\(C\times T\) is specified by its degree \(n\) on the normalized
\(\mathbb P^1\)-fiber and by a branch-gluing unit
\(\lambda(u)\in\Gamma(T,\mathcal O_T^*)\).  For
\(T=\mathbb A^1\), write \(\lambda=t\in k^*\), giving \((t,n)\).
For \(T=\mathbb G_m\), write uniquely
\[
\lambda(u)=tu^d,\qquad t\in k^*,\ d\in\mathbb Z,
\]
giving \((t,d,n)\).  Restriction from \(\mathbb A^1\) to
\(\mathbb G_m\) leaves a constant gluing unit constant, and hence is
\[
(t,n)\longmapsto(t,0,n).
\]

Multiplication by \(u\) on the normalized \(\mathbb P^1\) fixes
\(0,\infty\).  On a degree-\(n\) line bundle it changes the ratio of
the standard trivializations at those two points by \(u^n\).  Thus
pullback by the gluing automorphism replaces
\(\lambda(u)=tu^d\) by \(tu^{d+n}\), while preserving the degree.
With this consistent choice of branch order,
\[
\varphi^*(t,d,n)=(t,d+n,n).
\]

(d) Let an invertible sheaf \(\mathcal L\) on \(X\) have classes
\((t_0,n_0)\) and \((t_\infty,n_\infty)\) on \(X_0\) and
\(X_\infty\).  On the overlap, the gluing of \(X\) requires
\[
(t_0,0,n_0)=\varphi^*(t_\infty,0,n_\infty).
\]
Part (c) turns this into
\[
(t_0,0,n_0)=(t_\infty,n_\infty,n_\infty).
\]
Consequently
\[
n_0=n_\infty=0.
\tag{3}
\]
In particular, \(\mathcal L|_{C\times\{0\}}\) has degree zero.

If \(X\) were projective, it would possess a very ample invertible
sheaf.  Its restriction to the closed fiber \(C\times\{0\}\) would
still be very ample, but (3) makes its pullback to the normalization
\(\mathbb P^1\) a degree-zero line bundle, hence \(\mathcal O_{\mathbb
P^1}\).  All its pulled-back global sections are constant, so the
induced map on \(\mathbb P^1\) is constant, contradicting a closed
immersion of the nontrivial curve \(C\).  Thus \(X\) is not projective.
Finally, if \(\pi\) were projective, composing it with the projective
morphism \(\mathbb P^1\to\operatorname{Spec}k\) would make \(X\)
projective.  Hence \(\pi\) is not projective.
\end{solution}

\begin{exerciseintro}
% record: H-II-07-014
\exerciseheading{Exercise II.7.14}
\addcontentsline{toc}{subsection}{Exercise II.7.14}

\begin{context}
For part (b), condition \((\dagger)\) supplies a quasi-coherent graded
\(\mathcal O_X\)-algebra
\(\mathcal S=\bigoplus_{d\geq0}\mathcal S_d\), with
\(\mathcal S_0=\mathcal O_X\), generated in degree one by the coherent
sheaf \(\mathcal S_1\), and such that the associated
\(\mathcal O_P(1)\) on
\(P=\mathbf{Proj}_X\mathcal S\) is invertible.  Ampleness of
\(\mathcal L\) is understood relative to \(f:X\to Y\).
Throughout, very ample relative to a base means realized by an
immersion into a single finite-dimensional projective space over
that base.  In (b) the base is required to be affine and noetherian.
\end{context}

\begin{problem}
(a) Give a noetherian \(X\) and a finite locally free
\(\mathcal E\) for which
\(\mathcal O_{\mathbb P(\mathcal E)}(1)\) is not very ample relative
to \(X\).

(b) Let \(Y=\operatorname{Spec}A\) with \(A\) noetherian, let \(f:X\to Y\) be of finite type, let \(\mathcal L\) be
\(f\)-ample, and let \(\pi:P\to X\) be as in the context.  Prove that
\[
\mathcal O_P(1)\otimes\pi^*\mathcal L^n
\]
is very ample relative to \(Y\) for every sufficiently large \(n\).
\end{problem}
\end{exerciseintro}

\begin{solution}
(a) Take
\[
X=\mathbb P_k^1,\qquad
\mathcal E=\mathcal O_X(-1)^{\oplus2}.
\]
Twisting a vector bundle by a line bundle does not change its
projectivization, so
\[
\mathbb P(\mathcal E)
\cong\mathbb P(\mathcal O_X^{\oplus2})
\cong\mathbb P^1\times\mathbb P^1.
\]
Under this isomorphism, with the base as the first factor,
\[
\mathcal O_{\mathbb P(\mathcal E)}(1)
\cong\mathcal O_{\mathbb P^1\times\mathbb P^1}(-1,1).
\]
It has no nonzero global sections.  A relatively very ample sheaf is
the pullback of \(\mathcal O(1)\) under an immersion into some
\(\mathbb P_X^N\), and is therefore generated by its defining
sections.  Hence this \(\mathcal O(1)\) is not very ample relative to
\(X\).

(b) Choose \(a>0\) such that \(\mathcal L^a\) is very ample relative
to \(Y\).  Such an \(a\) exists by the ample-sheaf criterion.  By
ampleness again, for every sufficiently large \(q\), the coherent
sheaf
\[
\mathcal S_1\otimes\mathcal L^q
\]
is generated by finitely many sections.  A resulting surjection
\[
\mathcal O_X^{N+1}\twoheadrightarrow
\mathcal S_1\otimes\mathcal L^q
\]
extends, because \(\mathcal S\) is generated in degree one, to a
surjection onto the degreewise twisted graded algebra
\[
\bigoplus_{d\geq0}\mathcal S_d\otimes\mathcal L^{qd}.
\]
Relative Proj of this algebra is still \(P\), and its twisting sheaf
is
\[
\mathcal O_P(1)\otimes\pi^*\mathcal L^q.
\]
The graded surjection therefore gives a closed immersion
\[
P\hookrightarrow\mathbb P_X^N
\]
realizing this sheaf.  It is very ample relative to \(X\).

Now let \(n\geq a\) and put \(q=n-a\).  Once \(n\) is large, so is
\(q\), and the preceding paragraph applies.  The composition rule for
very ample sheaves, applied to
\(P\xrightarrow{\pi}X\xrightarrow{f}Y\),
says that
\[
\bigl(\mathcal O_P(1)\otimes\pi^*\mathcal L^q\bigr)
\otimes\pi^*(\mathcal L^a)
=\mathcal O_P(1)\otimes\pi^*\mathcal L^n
\]
is very ample relative to \(Y\).  This holds for every sufficiently
large \(n\).
The affine-base hypothesis cannot be dropped with this definition.
For \(f=\operatorname{id}_{\mathbb P^1}\),
\(\mathcal L=\mathcal O_{\mathbb P^1}\), and
\(\mathcal S=\operatorname{Sym}(\mathcal O(-1)^{\oplus2})\),
the sheaf in (b) is \(\mathcal O(-1,1)\) for every \(n\).
The identity makes \(\mathcal L\) relatively ample, but (a) proves
that this sheaf is never relatively very ample.
\end{solution}

\begin{exerciseintro}
\corpussection{II.8 Differentials}
\addcontentsline{toc}{section}{II.8 Differentials}

% record: H-II-08-001
\exerciseheading{Exercise II.8.1}
\addcontentsline{toc}{subsection}{Exercise II.8.1}

\begin{context}
Write \(\kappa(B)=B/\mathfrak m\) for the residue field of a local
ring \((B,\mathfrak m)\).  We may use the conormal sequence for
\(k\to B\to\kappa(B)\), the equality
\[
\dim_L\Omega_{L/k}=\operatorname{trdeg}_kL
\]
when the finitely generated field extension \(L/k\) is separably
generated, and the dimension formula for an integral algebra of
finite type over a field.  We also use the relative coefficient-field
theorem: if a complete equicharacteristic local ring contains \(k\)
and its residue field is separably generated over \(k\), it has a
coefficient field containing \(k\).
\end{context}

\begin{problem}
(a) Let \((B,\mathfrak m)\) be a local ring containing a field \(k\),
and suppose \(\kappa=\kappa(B)\) is separably generated over \(k\).
Prove that the conormal sequence is short exact:
\[
0\longrightarrow\mathfrak m/\mathfrak m^2
\longrightarrow\Omega_{B/k}\otimes_B\kappa
\longrightarrow\Omega_{\kappa/k}\longrightarrow0.
\]

(b) Suppose in addition that \(k\) is perfect and that \(B\) is a
localization of a finitely generated \(k\)-algebra.  Prove that \(B\)
is regular if and only if \(\Omega_{B/k}\) is free of rank
\[
\dim B+\operatorname{trdeg}_k\kappa.
\]

(c) Let \(X\) be an irreducible finite-type scheme of dimension
\(n\) over a perfect field.  For every point \(x\in X\), prove that
\(\mathcal O_{X,x}\) is regular exactly when
\((\Omega_{X/k})_x\) is free of rank \(n\).

(d) Deduce that the regular points of a variety over an algebraically
closed field form an open dense subset.
\end{problem}
\end{exerciseintro}

\begin{solution}
(a) Put \(\overline B=B/\mathfrak m^2\) and
\(\overline{\mathfrak m}=\mathfrak m/\mathfrak m^2\).  This is a
complete local ring whose maximal ideal is square-zero.  The
coefficient-field theorem supplies a subfield
\(K\subseteq\overline B\), containing the image of \(k\), such that
\(K\to\kappa\) is an isomorphism.  Consequently every element of
\(\overline B\) has a unique expression
\[
\lambda+u,\qquad \lambda\in K,\quad
u\in\overline{\mathfrak m}.
\tag{1}
\]

Let \(h:\mathfrak m/\mathfrak m^2\to\kappa\) be any
\(\kappa\)-linear map.  Define
\[
D:B\longrightarrow\kappa,
\qquad D(b)=h(u),
\]
where the image of \(b\) in \(\overline B\) is written as in (1).
This map vanishes on \(k\).  If
\(b=\lambda+u\) and \(c=\mu+v\) modulo \(\mathfrak m^2\), then
\(uv=0\), so
\[
D(bc)=h(\lambda v+\mu u)
=\overline bD(c)+\overline cD(b).
\]
Thus \(D\) is a \(k\)-derivation, and its restriction to
\(\mathfrak m/\mathfrak m^2\) is \(h\).

Dualizing the first map in the conormal sequence identifies its dual
with the restriction map
\[
\operatorname{Der}_k(B,\kappa)
\longrightarrow
\operatorname{Hom}_{\kappa}(\mathfrak m/\mathfrak m^2,\kappa).
\]
The construction shows that this map is surjective.  Linear
functionals separate vectors, hence the conormal map is injective.
The other two terms were already right exact, proving (a).

(b) Set
\[
r=\dim B+\operatorname{trdeg}_k\kappa.
\]
If \(\Omega_{B/k}\) is free of rank \(r\), part (a) and the field
differential formula give
\[
\dim_{\kappa}\mathfrak m/\mathfrak m^2
=r-\dim_{\kappa}\Omega_{\kappa/k}
=\dim B.
\]
This is precisely the regular-local-ring criterion.

Conversely, suppose \(B\) is regular.  Then \(B\) is a domain; write
\(K=\operatorname{Frac}(B)\).  Part (a) gives
\[
\dim_{\kappa}(\Omega_{B/k}\otimes_B\kappa)=r.
\tag{2}
\]
Because \(k\) is perfect, \(K/k\) is separably generated.  Localization
of differentials and the dimension formula give
\[
\dim_K(\Omega_{B/k}\otimes_BK)
=\operatorname{trdeg}_kK
=\dim B+\operatorname{trdeg}_k\kappa=r.
\tag{3}
\]
The module \(\Omega_{B/k}\) is finite.  Lift a basis in (2) and use
Nakayama's lemma to obtain a surjection
\(B^r\twoheadrightarrow\Omega_{B/k}\).  After tensoring with \(K\),
(3) makes this an isomorphism.  Its kernel therefore becomes zero
over \(K\).  That kernel lies in the torsion-free module \(B^r\), so
it is zero.  Hence \(\Omega_{B/k}\cong B^r\).

(c) Put \(B=\mathcal O_{X,x}\) and \(\kappa=\kappa(x)\).  Since
\(X_{\mathrm{red}}\) is integral of dimension \(n\) and has
the same local dimensions and residue fields as \(X\), the
dimension formula applied to \(X_{\mathrm{red}}\) says
\[
\dim B+\operatorname{trdeg}_k\kappa=n.
\]
Moreover, differentials commute with localization:
\[
(\Omega_{X/k})_x\cong\Omega_{B/k}.
\]
Part (b) now gives the asserted equivalence.

(d) By (c), the regular locus is the locus on which the coherent
sheaf \(\Omega_{X/k}\) is locally free of rank \(n=\dim X\).  This
locus is open.  At the generic point \(\eta\),
\[
(\Omega_{X/k})_\eta=\Omega_{k(X)/k}
\]
is an \(n\)-dimensional \(k(X)\)-vector space, because an algebraically
closed field is perfect.  Thus the open locus contains \(\eta\), and
is therefore dense.
\end{solution}

\begin{exerciseintro}
% record: H-II-08-002
\exerciseheading{Exercise II.8.2}
\addcontentsline{toc}{subsection}{Exercise II.8.2}

\begin{context}
A variety is integral and of finite type over an algebraically closed
field \(k\).  A family of global sections generates a locally free
sheaf if its evaluation map is surjective.  We may use that the image
of a finite-type morphism is constructible and that the dimension of
the closure of an image is at most the dimension of its source.
\end{context}

\begin{problem}
Let \(X\) be an \(n\)-dimensional variety, let \(\mathcal E\) be
locally free of rank \(r>n\), and let
\(V\subseteq\Gamma(X,\mathcal E)\) generate \(\mathcal E\).  Prove
that some \(s\in V\) is nonzero in every fiber
\(\mathcal E_x/\mathfrak m_x\mathcal E_x\).  Deduce an exact sequence
\[
0\longrightarrow\mathcal O_X\longrightarrow\mathcal E
\longrightarrow\mathcal E'\longrightarrow0
\]
in which \(\mathcal E'\) is locally free.
\end{problem}
\end{exerciseintro}

\begin{solution}
First replace \(V\) by a finite-dimensional subspace that still
generates \(\mathcal E\).  Indeed, at each point finitely many members
of \(V\) span the fiber.  The same remains true on an open
neighborhood, and quasi-compactness of \(X\) gives finitely many such
neighborhoods and sections.  Let \(m=\dim_kV\).

The evaluation map is a surjection of vector bundles
\[
V\otimes_k\mathcal O_X\twoheadrightarrow\mathcal E.
\]
Its kernel \(\mathcal K\) is locally free of rank \(m-r\), since the
surjection splits locally.  Inside \(X\times\mathbb A(V)\), consider
the incidence scheme
\[
Z=\{(x,s):s(x)=0\}.
\]
It is the total space of \(\mathcal K\), and hence
\[
\dim Z=n+m-r<m=\dim\mathbb A(V).
\]
For the projection \(p:Z\to\mathbb A(V)\), the closure of \(p(Z)\)
therefore has dimension strictly less than \(m\).  Choose a
\(k\)-rational point \(s\) outside that proper closed subset.  Such a
point exists because \(k\) is algebraically closed.  No point of
\(Z\) lies over \(s\), which says exactly that
\[
s_x\notin\mathfrak m_x\mathcal E_x
\quad\text{for every }x\in X.
\]

The section defines \(\mathcal O_X\to\mathcal E\).  At a point
\(x\), trivialize \(\mathcal E_x\cong\mathcal O_{X,x}^r\).  At least
one coordinate of \(s_x\) is a unit, so elementary changes of basis
send \(s_x\) to the first standard basis vector.  Consequently the
stalk map is a split injection and its cokernel is free of rank
\(r-1\).  These stalkwise cokernels form a locally free sheaf
\(\mathcal E'\), giving the displayed exact sequence.
\end{solution}

\begin{exerciseintro}
% record: H-II-08-003
\exerciseheading{Exercise II.8.3}
\addcontentsline{toc}{subsection}{Exercise II.8.3}

\begin{context}
For a nonsingular \(d\)-dimensional variety \(T\), its canonical
sheaf is \(\omega_T=\bigwedge^d\Omega_{T/k}\), and
\(p_g(T)=\dim_k\Gamma(T,\omega_T)\).  A smooth plane cubic has
\(\omega\cong\mathcal O\) and arithmetic genus one.  For projective
curves \(C,D\), the arithmetic-genus product formula specializes to
\[
p_a(C\times D)=p_a(C)p_a(D)-p_a(C)-p_a(D).
\]
\end{context}

\begin{problem}
(a) For \(S\)-schemes \(X,Y\), prove
\[
\Omega_{X\times_SY/S}
\cong p_1^*\Omega_{X/S}\oplus p_2^*\Omega_{Y/S}.
\]

(b) If \(X,Y\) are nonsingular varieties over \(k\), deduce
\[
\omega_{X\times_kY}
\cong p_1^*\omega_X\otimes p_2^*\omega_Y.
\]

(c) If \(Y\) is a nonsingular plane cubic and \(T=Y\times Y\), show
that \(p_g(T)=1\) whereas \(p_a(T)=-1\).
\end{problem}
\end{exerciseintro}

\begin{solution}
(a) The assertion is local on all three schemes.  For ring maps
\(A\to B\) and \(A\to C\), and for any \(B\otimes_AC\)-module
\(M\), restriction to the two factors gives a natural bijection
\[
\operatorname{Der}_A(B\otimes_AC,M)
\longrightarrow
\operatorname{Der}_A(B,M)\oplus\operatorname{Der}_A(C,M).
\]
Its inverse sends \((D_B,D_C)\) to
\[
D(b\otimes c)=cD_B(b)+bD_C(c).
\]
The formula is balanced over \(A\), satisfies the Leibniz rule, and
is forced on pure tensors.  The universal property of differentials
therefore yields
\[
\Omega_{(B\otimes_AC)/A}
\cong(\Omega_{B/A}\otimes_B(B\otimes_AC))
\oplus
(\Omega_{C/A}\otimes_C(B\otimes_AC)).
\]
Sheafifying gives the required formula.

(b) Put \(a=\dim X\) and \(b=\dim Y\).  Part (a) gives a locally
free sheaf of rank \(a+b\):
\[
\Omega_{X\times Y/k}
\cong p_1^*\Omega_{X/k}\oplus p_2^*\Omega_{Y/k}.
\]
The differential criterion shows that the product is nonsingular.
Taking its top exterior power and using
\(\det(E\oplus F)=\det E\otimes\det F\), we obtain
\[
\omega_{X\times Y}
\cong p_1^*\omega_X\otimes p_2^*\omega_Y.
\]

(c) Since \(\omega_Y\cong\mathcal O_Y\), part (b) gives
\(\omega_T\cong\mathcal O_T\).  The projective integral variety
\(T\) has only constant global regular functions, so
\[
p_g(T)=\dim_k\Gamma(T,\mathcal O_T)=1.
\]
On the other hand, \(p_a(Y)=1\), and the stated product formula gives
\[
p_a(T)=1\cdot1-1-1=-1.
\]
\end{solution}

\begin{exerciseintro}
% record: H-II-08-004
\exerciseheading{Exercise II.8.4}
\addcontentsline{toc}{subsection}{Exercise II.8.4}

\begin{context}
Let \(S=k[x_0,\ldots,x_n]\).  A closed subscheme
\(Y\subseteq\mathbb P_k^n\), with saturated homogeneous ideal
\(I_Y\), is a strict complete intersection if \(I_Y\) is generated
by \(r=\operatorname{codim}(Y,\mathbb P^n)\) elements.  We may use
the following facts: a regular sequence quotient of \(S\) is
Cohen--Macaulay and unmixed; Serre's conditions \(R_1\) and \(S_2\)
characterize normality; the homogeneous coordinate ring is normal
exactly when the embedding is projectively normal; and Bertini's
theorem applies to the complete linear systems of degree-\(d\)
hypersurfaces.  We may also use the arithmetic-genus formulas obtained
from the Hilbert polynomial for hypersurfaces and complete-intersection
curves.  Binomial coefficients \(\binom ab\) below are
understood to be zero when \(a<b\).
The ground field \(k\) is algebraically closed.
\end{context}

\begin{problem}
(a) If \(Y\subseteq\mathbb P^n\) has codimension \(r\), prove that it
is a strict complete intersection if and only if it is the
scheme-theoretic intersection of \(r\) hypersurfaces.

(b) Prove that a normal complete intersection of positive dimension
is projectively normal.

(c) Under the hypotheses of (b), prove that
\[
H^0(\mathbb P^n,\mathcal O_{\mathbb P^n}(\ell))
\longrightarrow H^0(Y,\mathcal O_Y(\ell))
\]
is surjective for every \(\ell\geq0\), and deduce that \(Y\) is
connected.

(d) Given \(d_1,\ldots,d_r\geq1\) with \(r<n\), construct smooth
hypersurfaces \(H_i\) of degrees \(d_i\) whose intersection is
smooth, irreducible, and of codimension \(r\).

(e) For the smooth complete intersection \(Y\) in (d), prove
\[
\omega_Y\cong
\mathcal O_Y(d_1+\cdots+d_r-n-1).
\]

(f) If \(n\geq2\) and \(Y\subseteq\mathbb P^n\) is a smooth
hypersurface of degree \(d\), prove
\[
p_g(Y)=\binom{d-1}{n}.
\]
Show that this equals the arithmetic genus.  In particular, recover
\((d-1)(d-2)/2\) for a smooth plane curve.

(g) If a smooth curve \(Y\subseteq\mathbb P^3\) is the intersection
of smooth surfaces of degrees \(d,e\), prove
\[
p_g(Y)=\frac12de(d+e-4)+1,
\]
and show again that the geometric and arithmetic genera agree.
\end{problem}
\end{exerciseintro}

\begin{solution}
(a) Suppose first that \(I_Y\) can be generated by \(r\) elements.
The graded Nakayama lemma says that a basis of the graded vector space
\(I_Y/S_+I_Y\) lifts to a minimal homogeneous generating set.  Its
cardinality is at most \(r\), while the height theorem says that at
least \(r=\operatorname{ht}I_Y\) generators are necessary.  Thus
there are homogeneous generators \(f_1,\ldots,f_r\).  The
hypersurfaces \(H_i=V_+(f_i)\) then satisfy
\[
\mathcal I_Y=\mathcal I_{H_1}+\cdots+\mathcal I_{H_r},
\]
so their scheme-theoretic intersection is \(Y\).

Conversely, suppose \(Y=H_1\cap\cdots\cap H_r\) as schemes.  A
hypersurface in \(\mathbb P^n\) has invertible ideal sheaf
\(\mathcal O_{\mathbb P^n}(-d_i)\), and its inclusion into
\(\mathcal O_{\mathbb P^n}\) is given by a homogeneous polynomial
\(f_i\) of degree \(d_i\).  Put \(J=(f_1,\ldots,f_r)\).  The
codimension assumption gives \(\operatorname{ht}J=r\); hence the
\(f_i\) form a regular sequence in the Cohen--Macaulay ring \(S\).
It follows that \(S/J\) is unmixed and every associated prime has
height \(r\).  The irrelevant ideal \(S_+\), of height \(n+1>r\), is
therefore not associated to \(S/J\).  Thus \(S/J\) has no
\(S_+\)-torsion, which is equivalent to \(J\) being saturated.  The
ideal sheaf equality in the hypothesis says
\(I_Y=J^{\mathrm{sat}}\), so in fact \(I_Y=J\).  Therefore \(Y\) is
a strict complete intersection.

(b) Let \(A=S/I_Y\), the homogeneous coordinate ring.  Part (a)
makes \(A\) a complete-intersection ring, so it satisfies \(S_2\).
The punctured affine cone
\[
\operatorname{Spec}A\setminus\{A_+\}
\longrightarrow Y
\]
is locally a product with \(\mathbb G_m\).  Since \(Y\) is normal,
the punctured cone is normal and in particular regular in codimension
one.  The vertex has codimension
\(\dim A=\dim Y+1\geq2\), so it introduces no codimension-one point.
Hence \(A\) satisfies \(R_1\) as well as \(S_2\), and is normal by
Serre's criterion.  Every irreducible component of this affine
cone contains its vertex.  Distinct components of a normal
noetherian scheme are disjoint, so there is only one.
Thus \(A\) is a domain, proving projective normality of the embedding.

(c) The projective-normality criterion identifies \(A\) with the
section ring of \(\mathcal O_Y(1)\).  Degree by degree, this says that
\[
S_\ell=H^0(\mathbb P^n,\mathcal O(\ell))
\longrightarrow H^0(Y,\mathcal O_Y(\ell))
\]
is surjective for every \(\ell\geq0\).  At \(\ell=0\) it gives
\(H^0(Y,\mathcal O_Y)=k\).  A disconnection of \(Y\) would produce a
nonconstant idempotent global function, so \(Y\) is connected.

(d) Begin with \(Y_0=\mathbb P^n\).  Suppose smooth irreducible
\(Y_{i-1}\) of dimension \(n-i+1\) has been constructed.  Because
\(i\leq r<n\), this dimension is at least two.  Apply Bertini to the
degree-\(d_i\) complete linear system, equivalently to hyperplanes
after the \(d_i\)-uple embedding.  There is a nonempty open set of
degree-\(d_i\) forms whose zero hypersurface is smooth in
\(\mathbb P^n\), and another nonempty open set whose restriction to
\(Y_{i-1}\) has smooth irreducible zero locus of codimension one.
Choose a form in their intersection and call its hypersurface
\(H_i\).  Then
\(Y_i=Y_{i-1}\cap H_i\) has all the required properties.  Induction
through \(i=r\) constructs \(Y=Y_r\).

(e) Let \(W\) be smooth and let \(Z=W\cap H\) be a smooth section by
a hypersurface of degree \(d\).  The conormal sequence is
\[
0\longrightarrow\mathcal O_Z(-d)
\longrightarrow\Omega_{W/k}|_Z
\longrightarrow\Omega_{Z/k}\longrightarrow0.
\]
Taking determinants gives the adjunction formula
\[
\omega_Z\cong(\omega_W\otimes\mathcal O_W(d))|_Z.
\tag{1}
\]
Starting with \(\omega_{\mathbb P^n}=\mathcal O_{\mathbb P^n}(-n-1)\)
and applying (1) successively to the hypersurfaces in (d) yields
\[
\omega_Y\cong\mathcal O_Y(d_1+\cdots+d_r-n-1).
\]

(f) Set \(q=d-n-1\).  By (e), \(\omega_Y\cong\mathcal O_Y(q)\).
If \(q<0\), an anti-ample line bundle on the positive-dimensional
projective integral variety \(Y\) has no nonzero section, so
\(p_g(Y)=0=\binom{d-1}{n}\).  If \(q\geq0\), part (c) makes the
restriction map in degree \(q\) surjective.  Its kernel is
\((I_Y)_q\), which is zero because \(I_Y\) is generated in degree
\(d>q\).  Therefore
\[
p_g(Y)=\dim_kS_q=\binom{q+n}{n}=\binom{d-1}{n}.
\]
The complete-intersection Hilbert-polynomial calculation from the
arithmetic-genus formula gives the same binomial coefficient.  Hence
\(p_g(Y)=p_a(Y)\).  For \(n=2\), their common value is
\(\binom{d-1}{2}=(d-1)(d-2)/2\).

(g) Adjunction gives
\[
\omega_Y\cong\mathcal O_Y(q),
\qquad q=d+e-4.
\]
If \(q<0\), then \((d,e)\), up to order, is \((1,1)\) or
\((1,2)\); both the asserted formula and \(p_g(Y)\) are zero.  Now
assume \(q\geq0\).  If \(F,G\) are the defining forms, their Koszul
resolution gives
\[
0\longrightarrow S(-d-e)
\longrightarrow S(-d)\oplus S(-e)
\longrightarrow I_Y\longrightarrow0.
\]
In degree \(q\), the left term is \(S_{-4}=0\), and hence
\[
\dim_k(I_Y)_q
=\binom{e-1}{3}+\binom{d-1}{3}.
\]
Using the surjectivity in (c), we obtain
\[
\begin{aligned}
p_g(Y)
&=\binom{d+e-1}{3}
-\binom{e-1}{3}-\binom{d-1}{3}\\
&=\frac12de(d+e-4)+1.
\end{aligned}
\]
This is also the arithmetic genus furnished by the Hilbert-polynomial
formula for a complete intersection of degrees \(d,e\), so the two
genera agree.
\end{solution}

\begin{exerciseintro}
% record: H-II-08-005
\exerciseheading{Exercise II.8.5}
\addcontentsline{toc}{subsection}{Exercise II.8.5}

\begin{context}
Let \(X\) be a nonsingular variety and let
\(Y\subseteq X\) be a nonsingular subvariety of codimension
\(r\geq2\).  Write
\[
\pi:\widetilde X=\operatorname{Bl}_YX\longrightarrow X,
\qquad E=\pi^{-1}(Y).
\]
The blowup is nonsingular, \(E\) is the projective bundle associated
to the conormal bundle of \(Y\), and
\(\mathcal O_{\widetilde X}(-E)|_{\mathbb P^{r-1}}\cong
\mathcal O_{\mathbb P^{r-1}}(1)\) on every fiber.  On a regular
scheme, the Picard and divisor class groups agree.
\end{context}

\begin{problem}
(a) Prove that pullback and the exceptional divisor induce an
isomorphism
\[
\operatorname{Pic}(X)\oplus\mathbb Z
\longrightarrow\operatorname{Pic}(\widetilde X),
\qquad(\mathcal L,m)\longmapsto
\pi^*\mathcal L\otimes\mathcal O_{\widetilde X}(mE).
\]

(b) Prove the canonical-sheaf formula
\[
\omega_{\widetilde X}
\cong\pi^*\omega_X\otimes
\mathcal O_{\widetilde X}((r-1)E).
\]
\end{problem}
\end{exerciseintro}

\begin{solution}
(a) Put \(U=X\setminus Y\).  The blowup identifies
\(\widetilde X\setminus E\) with \(U\).  Because \(Y\) has
codimension at least two, deleting it does not change the divisor
class group of the regular scheme \(X\), and hence
\[
\operatorname{Pic}(X)\cong\operatorname{Cl}(X)
\xrightarrow{\sim}\operatorname{Cl}(U).
\]
The localization sequence for divisor class groups on
\(\widetilde X\), together with the irreducibility of the projective
bundle \(E\to Y\), gives an exact sequence
\[
\mathbb Z[E]\longrightarrow
\operatorname{Pic}(\widetilde X)
\longrightarrow\operatorname{Pic}(U)
\longrightarrow0.
\tag{1}
\]
Under the identification \(\operatorname{Pic}(U)\cong
\operatorname{Pic}(X)\), the composite of \(\pi^*\) with the second
map in (1) is the identity.  Thus (1) is split and only injectivity of
its first map remains.

Restrict \(\mathcal O_{\widetilde X}(mE)\) to a fiber
\(F\cong\mathbb P^{r-1}\) of \(E\to Y\).  The blowup description
gives
\[
\mathcal O_{\widetilde X}(mE)|_F
\cong\mathcal O_F(-m).
\]
This line bundle is trivial only for \(m=0\).  Hence \([E]\) has
infinite order, the first map in (1) is injective, and the asserted
direct-sum decomposition follows.

(b) The formula may be checked etale-locally at a general point of
\(Y\).  Choose smooth coordinates \(x_1,\ldots,x_n\) in which
\(Y\) is given by \(x_1=\cdots=x_r=0\).  On the blowup chart in which
the first homogeneous coordinate is nonzero, use coordinates
\[
x_1,\quad u_2=\frac{x_2}{x_1},\ldots,
u_r=\frac{x_r}{x_1},\quad x_{r+1},\ldots,x_n.
\]
Here \(E\) is defined by \(x_1=0\), and direct differentiation gives,
up to sign,
\[
\pi^*(dx_1\wedge\cdots\wedge dx_n)
=x_1^{r-1}
dx_1\wedge du_2\wedge\cdots\wedge du_r
\wedge dx_{r+1}\wedge\cdots\wedge dx_n.
\tag{2}
\]
Thus the natural map \(\pi^*\omega_X\to\omega_{\widetilde X}\)
vanishes to order \(r-1\) along \(E\).  Away from \(E\), \(\pi\) is
an isomorphism, so (2) has no other zero divisor.  Therefore its
divisor is exactly \((r-1)E\), which is equivalent to
\[
\omega_{\widetilde X}
\cong\pi^*\omega_X\otimes
\mathcal O_{\widetilde X}((r-1)E).
\]
\end{solution}

\begin{exerciseintro}
% record: H-II-08-006
\exerciseheading{Exercise II.8.6}
\addcontentsline{toc}{subsection}{Exercise II.8.6}

\begin{context}
Let \(k\) be algebraically closed, let \(A\) be a finitely generated
\(k\)-algebra with \(\operatorname{Spec}A\) nonsingular, and let
\[
0\longrightarrow I\longrightarrow B'\longrightarrow B
\longrightarrow0
\]
be an exact sequence of \(k\)-algebras with \(I^2=0\).  The action
of \(B'\) on \(I\) factors through \(B\), so a map
\(f:A\to B\) makes \(I\) an \(A\)-module.  For a nonsingular affine
variety, \(\Omega_{A/k}\) is a finite projective \(A\)-module and the
conormal sequence of a polynomial presentation is left exact.
\end{context}

\begin{problem}
Given a \(k\)-algebra map \(f:A\to B\), prove that it lifts to a
\(k\)-algebra map \(g:A\to B'\).  Establish this in the following
steps.

(a) If \(g,g':A\to B'\) are two lifts, prove that \(g-g'\) is a
\(k\)-derivation \(A\to I\).  Conversely, show that adding any such
derivation to a lift produces another lift.

(b) Present \(A=P/J\), where \(P=k[x_1,\ldots,x_N]\).  Construct a
lift \(h:P\to B'\) of \(P\to A\xrightarrow{f}B\), and show that its
restriction induces an \(A\)-linear map
\[
\overline h:J/J^2\longrightarrow I.
\]

(c) Use the conormal sequence and (a)--(b) to modify \(h\) so that it
kills \(J\), and thereby obtain the desired lift from \(A\).
\end{problem}
\end{exerciseintro}

\begin{solution}
(a) Since the two maps have the same composite with \(B'\to B\),
their difference \(\theta=g-g'\) takes values in \(I\).  It is
additive and vanishes on \(k\).  For \(a,b\in A\),
\[
\begin{aligned}
\theta(ab)
&=g(a)g(b)-g'(a)g'(b)\\
&=g(a)\theta(b)+g'(b)\theta(a)\\
&=a\theta(b)+b\theta(a).
\end{aligned}
\]
The last equality uses the \(A\)-module structure on \(I\).  Thus
\(\theta\in\operatorname{Der}_k(A,I)\), equivalently
\(\theta\in\operatorname{Hom}_A(\Omega_{A/k},I)\).

Conversely, let \(\theta:A\to I\) be a \(k\)-derivation and set
\(g_\theta=g+\theta\).  Because \(I^2=0\),
\[
g_\theta(a)g_\theta(b)
=g(ab)+a\theta(b)+b\theta(a)
=g_\theta(ab).
\]
Also \(\theta(1)=0\), so \(g_\theta\) is a unital \(k\)-algebra map,
and it still lifts \(f\).

(b) Choose arbitrary lifts in \(B'\) of the elements
\(f(\overline{x}_i)\in B\).  The universal property of the
polynomial ring gives a homomorphism \(h:P\to B'\) lifting the
prescribed map \(P\to B\).  Since the latter kills \(J\), one has
\(h(J)\subseteq I\), and therefore \(h(J^2)\subseteq I^2=0\).  Thus
restriction induces an additive map
\(\overline h:J/J^2\to I\).

If \(p\in P\) and \(j\in J\), then
\(h(pj)=h(p)h(j)\).  On \(I\), multiplication by \(h(p)\) depends
only on the image of \(p\) in \(A\), because elements of \(I\)
annihilate \(I\).  This proves that \(\overline h\) is
\(A\)-linear.

(c) Nonsingularity gives the exact conormal sequence
\[
0\longrightarrow J/J^2
\longrightarrow\Omega_{P/k}\otimes_PA
\longrightarrow\Omega_{A/k}\longrightarrow0.
\tag{1}
\]
Since \(\Omega_{A/k}\) is projective, applying
\(\operatorname{Hom}_A(-,I)\) to (1) gives a surjection
\[
\operatorname{Hom}_A(\Omega_{P/k}\otimes_PA,I)
\twoheadrightarrow
\operatorname{Hom}_A(J/J^2,I).
\tag{2}
\]
The source of (2) is naturally
\(\operatorname{Hom}_P(\Omega_{P/k},I)\), the module of
\(k\)-derivations \(P\to I\).  Choose a derivation \(\theta\) whose
restriction in (2) is \(\overline h\).

Regard \(I\) as an ideal of \(B'\) and put \(h'=h-\theta\).  The
square-zero calculation in (a), applied to \(P\), shows that \(h'\)
is again a \(k\)-algebra homomorphism lifting \(P\to B\).  For every
\(j\in J\), the choice of \(\theta\) gives
\(\theta(j)=h(j)\), so \(h'(j)=0\).  Hence \(h'\) factors uniquely
through \(P/J=A\), producing the required lift
\(g:A\to B'\).
\end{solution}

\begin{exerciseintro}
% record: H-II-08-007
\exerciseheading{Exercise II.8.7}
\addcontentsline{toc}{subsection}{Exercise II.8.7}

\begin{context}
An infinitesimal extension of a finite-type \(k\)-scheme \(X\) by a
coherent sheaf \(\mathcal F\) is a scheme \(X'\) with a square-zero
ideal \(\mathcal I\) such that
\((X',\mathcal O_{X'}/\mathcal I)\cong
(X,\mathcal O_X)\) and \(\mathcal I\cong\mathcal F\) as
\(\mathcal O_X\)-modules.  Its trivial extension has sheaf of rings
\(\mathcal O_X\oplus\mathcal F\), with multiplication
\[
(a,u)(b,v)=(ab,av+bu).
\]
We may use the infinitesimal lifting property for smooth affine
\(k\)-schemes and exactness of the augmented Cech complex of a
quasi-coherent sheaf for a finite principal-open cover of an affine
scheme.
\end{context}

\begin{problem}
Let \(X\) be affine and smooth over \(k\), and let \(\mathcal F\) be
coherent on \(X\).  Prove that every infinitesimal extension of
\(X\) by \(\mathcal F\) is isomorphic, as an extension, to the
trivial one.
\end{problem}
\end{exerciseintro}

\begin{solution}
Write \(X=\operatorname{Spec}A\), and let
\((X',\mathcal I)\) be such an extension.  A square-zero closed
immersion is a homeomorphism on underlying spaces, so we identify the
spaces of \(X\) and \(X'\).

Choose a finite principal-open cover
\(U_i=D(f_i)\) of \(X\) such that the corresponding opens \(U_i'\)
of \(X'\) are affine.  To obtain one, start with affine opens of
\(X'\), refine their reductions by principal opens of the affine
scheme \(X\), and lift each defining function; its nonvanishing locus
in the affine thickening is affine.  On each member of this cover we
have a square-zero extension of rings
\[
0\longrightarrow\mathcal F(U_i)
\longrightarrow\mathcal O_{X'}(U_i')
\longrightarrow\mathcal O_X(U_i)
\longrightarrow0.
\]
Since \(U_i\) is smooth and affine, the infinitesimal lifting property
provides a \(k\)-algebra section \(s_i\) of the last arrow.

On \(U_{ij}=U_i\cap U_j\), part (a) of the lifting property identifies
\[
c_{ij}=s_i-s_j
\]
with a \(k\)-derivation
\(\mathcal O_X|_{U_{ij}}\to\mathcal F|_{U_{ij}}\).  Equivalently,
\(c_{ij}\) is a section on \(U_{ij}\) of
\[
\mathcal D=
\mathcal Hom_{\mathcal O_X}(\Omega_{X/k},\mathcal F).
\]
Because \(X\) is nonsingular, \(\Omega_{X/k}\) is finite locally
free, so \(\mathcal D\) is quasi-coherent.  The identities among the
differences make \((c_{ij})\) a Cech one-cocycle.  Exactness of the
Cech complex on the affine scheme \(X\) gives sections
\(d_i\in\mathcal D(U_i)\) such that
\[
c_{ij}=d_i-d_j.
\]
By the converse assertion in the lifting property,
\(s_i'=s_i-d_i\) is still a ring-theoretic splitting.  The displayed
identity says that the \(s_i'\) agree on all overlaps, so they glue to
a morphism of sheaves of \(k\)-algebras
\[
s:\mathcal O_X\longrightarrow\mathcal O_{X'}
\]
splitting \(\mathcal O_{X'}\twoheadrightarrow\mathcal O_X\).

Now define
\[
\Phi:\mathcal O_X\oplus\mathcal F
\longrightarrow\mathcal O_{X'},
\qquad (a,u)\longmapsto s(a)+u.
\]
It is an additive isomorphism because \(s\) splits the quotient and
\(\mathcal F=\mathcal I\) is its kernel.  Moreover, \(\mathcal I^2=0\)
and its \(\mathcal O_X\)-module structure is induced through the
quotient, so
\[
\Phi(a,u)\Phi(b,v)=\Phi(ab,av+bu).
\]
Thus \(\Phi\) is an isomorphism of sheaves of rings and identifies
\((X',\mathcal I)\) with the trivial extension.
\end{solution}

\begin{exerciseintro}
% record: H-II-08-008
\exerciseheading{Exercise II.8.8}
\addcontentsline{toc}{subsection}{Exercise II.8.8}

\begin{context}
For a nonsingular projective variety \(X\), define
\[
P_m(X)=\dim_kH^0(X,\omega_X^{\otimes m})\quad(m>0),
\qquad
h^{q,0}(X)=\dim_kH^0(X,\Omega_{X/k}^q),
\]
where \(\Omega_{X/k}^q=\bigwedge^q\Omega_{X/k}\).  We may use two
facts from the rational-map extension argument: a rational map from a
nonsingular variety to a projective variety is defined outside a
subset of codimension at least two; and sections of a locally free
sheaf on a nonsingular variety extend uniquely across such a subset.
\end{context}

\begin{problem}
Prove that every plurigenus \(P_m\) and every number \(h^{q,0}\) is a
birational invariant among nonsingular projective varieties.
\end{problem}
\end{exerciseintro}

\begin{solution}
Let \(X\) and \(X'\) be birational nonsingular projective varieties,
and choose dense opens \(U\subseteq X\) and \(U'\subseteq X'\) that
are isomorphic.  Regard this isomorphism as a rational map
\(\varphi:X\dashrightarrow X'\).  There is an open
\(V\subseteq X\), containing \(U\), such that
\(X\setminus V\) has codimension at least two and \(\varphi\) extends
to a morphism \(f:V\to X'\).

Functoriality of differentials gives
\[
f^*\Omega_{X'/k}^q\longrightarrow\Omega_{V/k}^q.
\]
Likewise, taking the top exterior power of the differential and then
the \(m\)-th tensor power gives
\[
f^*(\omega_{X'}^{\otimes m})
\longrightarrow\omega_V^{\otimes m}.
\]
Pulling back global sections and using extension across
\(X\setminus V\) therefore defines linear maps
\[
H^0(X',\Omega_{X'/k}^q)
\longrightarrow H^0(X,\Omega_{X/k}^q)
\tag{1}
\]
and
\[
H^0(X',\omega_{X'}^{\otimes m})
\longrightarrow H^0(X,\omega_X^{\otimes m}).
\tag{2}
\]
Both are injective.  Indeed, on \(U\) the map \(f\) is the chosen
isomorphism \(U\cong U'\).  A section in the kernel consequently
vanishes on the dense open \(U'\), and a section of a locally free
sheaf that vanishes on a dense open of an integral scheme is zero.
It follows from (1)--(2) that
\[
h^{q,0}(X')\leq h^{q,0}(X),
\qquad P_m(X')\leq P_m(X).
\]
Apply the same construction to the inverse birational map
\(X'\dashrightarrow X\).  The reverse inequalities follow, proving
\[
h^{q,0}(X')=h^{q,0}(X),
\qquad P_m(X')=P_m(X)
\]
for every permitted \(q\) and every \(m>0\).
\end{solution}

\begin{exerciseintro}
\corpussection{II.9 Formal Schemes}
\addcontentsline{toc}{section}{II.9 Formal Schemes}

% record: H-II-09-001
\exerciseheading{Exercise II.9.1}
\addcontentsline{toc}{subsection}{Exercise II.9.1}

\begin{context}
Let \(k\) be algebraically closed, let
\(X=\mathbb P_k^n\), and let \(Y\subseteq X\) be a connected,
nonsingular, positive-dimensional closed subvariety with ideal sheaf
\(\mathcal I\).  Its formal completion has underlying space \(Y\) and
structure sheaf
\[
\mathcal O_{\widehat X}
=\varprojlim_{r\geq1}\mathcal O_X/\mathcal I^r.
\]
We may use the Euler sequence on projective space, the exact conormal
sequence for the regular immersion \(Y\hookrightarrow X\), and
\[
\mathcal I^r/\mathcal I^{r+1}
\cong\operatorname{Sym}^r_{\mathcal O_Y}
(\mathcal I/\mathcal I^2).
\]
\end{context}

\begin{problem}
(a) Prove that there is an injection
\[
\mathcal I/\mathcal I^2
\hookrightarrow\mathcal O_Y(-1)^{\oplus(n+1)}.
\]

(b) Prove that
\[
H^0(Y,\mathcal I^r/\mathcal I^{r+1})=0
\qquad(r\geq1).
\]

(c) Using the successive infinitesimal-neighborhood sequences, prove
\[
H^0(Y,\mathcal O_X/\mathcal I^r)=k
\qquad(r\geq1).
\]

(d) Deduce that every formal-regular function along \(Y\) is constant:
\[
\Gamma(\widehat X,\mathcal O_{\widehat X})=k.
\]
\end{problem}
\end{exerciseintro}

\begin{solution}
(a) Since both \(X\) and \(Y\) are nonsingular, the closed immersion
is regular and its conormal sequence begins
\[
0\longrightarrow\mathcal I/\mathcal I^2
\longrightarrow\Omega_{X/k}|_Y.
\tag{1}
\]
The Euler sequence begins
\[
0\longrightarrow\Omega_{X/k}
\longrightarrow\mathcal O_X(-1)^{\oplus(n+1)}
\longrightarrow\mathcal O_X\longrightarrow0.
\]
Tensoring with \(\mathcal O_Y\) preserves injectivity on the left
because the last term \(\mathcal O_X\) is locally free.  Restricting
the resulting injection to \(Y\) and composing with (1) gives
\[
\mathcal I/\mathcal I^2
\hookrightarrow\mathcal O_Y(-1)^{\oplus(n+1)}.
\]

(b) Let \(\mathcal Q\) be the quotient of the injection in (a).
The conormal and Euler sequences give
\[
0\longrightarrow\Omega_{Y/k}\longrightarrow\mathcal Q
\longrightarrow\mathcal O_Y\longrightarrow0.
\]
Its end terms are locally free, so \(\mathcal Q\) is locally free.
It is this fact about the quotient that makes the injection in (a)
locally split.  Taking the \(r\)-th symmetric power therefore gives an
injection
\[
\mathcal I^r/\mathcal I^{r+1}
\cong\operatorname{Sym}^r(\mathcal I/\mathcal I^2)
\hookrightarrow
\operatorname{Sym}^r
(\mathcal O_Y(-1)^{\oplus(n+1)})
\cong\mathcal O_Y(-r)^{\oplus N_r}
\tag{2}
\]
for a positive integer \(N_r\).

It remains to show \(H^0(Y,\mathcal O_Y(-r))=0\).  A connected
nonsingular variety is integral: distinct irreducible components of
a regular scheme are disjoint, and connectedness leaves only one.
Choose a hyperplane \(H\) that does not contain \(Y\).  Since
\(\dim Y>0\), the projective dimension theorem gives
\(H\cap Y\neq\varnothing\).  Multiplication by an equation of
\(rH\) is an injection
\[
\mathcal O_Y(-r)\hookrightarrow\mathcal O_Y.
\]
A global section maps to a global regular function on the projective
integral variety \(Y\), hence to a constant in \(k\).  That function
vanishes on \(H\cap Y\), so the constant is zero.  The original
section is therefore zero.  Taking global sections in (2) proves (b).

(c) For every \(r\geq1\) there is an exact sequence on \(Y\)
\[
0\longrightarrow\mathcal I^r/\mathcal I^{r+1}
\longrightarrow\mathcal O_X/\mathcal I^{r+1}
\longrightarrow\mathcal O_X/\mathcal I^r
\longrightarrow0.
\tag{3}
\]
By (b), taking global sections makes the last map in (3) injective:
\[
H^0(Y,\mathcal O_X/\mathcal I^{r+1})
\hookrightarrow H^0(Y,\mathcal O_X/\mathcal I^r).
\]
For \(r=1\), the right side is
\(H^0(Y,\mathcal O_Y)=k\).  At every level the constant sections give
a copy of \(k\) in the group on the left.  Induction therefore
squeezes every group between \(k\) and \(k\), proving the assertion.

(d) Inverse limits of sheaves are computed sectionwise.  Hence
\[
\begin{aligned}
\Gamma(\widehat X,\mathcal O_{\widehat X})
&=\varprojlim_r
\Gamma(Y,\mathcal O_X/\mathcal I^r)\\
&=\varprojlim_r k=k.
\end{aligned}
\]
The transition maps on the final inverse system are the identity on
constant sections.
\end{solution}

\begin{exerciseintro}
% record: H-II-09-002
\exerciseheading{Exercise II.9.2}
\addcontentsline{toc}{subsection}{Exercise II.9.2}

\begin{context}
Retain the hypotheses of Exercise II.9.1: \(X=\mathbb P_k^n\) and
\(Y\subseteq X\) is connected, nonsingular, and positive-dimensional.
If \(\mathcal I\) is its ideal sheaf and
\[
Y_r=(Y,\mathcal O_X/\mathcal I^r),
\]
then
\(\Gamma(Y_r,\mathcal O_{Y_r})=k\) for every \(r\geq1\).
We may use the Krull intersection theorem for a noetherian domain.
\end{context}

\begin{problem}
Let \(f:X\to Z\) be a morphism of \(k\)-varieties.  If \(f(Y)\) is
the single point \(P\in Z\), prove that \(f(X)=\{P\}\).
\end{problem}
\end{exerciseintro}

\begin{solution}
Choose an affine neighborhood \(W=\operatorname{Spec}B\) of \(P\),
and let \(\mathfrak m_P\subseteq B\) be its maximal ideal.  The
underlying space of every \(Y_r\) is \(Y\), whose image is \(P\), so
the restriction \(Y_r\to Z\) factors through the open set \(W\).
It corresponds to a \(k\)-algebra map
\[
B\longrightarrow\Gamma(Y_r,\mathcal O_{Y_r})=k.
\]
After reduction to \(Y\), this map is evaluation at \(P\).  It is
therefore the quotient \(B\to B/\mathfrak m_P=k\).  Thus every
infinitesimal neighborhood \(Y_r\) maps scheme-theoretically to
\(P\).

The open set \(f^{-1}(W)\) contains \(Y\).  Choose a nonempty affine
open \(U=\operatorname{Spec}A\) inside it that meets \(Y\), and let
\(I\subsetneq A\) define \(Y\cap U\).  Since \(U\) is an open
subscheme of projective space, \(A\) is a noetherian domain.  If
\(\varphi:B\to A\) is the map induced by \(f|_U\), the conclusion for
all \(Y_r\) says
\[
\varphi(\mathfrak m_P)\subseteq I^r
\qquad\text{for every }r\geq1.
\]
The Krull intersection theorem gives
\[
\bigcap_{r\geq1}I^r=0,
\]
and consequently \(\varphi(\mathfrak m_P)=0\).  Hence \(f|_U\)
factors through the reduced point \(P\).

Let \(c_P:X\to Z\) be the constant morphism with value \(P\).  A
variety is separated, so the equalizer of \(f\) and \(c_P\) is closed
in \(X\).  It contains the nonempty open \(U\), which is dense because
\(X\) is integral.  The equalizer is therefore all of \(X\), and
\(f=c_P\).
\end{solution}

\begin{exerciseintro}
% record: H-II-09-003
\exerciseheading{Exercise II.9.3}
\addcontentsline{toc}{subsection}{Exercise II.9.3}

\begin{context}
Let \(\mathfrak X\) be an affine noetherian formal scheme with an
ideal of definition \(\mathfrak I\).  A coherent sheaf
\(\mathfrak G\) is complete:
\[
\mathfrak G\cong
\varprojlim_n\mathfrak G/\mathfrak I^n\mathfrak G,
\]
and on an affine formal open its quotient sections form the
surjective system \(M/\mathfrak I^nM\).  We may also use the affine
lifting lemma: in an exact sequence of sheaves on an affine scheme,
global sections are exact if the kernel is quasi-coherent, even when
the other two sheaves are only sheaves of modules.
\end{context}

\begin{problem}
Let
\[
0\longrightarrow\mathfrak F'
\longrightarrow\mathfrak F
\longrightarrow\mathfrak F''\longrightarrow0
\tag{1}
\]
be exact, and suppose \(\mathfrak F'\) is coherent.

(a) For an affine formal open
\(\mathfrak U\subseteq\mathfrak X\) and \(n>0\), prove exactness of
\[
\begin{aligned}
0\longrightarrow&
\Gamma(\mathfrak U,\mathfrak F'/\mathfrak I^n\mathfrak F')
\longrightarrow
\Gamma(\mathfrak U,\mathfrak F/\mathfrak I^n\mathfrak F')\\
\longrightarrow&
\Gamma(\mathfrak U,\mathfrak F'')
\longrightarrow0.
\end{aligned}
\tag{2}
\]

(b) Pass to the inverse limit, prove
\[
\mathfrak F\cong
\varprojlim_n\mathfrak F/\mathfrak I^n\mathfrak F',
\]
and deduce that applying global sections to (1) gives a short exact
sequence.
\end{problem}
\end{exerciseintro}

\begin{solution}
(a) Regard \(\mathfrak I^n\mathfrak F'\) as a subsheaf of
\(\mathfrak F\) through the first map in (1).  Quotienting gives an
exact sequence
\[
0\longrightarrow
\mathfrak F'/\mathfrak I^n\mathfrak F'
\longrightarrow
\mathfrak F/\mathfrak I^n\mathfrak F'
\longrightarrow\mathfrak F''\longrightarrow0.
\tag{3}
\]
The first term is coherent and is annihilated by
\(\mathfrak I^n\).  It is therefore a quasi-coherent sheaf on the
ordinary affine scheme
\[
U_n=(\mathfrak U,
\mathcal O_{\mathfrak U}/\mathfrak I^n).
\]
All sheaves in (3) have the same underlying topological space.
Only its kernel is asserted to be a module on \(U_n\).
The vanishing \(H^1(U_n,\mathfrak F'/\mathfrak I^n\mathfrak F')=0\)
therefore proves (2) by the long exact sequence of abelian sheaves.  Concretely,
local lifts of a section of \(\mathfrak F''\) differ by sections of
the quasi-coherent first term, and affine exactness removes that
Cech one-cocycle so that the lifts glue.

(b) For a fixed affine \(\mathfrak U\), abbreviate the three terms of
(2) by \(A_n,B_n,C\).  Coherence of \(\mathfrak F'\) identifies
\[
A_n=M/\mathfrak I^nM
\]
for a finite module \(M\), so the transition maps
\(A_{n+1}\to A_n\) are surjective.  Taking inverse limits of
\[
0\longrightarrow A_n\longrightarrow B_n\longrightarrow C
\longrightarrow0
\]
is consequently exact.  For surjectivity on the right, one may choose
a lift in \(B_1\) and recursively alter a lift in \(B_{n+1}\) by an
element of \(A_{n+1}\) so that it restricts to the lift already chosen
in \(B_n\).  Hence
\[
0\longrightarrow\Gamma(\mathfrak U,\mathfrak F')
\longrightarrow
\varprojlim_n
\Gamma(\mathfrak U,\mathfrak F/\mathfrak I^n\mathfrak F')
\longrightarrow
\Gamma(\mathfrak U,\mathfrak F'')
\longrightarrow0.
\tag{4}
\]

Set
\[
\mathfrak G=
\varprojlim_n\mathfrak F/\mathfrak I^n\mathfrak F'.
\]
Limits of sheaves are computed sectionwise, so (4), for all affine
formal opens in a basis, gives an exact sequence
\[
0\longrightarrow\mathfrak F'
\longrightarrow\mathfrak G
\longrightarrow\mathfrak F''\longrightarrow0.
\tag{5}
\]
There is a natural morphism from (1) to (5): it is the completeness
isomorphism on \(\mathfrak F'\), the map
\(\mathfrak F\to\mathfrak G\) in the middle, and the identity on
\(\mathfrak F''\).  On every stalk, the short five lemma makes the
middle map an isomorphism.  Thus
\[
\mathfrak F\cong
\varprojlim_n\mathfrak F/\mathfrak I^n\mathfrak F'.
\]
Finally take \(\mathfrak U=\mathfrak X\) in (4) and use this
isomorphism.  The resulting sequence is
\[
0\longrightarrow\Gamma(\mathfrak X,\mathfrak F')
\longrightarrow\Gamma(\mathfrak X,\mathfrak F)
\longrightarrow\Gamma(\mathfrak X,\mathfrak F'')
\longrightarrow0,
\]
as required.
\end{solution}

\begin{exerciseintro}
% record: H-II-09-004
\exerciseheading{Exercise II.9.4}
\addcontentsline{toc}{subsection}{Exercise II.9.4}

\begin{context}
On an affine noetherian formal scheme
\(\mathfrak U=\operatorname{Spf}A\), the functor
\[
M\longmapsto M^\Delta
\]
is an exact equivalence from finite \(A\)-modules to coherent
\(\mathcal O_{\mathfrak U}\)-modules, with quasi-inverse global
sections.  We may use Exercise II.9.3: global sections are exact when
the kernel in a short exact sequence is coherent.
\end{context}

\begin{problem}
Let \(\mathfrak X\) be a noetherian formal scheme.  If
\[
0\longrightarrow\mathfrak F'
\longrightarrow\mathfrak F
\longrightarrow\mathfrak F''\longrightarrow0
\]
is exact and both \(\mathfrak F'\) and \(\mathfrak F''\) are
coherent, prove that \(\mathfrak F\) is coherent.
\end{problem}
\end{exerciseintro}

\begin{solution}
Coherence is local, so take
\(\mathfrak X=\operatorname{Spf}A\).  Put
\[
M'=\Gamma(\mathfrak X,\mathfrak F'),\qquad
M=\Gamma(\mathfrak X,\mathfrak F),\qquad
M''=\Gamma(\mathfrak X,\mathfrak F'').
\]
Exercise II.9.3 gives a short exact sequence
\[
0\longrightarrow M'\longrightarrow M\longrightarrow M''
\longrightarrow0.
\tag{1}
\]
The outer modules are finite over \(A\), by the affine equivalence.
Lifting a finite generating set of \(M''\) to \(M\) and adjoining a
finite generating set of \(M'\) shows that \(M\) is finite.

Apply the exact functor \((\,\cdot\,)^\Delta\) to (1).  There is a
natural morphism from the resulting sequence
\[
0\longrightarrow(M')^\Delta
\longrightarrow M^\Delta
\longrightarrow(M'')^\Delta
\longrightarrow0
\]
to the original sheaf sequence.  The two outer vertical maps are
isomorphisms because \(\mathfrak F'\) and \(\mathfrak F''\) are
coherent.  The short five lemma, applied on stalks, makes
\[
M^\Delta\xrightarrow{\sim}\mathfrak F
\]
an isomorphism.  Since \(M\) is finite, \(\mathfrak F\) is coherent.
The same argument on each affine formal open proves the result on
\(\mathfrak X\).
\end{solution}

\begin{exerciseintro}
% record: H-II-09-005
\exerciseheading{Exercise II.9.5}
\addcontentsline{toc}{subsection}{Exercise II.9.5}

\begin{context}
A noetherian formal scheme is quasi-compact and has a finite affine
formal open cover.  On
\(\mathfrak U=\operatorname{Spf}A\), coherent sheaves are exactly the
sheaves \(M^\Delta\) associated to finite \(A\)-modules, and this
correspondence is faithful.
\end{context}

\begin{problem}
Let \(\mathfrak F\) be a coherent sheaf on a noetherian formal scheme.
If its global sections generate it, prove that finitely many of those
global sections already generate it.
\end{problem}
\end{exerciseintro}

\begin{solution}
Let \(\{s_\alpha\}_{\alpha\in\Lambda}\) be a generating family and
choose a finite affine formal open cover
\[
\mathfrak X=\mathfrak U_1\cup\cdots\cup\mathfrak U_t.
\]
Write
\[
\mathfrak F|_{\mathfrak U_i}\cong M_i^\Delta
\]
for a finite module \(M_i\) over
\(A_i=\Gamma(\mathfrak U_i,\mathcal O_{\mathfrak X})\).  The
restrictions \(s_\alpha|_{\mathfrak U_i}\) correspond to elements
\(m_{\alpha,i}\in M_i\).

Let \(N_i\subseteq M_i\) be the submodule generated by all the
\(m_{\alpha,i}\).  Since the original sections generate the sheaf,
the associated sheaf of \(M_i/N_i\) is zero.  The quotient is a
finite \(A_i\)-module, so faithfulness of the affine equivalence gives
\(M_i/N_i=0\).  Thus the elements \(m_{\alpha,i}\) generate \(M_i\).
Because \(M_i\) is finite, some finite subfamily, indexed by
\(\Lambda_i\subseteq\Lambda\), already generates it.

The union
\(\Lambda_1\cup\cdots\cup\Lambda_t\) is finite.  The corresponding
global sections generate \(\mathfrak F\) on every member of the
cover, and hence on all of \(\mathfrak X\).
\end{solution}

\begin{exerciseintro}
% record: H-II-09-006
\exerciseheading{Exercise II.9.6}
\addcontentsline{toc}{subsection}{Exercise II.9.6}

\begin{context}
Let \(\mathfrak X\) be a noetherian formal scheme with ideal of
definition \(\mathfrak I\), and put
\[
Y_n=(|\mathfrak X|,
\mathcal O_{\mathfrak X}/\mathfrak I^n),\qquad n\geq1.
\]
Assume that the inverse system of rings
\[
A_n=\Gamma(Y_n,\mathcal O_{Y_n})
\]
satisfies the Mittag--Leffler condition.  We may use the formal
coherence theorem: a compatible inverse system of coherent sheaves
\(\mathcal F_n\) on \(Y_n\), with
\(\mathcal F_{n+1}\otimes\mathcal O_{Y_n}\cong\mathcal F_n\), has a
coherent limit on \(\mathfrak X\) whose reduction to \(Y_n\) is
\(\mathcal F_n\).  We also use the elementary fact that a nonempty
Mittag--Leffler inverse sequence of sets has a nonempty inverse limit.
In (d), all the projective schemes \(Y_n\) and their transition maps are over the same field \(k\).
\end{context}

\begin{problem}
Prove that reduction induces an isomorphism
\[
\operatorname{Pic}(\mathfrak X)
\xrightarrow{\sim}
\varprojlim_n\operatorname{Pic}(Y_n)
\]
by establishing the following points.

(a) Show that the inverse system of unit groups \(A_n^*\) is
Mittag--Leffler.

(b) If a coherent \(\mathfrak F\) has
\(\mathfrak F/\mathfrak I^n\mathfrak F\cong\mathcal O_{Y_n}\) for
every \(n\), prove \(\mathfrak F\cong\mathcal O_{\mathfrak X}\).
Deduce injectivity of the Picard-group map.

(c) Starting from a compatible family of invertible sheaf classes on
the \(Y_n\), construct an invertible sheaf on \(\mathfrak X\) that
reduces to that family.

(d) Verify the ring Mittag--Leffler hypothesis when \(\mathfrak X\)
is affine, and also when every \(Y_n\) is projective over a field.
\end{problem}
\end{exerciseintro}

\begin{solution}
(a) For \(m\geq n\), let
\(\rho_{mn}:A_m\to A_n\) be restriction.  Its kernel consists of
sections lying in the nilpotent ideal
\(\mathfrak I^n/\mathfrak I^m\); hence that kernel is nilpotent.
More importantly, nilpotence is valid sheafwise: a section of
\(\mathcal O_{Y_m}\) whose restriction to \(Y_n\) is a unit is
locally a unit, and therefore has a unique global inverse.  It follows
that
\[
A_m^*=\rho_{mn}^{-1}(A_n^*)
\]
and consequently
\[
\operatorname{im}(A_m^*\to A_n^*)
=\operatorname{im}(A_m\to A_n)\cap A_n^*.
\tag{1}
\]
For fixed \(n\), the ring images on the right stabilize as
\(m\) increases by hypothesis.  Their intersections with \(A_n^*\)
then stabilize as well.  Thus \((A_n^*)\) is Mittag--Leffler.

(b) Put
\(\mathfrak F_n=\mathfrak F/\mathfrak I^n\mathfrak F\), and let
\[
T_n=\operatorname{Isom}_{\mathcal O_{Y_n}}
(\mathfrak F_n,\mathcal O_{Y_n}).
\]
Each \(T_n\) is nonempty and is a torsor under
\[
\operatorname{Aut}(\mathcal O_{Y_n})=A_n^*.
\]
Restriction gives an inverse system of the \(T_n\).  For fixed \(n\),
the image of \(T_m\to T_n\) is a coset under the image of
\(A_m^*\to A_n^*\).  By (a), the latter subgroups eventually
stabilize.  The former images are nested nonempty cosets, so they
stabilize too.  Hence \((T_n)\) is a nonempty Mittag--Leffler system
of sets and has a compatible element
\((\varphi_n)\in\varprojlim T_n\).

Coherence gives the completeness isomorphisms
\[
\mathfrak F\cong\varprojlim_n\mathfrak F_n,
\qquad
\mathcal O_{\mathfrak X}\cong
\varprojlim_n\mathcal O_{Y_n}.
\]
Taking the inverse limit of the compatible \(\varphi_n\) therefore
gives
\[
\mathfrak F\xrightarrow{\sim}\mathcal O_{\mathfrak X}.
\]
If an invertible sheaf on \(\mathfrak X\) maps to zero in every
\(\operatorname{Pic}(Y_n)\), its reductions satisfy the hypothesis
just considered, so it is trivial.  The natural Picard-group map is
injective.

(c) Let \((\xi_n)\) represent an element of the inverse limit of the
Picard groups.  Choose an invertible sheaf \(\mathcal L_n\) in the
class \(\xi_n\).  Compatibility of the classes permits choices of
isomorphisms
\[
\alpha_n:
\mathcal L_{n+1}\otimes_{\mathcal O_{Y_{n+1}}}
\mathcal O_{Y_n}\xrightarrow{\sim}\mathcal L_n.
\]
Compose the quotient map
\(\mathcal L_{n+1}\to
\mathcal L_{n+1}\otimes\mathcal O_{Y_n}\)
with \(\alpha_n\), and use composites to define
\(\mathcal L_m\to\mathcal L_n\) for every \(m\geq n\).  This is an
inverse system satisfying the hypotheses of the formal coherence
theorem.  Consequently
\[
\mathfrak L=\varprojlim_n\mathcal L_n
\]
is coherent and
\(\mathfrak L/\mathfrak I^n\mathfrak L\cong\mathcal L_n\).

We prove that \(\mathfrak L\) is invertible without allowing the
trivializing neighborhood to depend on \(n\).  Fix a point, choose
an affine formal neighborhood
\(\mathfrak U=\operatorname{Spf}A\), and write
\[
\mathfrak L|_{\mathfrak U}=M^\Delta
\]
for a finite \(A\)-module \(M\).  Let \(J\subseteq A\) be the ideal
of definition induced by \(\mathfrak I\).  Shrink \(\mathfrak U\)
once so that the line bundle
\(M/JM\) on \(\operatorname{Spec}(A/J)\) is free.  Choose a basis
\(\overline s\) and lift it to \(s\in M\).

For every \(n\), consider
\[
A/J^n\longrightarrow M/J^nM,\qquad 1\longmapsto s.
\tag{2}
\]
The target is the module of the line bundle
\(\mathcal L_n|_{\mathfrak U}\), so it is locally free of rank one.
In a local trivialization, (2) is multiplication by an element
\(c\in A/J^n\).  Its reduction modulo \(J\) is a unit because
\(\overline s\) is a basis.  Since \(J/J^n\) is nilpotent, \(c\) is
a unit: if \(cd=1-e\) with \(e\in J/J^n\), a finite geometric series
in \(e\) supplies an inverse.  Thus (2) is an isomorphism at every
point, and hence globally.

The adic ring \(A\) and its finite module \(M\) are complete, so
taking inverse limits in (2) gives
\[
A\xrightarrow{\sim}M.
\]
Therefore \(\mathfrak L\) is free of rank one on the fixed
neighborhood \(\mathfrak U\).  It is an invertible sheaf on
\(\mathfrak X\), and its reductions represent the prescribed
\(\xi_n\).  This proves surjectivity.

(d) If \(\mathfrak X=\operatorname{Spf}A\) with ideal of definition
\(J\), then
\[
Y_n=\operatorname{Spec}(A/J^n),\qquad
\Gamma(Y_n,\mathcal O_{Y_n})=A/J^n.
\]
The transition maps are surjective, so the system is
Mittag--Leffler.

If every \(Y_n\) is projective over a field \(k\), each
\(A_n=H^0(Y_n,\mathcal O_{Y_n})\) is a finite-dimensional
\(k\)-vector space.  For fixed \(n\), the images of \(A_m\to A_n\)
form a descending chain of subspaces of the finite-dimensional space
\(A_n\), so they stabilize.  This is the Mittag--Leffler condition.
\end{solution}

\corpuschapter{Chapter III}
\addcontentsline{toc}{chapter}{Chapter III}

\begin{exerciseintro}
\corpussection{III.2 Cohomology of Sheaves}
\addcontentsline{toc}{section}{III.2 Cohomology of Sheaves}

% record: H-III-02-001
\exerciseheading{Exercise III.2.1}
\addcontentsline{toc}{subsection}{Exercise III.2.1}

\begin{context}
For a locally closed subset \(T\subseteq X\), write
\(\mathbb Z_T\) for the constant sheaf \(\mathbb Z\) on \(T\),
extended by zero to \(X\).  A constant sheaf on an irreducible
topological space is flasque, as is its direct image under a closed
immersion.  A flasque resolution computes sheaf cohomology.
\end{context}

\begin{problem}
(a) Let \(X=\mathbb A_k^1\), where \(k\) is infinite.  For distinct
closed points \(P,Q\), put \(U=X\setminus\{P,Q\}\).  Prove
\[
H^1(X,\mathbb Z_U)\neq0.
\]

(b) More generally, in \(X=\mathbb A_k^n\) take \(n+1\) affine
hyperplanes in general position, let \(Y\) be their union, and put
\(U=X\setminus Y\).  Prove
\[
H^n(X,\mathbb Z_U)\neq0.
\]
\end{problem}
\end{exerciseintro}

\begin{solution}
(a) Let \(V=\{P,Q\}\).  There is an exact sequence
\[
0\longrightarrow\mathbb Z_U
\longrightarrow\mathbb Z_X
\longrightarrow\mathbb Z_V\longrightarrow0.
\tag{1}
\]
The affine line is irreducible, so \(\mathbb Z_X\) is flasque.
The sheaf \(\mathbb Z_V\) is a finite sum of skyscraper sheaves and
is also flasque.  Thus (1) is a flasque resolution.  On global
sections its nonzero map is the diagonal map, and hence
\[
H^1(X,\mathbb Z_U)
\cong\operatorname{coker}
(\mathbb Z\longrightarrow\mathbb Z\oplus\mathbb Z)
\cong\mathbb Z.
\]

(b) Write the hyperplanes as \(H_0,\ldots,H_n\), chosen so that every
intersection of \(r\leq n\) of them is a nonempty affine linear space
of codimension \(r\), while
\(\bigcap_{i=0}^nH_i=\varnothing\).  For a nonempty index set \(I\),
put \(H_I=\bigcap_{i\in I}H_i\).  Alternating restriction maps give
an exact complex
\[
\begin{aligned}
0\longrightarrow\mathbb Z_U\longrightarrow\mathbb Z_X
&\longrightarrow\bigoplus_{|I|=1}\mathbb Z_{H_I}
\longrightarrow\bigoplus_{|I|=2}\mathbb Z_{H_I}\\
&\longrightarrow\cdots
\longrightarrow\bigoplus_{|I|=n}\mathbb Z_{H_I}
\longrightarrow0.
\end{aligned}
\tag{2}
\]
Exactness may be checked at a point \(x\).  If \(x\in U\), the first
map is an isomorphism on the stalk.  Otherwise, the indices of the
hyperplanes through \(x\) form a nonempty set, and the stalk complex
is the augmented cochain complex of a simplex, hence exact.

Every nonempty \(H_I\) is irreducible.  All terms following
\(\mathbb Z_U\) in (2) are therefore flasque, so (2) computes
cohomology.  Taking global sections produces
\[
\mathbb Z\longrightarrow
\mathbb Z^{\binom{n+1}{1}}\longrightarrow\cdots
\longrightarrow\mathbb Z^{\binom{n+1}{n}}\longrightarrow0.
\tag{3}
\]
If one appends the final group \(\mathbb Z\), indexed by the full set
\(\{0,\ldots,n\}\), this is the exact augmented cochain complex of
the \(n\)-simplex.  Its last differential onto \(\mathbb Z\) is
surjective.  Consequently the cokernel of the penultimate map in
(3) is \(\mathbb Z\), and
\[
H^n(X,\mathbb Z_U)\cong\mathbb Z\neq0.
\]
\end{solution}

\begin{exerciseintro}
% record: H-III-02-002
\exerciseheading{Exercise III.2.2}
\addcontentsline{toc}{subsection}{Exercise III.2.2}

\begin{context}
Let \(X=\mathbb P_k^1\) over an algebraically closed field, let
\(\mathcal K\) be its sheaf of rational functions, and regard
\(\mathcal O_X\subseteq\mathcal K\).  We use the elementary
partial-fraction decomposition of rational functions.
\end{context}

\begin{problem}
Prove that
\[
0\longrightarrow\mathcal O_X\longrightarrow\mathcal K
\longrightarrow\mathcal K/\mathcal O_X\longrightarrow0
\]
is a flasque resolution.  Deduce that
\[
H^i(X,\mathcal O_X)=0\qquad(i>0).
\]
\end{problem}
\end{exerciseintro}

\begin{solution}
Every nonempty open subset of the integral curve \(X\) has the same
function field \(K=k(X)\).  Thus \(\mathcal K\) is the constant sheaf
with value \(K\), and it is flasque.

Taking the principal part at each closed point gives a sheaf
isomorphism
\[
\mathcal K/\mathcal O_X
\cong
\bigoplus_{p\in X}
i_{p*}(K/\mathcal O_{X,p}).
\tag{1}
\]
Indeed, this can be checked on stalks; a rational function has poles
at only finitely many points.  Each summand is a skyscraper sheaf.
On the noetherian space \(X\), sections of the direct sum over an
open set involve only finitely many summands, so restrictions remain
surjective.  Hence the sheaf in (1) is flasque.

The displayed sequence is therefore a length-one flasque resolution.
It gives \(H^i(X,\mathcal O_X)=0\) for \(i\geq2\), while
\[
H^1(X,\mathcal O_X)
\cong\operatorname{coker}\left(
K\longrightarrow
\bigoplus_{p\in X}K/\mathcal O_{X,p}\right).
\]
Choose an affine coordinate \(t\).  At a finite point \(t=a\),
write the specified principal part as a finite sum of
\(c_{a,j}/(t-a)^j\).  At infinity write it as a polynomial in
\(t\) with zero constant term.  Adding these finitely many
expressions realizes all the prescribed principal parts and has
no poles elsewhere.  The principal-parts map is therefore
surjective.  Its cokernel is zero, proving the assertion for every
\(i>0\).
\end{solution}

\begin{exerciseintro}
% record: H-III-02-003
\exerciseheading{Exercise III.2.3}
\addcontentsline{toc}{subsection}{Exercise III.2.3}

\begin{context}
Let \(Y\) be closed in a topological space \(X\).  For a sheaf
\(\mathcal F\) of abelian groups, define
\[
\Gamma_Y(X,\mathcal F)
=\{s\in\Gamma(X,\mathcal F):
\operatorname{supp}(s)\subseteq Y\}.
\]
Write \(H_Y^i(X,-)\) for the right derived functors of
\(\Gamma_Y(X,-)\).  Injective sheaves are flasque, restrictions of
flasque sheaves are flasque, and flasque resolutions compute ordinary
sheaf cohomology.
\end{context}

\begin{problem}
(a) Prove that \(\Gamma_Y(X,-)\) is left exact.

(b) If
\[
0\to\mathcal F'\to\mathcal F\to\mathcal F''\to0
\]
is exact and \(\mathcal F'\) is flasque, prove that applying
\(\Gamma_Y(X,-)\) gives a short exact sequence.

(c) Prove that every flasque sheaf is acyclic for
\(\Gamma_Y(X,-)\).

(d) If \(\mathcal F\) is flasque and \(U=X\setminus Y\), prove
\[
0\to\Gamma_Y(X,\mathcal F)\to\Gamma(X,\mathcal F)
\to\Gamma(U,\mathcal F|_U)\to0.
\]

(e) Construct the natural long exact sequence
\[
\begin{aligned}
0&\to H_Y^0(X,\mathcal F)\to H^0(X,\mathcal F)
\to H^0(U,\mathcal F|_U)\\
&\to H_Y^1(X,\mathcal F)\to H^1(X,\mathcal F)
\to H^1(U,\mathcal F|_U)\to\cdots.
\end{aligned}
\]

(f) If \(V\subseteq X\) is open and contains \(Y\), prove the
excision isomorphisms
\[
H_Y^i(X,\mathcal F)
\cong H_Y^i(V,\mathcal F|_V)
\]
functorially in \(\mathcal F\) and for all \(i\).
\end{problem}
\end{exerciseintro}

\begin{solution}
(a) The functor of global sections is left exact, and support can be
tested on germs.  Thus, for a morphism
\(\varphi:\mathcal F\to\mathcal G\),
\[
\begin{aligned}
\Gamma_Y(X,\ker\varphi)
&=\{s\in\Gamma(X,\mathcal F):
\varphi(s)=0,\ s|_{X\setminus Y}=0\}\\
&=\ker\bigl(
\Gamma_Y(X,\mathcal F)\to\Gamma_Y(X,\mathcal G)\bigr).
\end{aligned}
\]
This proves left exactness.

(b) Let \(s''\in\Gamma_Y(X,\mathcal F'')\).  A short exact sequence
with flasque kernel is exact on global sections, so choose
\(t\in\Gamma(X,\mathcal F)\) mapping to \(s''\).  On
\(U=X\setminus Y\), the restriction \(t|_U\) maps to zero and
therefore comes from some \(a'\in\Gamma(U,\mathcal F')\).
Flasqueness extends \(a'\) to
\(\widetilde a'\in\Gamma(X,\mathcal F')\).  Then
\[
t-\widetilde a'
\]
still maps to \(s''\) and vanishes on \(U\), so it is supported in
\(Y\).  This proves surjectivity; exactness at the other two terms
follows from (a).

(c) Embed a flasque sheaf \(\mathcal F\) into an injective sheaf
\(\mathcal I\), and let \(\mathcal G=\mathcal I/\mathcal F\).
Both \(\mathcal I\) and \(\mathcal F\) are flasque, and the usual
restriction-lifting argument shows that \(\mathcal G\) is flasque.
Part (b) makes
\[
\Gamma_Y(X,\mathcal I)\to\Gamma_Y(X,\mathcal G)
\]
surjective.  The derived-functor sequence therefore gives
\(H_Y^1(X,\mathcal F)=0\), and, for \(i\geq2\),
\[
H_Y^i(X,\mathcal F)\cong H_Y^{i-1}(X,\mathcal G).
\]
Repeating this argument and inducting on \(i\) proves
\(H_Y^i(X,\mathcal F)=0\) for every \(i>0\).

(d) The kernel of restriction
\(\Gamma(X,\mathcal F)\to\Gamma(U,\mathcal F|_U)\) consists exactly
of the sections supported in \(Y\).  The restriction map is
surjective because \(\mathcal F\) is flasque.  This gives the claimed
short exact sequence.

(e) Choose a flasque resolution
\(\mathcal F\to\mathcal I^\bullet\).  By (c), it computes cohomology
with supports; it also computes ordinary cohomology on \(X\) and,
after restriction, on \(U\).  Part (d), degree by degree, gives a
short exact sequence of complexes
\[
0\to\Gamma_Y(X,\mathcal I^\bullet)
\to\Gamma(X,\mathcal I^\bullet)
\to\Gamma(U,\mathcal I^\bullet|_U)\to0.
\]
The associated cohomology sequence is exactly the displayed long
exact sequence.

(f) Use the same flasque resolution.  Its restriction to \(V\) is a
flasque resolution of \(\mathcal F|_V\), and hence is acyclic for
sections supported in \(Y\) by (c).  Restriction induces a natural
isomorphism
\[
\Gamma_Y(X,\mathcal I^q)
\xrightarrow{\sim}
\Gamma_Y(V,\mathcal I^q|_V)
\tag{1}
\]
in every degree.  To see surjectivity in (1), glue a section on \(V\)
supported in \(Y\) with the zero section on \(X\setminus Y\); the two
agree on \(V\setminus Y\).  Injectivity is immediate.  Thus the two
complexes computing supported cohomology are isomorphic, and taking
cohomology proves excision and its functoriality.
\end{solution}

\begin{exerciseintro}
% record: H-III-02-004
\exerciseheading{Exercise III.2.4}
\addcontentsline{toc}{subsection}{Exercise III.2.4}

\begin{context}
For a closed subset \(Y\subseteq X\), write
\(H_Y^i(X,\mathcal F)\) for cohomology with support in \(Y\).
Flasque sheaves are acyclic for this functor, so flasque resolutions
may be used to compute it.
\end{context}

\begin{problem}
For closed subsets \(Y_1,Y_2\subseteq X\), construct the
Mayer--Vietoris sequence
\[
\begin{aligned}
\cdots\to&H_{Y_1\cap Y_2}^i(X,\mathcal F)
\to H_{Y_1}^i(X,\mathcal F)\oplus H_{Y_2}^i(X,\mathcal F)\\
\to&H_{Y_1\cup Y_2}^i(X,\mathcal F)
\to H_{Y_1\cap Y_2}^{i+1}(X,\mathcal F)\to\cdots.
\end{aligned}
\]
\end{problem}
\end{exerciseintro}

\begin{solution}
If \(\mathcal I\) is flasque, there is a short exact sequence
\[
\begin{aligned}
0\longrightarrow\Gamma_{Y_1\cap Y_2}(X,\mathcal I)
&\longrightarrow
\Gamma_{Y_1}(X,\mathcal I)\oplus
\Gamma_{Y_2}(X,\mathcal I)\\
&\longrightarrow
\Gamma_{Y_1\cup Y_2}(X,\mathcal I)
\longrightarrow0,
\end{aligned}
\tag{1}
\]
where the first map is \(t\mapsto(t,-t)\) and the second is addition.
Only surjectivity needs proof.  Let
\(s\) be supported in \(Y_1\cup Y_2\).  On
\[
W=(X\setminus Y_1)\cup(X\setminus Y_2)
=X\setminus(Y_1\cap Y_2)
\]
glue the zero section on \(X\setminus Y_1\) to
\(s|_{X\setminus Y_2}\).  They agree on
\(X\setminus(Y_1\cup Y_2)\).  Since \(\mathcal I\) is flasque, the
resulting section on \(W\) extends to a global section \(s_1\).
It vanishes off \(Y_1\), while \(s-s_1\) vanishes off \(Y_2\).
Thus \(s=s_1+(s-s_1)\), proving surjectivity in (1).

Choose a flasque resolution
\(\mathcal F\to\mathcal I^\bullet\).  Applying (1) in every degree
gives a short exact sequence of complexes.  Its long exact
cohomology sequence is the asserted Mayer--Vietoris sequence.
\end{solution}

\begin{exerciseintro}
% record: H-III-02-005
\exerciseheading{Exercise III.2.5}
\addcontentsline{toc}{subsection}{Exercise III.2.5}

\begin{context}
Let \(X\) be a Zariski space: it is noetherian, and every irreducible
closed subset has a unique generic point.  If \(P\in X\) is closed,
put
\[
X_P=\{Q\in X:P\in\overline{\{Q\}}\}
\]
with its subspace topology, and let \(j:X_P\to X\) be the inclusion.
For a sheaf \(\mathcal F\), write
\(\mathcal F_P=j^{-1}\mathcal F\).  The symbol \(P\) in
\(H_P^i\) denotes the closed support \(\{P\}\).
\end{context}

\begin{problem}
Prove, naturally for every sheaf \(\mathcal F\) and every \(i\), that
\[
H_P^i(X,\mathcal F)
\cong H_P^i(X_P,\mathcal F_P).
\]
\end{problem}
\end{exerciseintro}

\begin{solution}
We first record a topological description of inverse image in this
situation.  If \(W\subseteq X_P\) is open, then
\[
(j^{-1}\mathcal F)(W)
=\varinjlim_{\,W\subseteq V\cap X_P}\Gamma(V,\mathcal F),
\tag{1}
\]
where \(V\) runs through open subsets of \(X\).  In other words, the
usual presheaf defining \(j^{-1}\mathcal F\) is already a sheaf.

Here is the needed shrinking argument.  Every open subset of the
noetherian space \(X_P\) is quasi-compact, so gluing reduces to a
finite cover.  Suppose ambient representatives agree on an ambient
neighborhood \(N\) of an overlap.  The closed set \(X\setminus N\)
has finitely many irreducible components.  If \(\eta\) is the generic
point of one component and does not belong to an open
\(W\subseteq X_P\), then
\(\overline{\{\eta\}}\cap W=\varnothing\): otherwise \(\eta\) would
lie in \(X_P\), and stability of \(W\) under generization would put
\(\eta\) in \(W\).  For each component, remove it from one of the two
ambient neighborhoods whose corresponding opens do not contain its
generic point.  The shrunken ambient neighborhoods still contain the
given opens in \(X_P\), but now meet inside \(N\).  Repeating this for
the finitely many pairs lets the representatives glue.  The same
argument proves uniqueness, establishing (1).

If \(\mathcal I\) is flasque on \(X\), formula (1) shows that
\(j^{-1}\mathcal I\) is flasque on \(X_P\).  Indeed, a section on a
smaller open is represented on some ambient open \(V\); extend that
representative to \(X\) by flasqueness, then restrict it to the larger
open of \(X_P\).

We next compare the degree-zero functors.  Since an open subset of
\(X\) contains \(P\) exactly when it contains all of \(X_P\), (1)
gives
\[
\Gamma(X_P,\mathcal F_P)
=\varinjlim_{P\in V}\Gamma(V,\mathcal F)
=\mathcal F_P^{\mathrm{stalk}}.
\tag{2}
\]
We claim that the subgroups supported at \(P\) agree:
\[
\Gamma_P(X,\mathcal F)
\xrightarrow{\sim}
\Gamma_P(X_P,\mathcal F_P).
\tag{3}
\]
Injectivity is immediate.  For surjectivity, represent an element on
the right by \(s\in\Gamma(V,\mathcal F)\), where \(V\) contains
\(X_P\).  Let \(S\subseteq V\) be the support of \(s\).  The support
condition on \(X_P\) says
\[
S\cap X_P\subseteq\{P\}.
\]
The noetherian closed set \(S\) has finitely many irreducible
components.  If such a component contains \(P\), its generic point
lies in \(X_P\), so the component must be \(\{P\}\).  Remove all
other components from \(V\).  On the resulting neighborhood \(V'\)
of \(P\), the section is supported in \(\{P\}\).  It glues with zero
on \(X\setminus\{P\}\) to a global section on \(X\) supported at
\(P\).  This proves (3).

Finally, take an injective resolution
\(\mathcal F\to\mathcal I^\bullet\).  Inverse image is exact, so
\(\mathcal F_P\to j^{-1}\mathcal I^\bullet\) is a resolution.
Its terms are flasque by the preceding argument and hence acyclic for
sections supported at \(P\).  The isomorphisms (3), applied degree by
degree, identify the two complexes
\[
\Gamma_P(X,\mathcal I^\bullet)
\quad\text{and}\quad
\Gamma_P(X_P,j^{-1}\mathcal I^\bullet).
\]
Their cohomology groups are the two sides in the statement, which
proves the natural isomorphism.
\end{solution}

\begin{exerciseintro}
% record: H-III-02-006
\exerciseheading{Exercise III.2.6}
\addcontentsline{toc}{subsection}{Exercise III.2.6}

\begin{context}
For an open immersion \(u:U\to X\), let
\(\mathbb Z_U=u_!\mathbb Z\).  The sheaves \(\mathbb Z_U\), as \(U\)
runs through the open subsets of \(X\), form a family of generators
for sheaves of abelian groups.  We may use the corresponding Baer
criterion: a sheaf \(\mathcal I\) is injective if every map
\(\mathcal R\to\mathcal I\), for every subsheaf
\(\mathcal R\subseteq\mathbb Z_U\), extends across
\(\mathcal R\subseteq\mathbb Z_U\).
A space and its sobrification have the same lattice of open sets
and equivalent sheaf categories.  The sobrification consists of
nonempty irreducible closed subsets, with opens defined by meeting
an original open.  It is noetherian when the original space is,
and every irreducible closed subset has a generic point.
\end{context}

\begin{problem}
Let \(X\) be a noetherian topological space and
\((\mathcal I_\alpha)\) a filtered direct system of injective sheaves
of abelian groups.  Prove that
\[
\varinjlim_\alpha\mathcal I_\alpha
\]
is injective.
\end{problem}
\end{exerciseintro}

\begin{solution}
We use a finiteness consequence of the noetherian hypothesis: every
subsheaf of a finite sum of sheaves \(\mathbb Z_U\) is finitely
presented.
We can replace \(X\) by its sobrification; this leaves the sheaf
category and the assertion unchanged.  To prove the noetherian-object
assertion, let \(\mathcal R_1\subseteq\mathcal R_2\subseteq\cdots\)
be subsheaves of a fixed finite sum \(\mathcal G\) of the
\(\mathbb Z_U\).  On a dense open of each irreducible component,
disjoint from the other components, \(\mathcal G\) is a constant
sheaf of finite rank.  The union of the generic stalks of the
\(\mathcal R_i\) is a subgroup \(H\) of a finite-rank free abelian
group, so it is finitely generated and occurs at some finite stage.
Represent those generators on a smaller dense open.
On an irreducible open, sections of the ambient constant sheaf
inject into its generic stalk.  Therefore every later
\(\mathcal R_i\) equals the same constant subgroup \(H\) on this
open; generic generation suffices here because of the embedding
in the constant ambient sheaf.

The complement is a proper closed subset.  Restriction to it
sends \(\mathbb Z_U\) to extension by zero of the constant sheaf
on its intersection with \(U\), and is exact.
Noetherian induction on closed subsets therefore makes the
restricted chain stationary.  Combining the open and closed
parts proves that the original chain is stationary on \(X\).
Thus the finite sums \(\mathcal G\) are noetherian objects.
They are also finitely
presented: by adjunction,
\(\operatorname{Hom}(\mathbb Z_U,\mathcal G)=\Gamma(U,\mathcal G)\),
and sections over \(U\) commute with filtered colimits because all
opens of a noetherian space, and all their pairwise intersections,
are quasi-compact.  After choosing finitely many generators of a
subsheaf, its relation sheaf is again a subsheaf of a finite sum of
noetherian objects and is finitely generated.  This gives a finite
presentation and proves the assertion.  Equivalently,
for such a sheaf \(\mathcal R\),
\[
\operatorname{Hom}\left(
\mathcal R,\varinjlim_\alpha\mathcal G_\alpha\right)
\cong
\varinjlim_\alpha
\operatorname{Hom}(\mathcal R,\mathcal G_\alpha)
\tag{2}
\]
for every filtered system.

Put
\(\mathcal I=\varinjlim_\alpha\mathcal I_\alpha\).  To apply the Baer
criterion, take
\(\mathcal R\subseteq\mathbb Z_U\) and a morphism
\(f:\mathcal R\to\mathcal I\).  By this finiteness and (2), \(f\)
factors through
some stage:
\[
\mathcal R\xrightarrow{f_\alpha}\mathcal I_\alpha
\longrightarrow\mathcal I.
\]
The sheaf \(\mathcal I_\alpha\) is injective, so \(f_\alpha\) extends
to a morphism
\[
\mathbb Z_U\longrightarrow\mathcal I_\alpha.
\]
Composing with \(\mathcal I_\alpha\to\mathcal I\) extends \(f\).
The Baer criterion now shows that \(\mathcal I\) is injective.
\end{solution}

\begin{exerciseintro}
% record: H-III-02-007
\exerciseheading{Exercise III.2.7}
\addcontentsline{toc}{subsection}{Exercise III.2.7}

\begin{context}
Give \(S^1\) its usual topology.  Let \(\mathbb Z\) denote the
constant sheaf, let \(\mathcal R\) be the sheaf of germs of continuous
real-valued functions, and let \(\mathcal T\) be the sheaf of germs of
continuous \(S^1\)-valued functions.  Finite open covers of the circle
admit subordinate partitions of unity.
\end{context}

\begin{problem}
(a) Compute \(H^1(S^1,\mathbb Z)\) from the derived-functor
definition of sheaf cohomology.

(b) Prove that \(H^1(S^1,\mathcal R)=0\).
\end{problem}
\end{exerciseintro}

\begin{solution}
(b) Consider an exact
sequence
\[
0\longrightarrow\mathcal R\longrightarrow\mathcal I
\longrightarrow\mathcal Q\longrightarrow0
\tag{1}
\]
with \(\mathcal I\) injective.  Let
\(q\in\Gamma(S^1,\mathcal Q)\).  Choose a finite open cover
\((U_i)\) and lifts \(t_i\in\mathcal I(U_i)\).  On a double
intersection put
\[
v_{ij}=t_i-t_j\in\mathcal R(U_i\cap U_j).
\]
These sections satisfy
\[
v_{ij}-v_{ik}+v_{jk}=0.
\tag{2}
\]

Choose a partition of unity \((\varphi_k)\) subordinate to the cover.
On \(U_i\), set
\[
w_i=\sum_k\varphi_kv_{ik}.
\tag{3}
\]
The product in (3), initially defined on \(U_i\cap U_k\), extends by
zero across the complement of the support of \(\varphi_k\), so
\(w_i\) is a well-defined continuous function.  From (2) and
\(\sum_k\varphi_k=1\), one gets
\[
v_{ij}=w_i-w_j.
\]
Consequently the local sections \(t_i-w_i\) agree on overlaps and
glue to a global section of \(\mathcal I\) lifting \(q\).  Thus
\(\Gamma(S^1,\mathcal I)\to\Gamma(S^1,\mathcal Q)\) is surjective.
The long exact sequence of (1), together with injectivity of
\(\mathcal I\), now gives
\[
H^1(S^1,\mathcal R)=0.
\]
This proves (b).

(a) The local existence of a continuous argument gives an exact
exponential sequence of sheaves
\[
0\longrightarrow\mathbb Z
\longrightarrow\mathcal R
\xrightarrow{\,f\mapsto\exp(2\pi i f)\,}
\mathcal T\longrightarrow0.
\tag{4}
\]
Using (b), the long exact sequence of (4) identifies
\[
H^1(S^1,\mathbb Z)
\cong
\frac{C(S^1,S^1)}
{\{\exp(2\pi i f):f\in C(S^1,\mathbb R)\}}.
\tag{5}
\]
The winding-number map
\[
\deg:C(S^1,S^1)\longrightarrow\mathbb Z
\]
is surjective, since \(z\mapsto z^d\) has degree \(d\).  Its kernel
consists exactly of maps admitting a continuous lift
\(S^1\to\mathbb R\) through
\(\mathbb R\to S^1\), and these are precisely the maps in the
denominator of (5).  Therefore degree induces an isomorphism
\[
H^1(S^1,\mathbb Z)\cong\mathbb Z.
\]
\end{solution}

\begin{exerciseintro}
\corpussection{III.3 Cohomology of a Noetherian Affine Scheme}
\addcontentsline{toc}{section}{III.3 Cohomology of a Noetherian Affine Scheme}

% record: H-III-03-001
\exerciseheading{Exercise III.3.1}
\addcontentsline{toc}{subsection}{Exercise III.3.1}

\begin{context}
For a noetherian scheme \(X\), the reduction
\(i:X_{\mathrm{red}}\to X\) is defined by the nilradical ideal sheaf
\(\mathcal N\).  We may use Serre's criterion in the following form:
\(X\) is affine if and only if
\(H^1(X,\mathcal J)=0\) for every coherent ideal sheaf
\(\mathcal J\subseteq\mathcal O_X\).  Higher cohomology of a
quasi-coherent sheaf on a noetherian affine scheme vanishes.
\end{context}

\begin{problem}
Let \(X\) be a noetherian scheme.  Prove that \(X\) is affine if and
only if \(X_{\mathrm{red}}\) is affine.
\end{problem}
\end{exerciseintro}

\begin{solution}
If \(X=\operatorname{Spec}A\), then
\[
X_{\mathrm{red}}=\operatorname{Spec}(A/\sqrt{(0)}),
\]
so one implication is immediate.

Conversely, suppose that \(X_{\mathrm{red}}\) is affine.  The
nilradical \(\mathcal N\) is coherent and nilpotent.  Indeed, this is
true on every noetherian affine open, and a finite affine cover gives
one exponent \(r\) such that \(\mathcal N^r=0\).

Let \(\mathcal J\subseteq\mathcal O_X\) be any coherent ideal sheaf.
It has the finite filtration
\[
\mathcal J\supseteq\mathcal N\mathcal J\supseteq\cdots
\supseteq\mathcal N^r\mathcal J=0.
\tag{1}
\]
Each quotient
\[
\mathcal Q_q=
\mathcal N^q\mathcal J/\mathcal N^{q+1}\mathcal J
\]
is coherent and is annihilated by \(\mathcal N\).  Hence
\(\mathcal Q_q=i_*\mathcal G_q\) for a coherent sheaf
\(\mathcal G_q\) on \(X_{\mathrm{red}}\).  A closed immersion does
not change cohomology of a sheaf under direct image, and therefore
\[
H^1(X,\mathcal Q_q)
 =H^1(X_{\mathrm{red}},\mathcal G_q)=0.
\tag{2}
\]

Starting with the zero term of (1) and using the long exact sequence
for each successive quotient, (2) gives
\(H^1(X,\mathcal J)=0\).  Since this holds for every coherent ideal
sheaf, Serre's criterion shows that \(X\) is affine.
\end{solution}

\begin{exerciseintro}
% record: H-III-03-002
\exerciseheading{Exercise III.3.2}
\addcontentsline{toc}{subsection}{Exercise III.3.2}

\begin{context}
A noetherian scheme has finitely many irreducible components.  We may
use Serre's criterion: a noetherian scheme \(X\) is affine if and only
if \(H^1(X,\mathcal J)=0\) for every coherent ideal sheaf
\(\mathcal J\subseteq\mathcal O_X\).
\end{context}

\begin{problem}
Let \(X\) be a reduced noetherian scheme.  Prove that \(X\) is affine
if and only if every irreducible component of \(X\), with its reduced
induced scheme structure, is affine.
\end{problem}
\end{exerciseintro}

\begin{solution}
If \(X\) is affine, every irreducible component is the closed
subscheme defined by a minimal prime, and hence is affine.

For the converse, write the irreducible components as
\(C_1,\ldots,C_r\), and let \(\mathcal I_i\) be the ideal sheaf of
\(C_i\).  Since \(X\) is reduced,
\[
\mathcal I_1\cdots\mathcal I_r
 \subseteq\bigcap_{i=1}^r\mathcal I_i=0.
\tag{1}
\]
Put \(\mathcal K_0=\mathcal O_X\) and
\(\mathcal K_i=\mathcal I_1\cdots\mathcal I_i\).

Fix a coherent ideal sheaf \(\mathcal J\).  Equation (1) gives a
finite filtration
\[
0=\mathcal J\mathcal K_r\subseteq
\mathcal J\mathcal K_{r-1}\subseteq\cdots\subseteq
\mathcal J\mathcal K_0=\mathcal J.
\tag{2}
\]
The quotient
\[
\mathcal Q_i=
\mathcal J\mathcal K_{i-1}/\mathcal J\mathcal K_i
\]
is coherent and is annihilated by \(\mathcal I_i\).  If
\(j_i:C_i\to X\) denotes the closed immersion, then
\(\mathcal Q_i=(j_i)_*\mathcal G_i\) for a coherent sheaf
\(\mathcal G_i\) on \(C_i\).  By hypothesis \(C_i\) is affine, so
\[
H^1(X,\mathcal Q_i)=H^1(C_i,\mathcal G_i)=0.
\tag{3}
\]
The long exact sequences associated to (2), read successively from
left to right, now imply \(H^1(X,\mathcal J)=0\).  Serre's criterion
therefore proves that \(X\) is affine.
\end{solution}

\begin{exerciseintro}
% record: H-III-03-003
\exerciseheading{Exercise III.3.3}
\addcontentsline{toc}{subsection}{Exercise III.3.3}

\begin{context}
Let \(A\) be noetherian and \(\mathfrak a\subseteq A\) an ideal.  For
an \(A\)-module \(M\), set
\[
\Gamma_{\mathfrak a}(M)=
\{m\in M:\mathfrak a^n m=0\text{ for some }n\geq1\}.
\]
If \(X=\operatorname{Spec}A\) and \(Y=V(\mathfrak a)\), an injective
\(A\)-module \(I\) gives a flasque sheaf \(\widetilde I\) on \(X\).
\end{context}

\begin{problem}
(a) Prove that \(\Gamma_{\mathfrak a}\) is a left-exact functor, and
denote its right derived functors by \(H^i_{\mathfrak a}\).

(b) For every \(A\)-module \(M\), prove the natural isomorphism
\[
H^i_{\mathfrak a}(M)\cong H^i_Y(X,\widetilde M).
\]

(c) Prove that every local-cohomology module is
\(\mathfrak a\)-power torsion:
\[
\Gamma_{\mathfrak a}\bigl(H^i_{\mathfrak a}(M)\bigr)
=H^i_{\mathfrak a}(M).
\]
\end{problem}
\end{exerciseintro}

\begin{solution}
(a) The displayed rule defines a submodule and is functorial in
\(M\).  If \(u:M\to N\), then
\[
\Gamma_{\mathfrak a}(\ker u)
=\ker u\cap\Gamma_{\mathfrak a}(M)
=\ker\bigl(\Gamma_{\mathfrak a}(M)
 \longrightarrow\Gamma_{\mathfrak a}(N)\bigr).
\]
Thus \(\Gamma_{\mathfrak a}\) preserves kernels and is left exact.
Its right derived functors in the category of \(A\)-modules are, by
definition, the modules \(H^i_{\mathfrak a}(-)\).

(b) For every module \(L\) there is a natural equality
\[
\Gamma_{\mathfrak a}(L)=\Gamma_Y(X,\widetilde L).
\tag{1}
\]
Indeed, if \(\mathfrak a=(a_1,\ldots,a_s)\), the section defined by
\(\ell\in L\) vanishes on every \(D(a_j)\) if and only if suitable
powers of all the \(a_j\) kill \(\ell\); a sufficiently large power
of \(\mathfrak a\) then kills \(\ell\).

Choose an injective resolution \(M\to I^\bullet\).  Sheafification on
an affine scheme is exact, so
\(\widetilde M\to\widetilde{I^\bullet}\) is a resolution.  Its terms
are flasque, hence are acyclic for sections with support.  Applying
(1) degree by degree gives an isomorphism of complexes
\[
\Gamma_{\mathfrak a}(I^\bullet)
 \cong\Gamma_Y(X,\widetilde{I^\bullet}).
\]
Taking cohomology proves the asserted natural isomorphism.

(c) Each term \(\Gamma_{\mathfrak a}(I^q)\) is
\(\mathfrak a\)-power torsion.  The class of such modules is closed
under submodules and quotients: the exponent may depend on the
element, which is exactly what the definition allows.  Therefore
every cohomology module of the complex
\(\Gamma_{\mathfrak a}(I^\bullet)\), being a subquotient of one of
its terms, is again \(\mathfrak a\)-power torsion.  This is precisely
the displayed equality in (c).
\end{solution}

\begin{exerciseintro}
% record: H-III-03-004
\exerciseheading{Exercise III.3.4}
\addcontentsline{toc}{subsection}{Exercise III.3.4}

\begin{context}
Let \(A\) be noetherian, \(\mathfrak a\subseteq A\) an ideal, and
\(M\) an \(A\)-module.  The number
\(\operatorname{depth}_{\mathfrak a}M\) is the supremum of the
lengths of \(M\)-regular sequences whose terms lie in
\(\mathfrak a\).  For a finite module, the zero divisors on \(M\)
form the union of the finitely many primes in
\(\operatorname{Ass}_A(M)\).  Every \(H^i_{\mathfrak a}(M)\) is
\(\mathfrak a\)-power torsion.
For the finite modules in this exercise, require \(M\ne0\) and
\(\mathfrak aM\ne M\).  Regular sequences are required to leave a
nonzero quotient.  These conditions remove the degenerate empty-support
case from the depth comparison.
\end{context}

\begin{problem}
(a) Prove that \(\operatorname{depth}_{\mathfrak a}M\geq1\) implies
\(\Gamma_{\mathfrak a}(M)=0\).  Prove the converse when \(M\) is
finitely generated.

(b) If \(M\) is finitely generated, prove for every \(n\geq0\) that
the following conditions are equivalent:
\[
\operatorname{depth}_{\mathfrak a}M\geq n,
\qquad
H^i_{\mathfrak a}(M)=0\quad(0\leq i<n).
\]
\end{problem}
\end{exerciseintro}

\begin{solution}
(a) Choose an \(M\)-regular element \(x\in\mathfrak a\).  If
\(\mathfrak a^r m=0\), then \(x^r m=0\); injectivity of successive
multiplications by \(x\) gives \(m=0\).  Thus
\(\Gamma_{\mathfrak a}(M)=0\).

Conversely, suppose that \(M\) is finite and
\(\Gamma_{\mathfrak a}(M)=0\).  No associated prime of \(M\) can
contain \(\mathfrak a\): if
\(\mathfrak p=\operatorname{Ann}(m)\) contained \(\mathfrak a\),
then the nonzero element \(m\) would lie in
\(\Gamma_{\mathfrak a}(M)\).  Finite prime avoidance therefore gives
\[
x\in\mathfrak a\setminus
 \bigcup_{\mathfrak p\in\operatorname{Ass}_A(M)}\mathfrak p.
\]
Such an \(x\) is not a zero divisor on \(M\), so it is a regular
sequence of length one.

(b) We argue by induction on \(n\); the assertion for \(n=0\) is
empty, and the case \(n=1\) is (a), since
\(H^0_{\mathfrak a}(M)=\Gamma_{\mathfrak a}(M)\).

Assume the equivalence known with \(n-1\) in place of \(n\).  First
suppose that \(\operatorname{depth}_{\mathfrak a}M\geq n\).  Choose
a regular sequence beginning with \(x\in\mathfrak a\), and put
\(\overline M=M/xM\).  Then
\[
0\longrightarrow M\xrightarrow{\,x\,}M
 \longrightarrow\overline M\longrightarrow0
\tag{1}
\]
is exact, while both
\(\operatorname{depth}_{\mathfrak a}M\geq n-1\) and
\(\operatorname{depth}_{\mathfrak a}\overline M\geq n-1\).
The induction hypothesis gives
\[
H^i_{\mathfrak a}(M)=H^i_{\mathfrak a}(\overline M)=0
\quad(i<n-1).
\]
At the next degree, the long exact sequence of (1) shows that
multiplication by \(x\) on \(H^{n-1}_{\mathfrak a}(M)\) is
injective.  This module is \(\mathfrak a\)-power torsion, so every
one of its elements is killed by a power of \(x\).  An injective
nilpotent action is possible only on the zero module.  Hence
\(H^{n-1}_{\mathfrak a}(M)=0\) as well.

Conversely, suppose \(H^i_{\mathfrak a}(M)=0\) for \(i<n\).  When
\(n>0\), part (a) supplies an \(M\)-regular
\(x\in\mathfrak a\).  In the long exact sequence of (1), for
\(i<n-1\) the groups on either side of
\(H^i_{\mathfrak a}(\overline M)\) are
\(H^i_{\mathfrak a}(M)\) and
\(H^{i+1}_{\mathfrak a}(M)\), both zero.  Therefore
\[
H^i_{\mathfrak a}(\overline M)=0\quad(i<n-1).
\]
By induction, \(\overline M\) has an
\(\overline M\)-regular sequence of length \(n-1\) in
\(\mathfrak a\).  Prepending \(x\) gives an \(M\)-regular sequence
of length \(n\), completing the induction.
\end{solution}

\begin{exerciseintro}
% record: H-III-03-005
\exerciseheading{Exercise III.3.5}
\addcontentsline{toc}{subsection}{Exercise III.3.5}

\begin{context}
Let \(P\) be a closed point of a noetherian scheme \(X\).  We may use
the excision and localization isomorphisms for cohomology with support
at \(P\), as well as the exact sequence
\[
0\to H^0_P(W,\mathcal F)\to H^0(W,\mathcal F)
 \to H^0(W\setminus\{P\},\mathcal F)
 \to H^1_P(W,\mathcal F)\to H^1(W,\mathcal F).
\]
\end{context}

\begin{problem}
Prove that the following conditions are equivalent:

(i) \(\operatorname{depth}\mathcal O_{X,P}\geq2\).

(ii) For every open neighborhood \(W\) of \(P\), every regular
function on \(W\setminus\{P\}\) extends uniquely to a regular
function on \(W\).
\end{problem}
\end{exerciseintro}

\begin{solution}
Choose an affine neighborhood \(V=\operatorname{Spec}A\) of \(P\),
and let \(\mathfrak m\) be the corresponding maximal ideal.  By the
local-cohomology description of depth,
\[
\operatorname{depth}\mathcal O_{X,P}\geq2
\quad\Longleftrightarrow\quad
H^0_{\mathfrak mA_{\mathfrak m}}(A_{\mathfrak m})=
H^1_{\mathfrak mA_{\mathfrak m}}(A_{\mathfrak m})=0.
\tag{1}
\]
The comparison between algebraic local cohomology and cohomology with
supports, followed by localization to
\(\operatorname{Spec}A_{\mathfrak m}\), turns (1) into
\[
H^0_P(V,\mathcal O_V)=H^1_P(V,\mathcal O_V)=0.
\tag{2}
\]

Since \(V\) is affine, \(H^1(V,\mathcal O_V)=0\).  The support exact
sequence therefore shows that (2) is equivalent to the restriction
map
\[
\Gamma(V,\mathcal O_V)\longrightarrow
\Gamma(V\setminus\{P\},\mathcal O_V)
\tag{3}
\]
being an isomorphism.  Thus condition (i) is equivalent to (3) for
every sufficiently small affine neighborhood \(V\) of \(P\).

It remains only to pass from affine neighborhoods to arbitrary ones.
Let \(W\) be any neighborhood of \(P\), choose an affine
\(P\in V\subseteq W\), and take
\(s\in\Gamma(W\setminus\{P\},\mathcal O_X)\).  If (3) is an
isomorphism, \(s|_{V\setminus\{P\}}\) has a unique extension to
\(V\).  That extension and \(s\) agree on their overlap and hence
glue over the cover
\[
W=V\cup(W\setminus\{P\}).
\]
The same cover proves uniqueness.  Conversely, (ii) applies in
particular to every affine neighborhood \(V\), so (3) is an
isomorphism and (1) follows.  This proves the equivalence.
\end{solution}

\begin{exerciseintro}
% record: H-III-03-006
\exerciseheading{Exercise III.3.6}
\addcontentsline{toc}{subsection}{Exercise III.3.6}

\begin{context}
Let \(X\) be a noetherian scheme.  Choose a finite affine open cover
\(X=\bigcup_{i=1}^r U_i\), write \(U_i=\operatorname{Spec}A_i\), and
let \(j_i:U_i\to X\) be the inclusions.  In the standard construction,
a quasi-coherent sheaf \(\mathcal F\) is embedded into
\[
\mathcal G=\bigoplus_{i=1}^r(j_i)_*\widetilde I_i,
\]
where each \(I_i\) is an injective \(A_i\)-module containing the
module corresponding to \(\mathcal F|_{U_i}\).
\end{context}

\begin{problem}
(a) Prove that \(\mathcal G\) is injective in the category
\(\operatorname{Qco}(X)\).  Conclude that this category has enough
injectives.

(b) Prove that every injective object of \(\operatorname{Qco}(X)\)
is flasque as a sheaf.

(c) Prove that the derived functors of global sections computed in
\(\operatorname{Qco}(X)\) agree with ordinary sheaf cohomology.
\end{problem}
\end{exerciseintro}

\begin{solution}
(a) Because \(X\) is noetherian, each \(j_i\) is quasi-compact, so
\((j_i)_*\) carries quasi-coherent sheaves to quasi-coherent sheaves.
For \(\mathcal E\in\operatorname{Qco}(X)\), adjunction gives
\[
\operatorname{Hom}_X
 \bigl(\mathcal E,(j_i)_*\widetilde I_i\bigr)
\cong
\operatorname{Hom}_{U_i}
 \bigl(\mathcal E|_{U_i},\widetilde I_i\bigr).
\tag{1}
\]
Restriction to \(U_i\) is exact.  Under the equivalence
\(\operatorname{Qco}(U_i)\simeq\operatorname{Mod}(A_i)\), the sheaf
\(\widetilde I_i\) corresponds to the injective module \(I_i\).
Consequently the functor in (1) is exact, and
\((j_i)_*\widetilde I_i\) is injective in
\(\operatorname{Qco}(X)\).  The cover is finite, so its direct sum is
a finite biproduct of injectives and is injective.  Since the standard
construction embeds every quasi-coherent \(\mathcal F\) into such a
\(\mathcal G\), the category has enough injectives.

(b) Each \(\widetilde I_i\) is flasque on \(U_i\), and direct image
under an open immersion preserves flasqueness.  Hence the sheaf
\(\mathcal G\) above is flasque.  Apply the construction to an
injective object \(\mathcal I\in\operatorname{Qco}(X)\), obtaining
\[
0\longrightarrow\mathcal I\longrightarrow\mathcal G.
\]
Injectivity of \(\mathcal I\) splits this monomorphism inside
\(\operatorname{Qco}(X)\), so
\(\mathcal G\cong\mathcal I\oplus\mathcal K\).  A direct summand of
a flasque sheaf is flasque: projecting a lift in \(\mathcal G(W)\)
onto \(\mathcal I(W)\) lifts any section of \(\mathcal I(V)\) for
\(V\subseteq W\).  Thus \(\mathcal I\) is flasque.

(c) Resolve \(\mathcal F\in\operatorname{Qco}(X)\) by injectives in
that category.  By (b), every term of this resolution is flasque as
an ordinary \(\mathcal O_X\)-module.  Applying \(\Gamma(X,-)\) and
taking cohomology therefore computes both the derived functors in
\(\operatorname{Qco}(X)\) and ordinary sheaf cohomology.  The two
constructions consequently agree, naturally in \(\mathcal F\).
\end{solution}

\begin{exerciseintro}
% record: H-III-03-007
\exerciseheading{Exercise III.3.7}
\addcontentsline{toc}{subsection}{Exercise III.3.7}

\begin{context}
Let \(A\) be noetherian, \(X=\operatorname{Spec}A\),
\(\mathfrak a\subseteq A\), and
\(U=X\setminus V(\mathfrak a)\).  The transition map in the system
\(\operatorname{Hom}_A(\mathfrak a^n,M)\) is restriction along
\(\mathfrak a^{n+1}\subseteq\mathfrak a^n\).
\end{context}

\begin{problem}
(a) Prove Deligne's formula
\[
\Gamma(U,\widetilde M)
\cong\varinjlim_n\operatorname{Hom}_A(\mathfrak a^n,M)
\]
for every \(A\)-module \(M\).

(b) Use the formula to prove that \(\widetilde I\) is flasque whenever
\(I\) is an injective \(A\)-module.
\end{problem}
\end{exerciseintro}

\begin{solution}
(a) Since \(\widetilde{\mathfrak a^n}|_U=\mathcal O_U\), restriction
defines compatible maps
\[
\operatorname{Hom}_A(\mathfrak a^n,M)
\longrightarrow\Gamma(U,\widetilde M),
\]
and hence a natural map
\[
\Phi:\varinjlim_n\operatorname{Hom}_A(\mathfrak a^n,M)
\longrightarrow\Gamma(U,\widetilde M).
\tag{1}
\]

We first prove injectivity.  Suppose
\(f:\mathfrak a^n\to M\) restricts to zero on \(U\).  Its image
\(N\) is finite, and \(\widetilde N\) is supported on
\(V(\mathfrak a)\).  Thus \(\mathfrak a^rN=0\) for some \(r\).
It follows that the restriction of \(f\) to
\(\mathfrak a^{n+r}\) is zero, so the class of \(f\) vanishes in the
colimit.  Hence \(\Phi\) is injective.

For surjectivity, first suppose that \(M\) is finite.  Set
\[
T_j=\{m\in M:\mathfrak a^j m=0\},
\qquad T=\bigcup_jT_j.
\]
The ascending chain \((T_j)\) stabilizes, so
\(T=T_t\) for some \(t\).  Artin--Rees, applied to \(T\subseteq M\),
gives an integer \(c\) such that
\[
\mathfrak a^qM\cap T\subseteq\mathfrak a^{q-c}T
\quad(q\geq c).
\]
Choose \(q\geq c+t\) and put \(M'=\mathfrak a^qM\).  Then
\(M'\cap T=0\), so \(M'\) has no nonzero
\(\mathfrak a\)-power torsion, while
\(\widetilde{M'}|_U=\widetilde M|_U\).  We may therefore replace
\(M\) by \(M'\) and assume
\[
\Gamma_{\mathfrak a}(M)=0.
\tag{2}
\]

Write \(\mathfrak a=(f_1,\ldots,f_s)\), and take
\(z\in\Gamma(U,\widetilde M)\).  Condition (2) makes the natural map
\(M\to\Gamma(U,\widetilde M)\) injective.  On \(D(f_i)\), write
\[
z=\frac{m_i}{f_i^{e_i}}
\qquad(m_i\in M).
\]
The section \(m_i-f_i^{e_i}z\) on \(U\) vanishes on \(D(f_i)\).
The finite cover \(U=\bigcup_jD(f_j)\) then supplies an exponent
\(r_i\) for which
\[
f_i^{r_i}(m_i-f_i^{e_i}z)=0
\quad\text{on }U.
\]
Choose one \(L\geq1\) at least as large as every \(e_i+r_i\).
Then \(f_i^Lz\) belongs to the embedded copy of \(M\) for every
\(i\).  Every monomial of total degree \(sL\) in the \(f_i\) has
some exponent at least \(L\); hence
\[
\mathfrak a^{sL}z\subseteq M
\quad\text{inside }\Gamma(U,\widetilde M).
\tag{3}
\]
Multiplication by \(z\) in (3) defines an \(A\)-linear map
\[
\mathfrak a^{sL}\longrightarrow M,
\qquad b\longmapsto bz,
\]
whose image under \(\Phi\) is \(z\).  This proves surjectivity for
finite modules.

For arbitrary \(M\), write it as the filtered union of its finite
submodules.  The open set \(U\) is quasi-compact and
quasi-separated, so sections of quasi-coherent sheaves over \(U\)
commute with this filtered union.  Thus every
\(z\in\Gamma(U,\widetilde M)\) comes from
\(\Gamma(U,\widetilde N)\) for some finite submodule \(N\subseteq M\).
The finite case produces a map \(\mathfrak a^n\to N\); composing with
\(N\to M\) shows that \(z\) lies in the image of (1).  This completes
the proof of Deligne's formula.

(b) Let \(W'\subseteq W\subseteq X\) be open.  Since \(X\) is
noetherian, there are radical ideals \(\mathfrak b\subseteq\mathfrak a\)
such that
\[
W=X\setminus V(\mathfrak a),
\qquad W'=X\setminus V(\mathfrak b).
\]
By (a), a section of \(\widetilde I\) over \(W'\) is represented by
a map \(g:\mathfrak b^n\to I\).  The inclusion
\(\mathfrak b^n\subseteq\mathfrak a^n\) and injectivity of \(I\)
extend \(g\) to a map \(\mathfrak a^n\to I\).  The corresponding
section over \(W\) restricts to the given section.  Every restriction
map is therefore surjective, so \(\widetilde I\) is flasque.
\end{solution}

\begin{exerciseintro}
% record: H-III-03-008
\exerciseheading{Exercise III.3.8}
\addcontentsline{toc}{subsection}{Exercise III.3.8}

\begin{context}
Let
\[
A=k[x_0,x_1,x_2,\ldots]/(x_0^n x_n:n\geq1),
\]
and let \(I\) be an injective \(A\)-module containing \(A\) as a
submodule.  Equality of two fractions in \(I_{x_0}\) means that a
power of \(x_0\) kills their cross-multiplied difference.
\end{context}

\begin{problem}
Prove that the localization map
\[
I\longrightarrow I_{x_0}
\]
is not surjective.  This gives the stated failure, without a
noetherian hypothesis, of the flasqueness result for sheaves
associated to injective modules.
\end{problem}
\end{exerciseintro}

\begin{solution}
The element \(1/x_0\in I_{x_0}\) is well defined because
\(1\in A\subseteq I\).  Suppose it were the image of some
\(u\in I\).  Equality
\[
\frac{u}{1}=\frac{1}{x_0}
\quad\text{in }I_{x_0}
\]
would give an integer \(r\geq0\) such that
\[
x_0^r(x_0u-1)=0,
\qquad\text{or equivalently}\qquad
x_0^{r+1}u=x_0^r.
\tag{1}
\]
Multiply (1) by \(x_{r+1}\).  The defining relation
\(x_0^{r+1}x_{r+1}=0\) yields
\[
x_0^r x_{r+1}=x_0^{r+1}x_{r+1}u=0
\quad\text{in }I.
\tag{2}
\]
But \(x_0^r x_{r+1}\neq0\) in \(A\): the defining ideal is a
monomial ideal, and this monomial is not divisible by any generator
\(x_0^n x_n\).  Equation (2) contradicts the injectivity of
\(A\to I\).  Hence \(1/x_0\) is not in the image, and
\(I\to I_{x_0}\) is not surjective.

Indeed, this map is the restriction
\[
\Gamma(\operatorname{Spec}A,\widetilde I)
\longrightarrow
\Gamma(D(x_0),\widetilde I),
\]
so \(\widetilde I\) is not flasque.
\end{solution}

\begin{exerciseintro}
\corpussection{III.4 Cech Cohomology}
\addcontentsline{toc}{section}{III.4 Cech Cohomology}

% record: H-III-04-001
\exerciseheading{Exercise III.4.1}
\addcontentsline{toc}{subsection}{Exercise III.4.1}

\begin{context}
For a quasi-coherent sheaf on a noetherian separated scheme, an
affine open cover computes sheaf cohomology by its Čech complex.  The
inverse image of an affine open under an affine morphism is affine,
and the pushforward of a quasi-coherent sheaf under such a morphism is
quasi-coherent.
\end{context}

\begin{problem}
Let \(f:X\to Y\) be an affine morphism of noetherian separated
schemes.  Prove that for every quasi-coherent sheaf \(\mathcal F\) on
\(X\) there are natural isomorphisms
\[
H^i(X,\mathcal F)\cong H^i(Y,f_*\mathcal F)
\qquad(i\geq0).
\]
\end{problem}
\end{exerciseintro}

\begin{solution}
Choose an affine open cover \(\mathfrak U=(U_j)\) of \(Y\).  Because
\(Y\) is separated, every finite intersection
\(U_{j_0\ldots j_p}\) is affine.  Since \(f\) is affine, every
\(f^{-1}(U_{j_0\ldots j_p})\) is affine as well.  Thus both covers are
acyclic for the relevant quasi-coherent sheaves.

By the definition of direct image, for every finite intersection
there is a natural identity
\[
\Gamma(U_{j_0\ldots j_p},f_*\mathcal F)
=
\Gamma(f^{-1}(U_{j_0\ldots j_p}),\mathcal F).
\tag{1}
\]
The identities (1) commute with all restriction maps and therefore
give an isomorphism of augmented Čech complexes
\[
C^\bullet(\mathfrak U,f_*\mathcal F)
\cong
C^\bullet(f^{-1}\mathfrak U,\mathcal F).
\]
The Čech comparison theorem identifies the cohomology of the two
complexes with the two sheaf-cohomology groups in the statement.
Taking cohomology of this natural isomorphism proves the result in
every degree.
\end{solution}

\begin{exerciseintro}
% record: H-III-04-002
\exerciseheading{Exercise III.4.2}
\addcontentsline{toc}{subsection}{Exercise III.4.2}

\begin{context}
Finite morphisms are affine, closed, and stable under base change.
On a noetherian scheme, every quasi-coherent sheaf embeds in a flasque
quasi-coherent sheaf.  We may use generic freeness, the affine
morphism equivalence for quasi-coherent modules, and Serre's
cohomological criterion for affineness.
In part (b), retain the integral hypothesis of (a) when referring to the generic point.
\end{context}

\begin{problem}
Let \(f:X\to Y\) be a finite surjective morphism of noetherian
separated schemes, with \(X\) affine.

(a) If \(X\) and \(Y\) are integral, construct a coherent sheaf
\(\mathcal M\) on \(X\) and a morphism
\[
\alpha:\mathcal O_Y^r\longrightarrow f_*\mathcal M
\]
which is an isomorphism at the generic point of \(Y\).

(b) For every coherent \(\mathcal F\) on \(Y\), construct a coherent
\(\mathcal G\) on \(X\) and a morphism
\[
\beta:f_*\mathcal G\longrightarrow\mathcal F^r
\]
which is an isomorphism at the generic point.

(c) Prove Chevalley's theorem: \(Y\) is affine.
\end{problem}
\end{exerciseintro}

\begin{solution}
(a) Let \(V=\operatorname{Spec}B\subseteq Y\) be a nonempty affine
open and write
\(f^{-1}(V)=\operatorname{Spec}A\).  The map \(B\to A\) is finite and
injective.  After shrinking \(V\), generic freeness permits us to
assume that \(A\) is a free \(B\)-module with basis
\(a_1,\ldots,a_r\).

Let \(j:f^{-1}(V)\to X\) be the open immersion.  The sheaf
\(j_*\mathcal O_{f^{-1}(V)}\) is quasi-coherent; embed it into a
flasque quasi-coherent sheaf \(\mathcal I\).  Denote by
\(s\in\Gamma(f^{-1}(V),\mathcal I)\) the image of \(1\).  Flasqueness
allows each \(a_i s\) to extend to a section
\(\widetilde s_i\in\Gamma(X,\mathcal I)\).  Let \(\mathcal M\) be
the \(\mathcal O_X\)-submodule generated by these finitely many
sections.  It is coherent because \(X\) is noetherian.

On \(f^{-1}(V)\), the \(a_i\) generate \(A\), so
\(\mathcal M|_{f^{-1}(V)}=\mathcal O_{f^{-1}(V)}s\), a copy of the
structure sheaf.  The global sections \(\widetilde s_i\) define
\[
\alpha:\mathcal O_Y^r\longrightarrow f_*\mathcal M.
\]
Over \(V\), this is the map of \(B\)-modules
\[
B^r\longrightarrow A,
\qquad(b_1,\ldots,b_r)\longmapsto\sum_i b_i a_i,
\]
and hence is an isomorphism.  In particular it is an isomorphism at
the generic point.

(b) Apply the sheaf-Hom functor \(\mathcal Hom_Y(-,\mathcal F)\) to
\(\alpha\).  Precomposition gives
\[
\beta:\mathcal H:=
\mathcal Hom_Y(f_*\mathcal M,\mathcal F)
\longrightarrow\mathcal F^r.
\tag{1}
\]
The sheaf \(\mathcal H\) is coherent, and (1) is an isomorphism at
the generic point.  It also has a natural
\(f_*\mathcal O_X\)-module structure:
\[
(a\varphi)(m)=\varphi(am).
\]
The affine-morphism equivalence therefore gives a quasi-coherent
\(\mathcal O_X\)-module \(\mathcal G\) with
\(f_*\mathcal G\cong\mathcal H\).  It is coherent: finitely many
\(\mathcal O_Y\)-generators of \(\mathcal H\) also generate it as an
\(f_*\mathcal O_X\)-module, and finiteness can be checked on affine
opens.  With this identification, (1) is the required morphism.

(c) We first make two reductions.  Passing to reductions preserves
finiteness and surjectivity, \(X_{\mathrm{red}}\) is affine, and a
noetherian scheme is affine exactly when its reduction is affine.
Thus we may suppose that \(X\) and \(Y\) are reduced.  For each
irreducible component \(Y_j\), some irreducible component \(X_i\)
maps onto \(Y_j\): apply surjectivity to the generic point of
\(Y_j\), and use that a finite map is closed.  The induced map
\(X_i\to Y_j\) is finite and surjective, and \(X_i\) is affine.  If
the theorem is known in the integral case, every \(Y_j\) is affine;
the component criterion for a reduced noetherian scheme then makes
\(Y\) affine.  It remains to treat the integral case.

We use noetherian induction on the closed subsets of the integral
scheme \(Y\).  Assume the theorem has been proved over every proper
closed subset of \(Y\).  Let
\(\mathcal J\subseteq\mathcal O_Y\) be a coherent ideal sheaf.  Part
(b) supplies
\[
\beta:f_*\mathcal G\longrightarrow\mathcal J^r
\]
which is an isomorphism at the generic point.  Put
\(\mathcal K=\ker\beta\), \(\mathcal Q=\operatorname{im}\beta\),
and \(\mathcal C=\operatorname{coker}\beta\).  The coherent sheaves
\(\mathcal K\) and \(\mathcal C\) have proper closed support.

Give either support a closed scheme structure annihilating the
corresponding sheaf.  Its inverse image in \(X\) is affine and maps
to it finitely and surjectively.  The inductive hypothesis, together
with the reduction criterion, says that each such support scheme is
affine.  Consequently
\[
H^i(Y,\mathcal K)=H^i(Y,\mathcal C)=0
\qquad(i>0).
\tag{2}
\]
Moreover, the result for affine morphisms and the affineness of \(X\)
give
\[
H^i(Y,f_*\mathcal G)=H^i(X,\mathcal G)=0
\qquad(i>0).
\tag{3}
\]
Apply cohomology to
\[
0\to\mathcal K\to f_*\mathcal G\to\mathcal Q\to0,
\qquad
0\to\mathcal Q\to\mathcal J^r\to\mathcal C\to0.
\]
Equations (2) and (3) first give \(H^1(Y,\mathcal Q)=0\), and then
\(H^1(Y,\mathcal J^r)=0\).  Since \(r>0\), this implies
\(H^1(Y,\mathcal J)=0\).  Serre's criterion now shows that \(Y\) is
affine, completing the induction and the proof.
\end{solution}

\begin{exerciseintro}
% record: H-III-04-003
\exerciseheading{Exercise III.4.3}
\addcontentsline{toc}{subsection}{Exercise III.4.3}

\begin{context}
Let \(X=\mathbb A_k^2=\operatorname{Spec}k[x,y]\) and
\(U=X\setminus\{(0,0)\}\).  The cover
\(U=D(x)\cup D(y)\) is an affine cover with affine intersection, so
its Čech complex computes the cohomology of \(\mathcal O_U\).
\end{context}

\begin{problem}
Compute \(H^1(U,\mathcal O_U)\) and prove that it is the
infinite-dimensional \(k\)-vector space with basis
\[
\{x^i y^j:i<0,\ j<0\}.
\]
\end{problem}
\end{exerciseintro}

\begin{solution}
The degree-zero and degree-one terms of the Čech complex are
\[
k[x,x^{-1},y]\oplus k[x,y,y^{-1}]
\longrightarrow k[x,x^{-1},y,y^{-1}],
\tag{1}
\]
where the map is the difference of the two restrictions.  Therefore
\[
H^1(U,\mathcal O_U)
\cong
\frac{k[x,x^{-1},y,y^{-1}]}
{k[x,x^{-1},y]+k[x,y,y^{-1}]}.
\tag{2}
\]

The Laurent monomials \(x^i y^j\), \((i,j)\in\mathbb Z^2\), form a
\(k\)-basis of the numerator.  The first summand in the denominator
is spanned by those with \(j\geq0\), and the second by those with
\(i\geq0\).  Thus precisely the monomials with both exponents
negative survive in (2), and their residue classes remain linearly
independent.  Hence
\[
H^1(U,\mathcal O_U)
=\bigoplus_{i<0,\,j<0}k\,x^i y^j,
\]
which is infinite-dimensional.
\end{solution}

\begin{exerciseintro}
% record: H-III-04-004
\exerciseheading{Exercise III.4.4}
\addcontentsline{toc}{subsection}{Exercise III.4.4}

\begin{context}
Let \(\mathfrak U=(U_i)_{i\in I}\) be an open cover of a topological
space \(X\), and let \(\mathcal F\) be a sheaf of abelian groups.  We
use alternating Čech cochains; thus a cochain is extended to all
ordered index tuples by the sign rule and is zero when two indices
coincide.
\end{context}

\begin{problem}
(a) Construct the maps on Čech cohomology induced by a refinement and
the direct limit over all open covers.

(b) Prove that the comparison maps from Čech cohomology to derived
sheaf cohomology are compatible with refinements.

(c) Prove that the resulting natural map
\[
\varinjlim_{\mathfrak U}\check H^1(\mathfrak U,\mathcal F)
\longrightarrow H^1(X,\mathcal F)
\]
is an isomorphism.
\end{problem}
\end{exerciseintro}

\begin{solution}
(a) Let \(\mathfrak V=(V_j)_{j\in J}\) refine \(\mathfrak U\), and
choose a refinement function \(\lambda:J\to I\) with
\(V_j\subseteq U_{\lambda(j)}\).  Define
\[
(\lambda^*c)_{j_0\ldots j_p}
=c_{\lambda(j_0)\ldots\lambda(j_p)}
 \big|_{V_{j_0}\cap\cdots\cap V_{j_p}}.
\tag{1}
\]
The alternating convention makes (1) meaningful even when the
\(\lambda(j_q)\) are repeated or out of order.  Direct substitution
in the alternating differential shows that
\(d\lambda^*=\lambda^*d\), so (1) induces maps
\[
\lambda^p:\check H^p(\mathfrak U,\mathcal F)
\longrightarrow\check H^p(\mathfrak V,\mathcal F).
\]

These maps do not depend on the choice of refinement function.  In
fact, for two choices \(\lambda,\mu:J\to I\), the standard prism
operator
\[
(hc)_{j_0\ldots j_{p-1}}
=\sum_{q=0}^{p-1}(-1)^q
c_{\lambda(j_0)\ldots\lambda(j_q)
  \mu(j_q)\ldots\mu(j_{p-1})}
 \big|_{V_{j_0}\cap\cdots\cap V_{j_{p-1}}}
\]
satisfies \(dh+hd=\mu^*-\lambda^*\).  Thus the induced maps on
cohomology agree.  Any two covers have the common refinement
\((U_i\cap V_j)_{i,j}\); consequently these maps define the filtered
direct limit in the statement.

(b) Regard the augmented Čech complex as a resolution of
\(\mathcal F\) by sheaves and compare it with an injective resolution
\(\mathcal F\to\mathcal I^\bullet\).  A refinement gives a morphism
of augmented Čech resolutions which is the identity on
\(\mathcal F\).  By the comparison theorem for resolutions, its
composite with a lift to \(\mathcal I^\bullet\) is homotopic to the
lift used for the original cover.  After taking global sections and
cohomology, this says precisely that the comparison map for
\(\mathfrak V\), composed with \(\lambda^p\), is the comparison map
for \(\mathfrak U\).  The comparison maps therefore induce a map from
the direct limit.

(c) Choose a short exact sequence
\[
0\longrightarrow\mathcal F\longrightarrow\mathcal I
\longrightarrow\mathcal Q\longrightarrow0
\tag{2}
\]
with \(\mathcal I\) injective, hence flasque.  The associated long
exact sequence identifies
\[
H^1(X,\mathcal F)
\cong
\operatorname{coker}\bigl(
\Gamma(X,\mathcal I)\to\Gamma(X,\mathcal Q)\bigr).
\tag{3}
\]

Let a class in (3) be represented by
\(q\in\Gamma(X,\mathcal Q)\).  Since
\(\mathcal I\to\mathcal Q\) is an epimorphism of sheaves, there is an
open cover \((U_i)\) and lifts \(t_i\in\mathcal I(U_i)\).  On an
overlap,
\[
c_{ij}=t_j-t_i\in\mathcal F(U_i\cap U_j)
\]
is a Čech \(1\)-cocycle.  Its comparison image is the connecting
class of \(q\), so every element of \(H^1(X,\mathcal F)\) is reached
after choosing a suitable cover.  The direct-limit map is surjective.

It is also injective, in fact already on each cover.  Suppose a Čech
cocycle \(c\) for \(\mathfrak U\) maps to zero in
\(H^1(X,\mathcal F)\).  Viewed as an \(\mathcal I\)-valued cocycle,
it is a coboundary because a flasque sheaf has zero positive Čech
cohomology.  Thus
\(c_{ij}=t_j-t_i\) for sections \(t_i\in\mathcal I(U_i)\).  Their
images in \(\mathcal Q\) glue to a global section \(q\).  The
comparison image of \(c\) is the class of \(q\) in (3).  Its
vanishing supplies \(t\in\Gamma(X,\mathcal I)\) mapping to \(q\).
Then \(t_i-t|_{U_i}\) lies in \(\mathcal F(U_i)\), and these sections
exhibit \(c\) as an \(\mathcal F\)-valued coboundary.  Hence the
direct-limit map is injective and therefore an isomorphism.
\end{solution}

\begin{exerciseintro}
% record: H-III-04-005
\exerciseheading{Exercise III.4.5}
\addcontentsline{toc}{subsection}{Exercise III.4.5}

\begin{context}
For a ringed space \((X,\mathcal O_X)\), let \(\mathcal O_X^*\) be
the sheaf of multiplicative units and let \(\operatorname{Pic}X\) be
the group of isomorphism classes of invertible sheaves, with tensor
product as its operation.  First sheaf cohomology is the direct limit
of first Čech cohomology over open covers.
\end{context}

\begin{problem}
Prove the natural group isomorphism
\[
\operatorname{Pic}X\cong H^1(X,\mathcal O_X^*).
\]
\end{problem}
\end{exerciseintro}

\begin{solution}
Fix an open cover \(\mathfrak U=(U_i)\).  If an invertible sheaf
\(\mathcal L\) is trivial on every \(U_i\), choose generators
\(e_i\) of \(\mathcal L|_{U_i}\).  On overlaps there are unique
units \(g_{ij}\in\mathcal O_X^*(U_i\cap U_j)\) such that
\[
e_j=g_{ij}e_i.
\]
On triple intersections these satisfy
\(g_{ij}g_{jk}=g_{ik}\), so they form a Čech \(1\)-cocycle.  Replacing
\(e_i\) by \(u_i e_i\) multiplies \((g_{ij})\) by a Čech
coboundary.  Thus \(\mathcal L\), together with its local
triviality over this cover, determines a class in
\(\check H^1(\mathfrak U,\mathcal O_X^*)\).

Conversely, a cocycle \((g_{ij})\) glues the sheaves
\(\mathcal O_{U_i}\) by the isomorphisms given by multiplication by
\(g_{ij}\).  The cocycle identity is exactly the compatibility
needed for gluing, and the resulting sheaf is invertible.  A
coboundary changes the local generators and hence gives an
isomorphic sheaf.  These two constructions are inverse.

Every invertible sheaf is trivial on some open cover.  Passing to the
direct limit over refinements therefore gives a bijection
\[
\operatorname{Pic}X
\longrightarrow
\varinjlim_{\mathfrak U}
\check H^1(\mathfrak U,\mathcal O_X^*)
\cong H^1(X,\mathcal O_X^*).
\]
Under this bijection, tensoring line bundles multiplies their
transition functions.  Hence the bijection is an isomorphism of
groups.
\end{solution}

\begin{exerciseintro}
% record: H-III-04-006
\exerciseheading{Exercise III.4.6}
\addcontentsline{toc}{subsection}{Exercise III.4.6}

\begin{context}
Let \((X,\mathcal O_X)\) be a ringed space, let
\(\mathcal I\subseteq\mathcal O_X\) be an ideal sheaf with
\(\mathcal I^2=0\), and put
\(X_0=(X,\mathcal O_X/\mathcal I)\).  We identify Picard groups with
first cohomology of the corresponding sheaves of units.
\end{context}

\begin{problem}
Prove that \(a\mapsto1+a\) fits into an exact sequence of sheaves of
abelian groups
\[
0\longrightarrow\mathcal I\longrightarrow\mathcal O_X^*
\longrightarrow\mathcal O_{X_0}^*\longrightarrow0.
\]
Deduce the exact sequence
\[
\cdots\longrightarrow H^1(X,\mathcal I)
\longrightarrow\operatorname{Pic}X
\longrightarrow\operatorname{Pic}X_0
\longrightarrow H^2(X,\mathcal I)\longrightarrow\cdots.
\]
\end{problem}
\end{exerciseintro}

\begin{solution}
Because \(\mathcal I^2=0\),
\[
(1+a)(1+b)=1+a+b
\qquad(a,b\in\mathcal I),
\]
so \(a\mapsto1+a\) is a homomorphism from the additive sheaf
\(\mathcal I\) to the multiplicative sheaf \(\mathcal O_X^*\).

Exactness can be checked on stalks.  Let \(R=\mathcal O_{X,x}\) and
\(I=\mathcal I_x\).  The kernel of
\(R^*\to(R/I)^*\) consists exactly of the elements \(1+a\),
\(a\in I\), and this representation is unique.  The map on units is
surjective as well.  Indeed, if the class of \(r\in R\) is invertible
modulo \(I\), choose \(s\in R\) with
\(rs=1+c\), \(c\in I\).  Since
\((1+c)^{-1}=1-c\), the element \(r\) is a unit in \(R\).  This proves
the short exact sequence of sheaves.

Its long exact cohomology sequence contains
\[
H^1(X,\mathcal I)\longrightarrow H^1(X,\mathcal O_X^*)
\longrightarrow H^1(X,\mathcal O_{X_0}^*)
\longrightarrow H^2(X,\mathcal I).
\]
The identifications
\(H^1(X,\mathcal O_X^*)\cong\operatorname{Pic}X\) and
\(H^1(X,\mathcal O_{X_0}^*)\cong\operatorname{Pic}X_0\) turn this
into the required exact sequence.
\end{solution}

\begin{exerciseintro}
% record: H-III-04-007
\exerciseheading{Exercise III.4.7}
\addcontentsline{toc}{subsection}{Exercise III.4.7}

\begin{context}
Let \(X\subseteq\mathbb P_k^2\) be the hypersurface defined by a
homogeneous polynomial \(f(x_0,x_1,x_2)\) of degree \(d>0\), not
necessarily irreducible.  Assume \((1:0:0)\notin X\).  Standard
affine opens and all their intersections are acyclic for
\(\mathcal O_X\), so the associated Čech complex computes
cohomology.
\end{context}

\begin{problem}
Show that
\[
U=X\cap D_+(x_1),\qquad V=X\cap D_+(x_2)
\]
are affine and cover \(X\).  Compute their Čech complex explicitly
and prove
\[
\dim_k H^0(X,\mathcal O_X)=1,\qquad
\dim_k H^1(X,\mathcal O_X)=\frac{(d-1)(d-2)}2.
\]
\end{problem}
\end{exerciseintro}

\begin{solution}
The two sets cover \(X\), because a point outside both would have
\(x_1=x_2=0\) and hence would be \((1:0:0)\).  Each is closed in the
corresponding standard affine open of \(\mathbb P_k^2\), and is
therefore affine.

Put
\[
S=k[x_0,x_1,x_2]/(f).
\]
The rings in the Čech complex are
\[
\Gamma(U,\mathcal O_X)=S_{(x_1)},\qquad
\Gamma(V,\mathcal O_X)=S_{(x_2)},\qquad
\Gamma(U\cap V,\mathcal O_X)=S_{(x_1x_2)},
\tag{1}
\]
where parentheses denote the degree-zero part of a homogeneous
localization.

The hypothesis on \((1:0:0)\) says that the coefficient of \(x_0^d\)
in \(f\) is nonzero.  After scaling \(f\), it is monic as a polynomial
in \(x_0\).  Polynomial division therefore gives a unique normal form
in the last ring of (1) as a finite \(k\)-linear combination of
\[
x_0^a x_1^b x_2^c,
\qquad
0\leq a<d,\quad b,c\in\mathbb Z,\quad a+b+c=0.
\tag{2}
\]
The image of \(S_{(x_1)}\) is spanned by the normal monomials (2) with
\(c\geq0\), while the image of \(S_{(x_2)}\) is spanned by those with
\(b\geq0\).  Indeed, division by the monic polynomial \(f\) introduces
no negative power of a variable that was not already inverted.

The Čech differential is
\[
S_{(x_1)}\oplus S_{(x_2)}
\longrightarrow S_{(x_1x_2)},\qquad
(g,h)\longmapsto h-g.
\tag{3}
\]
For \(H^0\), the intersection of the two image subspaces has
\(b,c\geq0\).  Together with \(a\geq0\) and \(a+b+c=0\), this forces
\(a=b=c=0\).  Hence
\[
H^0(X,\mathcal O_X)=k.
\]

For \(H^1\), quotienting the target of (3) by the two images leaves
exactly the normal monomials for which \(b<0\) and \(c<0\).  Write
\(p=-b\) and \(q=-c\).  Then \(p,q\geq1\) and \(a=p+q\), with
\(a<d\).  Thus
\[
H^1(X,\mathcal O_X)
\cong
\bigoplus_{\substack{p,q\geq1\\ p+q<d}}
k\,x_0^{p+q}x_1^{-p}x_2^{-q}.
\]
For a fixed \(a=2,\ldots,d-1\), there are \(a-1\) pairs
\((p,q)\) with \(p+q=a\).  Consequently
\[
\dim_kH^1(X,\mathcal O_X)
=\sum_{a=2}^{d-1}(a-1)
=\frac{(d-1)(d-2)}2.
\]
\end{solution}

\begin{exerciseintro}
% record: H-III-04-008
\exerciseheading{Exercise III.4.8}
\addcontentsline{toc}{subsection}{Exercise III.4.8}

\begin{context}
For a noetherian separated scheme \(X\), define
\(\operatorname{cd}(X)\) to be the least \(n\) such that
\(H^i(X,\mathcal F)=0\) for every quasi-coherent \(\mathcal F\) and
every \(i>n\).  Cohomology on a noetherian space commutes with
filtered direct limits, and every quasi-coherent sheaf is the
filtered union of its coherent subsheaves.
For \(r=0\) in (e), an empty complement has zero cohomology in every degree.
\end{context}

\begin{problem}
(a) Show that coherent sheaves suffice in the definition of
\(\operatorname{cd}(X)\).

(b) If \(X\) is quasi-projective over a field, show that locally free
coherent sheaves suffice.

(c) If \(X\) is covered by \(r+1\) affine opens, prove
\(\operatorname{cd}(X)\leq r\).

(d) Prove that a quasi-projective scheme of dimension \(r\) over a
field can be covered by \(r+1\) affine opens.  Deduce
\(\operatorname{cd}(X)\leq\dim X\).

(e) If \(Y\subseteq\mathbb P_k^n\) is a set-theoretic complete
intersection of codimension \(r\), prove
\(\operatorname{cd}(\mathbb P_k^n\setminus Y)\leq r-1\).
\end{problem}
\end{exerciseintro}

\begin{solution}
(a) Write a quasi-coherent sheaf as a filtered colimit
\(\mathcal F=\varinjlim_\lambda\mathcal F_\lambda\) of coherent
subsheaves.  Then
\[
H^i(X,\mathcal F)
\cong\varinjlim_\lambda H^i(X,\mathcal F_\lambda).
\]
Thus vanishing above a fixed degree for all coherent sheaves implies
the same vanishing for every quasi-coherent sheaf.  The converse is
immediate.

(b) Suppose \(H^i(X,\mathcal E)=0\) for \(i>n\) and every locally
free coherent \(\mathcal E\).  The ample-sheaf theorem for a
quasi-projective scheme gives, for each coherent \(\mathcal F_j\), a
surjection from a finite direct sum of twists,
\[
0\longrightarrow\mathcal F_{j+1}\longrightarrow
\mathcal E_j\longrightarrow\mathcal F_j\longrightarrow0,
\qquad \mathcal F_0=\mathcal F,
\tag{1}
\]
with \(\mathcal E_j\) locally free and \(\mathcal F_{j+1}\) coherent.
For fixed \(i>n\), the long exact sequences of (1) give dimension
shifting
\[
H^i(X,\mathcal F)
\cong H^{i+t}(X,\mathcal F_t)
\tag{2}
\]
for every \(t\geq0\).  Choose \(t\) so that
\(i+t>\dim X\).  Grothendieck vanishing makes the right side of (2)
zero.  Hence all coherent sheaves have the required vanishing, and
part (a) finishes the claim.

(c) Let \(\mathfrak U=(U_0,\ldots,U_r)\) be such a cover.  Since
\(X\) is separated, every finite intersection of the \(U_i\) is
affine.  For a quasi-coherent \(\mathcal F\), this is an acyclic
cover, so Čech cohomology computes \(H^*(X,\mathcal F)\).  Its Čech
complex has no terms in degrees greater than \(r\).  Therefore
\[
H^i(X,\mathcal F)=0\quad(i>r),
\]
and \(\operatorname{cd}(X)\leq r\).

(d) Realize \(X\) as a closed subscheme of an open
\(W\subseteq\mathbb P_k^N\), and let
\(Z=\mathbb P_k^N\setminus W\).  We inductively choose homogeneous
polynomials \(f_0,\ldots,f_r\) vanishing on \(Z\).  Start with
\(T_0=X\).  If \(T_j\neq\varnothing\), none of its finitely many
generic points lies in \(Z\).  Homogeneous prime avoidance, after
raising candidates to a common degree, gives
\[
f_j\in I(Z)
\]
which vanishes at none of those generic points.  Set
\[
T_{j+1}=T_j\cap V_+(f_j).
\]
Every irreducible component drops in dimension, so
\(\dim T_{j+1}<\dim T_j\).  Since \(\dim X=r\), after \(r+1\) choices
we have \(T_{r+1}=\varnothing\).

Now \(D_+(f_j)\subseteq W\), because \(f_j\) vanishes on \(Z\), and
\(D_+(f_j)\) is affine.  Hence
\[
X_j=X\cap D_+(f_j)
\]
is closed in an affine scheme and is affine.  The equality
\(T_{r+1}=\varnothing\) says that the \(X_j\) cover \(X\).  Part (c)
then gives \(\operatorname{cd}(X)\leq r=\dim X\).

(e) Write \(Y\) set-theoretically as
\[
Y=V_+(f_1)\cap\cdots\cap V_+(f_r)
\]
for homogeneous polynomials \(f_i\).  Then
\[
\mathbb P_k^n\setminus Y
=D_+(f_1)\cup\cdots\cup D_+(f_r).
\]
These are \(r\) affine opens.  Applying (c) gives
\(\operatorname{cd}(\mathbb P_k^n\setminus Y)\leq r-1\).
\end{solution}

\begin{exerciseintro}
% record: H-III-04-009
\exerciseheading{Exercise III.4.9}
\addcontentsline{toc}{subsection}{Exercise III.4.9}

\begin{context}
Let
\[
X=\operatorname{Spec}k[x_1,x_2,x_3,x_4],
\quad
Y_1=V(x_1,x_2),
\quad
Y_2=V(x_3,x_4),
\quad
Y=Y_1\cup Y_2.
\]
An open cover by two affine subsets has cohomological dimension at
most one.
\end{context}

\begin{problem}
Prove that \(Y\) is not a set-theoretic complete intersection in
\(X\).  Deduce that its projective closure in \(\mathbb P_k^4\) is
not a set-theoretic complete intersection.
\end{problem}
\end{exerciseintro}

\begin{solution}
The closed set \(Y\) has codimension two.  If it were a
set-theoretic complete intersection, we could write
\[
Y=H_1\cap H_2
\]
with affine hypersurfaces \(H_i=V(f_i)\).  Its complement would then
be
\[
X\setminus Y=D(f_1)\cup D(f_2),
\]
a union of two affine opens.  This would imply
\(\operatorname{cd}(X\setminus Y)\leq1\).

We show instead that \(H^2(X\setminus Y,\mathcal O_X)\neq0\).  Put
\[
U_1=X\setminus Y_1,\qquad U_2=X\setminus Y_2.
\]
Then
\[
U_1\cap U_2=X\setminus Y,\qquad
U_1\cup U_2=X\setminus\{0\}.
\tag{1}
\]
Each \(U_i\) is covered by two principal affine opens, so
\(H^q(U_i,\mathcal O_X)=0\) for \(q\geq2\).  The Mayer--Vietoris
sequence for the cover in (1) therefore gives an isomorphism
\[
H^2(X\setminus Y,\mathcal O_X)
\cong H^3(X\setminus\{0\},\mathcal O_X).
\tag{2}
\]

Cover \(X\setminus\{0\}\) by the four opens \(D(x_i)\).  The cover is
acyclic, and the top Čech quotient is
\[
\frac{k[x_1^{\pm1},x_2^{\pm1},x_3^{\pm1},x_4^{\pm1}]}
{\displaystyle\sum_{i=1}^4
k[x_1,\ldots,x_4,x_j^{-1}:j\neq i]}.
\tag{3}
\]
The residue classes of Laurent monomials
\[
x_1^{a_1}x_2^{a_2}x_3^{a_3}x_4^{a_4},
\qquad a_i<0\text{ for every }i,
\]
form a basis of (3).  In particular, the class of
\((x_1x_2x_3x_4)^{-1}\) is nonzero.  Thus the right side of (2) is
nonzero, contradicting the cohomological-dimension consequence of a
set-theoretic complete-intersection presentation.

Finally, if the projective closure \(\overline Y\subseteq\mathbb P_k^4\)
were the set-theoretic intersection of two projective hypersurfaces,
restriction to the standard affine chart containing \(X=\mathbb A_k^4\)
would express \(Y\) as the intersection of two affine hypersurfaces.
This is impossible by what we have just proved.
\end{solution}

\begin{exerciseintro}
% record: H-III-04-010
\exerciseheading{Exercise III.4.10}
\addcontentsline{toc}{subsection}{Exercise III.4.10}

\begin{context}
Let \(X\) be a nonsingular variety over an algebraically closed field
\(k\), let \(\mathcal F\) be coherent, and let
\(\mathcal T_X=\mathcal Hom_{\mathcal O_X}
(\Omega_{X/k},\mathcal O_X)\).  An infinitesimal extension of \(X\)
by \(\mathcal F\) is a square-zero extension of its structure sheaf
whose ideal is identified with \(\mathcal F\).  Such an extension is
split on a nonsingular affine open.
\end{context}

\begin{problem}
Prove that the isomorphism classes of infinitesimal extensions of
\(X\) by \(\mathcal F\) are in natural one-to-one correspondence with
\[
H^1(X,\mathcal F\otimes_{\mathcal O_X}\mathcal T_X).
\]
\end{problem}
\end{exerciseintro}

\begin{solution}
Choose an affine open cover \(\mathfrak U=(U_i)\) of \(X\).  On each
\(U_i\), local splitting gives an isomorphism of square-zero
extensions with
\[
\mathcal O_{U_i}\oplus\mathcal F|_{U_i},
\qquad
(a,m)(b,n)=(ab,an+bm).
\tag{1}
\]
Choose such a trivialization \(\varphi_i\).

We first identify the possible transition automorphisms.  An
automorphism of (1) which induces the identity on both the quotient
\(\mathcal O_{U_i}\) and the ideal \(\mathcal F|_{U_i}\) has the form
\[
(a,m)\longmapsto(a,m+D(a)).
\tag{2}
\]
Compatibility with multiplication is equivalent to
\[
D(ab)=aD(b)+bD(a),
\]
so \(D\) is a \(k\)-derivation
\(\mathcal O_{U_i}\to\mathcal F|_{U_i}\).  Conversely every such
derivation defines (2), and composition of these automorphisms
corresponds to addition of derivations.

On \(U_i\cap U_j\), the two chosen trivializations therefore differ
by a derivation
\[
D_{ij}\in
\operatorname{Der}_k(\mathcal O_X,\mathcal F)(U_i\cap U_j).
\]
The compatibility on triple intersections is
\[
D_{ij}+D_{jk}=D_{ik},
\tag{3}
\]
so \((D_{ij})\) is a Čech \(1\)-cocycle.  Replacing
\(\varphi_i\) by a different local trivialization adds a local
derivation \(D_i\), and changes \((D_{ij})\) by its Čech
coboundary.  Isomorphic extensions give the same cohomology class.

Conversely, a cocycle satisfying (3) gives transition automorphisms
of the split algebras (1).  Gluing them produces a square-zero
extension of \(X\) by \(\mathcal F\), and cohomologous cocycles give
isomorphic extensions.  Thus the desired set of isomorphism classes
is
\[
\check H^1\bigl(\mathfrak U,
\operatorname{Der}_k(\mathcal O_X,\mathcal F)\bigr).
\tag{4}
\]

The universal property of differentials gives
\[
\operatorname{Der}_k(\mathcal O_X,\mathcal F)
\cong\mathcal Hom_{\mathcal O_X}(\Omega_{X/k},\mathcal F).
\]
Since \(X\) is nonsingular, \(\Omega_{X/k}\) is locally free of finite
rank, and hence
\[
\mathcal Hom_{\mathcal O_X}(\Omega_{X/k},\mathcal F)
\cong\mathcal F\otimes_{\mathcal O_X}\mathcal T_X.
\tag{5}
\]
The intersections in the affine cover are affine because \(X\) is
separated, so Čech cohomology of the quasi-coherent sheaf (5) equals
derived sheaf cohomology.  Combining (4) and (5) gives the asserted
correspondence.
\end{solution}

\begin{exerciseintro}
% record: H-III-04-011
\exerciseheading{Exercise III.4.11}
\addcontentsline{toc}{subsection}{Exercise III.4.11}

\begin{context}
Let \(X\) be a topological space, \(\mathcal F\) a sheaf of abelian
groups, and \(\mathfrak U=(U_i)\) an open cover.  Assume that for
every nonempty finite intersection
\[
V=U_{i_0}\cap\cdots\cap U_{i_p}
\]
one has \(H^q(V,\mathcal F|_V)=0\) for all \(q>0\).
\end{context}

\begin{problem}
Prove that every natural comparison map
\[
\check H^p(\mathfrak U,\mathcal F)
\longrightarrow H^p(X,\mathcal F)
\]
is an isomorphism.  Deduce the Čech comparison theorem for a
quasi-coherent sheaf on a noetherian separated scheme and an affine
open cover.
\end{problem}
\end{exerciseintro}

\begin{solution}
Choose an injective resolution
\(\mathcal F\to\mathcal I^\bullet\), and form the first-quadrant
double complex
\[
K^{p,q}
=
\prod_{i_0<\cdots<i_p}
\Gamma(U_{i_0}\cap\cdots\cap U_{i_p},\mathcal I^q).
\tag{1}
\]
The horizontal differential is the Čech differential and the
vertical differential comes from the resolution.

Compute the cohomology of the total complex in the horizontal
direction first.  Every injective sheaf is flasque, and a flasque
sheaf has zero positive Čech cohomology for every cover.  Hence the
augmented horizontal row
\[
0\longrightarrow\Gamma(X,\mathcal I^q)
\longrightarrow K^{0,q}\longrightarrow K^{1,q}
\longrightarrow\cdots
\]
is exact.  It follows that
\[
H^n(\operatorname{Tot}K^{\bullet,\bullet})
\cong H^n(\Gamma(X,\mathcal I^\bullet))
=H^n(X,\mathcal F).
\tag{2}
\]

Now compute vertically first.  Restricting
\(\mathcal I^\bullet\) to an open set gives a flasque resolution of
the restricted sheaf, so
\[
H^q\!\left(
\Gamma(U_{i_0}\cap\cdots\cap U_{i_p},\mathcal I^\bullet)
\right)
\cong
H^q(U_{i_0}\cap\cdots\cap U_{i_p},\mathcal F).
\tag{3}
\]
By hypothesis, (3) is zero for \(q>0\), while for \(q=0\) it is
\(\Gamma(U_{i_0}\cap\cdots\cap U_{i_p},\mathcal F)\).  Thus the
vertical cohomology of (1) is concentrated in its bottom row, and
that row is exactly the Čech complex \(C^\bullet(\mathfrak U,\mathcal
F)\).  Therefore
\[
H^n(\operatorname{Tot}K^{\bullet,\bullet})
\cong\check H^n(\mathfrak U,\mathcal F).
\tag{4}
\]
The identifications (2) and (4) are induced by the augmentations and
their composite is the natural comparison map.  Hence that map is an
isomorphism in every degree.

Finally, let \(X\) be noetherian and separated, let \(\mathcal F\) be
quasi-coherent, and let \(\mathfrak U\) be an affine open cover.
Every finite intersection of members of \(\mathfrak U\) is affine,
and the restriction of \(\mathcal F\) to it has no higher
cohomology.  The hypothesis above is satisfied, so the usual Čech
comparison theorem follows.
\end{solution}

\begin{exerciseintro}
\corpussection{III.5 The Cohomology of Projective Space}
\addcontentsline{toc}{section}{III.5 The Cohomology of Projective Space}

% record: H-III-05-001
\exerciseheading{Exercise III.5.1}
\addcontentsline{toc}{subsection}{Exercise III.5.1}

\begin{context}
Let \(X\) be projective over a field \(k\).  For a coherent sheaf
\(\mathcal F\), all \(H^i(X,\mathcal F)\) are finite-dimensional and
vanish for \(i\) sufficiently large, so
\[
\chi(\mathcal F)=\sum_i(-1)^i\dim_kH^i(X,\mathcal F)
\]
is well defined.
\end{context}

\begin{problem}
If
\[
0\longrightarrow\mathcal F'\longrightarrow\mathcal F
\longrightarrow\mathcal F''\longrightarrow0
\]
is a short exact sequence of coherent sheaves, prove
\[
\chi(\mathcal F)=\chi(\mathcal F')+\chi(\mathcal F'').
\]
\end{problem}
\end{exerciseintro}

\begin{solution}
The short exact sequence gives a long exact sequence of
finite-dimensional vector spaces
\[
0\to H^0(\mathcal F')\to H^0(\mathcal F)\to H^0(\mathcal F'')
\to H^1(\mathcal F')\to\cdots,
\]
which terminates because the cohomology groups vanish in high degree.
For every finite exact sequence of finite-dimensional vector spaces,
the alternating sum of the dimensions is zero.  Applying this fact
to the displayed sequence and collecting the terms belonging to each
sheaf gives
\[
\sum_i(-1)^i\dim H^i(\mathcal F)
=
\sum_i(-1)^i\dim H^i(\mathcal F')
+
\sum_i(-1)^i\dim H^i(\mathcal F''),
\]
which is the required identity.
\end{solution}

\begin{exerciseintro}
% record: H-III-05-002
\exerciseheading{Exercise III.5.2}
\addcontentsline{toc}{subsection}{Exercise III.5.2}

\begin{context}
Let \(X\) be projective over \(k\), let \(\mathcal O_X(1)\) be very
ample, and let \(\mathcal F\) be coherent.  After embedding
\(i:X\to\mathbb P_k^N\), the sheaf \(i_*\mathcal F\) admits a finite
resolution by finite direct sums of twists of
\(\mathcal O_{\mathbb P^N}\).  Also
\[
\chi(\mathbb P^N,\mathcal O_{\mathbb P^N}(m))
=\binom{m+N}{N}
\]
as a polynomial identity for every \(m\in\mathbb Z\).
The section module \(\Gamma_*(\mathcal F)\), indexed over all
integers, need not be finitely generated.  Its Hilbert polynomial
means that of a sufficiently high finite truncation, equivalently
the polynomial agreeing with its dimensions in large degrees.
\end{context}

\begin{problem}
(a) Prove that there is a polynomial \(P_{\mathcal F}(z)\in
\mathbb Q[z]\) such that
\[
\chi(\mathcal F(n))=P_{\mathcal F}(n)
\qquad(n\in\mathbb Z).
\]

(b) For \(X=\mathbb P_k^r\) and
\(M=\Gamma_*(\mathcal F)\), prove that \(P_{\mathcal F}\) is the
usual Hilbert polynomial of the graded module \(M\).
\end{problem}
\end{exerciseintro}

\begin{solution}
(a) Use the projective embedding defined by
\(\mathcal O_X(1)\).  Direct image under the closed immersion
preserves cohomology and satisfies
\[
i_*(\mathcal F(n))=(i_*\mathcal F)(n).
\]
Choose a finite resolution
\[
0\longrightarrow\mathcal E_m\longrightarrow\cdots
\longrightarrow\mathcal E_0\longrightarrow i_*\mathcal F
\longrightarrow0,
\]
where
\[
\mathcal E_j=\bigoplus_\ell
\mathcal O_{\mathbb P^N}(-a_{j\ell}).
\]
Twisting is exact.  Repeated additivity of Euler characteristic gives
\[
\chi(\mathcal F(n))
=
\sum_{j,\ell}(-1)^j
\chi\bigl(\mathcal O_{\mathbb P^N}(n-a_{j\ell})\bigr)
=
\sum_{j,\ell}(-1)^j
\binom{n-a_{j\ell}+N}{N}.
\tag{1}
\]
The right side is the value at \(n\) of a polynomial in
\(\mathbb Q[z]\), and (1) holds for every integer \(n\).  This defines
\(P_{\mathcal F}\); it is unique because two polynomials agreeing at
all integers are equal.

(b) By Serre vanishing, for \(n\) sufficiently large,
\[
H^i(\mathbb P^r,\mathcal F(n))=0\qquad(i>0).
\]
Consequently
\[
P_{\mathcal F}(n)=\chi(\mathcal F(n))
=\dim_kH^0(\mathbb P^r,\mathcal F(n))
=\dim_k M_n
\qquad(n\gg0).
\]
The graded-module Hilbert polynomial is characterized as the
polynomial eventually equal to \(\dim_kM_n\).  Hence it equals
\(P_{\mathcal F}\).
\end{solution}

\begin{exerciseintro}
% record: H-III-05-003
\exerciseheading{Exercise III.5.3}
\addcontentsline{toc}{subsection}{Exercise III.5.3}

\begin{context}
For a projective scheme \(X\) of dimension \(r\), define
\[
p_a(X)=(-1)^r\bigl(\chi(\mathcal O_X)-1\bigr).
\]
For a projective embedding, the Hilbert polynomial satisfies
\(P_X(n)=\chi(\mathcal O_X(n))\) for every integer \(n\).
\end{context}

\begin{problem}
(a) If \(X\) is integral and \(k\) is algebraically closed, prove
\(H^0(X,\mathcal O_X)=k\), and deduce
\[
p_a(X)=\sum_{j=0}^{r-1}(-1)^j
\dim_kH^{r-j}(X,\mathcal O_X).
\]
In particular, for a curve,
\(p_a(X)=\dim_kH^1(X,\mathcal O_X)\).

(b) Prove that for a closed projective variety this definition agrees
with the Hilbert-polynomial definition of arithmetic genus.

(c) Prove that arithmetic genus is a birational invariant of
nonsingular projective curves over an algebraically closed field.
Deduce that a nonsingular plane curve of degree \(d\geq3\) is not
rational.
\end{problem}
\end{exerciseintro}

\begin{solution}
(a) A global regular function gives a morphism
\(X\to\mathbb A_k^1\).  Since \(X\) is proper, its image is a proper
closed integral subscheme of \(\mathbb A_k^1\), hence a single closed
point.  The field is algebraically closed, so the function is a
constant in \(k\).  Thus \(H^0(X,\mathcal O_X)=k\).

Writing \(h^i=\dim_kH^i(X,\mathcal O_X)\), we obtain
\[
\chi(\mathcal O_X)-1
=\sum_{i=1}^r(-1)^i h^i.
\]
Multiplication by \((-1)^r\), followed by the substitution
\(i=r-j\), gives the formula in the statement.  When \(r=1\), only
\(h^1\) remains.

(b) If \(e=\dim X\), the earlier definition from the Hilbert
polynomial is
\[
(-1)^e\bigl(P_X(0)-1\bigr).
\]
But \(P_X(0)=\chi(\mathcal O_X)\).  Hence this expression is exactly
the present definition, and in particular does not depend on the
chosen projective embedding.

(c) A birational map between nonsingular projective curves extends
to an isomorphism.  Their structure-sheaf cohomology, and therefore
their arithmetic genera, are consequently equal.  A rational smooth
projective curve is isomorphic to \(\mathbb P^1\) and has arithmetic
genus zero.  On the other hand, the Čech computation for a plane
curve of degree \(d\) gives
\[
p_a(X)=\dim_kH^1(X,\mathcal O_X)
=\frac{(d-1)(d-2)}2.
\]
This is positive for \(d\geq3\), so such a nonsingular plane curve
cannot be rational.
\end{solution}

\begin{exerciseintro}
% record: H-III-05-004
\exerciseheading{Exercise III.5.4}
\addcontentsline{toc}{subsection}{Exercise III.5.4}

\begin{context}
Let \(K(X)\) be the Grothendieck group of coherent sheaves on a
noetherian scheme.  It is generated by symbols
\(\gamma(\mathcal F)\), with the relation
\[
\gamma(\mathcal F)=\gamma(\mathcal F')+\gamma(\mathcal F'')
\]
for every short exact sequence
\(0\to\mathcal F'\to\mathcal F\to\mathcal F''\to0\).
For an arbitrary choice of linear subspaces, first prove the result for a flag. Injectivity of the Hilbert-polynomial map then identifies any other subspace of the same dimension with its flag representative.
\end{context}

\begin{problem}
(a) For a projective scheme \(X\) with a fixed very ample
\(\mathcal O_X(1)\), prove that there is a unique additive map
\[
P:K(X)\longrightarrow\mathbb Q[z]
\]
sending \(\gamma(\mathcal F)\) to the Hilbert polynomial of
\(\mathcal F\).

(b) Let \(X=\mathbb P_k^r\), and for \(0\leq i\leq r\) let
\(L_i\subseteq X\) be a linear subspace of dimension \(i\).  Prove
that \(K(X)\) is freely generated by the classes
\(\gamma(\mathcal O_{L_i})\), and that \(P\) is injective.
\end{problem}
\end{exerciseintro}

\begin{solution}
(a) Hilbert polynomials are additive in short exact sequences because
Euler characteristics are additive after every twist.  Thus the
assignment
\[
\gamma(\mathcal F)\longmapsto P_{\mathcal F}
\]
respects all defining relations of \(K(X)\), and so induces an
additive homomorphism \(P\).  The classes of coherent sheaves
generate \(K(X)\), which also proves uniqueness.

(b) We first prove generation by induction on \(r\).  Let
\(H=\mathbb P_k^{r-1}\) be a hyperplane and
\(U=X\setminus H\cong\mathbb A_k^r\).  The localization sequence is
\[
K(H)\longrightarrow K(X)\longrightarrow K(U)\longrightarrow0.
\tag{1}
\]
Every finite module over \(k[x_1,\ldots,x_r]\) has a finite free
resolution.  Hence \(K(U)\cong\mathbb Z\), generated by
\(\gamma(\mathcal O_U)\); the rank map shows that this generator has
infinite order.  The class
\(\gamma(\mathcal O_{L_r})=\gamma(\mathcal O_X)\) restricts to this
generator.  By induction, \(K(H)\) is generated by the classes of
linear spaces \(L_0,\ldots,L_{r-1}\) lying in \(H\).  Exactness of
(1) therefore shows that
\[
\gamma(\mathcal O_{L_0}),\ldots,\gamma(\mathcal O_{L_r})
\]
generate \(K(X)\).

Their Hilbert polynomials are
\[
P\bigl(\gamma(\mathcal O_{L_i})\bigr)
=\binom{z+i}{i}
=\frac{(z+1)\cdots(z+i)}{i!}.
\tag{2}
\]
The polynomials in (2) have distinct degrees \(0,1,\ldots,r\), with
nonzero leading coefficients, and are therefore linearly independent
over \(\mathbb Q\).  Any integral relation among the classes would,
after applying \(P\), give a relation among these polynomials, so all
coefficients must vanish.  The displayed classes are consequently a
free basis.  Finally, every element of \(K(X)\) has a unique
expression in that basis; (2) shows that an element with zero Hilbert
polynomial has all coefficients zero.  Thus \(P\) is injective.
\end{solution}

\begin{exerciseintro}
% record: H-III-05-005
\exerciseheading{Exercise III.5.5}
\addcontentsline{toc}{subsection}{Exercise III.5.5}

\begin{context}
Let \(X=\mathbb P_k^r\), and let \(Y\subseteq X\) be a complete
intersection of dimension \(q\geq1\).  Write \(c=r-q\).  Its
homogeneous ideal is generated by a regular sequence
\(f_1,\ldots,f_c\), with \(d_i=\deg f_i\).  Put
\[
Y_i=V_+(f_1,\ldots,f_i),\qquad Y_0=X,\qquad Y_c=Y.
\]
\end{context}

\begin{problem}
(a) Prove that for every \(n\in\mathbb Z\), the restriction map
\[
H^0(X,\mathcal O_X(n))\longrightarrow
H^0(Y,\mathcal O_Y(n))
\]
is surjective.

(b) Prove that \(Y\) is connected.

(c) Prove
\[
H^j(Y,\mathcal O_Y(n))=0
\qquad(0<j<q,\ n\in\mathbb Z).
\]

(d) Prove
\[
p_a(Y)=\dim_kH^q(Y,\mathcal O_Y).
\]
\end{problem}
\end{exerciseintro}

\begin{solution}
(a) Regularity of the sequence gives, for \(1\leq i\leq c\), exact
sequences
\[
0\longrightarrow
\mathcal O_{Y_{i-1}}(n-d_i)
\xrightarrow{\,f_i\,}
\mathcal O_{Y_{i-1}}(n)
\longrightarrow\mathcal O_{Y_i}(n)
\longrightarrow0.
\tag{1}
\]
We prove simultaneously, by induction on \(i\), that
\[
\begin{aligned}
&H^0(X,\mathcal O_X(n))\longrightarrow
H^0(Y_i,\mathcal O_{Y_i}(n))
\quad\text{is surjective for every }n,\\
&H^j(Y_i,\mathcal O_{Y_i}(n))=0
\quad\text{for }0<j<\dim Y_i.
\end{aligned}
\tag{2}
\]
For \(i=0\), these assertions are the standard cohomology
calculation for projective space.

Assume (2) for \(Y_{i-1}\), whose dimension is
\(r-i+1\geq q+1\geq2\).  The \(H^1\)-vanishing in (2), applied to
\(\mathcal O_{Y_{i-1}}(n-d_i)\), shows from (1) that
\[
H^0(Y_{i-1},\mathcal O_{Y_{i-1}}(n))
\longrightarrow H^0(Y_i,\mathcal O_{Y_i}(n))
\]
is surjective.  Composing with the map from \(X\) proves the first
line of (2) for \(Y_i\).

Now let \(0<j<\dim Y_i=r-i\).  The portion of the long exact sequence
of (1) surrounding \(H^j(Y_i,\mathcal O_{Y_i}(n))\) has outer terms
\[
H^j(Y_{i-1},\mathcal O_{Y_{i-1}}(n))
\quad\text{and}\quad
H^{j+1}(Y_{i-1},\mathcal O_{Y_{i-1}}(n-d_i)).
\]
Both vanish by the inductive hypothesis, because
\(0<j<j+1<\dim Y_{i-1}\).  This proves the second line of (2).  Taking
\(i=c\) proves (a).

(b) Part (a) with \(n=0\) gives a surjection
\[
k=H^0(X,\mathcal O_X)\longrightarrow H^0(Y,\mathcal O_Y).
\]
Constants inject, so \(H^0(Y,\mathcal O_Y)=k\).  A decomposition of
\(Y\) into two nonempty open-and-closed subsets would give a
nontrivial idempotent in this ring.  Hence \(Y\) is connected.

(c) This is the second assertion of (2) for \(i=c\).

(d) By (b), \(H^0(Y,\mathcal O_Y)=k\), and by (c) all cohomology in
degrees \(1,\ldots,q-1\) vanishes.  Therefore
\[
\chi(\mathcal O_Y)=1+(-1)^q
\dim_kH^q(Y,\mathcal O_Y).
\]
The definition
\(p_a(Y)=(-1)^q(\chi(\mathcal O_Y)-1)\) gives the claimed equality.
\end{solution}

\begin{exerciseintro}
% record: H-III-05-006
\exerciseheading{Exercise III.5.6}
\addcontentsline{toc}{subsection}{Exercise III.5.6}

\begin{context}
Let \(Q=V(xy-zw)\subseteq\mathbb P_k^3\).  The Segre map identifies
\(Q\) with \(\mathbb P_k^1\times_k\mathbb P_k^1\), and under this
identification

\[
\mathcal O_Q(a,b)
=p_1^*\mathcal O_{\mathbb P^1}(a)\otimes
 p_2^*\mathcal O_{\mathbb P^1}(b),
\qquad
\mathcal O_Q(n)=\mathcal O_Q(n,n).
\]

For \(q>0\), unions of \(q\) distinct lines in either ruling give
exact sequences

\[
\begin{aligned}
0&\to\mathcal O_Q(a,b-q)\to\mathcal O_Q(a,b)
\to\mathcal O_{\mathbb P^1}(a)^{\oplus q}\to0,\\
0&\to\mathcal O_Q(a-q,b)\to\mathcal O_Q(a,b)
\to\mathcal O_{\mathbb P^1}(b)^{\oplus q}\to0.
\end{aligned}
\tag{1}
\]

For the cohomological assertions one may first extend \(k\) to an
algebraic closure: coherent cohomology commutes with field extension,
and the extension is faithfully flat.  We may use the cohomology of
\(\mathcal O_{\mathbb P^r}(m)\) and Bertini's theorem.
A smooth curve of type \((1,3)\) on a quadric has degree four and a degree-one projection to \(\mathbb P^1\), hence is rational.
\end{context}

\begin{problem}
(a) Prove the following assertions.

\begin{enumerate}
\item[(1)] If \(\lvert a-b\rvert\leq1\), then
\(H^1(Q,\mathcal O_Q(a,b))=0\).
\item[(2)] If \(a,b<0\), then
\(H^1(Q,\mathcal O_Q(a,b))=0\).
\item[(3)] If \(a\leq-2\), then
\(H^1(Q,\mathcal O_Q(a,0))\neq0\).
\end{enumerate}

(b) Let \(Y\subseteq Q\) be a locally principal closed subscheme of
type \((a,b)\).

\begin{enumerate}
\item[(1)] If \(a,b>0\), prove that \(Y\) is connected.
\item[(2)] If \(k\) is algebraically closed and \(a,b>0\), prove that
there is an irreducible nonsingular curve of type \((a,b)\).
\item[(3)] If \(Y\) is irreducible and nonsingular and \(a,b>0\), prove
that its given embedding in \(\mathbb P^3\) is projectively normal if
and only if \(\lvert a-b\rvert\leq1\).  In particular, type \((1,3)\)
gives a nonsingular rational quartic that is not projectively normal.
\end{enumerate}

(c) Prove
\[
p_a(Y)=ab-a-b+1.
\]
\end{problem}
\end{exerciseintro}

\begin{solution}
(a) The quadric sequence

\[
0\longrightarrow\mathcal O_{\mathbb P^3}(n-2)
\longrightarrow\mathcal O_{\mathbb P^3}(n)
\longrightarrow\mathcal O_Q(n,n)\longrightarrow0
\]

and the vanishing of the intermediate cohomology of line bundles on
\(\mathbb P^3\) give

\[
H^1(Q,\mathcal O_Q(n,n))=0
\qquad(n\in\mathbb Z).
\tag{2}
\]

Let \(L\) be a ruling line of type \((0,1)\).  The case \(q=1\) of
(1) gives

\[
0\to\mathcal O_Q(n,n-1)\to\mathcal O_Q(n,n)
\to\mathcal O_L(n)\to0.
\tag{3}
\]

If \(n<0\), then \(H^0(L,\mathcal O_L(n))=0\).  If \(n\geq0\), the
restriction map from \(\mathbb P^3\), and hence the one from \(Q\), to
\(L\) is surjective on degree-\(n\) forms.  The cohomology sequence
of (3) and (2) therefore show that
\(H^1(Q,\mathcal O_Q(n,n-1))=0\) for every \(n\).  Interchanging the
two rulings gives the same conclusion for
\(\mathcal O_Q(n-1,n)\).  Together with (2), this proves (1).

For (2), after the harmless field extension described in the context,
choose a reduced divisor \(D\) consisting of \(-a\) lines of one
ruling and \(-b\) lines of the other.  It has type \((-a,-b)\), and
it is connected because every line of one ruling meets every line of
the other.  The divisor sequence is

\[
0\to\mathcal O_Q(a,b)\to\mathcal O_Q\to\mathcal O_D\to0.
\]

Both \(Q\) and \(D\) are proper, reduced, and connected, so their
rings of global functions are \(k\).  Since
\(H^1(Q,\mathcal O_Q)=0\) by (2), the cohomology sequence gives
\(H^1(Q,\mathcal O_Q(a,b))=0\).

For (3), let \(D\) be the disjoint union of \(-a\) lines in one
ruling.  From

\[
0\to\mathcal O_Q(a,0)\to\mathcal O_Q\to\mathcal O_D\to0
\]

we obtain

\[
k\longrightarrow k^{-a}\longrightarrow
H^1(Q,\mathcal O_Q(a,0))\longrightarrow0.
\]

The first map is diagonal, so its cokernel has dimension
\(-a-1>0\).  Faithful flatness then gives the same nonvanishing over
the original field.

(b) For (1), the divisor sequence is

\[
0\longrightarrow\mathcal O_Q(-a,-b)
\longrightarrow\mathcal O_Q
\longrightarrow\mathcal O_Y\longrightarrow0.
\tag{4}
\]

Since \(-a,-b<0\), part (a2) shows that the restriction

\[
k=H^0(Q,\mathcal O_Q)\longrightarrow H^0(Y,\mathcal O_Y)
\]

is surjective.  It is injective because it preserves constants, so
\(H^0(Y,\mathcal O_Y)=k\).  This ring has no nontrivial idempotent;
hence \(Y\) is connected.

For (2), if \(a,b>0\), the bundle \(\mathcal O_Q(a,b)\) is very ample.
Bertini's theorem gives a nonsingular general member \(Y\) of its
complete linear system.  Every such member is connected by part (1).  A
nonsingular curve over an algebraically closed field is the disjoint
union of its irreducible components, so this connected \(Y\) is
irreducible.

For (3), since \(Y\) is nonsingular, it is normal.  Thus it is
projectively normal precisely when

\[
H^0(\mathbb P^3,\mathcal O_{\mathbb P^3}(n))
\longrightarrow H^0(Y,\mathcal O_Y(n))
\tag{5}
\]

is surjective for every \(n\geq0\).  The quadric sequence

\[
0\to\mathcal O_{\mathbb P^3}(n-2)
\to\mathcal O_{\mathbb P^3}(n)
\to\mathcal O_Q(n,n)\to0
\]

and the vanishing of \(H^1(\mathbb P^3,\mathcal O(m))\) show that the
restriction from \(\mathbb P^3\) to \(Q\) is always surjective.
Meanwhile, twisting (4) gives

\[
0\to\mathcal O_Q(n-a,n-b)\to\mathcal O_Q(n,n)
\to\mathcal O_Y(n)\to0.
\tag{6}
\]

Because \(H^1(Q,\mathcal O_Q(n,n))=0\), its cohomology sequence
identifies the cokernel of the restriction in (6) with
\(H^1(Q,\mathcal O_Q(n-a,n-b))\).  If
\(\lvert a-b\rvert\leq1\), part (a1) makes this group zero for every
\(n\), so (5) is surjective.

Conversely, suppose for instance that \(a\geq b+2\).  Taking \(n=b\)
in (6), part (a3) gives

\[
H^1(Q,\mathcal O_Q(b-a,0))\neq0.
\]

Thus restriction fails to be surjective in degree \(b\geq0\), and
\(Y\) is not projectively normal.  The case \(b\geq a+2\) is
symmetric.

For \((a,b)=(1,3)\), the criterion shows failure of projective
normality; this curve is the rational quartic from Exercise I.3.18.

(c) A ruling line and additivity of Euler characteristic give, for
all integers \(u,v\),

\[
\begin{aligned}
\chi(\mathcal O_Q(u,v))-\chi(\mathcal O_Q(u,v-1))&=u+1,\\
\chi(\mathcal O_Q(u,0))-\chi(\mathcal O_Q(u-1,0))&=1.
\end{aligned}
\]

The quadric sequence gives \(\chi(\mathcal O_Q)=1\).  The second
recurrence, followed by the first, therefore yields

\[
\chi(\mathcal O_Q(u,v))=(u+1)(v+1)
\qquad(u,v\in\mathbb Z).
\]

Taking Euler characteristics in (4),

\[
\chi(\mathcal O_Y)
=1-(1-a)(1-b)=a+b-ab.
\]

Since \(Y\) has dimension one,

\[
p_a(Y)=1-\chi(\mathcal O_Y)=ab-a-b+1.
\]
\end{solution}

\begin{exerciseintro}
% record: H-III-05-007
\exerciseheading{Exercise III.5.7}
\addcontentsline{toc}{subsection}{Exercise III.5.7}

\begin{context}
Let \(X\) and \(Y\) be proper schemes over a noetherian ring \(A\),
and let \(\mathcal L\) denote an invertible sheaf.  We may use both
standard criteria for ampleness: sufficiently high twists of every
coherent sheaf are globally generated, and, on a proper scheme,

\[
\mathcal L\text{ is ample}
\quad\Longleftrightarrow\quad
H^i(X,\mathcal F\otimes\mathcal L^n)=0
\quad(i>0,\ n\gg0)
\]

for every coherent \(\mathcal F\).
\end{context}

\begin{problem}
(a) If \(\mathcal L\) is ample on \(X\) and
\(i:Y\hookrightarrow X\) is a closed immersion, prove that
\(i^*\mathcal L\) is ample on \(Y\).

(b) Prove that \(\mathcal L\) is ample on \(X\) if and only if its
restriction \(\mathcal L_{\mathrm{red}}\) is ample on
\(X_{\mathrm{red}}\).

(c) Suppose that \(X\) is reduced, with irreducible components
\(X_1,\ldots,X_s\).  Prove that \(\mathcal L\) is ample on \(X\) if
and only if every \(\mathcal L|_{X_j}\) is ample.

(d) Let \(f:X\to Y\) be finite and surjective.  Prove that
\(\mathcal L\) is ample on \(Y\) if and only if \(f^*\mathcal L\) is
ample on \(X\).
\end{problem}
\end{exerciseintro}

\begin{solution}
(a) Let \(\mathcal F\) be coherent on \(Y\).  The sheaf
\(i_*\mathcal F\) is coherent, so for \(n\gg0\),

\[
i_*\mathcal F\otimes\mathcal L^n
\cong i_*(\mathcal F\otimes i^*\mathcal L^n)
\]

is generated by its global sections.  Its global sections are exactly
those of \(\mathcal F\otimes i^*\mathcal L^n\), and restricting the
evaluation map to \(Y\) shows that the latter sheaf is globally
generated.  Hence \(i^*\mathcal L\) is ample.

(b) The forward implication is (a).  Conversely, let
\(\mathcal N\subseteq\mathcal O_X\) be the nilradical.  Since \(X\) is
noetherian, \(\mathcal N^m=0\) for some \(m\).  For any coherent
\(\mathcal F\), the finite filtration

\[
\mathcal F\supseteq\mathcal N\mathcal F\supseteq\cdots
\supseteq\mathcal N^m\mathcal F=0
\tag{1}
\]

has coherent successive quotients annihilated by \(\mathcal N\), and
therefore coherent on \(X_{\mathrm{red}}\).  If
\(\mathcal L_{\mathrm{red}}\) is ample, Serre vanishing applies to
every quotient of (1).  The long exact cohomology sequences for the
filtration then give

\[
H^i(X,\mathcal F\otimes\mathcal L^n)=0
\qquad(i>0,\ n\gg0).
\]

The cohomological criterion proves that \(\mathcal L\) is ample.

(c) Necessity follows from (a).  For sufficiency, let
\(\mathcal I_j\) be the ideal of \(X_j\).  Reducedness gives
\(\bigcap_j\mathcal I_j=0\), hence
\(\mathcal I_1\cdots\mathcal I_s=0\).  Filter a coherent sheaf
\(\mathcal F\) by

\[
\mathcal F\supseteq\mathcal I_1\mathcal F\supseteq
\mathcal I_1\mathcal I_2\mathcal F\supseteq\cdots\supseteq0.
\tag{2}
\]

The \(j\)-th successive quotient is annihilated by
\(\mathcal I_j\), so it is a coherent sheaf on \(X_j\).  If every
\(\mathcal L|_{X_j}\) is ample, its higher cohomology vanishes after a
sufficiently high twist.  There are only finitely many quotients in
(2), so the long exact sequences and the cohomological criterion show
that \(\mathcal L\) is ample on \(X\).

(d) First suppose that \(\mathcal L\) is ample, and let
\(\mathcal F\) be coherent on \(X\).  The finite morphism \(f\) is
affine and \(f_*\mathcal F\) is coherent.  For \(n\gg0\), the
projection formula makes

\[
f_*(\mathcal F\otimes f^*\mathcal L^n)
\cong f_*\mathcal F\otimes\mathcal L^n
\]

globally generated.  Pull its evaluation map back to \(X\), and
compose with the counit

\[
f^*f_*(\mathcal F\otimes f^*\mathcal L^n)
\longrightarrow\mathcal F\otimes f^*\mathcal L^n.
\tag{3}
\]

The map (3) is surjective: over
\(\operatorname{Spec}B\subseteq Y\), it is the multiplication map
\(C\otimes_B M\to M\), and \(1\otimes m\mapsto m\).  Thus every high
twist by \(f^*\mathcal L\) is globally generated, proving that
\(f^*\mathcal L\) is ample.

Conversely, assume that \(f^*\mathcal L\) is ample.  We use noetherian
induction on closed subsets of \(Y\).  Parts (b) and (c), together
with base change of \(f\), reduce the inductive step to the case in
which \(Y\) is integral and \(\mathcal L\) is already known to be
ample on every proper closed subscheme of \(Y\).  If necessary, choose
an irreducible component of \(X_{\mathrm{red}}\) that dominates
\(Y\); its map to \(Y\) is still finite and surjective, and the
restriction of \(f^*\mathcal L\) is ample by (a).  We may therefore
also take \(X\) integral.

Let \(\mathcal F\) be coherent on \(Y\).  The finite-morphism lemma
used in the proof of Chevalley's theorem supplies a coherent sheaf
\(\mathcal G\) on \(X\), an integer \(r>0\), and a morphism

\[
\beta:f_*\mathcal G\longrightarrow\mathcal F^{\oplus r}
\]

which is an isomorphism at the generic point of \(Y\).  Put
\(\mathcal K=\ker\beta\), \(\mathcal Q=\operatorname{im}\beta\), and
\(\mathcal C=\operatorname{coker}\beta\).  Both \(\mathcal K\) and
\(\mathcal C\) are supported on proper closed subsets of \(Y\).  The
inductive hypothesis and Serre vanishing therefore give, for \(i>0\)
and \(n\gg0\),

\[
H^i(Y,\mathcal K\otimes\mathcal L^n)=
H^i(Y,\mathcal C\otimes\mathcal L^n)=0.
\tag{4}
\]

On the other hand, finiteness of \(f\), the projection formula, and
the assumed ampleness upstairs give

\[
\begin{split}
H^i(Y,f_*\mathcal G\otimes\mathcal L^n)
&\cong H^i(Y,f_*(\mathcal G\otimes f^*\mathcal L^n))\\
&\cong H^i(X,\mathcal G\otimes f^*\mathcal L^n)=0
\end{split}
\tag{5}
\]

for \(i>0\) and \(n\gg0\).  Apply cohomology first to

\[
0\to\mathcal K\to f_*\mathcal G\to\mathcal Q\to0
\]

and then to

\[
0\to\mathcal Q\to\mathcal F^{\oplus r}\to\mathcal C\to0.
\]

Equations (4) and (5) imply successively that
\(H^i(Y,\mathcal Q\otimes\mathcal L^n)=0\) and
\(H^i(Y,\mathcal F^{\oplus r}\otimes\mathcal L^n)=0\) for every
\(i>0\) and all sufficiently large \(n\).  Hence the same vanishing
holds for \(\mathcal F\).  The cohomological criterion proves that
\(\mathcal L\) is ample on \(Y\).
\end{solution}

\begin{exerciseintro}
% record: H-III-05-008
\exerciseheading{Exercise III.5.8}
\addcontentsline{toc}{subsection}{Exercise III.5.8}

\begin{context}
Let \(k\) be algebraically closed and let \(X\) be a one-dimensional
proper \(k\)-scheme.  We may use that a proper scheme over \(k\) with
an ample invertible sheaf is projective, and the ampleness results of
Exercise III.5.7.
We use the projectivity theorem for nonsingular proper curves, equivalently ampleness of divisors of positive degree on such a curve.
\end{context}

\begin{problem}
Prove that \(X\) is projective by treating successively the following
cases.

(a) \(X\) is irreducible and nonsingular.

(b) \(X\) is integral.  Use its normalization and descend an ample
invertible sheaf.

(c) \(X\) is reduced.  If \(X_1,\ldots,X_s\) are its irreducible
components, prove that

\[
\operatorname{Pic}X\longrightarrow
\bigoplus_{j=1}^s\operatorname{Pic}X_j
\]

is surjective, and use this to construct an ample invertible sheaf on
\(X\).

(d) Treat an arbitrary \(X\) by lifting an ample invertible sheaf
from \(X_{\mathrm{red}}\).
\end{problem}
\end{exerciseintro}

\begin{solution}
(a) An irreducible nonsingular proper curve has a positive effective
divisor, and the associated invertible sheaf is ample.  Equivalently,
this is the projectivity theorem for nonsingular proper curves.

(b) Let \(f:\widetilde X\to X\) be the normalization.  Normalization
is finite for an integral scheme of finite type over a field, so
\(\widetilde X\) is proper.  A normal noetherian curve is regular:
each local ring at a closed point is a one-dimensional noetherian
normal local domain, hence a discrete valuation ring.  Since \(k\) is
perfect, \(\widetilde X\) is nonsingular and is projective by (a).

Choose a very ample invertible sheaf \(\mathcal M\) on
\(\widetilde X\).  The finite birational map \(f\) is an isomorphism
away from a finite set \(S\subseteq\widetilde X\).  Embed
\(\widetilde X\) by \(\mathcal M\) and choose a hyperplane avoiding
\(S\), which is possible because \(k\) is infinite.  Its intersection
with \(\widetilde X\) is an effective Cartier divisor \(D\), supported
where \(f\) is an isomorphism, and

\[
\mathcal O_{\widetilde X}(D)\cong\mathcal M.
\tag{1}
\]

The divisor \(D\) therefore descends to an effective Cartier divisor
\(D_0\) on \(X\): near its support use the isomorphism furnished by
\(f\), and use the zero divisor elsewhere.  Put
\(\mathcal L_0=\mathcal O_X(D_0)\).  Then

\[
f^*\mathcal L_0\cong\mathcal O_{\widetilde X}(D)
\cong\mathcal M.
\]

Since \(f\) is finite and surjective, Exercise III.5.7(d) implies that
\(\mathcal L_0\) is ample.  Thus \(X\) is projective.

(c) Let \(i_j:X_j\hookrightarrow X\) be the inclusions.  Reducedness
gives an injection of sheaves of abelian groups

\[
1\longrightarrow\mathcal O_X^\times
\longrightarrow\prod_{j=1}^s i_{j*}\mathcal O_{X_j}^\times
\longrightarrow\mathcal Q\longrightarrow1.
\tag{2}
\]

Away from the intersections of distinct components, the middle map
is an isomorphism.  Those intersections form a finite set, so
\(\mathcal Q\) is supported at finitely many closed points and
\(H^1(X,\mathcal Q)=0\).  The cohomology sequence of (2) consequently
gives a surjection

\[
\operatorname{Pic}X=H^1(X,\mathcal O_X^\times)
\longrightarrow
\prod_{j=1}^sH^1(X_j,\mathcal O_{X_j}^\times)
=\bigoplus_{j=1}^s\operatorname{Pic}X_j.
\tag{3}
\]

Every one-dimensional component \(X_j\) is projective by (b), while a
zero-dimensional component is already projective.  Choose an ample
invertible sheaf \(\mathcal L_j\) on each component and lift their
tuple through (3) to an invertible sheaf \(\mathcal L\) on \(X\).
Exercise III.5.7(c) says that \(\mathcal L\) is ample, so \(X\) is
projective.

(d) Let \(\mathcal N\) be the nilradical.  It is nilpotent, say
\(\mathcal N^{m+1}=0\), and the schemes

\[
X_j=(X,\mathcal O_X/\mathcal N^j),
\qquad 1\leq j\leq m+1,
\]

form a chain from \(X_1=X_{\mathrm{red}}\) to \(X_{m+1}=X\).  The
ideal
\(\mathcal I_j=\mathcal N^j/\mathcal N^{j+1}\) of
\(X_j\subseteq X_{j+1}\) has square zero.  Hence there is an exact
sequence

\[
0\longrightarrow\mathcal I_j
\xrightarrow{\,a\mapsto1+a\,}
\mathcal O_{X_{j+1}}^\times
\longrightarrow\mathcal O_{X_j}^\times
\longrightarrow1.
\tag{4}
\]

Grothendieck vanishing gives
\(H^2(X_j,\mathcal I_j)=0\), because the underlying space has
dimension one.  The cohomology sequence of (4) therefore makes

\[
\operatorname{Pic}X_{j+1}\longrightarrow\operatorname{Pic}X_j
\]

surjective.  Iterating, every invertible sheaf on
\(X_{\mathrm{red}}\) lifts to \(X\).  By (c), choose an ample one and
lift it to \(\mathcal L\) on \(X\).  Exercise III.5.7(b) shows that
\(\mathcal L\) is ample, and therefore \(X\) is projective.
\end{solution}

\begin{exerciseintro}
% record: H-III-05-009
\exerciseheading{Exercise III.5.9}
\addcontentsline{toc}{subsection}{Exercise III.5.9}

\begin{context}
Let \(k\) be algebraically closed of characteristic zero, let
\(X=\mathbb P_k^2\), and put
\(\omega=\Omega_X^2\cong\mathcal O_X(-3)\).  Write
\(U_i=D_+(x_i)\) for the standard affine cover.  Under
\(\omega\otimes\mathcal T_X\cong\Omega_X^1\), the Cech cocycle

\[
\xi_{ij}=\frac{x_j}{x_i}\,d\!\left(\frac{x_i}{x_j}\right)
\in\Gamma(U_i\cap U_j,\Omega_X^1)
\tag{1}
\]

defines a square-zero extension \(X'\) of \(X\) by \(\omega\):

\[
0\longrightarrow\omega\longrightarrow\mathcal O_{X'}
\longrightarrow\mathcal O_X\longrightarrow0.
\]

The associated sequence of units yields

\[
H^1(X,\omega)\longrightarrow\operatorname{Pic}X'
\longrightarrow\operatorname{Pic}X
\xrightarrow{\delta}H^2(X,\omega).
\tag{2}
\]
\end{context}

\begin{problem}
Compute \(\delta(\mathcal O_X(1))\) in Cech cohomology and prove that
\(\delta\) is injective.  Deduce that
\(\operatorname{Pic}X'=0\), that \(X'\) has no ample invertible sheaf,
and hence that \(X'\) is not projective.
\end{problem}
\end{exerciseintro}

\begin{solution}
Choose local splittings

\[
\mathcal O_{X'}|_{U_i}\cong
\mathcal O_{U_i}\oplus\omega|_{U_i},
\]

where the second summand is square zero.  With the natural
identification used in (1), passage from the \(j\)-splitting to the
\(i\)-splitting on \(U_{ij}\) is

\[
(a,m)\longmapsto(a,m+D_{ij}(a)),
\qquad D_{ij}(a)=da\wedge\xi_{ij}.
\tag{3}
\]

The sign changes if one chooses the opposite convention for
\(\omega\otimes\mathcal T_X\cong\Omega_X^1\), but the nonvanishing
below is unaffected.

Represent \(\mathcal O_X(1)\) by

\[
g_{ij}=\frac{x_j}{x_i},
\qquad g_{01}g_{12}g_{20}=1.
\]

Lift \(g_{01},g_{12},g_{20}\) using respectively the splittings on
\(U_0,U_1,U_2\).  Expressing all three lifts in the \(U_0\)-splitting
and multiplying them with

\[
(a,m)(b,n)=(ab,an+bm)
\]

gives \((1,\eta)\), where

\[
\begin{split}
\eta
&=g_{20}^{-1}D_{02}(g_{20})
  +g_{12}^{-1}D_{01}(g_{12})\\
&=d\log(g_{20})\wedge\xi_{02}
  +d\log(g_{12})\wedge\xi_{01}.
\end{split}
\tag{4}
\]

Here \(d\log(h)\) is shorthand for \(h^{-1}dh\); no logarithm of a
regular function is being chosen.

The first term is zero because
\(\xi_{02}=d\log(x_0/x_2)=d\log(g_{20})\).  On \(U_0\), set

\[
u=\frac{x_1}{x_0},
\qquad v=\frac{x_2}{x_0}.
\]

Then the second term in (4) is

\[
d\log(v/u)\wedge d\log(1/u)
=\frac{du}{u}\wedge\frac{dv}{v}.
\tag{5}
\]

Thus \(\delta(\mathcal O_X(1))\) is represented, up to the harmless
overall sign just mentioned, by

\[
\frac{1}{uv}\,du\wedge dv
\in\Gamma(U_0\cap U_1\cap U_2,\omega).
\tag{6}
\]

We verify directly that this class is nonzero.  On the triple
intersection,

\[
\Gamma(U_{012},\omega)
=k[u^{\pm1},v^{\pm1}]\,du\wedge dv.
\]

The restrictions from \(U_{01}\) span the Laurent monomials
\(u^av^bdu\wedge dv\) with \(b\geq0\), and those from \(U_{02}\) span
the ones with \(a\geq0\).  On \(U_1\), use
\(s=x_0/x_1=u^{-1}\) and \(t=x_2/x_1=v/u\).  Since

\[
ds\wedge dt=-u^{-3}\,du\wedge dv,
\]

the restrictions from \(U_{12}\) span precisely the monomials with
\(a+b\leq-3\).  Modulo these three subspaces, the only remaining
Laurent monomial is
\(u^{-1}v^{-1}du\wedge dv\).  Consequently

\[
\check H^2(\{U_i\},\omega)
=k\left[\frac{du\wedge dv}{uv}\right]
\cong H^2(X,\omega),
\]

and (6) is a generator.  In particular,
\(\delta(\mathcal O_X(1))\neq0\).

Now
\(\operatorname{Pic}X\cong\mathbb Z[\mathcal O_X(1)]\), while
\(H^2(X,\omega)\cong(k,+)\).  Since \(\delta\) is a group
homomorphism,

\[
\delta(\mathcal O_X(n))=n\,\delta(\mathcal O_X(1)).
\]

Characteristic zero makes this zero only when \(n=0\), so \(\delta\)
is injective.  Also \(H^1(X,\omega)=H^1(X,\mathcal O_X(-3))=0\).
Exactness of (2) now gives
\(\operatorname{Pic}X'=0\).

If \(X'\) possessed an ample invertible sheaf, it would be trivial;
its restriction to the closed subscheme \(X\) would then be
\(\mathcal O_X\).  But the restriction of an ample sheaf is ample,
whereas \(\mathcal O_{\mathbb P^2}\) is not ample.  Thus \(X'\) has no
ample invertible sheaf and cannot be projective.
\end{solution}

\begin{exerciseintro}
% record: H-III-05-010
\exerciseheading{Exercise III.5.10}
\addcontentsline{toc}{subsection}{Exercise III.5.10}

\begin{context}
Let \(X\) be projective over a noetherian ring \(A\), with a fixed
very ample sheaf \(\mathcal O_X(1)\).  Kernels and images of morphisms
of coherent sheaves on \(X\) are coherent, and Serre vanishing holds
for every coherent sheaf.
\end{context}

\begin{problem}
Let

\[
\mathcal F^1\longrightarrow\mathcal F^2\longrightarrow\cdots
\longrightarrow\mathcal F^r
\]

be an exact sequence of coherent sheaves.  Prove that there is an
integer \(n_0\) such that, for every \(n\geq n_0\), the sequence

\[
\Gamma(X,\mathcal F^1(n))\longrightarrow
\Gamma(X,\mathcal F^2(n))\longrightarrow\cdots\longrightarrow
\Gamma(X,\mathcal F^r(n))
\]

is exact.
\end{problem}
\end{exerciseintro}

\begin{solution}
Write \(d_i:\mathcal F^i\to\mathcal F^{i+1}\), and put
\(\mathcal K_i=\ker d_i\).  Exactness of the given sequence says that

\[
\operatorname{im}d_{i-1}=\mathcal K_i
\qquad(2\leq i\leq r-1).
\]

Consequently, for each such \(i\) there is a short exact sequence of
coherent sheaves

\[
0\longrightarrow\mathcal K_{i-1}
\longrightarrow\mathcal F^{i-1}
\longrightarrow\mathcal K_i\longrightarrow0.
\tag{1}
\]

After twisting, the only obstruction to surjectivity of

\[
\Gamma(X,\mathcal F^{i-1}(n))
\longrightarrow\Gamma(X,\mathcal K_i(n))
\]

is \(H^1(X,\mathcal K_{i-1}(n))\).  By Serre vanishing this group is
zero for all sufficiently large \(n\).  Since there are only finitely
many indices, choose one \(n_0\) that works for every sequence (1).

For \(n\geq n_0\), left exactness of global sections gives

\[
\ker\Gamma(d_i(n))=\Gamma(X,\mathcal K_i(n)),
\]

while the preceding surjectivity gives

\[
\operatorname{im}\Gamma(d_{i-1}(n))
=\Gamma(X,\mathcal K_i(n)).
\]

Thus image equals kernel at every interior term, as required.  If
exactness at either endpoint is part of the original sequence, the
same argument gives surjectivity at the last term, while injectivity
at the first is preserved automatically by the left exact functor
\(\Gamma(X,-)\).
\end{solution}

\begin{exerciseintro}
\corpussection{III.6 Ext Groups and Sheaves}
\addcontentsline{toc}{section}{III.6 Ext Groups and Sheaves}

% record: H-III-06-001
\exerciseheading{Exercise III.6.1}
\addcontentsline{toc}{subsection}{Exercise III.6.1}

\begin{context}
Let \((X,\mathcal O_X)\) be a ringed space, and let
\(\mathcal A,\mathcal C\in\operatorname{Mod}(X)\).  An extension of
\(\mathcal C\) by \(\mathcal A\) is a short exact sequence

\[
0\longrightarrow\mathcal A\longrightarrow\mathcal E
\longrightarrow\mathcal C\longrightarrow0.
\]

Two extensions are equivalent if an isomorphism between their middle
terms induces the identity on \(\mathcal A\) and \(\mathcal C\).
\end{context}

\begin{problem}
For an extension as above, apply
\(\operatorname{Hom}(\mathcal C,-)\) and let

\[
\xi_{\mathcal E}
=\delta(1_{\mathcal C})
\in\operatorname{Ext}^1(\mathcal C,\mathcal A).
\]

Prove that
\(\mathcal E\mapsto\xi_{\mathcal E}\) is a bijection from equivalence
classes of extensions of \(\mathcal C\) by \(\mathcal A\) to
\(\operatorname{Ext}^1(\mathcal C,\mathcal A)\).
\end{problem}
\end{exerciseintro}

\begin{solution}
Choose an injective sheaf \(\mathcal I\) containing \(\mathcal A\),
and set \(\mathcal Q=\mathcal I/\mathcal A\).  From

\[
0\longrightarrow\mathcal A\longrightarrow\mathcal I
\xrightarrow{p}\mathcal Q\longrightarrow0
\tag{1}
\]

and the injectivity of \(\mathcal I\), we obtain an exact sequence

\[
\operatorname{Hom}(\mathcal C,\mathcal I)
\longrightarrow\operatorname{Hom}(\mathcal C,\mathcal Q)
\xrightarrow{\partial}
\operatorname{Ext}^1(\mathcal C,\mathcal A)
\longrightarrow0.
\tag{2}
\]

Given \(u:\mathcal C\to\mathcal Q\), form the pullback

\[
\mathcal E_u=\mathcal I\times_{\mathcal Q}\mathcal C
=\{(i,c):p(i)=u(c)\}.
\]

It fits into an exact sequence

\[
0\longrightarrow\mathcal A\longrightarrow\mathcal E_u
\longrightarrow\mathcal C\longrightarrow0.
\tag{3}
\]

The boundary map in (2) sends \(u\) to the pullback of (1) along
\(u\), namely (3).  Equivalently, the connecting map belonging to
(3) sends \(1_{\mathcal C}\) to \(\partial(u)\).

If \(u'=u+p\circ v\) for some \(v:\mathcal C\to\mathcal I\), then

\[
(i,c)\longmapsto(i+v(c),c)
\]

is an equivalence \(\mathcal E_u\cong\mathcal E_{u'}\).  Thus (2)
shows that the construction depends only on the class in
\(\operatorname{Ext}^1(\mathcal C,\mathcal A)\), and every such class
is represented by an extension.

It remains to show that no extensions have been identified
incorrectly.  Start with an arbitrary extension

\[
0\to\mathcal A\to\mathcal E\xrightarrow{\pi}\mathcal C\to0.
\tag{4}
\]

Injectivity of \(\mathcal I\) extends the inclusion
\(\mathcal A\hookrightarrow\mathcal I\) to a map
\(s:\mathcal E\to\mathcal I\).  The composite \(p\circ s\) vanishes
on \(\mathcal A\), so it descends uniquely through \(\pi\) to a map
\(u:\mathcal C\to\mathcal Q\).  There is then a commutative diagram
of short exact sequences

\[
\begin{array}{ccccccccc}
0&\to&\mathcal A&\to&\mathcal E&\xrightarrow{\pi}&\mathcal C&\to&0\\
 &&\Vert&&\downarrow&&\Vert\\
0&\to&\mathcal A&\to&\mathcal E_u&\to&\mathcal C&\to&0,
\end{array}
\]

where the middle vertical map is \(e\mapsto(s(e),\pi(e))\).
Since the outer vertical maps are identities, the middle one is an
isomorphism.  Hence (4) is equivalent to \(\mathcal E_u\).

A different choice of \(s\) changes \(u\) by an element in the image
of \(\operatorname{Hom}(\mathcal C,\mathcal I)\), so it has the same
class in (2).  The two constructions are therefore mutually inverse,
which proves the bijection.
\end{solution}

\begin{exerciseintro}
% record: H-III-06-002
\exerciseheading{Exercise III.6.2}
\addcontentsline{toc}{subsection}{Exercise III.6.2}

\begin{context}
Let \(X=\mathbb P_k^1\), where \(k\) is infinite.  Write
\(\operatorname{Mod}(X)\), \(\operatorname{Qco}(X)\), and
\(\operatorname{Coh}(X)\) for the categories of all
\(\mathcal O_X\)-modules, quasi-coherent modules, and coherent
modules, respectively.
\end{context}

\begin{problem}
(a) Prove that there is no projective object
\(\mathcal P\in\operatorname{Mod}(X)\) admitting an epimorphism
\(\mathcal P\twoheadrightarrow\mathcal O_X\).

(b) Prove the same assertion when projective means projective in
either \(\operatorname{Qco}(X)\) or \(\operatorname{Coh}(X)\).
\end{problem}
\end{exerciseintro}

\begin{solution}
(a) Suppose that
\(\alpha:\mathcal P\twoheadrightarrow\mathcal O_X\) exists.  Choose a
closed \(k\)-rational point \(x\).  Surjectivity on the stalk at \(x\)
gives a neighborhood \(U\) of \(x\) and
\(s\in\Gamma(U,\mathcal P)\) whose image has residue \(1\) at \(x\).
Because \(k\) is infinite, choose another \(k\)-rational point
\(y\in U\setminus\{x\}\), put \(V=X\setminus\{y\}\), and let
\(j:V\hookrightarrow X\).

There is a stalkwise surjection

\[
j_!\mathcal O_V\twoheadrightarrow i_{x*}k(x).
\tag{1}
\]

The composite
\(\mathcal P\xrightarrow{\alpha}\mathcal O_X\twoheadrightarrow
i_{x*}k(x)\) must lift through (1), since \(\mathcal P\) is
projective.  Let
\(\varphi:\mathcal P\to j_!\mathcal O_V\) be such a lift.  Then
\(\varphi(s)\in\Gamma(U,j_!\mathcal O_V)\) has nonzero germ at \(x\).

On the other hand,

\[
\Gamma(U,j_!\mathcal O_V)=0.
\tag{2}
\]

Indeed, a section on the left is a regular function on
\(U\cap V=U\setminus\{y\}\) whose support is closed in \(U\) and
avoids \(y\).  If it were nonzero, its support would contain the
generic point of the integral open set \(U\), and its closure in
\(U\) would be all of \(U\), including \(y\).  This contradicts the
defining support condition for extension by zero.  Equation (2)
contradicts the nonzero germ of \(\varphi(s)\).

(b) Suppose now that \(\mathcal P\) is projective in either of the two
stated subcategories and that
\(\alpha:\mathcal P\twoheadrightarrow\mathcal O_X\).  Fix a rational
point \(x\).  On some open neighborhood \(U\) there is
\(s\in\Gamma(U,\mathcal P)\) whose image in \(k(x)\) is nonzero.
After shrinking \(U\) to a principal open, write
\(U=X\setminus D\), where \(D\) is an effective Cartier divisor.
Let \(\mathcal L=\mathcal O_X(D)\) and let
\(f\in\Gamma(X,\mathcal L)\) be its canonical section, which is
invertible on \(U\).

The extension lemma for quasi-coherent sheaves says that for some
\(N\geq0\), the section \(f^Ns\) extends to

\[
t\in\Gamma(X,\mathcal P\otimes\mathcal L^N).
\tag{3}
\]

Its image in the fiber \(\mathcal L_x^N\otimes k(x)\) is nonzero.
Set

\[
\mathcal E=\mathcal L^{-N}\otimes\mathcal O_X(-1).
\]

After choosing an identification of its one-dimensional fiber at
\(x\) with \(k(x)\), there is an epimorphism in both
\(\operatorname{Coh}(X)\) and \(\operatorname{Qco}(X)\),

\[
\mathcal E\twoheadrightarrow i_{x*}k(x).
\tag{4}
\]

Projectivity of \(\mathcal P\) lifts
\(\mathcal P\to\mathcal O_X\to i_{x*}k(x)\) through (4), producing
\(\varphi:\mathcal P\to\mathcal E\).  Tensoring by
\(\mathcal L^N\), the section (3) maps to a global section of

\[
\mathcal E\otimes\mathcal L^N\cong\mathcal O_X(-1).
\]

That global section is zero.  Its value in the fiber at \(x\) is
therefore zero, whereas commutativity with (4) says that it is the
nonzero image of \(t\).  This contradiction proves the assertion in
both categories.
\end{solution}

\begin{exerciseintro}
% record: H-III-06-003
\exerciseheading{Exercise III.6.3}
\addcontentsline{toc}{subsection}{Exercise III.6.3}

\begin{context}
Let \(X\) be a noetherian scheme.  For sheaves
\(\mathcal F,\mathcal G\in\operatorname{Mod}(X)\), write
\(\mathcal E\!xt_X^i(\mathcal F,\mathcal G)\) for the sheaf Ext.
Sheaf Ext is local on \(X\), and a locally free resolution of finite
rank may be used to compute it.
\end{context}

\begin{problem}
(a) If \(\mathcal F\) and \(\mathcal G\) are coherent, prove that
\(\mathcal E\!xt_X^i(\mathcal F,\mathcal G)\) is coherent for every
\(i\geq0\).

(b) If \(\mathcal F\) is coherent and \(\mathcal G\) is
quasi-coherent, prove that
\(\mathcal E\!xt_X^i(\mathcal F,\mathcal G)\) is quasi-coherent for
every \(i\geq0\).
\end{problem}
\end{exerciseintro}

\begin{solution}
Both assertions are local.  On an affine open
\(U=\operatorname{Spec}A\), the coherent sheaf
\(\mathcal F|_U\) corresponds to a finite \(A\)-module.  Since \(A\)
is noetherian, repeatedly taking kernels of surjections from finite
free modules gives a resolution

\[
\cdots\longrightarrow\mathcal E_2\longrightarrow\mathcal E_1
\longrightarrow\mathcal E_0\longrightarrow\mathcal F|_U
\longrightarrow0
\tag{1}
\]

in which every \(\mathcal E_j\) is locally free of finite rank.
Applying \(\mathcal H\!om(-,\mathcal G|_U)\) gives the complex

\[
0\longrightarrow
\mathcal E_0^\vee\otimes\mathcal G|_U
\longrightarrow\mathcal E_1^\vee\otimes\mathcal G|_U
\longrightarrow\cdots,
\tag{2}
\]

whose \(i\)-th cohomology sheaf is
\(\mathcal E\!xt_X^i(\mathcal F,\mathcal G)|_U\).

(a) If \(\mathcal G\) is coherent, every term of (2) is coherent.
Kernels, images, and cokernels of maps of coherent sheaves are
coherent on a noetherian scheme, so every cohomology sheaf of (2) is
coherent.

(b) If \(\mathcal G\) is quasi-coherent, every term of (2) is
quasi-coherent.  Kernels and cokernels of morphisms of
quasi-coherent sheaves on a scheme are quasi-coherent.  Hence the
cohomology sheaves of (2), and therefore all the stated sheaf Exts,
are quasi-coherent.
\end{solution}

\begin{exerciseintro}
% record: H-III-06-004
\exerciseheading{Exercise III.6.4}
\addcontentsline{toc}{subsection}{Exercise III.6.4}

\begin{context}
Let \(X\) be noetherian, and suppose that every coherent sheaf is a
quotient of a locally free coherent sheaf.  In this situation we say
that \(\operatorname{Coh}(X)\) has enough locally free sheaves.  Fix
\(\mathcal G\in\operatorname{Mod}(X)\).
\end{context}

\begin{problem}
Prove that

\[
\bigl(\mathcal E\!xt_X^i(-,\mathcal G)\bigr)_{i\geq0}
\]

is a contravariant universal delta functor from
\(\operatorname{Coh}(X)\) to \(\operatorname{Mod}(X)\).
\end{problem}
\end{exerciseintro}

\begin{solution}
The long exact sheaf-Ext sequence makes the displayed family a
contravariant delta functor.  By the standard universality criterion,
it is enough to prove that its positive-degree terms are
coeffaceable: for each coherent \(\mathcal F\), there must be an
epimorphism \(\mathcal E\twoheadrightarrow\mathcal F\) such that the
induced maps

\[
\mathcal E\!xt_X^i(\mathcal F,\mathcal G)
\longrightarrow
\mathcal E\!xt_X^i(\mathcal E,\mathcal G)
\]

vanish for all \(i>0\).

By hypothesis, choose such an epimorphism with \(\mathcal E\) locally
free of finite rank.  Then

\[
\mathcal H\!om_X(\mathcal E,-)
\cong\mathcal E^\vee\otimes_{\mathcal O_X}-
\]

is exact.  Consequently

\[
\mathcal E\!xt_X^i(\mathcal E,\mathcal G)=0
\qquad(i>0),
\]

so all the maps above have zero target.  The delta functor is
coeffaceable and hence universal.
\end{solution}

\begin{exerciseintro}
% record: H-III-06-005
\exerciseheading{Exercise III.6.5}
\addcontentsline{toc}{subsection}{Exercise III.6.5}

\begin{context}
Let \(X\) be noetherian and assume that
\(\operatorname{Coh}(X)\) has enough locally free sheaves.  For a
coherent sheaf \(\mathcal F\), its homological dimension
\(\operatorname{hd}(\mathcal F)\) is the least length of a locally
free resolution, or \(+\infty\) if no finite resolution exists.
\end{context}

\begin{problem}
(a) Prove that \(\mathcal F\) is locally free if and only if

\[
\mathcal E\!xt_X^1(\mathcal F,\mathcal G)=0
\]

for every \(\mathcal G\in\operatorname{Mod}(X)\).

(b) Prove that \(\operatorname{hd}(\mathcal F)\leq n\) if and only if

\[
\mathcal E\!xt_X^i(\mathcal F,\mathcal G)=0
\qquad(i>n)
\]

for every \(\mathcal G\in\operatorname{Mod}(X)\).

(c) Prove

\[
\operatorname{hd}(\mathcal F)
=\sup_{x\in X}
\operatorname{pd}_{\mathcal O_{X,x}}\mathcal F_x.
\]
\end{problem}
\end{exerciseintro}

\begin{solution}
(a) If \(\mathcal F\) is locally free of finite rank, then
\(\mathcal H\!om(\mathcal F,-)\cong\mathcal F^\vee\otimes-\) is
exact, so every positive sheaf Ext vanishes.

Conversely, choose an exact sequence

\[
0\longrightarrow\mathcal K\longrightarrow\mathcal E
\xrightarrow{q}\mathcal F\longrightarrow0
\tag{1}
\]

with \(\mathcal E\) locally free.  Apply
\(\mathcal H\!om(\mathcal F,-)\).  The relevant part of the long
exact sequence is

\[
\mathcal H\!om(\mathcal F,\mathcal E)
\longrightarrow\mathcal H\!om(\mathcal F,\mathcal F)
\longrightarrow\mathcal E\!xt_X^1(\mathcal F,\mathcal K).
\]

The last sheaf is zero by hypothesis.  Hence the identity of
\(\mathcal F\) lifts locally through \(q\), so (1) is locally split.
Thus \(\mathcal F\) is locally a direct summand of a finite locally
free sheaf.  Such a coherent direct summand is locally free.

(b) If \(\operatorname{hd}(\mathcal F)\leq n\), apply
\(\mathcal H\!om(-,\mathcal G)\) to a length-\(n\) locally free
resolution.  The resulting complex has no terms in degrees greater
than \(n\), so its cohomology gives

\[
\mathcal E\!xt_X^i(\mathcal F,\mathcal G)=0
\qquad(i>n).
\]

For the converse, set \(\mathcal K_0=\mathcal F\).  Using enough
locally free sheaves, construct exact sequences

\[
0\longrightarrow\mathcal K_{j+1}\longrightarrow\mathcal E_j
\longrightarrow\mathcal K_j\longrightarrow0,
\qquad 0\leq j<n,
\tag{2}
\]

with each \(\mathcal E_j\) locally free and each
\(\mathcal K_j\) coherent.  Dimension shifting in (2) gives, for
every \(\mathcal G\),

\[
\mathcal E\!xt_X^1(\mathcal K_n,\mathcal G)
\cong
\mathcal E\!xt_X^{n+1}(\mathcal F,\mathcal G)=0.
\]

Part (a) shows that \(\mathcal K_n\) is locally free.  Splicing the
sequences (2) therefore gives a locally free resolution of
\(\mathcal F\) of length at most \(n\).  The case \(n=0\) is exactly
part (a).

(c) Any locally free resolution of \(\mathcal F\) remains a free
resolution after taking a stalk.  Therefore

\[
\sup_{x\in X}\operatorname{pd}_{\mathcal O_{X,x}}\mathcal F_x
\leq\operatorname{hd}(\mathcal F).
\tag{3}
\]

Let the supremum on the left be \(d<\infty\).  For a coherent first
argument, sheaf Ext has the stalk formula

\[
\bigl(\mathcal E\!xt_X^i(\mathcal F,\mathcal G)\bigr)_x
\cong
\operatorname{Ext}_{\mathcal O_{X,x}}^i
(\mathcal F_x,\mathcal G_x).
\tag{4}
\]

If \(i>d\), the right side of (4) vanishes for every \(x\) and every
\(\mathcal G\).  Hence the sheaf on the left is zero.  Part (b) gives
\(\operatorname{hd}(\mathcal F)\leq d\), which together with (3)
proves equality.  If the supremum is infinite, (3) already proves
the assertion.
\end{solution}

\begin{exerciseintro}
% record: H-III-06-006
\exerciseheading{Exercise III.6.6}
\addcontentsline{toc}{subsection}{Exercise III.6.6}

\begin{context}
Let \((A,\mathfrak m)\) be a regular local ring and let \(M\) be a
finite \(A\)-module.  We may use that \(A\) has finite global
dimension

\[
\operatorname{gldim}A=\dim A.
\]

In particular, every finite \(A\)-module has finite projective
dimension.
\end{context}

\begin{problem}
(a) Prove that \(M\) is projective if and only if

\[
\operatorname{Ext}_A^i(M,A)=0
\qquad(i>0).
\]

(b) For every \(n\geq0\), prove that

\[
\operatorname{pd}_A M\leq n
\quad\Longleftrightarrow\quad
\operatorname{Ext}_A^i(M,A)=0
\quad(i>n).
\]
\end{problem}
\end{exerciseintro}

\begin{solution}
(a) The forward implication is immediate.  Conversely, choose a
finite free presentation

\[
0\longrightarrow N\longrightarrow A^r\longrightarrow M
\longrightarrow0.
\tag{1}
\]

It is enough to prove that (1) splits, or equivalently that
\(\operatorname{Ext}_A^1(M,N)=0\).  Resolve \(N\) by finite free
modules, writing

\[
0\longrightarrow N_{j+1}\longrightarrow A^{r_j}
\longrightarrow N_j\longrightarrow0,
\qquad N_0=N.
\tag{2}
\]

The hypothesis gives
\(\operatorname{Ext}_A^i(M,A^{r_j})=0\) for every \(i>0\).
The long exact sequences from (2) therefore give isomorphisms

\[
\operatorname{Ext}_A^1(M,N)
\cong\operatorname{Ext}_A^2(M,N_1)
\cong\cdots\cong
\operatorname{Ext}_A^{d+1}(M,N_d),
\]

where \(d=\dim A\).  The last group is zero because
\(\operatorname{pd}_A M\leq d\).  Thus (1) splits and \(M\) is
projective.  Since \(A\) is local, it is in fact free.

(b) If \(\operatorname{pd}_A M\leq n\), the asserted Ext vanishing
follows by applying \(\operatorname{Hom}_A(-,A)\) to a projective
resolution of length \(n\).

Conversely, take the first \(n\) steps of a finite free resolution:

\[
0\longrightarrow K_n\longrightarrow F_{n-1}\longrightarrow\cdots
\longrightarrow F_0\longrightarrow M\longrightarrow0.
\tag{3}
\]

For \(j>0\), dimension shifting in (3) gives

\[
\operatorname{Ext}_A^j(K_n,A)
\cong\operatorname{Ext}_A^{j+n}(M,A)=0.
\]

Part (a) shows that \(K_n\) is projective.  Thus (3) is a projective
resolution of length \(n\), so \(\operatorname{pd}_A M\leq n\).
For \(n=0\), this is precisely part (a).
\end{solution}

\begin{exerciseintro}
% record: H-III-06-007
\exerciseheading{Exercise III.6.7}
\addcontentsline{toc}{subsection}{Exercise III.6.7}

\begin{context}
Let \(X=\operatorname{Spec}A\) be noetherian.  Let \(M,N\) be
\(A\)-modules, with \(M\) finite, and let
\(\widetilde M,\widetilde N\) be their associated quasi-coherent
sheaves.
\end{context}

\begin{problem}
Prove naturally, for every \(i\geq0\),

\[
\operatorname{Ext}_X^i(\widetilde M,\widetilde N)
\cong\operatorname{Ext}_A^i(M,N)
\]

and

\[
\mathcal E\!xt_X^i(\widetilde M,\widetilde N)
\cong\widetilde{\operatorname{Ext}_A^i(M,N)}.
\]
\end{problem}
\end{exerciseintro}

\begin{solution}
Because \(A\) is noetherian and \(M\) is finite, there is a free
resolution

\[
\cdots\longrightarrow P_2\longrightarrow P_1\longrightarrow P_0
\longrightarrow M\longrightarrow0
\tag{1}
\]

with every \(P_j\) finite free.  Sheafifying (1) gives a locally free
resolution of \(\widetilde M\).  Applying sheaf Hom into
\(\widetilde N\) produces

\[
\mathcal H\!om_X(\widetilde P_j,\widetilde N)
\cong\widetilde{\operatorname{Hom}_A(P_j,N)}
\tag{2}
\]

in every degree.  Sheafification is exact, so it commutes with the
cohomology of this complex.  Therefore

\[
\mathcal E\!xt_X^i(\widetilde M,\widetilde N)
\cong
\widetilde{H^i(\operatorname{Hom}_A(P_\bullet,N))}
=\widetilde{\operatorname{Ext}_A^i(M,N)}.
\tag{3}
\]

Now use the local-to-global Ext spectral sequence

\[
E_2^{p,q}
=H^p\!\left(
X,\mathcal E\!xt_X^q(\widetilde M,\widetilde N)
\right)
\Longrightarrow
\operatorname{Ext}_X^{p+q}(\widetilde M,\widetilde N).
\tag{4}
\]

The sheaves in (3) are quasi-coherent, and higher cohomology of a
quasi-coherent sheaf on the affine scheme \(X\) vanishes.  Thus (4)
has only its \(p=0\) column and degenerates.  Taking global sections in
(3) gives

\[
\operatorname{Ext}_X^i(\widetilde M,\widetilde N)
\cong
\Gamma\!\left(X,
\widetilde{\operatorname{Ext}_A^i(M,N)}\right)
\cong\operatorname{Ext}_A^i(M,N).
\]

Every construction used above is functorial in \(M\) and \(N\), so
the isomorphisms are natural.
\end{solution}

\begin{exerciseintro}
% record: H-III-06-008
\exerciseheading{Exercise III.6.8}
\addcontentsline{toc}{subsection}{Exercise III.6.8}

\begin{context}
Let \(X\) be a noetherian, integral, separated, locally factorial
scheme.  If \(\mathcal L\) is invertible and
\(s\in\Gamma(X,\mathcal L)\), write

\[
X_s=\{x\in X:s_x\text{ generates }\mathcal L_x\}.
\]

This is the complement of the zero locus of \(s\).
\end{context}

\begin{problem}
(a) Prove that the opens \(X_s\), as \(\mathcal L\) and \(s\) vary,
form a basis for the topology of \(X\).

(b) Deduce that every coherent sheaf on \(X\) is a quotient of a
locally free sheaf of finite rank.
\end{problem}
\end{exerciseintro}

\begin{solution}
(a) We first prove a lemma: if \(W\subseteq X\) is a nonempty affine
open, every irreducible component of \(X\setminus W\) has codimension
one.

Let \(Z\) be such a component, with generic point \(\eta\), and choose
an affine neighborhood \(V\) of \(\eta\).  Because \(X\) is
separated, \(V\cap W\) is affine.  Put
\(R=\mathcal O_{X,\eta}\), with maximal ideal \(\mathfrak m\).
The inverse image of \(V\cap W\) in \(\operatorname{Spec}R\) is

\[
\operatorname{Spec}R\setminus\{\mathfrak m\}.
\tag{1}
\]

Indeed, a point of the local spectrum other than its closed point
corresponds to a proper generalization of \(\eta\).  If that
generalization also lay in \(X\setminus W\), its closure would be an
irreducible subset of \(X\setminus W\) strictly containing \(Z\),
contrary to the maximality of \(Z\).  The open in (1) is affine,
because it is obtained by base change from the affine morphism
\(V\cap W\to V\).

The local ring \(R\) is a unique factorization domain, hence normal.
If \(\dim R\geq2\), the Hartogs property for a noetherian normal local
domain gives

\[
\Gamma(\operatorname{Spec}R\setminus\{\mathfrak m\},
\mathcal O)=R.
\tag{2}
\]

If the punctured spectrum were affine, (2) would identify it with
\(\operatorname{Spec}R\); under this identification its open
immersion into \(\operatorname{Spec}R\) would be induced by the
identity of \(R\), an impossibility because the closed point is
missing.  Hence \(\dim R<2\).  The point \(\eta\) is not the generic
point of \(X\), since every nonempty open contains that point, so
\(\dim R=1\).  This proves the lemma.

Now let \(x\in U\), with \(U\) open.  Choose an affine open
\(W\) such that \(x\in W\subseteq U\).  The closed set
\(X\setminus W\) has finitely many irreducible components
\(D_1,\ldots,D_r\), and the lemma says that each \(D_j\) is a prime
Weil divisor.  Since \(X\) is locally factorial, every \(D_j\) is
Cartier.  Thus

\[
D=D_1+\cdots+D_r
\]

is an effective Cartier divisor.  Its canonical section
\(s_D\in\Gamma(X,\mathcal O_X(D))\) has zero locus
\(\lvert D\rvert=X\setminus W\), and consequently

\[
X_{s_D}=W.
\]

This gives \(x\in X_{s_D}\subseteq U\), proving the basis assertion.

(b) Let \(\mathcal F\) be coherent.  By part (a), noetherian
quasi-compactness, and coherence, choose finitely many affine basis
opens

\[
U_i=X_{s_i},\qquad
s_i\in\Gamma(X,\mathcal L_i),
\]

together with finitely many sections
\(t_{ij}\in\Gamma(U_i,\mathcal F)\) that generate
\(\mathcal F|_{U_i}\).  The extension lemma for sections says that,
after replacing an exponent by a sufficiently large integer
\(n_{ij}\), the section

\[
s_i^{\,n_{ij}}t_{ij}
\in\Gamma(U_i,\mathcal F\otimes\mathcal L_i^{n_{ij}})
\]

extends to a global section.  It defines a morphism

\[
\mathcal L_i^{-n_{ij}}\longrightarrow\mathcal F.
\tag{3}
\]

On \(U_i\), multiplication by \(s_i^{n_{ij}}\) is a unit, so the
images of the maps (3) generate exactly when the \(t_{ij}\) do.
Therefore the combined map

\[
\bigoplus_{i,j}\mathcal L_i^{-n_{ij}}
\longrightarrow\mathcal F
\]

is surjective.  Its source is a finite direct sum of invertible
sheaves, hence is locally free of finite rank.
\end{solution}

\begin{exerciseintro}
% record: H-III-06-009
\exerciseheading{Exercise III.6.9}
\addcontentsline{toc}{subsection}{Exercise III.6.9}

\begin{context}
Let \(X\) be a noetherian, integral, separated, regular scheme.
Let \(K(X)\) be the Grothendieck group of coherent sheaves.  Define
\(K_1(X)\) from finite-rank locally free sheaves, with the relations

\[
[\mathcal E]_1=[\mathcal E']_1+[\mathcal E'']_1
\]

for every short exact sequence
\(0\to\mathcal E'\to\mathcal E\to\mathcal E''\to0\) of locally free
sheaves.  There is a natural homomorphism

\[
\epsilon:K_1(X)\longrightarrow K(X).
\]
The notation \(K_1(X)\) here is the book's group of vector
bundles with exact-sequence relations, usually denoted
\(K_0\) of vector bundles; it is not higher algebraic \(K_1\).
\end{context}

\begin{problem}
(a) Prove that every coherent sheaf on \(X\) has a finite locally
free resolution.

(b) If

\[
0\to\mathcal E_n\to\cdots\to\mathcal E_0\to\mathcal F\to0
\]

is such a resolution, set

\[
\delta(\mathcal F)=\sum_{i=0}^n(-1)^i[\mathcal E_i]_1.
\]

Prove that this class is independent of the resolution, is additive
in short exact sequences of coherent sheaves, and induces an inverse
\(\delta:K(X)\to K_1(X)\) to \(\epsilon\).
\end{problem}
\end{exerciseintro}

\begin{solution}
(a) A regular local ring is factorial, so \(X\) is locally factorial.
Exercise III.6.8 therefore shows that coherent sheaves on \(X\) have
enough locally free sheaves.  Starting with a coherent
\(\mathcal F=\mathcal K_0\), construct exact sequences

\[
0\longrightarrow\mathcal K_{i+1}\longrightarrow\mathcal E_i
\longrightarrow\mathcal K_i\longrightarrow0
\tag{1}
\]

with each \(\mathcal E_i\) locally free of finite rank.  All
\(\mathcal K_i\) are coherent.

Let

\[
C_i=\{x\in X:(\mathcal K_i)_x
\text{ is not a free }\mathcal O_{X,x}\text{-module}\}.
\]

The locally free locus of a coherent sheaf is open, so \(C_i\) is
closed.  If \((\mathcal K_i)_x\) is free, the stalk of (1) splits and
\((\mathcal K_{i+1})_x\) is a finite projective module over a local
ring, hence free.  Thus

\[
C_0\supseteq C_1\supseteq C_2\supseteq\cdots.
\tag{2}
\]

At every \(x\), the regular local ring \(\mathcal O_{X,x}\) has finite
global dimension.  The stalk resolution therefore reaches a free
syzygy, so \(x\notin C_i\) for some \(i\).  Hence
\(\bigcap_iC_i=\varnothing\).  Since the topological space of \(X\)
is noetherian, the descending chain (2) stabilizes; its stable value
must be empty.  Thus \(\mathcal K_N\) is locally free for some \(N\),
and (1) splices to the required finite resolution.

(b) We first record that a bounded exact complex of locally free
sheaves has alternating class zero in \(K_1(X)\).  Starting at one
end, its cycle sheaves are locally free: each is the kernel of a
surjection of locally free sheaves onto a locally free cycle, and
that sequence splits locally.  The short exact sequences involving
successive cycles then make all terms cancel in the alternating sum.

Next let \(\mathcal E_\bullet\to\mathcal F\) and
\(\mathcal E'_\bullet\to\mathcal F\) be two finite locally free
resolutions.  They admit a common locally free refinement
\(\mathcal A_\bullet\to\mathcal F\) with augmentation-preserving
chain maps

\[
\mathcal A_\bullet\longrightarrow\mathcal E_\bullet,
\qquad
\mathcal A_\bullet\longrightarrow\mathcal E'_\bullet.
\tag{3}
\]

Here is the construction.  In degree zero, choose a locally free
sheaf mapping onto the coherent fiber product
\(\mathcal E_0\times_{\mathcal F}\mathcal E'_0\).  Inductively, after
the lower degrees have been constructed, take the fiber product of
the two lifting problems over the current syzygy and choose a locally
free sheaf surjecting onto it.  This simultaneously continues the
resolution and both chain maps.  Part (a), applied to the coherent
syzygies through the same descending nonfree-locus argument, shows
that the current syzygy is locally free after finitely many stages.
Continue until both original finite resolutions have also ended,
and only then use the syzygy as the final term.  At this degree
its required maps to the original resolutions are zero, and the
chain identities already follow from exactness there.
This gives a finite truncation without assuming global lifting
from a locally free sheaf into an unfinished resolution.

Each map in (3) is a quasi-isomorphism, so its mapping cone is a
bounded exact complex of locally free sheaves.  The preceding
observation applied to the two cones gives

\[
\sum_i(-1)^i[\mathcal E_i]_1
=\sum_i(-1)^i[\mathcal A_i]_1
=\sum_i(-1)^i[\mathcal E'_i]_1.
\tag{4}
\]

Thus \(\delta(\mathcal F)\) is independent of all choices.

Now take a short exact sequence

\[
0\longrightarrow\mathcal F'\longrightarrow\mathcal F
\longrightarrow\mathcal F''\longrightarrow0.
\tag{5}
\]

A fiber-product version of the horseshoe construction gives bounded
locally free resolutions in a termwise short exact sequence

\[
0\longrightarrow\mathcal E'_\bullet
\longrightarrow\mathcal E_\bullet
\longrightarrow\mathcal E''_\bullet
\longrightarrow0
\tag{6}
\]

resolving (5).  For completeness, in degree zero choose a locally
free \(\mathcal E''_0\twoheadrightarrow\mathcal F''\), form
\(\mathcal F\times_{\mathcal F''}\mathcal E''_0\), and choose a
locally free \(\mathcal E_0\) surjecting onto that fiber product.
Then \(\mathcal E_0\to\mathcal E''_0\) is surjective and its kernel
\(\mathcal E'_0\) is locally free; the induced map
\(\mathcal E'_0\to\mathcal F'\) is surjective.  The snake lemma gives
a new short exact sequence of syzygies, and the construction repeats.
The nonfree-locus argument from part (a), applied simultaneously to
the three syzygies, makes the resulting resolutions finite.

Each degree of (6) gives

\[
[\mathcal E_i]_1=[\mathcal E'_i]_1+[\mathcal E''_i]_1.
\]

Taking alternating sums and using (4) yields

\[
\delta(\mathcal F)
=\delta(\mathcal F')+\delta(\mathcal F'').
\]

Therefore \(\delta\) descends to a homomorphism
\(K(X)\to K_1(X)\).

If \(\mathcal E\) is locally free, its length-zero resolution gives
\(\delta(\epsilon[\mathcal E]_1)=[\mathcal E]_1\), so
\(\delta\circ\epsilon\) is the identity.  Conversely, exactness of a
finite resolution gives in \(K(X)\)

\[
[\mathcal F]=\sum_i(-1)^i[\mathcal E_i].
\]

Hence \(\epsilon\circ\delta\) is also the identity, and \(\epsilon\)
is an isomorphism.
\end{solution}

\begin{exerciseintro}
% record: H-III-06-010
\exerciseheading{Exercise III.6.10}
\addcontentsline{toc}{subsection}{Exercise III.6.10}

\begin{context}
Let \(f:X\to Y\) be a finite morphism of noetherian schemes, and put

\[
\mathcal B=f_*\mathcal O_X.
\]

This is a coherent \(\mathcal O_Y\)-algebra.  Quasi-coherent
\(\mathcal O_X\)-modules are equivalent to quasi-coherent
\(\mathcal B\)-modules on \(Y\).
The derived comparison used here is the equivalence
between the derived category of quasi-coherent modules and the
subcategory of complexes with quasi-coherent cohomology on a
noetherian scheme (Stacks Project, Proposition 36.8.3, tag 09T1).
Thus one may resolve inside quasi-coherent modules and compare Ext
afterwards.
\end{context}

\begin{problem}
(a) For a quasi-coherent \(\mathcal O_Y\)-module \(\mathcal G\), show
that

\[
\mathcal H\!om_Y(\mathcal B,\mathcal G)
\]

is a quasi-coherent \(\mathcal B\)-module.  Denote the corresponding
quasi-coherent sheaf on \(X\) by \(f^!\mathcal G\).

(b) If \(\mathcal F\) is coherent on \(X\), prove the natural
isomorphism

\[
f_*\mathcal H\!om_X(\mathcal F,f^!\mathcal G)
\xrightarrow{\sim}
\mathcal H\!om_Y(f_*\mathcal F,\mathcal G).
\]

(c) Construct natural maps, for \(i\geq0\),

\[
\varphi_i:
\operatorname{Ext}_X^i(\mathcal F,f^!\mathcal G)
\longrightarrow
\operatorname{Ext}_Y^i(f_*\mathcal F,\mathcal G).
\]

(d) Assume in addition that \(X\) and \(Y\) are separated, that
\(\operatorname{Coh}(X)\) has enough locally free sheaves, and that
\(\mathcal B\) is locally free on \(Y\), equivalently that \(f\) is
flat.  Prove that every \(\varphi_i\) is an isomorphism.
\end{problem}
\end{exerciseintro}

\begin{solution}
(a) Work over an affine open \(U=\operatorname{Spec}A\subseteq Y\),
and write \(f^{-1}U=\operatorname{Spec}B\) and
\(\mathcal G|_U=\widetilde N\).  The \(A\)-module \(B\) is finite and
therefore finitely presented.  Consequently localization commutes
with \(\operatorname{Hom}_A(B,N)\), and

\[
\mathcal H\!om_Y(\mathcal B,\mathcal G)|_U
\cong\widetilde{\operatorname{Hom}_A(B,N)}.
\]

This proves quasi-coherence.  Give this Hom module the \(B\)-action

\[
(b\cdot\lambda)(b')=\lambda(bb').
\tag{1}
\]

These local actions glue to a \(\mathcal B\)-module structure, which
corresponds to the asserted sheaf \(f^!\mathcal G\) on \(X\).

(b) The assertion is the sheaf form of coinduction adjunction.  For
an \(A\)-algebra \(B\), a \(B\)-module \(M\), and an \(A\)-module
\(N\), there is a natural isomorphism

\[
\operatorname{Hom}_B(M,\operatorname{Hom}_A(B,N))
\xrightarrow{\sim}\operatorname{Hom}_A(M,N).
\tag{2}
\]

It sends \(\psi\) to \(m\mapsto\psi(m)(1)\); its inverse sends
\(h:M\to N\) to

\[
m\longmapsto\bigl(b\longmapsto h(bm)\bigr).
\]

Applying (2) over every affine open \(U\subseteq Y\), with
\(M=\Gamma(f^{-1}U,\mathcal F)\), gives compatible isomorphisms on
sections.  They sheafify to

\[
f_*\mathcal H\!om_X(\mathcal F,f^!\mathcal G)
\cong\mathcal H\!om_Y(f_*\mathcal F,\mathcal G).
\tag{3}
\]

(c) Resolve in the category of quasi-coherent modules and use the
derived comparison from the context to identify its Ext groups
with the sheaf Ext groups in the problem. Since \(f\) is affine,

\[
f_*:\operatorname{Qco}(X)\longrightarrow\operatorname{Qco}(Y)
\]

is exact, and (3) says that \(f^!\) is its right adjoint.  In
particular, \(f^!\) carries injective objects of
\(\operatorname{Qco}(Y)\) to injective objects of
\(\operatorname{Qco}(X)\).
The Grothendieck spectral sequence for

\[
\operatorname{Hom}_X(\mathcal F,-)\circ f^!
\cong\operatorname{Hom}_Y(f_*\mathcal F,-)
\]

is

\[
E_2^{p,q}
=\operatorname{Ext}_X^p
(\mathcal F,R^qf^!\mathcal G)
\Longrightarrow
\operatorname{Ext}_Y^{p+q}(f_*\mathcal F,\mathcal G).
\tag{4}
\]

The edge morphism from the bottom row of (4) is

\[
E_2^{i,0}
=\operatorname{Ext}_X^i(\mathcal F,f^!\mathcal G)
\longrightarrow
\operatorname{Ext}_Y^i(f_*\mathcal F,\mathcal G).
\]

Define this to be \(\varphi_i\).  The spectral-sequence construction
makes it natural, and for \(i=0\) it is exactly the adjunction (3) on
global Hom groups.

(d) If \(\mathcal B\) is locally free of finite rank, then

\[
f^!\mathcal G
\quad\leftrightarrow\quad
\mathcal H\!om_Y(\mathcal B,\mathcal G)
\cong\mathcal B^\vee\otimes\mathcal G.
\tag{5}
\]

Thus \(f^!\) is exact and \(R^qf^!=0\) for \(q>0\).  The spectral
sequence (4) has only its bottom row, so it degenerates and every edge
map \(\varphi_i\) is an isomorphism.

Equivalently, if \(\mathcal G\to\mathcal I^\bullet\) is an injective
resolution in \(\operatorname{Qco}(Y)\), exactness of \(f^!\)
and the adjunction with the exact
functor \(f_*\) make
\(f^!\mathcal I^\bullet\) an injective resolution of
\(f^!\mathcal G\).  Termwise adjunction then gives

\[
\operatorname{Hom}_X(\mathcal F,f^!\mathcal I^\bullet)
\cong
\operatorname{Hom}_Y(f_*\mathcal F,\mathcal I^\bullet),
\]

whose cohomology is the desired isomorphism.  In fact, this exact
adjunction argument proves a slightly stronger statement; the extra
separatedness and enough-locally-free hypotheses ensure the standard
comparison framework used in the exercise but are not otherwise
needed once (4) is available.
\end{solution}

\begin{exerciseintro}
\corpussection{III.7 The Serre Duality Theorem}
\addcontentsline{toc}{section}{III.7 The Serre Duality Theorem}

% record: H-III-07-001
\exerciseheading{Exercise III.7.1}
\addcontentsline{toc}{subsection}{Exercise III.7.1}

\begin{context}
Let \(X\) be an integral projective scheme of dimension at least one
over a field \(k\), and let \(\mathcal L\) be an ample invertible
sheaf on \(X\).
\end{context}

\begin{problem}
Prove that

\[
H^0(X,\mathcal L^{-1})=0.
\]
\end{problem}
\end{exerciseintro}

\begin{solution}
Suppose that \(0\neq s\in H^0(X,\mathcal L^{-1})\).  Because \(X\)
is integral, the corresponding map

\[
\mathcal O_X\xrightarrow{s}\mathcal L^{-1}
\]

is injective: after trivializing \(\mathcal L^{-1}\), it is
multiplication by a nonzero element of a domain.  Tensoring by
\(\mathcal L\) gives an injection
\(\mathcal L\hookrightarrow\mathcal O_X\), and taking tensor powers
gives

\[
\mathcal L^r\hookrightarrow\mathcal O_X
\qquad(r\geq1).
\tag{1}
\]

Thus

\[
\dim_kH^0(X,\mathcal L^r)
\leq\dim_kH^0(X,\mathcal O_X)
\tag{2}
\]

for every \(r\).  The number on the right is finite because \(X\) is
projective.

Choose \(m>0\) such that \(\mathcal L^m\) is very ample.  For
\(r\gg0\), the dimension of
\(H^0(X,\mathcal L^{mr})\) is the value of the Hilbert polynomial of
\(X\) in this embedding.  That polynomial has degree
\(\dim X\geq1\) and positive leading coefficient, so these dimensions
are unbounded as \(r\to\infty\).  This contradicts the uniform bound
(2).  Hence no nonzero section of \(\mathcal L^{-1}\) exists.
\end{solution}

\begin{exerciseintro}
% record: H-III-07-002
\exerciseheading{Exercise III.7.2}
\addcontentsline{toc}{subsection}{Exercise III.7.2}

\begin{context}
Let \(f:X\to Y\) be a finite morphism of projective \(k\)-schemes,
both of dimension \(n\).  Let \(\omega_Y^\circ\) be a dualizing sheaf
on \(Y\), with trace

\[
t_Y:H^n(Y,\omega_Y^\circ)\longrightarrow k.
\]

Use the functor \(f^!\) from Exercise III.6.10.
\end{context}

\begin{problem}
(a) Prove that \(f^!\omega_Y^\circ\) is a dualizing sheaf on \(X\).

(b) If \(X\) and \(Y\) are nonsingular and \(k\) is algebraically
closed, construct the natural trace morphism

\[
\operatorname{Tr}_f:f_*\omega_X\longrightarrow\omega_Y.
\]
\end{problem}
\end{exerciseintro}

\begin{solution}
(a) For every coherent \(\mathcal F\) on \(X\), finite duality
adjunction and Serre duality on \(Y\) give natural isomorphisms

\[
\begin{split}
\operatorname{Hom}_X(\mathcal F,f^!\omega_Y^\circ)
&\cong\operatorname{Hom}_Y(f_*\mathcal F,\omega_Y^\circ)\\
&\cong H^n(Y,f_*\mathcal F)^\vee.
\end{split}
\tag{1}
\]

Since \(f\) is finite, it is affine and has no higher direct images
on quasi-coherent sheaves.  The Leray spectral sequence therefore
gives

\[
H^n(Y,f_*\mathcal F)\cong H^n(X,\mathcal F).
\tag{2}
\]

Combining (1) and (2), we obtain functorially

\[
\operatorname{Hom}_X(\mathcal F,f^!\omega_Y^\circ)
\cong H^n(X,\mathcal F)^\vee.
\tag{3}
\]

Let \(t_X\in H^n(X,f^!\omega_Y^\circ)^\vee\) be the functional
corresponding under (3), with
\(\mathcal F=f^!\omega_Y^\circ\), to its identity endomorphism.
Naturality of (3) says that a map
\(\alpha:\mathcal F\to f^!\omega_Y^\circ\) corresponds exactly to

\[
H^n(X,\mathcal F)\xrightarrow{H^n(\alpha)}
H^n(X,f^!\omega_Y^\circ)\xrightarrow{t_X}k.
\]

Thus the pair \((f^!\omega_Y^\circ,t_X)\) has the defining
representing property of a dualizing sheaf on \(X\).

(b) On a nonsingular projective \(n\)-fold, the canonical sheaf
\(\omega=\Omega^n\), with its trace, is dualizing.  By the uniqueness
of a representing object with trace, part (a) gives a canonical
trace-preserving isomorphism

\[
\omega_X\xrightarrow{\sim}f^!\omega_Y.
\tag{4}
\]

The counit of the adjunction \(f_*\dashv f^!\) is

\[
f_*f^!\omega_Y\longrightarrow\omega_Y.
\tag{5}
\]

Affine-locally, for a finite map \(A\to B\), it is the evaluation map

\[
\operatorname{Hom}_A(B,\omega_A)\longrightarrow\omega_A,
\qquad \lambda\longmapsto\lambda(1).
\]

Composing \(f_*\) of (4) with (5) gives the natural morphism

\[
\operatorname{Tr}_f:f_*\omega_X\longrightarrow\omega_Y.
\]
\end{solution}

\begin{exerciseintro}
% record: H-III-07-003
\exerciseheading{Exercise III.7.3}
\addcontentsline{toc}{subsection}{Exercise III.7.3}

\begin{context}
Let \(X=\mathbb P_k^n\), and write
\(\Omega_X^p=\bigwedge^p\Omega_X^1\).  The Euler sequence is

\[
0\longrightarrow\Omega_X^1
\longrightarrow\mathcal O_X(-1)^{\oplus(n+1)}
\longrightarrow\mathcal O_X\longrightarrow0.
\]
\end{context}

\begin{problem}
For \(0\leq p,q\leq n\), prove

\[
H^q(X,\Omega_X^p)\cong
\begin{cases}
k,&p=q,\\
0,&p\neq q.
\end{cases}
\]
\end{problem}
\end{exerciseintro}

\begin{solution}
Because the quotient in the Euler sequence has rank one, its
\(p\)-th exterior-power sequence is short exact:

\[
0\longrightarrow\Omega_X^p
\longrightarrow
\bigwedge^p\!\left(\mathcal O_X(-1)^{\oplus(n+1)}\right)
\longrightarrow\Omega_X^{p-1}\longrightarrow0.
\]

Equivalently,

\[
0\longrightarrow\Omega_X^p
\longrightarrow
\mathcal O_X(-p)^{\oplus\binom{n+1}{p}}
\longrightarrow\Omega_X^{p-1}\longrightarrow0.
\tag{1}
\]

For \(1\leq p\leq n\), the line bundle
\(\mathcal O_X(-p)\) is acyclic.  Indeed it has no global sections,
all its intermediate cohomology vanishes, and by Serre duality

\[
H^n(X,\mathcal O_X(-p))
\cong H^0(X,\mathcal O_X(p-n-1))^\vee=0.
\]

The long exact sequence of (1) consequently gives

\[
H^q(X,\Omega_X^p)
\cong H^{q-1}(X,\Omega_X^{p-1})
\tag{2}
\]

with the convention that a negative cohomology group is zero.  The
initial case is

\[
H^q(X,\Omega_X^0)=H^q(X,\mathcal O_X)
\cong
\begin{cases}
k,&q=0,\\
0,&q>0.
\end{cases}
\]

Iterating (2) proves that the only nonzero group occurs when \(q=p\),
and in that case it is \(k\).
\end{solution}

\begin{exerciseintro}
% record: H-III-07-004
\exerciseheading{Exercise III.7.4}
\addcontentsline{toc}{subsection}{Exercise III.7.4}

\begin{context}
Let \(X\) be a nonsingular projective variety of dimension \(n\) over
an algebraically closed field \(k\), and let
\(i:Y\hookrightarrow X\) be a nonsingular subvariety of codimension
\(p\).  Restriction of differential forms and the trace on \(Y\)
give a functional

\[
H^{n-p}(X,\Omega_X^{n-p})
\xrightarrow{i^*}
H^{n-p}(Y,\omega_Y)
\xrightarrow{t_Y}k.
\]

Under the Serre-duality pairing

\[
H^p(X,\Omega_X^p)\times
H^{n-p}(X,\Omega_X^{n-p})
\xrightarrow{\smile,\,\wedge}
H^n(X,\omega_X)\xrightarrow{t_X}k,
\]

this functional corresponds to a unique class
\(\eta(Y)\in H^p(X,\Omega_X^p)\).  Thus

\[
t_X\bigl(\eta(Y)\smile\alpha\bigr)
=t_Y(i^*\alpha)
\tag{1}
\]

for every
\(\alpha\in H^{n-p}(X,\Omega_X^{n-p})\).
\end{context}

\begin{problem}
(a) If \(P\in X\) is a closed point, prove

\[
t_X(\eta(P))=1.
\]

(b) If \(X=\mathbb P_k^n\), use the identification
\(H^p(X,\Omega_X^p)\cong k\) from Exercise III.7.3, normalized
so that \(c(\mathcal O(1))^p\) maps to \(1\), to prove

\[
\eta(Y)=(\deg Y)\cdot1.
\]

(c) For a finite-type \(k\)-scheme \(Z\), show that

\[
d\log:\mathcal O_Z^\times\longrightarrow\Omega_Z^1,
\qquad f\longmapsto f^{-1}df,
\]

is a homomorphism of sheaves of abelian groups.  Denote the induced
homomorphism

\[
c:\operatorname{Pic}Z=H^1(Z,\mathcal O_Z^\times)
\longrightarrow H^1(Z,\Omega_Z^1)
\]

by \(c\).

(d) If \(p=1\), prove

\[
\eta(Y)=c(\mathcal O_X(Y)).
\]
\end{problem}
\end{exerciseintro}

\begin{solution}
(a) Here \(p=n\), so the functional defining \(\eta(P)\) is

\[
H^0(X,\mathcal O_X)\longrightarrow H^0(P,\mathcal O_P)
\xrightarrow{t_P}k.
\]

Because \(k\) is algebraically closed, \(P=\operatorname{Spec}k\)
and \(t_P\) is the identity.  Taking \(\alpha=1\) in (1) gives

\[
t_X(\eta(P))=t_P(1)=1.
\]

(c) The Leibniz rule gives

\[
d\log(fg)
=(fg)^{-1}(f\,dg+g\,df)
=f^{-1}df+g^{-1}dg.
\]

The construction commutes with restriction, so \(d\log\) is a
homomorphism of sheaves from the multiplicative group
\(\mathcal O_Z^\times\) to the additive group \(\Omega_Z^1\).
Applying \(H^1\) gives \(c\).  If an invertible sheaf has transition
functions \(g_{ij}\), its class \(c\) is represented by the Cech
cocycle \(d\log(g_{ij})\).  This description also shows that \(c\) is
additive under tensor products and compatible with pullback.

(d) Let \(Y\) be a nonsingular effective Cartier divisor.  We use the
residue sequence

\[
0\longrightarrow\omega_X
\longrightarrow\omega_X(Y)
\xrightarrow{\operatorname{res}}i_*\omega_Y
\longrightarrow0.
\tag{2}
\]

Write

\[
\partial:H^{n-1}(Y,\omega_Y)\longrightarrow H^n(X,\omega_X)
\tag{3}
\]

for its connecting map.  The residue normalization of the trace
satisfies

\[
t_X\circ\partial=t_Y.
\tag{4}
\]

There is also the standard divisor identity

\[
\partial(i^*\alpha)
=c(\mathcal O_X(Y))\smile\alpha
\qquad
\bigl(\alpha\in H^{n-1}(X,\Omega_X^{n-1})\bigr).
\tag{5}
\]

To see (5), choose local equations \(f_i\) for \(Y\).  The transition
functions of \(\mathcal O_X(Y)\) are, with compatible frame
conventions, \(g_{ij}=f_i/f_j\), so its \(c\)-class is represented by

\[
d\log(g_{ij})=\frac{df_i}{f_i}-\frac{df_j}{f_j}.
\]

Local logarithmic lifts in (2) have exactly these differences.
Computing the boundary in Cech cohomology therefore gives cup product
with this cocycle, which is (5).

Combining (4) and (5), for every \(\alpha\) we obtain

\[
t_X\!\left(c(\mathcal O_X(Y))\smile\alpha\right)
=t_Y(i^*\alpha).
\]

Comparison with (1), followed by perfection of the Serre-duality
pairing, proves
\(\eta(Y)=c(\mathcal O_X(Y))\).

(b) Put

\[
h=c(\mathcal O_{\mathbb P^n}(1))\in H^1(\mathbb P^n,\Omega^1).
\]

Under the recursive isomorphisms of Exercise III.7.3, \(h^r\) is the
element \(1\in H^r(\mathbb P^n,\Omega^r)\cong k\).  One way to fix
the normalization is to choose a transverse flag of hyperplanes.
Part (d), iterated through the residue sequences, gives

\[
h^r=\eta(\mathbb P^{n-r}).
\tag{6}
\]

For \(r=n\), part (a) consequently gives

\[
t_{\mathbb P^n}(h^n)=1.
\tag{7}
\]

We use the same residue calculation in the following general form:
if smooth divisors \(D_1,\ldots,D_r\) meet transversely on a smooth
projective variety, then

\[
\eta(D_1\cap\cdots\cap D_r)
=\prod_{j=1}^r c(\mathcal O(D_j)).
\tag{8}
\]

Indeed, (8) follows by iterating (5), or equivalently by the
projection formula for the successive residue maps.

Let \(m=\dim Y=n-p\).  Since
\(H^p(\mathbb P^n,\Omega^p)=k h^p\), write

\[
\eta(Y)=a h^p.
\]

Pairing with \(h^m\) and using (1) and (7) gives

\[
a
=t_{\mathbb P^n}(\eta(Y)\smile h^m)
=t_Y\bigl((i^*h)^m\bigr).
\tag{9}
\]

Choose \(m\) general hyperplanes meeting \(Y\) transversely.  Their
intersection with \(Y\) is a reduced zero-dimensional scheme
consisting of \(\deg Y\) points.  By (8) on \(Y\),
\((i^*h)^m\) is the sum of the cohomology classes of those points.
Part (a) says that the trace of each point class is \(1\), so the
right side of (9) is \(\deg Y\).  Hence

\[
\eta(Y)=(\deg Y)h^p=(\deg Y)\cdot1
\]

under the identification of Exercise III.7.3.
\end{solution}

\begin{exerciseintro}
\corpussection{III.8 Higher Direct Images of Sheaves}
\addcontentsline{toc}{section}{III.8 Higher Direct Images of Sheaves}

% record: H-III-08-001
\exerciseheading{Exercise III.8.1}
\addcontentsline{toc}{subsection}{Exercise III.8.1}

\begin{context}
Let \(f:X\to Y\) be a continuous map.  Direct image is a left-exact
functor on sheaves of abelian groups, and its right-derived functors
are denoted by \(R^qf_*\).  Injective sheaves are flasque, direct image
preserves flasque sheaves, and flasque resolutions compute sheaf
cohomology.
\end{context}

\begin{problem}
Let \(\mathcal F\) be a sheaf of abelian groups on \(X\), and suppose

\[
R^qf_*\mathcal F=0\qquad(q>0).
\]

Prove that, naturally in \(\mathcal F\),

\[
H^i(X,\mathcal F)\cong H^i(Y,f_*\mathcal F)
\qquad(i\geq0).
\]
\end{problem}
\end{exerciseintro}

\begin{solution}
Choose an injective resolution

\[
0\longrightarrow\mathcal F\longrightarrow\mathcal I^0
\longrightarrow\mathcal I^1\longrightarrow\cdots .
\tag{1}
\]

Set \(\mathcal K^0=\mathcal F\), and for \(q\geq0\) let
\(\mathcal K^{q+1}\) be the cokernel of
\(\mathcal K^q\hookrightarrow\mathcal I^q\).  Thus there are short
exact sequences

\[
0\longrightarrow\mathcal K^q\longrightarrow\mathcal I^q
\longrightarrow\mathcal K^{q+1}\longrightarrow0.
\tag{2}
\]

Because \(\mathcal I^q\) is \(f_*\)-acyclic, the long exact sequence
of higher direct images associated to (2) gives

\[
R^1f_*\mathcal K^q\cong R^{q+1}f_*\mathcal F.
\tag{3}
\]

Indeed, this follows by dimension shifting, starting with
\(\mathcal K^0=\mathcal F\).  The hypothesis makes the right-hand
side of (3) zero.  The beginning of the same long exact sequence is
therefore

\[
0\longrightarrow f_*\mathcal K^q\longrightarrow f_*\mathcal I^q
\longrightarrow f_*\mathcal K^{q+1}\longrightarrow0.
\]

Consequently, applying \(f_*\) to (1) produces an exact resolution

\[
0\longrightarrow f_*\mathcal F\longrightarrow f_*\mathcal I^0
\longrightarrow f_*\mathcal I^1\longrightarrow\cdots .
\tag{4}
\]

Every \(f_*\mathcal I^q\) is flasque: if \(V\subseteq U\subseteq Y\),
its restriction map is the restriction map

\[
\mathcal I^q(f^{-1}U)\longrightarrow\mathcal I^q(f^{-1}V),
\]

which is surjective.  Thus (4) is a flasque resolution and computes
\(H^i(Y,f_*\mathcal F)\).  On the other hand, (1) computes
\(H^i(X,\mathcal F)\), and in every degree

\[
\Gamma(Y,f_*\mathcal I^q)=\Gamma(X,\mathcal I^q).
\]

The two cohomology groups are therefore the cohomology groups of the
same complex.  The comparison isomorphism is independent of the
chosen resolution by the comparison theorem for resolutions, so it
is natural in \(\mathcal F\).
\end{solution}

\begin{exerciseintro}
% record: H-III-08-002
\exerciseheading{Exercise III.8.2}
\addcontentsline{toc}{subsection}{Exercise III.8.2}

\begin{context}
Let \(f:X\to Y\) be a morphism of schemes with \(X\) noetherian.
For a quasi-coherent sheaf \(\mathcal F\) on \(X\), higher direct
images may be computed locally on the target: on an open subset
\(U\subseteq Y\),

\[
(R^if_*\mathcal F)|_U
\cong R^i(f_U)_*(\mathcal F|_{f^{-1}U}),
\]

where \(f_U:f^{-1}U\to U\).  Equivalently, on an affine-open basis
they are obtained from the groups
\(H^i(f^{-1}U,\mathcal F)\).  Quasi-coherent sheaves have no positive
cohomology on noetherian affine schemes.
\end{context}

\begin{problem}
Suppose that \(f:X\to Y\) is affine, \(X\) is noetherian, and
\(\mathcal F\) is quasi-coherent.  Prove

\[
R^if_*\mathcal F=0\qquad(i>0),
\]

and deduce natural isomorphisms

\[
H^i(X,\mathcal F)\cong H^i(Y,f_*\mathcal F)
\qquad(i\geq0).
\]
\end{problem}
\end{exerciseintro}

\begin{solution}
Let \(U\subseteq Y\) be affine.  Since \(f\) is affine,
\(f^{-1}U\) is affine; it is also noetherian because it is an open
subscheme of the noetherian scheme \(X\).  The restriction of
\(\mathcal F\) to \(f^{-1}U\) is quasi-coherent, so affine vanishing
gives

\[
H^i(f^{-1}U,\mathcal F)=0\qquad(i>0).
\tag{1}
\]

The same assertion holds after replacing \(U\) by any affine open
subset of it.  By the local description of higher direct images,
(1) says that \((R^if_*\mathcal F)|_U=0\) for every affine open
\(U\subseteq Y\).  Such opens form a basis, hence

\[
R^if_*\mathcal F=0\qquad(i>0).
\]

Exercise III.8.1 now applies and yields

\[
H^i(X,\mathcal F)\cong H^i(Y,f_*\mathcal F)
\]

in every degree.  Naturality follows from the naturality of both
affine cohomology and the comparison in that exercise.
\end{solution}

\begin{exerciseintro}
% record: H-III-08-003
\exerciseheading{Exercise III.8.3}
\addcontentsline{toc}{subsection}{Exercise III.8.3}

\begin{context}
Let \(f:(X,\mathcal O_X)\to(Y,\mathcal O_Y)\) be a morphism of
ringed spaces.  Higher direct images commute with restriction to an
open subset of the target and with finite direct sums.  If
\(\mathcal E\) is finite locally free, tensoring with \(\mathcal E\)
is exact, and the ordinary projection formula is

\[
f_*(\mathcal F\otimes_{\mathcal O_X}f^*\mathcal E)
\cong f_*\mathcal F\otimes_{\mathcal O_Y}\mathcal E.
\]
\end{context}

\begin{problem}
For an \(\mathcal O_X\)-module \(\mathcal F\) and a finite locally
free \(\mathcal O_Y\)-module \(\mathcal E\), prove the higher
projection formula

\[
R^if_*(\mathcal F\otimes_{\mathcal O_X}f^*\mathcal E)
\cong R^if_*\mathcal F\otimes_{\mathcal O_Y}\mathcal E
\qquad(i\geq0).
\]
\end{problem}
\end{exerciseintro}

\begin{solution}
Choose an open subset
\(U\subseteq Y\) on which
\(\mathcal E|_U\cong\mathcal O_U^{\oplus r}\), and write
\(f_U:f^{-1}U\to U\).  Then

\[
(\mathcal F\otimes f^*\mathcal E)|_{f^{-1}U}
\cong(\mathcal F|_{f^{-1}U})^{\oplus r}.
\]

Additivity of higher direct images now gives an isomorphism

\[
R^if_{U*}((\mathcal F|_{f^{-1}U})^{\oplus r})
\cong
(R^if_{U*}(\mathcal F|_{f^{-1}U}))^{\oplus r}.
\tag{1}
\]

Via the chosen frame, the left side of (1) is the restriction of
\(R^if_*(\mathcal F\otimes f^*\mathcal E)\), and the right side is
the restriction of \(R^if_*\mathcal F\otimes\mathcal E\).

This local isomorphism does not depend on the frame.  Indeed, a change
of frame is represented by an invertible matrix with entries in
\(\mathcal O_U\).  Functoriality and \(\mathcal O_U\)-linearity of
\(R^if_{U*}\) make that matrix act in exactly the same way on the two
sides of (1).  The local isomorphisms therefore agree on overlaps and
glue to the required global isomorphism.  For \(i=0\), the resulting
map is precisely the ordinary projection formula.
\end{solution}

\begin{exerciseintro}
% record: H-III-08-004
\exerciseheading{Exercise III.8.4}
\addcontentsline{toc}{subsection}{Exercise III.8.4}

\begin{context}
Let \(Y\) be noetherian, let \(\mathcal E\) be a locally free
\(\mathcal O_Y\)-module of rank \(n+1\), where \(n\geq1\), and use
the convention

\[
X=\mathbb P(\mathcal E)
=\operatorname{Proj}_Y(\operatorname{Sym}\mathcal E).
\]

Write \(\pi:X\to Y\) for the projection and
\(\mathcal O_X(1)\) for the tautological invertible sheaf.  Thus the
convention is normalized by
\(\pi_*\mathcal O_X(l)=\operatorname{Sym}^l\mathcal E\) for
\(l\geq0\).  Put
\(\mathcal L=\bigwedge^{n+1}\mathcal E\).

When arithmetic and geometric genera occur below, assume in addition
that \(Y\) is a projective variety over a field \(k\); for the
geometric-genus assertion assume that \(Y\) is nonsingular.  Then
\(X\) is also nonsingular and projective.
\end{context}

\begin{problem}
(a) Compute the direct images of the twists sufficiently to prove

\[
\pi_*\mathcal O_X(l)=
\begin{cases}
\operatorname{Sym}^l\mathcal E,&l\geq0,\\
0,&l<0,
\end{cases}
\]

\[
R^i\pi_*\mathcal O_X(l)=0
\quad(0<i<n,\ l\in\mathbb Z),
\qquad
R^n\pi_*\mathcal O_X(l)=0
\quad(l>-n-1).
\]

(b) Prove the relative Euler sequence

\[
0\longrightarrow\Omega_{X/Y}
\longrightarrow\pi^*\mathcal E\otimes\mathcal O_X(-1)
\longrightarrow\mathcal O_X\longrightarrow0.
\]

Deduce

\[
\omega_{X/Y}:=\bigwedge^n\Omega_{X/Y}
\cong
\pi^*\mathcal L\otimes\mathcal O_X(-n-1),
\]

and construct a natural isomorphism
\(R^n\pi_*\omega_{X/Y}\cong\mathcal O_Y\).

(c) For every \(l\in\mathbb Z\), prove

\[
R^n\pi_*\mathcal O_X(l)
\cong
\bigl(\pi_*\mathcal O_X(-l-n-1)\bigr)^\vee
\otimes\mathcal L^\vee.
\]

(d) Prove

\[
p_a(X)=(-1)^n p_a(Y),
\qquad
p_g(X)=0.
\]

(e) Let \(Y\) be a nonsingular projective curve of genus \(g\), and
let \(\mathcal E\) have rank \(2\).  Show that \(X\) is a projective
surface with

\[
p_a(X)=-g,\qquad p_g(X)=0,\qquad q(X)=g.
\]
\end{problem}
\end{exerciseintro}

\begin{solution}
(a) All assertions are local on \(Y\).  Take an affine open
\(U=\operatorname{Spec}A\) on which
\(\mathcal E|_U\cong\mathcal O_U^{\oplus(n+1)}\).  Then

\[
X_U:=\pi^{-1}U\cong\mathbb P_A^n
=\operatorname{Proj}A[x_0,\ldots,x_n].
\]

The standard affine cover \(D_+(x_0),\ldots,D_+(x_n)\) and its Cech
complex give

\[
H^q(\mathbb P_A^n,\mathcal O(l))=0
\qquad(0<q<n)
\tag{1}
\]

for every \(l\), while

\[
H^0(\mathbb P_A^n,\mathcal O(l))=
\begin{cases}
A[x_0,\ldots,x_n]_l,&l\geq0,\\
0,&l<0.
\end{cases}
\tag{2}
\]

The only possible remaining positive cohomology is in degree \(n\).
The same Cech calculation identifies it with the free \(A\)-module
having basis

\[
x_0^{-a_0}\cdots x_n^{-a_n},
\qquad
a_j\geq1,\quad \sum_{j=0}^n a_j=-l.
\tag{3}
\]

In particular, this module is zero when \(l\geq-n\), equivalently
when \(l>-n-1\).  These computations commute with localization in
\(A\), so they describe the restrictions of the higher direct-image
sheaves to \(U\).  Equations (1)--(3) prove all the vanishing
statements.  Finally, (2) is compatible with change of frame in
\(\mathcal E\) and glues to

\[
\pi_*\mathcal O_X(l)=\operatorname{Sym}^l\mathcal E
\]

for \(l\geq0\); for \(l<0\) it glues to zero.

(b) The tautological quotient

\[
\pi^*\mathcal E\longrightarrow\mathcal O_X(1)
\]

becomes, after tensoring by \(\mathcal O_X(-1)\), a surjection

\[
\pi^*\mathcal E\otimes\mathcal O_X(-1)
\longrightarrow\mathcal O_X.
\tag{4}
\]

On every open set \(U\) used in part (a), (4) is the last map in the
usual Euler sequence on \(\mathbb P_A^n\); its kernel is
\(\Omega_{\mathbb P_A^n/A}\).  Since (4) is intrinsic, these local
kernel identifications agree on overlaps.  This proves the asserted
relative Euler sequence.

Taking determinants in that sequence gives

\[
\begin{aligned}
\omega_{X/Y}
&\cong
\det(\pi^*\mathcal E\otimes\mathcal O_X(-1))\\
&\cong
\pi^*(\det\mathcal E)\otimes\mathcal O_X(-n-1)\\
&=
\pi^*\mathcal L\otimes\mathcal O_X(-n-1).
\end{aligned}
\tag{5}
\]

There is a canonical local residue isomorphism

\[
R^n\pi_*\mathcal O_X(-n-1)\cong\mathcal L^\vee.
\tag{6}
\]

To see both the isomorphism and its canonicity, choose a frame
\(x_0,\ldots,x_n\) of \(\mathcal E|_U\).  Formula (3) says that the
left side of (6) is freely generated by the Cech class

\[
\left[\frac{1}{x_0\cdots x_n}\right].
\]

Pair it with \(x_0\wedge\cdots\wedge x_n\) by sending the displayed
pair to \(1\).  Under a change of frame the exterior product is
multiplied by the determinant, while the top Cech class is multiplied
by the inverse determinant modulo a coboundary.  Hence the pairing is
independent of the frame and glues to (6).

Using (5), the higher projection formula, and (6), we obtain

\[
\begin{aligned}
R^n\pi_*\omega_{X/Y}
&\cong
R^n\pi_*(\pi^*\mathcal L\otimes\mathcal O_X(-n-1))\\
&\cong
\mathcal L\otimes R^n\pi_*\mathcal O_X(-n-1)\\
&\cong
\mathcal L\otimes\mathcal L^\vee
\cong\mathcal O_Y.
\end{aligned}
\tag{7}
\]

This is the natural relative trace isomorphism.

(c) Multiplication of twists, followed by the trace in (7), gives a
pairing

\[
\pi_*\mathcal Hom_{\mathcal O_X}
   (\mathcal O_X(l),\omega_{X/Y})
\otimes R^n\pi_*\mathcal O_X(l)
\longrightarrow\mathcal O_Y.
\tag{8}
\]

It is perfect.  This may be checked on a trivializing affine open
\(U=\operatorname{Spec}A\).  Put \(d=-l-n-1\).  If \(d<0\), both
terms relevant to the asserted duality vanish by part (a).  If
\(d\geq0\), degree-\(d\) monomials

\[
x_0^{b_0}\cdots x_n^{b_n},
\qquad b_j\geq0,\quad \sum b_j=d,
\]

form a basis of \(H^0(\mathbb P_A^n,\mathcal O(d))\).  The basis
(3) of \(H^n(\mathbb P_A^n,\mathcal O(l))\) is indexed by
\(a_j=b_j+1\).  Multiplication followed by residue sends the
corresponding two basis elements to \(1\), and sends distinct paired
basis elements to \(0\).  Thus (8) is a perfect pairing over \(A\);
the residue construction shows that these local pairings glue.

Now (5) and the ordinary projection formula give

\[
\begin{aligned}
\pi_*\mathcal Hom(\mathcal O_X(l),\omega_{X/Y})
&\cong
\pi_*(\pi^*\mathcal L\otimes\mathcal O_X(-l-n-1))\\
&\cong
\mathcal L\otimes\pi_*\mathcal O_X(-l-n-1).
\end{aligned}
\]

Dualizing this identity through the perfect pairing (8) yields

\[
R^n\pi_*\mathcal O_X(l)
\cong
\bigl(\pi_*\mathcal O_X(-l-n-1)\bigr)^\vee
\otimes\mathcal L^\vee,
\]

as required.

(d) Taking \(l=0\) in part (a) gives

\[
\pi_*\mathcal O_X=\mathcal O_Y,
\qquad
R^i\pi_*\mathcal O_X=0\quad(i>0).
\tag{9}
\]

For \(i=n\), the second assertion follows because
\(0>-n-1\); for \(i>n\), it follows from the same affine Cech cover,
whose complex has length \(n\).  Exercise III.8.1 applied to (9)
therefore gives

\[
H^i(X,\mathcal O_X)\cong H^i(Y,\mathcal O_Y)
\qquad(i\geq0).
\tag{10}
\]

In particular, \(\chi(\mathcal O_X)=\chi(\mathcal O_Y)\).  If
\(m=\dim Y\), then \(\dim X=m+n\), and the definition

\[
p_a(Z)=(-1)^{\dim Z}\bigl(\chi(\mathcal O_Z)-1\bigr)
\]

implies

\[
p_a(X)=(-1)^{m+n}(\chi(\mathcal O_Y)-1)
=(-1)^n p_a(Y).
\]

For the geometric genus, smoothness of the projective bundle and the
exact sequence of relative differentials give

\[
\omega_X\cong\pi^*\omega_Y\otimes\omega_{X/Y}
\cong
\pi^*(\omega_Y\otimes\mathcal L)
\otimes\mathcal O_X(-n-1).
\]

By part (a) and the projection formula,

\[
\pi_*\omega_X
\cong
(\omega_Y\otimes\mathcal L)
\otimes\pi_*\mathcal O_X(-n-1)=0.
\]

Hence \(H^0(X,\omega_X)=0\), so \(p_g(X)=0\).

(e) Here \(n=1\), so \(X\) has dimension \(2\).  It is projective
over the projective curve \(Y\), and hence projective over \(k\).
Since \(p_a(Y)=g\), part (d) gives

\[
p_a(X)=-g,\qquad p_g(X)=0.
\]

Finally, (10) gives

\[
q(X):=h^1(X,\mathcal O_X)
=h^1(Y,\mathcal O_Y)=g.
\]

Thus \(X\) has all the stated invariants.
\end{solution}

\begin{exerciseintro}
\corpussection{III.9 Flat Morphisms}
\addcontentsline{toc}{section}{III.9 Flat Morphisms}

% record: H-III-09-001
\exerciseheading{Exercise III.9.1}
\addcontentsline{toc}{subsection}{Exercise III.9.1}

\begin{context}
A finite-type morphism of noetherian schemes has constructible image.
In a noetherian topological space, a constructible subset is open if
and only if it is stable under generization.  Flatness is preserved by
restriction and base change.  A flat module over a domain is
torsion-free.
\end{context}

\begin{problem}
Prove that a flat morphism \(f:X\to Y\) of finite type between
noetherian schemes is open.
\end{problem}
\end{exerciseintro}

\begin{solution}
It suffices to show that the image of \(X\) is open.  Indeed, if
\(W\subseteq X\) is open, then \(W\) is noetherian and therefore
quasi-compact.  The restriction \(W\to Y\) is again flat and of
finite type, so the same assertion applied to \(W\) makes \(f(W)\)
open.

The question is local on \(Y\), and a finite affine-open cover of
\(X\) reduces it to

\[
f:\operatorname{Spec}B\longrightarrow\operatorname{Spec}A,
\]

where \(B\) is a finitely generated flat \(A\)-algebra.  Its image is
constructible.  We prove that it is stable under generization.

Let \(\mathfrak q\subseteq\mathfrak p\) be primes of \(A\), and
suppose \(\mathfrak p\) is in the image.  Make the base change

\[
A'=A_{\mathfrak p}/\mathfrak qA_{\mathfrak p},
\qquad
B'=B\otimes_A A'.
\]

The ring \(A'\) is a local domain.  Its closed point is the point
induced by \(\mathfrak p\), and its generic point is induced by
\(\mathfrak q\).  Since the fiber over \(\mathfrak p\) was nonempty,

\[
B'\otimes_{A'}\kappa(\mathfrak p)\ne0.
\]

In particular \(B'\ne0\).  The \(A'\)-module \(B'\) is flat, so it is
torsion-free.  If \(K=\operatorname{Frac}(A')\), the localization map

\[
B'\longrightarrow B'\otimes_{A'}K
\]

is therefore injective.  Its source is nonzero, and hence the generic
fiber \(\operatorname{Spec}(B'\otimes_{A'}K)\) is nonempty.  This says
that \(\mathfrak q\) belongs to the image of
\(\operatorname{Spec}B\to\operatorname{Spec}A\).

Thus the constructible image is stable under generization, so it is
open.  Applying this argument to every open \(W\subseteq X\) proves
that \(f\) is an open map.
\end{solution}

\begin{exerciseintro}
% record: H-III-09-002
\exerciseheading{Exercise III.9.2}
\addcontentsline{toc}{subsection}{Exercise III.9.2}

\begin{context}
Let \(a\) be the coordinate on \(\mathbb A_k^1\), and use homogeneous
coordinates \(x,y,z,w\) on \(\mathbb P_k^3\).  The twisted cubic has
parametrization

\[
[u:v]\longmapsto[u^3:u^2v:uv^2:v^3]
\]

and ideal

\[
(xz-y^2,\;xw-yz,\;yw-z^2).
\]

Apply the one-parameter coordinate change used in the flat-limit
construction by replacing the third coordinate by \(a\) times that
coordinate, and take the scheme-theoretic closure over
\(\mathbb A^1\).
We use the monomial basis theorem and Buchberger criterion for a Groebner basis.
\end{context}

\begin{problem}
Compute the special fiber of this family and prove that it is the
plane cuspidal cubic together with an embedded point at its cusp.
\end{problem}
\end{exerciseintro}

\begin{solution}
Over \(a\ne0\), the transformed twisted cubic is parametrized by

\[
x=u^3,\qquad y=u^2v,\qquad z=auv^2,\qquad w=v^3.
\tag{1}
\]

Its ideal in \(k[a,a^{-1},x,y,z,w]\) is

\[
(xz-ay^2,\;axw-yz,\;a^2yw-z^2).
\tag{2}
\]

The ideal of the closure is the contraction of (2) to
\(k[a,x,y,z,w]\).  It is

\[
J=(xz-ay^2,\;axw-yz,\;a^2yw-z^2,\;x^2w-y^3).
\tag{3}
\]

Here is an explicit verification that (3) is the contraction.
Use lexicographic order \(z>y>x>w>a\) and write
\[
g_1=xz-ay^2,\quad g_2=yz-axw,\quad
g_3=z^2-a^2yw,\quad g_4=y^3-x^2w.
\]
The leading monomials are \(xz,yz,z^2,y^3\).
The six S-polynomials have the following reductions:
\[
\begin{array}{c|l}
(1,2)&-a g_4\\
(1,3)&-ay g_2\\
(1,4)&x^2w g_1-ay^2g_4\\
(2,3)&-aw g_1\\
(2,4)&xw g_1\\
(3,4)&x^2w g_3-a^2yw g_4.
\end{array}
\]
In each row successive leading-term division gives zero.
Buchberger's criterion therefore makes these polynomials a
Groebner basis.  Their leading monomials do not involve \(a\);
the standard monomials give a free \(k[a]\)-basis of the quotient.
Thus multiplication by \(a\) is injective modulo \(J\), and
\(J\) is saturated with respect to \(a\).
After inverting \(a\), its fourth generator follows from the first
two by
\[
x(axw-yz)+y(xz-ay^2)=a(x^2w-y^3).
\]
Hence \(J[a^{-1}]\) is (2), and saturation proves the required
contraction equality.  All leading coefficients are one, so this
verification works in every characteristic.

Setting \(a=0\) gives the special-fiber ideal

\[
J_0=(xz,\;yz,\;z^2,\;x^2w-y^3).
\tag{4}
\]

Its radical is

\[
\sqrt{J_0}=(z,\;x^2w-y^3),
\]

the plane cubic in the plane \(z=0\).  This cubic is cuspidal at
\(P=[0:0:0:1]\).

To see the embedded structure, work on the affine chart \(w=1\).
Then (4) becomes

\[
I=(xz,\;yz,\;z^2,\;x^2-y^3)\subseteq k[x,y,z].
\]

There is a primary decomposition

\[
I=(z,\;x^2-y^3)\cap(x,\;y,\;z^2).
\tag{5}
\]

For completeness, the inclusion from left to right is immediate.
Conversely, if
\(h=cz+d(x^2-y^3)\) lies in \((x,y,z^2)\), then
\(cz\in(x,y,z^2)\), whence
\(c\in(x,y,z)\); it follows that
\(h\in(xz,yz,z^2,x^2-y^3)\).

The first component in (5) is the cuspidal cubic.  The second is
primary to the maximal ideal \((x,y,z)\), and that maximal ideal
strictly contains the minimal prime of the cubic.  Hence it is an
embedded associated point, supported at \(P\), with its nilpotent
direction transverse to the plane \(z=0\).
\end{solution}

\begin{exerciseintro}
% record: H-III-09-003
\exerciseheading{Exercise III.9.3}
\addcontentsline{toc}{subsection}{Exercise III.9.3}

\begin{context}
For a finite module over a noetherian local ring, flatness is
equivalent to freeness.  A regular local ring has finite global
dimension, and the Auslander--Buchsbaum formula says

\[
\operatorname{pd}_A M+\operatorname{depth}_A M
=\operatorname{depth}A
\]

for every nonzero finite \(A\)-module \(M\).  Regular local rings are
Cohen--Macaulay.  Flatness can also be detected by vanishing of
\(\operatorname{Tor}_1\).
\end{context}

\begin{problem}
(a) Let \(f:X\to Y\) be a finite surjective morphism of nonsingular
varieties over an algebraically closed field.  Prove that \(f\) is
flat.

(b) Put \(A=k[x,y]\), and let

\[
B=\frac{k[x,y,z,w]}
{(z,w)\cap(x+z,y+w)}.
\]

Geometrically, \(\operatorname{Spec}B\) is the union of two planes
meeting at one point, and each plane maps isomorphically to
\(\operatorname{Spec}A\).  Prove that \(B\) is not flat over \(A\).

(c) Put

\[
C=\frac{A[z,w]}{(z^2,zw,w^2,xz-yw)}.
\]

Show that \((\operatorname{Spec}C)_{\mathrm{red}}\cong
\operatorname{Spec}A\), that \(\operatorname{Spec}C\) has no embedded
points, and that \(C\) is nevertheless not flat over \(A\).
\end{problem}
\end{exerciseintro}

\begin{solution}
(a) Fix \(y\in Y\), set \(A=\mathcal O_{Y,y}\), and use the finite
semilocal algebra
\[
B=(f_*\mathcal O_X)_y.
\]
Its maximal ideals correspond to the finitely many points \(x\)
above \(y\), and \(B_{\mathfrak n_x}=\mathcal O_{X,x}\).
The dimension formula for finite dominant morphisms of varieties
gives \(\dim B_{\mathfrak n_x}=\dim A=d\) for all these points.
Each \(B_{\mathfrak n_x}\) is regular, hence Cohen--Macaulay.

A regular system of parameters \(t_1,\ldots,t_d\) of \(A\)
generates a primary ideal in each \(B_{\mathfrak n_x}\), because the
finite fiber is zero-dimensional.  It is therefore a regular
sequence in each such localization.  Injectivity of multiplication
on the successive quotients can be tested at all maximal ideals,
so the same sequence is \(B\)-regular.  Thus
\(\operatorname{depth}_A B=d\).
Auslander--Buchsbaum applies to this finite \(A\)-module and gives
\(\operatorname{pd}_A B=0\).  Consequently \(B\) is free over \(A\).
Localizing \(B\) at \(\mathfrak n_x\) preserves flatness over \(A\),
so all the local maps of \(f\) are flat.

(b) After identifying the coordinate ring of either plane with
\(A\), the coordinate ring of their union is the fiber product

\[
B=\{(r,s)\in A\oplus A:r(0,0)=s(0,0)\}.
\]

Thus there is an exact sequence of \(A\)-modules

\[
0\longrightarrow B\longrightarrow A\oplus A
\xrightarrow{(r,s)\mapsto r(0,0)-s(0,0)}k
\longrightarrow0.
\tag{2}
\]

Let \(\mathfrak m=(x,y)\), so \(k=A/\mathfrak m\).  The Koszul
resolution of \(k\) is

\[
0\longrightarrow A\longrightarrow A^{\oplus2}
\longrightarrow A\longrightarrow k\longrightarrow0.
\]

In particular,
\(\operatorname{Tor}_2^A(k,k)\cong k\).  Applying
\(-\otimes_A k\) to (2), and using that \(A\oplus A\) is free, gives

\[
\operatorname{Tor}_1^A(B,k)
\cong\operatorname{Tor}_2^A(k,k)
\cong k\ne0.
\]

Therefore \(B\) is not flat over \(A\).  Equivalently, the generic
fiber has length \(2\), whereas the fiber over \(\mathfrak m\) has
length \(3\).

(c) Let \(N=(z,w)\subseteq C\).  Then \(N^2=0\) and
\(C/N\cong A\).  Since \(A\) is reduced, \(N\) is exactly the
nilradical, and hence

\[
(\operatorname{Spec}C)_{\mathrm{red}}\cong\operatorname{Spec}A.
\]

As an \(A\)-module, \(C\) has the presentation

\[
0\longrightarrow A
\xrightarrow{\,1\mapsto(0,x,-y)\,}
A^{\oplus3}
\xrightarrow{(a,b,c)\mapsto a+bz+cw}
C\longrightarrow0.
\tag{3}
\]

In particular,

\[
C\cong A\oplus
\frac{Az\oplus Aw}{A(xz-yw)}.
\]

The second summand is isomorphic to the ideal \((x,y)\subseteq A\):
send \(z\mapsto y\) and \(w\mapsto x\).  Thus \(C\), viewed as an
\(A\)-module, is torsion-free.

Every prime of \(C\) contains the nilradical \(N\), and primes of
\(C\) correspond to primes of \(A=C/N\).  If
\(\mathfrak p\in\operatorname{Ass}_C(C)\), then
\(\mathfrak p\cap A=(0)\), because every nonzero element of the
domain \(A\) acts injectively on the torsion-free \(A\)-module \(C\).
It follows that \(\mathfrak p=N\).  Hence the only associated prime
of \(C\) is its minimal prime \(N\), so
\(\operatorname{Spec}C\) has no embedded points.

Finally, tensor (3) with \(k=A/(x,y)\).  Its left-hand map becomes
the zero map, so

\[
\operatorname{Tor}_1^A(C,k)
\cong k\ne0.
\]

Therefore \(C\) is not flat over \(A\), despite being generically a
double structure without embedded points.  Its generic fiber has
length \(2\), while its fiber over the origin has length \(3\).
\end{solution}

\begin{exerciseintro}
% record: H-III-09-004
\exerciseheading{Exercise III.9.4}
\addcontentsline{toc}{subsection}{Exercise III.9.4}

\begin{context}
We use the local criterion for flatness in the following form.  Let
\((R,\mathfrak m)\to(S,\mathfrak n)\) be a local homomorphism of
noetherian local rings, let \(I\subseteq\mathfrak m\), and let \(M\)
be a finite \(S\)-module.  Then \(M\) is flat over \(R\) if and only
if

\[
M/IM\ \text{is flat over }R/I
\quad\text{and}\quad
\operatorname{Tor}_1^R(R/I,M)=0.
\]

We also use that Tor commutes with localization.
The topological generization criterion is applied to the noetherian sober space \(\operatorname{Spec}B\); generic points exist there.
\end{context}

\begin{problem}
Let \(f:X\to Y\) be a finite-type morphism of noetherian schemes.
Prove that the flat locus

\[
\{x\in X:f\text{ is flat at }x\}
\]

is open in \(X\).
\end{problem}
\end{exerciseintro}

\begin{solution}
We first prove the generic-freeness lemma needed below.  If \(R\) is
a noetherian domain, \(C\) is
a finite-type \(R\)-algebra, and \(M\) is a finite \(C\)-module, then
there is a nonzero \(r\in R\) such that \(M_r\) is a free
\(R_r\)-module.

Write \(C\) as a quotient of
\(R[x_1,\ldots,x_n]\), and regard \(M\) as a finite module over that
polynomial ring.  We argue by induction on \(n\).  If \(n=0\), choose
elements \(m_1,\ldots,m_s\in M\) whose images form a basis of
\(M\otimes_R K\), where \(K=\operatorname{Frac}R\).  After inverting
one nonzero element of \(R\), the \(m_i\) generate \(M\), and their
linear independence over \(K\) makes them a basis.

Suppose \(n>0\), and put
\(R'=R[x_1,\ldots,x_{n-1}]\).  Choose
\(R'[x_n]\)-generators of \(M\), and let \(N\) be their
\(R'\)-span.  Define

\[
F_j=N+x_nN+\cdots+x_n^{j-1}N
\quad(j\geq1),\qquad F_0=0.
\]

Then \(M=\bigcup_jF_j\).  Multiplication by \(x_n\) induces
surjections

\[
F_1/F_0\twoheadrightarrow F_2/F_1
\twoheadrightarrow F_3/F_2\twoheadrightarrow\cdots .
\tag{1}
\]

All these quotients are finite \(R'\)-modules.  Their kernels, viewed
as an increasing sequence of submodules of the noetherian module
\(F_1\), eventually stabilize.  Thus only finitely many quotients in
(1) occur.  By induction, after inverting one nonzero element
\(r\in R\), every quotient \((F_{j+1}/F_j)_r\) is free over \(R_r\).
Each exact sequence

\[
0\longrightarrow(F_j)_r\longrightarrow(F_{j+1})_r
\longrightarrow(F_{j+1}/F_j)_r\longrightarrow0
\]

then splits as a sequence of \(R_r\)-modules.  Taking the increasing
union expresses \(M_r\) as a direct sum of free \(R_r\)-modules, so
it is free.  This proves generic freeness.

We now prove a slightly more general affine statement.  Let
\(A\to B\) be a finite-type homomorphism of noetherian rings, and let
\(M\) be a finite \(B\)-module.  Set

\[
U=\left\{\mathfrak p\in\operatorname{Spec}B:
M_{\mathfrak p}\text{ is flat over }A_{\mathfrak q},
\ \mathfrak q=\mathfrak p\cap A\right\}.
\tag{2}
\]

First, \(U\) is stable under generization.  Indeed, suppose
\(\mathfrak p_1\subseteq\mathfrak p_2\), put
\(\mathfrak q_i=\mathfrak p_i\cap A\), and assume
\(\mathfrak p_2\in U\).  Base change from
\(A_{\mathfrak q_2}\) to \(A_{\mathfrak q_1}\), followed by
localization from \(\mathfrak p_2\) to \(\mathfrak p_1\), shows that
\(M_{\mathfrak p_1}\) is flat over \(A_{\mathfrak q_1}\).

Fix \(\mathfrak p\in U\) and put
\(\mathfrak q=\mathfrak p\cap A\).  We claim that \(U\cap
V(\mathfrak p)\) contains an open neighborhood of the generic point
\(\mathfrak p\) of \(V(\mathfrak p)\).

Consider the finite \(B\)-module

\[
T=\operatorname{Tor}_1^A(A/\mathfrak q,M).
\tag{3}
\]

It is finite because \(\mathfrak q\) is finitely generated and Tor
in degree one can be computed from a finite presentation of
\(A/\mathfrak q\).  Since \(M_{\mathfrak p}\) is flat,
\(T_{\mathfrak p}=0\).  Hence there is \(s\in B\setminus\mathfrak p\)
such that \(T_s=0\).

Apply generic freeness to the domain \(R=A/\mathfrak q\), the
finite-type \(R\)-algebra \(B/\mathfrak qB\), and the finite module

\[
\overline M=M/\mathfrak qM.
\]

There is \(a\in A\setminus\mathfrak q\) such that
\(\overline M_a\) is free over \((A/\mathfrak q)_a\).
Because \(\mathfrak p\cap A=\mathfrak q\), both \(s\) and \(a\) are
outside \(\mathfrak p\).  Therefore

\[
V(\mathfrak p)\cap D(s)\cap D(a)
\tag{4}
\]

is a nonempty open neighborhood of \(\mathfrak p\) in
\(V(\mathfrak p)\).

Let \(\mathfrak p'\) belong to (4), and put
\(\mathfrak q'=\mathfrak p'\cap A\).  Localization of (3) gives

\[
\operatorname{Tor}_1^{A_{\mathfrak q'}}
\left(A_{\mathfrak q'}/\mathfrak qA_{\mathfrak q'},
M_{\mathfrak p'}\right)=0.
\tag{5}
\]

Moreover,

\[
M_{\mathfrak p'}/\mathfrak qM_{\mathfrak p'}
\cong\overline M_{\mathfrak p'}
\]

is flat over
\(A_{\mathfrak q'}/\mathfrak qA_{\mathfrak q'}\): it is obtained
from the free module \(\overline M_a\) by module localization, which
is a filtered colimit and therefore preserves flatness.  The local
criterion for flatness, applied with the ideal
\(\mathfrak qA_{\mathfrak q'}\), now shows that
\(M_{\mathfrak p'}\) is flat over \(A_{\mathfrak q'}\).  Thus (4) is
contained in \(U\).

We finish with a noetherian topological argument.  In a noetherian
space, a subset \(U\) is open if

\[
\begin{gathered}
U\text{ is stable under generization, and}\\
U\cap\overline{\{u\}}\text{ contains a nonempty open subset of }
\overline{\{u\}}\text{ for every }u\in U.
\end{gathered}
\tag{6}
\]

Indeed, let \(Z\) be the complement of \(U\).  If an irreducible
component of \(\overline Z\) had generic point in \(U\), the second
condition in (6) would give a nonempty open part of that component
disjoint from \(Z\), contradicting the density of \(Z\) there.
Every generic point of \(\overline Z\) therefore lies in \(Z\).
The first condition says that \(Z\) is stable under specialization,
so every such irreducible component is contained in \(Z\).  Hence
\(Z=\overline Z\), proving that \(U\) is open.

The two conditions in (6) have just been verified for (2), so its
flat locus is open.  Finally, cover \(Y\) by affine opens and their
inverse images in \(X\) by affine opens.  On each such pair take
\(M=B\).  The set (2) is then precisely the locus where the local
ring of \(X\) is flat over the corresponding local ring of \(Y\).
These affine flat loci glue, proving the assertion for \(f\).
\end{solution}

\begin{exerciseintro}
% record: H-III-09-005
\exerciseheading{Exercise III.9.5}
\addcontentsline{toc}{subsection}{Exercise III.9.5}

\begin{context}
For a closed subscheme \(Z\subseteq\mathbb P_F^n\), let
\(I_Z\subseteq F[x_0,\ldots,x_n]\) be its saturated homogeneous ideal
and put

\[
h_Z(d)=\dim_F(F[x_0,\ldots,x_n]/I_Z)_d.
\]

The projective cone \(C(Z)\subseteq\mathbb P_F^{n+1}\) is defined by
the extension of \(I_Z\) to
\(F[x_0,\ldots,x_n,x_{n+1}]\).

We use the Hilbert-polynomial flatness criterion for projective
morphisms over an integral noetherian base.  We also use
cohomology and base change in the following form: for a projective
morphism and a coherent sheaf flat over the base, constancy of the
fiber dimensions of \(H^0\) makes the direct image locally free and
its formation compatible with base change.

In (d), the scheme-theoretic fibers themselves are assumed
projectively normal.  This implies
\(H^1(\mathbb P^n_{\kappa(t)},\mathcal I_t(d))=0\) for every \(d\geq0\).
Vanishing in this degree gives cohomology and base change for
\(p_*\mathcal I(d)\).  Over a DVR its torsion-free coherent direct
image is locally free.
\end{context}

\begin{problem}
(a) Give a flat family \(X_\lambda\subseteq\mathbb P^n\) for which
the fiberwise projective cones \(C(X_\lambda)\) cannot be the fibers of a flat projective family.

(b) Let \(T\) be an integral noetherian scheme and
\(X\subseteq\mathbb P_T^n\) a closed subscheme.  For
\(t\in T\), write

\[
S_t=\kappa(t)[x_0,\ldots,x_n]
\]

and let \(I_t\) be the saturated homogeneous ideal of the fiber
\(X_t\).  Call the family very flat if

\[
\dim_{\kappa(t)}(S_t/I_t)_d
\]

is independent of \(t\) for every \(d\geq0\).

(c) Prove that every very flat family is flat.  Prove also that its
family of projective cones is very flat, and therefore flat.

(d) Let \(T\) be an integral nonsingular curve over an algebraically
closed field, and let \(X\subseteq\mathbb P_T^n\) be flat over \(T\).
Assume that every scheme-theoretic fiber is projectively normal.
Prove that this is a very flat family.  Explain why normality of
the reduced members alone would not suffice.
\end{problem}
\end{exerciseintro}

\begin{solution}
(a) Let \(\lambda\) be the coordinate on \(\mathbb A_k^1\), and
consider

\[
\begin{aligned}
\Phi:\mathbb P_k^1\times\mathbb A_k^1
&\longrightarrow\mathbb P_k^4\times\mathbb A_k^1,\\
([u:v],\lambda)
&\longmapsto
([u^4:u^3v:\lambda u^2v^2:uv^3:v^4],\lambda).
\end{aligned}
\tag{1}
\]

This is an immersion.  On \(u\ne0\), the ratio of the second
coordinate to the first recovers \(v/u\), and on \(v\ne0\), the ratio
of the fourth coordinate to the fifth recovers \(u/v\); these two
calculations also identify the maps on local coordinate rings.
Since the source is proper over \(\mathbb A^1\), the immersion is
closed.  Its image \(X\) is isomorphic to
\(\mathbb P^1\times\mathbb A^1\), so \(X\to\mathbb A^1\) is flat.

For \(\lambda\ne0\), the fiber is projectively equivalent to the
rational normal quartic.  Its degree-\(d\) coordinate functions are
all monomials of degree \(4d\) in \(u,v\), and hence

\[
h_{X_\lambda}(d)=4d+1\qquad(d\geq0,\ \lambda\ne0).
\tag{2}
\]

For \(\lambda=0\), the degree-one coordinate functions have
\(v\)-exponents \(0,1,3,4\).  Thus

\[
h_{X_0}(0)=1,\qquad h_{X_0}(1)=4.
\]

Every integer from \(0\) through \(8\) is a sum of two elements of
\(\{0,1,3,4\}\).  Inductively, adding either \(0\) or \(4\) shows
that every integer from \(0\) through \(4d\) occurs for every
\(d\geq2\).  Consequently,

\[
h_{X_0}(d)=4d+1\qquad(d\geq2).
\tag{3}
\]

If \(R_Z\) is the homogeneous coordinate ring of \(Z\), then the
coordinate ring of its projective cone is \(R_Z[x_{n+1}]\), so

\[
h_{C(Z)}(m)=\sum_{d=0}^m h_Z(d).
\tag{4}
\]

Equations (2)--(4) give, for \(m\gg0\),

\[
\begin{aligned}
P_{C(X_\lambda)}(m)
&=\sum_{d=0}^m(4d+1)
=(2m+1)(m+1) &&(\lambda\ne0),\\
P_{C(X_0)}(m)
&=(2m+1)(m+1)-1.
\end{aligned}
\]

The Hilbert polynomials differ, so no flat projective family can
have exactly these fiberwise cones as its scheme fibers.  The
scheme-theoretic closure of the generic cones must have extra
structure in its special fiber; it cannot be identified with
the cone of \(X_0\).

(b) This is the stated definition.  Notice that it asks for
constancy of the entire Hilbert function, not merely the eventual
Hilbert polynomial.

(c) In a very flat family, the fiber Hilbert functions, and hence
their eventual Hilbert polynomials, are independent of \(t\).
The Hilbert-polynomial criterion therefore shows that \(X\to T\) is
flat.

We next construct the cone family with the correct scheme fibers.
Let

\[
p:\mathbb P_T^n\longrightarrow T,\qquad
\mathcal I=\mathcal I_X,
\]

and write

\[
\mathcal S=\operatorname{Sym}_{\mathcal O_T}
(\mathcal O_T^{\oplus(n+1)})
=\bigoplus_{d\geq0}\mathcal S_d.
\]

For every \(d\), the sequence

\[
0\longrightarrow\mathcal I(d)
\longrightarrow\mathcal O_{\mathbb P_T^n}(d)
\longrightarrow\mathcal O_X(d)\longrightarrow0
\tag{5}
\]

has its middle and right terms flat over \(T\), so
\(\mathcal I(d)\) is flat over \(T\).  On a fiber,

\[
\dim_{\kappa(t)}H^0(\mathbb P_{\kappa(t)}^n,\mathcal I_t(d))
=\dim_{\kappa(t)}(I_t)_d
=\dim_{\kappa(t)}(S_t)_d-h_{X_t}(d),
\]

which is constant by very flatness.  Cohomology and base change
applied to (5) therefore shows that

\[
\mathcal J_d:=p_*\mathcal I(d)\subseteq\mathcal S_d
\]

is locally free and satisfies

\[
\mathcal J_d\otimes\kappa(t)\cong(I_t)_d
\tag{6}
\]

for every \(t\).  Multiplication of homogeneous polynomials makes

\[
\mathcal J=\bigoplus_{d\geq0}\mathcal J_d
\]

a homogeneous ideal in \(\mathcal S\).

Introduce a new degree-one variable \(z\) and define

\[
\mathcal C=
\operatorname{Proj}_T
\left(\mathcal S[z]/\mathcal J\mathcal S[z]\right)
\subseteq\mathbb P_T^{n+1}.
\tag{7}
\]

By (6), formation of (7) commutes with every fiber and
\(\mathcal C_t=C(X_t)\) scheme-theoretically.  Moreover,

\[
\left(
\frac{S_t[z]}{I_tS_t[z]}
\right)_m
\cong
\bigoplus_{d=0}^m(S_t/I_t)_d z^{m-d}.
\tag{8}
\]

Its dimension is

\[
h_{C(X_t)}(m)=\sum_{d=0}^m h_{X_t}(d),
\]

which is independent of \(t\) for every \(m\).  Thus the cone family
is very flat.  Its Hilbert polynomial is consequently constant, so
the Hilbert-polynomial criterion also makes
\(\mathcal C\to T\) flat.

(d) Let \(p:\mathbb P_T^n\to T\), with ideal sheaf \(\mathcal I\).
Flatness of \(X\) makes \(\mathcal I(d)\) flat over \(T\).
Projective normality of each scheme fiber gives
\[
H^1(\mathbb P^n_{\kappa(t)},\mathcal I_t(d))=0
\qquad(d\geq0).
\]
Hence \(\mathcal J_d=p_*\mathcal I(d)\) commutes with passage to
the fibers.  This coherent sheaf is torsion-free on the nonsingular
curve, hence locally free.  Its inclusion into the free polynomial
piece \(\mathcal S_d\) is injective on every fiber, so its cokernel
\(\mathcal R_d\) is locally free as well.  Its fiber is
\((S_t/I_t)_d\).  Its rank is constant on the integral curve,
which proves constancy of every Hilbert-function value.

To see the need for the precise fiber assumption, in characteristic
different from two consider
\[
x_0^2-tx_1x_2=0
\quad\text{in }\mathbb P^2\times\mathbb A^1.
\]
This is a flat integral total family.  Its reduced members are
smooth conics for \(t\ne0\) and a line for \(t=0\), all projectively
normal.  Their Hilbert functions are \(2d+1\) and \(d+1\).
The scheme-theoretic special fiber is a double line.  Thus an
assertion about reduced members alone cannot yield very flatness.
\end{solution}

\begin{exerciseintro}
% record: H-III-09-006
\exerciseheading{Exercise III.9.6}
\addcontentsline{toc}{subsection}{Exercise III.9.6}

\begin{context}
Let \(k\) be algebraically closed.  A projective complete
intersection is arithmetically Cohen--Macaulay; this follows from its
Koszul resolution.  In particular, if
\(W\subseteq\mathbb P^N\) is a complete intersection of positive
dimension, then

\[
H^1(\mathbb P^N,\mathcal I_W(m))=0
\qquad(m\in\mathbb Z).
\]

A complete intersection has a Cohen--Macaulay homogeneous coordinate
ring and consequently has no embedded associated points.
\end{context}

\begin{problem}
Let \(Y\subseteq\mathbb P_k^n\) be a nonsingular variety with
\(\dim Y\geq2\).  Let \(H\cong\mathbb P_k^{n-1}\) be a hyperplane
which does not contain \(Y\), and suppose the scheme-theoretic
intersection

\[
Y'=Y\cap H
\]

is nonsingular.  Prove that \(Y\) is a complete intersection in
\(\mathbb P^n\) if and only if \(Y'\) is a complete intersection in
\(H\).
\end{problem}
\end{exerciseintro}

\begin{solution}
Write \(H=V(h)\), where \(h\) is linear.

Suppose first that \(Y\) is a complete intersection of codimension
\(r\).  Its ideal is generated by \(r\) homogeneous equations, and
their restrictions to \(H\) define the scheme \(Y\cap H=Y'\).
Since \(h\) is a nonzerodivisor on the coordinate ring of the
integral variety \(Y\),

\[
\operatorname{codim}_H Y'
=\operatorname{codim}_{\mathbb P^n}Y=r.
\]

Thus the \(r\) restricted equations generate an ideal of height
\(r\) in the Cohen--Macaulay coordinate ring of \(H\), so they form
a regular sequence.  Hence \(Y'\) is a complete intersection.

Conversely, assume \(Y'\) is a complete intersection of codimension
\(r\) in \(H\).  Let \(\mathcal I_Y\) and
\(\mathcal I_{Y'/H}\) denote the corresponding ideal sheaves.
Because \(h\) is a nonzerodivisor on \(\mathcal O_Y\), restriction to
\(H\) gives, for every \(m\), an exact sequence

\[
0\longrightarrow\mathcal I_Y(m-1)
\xrightarrow{\,h\,}\mathcal I_Y(m)
\longrightarrow\mathcal I_{Y'/H}(m)
\longrightarrow0.
\tag{1}
\]

The complete intersection \(Y'\) has dimension
\(\dim Y-1\geq1\), so its Koszul resolution gives

\[
H^1(H,\mathcal I_{Y'/H}(m))=0
\qquad(m\in\mathbb Z).
\tag{2}
\]

We claim that

\[
H^1(\mathbb P^n,\mathcal I_Y(m))=0
\qquad(m\in\mathbb Z).
\tag{3}
\]

For \(m<0\), both
\(H^0(\mathbb P^n,\mathcal O_{\mathbb P^n}(m))\) and
\(H^0(Y,\mathcal O_Y(m))\) vanish; the latter vanishing follows
because the inverse of an ample line bundle has no nonzero section on
a positive-dimensional integral projective variety.  Also
\(H^1(\mathbb P^n,\mathcal O_{\mathbb P^n}(m))=0\), since \(n\geq2\).
The ideal-sheaf sequence for \(Y\) therefore proves (3) for \(m<0\).

Taking cohomology in (1) and using (2) shows that multiplication by
\(h\) induces a surjection

\[
H^1(\mathbb P^n,\mathcal I_Y(m-1))
\twoheadrightarrow
H^1(\mathbb P^n,\mathcal I_Y(m))
\]

for every \(m\).  Starting with the negative-degree vanishing and
inducting upward proves (3) in all degrees.

Let

\[
I_{Y'}=(f_1,\ldots,f_r)
\subseteq k[x_0,\ldots,x_n]/(h)
\]

be the saturated homogeneous ideal of \(Y'\), with the \(f_i\)
forming a regular sequence.  The cohomology sequence of (1), together
with (3), makes

\[
H^0(\mathbb P^n,\mathcal I_Y(d))
\longrightarrow H^0(H,\mathcal I_{Y'/H}(d))
\tag{4}
\]

surjective for every \(d\).  We may therefore lift each \(f_i\) to a
homogeneous

\[
g_i\in H^0(\mathbb P^n,\mathcal I_Y(\deg f_i)).
\]

Put \(S=k[x_0,\ldots,x_n]\), \(J=(g_1,\ldots,g_r)\), and
\[
Z=V(J)\subseteq\mathbb P^n.
\]

Since the images of the \(g_i\) modulo \(h\) are the regular sequence
\(f_1,\ldots,f_r\),

\[
\operatorname{ht}(J,h)=r+1.
\]

Krull's height theorem gives \(\operatorname{ht}J=r\).  The
polynomial ring \(S\) is Cohen--Macaulay, so \(g_1,\ldots,g_r\) is a
regular sequence.  Thus \(Z\) is a complete intersection.  Moreover,
\(h\) is a nonzerodivisor on \(S/J\): otherwise it would lie in a
minimal prime of the unmixed ideal \(J\), contradicting
\(\operatorname{ht}(J,h)=r+1\).

The choices of the \(g_i\) give

\[
Y\subseteq Z,
\qquad
Z\cap H=Y'
\tag{5}
\]

scheme-theoretically.  Indeed, modulo \(h\), both hyperplane-section
ideals are generated by \(f_1,\ldots,f_r\).

Because \(h\) is a nonzerodivisor on both \(Z\) and \(Y\), (5) gives
exact sequences

\[
\begin{aligned}
0&\longrightarrow\mathcal O_Z(m-1)
\xrightarrow{\,h\,}\mathcal O_Z(m)
\longrightarrow\mathcal O_{Y'}(m)\longrightarrow0,\\
0&\longrightarrow\mathcal O_Y(m-1)
\xrightarrow{\,h\,}\mathcal O_Y(m)
\longrightarrow\mathcal O_{Y'}(m)\longrightarrow0.
\end{aligned}
\tag{6}
\]

Taking Hilbert polynomials in (6) shows

\[
P_Z(m)-P_Z(m-1)=P_Y(m)-P_Y(m-1).
\]

Hence \(P_Z-P_Y\) is constant.  From

\[
0\longrightarrow\mathcal K
\longrightarrow\mathcal O_Z
\longrightarrow\mathcal O_Y\longrightarrow0,
\qquad
\mathcal K=\mathcal I_{Y/Z},
\tag{7}
\]

we conclude that either \(\mathcal K=0\), or
\(\mathcal K\) is supported on finitely many closed points.

The second possibility cannot occur.  A nonzero finite-length
submodule of a local ring contains a nonzero socle element, so a
closed point in \(\operatorname{Supp}\mathcal K\) would be an
associated point of \(Z\).  It would be embedded because
\(\dim Z=\dim Y\geq2\).  But the complete intersection \(Z\) is
Cohen--Macaulay and has no embedded associated points.  Therefore
\(\mathcal K=0\), so \(Y=Z\).  Since \(Z\) is a complete
intersection, so is \(Y\).
\end{solution}

\begin{exerciseintro}
% record: H-III-09-007
\exerciseheading{Exercise III.9.7}
\addcontentsline{toc}{subsection}{Exercise III.9.7}

\begin{context}
Let \(D=k[\epsilon]/(\epsilon^2)\).  A \(D\)-module \(M\) is flat if
and only if multiplication by \(\epsilon\) induces an isomorphism

\[
M/\epsilon M\xrightarrow{\ \sim\ }\epsilon M.
\]

Let \(Y\hookrightarrow X\) be a closed immersion with ideal sheaf
\(\mathcal I\).  Its normal sheaf is

\[
\mathcal N_{Y/X}
=\mathcal Hom_{\mathcal O_Y}
(\mathcal I/\mathcal I^2,\mathcal O_Y).
\]
\end{context}

\begin{problem}
Assume \(X\) is of finite type over \(k\).  An embedded
infinitesimal deformation of \(Y\) in \(X\) is a closed subscheme

\[
Y'\subseteq X\times_k\operatorname{Spec}D
\]

which is flat over \(D\) and whose closed fiber is \(Y\).  Prove that
such deformations are naturally classified by

\[
H^0(Y,\mathcal N_{Y/X}).
\]
\end{problem}
\end{exerciseintro}

\begin{solution}
We first work on an affine open \(X=\operatorname{Spec}A\), with
\(Y\) defined by an ideal \(I\subseteq A\).  Put
\(A_D=A[\epsilon]/(\epsilon^2)\).  A deformation is determined by an
ideal \(I'\subseteq A_D\) such that

\[
B=A_D/I'
\]

is flat over \(D\) and the reduction of \(I'\) modulo \(\epsilon\)
is \(I\).

Flatness of \(B\) implies the useful identity

\[
I'\cap\epsilon A_D=\epsilon I.
\tag{1}
\]

Indeed, if \(\epsilon c\in I'\), then the class of \(c\) in \(B\) is
killed by \(\epsilon\).  The flatness criterion says that this class
lies in \(\epsilon B\), so its reduction in
\(B/\epsilon B\cong A/I\) is zero; hence \(c\in I\).  Conversely, if
\(c\in I\), choose \(c+\epsilon d\in I'\) and multiply by
\(\epsilon\) to obtain \(\epsilon c\in I'\).

For \(a\in I\), choose a lift

\[
a+\epsilon b\in I'
\]

and define

\[
\varphi_{I'}(a)=b\bmod I\in A/I.
\tag{2}
\]

This is well defined: two choices differ by
\(\epsilon(b-b')\in I'\), and (1) then gives \(b-b'\in I\).
Multiplying a lift by an element of \(A\) shows that (2) is
\(A\)-linear.  Since \(I\) annihilates \(A/I\), it kills \(I^2\) and
therefore factors uniquely as

\[
\varphi_{I'}:I/I^2\longrightarrow A/I.
\tag{3}
\]

Conversely, start with an \(A/I\)-linear homomorphism
\(\varphi:I/I^2\to A/I\).  Define

\[
I_\varphi=
\left\{
a+\epsilon b\in A_D:
a\in I,\quad b\bmod I=\varphi(a\bmod I^2)
\right\}.
\tag{4}
\]

This is an ideal.  In fact, for \(c+\epsilon d\in A_D\),

\[
(c+\epsilon d)(a+\epsilon b)
=ca+\epsilon(cb+da),
\]

and modulo \(I\) the second coefficient is
\(c\varphi(a)=\varphi(ca)\).
The reduction of \(I_\varphi\) is plainly \(I\), and (4) gives

\[
I_\varphi\cap\epsilon A_D=\epsilon I.
\tag{5}
\]

Let \(B_\varphi=A_D/I_\varphi\).  Multiplication by \(\epsilon\)
therefore induces

\[
B_\varphi/\epsilon B_\varphi
\cong A/I
\xrightarrow{\ \sim\ }
\epsilon B_\varphi,
\qquad
\overline a\longmapsto\epsilon a.
\]

Thus \(B_\varphi\) is flat over \(D\).  Formulas (2) and (4) are
inverse constructions, yielding a natural bijection

\[
\left\{
\begin{matrix}
\text{embedded deformations of }Y\\
\text{inside }\operatorname{Spec}A
\end{matrix}
\right\}
\longleftrightarrow
\operatorname{Hom}_{A/I}(I/I^2,A/I).
\tag{6}
\]

The construction commutes with localization in \(A\): both choosing
the coefficient of \(\epsilon\) in (2) and forming the ideal (4) do.
Consequently, the local maps obtained from a global deformation agree
on overlaps and define a morphism

\[
\mathcal I/\mathcal I^2\longrightarrow\mathcal O_Y.
\]

Conversely, a global morphism of this kind produces through (4)
compatible ideal sheaves on an affine cover, which glue to a unique
closed subscheme \(Y'\subseteq X\times\operatorname{Spec}D\).
Flatness is local, so the glued deformation is flat.  We have
therefore obtained the required natural bijection

\[
\{\text{embedded infinitesimal deformations of }Y\text{ in }X\}
\cong
\operatorname{Hom}_{\mathcal O_Y}
(\mathcal I/\mathcal I^2,\mathcal O_Y)
=H^0(Y,\mathcal N_{Y/X}).
\]
\end{solution}

\begin{exerciseintro}
% record: H-III-09-008
\exerciseheading{Exercise III.9.8}
\addcontentsline{toc}{subsection}{Exercise III.9.8}

\begin{context}
Let \(A\) be a finitely generated \(k\)-algebra, choose a
presentation

\[
0\longrightarrow J\longrightarrow P\longrightarrow A
\longrightarrow0,
\qquad
P=k[x_1,\ldots,x_n],
\]

and let \(D=k[\epsilon]/(\epsilon^2)\).  Embedded first-order
deformations of \(\operatorname{Spec}A\) in
\(\operatorname{Spec}P\) are classified by

\[
\operatorname{Hom}_A(J/J^2,A).
\]

The universal property of Kahler differentials identifies

\[
\operatorname{Hom}_A(\Omega_{P/k}\otimes_PA,A)
\cong\operatorname{Der}_k(P,A).
\]
\end{context}

\begin{problem}
Apply \(\operatorname{Hom}_A(-,A)\) to the conormal sequence

\[
J/J^2\longrightarrow\Omega_{P/k}\otimes_PA
\longrightarrow\Omega_{A/k}\longrightarrow0
\]

and define

\[
T^1(A)=
\operatorname{coker}\left(
\operatorname{Hom}_A(\Omega_{P/k}\otimes_PA,A)
\longrightarrow
\operatorname{Hom}_A(J/J^2,A)
\right).
\]

Prove that \(T^1(A)\) classifies isomorphism classes of flat
\(D\)-algebras \(A'\) equipped with an identification
\(A'/\epsilon A'\cong A\), where isomorphisms must induce the
identity on \(A\).  Deduce that \(T^1(A)\) is independent of the
chosen polynomial presentation.
\end{problem}
\end{exerciseintro}

\begin{solution}
For
\(\varphi\in\operatorname{Hom}_A(J/J^2,A)\), let
\(J_\varphi\subseteq P_D=P[\epsilon]/(\epsilon^2)\) be

\[
J_\varphi=
\left\{
j+\epsilon q:
j\in J,\quad q\bmod J=\varphi(j\bmod J^2)
\right\}.
\tag{1}
\]

Then

\[
A_\varphi=P_D/J_\varphi
\tag{2}
\]

is flat over \(D\) and has closed fiber \(A\).  Every deformation
embedded in \(\operatorname{Spec}P_D\) arises uniquely in this way.

We determine when (2) and \(A_\psi\) are isomorphic as abstract
deformations.  An element

\[
\delta\in
\operatorname{Hom}_A(\Omega_{P/k}\otimes_PA,A)
=\operatorname{Der}_k(P,A)
\]

is determined by \(a_i=\delta(x_i)\in A\).  Choose lifts
\(\widetilde a_i\in P\), and let

\[
\widetilde\delta
=\sum_{i=1}^n\widetilde a_i\frac{\partial}{\partial x_i}
:P\longrightarrow P.
\]

There is a \(D\)-algebra automorphism, reducing to the identity
modulo \(\epsilon\),

\[
\begin{aligned}
\Theta_\delta:P_D&\longrightarrow P_D,\\
p+\epsilon q&\longmapsto
p+\epsilon\bigl(q+\widetilde\delta(p)\bigr).
\end{aligned}
\tag{3}
\]

Its inverse is obtained by replacing
\(\widetilde\delta\) with \(-\widetilde\delta\).  For
\(j+\epsilon q\in J_\varphi\), formula (3) gives

\[
\Theta_\delta(j+\epsilon q)
=j+\epsilon(q+\widetilde\delta(j)).
\]

Modulo \(J\), the new coefficient of \(\epsilon\) is

\[
\varphi(j)+\delta(j).
\]

Consequently,

\[
\Theta_\delta(J_\varphi)=J_{\varphi+\delta|_J},
\tag{4}
\]

where \(\delta|_J\) is exactly the image of \(\delta\) under the map
appearing in the definition of \(T^1(A)\).  Thus elements of
\(\operatorname{Hom}_A(J/J^2,A)\) differing by that image give
isomorphic abstract deformations.

Conversely, suppose

\[
\alpha:A_\varphi\xrightarrow{\ \sim\ }A_\psi
\tag{5}
\]

is a \(D\)-algebra isomorphism inducing the identity on \(A\).
For each polynomial generator, choose \(a_i\in A\) such that

\[
\alpha([x_i])=[x_i+\epsilon a_i].
\tag{6}
\]

Such \(a_i\) exist because the difference in (6) belongs to
\(\epsilon A_\psi\).  Since the \(x_i\) generate \(P_D\) as a
\(D\)-algebra, (6) says that \(\alpha\) is induced by an
automorphism of the form (3), with
\(\delta(x_i)=a_i\).  The condition that this automorphism carry
\(J_\varphi\) to \(J_\psi\) gives

\[
\psi-\varphi=\delta|_J.
\tag{7}
\]

Hence the two embedded classes determine isomorphic abstract
deformations if and only if their difference lies in the image of
\(\operatorname{Der}_k(P,A)\).

It remains to see that every abstract deformation has an embedded
representative.  Let \(B\) be flat over \(D\), together with
\(B/\epsilon B\cong A\).  Lift the images of
\(x_1,\ldots,x_n\) to elements \(b_1,\ldots,b_n\in B\).  They define
a \(D\)-algebra homomorphism

\[
P_D\longrightarrow B.
\tag{8}
\]

This map is surjective.  Indeed, for \(b\in B\), first choose a
polynomial \(p\) such that \(b-p(b_1,\ldots,b_n)\in\epsilon B\).
Write this difference as \(\epsilon c\).  Choose a polynomial \(q\)
whose value modulo \(\epsilon\) is the reduction of \(c\).  Flatness
is not needed for this last equality: the difference
\(c-q(b_1,\ldots,b_n)\) lies in \(\epsilon B\), and hence
\(\epsilon(c-q(b_1,\ldots,b_n))=0\).  Thus

\[
\epsilon c=\epsilon q(b_1,\ldots,b_n).
\]

Thus \(b\) is the image of \(p+\epsilon q\) under (8).

Let \(K\) be the kernel of (8).  Tensoring
\(0\to K\to P_D\to B\to0\) with \(k=D/(\epsilon)\) remains exact on
the left because \(B\) is \(D\)-flat.  It follows that
\(K/\epsilon K\) maps isomorphically to the kernel \(J\) of
\(P\to A\).  Thus \(K\) is an embedded deformation ideal, so it is
\(J_\varphi\) for a unique
\(\varphi\in\operatorname{Hom}_A(J/J^2,A)\).  Combining this fact
with (4)--(7) produces a bijection

\[
T^1(A)
\longleftrightarrow
\left\{
\begin{matrix}
\text{isomorphism classes of flat \(D\)-algebras \(A'\)}\\
\text{equipped with \(A'/\epsilon A'\cong A\)}
\end{matrix}
\right\}.
\]

The right-hand side is intrinsic to \(A\).  Moreover, addition of
first-order coefficients in (1) agrees with addition in the
cokernel.  Therefore any two polynomial presentations give
canonically isomorphic \(k\)-vector spaces \(T^1(A)\), proving the
claimed independence.
\end{solution}

\begin{exerciseintro}
% record: H-III-09-009
\exerciseheading{Exercise III.9.9}
\addcontentsline{toc}{subsection}{Exercise III.9.9}

\begin{context}
For a presentation \(A=P/J\), where \(P\) is a polynomial
\(k\)-algebra, first-order abstract deformations of \(A\) are
classified by

\[
T^1(A)=
\operatorname{coker}\left(
\operatorname{Der}_k(P,A)
\longrightarrow
\operatorname{Hom}_A(J/J^2,A)
\right).
\]

The algebra \(A\) is rigid if \(T^1(A)=0\).
\end{context}

\begin{problem}
Let

\[
A=
\frac{k[x,y,z,w]}{(x,y)\cap(z,w)}.
\]

This is the union of two coordinate planes in \(\mathbb A_k^4\)
meeting at the origin.  Prove that \(A\) is rigid.
\end{problem}
\end{exerciseintro}

\begin{solution}
Put \(P=k[x,y,z,w]\).  Since the variables in the two ideals are
disjoint,

\[
J=(x,y)\cap(z,w)=(xz,xw,yz,yw).
\tag{1}
\]

Write \(e_{xz},e_{xw},e_{yz},e_{yw}\) for these four generators.
Their first syzygies are generated by

\[
\begin{aligned}
y e_{xz}-x e_{yz}&=0,&
y e_{xw}-x e_{yw}&=0,\\
w e_{xz}-z e_{xw}&=0,&
w e_{yz}-z e_{yw}&=0.
\end{aligned}
\tag{2}
\]

For example, this presentation follows by tensoring the two Koszul
presentations of the ideals \((x,y)\) and \((z,w)\).  Tensoring the
resulting presentation of \(J\) with \(A=P/J\) presents
\(J/J^2\), so (2) gives all the conditions on a homomorphism
\(J/J^2\to A\).

It is useful to write

\[
A=
\left\{
(F,G)\in k[x,y]\times k[z,w]:
F(0,0)=G(0,0)
\right\}.
\tag{3}
\]

Let

\[
\varphi\in\operatorname{Hom}_A(J/J^2,A),
\]

and denote the images of the four generators by
\(a_{xz},a_{xw},a_{yz},a_{yw}\).  The first two relations in (2),
applied to the \(k[x,y]\)-components in (3), imply that there are
\(f_z,f_w\in k[x,y]\) such that

\[
\begin{aligned}
(a_{xz})_1&=x f_z,&(a_{yz})_1&=y f_z,\\
(a_{xw})_1&=x f_w,&(a_{yw})_1&=y f_w.
\end{aligned}
\tag{4}
\]

Indeed, for example
\(y(a_{xz})_1=x(a_{yz})_1\), and \(x,y\) are relatively prime.
The last two relations in (2), applied to the \(k[z,w]\)-components,
similarly give \(g_x,g_y\in k[z,w]\) such that

\[
\begin{aligned}
(a_{xz})_2&=z g_x,&(a_{xw})_2&=w g_x,\\
(a_{yz})_2&=z g_y,&(a_{yw})_2&=w g_y.
\end{aligned}
\tag{5}
\]

Thus every \(\varphi\) has the form

\[
\begin{aligned}
a_{xz}&=(xf_z,z g_x),&
a_{xw}&=(xf_w,w g_x),\\
a_{yz}&=(yf_z,z g_y),&
a_{yw}&=(yf_w,w g_y).
\end{aligned}
\tag{6}
\]

We now construct a derivation inducing it.  Since \(P\) is
polynomial, prescribe

\[
\begin{aligned}
\delta(x)&=(g_x(0,0),g_x),&
\delta(y)&=(g_y(0,0),g_y),\\
\delta(z)&=(f_z,f_z(0,0)),&
\delta(w)&=(f_w,f_w(0,0)).
\end{aligned}
\tag{7}
\]

Each pair in (7) belongs to \(A\) by (3), and the assignments extend
uniquely to a \(k\)-derivation \(\delta:P\to A\).  The Leibniz rule,
together with the fact that every product of an element of
\((x,y)\) and an element of \((z,w)\) is zero in \(A\), gives

\[
\begin{aligned}
\delta(xz)&=(xf_z,z g_x),&
\delta(xw)&=(xf_w,w g_x),\\
\delta(yz)&=(yf_z,z g_y),&
\delta(yw)&=(yf_w,w g_y).
\end{aligned}
\]

By (6), the restriction of \(\delta\) to \(J/J^2\) is exactly
\(\varphi\).  Therefore

\[
\operatorname{Der}_k(P,A)
\longrightarrow\operatorname{Hom}_A(J/J^2,A)
\]

is surjective.  Its cokernel \(T^1(A)\) is zero, so \(A\) is rigid.
\end{solution}

\begin{exerciseintro}
% record: H-III-09-010
\exerciseheading{Exercise III.9.10}
\addcontentsline{toc}{subsection}{Exercise III.9.10}

\begin{context}
For a nonsingular \(k\)-scheme \(Z\), first-order deformations are
classified by \(H^1(Z,\mathcal T_Z)\).  Smoothness is open, and a
flat morphism of finite presentation is smooth at a point when its
geometric fiber is regular there.  A smooth morphism has a section
through a chosen rational point after an etale base change.

We also use cohomology and base change for a proper flat family of
curves: if a line bundle has constant \(h^0\) and vanishing \(h^1\)
on the fibers, its direct image is locally free, commutes with base
change, and its evaluation map may be checked fiberwise.
\end{context}

\begin{problem}
A scheme over a field is called rigid if it has no nontrivial
first-order deformations.

(a) Prove that \(\mathbb P_k^1\) is rigid.

(b) Let \(k\) be algebraically closed.  Construct a proper flat
morphism

\[
f:X\longrightarrow\mathbb A_k^2
\]

whose fiber \(X_0\) over the origin is isomorphic to
\(\mathbb P_k^1\), but for which there is no open neighborhood
\(U\) of the origin satisfying
\(X_U\cong U\times\mathbb P^1\).

(c) Let \(f:X\to T\) be flat and projective, where \(T\) is a
nonsingular curve over an algebraically closed field.  Suppose
\(X_t\cong\mathbb P_k^1\) for a closed point \(t\in T\).  Prove that
there are a nonsingular curve \(T'\) and a flat morphism
\(g:T'\to T\), whose image contains \(t\), such that

\[
X\times_TT'\cong\mathbb P_{T'}^1.
\]
\end{problem}
\end{exerciseintro}

\begin{solution}
(a) The tangent bundle of the projective line is

\[
\mathcal T_{\mathbb P^1}\cong\mathcal O_{\mathbb P^1}(2).
\]

Therefore

\[
H^1(\mathbb P^1,\mathcal T_{\mathbb P^1})
=H^1(\mathbb P^1,\mathcal O_{\mathbb P^1}(2))=0.
\]

The deformation classification in the context shows that
\(\mathbb P^1\) is rigid.

(b) Write \(a,b\) for the coordinates on \(\mathbb A^2\), and
\([x:y:z]\) for those on \(\mathbb P^2\).  If
\(\operatorname{char}k\ne2\), let

\[
X=V\bigl(x^2+(1+a)y^2+(1+b)z^2\bigr)
\subseteq\mathbb P_k^2\times\mathbb A_k^2.
\tag{1}
\]

If \(\operatorname{char}k=2\), use instead

\[
X=V\bigl(x^2+(1+a)xy+(1+a)y^2+(1+b)z^2\bigr).
\tag{2}
\]

In either case the defining equation is monic in \(x^2\).  Thus every
graded piece of its homogeneous coordinate ring is a free
\(k[a,b]\)-module, with representatives having \(x\)-degree at most
one.  Hence \(f:X\to\mathbb A^2\) is flat.  It is projective, and
therefore proper.

At \(a=b=0\), equation (1) is a nonsingular conic when
\(\operatorname{char}k\ne2\).  In characteristic \(2\), the partial
derivatives of the equation

\[
x^2+xy+y^2+z^2
\]

are \(y,x,0\); their only simultaneous affine zero has
\(x=y=0\), and the equation would then force \(z=0\), so the
projective conic is again nonsingular.  A nonsingular conic over an
algebraically closed field is isomorphic to \(\mathbb P^1\).

Let \(K=k(a,b)\).  We show that the generic conic has no \(K\)-point.
Suppose first that the characteristic is not \(2\).  Clearing
denominators in a hypothetical point gives relatively prime
polynomials \(X_1,Y_1,Z_1\in k[a,b]\), not all zero, satisfying

\[
X_1^2+(1+a)Y_1^2+(1+b)Z_1^2=0.
\tag{3}
\]

Set \(a=-1\).  Since \(-(1+b)\) is not a square in \(k(b)\), equation
(3) implies

\[
X_1(-1,b)=Z_1(-1,b)=0.
\]

Thus \(1+a\) divides \(X_1\) and \(Z_1\).  Substitution back into
(3) then shows that \(1+a\) divides \(Y_1^2\), and hence \(Y_1\).
This contradicts the choice of a relatively prime triple.

In characteristic \(2\), the same argument works with (2).  After
setting \(a=-1\), the equation is

\[
X_1^2+(1+b)Z_1^2=0.
\]

The element \(1+b\) is not a square in \(k(b)\), so \(1+a\) divides
\(X_1\) and \(Z_1\).  Substitution into

\[
X_1^2+(1+a)X_1Y_1+(1+a)Y_1^2+(1+b)Z_1^2=0
\]

then forces \(1+a\mid Y_1^2\), and hence
\(1+a\mid Y_1\), giving the same contradiction.

If \(X_U\) were isomorphic to \(U\times\mathbb P^1\) on any open
neighborhood \(U\) of the origin, then \(U\) would contain the
generic point of \(\mathbb A^2\), and the generic conic would have a
\(K\)-rational point.  The contradiction proves the required
nontriviality.

(c) We may replace \(T\) by an open neighborhood of \(t\).  Because
\(f\) is flat and the fiber \(X_t\) is smooth, \(f\) is smooth at
every point of \(X_t\).  The nonsmooth locus is closed in \(X\), and
its image is closed because \(f\) is projective.  After shrinking
\(T\), we may therefore assume that \(f\) is a smooth proper family
of curves.  The arithmetic genus and the number of geometric
connected components are locally constant in such a family, so,
after one more shrinking, its fibers are geometrically connected
curves of genus \(0\).

Choose \(x\in X_t(k)\).  Smooth morphisms admit sections etale-locally,
so there are an etale morphism

\[
g_0:T_0\longrightarrow T,
\]

a point \(t_0\in T_0\) above \(t\), and a section

\[
s:T_0\longrightarrow X_0:=X\times_TT_0
\]

through \(x\).  Replace \(T_0\) by a sufficiently small affine open
neighborhood of \(t_0\).  It remains a nonsingular curve, and
\(g_0\) is flat.

The image \(s(T_0)\) is an effective Cartier divisor on the smooth
relative curve \(X_0/T_0\).  Put

\[
\mathcal L=\mathcal O_{X_0}(s(T_0)).
\]

On every geometric fiber, \(\mathcal L\) has degree \(1\).  A
geometrically connected smooth genus-zero curve with a rational point
is a projective line, so

\[
\mathcal L|_{(X_0)_u}\cong\mathcal O_{\mathbb P^1}(1),
\qquad
h^0((X_0)_u,\mathcal L_u)=2,\quad
h^1((X_0)_u,\mathcal L_u)=0.
\]

Cohomology and base change now makes

\[
\mathcal E=(f_0)_*\mathcal L
\]

a locally free sheaf of rank \(2\), and the evaluation map
\[
f_0^*\mathcal E\longrightarrow\mathcal L
\]

is surjective.  It defines a \(T_0\)-morphism

\[
\Phi:X_0\longrightarrow\mathbb P_{T_0}(\mathcal E)
\tag{4}
\]

whose restriction to every geometric fiber is the isomorphism
defined by \(\mathcal O_{\mathbb P^1}(1)\).  The morphism \(\Phi\) is
proper, since \(X_0\) is proper over \(T_0\) and the target is
separated over \(T_0\).  It is quasi-finite because this may be
checked on the geometric fibers, so it is finite.  Locally on the
target, its finite algebra map becomes an isomorphism after tensoring
with every residue field of \(T_0\).  Surjectivity follows from
Nakayama's lemma; because both source and target are flat over
\(T_0\), formation of the kernel commutes with passage to a fiber,
and a second application of Nakayama shows that the kernel is zero.
Consequently (4) identifies
\[
X_0\cong\mathbb P_{T_0}(\mathcal E).
\]

Finally, shrink \(T_0\) once more around \(t_0\) so that the rank-two
bundle \(\mathcal E\) is trivial.  Calling this open curve \(T'\), we
obtain

\[
X\times_TT'
\cong\mathbb P_{T'}(\mathcal O_{T'}^{\oplus2})
\cong\mathbb P_{T'}^1.
\]

The composite \(g:T'\to T\) is flat and its image contains \(t\), as
required.
\end{solution}

\begin{exerciseintro}
% record: H-III-09-011
\exerciseheading{Exercise III.9.11}
\addcontentsline{toc}{subsection}{Exercise III.9.11}

\begin{context}
For an integral projective curve \(Z\) over an algebraically closed
field,

\[
p_a(Z)=1-\chi(\mathcal O_Z).
\]

A finite birational morphism from a normal curve to an integral curve
is its normalization.  If \(\nu:\widetilde C\to C\) is the
normalization of a projective curve, then

\[
0\longrightarrow\mathcal O_C
\longrightarrow\nu_*\mathcal O_{\widetilde C}
\longrightarrow\mathcal Q\longrightarrow0,
\]

where \(\mathcal Q\) has finite length.  We also use the usual
dimension formula for Grassmannians and the fact that a proper
quasi-finite morphism is finite.
\end{context}

\begin{problem}
Let \(Y\) be a nonsingular curve of degree \(d\) in
\(\mathbb P_k^n\), where \(k\) is algebraically closed.  Prove that

\[
0\leq p_a(Y)\leq\frac{1}{2}(d-1)(d-2).
\]
\end{problem}
\end{exerciseintro}

\begin{solution}
Because \(Y\) is integral and projective,
\(H^0(Y,\mathcal O_Y)=k\).  Hence

\[
p_a(Y)=1-\chi(\mathcal O_Y)=h^1(Y,\mathcal O_Y)\geq0.
\tag{1}
\]

For the upper bound, replace the ambient projective space by the
linear span of \(Y\), say \(\mathbb P^r\).  This does not change
the degree.  If \(r=1\), then \(Y=\mathbb P^1\), and the result is
immediate.  If \(r=2\), then \(Y\) is a plane curve of degree \(d\),
and the standard plane-curve computation already gives

\[
p_a(Y)=\frac{(d-1)(d-2)}{2}.
\]

Assume now that \(r\geq3\).  We first construct a suitable linear
projection.  Let \(V\) be a vector space of dimension \(r+1\), so
that \(\mathbb P^r=\mathbb P(V)\), and let

\[
G=\operatorname{Gr}(r-2,V).
\]

A point \(W\in G\) corresponds to a prospective center
\(\Lambda=\mathbb P(W)\cong\mathbb P^{r-3}\), and

\[
\dim G=3(r-2).
\tag{2}
\]

For a fixed line \(L\subseteq\mathbb P(V)\), the locus of
\(W\in G\) for which \(\mathbb P(W)\cap L\ne\varnothing\) has
dimension at most

\[
1+\dim\operatorname{Gr}(r-3,V/kq)
=1+3(r-3)=3r-8=\dim G-2.
\tag{3}
\]

Here \(q\) varies on \(L\), and the Grassmannian parametrizes the
subspaces \(W\) containing the corresponding one-dimensional
subspace \(kq\).

Apply (3) first to the embedded tangent lines \(T_yY\), with
\(y\in Y\).  The incidence variety of pairs \((y,W)\) for which
\(\mathbb P(W)\) meets \(T_yY\) has dimension at most
\(1+(\dim G-2)=\dim G-1\).  Thus a general center meets no tangent
line to \(Y\).  The analogous incidence for secant lines,

\[
\left\{(y,z,W):y\ne z,
\ \mathbb P(W)\cap\overline{yz}\ne\varnothing\right\},
\]

has dimension at most \(2+(\dim G-2)=\dim G\).  The fiber-dimension
theorem therefore shows that, for general \(W\), only finitely many
ordered pairs \((y,z)\) occur.  Finally, the centers meeting \(Y\)
form a proper closed subset of \(G\): for fixed \(y\), the locus of
centers containing \(y\) has dimension \(3(r-3)=\dim G-3\).

Choose a center \(\Lambda\) satisfying all three conclusions.  The
linear projection from \(\Lambda\) restricts to a morphism

\[
\pi:Y\longrightarrow C\subseteq\mathbb P^2.
\tag{4}
\]

The morphism is quasi-finite: distinct points in one fiber determine
a secant meeting \(\Lambda\), and there are only finitely many such
pairs.  Since it is also projective, it is finite.  Moreover, its
differential is nonzero at every point, because \(\Lambda\) meets no
embedded tangent line.  Thus the finite extension
\(k(C)\subseteq k(Y)\) is separable.  If its degree were greater than
one, then over a dense open subset of \(C\) every geometric fiber
would contain at least two distinct points, contradicting the
finiteness of the set of secant pairs above.  Consequently \(\pi\)
is birational.

The image \(C\) is therefore an integral plane curve.  Since

\[
\pi^*\mathcal O_C(1)=\mathcal O_Y(1),
\]

the degree formula for a finite map of curves gives

\[
\deg C=\deg\mathcal O_C(1)
=\deg\mathcal O_Y(1)=d.
\tag{5}
\]

Because \(Y\) is nonsingular, it is normal, so the finite birational
map (4) is the normalization of \(C\).  We consequently have an
exact sequence

\[
0\longrightarrow\mathcal O_C
\longrightarrow\pi_*\mathcal O_Y
\longrightarrow\mathcal Q\longrightarrow0
\tag{6}
\]

with \(\mathcal Q\) of finite length.  Taking Euler
characteristics, and using that a finite morphism has no higher
direct images, yields

\[
p_a(Y)=p_a(C)-\operatorname{length}\mathcal Q
\leq p_a(C).
\tag{7}
\]

Finally, an integral plane curve of degree \(d\) has ideal sheaf
\(\mathcal O_{\mathbb P^2}(-d)\).  The sequence

\[
0\longrightarrow\mathcal O_{\mathbb P^2}(-d)
\longrightarrow\mathcal O_{\mathbb P^2}
\longrightarrow\mathcal O_C\longrightarrow0
\]

gives

\[
p_a(C)=\frac{(d-1)(d-2)}{2}.
\tag{8}
\]

Combining (1), (7), and (8) proves both inequalities.
\end{solution}

\begin{exerciseintro}
\corpussection{III.10 Smooth Morphisms}
\addcontentsline{toc}{section}{III.10 Smooth Morphisms}

% record: H-III-10-001
\exerciseheading{Exercise III.10.1}
\addcontentsline{toc}{subsection}{Exercise III.10.1}

\begin{context}
A localization of a regular ring is regular.  A scheme of finite type
over a field is smooth only if it remains regular after every field
extension; equivalently, it is geometrically regular.
\end{context}

\begin{problem}
Let \(k_0\) be a field of characteristic \(p>0\), put
\(k=k_0(t)\), and let

\[
X=V(y^2-x^p+t)\subseteq\mathbb A_k^2.
\]

Show that every local ring of \(X\) is regular, but that \(X\) is
not smooth over \(k\).
\end{problem}
\end{exerciseintro}

\begin{solution}
Introduce an indeterminate \(T\) over \(k_0\), and map

\[
k_0[T]\longrightarrow k_0[x,y],
\qquad T\longmapsto x^p-y^2.
\]

This map is injective, and the coordinate ring of \(X\) can be
written as

\[
\begin{aligned}
A
&=k_0(T)[x,y]/(y^2-x^p+T)\\
&\cong k_0[x,y]\otimes_{k_0[T]}k_0(T).
\end{aligned}
\tag{1}
\]

Thus \(A\) is obtained from the polynomial ring \(k_0[x,y]\) by
inverting the nonzero elements of the subring
\(k_0[x^p-y^2]\).  Since a polynomial ring over a field is regular
and regularity is preserved by localization, \(A\) is a regular
ring.  Consequently every local ring of \(X=\operatorname{Spec}A\)
is regular.

Now let

\[
K=k(t^{1/p}),\qquad \alpha=t^{1/p}\in K.
\]

After this field extension the equation becomes

\[
y^2=(x-\alpha)^p.
\tag{2}
\]

At the \(K\)-rational point \((x,y)=(\alpha,0)\), put
\(u=x-\alpha\).  The local ring is the localization at \((u,y)\)
of

\[
K[u,y]/(y^2-u^p).
\]

It has Krull dimension \(1\), whereas its maximal ideal has two
independent generators modulo its square, because
\(y^2-u^p\in(u,y)^2\).  Its embedding dimension is therefore \(2\),
so this local ring is not regular.  Hence \(X_K\) is not regular.
Smoothness is preserved by base change and a smooth scheme over a
field is regular, so \(X\) cannot be smooth over \(k\).
\end{solution}

\begin{exerciseintro}
% record: H-III-10-002
\exerciseheading{Exercise III.10.2}
\addcontentsline{toc}{subsection}{Exercise III.10.2}

\begin{context}
For a flat morphism of finite presentation, a point of the source is
a smooth point of the morphism if and only if the corresponding
geometric fiber is regular there.  The smooth locus is open.  A
proper morphism is closed.
\end{context}

\begin{problem}
Let \(f:X\to Y\) be a proper flat morphism of varieties over \(k\).
Suppose that, for some \(y\in Y\), the fiber \(X_y\) is smooth over
\(k(y)\).  Show that there is an open neighborhood \(U\) of \(y\)
such that

\[
f^{-1}(U)\longrightarrow U
\]

is smooth.
\end{problem}
\end{exerciseintro}

\begin{solution}
Let \(X^{\mathrm{sm}}\subseteq X\) be the smooth locus of \(f\).
It is open.  Since \(f\) is flat and of finite presentation, the
fiberwise criterion for smoothness applies.  The hypothesis that
\(X_y\) is smooth over \(k(y)\) therefore gives

\[
X_y\subseteq X^{\mathrm{sm}}.
\tag{1}
\]

Set \(Z=X\setminus X^{\mathrm{sm}}\).  This is closed, and its image
\(f(Z)\) is closed because \(f\) is proper.  By (1), the point
\(y\) does not belong to \(f(Z)\).  Hence

\[
U=Y\setminus f(Z)
\]

is an open neighborhood of \(y\).  Its inverse image is contained
in the smooth locus, so the restricted morphism
\(f^{-1}(U)\to U\) is smooth.
\end{solution}

\begin{exerciseintro}
% record: H-III-10-003
\exerciseheading{Exercise III.10.3}
\addcontentsline{toc}{subsection}{Exercise III.10.3}

\begin{context}
For a morphism locally of finite type, formation of K\"ahler
differentials commutes with base change.  We use the following local
criterion.  If \(B\) is a local algebra essentially of finite type
over a field \(K\), with maximal ideal \(\mathfrak n\) and residue
field \(L\), then

\[
\Omega_{B/K}=0
\quad\Longleftrightarrow\quad
\mathfrak n=0\ \text{ and }\ L/K\ \text{is finite separable}.
\tag{*}
\]

For completeness, if the left side vanishes, then the quotient map
to \(L\) gives \(\Omega_{L/K}=0\); the separating-basis theorem says
that \(L/K\) is finite separable.  Separability provides a
\(K\)-algebra section

\[
L\longrightarrow B/\mathfrak n^2.
\]

Thus \(B/\mathfrak n^2=L\oplus\mathfrak n/\mathfrak n^2\), and
projection onto the second summand is a \(K\)-derivation.  But
\(\Omega_{(B/\mathfrak n^2)/K}=0\), as follows from the quotient
sequence for differentials, so this projection is zero.  Hence
\(\mathfrak n/\mathfrak n^2=0\), and Nakayama's lemma gives
\(\mathfrak n=0\).  The reverse implication is the familiar
vanishing of differentials for a finite separable field extension.
\end{context}

\begin{problem}
A morphism \(f:X\to Y\) of schemes of finite type over \(k\) is
called etale if it is smooth of relative dimension \(0\).  It is
called unramified if, for every \(x\in X\), with
\(y=f(x)\),

\[
\mathfrak m_y\mathcal O_{X,x}=\mathfrak m_x
\]

and \(k(x)/k(y)\) is a separable algebraic extension.  Prove that the
following are equivalent:

\begin{enumerate}
\item[(i)] \(f\) is etale;
\item[(ii)] \(f\) is flat and \(\Omega_{X/Y}=0\);
\item[(iii)] \(f\) is flat and unramified.
\end{enumerate}
\end{problem}
\end{exerciseintro}

\begin{solution}
Suppose first that \(f\) is etale.  A smooth morphism is flat, and
for a smooth morphism of relative dimension \(n\), the sheaf
\(\Omega_{X/Y}\) is locally free of rank \(n\).  Here \(n=0\), so
\(\Omega_{X/Y}=0\).  Thus (i) implies (ii).

Assume (ii), and fix \(x\in X\), with \(y=f(x)\).  The local ring of
the fiber at \(x\) is

\[
B=\mathcal O_{X,x}/\mathfrak m_y\mathcal O_{X,x},
\]

and its residue field is \(k(x)\).  Base change for differentials
gives

\[
\Omega_{B/k(y)}
\cong
(\Omega_{X/Y})_x\otimes_{\mathcal O_{X,x}}B=0.
\]

Criterion (*) shows both that

\[
\mathfrak m_x/\mathfrak m_y\mathcal O_{X,x}=0
\]

and that \(k(x)/k(y)\) is finite separable.  Hence \(f\) is
unramified, proving (iii).

Finally, assume (iii).  At the same point \(x\), the equality of
maximal ideals says that the local ring of the fiber is precisely

\[
B=k(x),
\]

and this is a finite separable extension of \(k(y)\).  Thus every
fiber is zero-dimensional and is locally the spectrum of a finite
separable residue-field extension.  After any extension of
\(k(y)\) to an algebraic closure, the fiber becomes a disjoint union
of reduced points.  In other words, all geometric fibers are regular
of dimension \(0\).

The morphism is flat by hypothesis and is of finite presentation
because the schemes are of finite type over \(k\).  The fiberwise
criterion for smoothness now shows that \(f\) is smooth of relative
dimension \(0\), hence etale.  This proves (i).
\end{solution}

\begin{exerciseintro}
% record: H-III-10-004
\exerciseheading{Exercise III.10.4}
\addcontentsline{toc}{subsection}{Exercise III.10.4}

\begin{context}
Completion of a noetherian local ring is faithfully flat, and
flatness descends through a faithfully flat algebra.  We use the
criterion that a finite-type morphism is etale exactly when it is
flat and unramified.

If \((R,\mathfrak m)\) is local and \(M\) is a finite flat
\(R\)-module, then \(M\) is free.  Its rank is the dimension of
\(M/\mathfrak mM\) over the residue field.
Equicharacteristic complete noetherian local rings admit coefficient
fields.  If the residue extension of a local homomorphism is finite
separable, a chosen coefficient field of the completed base has a
compatible extension in the completed source: lift a primitive
element by Hensel's lemma applied to its separable minimal polynomial.
No compatibility is asserted for arbitrary inseparable residue
extensions.
\end{context}

\begin{problem}
Let \(f:X\to Y\) be a morphism of schemes of finite type over \(k\).
For \(x\in X\), put \(y=f(x)\), and let
\(\widehat{\mathcal O}_{X,x}\) and
\(\widehat{\mathcal O}_{Y,y}\) be the completed local rings.
Prove that \(f\) is etale if and only if, for every \(x\), the
extension \(k(x)/k(y)\) is finite separable and, after choosing a
coefficient field in the completed base and its compatible finite
separable lift in the completed source, the natural map
\[
\widehat{\mathcal O}_{Y,y}\otimes_{k(y)}k(x)
\longrightarrow\widehat{\mathcal O}_{X,x}
\tag{1}
\]
is an isomorphism.
\end{problem}
\end{exerciseintro}

\begin{solution}
Fix \(x\in X\), set \(y=f(x)\), and abbreviate

\[
A=\mathcal O_{Y,y},\quad B=\mathcal O_{X,x},\quad
\kappa=k(y),\quad\lambda=k(x).
\]

Write \(\mathfrak m_A\) and \(\mathfrak m_B\) for the maximal
ideals.

Suppose first that \(\lambda/\kappa\) is separable algebraic and
that (1) is an isomorphism.  Since \(f\) is of finite type,
\(\lambda/\kappa\) is in fact finite.  The maximal ideal of
\(\widehat A\otimes_\kappa\lambda\) is
\(\widehat{\mathfrak m}_A\otimes_\kappa\lambda\).  Consequently (1)
gives

\[
\mathfrak m_A\widehat B=\mathfrak m_B\widehat B.
\]

The map \(B\to\widehat B\) is faithfully flat, so contraction to
\(B\) yields

\[
\mathfrak m_AB=\mathfrak m_B.
\tag{2}
\]

Together with the separability of \(\lambda/\kappa\), equation (2)
says that \(A\to B\) is unramified.

Moreover,

\[
\widehat A\longrightarrow
\widehat A\otimes_\kappa\lambda\cong\widehat B
\]

is flat.  Since \(A\to\widehat A\) is flat, the composite
\(A\to\widehat B\) is flat.  Faithfully flat descent along
\(B\to\widehat B\) now shows that \(A\to B\) is flat.  This holds
at every \(x\), so \(f\) is flat and unramified, hence etale.

Conversely, suppose that \(f\) is etale.  Then \(A\to B\) is flat
and unramified.  In particular,

\[
\mathfrak m_B=\mathfrak m_AB
\quad\text{and}\quad
\lambda/\kappa\text{ is finite separable}.
\tag{3}
\]

For each \(n\geq1\), put

\[
A_n=A/\mathfrak m_A^n,
\qquad
B_n=B/\mathfrak m_B^n.
\]

By (3), \(\mathfrak m_B^n=\mathfrak m_A^nB\), and hence

\[
B_n=B\otimes_AA_n.
\]

Thus \(B_n\) is flat over the local Artinian ring \(A_n\).  It is
also finite over \(A_n\): it is a zero-dimensional local algebra
essentially of finite type whose residue field \(\lambda\) is finite
over \(\kappa\).  Therefore \(B_n\) is finite free over \(A_n\), of
rank

\[
\dim_\kappa(B_n/\mathfrak m_AB_n)
=\dim_\kappa\lambda=[\lambda:\kappa].
\tag{4}
\]

The compatible coefficient fields define an \(A_n\)-linear map

\[
A_n\otimes_\kappa\lambda\longrightarrow B_n.
\tag{5}
\]

Modulo the maximal ideal of \(A_n\), map (5) is the identity of
\(\lambda\), so Nakayama's lemma makes it surjective.  By (4), its
source and target are finite free \(A_n\)-modules of the same rank;
therefore (5) is an isomorphism.

Taking inverse limits over \(n\), and using the finiteness of
\(\lambda/\kappa\), gives

\[
\widehat A\otimes_\kappa\lambda
\cong
\varprojlim_n(A_n\otimes_\kappa\lambda)
\cong
\varprojlim_n B_n
=\widehat B.
\]

This is precisely (1), and the proof is complete.
\end{solution}

\begin{exerciseintro}
% record: H-III-10-005
\exerciseheading{Exercise III.10.5}
\addcontentsline{toc}{subsection}{Exercise III.10.5}

\begin{context}
A flat local homomorphism of local rings is faithfully flat.  Flatness
of modules descends along faithfully flat maps, and a finitely
presented flat module is projective.  A finite projective module over
a local ring is free.
\end{context}

\begin{problem}
An etale neighborhood of a point \(x\) of a scheme \(X\) is an etale
morphism \(f:U\to X\), together with a point \(x'\in U\) satisfying
\(f(x')=x\).

Let \(\mathcal F\) be a coherent sheaf on \(X\).  Suppose every
point of \(X\) has an etale neighborhood \(f:U\to X\) for which
\(f^*\mathcal F\) is a free \(\mathcal O_U\)-module.  Prove that
\(\mathcal F\) is locally free on \(X\).
\end{problem}
\end{exerciseintro}

\begin{solution}
Fix \(x\in X\), and choose such a neighborhood \((U,x')\).  Set

\[
A=\mathcal O_{X,x},\qquad
B=\mathcal O_{U,x'},\qquad
M=\mathcal F_x.
\]

Because \(f\) is etale, the local homomorphism \(A\to B\) is flat.
It is a local homomorphism, and is therefore faithfully flat.  The
hypothesis gives

\[
M\otimes_AB
\cong(f^*\mathcal F)_{x'},
\]

which is a finite free \(B\)-module.  It is consequently flat over
\(B\), and hence over \(A\).  Faithfully flat descent now implies
that \(M\) is flat over \(A\).

The sheaf \(\mathcal F\) is coherent, so \(M\) is a finitely
presented \(A\)-module.  A finitely presented flat module is
projective, and over the local ring \(A\) it is therefore free, say
of rank \(r\).

Choose sections near \(x\) representing a basis of \(M\).  They
give a homomorphism

\[
\mathcal O_X^{\oplus r}\longrightarrow\mathcal F
\]

which is an isomorphism on the stalk at \(x\).  Its finite-type
cokernel vanishes after shrinking around \(x\), giving a surjection.
Since \(\mathcal F\) is finitely presented, the kernel of this
surjection from a finite free sheaf is of finite type.
Its zero stalk at \(x\) therefore makes it vanish after a further
shrinking.  Thus \(\mathcal F\) is free on
that neighborhood.  Since \(x\) was arbitrary, \(\mathcal F\) is
locally free.
\end{solution}

\begin{exerciseintro}
% record: H-III-10-006
\exerciseheading{Exercise III.10.6}
\addcontentsline{toc}{subsection}{Exercise III.10.6}

\begin{context}
Because the cubic in the problem is nodal, we assume
\(\operatorname{char}k\ne2\).  At an ordinary node the completed local
ring is \(k[[u,v]]/(uv)\).  A finite-type morphism inducing an
isomorphism on completed local rings, with identical residue fields,
is etale at the points in question.
In this record the nodal cubic is affine. For its projective closure, replace the two affine normalization lines by projective lines and perform the same crosswise gluing.
\end{context}

\begin{problem}
Let \(Y\) be the plane nodal cubic

\[
y^2=x^2(x+1).
\]

Construct a finite etale covering \(p:X\to Y\) of degree \(2\) such
that \(X\) is the union of two irreducible components, each
isomorphic to the normalization of \(Y\).
\end{problem}
\end{exerciseintro}

\begin{solution}
Put

\[
A=k[x,y]/(y^2-x^2(x+1)).
\]

The normalization \(\nu:N\to Y\) is

\[
N=\operatorname{Spec}k[t],
\qquad
x=t^2-1,\quad y=t(t^2-1).
\tag{1}
\]

Indeed, \(t\) is integral over \(A\), since \(t^2=x+1\), and away
from \(x=0\) it is the rational function \(y/x\).  The inverse image
of the node \(P=(0,0)\) consists of the two points

\[
P_+=(t-1),\qquad P_-=(t+1).
\]

Take two copies \(N_1,N_2\) of \(N\).  Glue \(P_+\in N_1\) to
\(P_-\in N_2\), and glue \(P_-\in N_1\) to \(P_+\in N_2\).  The
resulting affine curve can be written explicitly as

\[
X=\operatorname{Spec}B,
\]

where

\[
B=
\left\{
(f_1,f_2)\in k[t]\times k[t]:
\begin{array}{l}
f_1(1)=f_2(-1),\\
f_1(-1)=f_2(1)
\end{array}
\right\}.
\tag{2}
\]

The two projection maps \(B\to k[t]\) are surjective: prescribed
values at \(1\) and \(-1\) can be attained by a linear polynomial.
Their kernels are the two minimal primes of \(B\).  Thus \(X\) has
exactly two irreducible components, and each component is
isomorphic to \(N\).

The two normalization maps in (1) agree at every pair of points
identified in (2), so they induce a morphism

\[
p:X\longrightarrow Y.
\tag{3}
\]

Algebraically, \(A\to B\) is the diagonal map through
\(A\to k[t]\).  The \(A\)-module \(k[t]\) is finite, and \(B\) is an
\(A\)-submodule of \(k[t]\oplus k[t]\).  Since \(A\) is noetherian,
\(B\) is finite over \(A\); hence \(p\) is finite.  Over
\(Y\setminus\{P\}\), the normalization is an isomorphism, so (3) is
the disjoint union of two isomorphisms there.

It remains to check the two points of \(X\) over \(P\).  First,
\[
\widehat{\mathcal O}_{Y,P}
\cong
k[[x,y]]/(y^2-x^2(1+x))
\cong k[[r,s]]/(rs).
\tag{4}
\]

For the second isomorphism, Hensel's lemma gives a power series
\(h(x)\) with \(h(0)=1\) and \(h(x)^2=1+x\).  One may take

\[
r=y-xh(x),\qquad s=y+xh(x);
\]

the change of variables is invertible because \(2\) is a unit and
\(xh(x)\) is a local parameter.

At either of the two crosswise identifications in \(X\), the
completed local ring is the fiber product of the completed rings of
the two branches:

\[
k[[u]]\times_k k[[v]]
\cong k[[u,v]]/(uv).
\tag{5}
\]

Under \(p\), these two branches map respectively to the two distinct
branches of \(Y\) at \(P\).  The induced homomorphism from (4) to
(5) is therefore an isomorphism.  The residue fields are both \(k\),
so the completed-local-ring criterion shows that \(p\) is etale at
both points over \(P\).  We already saw that it is an isomorphism
locally away from those points.  Hence \(p\) is finite etale.

The generic fiber consists of one point on each of the two
components, so the finite etale cover has degree \(2\), as required.
\end{solution}

\begin{exerciseintro}
% record: H-III-10-007
\exerciseheading{Exercise III.10.7}
\addcontentsline{toc}{subsection}{Exercise III.10.7}

\begin{context}
The classification assertion in part (b), as printed, is incomplete:
there are also members consisting of a smooth conic and its tangent
line.  The solution proves all the valid assertions, gives an
explicit counterexample to the stated dichotomy, and proves the
corrected classification.

We use the standard classification of singular irreducible plane
cubics: a double point with two distinct tangent lines is a node,
whereas a double point with one double tangent line is a cusp.
\end{context}

\begin{problem}
Let \(k\) be an algebraically closed field of characteristic \(2\),
let

\[
S=\mathbb P^2(\mathbb F_2)=\{P_1,\ldots,P_7\}
\subseteq X=\mathbb P_k^2,
\]

and let \(\mathfrak d\) be the linear system of cubic curves through
the seven points.

(a) Show that \(\mathfrak d\) has dimension \(2\), has base locus
exactly \(S\), and determines an inseparable morphism of degree \(2\)

\[
X\setminus S\longrightarrow\mathbb P^2.
\]

(b) Show that every member of \(\mathfrak d\) is singular and that
sending a member to its singular point gives a bijection

\[
\mathfrak d\longleftrightarrow\mathbb P^2.
\]

Prove the following classification, writing \(\mathcal A\) for the
union of the seven \(\mathbb F_2\)-lines and \(P\) for the unique
singular point of a member \(C\):
\begin{enumerate}
\item if \(P\in S\), then \(C\) is the union of the three
  \(\mathbb F_2\)-lines through \(P\);
\item if \(P\in\mathcal A\setminus S\), then \(C\) is the union
  of an \(\mathbb F_2\)-line and a smooth conic tangent to it at \(P\);
\item if \(P\notin\mathcal A\), then \(C\) is an irreducible
  cuspidal cubic with cusp \(P\).
\end{enumerate}

\emph{Correction to the printed statement.} The original dichotomy
lists only three concurrent lines or an irreducible cuspidal cubic.
It omits the line-plus-tangent-conic case above, as the explicit
factorization in the solution demonstrates.
\end{problem}
\end{exerciseintro}

\begin{solution}
Use homogeneous coordinates \([x_0:x_1:x_2]\), and put

\[
\begin{aligned}
F_0&=x_1x_2(x_1+x_2),\\
F_1&=x_2x_0(x_2+x_0),\\
F_2&=x_0x_1(x_0+x_1).
\end{aligned}
\tag{1}
\]

(a) Let us first determine the linear system.
Let \(F\) be a homogeneous cubic vanishing on \(S\).  Vanishing at
the three coordinate points kills the coefficients of
\(x_0^3,x_1^3,x_2^3\).  Vanishing at
\([1:1:0]\), \([1:0:1]\), and \([0:1:1]\) identifies the two
coefficients involving each pair of variables.  Finally, vanishing
at \([1:1:1]\) kills the coefficient of \(x_0x_1x_2\).  Therefore

\[
H^0(\mathbb P^2,\mathcal I_S(3))
=kF_0\oplus kF_1\oplus kF_2.
\tag{2}
\]

Thus the projective linear system has dimension \(2\).

The equations \(F_0=F_1=F_2=0\) say that, for each pair of
coordinates, either one coordinate is zero or the two are equal.
After projective scaling, every common zero consequently has all
coordinates in \(\mathbb F_2\).  Conversely, the forms in (1)
vanish at every such point.  Hence the base locus is exactly \(S\),
and the system defines

\[
\phi:\mathbb P^2\setminus S\longrightarrow\mathbb P^2,
\qquad
[x_0:x_1:x_2]\longmapsto[F_0:F_1:F_2].
\tag{3}
\]

We compute its function-field degree.  On the chart \(x_1\ne0\), set

\[
q=\frac{x_0}{x_1},\qquad r=\frac{x_2}{x_1},
\]

and, on the target chart \(F_2\ne0\), set

\[
A=\frac{F_0}{F_2}
=\frac{r(r+1)}{q(q+1)},
\qquad
B=\frac{F_1}{F_2}
=\frac{r(r+q)}{q+1}.
\tag{4}
\]

The identities

\[
\frac{F_1(F_0+F_1)}{F_2(F_0+F_2)}
=r^2
\tag{5}
\]

and

\[
\frac{F_0+F_1+F_2}{F_2(1+r)}+1=\frac rq
\tag{6}
\]

show that

\[
k(q,r)=k(A,B)(r),
\qquad r^2\in k(A,B).
\tag{7}
\]

In particular, \(\phi\) is dominant and its degree is at most \(2\).
It is not generically separable.  Indeed, if

\[
D=q(q+1),\qquad N=r(r+1),
\]

then differentiation of (4) in characteristic \(2\) gives

\[
dA=\frac{D\,dr+N\,dq}{D^2},
\qquad
dB=\frac{D\,dr+N\,dq}{(q+1)^2}.
\tag{8}
\]

These two differentials are nonzero and proportional, so the
generic differential of \(\phi\) has rank \(1\), not \(2\).  The
extension in (7) is therefore nontrivial and purely inseparable.
Its degree is exactly \(2\).

(b) We begin with the singular-point correspondence.
Every element of \(k\) has a unique square root.  Thus a member of
\(\mathfrak d\) can be written uniquely, up to simultaneous scaling,
as

\[
C_{a,b,c}:
\quad
a^2F_0+b^2F_1+c^2F_2=0
\tag{9}
\]

for a point \([a:b:c]\in\mathbb P^2\).  Its partial derivatives are

\[
\begin{aligned}
\frac{\partial F}{\partial x_0}
  &=(b x_2+c x_1)^2,\\
\frac{\partial F}{\partial x_1}
  &=(a x_2+c x_0)^2,\\
\frac{\partial F}{\partial x_2}
  &=(a x_1+b x_0)^2.
\end{aligned}
\tag{10}
\]

The alternating matrix of the three linear forms inside the squares
has rank \(2\), and its kernel is spanned by \((a,b,c)\).
Consequently the three derivatives have the unique common zero

\[
P_{a,b,c}=[a:b:c].
\tag{11}
\]

Euler's identity gives \(F(P_{a,b,c})=0\), so this point lies on the
curve and is singular.  This proves both that every member is
singular and that

\[
C_{a,b,c}\longmapsto P_{a,b,c}
\]

is a bijection from \(\mathfrak d\) to \(\mathbb P^2\).

We prove these claims.  The group
\(\operatorname{PGL}_3(\mathbb F_2)\) preserves \(S\), the linear
system, and the singular-point correspondence.  It is therefore
enough, for a point on an \(\mathbb F_2\)-line, to consider the line
\(x_0=0\).  Write such a point as \(P=[0:b:c]\).  Its member factors
as

\[
C_{0,b,c}=V(x_0Q),
\qquad
Q=b^2x_2(x_2+x_0)+c^2x_1(x_0+x_1).
\tag{12}
\]

If \([0:b:c]\) is an \(\mathbb F_2\)-point, (12) splits into the
three \(\mathbb F_2\)-lines through that point.  For example, at
\([1:0:0]\) the corresponding equation is

\[
F_0=x_1x_2(x_1+x_2).
\]

Now suppose \(b,c,b+c\ne0\), which is precisely the case in which
\([0:b:c]\) lies on \(x_0=0\) but is not an
\(\mathbb F_2\)-point.  The derivatives of \(Q\) are

\[
Q_{x_0}=b^2x_2+c^2x_1,\qquad
Q_{x_1}=c^2x_0,\qquad
Q_{x_2}=b^2x_0.
\]

Their only common projective zero is
\([0:b^2:c^2]\), and

\[
Q(0,b^2,c^2)=b^2c^2(b+c)^2\ne0.
\]

Thus \(Q\) is a smooth conic.  Its restriction to \(x_0=0\) is

\[
Q|_{x_0=0}=(b x_2+c x_1)^2,
\]

so the line is tangent to the conic at \([0:b:c]\).  This proves the
first two cases.

Conversely, suppose a member is reducible.  It has a line component
\(L\).  If \(L\) were not defined over \(\mathbb F_2\), it could
contain at most one point of \(S\), so the residual conic would
contain at least six points of \(S\).  No nonzero conic does so.  To
see this, use an \(\mathbb F_2\)-projective transformation to make
the omitted point \([1:1:1]\).  A quadratic through the other six
points first has all three square coefficients zero by evaluating at
the coordinate points, and then has all three mixed coefficients
zero by evaluating at
\([1:1:0]\), \([1:0:1]\), and \([0:1:1]\).

Hence \(L\) is an \(\mathbb F_2\)-line.  Every intersection of \(L\)
with the residual conic is singular on the cubic.  Since (11) is the
unique singular point, \(P_{a,b,c}\in L\).  We have therefore proved
that a member is reducible exactly when its singular point lies in
\(\mathcal A\).

It remains to identify the singularity outside \(\mathcal A\).
After permuting coordinates, assume \(a\ne0\) and scale so that
\(a=1\).  On \(x_0=1\), translate the singular point by writing

\[
x_1=b+u,\qquad x_2=c+v.
\]

A direct expansion of (9) gives the local equation

\[
uv(u+v)+c(c+1)u^2+b(b+1)v^2=0.
\tag{13}
\]

Outside \(\mathcal A\), the quadratic part is nonzero, and over the
algebraically closed field \(k\) it is the square

\[
\left(
\sqrt{c(c+1)}\,u+\sqrt{b(b+1)}\,v
\right)^2.
\tag{14}
\]

Thus the singularity has multiplicity \(2\) and a single double
tangent.  The curve is irreducible by the criterion just proved, so
the standard classification of singular plane cubics shows that
(13) is a cusp.  This completes the corrected classification.

For an explicit counterexample to the dichotomy printed in the
exercise, choose \(\lambda\in k\setminus\mathbb F_2\).  The member
singular at \([0:1:\lambda]\) is

\[
x_0\left(
x_2(x_0+x_2)+\lambda^2x_1(x_0+x_1)
\right)=0,
\]

which is a line together with a smooth tangent conic, not either of
the two printed alternatives.
\end{solution}

\begin{exerciseintro}
% record: H-III-10-008
\exerciseheading{Exercise III.10.8}
\addcontentsline{toc}{subsection}{Exercise III.10.8}

\begin{context}
The assertion labelled any characteristic is false for the printed
equation in characteristic \(2\): the alleged conic becomes a double
line.  We prove the statement when
\(\operatorname{char}k\ne2\), explain the failure in characteristic
\(2\), and give a replacement that works in every characteristic.
There is a second qualification: the cones are defined only for
\(t\ne0\), then completed to a pencil by taking the limiting
quadric.  The literal cone from \(P_0\in C\) is the plane \(z=0\),
not that limiting quadric.
\end{context}

\begin{problem}
In affine \(3\)-space with coordinates \(x,y,z\), let \(C\) be the
conic

\[
(x-1)^2+y^2=1
\]

in the \(xy\)-plane, and let \(P_t=(0,0,t)\) be a point of the
\(z\)-axis.  For \(t\ne0\), let \(Y_t\) be the closure in \(\mathbb P^3\) of the
cone over \(C\) with vertex \(P_t\).  Show that, as \(t\) varies, the
surfaces \(Y_t\) extend to a linear system of dimension \(1\), with a
moving singularity at \(P_t\), and that its base locus is \(C\)
together with the \(z\)-axis.
\end{problem}
\end{exerciseintro}

\begin{solution}
Use homogeneous coordinates \([x:y:z:w]\) on \(\mathbb P^3\), and
assume first that \(\operatorname{char}k\ne2\).  The projective
closure of \(C\) lies in \(z=0\) and has equation

\[
Q=x^2+y^2-2xw=0.
\tag{1}
\]

For \(t\ne0\), project a point \([x:y:z:w]\) from
\(P_t=[0:0:t:1]\) to the plane \(z=0\).  The resulting point is

\[
[tx:ty:0:tw-z].
\]

Requiring this point to satisfy (1) gives

\[
tQ+2xz=0.
\tag{2}
\]

Thus the cones belong to the pencil

\[
Y_{[s:t]}=V(tQ+2s xz),
\qquad [s:t]\in\mathbb P^1.
\tag{3}
\]

The two quadrics \(Q\) and \(2xz\) are linearly independent, so (3)
is a linear system of projective dimension \(1\).  The completed affine family is the chart \(s=1\); at \(t=0\), equation (3) gives
its flat limiting member \(V(xz)\).  This is a prescribed limit,
not the literal cone from \(P_0\).  At \(t=0\), the lines joining
\(P_0\in C\) to the other points of \(C\) fill the plane \(z=0\).
The complete pencil also includes the member at \([s:t]=[0:1]\).

For \(s=1\), let

\[
G_t=t(x^2+y^2-2xw)+2xz.
\]

Its partial derivatives are

\[
\begin{aligned}
(G_t)_x&=2tx-2tw+2z,&
(G_t)_y&=2ty,\\
(G_t)_z&=2x,&
(G_t)_w&=-2tx.
\end{aligned}
\]

They all vanish at \(P_t=[0:0:t:1]\), so \(P_t\) is singular on
\(Y_t\).  As \(t\) varies, this singular point moves along the
\(z\)-axis.  For \(t\ne0\), it is the vertex of the cone.  The
limiting member \(Y_0=V(xz)\) is singular along the intersection of
its two plane components, which still contains \(P_0\).

The set-theoretic base locus of the pencil is

\[
V(Q,xz).
\]

On \(z=0\), this is precisely the projective closure of \(C\).  On
\(x=0\), equation \(Q=0\) becomes \(y^2=0\), whose underlying set is
the line \(x=y=0\), namely the projective \(z\)-axis.  Hence

\[
|V(Q,xz)|=C\cup V(x,y).
\tag{4}
\]

Scheme-theoretically, the base scheme has a nonreduced structure
along the axis; (4) is the base locus asserted in the exercise.

In characteristic \(2\), however,

\[
Q=x^2+y^2=(x+y)^2,\qquad 2xz=0.
\]

The curve called \(C\) is a double line, and all the equations (2)
with \(t\ne0\) define the same double plane.  Thus the displayed
family has dimension \(0\), not \(1\), and the statement as printed
cannot hold in characteristic \(2\).

For comparison, the intended phenomenon does have a
characteristic-free version.  Replace \(C\) by the smooth conic

\[
xw-y^2=0,\qquad z=0.
\]

For \(t\ne0\), the cones with vertex \(P_t\) have equations

\[
t(xw-y^2)-xz=0,
\]

and extend at \(t=0\) by the flat limit \(V(xz)\).
These equations span the pencil generated by \(xw-y^2\) and \(xz\) in every
characteristic.  The same derivative and base-locus calculations
give the moving vertex \(P_t\) and the base set consisting of this
conic and the \(z\)-axis.
\end{solution}

\begin{exerciseintro}
% record: H-III-10-009
\exerciseheading{Exercise III.10.9}
\addcontentsline{toc}{subsection}{Exercise III.10.9}

\begin{context}
We use the local criterion for flatness in the following form.  Let
\((A,\mathfrak m)\to(B,\mathfrak n)\) be a local homomorphism of
noetherian local rings.  If \(a\in\mathfrak m\) is a nonzerodivisor
on both \(A\) and \(B\), and \(B/aB\) is flat over \(A/aA\), then
\(B\) is flat over \(A\).

In a Cohen--Macaulay local ring, every system of parameters is a
regular sequence.
\end{context}

\begin{problem}
Let \(f:X\to Y\) be a morphism of varieties over \(k\).  Assume that
\(Y\) is regular, \(X\) is Cohen--Macaulay, and every fiber of \(f\)
has dimension

\[
\dim X-\dim Y.
\]

Prove that \(f\) is flat.
\end{problem}
\end{exerciseintro}

\begin{solution}
Put

\[
d=\dim X-\dim Y.
\]

The fiber-dimension theorem first shows that \(f\) is dominant:
if the closure of its image had dimension smaller than \(\dim Y\),
the generic fiber over that image would have dimension greater than
\(d\).  Moreover, every irreducible component of every fiber has
dimension at least \(d\).  Since every fiber has dimension exactly
\(d\), all fibers are pure of dimension \(d\).

Fix \(x\in X\), put \(y=f(x)\), and set

\[
A=\mathcal O_{Y,y},\qquad B=\mathcal O_{X,x}.
\]

The local ring of the fiber is

\[
B/\mathfrak m_A B=\mathcal O_{X_y,x}.
\tag{1}
\]

Dimension theory for varieties, together with the purity just
proved, gives

\[
\dim B
=\dim A+\dim(B/\mathfrak m_A B).
\tag{2}
\]

For clarity, this can be seen from transcendence degrees:
\[
\begin{aligned}
\dim B-\dim A
&=\bigl(\dim X-\operatorname{trdeg}_k k(x)\bigr)
 -\bigl(\dim Y-\operatorname{trdeg}_k k(y)\bigr)\\
&=d-\operatorname{trdeg}_{k(y)}k(x)
=\dim\mathcal O_{X_y,x}.
\end{aligned}
\]

The ring \(A\) is regular.  Choose a regular system of parameters
\[
a_1,\ldots,a_n,\qquad n=\dim A.
\]

By (2),

\[
\dim B/(a_1,\ldots,a_n)B=\dim B-n.
\tag{3}
\]

Since \(B\) is Cohen--Macaulay, equation (3) says that the images of
\(a_1,\ldots,a_n\) form part of a system of parameters of \(B\).
They are therefore a \(B\)-regular sequence.

We now prove that \(A\to B\) is flat by induction on \(n\).  If
\(n=0\), then \(A\) is a field and there is nothing to prove.  If
\(n>0\), the element \(a_1\) is a nonzerodivisor on both \(A\) and
\(B\).  The quotient \(A/a_1A\) is regular, the quotient
\(B/a_1B\) is Cohen--Macaulay, and the remaining images
\(a_2,\ldots,a_n\) satisfy the same dimension equality.  By
induction,

\[
B/a_1B\quad\text{is flat over}\quad A/a_1A.
\]

The local criterion for flatness from the context then implies that
\(B\) is flat over \(A\).  This holds for every \(x\in X\), so
\(f\) is flat.
\end{solution}

\begin{exerciseintro}
\corpussection{III.11 The Theorem on Formal Functions}
\addcontentsline{toc}{section}{III.11 The Theorem on Formal Functions}

% record: H-III-11-001
\exerciseheading{Exercise III.11.1}
\addcontentsline{toc}{subsection}{Exercise III.11.1}

\begin{context}
For a quasi-compact separated morphism of noetherian schemes, higher
direct images of quasi-coherent sheaves are quasi-coherent.  On an
affine base, the Leray spectral sequence therefore identifies their
global sections with the cohomology of the source.
\end{context}

\begin{problem}
Show that the fiber-dimension vanishing theorem for higher direct
images is false without the projective hypothesis.  More precisely,
let \(n\geq2\),

\[
X=\mathbb A_k^n,\qquad P=(0,\ldots,0),\qquad U=X\setminus\{P\},
\]

and let \(j:U\hookrightarrow X\) be the inclusion.  Show that every
fiber of \(j\) has dimension at most \(0\), but

\[
R^{n-1}j_*\mathcal O_U\ne0.
\]
\end{problem}
\end{exerciseintro}

\begin{solution}
An open immersion has either an empty fiber or a fiber consisting of
one reduced point, so all fibers have dimension at most \(0\).

Write

\[
A=k[x_1,\ldots,x_n].
\]

The principal opens \(D(x_1),\ldots,D(x_n)\) cover \(U\).  Their
intersections are affine, so this cover computes the cohomology of
\(\mathcal O_U\).  The last differential in its Cech complex has
cokernel

\[
H^{n-1}(U,\mathcal O_U)
\cong
\frac{A_{x_1\cdots x_n}}
{\displaystyle\sum_{i=1}^n
 A_{x_1\cdots\widehat{x_i}\cdots x_n}}.
\tag{1}
\]

Every Laurent monomial in a summand of the denominator has
nonnegative exponent in at least one variable.  Hence the class of

\[
\frac{1}{x_1\cdots x_n}
\]

is nonzero in (1).  More precisely, the quotient has \(k\)-basis

\[
\left\{
x_1^{-a_1}\cdots x_n^{-a_n}:a_1,\ldots,a_n\geq1
\right\}.
\tag{2}
\]

Since \(X\) is affine and the higher direct images are
quasi-coherent, Leray gives

\[
\Gamma\left(X,R^{n-1}j_*\mathcal O_U\right)
\cong H^{n-1}(U,\mathcal O_U).
\]

The right side is nonzero by (2), so
\(R^{n-1}j_*\mathcal O_U\ne0\), despite the zero-dimensional fibers.
\end{solution}

\begin{exerciseintro}
% record: H-III-11-002
\exerciseheading{Exercise III.11.2}
\addcontentsline{toc}{subsection}{Exercise III.11.2}

\begin{context}
If a projective morphism has fibers of dimension at most \(r\), then
\(R^if_*\mathcal F=0\) for every coherent \(\mathcal F\) and every
\(i>r\).  Serre's cohomological criterion says that a noetherian
scheme \(Z\) is affine if
\(H^i(Z,\mathcal F)=0\) for every coherent \(\mathcal F\) and every
\(i>0\).  An affine proper morphism is finite.
The schemes in this exercise are noetherian.
\end{context}

\begin{problem}
Prove that a projective morphism with finite fibers, equivalently a
projective quasi-finite morphism, is finite.
\end{problem}
\end{exerciseintro}

\begin{solution}
Let \(f:X\to Y\) be projective and quasi-finite.  It is enough to
prove that \(f^{-1}(V)\) is affine for every affine open
\(V\subseteq Y\).  Replace \(Y\) by such a \(V\), and let
\(\mathcal F\) be any coherent sheaf on \(X\).

Every fiber of \(f\) has dimension \(0\).  The projective
fiber-dimension theorem therefore gives

\[
R^if_*\mathcal F=0\qquad(i>0).
\]

These higher direct images are coherent, and the base is affine, so
the Leray spectral sequence yields

\[
H^i(X,\mathcal F)
=\Gamma(Y,R^if_*\mathcal F)=0
\qquad(i>0).
\]

Serre's criterion now shows that \(X\) is affine.  Thus \(f\) is an
affine morphism.  It is also proper because it is projective, and an
affine proper morphism is finite.  Hence \(f\) is finite.
\end{solution}

\begin{exerciseintro}
% record: H-III-11-003
\exerciseheading{Exercise III.11.3}
\addcontentsline{toc}{subsection}{Exercise III.11.3}

\begin{context}
Stein factorization writes a projective morphism as a morphism with
connected fibers followed by a finite morphism.  The pullback of an
ample line bundle by a finite morphism is ample.  On a normal
projective variety of dimension at least \(2\), the support of every
effective ample Cartier divisor is connected.
\end{context}

\begin{problem}
Let \(X\) be a normal projective variety over an algebraically closed
field \(k\).  Let \(\mathfrak d\) be a base-point-free linear system
of effective Cartier divisors.  Suppose that \(\mathfrak d\) is not
composite with a pencil: if

\[
f:X\longrightarrow\mathbb P_k^n
\]

is the morphism defined by \(\mathfrak d\), then
\(\dim f(X)\geq2\).  Prove that every divisor in
\(\mathfrak d\) is connected.
\end{problem}
\end{exerciseintro}

\begin{solution}
Take the Stein factorization of \(f\):

\[
X\xrightarrow{\,g\,}Z\xrightarrow{\,h\,}f(X)\subseteq\mathbb P^n,
\tag{1}
\]

where \(g\) has connected fibers, \(h\) is finite, and
\(g_*\mathcal O_X=\mathcal O_Z\).

The variety \(Z\) is normal.  Indeed, for an affine open
\(V\subseteq Z\),

\[
\Gamma(V,\mathcal O_Z)
=\Gamma(g^{-1}(V),\mathcal O_X).
\]

The latter ring is integrally closed in \(k(X)\), since
\(g^{-1}(V)\) is an open subset of the normal integral variety
\(X\).  Also, because \(h\) is finite,

\[
\dim Z=\dim f(X)\geq2.
\tag{2}
\]

The finite morphism \(h\) makes \(Z\) projective, and

\[
\mathcal L=h^*\mathcal O_{\mathbb P^n}(1)
\]

is ample on \(Z\).

Let \(D\in\mathfrak d\), and let \(H\subseteq\mathbb P^n\) be the
hyperplane corresponding to \(D\).  Put

\[
E=h^{-1}(H).
\]

The hyperplane does not contain \(f(X)\), so \(E\) is an effective
Cartier divisor associated to the ample line bundle \(\mathcal L\).
By the connectedness theorem and (2), \(E\) is connected.  Moreover,

\[
D=f^{-1}(H)=g^{-1}(E).
\tag{3}
\]

It remains only to pass connectedness through \(g\).  The restriction
\[
g^{-1}(E)\longrightarrow E
\]

is a proper surjection with connected fibers.  If its source were a
disjoint union of two nonempty open-and-closed subsets, their images
would be closed, would cover \(E\), and could not meet because each
fiber is connected.  This would disconnect \(E\), a contradiction.
Thus \(g^{-1}(E)\), and hence \(D\) by (3), is connected.
\end{solution}

\begin{exerciseintro}
% record: H-III-11-004
\exerciseheading{Exercise III.11.4}
\addcontentsline{toc}{subsection}{Exercise III.11.4}

\begin{context}
We use the degree-zero theorem on formal functions.  We also use the
following standard facts.  Idempotents, equivalently open-and-closed
decompositions, lift uniquely across nilpotent thickenings.  A finite
torsion-free module over a discrete valuation ring is free.  Finally,
if the geometric generic fiber of a finite-type morphism over an
integral curve is disconnected, then, after a finite extension of the
function field and normalization of the base, that disconnection
spreads out over a nonempty open subset.
\end{context}

\begin{problem}
Principle of connectedness.
Let \(\{X_t\}\) be a flat family of closed subschemes of
\(\mathbb P_k^n\), parametrized by an irreducible curve \(T\) of
finite type over the algebraically closed field \(k\).  Suppose that
there is a nonempty open subset \(U\subseteq T\) such that \(X_t\)
is connected for every closed point \(t\in U\).  Prove that \(X_t\)
is connected for every \(t\in T\).
\end{problem}
\end{exerciseintro}

\begin{solution}
The hypothesis also makes the generic fiber geometrically connected:
otherwise its disconnection spreads out after a finite extension of
the function field, producing a disconnected closed fiber over the
given open set \(U\).  This is impossible.  It remains to treat
the closed fibers.

Let
\[
p:X\longrightarrow T
\]
be the projective flat morphism representing the family.  First
replace \(T\) by its normalization \(\nu:T'\to T\) and \(X\) by
\(X\times_TT'\).  This does not change the question: \(\nu\) is
finite and surjective, and, because \(k\) is algebraically closed,
the fiber over a closed point \(t'\) above \(t\) is just \(X_t\).
The inverse image of \(U\) is again nonempty.  We may therefore
assume that \(T\) is regular.

Fix a closed point \(t\in T\), and put
\[
A=\mathcal O_{T,t},\qquad \mathfrak m=(\pi),\qquad K=\operatorname{Frac}(A).
\]
Thus \(A\) is a discrete valuation ring.  Set
\[
B=(p_*\mathcal O_X)_t.
\]
Properness makes \(B\) a finite \(A\)-module.  It is torsion-free:
multiplication by any nonzero element of \(A\) is injective on
\(\mathcal O_X\), because \(X\) is flat over \(T\), and remains
injective after applying the left-exact functor \(p_*\).  Hence
\(B\) is a finite free \(A\)-module.

Suppose, toward a contradiction, that \(X_t\) is disconnected.  For
\(r\geq0\), let
\[
X_r=X\times_T\operatorname{Spec}(A/\mathfrak m^{r+1}).
\]
Every \(X_r\) has the same underlying topological space as \(X_t\).
The nontrivial open-and-closed decomposition of \(X_t\) consequently
lifts uniquely to each \(X_r\), and these lifts give compatible
nontrivial idempotents
\[
e_r\in H^0(X_r,\mathcal O_{X_r}).
\]
The theorem on formal functions gives
\[
\widehat B\cong
\varprojlim_r H^0(X_r,\mathcal O_{X_r}),
\tag{1}
\]
so the \(e_r\) determine an idempotent
\(e\in\widehat B\).  Its reduction modulo \(\mathfrak m\) is
neither \(0\) nor \(1\).

Because \(B\) is free over \(A\), both \(\widehat B\) and its direct
summands \(e\widehat B\) and \((1-e)\widehat B\) are torsion-free.
Thus \(e\) remains nontrivial after inverting \(\pi\).  Flat base
change identifies
\[
\widehat B[1/\pi]
\cong H^0(X_{\operatorname{Frac}(\widehat A)},\mathcal O),
\]
and the surviving idempotent shows that
\(X_{\operatorname{Frac}(\widehat A)}\) is disconnected.  Therefore
the generic fiber \(X_K\) is not geometrically connected.

This contradicts geometric connectedness of the generic fiber,
established at the start. Thus \(X_t\) is connected.
Since \(t\) was arbitrary, all fibers are connected.
\end{solution}

\begin{exerciseintro}
% record: H-III-11-005
\exerciseheading{Exercise III.11.5}
\addcontentsline{toc}{subsection}{Exercise III.11.5}

\begin{context}
For a square-zero thickening \(Z\subseteq Z'\) with ideal sheaf
\(\mathcal J\), the sequence
\[
1\longrightarrow 1+\mathcal J\longrightarrow
\mathcal O_{Z'}^*\longrightarrow\mathcal O_Z^*\longrightarrow1,
\qquad 1+\mathcal J\cong\mathcal J,
\]
gives the exact segment
\[
H^1(Z,\mathcal J)\longrightarrow\operatorname{Pic}Z'
\longrightarrow\operatorname{Pic}Z
\longrightarrow H^2(Z,\mathcal J).
\]
We also use \(\operatorname{Pic}\widehat X\cong
\varprojlim_n\operatorname{Pic}Y_n\) for the infinitesimal
neighborhoods of \(Y\).
\end{context}

\begin{problem}
Let \(Y\) be a hypersurface in \(X=\mathbb P_k^N\), where
\(N\geq4\), and let \(\widehat X\) be the formal completion of
\(X\) along \(Y\).  Prove that the natural map
\[
\operatorname{Pic}\widehat X\longrightarrow\operatorname{Pic}Y
\]
is an isomorphism.
\end{problem}
\end{exerciseintro}

\begin{solution}
Let \(Y\) have degree \(d\), let \(\mathcal I=\mathcal O_X(-d)\)
be its ideal sheaf, and write
\[
Y_n=(|Y|,\mathcal O_X/\mathcal I^{n+1})
\qquad(n\geq0).
\]
Thus \(Y_0=Y\).  The ideal of \(Y_n\) in \(Y_{n+1}\) is
square-zero and is naturally
\[
\mathcal I^{n+1}/\mathcal I^{n+2}
\cong\mathcal O_Y(-(n+1)d).
\tag{1}
\]
Consequently the square-zero Picard sequence gives
\[
H^1\bigl(Y,\mathcal O_Y(-(n+1)d)\bigr)
\longrightarrow\operatorname{Pic}Y_{n+1}
\longrightarrow\operatorname{Pic}Y_n
\longrightarrow
H^2\bigl(Y,\mathcal O_Y(-(n+1)d)\bigr).
\tag{2}
\]

For every integer \(m\), tensoring the hypersurface sequence by
\(\mathcal O_X(m)\) gives
\[
0\longrightarrow\mathcal O_X(m-d)
\longrightarrow\mathcal O_X(m)
\longrightarrow\mathcal O_Y(m)
\longrightarrow0.
\tag{3}
\]
On \(\mathbb P^N\),
\[
H^q(\mathbb P^N,\mathcal O(m))=0
\qquad(0<q<N)
\]
for every \(m\).  Since \(N\geq4\), the long exact sequence of
(3) therefore yields
\[
H^1(Y,\mathcal O_Y(m))=H^2(Y,\mathcal O_Y(m))=0
\qquad\text{for every }m\in\mathbb Z.
\tag{4}
\]
By (2) and (4), every transition map
\[
\operatorname{Pic}Y_{n+1}\xrightarrow{\sim}
\operatorname{Pic}Y_n
\]
is an isomorphism.  Taking inverse limits gives
\[
\operatorname{Pic}\widehat X
\cong\varprojlim_n\operatorname{Pic}Y_n
\cong\operatorname{Pic}Y_0
=\operatorname{Pic}Y,
\]
as required.
\end{solution}

\begin{exerciseintro}
% record: H-III-11-006
\exerciseheading{Exercise III.11.6}
\addcontentsline{toc}{subsection}{Exercise III.11.6}

\begin{context}
Let \(Y\subseteq X=\mathbb P_k^N\) be a hypersurface of degree
\(d\), with ideal \(\mathcal I=\mathcal O_X(-d)\), and put
\(Y_n=(|Y|,\mathcal O_X/\mathcal I^{n+1})\).  Sections on the formal
completion are inverse limits of sections on the \(Y_n\).  Completion
is exact on coherent sheaves in this noetherian setting.

We use Serre vanishing and Serre duality on projective space.  We
also use that completion of a noetherian local ring along an ideal
contained in its maximal ideal is faithfully flat.
\end{context}

\begin{problem}
Again let \(Y\) be a hypersurface in \(X=\mathbb P_k^N\), this time
with \(N\geq2\).
(a) If \(\mathcal F\) is locally free on \(X\), prove that
the natural map
\[
H^0(X,\mathcal F)\longrightarrow
H^0(\widehat X,\widehat{\mathcal F})
\]
is an isomorphism.

(b) Prove that the following conditions are equivalent:
\begin{enumerate}
\item[(i)] every locally free sheaf \(\mathfrak F\) on
\(\widehat X\) is algebraizable, meaning that
\(\mathfrak F\cong\widehat{\mathcal F}\) for some coherent sheaf
\(\mathcal F\) on \(X\);

\item[(ii)] for every locally free sheaf \(\mathfrak F\) on
\(\widehat X\), there is an integer \(n_0\) such that
\(\mathfrak F(n)\) is generated by global sections for all
\(n\geq n_0\).
\end{enumerate}

(c) Show that these equivalent conditions imply that
\[
\operatorname{Pic}X\longrightarrow\operatorname{Pic}\widehat X
\]
is an isomorphism.
\end{problem}
\end{exerciseintro}

\begin{solution}
(a) Tensoring
\[
0\longrightarrow\mathcal O_X(-(n+1)d)
\longrightarrow\mathcal O_X
\longrightarrow\mathcal O_{Y_n}
\longrightarrow0
\]
with the locally free sheaf \(\mathcal F\) gives
\[
0\longrightarrow\mathcal F(-(n+1)d)
\longrightarrow\mathcal F
\longrightarrow\mathcal F|_{Y_n}
\longrightarrow0.
\tag{1}
\]
For \(m\gg0\), one has
\[
H^0(X,\mathcal F(-m))=H^1(X,\mathcal F(-m))=0.
\tag{2}
\]
The first vanishing is the standard negative-twist vanishing for a
vector bundle.  For the second, Serre duality gives
\[
H^1(X,\mathcal F(-m))^\vee
\cong
H^{N-1}\bigl(X,\mathcal F^\vee(m-N-1)\bigr),
\]
which vanishes for \(m\gg0\) by Serre vanishing; here \(N\geq2\).
It follows from (1) and (2) that, for all sufficiently large \(n\),
restriction is an isomorphism
\[
H^0(X,\mathcal F)\xrightarrow{\sim}
H^0(Y_n,\mathcal F|_{Y_n}).
\]
The inverse system is therefore eventually constant, and hence
\[
H^0(\widehat X,\widehat{\mathcal F})
=\varprojlim_nH^0(Y_n,\mathcal F|_{Y_n})
\cong H^0(X,\mathcal F).
\]

(b) Assume (i), and write
\(\mathfrak F\cong\widehat{\mathcal F}\) with \(\mathcal F\)
coherent on \(X\).  By Serre's theorem, \(\mathcal F(n)\) is
generated by global sections for all \(n\gg0\).  Completing a
surjection from a finite free sheaf shows that
\(\mathfrak F(n)\) is generated by global sections for the same
values of \(n\).  Thus (ii) holds.

Conversely, assume (ii), and let \(\mathfrak F\) be locally free on
\(\widehat X\).  Choose \(n\) so that \(\mathfrak F(n)\) is
globally generated.  Finitely many sections give a surjection
\[
\widehat{\mathcal E}_1=
\mathcal O_{\widehat X}(-n)^{\oplus r}
\twoheadrightarrow\mathfrak F.
\tag{3}
\]
Its kernel \(\mathfrak K\) is locally free: locally the target of
(3) is free, and a surjection onto a finite free module splits.
Apply (ii) to \(\mathfrak K\).  For a suitable \(m\), finitely many
global sections give a surjection
\[
\widehat{\mathcal E}_2=
\mathcal O_{\widehat X}(-m)^{\oplus s}
\twoheadrightarrow\mathfrak K.
\]
After composing with the inclusion in (3), we obtain a presentation
\[
\widehat{\mathcal E}_2
\xrightarrow{\widehat\varphi}
\widehat{\mathcal E}_1
\longrightarrow\mathfrak F\longrightarrow0,
\tag{4}
\]
where
\[
\mathcal E_1=\mathcal O_X(-n)^{\oplus r},
\qquad
\mathcal E_2=\mathcal O_X(-m)^{\oplus s}.
\]

The bundle \(\mathcal H om(\mathcal E_2,\mathcal E_1)\) is locally
free, so part (a) identifies
\[
\begin{aligned}
\operatorname{Hom}_X(\mathcal E_2,\mathcal E_1)
&=H^0\bigl(X,\mathcal H om(\mathcal E_2,\mathcal E_1)\bigr)\\
&\xrightarrow{\sim}
H^0\bigl(\widehat X,
\mathcal H om(\widehat{\mathcal E}_2,
\widehat{\mathcal E}_1)\bigr).
\end{aligned}
\tag{5}
\]
Thus \(\widehat\varphi\) lifts uniquely to a morphism
\(\varphi:\mathcal E_2\to\mathcal E_1\).  Put
\[
\mathcal F=\operatorname{coker}\varphi.
\]
Exactness of completion and (4) give
\[
\widehat{\mathcal F}\cong\mathfrak F.
\]
Hence \(\mathfrak F\) is algebraizable, proving (i).

(c) First consider injectivity.  Since
\(\operatorname{Pic}(\mathbb P^N)\cong\mathbb Z\), every line
bundle on \(X\) is \(\mathcal O_X(a)\).  If its completion is
trivial, then \(\mathcal O_Y(a)\) is trivial.  For \(a\ne0\), this
would make a positive power of the ample line bundle
\(\mathcal O_Y(1)\) trivial.  That is impossible on the
positive-dimensional projective scheme \(Y\).  Thus \(a=0\), and
the map on Picard groups is injective.

For surjectivity, let \(\mathfrak L\) be an invertible sheaf on
\(\widehat X\).  By (i), there is a coherent sheaf \(\mathcal F\)
on \(X\) with
\[
\widehat{\mathcal F}\cong\mathfrak L.
\tag{6}
\]
At every \(y\in Y\), the completion of the finite
\(\mathcal O_{X,y}\)-module \(\mathcal F_y\) is free of rank one.
Faithful flatness of completion therefore implies that
\(\mathcal F_y\) itself is free of rank one.  The locally free locus
is open, so \(\mathcal F\) is an invertible sheaf on some open
neighborhood \(V\) of \(Y\).

The complement \(X\setminus V\) is finite.  Indeed, a
positive-dimensional closed projective subvariety of \(X\) must
meet the ample hypersurface \(Y\); otherwise the defining section of
\(\mathcal O_X(d)\) would trivialize an ample line bundle on that
subvariety.  Hence \(X\setminus V\) has dimension zero.  Since
\(X=\mathbb P^N\) is regular and \(N\geq2\), restriction across a
closed subset of codimension at least two induces an isomorphism
\[
\operatorname{Pic}X\xrightarrow{\sim}\operatorname{Pic}V.
\]
Thus \(\mathcal F|_V\) extends to a line bundle \(\mathcal L\) on
\(X\).  Completion only sees a neighborhood of \(Y\), so (6) gives
\(\widehat{\mathcal L}\cong\mathfrak L\).  This proves
surjectivity and completes the proof.
\end{solution}

\begin{exerciseintro}
% record: H-III-11-007
\exerciseheading{Exercise III.11.7}
\addcontentsline{toc}{subsection}{Exercise III.11.7}

\begin{context}
There is a characteristic issue in part (a).  The kernel is naturally
an infinite-dimensional \(k\)-vector space when
\(\operatorname{char}k=0\), by the formal logarithm.  In positive
characteristic it is generally not a vector space under tensor
product.  The characteristic-free statement is that it has a
complete filtration whose associated graded pieces are finite
dimensional \(k\)-vector spaces, with infinitely many nonzero pieces.
We prove this precise statement, the stated characteristic-zero
form, and the conclusions of (b) and (c) in every characteristic.
\end{context}

\begin{problem}
Let \(Y\) be a curve in \(X=\mathbb P_k^2\), and let
\(\widehat X\) be the formal completion of \(X\) along \(Y\).
(a) Use the method of Exercise III.11.5 to show that
\[
\operatorname{Pic}\widehat X\longrightarrow\operatorname{Pic}Y
\]
is surjective and that its kernel is an infinite-dimensional vector
space over \(k\).

(b) Conclude that there is an invertible sheaf on
\(\widehat X\) which is not algebraizable.

(c) Conclude that there is a locally free sheaf
\(\mathfrak F\) on \(\widehat X\) such that no twist
\(\mathfrak F(n)\) is generated by global sections.
\end{problem}
\end{exerciseintro}

\begin{solution}
Let \(Y\) have degree \(d\), let
\(\mathcal I=\mathcal O_X(-d)\), and set
\[
Y_n=(|Y|,\mathcal O_X/\mathcal I^{n+1})
\qquad(n\geq0).
\]
Thus \(Y_0=Y\) and
\(\operatorname{Pic}\widehat X=\varprojlim_n\operatorname{Pic}Y_n\).

(a) The ideal of \(Y_n\) in \(Y_{n+1}\) is square-zero and equals
\[
\mathcal J_n=\mathcal I^{n+1}/\mathcal I^{n+2}
\cong\mathcal O_Y(-(n+1)d).
\tag{1}
\]
The exact sequence of units for this thickening gives
\[
H^1(Y,\mathcal J_n)\longrightarrow
\operatorname{Pic}Y_{n+1}\longrightarrow
\operatorname{Pic}Y_n\longrightarrow H^2(Y,\mathcal J_n).
\tag{2}
\]
The last group vanishes because \(Y\) is a curve.

We can also identify the kernel in (2).  From
\[
0\longrightarrow\mathcal O_X(-(n+1)d)
\longrightarrow\mathcal O_X
\longrightarrow\mathcal O_{Y_n}
\longrightarrow0
\]
and the vanishing of \(H^0(\mathbb P^2,\mathcal O(-a))\) for
\(a>0\) and of \(H^1(\mathbb P^2,\mathcal O(m))\) for every \(m\),
we obtain
\[
H^0(Y_n,\mathcal O_{Y_n})=k.
\]
Thus the restriction map on global units is \(k^*\to k^*\), hence
is an isomorphism.  The full units sequence therefore turns (2) into
the short exact sequence
\[
0\longrightarrow V_n\longrightarrow
\operatorname{Pic}Y_{n+1}\longrightarrow
\operatorname{Pic}Y_n\longrightarrow0,
\qquad
V_n=H^1\bigl(Y,\mathcal O_Y(-(n+1)d)\bigr).
\tag{3}
\]
In particular all transition maps are surjective, so every line
bundle on \(Y\) lifts successively to all \(Y_n\).  Hence
\[
\operatorname{Pic}\widehat X\longrightarrow\operatorname{Pic}Y
\]
is surjective.

The vector spaces \(V_n\) are nonzero for all sufficiently large
\(n\), and their dimensions tend to infinity.  Indeed, the
hypersurface sequence gives
\[
0\longrightarrow\mathcal O_{\mathbb P^2}(-(n+2)d)
\longrightarrow\mathcal O_{\mathbb P^2}(-(n+1)d)
\longrightarrow\mathcal O_Y(-(n+1)d)
\longrightarrow0.
\]
Using \(H^1(\mathbb P^2,\mathcal O(m))=0\) and
\(H^2(Y,-)=0\), one obtains, for \(n\gg0\),
\[
\dim_kV_n
=\binom{(n+2)d-1}{2}-\binom{(n+1)d-1}{2},
\tag{4}
\]
which is positive and grows linearly with \(n\).

Let
\[
K=\ker\bigl(\operatorname{Pic}\widehat X
\longrightarrow\operatorname{Pic}Y\bigr)
\]
and, for \(r\geq1\), define
\[
F^rK=\ker\bigl(\operatorname{Pic}\widehat X
\longrightarrow\operatorname{Pic}Y_{r-1}\bigr).
\]
Surjectivity in (3), including the ability to continue every
finite-level lift, gives
\[
F^rK/F^{r+1}K
\cong H^1(Y,\mathcal O_Y(-rd)).
\tag{5}
\]
Thus, in every characteristic, \(K\) is a complete filtered abelian
group with infinitely many nonzero \(k\)-vector-space graded pieces.

If \(\operatorname{char}k=0\), the convergent formal power series
\[
\log(1+x)=x-\frac{x^2}{2}+\frac{x^3}{3}-\cdots
\]
identifies the sheaf \(1+\widehat{\mathcal I}\) with the additive
sheaf \(\widehat{\mathcal I}\).  Since part (a) of Exercise III.11.6
gives
\(H^0(\widehat X,\mathcal O_{\widehat X})=k\), the global units on
\(\widehat X\) and on \(Y\) are both \(k^*\).  The units sequence
therefore identifies
\[
K\cong H^1(\widehat X,\widehat{\mathcal I})
\]
as a \(k\)-vector space.  Its filtration has the graded pieces (5),
so (4) proves that it is infinite-dimensional.  This proves the
assertion in the problem under its necessary characteristic-zero
hypothesis.

Here is why that qualification cannot simply be omitted.  Suppose
\(\operatorname{char}k=p>0\) and take \(Y\) to be a line.  The
class in
\[
H^1(\mathbb P^1,\mathcal O_{\mathbb P^1}(-2))
\]
represented on the standard overlap by \(t^{-1}\) lifts, using
(3), to an element \(L\in F^2K\).  Its leading class is nonzero.
Taking the \(p\)-th tensor power sends a transition function
\(1+a\) to \(1+a^p\).  Hence the leading class of \(L^{\otimes p}\)
is represented by \(t^{-p}\) in
\[
H^1(\mathbb P^1,\mathcal O_{\mathbb P^1}(-2p)),
\]
where it is still nonzero.  Thus \(L^{\otimes p}\ne\mathcal O\).
But the additive group of a \(k\)-vector space in characteristic
\(p\) is killed by \(p\).  Consequently the literal vector-space
claim is false in positive characteristic, while (3)--(5) give its
correct replacement.

(b) Choose a nontrivial \(\mathfrak L\in K\), which exists by (4) and
(5).  Suppose that \(\mathfrak L\) were algebraizable by a coherent
sheaf \(\mathcal F\) on \(X\).  The argument in Exercise III.11.6(c)
does not require algebraizability of every formal bundle: it shows
that \(\mathcal F\) is a line bundle near \(Y\), and then extends it
across the finite complement to a line bundle \(\mathcal M\) on
\(\mathbb P^2\).  Thus
\[
\mathfrak L\cong\widehat{\mathcal M}.
\]
Now \(\mathcal M\cong\mathcal O_{\mathbb P^2}(a)\) for some
\(a\in\mathbb Z\).  Since \(\mathfrak L|_Y\cong\mathcal O_Y\), we
have \(\mathcal O_Y(a)\cong\mathcal O_Y\).  Ampleness of
\(\mathcal O_Y(1)\) and \(\dim Y=1\) force \(a=0\).  This would
make \(\mathfrak L\) trivial, a contradiction.  Hence
\(\mathfrak L\) is not algebraizable.

(c) Suppose, to the contrary, that every formal locally free sheaf
has at least one globally generated twist.
Since \(\mathcal O_{\widehat X}(1)\) is globally generated, each
such sheaf would then have all sufficiently large twists globally
generated.  Condition (ii) of Exercise III.11.6(b) would hold for
every formal vector bundle, including the kernels needed in that
exercise's construction.  Its equivalence would make every such
bundle algebraizable, contradicting the line bundle in (b).
Therefore some formal locally free sheaf has no globally generated
twist.  This argument does not identify that sheaf with the
particular line bundle chosen in (b).
\end{solution}

\begin{exerciseintro}
% record: H-III-11-008
\exerciseheading{Exercise III.11.8}
\addcontentsline{toc}{subsection}{Exercise III.11.8}

\begin{context}
For a projective morphism of noetherian schemes, higher direct images
of coherent sheaves are coherent.  We use the theorem on formal
functions and Nakayama's lemma.
\end{context}

\begin{problem}
Let \(f:X\to Y\) be a projective morphism, let \(\mathcal F\) be a
coherent sheaf on \(X\) which is flat over \(Y\), and suppose that
\[
H^i(X_y,\mathcal F_y)=0
\]
for some \(i\) and some \(y\in Y\).  Prove that
\(R^if_*\mathcal F\) is zero in a neighborhood of \(y\).
\end{problem}
\end{exerciseintro}

\begin{solution}
Put
\[
A=\mathcal O_{Y,y},\qquad \mathfrak m=\mathfrak m_y,
\]
and, for \(n\geq0\), set
\[
A_n=A/\mathfrak m^{n+1},\qquad
X_n=X\times_Y\operatorname{Spec}A_n,\qquad
\mathcal F_n=\mathcal F|_{X_n}.
\]
Thus \(X_0=X_y\) and \(\mathcal F_0=\mathcal F_y\).

We first prove by induction that
\[
H^i(X_n,\mathcal F_n)=0
\qquad\text{for every }n\geq0.
\tag{1}
\]
The case \(n=0\) is the hypothesis.  For \(n\geq1\), the exact
sequence
\[
0\longrightarrow\mathfrak m^n/\mathfrak m^{n+1}
\longrightarrow A_n\longrightarrow A_{n-1}\longrightarrow0
\]
remains exact after tensoring with \(\mathcal F\), because
\(\mathcal F\) is flat over \(Y\).  It gives
\[
0\longrightarrow
\mathcal F_0\otimes_{k(y)}
(\mathfrak m^n/\mathfrak m^{n+1})
\longrightarrow\mathcal F_n
\longrightarrow\mathcal F_{n-1}
\longrightarrow0.
\tag{2}
\]
The left term in (2) is a finite direct sum of copies of
\(\mathcal F_0\), so its \(i\)-th cohomology vanishes.  If (1) holds
for \(n-1\), the degree-\(i\) part of the long exact cohomology
sequence of (2) reads
\[
0\longrightarrow H^i(X_n,\mathcal F_n)\longrightarrow0.
\]
This proves (1).

The theorem on formal functions now yields
\[
\widehat{(R^if_*\mathcal F)_y}
\cong\varprojlim_nH^i(X_n,\mathcal F_n)=0.
\tag{3}
\]
The stalk \((R^if_*\mathcal F)_y\) is a finite \(A\)-module.
Reducing its completion modulo \(\mathfrak m\), or equivalently using
faithful flatness of completion, (3) implies
\[
(R^if_*\mathcal F)_y=0.
\]
Finally, \(R^if_*\mathcal F\) is coherent, so the complement of its
support is open.  It therefore vanishes on an open neighborhood of
\(y\), as required.
\end{solution}

\begin{exerciseintro}
\corpussection{III.12 The Semicontinuity Theorem}
\addcontentsline{toc}{section}{III.12 The Semicontinuity Theorem}

% record: H-III-12-001
\exerciseheading{Exercise III.12.1}
\addcontentsline{toc}{subsection}{Exercise III.12.1}

\begin{context}
If \(\mathcal F\) is a coherent sheaf on a noetherian scheme, then
\[
y\longmapsto
\dim_{k(y)}\bigl(\mathcal F\otimes k(y)\bigr)
\]
is upper semicontinuous.
\end{context}

\begin{problem}
Let \(Y\) be a scheme of finite type over an algebraically closed
field \(k\).  Prove that the function
\[
\varphi(y)=\dim_k(\mathfrak m_y/\mathfrak m_y^2)
\]
is upper semicontinuous on the set of closed points of \(Y\).
\end{problem}
\end{exerciseintro}

\begin{solution}
The sheaf of differentials \(\Omega_{Y/k}\) is coherent.  If \(y\)
is a closed point, then \(k(y)=k\), and localization gives
\[
(\Omega_{Y/k})_y\cong\Omega_{\mathcal O_{Y,y}/k}.
\]
The residue homomorphism
\(\mathcal O_{Y,y}\to k\) is split by the structure map
\(k\to\mathcal O_{Y,y}\).  The cotangent-space sequence therefore
identifies
\[
\mathfrak m_y/\mathfrak m_y^2
\xrightarrow{\sim}
\Omega_{\mathcal O_{Y,y}/k}\otimes_{\mathcal O_{Y,y}}k
\cong \Omega_{Y/k}\otimes k(y).
\tag{1}
\]
Equivalently, both sides of the first isomorphism represent
\(k\)-derivations from \(\mathcal O_{Y,y}\) to a \(k\)-vector space.

It follows from (1) that, on closed points,
\[
\varphi(y)
=\dim_{k(y)}\bigl(\Omega_{Y/k}\otimes k(y)\bigr).
\]
This is the fiber-dimension function of the coherent sheaf
\(\Omega_{Y/k}\), so it is upper semicontinuous.
\end{solution}

\begin{exerciseintro}
% record: H-III-12-002
\exerciseheading{Exercise III.12.2}
\addcontentsline{toc}{subsection}{Exercise III.12.2}

\begin{context}
For every field \(K\), the dimensions of
\(H^q(\mathbb P_K^n,\mathcal O_{\mathbb P^n}(m))\) depend only on
\(n,m,q\), not on \(K\).  Also
\[
H^q(\mathbb P_K^n,\mathcal O_{\mathbb P^n})=0
\qquad(q>0).
\]
Assume \(n\geq1\).
\end{context}

\begin{problem}
Let \(\{X_t\}\) be a family of hypersurfaces of the same degree in
\(\mathbb P_k^n\).  Prove that, for every \(i\), the function
\[
t\longmapsto h^i(X_t,\mathcal O_{X_t})
\]
is constant.
\end{problem}
\end{exerciseintro}

\begin{solution}
Let the common degree be \(d>0\).  Over the residue field \(k(t)\),
the fiber \(X_t\) is cut out by a nonzero homogeneous polynomial
\(f_t\) of degree \(d\), so there is an exact sequence
\[
0\longrightarrow
\mathcal O_{\mathbb P_{k(t)}^n}(-d)
\xrightarrow{\,f_t\,}
\mathcal O_{\mathbb P_{k(t)}^n}
\longrightarrow\mathcal O_{X_t}
\longrightarrow0.
\tag{1}
\]
Since \(H^0(\mathbb P^n,\mathcal O(-d))=0\) and the positive
cohomology of \(\mathcal O_{\mathbb P^n}\) vanishes, the long exact
sequence of (1) gives
\[
0\longrightarrow k(t)\longrightarrow
H^0(X_t,\mathcal O_{X_t})
\longrightarrow
H^1(\mathbb P_{k(t)}^n,\mathcal O(-d))
\longrightarrow0
\tag{2}
\]
and, for every \(i\geq1\),
\[
H^i(X_t,\mathcal O_{X_t})
\cong
H^{i+1}(\mathbb P_{k(t)}^n,\mathcal O(-d)).
\tag{3}
\]
The dimensions of the groups on the right sides of (2) and (3)
depend only on \(n,d,i\).  Therefore all the numbers
\(h^i(X_t,\mathcal O_{X_t})\) are independent of \(t\).
\end{solution}

\begin{exerciseintro}
% record: H-III-12-003
\exerciseheading{Exercise III.12.3}
\addcontentsline{toc}{subsection}{Exercise III.12.3}

\begin{context}
The two numerical assertions in the problem are false for the
untwisted ideal sheaf: its zeroth and first cohomology vanish on every
fiber.  The asserted values are exactly correct for
\(\mathcal I(1)\).  We construct the requested family, verify the
failure as printed, and prove this standard corrected form.
The ground field is algebraically closed.  We use that all rational
normal quartics are projectively equivalent and that elementary
matrices generate \(\operatorname{SL}_5(k)\).
\end{context}

\begin{problem}
Let \(X_1\subseteq\mathbb P_k^4\) be the rational normal quartic,
the \(4\)-uple embedding of \(\mathbb P^1\).  Let
\(X_0\subseteq\mathbb P_k^3\subseteq\mathbb P_k^4\) be a nonsingular
rational quartic.  Construct a flat family \(\{X_a\}\) in
\(\mathbb P^4\), parametrized by \(T=\mathbb A^1\), with the given
fibers at \(a=1\) and \(a=0\).

Let \(\mathcal I\subseteq\mathcal O_{\mathbb P^4\times T}\) be the
ideal sheaf of the total family.  Show that \(\mathcal I\) is flat
over \(T\), and show that
\[
h^0(a,\mathcal I)=
\begin{cases}
0,&a\ne0,\\
1,&a=0,
\end{cases}
\qquad
h^1(a,\mathcal I)=
\begin{cases}
0,&a\ne0,\\
1,&a=0.
\end{cases}
\]
\end{problem}
\end{exerciseintro}

\begin{solution}
Work over the algebraically closed ground field \(k\).
Choose an isomorphism \(X_0\cong\mathbb P^1\), and choose coordinates
so its containing hyperplane is \(x_4=0\).
The prescribed embedding is defined by a basis
\(s_0,s_1,s_2,s_3\) of a four-dimensional subspace
\(V\subset H^0(\mathbb P^1,\mathcal O(4))\).
Extend it to a basis by a section \(s_4\), and set
\[
\Phi(p,a)=([s_0(p):s_1(p):s_2(p):s_3(p):a s_4(p)],a).
\tag{1}
\]
Projection away from \([0:0:0:0:1]\) is defined on the image.
There \(\Phi\) is the graph over the fixed closed embedding
defined by \(V\), hence is an immersion.  Since its source is proper
over \(\mathbb A^1\), it is a closed immersion into all of
\(\mathbb P^4\times\mathbb A^1\).  Its image is therefore flat,
isomorphic to \(\mathbb P^1\times\mathbb A^1\).
It has the prescribed fiber \(X_0\), and every nonzero fiber is a
rational normal quartic.

If the embedded endpoint \(X_1\) is prescribed exactly, choose
\(g\in\operatorname{PGL}_5(k)\) taking the fiber of (1) at \(1\)
to \(X_1\).  Scale a matrix representative to determinant one.
Write it as a product of elementary matrices
\(\prod_j(I+c_jE_{i_jk_j})\), where \(i_j\ne k_j\).
Then
\[
g(a)=\prod_j(I+a c_jE_{i_jk_j})
\]
is invertible over \(k[a]\), with \(g(0)=I\) and \(g(1)=g\).
Applying this automorphism of \(\mathbb P^4\times\mathbb A^1\)
to the family fixes \(X_0\) and gives precisely \(X_1\) at \(1\).
It preserves flatness and all the fiber cohomology computations.

In the exact sequence
\[
0\longrightarrow\mathcal I\longrightarrow
\mathcal O_{\mathbb P^4\times T}\longrightarrow
\mathcal O_X\longrightarrow0,
\tag{2}
\]
the middle and right terms are flat over \(T\).  The kernel
\(\mathcal I\) is therefore flat over \(T\), and restricting (2) to
a fiber remains exact.

We now encounter the error in the printed assertion.  On every fiber,
\[
0\longrightarrow\mathcal I_a\longrightarrow
\mathcal O_{\mathbb P^4}\longrightarrow
\mathcal O_{X_a}\longrightarrow0.
\tag{3}
\]
Both \(X_a\cong\mathbb P^1\) and \(\mathbb P^4\) have only the
constant global functions, and the map between those constants is an
isomorphism.  Since
\(H^1(\mathbb P^4,\mathcal O_{\mathbb P^4})=0\), sequence (3) gives
\[
H^0(\mathbb P^4,\mathcal I_a)
=H^1(\mathbb P^4,\mathcal I_a)=0
\qquad\text{for every }a.
\tag{4}
\]
Thus the requested jump does not occur for \(\mathcal I\).

For the corrected statement, twist (3) by
\(\mathcal O_{\mathbb P^4}(1)\).  Since
\(\mathcal O_{\mathbb P^4}(1)|_{X_a}\cong
\mathcal O_{\mathbb P^1}(4)\), cohomology gives
\[
0\longrightarrow H^0(\mathcal I_a(1))
\longrightarrow k^5
\xrightarrow{\rho_a}H^0(\mathbb P^1,\mathcal O(4))
\longrightarrow H^1(\mathcal I_a(1))
\longrightarrow0.
\tag{5}
\]
Before applying the harmless fiberwise coordinate change \(g(a)\),
the map \(\rho_a\) sends the coordinate sections to
\(s_0,s_1,s_2,s_3,a s_4\).
It has rank five for \(a\ne0\), and rank four for \(a=0\);
at zero its kernel is \(kx_4\) and its cokernel is represented by
\(s_4\).  Thus
\[
h^0(a,\mathcal I(1))=h^1(a,\mathcal I(1))
=
\begin{cases}
0,&a\ne0,\\
1,&a=0.
\end{cases}
\]
These are the displayed values intended by the exercise.
\end{solution}

\begin{exerciseintro}
% record: H-III-12-004
\exerciseheading{Exercise III.12.4}
\addcontentsline{toc}{subsection}{Exercise III.12.4}

\begin{context}
We use the following consequence of cohomology and base change.  If
\(f:X\to Y\) is projective, \(\mathcal F\) is coherent and flat over
\(Y\), with \(Y\) reduced and noetherian, and
\[
\dim_{k(y)}H^0(X_y,\mathcal F_y)
\]
is constant, then \(f_*\mathcal F\) is locally free and its formation
commutes with passage to the fibers.
\end{context}

\begin{problem}
Let \(Y\) be an integral scheme of finite type over an algebraically
closed field \(k\).  Let \(f:X\to Y\) be a flat projective morphism
whose fibers are integral schemes.  Let \(\mathcal L,\mathcal M\)
be invertible sheaves on \(X\), and suppose that
\[
\mathcal L_y\cong\mathcal M_y
\]
on \(X_y\) for every \(y\in Y\).  Prove that there is an invertible
sheaf \(\mathcal N\) on \(Y\) such that
\[
\mathcal L\cong\mathcal M\otimes f^*\mathcal N.
\]
\end{problem}
\end{exerciseintro}

\begin{solution}
Set
\[
\mathcal P=\mathcal L\otimes\mathcal M^{-1}.
\]
Then \(\mathcal P_y\cong\mathcal O_{X_y}\) for every \(y\).  In
particular, for every closed point \(y\), the fiber \(X_y\) is an
integral projective scheme over the algebraically closed field \(k\),
and hence
\[
H^0(X_y,\mathcal P_y)
\cong H^0(X_y,\mathcal O_{X_y})=k.
\tag{1}
\]
For an arbitrary point \(y\), the triviality of \(\mathcal P_y\)
gives at least one section.  If its space of sections had dimension
at least \(2\), upper semicontinuity would put \(y\) in the closed
locus where \(h^0\geq2\).  Any nonempty such locus contains a closed
point, contradicting (1).  Thus \(h^0(X_y,\mathcal P_y)=1\) for
every \(y\in Y\).

The sheaf \(\mathcal P\) is flat over \(Y\), because it is locally
free on the flat \(Y\)-scheme \(X\).  Cohomology and base change,
applied to this constant value, therefore shows that
\[
\mathcal N=f_*\mathcal P
\]
is invertible on \(Y\), and that
\[
\mathcal N\otimes k(y)
\xrightarrow{\sim}H^0(X_y,\mathcal P_y)
\tag{2}
\]
for every closed point \(y\).

Consider the adjunction evaluation map
\[
\varepsilon:f^*\mathcal N\longrightarrow\mathcal P.
\tag{3}
\]
After restricting (3) to \(X_y\) and using (2), it becomes
\[
H^0(X_y,\mathcal O_{X_y})\otimes_k\mathcal O_{X_y}
\longrightarrow\mathcal O_{X_y},
\]
the evaluation of constant functions.  This is an isomorphism.
Consequently (3) is an isomorphism at every point lying over a
closed point of \(Y\).  If its cokernel were nonzero, its support
would contain a closed point of \(X\), whose image in \(Y\) would
also be closed, a contradiction.  Thus (3) is surjective; being a
surjection between invertible sheaves, it is an isomorphism.

We have proved
\[
\mathcal L\otimes\mathcal M^{-1}
=\mathcal P\cong f^*\mathcal N,
\]
or equivalently
\[
\mathcal L\cong\mathcal M\otimes f^*\mathcal N.
\]
\end{solution}

\begin{exerciseintro}
% record: H-III-12-005
\exerciseheading{Exercise III.12.5}
\addcontentsline{toc}{subsection}{Exercise III.12.5}

\begin{context}
The assertion requires \(\operatorname{rank}\mathcal E\geq2\);
for a rank-one bundle, \(\mathbb P(\mathcal E)\cong Y\) and there is
no additional independent \(\mathbb Z\)-factor.

We use constancy of the Hilbert polynomial in a flat projective
family and Exercise III.12.4.
\end{context}

\begin{problem}
Let \(Y\) be an integral scheme of finite type over an algebraically
closed field \(k\).  Let \(\mathcal E\) be a locally free sheaf of
rank at least \(2\) on \(Y\), and put
\[
X=\mathbb P(\mathcal E).
\]
Prove that
\[
\operatorname{Pic}X\cong
\bigl(\operatorname{Pic}Y\bigr)\times\mathbb Z.
\]
\end{problem}
\end{exerciseintro}

\begin{solution}
Let
\[
\pi:X=\mathbb P(\mathcal E)\longrightarrow Y
\]
be the projection, and write
\(\operatorname{rank}\mathcal E=r+1\), where \(r\geq1\).  Define
\[
\Psi:\operatorname{Pic}Y\times\mathbb Z
\longrightarrow\operatorname{Pic}X,
\qquad
(\mathcal M,d)\longmapsto
\pi^*\mathcal M\otimes\mathcal O_X(d).
\tag{1}
\]

We first prove surjectivity.  Let \(\mathcal L\) be an invertible
sheaf on \(X\).  For every point \(y\in Y\),
\[
X_y\cong\mathbb P_{k(y)}^r,
\qquad
\mathcal L_y\cong\mathcal O_{\mathbb P^r}(d_y)
\]
for a unique integer \(d_y\).  The sheaf \(\mathcal L\) is flat over
\(Y\), and \(\mathcal O_X(1)\) is relatively ample.  Therefore the
fiber Hilbert polynomial
\[
\chi\bigl(X_y,\mathcal L_y(m)\bigr)
=\binom{m+d_y+r}{r}
\tag{2}
\]
is locally constant in \(y\).  Since \(Y\) is integral, it is
connected, so (2) is independent of \(y\).  Equality of these
polynomials forces all the integers \(d_y\) to have one common value
\(d\).

Put
\[
\mathcal L'=\mathcal L\otimes\mathcal O_X(-d).
\]
Its restriction to every fiber is trivial.  The morphism \(\pi\) is
flat and projective with integral fibers, so Exercise III.12.4 gives
an invertible sheaf \(\mathcal M\) on \(Y\) such that
\[
\mathcal L'\cong\pi^*\mathcal M.
\]
Hence
\[
\mathcal L\cong\pi^*\mathcal M\otimes\mathcal O_X(d),
\]
and \(\Psi\) is surjective.

For injectivity, suppose
\[
\pi^*\mathcal M\otimes\mathcal O_X(d)\cong\mathcal O_X.
\]
Restriction to any fiber gives
\(\mathcal O_{\mathbb P^r}(d)\cong\mathcal O_{\mathbb P^r}\), so
\(d=0\).  We are left with \(\pi^*\mathcal M\cong\mathcal O_X\).
Since \(\pi_*\mathcal O_X=\mathcal O_Y\), the projection formula
gives
\[
\mathcal M
\cong\mathcal M\otimes\pi_*\mathcal O_X
\cong\pi_*\pi^*\mathcal M
\cong\mathcal O_Y.
\]
Thus the kernel of \(\Psi\) is zero, and (1) is the desired
isomorphism.
\end{solution}

\begin{exerciseintro}
% record: H-III-12-006
\exerciseheading{Exercise III.12.6}
\addcontentsline{toc}{subsection}{Exercise III.12.6}

\begin{context}
We use Exercise III.11.8 and the cohomology-and-base-change theorem.
In particular, if \(p:Z\to S\) is projective, \(\mathcal A\) is
flat over \(S\), and
\[
H^1(Z_s,\mathcal A_s)=0,
\]
then, after shrinking around \(s\), the fiber section map
\[
(p_*\mathcal A)\otimes k(s)
\longrightarrow H^0(Z_s,\mathcal A_s)
\]
is surjective.

An integral projective scheme over an algebraically closed field has
only constant global functions.
Assume \(T\) is nonempty when choosing the reference point \(t_0\).
\end{context}

\begin{problem}
Let \(X\) be an integral projective scheme over an algebraically
closed field \(k\), and suppose that
\[
H^1(X,\mathcal O_X)=0.
\]
Let \(T\) be a connected scheme of finite type over \(k\).

(a) If \(\mathcal L\) is an invertible sheaf on \(X\times T\), prove
that the invertible sheaves
\[
\mathcal L_t=\mathcal L|_{X\times\{t\}}
\]
on \(X\) are isomorphic for all closed points \(t\in T\).

(b) Prove, without assuming that \(T\) is reduced, that
\[
\operatorname{Pic}(X\times T)
\cong\operatorname{Pic}X\times\operatorname{Pic}T.
\]
\end{problem}
\end{exerciseintro}

\begin{solution}
Write
\[
p:X\times T\longrightarrow T,
\qquad
q:X\times T\longrightarrow X
\]
for the two projections.

We first establish a local observation.  Suppose that \(t\in T\) is
a closed point and that an invertible sheaf \(\mathcal A\) on
\(X\times T\) satisfies
\[
\mathcal A_t\cong\mathcal O_X.
\tag{1}
\]
The sheaf \(\mathcal A\) is flat over \(T\), and
\[
H^1(X,\mathcal A_t)=H^1(X,\mathcal O_X)=0.
\]
Exercise III.11.8 makes \(R^1p_*\mathcal A\) zero near \(t\), and
cohomology and base change then makes
\[
(p_*\mathcal A)\otimes k(t)
\longrightarrow H^0(X,\mathcal A_t)
\tag{2}
\]
surjective.  After replacing \(T\) by an affine neighborhood \(U\)
of \(t\), lift the section \(1\in H^0(X,\mathcal O_X)\), using (1)
and (2), to
\[
\sigma\in H^0(X\times U,\mathcal A).
\]
Its zero locus \(Z(\sigma)\) is closed and misses \(X\times\{t\}\).
Since \(p\) is proper, \(p(Z(\sigma))\) is closed.  Replacing \(U\)
by \(U\setminus p(Z(\sigma))\), the section is nowhere zero and
therefore gives an isomorphism
\[
\mathcal O_{X\times U}\xrightarrow{\sim}
\mathcal A|_{X\times U}.
\tag{3}
\]
This argument is valid even when \(T\) is nonreduced.

(a) Fix a closed point \(t_0\in T\), and replace \(\mathcal L\) by
\[
\mathcal L\otimes q^*(\mathcal L_{t_0}^{-1}).
\]
Its fiber at \(t_0\) is trivial, so (3) shows that its fibers are
trivial on a neighborhood of \(t_0\).  Applying the same argument
at every closed point shows that
\[
t\longmapsto[\mathcal L_t]\in\operatorname{Pic}X
\]
is locally constant on the closed points of \(T\).  A scheme of
finite type over \(k\) is Jacobson.  Neighborhoods belonging to two
different isomorphism classes cannot meet, since a nonempty
intersection would contain a closed point belonging to both classes.
Their union is all of \(T\), since a nonempty closed complement would
also contain a closed point.  They therefore give a disjoint open
partition of \(T\).  Since \(T\) is connected, there is only one
part, and all \(\mathcal L_t\) are isomorphic.

(b) Consider the homomorphism
\[
\Theta:\operatorname{Pic}X\times\operatorname{Pic}T
\longrightarrow\operatorname{Pic}(X\times T),
\qquad
(\mathcal A,\mathcal N)\longmapsto
q^*\mathcal A\otimes p^*\mathcal N.
\tag{4}
\]
To prove surjectivity, let \(\mathcal L\) be a line bundle on
\(X\times T\), choose a closed point \(t_0\), and put
\[
\mathcal A=\mathcal L_{t_0},
\qquad
\mathcal P=\mathcal L\otimes q^*(\mathcal A^{-1}).
\]
Part (a) gives
\(\mathcal P_t\cong\mathcal O_X\) for every closed \(t\).  The local
observation (3), applied at every closed point, gives an open cover
\(\{U_i\}\) of \(T\) such that
\[
\mathcal P|_{X\times U_i}\cong\mathcal O_{X\times U_i}.
\tag{5}
\]
The opens cover all of \(T\), because a nonempty closed subset of a
finite-type \(k\)-scheme contains a closed point.

Now
\[
p_*\mathcal O_{X\times T}=\mathcal O_T,
\tag{6}
\]
since \(H^0(X,\mathcal O_X)=k\); this equality remains valid for a
nonreduced \(T\).  On \(X\times(U_i\cap U_j)\), the two
trivializations in (5) consequently differ by the pullback of a
unique unit on \(U_i\cap U_j\).  These units form a cocycle defining
an invertible sheaf \(\mathcal N\) on \(T\), and (5) glues to
\[
\mathcal P\cong p^*\mathcal N.
\]
Thus
\[
\mathcal L\cong q^*\mathcal A\otimes p^*\mathcal N,
\]
so \(\Theta\) is surjective.

Finally, suppose
\[
q^*\mathcal A\otimes p^*\mathcal N\cong\mathcal O_{X\times T}.
\]
Restriction to \(X\times\{t_0\}\) gives
\(\mathcal A\cong\mathcal O_X\).  Hence
\(p^*\mathcal N\cong\mathcal O_{X\times T}\).  Choose a closed
\(k\)-point \(x\in X\), and pull back along the section
\[
T\longrightarrow X\times T,\qquad t\longmapsto(x,t).
\]
This gives \(\mathcal N\cong\mathcal O_T\).  Therefore \(\Theta\)
is injective, and (4) is an isomorphism.
\end{solution}

\corpuschapter{Chapter IV}
\addcontentsline{toc}{chapter}{Chapter IV}

\begin{exerciseintro}
\corpussection{IV.1 Riemann--Roch Theorem}
\addcontentsline{toc}{section}{IV.1 Riemann--Roch Theorem}

% record: H-IV-01-001
\exerciseheading{Exercise IV.1.1}
\addcontentsline{toc}{subsection}{Exercise IV.1.1}

\begin{context}
Throughout, a curve is a complete nonsingular integral curve over an
algebraically closed field \(k\).  For a divisor \(D\) on a curve \(X\), set

\[
L(D)=H^0(X,\mathcal O_X(D))
=\{f\in K(X)^\times:\operatorname{div}(f)+D\geq 0\}\cup\{0\}
\]

and write \(\ell(D)=\dim_k L(D)\).  The Riemann--Roch theorem says

\[
\ell(D)-\ell(K-D)=\deg D+1-g,
\]

where \(K\) is a canonical divisor and \(g\) is the genus of \(X\).
A divisor of negative degree has no nonzero global section.
\end{context}

\begin{problem}
Let \(X\) be a curve and let \(P\in X\) be a point.  Prove that there is a
nonconstant rational function \(f\in K(X)\) that is regular on
\(X\setminus\{P\}\).
\end{problem}
\end{exerciseintro}

\begin{solution}
Choose an integer \(m\geq 2g+1\).  Since \(\deg K=2g-2\),

\[
\deg(K-mP)=2g-2-m<0.
\]

Thus \(\ell(K-mP)=0\), and Riemann--Roch gives

\[
\ell(mP)=m+1-g\geq g+2>1.
\]

The constants form a one-dimensional subspace of \(L(mP)\), so there is
an \(f\in L(mP)\) that is not constant.  The condition

\[
\operatorname{div}(f)+mP\geq0
\]

says that \(f\) has no pole away from \(P\).  Hence \(f\) is regular on
\(X\setminus\{P\}\), as required.  In fact, \(f\) must have a pole at
\(P\), because every global regular function on the complete integral
curve \(X\) is constant.
\end{solution}

\begin{exerciseintro}
% record: H-IV-01-002
\exerciseheading{Exercise IV.1.2}
\addcontentsline{toc}{subsection}{Exercise IV.1.2}

\begin{context}
For a divisor \(D\) on a complete nonsingular curve \(X\), write

\[
L(D)=H^0(X,\mathcal O_X(D)),
\qquad
\ell(D)=\dim_k L(D).
\]

If \(\deg D>2g-2\), Riemann--Roch and \(\deg K=2g-2\) give

\[
\ell(D)=\deg D+1-g.
\]

The ground field \(k\) is algebraically closed and therefore infinite; a
finite-dimensional vector space over an infinite field is not the union
of finitely many proper linear subspaces.
For an empty list of prescribed poles, take a constant function.
\end{context}

\begin{problem}
Let \(X\) be a curve, and let \(P_1,\ldots,P_r\in X\) be points.  Prove
that there is a rational function \(f\in K(X)\) having a pole of positive
order at every \(P_i\) and no poles elsewhere.
\end{problem}
\end{exerciseintro}

\begin{solution}
Discarding repetitions, we may suppose that \(P_1,\ldots,P_r\) are
distinct.  Choose \(m>0\) so large that

\[
mr-1>2g-2,
\qquad
D=m(P_1+\cdots+P_r).
\]

For each \(i\), both \(D\) and \(D-P_i\) have degree greater than
\(2g-2\).  Riemann--Roch therefore gives

\[
\ell(D)=mr+1-g,
\qquad
\ell(D-P_i)=mr-g.
\]

Consequently \(L(D-P_i)\) is a proper, codimension-one subspace of
\(L(D)\).  Since \(k\) is infinite, the finite union

\[
\bigcup_{i=1}^r L(D-P_i)
\]

does not exhaust \(L(D)\).  Choose

\[
f\in L(D)\setminus\bigcup_{i=1}^r L(D-P_i).
\]

Membership in \(L(D)\) shows that \(f\) is regular away from the
\(P_i\) and has pole order at most \(m\) at each \(P_i\).  If its pole
order at \(P_i\) were at most \(m-1\), then
\(f\in L(D-P_i)\), contrary to the choice of \(f\).  Thus \(f\) has pole
order exactly \(m\) at every \(P_i\) and has no other poles.
\end{solution}

\begin{exerciseintro}
% record: H-IV-01-003
\exerciseheading{Exercise IV.1.3}
\addcontentsline{toc}{subsection}{Exercise IV.1.3}

\begin{context}
We use the following standard facts about curves over an algebraically
closed field \(k\).

An integral separated regular scheme of dimension one and finite type over
\(k\) is an open subscheme of a complete nonsingular curve.  A rational
function on a complete nonsingular curve determines a morphism to
\(\mathbb P^1\); the inverse image of \(\infty\) is precisely the support
of its pole divisor.  Finally, a nonconstant morphism between complete
integral curves is finite, and a finite morphism is affine.

We may also use Exercise IV.1.2: given finitely many points on a curve,
there is a rational function having a pole at every one of those points
and no poles elsewhere.
\end{context}

\begin{problem}
Let \(X\) be an integral, separated, regular, one-dimensional scheme of
finite type over \(k\).  Assume that \(X\) is not proper over \(k\).  Prove
that \(X\) is affine.
\end{problem}
\end{exerciseintro}

\begin{solution}
Embed \(X\) as a dense open subscheme of a complete nonsingular curve
\(\overline X\).  Its complement is a proper closed subset of the
one-dimensional integral scheme \(\overline X\), hence is a finite set of
closed points:

\[
\overline X\setminus X=\{P_1,\ldots,P_r\}.
\]

This set is nonempty, because otherwise \(X=\overline X\) would be proper.

By Exercise IV.1.2, there is \(f\in K(\overline X)\) having a pole at every
\(P_i\) and no poles elsewhere.  The associated morphism

\[
\overline f:\overline X\longrightarrow\mathbb P^1
\]

therefore satisfies

\[
\overline f^{-1}(\infty)=\{P_1,\ldots,P_r\},
\qquad
\overline f^{-1}(\mathbb A^1)=X.
\]

Because \(f\) has a pole, it is nonconstant.  Thus \(\overline f\) is a
finite morphism of complete integral curves.  Its restriction

\[
X=\overline f^{-1}(\mathbb A^1)\longrightarrow\mathbb A^1
\]

is finite by base change.  A scheme finite over an affine scheme is
affine, so \(X\) is affine.
\end{solution}

\begin{exerciseintro}
% record: H-IV-01-004
\exerciseheading{Exercise IV.1.4}
\addcontentsline{toc}{subsection}{Exercise IV.1.4}

\begin{context}
We may use Exercise IV.1.3, which says that every nonproper integral
regular curve of finite type over \(k\) is affine, together with the
following consequences of the cohomological criterion for affineness.

First, a noetherian scheme is affine if and only if its reduction is
affine.  Second, a reduced noetherian scheme is affine if and only if
all of its irreducible components, with their reduced structures, are
affine.  Third, if \(Y\to Z\) is a finite surjective morphism of
noetherian schemes and \(Y\) is affine, then \(Z\) is affine.

The normalization of an integral curve of finite type over a field is
finite and is a regular curve.  Properness descends along a finite
surjective morphism: if the source is proper over \(k\), then so is the
target.
\end{context}

\begin{problem}
Let \(X\) be a separated, one-dimensional scheme of finite type over
\(k\).  Suppose that none of the irreducible components of \(X\) is
proper over \(k\).  Prove that \(X\) is affine.
\end{problem}
\end{exerciseintro}

\begin{solution}
Let \(X_1,\ldots,X_s\) be the irreducible components of
\(X_{\mathrm{red}}\), each with its reduced induced structure.  Every
\(X_i\) is an integral separated curve of finite type over \(k\), and by
hypothesis it is not proper.  Let

\[
\nu_i:\widetilde X_i\longrightarrow X_i
\]

be its normalization.  This morphism is finite and surjective, and
\(\widetilde X_i\) is regular.  Moreover, \(\widetilde X_i\) cannot be
proper: otherwise properness would descend through \(\nu_i\), contrary
to the hypothesis on \(X_i\).

Exercise IV.1.3 now shows that \(\widetilde X_i\) is affine.  Since
\(\nu_i\) is finite and surjective, its target \(X_i\) is affine as well.
Thus every irreducible component of the reduced noetherian scheme
\(X_{\mathrm{red}}\) is affine.  The component criterion makes
\(X_{\mathrm{red}}\) affine, and the reduction criterion then makes
\(X\) affine.  This last step also accounts for every nilpotent structure
on \(X\).
\end{solution}

\begin{exerciseintro}
% record: H-IV-01-005
\exerciseheading{Exercise IV.1.5}
\addcontentsline{toc}{subsection}{Exercise IV.1.5}

\begin{context}
For a divisor \(D\) on a curve \(X\), write

\[
L(D)=H^0(X,\mathcal O_X(D)),
\qquad
\ell(D)=\dim_kL(D).
\]

If \(D\) is effective, then its complete linear system has dimension
\(\dim|D|=\ell(D)-1\).  Riemann--Roch says

\[
\ell(D)-\ell(K-D)=\deg D+1-g,
\]

and \(\ell(K)=g\).  A nonconstant rational function on a complete
nonsingular curve defines a finite morphism to \(\mathbb P^1\), whose
degree is the degree of its pole divisor.
\end{context}

\begin{problem}
Let \(D\) be an effective divisor on a curve \(X\) of genus \(g\).
Prove that

\[
\dim|D|\leq\deg D.
\]

Prove also that equality holds if and only if \(D=0\) or \(g=0\).
\end{problem}
\end{exerciseintro}

\begin{solution}
Because \(D\geq0\), there is an inclusion

\[
L(K-D)\subseteq L(K).
\]

Hence \(\ell(K-D)\leq\ell(K)=g\).  Riemann--Roch gives

\[
\dim|D|
=\ell(D)-1
=\deg D-g+\ell(K-D)
\leq\deg D.
\tag{1}
\]

If \(D=0\), both sides of (1) are zero.  If \(g=0\), then
\(\ell(K-D)=0\), so Riemann--Roch gives
\(\ell(D)=\deg D+1\); equality again holds.

Conversely, suppose equality holds and \(g>0\).  We first show that the
canonical linear system has no base point.  Fix \(P\in X\).  If
\(\ell(P)\geq2\), then \(L(P)\) contains a nonconstant rational function
whose pole divisor has degree at most one.  Its morphism to
\(\mathbb P^1\) would have degree one, hence would be an isomorphism.
This would give \(g=0\), a contradiction.  Therefore \(\ell(P)=1\).
Applying Riemann--Roch to \(P\) gives

\[
1-\ell(K-P)=2-g,
\qquad\text{so}\qquad
\ell(K-P)=g-1.
\tag{2}
\]

Thus \(L(K-P)\) is a proper subspace of \(L(K)\) for every point \(P\).

Equality in (1) forces \(\ell(K-D)=g\), and hence
\(L(K-D)=L(K)\).  If \(D\neq0\), choose a point \(P\) in its support.
Then \(D\geq P\), so

\[
L(K-D)\subseteq L(K-P)\subsetneq L(K),
\]

contradicting \(L(K-D)=L(K)\).  Therefore \(D=0\).  This proves the
stated characterization of equality.
\end{solution}

\begin{exerciseintro}
% record: H-IV-01-006
\exerciseheading{Exercise IV.1.6}
\addcontentsline{toc}{subsection}{Exercise IV.1.6}

\begin{context}
For a divisor \(D\) on a curve \(X\), Riemann--Roch gives

\[
\ell(D)=\deg D+1-g+\ell(K-D).
\]

A nonconstant rational function \(f\in K(X)\) defines a finite morphism
\(f:X\to\mathbb P^1\).  Its degree is the degree of the pole divisor
\(f^*(\infty)\).
\end{context}

\begin{problem}
Let \(X\) be a curve of genus \(g\).  Prove that there is a finite
morphism

\[
f:X\longrightarrow\mathbb P^1
\]

of degree at most \(g+1\).
\end{problem}
\end{exerciseintro}

\begin{solution}
Choose a point \(P\in X\) and put \(D=(g+1)P\).  Riemann--Roch gives

\[
\ell(D)=\deg D+1-g+\ell(K-D)
=2+\ell(K-D)\geq2.
\]

Thus \(L(D)\) contains a nonconstant rational function \(f\).  Since
\(f\in L(D)\), its only possible pole is at \(P\), with order at most
\(g+1\).  It must have a pole, because a global regular function on the
complete integral curve \(X\) is constant.  Consequently the morphism

\[
f:X\longrightarrow\mathbb P^1
\]

is nonconstant and therefore finite, and

\[
\deg f=\deg f^*(\infty)\leq g+1.
\]
\end{solution}

\begin{exerciseintro}
% record: H-IV-01-007
\exerciseheading{Exercise IV.1.7}
\addcontentsline{toc}{subsection}{Exercise IV.1.7}

\begin{context}
A curve \(X\) is hyperelliptic if \(g(X)\geq2\) and it admits a finite
morphism of degree two to \(\mathbb P^1\).  On a curve of positive genus,
the canonical linear system is base-point-free: for every \(P\in X\),

\[
\ell(K-P)=g-1<g=\ell(K).
\]

A base-point-free linear system of dimension one defines a morphism to
\(\mathbb P^1\).  If the system is complete and corresponds to a divisor
\(D\), then the resulting morphism \(f\) satisfies
\(f^*\mathcal O_{\mathbb P^1}(1)\cong\mathcal O_X(D)\).

Identify a nonsingular quadric surface with
\(Q=\mathbb P^1\times\mathbb P^1\).  A curve of type \((a,b)\) is the zero
divisor of a section of \(\mathcal O_Q(a,b)\); its intersection with a
fiber of the first projection is \(b\).
Work over an algebraically closed field.
Adjunction on the quadric gives genus \((a-1)(b-1)\) for a smooth
curve of type \((a,b)\).  For \(a,b>0\), \(\mathcal O_Q(a,b)\) is
very ample and Bertini gives a smooth integral general member.
\end{context}

\begin{problem}
(a) Let \(X\) be a curve of genus \(2\).  Prove that its canonical divisor
defines a complete linear system \(|K|\) of degree \(2\) and dimension
\(1\), with no base points.  Deduce that \(X\) is hyperelliptic.

(b) For every \(g\geq2\), let \(X\) be a nonsingular irreducible curve of
type \((g+1,2)\) on
\(\mathbb P^1\times\mathbb P^1\).  Prove that \(X\) admits a morphism of
degree \(2\) to \(\mathbb P^1\).  Thus hyperelliptic curves exist in every
genus \(g\geq2\).
\end{problem}
\end{exerciseintro}

\begin{solution}
(a) Since \(g=2\),

\[
\deg K=2g-2=2,
\qquad
\ell(K)=g=2.
\]

Therefore \(|K|\) has degree \(2\) and projective dimension
\(\ell(K)-1=1\).  The canonical system has no base points, so it defines
a morphism

\[
f:X\longrightarrow\mathbb P^1
\]

with \(f^*\mathcal O_{\mathbb P^1}(1)\cong\mathcal O_X(K)\).  It follows
that

\[
\deg f=\deg f^*\mathcal O_{\mathbb P^1}(1)=\deg K=2.
\]

The morphism is nonconstant and hence finite.  Thus \(X\) is
hyperelliptic.

(b) Let \(\pi_1:Q\to\mathbb P^1\) be the first projection.  For a general
point \(t\in\mathbb P^1\), its fiber
\(F_t=\{t\}\times\mathbb P^1\) meets \(X\) in a zero-dimensional scheme
of length

\[
X\mathbin{\cdot}F_t=2,
\]

because \(X\) has type \((g+1,2)\).  Hence
\(\pi_1|_X:X\to\mathbb P^1\) is a nonconstant morphism of degree \(2\).
It is finite because \(X\) is complete.  Adjunction gives
\(g(X)=((g+1)-1)(2-1)=g\).
The system \(\mathcal O_Q(g+1,2)\) is very ample, so Bertini
supplies a smooth integral member.  Thus the construction yields
a hyperelliptic curve in each genus \(g\geq2\).
\end{solution}

\begin{exerciseintro}
% record: H-IV-01-008
\exerciseheading{Exercise IV.1.8}
\addcontentsline{toc}{subsection}{Exercise IV.1.8}

\begin{context}
Let \(X\) be an integral projective scheme of dimension one over an
algebraically closed field \(k\), and let
\(\nu:\widetilde X\to X\) be its normalization.  This is a finite
morphism.  There is an exact sequence

\[
0\longrightarrow\mathcal O_X
\longrightarrow\nu_*\mathcal O_{\widetilde X}
\longrightarrow
\bigoplus_{P\in X}\overline{\mathcal O}_{P}/\mathcal O_{X,P}
\longrightarrow0,
\tag{1}
\]

where the final terms are skyscraper sheaves and only finitely many are
nonzero.  Put

\[
\delta_P=\operatorname{length}_{\mathcal O_{X,P}}
(\overline{\mathcal O}_{P}/\mathcal O_{X,P}).
\]

For a projective curve \(Y\),
\(p_a(Y)=1-\chi(\mathcal O_Y)\).  A finite morphism has no higher direct
images on quasi-coherent sheaves, so
\(\chi(\nu_*\mathcal O_{\widetilde X})=
\chi(\mathcal O_{\widetilde X})\).

For a reduced curve singularity \(A\) of finite type over \(k\), completion
preserves the length of the normalization quotient:

\[
\operatorname{length}_A(\overline A/A)
=\dim_k(\overline{\widehat A}/\widehat A).
\tag{2}
\]

Thus this length depends only on the analytic isomorphism class of the
singularity.
\end{context}

\begin{problem}
With the notation above:

(a) Prove that

\[
p_a(X)=p_a(\widetilde X)+\sum_{P\in X}\delta_P.
\]

(b) Prove that if \(p_a(X)=0\), then \(X\) is nonsingular and
\(X\cong\mathbb P^1\).

(c) If \(P\) is a node or an ordinary cusp, prove that
\(\delta_P=1\).
\end{problem}
\end{exerciseintro}

\begin{solution}
(a) Taking Euler characteristics in (1) gives

\[
\chi(\mathcal O_X)
=\chi(\nu_*\mathcal O_{\widetilde X})
-\sum_{P\in X}\delta_P
=\chi(\mathcal O_{\widetilde X})
-\sum_{P\in X}\delta_P.
\]

Here a skyscraper sheaf of length \(n\) has \(n\)-dimensional space of
global sections and no higher cohomology.  Subtracting the displayed
identity from \(1\) on both sides yields

\[
p_a(X)=p_a(\widetilde X)+\sum_{P\in X}\delta_P.
\]

(b) The normalization \(\widetilde X\) is a complete nonsingular curve,
so \(p_a(\widetilde X)=g(\widetilde X)\geq0\); every
\(\delta_P\) is also nonnegative.  If \(p_a(X)=0\), part (a) therefore
forces

\[
g(\widetilde X)=0,
\qquad
\delta_P=0\quad\text{for every }P\in X.
\]

The vanishing of every normalization quotient says that
\(\mathcal O_X\to\nu_*\mathcal O_{\widetilde X}\) is an isomorphism.
Hence \(\nu\) is an isomorphism and \(X\) is nonsingular.  A complete
nonsingular curve of genus zero over an algebraically closed field is
isomorphic to \(\mathbb P^1\), so \(X\cong\mathbb P^1\).

(c) By (2), we may compute after replacing the singularity by its
standard analytic model.

For a node, the completed local ring and its normalization are

\[
A=k[[x,y]]/(xy),
\qquad
\overline A=k[[x]]\oplus k[[y]].
\]

The map \(A\to\overline A\) sends the class of \(h(x,y)\) to
\((h(x,0),h(0,y))\).  Its image consists exactly of the pairs
\((a(x),b(y))\) satisfying \(a(0)=b(0)\).  Therefore

\[
\overline A/A\cong k,
\]

with the quotient measured by \(a(0)-b(0)\), and the node has
\(\delta_P=1\).

For an ordinary cusp, the standard analytic model is

\[
A=k[[x,y]]/(y^2-x^3)
\cong k[[t^2,t^3]]\subseteq k[[t]]=\overline A,
\]

under \(x\mapsto t^2\) and \(y\mapsto t^3\).  Since every integer at
least two lies in the semigroup generated by \(2\) and \(3\),

\[
k[[t^2,t^3]]=k+t^2k[[t]].
\]

Thus \(\overline A/A\) is the one-dimensional vector space generated by
the class of \(t\).  Hence an ordinary cusp also has
\(\delta_P=1\).
\end{solution}

\begin{exerciseintro}
% record: H-IV-01-009
\exerciseheading{Exercise IV.1.9}
\addcontentsline{toc}{subsection}{Exercise IV.1.9}

\begin{context}
Let \(X\) be an integral projective scheme of dimension one over an
algebraically closed field \(k\), and let \(X_{\mathrm{reg}}\) be its
regular locus.  A divisor supported in \(X_{\mathrm{reg}}\) is a finite
sum \(D=\sum n_iP_i\); define \(\deg D=\sum n_i\).  Every such divisor is
Cartier and determines an invertible sheaf \(\mathcal O_X(D)\).

We use \(p_a(X)=1-\chi(\mathcal O_X)\), additivity of Euler
characteristic, Serre's theorem on twists of coherent sheaves, and the
standard properties of ample and very ample invertible sheaves.  The
singular locus of an integral curve is finite.

If \(X\) is a locally complete intersection in projective space, its
dualizing sheaf \(\omega_X\) is invertible, and Serre duality gives

\[
H^1(X,\mathcal L)^\vee
\cong H^0(X,\omega_X\otimes\mathcal L^{-1})
\]

for every invertible sheaf \(\mathcal L\).
\end{context}

\begin{problem}
(a) Let \(D=\sum n_iP_i\) be a divisor whose support lies in
\(X_{\mathrm{reg}}\).  Prove that

\[
\chi(\mathcal O_X(D))=\deg D+1-p_a(X).
\]

(b) Prove that every Cartier divisor on \(X\) is the difference of two
very ample Cartier divisors.

(c) Deduce that every invertible sheaf \(\mathcal L\) on \(X\) is
isomorphic to \(\mathcal O_X(D)\) for a divisor \(D\) supported in
\(X_{\mathrm{reg}}\).

(d) Suppose in addition that \(X\) is a locally complete intersection in
some projective space.  Choose a divisor \(K\), supported in
\(X_{\mathrm{reg}}\), such that
\(\mathcal O_X(K)\cong\omega_X\).  With
\(\ell(D)=h^0(X,\mathcal O_X(D))\), prove

\[
\ell(D)-\ell(K-D)=\deg D+1-p_a(X).
\]
\end{problem}
\end{exerciseintro}

\begin{solution}
(a) Let \(P\in X_{\mathrm{reg}}\).  Its ideal sheaf is the invertible
sheaf \(\mathcal O_X(-P)\), and tensoring its defining sequence by
\(\mathcal O_X(D+P)\) gives

\[
0\longrightarrow\mathcal O_X(D)
\longrightarrow\mathcal O_X(D+P)
\longrightarrow k(P)\longrightarrow0.
\tag{1}
\]

Since \(k\) is algebraically closed, the skyscraper sheaf \(k(P)\) has
Euler characteristic one.  Thus

\[
\chi(\mathcal O_X(D+P))=\chi(\mathcal O_X(D))+1.
\tag{2}
\]

Starting with \(D=0\), every divisor supported in the regular locus can
be reached by repeatedly adding or subtracting a point.  Equation (2)
therefore gives

\[
\chi(\mathcal O_X(D))
=\chi(\mathcal O_X)+\deg D
=\deg D+1-p_a(X).
\]

(b) Fix a very ample Cartier divisor \(H\), and let \(C\) be any Cartier
divisor.  For all sufficiently large \(n\), the invertible sheaf

\[
\mathcal O_X(C+(n-1)H)
\]

is generated by global sections.  Since \(\mathcal O_X(H)\) is very
ample, their tensor product \(\mathcal O_X(C+nH)\) is very ample.
Increasing \(n\) if necessary, \(\mathcal O_X(nH)\) is also very ample
(indeed, its embedding is obtained from that of \(H\) by a Veronese
embedding).  Hence

\[
C=(C+nH)-nH
\]

expresses \(C\) as the difference of two very ample Cartier divisors.

(c) Choose a nonzero rational section of \(\mathcal L\).  Its local
expressions in trivializations of \(\mathcal L\) define a Cartier divisor
\(C\) with

\[
\mathcal L\cong\mathcal O_X(C).
\]

By part (b), write \(C=C_1-C_2\), where each
\(\mathcal O_X(C_i)\) is very ample.  Use each of these very ample sheaves
to embed \(X\) in a projective space.  Since the singular locus is finite
and \(k\) is infinite, there is a hyperplane containing none of its
points.  Its pullback is an effective Cartier divisor \(E_i\), supported
entirely in \(X_{\mathrm{reg}}\), with \(E_i\sim C_i\).  Therefore

\[
\mathcal L
\cong\mathcal O_X(C)
\cong\mathcal O_X(E_1-E_2),
\]

and \(E_1-E_2\) has support in the regular locus.

(d) Apply part (a) to \(D\):

\[
\deg D+1-p_a(X)
=\chi(\mathcal O_X(D))
=h^0(X,\mathcal O_X(D))-h^1(X,\mathcal O_X(D)).
\]

Serre duality and the choice of \(K\) identify the final term as

\[
h^1(X,\mathcal O_X(D))
=h^0(X,\omega_X\otimes\mathcal O_X(-D))
=h^0(X,\mathcal O_X(K-D))
=\ell(K-D).
\]

The desired formula follows.
\end{solution}

\begin{exerciseintro}
% record: H-IV-01-010
\exerciseheading{Exercise IV.1.10}
\addcontentsline{toc}{subsection}{Exercise IV.1.10}

\begin{context}
Let \(X\) be an integral projective locally complete-intersection curve
over an algebraically closed field \(k\), and let \(X_{\mathrm{reg}}\)
denote its regular locus.  Every invertible sheaf is represented by a
divisor supported in \(X_{\mathrm{reg}}\), and its degree is the degree of
such a divisor.  Write \(\operatorname{Pic}^0X\) for the subgroup of
degree-zero invertible sheaves.

The dualizing sheaf \(\omega_X\) is invertible.  If \(K\) represents it
and \(D\) is supported in \(X_{\mathrm{reg}}\), singular Riemann--Roch
and Serre duality give

\[
\ell(D)-\ell(K-D)=\deg D+1-p_a(X).
\tag{1}
\]

A nonzero section of an invertible sheaf on an integral curve cuts out
an effective Cartier divisor \(E\).  Its length satisfies

\[
\operatorname{length}(E)
=\chi(\mathcal O_X(E))-\chi(\mathcal O_X),
\]

by the defining exact sequence.  An effective Cartier divisor of length
one is a single regular point: locally, length one means that its local
equation generates the maximal ideal, so the one-dimensional local ring
is regular.
\end{context}

\begin{problem}
Assume \(p_a(X)=1\), and fix \(P_0\in X_{\mathrm{reg}}\).  Prove that

\[
X_{\mathrm{reg}}\longrightarrow\operatorname{Pic}^0X,
\qquad
P\longmapsto\mathcal O_X(P-P_0),
\]

is a bijection.
\end{problem}
\end{exerciseintro}

\begin{solution}
First determine the canonical divisor.  Since \(X\) is integral and
projective,

\[
h^0(X,\mathcal O_X)=1,
\qquad
\chi(\mathcal O_X)=1-p_a(X)=0,
\]

so \(h^1(X,\mathcal O_X)=1\).  Serre duality gives
\(\ell(K)=h^0(X,\omega_X)=1\).  Applying (1) with \(D=K\) and using
\(\ell(0)=1\) yields

\[
\deg K=0.
\tag{2}
\]

Let \(\mathcal L\in\operatorname{Pic}^0X\).  Represent it by a divisor
\(D\) supported in \(X_{\mathrm{reg}}\), so that
\(\deg D=0\).  Apply (1) to \(D+P_0\).  By (2),

\[
\deg(K-D-P_0)=-1.
\]

An invertible sheaf of negative degree has no nonzero section, because
such a section would cut out an effective Cartier divisor of negative
degree.  Hence

\[
\ell(D+P_0)=\deg(D+P_0)+1-p_a(X)=1.
\]

Choose its nonzero section, and let \(E\) be the resulting effective
Cartier divisor.  Its defining exact sequence and the Euler-characteristic
formula for degree give

\[
\operatorname{length}(E)
=\chi(\mathcal O_X(D+P_0))-\chi(\mathcal O_X)
=1.
\]

Hence \(E\) is a single point \(P\in X_{\mathrm{reg}}\).  Thus

\[
\mathcal L(P_0)
\cong\mathcal O_X(D+P_0)
\cong\mathcal O_X(P),
\]

or \(\mathcal L\cong\mathcal O_X(P-P_0)\).  This proves surjectivity.

For uniqueness, let \(P\in X_{\mathrm{reg}}\).  Equation (2) gives
\(\deg(K-P)=-1\), so another application of (1) yields

\[
\ell(P)=1.
\]

Consequently the complete linear system \(|P|\) consists only of the
divisor \(P\).  If

\[
\mathcal O_X(P-P_0)\cong\mathcal O_X(Q-P_0),
\]

then \(P\sim Q\); since both are effective and \(Q\in|P|\), it follows
that \(Q=P\).  The displayed map is therefore injective as well as
surjective.
\end{solution}

\begin{exerciseintro}
\corpussection{IV.2 Hurwitz's Theorem}
\addcontentsline{toc}{section}{IV.2 Hurwitz's Theorem}

% record: H-IV-02-001
\exerciseheading{Exercise IV.2.1}
\addcontentsline{toc}{subsection}{Exercise IV.2.1}

\begin{context}
A scheme \(Y\) is simply connected if every finite étale cover of \(Y\)
is a disjoint union of copies of \(Y\).  We use the fact that
\(\mathbb P^1\) is simply connected.

We also use the connectedness theorem: on an integral normal projective
variety of dimension at least two, the support of an effective ample
divisor is connected.  Finite étale morphisms and their degrees are
preserved by base change.
All schemes are over an algebraically closed field \(k\).
\end{context}

\begin{problem}
Prove that \(\mathbb P^n\) is simply connected for every \(n\geq1\).
\end{problem}
\end{exerciseintro}

\begin{solution}
The assertion for \(n=1\) is part of the permitted context.  Let
\[
f:X\longrightarrow\mathbb P^n
\]
be a finite étale cover, and suppose first that \(X\) is connected.
Because \(X\) is étale over the smooth variety \(\mathbb P^n\), it is
regular.  A connected regular scheme has only one irreducible component,
so \(X\) is integral.

Choose a flag of linear subspaces
\[
\mathbb P^1=L_1\subset L_2\subset\cdots\subset L_n=\mathbb P^n
\]
and put \(X_i=f^{-1}(L_i)\).  We claim, by descending induction, that
every \(X_i\) is connected.  This is true for \(X_n=X\).  If \(i\geq2\)
and \(X_i\) is connected, then \(X_i\) is again regular and hence
integral and normal.  The subscheme
\[
X_{i-1}=f|_{X_i}^{-1}(L_{i-1})
\]
is the support of the pullback of a hyperplane, which is an effective
ample divisor on \(X_i\).  The connectedness theorem therefore makes
\(X_{i-1}\) connected.  The claim follows.

In particular,
\[
X_1\longrightarrow L_1\cong\mathbb P^1
\]
is a connected finite étale cover.  Since \(\mathbb P^1\) is simply
connected, this cover has degree one.  Degree is unchanged by the base
change to \(L_1\), so \(f\) itself has degree one and is an isomorphism.

For an arbitrary finite étale cover, apply this argument to each
connected component.  Every component maps isomorphically to
\(\mathbb P^n\), so the cover is a finite disjoint union of copies of
\(\mathbb P^n\).
\end{solution}

\begin{exerciseintro}
% record: H-IV-02-002
\exerciseheading{Exercise IV.2.2}
\addcontentsline{toc}{subsection}{Exercise IV.2.2}

\begin{context}
Let \(k\) be algebraically closed of characteristic different from \(2\).
Every curve \(X\) of genus \(2\) has a base-point-free canonical pencil,
which defines a finite morphism \(f:X\to\mathbb P^1\) of degree two.
Changing a basis of \(H^0(X,\omega_X)\) changes \(f\) only by an
automorphism of \(\mathbb P^1\).

For a separable finite morphism \(f:X\to Y\) of degree \(n\), Hurwitz's
formula is
\[
2g(X)-2=n(2g(Y)-2)+\deg R.
\]
In the tame case, a point of ramification index \(e_P\) contributes
\(e_P-1\) to \(R\).

A separable quadratic extension of \(k(x)\) has the form
\(k(x)(\sqrt h)\).  A place is ramified precisely when the valuation of
\(h\) there is odd.  Two such extensions are isomorphic over \(k(x)\)
if their defining elements differ by a square.
\end{context}

\begin{problem}
(a) Let \(X\) be a curve of genus \(2\), and let
\(f:X\to\mathbb P^1\) be its canonical morphism.  Prove that \(f\) is
ramified at exactly six points of \(X\), each with ramification index
\(2\).  Explain why \(X\) determines an unordered configuration of six
points of \(\mathbb P^1\), up to an automorphism of \(\mathbb P^1\).

(b) Given six distinct elements \(a_1,\ldots,a_6\in k\), let
\[
K=k(x,z),\qquad
z^2=\prod_{i=1}^6(x-a_i),
\]
and let \(f:X\to\mathbb P^1\) be the morphism of nonsingular projective
curves corresponding to \(k(x)\subset K\).  Prove that \(g(X)=2\), that
\(f\) is the canonical morphism, and that it is ramified precisely over
\(x=a_1,\ldots,a_6\).

(c) Prove that for any three distinct points \(P_1,P_2,P_3\) of
\(\mathbb P^1\), there is a unique
\(\varphi\in\operatorname{Aut}\mathbb P^1\) satisfying
\[
\varphi(P_1)=0,\qquad
\varphi(P_2)=1,\qquad
\varphi(P_3)=\infty.
\]
Deduce that an ordered branch configuration can be normalized to
\[
(0,1,\infty,\beta_1,\beta_2,\beta_3),
\]
where the \(\beta_i\) are pairwise distinct and avoid \(0,1\).

(d) Let \(S_6\) act on such triples as follows: reorder the six entries
\[
(0,1,\infty,\beta_1,\beta_2,\beta_3)
\]
and then apply the unique automorphism that sends the first three
entries to \(0,1,\infty\); the last three entries form the new triple.
Prove that this defines an action.

(e) Prove that isomorphism classes of genus-\(2\) curves over \(k\) are
in bijection with triples
\[
(\beta_1,\beta_2,\beta_3)\in
(\mathbb A^1\setminus\{0,1\})^3
\]
having pairwise distinct entries, modulo the action of \(S_6\) in
part (d).  In this sense, curves of genus \(2\) depend on three
parameters.
\end{problem}
\end{exerciseintro}

\begin{solution}
(a) Because \(\operatorname{char}k\neq2\), the degree-two map \(f\) is
separable and all of its ramification is tame.  Hurwitz's formula gives
\[
2
=2g(X)-2
=2(2g(\mathbb P^1)-2)+\deg R
=-4+\deg R.
\]
Thus \(\deg R=6\).  At a ramification point of a degree-two map the
ramification index must be \(2\), so each ramification point contributes
one to \(R\).  There are therefore exactly six ramification points.
Their six images are distinct: a fiber of degree two cannot contain two
different points that both have ramification index two.

The map \(f\) is defined by the complete canonical pencil.  A different
basis of that two-dimensional space of sections composes \(f\) with an
element of \(\operatorname{PGL}_2(k)\).  Hence \(X\) determines the
unordered set of six branch points up to
\(\operatorname{Aut}\mathbb P^1\).

(b) The polynomial
\[
h(x)=\prod_{i=1}^6(x-a_i)
\]
is square-free, so \(K/k(x)\) is a separable quadratic extension.  At
the place \(x=a_i\), one has \(v_{a_i}(h)=1\); hence there is one point
above \(a_i\), with ramification index two.  Every other finite place
has even valuation zero unless it is one of these six places.

At infinity, put \(u=x^{-1}\) and \(w=zu^3\).  The equation becomes
\[
w^2=\prod_{i=1}^6(1-a_i u).
\]
At \(u=0\) the right side equals \(1\), so, since
\(\operatorname{char}k\neq2\), there are two distinct unramified points
above infinity.  Thus the six \(a_i\) are exactly the branch points.

Hurwitz's formula now yields
\[
2g(X)-2=2(-2)+6=2,
\]
so \(g(X)=2\).  Let \(D=f^*(\infty)\).  It has degree two, and the
pullbacks of \(1\) and a coordinate on \(\mathbb P^1\) show that
\(\ell(D)\geq2\).  Riemann--Roch on the genus-two curve gives
\[
\ell(D)-\ell(K_X-D)=\deg D+1-g(X)=1.
\]
Hence \(\ell(K_X-D)\geq1\).  This divisor has degree zero, so a nonzero
section forces \(K_X-D\sim0\).  Therefore
\[
\mathcal O_X(D)\cong\omega_X,\qquad \ell(D)=2,
\]
and the pencil defining \(f\) is the complete canonical pencil.  Thus
\(f\) is the canonical morphism, up to the choice of coordinate on its
target.

(c) Choose an affine coordinate \(t\) in which the three points are
finite; after a preliminary change of coordinate this is always
possible.  Then
\[
\varphi(t)=
\frac{(t-P_1)(P_2-P_3)}
     {(t-P_3)(P_2-P_1)}
\]
has the required values.  This formula has the evident homogeneous
interpretation if one of the original points is infinity.

For uniqueness, an automorphism fixing \(0\), \(1\), and \(\infty\)
is represented by a diagonal fractional linear transformation fixing
\(1\), and is therefore the identity.  Applying the unique
\(\varphi\) to an ordering of the six branch points gives the displayed
normal form.  Distinctness of the branch points implies the asserted
conditions on the \(\beta_i\).

(d) Let
\[
\mathbf q=(0,1,\infty,\beta_1,\beta_2,\beta_3).
\]
Use the convention that \(\sigma\in S_6\) replaces this ordered
sextuple by
\((q_{\sigma^{-1}(1)},\ldots,q_{\sigma^{-1}(6)})\), and then normalize
its first three entries as in part (c).  The identity permutation
clearly fixes the triple.

For \(\sigma,\tau\in S_6\), performing the construction first with
\(\tau\) and then with \(\sigma\) has the same effect on the ordered
sextuple as the permutation \(\sigma\tau\), followed by a projective
automorphism.  In both constructions the final first three entries are
\(0,1,\infty\).  The uniqueness in part (c) therefore says that the two
final normalizations are identical.  Hence
\[
\sigma\mathbin{\cdot}
  \bigl(\tau\mathbin{\cdot}(\beta_1,\beta_2,\beta_3)\bigr)
=(\sigma\tau)\mathbin{\cdot}(\beta_1,\beta_2,\beta_3),
\]
which proves the action law.

(e) By part (a), a genus-two curve gives an unordered six-point branch
configuration modulo \(\operatorname{PGL}_2(k)\).  Ordering the points
and applying part (c) produces a triple
\((\beta_1,\beta_2,\beta_3)\).  Changing
the ordering changes this triple precisely by the action in part (d),
so the resulting \(S_6\)-orbit is intrinsic.

Conversely, a six-point configuration on \(\mathbb P^1\) determines a
quadratic cover branched at those points.  After moving infinity away
from the configuration it has the form in part (b); alternatively, if
the normalized configuration contains infinity, one may use
\[
z^2=x(x-1)(x-\beta_1)(x-\beta_2)(x-\beta_3),
\]
whose odd degree supplies ramification at infinity.  Part (b) shows
that its nonsingular projective model has genus two.

It remains to check uniqueness of the cover.  If \(h,h'\in k(x)^\times\)
give quadratic extensions with the same six branch places, then
\(\operatorname{div}(h/h')\) is even.  Since every degree-zero divisor
on \(\mathbb P^1\) is principal, there are \(r\in k(x)^\times\) and
\(c\in k^\times\) with
\[
h/h'=c\,r^2.
\]
The constant \(c\) is a square because \(k\) is algebraically closed,
so the two quadratic extensions are isomorphic.  Projectively
equivalent branch configurations therefore give isomorphic curves,
and an isomorphism of genus-two curves preserves their canonical maps
and hence their branch configurations.  This proves the claimed
bijection.  The triples form a nonempty open subset of
\(\mathbb A^3\), and the remaining quotient is by the finite group
\(S_6\), which explains the three-parameter statement.
\end{solution}

\begin{exerciseintro}
% record: H-IV-02-003
\exerciseheading{Exercise IV.2.3}
\addcontentsline{toc}{subsection}{Exercise IV.2.3}

\begin{context}
Let \(X\subset\mathbb P^2\) be a nonsingular plane curve of degree
\(d\geq2\) over an algebraically closed field \(k\) of characteristic
zero.  Its genus is
\[
g(X)=\frac{(d-1)(d-2)}2.
\]
For \(P\in X\), write \(T_PX\) for its tangent line.  The Gauss map
\[
\gamma:X\longrightarrow(\mathbb P^2)^*,
\qquad P\longmapsto[T_PX],
\]
has image the dual curve \(X^*\).

A point \(P\) is an inflection point if
\(r_P=I_P(X,T_PX)\geq3\), and it is counted with weight \(r_P-2\).
A line tangent at more than one point is a multiple tangent; a strict
bitangent is tangent at exactly two distinct points.

We use Hurwitz's formula, Bézout's theorem, and the formula
\[
p_a(C)=\frac{(m-1)(m-2)}2
\]
for a plane curve \(C\) of degree \(m\).  A node and an ordinary cusp
both have normalization defect one.
A finite birational map from a nonsingular curve is its normalization,
and the points of a normalization index the branches of a curve germ.
For formal power series of orders two and three in characteristic
zero, the corresponding plane branch is formally equivalent to
the ordinary cusp.
\end{context}

\begin{problem}
(a) Fix a line \(L\subset\mathbb P^2\) that is not tangent to \(X\), and
define
\[
\varphi:X\longrightarrow L,\qquad
\varphi(P)=T_PX\cap L.
\]
Prove that \(\varphi\) is ramified at \(P\) exactly when either
\(P\in L\) or \(P\) is an inflection point.  Deduce that \(X\) has only
finitely many inflection points.

(b) Let \(M\) be tangent to \(X\) at distinct points
\(P_1,\ldots,P_r\), which are all its tangency points, none of which is an inflection point.  Prove that
\([M]\in X^*\) is an ordinary \(r\)-fold point.  Deduce that \(X\) has
only finitely many multiple tangents.

(c) Let \(O\in\mathbb P^2\) lie neither on \(X\) nor on an inflectional
or multiple tangent, and project \(X\) from \(O\) to a line.  Prove that
the projection is ramified at \(P\) exactly when \(OP=T_PX\), and then
has ramification index \(2\).  Deduce that exactly \(d(d-1)\) tangents
to \(X\) pass through \(O\), and hence
\[
\deg X^*=d(d-1).
\]

(d) Prove that, outside a finite set of choices of \(O\in X\), exactly
\[
(d+1)(d-2)
\]
tangent lines other than \(T_OX\) pass through \(O\).

(e) Prove that the morphism \(\varphi\) in part (a) has degree
\(d(d-1)\).  Deduce that the weighted number of inflection points is
\[
3d(d-2).
\]
Prove that an ordinary inflection point corresponds to an ordinary
cuspidal branch of \(X^*\).  It is an ordinary cusp of the whole dual
curve if the tangent has no other tangency.

(f) Suppose that \(X^*\) has only nodes and ordinary cusps.  Prove that
the number of bitangents is
\[
\frac12d(d-2)(d-3)(d+3).
\]

(g) Deduce that a plane cubic has exactly nine inflection points, all
ordinary, and that the line through any two distinct inflection points
meets the cubic in a third inflection point.

(h) Deduce that a smooth plane quartic has \(28\) generalized
bitangents, where a tangent with fourfold contact is counted as a
degenerate bitangent.  Clarify what happens to \(X^*\) at such a
hyperflex.
\end{problem}
\end{exerciseintro}

\begin{solution}
We begin with a local calculation used throughout.  Choose affine
coordinates at \(P\) such that \(P=(0,0)\), \(T_PX\) is \(y=0\), and
\(t=x\) is a uniformizer on \(X\).  Then
\[
y(t)=a t^r+\text{terms of order greater than }r,
\qquad a\neq0,\qquad r=I_P(X,T_PX)\geq2.
\tag{1}
\]
The tangent line at \((t,y(t))\) has equation
\[
-y'(t)X+Y+t\,y'(t)-y(t)=0.
\tag{2}
\]
In the affine chart of the dual plane where the coefficient of \(Y\)
is one, the Gauss map is therefore
\[
t\longmapsto
\bigl(-y'(t),\,t\,y'(t)-y(t)\bigr)
=\bigl(-ar t^{r-1}+\cdots,\,
       a(r-1)t^r+\cdots\bigr).
\tag{3}
\]
Characteristic zero ensures that the displayed leading coefficients
are nonzero.

We will also use that \(\gamma\) is the normalization of \(X^*\).
It is nonconstant and finite onto its image.  At a general point,
characteristic zero makes its differential nonzero.  Differentiating
(3) shows that the tangent to its image is \(P^*\).
Two distinct points over a general smooth point of \(X^*\) would
therefore give two different tangents to that smooth point.
This is impossible; thus the generic degree is one.
The finite birational map from the nonsingular curve \(X\) is its
normalization.  In particular its fibers index all local branches.

(a) If \(P\notin L\), choose the projective coordinates above so that
\(L\) is the line at infinity.  The point \(T_{(t,y(t))}X\cap L\)
records the tangent direction, and a local coordinate for it is
\(y'(t)\).  By (1), its order is \(r-1\).  Thus
\[
e_P(\varphi)=r-1\qquad(P\notin L).
\tag{4}
\]
If \(P\in L\), then \(L\neq T_PX\); take \(L\) to be \(x=0\).  The
intersection of (2) with \(L\) has coordinate
\[
y(t)-t\,y'(t)=a(1-r)t^r+\cdots,
\]
so
\[
e_P(\varphi)=r\qquad(P\in L).
\tag{5}
\]
Equations (4) and (5) say exactly that \(\varphi\) ramifies at the
points of \(X\cap L\) and at the inflection points.  The morphism is
nonconstant, and in characteristic zero it is separable, so its
ramification locus is finite.  Hence there are only finitely many
inflection points.

(b) At a noninflectionary point, \(r=2\), and (3) has the form
\[
t\longmapsto(c_1t+\cdots,c_2t^2+\cdots),
\qquad c_1c_2\neq0.
\]
Thus the corresponding branch of \(X^*\) is nonsingular.  Its tangent
line in the dual plane is the line \(P^*\) parametrizing all lines in
\(\mathbb P^2\) through \(P\).  At \([M]\), the \(r\) points
\(P_1,\ldots,P_r\) therefore give \(r\) smooth branches whose tangent
directions are the distinct lines \(P_1^*,\ldots,P_r^*\).  Hence
\([M]\) is an ordinary \(r\)-fold point.

The singular locus of the integral plane curve \(X^*\) is finite.
Every multiple tangent not passing through an inflection point gives
one of these singular points, while the tangent lines at inflection
points form a finite set by part (a).  Thus only finitely many multiple
tangents exist.

(c) Let
\[
\psi:X\longrightarrow\mathbb P^1
\]
be projection from \(O\).  Its fiber through \(P\) is cut out by the
line \(OP\), and the local ramification index is
\[
e_P(\psi)=I_P(X,OP).
\tag{6}
\]
Consequently \(\psi\) ramifies precisely when \(OP=T_PX\).  The
assumptions on \(O\) exclude inflectional tangents, so the contact order
is then two.  They also exclude multiple tangents, so distinct
ramification points yield distinct tangent lines through \(O\).

The map \(\psi\) has degree \(d\).  Since
\[
2g(X)-2=d(d-3),
\]
Hurwitz gives
\[
\deg R_\psi
=2g(X)-2+2d
=d(d-1).
\]
Every ramification point contributes one, proving the tangent count.
The lines through a general point \(O\) form a line \(O^*\) in the dual
plane, and \(O^*\cap X^*\) parametrizes these tangents.  Therefore
\(\deg X^*=d(d-1)\).

(d) Exclude from \(X\) the inflection points and the intersections of
\(X\) with every inflectional or multiple tangent.  Parts (a) and (b)
show that this removes only finitely many points.  For a remaining
point \(O\), projection from \(O\) is defined by the pencil of lines
through \(O\) after its simple base point at \(O\) is removed.  It has
degree \(d-1\).  The local calculation gives ramification index
\(r_O-1=1\) at \(O\), so \(T_OX\) contributes no ramification.  Every
other ramification point has index two and corresponds to exactly one
tangent through \(O\).  Hurwitz now gives
\[
\deg R
=2g(X)-2+2(d-1)
=d(d-3)+2d-2
=(d+1)(d-2),
\]
which is the required number.

(e) For a general \(Q\in L\), the fiber \(\varphi^{-1}(Q)\) consists of
the points \(P\) for which \(T_PX\) passes through \(Q\).  Part (c)
shows that there are \(d(d-1)\) of them, so
\[
\deg\varphi=d(d-1).
\]
Hurwitz and the genus formula give
\[
\deg R_\varphi
=2g(X)-2+2d(d-1)
=d(3d-5).
\tag{7}
\]

Because \(L\) is not tangent to \(X\), it meets \(X\) transversely in
\(d\) points.  Formulas (4) and (5) show that every point of
\(X\cap L\) contributes a baseline amount one to \(R_\varphi\), while
an inflection point of contact order \(r\) contributes an additional
\(r-2\), whether or not it lies on \(L\).  Thus
\[
\deg R_\varphi
=d+\sum_{\text{\(P\) an inflection}}(r_P-2).
\]
Together with (7), this gives
\[
\sum_{\text{\(P\) an inflection}}(r_P-2)=3d(d-2).
\tag{8}
\]

At an ordinary inflection \(r=3\), formula (3) parametrizes the dual
branch by
\[
(u,v)=(c_1t^2+\cdots,c_2t^3+\cdots),
\qquad c_1c_2\neq0.
\]
After formal coordinate changes this branch is an ordinary cusp.
Another point of tangency on the same line contributes another
branch at its dual point.  Thus without the additional hypothesis
in (e), the entire singular point need not be an ordinary cusp.

(f) We have already identified \(\gamma\) with the normalization.
Put \(m=\deg X^*=d(d-1)\).  Under the stated hypothesis, the nodes of
\(X^*\) correspond exactly to bitangents, and its ordinary cusps
correspond exactly to ordinary inflection points.  Equation (8) says
that the number of cusps is
\[
\kappa=3d(d-2).
\]
If \(b\) is the number of nodes, the normalization-genus formula gives
\[
\frac{(m-1)(m-2)}2
=\frac{(d-1)(d-2)}2+b+\kappa.
\]
Solving and substituting \(m=d(d-1)\) and the value of \(\kappa\)
yields
\[
b
=\frac12d(d-2)(d-3)(d+3).
\]

(g) For \(d=3\), equation (8) gives nine inflection points counted with
weight.  Bézout bounds every tangent contact by three, so every
inflection is ordinary and the nine points are distinct.

Let \(H\) denote a line section.  If \(P\) is an inflection point, its
tangent line cuts out \(3P\), so \(3P\sim H\).  If a line through two
distinct inflection points \(P,Q\) cuts out \(P+Q+R\), then
\[
P+Q+R\sim H.
\]
Multiplying by three and using \(3P\sim3Q\sim H\) gives
\[
3R\sim H.
\]
The line-section system is complete: the hypersurface sequence
on \(\mathbb P^2\) gives \(h^0(X,\mathcal O_X(1))=3\).
Thus some line cuts out \(3R\), so \(R\) is also an inflection point.
The third point is distinct from \(P,Q\): otherwise a degree-zero
divisor such as \(P-Q\) would be principal.  The corresponding rational
function would have one simple pole and define a degree-one morphism
from the genus-one curve \(X\) to \(\mathbb P^1\), an impossibility.

(h) For \(d=4\), equation (8) says that the total inflectionary weight
is \(24\).  Let \(f\) be the number of ordinary inflections, \(h\) the
number of hyperflexes with contact order four, and \(b\) the number of
strict bitangents.  Then
\[
f+2h=24.
\tag{9}
\]
The dual of an ordinary inflection is an ordinary cusp of defect one,
and a strict bitangent gives a node of defect one.  At a hyperflex,
(3) has local parametrization
\[
(u,v)=(c_1t^3+\cdots,c_2t^4+\cdots).
\]
Write its complete local subring as \(A=k[[u(t),v(t)]]\).
Every integer \(m\geq6\) is \(3a+4b\) for nonnegative \(a,b\).
Thus a monomial in \(u,v\) has each possible leading order
\(m\geq6\).  Successively canceling leading terms, and using
completeness, gives \(t^6k[[t]]\subseteq A\).
Modulo \(t^6\), the ring \(A\) is spanned by \(1,u,v\), which are
independent because their leading orders are \(0,3,4\).
Consequently
\[
\dim_k(k[[t]]/A)=6-3=3.
\]
This computes the normalization defect without identifying the
subring in the original parameter with a monomial subring.
These exhaust the singularities of \(X^*\): a singular branch of the
Gauss image comes from contact order at least three, and two branches
have the same image exactly when their two points share a tangent line.
For a quartic, Bézout allows only the cases just listed.

Now \(\deg X^*=12\), so
\[
p_a(X^*)=\frac{11\cdot10}{2}=55,
\qquad
g(X)=3.
\]
The normalization formula therefore gives
\[
55=3+b+f+3h.
\tag{10}
\]
Combining (9) and (10) yields
\[
b+h=28.
\]
Thus there are \(28\) generalized bitangent lines: the \(b\) lines
tangent at two distinct points together with the \(h\) hyperflex lines,
each regarded as a bitangent whose two contact points coincide.  Under
the strict definition from part (b), the number of two-point bitangents
is \(28-h\).  The dual singularity at a hyperflex is the unibranch
\((3,4)\)-cusp computed above, not a tacnode.
\end{solution}

\begin{exerciseintro}
% record: H-IV-02-004
\exerciseheading{Exercise IV.2.4}
\addcontentsline{toc}{subsection}{Exercise IV.2.4}

\begin{context}
Let \(k\) be an algebraically closed field of characteristic \(3\), and
let
\[
X=V(F)\subset\mathbb P^2_k,\qquad
F(x,y,z)=x^3y+y^3z+z^3x.
\]
For a smooth plane curve, the Gauss map sends a point \(P\) to its
tangent line in the dual projective plane.  Its image has closure the
dual curve \(X^*\).

A point \(P\) is an inflection point when
\(I_P(X,T_PX)\geq3\).  We use the Jacobian criterion and the fact that
the relative Frobenius of a smooth curve in characteristic \(p\) is a
finite purely inseparable morphism of degree \(p\).
\end{context}

\begin{problem}
Prove that \(X\) is nonsingular, that every point of \(X\) is an
inflection point, and that \(X^*\) is isomorphic to \(X\), although the
natural map \(X\to X^*\) is purely inseparable.
\end{problem}
\end{exerciseintro}

\begin{solution}
In characteristic \(3\), the partial derivatives are
\[
F_x=z^3,\qquad F_y=x^3,\qquad F_z=y^3.
\]
They cannot vanish simultaneously at a projective point.  The Jacobian
criterion therefore shows that \(X\) is nonsingular.  In particular,
\(X\) is irreducible: two positive-degree components in the projective
plane would meet and give a singular point.

Let \(P=[a:b:c]\in X\).  Its tangent line is
\[
T_PX:\quad c^3x+a^3y+b^3z=0.
\tag{1}
\]
Choose a direction \(V=(u,v,w)\) on this line, so that
\[
c^3u+a^3v+b^3w=0.
\tag{2}
\]
Using \((\alpha+t\beta)^3=\alpha^3+t^3\beta^3\), expansion along the
line \(P+tV\) gives
\[
\begin{aligned}
F(P+tV)
={}&F(P)
 +t(c^3u+a^3v+b^3w)\\
&+t^3(u^3b+v^3c+w^3a)
 +t^4(u^3v+v^3w+w^3u).
\end{aligned}
\tag{3}
\]
The constant term vanishes because \(P\in X\), and the coefficient of
\(t\) vanishes by (2).  There is no quadratic term.  Since the tangent
line is not a component of \(X\), equation (3) proves
\[
I_P(X,T_PX)\geq3.
\]
Thus every point of \(X\) is an inflection point.

Write \([U:V:W]\) for coordinates in the dual plane, representing the
line \(Ux+Vy+Wz=0\).  Formula (1) shows that the Gauss map is
\[
\gamma:X\longrightarrow(\mathbb P^2)^*,
\qquad
[a:b:c]\longmapsto[c^3:a^3:b^3].
\tag{4}
\]
The image satisfies
\[
\begin{aligned}
U^3V+V^3W+W^3U
&=c^9a^3+a^9b^3+b^9c^3\\
&=(c^3a+a^3b+b^3c)^3\\
&=F(a,b,c)^3=0.
\end{aligned}
\tag{5}
\]
Consequently \(X^*\) is contained in the plane quartic
\[
C:\quad U^3V+V^3W+W^3U=0.
\]
This quartic is nonsingular by the same partial-derivative calculation
as above, hence integral.  The map (4) is nonconstant, so its image is
dense in \(C\).  Therefore \(X^*=C\).  The equation of \(C\) is the
same as that of \(X\) after merely renaming the coordinates, and hence
\[
X^*\cong X.
\]

Finally, (4) factors as
\[
X\xrightarrow{\,F_{X/k}\,}X^{(3)}
\xrightarrow{\ [A:B:C]\mapsto[C:A:B]\ }X^*,
\tag{6}
\]
where the first arrow is relative Frobenius; the equation has
coefficients in the prime field, so \(X^{(3)}\) is identified with the
same plane quartic.  The second arrow is an isomorphism, whereas
\(F_{X/k}\) is purely inseparable of degree \(3\).  Thus the natural
Gauss map \(X\to X^*\) is purely inseparable, even though its source
and target are isomorphic smooth curves.
\end{solution}

\begin{exerciseintro}
% record: H-IV-02-005
\exerciseheading{Exercise IV.2.5}
\addcontentsline{toc}{subsection}{Exercise IV.2.5}

\begin{context}
Let \(X\) be a nonsingular projective curve of genus \(g\geq2\) over
an algebraically closed field \(k\) of characteristic zero.  Assume the
known finiteness of
\[
G=\operatorname{Aut}(X),\qquad |G|=n.
\]
The fixed-field inclusion
\[
K(X)^G=K(Y)\subset K(X)
\]
defines a finite quotient morphism \(f:X\to Y\) of degree \(n\).

For a tame finite morphism, Hurwitz's formula is
\[
2g(X)-2
=\deg(f)\bigl(2g(Y)-2\bigr)
  +\sum_{P\in X}(e_P-1).
\]
\end{context}

\begin{problem}
(a) If \(P\in X\) is ramified and \(e_P=r\), prove that
\(f^{-1}(f(P))\) consists of exactly \(n/r\) points, all of
ramification index \(r\).  Choose one point \(P_i\) over each branch
point of \(Y\), and write \(r_i=e_{P_i}\).  Prove that
\[
\frac{2g-2}{n}
=2g(Y)-2+\sum_{i=1}^s\left(1-\frac1{r_i}\right).
\]

(b) Suppose \(h,s\geq0\) and \(r_i\geq2\) are integers for which
\[
2h-2+\sum_{i=1}^s\left(1-\frac1{r_i}\right)>0.
\]
Prove that the least possible value of this expression is
\(1/42\).  Deduce Hurwitz's bound
\[
|\operatorname{Aut}(X)|\leq84(g-1).
\]
\end{problem}
\end{exerciseintro}

\begin{solution}
(a) The extension \(K(X)/K(Y)\) is Galois with group \(G\).  The group
\(G\) acts transitively on the points of \(X\) above any given point
of \(Y\).  For \(P\in X\), its stabilizer
\[
G_P=\{\sigma\in G:\sigma(P)=P\}
\]
is the decomposition group of the corresponding discrete valuation.
All residue fields are \(k\), so the decomposition group equals the
inertia group.  Characteristic zero makes the ramification tame, and
hence
\[
|G_P|=e_P.
\]
If \(e_P=r\), the orbit--stabilizer formula now gives
\[
\#f^{-1}(f(P))=\#(G\mathbin{\cdot}P)=\frac nr.
\]
The indices at all points in this orbit are equal because their local
extensions are conjugate.

Let \(Q_1,\ldots,Q_s\) be the branch points and choose
\(P_i\in f^{-1}(Q_i)\), with \(r_i=e_{P_i}\).  The fiber over \(Q_i\)
contains \(n/r_i\) points, each contributing \(r_i-1\) to the
ramification divisor.  Consequently
\[
\deg R
=\sum_{i=1}^s\frac n{r_i}(r_i-1)
=n\sum_{i=1}^s\left(1-\frac1{r_i}\right).
\]
Substitution in Hurwitz's formula and division by \(n\) gives
\[
\frac{2g-2}{n}
=2g(Y)-2+\sum_{i=1}^s\left(1-\frac1{r_i}\right).
\tag{1}
\]

(b) Put
\[
E=2h-2+\sum_{i=1}^s\left(1-\frac1{r_i}\right)>0.
\]
We determine its smallest possible value.  If \(h\geq2\), then
\(E\geq2\).  If \(h=1\), positivity requires \(s\geq1\), and every
summand is at least \(1/2\), so \(E\geq1/2\).

It remains to take \(h=0\), when
\[
E=s-2-\sum_{i=1}^s\frac1{r_i}.
\tag{2}
\]
Positivity forces \(s\geq3\).  If \(s\geq5\), then
\[
E\geq s-2-\frac s2=\frac s2-2\geq\frac12.
\]
If \(s=4\), the tuple \((2,2,2,2)\) gives zero; every other tuple has
at least one entry at least \(3\), and therefore
\[
E\geq2-\left(\frac12+\frac12+\frac12+\frac13\right)
=\frac16.
\]

Finally suppose \(s=3\), and reorder the indices so that
\[
2\leq r_1\leq r_2\leq r_3.
\]
If \(r_1\geq3\), the zero case \((3,3,3)\) is excluded, and the largest
reciprocal sum below \(1\) is attained at \((3,3,4)\).  Hence
\[
E\geq1-\frac13-\frac13-\frac14=\frac1{12}.
\]
If \(r_1=2\), then \(r_2=2\) is impossible because the reciprocal sum
would exceed \(1\).  For \(r_2=3\), positivity forces \(r_3\geq7\), so
\[
E\geq1-\frac12-\frac13-\frac17=\frac1{42}.
\]
For \(r_2\geq4\), the zero tuple \((2,4,4)\) is excluded, and the
largest reciprocal sum compatible with positivity is obtained from
\((2,4,5)\).  Thus
\[
E\geq1-\frac12-\frac14-\frac15=\frac1{20}.
\]
The global minimum is therefore \(1/42\), attained for
\[
(h;s;r_1,r_2,r_3)=(0;3;2,3,7).
\]

The left side of (1) is positive because \(g\geq2\).  Applying the
minimum just proved to its right side yields
\[
\frac{2g-2}{n}\geq\frac1{42}.
\]
Therefore
\[
n\leq42(2g-2)=84(g-1),
\]
which is Hurwitz's bound.
\end{solution}

\begin{exerciseintro}
% record: H-IV-02-006
\exerciseheading{Exercise IV.2.6}
\addcontentsline{toc}{subsection}{Exercise IV.2.6}

\begin{context}
Let \(f:X\to Y\) be a finite morphism of degree \(n\) between
nonsingular projective curves over an algebraically closed field.
For
\[
D=\sum_P m_P P\in\operatorname{Div}X
\]
define its divisor pushforward by
\[
f_*^{\mathrm{div}}D=\sum_Pm_P f(P).
\]
The morphism \(f\) is finite flat, so for every invertible sheaf
\(\mathcal M\) on \(X\), its sheaf-theoretic direct image
\(f_*\mathcal M\) is locally free of rank \(n\).

For a rank-\(r\) locally free sheaf \(\mathcal E\), write
\(\det\mathcal E=\bigwedge^r\mathcal E\).  We use finite duality in the
form
\[
f_*\omega_X
\cong\mathcal Hom_Y(f_*\mathcal O_X,\omega_Y).
\]
In positive characteristic, \(R\) is the different divisor with its actual multiplicities, including wild ramification.
\end{context}

\begin{problem}
(a) Prove that for every divisor \(D\) on \(X\),
\[
\det\bigl(f_*\mathcal O_X(D)\bigr)
\cong
\det(f_*\mathcal O_X)\otimes
\mathcal O_Y\bigl(f_*^{\mathrm{div}}D\bigr).
\]

(b) Deduce that divisor pushforward respects linear equivalence and
therefore induces a homomorphism
\[
f_*:\operatorname{Pic}X\longrightarrow\operatorname{Pic}Y.
\]
Prove that \(f_*f^*\) is multiplication by \(n\) on
\(\operatorname{Pic}Y\).

(c) Prove that
\[
\det(f_*\Omega_X^1)
\cong
\det(f_*\mathcal O_X)^{-1}\otimes\omega_Y^{\otimes n}.
\]

(d) Suppose that \(f\) is separable, let \(R\) be its ramification
divisor, and put \(B=f_*^{\mathrm{div}}R\).  Prove that
\[
\det(f_*\mathcal O_X)^{\otimes2}
\cong\mathcal O_Y(-B).
\]
\end{problem}
\end{exerciseintro}

\begin{solution}
(a) We first record the effect of an elementary modification.  If
\(\mathcal E'\subset\mathcal E\) are rank-\(n\) vector bundles on \(Y\)
and
\[
0\longrightarrow\mathcal E'
\longrightarrow\mathcal E
\longrightarrow k(Q)\longrightarrow0
\tag{1}
\]
is exact, then
\[
\det\mathcal E
\cong\det\mathcal E'\otimes\mathcal O_Y(Q).
\tag{2}
\]
Indeed, over the discrete valuation ring \(\mathcal O_{Y,Q}\), suitable
bases put the inclusion in (1) in diagonal form
\(\operatorname{diag}(t,1,\ldots,1)\), where \(t\) is a uniformizer.
Taking top exterior powers gives (2).

For any divisor \(D\) and closed point \(P\in X\), there is an exact
sequence
\[
0\longrightarrow\mathcal O_X(D)
\longrightarrow\mathcal O_X(D+P)
\longrightarrow k(P)\longrightarrow0.
\tag{3}
\]
Finite direct image is exact.  If \(Q=f(P)\), then
\(f_*k(P)=k(Q)\), so (2) applied to the direct image of (3) gives
\[
\det f_*\mathcal O_X(D+P)
\cong
\det f_*\mathcal O_X(D)\otimes\mathcal O_Y(Q).
\tag{4}
\]
Starting with \(D=0\), repeated use of (4), in both the positive and
negative directions, proves
\[
\det f_*\mathcal O_X(D)
\cong
\det(f_*\mathcal O_X)\otimes
\mathcal O_Y\bigl(f_*^{\mathrm{div}}D\bigr)
\tag{5}
\]
for every divisor \(D\).

(b) If \(D\sim D'\), then
\(\mathcal O_X(D)\cong\mathcal O_X(D')\).  Their direct images and
determinants are therefore isomorphic.  Formula (5), after cancelling
\(\det(f_*\mathcal O_X)\), shows that
\[
\mathcal O_Y(f_*^{\mathrm{div}}D)
\cong\mathcal O_Y(f_*^{\mathrm{div}}D').
\]
Thus \(f_*^{\mathrm{div}}D\sim f_*^{\mathrm{div}}D'\), and divisor
pushforward induces the asserted map on Picard groups.

For a closed point \(Q\in Y\), one has
\[
f^*Q=\sum_{P\mapsto Q}e_P P.
\]
Because the residue fields are all \(k\), the degree formula for the
fiber says \(\sum_{P\mapsto Q}e_P=n\).  Hence
\[
f_*^{\mathrm{div}}f^*Q
=\sum_{P\mapsto Q}e_PQ
=nQ.
\]
Linearity gives \(f_*f^*[D]=n[D]\) for every
\([D]\in\operatorname{Pic}Y\).

(c) Put \(\mathcal E=f_*\mathcal O_X\).  Since \(X\) and \(Y\) are
nonsingular curves, \(\Omega_X^1=\omega_X\) and finite duality gives
\[
f_*\Omega_X^1
\cong\mathcal Hom_Y(\mathcal E,\omega_Y)
\cong\mathcal E^\vee\otimes\omega_Y.
\]
For a rank-\(n\) vector bundle \(\mathcal E\) and a line bundle
\(\mathcal L\),
\[
\det(\mathcal E^\vee\otimes\mathcal L)
\cong(\det\mathcal E)^{-1}\otimes\mathcal L^{\otimes n}.
\]
Taking \(\mathcal L=\omega_Y\) proves
\[
\det(f_*\Omega_X^1)
\cong
\det(f_*\mathcal O_X)^{-1}\otimes\omega_Y^{\otimes n}.
\tag{6}
\]

(d) Separability gives the canonical-sheaf form of the ramification
formula,
\[
\omega_X\cong f^*\omega_Y\otimes\mathcal O_X(R).
\]
The projection formula and then part (a) yield
\[
\begin{aligned}
\det(f_*\omega_X)
&\cong
\det\bigl(\omega_Y\otimes f_*\mathcal O_X(R)\bigr)\\
&\cong
\omega_Y^{\otimes n}\otimes
\det(f_*\mathcal O_X)\otimes\mathcal O_Y(B).
\end{aligned}
\tag{7}
\]
On the other hand, part (c) says
\[
\det(f_*\omega_X)
\cong
\det(f_*\mathcal O_X)^{-1}\otimes\omega_Y^{\otimes n}.
\tag{8}
\]
Comparing (7) and (8), and cancelling
\(\omega_Y^{\otimes n}\), gives
\[
\det(f_*\mathcal O_X)^{\otimes2}
\cong\mathcal O_Y(-B),
\]
as required.
\end{solution}

\begin{exerciseintro}
% record: H-IV-02-007
\exerciseheading{Exercise IV.2.7}
\addcontentsline{toc}{subsection}{Exercise IV.2.7}

\begin{context}
Let \(Y\) be a nonsingular projective curve over an algebraically
closed field \(k\) of characteristic different from \(2\).  Covers are
considered up to isomorphism over \(Y\), and line bundles up to
isomorphism.

For a finite locally free morphism \(f:X\to Y\), write
\(\mathcal A=f_*\mathcal O_X\).  Then
\[
X\cong\operatorname{Spec}_Y\mathcal A.
\]
If \(f\) is finite étale of degree \(2\), its branch divisor is zero,
so Exercise IV.2.6 gives
\[
\det(\mathcal A)^{\otimes2}\cong\mathcal O_Y.
\]
\end{context}

\begin{problem}
(a) Let \(f:X\to Y\) be a finite étale morphism of degree \(2\).
For the natural inclusion
\[
\mathcal O_Y\longrightarrow f_*\mathcal O_X,
\]
prove that the cokernel \(\mathcal L\) is invertible, that
\(\mathcal L\cong\det(f_*\mathcal O_X)\), and hence that
\(\mathcal L^{\otimes2}\cong\mathcal O_Y\).

(b) Conversely, given an invertible sheaf \(\mathcal L\) and an
isomorphism
\(\phi:\mathcal L\otimes\mathcal L\to\mathcal O_Y\), give
\(\mathcal O_Y\oplus\mathcal L\) the multiplication
\[
(a,b)(a',b')
=\bigl(aa'+\phi(b\otimes b'),\,ab'+a'b\bigr).
\]
Prove that its relative spectrum is a finite étale cover of \(Y\) of
degree \(2\).

(c) Prove that the constructions in (a) and (b) are inverse.  Conclude
that degree-two finite étale covers of \(Y\) correspond bijectively to
the \(2\)-torsion elements of \(\operatorname{Pic}Y\).
\end{problem}
\end{exerciseintro}

\begin{solution}
(a) Put \(\mathcal A=f_*\mathcal O_X\).  This is a rank-two locally
free \(\mathcal O_Y\)-algebra.  Its trace map satisfies
\[
\operatorname{Tr}_{\mathcal A/\mathcal O_Y}(1)=2.
\]
Because \(2\) is invertible, the map
\(\frac12\operatorname{Tr}:\mathcal A\to\mathcal O_Y\) splits the
inclusion of \(\mathcal O_Y\).  Thus its cokernel \(\mathcal L\) is
locally free of rank one, and the exact sequence
\[
0\longrightarrow\mathcal O_Y
\longrightarrow\mathcal A
\longrightarrow\mathcal L
\longrightarrow0
\tag{1}
\]
gives
\[
\det\mathcal A
\cong\det\mathcal O_Y\otimes\det\mathcal L
\cong\mathcal L.
\]
Since \(f\) is étale, its ramification and branch divisors vanish.
Exercise IV.2.6(d) therefore gives
\[
\mathcal L^{\otimes2}
\cong\det(\mathcal A)^{\otimes2}
\cong\mathcal O_Y.
\]

(b) Let
\[
\mathcal A_\phi=\mathcal O_Y\oplus\mathcal L
\]
with the displayed multiplication.  It is a commutative
\(\mathcal O_Y\)-algebra, locally free of rank two, so
\[
X_\phi=\operatorname{Spec}_Y\mathcal A_\phi\longrightarrow Y
\]
is finite flat of degree two.

To check that it is étale, work on an open set on which
\(\mathcal L\) has a generator \(e\).  Write
\(\phi(e\otimes e)=u\), where \(u\) is a unit.  Then
\[
\mathcal A_\phi\cong
\mathcal O_Y[T]/(T^2-u).
\]
In this algebra \(T\) is a unit, and \(2\) is a unit by hypothesis.
Thus the derivative \(2T\) is invertible, so
\[
\Omega_{\mathcal A_\phi/\mathcal O_Y}=0.
\]
The finite flat morphism is therefore étale.

The cover depends only on the isomorphism class of the \(2\)-torsion
line bundle, not on the chosen \(\phi\).  Indeed, two trivializations
of \(\mathcal L^{\otimes2}\) differ by a scalar
\(c\in k^\times\).  Choose \(\lambda\in k^\times\) with
\(c\lambda^2=1\).  Multiplication by \(\lambda\) on the
\(\mathcal L\)-summand gives an isomorphism between the resulting
algebras.

(c) Starting from \((\mathcal L,\phi)\), the natural inclusion
\[
\mathcal O_Y\longrightarrow
\mathcal O_Y\oplus\mathcal L
\]
has cokernel \(\mathcal L\).  Hence the construction in (a) immediately
recovers the original line bundle.

Conversely, begin with a degree-two étale cover and put
\(\mathcal A=f_*\mathcal O_X\).  The splitting in (a) identifies the
cokernel in (1) with the trace-zero line bundle
\[
\mathcal M=\ker(\operatorname{Tr}:\mathcal A\to\mathcal O_Y),
\qquad
\mathcal A=\mathcal O_Y\oplus\mathcal M.
\tag{2}
\]
There is a canonical involution
\[
\tau:\mathcal A\longrightarrow\mathcal A,\qquad
\tau(a)=\operatorname{Tr}(a)-a.
\tag{3}
\]
Étale locally on \(Y\), the algebra is
\(\mathcal O_Y\times\mathcal O_Y\), and \(\tau\) interchanges the two
factors.  It follows by descent that \(\tau\) is an algebra
automorphism, that its invariant subalgebra is \(\mathcal O_Y\), and
that it acts as \(-1\) on \(\mathcal M\).

Consequently multiplication in \(\mathcal A\) sends
\(\mathcal M\otimes\mathcal M\) to the invariant summand
\(\mathcal O_Y\).  This pairing is an isomorphism: after the same
étale localization, \(\mathcal M\) consists of the pairs \((a,-a)\),
and
\[
(a,-a)(b,-b)=(ab,ab).
\]
Thus (2), including its algebra structure, is exactly the algebra
constructed in part (b) from
\[
\mathcal M^{\otimes2}\xrightarrow{\ \sim\ }\mathcal O_Y.
\]
Taking relative spectra recovers \(X\to Y\).  The two constructions
are therefore inverse and give the asserted bijection.

The trivial class \(\mathcal L=\mathcal O_Y\) corresponds to the split
cover \(Y\bigsqcup Y\to Y\); if only connected covers are admitted, it is
the unique class to omit.
\end{solution}

\begin{exerciseintro}
\corpussection{IV.3 Embeddings in Projective Space}
\addcontentsline{toc}{section}{IV.3 Embeddings in Projective Space}

% record: H-IV-03-001
\exerciseheading{Exercise IV.3.1}
\addcontentsline{toc}{subsection}{Exercise IV.3.1}

\begin{context}
Let \(X\) be a nonsingular projective integral curve of genus \(2\)
over an algebraically closed field.  We use
Riemann--Roch and the facts that a divisor of degree at least
\(2g+1\) is very ample, and that a nonsingular plane curve of degree
\(m\) has genus
\[
\frac{(m-1)(m-2)}2.
\]
A very ample complete linear system on a positive-genus curve has at
least three independent sections.
For a line bundle of nonnegative degree \(d\), a nonzero section
identifies it with an effective divisor; removing its \(d\) points
one at a time gives \(h^0\leq d+1\).  A line bundle of negative degree
has no nonzero sections.
\end{context}

\begin{problem}
Prove that a divisor \(D\) on \(X\) is very ample if and only if
\[
\deg D\geq5.
\]
\end{problem}
\end{exerciseintro}

\begin{solution}
If \(\deg D\geq5=2g+1\), the standard very-ampleness theorem for
curves shows that \(D\) is very ample.

Conversely, suppose that \(D\) is very ample, and put \(d=\deg D\).
Then \(d>0\), and the complete linear system embeds \(X\) in
\(\mathbb P^{h^0(D)-1}\).  Since a positive-genus curve cannot embed
in \(\mathbb P^0\) or \(\mathbb P^1\), one must have
\[
h^0(D)\geq3.
\tag{1}
\]
For \(d\leq1\), the elementary estimate
\(h^0(D)\leq d+1\) contradicts (1).  If \(d=2\), Riemann--Roch gives
\[
h^0(D)=1+h^0(K_X-D)\leq2,
\]
because \(K_X-D\) has degree zero and a degree-zero line bundle has at
most one independent section.  Thus \(d\geq3\).

If \(d=3\), then \(\deg(K_X-D)=-1\), so Riemann--Roch gives
\[
h^0(D)=d+1-g=2,
\]
again contradicting (1).  If \(d=4\), the same calculation gives
\[
h^0(D)=3.
\]
Very ampleness would therefore identify \(X\) with a nonsingular
plane curve of degree \(4\).  Such a curve has genus
\[
\frac{(4-1)(4-2)}2=3,
\]
contrary to \(g(X)=2\).  Hence \(d\geq5\), proving the converse.
\end{solution}

\begin{exerciseintro}
% record: H-IV-03-002
\exerciseheading{Exercise IV.3.2}
\addcontentsline{toc}{subsection}{Exercise IV.3.2}

\begin{context}
Let \(X\subset\mathbb P^2\) be a nonsingular plane curve of degree
\(4\), over an algebraically closed field.  We use Riemann--Roch.
Then \(g(X)=3\), and adjunction gives
\[
\omega_X\cong\mathcal O_X(1).
\]
We use \(H^1(\mathbb P^2,\mathcal O_{\mathbb P^2}(m))=0\) for every
integer \(m\).  A curve is hyperelliptic precisely when it has a
degree-two divisor \(D\) with \(\dim|D|=1\).
\end{context}

\begin{problem}
(a) Prove that the effective canonical divisors on \(X\) are exactly
the divisors \(X\mathbin{\cdot}L\), where \(L\subset\mathbb P^2\) is
a line.

(b) If \(D\) is any effective divisor of degree \(2\) on \(X\), prove
that
\[
\dim|D|=0.
\]

(c) Deduce that \(X\) is not hyperelliptic.
\end{problem}
\end{exerciseintro}

\begin{solution}
(a) Restricting linear forms to \(X\) gives the exact sequence
\[
0\longrightarrow
\mathcal O_{\mathbb P^2}(-3)
\longrightarrow
\mathcal O_{\mathbb P^2}(1)
\longrightarrow
\mathcal O_X(1)
\longrightarrow0.
\]
Since
\[
H^0(\mathbb P^2,\mathcal O_{\mathbb P^2}(-3))=0
\quad\text{and}\quad
H^1(\mathbb P^2,\mathcal O_{\mathbb P^2}(-3))=0,
\]
the restriction map
\[
H^0(\mathbb P^2,\mathcal O_{\mathbb P^2}(1))
\xrightarrow{\ \sim\ }
H^0(X,\mathcal O_X(1))
=H^0(X,\omega_X)
\]
is an isomorphism.  The zero divisor of a nonzero restricted linear
form is the scheme-theoretic line section \(X\mathbin{\cdot}L\), and
every nonzero canonical section arises this way.

(b) Riemann--Roch, with \(g(X)=3\) and \(\deg D=2\), gives
\[
h^0(D)=h^0(K_X-D).
\tag{1}
\]
Suppose that \(h^0(D)\geq2\).  By (1), choose an effective divisor
\[
E\in|K_X-D|.
\]
It has degree two.  For every \(D'\in|D|\), the divisor
\[
D'+E
\]
is effective and canonical, so part (a) writes it as a line section.

There is at most one line whose intersection divisor with \(X\)
contains \(E\).  Indeed, if \(E=P+Q\) with \(P\neq Q\), that line must
be the unique line through \(P\) and \(Q\); if \(E=2P\), it must be
the unique tangent line \(T_PX\).  Thus every \(D'\in|D|\) gives the
same line section as \(D\), and
\[
D'+E=D+E.
\]
Cancellation of effective divisors gives \(D'=D\).  This contradicts
\(h^0(D)\geq2\), which would make \(|D|\) positive-dimensional.
Therefore \(h^0(D)=1\), or equivalently \(\dim|D|=0\).

(c) A hyperelliptic curve possesses a \(g^1_2\), that is, an effective
degree-two divisor \(D\) with \(\dim|D|=1\).  Part (b) rules out every
such divisor.  Hence the plane quartic \(X\) is not hyperelliptic.
\end{solution}

\begin{exerciseintro}
% record: H-IV-03-003
\exerciseheading{Exercise IV.3.3}
\addcontentsline{toc}{subsection}{Exercise IV.3.3}

\begin{context}
Work over an algebraically closed field.  Let \(X\subset\mathbb P^n\)
be a nonsingular integral curve, cut out scheme-theoretically by
hypersurfaces of degrees \(d_1,\ldots,d_{n-1}\).
We use the canonical degree and dimension formulas
\[
\deg\omega_X=2g(X)-2,\qquad h^0(X,\omega_X)=g(X).
\]
The adjunction formula gives
\[
\omega_X
\cong
\mathcal O_X\left(\sum_{i=1}^{n-1}d_i-n-1\right).
\]
The line bundle \(\mathcal O_X(1)\) is very ample, and every positive
tensor power of a very ample line bundle is very ample.
\end{context}

\begin{problem}
If \(g(X)\geq2\), prove that the canonical divisor \(K_X\) is very
ample.  Deduce that a curve of genus \(2\) cannot be a complete
intersection in any projective space.
\end{problem}
\end{exerciseintro}

\begin{solution}
Set
\[
m=\sum_{i=1}^{n-1}d_i-n-1.
\]
The adjunction formula says
\[
\omega_X\cong\mathcal O_X(m).
\tag{1}
\]
Taking degrees in (1) gives
\[
2g(X)-2=m\deg\mathcal O_X(1).
\]
The left side is positive because \(g(X)\geq2\), and
\(\deg\mathcal O_X(1)>0\).  Hence the integer \(m\) is positive.
Thus \(\omega_X\) is a positive tensor power of the very ample line
bundle \(\mathcal O_X(1)\), and is itself very ample.

Now suppose that \(g(X)=2\).  The canonical space has dimension
\[
h^0(X,\omega_X)=g(X)=2,
\]
so the canonical map has target \(\mathbb P^1\).  A genus-two curve
cannot be embedded in \(\mathbb P^1\): a closed integral
one-dimensional subscheme of \(\mathbb P^1\) is \(\mathbb P^1\)
itself, which has genus zero.  Hence \(K_X\) is not very ample.  This
contradicts the first paragraph if \(X\) is a complete intersection,
so no genus-two curve is a complete intersection in any
\(\mathbb P^n\).
\end{solution}

\begin{exerciseintro}
% record: H-IV-03-004
\exerciseheading{Exercise IV.3.4}
\addcontentsline{toc}{subsection}{Exercise IV.3.4}

\begin{context}
Work over an algebraically closed field; curves below are nonsingular,
projective, and integral.
Let
\[
\nu_d:\mathbb P^1\longrightarrow\mathbb P^d,\qquad
[s:t]\longmapsto
[s^d:s^{d-1}t:\cdots:st^{d-1}:t^d]
\]
be the \(d\)-uple embedding.  Its image is the rational normal curve
of degree \(d\).

We use the criterion that a normal projective variety is projectively
normal when all restriction maps
\[
H^0(\mathbb P^N,\mathcal O_{\mathbb P^N}(m))
\longrightarrow H^0(X,\mathcal O_X(m))
\]
are surjective.  For an effective divisor \(D\) of degree \(d\) on a
curve, successively removing its points gives \(h^0(D)\leq d+1\).
\end{context}

\begin{problem}
(a) Prove that the rational normal curve is projectively normal and
that its homogeneous ideal is generated by quadrics.

(b) Let \(X\subset\mathbb P^n\) be a curve of degree \(d\leq n\) that
is not contained in a hyperplane.  Prove that
\[
d=n,\qquad g(X)=0,
\]
and that, up to an automorphism of \(\mathbb P^d\), \(X\) is the
rational normal curve.

(c) Deduce that every curve of degree \(2\) in projective space is a
conic in some \(\mathbb P^2\).

(d) Prove that every curve of degree \(3\) in projective space is
either a plane cubic or the twisted cubic in \(\mathbb P^3\).
\end{problem}
\end{exerciseintro}

\begin{solution}
(a) Let \(S=k[X_0,\ldots,X_d]\), and consider the graded map
\[
\theta:S\longrightarrow
k[s,t]^{(d)}
=\bigoplus_{m\geq0}k[s,t]_{md},
\qquad
X_i\longmapsto s^{d-i}t^i.
\tag{1}
\]
Every monomial \(s^{md-q}t^q\), \(0\leq q\leq md\), is a product of
\(m\) monomials of degree \(d\).  Hence \(\theta\) is surjective in
every degree.  Equivalently,
\[
\operatorname{Sym}^m H^0(\mathbb P^1,\mathcal O_{\mathbb P^1}(d))
\longrightarrow
H^0(\mathbb P^1,\mathcal O_{\mathbb P^1}(md))
\]
is surjective for every \(m\).  Since \(\mathbb P^1\) is normal, this
proves projective normality.  Alternatively, the target of (1) is the
\(d\)-th Veronese subring of \(k[s,t]\), which is integrally closed.

It remains to identify the kernel.  Let \(J\) be generated by the
quadrics
\[
X_iX_j-X_{i+1}X_{j-1},
\qquad 0\leq i<j-1\leq d-1.
\tag{2}
\]
These are the \(2\times2\) minors of
\[
\begin{pmatrix}
X_0&X_1&\cdots&X_{d-1}\\
X_1&X_2&\cdots&X_d
\end{pmatrix},
\]
and their images under \(\theta\) vanish, so \(J\subseteq\ker\theta\).

Modulo (2), if two indices in a monomial differ by at least two, they
may be moved one step toward each other without changing their sum.
Each move strictly decreases the sum of the squares of the indices,
so the process terminates with all indices differing by at most one.
It follows that every degree-\(m\) monomial, for \(m\geq1\), is congruent to
\[
X_a^{\,m-r}X_{a+1}^{\,r},
\qquad q=ma+r,\quad 0\leq r<m,
\tag{3}
\]
where \(q\) is the sum of its indices; for \(q=md\), (3) means
\(X_d^m\).  Thus \((S/J)_m\) is spanned by at most \(md+1\) elements,
one for each \(q=0,\ldots,md\).  But the induced map
\[
(S/J)_m\longrightarrow k[s,t]_{md}
\]
is surjective, and its target has dimension \(md+1\).  It is therefore
an isomorphism for every \(m\geq1\); in degree zero it is the identity
of \(k\).  Hence \(J=\ker\theta\), proving that
the homogeneous ideal is generated by quadrics.

(b) Let \(H\) be a hyperplane-section divisor on \(X\).  Since \(X\)
is nondegenerate, the \(n+1\) coordinate linear forms restrict to
linearly independent sections, and therefore
\[
n+1\leq h^0(X,\mathcal O_X(H)).
\tag{4}
\]
Choose an effective divisor in \(|H|\), of degree \(d\).  The
point-removal estimate and the hypothesis \(d\leq n\) give
\[
n+1\leq h^0(H)\leq d+1\leq n+1.
\tag{5}
\]
All inequalities are equalities.  In particular,
\[
d=n,\qquad h^0(H)=d+1.
\tag{6}
\]

Write an effective representative as
\[
H=P_1+\cdots+P_d
\]
with repetitions allowed.  At each step, adding one point can increase
\(h^0\) by at most one.  Since the total increase from
\(h^0(\mathcal O_X)=1\) to \(h^0(H)=d+1\) is exactly \(d\), equality
holds at every step.  In particular,
\[
h^0(P_1)=2.
\]
The pencil \(|P_1|\) has no base point: otherwise
\(h^0(P_1)=h^0(\mathcal O_X)=1\).  It therefore defines a degree-one
morphism \(X\to\mathbb P^1\), which is an isomorphism.  Hence
\(g(X)=0\).

Under \(X\cong\mathbb P^1\), the line bundle
\(\mathcal O_X(H)\) has degree \(d\), so it is
\(\mathcal O_{\mathbb P^1}(d)\).  Equations (4)--(6) show that the
given embedding uses its complete space of sections.  Changing the
basis of that space only applies an automorphism of \(\mathbb P^d\).
Thus the embedding is the rational normal curve up to projective
automorphism.

(c) Let \(\mathbb P^r\) be the linear span of a degree-two curve.
Then the curve is nondegenerate in \(\mathbb P^r\).  It cannot have
\(r=1\), since the only integral curve in \(\mathbb P^1\) is
\(\mathbb P^1\) itself, of degree one.  If \(r\geq2\), part (b), with
\(2\leq r\), gives \(r=2\) and identifies the curve with the rational
normal conic.

(d) Let \(\mathbb P^r\) be the linear span of a degree-three curve.
Again \(r\neq1\).  If \(r=2\), the curve is a plane cubic.  If
\(r\geq3\), part (b) gives \(r=3\) and identifies it with the rational
normal curve of degree three, namely the twisted cubic.
\end{solution}

\begin{exerciseintro}
% record: H-IV-03-005
\exerciseheading{Exercise IV.3.5}
\addcontentsline{toc}{subsection}{Exercise IV.3.5}

\begin{context}
Let \(X\subset\mathbb P^3\) be a nonsingular curve not contained in a
plane, over an algebraically closed field, and let \(O\notin X\).
The generic-projection theorem supplies a nonempty open set of centers
outside \(X\) inducing birational morphisms to plane curves.
We use the vanishing of \(H^1(\mathbb P^2,\mathcal O(m))\) for all
integers \(m\), and the fact that finite birational morphisms of
nonsingular projective curves are isomorphisms.
Projection from \(O\) is defined by the
three-dimensional space of hyperplanes through \(O\), and its
associated line bundle on \(X\) is \(\mathcal O_X(1)\).

For an integral plane curve \(C\) of degree \(d\), its arithmetic
genus is
\[
p_a(C)=\frac{(d-1)(d-2)}2.
\]
If \(\widetilde C\to C\) is its normalization, then
\[
p_a(C)=g(\widetilde C)+
\sum_{Q\in C}\dim_k
\bigl((\nu_*\mathcal O_{\widetilde C}/\mathcal O_C)_Q\bigr).
\]
We also use constancy of the Hilbert polynomial in a flat projective
family.
\end{context}

\begin{problem}
(a) Suppose projection from \(O\) induces a birational morphism
\(\varphi:X\to C=\varphi(X)\subset\mathbb P^2\).  Prove that \(C\)
is singular.

(b) If \(X\) has degree \(d\) and genus \(g\), deduce that
\[
g<\frac{(d-1)(d-2)}2.
\]

(c) Let \(\{X_t\}\) be the flat family induced by the projection, with
\(X_1=X\) and with \(X_0\) supported on \(C\).  Prove that \(X_0\)
has nilpotent elements.
\end{problem}
\end{exerciseintro}

\begin{solution}
(a) Suppose, to the contrary, that \(C\) is nonsingular.  A birational
morphism between nonsingular projective curves is an isomorphism, so
\[
X\cong C
\quad\text{and}\quad
\mathcal O_X(1)\cong\varphi^*\mathcal O_C(1).
\tag{1}
\]
The restriction sequence for a plane curve shows that
\[
h^0(C,\mathcal O_C(1))=3
\tag{2}
\]
when \(\deg C\geq2\); the case that \(C\) is a line would give only
two sections and leads to the same contradiction.  From (1) and (2),
\[
h^0(X,\mathcal O_X(1))\leq3.
\]
On the other hand, because \(X\) is not contained in a plane, the four
coordinate linear forms on \(\mathbb P^3\) restrict to linearly
independent sections of \(\mathcal O_X(1)\).  Thus
\[
h^0(X,\mathcal O_X(1))\geq4,
\]
a contradiction.  Therefore \(C\) is singular.

(b) Choose a center giving a birational projection, as provided by the
generic-projection theorem.  The morphism \(\varphi:X\to C\) is finite
and birational, so it is
the normalization of \(C\).  Moreover,
\[
\varphi^*\mathcal O_C(1)\cong\mathcal O_X(1),
\]
and hence \(\deg C=d\).  The normalization-genus formula gives
\[
\frac{(d-1)(d-2)}2
=g+\sum_{Q\in C}\delta_Q.
\tag{3}
\]
Part (a) says that at least one \(\delta_Q\) is positive.  Equation
(3) therefore yields the strict inequality
\[
g<\frac{(d-1)(d-2)}2.
\]

(c) The Hilbert polynomial of a projective curve of degree \(d\) and
arithmetic genus \(p_a\) is
\[
P(m)=dm+1-p_a.
\]
Flatness of the family gives
\[
P_{X_0}(m)=P_X(m)=dm+1-g.
\tag{4}
\]
If \(X_0\) were reduced, then, because its support is the integral
curve \(C\), it would equal \(C\) with its reduced scheme structure.
Its Hilbert polynomial would then be
\[
P_C(m)
=dm+1-\frac{(d-1)(d-2)}2.
\tag{5}
\]
The constant terms in (4) and (5) differ by part (b), contradicting
flatness.  Hence \(X_0\) is nonreduced, or equivalently its structure
sheaf contains nonzero nilpotent elements.
\end{solution}

\begin{exerciseintro}
% record: H-IV-03-006
\exerciseheading{Exercise IV.3.6}
\addcontentsline{toc}{subsection}{Exercise IV.3.6}

\begin{context}
Let \(X\) be a nonsingular projective curve of degree \(4\), and let
\(\mathbb P^r\) be its linear span, over an algebraically closed field.
Thus \(X\) is nondegenerate in
\(\mathbb P^r\).

We use the preceding results that a nondegenerate curve of degree
\(d\leq r\) in \(\mathbb P^r\) is the rational normal curve with
\(d=r\), and that a nonplanar curve of degree \(d\) in
\(\mathbb P^3\) has genus strictly less than
\((d-1)(d-2)/2\).  A genus-two curve has no very ample divisor of
degree four.
We also use Riemann--Roch, the plane-curve genus formula, and the
following complete-intersection facts: two surfaces without a common
component meet in degree equal to the product of their degrees, and
their intersection is Cohen--Macaulay, with no embedded associated
points.
The existence examples are the rational normal quartic in
\(\mathbb P^4\), the smooth rational quartic in \(\mathbb P^3\),
a genus-one curve embedded by a degree-four line bundle, and a
smooth plane quartic.
\end{context}

\begin{problem}
(a) Prove that exactly the following possibilities can occur:

\[
\begin{array}{c|c|c}
g(X)&\text{linear span}&\text{description}\\
0&\mathbb P^4&\text{rational normal quartic}\\
0&\mathbb P^3&\text{rational quartic}\\
1&\mathbb P^3&\text{elliptic quartic}\\
3&\mathbb P^2&\text{nonsingular plane quartic}.
\end{array}
\]

(b) In the genus-one case, prove that \(X\subset\mathbb P^3\) is a
complete intersection of two irreducible quadric surfaces.
\end{problem}
\end{exerciseintro}

\begin{solution}
(a) If \(r\geq4\), then \(\deg X=4\leq r\).  The rational-normal-curve
result gives
\[
r=4,\qquad g(X)=0,
\]
and \(X\) is the rational normal quartic in \(\mathbb P^4\).

Suppose \(r=3\).  The genus bound for nonplanar degree-four curves
gives
\[
g(X)<\frac{(4-1)(4-2)}2=3,
\]
so \(g(X)\in\{0,1,2\}\).  The line bundle
\(\mathcal O_X(1)\) is very ample of degree four, which is impossible
when \(g(X)=2\).  Hence \(g(X)=0\) or \(1\).  In genus zero,
\(X\cong\mathbb P^1\) and
\(\mathcal O_X(1)\cong\mathcal O_{\mathbb P^1}(4)\); the four
coordinate sections give the rational quartic in \(\mathbb P^3\).
The remaining possibility is the elliptic quartic.

If \(r=2\), then \(X\) is a nonsingular plane quartic and therefore
\[
g(X)=\frac{(4-1)(4-2)}2=3.
\]
Finally, \(r=1\) is impossible: the only integral
one-dimensional closed subscheme of \(\mathbb P^1\) with its reduced
structure is \(\mathbb P^1\), which has degree one.  This exhausts the
possibilities.

(b) Now assume \(g(X)=1\), so \(X\subset\mathbb P^3\) is
nondegenerate.  Twist its ideal sequence by \(2\):
\[
0\longrightarrow\mathcal I_X(2)
\longrightarrow\mathcal O_{\mathbb P^3}(2)
\longrightarrow\mathcal O_X(2)
\longrightarrow0.
\tag{1}
\]
There are ten quadratic forms on \(\mathbb P^3\).  The line bundle
\(\mathcal O_X(2)\) has degree \(8\), and Riemann--Roch on the
genus-one curve gives
\[
h^0(X,\mathcal O_X(2))=8.
\]
Taking global sections in (1), without needing surjectivity of the
restriction map, therefore yields
\[
h^0(\mathbb P^3,\mathcal I_X(2))
\geq10-8=2.
\tag{2}
\]
Choose linearly independent quadrics \(Q_1,Q_2\) containing \(X\).

Each \(Q_i\) is irreducible.  Indeed, over an algebraically closed
field a reducible quadric is a union of two planes, possibly the same
plane twice.  Since \(X\) is integral, containment in such a union
would put \(X\) in one plane, contrary to nondegeneracy.

The irreducible quadrics \(Q_1\) and \(Q_2\) are not associates, so
they have no common surface component.  Their scheme-theoretic
intersection is consequently a complete-intersection curve of degree
\[
\deg Q_1\deg Q_2=4.
\]
It contains the integral degree-four curve \(X\).  A complete
intersection has no embedded components and is pure of dimension one;
equality of the degrees leaves no residual component or
multiplicity.  More explicitly, if \(Y=Q_1\cap Q_2\), then the
kernel of \(\mathcal O_Y\to\mathcal O_X\) is zero at the generic
point and is supported at finitely many closed points.  A nonzero such
submodule would have a closed associated point of \(\mathcal O_Y\),
contrary to the absence of embedded associated points.  Hence
\[
X=Q_1\cap Q_2
\]
scheme-theoretically, proving the assertion.
\end{solution}

\begin{exerciseintro}
% record: H-IV-03-007
\exerciseheading{Exercise IV.3.7}
\addcontentsline{toc}{subsection}{Exercise IV.3.7}

\begin{context}
Assume that \(k\) is algebraically closed of characteristic different
from \(2\).  A plane quartic has arithmetic genus \(3\), and
normalization at one ordinary node lowers the genus by one.
We use the Jacobian criterion and Bézout's theorem.  Projection here
means projection from a point outside the source curve.

If a nonsingular curve \(X\subset\mathbb P^3\) projects birationally
onto a plane curve \(C\), then \(X\) is the normalization of \(C\) and
\[
\mathcal O_X(1)\cong\nu^*\mathcal O_C(1).
\]
In particular, their degrees agree.  A genus-two curve has no very
ample divisor of degree four.
\end{context}

\begin{problem}
Show that the affine curve
\[
xy+x^4+y^4=0
\]
provides a plane curve with only nodes that is not the birational
projection of any nonsingular curve in \(\mathbb P^3\).
\end{problem}
\end{exerciseintro}

\begin{solution}
Let \(C\subset\mathbb P^2\) be the projective closure
\[
F(x,y,z)=xyz^2+x^4+y^4=0.
\]
Its partial derivatives are
\[
F_x=yz^2+4x^3,\qquad
F_y=xz^2+4y^3,\qquad
F_z=2xyz.
\]
At a singular point, \(F_z=0\), so \(xyz=0\).
If \(z=0\), the first two derivatives force \(x=y=0\), which is not a
projective point.  If \(x=0\) or \(y=0\), the equation \(F=0\) forces
both \(x\) and \(y\) to vanish.  Thus the unique singular point is
\[
P=[0:0:1].
\]
In the affine chart \(z=1\), the lowest-degree part of the local
equation at \(P\) is \(xy\).  It has two distinct tangent directions,
so \(P\) is an ordinary node.

The curve \(C\) is reduced: a repeated irreducible factor of \(F\)
would divide all its partial derivatives and produce a whole curve
of singular points, contradicting the computation above.
It is also irreducible.  Indeed, if it were reducible, distinct
components could meet only at the unique singular point \(P\).
The node has just two transverse branches, whose intersection
multiplicity is one.  But a decomposition of a quartic into components
of positive degrees \(a\) and \(b\), with \(a+b=4\), has
\(ab\geq3\), contradicting Bézout's theorem.  More than two components
are also impossible at a two-branch node.

Let \(\nu:\widetilde C\to C\) be the normalization.  Since a plane
quartic has arithmetic genus three and an ordinary node has
\(\delta\)-invariant one,
\[
g(\widetilde C)=3-1=2.
\tag{1}
\]
Suppose that \(C\) were the birational projection of a nonsingular
curve \(X\subset\mathbb P^3\).  Then \(X\cong\widetilde C\), and the
projection linear system satisfies
\[
\mathcal O_X(1)\cong\nu^*\mathcal O_C(1).
\]
Consequently \(\mathcal O_X(1)\) is a very ample line bundle of degree
four on the genus-two curve \(X\).  The very-ampleness bound in the
context says that every
very ample divisor on a genus-two curve has degree at least five.
This contradiction proves that \(C\) is not such a projection.
\end{solution}

\begin{exerciseintro}
% record: H-IV-03-008
\exerciseheading{Exercise IV.3.8}
\addcontentsline{toc}{subsection}{Exercise IV.3.8}

\begin{context}
Work over an algebraically closed field.
An integral curve \(C\subset\mathbb P^n\) is called strange if some
point \(A\in\mathbb P^n\) lies on every tangent line to \(C\) at its
nonsingular points.

For projection \(\pi_A\) from \(A\), the kernel of the differential at
a point \(P\neq A\) is the tangent direction of the line \(AP\).  In
characteristic zero, a nonconstant morphism between integral curves
is separable and has nonzero differential at the generic point.
\end{context}

\begin{problem}
(a) Over a field of characteristic \(p>0\), prove that the projective
closure of the parametrized affine curve
\[
x=t,\qquad y=t^p,\qquad z=t^{2p}
\]
is a singular strange curve.

(b) Prove that in characteristic zero every strange integral curve is
a line, hence isomorphic to \(\mathbb P^1\).  In particular, there are
no singular strange curves.
\end{problem}
\end{exerciseintro}

\begin{solution}
(a) Use homogeneous coordinates \([W:X:Y:Z]\) on \(\mathbb P^3\).
The projective closure \(C\) is parametrized by
\[
\nu:\mathbb P^1\longrightarrow\mathbb P^3,\qquad
[s:u]\longmapsto
[s^{2p}:us^{2p-1}:u^ps^p:u^{2p}].
\tag{1}
\]
On the affine chart \(s\neq0\), this is
\[
t\longmapsto[1:t:t^p:t^{2p}].
\]
Because \(t=x\), the parametrization is birational onto its integral
image, and this affine part is nonsingular.  The morphism is
nonconstant and proper with finite fibers, hence finite; since
\(\mathbb P^1\) is normal, it is the normalization of \(C\).

Formal differentiation in characteristic \(p\) gives
\[
\frac{d}{dt}(t,t^p,t^{2p})=(1,0,0).
\]
Thus the tangent line at every affine point has the \(x\)-direction.
After projective closure, all these lines pass through
\[
A=[0:1:0:0].
\]
The affine points are all the nonsingular points except possibly the
single point at infinity, which we now show is singular.

It remains to see what happens at infinity.  Formula (1) sends
\([0:1]\) to \(Q=[0:0:0:1]\).  On the chart \(u=1\), with local
parameter \(q=s/u\) on the normalization and \(Z=1\) on the target,
the three local coordinate functions at \(Q\) are
\[
W=q^{2p},\qquad X=q^{2p-1},\qquad Y=q^p.
\]
Their least order is \(p>1\).  Hence the branch has multiplicity
\(p\), so \(Q\) is singular.  This proves that \(C\) is a singular
strange curve.

(b) Let \(C\subset\mathbb P^n\) be strange with strange point \(A\),
and let \(\widetilde C\to C\) be its normalization.  Projection from
\(A\) gives a rational map
\[
C\dashrightarrow\mathbb P^{n-1},
\]
which induces a morphism from \(\widetilde C\) because the latter is a
nonsingular projective curve.

At every nonsingular \(P\in C\) with \(P\neq A\), the tangent line
\(T_PC\) is the line \(AP\).  It is precisely the direction killed by
the differential of projection.  The induced morphism on
\(\widetilde C\) therefore has zero differential on a dense open
set, hence at the generic point.  In characteristic zero a
nonconstant morphism of curves is separable and cannot have
identically zero differential.  The projected map is consequently
constant.

Thus all points of \(C\) lie on one fixed line through \(A\).  Since
\(C\) is an integral closed curve, it equals that line.  Therefore
\(C\cong\mathbb P^1\), and in particular no singular strange curve
exists in characteristic zero.
\end{solution}

\begin{exerciseintro}
% record: H-IV-03-009
\exerciseheading{Exercise IV.3.9}
\addcontentsline{toc}{subsection}{Exercise IV.3.9}

\begin{context}
Let \(X\subset\mathbb P^3\) be a nonsingular curve of degree \(d\)
that is not contained in a plane, over an algebraically closed field
of arbitrary characteristic.  The family of secant lines is the
image of the two-dimensional parameter space \(X^{(2)}\) in the
Grassmannian of lines.

We use the multisecant and strange-curve results of the
generic-projection theorem (Hartshorne IV, Proposition 3.8 and
Theorem 3.9), in the following form:
if a general secant of a nonplanar curve is a multisecant, then the
curve is strange.  A nonsingular strange curve is a line, or in
characteristic \(2\) a plane conic.  Thus a nonplanar nonsingular
curve is not strange.
\end{context}

\begin{problem}
Prove the following Bertini lemma: for a nonempty open subset of
planes \(H\subset\mathbb P^3\), the intersection \(X\cap H\) consists
of exactly \(d\) distinct points, no three of which are collinear.
\end{problem}
\end{exerciseintro}

\begin{solution}
First consider transversality.  A plane \(H\) fails to meet \(X\)
transversely at \(P\) precisely when it contains the tangent line
\(T_PX\).  For a fixed \(P\), the planes containing \(T_PX\) form a
projective line.  Hence the incidence correspondence
\[
\{(P,H):T_PX\subset H\}
\subset X\times(\mathbb P^3)^*
\]
has dimension two.  Its image is therefore a proper closed subset of
the three-dimensional dual projective space.  A plane outside this
subset cuts \(X\) transversely, and Bézout's theorem then says that
\(X\cap H\) consists of \(d\) distinct points.

It remains to exclude three collinear points.  Let \(\mathcal T\) be
the family of lines meeting \(X\) in at least three distinct points.
We claim that
\[
\dim\mathcal T\leq1.
\tag{1}
\]
Indeed, the full secant family has dimension two.  Its map from
\(X^{(2)}\) has finite general fiber, since a line not contained in
\(X\) meets \(X\) in only finitely many points.  If
\(\dim\mathcal T=2\), a general secant would be a multisecant.  The
multisecant argument would then make \(X\) strange.  This is
impossible for a nonplanar nonsingular curve, proving (1).

Now form the incidence correspondence
\[
\mathcal I=
\{(L,H):L\in\mathcal T,\ L\subset H\}.
\]
For each line \(L\), the planes containing it form a projective line.
Thus (1) gives \(\dim\mathcal I\leq2\).  The closure of its image in
\((\mathbb P^3)^*\) is again a proper closed subset.  A plane \(H\)
outside it contains no trisecant of \(X\).

Choose \(H\) outside both exceptional closed subsets.  It meets \(X\)
in \(d\) distinct points.  If three of them were collinear, their line
would be a trisecant contained in \(H\), contrary to the choice of
\(H\).  These good planes form the required nonempty open subset.
\end{solution}

\begin{exerciseintro}
% record: H-IV-03-010
\exerciseheading{Exercise IV.3.10}
\addcontentsline{toc}{subsection}{Exercise IV.3.10}

\begin{context}
Let \(X\subset\mathbb P^n\) be an integral nondegenerate curve over an
algebraically closed field of characteristic zero.  Points are always
chosen in a suitable dense open subset of the nonsingular locus.

We use the characteristic-zero trisecant lemma in the following precise
form: for an integral projective curve \(C\subset\mathbb P^m\) not
contained in a plane, the closure in the Grassmannian of the lines
meeting \(C\) in three distinct points has dimension at most one.
The curve may be singular and \(m\) may exceed three.
We also use normalization of projective curves, separability in
characteristic zero, and the dimension theorem for fibers.
\end{context}

\begin{problem}
Assume \(n\geq2\).  Prove that for general points
\[
P_1,\ldots,P_{n-1}\in X,
\]
their linear span is an \((n-2)\)-plane and contains no other point of
\(X\).
\end{problem}
\end{exerciseintro}

\begin{solution}
We prove a slightly stronger statement by induction on \(n\), allowing
the curve to be singular but choosing all points in its nonsingular
locus.  This is useful because an intermediate projected image need
not be nonsingular.

For \(n=2\), one general point spans a point and plainly contains no
other point of the reduced integral curve.  Now let \(n\geq3\), choose
a general point \(P_1\in X\), and project from it:
\[
\pi_{P_1}:X\dashrightarrow X'\subset\mathbb P^{n-1}.
\]
Here \(X\) is not planar, since it is nondegenerate and \(n\geq3\).
Let \(\mathcal T\) denote its trisecant family.  The incidence of
pairs \((P,\ell)\) with \(P\in X\cap\ell\) and \(\ell\in\mathcal T\)
has dimension at most one: its fibers over \(\mathcal T\) are finite,
since \(X\) is not a line.  Consequently, only finitely many
trisecants pass through a general \(P_1\).  Projection from \(P_1\)
is therefore one-to-one away from finitely many points of \(X\).
On the normalization this rational map extends to a nonconstant
morphism.  It is finite onto \(X'\), and in characteristic zero
generic injectivity implies degree one, hence birationality.
The image \(X'\) is nondegenerate: if it lay in a
hyperplane of \(\mathbb P^{n-1}\), then \(X\) would lie in the
corresponding hyperplane of \(\mathbb P^n\) through \(P_1\).

Write \(\nu:\widetilde X\to X\) for normalization and
\(f:\widetilde X\to X'\) for the finite birational morphism just
obtained.  Remove from \(X'\) its singular points, the images under
\(f\) of \(\nu^{-1}(X_{\mathrm{sing}})\), and
\(f(\nu^{-1}(P_1))\).  This removes only finitely many points and
leaves a dense open \(U'\) for which the entire inverse image
\(f^{-1}(U')\to U'\) is an isomorphism.  Indeed, a finite
birational map onto a normal curve is an isomorphism.
Normalization identifies \(f^{-1}(U')\) with an open \(U\subset X\).
Thus we have an isomorphism
\[
U\xrightarrow{\ \sim\ }U'
\]
and no additional points of \(X\setminus\{P_1\}\) map into \(U'\).
Choose the remaining
general points \(P_2,\ldots,P_{n-1}\) in \(U\), and put
\[
P_i'=\pi_{P_1}(P_i).
\]
They are general points of \(X'\).  By the induction hypothesis,
\[
L'=\langle P_2',\ldots,P_{n-1}'\rangle
\subset\mathbb P^{n-1}
\]
has dimension \(n-3\) and contains no other point of \(X'\).

The inverse linear span of \(L'\) under projection is
\[
L=\langle P_1,P_2,\ldots,P_{n-1}\rangle,
\]
and
\[
\dim L=\dim L'+1=n-2.
\]
Suppose \(Q\in X\cap L\).  If \(Q=P_1\), it is already listed.
Otherwise,
\[
\pi_{P_1}(Q)\in X'\cap L',
\]
so the induction hypothesis gives
\[
\pi_{P_1}(Q)=P_i'
\]
for some \(i\geq2\).  The points \(P_i'\) were chosen in the open set
over which the birational projection is an isomorphism, and hence
\(Q=P_i\).  Thus \(L\) contains exactly the chosen \(n-1\) points of
\(X\), completing the induction.  These conditions hold on a dense
open subset of the parameter space of tuples: linear independence
is open, while tuples admitting an extra distinct point in their
span form a constructible set by the incidence correspondence.
The argument excludes that set at the generic tuple, so its closure
is proper.  This gives the asserted meaning of general points.
\end{solution}

\begin{exerciseintro}
% record: H-IV-03-011
\exerciseheading{Exercise IV.3.11}
\addcontentsline{toc}{subsection}{Exercise IV.3.11}

\begin{context}
Work over an algebraically closed field of arbitrary characteristic.
A variety in projective space is taken to be closed and integral.
For a smooth projective variety, a base-point-free space of sections
defines a closed immersion if it separates distinct points and
nonzero tangent vectors.

The secant variety \(S(X)\) is the closure of the union of lines
joining distinct points of \(X\).  The embedded tangent space
\(T_PX\) is a projective linear space of dimension \(\dim X\).
We use dimensions of incidence varieties and their fibers.
\end{context}

\begin{problem}
(a) Let \(X\subset\mathbb P^n\) be a nonsingular projective variety of dimension
\(r\), with \(n>2r+1\).  Prove that projection from a suitable point
\(O\notin X\) induces a closed immersion \(X\to\mathbb P^{n-1}\).

(b) Let \(X=\nu_2(\mathbb P^2)\subset\mathbb P^5\) be the Veronese
surface.  Prove that every point on a secant line lies on infinitely
many secant lines.  Show that \(\dim S(X)=4\), and deduce that
projection from a suitable point embeds \(X\) into \(\mathbb P^4\).
\end{problem}
\end{exerciseintro}

\begin{solution}
(a) If \(r=0\), then \(X\) is a single reduced point, and any center
outside \(X\) gives the required embedding.  Assume \(r\geq1\).
The incidence of triples
\[
(P,Q,O),\qquad P,Q\in X,\quad P\ne Q,\quad O\in\langle P,Q\rangle
\]
has dimension \(2r+1\).  Thus its image has closure of dimension at
most \(2r+1\), and
\[
\dim S(X)\leq2r+1.
\]
The incidence of pairs \((P,O)\) with \(O\in T_PX\) has dimension
\(2r\).  Its image is closed, since it is the image of a projective
bundle over the projective variety \(X\).  As \(n>2r+1\), we can
choose \(O\) outside \(X\), \(S(X)\), and all embedded tangent spaces.

Projection from \(O\) is defined everywhere on \(X\).  Two distinct
points could have the same image only if their joining line passed
through \(O\), which was excluded.  Its differential could kill a
nonzero tangent vector at \(P\) only if \(O\in T_PX\), also excluded.
The separation criterion therefore proves that the projection is a
closed immersion.

(b) Write the Veronese coordinates as
\[
\nu_2([x:y:z])=[x^2:y^2:z^2:xy:xz:yz].
\]
For a line \(L\subset\mathbb P^2\), its image is a nonsingular conic
\(C_L\) spanning a plane \(\Pi_L\subset\mathbb P^5\).
If \(P,Q\in X\) are distinct, let \(L\) join their preimages.
Then \(\langle P,Q\rangle\subset\Pi_L\).

Let \(R\) lie on this secant.  If \(R\in C_L\), the lines joining
\(R\) to the other points of \(C_L\) already give infinitely many
secants through \(R\).  If \(R\notin C_L\), projection of the conic
from \(R\) is a finite morphism of degree two to the pencil of lines
through \(R\).  The given secant meets the conic in two distinct
points, so this morphism cannot be purely inseparable.  It is
therefore separable; a nonempty open subset of the pencil consists
of lines meeting the conic in two distinct points.
These are infinitely many secants of \(X\) through \(R\).
This argument also works in characteristic two: the nucleus of a
conic lies on no secant with two distinct endpoints and hence is
not a possible \(R\) in this case.

We compute the dimension without assuming that every point of a
conic plane is on an actual secant.  The closure of the secants of
a nonsingular plane conic is its entire plane.  Indeed, for fixed
\(P\) on the conic, the lines \(PQ\), as \(Q\) varies, are dense in
the pencil through \(P\).  Consequently,
\[
S(X)=\bigcup_{L\in(\mathbb P^2)^*}\Pi_L.
\tag{1}
\]
The right side is closed: it is the image of the projective
plane bundle whose fiber at \(L\) is \(\Pi_L\).  This bundle has
dimension \(2+2=4\), so \(\dim S(X)\leq4\).

Identify the coordinates of \(\mathbb P^5\) with nonzero symmetric
\(3\times3\) matrices up to scalar, so that a point \([v]\) of
\(\mathbb P^2\) maps to \([vv^{\mathrm t}]\).
If \(L=\mathbb P(W)\), where \(W\subset k^3\) has dimension two,
every matrix represented in \(\Pi_L\) has column space contained
in \(W\).  The plane \(\Pi_L\) contains rank-two matrices, for
example \(uu^{\mathrm t}+vv^{\mathrm t}\) for a basis \(u,v\) of
\(W\).  A rank-two matrix in \(\Pi_L\) determines \(W\) uniquely
as its column space.  Thus the map from the bundle in (1) has
one-point fibers on its nonempty rank-two open subset.
The dimension theorem gives \(\dim S(X)=4\).

Finally every tangent line to \(X\) at \(\nu_2(p)\) is the tangent
line to \(\nu_2(L)\) for the line \(L\) through \(p\) in the
specified tangent direction.  It is contained in \(\Pi_L\),
hence in \(S(X)\).  Therefore \(S(X)\) contains all embedded
tangent planes as well as all secants.
A point of \(\mathbb P^5\setminus S(X)\) separates points and
tangent vectors under projection.  The criterion in the context
gives the required closed immersion into \(\mathbb P^4\).
\end{solution}

\begin{exerciseintro}
% record: H-IV-03-012
\exerciseheading{Exercise IV.3.12}
\addcontentsline{toc}{subsection}{Exercise IV.3.12}

\begin{context}
Work over an algebraically closed field of arbitrary characteristic.
A node is a double point whose quadratic tangent cone is a product
of two distinct linear forms.  An integral plane curve of degree
\(d\) with exactly \(r\) nodes and no other singularities has
normalization genus
\[
g=\frac{(d-1)(d-2)}2-r.
\]
We use the Jacobian criterion, the dimension theorem for incidence
varieties, and Bézout's theorem.  The first jet of a local equation
at \(Q\) is its residue modulo \(\mathfrak m_Q^2\); prescribing it
imposes three independent linear conditions if the available
sections surject onto this three-dimensional space.

We also use Riemann--Roch and \(\deg K_Y=2g-2\).
On a nonsingular projective curve of genus \(g\), every divisor of
degree at least \(2g+1\) is very ample.
The generic-projection theorems for curves give embeddings into
\(\mathbb P^3\) by successive projections from larger projective
spaces, and a birational projection of a smooth nonplanar space
curve onto a plane curve having only nodes.
All these projections preserve degree.
\end{context}

\begin{problem}
For every \(d\in\{2,3,4,5\}\) and every integer
\[
0\leq r\leq\frac{(d-1)(d-2)}2,
\]
construct an irreducible plane curve of degree \(d\) with exactly
\(r\) nodes and no other singularities.
\end{problem}
\end{exerciseintro}

\begin{solution}
We first construct curves with a small number of prescribed nodes.
The remaining cases will come from projecting smooth curves.

Let \(A\) be a selected subset of the three coordinate vertices of
\(\mathbb P^2\).  Let \(V_d(A)\) be the vector space of homogeneous
degree-\(d\) forms vanishing to order at least two at every point
of \(A\).  We claim that a general member has precisely the nodes
in \(A\), and is irreducible, in the cases
\[
(d,|A|)=(3,1),\ (4,1),\ (5,1),\ (5,2),\ (5,3).
\tag{1}
\]

First we show that \(V_d(A)\) separates first jets at every
\(Q\notin A\).  It suffices to find a degree-\((d-1)\) form \(B\)
vanishing doubly on \(A\) with \(B(Q)\ne0\): the forms \(B\ell\),
as \(\ell\) ranges over linear forms, then give all first jets at
\(Q\), since multiplication by the local unit \(B\) is invertible.

If \(A=\{[1:0:0]\}\), one of \(y,z\) is nonzero at \(Q\), so one
of \(y^{d-1},z^{d-1}\) works; here \(d-1\geq2\).
For \(d=5\) and two or three selected vertices, use the quartics
\[
x^2y^2,\qquad x^2z^2,\qquad y^2z^2.
\]
They vanish doubly at all three vertices and do not vanish
simultaneously away from those vertices.  If \(Q\) is an unselected
vertex, its corresponding pure fourth power is nonzero there and
vanishes doubly at every selected vertex.  This proves the claim
about first jets in all cases of (1).

Write \(N=\dim\mathbb P(V_d(A))\).
For fixed \(Q\notin A\), being singular at \(Q\) imposes three
independent linear conditions.  The incidence of such pairs
\((Q,[F])\) consequently has dimension at most
\[
2+(N-3)=N-1.
\]
The closure of its image in \(\mathbb P(V_d(A))\) is proper.
A member outside this closure has no singularity away from \(A\),
including on the line at infinity.

At \(P=[1:0:0]\in A\), the forms
\[
x^{d-2}y^2,\qquad x^{d-2}yz,\qquad x^{d-2}z^2
\]
belong to \(V_d(A)\) in all cases of (1) and realize arbitrary
quadratic terms in the local coordinates \(y/x,z/x\).
The same holds at each selected vertex by permuting coordinates.
The condition that this quadratic term have two distinct linear
factors is a nonempty open condition, also in characteristic two;
the quadratic \(yz\) is one example.  The parameter space is
irreducible, so these finitely many nonempty open conditions can
be imposed simultaneously with absence of other singularities.
This gives a curve with exactly \(|A|\) nodes.

Such a curve is reduced, since a repeated component would make
every point of that component singular.  If it were reducible,
divide its components into two nonempty groups of total degrees
\(a\) and \(d-a\).  Their intersections must all occur among the
nodes, and at a node belonging to both groups the intersection
multiplicity is one.  Bézout would give
\[
a(d-a)\leq |A|.
\]
But \(a(d-a)\geq d-1>|A|\) in every case of (1).
This contradiction proves irreducibility.

The same first-jet incidence argument, now in the full space of
degree-\(d\) forms, gives a nonsingular plane curve for every
\(d\geq2\).  It is irreducible: two positive-degree components
would meet by Bézout and produce a singular point.
Thus the case \(r=0\) is available for each required \(d\).

We now produce the remaining curves using the following construction.
Take a nonsingular projective curve \(Y\) of genus \(g\) and an
effective divisor \(D\) of degree \(d\geq2g+1\).
Its complete linear system embeds \(Y\) with degree \(d\).
If the ambient dimension exceeds three, project successively to
an embedding in \(\mathbb P^3\).  A birational generic projection
to \(\mathbb P^2\) then has integral image of degree \(d\), with
only nodes, and normalization \(Y\).  Its number of nodes is
therefore exactly
\[
r=\frac{(d-1)(d-2)}2-g.
\tag{2}
\]
In the applications below the complete embedding has dimension
at least three and is nonplanar, as required for this projection
theorem.

For \(g=0\), take \(Y=\mathbb P^1\) and \(D\) of degree
\(d=3,4,5\).  The complete embedding is the rational normal curve
in \(\mathbb P^d\), and (2) gives \(r=1,3,6\), respectively.

For \(g=1\), take \(Y\) to be any nonsingular plane cubic, whose
existence was just proved.  Divisors of degrees \(4\) and \(5\)
are very ample.  Riemann--Roch gives \(h^0(D)=d\), so the complete
embeddings lie in \(\mathbb P^3\) and \(\mathbb P^4\), respectively.
Formula (2) gives a quartic with \(r=2\) and a quintic with \(r=5\).

For \(g=2\), normalize the one-node plane quartic constructed in
(1).  The normalization has genus \(3-1=2\).
A divisor of degree \(5=2g+1\) is very ample and has
\(h^0(D)=5+1-2=4\).  Its embedding into \(\mathbb P^3\), followed
by generic projection, gives a quintic with \(r=6-2=4\).

Together the constructions cover exactly the following ranges:
\[
\begin{array}{c|c|c}
d&\text{direct plane construction}&\text{projection construction}\\
2&r=0&\\
3&r=0,1&r=1\\
4&r=0,1&r=2,3\\
5&r=0,1,2,3&r=4,5,6.
\end{array}
\]
All the curves are integral and have no singularities other than
the stated nodes.  No restriction on the characteristic was used.
\end{solution}

\begin{exerciseintro}
\corpussection{IV.4 Elliptic Curves}
\addcontentsline{toc}{section}{IV.4 Elliptic Curves}

% record: H-IV-04-001
\exerciseheading{Exercise IV.4.1}
\addcontentsline{toc}{subsection}{Exercise IV.4.1}

\begin{context}
Let \(X\) be a nonsingular projective integral curve of genus one over
an algebraically closed field \(k\) of characteristic different from
two, and let \(P\in X(k)\).  Write
\[
L(nP)=\{f\in k(X):\operatorname{div}(f)+nP\geq0\}\cup\{0\}.
\]
Riemann--Roch gives \(h^0(nP)=n\) for \(n\geq1\) and \(h^0(0)=1\).
The Legendre normal form with \(P\) at infinity supplies functions
\(u,v\) with their only poles at \(P\), of orders two and three,
and an equation
\[
v^2=u(u-1)(u-\lambda),\qquad \lambda\in k\setminus\{0,1\}.
\]
\end{context}

\begin{problem}
Show that the section ring
\[
R=\bigoplus_{n\geq0}H^0(X,\mathcal O_X(nP))
\]
admits a graded presentation
\[
R\cong k[t,x,y]/\bigl(y^2-x(x-t^2)(x-\lambda t^2)\bigr),
\qquad \deg t=1,\quad\deg x=2,\quad\deg y=3.
\]
\end{problem}
\end{exerciseintro}

\begin{solution}
Represent the direct sum as the subring
\(\bigoplus_{n\geq0}L(nP)T^n\subset k(X)[T]\), where \(T\) is a
formal variable.  Define
\[
t\longmapsto T,\qquad x\longmapsto uT^2,\qquad
y\longmapsto vT^3.
\]
The Legendre equation makes the indicated weighted homogeneous
polynomial vanish, so this defines a graded homomorphism from the
claimed quotient to \(R\).

Since the relation is monic of degree two in \(y\), the quotient
is a free \(k[t,x]\)-module with basis \(1,y\).  Its degree-\(n\)
part therefore has basis
\[
t^{n-2a}x^a\quad(0\leq2a\leq n),\qquad
t^{n-3-2a}x^ay\quad(0\leq3+2a\leq n).
\]
Their images, after removing \(T^n\), are \(u^a\) and \(u^av\).
Their pole orders at \(P\) are respectively \(2a\) and \(2a+3\).
These orders are all distinct.  A linear combination cannot vanish:
its term of largest pole order cannot be canceled by any other term.
Thus the map is injective on each homogeneous part.

For \(n\geq1\), the displayed orders are precisely
\[
0,2,3,\ldots,n,
\]
with only \(0\) present when \(n=1\).  There are \(n\) of them,
equal to \(\dim L(nP)\).  They therefore form a basis of \(L(nP)\).
In degree zero both sides are \(k\).  The homomorphism is consequently
an isomorphism in every degree, proving the presentation.
\end{solution}

\begin{exerciseintro}
% record: H-IV-04-002
\exerciseheading{Exercise IV.4.2}
\addcontentsline{toc}{subsection}{Exercise IV.4.2}

\begin{context}
Let \(X\) be a nonsingular projective integral genus-one curve over
an algebraically closed field.  Its canonical bundle is trivial;
for a positive-degree line bundle \(M\), Riemann--Roch gives
\(h^0(M)=\deg M\) and \(H^1(X,M)=0\).
A line bundle of degree at least three is very ample.
A complete embedding of a normal variety by \(L\) is projectively
normal if all multiplication maps
\[
\operatorname{Sym}^mH^0(X,L)\longrightarrow H^0(X,L^m)
\]
are surjective for \(m\geq0\).
We use \(H^1(\mathbb P^2,\mathcal O(q))=0\) for every integer \(q\).
\end{context}

\begin{problem}
If \(D\) is a divisor of degree at least three on \(X\), prove that
the embedding defined by its complete linear system is projectively
normal.
\end{problem}
\end{exerciseintro}

\begin{solution}
Put \(L=\mathcal O_X(D)\) and \(d=\deg L\).
We first prove surjectivity of multiplication in degree two by
induction on \(d\).
For \(d=3\), the complete embedding is a smooth plane cubic.
Its restriction sequence
\[
0\longrightarrow\mathcal O_{\mathbb P^2}(m-3)
\longrightarrow\mathcal O_{\mathbb P^2}(m)
\longrightarrow L^m\longrightarrow0
\]
and the stated \(H^1\)-vanishing give surjectivity of restriction
in every degree, in particular degree two.

Let \(d\geq4\), choose \(P\in X\), and put \(L'=L(-P)\).
By induction, products of sections of \(L'\) span
\[
H^0((L')^2)=H^0(L^2(-2P)).
\]
Choose \(s\in H^0(L)\) not vanishing at \(P\).
Also choose \(t\in H^0(L(-P))\) vanishing to order exactly one
at \(P\); this is possible because
\[
h^0(L(-P))-h^0(L(-2P))=(d-1)-(d-2)=1.
\]
The quotient \(H^0(L^2)/H^0(L^2(-2P))\) has dimension two.
In a local trivialization at \(P\), the classes of \(s^2\) and \(st\)
have vanishing orders zero and one, so they are linearly independent
in this quotient.  Together with the products from \(L'\), they
span \(H^0(L^2)\).  This proves the induction step.

To pass to higher degrees, choose two sections \(a,b\in H^0(L)\)
with no common zero.  One can first choose nonzero \(a\), then
choose \(b\) avoiding its finite zero set, since \(L\) is globally
generated and \(k\) is infinite.  Their evaluation sequence is
\[
0\longrightarrow L^{-1}
\longrightarrow\mathcal O_X^{\,2}
\xrightarrow{(a,b)}L\longrightarrow0.
\]
Tensor by \(L^m\), where \(m\geq2\).
Since \(H^1(L^{m-1})=0\), the resulting cohomology sequence gives
a surjection
\[
H^0(L^m)^{\,2}\longrightarrow H^0(L^{m+1}),
\qquad (u,v)\longmapsto au+bv.
\]
Starting from degree two, induction on \(m\) now shows that every
section of \(L^m\) is a sum of products of \(m\) sections of \(L\).
Degrees zero and one require no argument.  The normal-generation
criterion proves projective normality.
\end{solution}

\begin{exerciseintro}
% record: H-IV-04-003
\exerciseheading{Exercise IV.4.3}
\addcontentsline{toc}{subsection}{Exercise IV.4.3}

\begin{context}
Let \(k\) be algebraically closed with \(\operatorname{char}k\ne2\),
and let \(E\) be the smooth projective completion of
\[
y^2=F(x)=x(x-1)(x-\lambda),\qquad \lambda\ne0,1,
\]
with origin \(O=[0:1:0]\).
The functions \(x,y\) have poles of orders two and three at \(O\).
Riemann--Roch gives
\[
L(2O)=\langle1,x\rangle,\qquad L(3O)=\langle1,x,y\rangle.
\]
An automorphism fixing \(O\) is a group automorphism, so commutes
with inversion \((x,y)\mapsto(x,-y)\).
The invariant is
\[
j=256\frac{(\lambda^2-\lambda+1)^3}{\lambda^2(\lambda-1)^2}.
\]
The Legendre parameters for a fixed \(j\) are the orbit under
\(\lambda\mapsto1-\lambda\) and \(\lambda\mapsto1/\lambda\).
\end{context}

\begin{problem}
Prove that every automorphism of \(E\) fixing \(O\) is induced by a
projective linear transformation whose affine expression is
\[
x'=ax+b,\qquad y'=cy.
\]
Describe these transformations and determine the group
\(G=\operatorname{Aut}(E,O)\) in the four cases
\[
j\ne0,1728;\quad
j=1728,\ \operatorname{char}k\ne3;\quad
j=0,\ \operatorname{char}k\ne3;\quad
j=0,\ \operatorname{char}k=3.
\]
\end{problem}
\end{exerciseintro}

\begin{solution}
An automorphism \(\sigma\) preserving \(O\) preserves each pole-order
space, hence
\[
\sigma^*x=ax+b,\qquad
\sigma^*y=cy+dx+e,
\qquad a,c\ne0.
\]
The nonzero leading coefficients follow from the exact pole orders.
Commuting with inversion forces \(d=e=0\), since the characteristic
is not two.  Thus \(\sigma\) extends to
\[
[X:Y:Z]\longmapsto[aX+bZ:cY:Z]
\]
in the plane.  Such a transformation preserves \(E\) precisely when
\[
F(ax+b)=c^2F(x).
\tag{1}
\]
Equivalently, the affine map \(x\mapsto ax+b\) permutes
\(\{0,1,\lambda\}\) and \(c^2=a^3\).
Each allowed affine permutation has exactly two lifts, differing
by inversion.  These conditions describe the transformations for
every original choice of \(\lambda\).

For completeness, a nonidentity affine permutation of the three
roots either exchanges two roots and fixes the third, or cyclically
permutes them.  In the first case the fixed root is the midpoint
of the other two, giving \(\lambda\in\{-1,2,1/2\}\).
In the second case, after ordering the roots, the map is
\[
x\longmapsto(\lambda-1)x+1,
\]
and its sending \(\lambda\) to zero is equivalent to
\(\lambda^2-\lambda+1=0\).
Thus there are no other exceptional parameters.

If \(j\ne0,1728\), neither exception occurs, so \(a=1,b=0,c=\pm1\).
The group is cyclic of order two.

If \(j=1728\) and the characteristic is not three, an affine
relabeling and rescaling puts the equation in the form
\(y^2=x^3-x\), corresponding to \(\lambda=-1\).
Its automorphisms are
\[
(x,y)\longmapsto(u^2x,u^3y),\qquad u^4=1.
\]
These four transformations form a cyclic group of order four.
Conjugating by the change of coordinates gives the transformations
for \(\lambda=2\) or \(1/2\).

If \(j=0\) and the characteristic is not three, then
\(\lambda^2-\lambda+1=0\).
The transformation
\[
\rho(x,y)=((\lambda-1)x+1,y)
\]
cycles the three roots and has order three; indeed
\((\lambda-1)^3=1\).
Together with \(\iota(x,y)=(x,-y)\), it gives all six
automorphisms \(\rho^a\iota^b\), \(0\leq a<3\), \(0\leq b<2\).
The generators commute, so \(G\) is cyclic of order six.

Finally, in characteristic three, \(j=0\) gives \(\lambda=-1\)
and the equation \(y^2=x^3-x\).
The affine permutations of its root set \(\mathbb F_3\) are
\[
x\longmapsto ax+b,\qquad a=\pm1,\quad b\in\mathbb F_3.
\]
Equation (1) says \(c^2=a\).  Hence the twelve automorphisms are
\[
(x,y)\longmapsto(c^2x+b,cy),\qquad c^4=1,\quad b\in\mathbb F_3.
\]
Choose \(i\in k\) with \(i^2=-1\).  The maps
\[
\tau(x,y)=(x+1,y),\qquad \sigma(x,y)=(-x,iy)
\]
generate \(G\) and satisfy
\[
\tau^3=1,\qquad \sigma^4=1,\qquad
\sigma\tau\sigma^{-1}=\tau^{-1}.
\]
Thus \(G\) is the semidirect product of a cyclic group of order
three by a cyclic group of order four, with the latter acting
by inversion.  In particular this group is not cyclic.
\end{solution}

\begin{exerciseintro}
% record: H-IV-04-004
\exerciseheading{Exercise IV.4.4}
\addcontentsline{toc}{subsection}{Exercise IV.4.4}

\begin{context}
As in the elliptic-curve section, work over an algebraically closed
field \(k\) of characteristic different from two.
For a Legendre equation \(y^2=x(x-1)(x-\lambda)\), the invariant is
\[
j=256\frac{(\lambda^2-\lambda+1)^3}{\lambda^2(\lambda-1)^2}.
\]
A monic cubic with coefficients \(A,B,C\) has discriminant
\[
\delta=A^2B^2-4B^3-4A^3C-27C^2+18ABC.
\]
This is the product of the squared pairwise differences of its
roots; it is nonzero precisely when the roots are distinct.
The projective Weierstrass cubic \(y^2=x^3+Ax^2+Bx+C\) is smooth
if and only if \(\delta\ne0\).
\end{context}

\begin{problem}
Let \(E\) be the nonsingular projective cubic with affine equation
\[
y^2+a_1xy+a_3y=x^3+a_2x^2+a_4x+a_6.
\]
Express \(j(E)\) as a rational function of the coefficients with
rational coefficients.  Prove that \(j(E)\in k_0\) if all \(a_i\)
belong to a subfield \(k_0\subset k\).
Conversely, for each \(j_0\in k_0\), construct an elliptic curve
over \(k_0\) with invariant \(j_0\).
\end{problem}
\end{exerciseintro}

\begin{solution}
Complete the square by setting \(Y=y+(a_1x+a_3)/2\).
The equation becomes \(Y^2=x^3+Ax^2+Bx+C\), where
\[
A=a_2+\frac{a_1^2}{4},\qquad
B=a_4+\frac{a_1a_3}{2},\qquad
C=a_6+\frac{a_3^2}{4}.
\]
We claim that
\[
j(E)=256\frac{(A^2-3B)^3}
{A^2B^2-4B^3-4A^3C-27C^2+18ABC}.
\tag{1}
\]
To see this, let the distinct roots of the cubic be
\(\alpha,\beta,\gamma\), put \(h=\beta-\alpha\), and set
\(\lambda=(\gamma-\alpha)/h\).
After \(x=\alpha+hu\) and rescaling \(Y\), the equation is
Legendre form.  Directly from the symmetric expressions in the
roots,
\[
A^2-3B=h^2(1-\lambda+\lambda^2),\qquad
\delta=h^6\lambda^2(1-\lambda)^2.
\]
Substitution into the Legendre formula gives (1).
No division by three is used, so the formula also holds in
characteristic three.

Formula (1), with the displayed expressions for \(A,B,C\), is
the required rational function over \(\mathbb Q\).  Its only
constant denominators are powers of two, and its discriminant
denominator is nonzero because \(E\) is smooth.  If the original
coefficients lie in \(k_0\), the formula therefore gives \(j(E)\in k_0\).

We now construct the converse over \(k_0\) itself.
First suppose its characteristic is different from both two and
three.  For \(j_0=0\), take \(y^2=x^3+1\); for \(j_0=1728\),
take \(y^2=x^3+x\).
For \(j_0\notin\{0,1728\}\), put
\[
t=\frac{j_0}{1728-j_0},\qquad
E:\ y^2=x^3+3tx+2t.
\]
Here \(t\ne0,-1\), and the cubic discriminant is
\(-108t^2(t+1)\ne0\).
Formula (1) gives
\[
j(E)=1728\frac{t}{t+1}=j_0.
\]

In characteristic three, take \(y^2=x^3-x\) when \(j_0=0\).
Its discriminant is \(4=1\), and (1) gives \(j=0\).
When \(j_0\ne0\), take
\[
E:\ y^2=x^3+x^2-\frac1{j_0}.
\]
Now \(A=1,B=0,C=-1/j_0\), so
\(\delta=1/j_0\ne0\) and \(256(A^2-3B)^3=1\) in \(k_0\).
Again (1) gives \(j(E)=j_0\).
Every constructed cubic is smooth and has the \(k_0\)-rational
point at infinity, hence is an elliptic curve over \(k_0\).
\end{solution}

\begin{exerciseintro}
% record: H-IV-04-005
\exerciseheading{Exercise IV.4.5}
\addcontentsline{toc}{subsection}{Exercise IV.4.5}

\begin{context}
Let \(k\) be algebraically closed of characteristic different from
two.  A nonconstant morphism between elliptic curves is finite.
An origin-preserving morphism is a group homomorphism, and a
separable such morphism is unramified by Hurwitz.
Every degree-two morphism from an elliptic curve to \(\mathbb P^1\)
is the quotient by an involution \(Q\mapsto A-Q\).
If it is ramified at the origin, its involution is inversion.
A double cover of \(\mathbb P^1\) is determined, up to isomorphism,
by its branch set; four branch points give genus one.
The Legendre invariant is
\[
j(\lambda)=256\frac{(\lambda^2-\lambda+1)^3}
{\lambda^2(\lambda-1)^2},
\]
and two elliptic curves over \(k\) are isomorphic if and only if
their invariants agree.
\end{context}

\begin{problem}
Let \(f:E\to E\) be an endomorphism of degree two, and let
\(\pi:E\to\mathbb P^1\) have degree two and be ramified at the origin.

(a) Find a degree-two map \(\pi':E\to\mathbb P^1\) and a
degree-two map \(g:\mathbb P^1\to\mathbb P^1\) such that
\(\pi f=g\pi'\).

(b) Choose coordinates so that \(g(u)=u^2\).

(c) Relate the branch points of \(\pi,\pi'\), and \(g\), and derive
an equation for a Legendre parameter of \(E\).

(d) Determine the possible \(j\)-invariants and corresponding
parameters.  In characteristic zero there are three distinct
invariants; in positive characteristic allow their reductions
to coincide.
\end{problem}
\end{exerciseintro}

\begin{solution}
(a) Since \(f\) is a homomorphism, it commutes with inversion.
Use the inversion quotient also on its source as \(\pi'\).
Then \(\pi f\) descends to a map \(g\) of the two quotient lines.
The equality of degrees in \(\pi f=g\pi'\) gives \(\deg g=2\).
Since the characteristic is not two, all these double covers
are separable, and \(f\) is unramified.

(b) Hurwitz shows that \(g\) has two ramification points, with
distinct images.  Send its ramification points to \(0,\infty\)
in the source and their images to \(0,\infty\) in the target.
The resulting rational function has zero divisor \(2(0)\) and
pole divisor \(2(\infty)\), hence is \(cu^2\).
Rescale a coordinate to obtain \(g(u)=u^2\).

(c) For \(P\in E\), compare ramification indices in the square:
\[
e_\pi(f(P))=e_g(\pi'(P))e_{\pi'}(P).
\tag{1}
\]
The left side is at most two.  Above either ramification point
of \(g\), formula (1) forces \(e_{\pi'}=1\) and \(e_\pi=2\).
Thus the two branch values of \(g\) belong to the branch set
of \(\pi\).  Above each of the other two branch values of
\(\pi\), the map \(g\) is unramified; (1) then forces
\(\pi'\) to ramify at each of their four preimages.
These are all four branch points of \(\pi'\).

Scale coordinates so that the branch set of \(\pi\) is
\(\{0,\infty,1,\lambda\}\), where \(\lambda\ne0,1\), and that
\(g(u)=u^2\) still holds.  Write \(s^2=\lambda\).
The branch set of \(\pi'\) is
\[
\{1,-1,s,-s\}.
\]
Its cross-ratio can be taken to be
\(\mu=((s-1)/(s+1))^2\).
Substitution into the Legendre formula gives
\[
j(\mu)=16\frac{(\lambda^2+14\lambda+1)^3}
{\lambda(\lambda-1)^4}.
\]
Both covers have source isomorphic to \(E\), so \(j(\lambda)=j(\mu)\).
Clearing the nonzero denominators yields
\[
16(\lambda^2-\lambda+1)^3(\lambda-1)^2
-\lambda(\lambda^2+14\lambda+1)^3=0.
\tag{2}
\]

(d) The polynomial in (2) factors as
\[
(\lambda+1)^2(\lambda^2-6\lambda+1)
(\lambda^2-\lambda+16)(16\lambda^2-\lambda+1).
\tag{3}
\]
The resulting possibilities, with this choice of branch ordering, are
\[
\begin{array}{c|c}
\text{condition on }\lambda&j\\
\lambda=-1&1728\\
\lambda^2-6\lambda+1=0&8000\\
\lambda^2-\lambda+16=0\ \text{or }16\lambda^2-\lambda+1=0&-3375.
\end{array}
\tag{4}
\]
The table follows by substituting each quadratic relation in
the formula for \(j(\lambda)\).
In characteristic zero its parameters are
\[
-1,\quad 3\pm2\sqrt2,\quad
\frac{1\pm3\sqrt{-7}}2,\quad
\frac{1\pm3\sqrt{-7}}{32}.
\]
Other orderings of the four branch points give the usual six
Legendre transforms of these parameters.

We verify sufficiency as well.  For any root \(\lambda\) of (3),
take the smooth completion \(D\) of
\[
v^2=(u^2-1)(u^2-\lambda).
\]
None of these roots is zero or one: the left side of (2) takes
the nonzero values \(16\) and \(-4096\) there.
Thus \(D\) is a smooth genus-one double cover with branch set
\(\{\pm1,\pm s\}\).  The formulas
\[
x=u^2,\qquad y=uv
\]
give a degree-two morphism \(D\to E_\lambda\), where
\(E_\lambda\) is the Legendre curve.
Its function-field extension has degree at most two, since
\(u^2=x\) and \(v=y/u\).  It has degree exactly two because the
nontrivial involution \((u,v)\mapsto(-u,-v)\) fixes \(x,y\).
By (2), \(D\) and \(E_\lambda\) have equal invariant, hence are
isomorphic.  Choose as origin of \(D\) a point above the origin
of \(E_\lambda\); composing an isomorphism with a translation
makes it preserve these origins.
The displayed map then gives a degree-two endomorphism of
\(E_\lambda\).

Consequently the possible invariants are precisely the reductions
of \(1728,8000,-3375\).  They are distinct in characteristic zero.
In odd positive characteristic collisions occur exactly in
characteristics \(3,5,7,13\), as is seen by factoring their
pairwise differences.  The statement about three distinct
invariants must be understood with this qualification.
\end{solution}

\begin{exerciseintro}
% record: H-IV-04-006
\exerciseheading{Exercise IV.4.6}
\addcontentsline{toc}{subsection}{Exercise IV.4.6}

\begin{context}
Work over an algebraically closed field of characteristic zero.
For a base-point-free linear series \(V\subset H^0(X,L)\) of
dimension \(n+1\) on a smooth projective curve, its vanishing
sequence at \(P\) is
\[
0=a_0(P)<a_1(P)<\cdots<a_n(P),
\]
the orders of vanishing of a basis adapted to \(P\).
Its ramification weight is
\[
w(P)=\sum_{i=0}^n(a_i(P)-i).
\]
We use \(\deg\omega_X=2g-2\), Riemann--Roch, and the
identification \(P\mapsto\mathcal O_X(P-O)\) of an elliptic curve
with its degree-zero Picard group.
For an immersed plane branch, the vanishing sequence of line
sections is \((0,1,m)\), where \(m\) is its tangent-line contact.
\end{context}

\begin{problem}
(a) Let a smooth genus-\(g\) curve \(X\) be the normalization of
an integral plane curve of degree \(d\geq2\) whose only singularities
are nodes.  Establish the inflection formula
\(6(g-1)+3d\), with ramification multiplicities, and explain the
qualification needed if points lying over nodes are excluded.

(b) Let \(X\subset\mathbb P^n\), \(n\geq3\), be a nondegenerate
smooth curve of degree \(d\).  A hyperplane with contact at least
\(n\) is osculating; contact at least \(n+1\) is hyperosculating.
Show that every point has an osculating hyperplane and that the
total hyperosculation weight is
\[
n(n+1)(g-1)+(n+1)d.
\]

(c) For an elliptic curve \(E\) and \(d\geq3\), apply the formula
to the embedding by \(|dO|\) to prove \(\#E[d](k)=d^2\).
Here \(E[d]\) means points killed by \(d\), rather than points of
exact order \(d\).
\end{problem}
\end{exerciseintro}

\begin{solution}
We first prove the ramification formula for any such linear series.
On an open set choose a local parameter \(z\) and a frame \(e\)
for \(L\), and write a basis of \(V\) as \(s_j=f_je\).
Form the Wronskian
\[
W=\det\left(\frac{d^if_j}{dz^i}\right)_{0\leq i,j\leq n}.
\]
The expressions
\[
W e^{\,n+1}(dz)^{n(n+1)/2}
\]
glue to a global section of
\[
L^{n+1}\otimes\omega_X^{\,n(n+1)/2}.
\tag{1}
\]
Indeed, replacing a frame or a parameter changes successive
derivatives by triangular row operations; their diagonal factors
are exactly canceled by those of the displayed tensor factors.
Changing the basis of \(V\) multiplies the section by a nonzero
constant, which leaves its zero divisor unchanged.

At \(P\), take an adapted basis with
\(f_j=c_jz^{a_j}+\text{higher terms}\), where \(c_j\ne0\).
The lowest coefficient of the determinant is
\[
\left(\prod_jc_j\right)\prod_{i<j}(a_j-a_i),
\]
up to the fixed nonzero sign determined by the basis ordering.
It is nonzero in characteristic zero.  Thus the section in (1)
is nonzero, and its order at \(P\) is
\(\sum_j a_j-n(n+1)/2=w(P)\).
Taking degrees gives
\[
\sum_{P\in X}w(P)=(n+1)d+n(n+1)(g-1).
\tag{2}
\]

(a) Pull back the plane line sections to the normalization.
They form a base-point-free three-dimensional series of degree
\(d\).  With \(n=2\), formula (2) gives \(3d+6(g-1)\).
At a smooth point of the plane curve the contribution is \(m-2\),
where \(m\) is tangent contact.  In particular an ordinary flex
has weight one and a higher flex has larger weight.

Each branch over a node is also immersed, with sequence
\((0,1,m)\).  Therefore, if the nodes themselves are excluded,
the precise weighted count of smooth inflection points is
\[
3d+6(g-1)-\sum_{\nu(P)\text{ a node}}(m_P-2).
\tag{3}
\]
When neither branch of any node has tangent contact above two,
the correction vanishes and the claimed formula holds.
A node contributes no automatic weight, but an inflectionary
branch above it does contribute to the Wronskian.

This qualification is necessary.  For instance, the quartic
\[
xyz^2+x^4+y^4=0
\]
has a unique node at \([0:0:1]\), as its three partial derivatives
show.  It is reduced, since a repeated component would give
infinitely many singular points.  If it were reducible, its
components could meet only at that node, with total intersection
multiplicity one, contradicting Bézout for degree four.
Its normalization thus has genus two.
The two local branches have expansions
\(y=-x^3+\text{higher terms}\) and
\(x=-y^3+\text{higher terms}\).
Each has tangent contact three and weight one.  Consequently
the weighted number of smooth flexes is \(18-2=16\), although
the total ramification weight on the normalization is \(18\).

(b) A basis adapted to \(P\) has \(a_n(P)\geq n\).
Its last section defines an osculating hyperplane.
There is a hyperosculating hyperplane exactly when \(a_n(P)>n\),
equivalently \(w(P)>0\), since the increasing nonnegative sequence
equals \(0,1,\ldots,n\) precisely when its last term is \(n\).
Formula (2) proves the assertion, counting these points with
weights \(w(P)\); the unweighted count can be smaller.

(c) Put \(L=\mathcal O_E(dO)\).  It has \(h^0(L)=d\) and is
very ample.  For \(0\leq m<d\), Riemann--Roch gives
\[
h^0(L(-mP))=d-m.
\]
For \(m=d\), the dimension is one if \(L\cong\mathcal O_E(dP)\)
and zero otherwise, since a degree-zero line bundle has a
nonzero section exactly when it is trivial.
Thus the vanishing sequence is \(0,1,\ldots,d-1\), except that
its last term is \(d\) when \(d(P-O)=0\) in \(\operatorname{Pic}^0(E)\).
The exceptional points are precisely \(E[d](k)\), and each has
weight one.  Formula (2), with \(g=1,n=d-1\), gives
\[
\#E[d](k)=d^2.
\]
This is a count of distinct points because all their weights
have just been shown to be one.
\end{solution}

\begin{exerciseintro}
% record: H-IV-04-007
\exerciseheading{Exercise IV.4.7}
\addcontentsline{toc}{subsection}{Exercise IV.4.7}

\begin{context}
Let \(E,E'\) be elliptic curves over an algebraically closed field
of characteristic different from two, with origins \(O,O'\).
Use \(Q\mapsto\mathcal O_E(Q-O)\) to identify
\(E\cong\operatorname{Pic}^0(E)\), including the functorial
identification for algebraic families.
Write \([n]_E\) for multiplication by \(n\).
An origin-preserving morphism is a group homomorphism; if
nonconstant, it is finite.  Its separable part is unramified,
and its purely inseparable part has one geometric point in each
fiber with multiplicity equal to its degree.

We use the following form of the see-saw principle.
If a line bundle \(N\) on a product of smooth projective curves
restricts trivially to every fiber of the first projection, then
\(N\) is pulled back from the first factor.
Indeed, cohomology and base change give a rank-one direct image,
and the evaluation from its pullback is an isomorphism on every
fiber, hence an isomorphism.  Restriction to a horizontal section
identifies the line bundle on the first factor.
Finally \([2]_E\) is separable of degree four: its differential
is multiplication by two, and its kernel consists of the origin
and the three finite branch points of the Legendre map.
\end{context}

\begin{problem}
(a) Show that pullback by a morphism \(f:E\to E'\) induces a
homomorphism of elliptic curves
\(\widehat f:E'\to E\), called its dual.

(b) Prove \(\widehat{g\circ f}=\widehat f\circ\widehat g\).

(c) If \(f\) preserves the origins and has degree \(n>0\), prove
\[
\widehat f\circ f=[n]_E,\qquad
f\circ\widehat f=[n]_{E'}.
\]
Explain the separable and inseparable cases.

(d) For origin-preserving \(f,g:E\to E'\), prove
\(\widehat{f+g}=\widehat f+\widehat g\).

(e) Prove \(\widehat{[n]_E}=[n]_E\) for every integer \(n\), and
\(\deg[n]_E=n^2\) for \(n\ne0\).

(f) Prove \(\deg\widehat f=\deg f\) for nonconstant \(f\).
\end{problem}
\end{exerciseintro}

\begin{solution}
(a) Pullback takes degree-zero line bundles to degree-zero line
bundles and respects tensor products.
For a constant \(f\), its pullback on \(\operatorname{Pic}^0\) is
zero.  Otherwise degrees of divisors are multiplied by \(\deg f\),
which also proves preservation of degree zero.
The functorial Picard identification makes this pullback an
algebraic morphism of the representing elliptic curves.
It preserves the identity and is a homomorphism; explicitly,
\[
\mathcal O_E(\widehat f(Q')-O)
\cong f^*\mathcal O_{E'}(Q'-O').
\tag{1}
\]

(b) The identity \((g\circ f)^*=f^*g^*\) on line bundles and
(1) give the asserted reversal of composition.

(c) First suppose \(f\) is separable.
Its kernel \(K\) has \(n\) points, and translation identifies
the fiber over \(f(Q)\) with \(Q+K\).
Consequently
\[
f^*((f(Q))-(O'))=\sum_{T\in K}((Q+T)-(T)).
\]
The degree-zero divisor on the right represents \(nQ\) in the
elliptic group law.  Equation (1) gives \(\widehat f(f(Q))=nQ\).
If \(f\) is purely inseparable, its fibers over \(f(Q)\) and \(O'\)
are respectively \(n(Q)\) and \(n(O)\), giving the same result.
In general let the inseparable degree be \(e\) and the separable
degree be \(s\).  There are \(s\) kernel points, each fiber point
has multiplicity \(e\), and the same computation gives
\(\widehat f(f(Q))=esQ=nQ\).
Thus \(\widehat f f=[n]_E\).
Since \(f\) is surjective,
\[
(f\widehat f)f=f[n]_E=[n]_{E'}f
\]
implies \(f\widehat f=[n]_{E'}\).
Equality on geometric points here is equality of morphisms,
since the curves are reduced over an algebraically closed field.

(d) A degree-zero line bundle \(M\) on \(E'\) is invariant under
translation.  To verify this, represent it by a degree-zero
divisor \(\sum_i a_i(Q_i)\).  Translating its pullback shifts all
the \(Q_i\) by the same point, and their group sum changes by
that point times \(\sum_i a_i=0\).
The degree-zero Picard identification therefore leaves its class
unchanged.

Let \(m:E'\times E'\to E'\) be addition and \(p_1,p_2\) the
projections.  Set
\[
N=m^*M\otimes(p_1^*M)^{-1}\otimes(p_2^*M)^{-1}.
\]
Translation invariance makes \(N\) trivial on every fiber of
\(p_1\).  Its restriction to \(E'\times\{O'\}\) is also trivial.
The see-saw principle gives \(N\cong\mathcal O_{E'\times E'}\).
Pulling back along \((f,g):E\to E'\times E'\) yields
\[
(f+g)^*M\cong f^*M\otimes g^*M.
\]
Under (1), this is exactly the additivity of the dual.

(e) Duality sends the identity to itself and the zero homomorphism
to zero.  Part (d), including additive inverses, gives
\(\widehat{[n]}=[n]\) for every integer \(n\).
We check that no nonzero integer acts as zero, without assuming
the degree formula to be proved.
The separable map \([2^r]\) has degree \(4^r\), so
\(\#E[2^r](k)=4^r\).
If a nonzero integer \(n=2^vu\), \(u\) odd, killed all points of
\(E\), then Bezout's identity would imply
\(E[2^r](k)\subset E[2^v](k)\) for every \(r>v\).
This contradicts their respective sizes.
Hence \([n]\ne0\) for \(n\ne0\), and the homomorphism
\(\mathbb Z\to\operatorname{End}(E)\) is injective.

Let \(a=\deg[n]\).  Apply part (c) and self-duality:
\[
[a]=\widehat{[n]}[n]=[n]^2=[n^2].
\]
Injectivity of the integer action gives \(a=n^2\).

(f) First suppose \(f\) preserves the origins, and put
\(n=\deg f>0\).  From \(f\widehat f=[n]\), both maps are
nonconstant and
\[
(\deg f)(\deg\widehat f)=\deg[n]=n^2.
\]
Thus \(\deg\widehat f=n\).
An arbitrary nonconstant \(f\) is an origin-preserving
homomorphism followed by translation.  Translation has degree
one and acts trivially on \(\operatorname{Pic}^0\), so it does
not affect the dual or the degree.  The same conclusion follows.
\end{solution}

\begin{exerciseintro}
% record: H-IV-04-008
\exerciseheading{Exercise IV.4.8}
\addcontentsline{toc}{subsection}{Exercise IV.4.8}

\begin{context}
Let \(E\) be an elliptic curve over an algebraically closed field
\(k\) of characteristic zero or \(p>2\), with origin \(O\).
Its algebraic fundamental group is the inverse limit of the
deck groups of connected pointed finite etale Galois covers.
Write \(\mathbb Z_\ell=\varprojlim_r\mathbb Z/\ell^r\mathbb Z\).

We use Hurwitz, and the following isogeny results:
\(\deg[n]=n^2\), the differential of \([n]\) is multiplication
by \(n\), and an isogeny \(f\) has a dual satisfying
\(f\widehat f=\widehat f f=[\deg f]\).
A finite map of smooth curves factors uniquely through a purely
inseparable map followed by a separable map; for an isogeny the
separable factor is etale.
The tangent space of \(\operatorname{Pic}^0(E)\) at the identity
is \(H^1(E,\mathcal O_E)\), and the differential of a pullback
map on Picard varieties is pullback on this cohomology.
The Hasse invariant \(h\in\{0,1\}\) records whether Frobenius on
this one-dimensional space is zero or nonzero.
\end{context}

\begin{problem}
Compute the algebraic fundamental group of \(E\), and prove
\[
\pi_1(E,O)\cong
\begin{cases}
\displaystyle\prod_{\ell\ \text{prime}}\mathbb Z_\ell^2,
 &\operatorname{char}k=0,\\
\displaystyle\prod_{\ell\ne p}\mathbb Z_\ell^2,
 &\operatorname{char}k=p,\ h=0,\\
\displaystyle\mathbb Z_p\times\prod_{\ell\ne p}\mathbb Z_\ell^2,
 &\operatorname{char}k=p,\ h=1.
\end{cases}
\]
\end{problem}
\end{exerciseintro}

\begin{solution}
Let \(f:Y\to E\) be a connected finite etale cover.
Hurwitz gives \(2g(Y)-2=(\deg f)(2g(E)-2)=0\).
Choose a point of \(Y\) over \(O\) as origin.  Then \(Y\) is an
elliptic curve and \(f\) is an isogeny.
Its fibers are translates of its kernel, so translation by the
kernel gives all deck transformations.  In particular every
connected finite etale cover is Galois and its deck group is
abelian.

For every positive integer \(n\), factor multiplication as
\[
E\xrightarrow{q_n}E_n\xrightarrow{h_n}E,
\qquad [n]=h_nq_n,
\]
with \(q_n\) purely inseparable and \(h_n\) separable, hence etale.
In characteristic zero \(q_n\) is the identity.
These covers are cofinal: if \(f:Y\to E\) has degree \(n\), then
\([n]=f\widehat f\).  The function field of \(Y\) is a separable
intermediate extension of that defined by \([n]\), hence is
contained in its maximal separable subextension \(k(E_n)\).
Thus \(h_n\) factors through \(f\).

The map \(q_n\) is bijective on geometric points and induces a
group isomorphism
\[
E[n](k)\cong\ker(h_n)(k).
\]
Indeed a point belongs to the kernel on the right exactly when
its unique preimage is killed by \([n]\).
If \(n\mid m\), the identity \([m]=[n]\circ[m/n]\) induces the
transition on these groups given by multiplication by \(m/n\).
It follows that
\[
\pi_1(E,O)\cong\varprojlim_n E[n](k).
\tag{1}
\]

For a prime \(\ell\) different from the characteristic,
\([\ell^r]\) is separable of degree \(\ell^{2r}\), so its kernel
has that many points.
This finite abelian group is killed by \(\ell^r\), and its
subgroup killed by \(\ell\) has order \(\ell^2\).
The classification of finite abelian groups consequently gives
\[
E[\ell^r](k)\cong(\mathbb Z/\ell^r\mathbb Z)^2.
\]
Multiplication by \(\ell\) from the next level is surjective:
lift a point along the surjective isogeny \([\ell]\), and note
that the lift is killed by \(\ell^{r+1}\).
One may therefore choose compatible bases noncanonically, giving inverse limit
\(\mathbb Z_\ell^2\).

In characteristic \(p\), consider relative Frobenius
\(F:E\to E^{(p)}\) and its dual \(V:E^{(p)}\to E\).
They have degree \(p\) and \(VF=[p]\).
The differential of \(V\), viewed as pullback on Picard
varieties, is the Frobenius pullback on \(H^1(\mathcal O)\),
with the scalar twist understood.  Thus it is nonzero exactly
when \(h=1\).
If \(h=1\), \(V\) is separable of degree \(p\);
if \(h=0\), \(V\) is purely inseparable.
Since \(F\) is bijective on points, this proves
\[
\#E[p](k)=p\quad(h=1),\qquad
\#E[p](k)=1\quad(h=0).
\]
Separable degrees multiply, so in the first case
\(\#E[p^r](k)=p^r\).  There is only one cyclic factor, since
the subgroup killed by \(p\) has order \(p\).  Hence
\(E[p^r](k)\cong\mathbb Z/p^r\mathbb Z\), compatibly under
multiplication by \(p\), with inverse limit \(\mathbb Z_p\).
In the second case \([p]\) and all its iterates are purely
inseparable, and every \(E[p^r](k)\) is trivial.

Finally, primary decomposition of finite torsion groups and the
Chinese remainder theorem decompose the inverse limit in (1)
as the product of these prime-power inverse limits.
This gives exactly the three groups stated.
\end{solution}

\begin{exerciseintro}
% record: H-IV-04-009
\exerciseheading{Exercise IV.4.9}
\addcontentsline{toc}{subsection}{Exercise IV.4.9}

\begin{context}
Work over an algebraically closed field of characteristic different
from two.  Two elliptic curves are isogenous if there is a finite
surjective morphism between them.
After composing with a translation, such a morphism preserves
origins and is a homomorphism.
A nonconstant homomorphism \(f\) has a dual isogeny \(\widehat f\)
with \(f\widehat f=[\deg f]\).
Multiplication by \(n\ne0\) has degree \(n^2\), so \(E[n](k)\)
is finite.
In characteristic \(p\), any isogeny factors as an iterated
relative Frobenius \(E\to E^{(p^r)}\), followed by a separable
isogeny.  A separable isogeny is the quotient by translations
by its finite kernel; its target function field is the invariant
subfield.
\end{context}

\begin{problem}
(a) Prove that isogeny is an equivalence relation on elliptic curves.

(b) For a fixed elliptic curve \(E\), prove that only countably
many isomorphism classes of elliptic curves are isogenous to \(E\).
\end{problem}
\end{exerciseintro}

\begin{solution}
(a) The identity proves reflexivity.  The composite of two finite
surjective morphisms is again finite and surjective, proving
transitivity.
For symmetry, translate a given finite map to an isogeny
preserving origins and take its dual.  This is a finite map in
the opposite direction.

(b) First count the targets of separable isogenies from a fixed
elliptic curve \(A\).  The torsion group
\[
A(k)_{\mathrm{tors}}=\bigcup_{n\geq1}A[n](k)
\]
is countable, being a countable union of finite sets.
It has only countably many finite subsets, hence only countably
many finite subgroups.
Every separable isogeny \(A\to B\) has one of these groups \(H\)
as its kernel, and \(B\) is determined up to isomorphism by \(H\):
its function field is \(k(A)^H\), and a one-variable function
field determines its smooth projective curve.
Thus there are only countably many such \(B\).

In characteristic zero all isogenies are separable, so this
already proves the assertion.
In characteristic \(p\), every isogeny from \(E\) first passes
through one of the countably many curves \(E^{(p^r)}\), \(r\geq0\),
and then is separable.
Each such twist has only countably many possible separable
quotients by the argument just given.
A countable union of countable sets is countable.
By symmetry from (a), considering maps out of \(E\) includes its
whole isogeny class.
\end{solution}

\begin{exerciseintro}
% record: H-IV-04-010
\exerciseheading{Exercise IV.4.10}
\addcontentsline{toc}{subsection}{Exercise IV.4.10}

\begin{context}
Let \(E\) be an elliptic curve over an algebraically closed field,
with origin \(O\), and let \(p_1,p_2:E\times E\to E\) be the
projections.
The identification \(E\cong\operatorname{Pic}^0(E)\) represents
algebraic families of degree-zero line bundles normalized along
the origin.  Its normalized Poincare bundle is
\[
\mathcal P=\mathcal O_{E\times E}(\Delta)
\otimes p_1^*\mathcal O_E(-O)\otimes p_2^*\mathcal O_E(-O).
\]
Its restriction to \(\{Q\}\times E\) represents
\(\mathcal O_E(Q-O)\), and its restrictions to both axes through
\((O,O)\) are trivial.
A morphism \(E\to E\) fixing \(O\) is a group homomorphism.
Degree is constant in a connected family of line bundles on a
projective curve.  The see-saw principle says that a line bundle
trivial on every fiber of \(p_1\) is pulled back from the first
factor.
\end{context}

\begin{problem}
Writing \(R=\operatorname{End}(E,O)\) with its additive group
structure, construct and prove exactness of
\[
0\longrightarrow\operatorname{Pic}(E)\oplus\operatorname{Pic}(E)
\xrightarrow{(A,B)\mapsto p_1^*A\otimes p_2^*B}
\operatorname{Pic}(E\times E)
\longrightarrow R\longrightarrow0.
\]
\end{problem}
\end{exerciseintro}

\begin{solution}
For \(L\in\operatorname{Pic}(E\times E)\), put
\[
A=L|_{E\times\{O\}},\qquad B=L|_{\{O\}\times E},
\qquad
L_0=L\otimes p_1^*A^{-1}\otimes p_2^*B^{-1}.
\]
Both restrictions of \(L_0\) to the axes are trivial as line
bundles.  Their possible constant one-dimensional fiber factors
do not affect their Picard classes.
The restriction of \(L_0\) to every vertical fiber has degree
zero, because its degree is constant and it is zero on the
fiber over \(O\).

The family of restrictions therefore defines a morphism
\[
\Phi_L:E\longrightarrow\operatorname{Pic}^0(E)\cong E,
\qquad
Q\longmapsto[L_0|_{\{Q\}\times E}].
\]
At \(Q=O\) this is the trivial class, so \(\Phi_L(O)=O\).
Thus \(\Phi_L\in R\).
Normalization and restriction commute with tensor product,
which shows that \(L\mapsto\Phi_L\) is a group homomorphism.

If \(L=p_1^*A\otimes p_2^*B\), its normalized bundle is trivial,
so \(\Phi_L=0\).  Conversely, if \(\Phi_L=0\), every vertical
restriction of \(L_0\) is trivial.
By see-saw, \(L_0=p_1^*C\) for some \(C\in\operatorname{Pic}(E)\).
Restriction to \(E\times\{O\}\) gives \(C\cong\mathcal O_E\).
Hence \(L=p_1^*A\otimes p_2^*B\), proving exactness at the middle
term.

Given \(f\in R\), define
\[
L_f=(f\times\operatorname{id}_E)^*\mathcal P.
\]
It is normalized on both axes, and its restriction to
\(\{Q\}\times E\) represents \(\mathcal O_E(f(Q)-O)\).
Consequently \(\Phi_{L_f}=f\), proving surjectivity.

Finally, if \(p_1^*A\otimes p_2^*B\) is trivial, restricting
successively to the two axes proves \(A\) and \(B\) trivial.
This proves injectivity of the first nonzero map and completes
exactness of the sequence.
\end{solution}

\begin{exerciseintro}
% record: H-IV-04-011
\exerciseheading{Exercise IV.4.11}
\addcontentsline{toc}{subsection}{Exercise IV.4.11}

\begin{context}
Let \(E=\mathbb C/\Lambda\), where
\(\Lambda=\mathbb Z+\mathbb Z\tau\) and \(\operatorname{Im}\tau>0\).
Endomorphisms of the elliptic curve preserving the origin
correspond to complex numbers \(\alpha\) satisfying
\(\alpha\Lambda\subseteq\Lambda\).
A nonzero endomorphism is an unramified covering; its degree is
the number of points in its kernel.
Its dual satisfies \(\widehat f f=[\deg f]\).
A quadratic algebraic integer satisfies a monic quadratic
polynomial over \(\mathbb Z\).
\end{context}

\begin{problem}
(a) If a nonzero endomorphism \(f\) corresponds to multiplication
by \(\alpha\), prove \(\deg f=|\alpha|^2\).

(b) Show that its dual corresponds to \(\overline\alpha\).

(c) If \(\tau\) is an algebraic integer in an imaginary quadratic
field, prove
\(\operatorname{End}(E)=\mathbb Z[\tau]\).
\end{problem}
\end{exerciseintro}

\begin{solution}
(a) The kernel of multiplication by \(\alpha\) is
\(\alpha^{-1}\Lambda/\Lambda\).
Multiplication by \(\alpha\) identifies this finite group with
\(\Lambda/\alpha\Lambda\), so
\[
\deg f=[\Lambda:\alpha\Lambda].
\]
The index is the ratio of the areas of fundamental
parallelograms for \(\alpha\Lambda\) and \(\Lambda\).
As a real-linear transformation of \(\mathbb C\), multiplication
by \(\alpha\) has determinant \(|\alpha|^2\).
Therefore the ratio, and hence \(\deg f\), equals \(|\alpha|^2\).

(b) Write \(\beta\) for the complex scalar corresponding to
\(\widehat f\).  The identity \(\widehat f f=[\deg f]\) gives
\[
\beta\alpha=\deg f=\alpha\overline\alpha.
\]
Since \(\alpha\ne0\), cancellation yields
\(\beta=\overline\alpha\).
The zero endomorphism and its dual both correspond to zero as
well.

(c) Any scalar \(\alpha\) preserving \(\Lambda\) sends \(1\) into
\(\Lambda\), so \(\alpha\in\mathbb Z+\mathbb Z\tau\).
Conversely, integrality of the quadratic number \(\tau\) gives
\[
\tau^2=a\tau+b,\qquad a,b\in\mathbb Z.
\]
Thus multiplication by \(\tau\) preserves \(\Lambda\), and so
does multiplication by every element of
\(\mathbb Z+\mathbb Z\tau=\mathbb Z[\tau]\).
These two inclusions prove the equality.
\end{solution}

\begin{exerciseintro}
% record: H-IV-04-012
\exerciseheading{Exercise IV.4.12}
\addcontentsline{toc}{subsection}{Exercise IV.4.12}

\begin{context}
Write \(E_\tau=\mathbb C/(\mathbb Z+\mathbb Z\tau)\), with the unique
representative satisfying \(\operatorname{Im}\tau>0\),
\(-1/2\leq\operatorname{Re}\tau<1/2\), \(|\tau|\geq1\), and
\(\operatorname{Re}\tau\leq0\) when \(|\tau|=1\).
Lattices define isomorphic curves exactly when homothetic.
Endomorphisms are lattice-preserving scalars \(\alpha\), of degree
\(|\alpha|^2\), with integral trace and determinant.
We use the established facts that degree-two endomorphisms occur
exactly at \(j=1728,8000,-3375\), and order-four origin-preserving
automorphisms exactly at \(j=1728\).
The Weierstrass function's nonzero half-period values are the three
cubic roots; for a conjugation-invariant lattice, conjugating its
argument conjugates its value.
\end{context}

\begin{problem}
(a) Prove that \(E_\tau\) has an origin-preserving automorphism
other than \(1,-1\) only when
\[
\tau=i\quad\text{or}\quad\tau=\frac{-1+\sqrt{-3}}2.
\]

(b) Prove that \(E_\tau\) has an endomorphism of degree two
exactly for
\[
\tau=i,\qquad \tau=\sqrt{-2},\qquad
\tau=\frac{-1+\sqrt{-7}}2,
\]
and match these with the invariants \(1728,8000,-3375\).
\end{problem}
\end{exerciseintro}

\begin{solution}
We first give an elementary lattice observation.
Each of the rings
\[
\mathbb Z[i],\quad \mathbb Z[\sqrt{-2}],\quad
\mathbb Z\left[\frac{-1+\sqrt{-3}}2\right],\quad
\mathbb Z\left[\frac{-1+\sqrt{-7}}2\right]
\tag{1}
\]
is Euclidean for the norm \(|z|^2\).
For the first two, a nearest rectangular-grid point leaves
squared distance at most \(1/2\) and \(3/4\), respectively.
For either of the last two, first approximate the imaginary
coordinate by a multiple of \(\sqrt{m}/2\), then approximate the
real coordinate in that row by an integer, where \(m=3,7\).
The squared error is at most
\[
\frac14+\frac{m}{16}<1.
\]
Approximating a quotient of two ring elements and multiplying
by the denominator gives Euclidean division with smaller norm.
Thus all four rings are principal ideal domains.

If a lattice \(\Lambda\subset\mathbb C\) is stable under one of
these rings \(A=\mathbb Z[\alpha]\), choose \(0\ne\omega\in\Lambda\).
The two vectors \(\omega,\alpha\omega\) span
\(\Lambda\otimes_{\mathbb Z}\mathbb Q\) over \(\mathbb Q\),
since \(\alpha\) is nonreal.
It follows that \(\omega^{-1}\Lambda\) lies in
\(\mathbb Q(\alpha)\).  After clearing the denominators of a
lattice basis, it is a fractional ideal of \(A\).
Since \(A\) is principal, for some \(c\in\mathbb C^\times\) we have
\[
\Lambda=cA.
\tag{2}
\]

(a) An automorphism corresponds to a scalar \(\alpha\) of
absolute value one.  If \(\alpha\) is real, preservation of
\(\Lambda=\mathbb Z+\mathbb Z\tau\) forces
\(\alpha\in\mathbb Z\), so \(\alpha=\pm1\).
Otherwise its integral trace satisfies
\[
\alpha+\overline\alpha=t,\qquad t\in\{-1,0,1\},
\]
and its norm is one.
Thus \(\alpha\) is a primitive fourth, third, or sixth root
of unity.  Its ring \(\mathbb Z[\alpha]\) is respectively
\(\mathbb Z[i]\) or
\(\mathbb Z[(-1+\sqrt{-3})/2]\).
By (2), these give the only two lattice classes.
Their representatives in the specified fundamental region are
the two values stated; each indeed has the indicated extra
root-of-unity automorphisms.

(b) A degree-two endomorphism has norm two.  Its scalar is
nonreal, since a real scalar preserving \(\Lambda\) is an
integer and has square degree.
Its trace is an integer with \(|t|<2\sqrt2\), so
\[
t\in\{-2,-1,0,1,2\},\qquad
\alpha^2-t\alpha+2=0.
\]
Up to replacing \(\alpha\) by its negative or conjugate, the
possibilities are
\[
1+i,\qquad \sqrt{-2},\qquad
\frac{1+\sqrt{-7}}2.
\]
Their rings are the first, second, and fourth rings of (1).
By (2), the corresponding lattices have exactly the three
claimed representatives.  Conversely each displayed scalar
preserves its ring lattice and has norm two, so each curve
does have such an endomorphism.

The square lattice has an automorphism of order four and
therefore \(j=1728\).
For \(\tau=\sqrt{-2}\), the rectangular lattice is invariant
under conjugation.  Each of its nonzero half-periods is fixed
by conjugation modulo the lattice and sign, so the three
Weierstrass branch values are real and distinct.
Ordering them gives a Legendre parameter \(0<\lambda<1\).
In particular \(j(\lambda)>0\).
This lattice class differs from the square one, so its
invariant is not \(1728\); among the three possible invariants
it must be \(8000\).
The remaining, distinct lattice class
\(\tau=(-1+\sqrt{-7})/2\) therefore has \(j=-3375\).
\end{solution}

\begin{exerciseintro}
% record: H-IV-04-013
\exerciseheading{Exercise IV.4.13}
\addcontentsline{toc}{subsection}{Exercise IV.4.13}

\begin{context}
Work over an algebraically closed field of characteristic thirteen.
Every elliptic curve is isomorphic to a Legendre curve
\(y^2=x(x-1)(x-\lambda)\), with \(\lambda\ne0,1\).
For an equation \(y^2=f(x)\) in odd characteristic \(p\), the
Hasse invariant vanishes exactly when the coefficient of
\(x^{p-1}\) in \(f(x)^{(p-1)/2}\) is zero.
The invariant of a Legendre curve is
\[
j=256\frac{(\lambda^2-\lambda+1)^3}{\lambda^2(\lambda-1)^2}.
\]
\end{context}

\begin{problem}
Find all \(j\)-invariants of elliptic curves in characteristic
thirteen whose Hasse invariant is zero.
\end{problem}
\end{exerciseintro}

\begin{solution}
The desired coefficient is that of \(x^{12}\) in
\(x^6(x-1)^6(x-\lambda)^6\).
Expanding the last two factors, it equals
\[
H(\lambda)=\sum_{i=0}^6\binom6i^2\lambda^i.
\]
Modulo thirteen this is
\[
H(\lambda)=
\lambda^6+10\lambda^5+4\lambda^4+10\lambda^3
+4\lambda^2+10\lambda+1.
\]
A direct expansion gives the useful identity
\[
H(\lambda)=(\lambda^2-\lambda+1)^3
-2\lambda^2(\lambda-1)^2.
\]
Since \(256=9\) and \(9\cdot2=5\) in characteristic thirteen,
the Legendre formula implies
\[
j-5=\frac{9H(\lambda)}{\lambda^2(\lambda-1)^2}.
\]
The denominator is nonzero, so the Hasse invariant vanishes
if and only if \(j=5\).

There do exist such curves: the nonconstant polynomial \(H\)
has a root in the algebraically closed field, and
\(H(0)=H(1)=1\), so its roots are allowed Legendre parameters.
Thus \(j=5\) is exactly the one required value.
\end{solution}

\begin{exerciseintro}
% record: H-IV-04-014
\exerciseheading{Exercise IV.4.14}
\addcontentsline{toc}{subsection}{Exercise IV.4.14}

\begin{context}
For an elliptic curve in characteristic \(p\), the Hasse invariant
is zero precisely when absolute Frobenius acts as zero on
\(H^1(E,\mathcal O_E)\).
For a smooth plane cubic \(F=0\), the restriction sequence
identifies
\[
H^1(E,\mathcal O_E)\cong
H^2(\mathbb P^2,\mathcal O_{\mathbb P^2}(-3)).
\]
In the standard three-chart Cech computation the latter space
is one-dimensional, represented by \(1/(xyz)\).
A homogeneous Laurent monomial with a nonnegative exponent
in any coordinate represents zero in this quotient.
We use the prime-progression density theorem: each reduced
residue class modulo three has Dirichlet density \(1/2\).
\end{context}

\begin{problem}
For each prime \(p\ne3\), consider the Fermat elliptic curve
\[
E_p:\ x^3+y^3=z^3
\]
over an algebraic closure of \(\mathbb F_p\).
Determine the primes for which its Hasse invariant is zero,
and determine their density.
\end{problem}
\end{exerciseintro}

\begin{solution}
The partial derivatives of \(F=x^3+y^3-z^3\) have no common
projective zero when \(p\ne3\), so \(E_p\) is smooth.
It has the rational point \([1:0:1]\), hence is elliptic.
This also applies when \(p=2\).

Under the stated cohomology identification, Frobenius sends
the class of \(1/(xyz)\) to the class of
\[
\frac{F^{p-1}}{x^py^pz^p}.
\tag{1}
\]
Indeed Frobenius first raises the Cech representative to the
\(p\)-th power, putting it in degree \(-3p\).  Multiplication
by \(F^{p-1}\) returns to degree \(-3\).
This multiplication is the map on the ideal term of the
restriction sequence that makes the Frobenius square commute,
since \(F\cdot F^{p-1}=F^p\).

Only the coefficient of \(x^{p-1}y^{p-1}z^{p-1}\) in \(F^{p-1}\)
survives in (1): to have three negative exponents after division,
each numerator exponent must be at most \(p-1\), and their sum
is \(3(p-1)\), forcing equality for all three.
Every monomial of \(F^{p-1}\) has exponents divisible by three.
Thus the surviving coefficient is zero unless \(p\equiv1\pmod3\).

If \(p=3m+1\), that coefficient is
\[
(-1)^m\frac{(p-1)!}{(m!)^3},
\]
which is nonzero modulo \(p\), since every factorial involved
has argument less than \(p\).
It follows that the Hasse invariant vanishes exactly for
\[
p\equiv2\pmod3,
\]
including \(p=2\).
The density theorem for the two reduced classes modulo three
gives this set Dirichlet density \(1/2\).
Its natural density among primes is also \(1/2\) if one uses
the prime number theorem for arithmetic progressions.
\end{solution}

\begin{exerciseintro}
% record: H-IV-04-015
\exerciseheading{Exercise IV.4.15}
\addcontentsline{toc}{subsection}{Exercise IV.4.15}

\begin{context}
Let \(E\) be an elliptic curve over an algebraically closed field
of characteristic \(p>2\), as in the section's characteristic
convention.
Let \(E_p\) denote the inverse Frobenius twist, and
\(F':E_p\to E\) the \(k\)-linear Frobenius, of degree \(p\).
It is purely inseparable and bijective on geometric points.
Write \(V=\widehat{F'}:E\to E_p\).
Duality gives \(\deg V=p\) and \(F'V=[p]_E\).

The tangent space of \(\operatorname{Pic}^0(E)\) is
\(H^1(E,\mathcal O_E)\); the differential of a map induced by
pullback is the corresponding pullback on this cohomology.
The Hasse invariant \(h\in\{0,1\}\) is zero or one according
as the semilinear absolute Frobenius on this one-dimensional
space is zero or nonzero.
A homomorphism of elliptic curves is separable exactly when
its differential at the origin is nonzero.
\end{context}

\begin{problem}
Prove that \(V\) is separable if and only if \(h=1\).
Deduce that the group of geometric points killed by \(p\) is
\[
E[p](k)\cong
\begin{cases}
\mathbb Z/p\mathbb Z,&h=1,\\
0,&h=0.
\end{cases}
\]
\end{problem}
\end{exerciseintro}

\begin{solution}
By definition of the dual, \(V\) represents the pullback
\[
(F')^*:\operatorname{Pic}^0(E)\longrightarrow
\operatorname{Pic}^0(E_p).
\]
Its tangent map is therefore
\[
(F')^*:H^1(E,\mathcal O_E)\longrightarrow
H^1(E_p,\mathcal O_{E_p}).
\]
Undoing the scalar twist identifies its vanishing or
nonvanishing with that of absolute Frobenius on
\(H^1(E,\mathcal O_E)\).
Thus \(dV\ne0\) precisely when \(h=1\), proving the first
assertion by the separability criterion.

If \(h=1\), the degree-\(p\) homomorphism \(V\) is etale:
translation carries its nonzero differential at the origin
to every point.  Its kernel consists of \(p\) geometric
points and is a group of prime order, hence cyclic.
Since \(F'\) has only the origin in its geometric kernel,
\[
[p](Q)=F'(V(Q))=O\quad\Longleftrightarrow\quad V(Q)=O_p.
\]
Thus \(E[p](k)=\ker V(k)\cong\mathbb Z/p\mathbb Z\).

If \(h=0\), \(V\) is inseparable.  Its degree is the prime \(p\),
so it is purely inseparable and has a one-point geometric
kernel.  The same equivalence then gives \(E[p](k)=0\).
This statement concerns geometric points; the finite group
scheme \(E[p]\) still has length \(p^2\).
\end{solution}

\begin{exerciseintro}
% record: H-IV-04-016
\exerciseheading{Exercise IV.4.16}
\addcontentsline{toc}{subsection}{Exercise IV.4.16}

\begin{context}
Let \(E\) be an elliptic curve defined over \(\mathbb F_q\),
where \(q=p^r\), \(p>2\), and choose its origin in \(E(\mathbb F_q)\).
We work over an algebraic closure \(k\).
Let \(F':E_q\to E\) denote \(k\)-linear Frobenius from the
inverse \(q\)-power twist.
For an isogeny \(u\), duality satisfies
\[
u\widehat u=[\deg u],\qquad
\widehat{u+v}=\widehat u+\widehat v,\qquad
\widehat{[m]}=[m],\qquad\deg[m]=m^2.
\]
The action of \(\mathbb Z\) on \(E\) is faithful.
The degree of the zero homomorphism is defined to be zero.
A separable isogeny has degree equal to the number of its
geometric kernel points.
The dual of degree-\(p\) Frobenius is separable exactly when
the Hasse invariant is one.
\end{context}

\begin{problem}
(a) Identify \(E_q\) with \(E\) over \(k\) and describe \(F'\)
on projective coordinates.

(b) Prove that \(1-F'\) is separable and that its kernel is
\(E(\mathbb F_q)\).

(c) Write \(N=\#E(\mathbb F_q)\).  Prove that
\[
F'+\widehat{F'}=[a],\qquad N=q+1-a
\]
for an integer \(a\).

(d) Use nonnegativity of \(\deg([m]+[n]F')\) to prove
\(|a|\leq2\sqrt q\).

(e) When \(q=p\), show that the Hasse invariant is zero if and
only if \(a\equiv0\pmod p\).  Deduce that for \(p\geq5\) this
is equivalent to \(N=p+1\).
\end{problem}
\end{exerciseintro}

\begin{solution}
(a) Forming the inverse \(q\)-power twist replaces the
coefficients of a defining equation by their \(q\)-th roots.
Coefficients in \(\mathbb F_q\) are unchanged.
Thus \(E_q\cong E\), and \(F'\), henceforth denoted \(F\), is
\[
[X:Y:Z]\longmapsto[X^q:Y^q:Z^q].
\]
It fixes the chosen rational origin, has degree \(q\), and has
zero differential.

(b) The differential of \(1-F\) is the identity, so this
homomorphism is nonconstant and separable.
Its kernel consists of the points fixed by \(F\).
In any projective chart with one nonzero coordinate normalized
to one, fixedness says that the other coordinates satisfy
\(u^q=u\), hence lie in \(\mathbb F_q\).
Therefore the kernel is precisely \(E(\mathbb F_q)\), and
\(\deg(1-F)=N\).

(c) Put \(V=\widehat F\); then \(FV=VF=[q]\).
Duality and part (b) give
\[
[N]=(1-F)(1-V)=[q+1]-(F+V).
\]
It follows that \(F+V=[q+1-N]\).
Taking \(a=q+1-N\in\mathbb Z\) proves both assertions.

(d) For integers \(m,n\), additivity of duality gives
\[
\begin{split}
([m]+[n]F)([m]+[n]V)
&=[m^2]+[mn](F+V)+[n^2]FV\\
&=[m^2+amn+qn^2].
\end{split}
\]
Comparison with \(u\widehat u=[\deg u]\), including \(u=0\),
and faithfulness of integer multiplication imply
\[
\deg([m]+[n]F)=m^2+amn+qn^2\geq0.
\tag{1}
\]
If \(a^2>4q\), the real quadratic \(t^2+at+q\) is negative
on a nonempty open interval between its roots.
Choose a rational \(m/n\) in that interval.  Multiplying by
\(n^2>0\) contradicts (1).
Hence \(a^2\leq4q\), or \(|a|\leq2\sqrt q\).
Nonnegativity is the necessary assertion; strict positivity
cannot hold for the zero homomorphism.

(e) Suppose \(q=p\).
Differentiating \(F+V=[a]\) gives \(dV=a\), since \(dF=0\).
The Hasse invariant is zero exactly when \(V\) is inseparable,
equivalently when \(dV=0\), hence exactly when \(p\mid a\).
For \(p\geq5\), the strict inequality \(2\sqrt p<p\), together
with the bound in (d), shows that the only possible multiple
of \(p\) is \(a=0\).
By (c), this is equivalent to \(N=p+1\).
\end{solution}

\begin{exerciseintro}
% record: H-IV-04-017
\exerciseheading{Exercise IV.4.17}
\addcontentsline{toc}{subsection}{Exercise IV.4.17}

\begin{context}
Let \(E/\mathbb Q\) have equation \(y^2+y=x^3-x\), with origin
the point \(O=[0:1:0]\) at infinity.
Three intersections of a line with this cubic, counted with
multiplicity, sum to \(O\).
Inversion is \((x,y)\mapsto(x,-1-y)\).
The projective equation over any field is
\[
F(X,Y,Z)=Y^2Z+YZ^2-X^3+XZ^2=0.
\]
We use the Jacobian criterion to check smoothness.
\end{context}

\begin{problem}
(a) Put \(P=(0,0)\).  Compute \(P+Q\) for \(Q=(a,b)\) and
list the coordinates of \(nP\) for \(1\leq n\leq10\), including
the exceptional cases in the addition formula.

(b) Prove that the displayed integral equation has nonsingular
reduction modulo every prime other than \(37\).
\end{problem}
\end{exerciseintro}

\begin{solution}
(a) If \(a\ne0\), the line through \(P\) and \(Q\) has equation
\(y=mx\), where \(m=b/a\).
Substitution into the cubic gives
\[
x\bigl(x^2-m^2x-(m+1)\bigr)=0.
\]
The third intersection has \(x\)-coordinate \(m^2-a\).
Negating it gives
\[
P+Q=\left(m^2-a,\ -1-m(m^2-a)\right),
\qquad m=\frac ba.
\tag{1}
\]
If \(Q=-P=(0,-1)\), then \(P+Q=O\).
If \(Q=P\), the tangent slope is
\[
m=\frac{3x^2-1}{2y+1}\bigg|_{(0,0)}=-1,
\]
and the tangent computation gives \(2P=(1,0)\).
Finally \(P+O=P\).
These exhaust the cases with \(a=0\) or \(Q\) at infinity.

Iterating (1), starting with the stated tangent value, yields
\[
\begin{array}{c|r|r}
n&x(nP)&y(nP)\\
1&0&0\\
2&1&0\\
3&-1&-1\\
4&2&-3\\
5&1/4&-5/8\\
6&6&14\\
7&-5/9&8/27\\
8&21/25&-69/125\\
9&-20/49&-435/343\\
10&161/16&-2065/64.
\end{array}
\]
In particular the sixth multiple is the claimed \((6,14)\).
Each row is obtained from the preceding one by the rational
operations in (1), so the table also verifies membership in
the subgroup generated by \(P\).

(b) At infinity \(Z=0\) forces \(X=0\), so the only point is
\(O=[0:1:0]\).  There \(F_Z=1\), and hence it is smooth in
every characteristic.
In the affine chart the partial derivatives of
\(y^2+y-x^3+x\) are
\[
1-3x^2,\qquad 2y+1.
\]
In characteristic two the second is one; in characteristic
three the first is one.  Thus both these reductions are smooth.
For a prime other than two or three, a singular point must have
\[
y=-\frac12,\qquad x^2=\frac13.
\]
The original equation then becomes
\(-1/4=-2x/3\), so \(x=3/8\).
Combining with \(x^2=1/3\) gives \(27=64\), which holds only
in characteristic \(37\).
This proves smoothness at every other prime.
Conversely in characteristic \(37\), the point
\((x,y)=(5,18)\) satisfies the equation and both derivatives,
so that prime really is exceptional.
\end{solution}

\begin{exerciseintro}
% record: H-IV-04-018
\exerciseheading{Exercise IV.4.18}
\addcontentsline{toc}{subsection}{Exercise IV.4.18}

\begin{context}
On the elliptic curve \(E/\mathbb Q\) given by
\(y^2=x^3-7x+10\), the origin is the point at infinity and
\(-(x,y)=(x,-y)\).
For finite points \(A=(x_1,y_1)\), \(B=(x_2,y_2)\) with
\(B\ne-A\), define
\[
m=\frac{y_2-y_1}{x_2-x_1}\quad(x_1\ne x_2),\qquad
m=\frac{3x_1^2-7}{2y_1}\quad(A=B).
\]
Then the sum has coordinates
\[
x_3=m^2-x_1-x_2,\qquad
y_3=-y_1+m(x_1-x_3).
\]
\end{context}

\begin{problem}
Find at least twenty-six distinct integral points on \(E\).
For each, verify membership in the subgroup of \(E(\mathbb Q)\)
generated by \(P=(1,2)\) and \(Q=(2,2)\).
\end{problem}
\end{exerciseintro}

\begin{solution}
The following table lists thirteen points with positive
\(y\)-coordinate and expressions for them in terms of \(P,Q\):
\[
\begin{array}{r|r|l}
x&y&\text{group expression}\\
-3&2&-P-Q\\
-2&4&4P+Q\\
-1&4&-2P\\
1&2&P\\
2&2&Q\\
3&4&-2P-Q\\
5&10&3P+Q\\
9&26&-3P\\
13&46&P-Q\\
31&172&2P+2Q\\
41&262&-5P-2Q\\
67&548&6P+Q\\
302&5248&-4P+Q.
\end{array}
\tag{1}
\]
Negating these expressions supplies thirteen further points
with the same \(x\)-coordinates and opposite \(y\)-coordinates.
All twenty-six points are distinct, since the listed
\(x\)-coordinates are distinct and none of the listed \(y\)'s
is zero.

Here are intermediate additions verifying the larger entries
without an integral-point search assumption.
The addition formulas give
\[
2P=(-1,-4),\quad 3P=(9,-26),\quad
4P=(9/4,19/8),\quad P+Q=(-3,-2),\quad P-Q=(13,46).
\]
Continuing with the same formulas gives
\[
2P+Q=(3,-4),\qquad 3P+Q=(5,10),\qquad
4P+Q=(-2,4).
\]
Doubling \(P+Q\) gives
\[
2P+2Q=(31,172).
\]
Adding \(P\) to \(4P+Q=(-2,4)\) first yields
\[
5P+Q=\left(\frac{13}{9},-\frac{46}{27}\right).
\]
Adding \(Q\) to this point gives \((41,-262)\), so
\[
5P+2Q=(41,-262).
\tag{2}
\]
Adding \(P\) instead gives
\[
6P+Q=(67,548).
\]
Finally, \(4P=(9/4,19/8)\) gives
\(-4P+Q=(302,5248)\) by direct addition with \(Q\).

Negating (2) gives the positive-\(y\) entry at \(x=41\):
\[
(41,-262)=5P+2Q,\qquad (41,262)=-5P-2Q.
\]
The other expressions in (1) follow from the displayed
intermediate additions or their negatives.
Substitution of each pair into \(y^2=x^3-7x+10\) is an
independent integer check of its coordinates.
This constructs the requested twenty-six points; it does not
assert that no other integral or rational points exist.
\end{solution}

\begin{exerciseintro}
% record: H-IV-04-019
\exerciseheading{Exercise IV.4.19}
\addcontentsline{toc}{subsection}{Exercise IV.4.19}

\begin{context}
Let \(E/\mathbb Q\) be an elliptic curve with a proper integral
plane model, and let \(T\subset\operatorname{Spec}\mathbb Z\)
be its open set of good reduction with the prime two removed.
Write \(\mathcal E\to T\) for the resulting smooth proper
genus-one model.
Its rational origin extends by the valuative criterion of
properness.  A smooth proper genus-one family with a section
has a unique elliptic-scheme group law compatible with that
section; it is obtained from its relative degree-zero Picard
functor.

We use these facts: multiplication by a nonzero integer on a
geometric elliptic curve is finite of degree \(n^2\);
a proper quasi-finite morphism is finite; a finite dominant
morphism between regular equidimensional schemes of the same
dimension is flat (the finite case of miracle flatness).
A finite flat morphism with etale geometric fibers is etale.
Over the complete local ring \(\mathbb Z_p\), sections of a
finite etale scheme are uniquely determined by their reductions.
\end{context}

\begin{problem}
For \(n\ne0\), prove that \([n]:E\to E\) extends over \(T\)
to a finite flat morphism \([n]:\mathcal E\to\mathcal E\),
and that its kernel is flat over \(T\).
For a prime \(p\in T\) with \((n,p)=1\), prove that reduction
injects \(E(\mathbb Q)[n]\) into \(\mathcal E_p(\mathbb F_p)\).
Apply this to prove that the rational torsion subgroups of
\[
y^2+y=x^3-x,\qquad y^2=x^3-7x+10
\]
are trivial.
\end{problem}
\end{exerciseintro}

\begin{solution}
The group law of \(\mathcal E/T\) defines \([n]\) by repeated
addition and inversion, and it restricts to the original
multiplication map on the generic fiber.
For \(n\ne0\), every geometric fiber of this morphism is finite.
Since the source is proper over \(T\) and the target separated
over \(T\), the morphism is proper.  It is therefore finite.

The scheme \(\mathcal E\) is regular of dimension two, being
smooth over the regular one-dimensional base \(T\).
The multiplication map is surjective and has finite fibers.
Miracle flatness applies to its source and target local rings,
so \([n]\) is flat.  It is consequently finite locally free;
its rank is \(n^2\), as may be read on any geometric fiber.
The kernel is its base change along the identity section:
\[
\mathcal E[n]=\mathcal E\times_{\mathcal E,[n]}T.
\]
It is finite locally free of rank \(n^2\), hence flat over \(T\).

Fix \(p\in T\) with \(p\nmid n\).
On the base change to \(\mathbb Z_p\), the differential of
\([n]\) on each fiber is multiplication by the invertible
scalar \(n\).
Its geometric fibers are etale, so \(\mathcal E[n]\) is
finite etale over \(\mathbb Z_p\).
Every rational point of \(E\) extends to a \(\mathbb Z_p\)-point
by properness.  Its reduction is compatible with the group law.
If two rational \(n\)-torsion points have equal reduction,
their extended sections of \(\mathcal E[n]\) coincide by
uniqueness of finite-etale lifting over \(\mathbb Z_p\).
Their generic points, and hence the original rational points,
are equal.  This proves injectivity.

For the first displayed curve, direct counting gives
\[
\#E(\mathbb F_3)=7,\qquad \#E(\mathbb F_5)=8.
\]
For the second, the corresponding counts are \(7\) and \(10\).
Both curves have good reduction at three and five, as the
Jacobian criterion or their discriminants show.
For a prime \(\ell\ne3,5\), rational \(\ell\)-torsion injects
into both reduction groups; their relatively prime orders
exclude such torsion.
For \(\ell=3\), reduction at five excludes it, since neither
eight nor ten is divisible by three.
For \(\ell=5\), reduction at three excludes it, since seven is
not divisible by five.
Thus neither curve has a rational point of any prime order.
Any nonzero torsion point would have a multiple of prime order,
so both rational torsion subgroups are trivial.

The restriction \(n\ne0\) is necessary for flatness of the
multiplication morphism: \([0]\) maps each elliptic fiber to its
origin and is not flat.  Its kernel is the whole family, which
is nevertheless flat over \(T\).
\end{solution}

\begin{exerciseintro}
% record: H-IV-04-020
\exerciseheading{Exercise IV.4.20}
\addcontentsline{toc}{subsection}{Exercise IV.4.20}

\begin{context}
Let \(E\) be an elliptic curve over an algebraically closed field
of characteristic \(p>2\), and put \(R=\operatorname{End}(E,O)\).
The inverse Frobenius twist \(E_p\) is obtained by replacing
coefficients of an equation by their \(p\)-th roots.
The \(k\)-linear Frobenius \(F':E_p\to E\) is purely inseparable
of degree \(p\).
Dual isogenies have the same degree, and
\(f\widehat f=\widehat f f=[\deg f]\).
The dual of Frobenius is separable if and only if the Hasse
invariant is one.  In particular, \([p]\) is purely inseparable
if the invariant is zero and has separable degree \(p\) otherwise.

For a smooth curve over a perfect field, its purely inseparable
degree-\(p^r\) quotient is uniquely its \(r\)-fold Frobenius
twist, up to isomorphism.  This follows from
\([k(E):k(E)^{p^r}]=p^r\).
The invariant \(j\) is a rational function of Weierstrass
coefficients over the prime field and classifies elliptic
curves over \(k\).
The tangent space of \(\operatorname{Pic}^0(E)\) is
\(H^1(E,\mathcal O_E)\).
\end{context}

\begin{problem}
(a) Prove \(j(E_p)=j(E)^{1/p}\), and conclude
\(E\cong E_p\) if and only if \(j(E)\in\mathbb F_p\).

(b) Show that \([p]\) factors as \(\pi\widehat\pi\) for an
endomorphism \(\pi\) of degree \(p\) if and only if
\(E\cong E_p\).  In this case show that the Hasse invariant is
zero if and only if \(\pi,\widehat\pi\) differ by an automorphism.

(c) If the Hasse invariant is zero, prove
\(j(E)\in\mathbb F_{p^2}\).

(d) Let \(\lambda_f\) be the scalar by which \(f^*\) acts on
\(H^1(E,\mathcal O_E)\).
Show that \(f\mapsto\lambda_f\) is a ring homomorphism
\(R\to k\), and that its image is \(\mathbb F_p\) when the
Hasse invariant is nonzero.  Deduce that \(R\) has a prime
ideal with quotient \(\mathbb F_p\).
\end{problem}
\end{exerciseintro}

\begin{solution}
(a) Taking \(p\)-th roots of the coefficients commutes with
the rational formula for \(j\), because its constants lie in
the prime field.  Thus \(j(E_p)=j(E)^{1/p}\).
The two curves are isomorphic exactly when these invariants
are equal, equivalently \(j(E)^p=j(E)\), which means
\(j(E)\in\mathbb F_p\).

(b) If \(E\cong E_p\), compose such an origin-preserving
isomorphism with \(F'\) to obtain a degree-\(p\)
endomorphism \(\pi\) of \(E\).
Duality gives \(\pi\widehat\pi=[p]\).

Conversely, in a factorization into two degree-\(p\) maps,
at least one factor is inseparable, since \([p]\) is inseparable.
An inseparable degree-\(p\) map is purely inseparable.
The uniqueness of the purely inseparable quotient therefore
implies \(E\cong E^{(p)}\), equivalently \(E\cong E_p\).

If the Hasse invariant is zero, \([p]\) is purely inseparable,
so both factors are purely inseparable of degree \(p\).
As maps with the same source, they have the same function-field
image \(k(E)^p\), and differ by an isomorphism of their targets.
Thus \(\pi\) and \(\widehat\pi\) differ by a unit of \(R\).
If the invariant is one, the separable degree of \([p]\) is
\(p\), so one factor is separable and the other is purely
inseparable.  Composing with an automorphism cannot change
separability, so they cannot be associates.

(c) When the invariant is zero, the degree-\(p^2\) map \([p]\)
is purely inseparable.  Uniqueness of this quotient gives
\(E\cong E^{(p^2)}\).
Their invariants satisfy \(j(E)=j(E)^{p^2}\).
The roots of \(T^{p^2}-T\) in \(k\) are exactly
\(\mathbb F_{p^2}\), proving the claim.

(d) Pullback reverses composition on \(H^1\), but the resulting
endomorphisms of its one-dimensional vector space are scalars
and commute.  Thus
\(\lambda_{fg}=\lambda_f\lambda_g\).
Additivity follows from the Picard interpretation:
\(\widehat{f+g}=\widehat f+\widehat g\), and the differential of
addition of group homomorphisms is the sum of their differentials.
The differential of \(\widehat f\) is exactly \(f^*\) on \(H^1\).
Also \(\lambda_1=1\), so this is a unital ring homomorphism.

Absolute Frobenius \(F:E\to E\), although not \(k\)-linear,
commutes with every morphism of schemes by functoriality.
On cohomology this gives \(f^*F^*=F^*f^*\).
Choose a nonzero vector \(v\) and write \(F^*v=cv\).
If the Hasse invariant is nonzero, then \(c\ne0\), and
semilinearity gives
\[
c\lambda_fv=f^*F^*v=F^*f^*v=\lambda_f^pcv.
\]
Hence \(\lambda_f^p=\lambda_f\), so every \(\lambda_f\) lies in
\(\mathbb F_p\).
The image contains \(\mathbb F_p\) because the homomorphism
is unital.  Its kernel is consequently a maximal, hence prime,
two-sided ideal with quotient \(\mathbb F_p\).
\end{solution}

\begin{exerciseintro}
% record: H-IV-04-021
\exerciseheading{Exercise IV.4.21}
\addcontentsline{toc}{subsection}{Exercise IV.4.21}

\begin{context}
Let \(K\) be a quadratic number field and \(\mathcal O_K\) its
ring of integers.
It has an integral basis \(1,\omega\), so
\(\mathcal O_K=\mathbb Z\oplus\mathbb Z\omega\), with
\(\omega^2=a\omega+b\) for integers \(a,b\).
Subrings are required to contain the identity, hence contain
\(\mathbb Z\).
Every nonzero additive subgroup of \(\mathbb Z\) is \(f\mathbb Z\)
for a unique positive integer \(f\).
\end{context}

\begin{problem}
Prove that every subring \(R\subseteq\mathcal O_K\), \(R\ne\mathbb Z\),
has the form
\[
R=\mathbb Z+f\mathcal O_K
\]
for a unique integer \(f\geq1\), called its conductor.
\end{problem}
\end{exerciseintro}

\begin{solution}
Project the additive group of \(\mathcal O_K\) onto
\(\mathcal O_K/\mathbb Z\cong\mathbb Z\), using the coefficient
of \(\omega\).
The image of \(R\) is nonzero because \(R\ne\mathbb Z\).
It is therefore \(f\mathbb Z\) for a unique \(f\geq1\).
Choose \(c+f\omega\in R\) lifting its generator.
Since \(c\in\mathbb Z\subset R\), subtraction shows \(f\omega\in R\).
Thus
\[
\mathbb Z+f\mathbb Z\omega\subseteq R.
\]
The reverse inclusion follows from the definition of the
image \(f\mathbb Z\): every element of \(R\) has its
\(\omega\)-coefficient divisible by \(f\).
Hence
\[
R=\mathbb Z+f\mathbb Z\omega=\mathbb Z+f\mathcal O_K.
\]
The subgroup image determines \(f\) uniquely.

Conversely this additive subgroup is a ring, since
\((f\omega)^2=f^2b+fa(f\omega)\) belongs to it.
Thus all positive \(f\) do produce subrings.
The usual conductor ideal also equals \(f\mathcal O_K\):
if \(r=c+df\omega\in R\) satisfies \(r\mathcal O_K\subset R\),
then the coefficient of \(\omega\) in
\(r\omega=dfb+(c+dfa)\omega\) must be divisible by \(f\).
This forces \(f\mid c\), hence \(r\in f\mathcal O_K\).
The reverse inclusion is immediate.
\end{solution}

\begin{exerciseintro}
% record: H-IV-04-022
\exerciseheading{Exercise IV.4.22}
\addcontentsline{toc}{subsection}{Exercise IV.4.22}

\begin{context}
A family of elliptic curves with section over
\(S=\mathbb A^1_{\mathbb C}\) means a smooth proper morphism
\(\mathcal E\to S\) with geometrically connected genus-one
fibers and a chosen section \(O\).
It carries an elliptic-scheme group law.
The kernel \(\mathcal E[2]\) is finite etale of rank four.
The quotient by inversion is a projective-line bundle \(Y\to S\);
the quotient map \(\mathcal E\to Y\) is finite flat of degree two
and branched at the images of the four two-torsion points.
It can be constructed by the relative linear system \(|2O|\).

We use Hurwitz for finite maps of smooth projective complex
curves and \(\operatorname{Pic}(\mathbb P^1\times S)\cong\mathbb Z\),
generated by the degree-one bundle on \(\mathbb P^1\).
For a finite flat double cover in characteristic different from
two, the trace decomposition of its algebra is
\(\mathcal O_Y\oplus M^{-1}\); multiplication on the second
summand is a section of \(M^2\), whose zero divisor is the
branch divisor.
\end{context}

\begin{problem}
Prove that every family of elliptic curves with a section over
\(\mathbb A^1_{\mathbb C}\) is trivial: it is isomorphic, with
section, to \(E_0\times\mathbb A^1_{\mathbb C}\) for a fixed
complex elliptic curve \(E_0\).
\end{problem}
\end{exerciseintro}

\begin{solution}
First, every connected finite etale cover \(Z\to\mathbb A^1_{\mathbb C}\)
has degree one.
To prove this, compactify it to a finite map
\(\overline Z\to\mathbb P^1\) of degree \(d\), with
\(\overline Z\) smooth and projective.
All ramification is over infinity.
If there are \(r\geq1\) points over infinity, their total
ramification contribution is \(d-r\), since the characteristic
is zero.  Hurwitz gives
\[
2g(\overline Z)-2=-2d+(d-r)=-d-r.
\]
As \(g(\overline Z)\geq0\), this forces \(d+r\leq2\), hence
\(d=r=1\).

It follows that \(\mathcal E[2]\) is the disjoint union of four
sections.  Label them \(O,T_1,T_2,T_3\).
Their images in the quotient line bundle \(Y\) are disjoint
sections.
The first three determine a unique trivialization
\[
Y\cong\mathbb P^1\times S
\]
carrying their images to \(\infty,0,1\), respectively.
Indeed the unique projective linear transformation carrying
three distinct points to these values varies regularly in a
local trivialization, and uniqueness makes these
transformations glue.
The fourth image is consequently a morphism
\[
\lambda:S\longrightarrow\mathbb A^1\setminus\{0,1\}.
\]
It is a polynomial with no zero on \(\mathbb C\), so it is
constant.  Write its value as \(\lambda_0\ne0,1\).
The branch divisor is therefore the fixed divisor
\[
B=(\infty)+(0)+(1)+(\lambda_0)
\]
pulled back from \(\mathbb P^1\).

It remains to rule out a varying double-cover twist, rather
than merely conclude that the fibers are isomorphic.
By the trace decomposition, the double cover is described by
a line bundle \(M\) on \(\mathbb P^1\times S\) and a section
\(b\in H^0(M^2)\) with divisor \(B\times S\).
In particular \(M^2\cong\mathcal O(4)\).
Since the Picard group of this product is \(\mathbb Z\),
we have \(M\cong\mathcal O(2)\), pulled back from \(\mathbb P^1\).
Under this identification, \(b\) differs from the fixed binary
quartic defining \(B\) by a global unit.
The global units on \(\mathbb P^1\times\mathbb A^1_{\mathbb C}\)
are \(\mathbb C^\times\), so that unit is a constant \(c\).
Choose a square root of \(c\) and rescale the trace-zero
generator.  This identifies the double-cover algebra with
the pullback of the fixed double-cover algebra on \(\mathbb P^1\).

Thus \(\mathcal E\cong E_0\times S\), where \(E_0\) is the
double cover branched at \(0,1,\lambda_0,\infty\), equivalently
the Legendre elliptic curve with parameter \(\lambda_0\).
The chosen section \(O\) is the unique ramification section
above infinity, so this isomorphism preserves the section as
well as the family.
\end{solution}

\begin{exerciseintro}
\corpussection{IV.5 The Canonical Embedding}
\addcontentsline{toc}{section}{IV.5 The Canonical Embedding}

% record: H-IV-05-001
\exerciseheading{Exercise IV.5.1}
\addcontentsline{toc}{subsection}{Exercise IV.5.1}

\begin{context}
Work over an algebraically closed field. A hyperelliptic curve is
a nonsingular projective integral curve of genus \(g\geq2\) with
a degree-two morphism to \(\mathbb P^1\); its canonical map factors
through that morphism and hence is not an embedding.
For a nonsingular complete intersection of hypersurfaces of degrees
\(d_1,\ldots,d_{n-1}\) in \(\mathbb P^n\), adjunction gives
\(\omega_X=\mathcal O_X(\sum d_i-n-1)\).
The canonical degree is \(2g-2\), and a positive tensor power of a
very ample line bundle is very ample.
\end{context}

\begin{problem}
Prove that a hyperelliptic curve is not a complete intersection in
any projective space.
\end{problem}
\end{exerciseintro}

\begin{solution}
Suppose an embedding made \(X\) such a complete intersection.
Set \(m=\sum d_i-n-1\), so adjunction gives
\(\omega_X=\mathcal O_X(m)\). Taking degrees yields
\[
2g-2=m\deg\mathcal O_X(1).
\]
Both the left side and \(\deg\mathcal O_X(1)\) are positive,
so \(m>0\). Consequently \(\omega_X\) is very ample.
Its complete linear system would embed \(X\), contradicting the
degree-two factorization of the canonical map of a hyperelliptic
curve. Since the argument applies to every embedding and every
ambient projective dimension, no complete-intersection embedding
exists.
\end{solution}

\begin{exerciseintro}
% record: H-IV-05-002
\exerciseheading{Exercise IV.5.2}
\addcontentsline{toc}{subsection}{Exercise IV.5.2}

\begin{context}
Let \(X\) be a nonsingular projective integral curve over an
algebraically closed field, of genus \(g\geq2\).
A line bundle of degree at least \(2g+1\) is very ample.
The tangent sheaf is \(T_X=\omega_X^{-1}\), and a negative-degree
line bundle has no nonzero global section.
The subgroup scheme of \(\operatorname{PGL}_{N+1}\) preserving
a fixed projective subscheme is closed and of finite type:
preservation of finitely many homogeneous ideal generators gives
its defining equations.
Its tangent vectors at the identity induce derivations on that
subscheme. We use that local dimension is at most tangent-space
dimension and that a zero-dimensional finite-type scheme has
finitely many geometric points.
\end{context}

\begin{problem}
In characteristic zero, prove that \(\operatorname{Aut}(X)\) is
finite. In fact, prove this assertion in arbitrary characteristic.
\end{problem}
\end{exerciseintro}

\begin{solution}
The line bundle \(L=\omega_X^{\otimes3}\) has degree \(6g-6\),
which is at least \(2g+1\) for every \(g\geq2\).
Embed \(X\) by the complete system \(|L|\) into \(\mathbb P^N\).
Every automorphism acts functorially on \(H^0(X,L)\), hence is
induced by a projective transformation. Conversely, a projective
transformation preserving \(X\) restricts to an automorphism.
These restrictions are faithful. Indeed, if a matrix fixes every
point of \(X\), each such point is represented by an eigenvector.
The finitely many eigenspaces cover the integral curve, so one
contains it. Since the complete embedding is nondegenerate,
that eigenspace is the whole vector space and the matrix is scalar.
Thus \(\operatorname{Aut}(X)\) is the group of geometric points
of the closed stabilizer \(G\subset\operatorname{PGL}_{N+1}\).

The tangent action at the identity gives a linear map
\[
T_1G\longrightarrow H^0(X,T_X).
\]
It is injective by the same eigenvector argument: a tangent vector
is represented by a matrix \(A\) modulo scalar matrices, and its
projective vector field is zero at \([v]\) exactly when
\(Av\in kv\). If its restriction to \(X\) is zero, the preceding
argument forces \(A\) scalar.
But \(\deg T_X=2-2g<0\), so \(H^0(X,T_X)=0\).
Hence \(T_1G=0\), and the local dimension of \(G\) at the identity
is zero. Translation gives the same conclusion at every geometric
point. Therefore \(G\) is zero-dimensional and of finite type,
and has finitely many geometric points. This proves the assertion
without a characteristic restriction.
\end{solution}

\begin{exerciseintro}
% record: H-IV-05-003
\exerciseheading{Exercise IV.5.3}
\addcontentsline{toc}{subsection}{Exercise IV.5.3}

\begin{context}
Work over an algebraically closed field, in any characteristic.
We measure a family's dimension in the moduli space of smooth curves.
We use finite automorphism groups in genus at least two, the
dimension-of-fibers theorem, and the uniqueness of the
hyperelliptic pencil.
A nonhyperelliptic genus-four canonical curve is the complete
intersection of a unique quadric \(Q\subset\mathbb P^3\) and a
cubic. The quadric is smooth or a cone; its rulings give exactly
two or one degree-three pencils, respectively.
Finite morphisms between smooth curves are flat. On
\(\mathbb P^1\), line bundles are classified by degree and
\(H^1(\mathcal O(m))=0\) for \(m\geq-1\).
We use smoothness as an open condition and Bertini for a very
ample system.
\end{context}

\begin{problem}
Show that the hyperelliptic genus-four curves form an irreducible
family of dimension seven, the nonhyperelliptic ones an
irreducible family of dimension nine, and those with just one
\(g^1_3\) an irreducible family of dimension eight.
\end{problem}
\end{exerciseintro}

\begin{solution}
First parametrize double covers without excluding characteristic
two. For a smooth hyperelliptic curve with covering map
\(\pi:X\to\mathbb P^1\), the quotient of
\(\pi_*\mathcal O_X\) by the unit subbundle is a line bundle:
locally the nonzero unit is part of a basis of this rank-two
bundle. Euler characteristics show that the quotient is
\(\mathcal O(-5)\). The extension splits since
\(H^1(\mathcal O(5))=0\).
Consequently the algebra has the form
\[
\mathcal O\oplus\mathcal O(-5),\qquad
w^2+a w=b,
\quad
a\in H^0(\mathcal O(5)),\quad
b\in H^0(\mathcal O(10)).
\tag{1}
\]
Conversely, (1) defines a finite flat double cover with
\(\chi(\mathcal O_X)=1-4=-3\), so a smooth connected such
cover has genus four. Smooth covers cannot be disconnected:
then their two degree-one components would both be \(\mathbb P^1\),
contradicting this Euler characteristic.
The coefficient space is affine of dimension \(6+11=17\).
Its smooth locus is open and nonempty. In odd characteristic,
take \(a=0\) and \(b\) a binary form with ten distinct roots.
In characteristic two, take homogeneous base coordinates
\([s:t]\), \(a=s^5\), and \(b=s t^9\).
On \(s\ne0\) the derivative with respect to \(w\) is one;
on \(t\ne0\), the equation is \(v^2+z^5v=z\), whose derivative
with respect to \(z\) is nonzero at the only point over \(z=0\).
Thus this example also is smooth.

Changing the algebra generator over a fixed base has the form
\(w\mapsto u w+c\), where \(u\in k^\times\) and
\(c\in H^0(\mathcal O(5))\); this group has dimension seven.
Allowing a base automorphism adds three dimensions.
Equivalently, the group of these changes fits over
\(\operatorname{PGL}_2\) with the preceding seven-dimensional
kernel; lifts of a base transformation to \(\mathcal O(5)\)
exist and differ by a scalar.
Uniqueness of the hyperelliptic pencil implies that two
parameters define isomorphic curves exactly when related by
these changes. The stabilizers are finite, since they act
faithfully as curve automorphisms. Thus every moduli fiber
has dimension ten, and the image is irreducible of dimension
\[
17-10=7.
\]

For nonhyperelliptic curves, the quadrics that are smooth or
cones form an irreducible open subset of the projective space
\(\mathbb P^9\) of quadratic forms.
For each such \(Q\), the restriction sequence gives
\[
h^0(Q,\mathcal O_Q(3))
=h^0(\mathbb P^3,\mathcal O(3))
 -h^0(\mathbb P^3,\mathcal O(1))
=20-4=16.
\]
The cubic sections therefore form a projective bundle with
fiber \(\mathbb P^{15}\). Its open subset of smooth complete
intersections is nonempty over both types of quadric:
on a cone require the cubic to avoid its vertex, and apply
Bertini on the smooth locus.
These curves are connected complete intersections, and
adjunction gives \(\omega_X=\mathcal O_X(1)\) and genus four.
Thus they are canonical curves, and they parametrize exactly
the nonhyperelliptic locus.
The parameter space has dimension \(9+15=24\).
Uniqueness of the quadric and completeness of the canonical
embedding identify each moduli fiber with a
\(\operatorname{PGL}_4\)-orbit. Its dimension is fifteen,
because its stabilizer is finite. The image therefore is
irreducible of dimension \(24-15=9\).

Finally, the cone quadrics form a single projective orbit of
dimension eight. One can also count them by choosing their
vertex (three parameters) and a smooth conic in the quotient
plane (five parameters). This description works in
characteristic two as well.
Restricting the cubic-section bundle to this irreducible locus
gives dimension \(8+15=23\), hence moduli dimension
\(23-15=8\). The canonical classification in the context
identifies its image exactly with the curves having one
\(g^1_3\).
\end{solution}

\begin{exerciseintro}
% record: H-IV-05-004
\exerciseheading{Exercise IV.5.4}
\addcontentsline{toc}{subsection}{Exercise IV.5.4}

\begin{context}
Work over an algebraically closed field. A nonhyperelliptic
genus-four canonical curve is a degree-six complete intersection
of a cubic and a quadric; a smooth quadric gives two \(g^1_3\)'s
and a cone gives one. A smooth quadric is
\(\mathbb P^1\times\mathbb P^1\), and its curve has type \((3,3)\).
Resolving a quadric cone gives the ruled surface \(\mathbb F_2\),
with section \(E\), fiber \(F\), \(E^2=-2\), \(E F=1\),
and hyperplane class \(E+2F\).
We use blowup and blowdown charts, the plane genus formula,
and the generic-projection theorem giving a nodal plane model
of the same degree from a smooth space curve.
A purely inseparable degree-three cover of \(\mathbb P^1\)
over this field is rational.
A node has delta invariant one. A tacnode here means two smooth
branches with intersection multiplicity two; it has delta invariant
two, including in characteristic two.
\end{context}

\begin{problem}
Let \(X\) be nonhyperelliptic of genus four.
(a) Prove that it has two \(g^1_3\)'s exactly when it has a plane
quintic model with two nodes and no other singularities.

(b) If it has one \(g^1_3\), construct a plane quintic model
with one tacnode and no other singularities, and prove that the
least degree of a nodal plane model is six.
\end{problem}
\end{exerciseintro}

\begin{solution}
(a) Put the canonical curve on its smooth quadric \(Q\).
Choose \(P\in X\) so that each of the two ruling fibers through
\(P\) meets \(X\) in three distinct points. Such choices form
a nonempty open subset: both degree-three maps to \(\mathbb P^1\)
are separable. Indeed, an inseparable degree-three map could
occur only in characteristic three and would be purely
inseparable, forcing \(X\) to be rational.

Project \(Q\) from \(P\) to \(\mathbb P^2\). On the blowup of
\(Q\) at \(P\), this is a birational morphism contracting exactly
the strict transforms of the two ruling fibers through \(P\);
they are disjoint curves of self-intersection minus one.
This description follows as well by writing the quadric in
Segre coordinates with \(P=[1:0:0:0]\) and projecting onto the
last three coordinates.
The strict transform of \(X\) meets each contracted curve in
two distinct points, transversely. Blowing down such a curve
therefore produces a node: its two intersection points specify
two distinct tangent directions at the image point.
Away from these curves the surface morphism is an isomorphism.
The restricted plane system is \(K_X-P\), of degree five.
Consequently the image is a plane quintic with exactly the
two stated nodes.

Conversely, projection from either node of a plane quintic
gives, on its normalization, a base-point-free degree-three
pencil. These pencils are distinct. In fact, the pair of
projections from two distinct plane points is birational on
the plane: outside their joining line, a point is recovered
as the intersection of its two projection lines.
Thus their restrictions jointly generate the function field
of the normalization. If the pencils were equal, their pair
would factor through a degree-three map to \(\mathbb P^1\),
which is impossible for a birational map. Hence \(X\) has
two \(g^1_3\)'s.

(b) Now the canonical quadric is a cone. Its vertex does not
lie on \(X\): if it did, the quadric gradient there would
vanish and the complete intersection would be singular.
On \(\mathbb F_2\), the curve therefore has class
\(3E+6F\) and is disjoint from \(E\).
Choose a fiber meeting \(X\) in three distinct points
\(P,A,B\), and blow up \(P\). The strict transform \(F'\)
of that fiber has self-intersection minus one.
Contract \(F'\); this changes \(\mathbb F_2\) to
\(\mathbb F_1\). The image \(E_1\) of \(E\) now has
self-intersection minus one, so contracting \(E_1\) gives
\(\mathbb P^2\). This composite is the resolution of projection
of the cone from \(P\), and restricts to the system \(K_X-P\).

The first contraction identifies \(A,B\) to a node \(R\) on
\(E_1\). Both branches are transverse to \(E_1\), since their
points on \(F'\) differ from \(E\cap F'\).
The second contraction makes these branches tangent.
For precision, take coordinates \((u,v)\) at \(R\) with
\(E_1=(u=0)\). The branches have expansions
\(v=a_i u+O(u^2)\), with \(a_1\ne a_2\).
The blowdown has coordinates \(x=u,\ y=uv\), so the images
are \(y=a_i x^2+O(x^3)\). Each branch is smooth and their
intersection multiplicity is exactly two.
Thus the image has a tacnode. The surface map is an
isomorphism elsewhere along the curve, so there are no other
singularities. Its degree is five.

A plane curve of degree at most four has arithmetic genus
at most three, so cannot normalize to \(X\).
A nodal quintic with normalization genus four would have
exactly two nodes, and by (a) would give two \(g^1_3\)'s.
Thus no nodal model of degree at most five exists in the
one-pencil case. Finally, generic projection of the smooth
degree-six canonical space curve gives a nodal sextic.
The least possible degree is therefore six.
\end{solution}

\begin{exerciseintro}
% record: H-IV-05-005
\exerciseheading{Exercise IV.5.5}
\addcontentsline{toc}{subsection}{Exercise IV.5.5}

\begin{context}
Work over an algebraically closed field. We use Riemann--Roch,
Clifford's theorem with its equality cases, adjunction,
Bertini, and finite automorphism groups for curves of genus
at least two. A plane curve of degree five has arithmetic
genus six; a reduced plane singularity with delta invariant
one is a node or a cusp. Here cusp means the unibranch
delta-one singularity, also in positive characteristic.
On \(S=\operatorname{Bl}_p\mathbb P^2\), write \(H\) for the
pullback of a line and \(E\) for the exceptional curve.
Then \(H^2=1\), \(H E=0\), \(E^2=-1\), and
\(K_S=-3H+E\). The conics through \(p\), of class \(2H-E\),
embed \(S\) in \(\mathbb P^4\). Curves of type \((a,b)\)
on \(\mathbb P^1\times\mathbb P^1\) have arithmetic genus
\((a-1)(b-1)\). We use the dimension-of-fibers theorem.
\end{context}

\begin{problem}
Let \(X\) be a nonhyperelliptic curve of genus five.
(a) Show that the canonical curves which are complete
intersections of three quadrics form a family of dimension twelve.

(b) Prove the corrected statement: \(X\) has a \(g^1_3\)
if and only if it has a plane quintic model with one node
or one cusp and no other singularities. These curves form
an irreducible family of dimension eleven.
The source's assertion with only a node omits the cuspidal case.

(c) Show that the conics through the singular point induce
the canonical embedding and place \(X\) on a smooth cubic
scroll \(V\cong\operatorname{Bl}_p\mathbb P^2\).
Prove that its ruling lines are the trisecants from the
degree-three pencil, that every quadric containing \(X\)
contains \(V\), and that both \(V\) and the pencil are unique.
\end{problem}
\end{exerciseintro}

\begin{solution}
(a) Quadratic forms in five variables form a vector space
of dimension fifteen. The three-dimensional subspaces form
\(\operatorname{Gr}(3,15)\), irreducible of dimension
\(3(15-3)=36\).
The locus whose common zero scheme is a smooth curve
complete intersection is nonempty and open: successive
general quadrics give examples by Bertini.
Adjunction gives \(K_X=\mathcal O_X(1)\);
the degree is eight, so the genus is five.
These embeddings are complete canonical embeddings:
restriction of linear forms is injective, and both spaces
of sections have dimension five.
The degree-two part of the ideal has dimension three,
as follows from its Koszul resolution, so it recovers
the original point of the Grassmannian.
Consequently moduli fibers are precisely the
\(\operatorname{PGL}_5\)-orbits. Stabilizers are finite,
so these fibers have dimension \(25-1=24\).
The resulting family is irreducible of dimension
\(36-24=12\).

(b) Let \(D\) be the divisor class of a \(g^1_3\).
It has no base points, since otherwise \(X\) would have
a degree-one or degree-two map to \(\mathbb P^1\).
Clifford's theorem gives \(h^0(D)=2\), and
Riemann--Roch then gives
\[
\deg(K_X-D)=5,\qquad h^0(K_X-D)=3.
\tag{1}
\]
This residual system is base-point-free. A base point \(P\)
would give a degree-four divisor \(K_X-D-P\) with at least
three sections. It is special by Riemann--Roch; equality
in Clifford's theorem would force hyperellipticity,
since this divisor is neither zero nor canonical.

Thus (1) maps \(X\) to a nondegenerate plane curve.
The product of the degree of the map and the degree of
its image is five. Since the image is not a line, this
prime degree forces the map to be birational and the
image to be a quintic. The difference between arithmetic
genus six and normalization genus five is one.
There is therefore exactly one singularity, a node or
a cusp, and no others.
Conversely, projection from a multiplicity-two singularity
of either type gives a degree-three pencil on the
normalization. Its degree is \(5-2=3\), because the two
fixed intersections at the center are removed on normalization.

We will need uniqueness. Suppose there are two distinct
degree-three pencils. The pair of maps to
\(\mathbb P^1\times\mathbb P^1\) has degree dividing three
onto its image. If that degree were three, both projections
of the image would have degree one; the pencils would
then differ by an automorphism of \(\mathbb P^1\).
Otherwise the pair is birational and its image has type
\((3,3)\), of arithmetic genus four, which cannot have
normalization genus five. Thus the pencil is unique.

For the parameter count, fix a marked point \(p\) of the
plane. Vanishing to order at least two at \(p\) imposes
three independent linear conditions on the twenty-one
quintic coefficients. Allowing \(p\) to vary gives an
irreducible projective bundle of dimension
\[
2+(20-3)=19.
\]
The locus where the quintic is integral, its strict transform
on the blowup at \(p\) is smooth, and the multiplicity at \(p\)
is exactly two, is nonempty and open.
The strict transform has class \(5H-2E\) and genus five
by adjunction; its intersection of length two with \(E\)
gives a node if reduced and a cusp if supported at one point.
There are no other singularities.
Nonemptiness follows from Bertini for the very ample system
\(5H-2E\), whose general restriction to \(E\) has two
distinct zeros.

Each resulting normalization is nonhyperelliptic: a
degree-two map together with its degree-three map would
be birational onto a curve of type \((2,3)\), with
arithmetic genus two. On the blowup, adjunction and
projection from \(p\) give, on the normalization,
\[
K_X=(2H-E)|_X,\qquad D=(H-E)|_X.
\tag{2}
\]
Thus its plane system is \(H|_X=K_X-D\), uniquely determined
by the unique pencil. Two of our plane models of the same
curve are therefore projectively equivalent. The moduli
fibers are \(\operatorname{PGL}_3\)-orbits of dimension eight,
so the image is irreducible of dimension \(19-8=11\).

The omitted cuspidal case really occurs. On the blowup,
impose a double intersection at a specified point of \(E\)
and take a general member of \(|5H-2E|\) with that condition.
Its local equation can be chosen to have the form
\(v^2-u+\text{higher terms}=0\), where \(E=(u=0)\).
It is smooth there and tangent to \(E\); Bertini away
from that point makes it smooth elsewhere.
For example these conditions come from plane quintics
with quadratic term \(y^2\) and nonzero \(x^3\) term at \(p\);
the freely varying terms of degree at least three remove
base points away from \(p\).
The smooth strict transform is connected (its class is
ample), hence integral, and blowing down gives the cusp.
By the uniqueness of \(D\) and \(K_X-D\), this normalization
cannot also have a nodal plane quintic model.

(c) Formula (2) shows that the conics through \(p\) restrict
to canonical divisors. Their vector space has dimension
five, and restriction is injective because a conic cannot
contain an integral quintic. Hence it is the full canonical
space \(H^0(X,K_X)\).
The conic embedding of \(S\) has degree
\((2H-E)^2=3\); call its image \(V\).
The fibers \(H-E\) of the pencil of lines through \(p\)
map to lines, since
\[
(2H-E)(H-E)=1,
\]
and each meets \(X\) in degree
\((5H-2E)(H-E)=3\).
Thus \(V\) is the union of the ruling trisecants of
the degree-three pencil.

Every quadric containing \(X\) contains each of these
lines: its restriction to a line has degree two but
vanishes on the length-three intersection with \(X\).
It consequently contains their union \(V\).
This already shows that these trigonal canonical curves
are not complete intersections of three quadrics.
For uniqueness of the surface, the unique pencil recovers
each ruling line as the span of its fiber in the
canonical embedding, so recovers their union.
Indeed, a length-three fiber has canonical span a line:
Riemann--Roch gives \(h^0(K_X-D)=3\), hence the span has
dimension \(h^0(K_X)-h^0(K_X-D)-1=1\).
This proves uniqueness of both the pencil and the scroll.
\end{solution}

\begin{exerciseintro}
% record: H-IV-05-006
\exerciseheading{Exercise IV.5.6}
\addcontentsline{toc}{subsection}{Exercise IV.5.6}

\begin{context}
Work over an algebraically closed field. A smooth plane
quintic has genus six and canonical bundle \(\mathcal O_X(2)\);
its canonical sections are restrictions of plane conics.
We use Riemann--Roch and that a curve with a purely inseparable
map of prime degree to \(\mathbb P^1\) is rational.
A curve of bidegree \((a,b)\) on
\(\mathbb P^1\times\mathbb P^1\) has arithmetic genus
\((a-1)(b-1)\), and its normalization has no larger genus.
Bertini supplies smooth integral members of the very ample
system \(\mathcal O(a,b)\) for positive \(a,b\).
\end{context}

\begin{problem}
Prove that a nonsingular plane quintic has no \(g^1_3\).
Construct a nonhyperelliptic genus-six curve that has no
nonsingular plane quintic model.
\end{problem}
\end{exerciseintro}

\begin{solution}
Let \(X\) be a smooth plane quintic. Its canonical bundle
\(\mathcal O_X(2)\) is very ample, so \(X\) is nonhyperelliptic.
A degree-three pencil on it could have no fixed part:
removing that part would give a map of degree at most two
to \(\mathbb P^1\), contradicting either genus six or
nonhyperellipticity.
Suppose such a pencil exists. Its map is separable, since an
inseparable degree-three map would force rationality.
A general fiber is therefore \(D=P+Q+R\) with three distinct
points. These points impose three independent conditions on
plane conics. To see independence at \(P\), choose a line
through \(Q\) avoiding \(P\) and a line through \(R\)
avoiding \(P\); their product vanishes at \(Q,R\) but not
at \(P\), and similarly for the other two points.
Since the six-dimensional conic space maps isomorphically
to \(H^0(X,K_X)\), it follows that
\[
h^0(X,K_X-D)=6-3=3.
\]
But Riemann--Roch and \(h^0(X,D)\geq2\) give
\[
h^0(X,K_X-D)=h^0(X,D)-3-1+6\geq4,
\]
a contradiction.

For the example, choose a smooth integral curve \(C\) of
bidegree \((3,4)\) on \(\mathbb P^1\times\mathbb P^1\).
It has genus \((3-1)(4-1)=6\), and one projection gives a
degree-three map. It is not hyperelliptic. Indeed, if it
also had a degree-two map, the pair of maps would have
degree dividing both two and three onto its image, and
therefore would be birational. The image would have
bidegree \((2,3)\), with arithmetic genus two, contradicting
the genus of its normalization \(C\).
If \(C\) admitted a smooth plane quintic model, its
degree-three pencil would contradict the first part.
Thus \(C\) is the required example.
\end{solution}

\begin{exerciseintro}
% record: H-IV-05-007
\exerciseheading{Exercise IV.5.7}
\addcontentsline{toc}{subsection}{Exercise IV.5.7}

\begin{context}
Work over an algebraically closed field \(k\).
Nonhyperelliptic genus-three curves embed canonically as
smooth plane quartics, forming an open dense locus in genus-three
moduli. Automorphism groups of curves of genus at least two
are finite. We use the tame Hurwitz bound \(84(g-1)\).
For a cyclic group of prime order \(\ell\), Hurwitz reads
\(2g-2=\ell(2h-2)+\sum d_P\); in the tame case each
ramified point contributes \(\ell-1\), and in characteristic
\(\ell\) it contributes \((\ell-1)(m_P+1)\), with \(m_P\geq1\).

We explicitly permit one finite-group representation fact:
over \(\mathbb C\), the matrices \(S,T,U\) displayed in part
(b) below generate the faithful Klein representation of
the simple group \(\operatorname{PSL}_2(\mathbb F_7)\), of
order \(168\).
This is the matrix representation lemma, not an assumption
about the full automorphism group of a curve; see Elkies,
\emph{The Klein Quartic in Number Theory}, Section 1.1,
equations (1.2)--(1.3).
We also use Jordan normal form, Sylow's and Cauchy's
theorems, and dimension of fibers.
\end{context}

\begin{problem}
(a) Prove that every automorphism of a nonhyperelliptic
genus-three curve is induced by a projective transformation
of its canonical plane. Explain the missing hypothesis in
the source statement.

(b) For
\[
X:\quad x^3y+y^3z+z^3x=0,
\]
prove that, in characteristic different from three and seven,
\(\operatorname{Aut}(X)\cong\operatorname{PSL}_2(\mathbb F_7)\)
has order \(168\). Explain the necessary exclusion of
characteristic seven and the scope of the Hurwitz bound.

(c) Show that a general genus-three curve has no nonidentity
automorphism, in any characteristic.
\end{problem}
\end{exerciseintro}

\begin{solution}
(a) An automorphism pulls back canonical differentials,
giving an invertible linear map on the three-dimensional
space \(H^0(X,\omega_X)\). The canonical map is equivariant
for its projectivization, so an automorphism is induced
by the corresponding projective transformation.
The canonical embedding makes this action faithful.
For a hyperelliptic genus-three curve the canonical map
has degree two onto a conic and is not an embedding.
Its hyperelliptic involution acts trivially on the canonical
image, so the assertion as phrased in the source needs
the nonhyperelliptic hypothesis.

(b) Write \(F=x^3y+y^3z+z^3x\).
First check smoothness. At a coordinate point on \(X\),
one partial derivative is nonzero.
If all coordinates are nonzero and all partial derivatives
vanish, put \(A=x^3y,\ B=y^3z,\ C=z^3x\).
Multiplying the three derivative equations by \(x,y,z\)
gives \(3A+C=A+3B=B+3C=0\).
Hence \(A=9C,\ B=-3C\), and \(F=7C\).
Thus a singular point on the curve forces characteristic
seven. In that characteristic \([1:2:4]\) is indeed on
the curve and all three derivatives vanish.
For every other characteristic the curve is smooth;
adjunction makes it a canonical genus-three curve.

Choose a primitive seventh root \(\zeta\), and set
\[
a=\zeta-\zeta^6,\quad b=\zeta^2-\zeta^5,\quad
c=\zeta^4-\zeta^3,\quad \delta=a+b+c.
\]
Then \(\delta^2=-7\). The promised matrices are
\[
S=\begin{pmatrix}\zeta^4&0&0\\ 0&\zeta^2&0\\ 0&0&\zeta\end{pmatrix},
\quad
T=\begin{pmatrix}0&1&0\\ 0&0&1\\ 1&0&0\end{pmatrix},
\quad
U=-\delta^{-1}
\begin{pmatrix}a&b&c\\ b&c&a\\ c&a&b\end{pmatrix}.
\tag{1}
\]
They preserve \(F\). For \(S,T\) this follows term by term;
for \(U\), writing \(M=-\delta U\), expansion using
\(1+\zeta+\cdots+\zeta^6=0\) gives
\[
M^2=-7I,\qquad F(M(x,y,z))=49F(x,y,z).
\tag{2}
\]
Since \(\delta^4=49\), (2) proves the assertion.

By the permitted representation lemma, in characteristic
zero they generate a copy of the simple group of order
168. All matrices and their inverses lie in
\(\mathbb Z[\zeta,1/7]\), since \(\delta^{-1}=-\delta/7\).
Thus all multiplication relations reduce to characteristic
\(p\ne7\), yielding a homomorphism from that simple group
to \(\operatorname{PGL}_3(k)\).
Its image is nontrivial, since \(S\) still has projective
order seven. Simplicity makes this homomorphism injective.
Consequently \(X\) has a subgroup \(G\) of 168 automorphisms
in every characteristic different from seven.

\emph{Characteristic zero and tame characteristic.}
In characteristic zero, Hurwitz immediately gives
\(\operatorname{Aut}(X)=G\).
The same holds when \(p\geq5\), \(p\ne7\), because no
genus-three curve has an automorphism of order \(p\)
in those characteristics. Indeed, for a cyclic action of
order \(p\), the quotient genus \(h\) is at most one.
If \(h=1\), the different would have degree four, whereas
each ramified point contributes at least \(2(p-1)>4\).
If \(h=0\), Hurwitz requires \(p-1\) to divide
\(4+2p\), hence to divide six, forcing \(p=7\) among
primes \(p\geq5\).
Cauchy's theorem therefore shows that the entire finite
automorphism group has order prime to \(p\). It is tame,
so its order is at most 168.

\emph{Characteristic-two bitangents.}
It remains to give an upper bound in characteristic two,
where a tame-bound argument would be invalid.
Consider lines whose intersection divisor with \(X\)
is twice a degree-two divisor.
Such a line has equation \(z=\alpha x+\beta y\):
no line with zero \(z\)-coefficient has this property,
as direct substitution into \(F\) shows.
The coefficients of \(x^3y\) and \(xy^3\) in the
restriction are respectively
\[
1+\alpha^2\beta,\qquad \alpha+\beta^3.
\]
The restriction is a square precisely when both vanish,
so these seven lines have
\[
\alpha=\beta^3,\qquad \beta^7=1.
\tag{3}
\]
They are ordinary bitangents: the coefficient of
\(x^2y^2\) is \(\alpha\beta^2=\beta^5\ne0\), so the
quadratic square root has two distinct zeros.

In the dual plane, (3) gives the points
\((\beta^3:\beta:1)=(\beta^4:\beta^2:\beta)\).
As \(\beta\) runs over \(\mathbb F_8^\times\), these form
a projective copy of the seven points of
\(\mathbb P^2(\mathbb F_2)\).
To verify this, the map
\(t\mapsto(t^4,t^2,t)\) is \(\mathbb F_2\)-linear,
and the images of an \(\mathbb F_2\)-basis of \(\mathbb F_8\)
are independent over \(k\): a dependence among the rows
would give a nonzero polynomial of degree at most four
vanishing on all eight elements of \(\mathbb F_8\).
Their projective incidence is therefore the Fano plane.
Its collineation group is \(\operatorname{GL}_3(\mathbb F_2)\),
of order \((8-1)(8-2)(8-4)=168\).
For example a collineation is determined by the images
of a projective frame, and the Fano incidence forces the
remaining points to be the corresponding sums.
Every projective automorphism of \(X\) permutes these
seven bitangents, faithfully because they contain a dual
projective frame. Hence \(|\operatorname{Aut}(X)|\leq168\).
Together with \(G\), this proves equality also in
characteristic two. The value agrees with \(84(g-1)\),
but here that equality is a computation, not an application
of the tame Hurwitz bound.

(c) \emph{General quartics.} Smooth plane quartics form an open subset of
\(\mathbb P^{14}\). We bound the locus with a nonidentity
automorphism by dimension twelve.
Every such finite group contains an element of prime
order \(\ell\). Hurwitz implies \(\ell\leq7\).
For completeness, if \(\ell\geq11\) the quotient genus
is at most one. Quotient genus one would require a
different of degree four, impossible since every
nonzero contribution is at least \(\ell-1\).
For quotient genus zero, all different exponents are
multiples of \(\ell-1\), so
\(\ell-1\mid4+2\ell\), hence \(\ell-1\mid6\), again
impossible. This argument includes wild cyclic actions.
Only the primes \(2,3,5,7\) need be considered.

If \(\ell\ne\operatorname{char}k\), the projective
transformation is diagonalizable. Up to projective
conjugacy there are finitely many possibilities, with
eigenvalues normalized to \(\ell\)-th roots of unity.
With one repeated eigenvalue its conjugacy class has
dimension four; with three distinct eigenvalues it has
dimension six.
A preserved quartic form is in an eigenspace of the
substitution action on degree-four forms.
The largest possible eigenspace dimensions, and hence
bounds for the corresponding families, are
\[
\begin{array}{c|c|c|c}
\ell&\text{eigenvalue type}&
\text{eigenspace dimension}&\text{family dimension}\\ \hline
2&\text{repeated}&9&4+8=12\\
3&\text{repeated}&7&4+6=10\\
3&\text{distinct}&5&6+4=10\\
5,7&\text{repeated}&5&4+4=8\\
5,7&\text{distinct}&\leq5&\leq6+4=10
\end{array}
\tag{4}
\]
These counts come from the fifteen monomials \(x^iy^jz^k\)
with \(i+j+k=4\).
For repeated eigenvalues, group by the exponent of the
exceptional variable modulo \(\ell\).
For three distinct cube roots each weight has five
monomials. For \(\ell\geq5\) and distinct eigenvalues,
fixing one exponent leaves at most one monomial of
each weight among the at most five possibilities for a
second exponent, giving the stated bound.

In characteristic \(\ell\), an element of projective
order \(\ell\) is unipotent. Its nontrivial Jordan type
is either \((2,1)\), with conjugacy-class dimension four,
or \((3)\), with dimension six. The latter has order
four in characteristic two and so does not enter the
prime-order count there.
For type \((2,1)\), use \(x\mapsto x+y\) with \(y,z\)
fixed. The invariant polynomial ring is
\[
k[y,z,x^\ell-xy^{\ell-1}].
\]
Its degree-four part has dimension nine, seven, or five
as \(\ell=2,3,\) or \(5,7\), respectively, giving the
same bounds as the repeated-eigenvalue rows of (4).
The ring formula also follows by division by the monic
invariant \(x^\ell-xy^{\ell-1}\): an invariant remainder
of \(x\)-degree less than \(\ell\) must be independent of \(x\).

For type \((3)\), use
\((x,y,z)\mapsto(x+y,y+z,z)\).
On degree-four forms, the kernel of substitution minus
identity has dimension five in characteristic three
and three in characteristics five and seven.
Here is an explicit coefficient computation of these
small ranks. Index rows and columns by triples of
nonnegative integers summing to four. In column
\((a,b,c)\), expand
\[
(x+y)^a(y+z)^b z^c-x^a y^b z^c
=\sum_{i=0}^a\sum_{j=0}^b
\binom ai\binom bj x^i y^{a-i+j}z^{b-j+c}
-x^a y^b z^c.
\]
The resulting fifteen-by-fifteen matrices have ranks
ten, twelve, twelve modulo three, five, seven.
Thus the projective fixed-form loci, with their
conjugacy classes, have dimensions at most
\(6+4=10\) or \(6+2=8\).

There are only finitely many conjugacy types and
eigenspaces in this argument. For each, the family of
preserved quartics is the image of the conjugacy class
with the projective eigenspace over it; its dimension
has just been bounded by twelve.
The closures of their images therefore form a proper
closed subset of \(\mathbb P^{14}\).
A smooth quartic outside that subset has no prime-order
automorphism, and hence no nonidentity automorphism.
This property is invariant under projective equivalence.
Since the nonhyperelliptic locus is open dense in
genus-three moduli, it proves the assertion for a
general genus-three curve.
\end{solution}

\begin{exerciseintro}
\corpussection{IV.6 Classification of Curves in Projective 3-Space}
\addcontentsline{toc}{section}{IV.6 Classification of Curves in Projective 3-Space}

% record: H-IV-06-001
\exerciseheading{Exercise IV.6.1}
\addcontentsline{toc}{subsection}{Exercise IV.6.1}

\begin{context}
Work over an algebraically closed field. Curves are nonsingular,
projective, and integral. A plane curve of degree \(d\) has genus
\((d-1)(d-2)/2\). On \(\mathbb P^1\),
\(h^0(\mathcal O(m))=m+1\) for \(m\geq0\).
We use Bezout and the complete-intersection genus formula
\(g=1+ab(a+b-4)/2\).
A smooth curve on a quadric cone of even degree \(2a\)
has genus \((a-1)^2\); one of odd degree \(2a+1\)
has genus \(a(a-1)\).
These cone formulas are the classification recalled in
Hartshorne IV, Remark 6.4.1(d).
\end{context}

\begin{problem}
Show that a rational quartic curve in \(\mathbb P^3\) lies
on a unique quadric, which is nonsingular.
\end{problem}
\end{exerciseintro}

\begin{solution}
Let \(L=\mathcal O_X(1)\).
Since \(X\cong\mathbb P^1\) and \(\deg L=4\),
\(h^0(X,L^2)=9\), whereas the space of ambient quadratic
forms has dimension ten. The restriction map therefore
has a nonzero kernel, giving a quadric \(Q\) through \(X\).
The curve cannot be planar: a nonsingular plane quartic
has genus three. Thus no quadric containing it can be
a union or double of planes, since integrality would
put \(X\) in one of those planes.

If two distinct quadrics contained \(X\), their irreducibility
would make their intersection a complete-intersection curve
of degree four. Since it contains the degree-four integral
curve \(X\), it equals \(X\) scheme-theoretically.
To justify the last assertion, a complete intersection
is Cohen--Macaulay with no embedded points, and equality
of degrees forces generic multiplicity one and leaves
no other component. Its ideal is consequently the
radical ideal of \(X\).
The complete-intersection genus would be one, a contradiction.
Hence \(Q\) is unique.

Finally, if \(Q\) were a cone, the even-degree cone formula
with \(2a=4\) would again give genus one.
Thus \(Q\) is nonsingular.
\end{solution}

\begin{exerciseintro}
% record: H-IV-06-002
\exerciseheading{Exercise IV.6.2}
\addcontentsline{toc}{subsection}{Exercise IV.6.2}

\begin{context}
Work over an algebraically closed field, in any characteristic.
A rational nonsingular projective curve is \(\mathbb P^1\),
and \(h^0(\mathbb P^1,\mathcal O(m))=m+1\) for \(m\geq0\).
An injective morphism of a smooth projective curve that
separates tangent vectors is a closed immersion.
There are twenty cubic forms and ten quadratic forms
in four variables.
\end{context}

\begin{problem}
Show that every rational quintic in \(\mathbb P^3\) is
contained in a cubic surface, and construct one contained
in no quadric surface.
\end{problem}
\end{exerciseintro}

\begin{solution}
For a rational quintic \(X\), the restriction of cubic
forms has target
\[
H^0(X,\mathcal O_X(3))\cong
H^0(\mathbb P^1,\mathcal O(15)),
\]
of dimension sixteen. Its kernel in the twenty-dimensional
space of cubic forms is therefore nonzero. Any nonzero
element defines the required cubic surface.

For the second assertion, consider
\[
[s:t]\longmapsto
[s^5:s^4t:s^3t^2+st^4:t^5].
\tag{1}
\]
The coordinates have no common zero. On \(s\ne0\),
put \(u=t/s\); the map is
\([1:u:u^2+u^4:u^5]\), so it separates points and tangent
vectors there. The point \(s=0\) has image \([0:0:0:1]\),
which is not an image of a finite \(u\).
Near it put \(q=s/t\). In the chart where the last
coordinate is one, the third coordinate is \(q^3+q\),
with derivative one at \(q=0\).
Thus (1) is a closed immersion, and its image is a
smooth rational curve of degree five.

The ten quadratic coordinate products, restricted to
\(s=1\), are
\[
1,\ u,\ u^2+u^4,\ u^5,\ u^2,\ u^3+u^5,\
u^6,\ u^4+2u^6+u^8,\ u^7+u^9,\ u^{10}.
\tag{2}
\]
They are linearly independent over every field.
In a relation, successively compare coefficients of
\(u^{10},u^9,u^8,u^6\) to remove the last, next-to-last,
eighth, and seventh terms. Then the coefficient of
\(u^4\) removes the third term, that of \(u^3\) the sixth,
and that of \(u^5\) the fourth. The remaining
\(1,u,u^2\) are independent. This reasoning does not
divide by two. Therefore no nonzero quadratic form
vanishes on the image of (1), as required.
\end{solution}

\begin{exerciseintro}
% record: H-IV-06-003
\exerciseheading{Exercise IV.6.3}
\addcontentsline{toc}{subsection}{Exercise IV.6.3}

\begin{context}
Work over an algebraically closed field.
On a genus-two curve \(X\), the canonical pencil \(|K_X|\)
is the unique degree-two pencil, and
\(\deg K_X=2\). Every line bundle of degree at least five
is very ample, and Riemann--Roch applies.
The variety \(\operatorname{Pic}^5(X)\) is irreducible
of dimension two; the image of a morphism from \(X\)
has dimension at most one.
We use Bezout and the quadric-cone resolution
as well as the plane genus formula.
\(\mathbb F_2\to Q\) denotes the resolution,
with \(E^2=-2,\ E F=1,\ F^2=0\)
and hyperplane class \(E+2F\).
A smooth degree-five curve on the cone has strict
transform \(2E+5F\), meeting \(E\) transversely once;
the restriction of \(F\) is its degree-two pencil.
\end{context}

\begin{problem}
Prove that a degree-five genus-two curve in \(\mathbb P^3\)
lies on a unique quadric.
For every abstract genus-two curve, construct degree-five
embeddings whose quadric is smooth, and others whose
quadric is singular.
\end{problem}
\end{exerciseintro}

\begin{solution}
For an embedded curve, put \(L=\mathcal O_X(1)\).
The curve is not planar, since a smooth plane quintic
has genus six. Riemann--Roch gives
\(h^0(L^2)=10+1-2=9\); hence some quadric contains \(X\).
It is irreducible, since otherwise \(X\) would lie in
a plane. Two distinct such quadrics would intersect
in a degree-four curve and could not contain the
degree-five curve \(X\), by Bezout. This proves uniqueness.

Now start with an arbitrary genus-two curve.
Every degree-five line bundle \(L\) is very ample,
and \(h^0(L)=5+1-2=4\). Thus its complete system embeds
\(X\) with the required degree into \(\mathbb P^3\).
We claim that its unique quadric is a cone exactly when
\[
L\cong\omega_X^{\otimes2}(P)
\quad\text{for some }P\in X.
\tag{1}
\]
For this choice of \(L\), let \(u,v\) be a basis of
\(H^0(X,\omega_X)\). The three independent sections
\(u^2,uv,v^2\) form a basis of \(H^0(X,\omega_X^{\otimes2})\):
their independence follows by pulling back independent
quadratics on the canonical \(\mathbb P^1\).
Include them in \(H^0(X,L)\) using the canonical section
of \(\mathcal O_X(P)\), and extend to a basis with one
further section. The resulting image satisfies
\(x_0x_2-x_1^2=0\), an irreducible quadric cone
(including in characteristic two). This proves one direction.

Conversely, for a curve on a cone, use its strict transform
on \(\mathbb F_2\). The section \(E\) cuts out one point \(P\),
and the ruling cuts out the canonical degree-two pencil.
Restricting the hyperplane class \(E+2F\) therefore gives
\(L\cong\mathcal O_X(P+2K_X)\), proving (1).

The bundles in (1) form the image of the proper curve
\(X\to\operatorname{Pic}^5(X)\), \(P\mapsto2K_X+P\).
This is a closed subset of dimension at most one,
so it does not fill the two-dimensional Picard variety.
Choose \(L\) outside it: its unique irreducible quadric
cannot be a cone, and hence is smooth.
Choosing any \(P\) in (1) gives the singular-quadric
embedding. Both kinds therefore exist on every \(X\).
\end{solution}

\begin{exerciseintro}
% record: H-IV-06-004
\exerciseheading{Exercise IV.6.4}
\addcontentsline{toc}{subsection}{Exercise IV.6.4}

\begin{context}
Work over an algebraically closed field.
We use Riemann--Roch, the plane genus formula, and
Clifford's theorem. A degree-two divisor with two
sections on a curve of positive genus is its
hyperelliptic pencil \(A\); then \(K_X\sim(g-1)A\)
and \(h^0(nA)=n+1\) for \(0\leq n\leq g\).
A curve of type \((a,b)\) on a smooth quadric has
degree \(a+b\) and genus \((a-1)(b-1)\).
A smooth curve of degree \(2a+1\) on a quadric cone
has genus \(a(a-1)\).
\end{context}

\begin{problem}
Prove that there is no nonsingular integral curve of
degree nine and genus eleven in \(\mathbb P^3\).
\end{problem}
\end{exerciseintro}

\begin{solution}
Suppose \(X\) were such a curve and let \(D\) be its
hyperplane divisor. It is not planar, since a smooth
plane curve of degree nine has genus twenty-eight.
Riemann--Roch gives
\[
h^0(2D)=18+1-11+h^0(K_X-2D)
=8+h^0(K_X-2D).
\tag{1}
\]
The residual divisor has degree two. We claim it has
at most one section. Otherwise it is a hyperelliptic
pencil \(A\), and \(K_X\sim10A\).
Then \(2D\sim9A\), whose complete linear system is
pulled back from \(|\mathcal O_{\mathbb P^1}(9)|\):
pullback supplies ten sections, and \(h^0(9A)=10\).
It cannot separate the degree-two fibers.
But \(D\) is very ample, so \(2D\) is very ample as well,
a contradiction. This proves the claim.

By (1), \(h^0(2D)\leq9\), whereas ambient quadratic forms
have dimension ten. Thus \(X\) lies on a quadric \(Q\).
The quadric is irreducible, because \(X\) is nonplanar.
If \(Q\) is smooth, writing its curve class as \((a,b)\)
gives
\[
a+b=9,\qquad ab-a-b+1=11,
\]
so \(ab=19\). The roots of \(t^2-9t+19\) are not
integers, since its discriminant is five.
If \(Q\) is a cone, degree nine means \(2a+1=9\),
so the cone genus formula gives genus \(4\cdot3=12\),
again a contradiction. These exhaust the possibilities.
\end{solution}

\begin{exerciseintro}
% record: H-IV-06-005
\exerciseheading{Exercise IV.6.5}
\addcontentsline{toc}{subsection}{Exercise IV.6.5}

\begin{context}
Work over an algebraically closed field.
For a scheme-theoretic complete intersection of two
hypersurfaces of degrees \(a,b\) in \(\mathbb P^3\),
the Koszul resolution of its ideal sheaf is
\[
0\to\mathcal O(-a-b)\to
\mathcal O(-a)\oplus\mathcal O(-b)\to\mathcal I_X\to0.
\]
On \(\mathbb P^3\), \(H^1(\mathcal O(t))=0\) for
all integers \(t\), and \(H^0(\mathcal O(t))=0\)
when \(t<0\).
\end{context}

\begin{problem}
If \(X\subset\mathbb P^3\) is a complete intersection
of surfaces of degrees \(a,b\), prove that it lies
on no surface of degree less than \(\min(a,b)\).
\end{problem}
\end{exerciseintro}

\begin{solution}
A surface of degree \(m\) containing \(X\)
scheme-theoretically corresponds to a nonzero element
of \(H^0(\mathbb P^3,\mathcal I_X(m))\).
Twist the Koszul resolution by \(\mathcal O(m)\).
Its cohomology gives
\[
H^0(\mathcal O(m-a))\oplus H^0(\mathcal O(m-b))
\longrightarrow H^0(\mathcal I_X(m))
\longrightarrow H^1(\mathcal O(m-a-b)).
\]
If \(m<\min(a,b)\), the first group is zero by
negativity and the last is zero by projective-space
cohomology. Exactness forces \(H^0(\mathcal I_X(m))=0\),
which proves the assertion.
For the nonsingular curves of the section, set-theoretic
containment also gives scheme-theoretic containment,
since their defining ideals are radical.
\end{solution}

\begin{exerciseintro}
% record: H-IV-06-006
\exerciseheading{Exercise IV.6.6}
\addcontentsline{toc}{subsection}{Exercise IV.6.6}

\begin{context}
Work over an algebraically closed field.
For a smooth projectively normal curve in projective space,
the restriction maps
\(H^0(\mathbb P^3,\mathcal O(n))\to
H^0(X,\mathcal O_X(n))\) are surjective for all \(n\geq0\).
A nonplanar degree-\(d\) space curve satisfies Castelnuovo's
bound \(g\leq\lfloor(d-2)^2/4\rfloor\).
We use Riemann--Roch and the dimensions four and ten
of the spaces of linear and quadratic forms on \(\mathbb P^3\).
\end{context}

\begin{problem}
Let \(X\subset\mathbb P^3\) be a nonsingular projectively
normal curve not contained in a plane.
Prove that degree six implies \(g=3\) or \(4\),
and degree seven implies \(g=5\) or \(6\).
\end{problem}
\end{exerciseintro}

\begin{solution}
Write \(D\) for the hyperplane divisor.
Nonplanarity makes restriction of linear forms injective,
and projective normality makes it surjective.
Hence \(h^0(D)=4\). Riemann--Roch yields
\[
4=d+1-g+h^1(D),\qquad g=d-3+h^1(D)\geq d-3.
\]
For \(d=6\), this gives \(g\geq3\), while Castelnuovo's
bound gives \(g\leq4\), proving the first assertion.

For \(d=7\), the same arguments give \(4\leq g\leq6\).
If \(g=4\), then Riemann--Roch gives
\[
h^0(2D)=14+1-4+h^1(2D)\geq11.
\]
But projective normality requires the ten-dimensional
space of quadrics to surject onto this space, which is
impossible. Thus \(g=5\) or \(6\).
\end{solution}

\begin{exerciseintro}
% record: H-IV-06-007
\exerciseheading{Exercise IV.6.7}
\addcontentsline{toc}{subsection}{Exercise IV.6.7}

\begin{context}
Work over an algebraically closed field. A multisecant
of a nonsingular integral curve \(X\subset\mathbb P^3\)
means a line not contained in \(X\) whose scheme-theoretic
intersection with \(X\) has length at least three.
Thus tangential intersections count with multiplicity.
A morphism from a smooth projective curve is an embedding
if its defining system separates every length-two subscheme.
We use the plane genus formula and Castelnuovo's
bound \(g\leq\lfloor(d-2)^2/4\rfloor\) for nonplanar curves.
A nonplanar smooth degree-three curve is a twisted cubic,
whose ideal is generated by quadrics. A smooth genus-one
degree-four space curve is a complete intersection of
two quadrics.
\end{context}

\begin{problem}
Prove that a line, conic, twisted cubic, and elliptic
quartic have no multisecants, whereas every other
nonsingular integral curve in \(\mathbb P^3\) has
infinitely many multisecants.
\end{problem}
\end{exerciseintro}

\begin{solution}
For a line the assertion follows from intersection
of two distinct lines. For a conic, a line outside
its plane meets it in length at most one, and a line
in the plane meets it in length at most two.
For the twisted cubic and elliptic quartic, any line
meeting the curve in length at least three would be
contained in every defining quadric: a polynomial of
degree two on the line cannot have three zeros counted
with multiplicity unless it is zero. Since their ideals
are generated by quadrics, such a line would be contained
in the curve, which is impossible.

A plane curve of degree at least three has infinitely
many multisecants, namely its general plane lines.
It remains to treat nonplanar curves.
Suppose there were only finitely many multisecants.
Choose \(P\in X\) on none of them, which is possible
because each has finite intersection with \(X\).
Projection from \(P\), after removing its fixed point,
is defined everywhere on \(X\) by the three-dimensional
space of hyperplanes through \(P\), a subsystem of
\(|D-P|\), where \(D\) is the hyperplane divisor.
It is base-point-free since the original embedding
separates length-two subschemes.

This projection is an embedding into \(\mathbb P^2\).
Indeed, failure to separate a length-two subscheme
\(Z\subset X\) would mean that the length-three divisor
\(P+Z\) fails to impose three independent conditions
on ambient linear forms. Its projective span would
then be a line, which is a multisecant through \(P\).
This also covers \(P\in Z\) and nonreduced \(Z\):
the divisors \(2P+Q\) and \(3P\) record tangent and
osculating failures. Such a multisecant was excluded
by the choice of \(P\).

The image is thus a smooth plane curve of degree \(d-1\),
so
\[
g=\frac{(d-2)(d-3)}2
\leq\left\lfloor\frac{(d-2)^2}{4}\right\rfloor.
\]
For \(d\geq5\), even the inequality without the floor
would give \(2(d-3)\leq d-2\), or \(d\leq4\),
a contradiction.
Thus \(d=3\) or \(4\), since a nonplanar curve has
degree at least three. The displayed genus formula
gives respectively \(g=0\) or \(1\).
These are precisely the twisted cubic and elliptic
quartic. Therefore any curve outside the four listed
classes has infinitely many multisecants.
\end{solution}

\begin{exerciseintro}
% record: H-IV-06-008
\exerciseheading{Exercise IV.6.8}
\addcontentsline{toc}{subsection}{Exercise IV.6.8}

\begin{context}
Work over an algebraically closed field.
A divisor \(D\) is nonspecial if \(H^1(X,\mathcal O_X(D))=0\).
Base-point-free means that its global sections generate
the line bundle, so an empty linear system does not qualify.
We use Riemann--Roch, Serre duality, and the irreducibility
and dimension \(g\) of \(\operatorname{Pic}^d(X)\).
Effective divisors of degree \(e\geq0\) are parametrized
by the proper \(e\)-dimensional variety \(X^{(e)}\);
there are none of negative degree.
On \(\mathbb P^1\), \(\mathcal O(d)\) is generated by
global sections exactly when \(d\geq0\).
\end{context}

\begin{problem}
For \(g\geq1\), prove that a genus-\(g\) curve admits a
nonspecial base-point-free divisor of degree \(d\)
exactly when \(d\geq g+1\).
For \(g=0\), identify the additional case omitted by
the literal statement in the source.
\end{problem}
\end{exerciseintro}

\begin{solution}
Suppose \(L=\mathcal O_X(D)\) is nonspecial and generated
by its sections. Then Riemann--Roch gives
\(h^0(L)=d+1-g\geq1\).
If \(h^0(L)=1\), its sole section must vanish nowhere,
so \(L\cong\mathcal O_X\).
Nonspeciality would then say
\(g=h^1(\mathcal O_X)=0\), and necessarily \(d=0\).
Apart from this case, \(h^0(L)\geq2\), so \(d\geq g+1\).
This proves necessity and identifies the possible exception.

Now let \(g\geq1\) and \(d\geq g+1\).
The special bundles in \(\operatorname{Pic}^d(X)\)
have the form
\[
L\cong\omega_X(-E),\qquad
E\text{ effective of degree }2g-2-d.
\]
When this degree is nonnegative their locus is a
proper closed image of dimension at most
\(2g-2-d\leq g-3<g\); otherwise it is empty.

For a nonspecial \(L\), the exact sequence at a point \(P\)
shows that \(P\) is a base point exactly when
\(H^1(L(-P))\ne0\). In detail,
\[
H^0(L)\to L|_P\to H^1(L(-P))\to H^1(L)=0;
\]
the first map is zero or surjective because its target
has dimension one.
By Serre duality, this exceptional condition gives
\[
L\cong\omega_X(P-E),\qquad
E\text{ effective of degree }2g-1-d.
\]
Allowing \(P\) and \(E\) to vary gives a closed image
of dimension at most
\(1+(2g-1-d)\leq g-1<g\), when the degree is
nonnegative, and an empty locus otherwise.
Irreducibility of the \(g\)-dimensional Picard variety
leaves a bundle outside both proper closed loci.
It is nonspecial and has no base points.
A line bundle on a smooth curve is represented by a
divisor, so it supplies the requested \(D\).

Finally, when \(g=0\), the curve is \(\mathbb P^1\).
The bundles \(\mathcal O(d)\) for all \(d\geq0\)
are generated and have \(H^1=0\), whereas no negative
degree bundle is generated.
Thus the correct condition is \(d\geq0\).
The source's inequality \(d\geq g+1\) omits precisely
the valid case \(g=0,d=0\).
\end{solution}

\begin{exerciseintro}
% record: H-IV-06-009
\exerciseheading{Exercise IV.6.9}
\addcontentsline{toc}{subsection}{Exercise IV.6.9}

\begin{context}
Let \(X\subset\mathbb P^3_k\) be a nonsingular integral
curve over an algebraically closed field.
Its ideal sheaf \(\mathcal I_X(t)\) is globally generated
for all sufficiently large \(t\), by Serre's theorem.
The conormal bundle \(\mathcal I_X/\mathcal I_X^2\)
has rank two. A hypersurface through a point of a smooth
ambient variety is singular there exactly when its
equation has zero first-order jet.
We use the dimension-of-fibers theorem and the facts
\(H^1(\mathbb P^3,\mathcal O(t))=0\) for all \(t\)
and \(H^0(\mathbb P^3,\mathcal O)=k\).
\end{context}

\begin{problem}
Prove that for every sufficiently large integer \(m\)
there is a nonsingular surface of degree \(m\) containing \(X\).
\end{problem}
\end{exerciseintro}

\begin{solution}
Choose \(m\) so large that both \(\mathcal I_X(m-1)\)
and \(\mathcal I_X(m)\) are globally generated, and put
\(W=H^0(\mathbb P^3,\mathcal I_X(m))\).
We exhibit two proper bad loci in the affine space \(W\).

For \(P\in X\), taking the first normal derivative gives
a surjection
\[
W\longrightarrow
(\mathcal I_X/\mathcal I_X^2)(m)|_P,
\]
whose target has dimension two.
Surjectivity follows from global generation and Nakayama's
lemma: the target is \(\mathcal I_{X,P}/\mathfrak m_P
\mathcal I_{X,P}\), with its twist.
A section vanishes on \(X\), so its tangential derivative
along \(X\) already vanishes. Its hypersurface is singular
at \(P\) exactly when its normal derivative is also zero.
The incidence of such pairs \((s,P)\in W\times X\)
therefore has dimension
\[
(\dim W-2)+1=\dim W-1.
\]
Since \(X\) is proper, its image in \(W\) is closed
and proper.

For \(P\notin X\), global generation of
\(\mathcal I_X(m-1)\) supplies a section \(h\)
with \(h(P)\ne0\).
The products \(h\ell\), for ambient linear forms \(\ell\),
belong to \(W\). Linear forms give all first-order
jets at \(P\), and multiplication by the unit \(h\)
is invertible on the local first-jet space.
Hence
\[
W\longrightarrow\mathcal O_{\mathbb P^3}(m)
\otimes(\mathcal O_{\mathbb P^3,P}/\mathfrak m_P^2)
\]
is surjective; its target has dimension four.
The incidence of a section with zero first-order jet
at some \(P\in\mathbb P^3\setminus X\) has dimension
\[
(\dim W-4)+3=\dim W-1.
\]
The closure of its image in \(W\) is again proper.

Choose \(s\) outside the union of these two proper
closed subsets. Its hypersurface \(F\) is smooth at
every point, contains \(X\), and has degree \(m\).
The choice is nonzero because the zero section is in
both bad loci.
For completeness \(F\) is integral: the hypersurface
sequence and \(H^1(\mathcal O(-m))=0\) give
\(H^0(F,\mathcal O_F)=k\), so it is connected.
Distinct components of a smooth scheme cannot intersect,
and smoothness also implies reducedness. Thus this
connected smooth hypersurface has just one integral
component. The construction works for every \(m\)
beyond the fixed global-generation threshold.
\end{solution}

\corpuschapter{Chapter V}
\addcontentsline{toc}{chapter}{Chapter V}

\begin{exerciseintro}
\corpussection{V.1 Geometry on a Surface}
\addcontentsline{toc}{section}{V.1 Geometry on a Surface}

% record: H-V-01-001
\exerciseheading{Exercise V.1.1}
\addcontentsline{toc}{subsection}{Exercise V.1.1}

\begin{context}
Let \(X\) be a nonsingular projective integral surface over
an algebraically closed field, with canonical divisor \(K\).
For a divisor \(A\), surface Riemann--Roch says
\[
\chi(\mathcal O_X(A))=\chi(\mathcal O_X)+\frac{A(A-K)}2.
\]
Intersection is symmetric, bilinear, and invariant under
linear equivalence. The Euler characteristic is the
alternating sum of coherent-cohomology dimensions.
\end{context}

\begin{problem}
For divisors \(C,D\), put \(L=\mathcal O_X(C)\) and
\(M=\mathcal O_X(D)\). Prove
\[
C D=\chi(\mathcal O_X)-\chi(L^{-1})-\chi(M^{-1})
+\chi(L^{-1}\otimes M^{-1}).
\]
\end{problem}
\end{exerciseintro}

\begin{solution}
Apply Riemann--Roch to \(-C,-D,-C-D\).
The four copies of \(\chi(\mathcal O_X)\) on the
right side cancel, leaving
\[
\frac{(C+D)(C+D+K)-C(C+K)-D(D+K)}2.
\]
By bilinearity the numerator is \(2CD\), proving the
formula. The calculation concerns integers and rational
numbers in the Riemann--Roch formula, so it is valid also
when the ground field has characteristic two.
\end{solution}

\begin{exerciseintro}
% record: H-V-01-002
\exerciseheading{Exercise V.1.2}
\addcontentsline{toc}{subsection}{Exercise V.1.2}

\begin{context}
Let \(X\) be a nonsingular projective integral surface
over an algebraically closed field, and \(H\) a very
ample divisor defining an embedding into \(\mathbb P^N\).
We use
\(\chi(\mathcal O_X(A))=\chi(\mathcal O_X)+A(A-K_X)/2\),
adjunction \(2g(H)-2=H(H+K_X)\) for a smooth member,
and Bertini's existence of such a member.
The Hilbert polynomial equals \(\chi(\mathcal O_X(nH))\).
For a dimension-\(r\) projective scheme its leading
coefficient is its degree divided by \(r!\).
The arithmetic genus of a surface is
\(p_a(X)=\chi(\mathcal O_X)-1\).
Here a curve means an integral one-dimensional closed subscheme;
on a nonsingular surface it is an effective Cartier divisor.
\end{context}

\begin{problem}
Write the Hilbert polynomial as
\(F(z)=az^2/2+bz+c\), and let \(\pi\) be the genus
of a smooth curve in \(|H|\).
Show \(a=H^2\), \(b=H^2/2+1-\pi\), and
\(c=1+p_a(X)\). Deduce \(\deg X=H^2\) and
\(\deg C=CH\) for every curve \(C\subset X\).
\end{problem}
\end{exerciseintro}

\begin{solution}
Riemann--Roch gives the polynomial identity
\[
F(z)=\frac{H^2}{2}z^2-\frac{HK_X}{2}z+
\chi(\mathcal O_X).
\]
Thus \(a=H^2\) and \(c=1+p_a(X)\).
Adjunction gives \(HK_X=2\pi-2-H^2\), so
\(b=-HK_X/2=H^2/2+1-\pi\).
By the leading-coefficient definition of degree,
\(\deg X=2!(H^2/2)=H^2\).

Every curve on a nonsingular surface is a Cartier
divisor. Its exact sequence, twisted by \(nH\), gives
\[
\chi(\mathcal O_C(nH))
=\chi(\mathcal O_X(nH))-\chi(\mathcal O_X(nH-C)).
\]
Substitution in Riemann--Roch makes the right side
\[
nHC-\frac{C^2+CK_X}{2}.
\]
This is the Hilbert polynomial of \(C\), whose linear
coefficient is its projective degree. Hence \(\deg C=HC\),
without any smoothness assumption on \(C\).
\end{solution}

\begin{exerciseintro}
% record: H-V-01-003
\exerciseheading{Exercise V.1.3}
\addcontentsline{toc}{subsection}{Exercise V.1.3}

\begin{context}
Let \(X\) be a nonsingular projective integral surface
over an algebraically closed field, with canonical divisor
\(K\). For a nonzero effective divisor \(D\),
\(p_a(D)=1-\chi(\mathcal O_D)\), whether or not \(D\)
is reduced or connected.
We use the exact sequence
\(0\to\mathcal O_X(-D)\to\mathcal O_X\to\mathcal O_D\to0\)
and surface Riemann--Roch.
Intersection depends only on linear equivalence.
For the empty divisor, the convention
\(1-\chi(\mathcal O_\varnothing)=1\) will agree with
the virtual genus below.
\end{context}

\begin{problem}
(a) Prove \(2p_a(D)-2=D(D+K)\) for an effective divisor.
(b) Deduce invariance under linear equivalence.
(c) Define the virtual arithmetic genus of any divisor
by this same formula and prove
\[
p_a(-D)=D^2-p_a(D)+2,\qquad
p_a(C+D)=p_a(C)+p_a(D)+CD-1.
\]
\end{problem}
\end{exerciseintro}

\begin{solution}
(a) Additivity of Euler characteristic gives
\[
\chi(\mathcal O_D)
=\chi(\mathcal O_X)-\chi(\mathcal O_X(-D))
=-\frac{D(D+K)}2.
\]
Since \(p_a(D)=1-\chi(\mathcal O_D)\), this is
exactly the desired formula. It also yields
\(p_a(0)=1\) under the stated empty-divisor convention.

(b) Linearly equivalent divisors have the same
self-intersection and the same intersection with \(K\).
The formula from (a) therefore gives the same arithmetic
genus for linearly equivalent effective divisors.

(c) Set \(p_a(A)=1+(A^2+AK)/2\) for every divisor \(A\).
This is an integer: Riemann--Roch applied to \(-A\)
shows that \(A(A+K)\) is even.
Then
\[
p_a(-D)+p_a(D)=2+D^2,
\]
which is the first identity. Similarly,
\[
\begin{split}
p_a(C+D)
&=1+\frac{C^2+D^2+2CD+CK+DK}{2}\\
&=p_a(C)+p_a(D)+CD-1.
\end{split}
\]
Thus the virtual extension agrees with the ordinary
genus on effective divisors and has both stated identities.
\end{solution}

\begin{exerciseintro}
% record: H-V-01-004
\exerciseheading{Exercise V.1.4}
\addcontentsline{toc}{subsection}{Exercise V.1.4}

\begin{context}
Work over an algebraically closed field.
For a nonsingular degree-\(d\) surface
\(X\subset\mathbb P^3\), \(K_X=(d-4)H\), where \(H\)
is the hyperplane class. Adjunction on a nonsingular
curve \(C\subset X\) gives \(2g(C)-2=C(C+K_X)\).
A line has genus zero and intersects \(H\) in degree one.
We use the projective hypersurface Jacobian criterion.
\end{context}

\begin{problem}
(a) If a nonsingular degree-\(d\) surface contains a line
\(C\), prove \(C^2=2-d\).
(b) In characteristic zero, construct such a surface
for every \(d\geq1\), containing the line \(x=y=0\).
\end{problem}
\end{exerciseintro}

\begin{solution}
(a) Adjunction gives
\[
-2=C^2+(d-4)CH=C^2+d-4,
\]
so \(C^2=2-d\).

(b) For \(d=1\), take the plane \(x=0\).
For \(d\geq2\), use homogeneous coordinates
\([x:y:z:w]\) and take
\[
F=xz^{d-1}+yw^{d-1}+x^d+y^d=0.
\tag{1}
\]
It contains \(x=y=0\).
When \(d=2\), the four partial derivatives are
\(z+2x,\ w+2y,\ x,\ y\); they cannot vanish at a
projective point.
When \(d\geq3\), vanishing of \(F_z=(d-1)xz^{d-2}\)
forces \(x=0\) or \(z=0\).
In either case vanishing of
\(F_x=z^{d-1}+dx^{d-1}\) forces both \(x=z=0\).
The equations for \(y,w\) likewise force \(y=w=0\).
Characteristic zero ensures that \(d\) and \(d-1\)
are nonzero. Thus the gradient never vanishes at a
projective point, and (1) defines a nonsingular surface.
A reducible hypersurface in \(\mathbb P^3\) would have
intersecting components and would be singular there.
Thus this surface is integral as well.
\end{solution}

\begin{exerciseintro}
% record: H-V-01-005
\exerciseheading{Exercise V.1.5}
\addcontentsline{toc}{subsection}{Exercise V.1.5}

\begin{context}
Work over an algebraically closed field.
For a nonsingular degree-\(d\) surface in \(\mathbb P^3\),
adjunction gives \(K_X=(d-4)H\), and \(H^2=d\).
For smooth curves \(C,C'\), the cotangent bundle of
their product is
\(p_1^*\Omega_C\oplus p_2^*\Omega_{C'}\).
A curve's canonical degree is \(2g-2\).
Fibers \(l=C\times\{P'\}\), \(m=\{P\}\times C'\)
satisfy \(l^2=m^2=0,\ lm=1\).
Pullbacks of divisors from either curve have intersections
determined by their degrees.
\end{context}

\begin{problem}
(a) For a nonsingular surface of degree \(d\) in
\(\mathbb P^3\), show \(K_X^2=d(d-4)^2\).
(b) For \(X=C\times C'\), with genera \(g,g'\),
show \(K_X^2=8(g-1)(g'-1)\).
\end{problem}
\end{exerciseintro}

\begin{solution}
(a) Bilinearity and adjunction give
\[
K_X^2=((d-4)H)^2=(d-4)^2H^2=d(d-4)^2.
\]

(b) Taking determinants of the product cotangent bundle
gives
\[
\omega_X=p_1^*\omega_C\otimes p_2^*\omega_{C'}.
\]
Thus the canonical divisor is the sum of the two
canonical pullbacks. Each pullback has self-intersection
zero, since it is numerically a multiple of a fiber.
Their mutual intersection is the product of their
degrees, \((2g-2)(2g'-2)\).
Therefore
\[
K_X^2=2(2g-2)(2g'-2)=8(g-1)(g'-1).
\]
The same calculation includes genus zero and genus one.
\end{solution}

\begin{exerciseintro}
% record: H-V-01-006
\exerciseheading{Exercise V.1.6}
\addcontentsline{toc}{subsection}{Exercise V.1.6}

\begin{context}
Let \(C\) be a nonsingular projective integral curve
of genus \(g\) over an algebraically closed field.
For the diagonal \(\Delta\subset C\times C\),
its conormal sheaf is
\(\mathcal I_\Delta/\mathcal I_\Delta^2\cong\Omega_C\).
The self-intersection of a smooth divisor is the degree
of its normal bundle.
Numerical equivalence means equality of intersections
with all curves; \(\operatorname{Num}(X)\) is the group
of divisor classes modulo this relation.
For the fibers \(l=C\times\{P\}\), \(m=\{P\}\times C\),
\(l^2=m^2=0\) and \(lm=1\).
A pullback of a degree-\(d\) line bundle from one factor
is numerically \(d\) times the corresponding fiber.
\end{context}

\begin{problem}
(a) Prove \(\Delta^2=2-2g\).
(b) For \(g\geq1\), prove that \(l,m,\Delta\) are
linearly independent in \(\operatorname{Num}(C\times C)\).
Deduce that the numerical group has rank at least three
and that pullbacks from the two factors do not generate
\(\operatorname{Pic}(C\times C)\).
\end{problem}
\end{exerciseintro}

\begin{solution}
(a) Dualizing the conormal identification gives
\(N_{\Delta/(C\times C)}\cong T_C\).
Its degree is \(-\deg\Omega_C=2-2g\), so this is
\(\Delta^2\).

(b) Each fiber meets the diagonal once, transversely.
The intersection matrix on \(l,m,\Delta\) is therefore
\[
\begin{pmatrix}
0&1&1\\
1&0&1\\
1&1&2-2g
\end{pmatrix},
\]
with determinant \(2g\).
For \(g\geq1\), this integer is nonzero.
If an integral linear combination of these classes
were numerically zero, pairing it with each of the
three classes would give a homogeneous linear system
with this invertible matrix over \(\mathbb Q\).
All coefficients must vanish. This proves independence,
regardless of the characteristic of the ground field:
the intersection matrix has integer entries, not entries
reduced modulo that characteristic.

If every line bundle on \(C\times C\) were a tensor
product of pullbacks, \(\mathcal O(\Delta)\) would
have the form \(p_1^*L\otimes p_2^*M\).
Its numerical class would then lie in the span of
\(m\) and \(l\), contradicting the independence just
proved. Hence these pullbacks do not generate the
Picard group.
\end{solution}

\begin{exerciseintro}
% record: H-V-01-007
\exerciseheading{Exercise V.1.7}
\addcontentsline{toc}{subsection}{Exercise V.1.7}

\begin{context}
Let \(X\) be a nonsingular projective integral surface
over an algebraically closed field.
Two effective divisors are prealgebraically equivalent
if they occur as fibers of an effective Cartier divisor
on \(X\times T\), flat over a nonsingular connected
curve \(T\). For arbitrary divisors, extend this by
differences of such pairs; algebraic equivalence is
the equivalence generated by finite chains of these pairs.
A relative effective Cartier divisor on the flat
family \(X\times T\) is flat when each fiber is Cartier.
Hilbert polynomials are constant in a projective flat
family over connected \(T\).
For very ample \(H\), the linear coefficient of a
divisor's Hilbert polynomial is its intersection with \(H\).
Every divisor becomes very ample after adding a
sufficiently large multiple of a fixed very ample divisor.
\end{context}

\begin{problem}
(a) Prove that the divisors algebraically equivalent
to zero form a subgroup of \(\operatorname{Div}(X)\).
(b) Show that linear equivalence implies algebraic
equivalence.
(c) Show that algebraic equivalence implies numerical
equivalence.
\end{problem}
\end{exerciseintro}

\begin{solution}
(a) Constant families show that every effective divisor
is prealgebraically equivalent to itself.
Adding a fixed effective divisor to a family gives a
new flat family: its fiber equation is the product
of the old Cartier equation and that of the fixed
divisor, again a nonzerodivisor.
Write any fixed divisor as \(A-B\) with \(A,B\)
effective. By adding the corresponding constant
families to the positive and negative pairs, translation
by \(A-B\) preserves prealgebraic equivalence.
Interchanging positive and negative pairs shows
compatibility with negation. These assertions also
hold for finite chains.

In particular, zero is algebraically equivalent to
itself. If \(D\) and \(E\) are each equivalent to zero,
translate a chain from \(E\) to zero by \(D\), then
append a chain from \(D\) to zero. This connects
\(D+E\) to zero. Negating a chain connects \(-D\)
to zero. Hence the zero class is a subgroup.

(b) Suppose \(D\sim D'\). Add a sufficiently positive
fixed divisor \(A\) so that \(D+A,D'+A\) are effective.
They are zero divisors of sections \(s_0,s_1\) of the
same line bundle \(L\).
If the sections are proportional, the effective
divisors coincide. Otherwise the pencil
\[
u s_0+v s_1,\qquad [u:v]\in\mathbb P^1,
\]
defines an effective Cartier divisor on
\(X\times\mathbb P^1\).
No linear combination at a parameter point is the
zero section, so each fiber is an effective Cartier
divisor on \(X\); the total divisor is flat over
\(\mathbb P^1\). Its fibers at \([1:0]\) and \([0:1]\)
are \(D+A\) and \(D'+A\).
Thus these are prealgebraically equivalent, and
subtracting the constant divisor \(A\) proves
\(D\) algebraically equivalent to \(D'\).

(c) Fix a very ample divisor \(H\).
In a flat family of effective divisors, constancy
of the Hilbert polynomial gives constancy of its
linear coefficient, hence of intersection with \(H\).
Taking differences and chains shows
\[
DH=D'H
\]
whenever \(D,D'\) are algebraically equivalent.
This holds for every very ample \(H\).

Put \(B=D-D'\) and let \(C\) be any curve on \(X\).
Choose \(n\) so large that both \(C+nH\) and \(nH\)
are very ample. Then
\[
BC=B(C+nH)-B(nH)=0-0=0.
\]
Thus \(B\) intersects every curve in zero, which is
exactly numerical equivalence of \(D\) and \(D'\).
\end{solution}

\begin{exerciseintro}
% record: H-V-01-008
\exerciseheading{Exercise V.1.8}
\addcontentsline{toc}{subsection}{Exercise V.1.8}

\begin{context}
Let \(X\) be a nonsingular projective integral surface
over an algebraically closed field \(k\).
Applying \(H^1\) to \(d\log:\mathcal O_X^\times\to\Omega_X^1\)
defines \(c:\operatorname{Pic}(X)\to H^1(X,\Omega_X^1)\).
It is additive and compatible with pullback.
The pairing on \(H^1(X,\Omega_X^1)\) is cup product,
followed by wedge of forms and the Serre trace on
\(H^2(X,\omega_X)\).

We use the divisor-class identity from Hartshorne III.7.4:
if \(i:Y\hookrightarrow X\) is a smooth divisor, then
\[
\operatorname{tr}_X(c(Y)\smile\alpha)
=\operatorname{tr}_Y(i^*\alpha).
\]
On a smooth projective curve,
\(\operatorname{tr}(c(\mathcal O(P)))=1\).
Every divisor is linearly equivalent to a difference
of smooth curves, by twisting to very ample systems
and applying Bertini.
Intersection with a smooth curve is the degree of
the restricted line bundle.
\end{context}

\begin{problem}
(a) Prove \((c(D),c(E))=(DE)\cdot1\) in \(k\).
(b) In characteristic zero, deduce that
\(\operatorname{Num}(X)\) is a finitely generated
free abelian group from finite dimensionality of
\(H^1(X,\Omega_X^1)\).
\end{problem}
\end{exerciseintro}

\begin{solution}
(a) First suppose \(D\) is a smooth curve.
The trace identity in the context and functoriality
of \(d\log\) give
\[
(c(D),c(E))
=\operatorname{tr}_D(c(\mathcal O_X(E)|_D)).
\]
Write the restricted line bundle as the line bundle
of a divisor \(\sum n_iP_i\) on \(D\).
Additivity of \(c\), and the trace-one normalization
for a closed point, make the right side
\[
\sum n_i\cdot1=\deg_D(\mathcal O_X(E)|_D)\cdot1
=(DE)\cdot1.
\]
For arbitrary \(D\), express it up to linear equivalence
as \(D_1-D_2\) with both \(D_i\) smooth and apply
bilinearity. No smoothness condition on \(E\) is needed.
This proves the formula in every characteristic,
where the integer on the right is mapped into \(k\).

(b) Let \(N=\operatorname{Num}(X)\) and
\(V=N\otimes_{\mathbb Z}\mathbb Q\).
The group \(N\) is torsion-free: if \(nD\) is
numerically zero with \(n\ne0\), its integer
intersections show that \(D\) is numerically zero.
The intersection pairing on \(V\) has zero radical
by the definition of numerical equivalence.

Put \(r=\dim_k H^1(X,\Omega_X^1)\).
We first prove \(\dim_{\mathbb Q}V\leq r\), without
assuming any finite generation.
Take \(n\) independent divisor classes \(D_1,\ldots,D_n\)
in \(V\). The functions
\[
E\longmapsto D_iE
\]
are independent rational linear functionals: a
dependence would put a nonzero combination of the
\(D_i\) in the radical.
Hence there exist divisor classes \(E_1,\ldots,E_n\)
such that the matrix \((D_iE_j)\) is invertible
over \(\mathbb Q\). Denominators can be cleared so
all chosen classes are represented by integral divisors.

In characteristic zero, its nonzero integer determinant
remains nonzero in \(k\).
By (a) it is the pairing matrix of
\(c(D_1),\ldots,c(D_n)\) with \(c(E_1),\ldots,c(E_n)\).
Thus the \(c(D_i)\) are independent over \(k\),
and \(n\leq r\). This proves the finite-rank assertion.

Now choose integral divisors \(A_1,\ldots,A_s\)
whose numerical classes form a rational basis of \(V\).
The homomorphism
\[
N\longrightarrow\mathbb Z^s,\qquad
[D]\longmapsto(DA_1,\ldots,DA_s)
\]
is injective: vanishing of these pairings forces
vanishing against all of \(V\), hence against every
curve. A subgroup of the finitely generated free
abelian group \(\mathbb Z^s\) is itself finitely
generated and free. This proves the assertion.
The integral-pairing step is essential; finite
rational rank alone would not imply finite generation.
\end{solution}

\begin{exerciseintro}
% record: H-V-01-009
\exerciseheading{Exercise V.1.9}
\addcontentsline{toc}{subsection}{Exercise V.1.9}

\begin{context}
Let \(X\) be a nonsingular projective integral surface
over an algebraically closed field.
An ample divisor \(H\) has \(H^2>0\).
The Hodge index theorem says that \(A H=0\) implies
\(A^2\leq0\) for a rational divisor class \(A\).
Numerical equivalence means equality of intersections
with all divisors.
On \(C\times C'\), the fibers
\(l=C\times\{P'\}\), \(m=\{P\}\times C'\)
satisfy \(l^2=m^2=0,\ lm=1\).
Positive-degree line bundles on curves are ample,
and the exterior tensor product of ample line bundles
is ample, as follows by taking very ample powers
and using the Segre embedding.
\end{context}

\begin{problem}
(a) Prove \(D^2H^2\leq(DH)^2\) for ample \(H\).
(b) On \(X=C\times C'\), say \(D\) has type
\((a,b)\) if \(Dl=a,\ Dm=b\).
Show \(D^2\leq2ab\), with equality exactly when
\(D\) is numerically equivalent to \(bl+am\).
\end{problem}
\end{exerciseintro}

\begin{solution}
(a) Put \(A=D-(DH/H^2)H\).
Then \(AH=0\), so the Hodge index theorem gives
\[
0\geq A^2=D^2-\frac{(DH)^2}{H^2}.
\]
Multiplying by \(H^2>0\) proves the inequality.

We also record its equality consequence.
If \(AH=0\) and \(A^2=0\), take any rational divisor
class \(B\) orthogonal to \(H\).
For every rational \(t\), Hodge index gives
\[
(A+tB)^2=2tAB+t^2B^2\leq0.
\]
For both signs of sufficiently small nonzero \(t\),
this is possible only if \(AB=0\).
An arbitrary divisor is the sum of a rational
multiple of \(H\) and a class orthogonal to \(H\).
Thus \(A\) intersects every divisor in zero; it is
numerically trivial. Conversely, a numerically
trivial \(A\) has \(A^2=0\).

(b) The divisor \(H=l+m\) is ample by the exterior
product criterion in the context.
Set
\[
R=D-bl-am.
\]
Using the definitions of \(a,b\) gives \(Rl=Rm=0\),
hence \(RH=0\). Its square is
\[
R^2=D^2-2ab.
\]
Hodge index therefore gives \(D^2\leq2ab\).
The equality consequence proved in (a) says that
\(R^2=0\) exactly when \(R\) is numerically trivial,
which is exactly \(D\equiv bl+am\).
\end{solution}

\begin{exerciseintro}
% record: H-V-01-010
\exerciseheading{Exercise V.1.10}
\addcontentsline{toc}{subsection}{Exercise V.1.10}

\begin{context}
Let \(C_0\) be a smooth projective geometrically integral
curve of genus \(g\) over \(\mathbb F_q\), and put
\(C=C_0\times_{\mathbb F_q}\overline{\mathbb F}_q\).
Base change of the \(q\)-power Frobenius of \(C_0\) defines an
\(\overline{\mathbb F}_q\)-morphism \(f:C\to C\) of degree \(q\),
with zero differential
and fixed points precisely \(C_0(\mathbb F_q)\).
For a smooth divisor its square is the degree of
its normal bundle. A transverse intersection contributes one.
The diagonal in \(C\times C\) has square \(2-2g\).
For \(l=C\times\{P\}\), \(m=\{P\}\times C\), the
product form of Hodge index says
\[
D^2\leq2(Dl)(Dm)
\]
for every divisor \(D\). We use \(\deg T_C=2-2g\).
\end{context}

\begin{problem}
Let \(N=\#C_0(\mathbb F_q)\), \(\Gamma\) be the graph
of \(f\), and \(\Delta\) the diagonal.
Show \(\Gamma^2=q(2-2g)\), \(\Gamma\Delta=N\), and
derive
\[
N=1-a+q,\qquad |a|\leq2g\sqrt q.
\]
\end{problem}
\end{exerciseintro}

\begin{solution}
Identify \(\Gamma\) with \(C\) by its first projection.
The tangent sequence of its graph identifies its normal
bundle as the quotient of \(T_C\oplus f^*T_C\)
by vectors \((v,df(v))\).
The map \((v,w)\mapsto w-df(v)\) identifies that
quotient with \(f^*T_C\).
Therefore
\[
\Gamma^2=\deg f^*T_C=q(2-2g).
\]
The graph and diagonal intersect at exactly the fixed
points of \(f\), namely the \(\mathbb F_q\)-points.
At such a point their tangent lines are
\(\{(v,0)\}\) and \(\{(v,v)\}\), since \(df=0\).
They are distinct and span the surface tangent space,
so every intersection is transverse. Hence
\(\Gamma\Delta=N\).

The intersections with the fibers are
\[
\Gamma l=q,\quad\Gamma m=1,\quad
\Delta l=\Delta m=1.
\]
For arbitrary integers \(r,s\), apply the product
inequality to \(D=r\Gamma+s\Delta\). It gives
\[
r^2q(2-2g)+2rsN+s^2(2-2g)
\leq2(rq+s)(r+s).
\]
Set \(a=q+1-N\) and rearrange:
\[
gq\,r^2+a rs+g s^2\geq0
\quad\text{for all integers }r,s.
\tag{1}
\]
By homogeneity it holds for rational \(r,s\), and
by continuity for real \(r,s\).
If \(g>0\), the discriminant of this nonnegative
quadratic form is nonpositive, so
\(a^2-4g^2q\leq0\).
If \(g=0\), (1) says \(a rs\geq0\) for both signs
of \(rs\), forcing \(a=0\).
Thus in every genus \(N=q+1-a\) with
\(|a|\leq2g\sqrt q\).
\end{solution}

\begin{exerciseintro}
% record: H-V-01-011
\exerciseheading{Exercise V.1.11}
\addcontentsline{toc}{subsection}{Exercise V.1.11}

\begin{context}
Let \(X\) be a nonsingular projective integral surface
over an algebraically closed field. Assume its numerical
divisor group \(N\) is finitely generated.
It is torsion-free, so its image in \(N_{\mathbb R}\)
is a lattice. For ample \(H\), intersection is
negative definite on \(H^\perp\).
An integral curve \(C\) has
\(C^2+CK_X=2p_a(C)-2\geq-2\).
Distinct integral curves have nonnegative intersection.
On a product \(C\times C\), the fibers \(l,m\)
have \(l+m\) ample and the diagonal has square \(2-2g(C)\).
We use the degree and intersection rules for graphs
of morphisms of curves.
\end{context}

\begin{problem}
(a) For ample \(H\) and fixed integer \(d\), prove
that effective divisors \(D\) with \(DH=d\) have only
finitely many numerical classes.
(b) Assuming this finite-generation hypothesis for
\(C\times C\), deduce that a curve \(C\) of genus
at least two has finite automorphism group.
\end{problem}
\end{exerciseintro}

\begin{solution}
(a) First bound integral curves of fixed degree
\(e=CH>0\). Set \(h=H^2>0,\ k=K_XH\), and
write in \(N_{\mathbb R}\)
\[
C=\frac e h H+v,\qquad
K_X=\frac k h H+w,\qquad v,w\in H^\perp.
\]
On \(H^\perp\) use the positive-definite inner
product \(\langle v,w\rangle=-vw\).
Adjunction yields
\[
-2\leq C^2+CK_X
=\frac{e^2+ek}{h}-\|v\|^2-\langle v,w\rangle.
\]
Completing the square gives
\[
\left\|v+\frac w2\right\|^2
\leq2+\frac{e^2+ek}{h}+\frac{\|w\|^2}{4}.
\tag{1}
\]
For fixed \(e\), this is a bounded closed ball
(or an empty set if the right side is negative).
Since the numerical classes of the corresponding
\(C\)'s lie in the affine hyperplane \(CH=e\)
inside the lattice \(N\), there are only finitely many.
This remains true when \(H^\perp\) has dimension zero.

Now let \(D=\sum n_iC_i\) be an effective divisor
of degree \(d>0\). Each \(C_iH\) is a positive
integer at most \(d\), and
\(\sum n_i\leq d\).
By the first step there are finitely many possible
classes for all \(C_i\), and the multiplicity bound
leaves finitely many sums.
For \(d<0\) there are no effective divisors;
for \(d=0\) only the zero divisor occurs, by ampleness.
This proves (a).

(b) For an automorphism \(\alpha\) of \(C\), let
\(\Gamma_\alpha\subset C\times C\) be its graph.
Both projections have degree one, so
\(\Gamma_\alpha(l+m)=2\).
Part (a) gives finitely many numerical classes of
such graphs.

These classes determine the graphs themselves.
First, if an integral graph \(\Gamma\) is numerically
equivalent to \(\Delta\), then
\[
\Gamma\Delta=\Delta^2=2-2g<0.
\]
Distinct integral curves have nonnegative intersection,
so \(\Gamma=\Delta\).
More generally, if
\(\Gamma_\alpha\equiv\Gamma_\beta\), applying the
surface automorphism \(\operatorname{id}\times\alpha^{-1}\)
preserves numerical equivalence and takes the graphs
to \(\Delta\) and \(\Gamma_{\alpha^{-1}\beta}\).
The preceding argument gives \(\alpha^{-1}\beta=\operatorname{id}\).
Thus distinct automorphisms have distinct numerical graph
classes. Finiteness of those classes proves that
\(\operatorname{Aut}(C)\) is finite.
\end{solution}

\begin{exerciseintro}
% record: H-V-01-012
\exerciseheading{Exercise V.1.12}
\addcontentsline{toc}{subsection}{Exercise V.1.12}

\begin{context}
Work over an algebraically closed field.
For a nonsingular projective surface, Nakai--Moishezon
says that a divisor \(D\) is ample exactly when
\(D^2>0\) and \(DE>0\) for every integral curve \(E\).
The exterior tensor product of very ample line bundles
is very ample. Restricting a very ample line bundle
to a closed subvariety gives a very ample line bundle.
A smooth plane quartic has genus three and canonical
bundle \(\mathcal O_C(1)\).
Bertini gives smooth plane quartics in every characteristic.
We use curve Riemann--Roch and that a nontrivial
degree-zero line bundle has no nonzero section.
\end{context}

\begin{problem}
Prove that numerical equivalence preserves ampleness
of divisors on a nonsingular projective surface.
Give an example showing that it does not preserve
very ampleness.
\end{problem}
\end{exerciseintro}

\begin{solution}
If \(D'\equiv D\), their intersections with every
curve agree. Bilinearity also gives
\[
(D')^2=D^2,
\]
since \((D'-D)(D'+D)=0\).
Thus both conditions of Nakai--Moishezon hold for
\(D'\) whenever they hold for \(D\), proving the
ampleness assertion.

For the example, take a smooth plane quartic \(C\)
and distinct points \(P,Q\in C\).
The degree-zero bundle \(\mathcal O_C(Q-P)\) is
nontrivial: a rational function with divisor \(Q-P\)
would give a degree-one map to \(\mathbb P^1\),
contradicting genus three.
Let
\[
L=\omega_C,\qquad L'=\omega_C(P-Q).
\]
Both have degree four, but Riemann--Roch gives
\[
h^0(C,L')=4+1-3+h^0(C,\mathcal O_C(Q-P))=2.
\]
Thus \(L'\) cannot be very ample: its complete
system could map only to \(\mathbb P^1\), where
a genus-three curve cannot embed.
In contrast, \(L=\mathcal O_C(1)\) is very ample.

On the smooth surface \(X=C\times\mathbb P^1\),
with projections \(p_1,p_2\), set
\[
A=p_1^*L\otimes p_2^*\mathcal O_{\mathbb P^1}(1),
\qquad
A'=p_1^*L'\otimes p_2^*\mathcal O_{\mathbb P^1}(1).
\]
The bundle \(A\) is very ample by the Segre embedding.
Their quotient is \(p_1^*\mathcal O_C(P-Q)\),
which is numerically trivial on \(X\).
Indeed, on the normalization of any integral curve
in \(X\), its degree is zero: projection to \(C\)
is either constant, giving a trivial pullback, or
finite, multiplying the degree-zero divisor degree.
Hence the divisors of \(A,A'\) are numerically equivalent.

Finally, \(A'\) restricts on \(C\times\{t\}\)
to \(L'\), which is not very ample.
Therefore \(A'\) is not very ample, giving the
required counterexample.
\end{solution}

\begin{exerciseintro}
\corpussection{V.2 Ruled Surfaces}
\addcontentsline{toc}{section}{V.2 Ruled Surfaces}

% record: H-V-02-001
\exerciseheading{Exercise V.2.1}
\addcontentsline{toc}{subsection}{Exercise V.2.1}

\begin{context}
Work over an algebraically closed field.
A birationally ruled surface is birational to
\(C\times\mathbb P^1\), with \(C\) a smooth projective
integral curve.
We use Luroth's theorem: a curve dominated by
\(\mathbb P^1\) has rational function field, also
for an inseparable map.
Rational maps from a smooth projective curve to a
projective curve extend to morphisms, and birational
smooth projective curves are isomorphic.
A rational map from \(C\times\mathbb P^1\) constant
on the generic ruling fiber factors rationally
through \(C\): constancy on the generic fiber makes the pulled-back
functions algebraic over \(k(C)\), and \(k(C)\) is algebraically
closed in the purely transcendental extension \(k(C)(t)\).
Thus those functions lie in \(k(C)\).
\end{context}

\begin{problem}
Prove that the base curve of a birationally ruled
surface is unique up to isomorphism.
\end{problem}
\end{exerciseintro}

\begin{solution}
Suppose \(C\times\mathbb P^1\) is birational to
\(C'\times\mathbb P^1\).
If \(C'\) has positive genus, compose the birational
map with projection to \(C'\).
On the generic ruling fiber this is a rational
map from \(\mathbb P^1_{k(C)}\) to the base change
of \(C'\). It must be constant by Luroth's theorem,
since the latter curve has positive genus.
Thus the composite factors through a dominant
rational map \(u:C\dashrightarrow C'\).

If both curves have positive genus, the inverse
birational map similarly gives \(v:C'\dashrightarrow C\).
Composing with the original birational maps shows
\(vu=\operatorname{id}_C\) and
\(uv=\operatorname{id}_{C'}\) as rational maps:
both are the maps induced on the respective bases
by the identity map of the product surface.
Therefore \(u,v\) are inverse birational maps of
smooth projective curves, hence isomorphisms.

It remains to handle a rational base.
If \(C\cong\mathbb P^1\) and \(C'\) had positive
genus, the same composite
\(\mathbb P^1\times\mathbb P^1\dashrightarrow C'\)
would be constant on each ruling. In function fields,
its pullback would lie in both \(k(s)\) and \(k(t)\)
inside \(k(s,t)\), whose intersection is \(k\).
It would be constant, contradicting dominance.
Thus both bases have genus zero, and over the
algebraically closed field both are \(\mathbb P^1\).
This completes the uniqueness proof.
\end{solution}

\begin{exerciseintro}
% record: H-V-02-002
\exerciseheading{Exercise V.2.2}
\addcontentsline{toc}{subsection}{Exercise V.2.2}

\begin{context}
Let \(\mathcal E\) be a rank-two vector bundle on
a smooth projective integral curve \(C\) over an
algebraically closed field.
Use the quotient convention for \(X=\mathbb P(\mathcal E)\):
a section corresponds to a line-bundle quotient
\(\mathcal E\twoheadrightarrow L\).
Two such quotient maps represent the same point
of a fiber precisely when their one-dimensional
kernels agree.
A morphism of vector bundles that is an isomorphism
on every fiber is an isomorphism, by Nakayama's lemma.
\end{context}

\begin{problem}
Prove that \(\mathcal E\) is decomposable exactly
when \(X\to C\) has two disjoint sections.
\end{problem}
\end{exerciseintro}

\begin{solution}
If \(\mathcal E=L_1\oplus L_2\), its two projection
quotients give sections. In every fiber their
kernels are the two different direct summands,
so the sections have no common point.

Conversely, let two disjoint sections correspond
to quotients \(q_i:\mathcal E\twoheadrightarrow L_i\).
The map
\[
(q_1,q_2):\mathcal E\longrightarrow L_1\oplus L_2
\]
is an isomorphism on every fiber: its two quotient
functionals have distinct kernels and hence are
independent. By the fiberwise criterion it is
an isomorphism of bundles. Thus \(\mathcal E\)
decomposes into the two line bundles.
\end{solution}

\begin{exerciseintro}
% record: H-V-02-003
\exerciseheading{Exercise V.2.3}
\addcontentsline{toc}{subsection}{Exercise V.2.3}

\begin{context}
Work over an algebraically closed field.
On a nonsingular curve, a coherent torsion-free
sheaf is locally free, since its local modules
are torsion-free over discrete valuation rings.
Saturating a subsheaf means enlarging it until
the quotient is torsion-free.
On \(\mathbb P^2\), every line bundle is
\(\mathcal O(a)\), \(H^1(\mathcal O(a))=0\) for
every integer \(a\), and
\[
0\to\Omega_{\mathbb P^2}^1(1)\to
\mathcal O^{\oplus3}\to\mathcal O(1)\to0
\]
is the Euler sequence.
Extensions of line bundles are classified by
\(\operatorname{Ext}^1(\mathcal O(b),\mathcal O(a))
=H^1(\mathcal O(a-b))\).
\end{context}

\begin{problem}
(a) Prove that every rank-\(r\) vector bundle on a
nonsingular curve has a filtration
\(0=\mathcal E_0\subset\mathcal E_1\subset\cdots
\subset\mathcal E_r=\mathcal E\)
whose successive quotients are line bundles.
(b) Show that this fails in dimension at least two;
in particular, \(\Omega_{\mathbb P^2}^1\) is not
an extension of line bundles.
\end{problem}
\end{exerciseintro}

\begin{solution}
(a) Choose a one-dimensional subspace of the
generic fiber of \(\mathcal E\). Extending a
generic generator after allowing poles, then
saturating its image in \(\mathcal E\), gives a
coherent rank-one subsheaf \(L\) with torsion-free
quotient. Locally this saturation is the intersection
of the generic line with the free module \(\mathcal E\);
over a discrete valuation ring it is a rank-one
direct summand. Hence \(L\) is a line bundle and
\(\mathcal E/L\) is a vector bundle of rank \(r-1\).
Induct on the rank, and pull the quotient's
filtration back to \(\mathcal E\).
For rank zero the empty filtration suffices.

(b) An extension
\(0\to\mathcal O(a)\to\Omega_{\mathbb P^2}^1
\to\mathcal O(b)\to0\)
would split, by the stated \(H^1\)-vanishing.
Taking determinants in the Euler sequence gives
\(\det\Omega_{\mathbb P^2}^1=\mathcal O(-3)\),
so \(a+b=-3\).
But the map on global sections
\(H^0(\mathcal O^{\oplus3})\to H^0(\mathcal O(1))\)
is an isomorphism, and therefore
\(H^0(\Omega_{\mathbb P^2}^1(1))=0\).
A splitting would force \(a+1<0\) and \(b+1<0\),
hence \(a,b\leq-2\), contradicting \(a+b=-3\).
Thus no such extension exists.

For each dimension \(n>2\), pull this rank-two
bundle back to \(\mathbb P^2\times\mathbb P^{n-2}\).
A filtration with line-bundle quotients would
restrict, still exactly because those quotients
are locally free, to such a filtration on a
fiber \(\mathbb P^2\). This is impossible by
the preceding argument. Hence the failure occurs
in every dimension at least two.
\end{solution}

\begin{exerciseintro}
% record: H-V-02-004
\exerciseheading{Exercise V.2.4}
\addcontentsline{toc}{subsection}{Exercise V.2.4}

\begin{context}
Let \(C\) be a smooth projective integral genus-\(g\)
curve over an algebraically closed field.
A section of \(C\times\mathbb P^1\to C\) is the
graph of a morphism \(h:C\to\mathbb P^1\).
The normal bundle of a graph is \(h^*T_{\mathbb P^1}\);
its degree is the self-intersection.
We use \(T_{\mathbb P^1}=\mathcal O(2)\) and the
existence, for \(r\geq g+1\), of a nonspecial
base-point-free degree-\(r\) line bundle.
A nonhyperelliptic genus-three curve is a canonical
plane quartic. A curve of type \((a,b)\) on
\(\mathbb P^1\times\mathbb P^1\) has arithmetic
genus \((a-1)(b-1)\).
\end{context}

\begin{problem}
For a section \(D\) of \(C\times\mathbb P^1\),
show \(D^2=2r\) with \(r\geq0\).
(a) Prove that \(r=0\) and every \(r\geq g+1\)
are possible.
(b) If \(g=3\), show that \(r=1\) is impossible,
and exactly one of \(r=2,3\) is possible, according
as \(C\) is hyperelliptic or nonhyperelliptic.
\end{problem}
\end{exerciseintro}

\begin{solution}
For the graph of \(h\), the normal-bundle formula
gives \(D^2=\deg h^*\mathcal O(2)=2r\), where
\(r=\deg h^*\mathcal O(1)\geq0\).
For nonconstant \(h\) this is its degree; for
constant \(h\) it is zero.

(a) Constant maps provide \(r=0\).
For \(r\geq g+1\), choose a generated nonspecial
degree-\(r\) line bundle \(L\).
Riemann--Roch gives \(h^0(L)=r+1-g\geq2\).
Choose a nonzero section \(s\); its zeros are finite.
Because \(L\) is generated, there is a second
section \(t\) nonzero at all of these zeros:
each excluded vanishing condition is a proper
linear subspace, and the field is infinite.
Thus \(s,t\) have no common zero and give a map
\(h:C\to\mathbb P^1\) with \(h^*\mathcal O(1)=L\).
It has degree \(r\), as desired.

(b) A degree-one map would make \(C\cong\mathbb P^1\),
so \(r=1\) is impossible.
If \(C\) is hyperelliptic its degree-two pencil
gives \(r=2\). It cannot also have a degree-three
map: the pair of maps would be birational onto
a curve of type \((2,3)\), since its degree onto
the image divides both two and three.
That image has arithmetic genus two, less than
the genus three of its normalization.

If \(C\) is nonhyperelliptic, \(r=2\) is impossible
by definition. Its canonical plane quartic
admits a degree-three map by projection from
any point on the curve: the lines through that
point give a pencil of hyperplane sections
with that one fixed point removed.
Smoothness makes the residual pencil base-point-free,
of degree \(4-1=3\). Thus \(r=3\) is possible.
\end{solution}

\begin{exerciseintro}
% record: H-V-02-005
\exerciseheading{Exercise V.2.5}
\addcontentsline{toc}{subsection}{Exercise V.2.5}

\begin{context}
Let \(C\) be a smooth projective integral genus-\(g\)
curve over an algebraically closed field, with \(g\geq1\).
A rank-two bundle is normalized if it has a section
and no positive-degree line subbundle; its ruled
surface invariant is \(e=-\deg\mathcal E\).
A nonzero section of a normalized bundle is nowhere
zero. A normalized decomposable bundle is
\(\mathcal O_C\oplus L\) with \(\deg L\leq0\).
We use Riemann--Roch, Serre duality, and
\(\operatorname{Ext}^1(L,\mathcal O_C)=H^1(L^{-1})\).
The curve symmetric power \(C^{(n)}\) has dimension \(n\).
Degree at least \(2g+1\) line bundles are very ample.
Projection of an embedded smooth curve from a point
is an embedding if that point lies on none of its
secant or tangent lines, interpreted by length-two
subschemes. A smooth plane quintic has genus six.
\end{context}

\begin{problem}
(a) For every \(0\leq e\leq2g-2\), construct an
indecomposable normalized bundle of invariant \(e\).
(b) Let \(d=-e>0\), \(\deg D=d\), and let
\(0\ne\xi\in H^1(\mathcal O_C(-D))\) define
\(0\to\mathcal O_C\to\mathcal E\to\mathcal O_C(D)\to0\).
Regard \(\xi\) as a functional on \(H^0(K_C+D)\)
and write \(H=\mathbb P(\ker\xi)\).
For each effective \(E\) of degree \(d-1\), put
\(L_E=|K_C+D-E|+E\).
Show that \(\mathcal E\) is normalized exactly when
no \(L_E\) is contained in \(H\).
(c) Deduce existence of ruled surfaces with
\(-g\leq e<0\).
(d) If \(g=2\), prove the necessary bound \(e\geq-2\).
\end{problem}
\end{exerciseintro}

\begin{solution}
(a) \emph{Construction.} Choose an effective divisor \(A\) of degree \(e\)
contained in a canonical divisor. Such a choice exists
when \(g\geq2\); for \(g=1\) the only case is \(e=0\)
and take \(A=0\).
Serre duality gives
\[
H^1(\mathcal O_C(A))^*
=H^0(\mathcal O_C(K_C-A))\ne0.
\]
Choose a nonzero extension class to obtain
\[
0\to\mathcal O_C\to\mathcal E
\to\mathcal O_C(-A)\to0.
\tag{1}
\]
A positive-degree line bundle can map nontrivially
to neither the first nor last term, so none can
be a subbundle of \(\mathcal E\).
Thus (1) is normalized and \(-\deg\mathcal E=e\).
If it decomposed, normalization would give
\(\mathcal E\cong\mathcal O_C\oplus\mathcal O_C(-A)\).
For \(e>0\) its only global sections come from
the first summand, so the displayed extension would
split. For \(e=0,A=0\), a nowhere-zero constant
section of \(\mathcal O_C^{\oplus2}\) also splits
off a summand. Both contradict the nonzero class.
Hence \(\mathcal E\) is indecomposable.

(b) \emph{The normalization criterion.}
The bundle already has a nonzero section.
It fails normalization precisely when some
positive-degree line bundle maps nontrivially into it.
By replacing such a bundle by its twist down by
an effective divisor, it is enough to test degree-one
line bundles \(M\).
A nonzero map \(M\to\mathcal E\) has nonzero composite
with \(\mathcal O_C(D)\), since
\(\operatorname{Hom}(M,\mathcal O_C)=0\).
Such a composite identifies
\[
M=\mathcal O_C(D-E)
\]
for an effective divisor \(E\) of degree \(d-1\).
It lifts to \(\mathcal E\) exactly when the
pulled-back extension class is zero in
\(H^1(\mathcal O_C(E-D))\).

By naturality of Serre duality, the dual of this
map on extension classes is the inclusion
\[
H^0(K_C+D-E)\longrightarrow H^0(K_C+D),
\]
multiplication by the section of \(E\).
Thus its vanishing on \(\xi\) says precisely
that \(\xi\) annihilates this subspace, or
\(L_E\subset H\). Excluding every such \(E\)
is exactly normalization.

(c) \emph{The incidence dimension bound.}
Let \(1\leq d\leq g\) and fix any degree-\(d\)
divisor \(D\).
The extension space
\(W=H^1(\mathcal O_C(-D))\) has dimension \(g+d-1\).
For any effective \(E\) of degree \(d-1\),
the divisor \(K_C+D-E\) has degree \(2g-1\),
so its section space has dimension \(g\).
Consequently its annihilator in \(W\) has
dimension \(d-1\).
The incidence of all such bad pairs \((\xi,E)\)
has dimension
\[
(d-1)+(d-1)=2d-2<g+d-1=\dim W.
\]
The section spaces have constant dimension and
vary algebraically with \(E\), so these annihilators
form a vector bundle over \(C^{(d-1)}\);
in particular the stated dimension bound applies.
Their image cannot fill \(W\).
Choose \(\xi\) outside its closure.
It is nonzero, since zero belongs to every annihilator.
By (b) its extension is normalized, giving invariant
\(e=-d\). This proves all the claimed negative values.

(d) \emph{The genus-two bound.}
Suppose \(g=2\) and a normalized bundle has
positive degree \(d=-e\).
It has exactly one independent global section.
Indeed, every nonzero section is nowhere zero.
Two independent sections cannot generically span
a line, since that saturated line would have degree
at most zero and at most one independent section.
If they generically span rank two, their determinant
is a nonzero section of a positive-degree line bundle;
at one of its zeros a nontrivial constant combination
of the two sections vanishes. This contradicts
normalization.
Riemann--Roch now gives
\(1=h^0(\mathcal E)\geq\chi(\mathcal E)=d-2\),
so \(d\leq3\).

It remains to exclude \(d=3\).
A nonzero section then gives an extension as in (b)
with \(\deg D=3\) and nonzero class \(\xi\).
The very ample system \(|K_C+D|\) embeds \(C\)
as a degree-five curve in
\(\mathbb P(H^0(K_C+D)^*)=\mathbb P^3\).
For a length-two effective divisor \(E\), the
projective annihilator of \(H^0(K_C+D-E)\)
is exactly its secant or tangent line.
Part (b) would put the point \([\xi]\) on none
of these lines. Projection from \([\xi]\)
would therefore embed \(C\) as a smooth plane
quintic, without changing its degree.
Such a quintic has genus six, contradicting \(g=2\).
Thus \(d\leq2\), or \(e\geq-2\).
\end{solution}

\begin{exerciseintro}
% record: H-V-02-006
\exerciseheading{Exercise V.2.6}
\addcontentsline{toc}{subsection}{Exercise V.2.6}

\begin{context}
Work over an algebraically closed field.
Every line bundle on \(\mathbb P^1\) is \(\mathcal O(a)\),
and \(H^1(\mathcal O(a))=0\) for \(a\geq-1\).
Vector bundles on a nonsingular curve admit line
subbundles with locally free quotient, by saturation.
For a vector bundle \(\mathcal E\), Serre generation
of \(\mathcal E^*(M)\) for large \(M\) gives an
injection \(\mathcal E\hookrightarrow\mathcal O(M)^{\oplus N}\).
Extensions of line bundles are classified by their
\(H^1\) groups.
\end{context}

\begin{problem}
Prove that every finite-rank vector bundle on
\(\mathbb P^1\) is a direct sum of line bundles.
\end{problem}
\end{exerciseintro}

\begin{solution}
Induct on the rank; ranks zero and one are immediate.
Choose a line subbundle \(\mathcal O(a)\subset\mathcal E\)
of maximal degree. Such a maximum exists: there
is a line subbundle by saturation, and the injection
into \(\mathcal O(M)^{\oplus N}\) bounds its degree
above by \(M\).
Its quotient \(\mathcal F\) is locally free.
By induction,
\(\mathcal F=\bigoplus_i\mathcal O(b_i)\).

We claim \(b_i\leq a\) for every \(i\).
Pull back the extension along the summand
\(\mathcal O(b_i)\subset\mathcal F\) and twist by
\(\mathcal O(-a-1)\), obtaining
\[
0\to\mathcal O(-1)\to\mathcal E_i(-a-1)
\to\mathcal O(b_i-a-1)\to0.
\]
Since \(H^1(\mathcal O(-1))=0\), a nonzero
section of the last term lifts to the middle.
If \(b_i\geq a+1\), this gives a nonzero map
\(\mathcal O(a+1)\to\mathcal E\).
Saturating its image yields a line subbundle
of degree at least \(a+1\), contradicting
maximality of \(a\). Hence the claim holds.

It follows that
\[
\operatorname{Ext}^1(\mathcal F,\mathcal O(a))
=\bigoplus_i H^1(\mathcal O(a-b_i))=0.
\]
The original extension therefore splits:
\(\mathcal E\cong\mathcal O(a)\oplus
\bigoplus_i\mathcal O(b_i)\).
This completes the induction.
\end{solution}

\begin{exerciseintro}
% record: H-V-02-007
\exerciseheading{Exercise V.2.7}
\addcontentsline{toc}{subsection}{Exercise V.2.7}

\begin{context}
Let \(C\) be an elliptic curve over an algebraically
closed field and \(\mathcal E\) the normalized
nonsplit extension
\(0\to\mathcal O_C\to\mathcal E\to\mathcal O_C(P_0)\to0\).
Thus \(\deg\mathcal E=1\), it has no positive-degree
line subbundle, and \(X=\mathbb P(\mathcal E)\)
is the elliptic ruled surface of invariant \(-1\).
A section with kernel \(M\subset\mathcal E\)
has divisor bundle
\(\mathcal O_X(1)\otimes\pi^*M^{-1}\) and square
\(\deg\mathcal E-2\deg M\).
We use Riemann--Roch, Serre duality with
\(\omega_C=\mathcal O_C\), and
\(\operatorname{Pic}^0(C)\cong C\).
A normalized Poincare bundle represents its
degree-zero line bundles; constant cohomology
dimensions give locally free pushforward and
base change by Grauert's theorem.
\end{context}

\begin{problem}
Show that the sections of \(X\to C\) with
self-intersection one form a one-dimensional
algebraic family parametrized by \(C\), and
that no two of them are linearly equivalent.
\end{problem}
\end{exerciseintro}

\begin{solution}
Fix \(M\in\operatorname{Pic}^0(C)\).
By Serre duality,
\[
H^1(C,\mathcal E\otimes M^{-1})
\cong\operatorname{Hom}(\mathcal E,M)^*.
\]
The Hom space is zero. A nonzero map
\(\mathcal E\to M\) would have image a line bundle
of degree at most zero; its kernel would be
a line subbundle of degree at least one,
contrary to normalization.
Riemann--Roch therefore gives
\[
h^0(C,\mathcal E\otimes M^{-1})=1.
\]
Its unique nonzero map \(M\to\mathcal E\), up to
scalar, is nowhere zero: a zero would make its
saturation a positive-degree line subbundle.
Its quotient is a line bundle, giving a unique
section \(D_M\) of \(X\).
The self-intersection formula gives \(D_M^2=1\).
Conversely, a section of square one has
\(\deg M=0\), so every such section is obtained.

These sections vary algebraically.
Let \(B=\operatorname{Pic}^0(C)\), and let
\(\mathcal P\) be the Poincare bundle on \(C\times B\).
The pushforward to \(B\) of
\(\operatorname{pr}_C^*\mathcal E\otimes\mathcal P^{-1}\)
is a line bundle, by the cohomology calculation
and Grauert's theorem.
Its evaluation, tensored with \(\mathcal P\),
gives a line subbundle of
\(\operatorname{pr}_C^*\mathcal E\).
The quotient defines an algebraic family of the
sections \(D_M\). Its parameter space \(B\cong C\)
has dimension one, and the preceding uniqueness
shows it parametrizes exactly the claimed sections.

Finally, if \(D_M\sim D_{M'}\), their divisor
bundles give \(\pi^*(M^{-1}\otimes M')\cong\mathcal O_X\).
Pullback from \(C\) is injective on Picard groups,
since restricting to any fixed section of \(\pi\)
is a left inverse. Hence \(M=M'\), and the
uniqueness already proved gives \(D_M=D_{M'}\).
Distinct sections in the family are therefore
not linearly equivalent.
\end{solution}

\begin{exerciseintro}
% record: H-V-02-008
\exerciseheading{Exercise V.2.8}
\addcontentsline{toc}{subsection}{Exercise V.2.8}

\begin{context}
Let \(C\) be a smooth projective integral curve of
genus \(g\) over an algebraically closed field.
For a vector bundle \(E\), put
\(\mu(E)=\deg E/\operatorname{rank}E\).
Stability requires every nonzero proper vector-bundle
quotient to have slope greater than \(\mu(E)\);
semistability replaces greater than by at least.
In rank two this is equivalent to every line
subbundle having degree less than, or at most,
\(\deg E/2\), respectively.
A normalized rank-two bundle has a nowhere-zero
section and maximal line-subbundle degree zero.
We use
\(\operatorname{Ext}^1(L,\mathcal O_C)=H^1(L^{-1})\)
and Serre duality.
Normalization is the convention of the section's
Notation 2.8.1; twisting by a line bundle preserves
stability, semistability, and indecomposability.
\end{context}

\begin{problem}
(a) Show that a decomposable bundle is not stable.
(b) For normalized rank-two \(E\), prove that it is
stable exactly when \(\deg E>0\), and semistable
exactly when \(\deg E\geq0\).
(c) Classify normalized indecomposable rank-two
bundles that are not semistable by
\(0<e\leq2g-2\), a line bundle \(L\) of degree
\(-e\), and a nonzero class in \(H^1(L^{-1})\)
up to scalar. Explain the additional twist needed
if normalization is not imposed.
\end{problem}
\end{exerciseintro}

\begin{solution}
(a) If \(E=E_1\oplus E_2\) with both summands nonzero,
then \(\mu(E)\) is a weighted average of
\(\mu(E_1),\mu(E_2)\).
At least one summand, viewed as a proper quotient,
has slope at most \(\mu(E)\). This violates stability.

(b) All line subbundles have degree at most zero,
and the degree zero is attained by the nowhere-zero
section \(\mathcal O_C\subset E\).
Thus all these degrees are strictly below
\(\deg E/2\) exactly when \(\deg E>0\), and are
at most \(\deg E/2\) exactly when \(\deg E\geq0\).
These are precisely the two slope criteria.

(c) Let \(E\) be normalized, indecomposable, and
not semistable. By (b), \(\deg E=-e<0\).
Its section gives
\[
0\to\mathcal O_C\to E\to L\to0,\qquad
L=\det E,\quad\deg L=-e.
\tag{1}
\]
Indecomposability makes its class
\(\xi\in H^1(L^{-1})\) nonzero.
Serre duality gives
\(H^1(L^{-1})^*=H^0(\omega_C\otimes L)\ne0\),
so \(2g-2-e\geq0\), proving \(e\leq2g-2\).

Conversely, any nonzero extension (1) with
\(\deg L=-e<0\) is normalized:
a positive-degree line bundle maps to neither
\(\mathcal O_C\) nor \(L\), and hence cannot be
a subbundle of \(E\).
Also \(h^0(E)=1\), since \(h^0(L)=0\).
If \(E\) decomposed, normalization would force
\(E=\mathcal O_C\oplus L\); its unique section
would then split (1), contrary to \(\xi\ne0\).
Thus it is indecomposable and, by (b), not semistable.

The unique section up to scalar makes the subbundle
\(\mathcal O_C\) in (1) intrinsic.
An isomorphism of middle bundles must preserve it,
and induces automorphisms of that line and of \(L\).
These are nonzero scalars. The extension classes
therefore give isomorphic bundles exactly when
they differ by a nonzero scalar, proving the
stated classification.

For an arbitrary unnormalized indecomposable bundle
that is not semistable, let \(M\) be its maximal-degree
line subbundle. Its degree is greater than half
the bundle degree, and it is unique: a different
such line could not map nontrivially into the
quotient by \(M\), whose degree is smaller.
After twisting by \(M^{-1}\) one obtains the
normalized bundle classified above.
Thus in the unnormalized formulation one must
also specify \(M\), and the bundle is \(M\otimes E\)
with \(E\) as in (1).
\end{solution}

\begin{exerciseintro}
% record: H-V-02-009
\exerciseheading{Exercise V.2.9}
\addcontentsline{toc}{subsection}{Exercise V.2.9}

\begin{context}
Work over an algebraically closed field.
The resolution of a quadric cone \(Q\subset\mathbb P^3\)
is \(\pi:\mathbb F_2\to Q\), contracting a section \(E\).
With \(F\) a ruling fiber,
\(E^2=-2,\ EF=1,\ F^2=0\), the hyperplane pullback
is \(H=E+2F\), and \(K_{\mathbb F_2}=-2E-4F\).
Also \(\pi^{-1}\mathcal I_{\mathrm{vertex}}\cdot
\mathcal O_{\mathbb F_2}=\mathcal O(-E)\) and
\(\pi_*\mathcal O_{\mathbb F_2}=\mathcal O_Q\).
We use adjunction, the universal property of blowing up,
and the projection formula. Restriction of degree-\(a\)
forms from \(\mathbb P^3\) to \(Q\) is surjective:
the quadric ideal is \(\mathcal O(-2)\), and
\(H^1(\mathbb P^3,\mathcal O(t))=0\).
\end{context}

\begin{problem}
Let \(Y\) be a nonsingular integral curve on \(Q\).
Show that either it is a complete intersection of
\(Q\) with a degree-\(a\) surface, with
\(\deg Y=2a\) and \(g(Y)=(a-1)^2\), or its
degree is odd, \(2a+1\), and \(g(Y)=a(a-1)\).
\end{problem}
\end{exerciseintro}

\begin{solution}
Let \(\widetilde Y\) be the strict transform.
It is isomorphic to \(Y\).
This is immediate if \(Y\) avoids the vertex.
If it passes through the vertex \(P\), the vertex
ideal restricts to the invertible ideal
\(\mathfrak m_{Y,P}\) on the smooth curve \(Y\).
The universal property of blowing up therefore
lifts \(Y\) to the resolution, giving its strict
transform isomorphically.
Since the exceptional ideal restricts to
\(\mathfrak m_{Y,P}\), the intersection with \(E\)
has length exactly one. Consequently
\[
\widetilde Y E=\delta,\qquad
\delta=0\text{ or }1.
\]

Write \(\widetilde Y\sim aE+bF\).
Then
\[
\deg Y=H\widetilde Y=b,\qquad b-2a=\delta.
\tag{1}
\]
Here \(a=\widetilde YF\geq0\); when \(\delta=0\),
positivity of the degree forces \(a\geq1\).
Adjunction gives
\[
g(Y)=1+\frac{\widetilde Y^2+
K_{\mathbb F_2}\widetilde Y}{2}
=1-a^2+ab-b.
\]
If \(\delta=0\), then \(b=2a\) and this becomes
\((a-1)^2\).
If \(\delta=1\), then \(b=2a+1\) and it becomes
\(a(a-1)\). The latter includes a ruling line,
where \(a=0\).

In the even case, \(\widetilde Y\sim aH\) and
it avoids \(E\).
The section cutting it out comes from a section
of \(\mathcal O_Q(a)\), by
\(\pi_*\mathcal O_{\mathbb F_2}=\mathcal O_Q\)
and the projection formula.
Its value at the vertex is nonzero, since its
pullback does not vanish on \(E\).
Its divisor on \(Q\) is therefore exactly \(Y\).
Lift this section to a degree-\(a\) form on
\(\mathbb P^3\). The corresponding surface cuts
out \(Y\) scheme-theoretically on \(Q\), proving
the complete-intersection assertion.
\end{solution}

\begin{exerciseintro}
% record: H-V-02-010
\exerciseheading{Exercise V.2.10}
\addcontentsline{toc}{subsection}{Exercise V.2.10}

\begin{context}
Work over an algebraically closed field.
On \(\mathbb F_e\), \(e\geq0\), use
\(C_0^2=-e,\ C_0f=1,\ f^2=0\) and
\(K=-2C_0-(e+2)f\).
For \(n>e\), \(H=C_0+nf\) embeds the surface
as a degree-\(d=2n-e\) scroll in \(\mathbb P^{d+1}\),
with \(h^0(H)=d+2\).
A divisor \(aC_0+bf\), \(a>0\), is very ample
if \(b>ae\), and is generated by sections if
\(b\geq ae\).
For \(e>0\), \(|C_0+ef|\) contracts only \(C_0\)
and is an isomorphism off it onto the complement
of the vertex of the rational cone.
We use Bertini and adjunction for smooth curves.
\end{context}

\begin{problem}
If \(n>e\geq0\) and \(n\geq2e-2\), prove that
the scroll contains a smooth curve \(Y\) of
genus \(d+2\) canonically embedded by its ambient
embedding. Deduce that every \(g\geq4\) occurs
for a nonhyperelliptic curve with a \(g^1_3\).
\end{problem}
\end{exerciseintro}

\begin{solution}
Take a general member of
\[
|H-K|=|3C_0+(n+e+2)f|.
\tag{1}
\]
The assumption \(n\geq2e-2\) says
\(n+e+2\geq3e\).
If this is strict, the system is very ample,
so Bertini gives a smooth integral member.
In the equality case, (1) is
\(|3(C_0+ef)|\). Here \(n>e\) forces \(e>2\).
It is the pullback of the cubic hyperplane
system on the rational cone. A general member
avoids the vertex and is smooth and integral
by Bertini on that cone away from its vertex.
Its pullback is likewise smooth and integral,
and avoids \(C_0\).
Thus in all cases we obtain the required smooth \(Y\).

Adjunction gives
\[
K_Y=(K+Y)|_Y=H|_Y,
\]
and its degree is
\[
HY=(C_0+nf)(3C_0+(n+e+2)f)
=4n-2e+2=2d+2.
\]
Hence \(2g(Y)-2=2d+2\), or \(g(Y)=d+2\).
The restriction of sections of \(H\) to \(Y\)
is injective, because \(H-Y=K\) has degree
minus two on a ruling fiber and has no global
section. Both its source and
\(H^0(Y,K_Y)\) have dimension \(d+2\), so
restriction is an isomorphism.
Thus the scroll embedding restricts to the
complete canonical embedding of \(Y\).
Also \(Yf=3\), so the ruling restricts to a
degree-three pencil. In particular \(Y\) is
nonhyperelliptic.

Finally, for any \(g\geq4\), put \(d=g-2\).
If \(d\) is even, choose \(e=0,\ n=d/2\).
If \(d\) is odd, choose \(e=1,\ n=(d+1)/2\).
In both cases \(n>e\) and \(n\geq2e-2\).
The preceding construction produces the desired
canonical trigonal genus-\(g\) curve.
\end{solution}

\begin{exerciseintro}
% record: H-V-02-011
\exerciseheading{Exercise V.2.11}
\addcontentsline{toc}{subsection}{Exercise V.2.11}

\begin{context}
Let \(\pi:X=\mathbb P(\mathcal E)\to C\) be a
ruled surface over a smooth projective curve.
Use the normalized extension
\(0\to\mathcal O_C\to\mathcal E\to\mathcal O_C(\epsilon)\to0\)
and \(\mathcal O_X(C_0)=\mathcal O_X(1)\).
For a divisor \(b\) on \(C\), write \(bf=\pi^*b\).
The sections of \(\mathcal O_X(C_0+bf)\) are
\(H^0(C,\mathcal E(b))\).
We use curve Riemann--Roch and the criterion that
a linear system is very ample if it separates
all length-two subschemes.
A general section of a globally generated
rank-two bundle on a curve is nowhere zero:
vanishing at a point imposes two linear conditions
over a one-dimensional parameter space.
\end{context}

\begin{problem}
(a) If \(|b|\) and \(|b+\epsilon|\) have no base
points and \(b\) is nonspecial, show that
\(|C_0+bf|\) has no base points and contains a section.
(b) If \(b,b+\epsilon\) are very ample and
\(b-P,b+\epsilon-P\) are nonspecial for every
\(P\in C\), prove \(C_0+bf\) very ample.
\end{problem}
\end{exerciseintro}

\begin{solution}
(a) Tensoring the normalized extension by
\(\mathcal O_C(b)\) gives
\[
0\to\mathcal O_C(b)\to\mathcal E(b)
\to\mathcal O_C(b+\epsilon)\to0.
\tag{1}
\]
Since \(H^1(\mathcal O_C(b))=0\), global sections
surject onto the quotient's sections.
At any point, sections of \(\mathcal O_C(b)\)
span the subline and lifts of sections of
\(\mathcal O_C(b+\epsilon)\) span the quotient.
Thus \(\mathcal E(b)\) is globally generated.
Its tautological quotient then shows that
\(\mathcal O_X(C_0+bf)\) is globally generated.

Choose a nowhere-zero section of \(\mathcal E(b)\).
Composing its pullback with the tautological quotient
defines a divisor in \(|C_0+bf|\).
In each fiber \(\mathbb P^1\) its equation is a
nonzero linear form, so its zero scheme has one
reduced point. Locally trivializing the nowhere-zero
subline shows that this divisor is a section of \(\pi\).

(b) Let \(Z=P+Q\) be any length-two divisor on \(C\),
allowing \(P=Q\).
For \(L=\mathcal O_C(b)\) and
\(L=\mathcal O_C(b+\epsilon)\), very ampleness
gives
\(h^0(L(-P))=h^0(L)-1\) and
\(h^0(L(-Z))=h^0(L)-2\).
Together with \(H^1(L(-P))=0\), Riemann--Roch
then gives \(H^1(L(-Z))=0\).
Twisting (1) by \(-Z\) consequently gives
\(H^1(\mathcal E(b)(-Z))=0\).
Hence
\[
H^0(C,\mathcal E(b))\longrightarrow
H^0(Z,\mathcal E(b)|_Z)
\]
is surjective for every such \(Z\).

This separates every length-two subscheme of \(X\).
For two points over distinct base points, it
prescribes the required quotient values independently.
For a vertical length-two subscheme in one fiber,
global generation supplies every linear form on
that fiber, and \(\mathcal O_{\mathbb P^1}(1)\)
separates it. For a nonvertical tangent vector,
use \(Z=2P\): local trivialization of \(\mathcal E(b)\)
identifies the two freely prescribed base jets with
all first-order linear forms on the local projective
bundle. They separate that tangent vector.
These cases exhaust length-two subschemes on
the smooth surface. Thus \(C_0+bf\) is very ample.
\end{solution}

\begin{exerciseintro}
% record: H-V-02-012
\exerciseheading{Exercise V.2.12}
\addcontentsline{toc}{subsection}{Exercise V.2.12}

\begin{context}
Let \(X=\mathbb P(\mathcal E)\to C\) be a ruled
surface over an elliptic curve, with \(\mathcal E\)
normalized and \(e=-\deg\mathcal E\geq-1\).
Write \(\det\mathcal E=\mathcal O_C(\epsilon)\),
\(\mathcal O_X(C_0)=\mathcal O_X(1)\), and
\(n=\deg b\), with \(C_0^2=-e\).
Line bundles on an elliptic curve are nonspecial
in positive degree, generated in degree at least
two, and very ample exactly in degree at least three.
We use the two generation and very-ampleness
criteria proved in V.2.11.
For \(e=-1\), \(\mathcal E\) is stable of degree one;
twists remain stable. On an elliptic curve a stable
positive-degree bundle has \(H^1=0\), by Serre
duality and the quotient slope criterion.
A stable negative-degree bundle has no section.
\end{context}

\begin{problem}
(a) If \(n\geq e+2\), prove that
\(|C_0+bf|\) is base-point-free and contains a section.
(b) Prove that \(C_0+bf\) is very ample exactly when
\(n\geq e+3\), including \(e=-1\).
\end{problem}
\end{exerciseintro}

\begin{solution}
(a) First suppose \(e\geq0\).
If \(n\geq e+2\), the degrees of \(b,b+\epsilon\)
are \(n\geq2\) and \(n-e\geq2\), so both are
generated and \(b\) is nonspecial.
The first criterion from V.2.11 proves (a).
Now suppose \(e=-1\).
The bundle \(V=\mathcal E(b)\) is stable of
degree \(2n+1\).
If \(n\geq1\), then for every point \(P\),
\(V(-P)\) is stable of positive degree \(2n-1\),
and has \(H^1=0\).
The evaluation \(H^0(V)\to V|_P\) is therefore
surjective. Thus \(V\) is globally generated.
As in V.2.11(a), its tautological system is generated
and a nowhere-zero general section gives a
section divisor. This proves (a) at \(n\geq1=e+2\).

(b) For \(e\geq0\) and \(n\geq e+3\), the degrees
of \(b,b+\epsilon\) are at least three,
and subtracting any point leaves positive degree.
Thus both are very ample and the required
point-subtractions are nonspecial.
The second criterion proves sufficiency.
Conversely, restriction of \(C_0+bf\) to \(C_0\)
has degree \(n-e\). Very ampleness on the surface
forces very ampleness on this elliptic section,
hence \(n-e\geq3\). This proves necessity for \(e\geq0\).

For \(e=-1\), use \(V=\mathcal E(b)\) as above.
If \(n\geq2\), then for every length-two divisor \(Z\),
\(V(-Z)\) is stable of positive degree \(2n-3\);
again \(H^1=0\).
Evaluation to \(V|_Z\) is surjective, so the
length-two argument of V.2.11(b) makes
\(C_0+bf\) very ample.

Finally, if \(n<0\), stability of the negative-degree
\(V\) gives no sections.
If \(n=0\) or \(1\), positive-degree stability
and Riemann--Roch give \(h^0(V)=2n+1\), respectively
one or three.
Such a complete system cannot embed this surface:
for \(n=0\) it maps to a point, and for \(n=1\)
a surface embedding into \(\mathbb P^2\) would
identify \(X\) with \(\mathbb P^2\).
But \(X\) contains an elliptic section of square
one; its image in \(\mathbb P^2\) would have
degree one, since a degree-\(d\) plane curve has
square \(d^2\). A line is not elliptic.
Thus \(n\geq2=e+3\) is necessary, completing (b).
\end{solution}

\begin{exerciseintro}
% record: H-V-02-013
\exerciseheading{Exercise V.2.13}
\addcontentsline{toc}{subsection}{Exercise V.2.13}

\begin{context}
Over an algebraically closed field, elliptic ruled
surfaces exist with every invariant \(e\geq-1\):
use \(\mathcal O_C\oplus L\) with \(\deg L=-e\)
for \(e\geq0\), and the normalized nonsplit
degree-one bundle for \(e=-1\).
For a normalized bundle \(\mathcal E\),
\(C_0^2=-e,\ C_0f=1,\ f^2=0\), and
\(H^0(X,\mathcal O_X(C_0+bf))=H^0(C,\mathcal E(b))\).
A divisor \(C_0+bf\) is very ample if
\(\deg b\geq e+3\), by V.2.12.
We use elliptic Riemann--Roch and the fact that
a stable positive-degree elliptic bundle has \(H^1=0\).
\end{context}

\begin{problem}
For \(e\geq-1\) and \(n\geq e+3\), construct an
elliptic scroll of degree \(d=2n-e\) in
\(\mathbb P^{d-1}\). In particular, construct
one of degree five in \(\mathbb P^4\).
\end{problem}
\end{exerciseintro}

\begin{solution}
Choose the indicated elliptic ruled surface,
and a divisor \(b\) of degree \(n\) on its base.
The divisor \(H=C_0+bf\) is very ample, and
\[
H^2=-e+2n=d.
\]
It has degree one on each ruling fiber, so the
embedding maps those fibers to lines.

We compute the dimension of its complete system.
For \(e\geq0\), tensoring
\(0\to\mathcal O_C\to\mathcal E\to
\mathcal O_C(\epsilon)\to0\)
by \(\mathcal O_C(b)\) gives line terms of
degrees \(n\) and \(n-e\), both positive.
Their \(H^1\)'s vanish, so \(H^1(\mathcal E(b))=0\).
For \(e=-1\), the bundle \(\mathcal E(b)\) is
stable of positive degree \(2n+1\), and has
the same vanishing.
In either case elliptic Riemann--Roch gives
\[
h^0(\mathcal E(b))=\deg\mathcal E(b)=2n-e=d.
\]
Thus the complete embedding is into
\(\mathbb P^{d-1}\), with surface degree \(H^2=d\).
Taking \(e=-1,\ n=2\) gives \(d=5\) and an
elliptic scroll in \(\mathbb P^4\).
\end{solution}

\begin{exerciseintro}
% record: H-V-02-014
\exerciseheading{Exercise V.2.14}
\addcontentsline{toc}{subsection}{Exercise V.2.14}

\begin{context}
Let \(X=\mathbb P(\mathcal E)\to C\) be a ruled
surface over an algebraically closed field of characteristic
\(p>0\). The base has genus \(g\geq2\), the bundle is
normalized, and \(e=-\deg\mathcal E<0\).
Numerically, \(C_0^2=-e,\ C_0f=1,\ f^2=0\) and
\(K_X\equiv-2C_0+(2g-2-e)f\).
Normalization means that \(\mathcal E\) has a section
and no line subbundle of positive degree.
We use adjunction, the Hurwitz inequality for a separable
map of smooth curves, and the Nakai--Moishezon criterion.
A purely inseparable map of smooth curves over a perfect
field preserves genus: it is an iterate of relative
Frobenius followed by an isomorphism.
\end{context}

\begin{problem}
(a) For an integral curve \(Y\equiv aC_0+bf\), other than
\(C_0\) or a fiber, prove the following alternatives:
\[
\begin{array}{ll}
a=1:& b\geq0,\\
2\leq a<p:& b\geq ae/2,\\
a\geq p:& b\geq ae/2+1-g.
\end{array}
\]
(b) Show that \(a>0\) and
\(b>a(e/2+(g-1)/p)\) imply ampleness of any divisor
\(D\equiv aC_0+bf\). Show conversely that ampleness
implies \(a>0,\ b>ae/2\).
\end{problem}
\end{exerciseintro}

\begin{solution}
(a) The projection of \(Y\) is finite and surjective;
its degree is \(a=Yf>0\).
Let \(\widetilde Y\) be its normalization, of genus \(h\).
Adjunction and normalization give
\[
2h-2\leq 2p_a(Y)-2
=(a-1)(2b-ae)+a(2g-2).
\tag{1}
\]
If \(a=1\), the finite degree-one map \(Y\to C\)
is an isomorphism: on affine opens its integral
coordinate ring lies in \(K(C)\), and the base ring
is integrally closed. Thus \(Y\) is a section.
Writing its line bundle as
\(\mathcal O_X(C_0)\otimes\pi^*B\), with \(\deg B=b\),
its defining section gives a nonzero map
\(B^{-1}\to\mathcal E\). Saturation and normalization
imply \(\deg B^{-1}\leq0\), hence \(b\geq0\).

If \(2\leq a<p\), the map
\(\widetilde Y\to C\) is separable, since a nontrivial
inseparable degree is divisible by \(p\).
Hurwitz gives \(2h-2\geq a(2g-2)\).
Together with (1), this yields \(2b-ae\geq0\).

For arbitrary \(a\), factor the function-field extension
into a purely inseparable part followed by a separable
part. The former preserves genus, and Hurwitz for the
latter gives \(h\geq g\), since \(g\geq2\).
For \(a>1\), (1) therefore implies
\[
0\leq(a-1)(2b-ae+2g-2),
\]
which is the last bound.

(b) Assume the sufficient inequalities.
Then \(Df=a>0\) and
\(D^2=a(2b-ae)>0\).
Also \(DC_0=b-ae>0\), because \(e<0\).
Consider any other integral curve
\(Y\equiv rC_0+sf\).
If \(r=1\), part (a) gives
\[
DY=as+b-ae\geq b-ae>0.
\]
If \(2\leq r<p\), it gives
\[
DY=as+rb-are\geq r(b-ae/2)>0.
\]
If \(r\geq p\), it gives
\[
DY\geq r(b-ae/2)-a(g-1)
>a(g-1)(r/p-1)\geq0.
\]
Thus \(D\) has positive square and positive intersection
with every integral curve, so Nakai--Moishezon makes
it ample.

Conversely, an ample divisor has \(Df=a>0\) and
\(D^2=a(2b-ae)>0\), whence \(b>ae/2\).
\end{solution}

\begin{exerciseintro}
% record: H-V-02-015
\exerciseheading{Exercise V.2.15}
\addcontentsline{toc}{subsection}{Exercise V.2.15}

\begin{context}
Work over an algebraically closed field \(k\) of
characteristic three, and put
\[
C=V(F)\subset\mathbb P^2,\qquad
F=x^3y+y^3z+z^3x.
\]
Its partial derivatives are \(z^3,x^3,y^3\), so
\(C\) is smooth and has genus three.
We use the projective-space Cech calculation
\[
H^1(C,\mathcal O_C)\simeq H^2(\mathbb P^2,\mathcal O(-4)),
\]
where a Laurent monomial represents zero if any
of its three exponents is nonnegative.
Extensions of a line bundle \(L\) by \(\mathcal O_C\)
are classified by \(H^1(C,L^{-1})\).
For a normalized ruled surface,
\(C_0^2=-e,\ C_0f=1,\ f^2=0\), and
\(K_X\equiv-2C_0+(2g-2-e)f\).
We use adjunction and the fact that a degree-\(p\)
purely inseparable map of smooth projective curves
over a perfect field is relative Frobenius, up to
isomorphism, and preserves genus.

Write \(F_C\) for absolute Frobenius.
It fixes points and satisfies \(F_C^*L=L^{\otimes3}\);
its cohomology action is semilinear.
The coordinate-cubing endomorphism of this curve,
which is defined over \(\mathbb F_3\), is \(k\)-linear.
Its cohomology matrix is the same, but it need not
fix an arbitrary \(k\)-point. We use absolute Frobenius
for the expressions involving \(P\) and \(3P\).
\end{context}

\begin{problem}
(a) Prove that Frobenius acts as zero on
\(H^1(C,\mathcal O_C)\).
(b) For any \(P\in C\), find a nonzero
\(\xi\in H^1(C,\mathcal O_C(-P))\) killed by
pullback to \(H^1(C,\mathcal O_C(-3P))\).
(c) Let
\[
0\longrightarrow\mathcal O_C\longrightarrow\mathcal E
\longrightarrow\mathcal O_C(P)\longrightarrow0
\]
be its extension, and put \(X=\mathbb P(\mathcal E)\).
Construct an integral curve \(Y\equiv3C_0-3f\)
whose projection to \(C\) is purely inseparable,
and use it to show that \(2C_0\) is not ample despite
satisfying the numerical inequalities in Hartshorne's
Proposition V.2.21(b), a textbook result rather than an exercise
in this collection.
Check the textbook's additional assertion that \(Y\)
is nonsingular.
\end{problem}
\end{exerciseintro}

\begin{solution}
(a) A basis for the indicated \(H^2\) consists of
\[
\frac1{x^2yz},\qquad \frac1{xy^2z},\qquad
\frac1{xyz^2}.
\]
Under this identification Frobenius sends a class
\(u\) to \(F^2u^3\). Indeed, cubing the exact sequence
defined by multiplication by \(F\) first gives the
sequence defined by \(F^3\); multiplication by \(F^2\)
then compares its degree-\(-12\) term to the original
degree-\(-4\) term.
In characteristic three,
\[
F^2=x^6y^2+y^6z^2+z^6x^2
+2x^3y^4z+2xy^3z^4+2x^4yz^3.
\]
After division by \(x^6y^3z^3\), every summand has
at least one nonnegative exponent, and hence vanishes
in Cech cohomology. Cyclic permutation proves the
same for the other two basis vectors. Thus the
semilinear action, and also the \(k\)-linear
coordinate-Frobenius action, is zero.

(b) From
\(0\to\mathcal O_C(-P)\to\mathcal O_C\to k(P)\to0\)
and the isomorphism on constants we get
\[
H^1(\mathcal O_C(-P))\simeq H^1(\mathcal O_C),
\qquad h^1(\mathcal O_C(-P))=3.
\]
The composite of Frobenius pullback with
\[
H^1(\mathcal O_C(-3P))\longrightarrow H^1(\mathcal O_C)
\]
is zero by (a).
The kernel of this last arrow has dimension two:
it is the cokernel of the constants in
\(H^0(\mathcal O_{3P})\), a three-dimensional space.
Thus a semilinear map from a three-dimensional space
has image in a two-dimensional space. Since \(k\)
is perfect, semilinear rank-nullity applies, and
there is a nonzero \(\xi\) in its kernel.

(c) Set \(L=\mathcal O_C(P)\).
The nonzero extension class makes \(\mathcal E\)
nonsplit. It is normalized. To see this, a
positive-degree line subbundle \(M\) cannot map
nontrivially to \(\mathcal O_C\), so it maps
nontrivially to \(L\).
Its degree is at most one. If it is positive, that
map has degree-zero divisor and is an isomorphism;
it would split the extension. This is impossible.
The given inclusion of \(\mathcal O_C\) supplies
the required section. Consequently \(e=-1\).

Since \(F_C^*\xi=0\),
\[
F_C^*\mathcal E\simeq\mathcal O_C\oplus L^3.
\]
The local cube map is an \(\mathcal O_C\)-linear
inclusion \(F_C^*\mathcal E\to\operatorname{Sym}^3\mathcal E\),
sending \(a\otimes v\) to \(av^3\).
Compose it with the splitting inclusion of \(L^3\).
This gives a section of
\(\operatorname{Sym}^3\mathcal E\otimes L^{-3}\),
or equivalently a divisor
\[
Y\in|\mathcal O_X(3)\otimes\pi^*L^{-3}|,
\qquad Y\equiv3C_0-3f.
\]
With a local frame \(e_0,e_1\) adapted to the original
extension, its equation is \(e_1^3+c e_0^3=0\).
The coefficient of \(e_1^3\) is a unit.
Therefore \(Y\) is finite flat of degree three over
\(C\), has no vertical component, and is disjoint
from the section \(e_0=0\).

Over \(K(C)\), this cubic is either irreducible or
the cube of a linear polynomial.
In the latter case its reduced support extends to a
section \(S\) of the proper ruled surface, and
\(Y=3S\), since there are no vertical components.
Then \(S\equiv C_0-f\).
But a section of this numerical class gives a
nonzero map from a line bundle of degree one into
\(\mathcal E\), contradicting normalization as above.
Thus the generic polynomial is irreducible.
The finite flat algebra of \(Y\) embeds into its
generic fiber, a field, so \(Y\) is integral.
Its function-field extension over \(C\) is purely
inseparable of degree three.

Now \(D=2C_0\) has \(a=2>0\) and \(b=0>ae/2=-1\),
the numerical inequalities in question. However,
\[
DY=2C_0(3C_0-3f)=0,
\]
so \(D\) is not ample. The characteristic-zero or
small-genus hypothesis of that textbook proposition is absent here.

Finally, the smoothness assertion is false.
Here \(K_X\equiv-2C_0+5f\), whence
\[
Y^2=-9,\qquad K_XY=15,\qquad
p_a(Y)=1+\frac{-9+15}{2}=4.
\]
The normalization of \(Y\) maps purely inseparably
of degree three to the smooth genus-three curve \(C\);
it therefore has genus three.
The difference between arithmetic and geometric
genus is one, so \(Y\) is singular, with total
\(\delta\)-invariant one. This correction leaves
the non-ampleness example intact.
\end{solution}

\begin{exerciseintro}
% record: H-V-02-016
\exerciseheading{Exercise V.2.16}
\addcontentsline{toc}{subsection}{Exercise V.2.16}

\begin{context}
Let \(C\) be a nonsingular affine curve over an
algebraically closed field.
Coherent sheaves on an affine scheme are generated by
finitely many global sections and have no higher cohomology.
On a nonsingular curve, torsion-free coherent sheaves
are locally free.
The Grothendieck group \(K(C)\) has the relations
\([\mathcal F]=[\mathcal F']+[\mathcal F'']\) for
short exact sequences of coherent sheaves.
Rank and determinant are additive and multiplicative,
respectively, in such sequences, so equality of classes
of vector bundles implies equality of their ranks and
determinants. One can equivalently compute this group
using vector bundles, since a coherent sheaf on a
regular affine curve has a finite locally free resolution.
\end{context}

\begin{problem}
Prove that vector bundles of the same rank on \(C\)
are isomorphic if and only if their classes in \(K(C)\)
coincide. Give a counterexample for a projective curve.
\end{problem}
\end{exerciseintro}

\begin{solution}
We prove that a rank-\(r\) bundle \(E\), for \(r\geq1\),
satisfies
\[
E\simeq\mathcal O_C^{\,r-1}\oplus\det E.
\tag{1}
\]
Suppose \(r\geq2\), and choose a finite-dimensional
space \(W\) of global sections generating \(E\).
For each point \(P\), the subspace of sections vanishing
at \(P\) has codimension \(r\) in \(W\).
Consequently the incidence variety of pairs
\((P,s)\) with \(s(P)=0\) has dimension
\(1+\dim W-r<\dim W\).
Its image has closure of dimension smaller than
\(\dim W\). A section outside that closure vanishes
nowhere, and gives an exact sequence
\[
0\longrightarrow\mathcal O_C
\longrightarrow E\longrightarrow Q\longrightarrow0
\]
with \(Q\) locally free of rank \(r-1\).
The sequence splits, because its extension group
is \(H^1(C,Q^\vee)=0\).
Induction, starting with rank one, proves (1);
the determinant of the remaining line summand must
be \(\det E\).

If \([E]=[E']\), their determinants and ranks agree,
and (1) gives an isomorphism. The converse follows
immediately from the definition of \(K(C)\).
The zero-rank case is immediate as well.

On \(\mathbb P^1\), the Euler sequence
\[
0\longrightarrow\mathcal O(-1)
\longrightarrow\mathcal O^{\,2}
\longrightarrow\mathcal O(1)\longrightarrow0
\]
shows that \(\mathcal O^{\,2}\) and
\(\mathcal O(-1)\oplus\mathcal O(1)\) have equal
classes in \(K(\mathbb P^1)\).
They are not isomorphic: after twisting by
\(\mathcal O(-1)\), their spaces of sections have
dimensions zero and one, respectively.
\end{solution}

\begin{exerciseintro}
% record: H-V-02-017
\exerciseheading{Exercise V.2.17}
\addcontentsline{toc}{subsection}{Exercise V.2.17}

\begin{context}
Work over an algebraically closed field.
For a smooth curve in a smooth threefold, the conormal
sheaf \(I/I^2\) is locally free of rank two.
We use the splitting theorem for bundles on
\(\mathbb P^1\), and the calculations
\[
h^0(\mathcal O(d))=\max(d+1,0),\qquad
H^1(\mathcal O(-2))\simeq k.
\]
On the two standard affine charts with coordinate
\(\zeta\) and \(\zeta^{-1}\), the cocycle \(\zeta^{-1}\)
generates the last group.
\end{context}

\begin{problem}
(a) Determine the splitting of the conormal bundle
of the twisted cubic, pulled back to \(\mathbb P^1\).
(b) Do the same for the rational quartic parametrized
by \([s:t]\mapsto[s^4:s^3t:st^3:t^4]\), keeping
track of the characteristic.
\end{problem}
\end{exerciseintro}

\begin{solution}
Treat both curves with
\[
\phi_n([s:t])=[s^{n+1}:s^nt:st^n:t^{n+1}],
\qquad n=2,3.
\]
The two charts \(X_0\ne0\) and \(X_3\ne0\)
cover the curve.
On the first, write
\(x=X_1/X_0,\ y=X_2/X_0,\ z=X_3/X_0\).
The parameter is \(x=t/s\), and the ideal is
generated by
\[
u=y-x^n,\qquad v=z-xy.
\]
On the second put
\(A=X_0/X_3,\ B=X_1/X_3,\ C=X_2/X_3\).
The ideal is generated by
\[
u'=B-C^n,\qquad v'=A-BC.
\]
These equations also show directly that the
parametrizations are closed immersions.

On the overlap let \(\zeta=x\).
Substitute \(y=\zeta^n+u\) and
\(z=\zeta^{n+1}+\zeta u+v\), and discard terms in
\((u,v)^2\). The transition is
\[
\begin{split}
v'&=\zeta^{-2n-2}v,\\
u'&=-\zeta^{-2n}u+(n-1)\zeta^{-2n-1}v.
\end{split}
\tag{1}
\]
For the first identity use
\(A-BC=(z-xy)/z^2\).
For the second expand
\(x/z-y^n/z^n\) to first order; the coefficient
of \(u\) is \(-\zeta^{-2n}\), and that of \(v\)
is \((n-1)\zeta^{-2n-1}\).
This calculation is valid in every characteristic.

If \(N^*=\phi_n^*(I/I^2)\), equation (1) gives
\[
0\longrightarrow\mathcal O(-2n-2)
\longrightarrow N^*
\longrightarrow\mathcal O(-2n)\longrightarrow0.
\tag{2}
\]
After twisting by \(\mathcal O(2n)\), the difference
between the two local lifts of the quotient's
constant section is, up to sign,
\((n-1)\zeta^{-1}\).
Thus (2) has extension class \(n-1\), up to a
nonzero scalar, in \(H^1(\mathcal O(-2))\).

If \(n-1=0\) in \(k\), (2) splits.
Otherwise put \(G=N^*(2n)\).
The boundary map
\(H^0(\mathcal O)\to H^1(\mathcal O(-2))\)
is an isomorphism, so \(H^0(G)=0\).
Its degree is \(-2\). Writing
\(G=\mathcal O(a)\oplus\mathcal O(b)\),
the vanishing gives \(a,b\leq-1\);
their sum is \(-2\), so both equal \(-1\).

(a) For \(n=2\), the extension class is nonzero
in every characteristic. Therefore
\[
\phi_2^*(I/I^2)=\mathcal O(-5)\oplus\mathcal O(-5).
\]
(b) For \(n=3\), the class is nonzero exactly when
the characteristic is different from two.
The answer is consequently
\[
\phi_3^*(I/I^2)\simeq
\begin{cases}
\mathcal O(-7)\oplus\mathcal O(-7),&\operatorname{char}k\ne2,\\
\mathcal O(-6)\oplus\mathcal O(-8),&\operatorname{char}k=2.
\end{cases}
\]
\end{solution}

\begin{exerciseintro}
\corpussection{V.3 Monoidal Transformations}
\addcontentsline{toc}{section}{V.3 Monoidal Transformations}

% record: H-V-03-001
\exerciseheading{Exercise V.3.1}
\addcontentsline{toc}{subsection}{Exercise V.3.1}

\begin{context}
Work over an algebraically closed field.
If \(X\) is smooth and \(Y\subset X\) is a smooth
closed subvariety of positive codimension \(r\),
then locally in the etale topology the pair is
\((\mathbb A^n,\mathbb A^{n-r})\), with the second
space a coordinate subspace.
Blowing up commutes with flat base change, as do
higher direct images of quasi-coherent sheaves
under quasi-compact separated morphisms.
The blow-up of \(\mathbb A^r\) at the origin is
the total space of \(\mathcal O_{\mathbb P^{r-1}}(-1)\).
We use the cohomology of \(\mathcal O(m)\) on projective
space, the Leray spectral sequence, and
\(p_a(X)=(-1)^{\dim X}(\chi(\mathcal O_X)-1)\).
\end{context}

\begin{problem}
For a smooth projective variety \(X\) and a smooth
proper closed center \(Y\), prove that its blow-up
\(\pi:\widetilde X\to X\) has the same arithmetic genus.
\end{problem}
\end{exerciseintro}

\begin{solution}
We show
\[
\pi_*\mathcal O_{\widetilde X}=\mathcal O_X,
\qquad R^i\pi_*\mathcal O_{\widetilde X}=0
\quad(i>0).
\tag{1}
\]
This can be checked on the stated etale local models.
For the model at the origin of \(\mathbb A^r\),
write \(Z=\operatorname{Tot}(\mathcal O_{\mathbb P^{r-1}}(-1))\).
Its projection \(q\) to \(\mathbb P^{r-1}\) is affine,
with
\[
q_*\mathcal O_Z=\bigoplus_{m\geq0}\mathcal O(m).
\]
A finite standard affine cover computes cohomology
and commutes with this direct sum.
For \(i>0\) all \(H^i(\mathbb P^{r-1},\mathcal O(m))\)
with \(m\geq0\) vanish. Hence \(H^i(Z,\mathcal O_Z)=0\).
Also
\[
H^0(Z,\mathcal O_Z)
=\bigoplus_{m\geq0}H^0(\mathbb P^{r-1},\mathcal O(m))
=k[x_1,\ldots,x_r]
\]
as a graded algebra, compatibly with the blow-up map.
Since its target is affine, these calculations give
(1) for that model.
The codimension-one case is simply the identity
blow-up and satisfies (1) as well.

Taking a product with \(\mathbb A^{n-r}\) and
then the etale base changes defining the pair
\((X,Y)\) proves (1) in general.
The assertion is already the identity away from
the center, so these local checks cover \(X\).
Leray now yields
\[
H^i(\widetilde X,\mathcal O_{\widetilde X})
\simeq H^i(X,\mathcal O_X)
\quad\text{for every }i.
\]
Thus the Euler characteristics agree.
The dimensions agree because the blow-up is
birational, and the formula for \(p_a\) proves
the claim. An empty center also gives the identity.
\end{solution}

\begin{exerciseintro}
% record: H-V-03-002
\exerciseheading{Exercise V.3.2}
\addcontentsline{toc}{subsection}{Exercise V.3.2}

\begin{context}
Let \(X\) be a smooth projective surface over an
algebraically closed field, and let
\(\pi:\widetilde X\to X\) blow up a point \(P\),
with exceptional curve \(E\).
For a curve \(C\) of multiplicity \(m_P(C)\),
\[
\widetilde C=\pi^*C-m_P(C)E,\qquad
E^2=-1,\qquad \pi^*C\,E=0.
\]
Pullback preserves intersections of divisors.
Curves with no common component have nonnegative
intersection, equal to the sum of their positive
local intersection lengths at intersection points.
An infinitely near point is a point on a surface
obtained by successive point blow-ups.
\end{context}

\begin{problem}
For curves \(C,D\) with no common component,
prove
\[
\widetilde C\,\widetilde D=CD-m_P(C)m_P(D).
\]
Deduce that \(CD\) is the sum of the products of
the multiplicities at all common ordinary and
infinitely near points of the successive strict
transforms.
\end{problem}
\end{exerciseintro}

\begin{solution}
Put \(r=m_P(C)\) and \(s=m_P(D)\).
Bilinearity gives
\[
(\pi^*C-rE)(\pi^*D-sE)
=CD+rsE^2=CD-rs.
\]
The mixed terms vanish by the projection formula.

Whenever the current strict transforms meet,
blow up one of their intersection points.
Their intersection number decreases by the
positive integer equal to the product of their
multiplicities there.
The strict transforms still have no common
component, so the new intersection number
remains nonnegative.
This process must terminate after at most
\(CD\) blow-ups. At termination the strict
transforms are disjoint: otherwise another
positive decrease would be possible.
Their final intersection is therefore zero.

Telescope the displayed formula through this
sequence. It gives
\[
CD=\sum_Q m_Q(C_Q)m_Q(D_Q),
\]
where \(C_Q,D_Q\) denote the strict transforms
at the stage containing the center \(Q\).
The centers include each original intersection
and every subsequent common infinitely near
point until the branches separate. Thus this
is exactly the asserted sum, and it is finite.
\end{solution}

\begin{exerciseintro}
% record: H-V-03-003
\exerciseheading{Exercise V.3.3}
\addcontentsline{toc}{subsection}{Exercise V.3.3}

\begin{context}
Let \(\pi:\widetilde X\to X\) blow up a point
\(P\) on a smooth projective surface over an
algebraically closed field, with exceptional curve
\(E\). We use \(E^2=-1\), \(\pi^*D\,E=0\),
and \(\widetilde C=\pi^*C-m_P(C)E\).
For an integral projective curve, a general
hyperplane through a point cuts there with
intersection multiplicity equal to the curve's
multiplicity; the total hyperplane intersection
is its degree.
The Nakai--Moishezon criterion on a surface says
that positive square and positive intersection
with every integral curve imply ampleness.
\end{context}

\begin{problem}
If \(D\) is very ample on \(X\), prove that
\(2\pi^*D-E\) is ample on \(\widetilde X\).
\end{problem}
\end{exerciseintro}

\begin{solution}
Use \(D\) to embed \(X\) into projective space.
For every integral curve \(C\) through \(P\),
a general hyperplane through \(P\), not containing
\(C\), gives
\[
m_P(C)\leq DC.
\tag{1}
\]
For curves avoiding \(P\), put \(m_P(C)=0\);
then (1) holds as well.

Set \(A=2\pi^*D-E\).
Since \(D^2\) is the positive integer degree of
the embedded surface,
\[
A^2=4D^2-1>0,\qquad AE=1>0.
\]
Every other integral curve of \(\widetilde X\)
is the strict transform of an integral curve \(C\)
on \(X\). Its intersection is
\[
A\widetilde C=2DC-m_P(C)\geq DC>0
\]
by (1) and very ampleness.
Thus all conditions of the surface
Nakai--Moishezon criterion hold, and \(A\) is ample.
\end{solution}

\begin{exerciseintro}
% record: H-V-03-004
\exerciseheading{Exercise V.3.4}
\addcontentsline{toc}{subsection}{Exercise V.3.4}

\begin{context}
Let \((A,\mathfrak m,\kappa)\) be a nonzero noetherian
local ring and write \(\ell\) for module length.
The associated graded ring
\[
G=\operatorname{gr}_{\mathfrak m}A
=\bigoplus_{j\geq0}\mathfrak m^j/\mathfrak m^{j+1}
\]
is a finitely generated standard graded
\(\kappa\)-algebra.
For such an algebra of dimension \(d\), the sum
of its graded dimensions below \(l\) is eventually
a polynomial of degree \(d\), with positive leading
coefficient; this is the cumulative form of the
Hilbert-polynomial theorem.
We use homogeneous systems of parameters for graded
algebras, systems of parameters in local rings,
Krull's height theorem, and Nakayama's lemma.
For a smooth surface local ring \(R\) at a closed
point, \(\operatorname{gr}_{\mathfrak m}R=\kappa[u,v]\).
An \(r\)-dimensional projective variety of degree
\(d\) has Hilbert polynomial with leading term
\(dt^r/r!\).
\end{context}

\begin{problem}
(a) Show that \(\ell(A/\mathfrak m^l)\) is eventually
a polynomial \(P_A(l)\) with rational coefficients.
(b) Prove \(\deg P_A=\dim A\).
(c) Define the multiplicity as \((\dim A)!\) times
the leading coefficient.
(d) At a closed point of a curve on a smooth surface,
show that this equals the order of a local equation.
(e) Show that the vertex of the affine cone over
a degree-\(d\) projective variety has multiplicity \(d\).
\end{problem}
\end{exerciseintro}

\begin{solution}
(a) The filtration by powers of \(\mathfrak m\) gives
\[
\ell(A/\mathfrak m^l)
=\sum_{j=0}^{l-1}\dim_\kappa G_j.
\]
The cumulative Hilbert-polynomial theorem therefore
gives the required polynomial, of degree
\(d=\dim G\).

(b) We verify \(d=\dim A\), rather than assuming it.
Put \(n=\dim A\), and choose a parameter ideal
\(\mathfrak q=(a_1,\ldots,a_n)\).
Since \(\mathfrak q\subseteq\mathfrak m\), and
\(\mathfrak q^j/\mathfrak q^{j+1}\) is generated
over \(A/\mathfrak q\) by the degree-\(j\)
monomials in its \(n\) generators,
\[
\ell(A/\mathfrak m^l)
\leq\ell(A/\mathfrak q^l)
\leq\ell(A/\mathfrak q)\binom{l+n-1}{n}.
\]
Thus \(d\leq n\).
If \(n=0\), the ring is artinian and the eventual
polynomial is the positive constant \(\ell(A)\),
which already proves the assertion in that case.

For the reverse inequality choose homogeneous
parameters \(g_1,\ldots,g_d\) in \(G\).
Lift them to elements \(b_i\in\mathfrak m\) of
their respective orders, and put \(J=(b_1,\ldots,b_d)\).
The quotient \(G/(g_1,\ldots,g_d)\) is finite
dimensional. In sufficiently high degree \(N\),
every leading form lies in the ideal of the \(g_i\);
lifting its expression gives
\[
\mathfrak m^N\subseteq J+\mathfrak m^{N+1}.
\]
In \(A/J\), this says
\(\overline{\mathfrak m}^{\,N}
=\overline{\mathfrak m}\,\overline{\mathfrak m}^{\,N}\).
Nakayama gives \(\overline{\mathfrak m}^{\,N}=0\).
Hence \(J\) is \(\mathfrak m\)-primary.
The height theorem, applied to its \(d\) generators,
gives \(n\leq d\). Together with the first bound,
this proves \(\deg P_A=n\).

(c) If \(P_A(l)=cl^n+\) lower terms, define
\[
\mu(A)=n!c.
\]
For a scheme point \(P\), put
\(\mu_P(X)=\mu(\mathcal O_{X,P})\).
The leading coefficient is positive by (a).
It gives a positive integer: the \(n\)-th finite
difference of an integer-valued eventual polynomial
is the constant \(n!c\).
In dimension zero this is the length of the ring.

(d) Let \(R\) be the regular local ring of the
surface, and \(f\in R\) a local equation of the
curve, of order \(r\).
Its initial form \(f_r\) is homogeneous of degree
\(r\) in \(\kappa[u,v]\).
Since this graded ring is a domain, orders add
under multiplication, so the initial ideal of
\((f)\) is exactly \((f_r)\).
Consequently
\[
\operatorname{gr}(R/(f))=\kappa[u,v]/(f_r).
\]
Multiplication by \(f_r\) is injective in the
polynomial ring. The degree-\(j\) component of
the quotient has dimension
\((j+1)-\max(j-r+1,0)=\min(j+1,r)\).
For \(l\geq r\), summation yields
\[
\ell\bigl((R/(f))/\mathfrak m^l(R/(f))\bigr)
=rl-\frac{r(r-1)}2.
\]
Its one-dimensional multiplicity is \(r\),
which is exactly the multiplicity used for
curves in the surface.

(e) Let \(S\) be the homogeneous coordinate ring
of the projective variety, of dimension \(r+1\),
and let \(\mathfrak n=S_+\) be the vertex ideal.
Because \(S\) is standard graded,
\[
\operatorname{gr}_{\mathfrak nS_{\mathfrak n}}
(S_{\mathfrak n})\simeq S.
\]
Indeed, \(\mathfrak n^j/\mathfrak n^{j+1}=S_j\),
and localization does not change these vector
spaces, which are annihilated by \(\mathfrak n\).
Thus the vertex Hilbert--Samuel function is
\(\sum_{j<l}\dim_k S_j\).
The summands have leading term \(dj^r/r!\),
so the sum has leading term \(dl^{r+1}/(r+1)!\).
Its multiplicity is therefore \(d\).
\end{solution}

\begin{exerciseintro}
% record: H-V-03-005
\exerciseheading{Exercise V.3.5}
\addcontentsline{toc}{subsection}{Exercise V.3.5}

\begin{context}
Work first over an algebraically closed field of
characteristic different from two.
For a reduced plane curve germ,
\(\delta\) drops by \(\binom{m}{2}\) when a
point of multiplicity \(m\) is blown up; the
remaining contributions are those of the strict
transform over that point.
A germ \(u^2=v^n\) has \(\delta=\lfloor n/2\rfloor\):
successive blow-ups replace \(n\) by \(n-2\),
ending in a smooth branch or a node.
A unit with constant term one has a square root
in \(k[[v]]\) in this characteristic.
For a degree-\(r\) integral plane curve,
\[
g(\widetilde C)=\frac{(r-1)(r-2)}2-\sum_P\delta_P.
\]
\end{context}

\begin{problem}
Let \(a_1,\ldots,a_r\) be distinct, with \(r\geq5\),
and let \(C\) be the projective closure of
\[
y^2=\prod_{i=1}^r(x-a_i).
\]
Describe its singular point at infinity, calculate
its \(\delta\)-invariant and normalization genus,
and obtain hyperelliptic curves of every genus
at least two. Explain why characteristic two
must be excluded from this construction.
\end{problem}
\end{exerciseintro}

\begin{solution}
The homogeneous equation is
\[
Y^2Z^{r-2}=\prod_{i=1}^r(X-a_iZ).
\]
Its only point at infinity is \(P=[0:1:0]\).
On the chart \(Y=1\), write \(u=X/Y,\ v=Z/Y\).
The equation is
\[
v^{r-2}=\prod_i(u-a_iv).
\]
Its lowest degree is \(r-2\), so \(P\) has
multiplicity \(r-2\geq3\) and is singular.

Blow up using \(v=uw\). After removing the
exceptional factor \(u^{r-2}\), the strict transform
has equation
\[
w^{r-2}=u^2h(w),\qquad h(w)=\prod_i(1-a_iw).
\]
It meets the exceptional curve \(u=0\) only at
\(Q=(0,0)\).
The formal change \(u'=u\sqrt{h(w)}\) identifies
its germ with \((u')^2=w^{r-2}\).
Therefore
\[
\delta_P=\binom{r-2}{2}
+\left\lfloor\frac{r-2}{2}\right\rfloor.
\tag{1}
\]
The affine curve is smooth: at a zero of the
right side, squarefreeness makes its derivative
nonzero, and elsewhere the derivative \(2y\)
is nonzero. Thus (1) accounts for every singularity.
Subtracting it from the plane arithmetic genus gives
\[
g(\widetilde C)=\left\lfloor\frac{r-1}{2}\right\rfloor.
\]
Equivalently, \(\delta_P=2(m-1)^2\) for \(r=2m\),
and \(\delta_P=2m(m-1)\) for \(r=2m+1\).

The function \(x\) defines a separable degree-two
map of the normalization to \(\mathbb P^1\),
because the squarefree polynomial is not a square
in \(k(x)\). Hence the normalization is hyperelliptic.
For any \(g\geq2\), choosing \(r=2g+1\), or \(2g+2\),
gives the required genus.

In characteristic two the same assertion is false.
If \(f(x)=\sum c_ix^i\), choose square roots of the
coefficients and put \(h(t)=\sum\sqrt{c_i}\,t^i\).
Then \(x=t^2,\ y=h(t)\) parametrizes \(y^2=f(x)\).
Since \(f\) is squarefree, \(h(t)\notin k(t^2)\).
The extension \(k(t)/k(t^2)\) has degree two, so
\[
k(t^2,h(t))=k(t).
\]
Thus the normalization is rational, rather than
of the stated positive genus. The numerical
\(\delta\)-formula need not survive either:
for \(f=x^6+x+1\), the derivative is one, so the
affine curve is smooth and its unique point at
infinity has \(\delta=10\), the full plane
arithmetic genus. Formula (1) would give eight.
The characteristic restriction is therefore essential.
\end{solution}

\begin{exerciseintro}
% record: H-V-03-006
\exerciseheading{Exercise V.3.6}
\addcontentsline{toc}{subsection}{Exercise V.3.6}

\begin{context}
For plane curve germs over an algebraically closed
field, analytic isomorphism means isomorphism of
their completed local \(k\)-algebras.
The discrete resolution equivalence of V.3.9.4
records the local branches, exceptional components,
and their multiplicities and incidences at each
point blow-up, without recording cross-ratios of
intersection points on an exceptional line.
Blow-up charts are obtained by substitutions
\(y=xv\) or \(x=yu\), followed by removing the
exceptional factor.
A formal coordinate substitution with invertible
linear part is an automorphism of \(k[[x,y]]\).
Four distinct points on \(\mathbb P^1\) are
projectively equivalent as unordered sets only
if their cross-ratios differ by one of the six
transformations generated by
\(\lambda\mapsto1-\lambda\) and
\(\lambda\mapsto1/\lambda\).
\end{context}

\begin{problem}
Show that analytically isomorphic plane curve
singularities have equivalent embedded resolution
data. Show that the converse fails.
\end{problem}
\end{exerciseintro}

\begin{solution}
Let the singular completed local rings be
\(k[[x,y]]/(f)\) and \(k[[u,v]]/(g)\).
An isomorphism between them is an isomorphism
on their two-dimensional cotangent spaces.
Lift the images of \(x,y\) to \(k[[u,v]]\).
Their linear terms are independent, so the resulting
formal substitution is an automorphism of the
two-variable power series rings.
Its induced quotient map is the given isomorphism,
and consequently it maps \((f)\) onto \((g)\).
In particular it preserves orders, tangent cones,
and the branches of the completed curve.

Blowing up the maximal ideal is functorial under
this formal isomorphism. In the explicit charts,
the corresponding tangent directions give
isomorphic completed local rings of the strict
transforms, together with their exceptional ideals.
Thus multiplicities, branch intersections, and
incidences with exceptional components agree.
Repeat at corresponding centers of an embedded
resolution. The same argument at each step
identifies the entire discrete resolution data,
which proves the asserted implication.
Smooth germs cause no exception; their local
resolution data are already trivial.

For the converse, over \(\mathbb C\) consider
\[
f_\lambda=xy(x-y)(x-\lambda y),
\qquad \lambda\notin\{0,1\}.
\]
Each germ is an ordinary quadruple point.
One blow-up separates its four smooth branches:
they meet the exceptional line transversely at
four distinct points. The multiplicities and
incidence data are identical for every \(\lambda\).
Thus all these germs are equivalent in the
discrete sense.

If two were analytically isomorphic, the linear
part of the ambient formal isomorphism constructed
above would identify their tangent cones.
It would therefore carry the four tangent directions
of one to those of the other by a projective
linear transformation of \(\mathbb P^1\).
For \(\lambda=2\), the unordered cross-ratio orbit is
\(\{2,1/2,-1\}\), whereas \(3\) is outside that orbit.
Hence the germs for \(\lambda=2\) and \(\lambda=3\)
are not analytically isomorphic, although their
embedded resolutions are equivalent.
\end{solution}

\begin{exerciseintro}
% record: H-V-03-007
\exerciseheading{Exercise V.3.7}
\addcontentsline{toc}{subsection}{Exercise V.3.7}

\begin{context}
Work over an algebraically closed field.
The strict transform under a point blow-up is
obtained in a chart by substitution and division
by the exceptional factor of order equal to the
curve's multiplicity.
An embedded resolution requires the reduced total
transform to have smooth components and only
transverse double intersections.
In particular, an ordinary triple intersection
requires one more blow-up.
For a reduced plane curve germ,
\[
\delta=\sum_Q\binom{m_Q}{2},
\]
where \(m_Q\) is the multiplicity of the strict
transform at each ordinary or infinitely near
singular point.
Resolution equivalence preserves the branches,
multiplicities, and incidence with exceptional
components at every step.
\end{context}

\begin{problem}
Give embedded resolutions, calculate \(\delta\),
and determine the equivalence classes of:

(a) \(x^3+y^5=0\).

(b) \(x^3+x^4+y^5=0\).

(c) \(x^3+y^4+y^5=0\).

(d) \(x^3+y^5+y^6=0\).

(e) \(x^3+xy^3+y^5=0\).
\end{problem}
\end{exerciseintro}

\begin{solution}
Let \(E_i\) denote the exceptional component of
the \(i\)-th blow-up in each calculation.
All five curves have multiplicity three at the
origin and unique tangent line \(x=0\).
The first chart is \(x=yu\), with exceptional
line \(E_1:y=0\).
The strict-transform equations are
\[
\begin{array}{c|l}
\text{case}&\text{equation after the first blow-up}\\ \hline
(a)&u^3+y^2\\
(b)&u^3+yu^4+y^2\\
(c)&u^3+y+y^2\\
(d)&u^3+y^2+y^3\\
(e)&u^3+yu+y^2.
\end{array}
\tag{1}
\]
Each meets \(E_1\) only at its origin.

(a) At that point the strict transform has
multiplicity two and tangent \(y=0\).
Blow up with \(y=uv\). Its new equation is
\(u+v^2=0\), with \(E_2:u=0\) and
the strict transform of \(E_1\) given by \(v=0\).
The curve is smooth but tangent to \(E_2\)
at the intersection of these two exceptional curves.
Blow up with \(u=vw\). The curve becomes
\(w+v=0\), while the exceptional components
through the point are \(E_3:v=0\) and \(E_2:w=0\).
These form three distinct tangent lines together
with the curve. One final blow-up separates
that ordinary triple point and gives normal crossings.
The strict-curve multiplicities at the four centers
are \(3,2,1,1\). Hence
\[
\delta_a=\binom32+\binom22=4.
\]

(b) Use the same successive charts.
After the second and third blow-ups, the equations are
\[
u+u^3v+v^2=0,\qquad
w+v+v^3w^3=0.
\]
Their relevant linear terms and tangencies are
the same as in (a): the second transform is
tangent to \(E_2\), and the third forms an
ordinary triple point with \(E_2,E_3\).
A fourth blow-up resolves it.
Thus the same multiplicities and incidences occur,
\(\delta_b=4\), and (b) is equivalent to (a).

(c) The first equation in (1) for this case is
already smooth, but it has contact order three
with \(E_1:y=0\).
Blow up with \(y=uv\); the strict transform is
\[
u^2+v+uv^2=0.
\]
It is smooth and tangent to \(E_1:v=0\) at the
intersection with \(E_2:u=0\).
Next set \(v=uw\); the equation becomes
\[
u+w+u^2w^2=0.
\]
This is an ordinary triple intersection with
\(E_3:u=0\) and \(E_1:w=0\).
A fourth blow-up gives normal crossings.
The strict-curve multiplicities at the centers
are \(3,1,1,1\), so
\[
\delta_c=\binom32=3.
\]
In particular (c) is not equivalent to (a) or (b).

(d) The same charts as in (a) give, after the
second and third blow-ups,
\[
u+v^2+uv^3=0,\qquad
w+v+v^3w=0.
\]
Again the second transform is smooth and tangent
to \(E_2\), and the third gives the same
ordinary triple point. The fourth blow-up
resolves it. Thus \(\delta_d=4\), and (d)
is equivalent to (a) and (b).

(e) The first transform has an ordinary node:
its quadratic tangent cone is \(y(u+y)\),
with two distinct tangent directions in every
characteristic.
The line \(E_1:y=0\) is tangent to one branch.
Blow up with \(y=uv\). After division by \(u^2\)
the equation is
\[
u+v+v^2=0.
\]
On \(E_2:u=0\), it has two points, \(v=0\)
and \(v=-1\).
At \(v=-1\) it is smooth and transverse to
\(E_2\); no other exceptional curve passes there.
At \(v=0\) it forms an ordinary triple point
with \(E_2:u=0\) and \(E_1:v=0\), since
its tangent is \(u+v=0\).
One further blow-up at that point gives normal
crossings. The multiplicities at the centers
are \(3,2,1\), giving
\[
\delta_e=\binom32+\binom22=4.
\]
This germ has two branches, whereas each of
(a)--(d) has one, as their resolutions show.
It is therefore inequivalent to all four.

The equivalence classes are consequently
\[
\{(a),(b),(d)\},\qquad\{(c)\},\qquad\{(e)\},
\]
with \(\delta\)-values \(4,3,4\), respectively.
All the chart calculations use unit linear
coefficients and distinct coordinate directions;
they remain valid in every characteristic.
\end{solution}

\begin{exerciseintro}
% record: H-V-03-008
\exerciseheading{Exercise V.3.8}
\addcontentsline{toc}{subsection}{Exercise V.3.8}

\begin{context}
Work over an algebraically closed field.
The multiplicity of a reduced plane germ is
the least degree in its local equation.
After a point blow-up, remove that power of
the exceptional equation to get the strict transform.
The delta invariant is the sum of
\(\binom{m_Q}{2}\) over ordinary and infinitely
near singular points of the strict transforms.
Embedded resolution equivalence additionally
preserves individual branch multiplicities and
their incidences with exceptional components.
The germ \(x^a-y^b\), with coprime positive
integers \(a,b\), has one branch, parametrized
by \(x=t^b,y=t^a\); this holds in every characteristic.
\end{context}

\begin{problem}
Check the textbook's claim that the following
two germs have the same multiplicities and
configuration of infinitely near singular points,
but are not equivalent:

(a) \(x^4-xy^4=0\).

(b) \(x^4-x^2y^3-x^2y^5+y^8=0\).

If the claimed numerical agreement fails,
identify it explicitly and give a corrected pair
that illustrates the intended distinction.
\end{problem}
\end{exerciseintro}

\begin{solution}
(a) The germ is \(x(x^3-y^4)=0\), of multiplicity
four, with a smooth branch and a branch of
multiplicity three.
Blow up with \(x=yu\). Removing \(y^4\) gives
\[
u(u^3-y)=0.
\]
Its only point on the exceptional line \(y=0\)
is the origin. The strict transform has
multiplicity two there, with two smooth branches
having distinct tangents \(u=0\) and \(y=0\).
Blowing up that node separates the branches and
makes the strict curve nonsingular.
Consequently its singular-point multiplicities
are \(4,2\), and
\[
\delta_a=\binom42+\binom22=7.
\tag{1}
\]
Extra blow-ups needed for normal crossings with
the exceptional components do not change this
delta calculation.

(b) Factor the second polynomial:
\[
(x^2-y^3)(x^2-y^5).
\]
It has two distinct branches, both of multiplicity
two, and total multiplicity four.
Under \(x=yu\), the strict transform is
\[
(u^2-y)(u^2-y^3)=0.
\]
At its only exceptional point, the lowest term
is \(-u^2y\), of degree three.
Thus its infinitely near multiplicity is three,
already different from (a).
One branch, \(u^2-y=0\), is smooth with tangent
\(y=0\); the other, \(u^2-y^3=0\), is a cusp
with tangent \(u=0\).
Blow up this point. In the chart \(y=uv\), the
first branch becomes \(u-v=0\), and the second
has no point at that chart's origin.
In the other chart \(u=yw\), the second becomes
\(w^2-y=0\), while the first has no point at
that origin. Both strict branches are now smooth
and disjoint.
The singular-point multiplicities are therefore
\(4,3\), and
\[
\delta_b=\binom42+\binom32=9.
\tag{2}
\]
Equations (1)--(2) disprove the printed agreement.
The printed germs are indeed inequivalent, but
even their delta invariants differ.

The intended example is recovered by replacing
the exponent four in the term \(xy^4\) of (a)
by five:
\[
C_a':\quad x^4-xy^5=x(x^3-y^5)=0.
\]
After \(x=yu\), its strict transform is
\(u(u^3-y^2)=0\).
It has multiplicity three at the origin,
with a smooth branch tangent to \(u=0\) and
a cusp tangent to \(y=0\).
Blowing up separates the two branches; the
cusp becomes smooth, since in the chart
\(y=uv\) its equation is \(u-v^2=0\).
The other branch is smooth in the complementary
chart.
Thus \(C_a'\) and the germ in (b) each have
exactly one infinitely near singular point:
their multiplicities are \(4,3\), giving
\(\delta=9\) in both cases.

Nevertheless they are inequivalent.
The original branches of \(C_a'\) have
multiplicities \(1,3\), while those of (b)
have multiplicities \(2,2\).
No correspondence preserving individual branch
multiplicities can identify their embedded
resolution data. The corrected pair therefore
demonstrates exactly why the total multiplicities
at infinitely near singular points are weaker
than the full embedded resolution invariant.
\end{solution}

\begin{exerciseintro}
\corpussection{V.4 The Cubic Surface in Projective 3-Space}
\addcontentsline{toc}{section}{V.4 The Cubic Surface in Projective 3-Space}

% record: H-V-04-001
\exerciseheading{Exercise V.4.1}
\addcontentsline{toc}{subsection}{Exercise V.4.1}

\begin{context}
Work over an algebraically closed field.
The closure of the graph of a rational map
defined by generating sections of an ideal sheaf
is the blow-up of that base ideal.
Multiplying the ideal by an invertible sheaf
does not change its blow-up.
The quadric \(AD=BC\) in \(\mathbb P^3\)
is smooth, and is the Segre image of
\(\mathbb P^1\times\mathbb P^1\).
Inverse birational maps have the same graph
closure, with its two projections interchanged.
\end{context}

\begin{problem}
Blow up two distinct points of \(\mathbb P^2\).
Show that the conics through those points define
a morphism to a smooth quadric, and that this
morphism identifies the blown-up plane with the
quadric blown up at one point.
\end{problem}
\end{exerciseintro}

\begin{solution}
Choose the points to be
\(P_1=[1:0:0]\) and \(P_2=[0:1:0]\).
The complete vector space of conics through them
has basis \(xy,xz,yz,z^2\), so its rational map is
\[
[x:y:z]\dashrightarrow[xy:xz:yz:z^2].
\tag{1}
\]
Near \(P_1\), where \(x=1\), its base ideal
is \((y,z)\); near \(P_2\) it is \((x,z)\).
There are no other base points.
Consequently its graph closure is exactly
\(X'=\operatorname{Bl}_{P_1,P_2}\mathbb P^2\),
and the second projection is a morphism on \(X'\).

The coordinates in (1) satisfy \(AD=BC\).
On \(z\ne0\), (1) is
\[
(u,v)\longmapsto[uv:u:v:1],
\]
an isomorphism onto the open subset \(D\ne0\)
of that quadric. Thus the morphism is birational
onto the entire smooth quadric \(Q\).
The inverse rational map is
\[
Q\dashrightarrow\mathbb P^2,\qquad
[A:B:C:D]\longmapsto[B:C:D].
\tag{2}
\]
Its base locus is the single point
\(q=[1:0:0:0]\).
In the chart \(A=1\), the equation is \(D=BC\),
so the base ideal \((B,C,D)\) equals the maximal
ideal \((B,C)\) of that smooth point.
Elsewhere the sections in (2) have no common zero.
Hence the graph closure of (2) is
\(\operatorname{Bl}_qQ\).

The two maps are birational inverses, so these
graph closures are the same closed subscheme
of \(\mathbb P^2\times Q\).
It follows that
\[
X'\simeq\operatorname{Bl}_qQ,
\]
compatibly with the conic-system morphism to \(Q\).
Its exceptional curve is the strict transform
of the line \(z=0\) joining \(P_1,P_2\),
which (1) sends to \(q\).
\end{solution}

\begin{exerciseintro}
% record: H-V-04-002
\exerciseheading{Exercise V.4.2}
\addcontentsline{toc}{subsection}{Exercise V.4.2}

\begin{context}
Let \(P_1,P_2,P_3\) be noncollinear points
of \(\mathbb P^2\) over an algebraically closed field.
Their quadratic transformation is resolved on
\(S=\operatorname{Bl}_{P_1,P_2,P_3}\mathbb P^2\).
Write \(l\) for the pullback of a line and \(e_i\)
for the exceptional classes.
Then \(l^2=1,\ e_i^2=-1,\ le_i=e_ie_j=0\) for
\(i\ne j\).
The other projection to \(\mathbb P^2\) has
line pullback \(l'=2l-e_1-e_2-e_3\) and
exceptional classes
\(e_i'=l-e_j-e_k\), where \(\{i,j,k\}=\{1,2,3\}\).
The curve \(e_i'\) contracts to the inverse
base point \(Q_i\).
A strict transform meets the exceptional curve
of a point blow-up with intersection equal to
its original multiplicity at that point.
\end{context}

\begin{problem}
For an integral degree-\(d\) plane curve \(C\)
of multiplicities \(r_i\) at \(P_i\), not one
of the three contracted joining lines, prove
that its transformed curve has
\[
d'=2d-r_1-r_2-r_3,\qquad
r_i'=d-r_j-r_k
\]
at the inverse base points.
Explain the excluded joining-line cases.
\end{problem}
\end{exerciseintro}

\begin{solution}
The strict transform of \(C\) on \(S\) has class
\[
\widetilde C=dl-r_1e_1-r_2e_2-r_3e_3.
\]
By assumption it is not exceptional for the
second projection. Its image \(C'\) is therefore
an integral curve, with the same strict transform
\(\widetilde C\) on \(S\).
Intersecting with the pullback of a target line gives
\[
d'=\widetilde C\,l'=2d-r_1-r_2-r_3.
\]
Intersecting instead with the target exceptional
curve gives
\[
r_i'=\widetilde C\,e_i'
=\widetilde C(l-e_j-e_k)=d-r_j-r_k.
\]
No ordinary-singularity assumption was used:
the strict-transform and intersection formulas
apply to arbitrary singularities.
The nonnegativity of each \(r_i'\) also follows
from Bezout with the line \(P_jP_k\), which is
not a component of \(C\).

If \(C=P_jP_k\), its strict transform is \(e_i'\)
and is contracted to \(Q_i\).
It has no transformed curve whose point
multiplicities could satisfy the displayed formulas.
Indeed, those expressions give \(d'=0\) and
\(r_i'=-1\) in this case.
Thus the curve formulation implicitly requires
excluding the three contracted lines; their
actual images are the corresponding points.
\end{solution}

\begin{exerciseintro}
% record: H-V-04-003
\exerciseheading{Exercise V.4.3}
\addcontentsline{toc}{subsection}{Exercise V.4.3}

\begin{context}
Work over an algebraically closed field.
A plane curve point is ordinary if its branches
are smooth with distinct tangent directions.
A quadratic transformation at three noncollinear
points blows up those points and blows down
their three joining lines.
Blowing down an exceptional curve meeting the
strict curve transversely at distinct smooth
points creates an ordinary multiple point.
At a multiplicity-\(m\) point, the sum of the
delta invariants of the strict transform above
it is \(\delta-\binom m2\).
An integral plane curve has only finitely many
singular points.

Call a point \(P\) a strange center for a curve
if every tangent at a general smooth point
passes through \(P\).
Projection from \(P\), on the normalization,
is separable exactly when \(P\) is not a strange
center: the differential of the pencil coordinate
vanishes exactly when its line is tangent.
For a separable projection, a general line
through \(P\) meets the curve away from \(P\)
transversely at smooth points.
General lines without a prescribed center also
have this transversality property.
\end{context}

\begin{problem}
Show that any integral plane curve can be
transformed, by finitely many quadratic
transformations, to a curve with only ordinary
singularities.
\end{problem}
\end{exerciseintro}

\begin{solution}
\emph{Preparatory lemma.} We first explain how to avoid a possible
inseparability obstruction to using a singular
point as a center.

Given finitely many points \(P\) on an integral
curve \(C\), a suitably general quadratic
transformation with its three centers outside
\(C\) has the following properties:
it is a local isomorphism at every singular
point of \(C\), creates only ordinary new
singularities, and makes the images of all
the prescribed points non-strange centers.
The first two assertions follow by requiring
the three contracted joining lines to avoid
the singular points and meet \(C\) transversely.
These are nonempty open conditions on the
ordered triples of centers.

Here is why the last assertion is also a
nonempty open condition for any fixed \(P\).
Choose a smooth point \(A\ne P\) and let \(T\)
be its tangent line.
If \(P\in T\), choose the three centers off
\(T\). The strict image of \(T\) is a smooth
conic: the quadratic transformation restricts
to its degree-two embedding.
The images of \(A,P\) are distinct points
on that conic, and the tangent at the first
does not pass through the second.
The image curve has that same tangent at
the image of \(A\), so its image of \(P\)
is not a strange center.
If \(P\notin T\), choose one center on \(T\),
outside \(C\), and the other two generically.
Then the strict image of \(T\) is a line.
Arrange that \(A,P\) avoid the contracted
lines. Their local images show that the
image tangent at \(A\) does not contain the
image of \(P\).
In both cases this proves the desired
nonempty open condition, since nonincidence
of that tangent and point is open.
The parameter space of noncollinear triples
is irreducible, so the finitely many conditions
can be satisfied simultaneously.

\emph{The local improvement.}
Now let \(P\) be a nonordinary singular point.
Apply such a preparatory transformation if
necessary. It preserves all existing
nonordinary germs and adds only ordinary
ones, and we may suppose \(P\) is not a
strange center.
Choose two further centers \(Q,R\) outside
the curve so that the lines \(PQ,PR\) have
directions different from all branch tangents
at \(P\), and meet the curve elsewhere
transversely at smooth points.
Also require \(QR\) to meet the curve
transversely at smooth points and all three
lines to avoid the other singularities.
Separability of projection from \(P\) and
generic transversality make these simultaneous
open conditions possible.

Resolve the quadratic transformation at
\(P,Q,R\).
Since \(Q,R\) are off the curve, the only
change to its singular germs on the blown-up
surface is the point blow-up at \(P\).
The strict transforms of \(PQ,PR\) do not
meet the strict curve above \(P\), because
their directions there were excluded.
Each of the three lines to be contracted
therefore meets the strict curve transversely
at distinct smooth points.
Blowing them down creates only ordinary
singularities.
All remaining germs of the strict curve,
including its points above \(P\), are
unchanged by these contractions.

\emph{Termination.} Use the nonnegative integer
\[
M(C)=\sum_{\substack{P\text{ singular}\\
                    P\text{ not ordinary}}}\delta_P.
\]
A preparatory transformation leaves \(M\)
unchanged. The subsequent transformation
at a nonordinary point of multiplicity \(m\)
decreases \(M\) by at least \(\binom m2>0\):
the sum of all new delta invariants above
that point has this decrease, and any of
them that become ordinary are no longer
counted. All other new points are ordinary.
After finitely many repetitions \(M=0\),
which means that every remaining singularity
is ordinary. Each step was a quadratic
transformation at three suitable ordinary
points of the ambient plane, as required.
\end{solution}

\begin{exerciseintro}
% record: H-V-04-004
\exerciseheading{Exercise V.4.4}
\addcontentsline{toc}{subsection}{Exercise V.4.4}

\begin{context}
Work over an algebraically closed field.
The nine-point theorem (Theorem V.4.5 in Hartshorne,
not Exercise V.4.5 in this collection) says that eight
plane points, no four collinear and no seven on
a conic, determine a ninth base point for all
cubics through them, allowing the stated
infinitely near versions.
On an integral plane cubic, write \(A\circ B\)
for the third intersection with the chord \(AB\),
using the tangent if \(A=B\).
On the regular locus this construction extends
algebraically to the tangency cases, and
\(A\circ(A\circ B)=B\), counting intersection
multiplicities.
Collinearity identities established for distinct
intersections therefore extend to all cases
where the chord construction is defined.
\end{context}

\begin{problem}
(a) Suppose lines cut an integral cubic in
\(P,Q,R\) and \(P',Q',R'\), respectively.
Prove that \(P\circ P',Q\circ Q',R\circ R'\)
are collinear.
(b) If \(O\) is a regular inflection point, use
(a) to prove associativity of the operation
\[
A+B=O\circ(A\circ B)
\]
on the regular points.
\end{problem}
\end{exerciseintro}

\begin{solution}
(a) First take the intersections distinct and
in general position for these constructions.
The union of the three chords \(PP',QQ',RR'\)
meets the cubic in the six original points
and in
\[
P''=P\circ P',\quad Q''=Q\circ Q',\quad
R''=R\circ R'.
\]
Let \(M\) be the line through \(P'',Q''\).
The cubic which is the union of the two original
lines and \(M\) passes through eight of these
nine intersection points.
No four of those eight are collinear and no
seven lie on a conic: either event would
contradict Bezout with the integral cubic.
The nine-point theorem forces this reducible
cubic through \(R''\) as well.
In the general situation \(R''\) is on neither
original line, so it lies on \(M\).
The tangent and coincident-point versions follow
by the algebraic extension of the chord
construction, or by its infinitely near form.

(b) Part (a), applied to the triples
\((A,B,A\circ B)\) and \((C,D,C\circ D)\),
gives the medial identity
\[
(A\circ B)\circ(C\circ D)
=(A\circ C)\circ(B\circ D).
\tag{1}
\]
This uses the fact that the third point on
the line through \(A\circ C,B\circ D\) is
\((A\circ B)\circ(C\circ D)\).

Put \(j(A)=O\circ A\).
The chord rule gives \(j^2=\operatorname{id}\);
the inflection hypothesis gives \(O\circ O=O\),
so \(j(O)=O\).
Taking \(C=D=O\) in (1) yields
\[
j(A\circ B)=j(A)\circ j(B).
\tag{2}
\]
Thus \(j\) preserves the chord operation.
Since \(A+B=j(A\circ B)\), applying (2) twice
and then (1) shows
\[
\begin{split}
(A+B)+(C+D)
&=(A\circ B)\circ(C\circ D)\\
&=(A\circ C)\circ(B\circ D)\\
&=(A+C)+(B+D).
\end{split}
\tag{3}
\]
Also \(A+O=j(A\circ O)=j^2(A)=A\).
Set \(C=O\) in (3) and replace \(D\) by \(C\).
The resulting equality is
\[
(A+B)+C=A+(B+C),
\]
which proves associativity directly from (a).
The same construction is commutative, and
\(A+j(A)=O\), so it gives the asserted group law.
\end{solution}

\begin{exerciseintro}
% record: H-V-04-005
\exerciseheading{Exercise V.4.5}
\addcontentsline{toc}{subsection}{Exercise V.4.5}

\begin{context}
Work over an algebraically closed field.
The nine-point theorem for cubics says that,
in a transverse complete intersection of two
cubics with the usual general-position conditions,
every cubic through eight of the nine points
passes through the ninth.
Collinearity is a closed algebraic condition.
For six points on a fixed smooth conic,
the parameter space is irreducible, and the
intersection of two distinct joining lines
varies algebraically wherever it is defined.
\end{context}

\begin{problem}
For six points \(A,B,C,A',B',C'\) on a smooth
conic, prove Pascal's theorem:
\[
P=AB'\cap A'B,\quad
Q=AC'\cap A'C,\quad
R=BC'\cap B'C
\]
are collinear, whenever these joining lines
and their intersections are defined.
\end{problem}
\end{exerciseintro}

\begin{solution}
First choose the six points generically.
Consider the two reducible cubics
\[
T_1=AB'\cup BC'\cup CA',\qquad
T_2=A'B\cup B'C\cup C'A.
\]
Their nine intersection points are exactly
the six conic points and \(P,Q,R\).
For example, the three intersections on
\(AB'\) are \(P,B',A\);
those on \(BC'\) are \(B,R,C'\);
those on \(CA'\) are \(A',C,Q\).

For a generic choice, these intersections are transverse and
no four of the eight points consisting of the six conic points
and \(P,Q\) are collinear. The unwanted incidences are excluded
by the generic choice. No seven are on a conic: seven of these
eight include at least five original conic points, and a conic
through those five must be the original one by Bezout; neither
\(P\) nor \(Q\) lies on it. These are the general-position
hypotheses of the nine-point theorem.

Let \(M=PQ\).
The union of the original conic and \(M\)
is another cubic, through eight of the nine
points. The nine-point theorem makes it pass
through \(R\).
In the generic configuration \(R\) does not
lie on the conic, so it lies on \(M\).
This proves collinearity for a dense open set
of configurations.

The sixfold product of the conic is irreducible,
and the displayed line intersections are
algebraic maps on their domain of definition.
The determinant expressing collinearity
vanishes on that dense open set, and therefore
vanishes throughout its domain.
Thus the conclusion holds for every configuration
for which \(P,Q,R\) are defined; tangent versions
are obtained by the same limiting construction.
\end{solution}

\begin{exerciseintro}
% record: H-V-04-006
\exerciseheading{Exercise V.4.6}
\addcontentsline{toc}{subsection}{Exercise V.4.6}

\begin{context}
Work over an algebraically closed field.
The vector space of plane quartic forms has
dimension fifteen.
Two quartics without a common component form
a complete intersection of length sixteen.
Its ideal is generated by the two equations;
in degree four its space of equations is exactly
their two-dimensional span.
The base scheme of a linear system is the common
zero scheme of all its sections.
\end{context}

\begin{problem}
Explain under what hypotheses thirteen prescribed
plane points force three further points through
which every quartic containing the thirteen must
pass. State what can go wrong without those
hypotheses.
\end{problem}
\end{exerciseintro}

\begin{solution}
Let \(Z=\{P_1,\ldots,P_{13}\}\) be thirteen
distinct points and put \(V=H^0(\mathcal I_Z(4))\).
The precise hypotheses for a pencil with three
distinct additional ordinary base points are:
\[
\dim V=2,\quad
\text{the pencil has no fixed curve component},\quad
\text{its base scheme is reduced}.
\tag{1}
\]
The first condition says that the thirteen
point conditions are independent.
Choose a basis \(F,G\) of \(V\).
The second says that they have no common component.
Bezout then makes \(B=V(F,G)\) a scheme of
length sixteen.
By the third condition it consists of sixteen
distinct points. Thirteen are the assigned
points, so the remaining three are uniquely
determined:
\[
B=Z\sqcup\{P_{14},P_{15},P_{16}\}.
\]
Every quartic through \(Z\) has equation
\(\lambda F+\mu G\), so it contains all three
residual points.
The base scheme is independent of the basis
of \(V\), proving their intrinsic determination.

These conditions are also necessary if the
common base scheme is to consist of precisely
sixteen distinct ordinary points.
Finiteness rules out a fixed curve.
Choose two members with no common component;
their intersection has length sixteen and
contains the sixteen base points, so it is
exactly that reduced scheme.
Any other quartic in \(V\) belongs to its
degree-four ideal, which is the span of the two
chosen equations. Hence \(\dim V=2\).

If \(\dim V>2\) but there is no fixed component,
a third independent member cannot contain the
entire intersection of the first two: otherwise
it would belong to their two-dimensional span.
Thus the asserted three additional common
ordinary points need not exist.
If there is a fixed curve, the base locus is
positive dimensional; for example, thirteen
points on a conic force every such quartic to
contain that conic by Bezout.
Finally, without reducedness, the length-sixteen
intersection can have coincident or infinitely
near base points. One still has an intersection
scheme, but cannot assert three distinct new
ordinary points. This is the appropriate
scheme-theoretic qualification of the statement.
\end{solution}

\begin{exerciseintro}
% record: H-V-04-007
\exerciseheading{Exercise V.4.7}
\addcontentsline{toc}{subsection}{Exercise V.4.7}

\begin{context}
Let \(X\) be a smooth cubic surface over an algebraically closed
field, presented as the blow-up of six plane points in general
position.  Write \(l\) for the pullback of a line and \(e_i\) for
the exceptional classes.  Then
\[
l^2=1,\quad e_i^2=-1,\quad l e_i=e_i e_j=0\ (i\ne j),
\qquad H=-K_X=3l-\sum_{i=1}^6e_i.
\]
The 27 lines have classes \(e_i\), \(l-e_i-e_j\), and
\(2l-\sum_{j\ne i}e_j\).  They are the only integral curves of
negative self-intersection.  The arithmetic genus of a divisor is
\(p_a(D)=1+(D^2-DH)/2\).

The Hodge index theorem gives \(D^2H^2\leq(DH)^2\).
We use the smooth-member criterion proved in Exercise V.4.8:
a class nonnegative on all 27 lines and of positive square contains
a smooth integral curve.  The class \(H-L\) for a line \(L\)
is a conic pencil with a smooth general member.
\end{context}

\begin{problem}
For every divisor \(D\) of degree \(d=DH\), prove
\[
p_a(D)\leq
\begin{cases}
(d-1)(d-2)/6,&d\equiv1,2\pmod3,\\
(d-1)(d-2)/6+1/3,&d\equiv0\pmod3.
\end{cases}
\]
For each positive integer \(d\), exhibit a smooth integral curve
attaining this bound.
\end{problem}
\end{exerciseintro}

\begin{solution}
Since \(H^2=3\), Hodge index and adjunction give
\[
p_a(D)\leq 1+\frac{d^2/3-d}{2}
=\frac{(d-1)(d-2)}6+\frac13.
\]
The arithmetic genus is an integer.  If \(d\equiv1,2\pmod3\),
the first term on the right is an integer, so rounding down
removes \(1/3\).  If \(3\mid d\), the whole right side is an
integer.  This proves the asserted bounds, including for divisors
of nonpositive degree.

Fix a line \(L\), and put \(F=H-L\).  For positive degrees take
\[
\begin{array}{c|c|c}
d&\text{class}&\text{square}\\ \hline
3n&nH&3n^2\\
3n+1&nH+L&3n^2+2n-1\\
3n+2&nH+F&3n^2+4n .
\end{array}
\]
Here \(n\geq1\) in the first row and \(n\geq0\) in the others.
For \(n=0\), use the line or a smooth conic, both of genus zero.
For \(n\geq1\), all three squares are positive.
The first class is ample.  The second meets \(L\) in \(n-1\)
and every other line nonnegatively.
For a line \(M\ne L\), \(F M=1-LM\geq0\), while \(FL=2\);
thus the third class is nonnegative on all lines as well.
The smooth-member criterion provides a smooth integral curve in
each class.  Substituting the displayed squares into adjunction
gives equality in the corresponding bound.
\end{solution}

\begin{exerciseintro}
% record: H-V-04-008
\exerciseheading{Exercise V.4.8}
\addcontentsline{toc}{subsection}{Exercise V.4.8}

\begin{context}
Let \(X\) be a smooth cubic surface over an algebraically closed
field, presented as the blow-up of six plane points in general
position.  Write \(l\) for the pullback of a line and \(e_i\) for
the exceptional classes.  Then
\[
l^2=1,\quad e_i^2=-1,\quad l e_i=e_i e_j=0\ (i\ne j),
\qquad H=-K_X=3l-\sum_{i=1}^6e_i.
\]
The 27 lines have classes \(e_i\), \(l-e_i-e_j\), and
\(2l-\sum_{j\ne i}e_j\).  They are the only integral curves of
negative self-intersection.  The arithmetic genus of a divisor is
\(p_a(D)=1+(D^2-DH)/2\).

A conic class means a ruling class \(H-L\), where \(L\) is one
of the 27 lines; its general member is a smooth integral curve
of degree two and square zero.

Quadratic transformations centered at three marked points give
new blow-up presentations of the same cubic surface.
For \(D=al-\sum b_ie_i\), the new coefficients are
\[
a'=2a-b_1-b_2-b_3,\qquad
b'_i=a-b_j-b_k\quad(\{i,j,k\}=\{1,2,3\}),
\]
with the other coefficients unchanged.
Blowing up three, four, or five points in general position has
very ample anticanonical class.  The conics through two points
give the morphism to a quadric in Exercise V.4.1.
We use Bertini irreducibility for a birational linear system with
two-dimensional image, and Bertini smoothness where its associated
morphism is an immersion.  These forms apply in every characteristic.
\end{context}

\begin{problem}
Prove that a divisor class \(D\) on \(X\) contains an integral
curve if and only if it contains a smooth integral curve, if and
only if it is one of the following:
(a) a line class;
(b) a conic class of square zero;
(c) a class of positive square with \(DL\geq0\) for all 27 lines.
\end{problem}
\end{exerciseintro}

\begin{solution}
We first analyze any class nonnegative on the 27 lines.
Write \(D=al-\sum b_ie_i\), and reorder so
\(b_1\geq\cdots\geq b_6\).  The line inequalities imply
\[
b_i\geq0,\qquad a\geq b_i+b_j,\qquad
2a\geq\sum_{j\ne i}b_j.
\tag{1}
\]
In particular \(a\geq0\).
If \(a<b_1+b_2+b_3\), perform the quadratic transformation
at the first three points.  Its new \(b_i'\) are nonnegative
by (1), and \(0\leq a'<a\).
Reordering and repeating must terminate.  Thus after a new marking
we have
\[
a\geq b_1+b_2+b_3,\qquad b_1\geq\cdots\geq b_6\geq0.
\tag{2}
\]

Define
\[
G_0=l,\quad G_1=l-e_1,\quad G_2=2l-e_1-e_2,\quad
G_j=3l-e_1-\cdots-e_j\quad(3\leq j\leq6).
\]
The first is a pullback from the plane, the second a pencil,
and the third the conic system giving a quadric.
For \(3\leq j\leq5\), \(G_j\) is the pullback of the very ample
anticanonical system on the blow-up of \(j\) points; \(G_6=H\).
Consequently all these classes are base point free and nef.
There is the explicit decomposition
\[
\begin{split}
D={}&(a-b_1-b_2-b_3)G_0+(b_1-b_2)G_1\\
 &+(b_2-b_3)G_2+(b_3-b_4)G_3\\
 &+(b_4-b_5)G_4+(b_5-b_6)G_5+b_6G_6 .
\end{split}
\tag{3}
\]
All coefficients are nonnegative integers.  Thus \(D\) itself
is base point free and nef.

Only \(G_1\) has square zero; all the other \(G_i\) have
positive square.  Since pairwise intersections of nef classes
are nonnegative, a nonzero \(D\) of square zero must be
\(mG_1\) for some \(m>0\).
The morphism defined by \(G_1\) is the connected-fiber ruling
to \(\mathbb P^1\).  Hence its full system \(|mG_1|\) is pulled
back from \(|\mathcal O_{\mathbb P^1}(m)|\).  It has an integral
member exactly when \(m=1\), and its general fiber is a smooth
conic on the cubic surface.

If \(D^2>0\), some generator \(B\ne G_1\) occurs in (3).
Write \(D=B+N\), with \(N\) base point free.
The morphism for \(B\) is birational and is an isomorphism onto
its image off finitely many exceptional lines.
They are the \(e_i\) for \(B=G_0\), the lines
\(l-e_1-e_2,e_3,\ldots,e_6\) for \(B=G_2\), and the
\(e_i\) with \(i>j\) for \(B=G_j\), \(3\leq j\leq6\).
On their complement, multiply sections of \(B\) by a section
of \(N\) nonvanishing at a given point.  This shows that the
system \(D\) is immersive there and is birational onto its image.
Bertini makes its general member integral and smooth on this open.

Smoothness along the omitted lines requires a separate check.
Every generator in (3) restricts to each such line \(E\)
with degree zero or one.  A base point free degree-one system
on \(E\cong\mathbb P^1\) is complete.  Multiplication of their
sections therefore shows that the restriction of \(|D|\) to \(E\)
contains the complete system
\(|\mathcal O_E(DE)|\).
If \(DE=0\), a general member avoids \(E\).
If \(DE>0\), it meets \(E\) in distinct transverse points and is
smooth at them.  There are only finitely many such lines,
so a general member is smooth everywhere.

Now let \(C\) be any integral member of a class \(D\).
If \(C^2<0\), it is a line, giving (a).
Otherwise \(C\) is distinct from every line, so its intersections
with all lines are nonnegative.
The preceding analysis gives (b) when \(C^2=0\), and (c) when
\(C^2>0\).  Conversely:

(a) Every line is smooth and integral.

(b) A square-zero conic class is the ruling class \(G_1\)
after a marking, and has a smooth integral general member.

(c) The argument following (3) gives a smooth integral member.

Thus each numerical alternative implies a smooth integral member,
and such a member is in particular integral.  All three conditions
are equivalent.
\end{solution}

\begin{exerciseintro}
% record: H-V-04-009
\exerciseheading{Exercise V.4.9}
\addcontentsline{toc}{subsection}{Exercise V.4.9}

\begin{context}
Let \(X\) be a smooth cubic surface over an algebraically closed
field, presented as the blow-up of six plane points in general
position.  Write \(l\) for the pullback of a line and \(e_i\) for
the exceptional classes.  Then
\[
l^2=1,\quad e_i^2=-1,\quad l e_i=e_i e_j=0\ (i\ne j),
\qquad H=-K_X=3l-\sum_{i=1}^6e_i.
\]
The 27 lines have classes \(e_i\), \(l-e_i-e_j\), and
\(2l-\sum_{j\ne i}e_j\).  They are the only integral curves of
negative self-intersection.  The arithmetic genus of a divisor is
\(p_a(D)=1+(D^2-DH)/2\).

Exercise V.4.8 proves that every class nonnegative on all lines
can be re-marked to satisfy
\(a\geq b_1+b_2+b_3\), \(b_1\geq\cdots\geq b_6\geq0\);
such a class is nef.  A positive-square class with these line
inequalities contains a smooth integral member.
\end{context}

\begin{problem}
Let \(C\) be a smooth integral curve of degree \(d\) and positive
genus \(g\) on \(X\).  Show
\[
g\geq
\begin{cases}
(d-6)/2,&d\text{ even},\ d\geq8,\\
(d-5)/2,&d\text{ odd},\ d\geq13.
\end{cases}
\]
Prove that every indicated lower bound is attained.
\end{problem}
\end{exerciseintro}

\begin{solution}
The curve is not a line, so it is nonnegative on all lines.
Choose a standard marking
\[
C=al-\sum b_ie_i,\qquad
a\geq b_1+b_2+b_3,\quad b_1\geq\cdots\geq b_6\geq0.
\]
Put \(D=C-H\).  Adjunction gives
\[
D^2=2g-d+1.
\tag{1}
\]
We show that \(D^2\geq-5\), with the sole possible exception
\(C=3l\), which has degree nine.

Let \(r\) be the number of positive \(b_i\), and write
\[
D=D_0+\sum_{i>r}e_i,\qquad
D_0=(a-3)l-\sum_{i\leq r}(b_i-1)e_i.
\]
The terms in this decomposition are orthogonal, so
\(D^2=D_0^2-(6-r)\).
If \(r\geq3\), \(D_0\) is standard and nef; hence
\(D^2\geq-(6-r)\geq-3\).
If \(r=2\) and \(a\geq b_1+b_2+1\), it is again standard,
giving \(D^2\geq-4\).
In the remaining \(r=2\) case, \(a=b_1+b_2\).
The genus formula becomes \(g=(b_1-1)(b_2-1)>0\), so both
\(b_i\geq2\).  Then
\[
D_0=(l-e_1-e_2)
 +(b_1-2)(l-e_1)+(b_2-2)(l-e_2).
\]
The last two terms form a nef class orthogonal to \(l-e_1-e_2\).
Thus \(D_0^2\geq-1\), giving \(D^2\geq-5\).

If \(r=1\), positive genus forces \(a\geq b_1+2\), and
\[
D_0=(a-b_1-2)l+(b_1-1)(l-e_1)
\]
is nef.  Thus \(D^2\geq-5\).
If \(r=0\), positive genus forces \(a\geq3\), and
\[
D^2=(a-3)^2-6.
\]
This is at least \(-5\) except when \(a=3\), the asserted
degree-nine exception.

For even \(d\geq8\), (1) now gives \(g\geq(d-6)/2\).
For odd \(d\geq13\), the exception is absent and (1) is even.
Thus \(D^2\geq-4\), giving \(g\geq(d-5)/2\).

For equality, take the following classes:
\[
\begin{array}{c|c|c}
C&d&g\\ \hline
al-(a-2)e_1&2a+2&a-2\quad(a\geq3)\\
al-(a-2)e_1-e_2&2a+1&a-2\quad(a\geq6).
\end{array}
\]
They have squares \(4a-4\) and \(4a-5\), respectively,
and meet every one of the 27 lines nonnegatively.
Exercise V.4.8 supplies smooth integral members.
The table covers every degree in the requested ranges and
attains exactly the corresponding lower bound.
\end{solution}

\begin{exerciseintro}
% record: H-V-04-010
\exerciseheading{Exercise V.4.10}
\addcontentsline{toc}{subsection}{Exercise V.4.10}

\begin{context}
A convex function on a compact polytope achieves a maximum at a
vertex: on any line segment its value does not exceed the larger
endpoint value, and iterating on faces reduces to vertices.
A vertex in six-dimensional space has six linearly independent
active defining constraints.
\end{context}

\begin{problem}
Let \(a,b_1,\ldots,b_6\) be real numbers such that
\[
b_i>0,\qquad a>b_i+b_j\ (i\ne j),\qquad
2a>\sum_{i\ne j}b_i\quad\text{for each }j.
\]
Prove directly that \(a^2>\sum_i b_i^2\).
\end{problem}
\end{exerciseintro}

\begin{solution}
Since \(a>0\), divide all \(b_i\) by \(a\) and reorder them.
It suffices to consider \(a=1\) and
\(b_1\geq\cdots\geq b_6>0\).
The strongest pair and five-term inequalities are
\(b_1+b_2<1\) and \(b_1+\cdots+b_5<2\).
Consider their closed relaxation
\[
P=\{b_1\geq\cdots\geq b_6\geq0,\
b_1+b_2\leq1,\ b_1+\cdots+b_5\leq2\}.
\]
It is compact.  The convex function \(q(b)=\sum b_i^2\)
is bounded by its values at the vertices.

Here are all its vertices and their values:
\[
\begin{array}{c|c}
(b_1,b_2,b_3,b_4,b_5,b_6)&q\\ \hline
(0,0,0,0,0,0)&0\\
(1,0,0,0,0,0)&1\\
(1/2,1/2,0,0,0,0)&1/2\\
(1/2,1/2,1/2,0,0,0)&3/4\\
(1/2,1/2,1/2,1/2,0,0)&1\\
(2/5,2/5,2/5,2/5,2/5,0)&4/5\\
(2/5,2/5,2/5,2/5,2/5,2/5)&24/25\\
(2/3,1/3,1/3,1/3,1/3,0)&8/9\\
(2/3,1/3,1/3,1/3,1/3,1/3)&1\\
(1/2,1/2,1/3,1/3,1/3,0)&5/6\\
(1/2,1/2,1/3,1/3,1/3,1/3)&17/18\\
(1/2,1/2,1/2,1/4,1/4,0)&7/8\\
(1/2,1/2,1/2,1/4,1/4,1/4)&15/16 .
\end{array}
\]
Here is a direct two-block verification of completeness.
Set \(b_7=0\) and \(x_i=b_i-b_{i+1}\) for \(1\leq i\leq6\).
This invertible linear change writes the constraints as
\[
x_i\geq0,\qquad A\cdot x\leq1,\qquad B\cdot x\leq2,
\quad A=(1,2,2,2,2,2),\quad B=(1,2,3,4,5,5).
\]
At a vertex at most two \(x_i\) can be positive: with three
positive coordinates there are at most three active coordinate
constraints and two active sum constraints, fewer than six.
The zero vector is one vertex. With one positive coordinate,
its value is \(\min(1/A_i,2/B_i)\), giving the next six rows.
With two positive coordinates \(x_i,x_j\), both sum constraints
must be equalities. Their weighted average ratio is two, so
\(B_i/A_i<2<B_j/A_j\), after relabeling.
The ratios are \(1,1,3/2,2,5/2,5/2\); hence exactly
\[
i\in\{1,2,3\},\qquad j\in\{5,6\}
\]
are possible. Solving the two equations gives the final six rows.
Thus the list is complete.
Every displayed value is at most one, so \(q\leq1\) on \(P\).

For the original strict inequalities, choose \(\epsilon>0\)
small enough that \((1+\epsilon)b\) still satisfies both weak
sum inequalities.  It lies in \(P\), so
\[
(1+\epsilon)^2\sum b_i^2\leq1.
\]
Hence \(\sum b_i^2<1\).  Undoing normalization proves
\(a^2>\sum b_i^2\).
\end{solution}

\begin{exerciseintro}
% record: H-V-04-011
\exerciseheading{Exercise V.4.11}
\addcontentsline{toc}{subsection}{Exercise V.4.11}

\begin{context}
Let \(X\) be a smooth cubic surface over an algebraically closed
field, presented as the blow-up of six plane points in general
position.  Write \(l\) for the pullback of a line and \(e_i\) for
the exceptional classes.  Then
\[
l^2=1,\quad e_i^2=-1,\quad l e_i=e_i e_j=0\ (i\ne j),
\qquad H=-K_X=3l-\sum_{i=1}^6e_i.
\]
The 27 lines have classes \(e_i\), \(l-e_i-e_j\), and
\(2l-\sum_{j\ne i}e_j\).  They are the only integral curves of
negative self-intersection.  The arithmetic genus of a divisor is
\(p_a(D)=1+(D^2-DH)/2\).

The automorphism group \(G\) of the line configuration is generated
by permutations of the six exceptional classes and the quadratic
transformation at the first three marked points.

We use the elementary chamber theorem for a finite real reflection
group.  Reflections in the walls of a chamber have the presentation
\(s_i^2=1\), \((s_is_j)^{m_{ij}}=1\), where the angle between the
walls is \(\pi/m_{ij}\).  The stabilizer of a vector is generated
by the reflections whose walls contain that vector.
One proves the presentation by following paths through adjacent
chambers: a general deformation of a path changes its wall word
only by backtracking or by passing a codimension-two wall
intersection, giving precisely these relations.
The same argument in a small neighborhood of a vector gives the
stabilizer assertion.
\end{context}

\begin{problem}
For a graph, associate one involution to each vertex, imposing
\((st)^3=1\) for joined vertices and \((st)^2=1\) for unjoined
distinct vertices.

(a) For a chain of \(n-1\) vertices, prove that this group is
\(S_n\), under the map to adjacent transpositions.

(b) For a chain \(x_1,\ldots,x_5\) and a sixth vertex \(y\)
joined to \(x_3\), construct a surjection of the associated
group \(E_6\) onto \(G\), taking \(x_i\) to the adjacent
permutation of exceptional curves and \(y\) to the quadratic
transformation.

(c) Determine the order of \(E_6\), and prove this surjection
is an isomorphism.
\end{problem}
\end{exerciseintro}

\begin{solution}
(a) Let the chain generators be \(s_1,\ldots,s_{n-1}\).
The adjacent transpositions satisfy the relations and generate
\(S_n\), giving a surjection.  For the reverse order bound,
let \(H\) be generated by \(s_1,\ldots,s_{n-2}\), and put
\[
r_k=s_{n-1}s_{n-2}\cdots s_k\quad(1\leq k<n),
\qquad r_n=1.
\]
The union of the \(n\) right cosets \(Hr_k\) is stable under
right multiplication by every generator.
Indeed, \(r_ks_i=s_ir_k\) for \(i<k-1\);
\(r_ks_{k-1}=r_{k-1}\), \(r_ks_k=r_{k+1}\);
and the braid relation gives \(r_ks_i=s_{i-1}r_k\) for \(i>k\).
For \(r_n\) the assertion is immediate.
The union contains the identity, hence the whole group.
Induction gives \(|H|\leq(n-1)!\), so the group has at most
\(n!\) elements.  Its surjection to \(S_n\) is therefore an
isomorphism.

(b) On the real space \(H^\perp\), use the positive definite
form \(\langle u,v\rangle=-uv\), and set
\[
\alpha_i=e_i-e_{i+1}\quad(1\leq i\leq5),\qquad
\beta=l-e_1-e_2-e_3.
\]
Each has square \(-2\).  Its reflection on the Picard lattice is
\[
s_\rho(D)=D+(D\rho)\rho.
\tag{1}
\]
The first five reflections permute neighboring \(e_i\).
The last sends \(l\) to \(2l-e_1-e_2-e_3\) and \(e_i\) to
\(l-e_j-e_k\) for the first three indices, so it is the
quadratic transformation.  Direct pairings give the chain
relations and the single extra edge from \(\beta\) to \(\alpha_3\).
All these operations preserve the line configuration and generate
\(G\), proving the requested surjection.

(c) We compute the finite reflection group without assuming a
classification of root systems.
The integral classes \(\rho\) with \(\rho H=0\), \(\rho^2=-2\)
are exactly
\[
e_i-e_j\ (i\ne j),\qquad
\pm(l-e_i-e_j-e_k)\ (i<j<k),\qquad
\pm\theta,\quad \theta=2l-\sum e_i.
\tag{2}
\]
There are \(30+40+2=72\) of them.
To verify the list, write \(\rho=al-\sum b_ie_i\).
Then \(\sum b_i=3a\), \(\sum b_i^2=a^2+2\).
Cauchy gives \(|a|\leq2\).
For \(a=0\) there are two nonzero coefficients, \(1,-1\).
For \(a=1\), \(\sum(b_i^2-b_i)=0\) forces three ones.
For \(a=2\), equality in Cauchy forces all six coefficients
to be one.  Negation handles negative \(a\).

The six roots in (b) form a basis of \(H^\perp\).
Every root in (2) has coefficients all nonnegative or all
nonpositive in this basis.  For \(e_i-e_j\), \(i<j\), its
expression is \(\alpha_i+\cdots+\alpha_{j-1}\).
For \(l-e_i-e_j-e_k\), the coefficient of \(\beta\) is one
and that of \(\alpha_m\) is
\(\min(m,3)-|\{i,j,k\}\cap\{1,\ldots,m\}|\geq0\).
Finally
\[
\theta=\alpha_1+2\alpha_2+3\alpha_3+2\alpha_4+\alpha_5+2\beta.
\]
Consequently the inequalities
\(\langle v,\alpha_i\rangle>0,\langle v,\beta\rangle>0\)
define a chamber of the root hyperplanes.

Let \(W\) be generated by its wall reflections.
They preserve the finite spanning root set, so \(W\) is finite.
The chamber theorem gives exactly the \(E_6\) presentation:
\(E_6\cong W\).
The roots form a single \(W\)-orbit.  Indeed, for a positive
root \(\rho=\sum c_i\rho_i\),
\[
2=\langle\rho,\rho\rangle
=\sum c_i\langle\rho,\rho_i\rangle
\]
forces a positive pairing with a simple root.
Unless \(\rho\) is that simple root, the pairing is one,
and reflection subtracts that simple root, reducing height.
This reaches a simple root.  Adjacent simple roots are conjugate
under their two reflections, the diagram is connected, and
negation is achieved by reflection in the root itself.

The roots perpendicular to \(\theta\) are exactly \(e_i-e_j\):
for a root as above, \(\theta\rho=-a\).
The stabilizer theorem identifies \(W_\theta\) with the group
generated by their reflections, which is \(S_6\) of order \(720\).
Orbit--stabilizer gives
\[
|E_6|=|W|=72\cdot720=51840.
\]
Finally the action of \(W\) on the 27 line classes is faithful:
these classes span the Picard space, since they include all
\(e_i\) and \(l-e_1-e_2\).  Thus the already surjective
map \(W\to G\) is injective as well.  Hence \(E_6\cong G\).
\end{solution}

\begin{exerciseintro}
% record: H-V-04-012
\exerciseheading{Exercise V.4.12}
\addcontentsline{toc}{subsection}{Exercise V.4.12}

\begin{context}
Let \(X\) be a smooth cubic surface over an algebraically closed
field, presented as the blow-up of six plane points in general
position.  Write \(l\) for the pullback of a line and \(e_i\) for
the exceptional classes.  Then
\[
l^2=1,\quad e_i^2=-1,\quad l e_i=e_i e_j=0\ (i\ne j),
\qquad H=-K_X=3l-\sum_{i=1}^6e_i.
\]
The 27 lines have classes \(e_i\), \(l-e_i-e_j\), and
\(2l-\sum_{j\ne i}e_j\).  They are the only integral curves of
negative self-intersection.  The arithmetic genus of a divisor is
\(p_a(D)=1+(D^2-DH)/2\).

Every ample divisor on a smooth cubic surface is very ample
(Theorem V.4.11).  A general member is smooth and integral.
The rational surface \(X\) satisfies \(H^1(X,\mathcal O_X)=0\);
both \(X\) and an integral projective curve on it have only
constant global regular functions.
\end{context}

\begin{problem}
Prove that \(H^1(X,\mathcal O_X(-D))=0\) for every ample divisor
\(D\) on a smooth cubic surface.
\end{problem}
\end{exerciseintro}

\begin{solution}
Choose a smooth integral \(C\in|D|\).  Its ideal sequence is
\[
0\longrightarrow\mathcal O_X(-D)\longrightarrow\mathcal O_X
\longrightarrow\mathcal O_C\longrightarrow0.
\]
The resulting cohomology sequence contains
\[
H^0(X,\mathcal O_X)\longrightarrow H^0(C,\mathcal O_C)
\longrightarrow H^1(X,\mathcal O_X(-D))
\longrightarrow H^1(X,\mathcal O_X).
\]
The first map is the identity on \(k\), and the last group is
zero.  Exactness therefore forces the middle first-cohomology
group to vanish.  No restriction on the characteristic is used.
\end{solution}

\begin{exerciseintro}
% record: H-V-04-013
\exerciseheading{Exercise V.4.13}
\addcontentsline{toc}{subsection}{Exercise V.4.13}

\begin{context}
Let \(X\) be the blow-up of five plane points in general position
over an algebraically closed field, embedded in \(\mathbb P^4\)
by the complete anticanonical system
\(H=3l-e_1-\cdots-e_5\).
Thus \(H^2=4\), \(K_X=-H\), \(l^2=1\), \(e_i^2=-1\),
and all the mixed basis intersections are zero.
The surface is rational, so \(H^1(\mathcal O_X)=H^2(\mathcal O_X)=0\).
We use adjunction, Bertini, Riemann--Roch on a genus-one curve,
and the fact that a complete intersection of two hypersurfaces
is Cohen--Macaulay and has the product of their degrees.
\end{context}

\begin{problem}
(a) Prove that \(X\) contains exactly sixteen lines.
(b) Prove that it is the scheme-theoretic complete intersection
of two quadrics in \(\mathbb P^4\).
\end{problem}
\end{exerciseintro}

\begin{solution}
(a) The exceptional curves give five lines.
Any other line has class \(C=al-\sum b_ie_i\) with
\(a\geq1\), \(b_i\geq0\).  Since \(C\cong\mathbb P^1\)
and \(HC=1\), adjunction gives \(C^2=-1\).  Thus
\[
\sum b_i=3a-1,\qquad \sum b_i^2=a^2+1.
\]
Cauchy gives \((3a-1)^2\leq5(a^2+1)\), hence \(a\leq2\).
For \(a=1\), the equations force two coefficients one and the
others zero: these are the ten transforms of joining lines.
For \(a=2\), all five coefficients are one: this is the unique
conic through the five points.  General position makes these
curves irreducible and smooth; each has anticanonical degree one,
so each is a line in the embedding.
There are exactly \(5+10+1=16\).

(b) Choose a smooth hyperplane section \(E\in|H|\).
Adjunction gives genus one.  The restrictions of \(H\) and
\(2H\) to \(E\) have degrees four and eight, hence have respectively
four and eight sections and zero first cohomology.
The sequences for adding \(E\), together with
\(H^1(\mathcal O_X)=H^2(\mathcal O_X)=0\), give
\[
h^0(\mathcal O_X(H))=5,\quad H^1(\mathcal O_X(H))=0,\quad
h^0(\mathcal O_X(2H))=5+8=13.
\]
There are fifteen quadrics in \(\mathbb P^4\), so the kernel
of restriction to \(X\) has dimension at least two.
Choose independent quadrics \(Q_1,Q_2\) in that kernel.
Each is irreducible: otherwise the irreducible, nondegenerate
surface \(X\) would lie in one of its hyperplane factors.
Independence therefore ensures that the two quadrics have no
common hypersurface component.

Their intersection \(Y\) is a complete intersection of pure
dimension two and degree four, containing \(X\), also of degree
four.  Hence \(X\) is its sole component, with generic
multiplicity one.  Since \(Y\) is Cohen--Macaulay, it has no
embedded associated components; generic reducedness then makes
it reduced everywhere.  Thus \(Y=X\) scheme-theoretically.
\end{solution}

\begin{exerciseintro}
% record: H-V-04-014
\exerciseheading{Exercise V.4.14}
\addcontentsline{toc}{subsection}{Exercise V.4.14}

\begin{context}
Let \(X\) be a smooth cubic surface over an algebraically closed
field, presented as the blow-up of six plane points in general
position.  Write \(l\) for the pullback of a line and \(e_i\) for
the exceptional classes.  Then
\[
l^2=1,\quad e_i^2=-1,\quad l e_i=e_i e_j=0\ (i\ne j),
\qquad H=-K_X=3l-\sum_{i=1}^6e_i.
\]
The 27 lines have classes \(e_i\), \(l-e_i-e_j\), and
\(2l-\sum_{j\ne i}e_j\).  They are the only integral curves of
negative self-intersection.  The arithmetic genus of a divisor is
\(p_a(D)=1+(D^2-DH)/2\).

A class \(D=al-\sum b_ie_i\) is ample if
\[
b_i>0,\qquad a>b_i+b_j,\qquad
2a>\sum_{j\ne i}b_j
\]
for all indicated distinct indices.  An ample class on \(X\)
is very ample and has a smooth integral general member.
\end{context}

\begin{problem}
Construct smooth integral curves in \(\mathbb P^3\) with
\[
(d,g)=(8,6),(8,7),(9,7),(9,8),(9,9),
(10,8),(10,9),(10,10),(10,11).
\]
\end{problem}
\end{exerciseintro}

\begin{solution}
Use the following divisor classes on a fixed smooth cubic surface:
\[
\begin{array}{cc|c|c}
d&g&a&(b_1,b_2,b_3,b_4,b_5,b_6)\\ \hline
8&6&6&(2,2,2,2,1,1)\\
8&7&7&(3,2,2,2,2,2)\\
9&7&6&(2,2,2,1,1,1)\\
9&8&7&(3,2,2,2,2,1)\\
9&9&7&(2,2,2,2,2,2)\\
10&8&6&(2,2,1,1,1,1)\\
10&9&7&(3,2,2,2,1,1)\\
10&10&7&(2,2,2,2,2,1)\\
10&11&8&(3,3,2,2,2,2).
\end{array}
\]
The coefficients are sorted.  Thus to check ampleness it suffices
to check positivity of \(b_6\), \(a-b_1-b_2\), and
\(2a-\sum b_i+b_6\), which holds in each row.
Each class consequently contains a smooth integral curve.
Its degree and genus in the cubic embedding are
\[
d=3a-\sum b_i,\qquad
g=1+\frac{a^2-\sum b_i^2-d}{2}.
\]
Substituting the entries gives the first two columns exactly.
These are therefore the requested curves in \(\mathbb P^3\).
\end{solution}

\begin{exerciseintro}
% record: H-V-04-015
\exerciseheading{Exercise V.4.15}
\addcontentsline{toc}{subsection}{Exercise V.4.15}

\begin{context}
Work over an algebraically closed field \(k\). A configuration consists of
distinct ordinary points of \(\mathbb P^2\). An admissible quadratic
transformation has three noncollinear configuration points as centers;
replace those centers by the three inverse centers and carry the other
points along. Call a configuration Cremona general if every finite
sequence of these operations leaves distinct ordinary points with no
three collinear.

On the blow-up write \(H,E_1,\ldots,E_r\) for the line and exceptional
classes. Their intersection matrix is \(\operatorname{diag}(1,-1,\ldots,-1)\),
and \(K=-3H+\sum E_i\). Adjunction gives
\(2p_a(C)-2=C^2+KC\). Distinct integral curves have nonnegative
intersection. An integral projective curve of arithmetic genus zero
over \(k\) is a smooth rational curve.

An admissible quadratic transformation lifts to an isomorphism between
the surfaces obtained by blowing up the two full configurations.
For the first three centers its action on divisor coordinates is
\[
d'=2d-m_1-m_2-m_3,\qquad
m'_1=d-m_2-m_3,\quad
m'_2=d-m_1-m_3,\quad
m'_3=d-m_1-m_2,
\]
with the other multiplicities unchanged. In particular this is the
reflection \(s_\gamma(D)=D+(D\gamma)\gamma\), where
\(\gamma=H-E_1-E_2-E_3\). Permuting exceptional labels is also allowed.
Polynomial equations in point coordinates describe incidence and
nonvanishing conditions for every fixed finite sequence.
The product \((\mathbb P^2)^m\) is irreducible, and a birational
map restricts to an isomorphism between dense open subsets.
Noether normalization gives a finite surjective morphism from an
affine integral variety of dimension \(n\) to \(\mathbb A^n\).
We may use Bezout's theorem for plane curves. Vanishing and prescribed
first-order vanishing at fixed points are linear conditions on
homogeneous forms; degree-\(d\) forms form a vector space of dimension
\(\binom{d+2}{2}\).
\end{context}

\begin{problem}
(a) For six points, prove that Cremona general position is equivalent
to having no three collinear and not having all six on a conic.

(b) Prove that an admissible sequence preserves Cremona general position.

(c) If \(k\) is uncountable and \(P_1,\ldots,P_r\) are fixed in
Cremona general position, prove that a dense subset of choices
of \(P_{r+1}\) extends this configuration to one in Cremona general
position. Prove the countable-union lemma suggested by this assertion.

(d) For \(r=7,8\), count the smooth rational curves of square \(-1\)
on the blow-up, obtaining respectively \(56,240\), and show that
there are no other negative integral curves.

(e) For \(r=9\), prove that there are infinitely many such curves.
\end{problem}
\end{exerciseintro}

\begin{solution}
(a) Normalize the centers to the coordinate vertices. The quadratic
map is \([x:y:z]\mapsto[yz:xz:xy]\). A conic through the three
vertices has equation \(axy+bxz+cyz=0\); its transform is a line.
Thus six points on a conic would give three collinear transformed
points.

Conversely, assume the two stated conditions. The three points
not used as centers have all coordinates nonzero. A line through
two inverse centers therefore cannot contain any of their images.
A line through one inverse center and two of their images pulls
back, after removing a coordinate-line factor, to a line through
one original center and the two corresponding original points.
A line through the three images pulls back to a conic through all
six original points. Each possibility is excluded. Finally a conic
through the three inverse centers and the three images pulls back
to \(xyz\) times a linear equation through the three original
noncenters, also impossible. The transformed configuration satisfies
the same two conditions. Induction proves the assertion for every
finite sequence.

(b) Concatenate any further sequence with the sequence already
performed. The defining condition for the original configuration
then gives the defining condition for the resulting one.

(c) \emph{The countable-union lemma.}
First prove the lemma, including density of
the complement. A countable family of proper closed subsets of
\(\mathbb A^1(k)\) is countable and cannot cover an uncountable field.
Inductively, for proper closed sets \(Z_j\subset\mathbb A^n\),
choose a nonzero polynomial \(f_j\) vanishing on each \(Z_j\).
View \(f_j\) as a polynomial in the last coordinate. Its coefficients
do not vanish simultaneously outside a proper closed subset of
\(\mathbb A^{n-1}\). Choose a point outside all these exceptional
sets by induction. On the resulting vertical line each \(f_j\)
has only finitely many zeros, so a point avoiding all of them exists.
Adding one further proper closed set to the family proves density.
For an arbitrary integral variety, restrict to a nonempty affine
open and use a finite surjective Noether-normalization map.
The image of each proper closed subset is closed of smaller
dimension. The affine-space case applies to these images.
The same argument on every nonempty affine open proves density.

\emph{Reduction with the old configuration fixed.}
Here is the precise remaining reduction for the extension assertion.
Hold \(P_1,\ldots,P_r\) fixed and let \(q\) vary in \(\mathbb P^2\).
For each finite word \(w\) of center triples, let \(U_w\) be the
set of \(q\) for which all successive configurations are defined,
ordinary, distinct, and have no collinear triple. Coordinate
formulas and nonvanishing determinants make \(U_w\) open.
There are only countably many words, and the desired set is
\[
V=\bigcap_w U_w.
\]
The lemma proves density if every \(U_w\) is nonempty. For words
whose centers never involve the additional label, this nonemptiness
follows by pulling back the finitely many forbidden lines under
fixed birational maps. For arbitrary words involving the additional
label, nonemptiness on the fixed-configuration fiber is an additional
condition that this argument does not prove. Nonemptiness of the
corresponding open subset of \((\mathbb P^2)^{r+1}\) does not suffice:
an open subset of a product can miss an entire specified fiber.
\begin{center}
\fbox{\begin{minipage}{0.9\linewidth}
\emph{Remaining gap in (c).}
For every fixed Cremona general configuration
\(P_1,\ldots,P_r\), one must prove \(U_w\ne\varnothing\)
for every word \(w\), including words whose centers involve
the new point. This fixed-fiber nonemptiness has not been
established here, so the original extension assertion remains unproved.
\end{minipage}}
\end{center}

\emph{A proved weaker statement: a very general full configuration.}
If all \(m=r+1\) points are allowed to vary, the Cremona general
configurations form a dense subset of \(S=(\mathbb P^2)^m\).
For \(m<3\) this follows just by excluding coincident points.
For \(m\geq3\), let \(S^\circ\) denote the nonempty open set of
distinct configurations with no collinear triple.

For each triple of labels, we construct a birational involution
of \(S\) representing the corresponding quadratic operation.
On a dense coordinate chart, choose normalized column vectors
for its three noncollinear centers and let \(A\) be their
invertible matrix. Conjugate
\(\sigma([x:y:z])=[yz:xz:xy]\) by \(A\).
Keep the three center labels at their original points: they
are precisely the inverse centers for \(A\sigma A^{-1}\).
Apply this map to all other points. On the dense open set
where those points avoid the three joining lines, the
construction is defined. Its second application uses the
same matrix \(A\), since the centers were retained, and
\(\sigma^2\) is the identity as a rational map.
Thus this operation is a birational involution of the full
parameter space.

Fix a finite word \(w\). Every prefix gives a birational
map of \(S\). Requiring each operation to be defined and
each prefix image to belong to \(S^\circ\) gives a dense
open subset \(G_w\): it is a finite intersection of domains
of definition and inverse images of dense opens under
birational maps. There are countably many words. The
countable-union lemma therefore makes
\(\bigcap_w G_w\) dense in \(S\), and every configuration
in this intersection is Cremona general.
This proves the weaker whole-configuration statement.
It supplies no nonemptiness assertion on an arbitrarily
specified fiber with the first \(r\) points fixed.

(d) Every subset of a Cremona general configuration is Cremona
general: a word on its labels can be performed on the full
configuration. In particular (a) excludes six points on a conic.
For eight points, a cubic through all of them singular at one
would have an effective divisor of class
\[
3H-2E_1-E_2-\cdots-E_8.
\]
Apply a quadratic transformation at labels \(1,2,3\).
The new class is that of a conic through six points.
A transformation at three of those six gives a line through
three points, contradicting general position. This argument uses
effectivity of divisor classes under the lifted isomorphisms,
and also covers multiplicities greater than the specified minima.

Put \(A=-K\). We first exhibit an integral member of \(|A|\).
For eight points, cubic equations through them form a nonzero
vector space. A reducible cubic is a line plus a conic, or three
lines; the restrictions on collinearity and conics allow it to
contain at most seven of the points. Such a cubic is therefore
integral. The preceding paragraph excludes singularities at the
eight points, so its strict transform has class \(A\).
For seven points, the projective space of cubics through them
has dimension at least two. A reducible member must be the union
of a line through two of the points and the conic through the
other five. There are only finitely many such members.
Five points with no three collinear determine a unique smooth
conic, by the elementary conic equations and Bezout.
For each specified point, one of these line-plus-conic members
is smooth there: put that point on the line, since the conic
cannot contain a sixth point. Thus singularity at any one of the
seven points is a proper linear condition. An integral member
smooth at all seven exists, again giving an integral divisor
of class \(A\).

Since \(A^2=9-r>0\), intersection with this member shows that
\(AC\geq0\) for every integral curve \(C\), including the member
itself. If \(C^2<0\), adjunction gives
\[
0\leq p_a(C)=1+\frac{C^2-AC}{2}.
\]
Consequently either \(C^2=-1,\ AC=1,\ p_a(C)=0\), or
\(C^2=-2,\ AC=0,\ p_a(C)=0\).
The latter case cannot be exceptional. Write its plane class as
\(dH-\sum m_iE_i\), with \(d>0,\ m_i\geq0\). Then
\[
\sum m_i=3d,\qquad \sum m_i^2=d^2+2,\qquad
(9-r)d^2\leq2r.
\]
For seven points \(d\leq2\), yielding only a line through three
points or a conic through six. For eight points \(d\leq4\);
equality \(d=4\) in Cauchy--Schwarz would require all \(m_i=3/2\).
The only additional case is \(d=3\), with multiplicities
\((2,1,1,1,1,1,1,1)\). These are precisely the excluded cases.

For a nonexceptional \((-1)\)-curve the equations instead are
\[
\sum m_i=3d-1,\qquad \sum m_i^2=d^2+1.
\]
Cauchy--Schwarz bounds \(d\leq3\) for seven points and \(d\leq7\)
for eight. In the second case \(d=7\) would require every
\(m_i=5/2\), so actually \(d\leq6\).
Enumerating nonnegative integer partitions of \(3d-1\) subject to
\(\sum m_i(m_i-1)=(d-1)(d-2)\) gives the following table.
Exponents indicate repeated entries; omitted entries are zero.
\[
\begin{array}{c|c|r|r}
\text{class or degree}&\text{multiplicities}&r=7&r=8\\ \hline
E_i& &7&8\\
1&1^2&21&28\\
2&1^5&21&56\\
3&2,1^6&7&56\\
4&2^3,1^5&0&56\\
5&2^6,1^2&0&28\\
6&3,2^7&0&8
\end{array}
\]
All listed classes occur. For each positive-degree row, perform a
quadratic transformation at the three largest multiplicities,
using a zero entry as the third center in degree one.
The successive degree reductions are \(1\) to an exceptional class,
\(2\) to \(1\), \(3\) to \(2\), \(4\) to \(2\), \(5\) to \(4\),
and \(6\) to \(5\), with the lower-degree patterns in the table.
General position makes every transformation admissible.
Pulling back the final exceptional curve gives a smooth rational
curve in the original class. Two different integral curves in
one such class would have intersection \(-1\), impossible.
The column sums are \(56\) and \(240\), as required.

(e) On the nine-point blow-up set
\[
\delta=3H-\sum_{i=1}^9E_i,\qquad \alpha=E_1-E_2.
\]
Thus \(\delta^2=\delta\alpha=0\) and \(\alpha^2=-2\).
The reflection \(s_\alpha\) exchanges labels \(1,2\).
The class \(\beta=\delta+\alpha\) has degree three and
multiplicities \((0,2,1,1,1,1,1,1,1)\).
Reflecting at its double point and two simple points reduces
it to a conic-through-six class, then to a line-through-three
class. Hence \(s_\beta\) is a conjugate of a quadratic reflection
by a word of permutations and quadratic reflections.

Set \(T=s_\beta s_\alpha\). Direct substitution gives
\[
T(D)=D+(D\delta)\alpha+(D\delta-D\alpha)\delta.
\]
Since \(E_9\delta=1,\ E_9\alpha=0\), induction yields
\[
T^n(E_9)=E_9+n\alpha+n^2\delta\qquad(n\geq1).
\]
Every word \(T^n\) can be realized by admissible transformations,
by general position. The lifted isomorphisms carry an exceptional
curve to a smooth rational \((-1)\)-curve in this class.
Its plane degree is \(3n^2\), and its multiplicities are
\[
(n^2-n,\ n^2+n,\ n^2,n^2,n^2,n^2,n^2,n^2,\ n^2-1).
\]
In particular these curves are distinct and their degrees are
unbounded. The line through two original points is also in the
reflection orbit of an exceptional class, as the degree-one
reduction in (d) shows, so this supplies the orbit required by
the source's hint.
\end{solution}

\begin{exerciseintro}
% record: H-V-04-016
\exerciseheading{Exercise V.4.16}
\addcontentsline{toc}{subsection}{Exercise V.4.16}

\begin{context}
Work over an algebraically closed field \(k\) with
\(\operatorname{char}k\ne3\).
A smooth cubic surface has 27 lines, and its anticanonical
embedding is its given embedding in \(\mathbb P^3\).
Thus every automorphism extends to a projective linear map.
In characteristic two, we use Hermitian linear algebra over
\(\mathbb F_4\), with conjugation \(u\mapsto u^2\).
A nondegenerate Hermitian space has an orthonormal basis.
\end{context}

\begin{problem}
For \(X=V(x_0^3+x_1^3+x_2^3+x_3^3)\), find all its lines
and their incidences.  Determine \(\operatorname{Aut}_k(X)\),
including the exceptional characteristic-two case.
Explain the exclusion of characteristic three.
\end{problem}
\end{exerciseintro}

\begin{solution}
Let every parameter below run through the three roots of \(-1\).
The following three sets give 27 distinct lines:
\[
\begin{array}{ll}
A(\zeta,\eta):&x_0=\zeta x_1,\quad x_2=\eta x_3,\\
B(\alpha,\beta):&x_0=\alpha x_2,\quad x_1=\beta x_3,\\
C(\gamma,\delta):&x_0=\gamma x_3,\quad x_1=\delta x_2.
\end{array}
\tag{1}
\]
Their equations make the cubic vanish identically.
They are distinct within each family; lines from different
families have different two-dimensional coordinate solution
spaces.  Since \(X\) is smooth and has 27 lines, the list is complete.

Within a family, two distinct lines meet exactly when one
parameter agrees.  Between families, substitution in (1) gives
\[
\begin{array}{c|c}
\text{pair}&\text{condition for intersection}\\ \hline
A(\zeta,\eta),B(\alpha,\beta)&\zeta\beta=\alpha\eta\\
A(\zeta,\eta),C(\gamma,\delta)&\gamma=\zeta\eta\delta\\
B(\alpha,\beta),C(\gamma,\delta)&\alpha\beta=\gamma\delta .
\end{array}
\tag{2}
\]
When the indicated equality holds the common solution space is
one-dimensional; otherwise it is zero.  Thus every line meets
four lines of its own family and three from each other family:
ten in total.  Equations (1)--(2) specify every pairwise incidence.

Suppose first that \(\operatorname{char}k\ne2,3\).
The Hessian determinant of \(F=\sum x_i^3\) is
\[
6^4 x_0x_1x_2x_3.
\]
The chain rule for a linear change of coordinates shows that
an automorphism preserving \(F\) up to scalar preserves this
zero set.  It must therefore permute the four coordinate
hyperplanes, and its matrix is monomial.
Conversely a permutation of coordinates preserves \(F\).
A diagonal matrix preserves \(F\) up to scalar exactly when
the cubes of its diagonal entries agree.
Modulo scalar matrices this gives
\[
\operatorname{Aut}(X)\cong\mu_3^3\rtimes S_4,
\qquad |\operatorname{Aut}(X)|=27\cdot24=648.
\tag{3}
\]

Now let \(\operatorname{char}k=2\).
For a matrix \(A\), expansion of \(F(Ax)\) shows that preservation
up to scalar is exactly
\[
A^{(2)T}A=\lambda I,
\tag{4}
\]
where \((2)\) squares every entry.  Rescale \(A\) to make
\(\lambda=1\).  Taking squared transpose of (4) gives
\(A^{(2)T}A^{(4)}=I\), so comparison with (4) yields
\(A^{(4)}=A\).
Every entry is therefore in \(\mathbb F_4\).
Conversely each unitary matrix over \(\mathbb F_4\) satisfies
(4) and preserves \(X\).
Two normalized matrices give the same projective map precisely
when they differ by a scalar in \(\mu_3\).
Consequently
\[
\operatorname{Aut}(X)=U_4(2)/\mu_3=\operatorname{PGU}_4(2).
\tag{5}
\]

For completeness, a vector in the standard Hermitian
\(\mathbb F_4^n\) has norm one exactly when it has an odd
number of nonzero coordinates.  There are therefore
\[
N_n=\sum_{j\text{ odd}}\binom nj3^j
=\frac{4^n-(-2)^n}{2}
\]
such vectors, where the counting takes place in the integers.
Choosing orthonormal basis vectors successively gives
\[
|U_4(2)|=N_4N_3N_2N_1=120\cdot36\cdot6\cdot3=77760.
\]
Dividing by the three scalar matrices gives
\(|\operatorname{Aut}(X)|=25920\).
These are automorphisms over \(k\); no field-semilinear maps
are included.

Finally, in characteristic three,
\(F=(x_0+x_1+x_2+x_3)^3\).
It defines a triple plane, so the smooth cubic-surface question
and its 27-line description do not apply.
\end{solution}

\begin{exerciseintro}
\corpussection{V.5 Birational Transformations}
\addcontentsline{toc}{section}{V.5 Birational Transformations}

% record: H-V-05-001
\exerciseheading{Exercise V.5.1}
\addcontentsline{toc}{subsection}{Exercise V.5.1}

\begin{context}
Work over an algebraically closed field \(k\). A surface is smooth,
integral, and projective over \(k\). For the blow-up
\(\pi:\widetilde X\to X\) at a closed point, with exceptional curve \(E\),
\[
E\cong\mathbb P^1,\quad E^2=-1,\quad
(\pi^*D)E=0,\quad(\pi^*D)(\pi^*D')=DD'.
\]
If an effective divisor \(D\) has multiplicity \(m\) at the center,
its strict transform has class \(\pi^*D-mE\).
Intersection multiplicities of effective divisors with no common
component are nonnegative.

Regular local rings of a smooth surface are factorial. Thus a rational
function with no polar component through a point is regular there.
\end{context}

\begin{problem}
For \(f\in k(X)\), construct a finite sequence of point blow-ups
\(g:X'\to X\) such that the rational map defined by \(f\) extends to a
morphism \(X'\to\mathbb P^1\).
\end{problem}
\end{exerciseintro}

\begin{solution}
The zero function and all constant functions already define morphisms.
Otherwise write
\[
\operatorname{div}(f)=Z-P
\]
with effective divisors \(Z,P\) having no common component.
Put \(N=ZP\), a nonnegative integer.
If their supports meet at \(x\), let \(m,n>0\) be their respective
multiplicities at \(x\), and blow up \(x\).
Write \(\widetilde Z,\widetilde P\) for their strict transforms.
The new divisor of \(f\) is
\[
\widetilde Z-\widetilde P+(m-n)E.
\]
When \(m\geq n\), the new zero and pole divisors are
\(\widetilde Z+(m-n)E\) and \(\widetilde P\). Their intersection is
\[
(ZP-mn)+(m-n)n=ZP-n^2.
\]
When \(n\geq m\), the analogous calculation gives \(ZP-m^2\).
In either case it decreases by \(\min(m,n)^2>0\).
The new zero and pole divisors again have no common component, so their
intersection remains nonnegative. Consequently this operation can occur
only finitely often.

At termination the supports of the zeros and poles are disjoint.
Off the poles, \(f\) is regular; off the zeros, \(f^{-1}\) is regular.
These two opens cover \(X'\). On them define the map to \(\mathbb P^1\)
by \([f:1]\) and \([1:f^{-1}]\), respectively. On the overlap \(f\) is
a unit and the expressions agree. They therefore glue to the required
morphism.
\end{solution}

\begin{exerciseintro}
% record: H-V-05-002
\exerciseheading{Exercise V.5.2}
\addcontentsline{toc}{subsection}{Exercise V.5.2}

\begin{context}
Work over an algebraically closed field \(k\). A surface is smooth,
integral, and projective over \(k\). For the blow-up
\(\pi:\widetilde X\to X\) at a closed point, with exceptional curve \(E\),
\[
E\cong\mathbb P^1,\quad E^2=-1,\quad
(\pi^*D)E=0,\quad(\pi^*D)(\pi^*D')=DD'.
\]
If an effective divisor \(D\) has multiplicity \(m\) at the center,
its strict transform has class \(\pi^*D-mE\).
Intersection multiplicities of effective divisors with no common
component are nonnegative.

We use Serre vanishing, the exact sequence for adding an effective
Cartier divisor, and \(H^1(\mathbb P^1,\mathcal O(d))=0\) for \(d\geq-1\).
A globally generated line bundle defines a morphism to projective space.
The normalization of a projective variety is finite and projective;
a dominant morphism from a normal variety factors through it.
\end{context}

\begin{problem}
Let \(Y\subset X\) be a curve isomorphic to \(\mathbb P^1\), with
\(Y^2<0\). Construct a birational morphism to a normal projective
surface \(f:X\to X_0\) which contracts \(Y\) to one point and is an
isomorphism away from \(Y\).
\end{problem}
\end{exerciseintro}

\begin{solution}
Put \(n=-Y^2>0\). Choose a sufficiently large multiple \(H\) of a
very ample divisor so that
\[
H^1(X,\mathcal O_X(H))=0,\qquad HY=kn
\]
for a positive integer \(k\). Divisibility is obtained by taking the
multiple divisible by \(n\), and vanishing by Serre's theorem.
For \(1\leq j\leq k\), the restriction of \(H+jY\) to \(Y\) has
degree \((k-j)n\geq0\). The sequences
\[
0\to\mathcal O_X(H+(j-1)Y)\to\mathcal O_X(H+jY)
\to\mathcal O_Y((H+jY)|_Y)\to0
\]
therefore show inductively that
\(H^1(X,\mathcal O_X(H+jY))=0\).

Set \(A=H+kY\). Multiplication by the canonical section of
\(\mathcal O_X(kY)\) supplies enough sections of \(A\) to generate it
away from \(Y\). At \(Y\), the last displayed sequence gives a
surjection
\[
H^0(X,\mathcal O_X(A))\longrightarrow
H^0(Y,\mathcal O_Y(A))=H^0(Y,\mathcal O_Y)=k,
\]
since the preceding \(H^1\) vanishes. A section lifting \(1\) generates
along \(Y\). Thus \(A\) is globally generated.

Let \(h:X\to X_1\subset\mathbb P^N\) be its morphism onto its image.
Every section restricts to a constant on \(Y\), so \(h(Y)\) is a point
\(p\). The subsystem obtained from \(H\) vanishes on \(Y\), while at
each point outside \(Y\) some member is nonzero. Thus \(h^{-1}(p)=Y\)
set-theoretically.
On a chart where such a member is nonzero, ratios of members of this
subsystem are exactly the ratios giving the very ample embedding by
\(H\). They recover both points and local coordinates. Hence
\(h:X\setminus Y\to X_1\setminus\{p\}\) is an isomorphism, and \(h\)
is birational.

Normalize \(X_1\), obtaining \(X_0\), and factor \(h\) as
\(X\xrightarrow{f}X_0\to X_1\).
The factor \(f\) is proper, since \(X\) is projective and \(X_0\)
is separated, and its image is dense by birationality.
Its closed image is therefore all of \(X_0\).
The finite set over \(p\) has only one point: the surjective map \(f\)
sends the connected curve \(Y\) onto that set. Off this point the
normalization is an isomorphism because \(X_1\setminus\{p\}\)
is smooth. Thus \(f\) has exactly the required properties.
The contracted point is allowed to be singular.
\end{solution}

\begin{exerciseintro}
% record: H-V-05-003
\exerciseheading{Exercise V.5.3}
\addcontentsline{toc}{subsection}{Exercise V.5.3}

\begin{context}
Work over an algebraically closed field \(k\). A surface is smooth,
integral, and projective over \(k\). For the blow-up
\(\pi:\widetilde X\to X\) at a closed point, with exceptional curve \(E\),
\[
E\cong\mathbb P^1,\quad E^2=-1,\quad
(\pi^*D)E=0,\quad(\pi^*D)(\pi^*D')=DD'.
\]
If an effective divisor \(D\) has multiplicity \(m\) at the center,
its strict transform has class \(\pi^*D-mE\).
Intersection multiplicities of effective divisors with no common
component are nonnegative.

Write \(\Omega_X=\Omega^1_{X/k}\). For a point blow-up,
\(\pi_*\mathcal O_{\widetilde X}=\mathcal O_X\) and
\(R^q\pi_*\mathcal O_{\widetilde X}=0\) for \(q>0\).
Projection formula and Leray then give
\(H^q(\widetilde X,\pi^*F)=H^q(X,F)\) for locally free \(F\).
Regular sections of a locally free sheaf on a smooth surface extend
across a missing closed point.
Serre duality gives
\(H^2(X,\Omega_X)^\vee=H^0(X,\Omega_X^\vee\otimes\omega_X)\);
the wedge pairing identifies \(\Omega_X^\vee\otimes\omega_X\)
with \(\Omega_X\).
We use \(H^0(\mathbb P^1,\Omega_{\mathbb P^1})=0\) and
\(H^1(\mathbb P^1,\Omega_{\mathbb P^1})=k\).
The class \(c_1(L)\in H^1(\Omega_X)\) is represented by logarithmic
differentials of transition functions; on a smooth projective curve
its trace is \(\deg L\), viewed in \(k\).
\end{context}

\begin{problem}
For the blow-up of a surface at one point, prove
\[
H^1(\widetilde X,\Omega_{\widetilde X})
\cong H^1(X,\Omega_X)\oplus k.
\]
Deduce that every sequence of contractions of exceptional curves
of the first kind has bounded length.
\end{problem}
\end{exerciseintro}

\begin{solution}
Let \(i:E\hookrightarrow\widetilde X\). There is an exact sequence
\[
0\to\pi^*\Omega_X\to\Omega_{\widetilde X}
\to i_*\Omega_E\to0.
\tag{1}
\]
To check it, take local smooth coordinates \(x,y\) at the center.
On the chart \(y=xt\), the pullbacks of \(dx,dy\) are
\(dx,t\,dx+x\,dt\). The map is injective, and its cokernel is
\((\mathcal O/(x))dt\), the differentials on that chart of \(E\).
The other chart gives the same quotient, with the usual transition
for \(\Omega_{\mathbb P^1}\). These calculations can be made after
an etale coordinate map, which suffices to verify (1).

The cohomology sequence and projection formula give
\[
\begin{split}
0\to H^1(X,\Omega_X)&\to H^1(\widetilde X,\Omega_{\widetilde X})
\to k\\
&\to H^2(X,\Omega_X)
\to H^2(\widetilde X,\Omega_{\widetilde X})\to0.
\end{split}
\tag{2}
\]
The initial zero follows from \(H^0(E,\Omega_E)=0\), and the final
zero from the absence of second cohomology on \(E\).
Moreover pullback identifies the spaces of global one-forms on
\(X\) and \(\widetilde X\): a form upstairs restricts to
\(X\setminus\{p\}\), extends across \(p\), and pulls back to the
original form because they agree on a dense open.
Serre duality now shows that the last two nonzero terms of (2)
have the same finite dimension. The displayed surjection between
them is therefore an isomorphism. Hence the preceding boundary map
from \(k\) is zero, and (2) reduces to a short exact sequence.

It splits explicitly: \(c_1(\mathcal O_{\widetilde X}(E))\)
restricts to the class of \(\mathcal O_E(-1)\), whose trace is
\(-1\ne0\), in every characteristic. Its span supplies a complement
to the pullback of \(H^1(X,\Omega_X)\).
This proves the asserted decomposition.

Each contraction to a smooth surface reverses one point blow-up
and lowers \(h^1(\Omega)\) by one. Thus a sequence starting from
\(X\) has length at most \(h^1(X,\Omega_X)\), since the final
dimension cannot be negative.
\end{solution}

\begin{exerciseintro}
% record: H-V-05-004
\exerciseheading{Exercise V.5.4}
\addcontentsline{toc}{subsection}{Exercise V.5.4}

\begin{context}
Work over an algebraically closed field \(k\). A surface is smooth,
integral, and projective over \(k\). For the blow-up
\(\pi:\widetilde X\to X\) at a closed point, with exceptional curve \(E\),
\[
E\cong\mathbb P^1,\quad E^2=-1,\quad
(\pi^*D)E=0,\quad(\pi^*D)(\pi^*D')=DD'.
\]
If an effective divisor \(D\) has multiplicity \(m\) at the center,
its strict transform has class \(\pi^*D-mE\).
Intersection multiplicities of effective divisors with no common
component are nonnegative.

Every birational morphism between smooth projective surfaces factors
as finitely many point blow-ups. Blowing up a smooth point of a
smooth curve leaves its strict transform isomorphic to that curve,
and decreases its self-intersection by one.
\end{context}

\begin{problem}
Let \(f:X\to X'\) be a birational morphism of smooth projective surfaces.

(a) If an integral curve \(Y\subset X\) is contracted by \(f\), prove
that \(Y\cong\mathbb P^1\) and \(Y^2<0\).

(b) For a point \(p\in X'\), let \(Y_1,\ldots,Y_r\) be all the
curve components of its fiber. Prove that their intersection
matrix is negative definite when \(r>0\).
\end{problem}
\end{exerciseintro}

\begin{solution}
(a) Factor \(f\) into point blow-ups, starting from \(X'\). Each new
exceptional curve is a copy of \(\mathbb P^1\) of square \(-1\).
Subsequent blow-ups preserve the isomorphism type of its strict
transform, and decrease its square only when their centers lie on it.
Every contracted curve is one of these strict transforms: away from
their union the composite is an isomorphism. This proves (a).

(b) Let \(E_1,\ldots,E_s\) be the total transforms on \(X\) of
the successive new exceptional curves, ordered by their creation.
The blow-up intersection formulas imply
\[
E_iE_j=0\quad(i\ne j),\qquad E_i^2=-1.
\]
Thus their real span has negative definite form, with
\((\sum a_iE_i)^2=-\sum a_i^2\).
The strict transform of the \(i\)-th exceptional curve equals
\(E_i\) minus an integer combination of \(E_j\) with \(j>i\).
Indeed each later blow-up subtracts its new exceptional class
exactly when its center lies on the then-current strict transform.
Consequently the change from total to strict transforms is
triangular with diagonal entries one.
All strict exceptional classes are therefore linearly independent
and span the same negative definite space.

The curves \(Y_1,\ldots,Y_r\) form a subset of these strict transforms.
For any nonzero real vector \((a_1,\ldots,a_r)\), their linear
independence makes \(D=\sum a_iY_i\) a nonzero vector in that space,
so
\[
\sum_{i,j}a_i a_j(Y_iY_j)=D^2<0.
\]
This is precisely negative definiteness of the requested matrix.
\end{solution}

\begin{exerciseintro}
% record: H-V-05-005
\exerciseheading{Exercise V.5.5}
\addcontentsline{toc}{subsection}{Exercise V.5.5}

\begin{context}
Let \(C\) be a smooth projective integral curve over an algebraically
closed field. A geometrically ruled surface over \(C\) is
\(\mathbb P_C(F)\) for a rank-two vector bundle \(F\), using the convention
that points are one-dimensional quotients. An elementary transformation
at a point of a fiber blows up that point and contracts the strict
transform of the fiber. A smooth rational curve of square \(-1\)
can be contracted to a smooth point.
Local rings of \(C\) at closed points are discrete valuation rings.
Every torsion-free coherent sheaf on \(C\) is locally free, and a
torsion sheaf has finite length and a filtration with quotients \(k(p)\).
\end{context}

\begin{problem}
Show that any two geometrically ruled surfaces over the same curve
\(C\) are related, over \(C\), by a finite sequence of elementary
transformations.
\end{problem}
\end{exerciseintro}

\begin{solution}
First consider a surjection \(F\to k(p)\) and its kernel \(F_1\).
The kernel is a rank-two vector bundle. Locally at \(p\), choose a
parameter \(t\) and a basis \(e_1,e_2\) so that the quotient sends
\(e_1\) to zero and \(e_2\) to one. Then
\[
F_1=\mathcal O e_1\oplus\mathcal O(t e_2).
\]
Restriction of a quotient of \(F\) to \(F_1\) gives the rational map
\[
(t,[u:v])\longmapsto(t,[u:tv])
\]
from \(\mathbb P(F)\) to \(\mathbb P(F_1)\).
Its unique base point is \(t=u=0\), the chosen quotient point.
On the chart \(v=1\), blowing up the ideal \((u,t)\) resolves the
pair \([u:t]\). On the other chart the pair was already regular.
The strict transform of \(t=0\) is mapped to \([1:0]\), while the
exceptional curve parametrizes the new fiber.
The inverse rational map is \([a:b]\mapsto[ta:b]\), with its single
base point at \(t=b=0\); the same two blow-up charts resolve it.
Thus the common resolution is a point blow-up on either side,
and \(\mathbb P(F_1)\) is exactly the elementary transformation
of \(\mathbb P(F)\). This also shows that the inverse operation
is an elementary transformation.

Now write the given surfaces as \(\mathbb P(F)\) and \(\mathbb P(G)\).
Identify their generic fibers as two-dimensional vector spaces over
\(k(C)\), and view \(F,G\) as lattices in that common rational vector
space. They agree away from finitely many points after this
identification. Their intersection \(L=F\cap G\) is coherent:
locally it is the intersection of two finite lattices over a
discrete valuation ring, hence again a finite free lattice.
It is therefore a rank-two vector bundle, and both \(F/L\) and
\(G/L\) are torsion sheaves of finite length.

Choose composition series of these two quotients and take
preimages. They give chains from \(F\) to \(L\) and from \(G\) to
\(L\), each step the kernel of a surjection onto some \(k(p)\).
The local calculation realizes every step by an elementary
transformation. Follow the chain from \(\mathbb P(F)\) to
\(\mathbb P(L)\), then reverse the chain from \(\mathbb P(G)\).
This is the required finite sequence over \(C\).
\end{solution}

\begin{exerciseintro}
% record: H-V-05-006
\exerciseheading{Exercise V.5.6}
\addcontentsline{toc}{subsection}{Exercise V.5.6}

\begin{context}
Work over an algebraically closed field \(k\). A surface is smooth,
integral, and projective over \(k\). For the blow-up
\(\pi:\widetilde X\to X\) at a closed point, with exceptional curve \(E\),
\[
E\cong\mathbb P^1,\quad E^2=-1,\quad
(\pi^*D)E=0,\quad(\pi^*D)(\pi^*D')=DD'.
\]
If an effective divisor \(D\) has multiplicity \(m\) at the center,
its strict transform has class \(\pi^*D-mE\).
Intersection multiplicities of effective divisors with no common
component are nonnegative.

A valuation ring \(R\) of \(K/k\) contains \(k\), has fraction field
\(K\), and is trivial on \(k\). It has a unique center on any proper
model, and dominates the local ring at that center. A valuation
overring of a discrete valuation ring in its fraction field is
either that ring or the field.
Every rational function on a smooth projective surface becomes a
morphism to \(\mathbb P^1\) after finitely many point blow-ups.
The center of a valuation lifts uniquely under every proper blow-up.
\end{context}

\begin{problem}
Let \(K=k(X)\). Classify the nontrivial valuation rings \(R\ne K\)
of \(K/k\) as follows:

(1) the local ring at the generic point of a curve on \(X\);

(2) the local ring at the generic point of an exceptional curve
on a surface obtained from \(X\) by finitely many point blow-ups;

(3) the increasing union of local rings along an infinite sequence
of closed valuation centers, each obtained by blowing up the
preceding center.

In case (3), prove that the union equals \(R\).
\end{problem}
\end{exerciseintro}

\begin{solution}
Properness supplies a center on \(X\). It is not the generic
point, since domination there would force \(R=K\).
If it is the generic point of a curve, its local ring is a
discrete valuation ring. Domination and the overring assertion
force equality, giving (1).

Otherwise the center \(x_0\) is closed. Blow it up and take the
new center. If this is a curve's generic point, that curve is
exceptional, and the same argument gives (2).
Repeat whenever the center remains closed.
If this process never stops, write
\[
X_0=X,\qquad X_{i+1}=\operatorname{Bl}_{x_i}X_i,\qquad
A_i=\mathcal O_{X_i,x_i}.
\]
The maps give local inclusions \(A_i\subset A_{i+1}\subset R\),
so their union \(R_0\) is contained in \(R\).

We verify the opposite inclusion. Fix \(f\in R\), with \(f\ne0\).
Choose a finite sequence of point blow-ups \(Z\to X\) on which
\(f\) defines a morphism to \(\mathbb P^1\).
The infinite sequence of centers eventually maps, in a neighborhood
of its center, to \(Z\).
Here is the elementary lifting argument. Process the finite
sequence defining \(Z\) one blow-up at a time. If its next center
differs from the current valuation center, that blow-up is an
isomorphism near the latter and needs no new step in our sequence.
If they agree, perform the next blow-up in our sequence: it is
the same local blow-up and its valuation center corresponds to
the center on the chosen model. Induction through the finite
list gives the asserted local map for some \(X_N\).

Let \(z\) be the center on \(Z\).
Because \(f\in R\), the map \(\operatorname{Spec}R\to\mathbb P^1\)
defined by \([f:1]\) has closed-point image in the affine chart
with coordinate \(f\). It agrees generically with the map induced
from \(Z\), hence everywhere by separatedness of \(\mathbb P^1\).
Thus \(z\) maps to this affine chart and \(f\in\mathcal O_{Z,z}\).
The local map just constructed gives \(f\in A_N\subset R_0\).
The zero element is already in every \(A_i\), so \(R\subset R_0\).
Therefore \(R=R_0\), establishing (3).

These alternatives exhaust the construction. If trivial valuation
rings are included in the terminology, \(K\) itself must be added
as a separate case.
\end{solution}

\begin{exerciseintro}
% record: H-V-05-007
\exerciseheading{Exercise V.5.7}
\addcontentsline{toc}{subsection}{Exercise V.5.7}

\begin{context}
Work over an algebraically closed field \(k\). A surface is smooth,
integral, and projective over \(k\). For the blow-up
\(\pi:\widetilde X\to X\) at a closed point, with exceptional curve \(E\),
\[
E\cong\mathbb P^1,\quad E^2=-1,\quad
(\pi^*D)E=0,\quad(\pi^*D)(\pi^*D')=DD'.
\]
If an effective divisor \(D\) has multiplicity \(m\) at the center,
its strict transform has class \(\pi^*D-mE\).
Intersection multiplicities of effective divisors with no common
component are nonnegative.

For a surjective generically finite morphism \(f:X\to X_0\) of
projective integral surfaces and Cartier divisors \(H,H'\) on \(X_0\),
\((f^*H)(f^*H')=\deg(f)(HH')\).
A pulled-back divisor has intersection zero with every curve
contracted by \(f\).
The Hodge index theorem says that the real numerical divisor space
has signature \((1,\rho-1)\); in particular, the orthogonal
complement of any class of positive square is negative definite.
An effective nonzero curve has positive intersection with an ample
divisor, and therefore is not numerically zero.
\end{context}

\begin{problem}
Let \(f:X\to X_0\) be a dominant morphism to a projective integral surface,
and suppose that an integral curve \(Y\subset X\) is contracted
to a point \(p\). Prove that \(Y^2<0\).
In particular this holds when \(Y\) is the whole set-theoretic fiber.
\end{problem}
\end{exerciseintro}

\begin{solution}
Since \(X\) is projective, its image is closed. Dominance therefore
makes \(f\) surjective, and equality of dimensions makes it
generically finite. For a very ample Cartier divisor \(H\)
on \(X_0\), put \(D=f^*H\). The projection formulas give
\[
D^2=\deg(f)H^2>0,\qquad DY=0.
\]
The class of \(Y\) is nonzero numerically since an ample divisor
on \(X\) has positive intersection with it. The Hodge index
theorem, applied to the orthogonal complement of \(D\), yields
\(Y^2<0\).
\end{solution}

\begin{exerciseintro}
% record: H-V-05-008
\exerciseheading{Exercise V.5.8}
\addcontentsline{toc}{subsection}{Exercise V.5.8}

\begin{context}
Let \(k\) be algebraically closed, of arbitrary characteristic.
The blow-up of the origin in affine three-space has, for example,
the \(z\)-chart \(x=zu,\ y=zv\). For a hypersurface of multiplicity
\(m\), substitute these expressions and divide the equation by \(z^m\)
to obtain its strict transform. Restriction gives the blow-up of the
integral hypersurface itself.

We use the Jacobian criterion, and Nagata's factoriality criterion:
if \(B\) is a noetherian domain, \(p\) is a prime element, and
\(B[p^{-1}]\) is a unique factorization domain, then \(B\) is one too.
The binomial \(x^2+y^3\) is irreducible in \(k[x,y]\): the homomorphism
\(x\mapsto s^3,\ y\mapsto-s^2\) has image \(k[s^2,s^3]\) and kernel
\((x^2+y^3)\). Polynomial division by the monic quadratic in \(x\)
verifies the kernel assertion, since the remaining even powers of \(s\)
and the powers \(s^{3+2j}\) are disjoint.

A hypersurface double point with smooth projective tangent conic is
resolved by one point blow-up. Indeed, on the exceptional divisor the
strict-transform equation is that conic; its nonzero derivatives in
projective tangent directions give smoothness of the transform there.
The exceptional curve is the smooth conic, hence isomorphic to
\(\mathbb P^1\). In particular this applies to tangent cone
\(u^2+st=0\), also in characteristic two.
\end{context}

\begin{problem}
Let
\[
X=V(x^2+y^3+z^5)\subset\mathbb A^3_k,\qquad
A=k[x,y,z]/(x^2+y^3+z^5).
\]

(a) Prove that \(A\) is a unique factorization domain.

(b) Resolve its singularity at the origin by eight successive
point blow-ups. Show that the reduced exceptional fiber consists of
eight copies of \(\mathbb P^1\), with intersection graph \(E_8\):
a chain of seven vertices with an additional vertex attached at the
third vertex from one end.
\end{problem}
\end{exerciseintro}

\begin{solution}
(a) The element \(z\) is a nonzerodivisor in \(A\): in the polynomial
ring, \(z\) is relatively prime to \(x^2+y^3+z^5\), so
\(zg\in(x^2+y^3+z^5)\) implies \(g\) is in that ideal.
Thus \(A\) injects into \(A[z^{-1}]\).
Set
\[
t=z^{-1},\qquad u=t^3x,\qquad v=t^2y.
\]
Multiplying the defining equation by \(t^6\) gives
\(u^2+v^3+t=0\). Conversely
\(x=u/t^3,\ y=v/t^2,\ z=1/t\).
These substitutions yield inverse isomorphisms
\[
A[z^{-1}]
\cong k[u,v,t,t^{-1}]/(u^2+v^3+t)
\cong k[u,v,(u^2+v^3)^{-1}].
\]
This is a localization of a polynomial ring, hence a domain and
a unique factorization domain. The injection already proves that
\(A\) is a domain. Furthermore
\[
A/(z)\cong k[x,y]/(x^2+y^3)
\]
is a domain by the kernel calculation in the context, so \(z\)
is prime in \(A\). Nagata's criterion proves that \(A\) is a
unique factorization domain.

(b) The Jacobian criterion shows that the only singular point is the
origin. In characteristics two, three, or five, the two nonzero
partial derivatives force the corresponding two coordinates to vanish,
and the equation forces the remaining coordinate to vanish as well.
We give all charts that can contain singular points above the origin;
away from the exceptional locus every blow-up is an isomorphism.
Denote the new exceptional curves, with their reduced structures,
by \(E_1,\ldots,E_8\), in order of creation.

Blow up the origin first. In the \(z\)-chart, put \(x=zu,\ y=zv\);
the equation becomes
\[
u^2+zv^3+z^3=0.
\tag{1}
\]
Its exceptional curve is \(z=u=0\). Along this curve the only
possible singular point is \(v=0\), because the derivative with
respect to \(z\) restricts to \(v^3\).
In the \(y\)-chart the equation is
\(u^2+y+y^3v^5=0\), which is smooth along \(y=0\).
The \(x\)-chart has no exceptional points, since its equation
restricts to \(1=0\) there.
Thus (1) has exactly one singular point above the origin.

Blow up that point. In its \(v\)-chart, put \(u=va,\ z=vb\)
and write \(s=v,\ t=b\). The equation is
\[
a^2+s^2t+st^3=0.
\tag{2}
\]
Here \(E_2\) is \(s=a=0\), and the strict transform of \(E_1\)
is \(t=a=0\).
The derivative with respect to \(s\) along \(E_2\) is \(t^3\),
so the only singular point in this chart is its origin.
The other relevant chart, \(u=za,\ v=zb\), has equation
\(a^2+z+z^2b^3=0\), smooth on \(z=0\).
Again the chart of the squared variable has no exceptional points.

The third blow-up is at the origin of (2).
Its two relevant charts are
\[
\begin{array}{ll}
a=su,\ t=sv:&u^2+sv+s^2v^3=0,\\
a=tu,\ s=tv:&u^2+tv^2+t^2v=0.
\end{array}
\tag{3}
\]
On \(E_3\), the first chart has one singular point, \(u=s=v=0\);
its quadratic tangent cone is \(u^2+sv=0\).
The second chart has one singular point, \(u=t=v=0\),
with exact equation
\[
u^2+tv(v+t)=0.
\tag{4}
\]
These are different points of \(E_3\): they correspond to the
directions \([s:t]=[1:0]\) and \([0:1]\).
They are all the singular points, since the relevant derivatives
on \(E_3\) are respectively \(v\) and \(v^2\).

Blow up the first point in (3) as the fourth blow-up.
Its tangent conic is smooth in every characteristic, so this resolves
that point and introduces the smooth conic \(E_4\).
The strict transforms of \(E_1\) and \(E_3\) meet \(E_4\) at two
distinct points, corresponding to their two coordinate-axis
directions in that conic. Their intersections are transverse:
in a chart at either axis direction, the transformed equation
solves for one coordinate, and the old curve and exceptional
curve are the two remaining coordinate axes.

The fifth blow-up is at the point (4).
In the \(t\)-chart, write \(u=tw,\ v=tb\); the equation is
\[
w^2+tb(b+1)=0.
\tag{5}
\]
Its exceptional curve is \(t=w=0\), and its singular points
on that curve are \(b=0,-1\).
The \(v\)-chart gives \(w^2+vb(b+1)=0\) and supplies one more
point, the direction missing from the \(t\)-chart.
Thus \(E_5\) contains exactly three singular points, corresponding
to the directions \(v=0,\ t=0,\ v+t=0\) in (4).
These directions are distinct even in characteristic two.

At each of these points the tangent cone has form
\(w^2+c\,rq=0\) for \(c=1\) or \(-1\).
For example, at \(b=0\) in (5) it is \(w^2+tb\);
at \(b=-1\), put \(b=-1+q\) and obtain \(w^2-tq\).
Each projective tangent conic is smooth. Blow up these three
points separately as blow-ups six, seven, and eight.
They are disjoint centers, and each is resolved by its one blow-up.
Name the resulting curves \(E_6,E_7,E_8\) according to the
directions \(v=0,\ t=0,\ v+t=0\), respectively.
The surface is now smooth.

It remains to record the graph, including the older curves.
After the second blow-up, \(E_1,E_2\) pass through the point (2).
The third blow-up separates their directions:
\(E_1\) meets \(E_3\) at the first point in (3), and \(E_2\)
meets \(E_3\) at (4).
The fourth inserts \(E_4\) between \(E_1\) and \(E_3\).
At (4), \(E_2\) has \(v=u=0\), whereas \(E_3\) has \(t=u=0\).
After the fifth blow-up they meet \(E_5\) in the first two
singular directions just listed.
The sixth and seventh blow-ups insert \(E_6\) between \(E_2,E_5\)
and \(E_7\) between \(E_3,E_5\).
The third singular direction contains neither older curve, so the
eighth adds \(E_8\) meeting only \(E_5\).
Consequently the edges are exactly
\begin{center}
\begin{picture}(280,80)
\qbezier(20,20)(40,20)(60,20)
\qbezier(60,20)(80,20)(100,20)
\qbezier(100,20)(120,20)(140,20)
\qbezier(140,20)(160,20)(180,20)
\qbezier(180,20)(200,20)(220,20)
\qbezier(220,20)(240,20)(260,20)
\qbezier(180,20)(180,40)(180,60)
\put(20,20){\makebox(0,0){\(\bullet\)}}
\put(60,20){\makebox(0,0){\(\bullet\)}}
\put(100,20){\makebox(0,0){\(\bullet\)}}
\put(140,20){\makebox(0,0){\(\bullet\)}}
\put(180,20){\makebox(0,0){\(\bullet\)}}
\put(220,20){\makebox(0,0){\(\bullet\)}}
\put(260,20){\makebox(0,0){\(\bullet\)}}
\put(180,60){\makebox(0,0){\(\bullet\)}}
\put(20,5){\makebox(0,0){\(E_1\)}}
\put(60,5){\makebox(0,0){\(E_4\)}}
\put(100,5){\makebox(0,0){\(E_3\)}}
\put(140,5){\makebox(0,0){\(E_7\)}}
\put(180,5){\makebox(0,0){\(E_5\)}}
\put(220,5){\makebox(0,0){\(E_6\)}}
\put(260,5){\makebox(0,0){\(E_2\)}}
\put(180,75){\makebox(0,0){\(E_8\)}}
\end{picture}
\end{center}
Thus the chain is \(E_1-E_4-E_3-E_7-E_5-E_6-E_2\), with
\(E_8\) attached only to \(E_5\).
Every indicated intersection is transverse, by the same coordinate
calculation at the two axes of a resolved quadratic cone; the
three points on \(E_5\) are distinct.
All other pairs are disjoint.
The exceptional curves of blow-ups one, two, three, and five
are reduced projective lines; those of the other four blow-ups
are smooth conics. Their later strict transforms remain isomorphic
to \(\mathbb P^1\).
Thus there are exactly eight exceptional curves isomorphic to
\(\mathbb P^1\), with the stated \(E_8\) graph; this does not
assert that they are lines in a chosen projective embedding.
\end{solution}

\begin{exerciseintro}
\corpussection{V.6 Classification of Surfaces}
\addcontentsline{toc}{section}{V.6 Classification of Surfaces}

% record: H-V-06-001
\exerciseheading{Exercise V.6.1}
\addcontentsline{toc}{subsection}{Exercise V.6.1}

\begin{context}
Work over an algebraically closed field. A surface here is smooth,
integral, and projective. Write \(H\) for its hyperplane divisor.
Adjunction for a smooth complete intersection of degrees
\(d_1,\ldots,d_c\) in \(\mathbb P^n\) gives
\[
K_X=\left(\sum_i d_i-n-1\right)H.
\]
Its structure sheaf has the Koszul resolution with \(j\)-th term
\[
Q_j=\bigoplus_{|J|=j}\mathcal O_{\mathbb P^n}
\left(-\sum_{i\in J}d_i\right).
\]
We use \(H^i(\mathbb P^n,\mathcal O(t))=0\) for \(0<i<n\),
Serre duality, and Castelnuovo's criterion: a smooth projective
surface with \(p_a=0\) and \(h^0(2K)=0\) is rational.
A surface with ample canonical divisor has Kodaira dimension two,
hence is of general type. A K3 surface is a smooth projective
surface with trivial canonical bundle and \(h^1(\mathcal O)=0\).
\end{context}

\begin{problem}
Let \(X\subset\mathbb P^n\), \(n\geq3\), be a smooth complete
intersection surface of degrees \(d_1,\ldots,d_{n-2}\), each at
least two. Determine exactly which degree lists fail to give a
surface of general type, and classify the exceptional surfaces.
Regard permutations of a degree list as equivalent.
\end{problem}
\end{exerciseintro}

\begin{solution}
First \(q(X)=h^1(X,\mathcal O_X)=0\).
For completeness, let \(c=n-2\), and let \(B_j\) be the image of
\(Q_j\to Q_{j-1}\), with \(Q_0=\mathcal O_{\mathbb P^n}\).
Thus \(B_1=\mathcal I_X\), \(B_c\cong Q_c\), and
\[
0\to B_{j+1}\to Q_j\to B_j\to0
\]
is exact for \(j<c\).
The sequence for \(\mathcal I_X\) gives
\(H^1(\mathcal O_X)\cong H^2(B_1)\).
Successive dimension shifts give
\[
H^2(B_1)\cong H^3(B_2)\cong\cdots
\cong H^{c+1}(Q_c)=H^{n-1}(Q_c)=0.
\]
Each shift uses only the stated intermediate cohomology vanishing;
when \(c=1\) the first group is already \(H^2(Q_1)=0\).

Put
\[
m=\sum_i d_i-n-1=\sum_i(d_i-1)-3.
\]
If \(m>0\), then \(K_X=mH\) is ample and \(X\) is of general type.
If \(m=0\), then \(K_X\) is trivial; together with \(q=0\)
this makes \(X\) a K3 surface, with \(p_g=1\) and Kodaira
dimension zero.
If \(m<0\), every positive multiple of \(K_X\) has negative
intersection with \(H\), so it has no effective representative.
In particular \(p_g=h^0(K_X)=0\) and \(h^0(2K_X)=0\).
Since \(p_a=p_g-q=0\), Castelnuovo's criterion makes \(X\)
rational, with Kodaira dimension \(-1\) in the convention of this
chapter.

Every \(d_i-1\) is a positive integer. Thus \(m\leq0\) is
equivalent to a partition of one, two, or three into \(c=n-2\)
positive parts. The complete list is
\[
\begin{array}{c|c|c|c}
n&(d_i)&m&\text{classification}\\ \hline
3&(2)&-2&\text{rational quadric}\\
3&(3)&-1&\text{rational cubic surface}\\
4&(2,2)&-1&\text{rational degree-four del Pezzo surface}\\
3&(4)&0&\text{K3 surface}\\
4&(2,3)&0&\text{K3 surface}\\
5&(2,2,2)&0&\text{K3 surface}.
\end{array}
\]
There are no further exceptional lists, since a sum of \(c\)
positive integers at most three has \(c\leq3\).
The labels del Pezzo in the negative cases agree with the
ampleness of \(-K_X\). All other smooth complete intersections
in the problem have ample canonical bundle.
\end{solution}

\begin{exerciseintro}
% record: H-V-06-002
\exerciseheading{Exercise V.6.2}
\addcontentsline{toc}{subsection}{Exercise V.6.2}

\begin{context}
Work over \(\mathbb C\). Let \(X\subset\mathbb P^{n+1}\) be a
smooth integral projective surface, and let \(H\) be its hyperplane
divisor. A general \(C\in|H|\) is smooth and integral, with
\(K_C=(K_X+H)|_C\).
Clifford's theorem says that a special line bundle \(L\) with
sections on a smooth curve satisfies
\(h^0(L)\leq\deg(L)/2+1\).
For positive degree, equality occurs only for the canonical
bundle or for a multiple of the degree-two pencil of a hyperelliptic
curve; in the latter noncanonical case the complete system factors
through the degree-two map and is not very ample.

We use Serre duality, surface Riemann--Roch
\[
\chi(\mathcal O_X(D))=\chi(\mathcal O_X)+\frac{D(D-K_X)}2,
\]
and Kodaira vanishing \(H^i(X,\mathcal O_X(K_X+H))=0\) for
\(i>0\). Write \(p_g=h^0(K_X)\), \(q=h^1(\mathcal O_X)\);
then \(\chi(\mathcal O_X)=1-q+p_g\).
A K3 surface has trivial canonical bundle and \(q=0\).
\end{context}

\begin{problem}
Let \(X\subset\mathbb P^{n+1}_{\mathbb C}\) be a smooth integral
surface not contained in a hyperplane, and let \(d=\deg X\).
Prove that \(d<2n\) implies \(p_g=0\).
If \(d=2n\), prove that either \(p_g=0\), or \(p_g=1\) and
\(X\) is a K3 surface.
\end{problem}
\end{exerciseintro}

\begin{solution}
For \(n=1\), a closed integral surface in \(\mathbb P^2\) is
\(\mathbb P^2\) itself; it has degree one and \(p_g=0\).
Assume \(n\geq2\), and suppose \(p_g>0\).
Choose a nonzero canonical section, with effective zero divisor \(D\),
and a general smooth hyperplane section \(C\) not contained in \(D\).
Its restriction to \(C\) is a nonzero section of \(K_X|_C\).
Set \(L=\mathcal O_C(H)\). Adjunction and duality on \(C\) give
\[
h^1(C,L)=h^0(C,K_C\otimes L^{-1})
=h^0(C,K_X|_C)>0,
\]
so \(L\) is special.

The linear forms on \(\mathbb P^{n+1}\) inject into
\(H^0(X,\mathcal O_X(H))\), by nondegeneracy.
Their restrictions to \(C\) have a one-dimensional kernel:
the exact sequence
\[
0\to\mathcal O_X\to\mathcal O_X(H)\to\mathcal O_C(H)\to0
\]
identifies the kernel of restriction on all sections with the
one-dimensional space generated by the equation of \(C\).
Consequently \(h^0(C,L)\geq n+1\).
Clifford's theorem yields
\[
n+1\leq h^0(C,L)\leq\frac d2+1.
\]
Hence \(d\geq2n\), proving the first assertion.

Now suppose \(d=2n\) and \(p_g>0\). Equality holds in Clifford.
The line bundle \(L\) has positive degree and is very ample,
since it gives the hyperplane embedding of \(C\).
The noncanonical hyperelliptic equality cases are therefore
impossible, and \(L\cong K_C\).
Adjunction implies \(K_X|_C\cong\mathcal O_C\), so \(K_XH=0\).
The effective divisor \(D\) of the chosen canonical section satisfies
\(DH=0\). Ampleness of \(H\) forces \(D=0\).
Thus that section never vanishes,
\[
\mathcal O_X(K_X)\cong\mathcal O_X,\qquad p_g=1.
\]

It remains to prove \(q=0\). Kodaira vanishing now gives
\(H^i(X,\mathcal O_X(H))=0\) for \(i>0\).
Since \(H^2=d=2n\), Riemann--Roch becomes
\[
h^0(X,\mathcal O_X(H))
=\chi(\mathcal O_X(H))
=2-q+n.
\]
Nondegeneracy supplies at least \(n+2\) independent linear-form
sections. Hence \(2-q+n\geq n+2\), giving \(q\leq0\).
As \(q\) is a dimension, it equals zero.
Together with the trivial canonical bundle this says exactly
that \(X\) is a K3 surface.
\end{solution}

\end{document}
